\section{Bench Section 0}
% synthetic bench section 0
% This is a sample LaTeX input file.  (Version of 12 August 2004.)
%
% A '%' character causes TeX to ignore all remaining text on the line,
% and is used for comments like this one.

\documentclass{article}      % Specifies the document class

                             % The preamble begins here.
\title{An Example Document}  % Declares the document's title.
\author{Leslie Lamport}      % Declares the author's name.
\date{January 21, 1994}      % Deleting this command produces today's date.

\newcommand{\ip}[2]{(#1, #2)}
                             % Defines \ip{arg1}{arg2} to mean
                             % (arg1, arg2).

%\newcommand{\ip}[2]{\langle #1 | #2\rangle}
                             % This is an alternative definition of
                             % \ip that is commented out.

\begin{document}             % End of preamble and beginning of text.

\maketitle                   % Produces the title.

This is an example input file.  Comparing it with
the output it generates can show you how to
produce a simple document of your own.

\section{Ordinary Text}      % Produces section heading.  Lower-level
                             % sections are begun with similar
                             % \subsection and \subsubsection commands.

The ends  of words and sentences are marked
  by   spaces. It  doesn't matter how many
spaces    you type; one is as good as 100.  The
end of   a line counts as a space.

One   or more   blank lines denote the  end
of  a paragraph.

Since any number of consecutive spaces are treated
like a single one, the formatting of the input
file makes no difference to
      \LaTeX,                % The \LaTeX command generates the LaTeX logo.
but it makes a difference to you.  When you use
\LaTeX, making your input file as easy to read
as possible will be a great help as you write
your document and when you change it.  This sample
file shows how you can add comments to your own input
file.

Because printing is different from typewriting,
there are a number of things that you have to do
differently when preparing an input file than if
you were just typing the document directly.
Quotation marks like
       ``this''
have to be handled specially, as do quotes within
quotes:
       ``\,`this'            % \, separates the double and single quote.
        is what I just
        wrote, not  `that'\,''.

Dashes come in three sizes: an
       intra-word
dash, a medium dash for number ranges like
       1--2,
and a punctuation
       dash---like
this.

A sentence-ending space should be larger than the
space between words within a sentence.  You
sometimes have to type special commands in
conjunction with punctuation characters to get
this right, as in the following sentence.
       Gnats, gnus, etc.\ all  % `\ ' makes an inter-word space.
       begin with G\@.         % \@ marks end-of-sentence punctuation.
You should check the spaces after periods when
reading your output to make sure you haven't
forgotten any special cases.  Generating an
ellipsis
       \ldots\               % `\ ' is needed after `\ldots' because TeX
                             % ignores spaces after command names like \ldots
                             % made from \ + letters.
                             %
                             % Note how a `%' character causes TeX to ignore
                             % the end of the input line, so these blank lines
                             % do not start a new paragraph.
                             %
with the right spacing around the periods requires
a special command.

\LaTeX\ interprets some common characters as
commands, so you must type special commands to
generate them.  These characters include the
following:
       \$ \& \% \# \{ and \}.

In printing, text is usually emphasized with an
       \emph{italic}
type style.

\begin{em}
   A long segment of text can also be emphasized
   in this way.  Text within such a segment can be
   given \emph{additional} emphasis.
\end{em}

It is sometimes necessary to prevent \LaTeX\ from
breaking a line where it might otherwise do so.
This may be at a space, as between the ``Mr.''\ and
``Jones'' in
       ``Mr.~Jones'',        % ~ produces an unbreakable interword space.
or within a word---especially when the word is a
symbol like
       \mbox{\emph{itemnum}}
that makes little sense when hyphenated across
lines.

Footnotes\footnote{This is an example of a footnote.}
pose no problem.

\LaTeX\ is good at typesetting mathematical formulas
like
       \( x-3y + z = 7 \)
or
       \( a_{1} > x^{2n} + y^{2n} > x' \)
or
       \( \ip{A}{B} = \sum_{i} a_{i} b_{i} \).
The spaces you type in a formula are
ignored.  Remember that a letter like
       $x$                   % $ ... $  and  \( ... \)  are equivalent
is a formula when it denotes a mathematical
symbol, and it should be typed as one.

\section{Displayed Text}

Text is displayed by indenting it from the left
margin.  Quotations are commonly displayed.  There
are short quotations
\begin{quote}
   This is a short quotation.  It consists of a
   single paragraph of text.  See how it is formatted.
\end{quote}
and longer ones.
\begin{quotation}
   This is a longer quotation.  It consists of two
   paragraphs of text, neither of which are
   particularly interesting.

   This is the second paragraph of the quotation.  It
   is just as dull as the first paragraph.
\end{quotation}
Another frequently-displayed structure is a list.
The following is an example of an \emph{itemized}
list.
\begin{itemize}
   \item This is the first item of an itemized list.
         Each item in the list is marked with a ``tick''.
         You don't have to worry about what kind of tick
         mark is used.

   \item This is the second item of the list.  It
         contains another list nested inside it.  The inner
         list is an \emph{enumerated} list.
         \begin{enumerate}
            \item This is the first item of an enumerated
                  list that is nested within the itemized list.

            \item This is the second item of the inner list.
                  \LaTeX\ allows you to nest lists deeper than
                  you really should.
         \end{enumerate}
         This is the rest of the second item of the outer
         list.  It is no more interesting than any other
         part of the item.
   \item This is the third item of the list.
\end{itemize}
You can even display poetry.
\begin{verse}
   There is an environment
    for verse \\             % The \\ command separates lines
   Whose features some poets % within a stanza.
   will curse.

                             % One or more blank lines separate stanzas.

   For instead of making\\
   Them do \emph{all} line breaking, \\
   It allows them to put too many words on a line when they'd rather be
   forced to be terse.
\end{verse}

Mathematical formulas may also be displayed.  A
displayed formula
is
one-line long; multiline
formulas require special formatting instructions.
   \[  \ip{\Gamma}{\psi'} = x'' + y^{2} + z_{i}^{n}\]
Don't start a paragraph with a displayed equation,
nor make one a paragraph by itself.

\end{document}               % End of document.

\section{Bench Section 1}
% synthetic bench section 1
% This is a sample LaTeX input file.  (Version of 12 August 2004.)
%
% A '%' character causes TeX to ignore all remaining text on the line,
% and is used for comments like this one.

\documentclass{article}      % Specifies the document class

                             % The preamble begins here.
\title{An Example Document}  % Declares the document's title.
\author{Leslie Lamport}      % Declares the author's name.
\date{January 21, 1994}      % Deleting this command produces today's date.

\newcommand{\ip}[2]{(#1, #2)}
                             % Defines \ip{arg1}{arg2} to mean
                             % (arg1, arg2).

%\newcommand{\ip}[2]{\langle #1 | #2\rangle}
                             % This is an alternative definition of
                             % \ip that is commented out.

\begin{document}             % End of preamble and beginning of text.

\maketitle                   % Produces the title.

This is an example input file.  Comparing it with
the output it generates can show you how to
produce a simple document of your own.

\section{Ordinary Text}      % Produces section heading.  Lower-level
                             % sections are begun with similar
                             % \subsection and \subsubsection commands.

The ends  of words and sentences are marked
  by   spaces. It  doesn't matter how many
spaces    you type; one is as good as 100.  The
end of   a line counts as a space.

One   or more   blank lines denote the  end
of  a paragraph.

Since any number of consecutive spaces are treated
like a single one, the formatting of the input
file makes no difference to
      \LaTeX,                % The \LaTeX command generates the LaTeX logo.
but it makes a difference to you.  When you use
\LaTeX, making your input file as easy to read
as possible will be a great help as you write
your document and when you change it.  This sample
file shows how you can add comments to your own input
file.

Because printing is different from typewriting,
there are a number of things that you have to do
differently when preparing an input file than if
you were just typing the document directly.
Quotation marks like
       ``this''
have to be handled specially, as do quotes within
quotes:
       ``\,`this'            % \, separates the double and single quote.
        is what I just
        wrote, not  `that'\,''.

Dashes come in three sizes: an
       intra-word
dash, a medium dash for number ranges like
       1--2,
and a punctuation
       dash---like
this.

A sentence-ending space should be larger than the
space between words within a sentence.  You
sometimes have to type special commands in
conjunction with punctuation characters to get
this right, as in the following sentence.
       Gnats, gnus, etc.\ all  % `\ ' makes an inter-word space.
       begin with G\@.         % \@ marks end-of-sentence punctuation.
You should check the spaces after periods when
reading your output to make sure you haven't
forgotten any special cases.  Generating an
ellipsis
       \ldots\               % `\ ' is needed after `\ldots' because TeX
                             % ignores spaces after command names like \ldots
                             % made from \ + letters.
                             %
                             % Note how a `%' character causes TeX to ignore
                             % the end of the input line, so these blank lines
                             % do not start a new paragraph.
                             %
with the right spacing around the periods requires
a special command.

\LaTeX\ interprets some common characters as
commands, so you must type special commands to
generate them.  These characters include the
following:
       \$ \& \% \# \{ and \}.

In printing, text is usually emphasized with an
       \emph{italic}
type style.

\begin{em}
   A long segment of text can also be emphasized
   in this way.  Text within such a segment can be
   given \emph{additional} emphasis.
\end{em}

It is sometimes necessary to prevent \LaTeX\ from
breaking a line where it might otherwise do so.
This may be at a space, as between the ``Mr.''\ and
``Jones'' in
       ``Mr.~Jones'',        % ~ produces an unbreakable interword space.
or within a word---especially when the word is a
symbol like
       \mbox{\emph{itemnum}}
that makes little sense when hyphenated across
lines.

Footnotes\footnote{This is an example of a footnote.}
pose no problem.

\LaTeX\ is good at typesetting mathematical formulas
like
       \( x-3y + z = 7 \)
or
       \( a_{1} > x^{2n} + y^{2n} > x' \)
or
       \( \ip{A}{B} = \sum_{i} a_{i} b_{i} \).
The spaces you type in a formula are
ignored.  Remember that a letter like
       $x$                   % $ ... $  and  \( ... \)  are equivalent
is a formula when it denotes a mathematical
symbol, and it should be typed as one.

\section{Displayed Text}

Text is displayed by indenting it from the left
margin.  Quotations are commonly displayed.  There
are short quotations
\begin{quote}
   This is a short quotation.  It consists of a
   single paragraph of text.  See how it is formatted.
\end{quote}
and longer ones.
\begin{quotation}
   This is a longer quotation.  It consists of two
   paragraphs of text, neither of which are
   particularly interesting.

   This is the second paragraph of the quotation.  It
   is just as dull as the first paragraph.
\end{quotation}
Another frequently-displayed structure is a list.
The following is an example of an \emph{itemized}
list.
\begin{itemize}
   \item This is the first item of an itemized list.
         Each item in the list is marked with a ``tick''.
         You don't have to worry about what kind of tick
         mark is used.

   \item This is the second item of the list.  It
         contains another list nested inside it.  The inner
         list is an \emph{enumerated} list.
         \begin{enumerate}
            \item This is the first item of an enumerated
                  list that is nested within the itemized list.

            \item This is the second item of the inner list.
                  \LaTeX\ allows you to nest lists deeper than
                  you really should.
         \end{enumerate}
         This is the rest of the second item of the outer
         list.  It is no more interesting than any other
         part of the item.
   \item This is the third item of the list.
\end{itemize}
You can even display poetry.
\begin{verse}
   There is an environment
    for verse \\             % The \\ command separates lines
   Whose features some poets % within a stanza.
   will curse.

                             % One or more blank lines separate stanzas.

   For instead of making\\
   Them do \emph{all} line breaking, \\
   It allows them to put too many words on a line when they'd rather be
   forced to be terse.
\end{verse}

Mathematical formulas may also be displayed.  A
displayed formula
is
one-line long; multiline
formulas require special formatting instructions.
   \[  \ip{\Gamma}{\psi'} = x'' + y^{2} + z_{i}^{n}\]
Don't start a paragraph with a displayed equation,
nor make one a paragraph by itself.

\end{document}               % End of document.

\section{Bench Section 2}
% synthetic bench section 2
% This is a sample LaTeX input file.  (Version of 12 August 2004.)
%
% A '%' character causes TeX to ignore all remaining text on the line,
% and is used for comments like this one.

\documentclass{article}      % Specifies the document class

                             % The preamble begins here.
\title{An Example Document}  % Declares the document's title.
\author{Leslie Lamport}      % Declares the author's name.
\date{January 21, 1994}      % Deleting this command produces today's date.

\newcommand{\ip}[2]{(#1, #2)}
                             % Defines \ip{arg1}{arg2} to mean
                             % (arg1, arg2).

%\newcommand{\ip}[2]{\langle #1 | #2\rangle}
                             % This is an alternative definition of
                             % \ip that is commented out.

\begin{document}             % End of preamble and beginning of text.

\maketitle                   % Produces the title.

This is an example input file.  Comparing it with
the output it generates can show you how to
produce a simple document of your own.

\section{Ordinary Text}      % Produces section heading.  Lower-level
                             % sections are begun with similar
                             % \subsection and \subsubsection commands.

The ends  of words and sentences are marked
  by   spaces. It  doesn't matter how many
spaces    you type; one is as good as 100.  The
end of   a line counts as a space.

One   or more   blank lines denote the  end
of  a paragraph.

Since any number of consecutive spaces are treated
like a single one, the formatting of the input
file makes no difference to
      \LaTeX,                % The \LaTeX command generates the LaTeX logo.
but it makes a difference to you.  When you use
\LaTeX, making your input file as easy to read
as possible will be a great help as you write
your document and when you change it.  This sample
file shows how you can add comments to your own input
file.

Because printing is different from typewriting,
there are a number of things that you have to do
differently when preparing an input file than if
you were just typing the document directly.
Quotation marks like
       ``this''
have to be handled specially, as do quotes within
quotes:
       ``\,`this'            % \, separates the double and single quote.
        is what I just
        wrote, not  `that'\,''.

Dashes come in three sizes: an
       intra-word
dash, a medium dash for number ranges like
       1--2,
and a punctuation
       dash---like
this.

A sentence-ending space should be larger than the
space between words within a sentence.  You
sometimes have to type special commands in
conjunction with punctuation characters to get
this right, as in the following sentence.
       Gnats, gnus, etc.\ all  % `\ ' makes an inter-word space.
       begin with G\@.         % \@ marks end-of-sentence punctuation.
You should check the spaces after periods when
reading your output to make sure you haven't
forgotten any special cases.  Generating an
ellipsis
       \ldots\               % `\ ' is needed after `\ldots' because TeX
                             % ignores spaces after command names like \ldots
                             % made from \ + letters.
                             %
                             % Note how a `%' character causes TeX to ignore
                             % the end of the input line, so these blank lines
                             % do not start a new paragraph.
                             %
with the right spacing around the periods requires
a special command.

\LaTeX\ interprets some common characters as
commands, so you must type special commands to
generate them.  These characters include the
following:
       \$ \& \% \# \{ and \}.

In printing, text is usually emphasized with an
       \emph{italic}
type style.

\begin{em}
   A long segment of text can also be emphasized
   in this way.  Text within such a segment can be
   given \emph{additional} emphasis.
\end{em}

It is sometimes necessary to prevent \LaTeX\ from
breaking a line where it might otherwise do so.
This may be at a space, as between the ``Mr.''\ and
``Jones'' in
       ``Mr.~Jones'',        % ~ produces an unbreakable interword space.
or within a word---especially when the word is a
symbol like
       \mbox{\emph{itemnum}}
that makes little sense when hyphenated across
lines.

Footnotes\footnote{This is an example of a footnote.}
pose no problem.

\LaTeX\ is good at typesetting mathematical formulas
like
       \( x-3y + z = 7 \)
or
       \( a_{1} > x^{2n} + y^{2n} > x' \)
or
       \( \ip{A}{B} = \sum_{i} a_{i} b_{i} \).
The spaces you type in a formula are
ignored.  Remember that a letter like
       $x$                   % $ ... $  and  \( ... \)  are equivalent
is a formula when it denotes a mathematical
symbol, and it should be typed as one.

\section{Displayed Text}

Text is displayed by indenting it from the left
margin.  Quotations are commonly displayed.  There
are short quotations
\begin{quote}
   This is a short quotation.  It consists of a
   single paragraph of text.  See how it is formatted.
\end{quote}
and longer ones.
\begin{quotation}
   This is a longer quotation.  It consists of two
   paragraphs of text, neither of which are
   particularly interesting.

   This is the second paragraph of the quotation.  It
   is just as dull as the first paragraph.
\end{quotation}
Another frequently-displayed structure is a list.
The following is an example of an \emph{itemized}
list.
\begin{itemize}
   \item This is the first item of an itemized list.
         Each item in the list is marked with a ``tick''.
         You don't have to worry about what kind of tick
         mark is used.

   \item This is the second item of the list.  It
         contains another list nested inside it.  The inner
         list is an \emph{enumerated} list.
         \begin{enumerate}
            \item This is the first item of an enumerated
                  list that is nested within the itemized list.

            \item This is the second item of the inner list.
                  \LaTeX\ allows you to nest lists deeper than
                  you really should.
         \end{enumerate}
         This is the rest of the second item of the outer
         list.  It is no more interesting than any other
         part of the item.
   \item This is the third item of the list.
\end{itemize}
You can even display poetry.
\begin{verse}
   There is an environment
    for verse \\             % The \\ command separates lines
   Whose features some poets % within a stanza.
   will curse.

                             % One or more blank lines separate stanzas.

   For instead of making\\
   Them do \emph{all} line breaking, \\
   It allows them to put too many words on a line when they'd rather be
   forced to be terse.
\end{verse}

Mathematical formulas may also be displayed.  A
displayed formula
is
one-line long; multiline
formulas require special formatting instructions.
   \[  \ip{\Gamma}{\psi'} = x'' + y^{2} + z_{i}^{n}\]
Don't start a paragraph with a displayed equation,
nor make one a paragraph by itself.

\end{document}               % End of document.

\section{Bench Section 3}
% synthetic bench section 3
% This is a sample LaTeX input file.  (Version of 12 August 2004.)
%
% A '%' character causes TeX to ignore all remaining text on the line,
% and is used for comments like this one.

\documentclass{article}      % Specifies the document class

                             % The preamble begins here.
\title{An Example Document}  % Declares the document's title.
\author{Leslie Lamport}      % Declares the author's name.
\date{January 21, 1994}      % Deleting this command produces today's date.

\newcommand{\ip}[2]{(#1, #2)}
                             % Defines \ip{arg1}{arg2} to mean
                             % (arg1, arg2).

%\newcommand{\ip}[2]{\langle #1 | #2\rangle}
                             % This is an alternative definition of
                             % \ip that is commented out.

\begin{document}             % End of preamble and beginning of text.

\maketitle                   % Produces the title.

This is an example input file.  Comparing it with
the output it generates can show you how to
produce a simple document of your own.

\section{Ordinary Text}      % Produces section heading.  Lower-level
                             % sections are begun with similar
                             % \subsection and \subsubsection commands.

The ends  of words and sentences are marked
  by   spaces. It  doesn't matter how many
spaces    you type; one is as good as 100.  The
end of   a line counts as a space.

One   or more   blank lines denote the  end
of  a paragraph.

Since any number of consecutive spaces are treated
like a single one, the formatting of the input
file makes no difference to
      \LaTeX,                % The \LaTeX command generates the LaTeX logo.
but it makes a difference to you.  When you use
\LaTeX, making your input file as easy to read
as possible will be a great help as you write
your document and when you change it.  This sample
file shows how you can add comments to your own input
file.

Because printing is different from typewriting,
there are a number of things that you have to do
differently when preparing an input file than if
you were just typing the document directly.
Quotation marks like
       ``this''
have to be handled specially, as do quotes within
quotes:
       ``\,`this'            % \, separates the double and single quote.
        is what I just
        wrote, not  `that'\,''.

Dashes come in three sizes: an
       intra-word
dash, a medium dash for number ranges like
       1--2,
and a punctuation
       dash---like
this.

A sentence-ending space should be larger than the
space between words within a sentence.  You
sometimes have to type special commands in
conjunction with punctuation characters to get
this right, as in the following sentence.
       Gnats, gnus, etc.\ all  % `\ ' makes an inter-word space.
       begin with G\@.         % \@ marks end-of-sentence punctuation.
You should check the spaces after periods when
reading your output to make sure you haven't
forgotten any special cases.  Generating an
ellipsis
       \ldots\               % `\ ' is needed after `\ldots' because TeX
                             % ignores spaces after command names like \ldots
                             % made from \ + letters.
                             %
                             % Note how a `%' character causes TeX to ignore
                             % the end of the input line, so these blank lines
                             % do not start a new paragraph.
                             %
with the right spacing around the periods requires
a special command.

\LaTeX\ interprets some common characters as
commands, so you must type special commands to
generate them.  These characters include the
following:
       \$ \& \% \# \{ and \}.

In printing, text is usually emphasized with an
       \emph{italic}
type style.

\begin{em}
   A long segment of text can also be emphasized
   in this way.  Text within such a segment can be
   given \emph{additional} emphasis.
\end{em}

It is sometimes necessary to prevent \LaTeX\ from
breaking a line where it might otherwise do so.
This may be at a space, as between the ``Mr.''\ and
``Jones'' in
       ``Mr.~Jones'',        % ~ produces an unbreakable interword space.
or within a word---especially when the word is a
symbol like
       \mbox{\emph{itemnum}}
that makes little sense when hyphenated across
lines.

Footnotes\footnote{This is an example of a footnote.}
pose no problem.

\LaTeX\ is good at typesetting mathematical formulas
like
       \( x-3y + z = 7 \)
or
       \( a_{1} > x^{2n} + y^{2n} > x' \)
or
       \( \ip{A}{B} = \sum_{i} a_{i} b_{i} \).
The spaces you type in a formula are
ignored.  Remember that a letter like
       $x$                   % $ ... $  and  \( ... \)  are equivalent
is a formula when it denotes a mathematical
symbol, and it should be typed as one.

\section{Displayed Text}

Text is displayed by indenting it from the left
margin.  Quotations are commonly displayed.  There
are short quotations
\begin{quote}
   This is a short quotation.  It consists of a
   single paragraph of text.  See how it is formatted.
\end{quote}
and longer ones.
\begin{quotation}
   This is a longer quotation.  It consists of two
   paragraphs of text, neither of which are
   particularly interesting.

   This is the second paragraph of the quotation.  It
   is just as dull as the first paragraph.
\end{quotation}
Another frequently-displayed structure is a list.
The following is an example of an \emph{itemized}
list.
\begin{itemize}
   \item This is the first item of an itemized list.
         Each item in the list is marked with a ``tick''.
         You don't have to worry about what kind of tick
         mark is used.

   \item This is the second item of the list.  It
         contains another list nested inside it.  The inner
         list is an \emph{enumerated} list.
         \begin{enumerate}
            \item This is the first item of an enumerated
                  list that is nested within the itemized list.

            \item This is the second item of the inner list.
                  \LaTeX\ allows you to nest lists deeper than
                  you really should.
         \end{enumerate}
         This is the rest of the second item of the outer
         list.  It is no more interesting than any other
         part of the item.
   \item This is the third item of the list.
\end{itemize}
You can even display poetry.
\begin{verse}
   There is an environment
    for verse \\             % The \\ command separates lines
   Whose features some poets % within a stanza.
   will curse.

                             % One or more blank lines separate stanzas.

   For instead of making\\
   Them do \emph{all} line breaking, \\
   It allows them to put too many words on a line when they'd rather be
   forced to be terse.
\end{verse}

Mathematical formulas may also be displayed.  A
displayed formula
is
one-line long; multiline
formulas require special formatting instructions.
   \[  \ip{\Gamma}{\psi'} = x'' + y^{2} + z_{i}^{n}\]
Don't start a paragraph with a displayed equation,
nor make one a paragraph by itself.

\end{document}               % End of document.

\section{Bench Section 4}
% synthetic bench section 4
% This is a sample LaTeX input file.  (Version of 12 August 2004.)
%
% A '%' character causes TeX to ignore all remaining text on the line,
% and is used for comments like this one.

\documentclass{article}      % Specifies the document class

                             % The preamble begins here.
\title{An Example Document}  % Declares the document's title.
\author{Leslie Lamport}      % Declares the author's name.
\date{January 21, 1994}      % Deleting this command produces today's date.

\newcommand{\ip}[2]{(#1, #2)}
                             % Defines \ip{arg1}{arg2} to mean
                             % (arg1, arg2).

%\newcommand{\ip}[2]{\langle #1 | #2\rangle}
                             % This is an alternative definition of
                             % \ip that is commented out.

\begin{document}             % End of preamble and beginning of text.

\maketitle                   % Produces the title.

This is an example input file.  Comparing it with
the output it generates can show you how to
produce a simple document of your own.

\section{Ordinary Text}      % Produces section heading.  Lower-level
                             % sections are begun with similar
                             % \subsection and \subsubsection commands.

The ends  of words and sentences are marked
  by   spaces. It  doesn't matter how many
spaces    you type; one is as good as 100.  The
end of   a line counts as a space.

One   or more   blank lines denote the  end
of  a paragraph.

Since any number of consecutive spaces are treated
like a single one, the formatting of the input
file makes no difference to
      \LaTeX,                % The \LaTeX command generates the LaTeX logo.
but it makes a difference to you.  When you use
\LaTeX, making your input file as easy to read
as possible will be a great help as you write
your document and when you change it.  This sample
file shows how you can add comments to your own input
file.

Because printing is different from typewriting,
there are a number of things that you have to do
differently when preparing an input file than if
you were just typing the document directly.
Quotation marks like
       ``this''
have to be handled specially, as do quotes within
quotes:
       ``\,`this'            % \, separates the double and single quote.
        is what I just
        wrote, not  `that'\,''.

Dashes come in three sizes: an
       intra-word
dash, a medium dash for number ranges like
       1--2,
and a punctuation
       dash---like
this.

A sentence-ending space should be larger than the
space between words within a sentence.  You
sometimes have to type special commands in
conjunction with punctuation characters to get
this right, as in the following sentence.
       Gnats, gnus, etc.\ all  % `\ ' makes an inter-word space.
       begin with G\@.         % \@ marks end-of-sentence punctuation.
You should check the spaces after periods when
reading your output to make sure you haven't
forgotten any special cases.  Generating an
ellipsis
       \ldots\               % `\ ' is needed after `\ldots' because TeX
                             % ignores spaces after command names like \ldots
                             % made from \ + letters.
                             %
                             % Note how a `%' character causes TeX to ignore
                             % the end of the input line, so these blank lines
                             % do not start a new paragraph.
                             %
with the right spacing around the periods requires
a special command.

\LaTeX\ interprets some common characters as
commands, so you must type special commands to
generate them.  These characters include the
following:
       \$ \& \% \# \{ and \}.

In printing, text is usually emphasized with an
       \emph{italic}
type style.

\begin{em}
   A long segment of text can also be emphasized
   in this way.  Text within such a segment can be
   given \emph{additional} emphasis.
\end{em}

It is sometimes necessary to prevent \LaTeX\ from
breaking a line where it might otherwise do so.
This may be at a space, as between the ``Mr.''\ and
``Jones'' in
       ``Mr.~Jones'',        % ~ produces an unbreakable interword space.
or within a word---especially when the word is a
symbol like
       \mbox{\emph{itemnum}}
that makes little sense when hyphenated across
lines.

Footnotes\footnote{This is an example of a footnote.}
pose no problem.

\LaTeX\ is good at typesetting mathematical formulas
like
       \( x-3y + z = 7 \)
or
       \( a_{1} > x^{2n} + y^{2n} > x' \)
or
       \( \ip{A}{B} = \sum_{i} a_{i} b_{i} \).
The spaces you type in a formula are
ignored.  Remember that a letter like
       $x$                   % $ ... $  and  \( ... \)  are equivalent
is a formula when it denotes a mathematical
symbol, and it should be typed as one.

\section{Displayed Text}

Text is displayed by indenting it from the left
margin.  Quotations are commonly displayed.  There
are short quotations
\begin{quote}
   This is a short quotation.  It consists of a
   single paragraph of text.  See how it is formatted.
\end{quote}
and longer ones.
\begin{quotation}
   This is a longer quotation.  It consists of two
   paragraphs of text, neither of which are
   particularly interesting.

   This is the second paragraph of the quotation.  It
   is just as dull as the first paragraph.
\end{quotation}
Another frequently-displayed structure is a list.
The following is an example of an \emph{itemized}
list.
\begin{itemize}
   \item This is the first item of an itemized list.
         Each item in the list is marked with a ``tick''.
         You don't have to worry about what kind of tick
         mark is used.

   \item This is the second item of the list.  It
         contains another list nested inside it.  The inner
         list is an \emph{enumerated} list.
         \begin{enumerate}
            \item This is the first item of an enumerated
                  list that is nested within the itemized list.

            \item This is the second item of the inner list.
                  \LaTeX\ allows you to nest lists deeper than
                  you really should.
         \end{enumerate}
         This is the rest of the second item of the outer
         list.  It is no more interesting than any other
         part of the item.
   \item This is the third item of the list.
\end{itemize}
You can even display poetry.
\begin{verse}
   There is an environment
    for verse \\             % The \\ command separates lines
   Whose features some poets % within a stanza.
   will curse.

                             % One or more blank lines separate stanzas.

   For instead of making\\
   Them do \emph{all} line breaking, \\
   It allows them to put too many words on a line when they'd rather be
   forced to be terse.
\end{verse}

Mathematical formulas may also be displayed.  A
displayed formula
is
one-line long; multiline
formulas require special formatting instructions.
   \[  \ip{\Gamma}{\psi'} = x'' + y^{2} + z_{i}^{n}\]
Don't start a paragraph with a displayed equation,
nor make one a paragraph by itself.

\end{document}               % End of document.

\section{Bench Section 5}
% synthetic bench section 5
% This is a sample LaTeX input file.  (Version of 12 August 2004.)
%
% A '%' character causes TeX to ignore all remaining text on the line,
% and is used for comments like this one.

\documentclass{article}      % Specifies the document class

                             % The preamble begins here.
\title{An Example Document}  % Declares the document's title.
\author{Leslie Lamport}      % Declares the author's name.
\date{January 21, 1994}      % Deleting this command produces today's date.

\newcommand{\ip}[2]{(#1, #2)}
                             % Defines \ip{arg1}{arg2} to mean
                             % (arg1, arg2).

%\newcommand{\ip}[2]{\langle #1 | #2\rangle}
                             % This is an alternative definition of
                             % \ip that is commented out.

\begin{document}             % End of preamble and beginning of text.

\maketitle                   % Produces the title.

This is an example input file.  Comparing it with
the output it generates can show you how to
produce a simple document of your own.

\section{Ordinary Text}      % Produces section heading.  Lower-level
                             % sections are begun with similar
                             % \subsection and \subsubsection commands.

The ends  of words and sentences are marked
  by   spaces. It  doesn't matter how many
spaces    you type; one is as good as 100.  The
end of   a line counts as a space.

One   or more   blank lines denote the  end
of  a paragraph.

Since any number of consecutive spaces are treated
like a single one, the formatting of the input
file makes no difference to
      \LaTeX,                % The \LaTeX command generates the LaTeX logo.
but it makes a difference to you.  When you use
\LaTeX, making your input file as easy to read
as possible will be a great help as you write
your document and when you change it.  This sample
file shows how you can add comments to your own input
file.

Because printing is different from typewriting,
there are a number of things that you have to do
differently when preparing an input file than if
you were just typing the document directly.
Quotation marks like
       ``this''
have to be handled specially, as do quotes within
quotes:
       ``\,`this'            % \, separates the double and single quote.
        is what I just
        wrote, not  `that'\,''.

Dashes come in three sizes: an
       intra-word
dash, a medium dash for number ranges like
       1--2,
and a punctuation
       dash---like
this.

A sentence-ending space should be larger than the
space between words within a sentence.  You
sometimes have to type special commands in
conjunction with punctuation characters to get
this right, as in the following sentence.
       Gnats, gnus, etc.\ all  % `\ ' makes an inter-word space.
       begin with G\@.         % \@ marks end-of-sentence punctuation.
You should check the spaces after periods when
reading your output to make sure you haven't
forgotten any special cases.  Generating an
ellipsis
       \ldots\               % `\ ' is needed after `\ldots' because TeX
                             % ignores spaces after command names like \ldots
                             % made from \ + letters.
                             %
                             % Note how a `%' character causes TeX to ignore
                             % the end of the input line, so these blank lines
                             % do not start a new paragraph.
                             %
with the right spacing around the periods requires
a special command.

\LaTeX\ interprets some common characters as
commands, so you must type special commands to
generate them.  These characters include the
following:
       \$ \& \% \# \{ and \}.

In printing, text is usually emphasized with an
       \emph{italic}
type style.

\begin{em}
   A long segment of text can also be emphasized
   in this way.  Text within such a segment can be
   given \emph{additional} emphasis.
\end{em}

It is sometimes necessary to prevent \LaTeX\ from
breaking a line where it might otherwise do so.
This may be at a space, as between the ``Mr.''\ and
``Jones'' in
       ``Mr.~Jones'',        % ~ produces an unbreakable interword space.
or within a word---especially when the word is a
symbol like
       \mbox{\emph{itemnum}}
that makes little sense when hyphenated across
lines.

Footnotes\footnote{This is an example of a footnote.}
pose no problem.

\LaTeX\ is good at typesetting mathematical formulas
like
       \( x-3y + z = 7 \)
or
       \( a_{1} > x^{2n} + y^{2n} > x' \)
or
       \( \ip{A}{B} = \sum_{i} a_{i} b_{i} \).
The spaces you type in a formula are
ignored.  Remember that a letter like
       $x$                   % $ ... $  and  \( ... \)  are equivalent
is a formula when it denotes a mathematical
symbol, and it should be typed as one.

\section{Displayed Text}

Text is displayed by indenting it from the left
margin.  Quotations are commonly displayed.  There
are short quotations
\begin{quote}
   This is a short quotation.  It consists of a
   single paragraph of text.  See how it is formatted.
\end{quote}
and longer ones.
\begin{quotation}
   This is a longer quotation.  It consists of two
   paragraphs of text, neither of which are
   particularly interesting.

   This is the second paragraph of the quotation.  It
   is just as dull as the first paragraph.
\end{quotation}
Another frequently-displayed structure is a list.
The following is an example of an \emph{itemized}
list.
\begin{itemize}
   \item This is the first item of an itemized list.
         Each item in the list is marked with a ``tick''.
         You don't have to worry about what kind of tick
         mark is used.

   \item This is the second item of the list.  It
         contains another list nested inside it.  The inner
         list is an \emph{enumerated} list.
         \begin{enumerate}
            \item This is the first item of an enumerated
                  list that is nested within the itemized list.

            \item This is the second item of the inner list.
                  \LaTeX\ allows you to nest lists deeper than
                  you really should.
         \end{enumerate}
         This is the rest of the second item of the outer
         list.  It is no more interesting than any other
         part of the item.
   \item This is the third item of the list.
\end{itemize}
You can even display poetry.
\begin{verse}
   There is an environment
    for verse \\             % The \\ command separates lines
   Whose features some poets % within a stanza.
   will curse.

                             % One or more blank lines separate stanzas.

   For instead of making\\
   Them do \emph{all} line breaking, \\
   It allows them to put too many words on a line when they'd rather be
   forced to be terse.
\end{verse}

Mathematical formulas may also be displayed.  A
displayed formula
is
one-line long; multiline
formulas require special formatting instructions.
   \[  \ip{\Gamma}{\psi'} = x'' + y^{2} + z_{i}^{n}\]
Don't start a paragraph with a displayed equation,
nor make one a paragraph by itself.

\end{document}               % End of document.

\section{Bench Section 6}
% synthetic bench section 6
% This is a sample LaTeX input file.  (Version of 12 August 2004.)
%
% A '%' character causes TeX to ignore all remaining text on the line,
% and is used for comments like this one.

\documentclass{article}      % Specifies the document class

                             % The preamble begins here.
\title{An Example Document}  % Declares the document's title.
\author{Leslie Lamport}      % Declares the author's name.
\date{January 21, 1994}      % Deleting this command produces today's date.

\newcommand{\ip}[2]{(#1, #2)}
                             % Defines \ip{arg1}{arg2} to mean
                             % (arg1, arg2).

%\newcommand{\ip}[2]{\langle #1 | #2\rangle}
                             % This is an alternative definition of
                             % \ip that is commented out.

\begin{document}             % End of preamble and beginning of text.

\maketitle                   % Produces the title.

This is an example input file.  Comparing it with
the output it generates can show you how to
produce a simple document of your own.

\section{Ordinary Text}      % Produces section heading.  Lower-level
                             % sections are begun with similar
                             % \subsection and \subsubsection commands.

The ends  of words and sentences are marked
  by   spaces. It  doesn't matter how many
spaces    you type; one is as good as 100.  The
end of   a line counts as a space.

One   or more   blank lines denote the  end
of  a paragraph.

Since any number of consecutive spaces are treated
like a single one, the formatting of the input
file makes no difference to
      \LaTeX,                % The \LaTeX command generates the LaTeX logo.
but it makes a difference to you.  When you use
\LaTeX, making your input file as easy to read
as possible will be a great help as you write
your document and when you change it.  This sample
file shows how you can add comments to your own input
file.

Because printing is different from typewriting,
there are a number of things that you have to do
differently when preparing an input file than if
you were just typing the document directly.
Quotation marks like
       ``this''
have to be handled specially, as do quotes within
quotes:
       ``\,`this'            % \, separates the double and single quote.
        is what I just
        wrote, not  `that'\,''.

Dashes come in three sizes: an
       intra-word
dash, a medium dash for number ranges like
       1--2,
and a punctuation
       dash---like
this.

A sentence-ending space should be larger than the
space between words within a sentence.  You
sometimes have to type special commands in
conjunction with punctuation characters to get
this right, as in the following sentence.
       Gnats, gnus, etc.\ all  % `\ ' makes an inter-word space.
       begin with G\@.         % \@ marks end-of-sentence punctuation.
You should check the spaces after periods when
reading your output to make sure you haven't
forgotten any special cases.  Generating an
ellipsis
       \ldots\               % `\ ' is needed after `\ldots' because TeX
                             % ignores spaces after command names like \ldots
                             % made from \ + letters.
                             %
                             % Note how a `%' character causes TeX to ignore
                             % the end of the input line, so these blank lines
                             % do not start a new paragraph.
                             %
with the right spacing around the periods requires
a special command.

\LaTeX\ interprets some common characters as
commands, so you must type special commands to
generate them.  These characters include the
following:
       \$ \& \% \# \{ and \}.

In printing, text is usually emphasized with an
       \emph{italic}
type style.

\begin{em}
   A long segment of text can also be emphasized
   in this way.  Text within such a segment can be
   given \emph{additional} emphasis.
\end{em}

It is sometimes necessary to prevent \LaTeX\ from
breaking a line where it might otherwise do so.
This may be at a space, as between the ``Mr.''\ and
``Jones'' in
       ``Mr.~Jones'',        % ~ produces an unbreakable interword space.
or within a word---especially when the word is a
symbol like
       \mbox{\emph{itemnum}}
that makes little sense when hyphenated across
lines.

Footnotes\footnote{This is an example of a footnote.}
pose no problem.

\LaTeX\ is good at typesetting mathematical formulas
like
       \( x-3y + z = 7 \)
or
       \( a_{1} > x^{2n} + y^{2n} > x' \)
or
       \( \ip{A}{B} = \sum_{i} a_{i} b_{i} \).
The spaces you type in a formula are
ignored.  Remember that a letter like
       $x$                   % $ ... $  and  \( ... \)  are equivalent
is a formula when it denotes a mathematical
symbol, and it should be typed as one.

\section{Displayed Text}

Text is displayed by indenting it from the left
margin.  Quotations are commonly displayed.  There
are short quotations
\begin{quote}
   This is a short quotation.  It consists of a
   single paragraph of text.  See how it is formatted.
\end{quote}
and longer ones.
\begin{quotation}
   This is a longer quotation.  It consists of two
   paragraphs of text, neither of which are
   particularly interesting.

   This is the second paragraph of the quotation.  It
   is just as dull as the first paragraph.
\end{quotation}
Another frequently-displayed structure is a list.
The following is an example of an \emph{itemized}
list.
\begin{itemize}
   \item This is the first item of an itemized list.
         Each item in the list is marked with a ``tick''.
         You don't have to worry about what kind of tick
         mark is used.

   \item This is the second item of the list.  It
         contains another list nested inside it.  The inner
         list is an \emph{enumerated} list.
         \begin{enumerate}
            \item This is the first item of an enumerated
                  list that is nested within the itemized list.

            \item This is the second item of the inner list.
                  \LaTeX\ allows you to nest lists deeper than
                  you really should.
         \end{enumerate}
         This is the rest of the second item of the outer
         list.  It is no more interesting than any other
         part of the item.
   \item This is the third item of the list.
\end{itemize}
You can even display poetry.
\begin{verse}
   There is an environment
    for verse \\             % The \\ command separates lines
   Whose features some poets % within a stanza.
   will curse.

                             % One or more blank lines separate stanzas.

   For instead of making\\
   Them do \emph{all} line breaking, \\
   It allows them to put too many words on a line when they'd rather be
   forced to be terse.
\end{verse}

Mathematical formulas may also be displayed.  A
displayed formula
is
one-line long; multiline
formulas require special formatting instructions.
   \[  \ip{\Gamma}{\psi'} = x'' + y^{2} + z_{i}^{n}\]
Don't start a paragraph with a displayed equation,
nor make one a paragraph by itself.

\end{document}               % End of document.

\section{Bench Section 7}
% synthetic bench section 7
% This is a sample LaTeX input file.  (Version of 12 August 2004.)
%
% A '%' character causes TeX to ignore all remaining text on the line,
% and is used for comments like this one.

\documentclass{article}      % Specifies the document class

                             % The preamble begins here.
\title{An Example Document}  % Declares the document's title.
\author{Leslie Lamport}      % Declares the author's name.
\date{January 21, 1994}      % Deleting this command produces today's date.

\newcommand{\ip}[2]{(#1, #2)}
                             % Defines \ip{arg1}{arg2} to mean
                             % (arg1, arg2).

%\newcommand{\ip}[2]{\langle #1 | #2\rangle}
                             % This is an alternative definition of
                             % \ip that is commented out.

\begin{document}             % End of preamble and beginning of text.

\maketitle                   % Produces the title.

This is an example input file.  Comparing it with
the output it generates can show you how to
produce a simple document of your own.

\section{Ordinary Text}      % Produces section heading.  Lower-level
                             % sections are begun with similar
                             % \subsection and \subsubsection commands.

The ends  of words and sentences are marked
  by   spaces. It  doesn't matter how many
spaces    you type; one is as good as 100.  The
end of   a line counts as a space.

One   or more   blank lines denote the  end
of  a paragraph.

Since any number of consecutive spaces are treated
like a single one, the formatting of the input
file makes no difference to
      \LaTeX,                % The \LaTeX command generates the LaTeX logo.
but it makes a difference to you.  When you use
\LaTeX, making your input file as easy to read
as possible will be a great help as you write
your document and when you change it.  This sample
file shows how you can add comments to your own input
file.

Because printing is different from typewriting,
there are a number of things that you have to do
differently when preparing an input file than if
you were just typing the document directly.
Quotation marks like
       ``this''
have to be handled specially, as do quotes within
quotes:
       ``\,`this'            % \, separates the double and single quote.
        is what I just
        wrote, not  `that'\,''.

Dashes come in three sizes: an
       intra-word
dash, a medium dash for number ranges like
       1--2,
and a punctuation
       dash---like
this.

A sentence-ending space should be larger than the
space between words within a sentence.  You
sometimes have to type special commands in
conjunction with punctuation characters to get
this right, as in the following sentence.
       Gnats, gnus, etc.\ all  % `\ ' makes an inter-word space.
       begin with G\@.         % \@ marks end-of-sentence punctuation.
You should check the spaces after periods when
reading your output to make sure you haven't
forgotten any special cases.  Generating an
ellipsis
       \ldots\               % `\ ' is needed after `\ldots' because TeX
                             % ignores spaces after command names like \ldots
                             % made from \ + letters.
                             %
                             % Note how a `%' character causes TeX to ignore
                             % the end of the input line, so these blank lines
                             % do not start a new paragraph.
                             %
with the right spacing around the periods requires
a special command.

\LaTeX\ interprets some common characters as
commands, so you must type special commands to
generate them.  These characters include the
following:
       \$ \& \% \# \{ and \}.

In printing, text is usually emphasized with an
       \emph{italic}
type style.

\begin{em}
   A long segment of text can also be emphasized
   in this way.  Text within such a segment can be
   given \emph{additional} emphasis.
\end{em}

It is sometimes necessary to prevent \LaTeX\ from
breaking a line where it might otherwise do so.
This may be at a space, as between the ``Mr.''\ and
``Jones'' in
       ``Mr.~Jones'',        % ~ produces an unbreakable interword space.
or within a word---especially when the word is a
symbol like
       \mbox{\emph{itemnum}}
that makes little sense when hyphenated across
lines.

Footnotes\footnote{This is an example of a footnote.}
pose no problem.

\LaTeX\ is good at typesetting mathematical formulas
like
       \( x-3y + z = 7 \)
or
       \( a_{1} > x^{2n} + y^{2n} > x' \)
or
       \( \ip{A}{B} = \sum_{i} a_{i} b_{i} \).
The spaces you type in a formula are
ignored.  Remember that a letter like
       $x$                   % $ ... $  and  \( ... \)  are equivalent
is a formula when it denotes a mathematical
symbol, and it should be typed as one.

\section{Displayed Text}

Text is displayed by indenting it from the left
margin.  Quotations are commonly displayed.  There
are short quotations
\begin{quote}
   This is a short quotation.  It consists of a
   single paragraph of text.  See how it is formatted.
\end{quote}
and longer ones.
\begin{quotation}
   This is a longer quotation.  It consists of two
   paragraphs of text, neither of which are
   particularly interesting.

   This is the second paragraph of the quotation.  It
   is just as dull as the first paragraph.
\end{quotation}
Another frequently-displayed structure is a list.
The following is an example of an \emph{itemized}
list.
\begin{itemize}
   \item This is the first item of an itemized list.
         Each item in the list is marked with a ``tick''.
         You don't have to worry about what kind of tick
         mark is used.

   \item This is the second item of the list.  It
         contains another list nested inside it.  The inner
         list is an \emph{enumerated} list.
         \begin{enumerate}
            \item This is the first item of an enumerated
                  list that is nested within the itemized list.

            \item This is the second item of the inner list.
                  \LaTeX\ allows you to nest lists deeper than
                  you really should.
         \end{enumerate}
         This is the rest of the second item of the outer
         list.  It is no more interesting than any other
         part of the item.
   \item This is the third item of the list.
\end{itemize}
You can even display poetry.
\begin{verse}
   There is an environment
    for verse \\             % The \\ command separates lines
   Whose features some poets % within a stanza.
   will curse.

                             % One or more blank lines separate stanzas.

   For instead of making\\
   Them do \emph{all} line breaking, \\
   It allows them to put too many words on a line when they'd rather be
   forced to be terse.
\end{verse}

Mathematical formulas may also be displayed.  A
displayed formula
is
one-line long; multiline
formulas require special formatting instructions.
   \[  \ip{\Gamma}{\psi'} = x'' + y^{2} + z_{i}^{n}\]
Don't start a paragraph with a displayed equation,
nor make one a paragraph by itself.

\end{document}               % End of document.

\section{Bench Section 8}
% synthetic bench section 8
% This is a sample LaTeX input file.  (Version of 12 August 2004.)
%
% A '%' character causes TeX to ignore all remaining text on the line,
% and is used for comments like this one.

\documentclass{article}      % Specifies the document class

                             % The preamble begins here.
\title{An Example Document}  % Declares the document's title.
\author{Leslie Lamport}      % Declares the author's name.
\date{January 21, 1994}      % Deleting this command produces today's date.

\newcommand{\ip}[2]{(#1, #2)}
                             % Defines \ip{arg1}{arg2} to mean
                             % (arg1, arg2).

%\newcommand{\ip}[2]{\langle #1 | #2\rangle}
                             % This is an alternative definition of
                             % \ip that is commented out.

\begin{document}             % End of preamble and beginning of text.

\maketitle                   % Produces the title.

This is an example input file.  Comparing it with
the output it generates can show you how to
produce a simple document of your own.

\section{Ordinary Text}      % Produces section heading.  Lower-level
                             % sections are begun with similar
                             % \subsection and \subsubsection commands.

The ends  of words and sentences are marked
  by   spaces. It  doesn't matter how many
spaces    you type; one is as good as 100.  The
end of   a line counts as a space.

One   or more   blank lines denote the  end
of  a paragraph.

Since any number of consecutive spaces are treated
like a single one, the formatting of the input
file makes no difference to
      \LaTeX,                % The \LaTeX command generates the LaTeX logo.
but it makes a difference to you.  When you use
\LaTeX, making your input file as easy to read
as possible will be a great help as you write
your document and when you change it.  This sample
file shows how you can add comments to your own input
file.

Because printing is different from typewriting,
there are a number of things that you have to do
differently when preparing an input file than if
you were just typing the document directly.
Quotation marks like
       ``this''
have to be handled specially, as do quotes within
quotes:
       ``\,`this'            % \, separates the double and single quote.
        is what I just
        wrote, not  `that'\,''.

Dashes come in three sizes: an
       intra-word
dash, a medium dash for number ranges like
       1--2,
and a punctuation
       dash---like
this.

A sentence-ending space should be larger than the
space between words within a sentence.  You
sometimes have to type special commands in
conjunction with punctuation characters to get
this right, as in the following sentence.
       Gnats, gnus, etc.\ all  % `\ ' makes an inter-word space.
       begin with G\@.         % \@ marks end-of-sentence punctuation.
You should check the spaces after periods when
reading your output to make sure you haven't
forgotten any special cases.  Generating an
ellipsis
       \ldots\               % `\ ' is needed after `\ldots' because TeX
                             % ignores spaces after command names like \ldots
                             % made from \ + letters.
                             %
                             % Note how a `%' character causes TeX to ignore
                             % the end of the input line, so these blank lines
                             % do not start a new paragraph.
                             %
with the right spacing around the periods requires
a special command.

\LaTeX\ interprets some common characters as
commands, so you must type special commands to
generate them.  These characters include the
following:
       \$ \& \% \# \{ and \}.

In printing, text is usually emphasized with an
       \emph{italic}
type style.

\begin{em}
   A long segment of text can also be emphasized
   in this way.  Text within such a segment can be
   given \emph{additional} emphasis.
\end{em}

It is sometimes necessary to prevent \LaTeX\ from
breaking a line where it might otherwise do so.
This may be at a space, as between the ``Mr.''\ and
``Jones'' in
       ``Mr.~Jones'',        % ~ produces an unbreakable interword space.
or within a word---especially when the word is a
symbol like
       \mbox{\emph{itemnum}}
that makes little sense when hyphenated across
lines.

Footnotes\footnote{This is an example of a footnote.}
pose no problem.

\LaTeX\ is good at typesetting mathematical formulas
like
       \( x-3y + z = 7 \)
or
       \( a_{1} > x^{2n} + y^{2n} > x' \)
or
       \( \ip{A}{B} = \sum_{i} a_{i} b_{i} \).
The spaces you type in a formula are
ignored.  Remember that a letter like
       $x$                   % $ ... $  and  \( ... \)  are equivalent
is a formula when it denotes a mathematical
symbol, and it should be typed as one.

\section{Displayed Text}

Text is displayed by indenting it from the left
margin.  Quotations are commonly displayed.  There
are short quotations
\begin{quote}
   This is a short quotation.  It consists of a
   single paragraph of text.  See how it is formatted.
\end{quote}
and longer ones.
\begin{quotation}
   This is a longer quotation.  It consists of two
   paragraphs of text, neither of which are
   particularly interesting.

   This is the second paragraph of the quotation.  It
   is just as dull as the first paragraph.
\end{quotation}
Another frequently-displayed structure is a list.
The following is an example of an \emph{itemized}
list.
\begin{itemize}
   \item This is the first item of an itemized list.
         Each item in the list is marked with a ``tick''.
         You don't have to worry about what kind of tick
         mark is used.

   \item This is the second item of the list.  It
         contains another list nested inside it.  The inner
         list is an \emph{enumerated} list.
         \begin{enumerate}
            \item This is the first item of an enumerated
                  list that is nested within the itemized list.

            \item This is the second item of the inner list.
                  \LaTeX\ allows you to nest lists deeper than
                  you really should.
         \end{enumerate}
         This is the rest of the second item of the outer
         list.  It is no more interesting than any other
         part of the item.
   \item This is the third item of the list.
\end{itemize}
You can even display poetry.
\begin{verse}
   There is an environment
    for verse \\             % The \\ command separates lines
   Whose features some poets % within a stanza.
   will curse.

                             % One or more blank lines separate stanzas.

   For instead of making\\
   Them do \emph{all} line breaking, \\
   It allows them to put too many words on a line when they'd rather be
   forced to be terse.
\end{verse}

Mathematical formulas may also be displayed.  A
displayed formula
is
one-line long; multiline
formulas require special formatting instructions.
   \[  \ip{\Gamma}{\psi'} = x'' + y^{2} + z_{i}^{n}\]
Don't start a paragraph with a displayed equation,
nor make one a paragraph by itself.

\end{document}               % End of document.

\section{Bench Section 9}
% synthetic bench section 9
% This is a sample LaTeX input file.  (Version of 12 August 2004.)
%
% A '%' character causes TeX to ignore all remaining text on the line,
% and is used for comments like this one.

\documentclass{article}      % Specifies the document class

                             % The preamble begins here.
\title{An Example Document}  % Declares the document's title.
\author{Leslie Lamport}      % Declares the author's name.
\date{January 21, 1994}      % Deleting this command produces today's date.

\newcommand{\ip}[2]{(#1, #2)}
                             % Defines \ip{arg1}{arg2} to mean
                             % (arg1, arg2).

%\newcommand{\ip}[2]{\langle #1 | #2\rangle}
                             % This is an alternative definition of
                             % \ip that is commented out.

\begin{document}             % End of preamble and beginning of text.

\maketitle                   % Produces the title.

This is an example input file.  Comparing it with
the output it generates can show you how to
produce a simple document of your own.

\section{Ordinary Text}      % Produces section heading.  Lower-level
                             % sections are begun with similar
                             % \subsection and \subsubsection commands.

The ends  of words and sentences are marked
  by   spaces. It  doesn't matter how many
spaces    you type; one is as good as 100.  The
end of   a line counts as a space.

One   or more   blank lines denote the  end
of  a paragraph.

Since any number of consecutive spaces are treated
like a single one, the formatting of the input
file makes no difference to
      \LaTeX,                % The \LaTeX command generates the LaTeX logo.
but it makes a difference to you.  When you use
\LaTeX, making your input file as easy to read
as possible will be a great help as you write
your document and when you change it.  This sample
file shows how you can add comments to your own input
file.

Because printing is different from typewriting,
there are a number of things that you have to do
differently when preparing an input file than if
you were just typing the document directly.
Quotation marks like
       ``this''
have to be handled specially, as do quotes within
quotes:
       ``\,`this'            % \, separates the double and single quote.
        is what I just
        wrote, not  `that'\,''.

Dashes come in three sizes: an
       intra-word
dash, a medium dash for number ranges like
       1--2,
and a punctuation
       dash---like
this.

A sentence-ending space should be larger than the
space between words within a sentence.  You
sometimes have to type special commands in
conjunction with punctuation characters to get
this right, as in the following sentence.
       Gnats, gnus, etc.\ all  % `\ ' makes an inter-word space.
       begin with G\@.         % \@ marks end-of-sentence punctuation.
You should check the spaces after periods when
reading your output to make sure you haven't
forgotten any special cases.  Generating an
ellipsis
       \ldots\               % `\ ' is needed after `\ldots' because TeX
                             % ignores spaces after command names like \ldots
                             % made from \ + letters.
                             %
                             % Note how a `%' character causes TeX to ignore
                             % the end of the input line, so these blank lines
                             % do not start a new paragraph.
                             %
with the right spacing around the periods requires
a special command.

\LaTeX\ interprets some common characters as
commands, so you must type special commands to
generate them.  These characters include the
following:
       \$ \& \% \# \{ and \}.

In printing, text is usually emphasized with an
       \emph{italic}
type style.

\begin{em}
   A long segment of text can also be emphasized
   in this way.  Text within such a segment can be
   given \emph{additional} emphasis.
\end{em}

It is sometimes necessary to prevent \LaTeX\ from
breaking a line where it might otherwise do so.
This may be at a space, as between the ``Mr.''\ and
``Jones'' in
       ``Mr.~Jones'',        % ~ produces an unbreakable interword space.
or within a word---especially when the word is a
symbol like
       \mbox{\emph{itemnum}}
that makes little sense when hyphenated across
lines.

Footnotes\footnote{This is an example of a footnote.}
pose no problem.

\LaTeX\ is good at typesetting mathematical formulas
like
       \( x-3y + z = 7 \)
or
       \( a_{1} > x^{2n} + y^{2n} > x' \)
or
       \( \ip{A}{B} = \sum_{i} a_{i} b_{i} \).
The spaces you type in a formula are
ignored.  Remember that a letter like
       $x$                   % $ ... $  and  \( ... \)  are equivalent
is a formula when it denotes a mathematical
symbol, and it should be typed as one.

\section{Displayed Text}

Text is displayed by indenting it from the left
margin.  Quotations are commonly displayed.  There
are short quotations
\begin{quote}
   This is a short quotation.  It consists of a
   single paragraph of text.  See how it is formatted.
\end{quote}
and longer ones.
\begin{quotation}
   This is a longer quotation.  It consists of two
   paragraphs of text, neither of which are
   particularly interesting.

   This is the second paragraph of the quotation.  It
   is just as dull as the first paragraph.
\end{quotation}
Another frequently-displayed structure is a list.
The following is an example of an \emph{itemized}
list.
\begin{itemize}
   \item This is the first item of an itemized list.
         Each item in the list is marked with a ``tick''.
         You don't have to worry about what kind of tick
         mark is used.

   \item This is the second item of the list.  It
         contains another list nested inside it.  The inner
         list is an \emph{enumerated} list.
         \begin{enumerate}
            \item This is the first item of an enumerated
                  list that is nested within the itemized list.

            \item This is the second item of the inner list.
                  \LaTeX\ allows you to nest lists deeper than
                  you really should.
         \end{enumerate}
         This is the rest of the second item of the outer
         list.  It is no more interesting than any other
         part of the item.
   \item This is the third item of the list.
\end{itemize}
You can even display poetry.
\begin{verse}
   There is an environment
    for verse \\             % The \\ command separates lines
   Whose features some poets % within a stanza.
   will curse.

                             % One or more blank lines separate stanzas.

   For instead of making\\
   Them do \emph{all} line breaking, \\
   It allows them to put too many words on a line when they'd rather be
   forced to be terse.
\end{verse}

Mathematical formulas may also be displayed.  A
displayed formula
is
one-line long; multiline
formulas require special formatting instructions.
   \[  \ip{\Gamma}{\psi'} = x'' + y^{2} + z_{i}^{n}\]
Don't start a paragraph with a displayed equation,
nor make one a paragraph by itself.

\end{document}               % End of document.

\section{Bench Section 10}
% synthetic bench section 10
% This is a sample LaTeX input file.  (Version of 12 August 2004.)
%
% A '%' character causes TeX to ignore all remaining text on the line,
% and is used for comments like this one.

\documentclass{article}      % Specifies the document class

                             % The preamble begins here.
\title{An Example Document}  % Declares the document's title.
\author{Leslie Lamport}      % Declares the author's name.
\date{January 21, 1994}      % Deleting this command produces today's date.

\newcommand{\ip}[2]{(#1, #2)}
                             % Defines \ip{arg1}{arg2} to mean
                             % (arg1, arg2).

%\newcommand{\ip}[2]{\langle #1 | #2\rangle}
                             % This is an alternative definition of
                             % \ip that is commented out.

\begin{document}             % End of preamble and beginning of text.

\maketitle                   % Produces the title.

This is an example input file.  Comparing it with
the output it generates can show you how to
produce a simple document of your own.

\section{Ordinary Text}      % Produces section heading.  Lower-level
                             % sections are begun with similar
                             % \subsection and \subsubsection commands.

The ends  of words and sentences are marked
  by   spaces. It  doesn't matter how many
spaces    you type; one is as good as 100.  The
end of   a line counts as a space.

One   or more   blank lines denote the  end
of  a paragraph.

Since any number of consecutive spaces are treated
like a single one, the formatting of the input
file makes no difference to
      \LaTeX,                % The \LaTeX command generates the LaTeX logo.
but it makes a difference to you.  When you use
\LaTeX, making your input file as easy to read
as possible will be a great help as you write
your document and when you change it.  This sample
file shows how you can add comments to your own input
file.

Because printing is different from typewriting,
there are a number of things that you have to do
differently when preparing an input file than if
you were just typing the document directly.
Quotation marks like
       ``this''
have to be handled specially, as do quotes within
quotes:
       ``\,`this'            % \, separates the double and single quote.
        is what I just
        wrote, not  `that'\,''.

Dashes come in three sizes: an
       intra-word
dash, a medium dash for number ranges like
       1--2,
and a punctuation
       dash---like
this.

A sentence-ending space should be larger than the
space between words within a sentence.  You
sometimes have to type special commands in
conjunction with punctuation characters to get
this right, as in the following sentence.
       Gnats, gnus, etc.\ all  % `\ ' makes an inter-word space.
       begin with G\@.         % \@ marks end-of-sentence punctuation.
You should check the spaces after periods when
reading your output to make sure you haven't
forgotten any special cases.  Generating an
ellipsis
       \ldots\               % `\ ' is needed after `\ldots' because TeX
                             % ignores spaces after command names like \ldots
                             % made from \ + letters.
                             %
                             % Note how a `%' character causes TeX to ignore
                             % the end of the input line, so these blank lines
                             % do not start a new paragraph.
                             %
with the right spacing around the periods requires
a special command.

\LaTeX\ interprets some common characters as
commands, so you must type special commands to
generate them.  These characters include the
following:
       \$ \& \% \# \{ and \}.

In printing, text is usually emphasized with an
       \emph{italic}
type style.

\begin{em}
   A long segment of text can also be emphasized
   in this way.  Text within such a segment can be
   given \emph{additional} emphasis.
\end{em}

It is sometimes necessary to prevent \LaTeX\ from
breaking a line where it might otherwise do so.
This may be at a space, as between the ``Mr.''\ and
``Jones'' in
       ``Mr.~Jones'',        % ~ produces an unbreakable interword space.
or within a word---especially when the word is a
symbol like
       \mbox{\emph{itemnum}}
that makes little sense when hyphenated across
lines.

Footnotes\footnote{This is an example of a footnote.}
pose no problem.

\LaTeX\ is good at typesetting mathematical formulas
like
       \( x-3y + z = 7 \)
or
       \( a_{1} > x^{2n} + y^{2n} > x' \)
or
       \( \ip{A}{B} = \sum_{i} a_{i} b_{i} \).
The spaces you type in a formula are
ignored.  Remember that a letter like
       $x$                   % $ ... $  and  \( ... \)  are equivalent
is a formula when it denotes a mathematical
symbol, and it should be typed as one.

\section{Displayed Text}

Text is displayed by indenting it from the left
margin.  Quotations are commonly displayed.  There
are short quotations
\begin{quote}
   This is a short quotation.  It consists of a
   single paragraph of text.  See how it is formatted.
\end{quote}
and longer ones.
\begin{quotation}
   This is a longer quotation.  It consists of two
   paragraphs of text, neither of which are
   particularly interesting.

   This is the second paragraph of the quotation.  It
   is just as dull as the first paragraph.
\end{quotation}
Another frequently-displayed structure is a list.
The following is an example of an \emph{itemized}
list.
\begin{itemize}
   \item This is the first item of an itemized list.
         Each item in the list is marked with a ``tick''.
         You don't have to worry about what kind of tick
         mark is used.

   \item This is the second item of the list.  It
         contains another list nested inside it.  The inner
         list is an \emph{enumerated} list.
         \begin{enumerate}
            \item This is the first item of an enumerated
                  list that is nested within the itemized list.

            \item This is the second item of the inner list.
                  \LaTeX\ allows you to nest lists deeper than
                  you really should.
         \end{enumerate}
         This is the rest of the second item of the outer
         list.  It is no more interesting than any other
         part of the item.
   \item This is the third item of the list.
\end{itemize}
You can even display poetry.
\begin{verse}
   There is an environment
    for verse \\             % The \\ command separates lines
   Whose features some poets % within a stanza.
   will curse.

                             % One or more blank lines separate stanzas.

   For instead of making\\
   Them do \emph{all} line breaking, \\
   It allows them to put too many words on a line when they'd rather be
   forced to be terse.
\end{verse}

Mathematical formulas may also be displayed.  A
displayed formula
is
one-line long; multiline
formulas require special formatting instructions.
   \[  \ip{\Gamma}{\psi'} = x'' + y^{2} + z_{i}^{n}\]
Don't start a paragraph with a displayed equation,
nor make one a paragraph by itself.

\end{document}               % End of document.

\section{Bench Section 11}
% synthetic bench section 11
% This is a sample LaTeX input file.  (Version of 12 August 2004.)
%
% A '%' character causes TeX to ignore all remaining text on the line,
% and is used for comments like this one.

\documentclass{article}      % Specifies the document class

                             % The preamble begins here.
\title{An Example Document}  % Declares the document's title.
\author{Leslie Lamport}      % Declares the author's name.
\date{January 21, 1994}      % Deleting this command produces today's date.

\newcommand{\ip}[2]{(#1, #2)}
                             % Defines \ip{arg1}{arg2} to mean
                             % (arg1, arg2).

%\newcommand{\ip}[2]{\langle #1 | #2\rangle}
                             % This is an alternative definition of
                             % \ip that is commented out.

\begin{document}             % End of preamble and beginning of text.

\maketitle                   % Produces the title.

This is an example input file.  Comparing it with
the output it generates can show you how to
produce a simple document of your own.

\section{Ordinary Text}      % Produces section heading.  Lower-level
                             % sections are begun with similar
                             % \subsection and \subsubsection commands.

The ends  of words and sentences are marked
  by   spaces. It  doesn't matter how many
spaces    you type; one is as good as 100.  The
end of   a line counts as a space.

One   or more   blank lines denote the  end
of  a paragraph.

Since any number of consecutive spaces are treated
like a single one, the formatting of the input
file makes no difference to
      \LaTeX,                % The \LaTeX command generates the LaTeX logo.
but it makes a difference to you.  When you use
\LaTeX, making your input file as easy to read
as possible will be a great help as you write
your document and when you change it.  This sample
file shows how you can add comments to your own input
file.

Because printing is different from typewriting,
there are a number of things that you have to do
differently when preparing an input file than if
you were just typing the document directly.
Quotation marks like
       ``this''
have to be handled specially, as do quotes within
quotes:
       ``\,`this'            % \, separates the double and single quote.
        is what I just
        wrote, not  `that'\,''.

Dashes come in three sizes: an
       intra-word
dash, a medium dash for number ranges like
       1--2,
and a punctuation
       dash---like
this.

A sentence-ending space should be larger than the
space between words within a sentence.  You
sometimes have to type special commands in
conjunction with punctuation characters to get
this right, as in the following sentence.
       Gnats, gnus, etc.\ all  % `\ ' makes an inter-word space.
       begin with G\@.         % \@ marks end-of-sentence punctuation.
You should check the spaces after periods when
reading your output to make sure you haven't
forgotten any special cases.  Generating an
ellipsis
       \ldots\               % `\ ' is needed after `\ldots' because TeX
                             % ignores spaces after command names like \ldots
                             % made from \ + letters.
                             %
                             % Note how a `%' character causes TeX to ignore
                             % the end of the input line, so these blank lines
                             % do not start a new paragraph.
                             %
with the right spacing around the periods requires
a special command.

\LaTeX\ interprets some common characters as
commands, so you must type special commands to
generate them.  These characters include the
following:
       \$ \& \% \# \{ and \}.

In printing, text is usually emphasized with an
       \emph{italic}
type style.

\begin{em}
   A long segment of text can also be emphasized
   in this way.  Text within such a segment can be
   given \emph{additional} emphasis.
\end{em}

It is sometimes necessary to prevent \LaTeX\ from
breaking a line where it might otherwise do so.
This may be at a space, as between the ``Mr.''\ and
``Jones'' in
       ``Mr.~Jones'',        % ~ produces an unbreakable interword space.
or within a word---especially when the word is a
symbol like
       \mbox{\emph{itemnum}}
that makes little sense when hyphenated across
lines.

Footnotes\footnote{This is an example of a footnote.}
pose no problem.

\LaTeX\ is good at typesetting mathematical formulas
like
       \( x-3y + z = 7 \)
or
       \( a_{1} > x^{2n} + y^{2n} > x' \)
or
       \( \ip{A}{B} = \sum_{i} a_{i} b_{i} \).
The spaces you type in a formula are
ignored.  Remember that a letter like
       $x$                   % $ ... $  and  \( ... \)  are equivalent
is a formula when it denotes a mathematical
symbol, and it should be typed as one.

\section{Displayed Text}

Text is displayed by indenting it from the left
margin.  Quotations are commonly displayed.  There
are short quotations
\begin{quote}
   This is a short quotation.  It consists of a
   single paragraph of text.  See how it is formatted.
\end{quote}
and longer ones.
\begin{quotation}
   This is a longer quotation.  It consists of two
   paragraphs of text, neither of which are
   particularly interesting.

   This is the second paragraph of the quotation.  It
   is just as dull as the first paragraph.
\end{quotation}
Another frequently-displayed structure is a list.
The following is an example of an \emph{itemized}
list.
\begin{itemize}
   \item This is the first item of an itemized list.
         Each item in the list is marked with a ``tick''.
         You don't have to worry about what kind of tick
         mark is used.

   \item This is the second item of the list.  It
         contains another list nested inside it.  The inner
         list is an \emph{enumerated} list.
         \begin{enumerate}
            \item This is the first item of an enumerated
                  list that is nested within the itemized list.

            \item This is the second item of the inner list.
                  \LaTeX\ allows you to nest lists deeper than
                  you really should.
         \end{enumerate}
         This is the rest of the second item of the outer
         list.  It is no more interesting than any other
         part of the item.
   \item This is the third item of the list.
\end{itemize}
You can even display poetry.
\begin{verse}
   There is an environment
    for verse \\             % The \\ command separates lines
   Whose features some poets % within a stanza.
   will curse.

                             % One or more blank lines separate stanzas.

   For instead of making\\
   Them do \emph{all} line breaking, \\
   It allows them to put too many words on a line when they'd rather be
   forced to be terse.
\end{verse}

Mathematical formulas may also be displayed.  A
displayed formula
is
one-line long; multiline
formulas require special formatting instructions.
   \[  \ip{\Gamma}{\psi'} = x'' + y^{2} + z_{i}^{n}\]
Don't start a paragraph with a displayed equation,
nor make one a paragraph by itself.

\end{document}               % End of document.

\section{Bench Section 12}
% synthetic bench section 12
% This is a sample LaTeX input file.  (Version of 12 August 2004.)
%
% A '%' character causes TeX to ignore all remaining text on the line,
% and is used for comments like this one.

\documentclass{article}      % Specifies the document class

                             % The preamble begins here.
\title{An Example Document}  % Declares the document's title.
\author{Leslie Lamport}      % Declares the author's name.
\date{January 21, 1994}      % Deleting this command produces today's date.

\newcommand{\ip}[2]{(#1, #2)}
                             % Defines \ip{arg1}{arg2} to mean
                             % (arg1, arg2).

%\newcommand{\ip}[2]{\langle #1 | #2\rangle}
                             % This is an alternative definition of
                             % \ip that is commented out.

\begin{document}             % End of preamble and beginning of text.

\maketitle                   % Produces the title.

This is an example input file.  Comparing it with
the output it generates can show you how to
produce a simple document of your own.

\section{Ordinary Text}      % Produces section heading.  Lower-level
                             % sections are begun with similar
                             % \subsection and \subsubsection commands.

The ends  of words and sentences are marked
  by   spaces. It  doesn't matter how many
spaces    you type; one is as good as 100.  The
end of   a line counts as a space.

One   or more   blank lines denote the  end
of  a paragraph.

Since any number of consecutive spaces are treated
like a single one, the formatting of the input
file makes no difference to
      \LaTeX,                % The \LaTeX command generates the LaTeX logo.
but it makes a difference to you.  When you use
\LaTeX, making your input file as easy to read
as possible will be a great help as you write
your document and when you change it.  This sample
file shows how you can add comments to your own input
file.

Because printing is different from typewriting,
there are a number of things that you have to do
differently when preparing an input file than if
you were just typing the document directly.
Quotation marks like
       ``this''
have to be handled specially, as do quotes within
quotes:
       ``\,`this'            % \, separates the double and single quote.
        is what I just
        wrote, not  `that'\,''.

Dashes come in three sizes: an
       intra-word
dash, a medium dash for number ranges like
       1--2,
and a punctuation
       dash---like
this.

A sentence-ending space should be larger than the
space between words within a sentence.  You
sometimes have to type special commands in
conjunction with punctuation characters to get
this right, as in the following sentence.
       Gnats, gnus, etc.\ all  % `\ ' makes an inter-word space.
       begin with G\@.         % \@ marks end-of-sentence punctuation.
You should check the spaces after periods when
reading your output to make sure you haven't
forgotten any special cases.  Generating an
ellipsis
       \ldots\               % `\ ' is needed after `\ldots' because TeX
                             % ignores spaces after command names like \ldots
                             % made from \ + letters.
                             %
                             % Note how a `%' character causes TeX to ignore
                             % the end of the input line, so these blank lines
                             % do not start a new paragraph.
                             %
with the right spacing around the periods requires
a special command.

\LaTeX\ interprets some common characters as
commands, so you must type special commands to
generate them.  These characters include the
following:
       \$ \& \% \# \{ and \}.

In printing, text is usually emphasized with an
       \emph{italic}
type style.

\begin{em}
   A long segment of text can also be emphasized
   in this way.  Text within such a segment can be
   given \emph{additional} emphasis.
\end{em}

It is sometimes necessary to prevent \LaTeX\ from
breaking a line where it might otherwise do so.
This may be at a space, as between the ``Mr.''\ and
``Jones'' in
       ``Mr.~Jones'',        % ~ produces an unbreakable interword space.
or within a word---especially when the word is a
symbol like
       \mbox{\emph{itemnum}}
that makes little sense when hyphenated across
lines.

Footnotes\footnote{This is an example of a footnote.}
pose no problem.

\LaTeX\ is good at typesetting mathematical formulas
like
       \( x-3y + z = 7 \)
or
       \( a_{1} > x^{2n} + y^{2n} > x' \)
or
       \( \ip{A}{B} = \sum_{i} a_{i} b_{i} \).
The spaces you type in a formula are
ignored.  Remember that a letter like
       $x$                   % $ ... $  and  \( ... \)  are equivalent
is a formula when it denotes a mathematical
symbol, and it should be typed as one.

\section{Displayed Text}

Text is displayed by indenting it from the left
margin.  Quotations are commonly displayed.  There
are short quotations
\begin{quote}
   This is a short quotation.  It consists of a
   single paragraph of text.  See how it is formatted.
\end{quote}
and longer ones.
\begin{quotation}
   This is a longer quotation.  It consists of two
   paragraphs of text, neither of which are
   particularly interesting.

   This is the second paragraph of the quotation.  It
   is just as dull as the first paragraph.
\end{quotation}
Another frequently-displayed structure is a list.
The following is an example of an \emph{itemized}
list.
\begin{itemize}
   \item This is the first item of an itemized list.
         Each item in the list is marked with a ``tick''.
         You don't have to worry about what kind of tick
         mark is used.

   \item This is the second item of the list.  It
         contains another list nested inside it.  The inner
         list is an \emph{enumerated} list.
         \begin{enumerate}
            \item This is the first item of an enumerated
                  list that is nested within the itemized list.

            \item This is the second item of the inner list.
                  \LaTeX\ allows you to nest lists deeper than
                  you really should.
         \end{enumerate}
         This is the rest of the second item of the outer
         list.  It is no more interesting than any other
         part of the item.
   \item This is the third item of the list.
\end{itemize}
You can even display poetry.
\begin{verse}
   There is an environment
    for verse \\             % The \\ command separates lines
   Whose features some poets % within a stanza.
   will curse.

                             % One or more blank lines separate stanzas.

   For instead of making\\
   Them do \emph{all} line breaking, \\
   It allows them to put too many words on a line when they'd rather be
   forced to be terse.
\end{verse}

Mathematical formulas may also be displayed.  A
displayed formula
is
one-line long; multiline
formulas require special formatting instructions.
   \[  \ip{\Gamma}{\psi'} = x'' + y^{2} + z_{i}^{n}\]
Don't start a paragraph with a displayed equation,
nor make one a paragraph by itself.

\end{document}               % End of document.

\section{Bench Section 13}
% synthetic bench section 13
% This is a sample LaTeX input file.  (Version of 12 August 2004.)
%
% A '%' character causes TeX to ignore all remaining text on the line,
% and is used for comments like this one.

\documentclass{article}      % Specifies the document class

                             % The preamble begins here.
\title{An Example Document}  % Declares the document's title.
\author{Leslie Lamport}      % Declares the author's name.
\date{January 21, 1994}      % Deleting this command produces today's date.

\newcommand{\ip}[2]{(#1, #2)}
                             % Defines \ip{arg1}{arg2} to mean
                             % (arg1, arg2).

%\newcommand{\ip}[2]{\langle #1 | #2\rangle}
                             % This is an alternative definition of
                             % \ip that is commented out.

\begin{document}             % End of preamble and beginning of text.

\maketitle                   % Produces the title.

This is an example input file.  Comparing it with
the output it generates can show you how to
produce a simple document of your own.

\section{Ordinary Text}      % Produces section heading.  Lower-level
                             % sections are begun with similar
                             % \subsection and \subsubsection commands.

The ends  of words and sentences are marked
  by   spaces. It  doesn't matter how many
spaces    you type; one is as good as 100.  The
end of   a line counts as a space.

One   or more   blank lines denote the  end
of  a paragraph.

Since any number of consecutive spaces are treated
like a single one, the formatting of the input
file makes no difference to
      \LaTeX,                % The \LaTeX command generates the LaTeX logo.
but it makes a difference to you.  When you use
\LaTeX, making your input file as easy to read
as possible will be a great help as you write
your document and when you change it.  This sample
file shows how you can add comments to your own input
file.

Because printing is different from typewriting,
there are a number of things that you have to do
differently when preparing an input file than if
you were just typing the document directly.
Quotation marks like
       ``this''
have to be handled specially, as do quotes within
quotes:
       ``\,`this'            % \, separates the double and single quote.
        is what I just
        wrote, not  `that'\,''.

Dashes come in three sizes: an
       intra-word
dash, a medium dash for number ranges like
       1--2,
and a punctuation
       dash---like
this.

A sentence-ending space should be larger than the
space between words within a sentence.  You
sometimes have to type special commands in
conjunction with punctuation characters to get
this right, as in the following sentence.
       Gnats, gnus, etc.\ all  % `\ ' makes an inter-word space.
       begin with G\@.         % \@ marks end-of-sentence punctuation.
You should check the spaces after periods when
reading your output to make sure you haven't
forgotten any special cases.  Generating an
ellipsis
       \ldots\               % `\ ' is needed after `\ldots' because TeX
                             % ignores spaces after command names like \ldots
                             % made from \ + letters.
                             %
                             % Note how a `%' character causes TeX to ignore
                             % the end of the input line, so these blank lines
                             % do not start a new paragraph.
                             %
with the right spacing around the periods requires
a special command.

\LaTeX\ interprets some common characters as
commands, so you must type special commands to
generate them.  These characters include the
following:
       \$ \& \% \# \{ and \}.

In printing, text is usually emphasized with an
       \emph{italic}
type style.

\begin{em}
   A long segment of text can also be emphasized
   in this way.  Text within such a segment can be
   given \emph{additional} emphasis.
\end{em}

It is sometimes necessary to prevent \LaTeX\ from
breaking a line where it might otherwise do so.
This may be at a space, as between the ``Mr.''\ and
``Jones'' in
       ``Mr.~Jones'',        % ~ produces an unbreakable interword space.
or within a word---especially when the word is a
symbol like
       \mbox{\emph{itemnum}}
that makes little sense when hyphenated across
lines.

Footnotes\footnote{This is an example of a footnote.}
pose no problem.

\LaTeX\ is good at typesetting mathematical formulas
like
       \( x-3y + z = 7 \)
or
       \( a_{1} > x^{2n} + y^{2n} > x' \)
or
       \( \ip{A}{B} = \sum_{i} a_{i} b_{i} \).
The spaces you type in a formula are
ignored.  Remember that a letter like
       $x$                   % $ ... $  and  \( ... \)  are equivalent
is a formula when it denotes a mathematical
symbol, and it should be typed as one.

\section{Displayed Text}

Text is displayed by indenting it from the left
margin.  Quotations are commonly displayed.  There
are short quotations
\begin{quote}
   This is a short quotation.  It consists of a
   single paragraph of text.  See how it is formatted.
\end{quote}
and longer ones.
\begin{quotation}
   This is a longer quotation.  It consists of two
   paragraphs of text, neither of which are
   particularly interesting.

   This is the second paragraph of the quotation.  It
   is just as dull as the first paragraph.
\end{quotation}
Another frequently-displayed structure is a list.
The following is an example of an \emph{itemized}
list.
\begin{itemize}
   \item This is the first item of an itemized list.
         Each item in the list is marked with a ``tick''.
         You don't have to worry about what kind of tick
         mark is used.

   \item This is the second item of the list.  It
         contains another list nested inside it.  The inner
         list is an \emph{enumerated} list.
         \begin{enumerate}
            \item This is the first item of an enumerated
                  list that is nested within the itemized list.

            \item This is the second item of the inner list.
                  \LaTeX\ allows you to nest lists deeper than
                  you really should.
         \end{enumerate}
         This is the rest of the second item of the outer
         list.  It is no more interesting than any other
         part of the item.
   \item This is the third item of the list.
\end{itemize}
You can even display poetry.
\begin{verse}
   There is an environment
    for verse \\             % The \\ command separates lines
   Whose features some poets % within a stanza.
   will curse.

                             % One or more blank lines separate stanzas.

   For instead of making\\
   Them do \emph{all} line breaking, \\
   It allows them to put too many words on a line when they'd rather be
   forced to be terse.
\end{verse}

Mathematical formulas may also be displayed.  A
displayed formula
is
one-line long; multiline
formulas require special formatting instructions.
   \[  \ip{\Gamma}{\psi'} = x'' + y^{2} + z_{i}^{n}\]
Don't start a paragraph with a displayed equation,
nor make one a paragraph by itself.

\end{document}               % End of document.

\section{Bench Section 14}
% synthetic bench section 14
% This is a sample LaTeX input file.  (Version of 12 August 2004.)
%
% A '%' character causes TeX to ignore all remaining text on the line,
% and is used for comments like this one.

\documentclass{article}      % Specifies the document class

                             % The preamble begins here.
\title{An Example Document}  % Declares the document's title.
\author{Leslie Lamport}      % Declares the author's name.
\date{January 21, 1994}      % Deleting this command produces today's date.

\newcommand{\ip}[2]{(#1, #2)}
                             % Defines \ip{arg1}{arg2} to mean
                             % (arg1, arg2).

%\newcommand{\ip}[2]{\langle #1 | #2\rangle}
                             % This is an alternative definition of
                             % \ip that is commented out.

\begin{document}             % End of preamble and beginning of text.

\maketitle                   % Produces the title.

This is an example input file.  Comparing it with
the output it generates can show you how to
produce a simple document of your own.

\section{Ordinary Text}      % Produces section heading.  Lower-level
                             % sections are begun with similar
                             % \subsection and \subsubsection commands.

The ends  of words and sentences are marked
  by   spaces. It  doesn't matter how many
spaces    you type; one is as good as 100.  The
end of   a line counts as a space.

One   or more   blank lines denote the  end
of  a paragraph.

Since any number of consecutive spaces are treated
like a single one, the formatting of the input
file makes no difference to
      \LaTeX,                % The \LaTeX command generates the LaTeX logo.
but it makes a difference to you.  When you use
\LaTeX, making your input file as easy to read
as possible will be a great help as you write
your document and when you change it.  This sample
file shows how you can add comments to your own input
file.

Because printing is different from typewriting,
there are a number of things that you have to do
differently when preparing an input file than if
you were just typing the document directly.
Quotation marks like
       ``this''
have to be handled specially, as do quotes within
quotes:
       ``\,`this'            % \, separates the double and single quote.
        is what I just
        wrote, not  `that'\,''.

Dashes come in three sizes: an
       intra-word
dash, a medium dash for number ranges like
       1--2,
and a punctuation
       dash---like
this.

A sentence-ending space should be larger than the
space between words within a sentence.  You
sometimes have to type special commands in
conjunction with punctuation characters to get
this right, as in the following sentence.
       Gnats, gnus, etc.\ all  % `\ ' makes an inter-word space.
       begin with G\@.         % \@ marks end-of-sentence punctuation.
You should check the spaces after periods when
reading your output to make sure you haven't
forgotten any special cases.  Generating an
ellipsis
       \ldots\               % `\ ' is needed after `\ldots' because TeX
                             % ignores spaces after command names like \ldots
                             % made from \ + letters.
                             %
                             % Note how a `%' character causes TeX to ignore
                             % the end of the input line, so these blank lines
                             % do not start a new paragraph.
                             %
with the right spacing around the periods requires
a special command.

\LaTeX\ interprets some common characters as
commands, so you must type special commands to
generate them.  These characters include the
following:
       \$ \& \% \# \{ and \}.

In printing, text is usually emphasized with an
       \emph{italic}
type style.

\begin{em}
   A long segment of text can also be emphasized
   in this way.  Text within such a segment can be
   given \emph{additional} emphasis.
\end{em}

It is sometimes necessary to prevent \LaTeX\ from
breaking a line where it might otherwise do so.
This may be at a space, as between the ``Mr.''\ and
``Jones'' in
       ``Mr.~Jones'',        % ~ produces an unbreakable interword space.
or within a word---especially when the word is a
symbol like
       \mbox{\emph{itemnum}}
that makes little sense when hyphenated across
lines.

Footnotes\footnote{This is an example of a footnote.}
pose no problem.

\LaTeX\ is good at typesetting mathematical formulas
like
       \( x-3y + z = 7 \)
or
       \( a_{1} > x^{2n} + y^{2n} > x' \)
or
       \( \ip{A}{B} = \sum_{i} a_{i} b_{i} \).
The spaces you type in a formula are
ignored.  Remember that a letter like
       $x$                   % $ ... $  and  \( ... \)  are equivalent
is a formula when it denotes a mathematical
symbol, and it should be typed as one.

\section{Displayed Text}

Text is displayed by indenting it from the left
margin.  Quotations are commonly displayed.  There
are short quotations
\begin{quote}
   This is a short quotation.  It consists of a
   single paragraph of text.  See how it is formatted.
\end{quote}
and longer ones.
\begin{quotation}
   This is a longer quotation.  It consists of two
   paragraphs of text, neither of which are
   particularly interesting.

   This is the second paragraph of the quotation.  It
   is just as dull as the first paragraph.
\end{quotation}
Another frequently-displayed structure is a list.
The following is an example of an \emph{itemized}
list.
\begin{itemize}
   \item This is the first item of an itemized list.
         Each item in the list is marked with a ``tick''.
         You don't have to worry about what kind of tick
         mark is used.

   \item This is the second item of the list.  It
         contains another list nested inside it.  The inner
         list is an \emph{enumerated} list.
         \begin{enumerate}
            \item This is the first item of an enumerated
                  list that is nested within the itemized list.

            \item This is the second item of the inner list.
                  \LaTeX\ allows you to nest lists deeper than
                  you really should.
         \end{enumerate}
         This is the rest of the second item of the outer
         list.  It is no more interesting than any other
         part of the item.
   \item This is the third item of the list.
\end{itemize}
You can even display poetry.
\begin{verse}
   There is an environment
    for verse \\             % The \\ command separates lines
   Whose features some poets % within a stanza.
   will curse.

                             % One or more blank lines separate stanzas.

   For instead of making\\
   Them do \emph{all} line breaking, \\
   It allows them to put too many words on a line when they'd rather be
   forced to be terse.
\end{verse}

Mathematical formulas may also be displayed.  A
displayed formula
is
one-line long; multiline
formulas require special formatting instructions.
   \[  \ip{\Gamma}{\psi'} = x'' + y^{2} + z_{i}^{n}\]
Don't start a paragraph with a displayed equation,
nor make one a paragraph by itself.

\end{document}               % End of document.

\section{Bench Section 15}
% synthetic bench section 15
% This is a sample LaTeX input file.  (Version of 12 August 2004.)
%
% A '%' character causes TeX to ignore all remaining text on the line,
% and is used for comments like this one.

\documentclass{article}      % Specifies the document class

                             % The preamble begins here.
\title{An Example Document}  % Declares the document's title.
\author{Leslie Lamport}      % Declares the author's name.
\date{January 21, 1994}      % Deleting this command produces today's date.

\newcommand{\ip}[2]{(#1, #2)}
                             % Defines \ip{arg1}{arg2} to mean
                             % (arg1, arg2).

%\newcommand{\ip}[2]{\langle #1 | #2\rangle}
                             % This is an alternative definition of
                             % \ip that is commented out.

\begin{document}             % End of preamble and beginning of text.

\maketitle                   % Produces the title.

This is an example input file.  Comparing it with
the output it generates can show you how to
produce a simple document of your own.

\section{Ordinary Text}      % Produces section heading.  Lower-level
                             % sections are begun with similar
                             % \subsection and \subsubsection commands.

The ends  of words and sentences are marked
  by   spaces. It  doesn't matter how many
spaces    you type; one is as good as 100.  The
end of   a line counts as a space.

One   or more   blank lines denote the  end
of  a paragraph.

Since any number of consecutive spaces are treated
like a single one, the formatting of the input
file makes no difference to
      \LaTeX,                % The \LaTeX command generates the LaTeX logo.
but it makes a difference to you.  When you use
\LaTeX, making your input file as easy to read
as possible will be a great help as you write
your document and when you change it.  This sample
file shows how you can add comments to your own input
file.

Because printing is different from typewriting,
there are a number of things that you have to do
differently when preparing an input file than if
you were just typing the document directly.
Quotation marks like
       ``this''
have to be handled specially, as do quotes within
quotes:
       ``\,`this'            % \, separates the double and single quote.
        is what I just
        wrote, not  `that'\,''.

Dashes come in three sizes: an
       intra-word
dash, a medium dash for number ranges like
       1--2,
and a punctuation
       dash---like
this.

A sentence-ending space should be larger than the
space between words within a sentence.  You
sometimes have to type special commands in
conjunction with punctuation characters to get
this right, as in the following sentence.
       Gnats, gnus, etc.\ all  % `\ ' makes an inter-word space.
       begin with G\@.         % \@ marks end-of-sentence punctuation.
You should check the spaces after periods when
reading your output to make sure you haven't
forgotten any special cases.  Generating an
ellipsis
       \ldots\               % `\ ' is needed after `\ldots' because TeX
                             % ignores spaces after command names like \ldots
                             % made from \ + letters.
                             %
                             % Note how a `%' character causes TeX to ignore
                             % the end of the input line, so these blank lines
                             % do not start a new paragraph.
                             %
with the right spacing around the periods requires
a special command.

\LaTeX\ interprets some common characters as
commands, so you must type special commands to
generate them.  These characters include the
following:
       \$ \& \% \# \{ and \}.

In printing, text is usually emphasized with an
       \emph{italic}
type style.

\begin{em}
   A long segment of text can also be emphasized
   in this way.  Text within such a segment can be
   given \emph{additional} emphasis.
\end{em}

It is sometimes necessary to prevent \LaTeX\ from
breaking a line where it might otherwise do so.
This may be at a space, as between the ``Mr.''\ and
``Jones'' in
       ``Mr.~Jones'',        % ~ produces an unbreakable interword space.
or within a word---especially when the word is a
symbol like
       \mbox{\emph{itemnum}}
that makes little sense when hyphenated across
lines.

Footnotes\footnote{This is an example of a footnote.}
pose no problem.

\LaTeX\ is good at typesetting mathematical formulas
like
       \( x-3y + z = 7 \)
or
       \( a_{1} > x^{2n} + y^{2n} > x' \)
or
       \( \ip{A}{B} = \sum_{i} a_{i} b_{i} \).
The spaces you type in a formula are
ignored.  Remember that a letter like
       $x$                   % $ ... $  and  \( ... \)  are equivalent
is a formula when it denotes a mathematical
symbol, and it should be typed as one.

\section{Displayed Text}

Text is displayed by indenting it from the left
margin.  Quotations are commonly displayed.  There
are short quotations
\begin{quote}
   This is a short quotation.  It consists of a
   single paragraph of text.  See how it is formatted.
\end{quote}
and longer ones.
\begin{quotation}
   This is a longer quotation.  It consists of two
   paragraphs of text, neither of which are
   particularly interesting.

   This is the second paragraph of the quotation.  It
   is just as dull as the first paragraph.
\end{quotation}
Another frequently-displayed structure is a list.
The following is an example of an \emph{itemized}
list.
\begin{itemize}
   \item This is the first item of an itemized list.
         Each item in the list is marked with a ``tick''.
         You don't have to worry about what kind of tick
         mark is used.

   \item This is the second item of the list.  It
         contains another list nested inside it.  The inner
         list is an \emph{enumerated} list.
         \begin{enumerate}
            \item This is the first item of an enumerated
                  list that is nested within the itemized list.

            \item This is the second item of the inner list.
                  \LaTeX\ allows you to nest lists deeper than
                  you really should.
         \end{enumerate}
         This is the rest of the second item of the outer
         list.  It is no more interesting than any other
         part of the item.
   \item This is the third item of the list.
\end{itemize}
You can even display poetry.
\begin{verse}
   There is an environment
    for verse \\             % The \\ command separates lines
   Whose features some poets % within a stanza.
   will curse.

                             % One or more blank lines separate stanzas.

   For instead of making\\
   Them do \emph{all} line breaking, \\
   It allows them to put too many words on a line when they'd rather be
   forced to be terse.
\end{verse}

Mathematical formulas may also be displayed.  A
displayed formula
is
one-line long; multiline
formulas require special formatting instructions.
   \[  \ip{\Gamma}{\psi'} = x'' + y^{2} + z_{i}^{n}\]
Don't start a paragraph with a displayed equation,
nor make one a paragraph by itself.

\end{document}               % End of document.

\section{Bench Section 16}
% synthetic bench section 16
% This is a sample LaTeX input file.  (Version of 12 August 2004.)
%
% A '%' character causes TeX to ignore all remaining text on the line,
% and is used for comments like this one.

\documentclass{article}      % Specifies the document class

                             % The preamble begins here.
\title{An Example Document}  % Declares the document's title.
\author{Leslie Lamport}      % Declares the author's name.
\date{January 21, 1994}      % Deleting this command produces today's date.

\newcommand{\ip}[2]{(#1, #2)}
                             % Defines \ip{arg1}{arg2} to mean
                             % (arg1, arg2).

%\newcommand{\ip}[2]{\langle #1 | #2\rangle}
                             % This is an alternative definition of
                             % \ip that is commented out.

\begin{document}             % End of preamble and beginning of text.

\maketitle                   % Produces the title.

This is an example input file.  Comparing it with
the output it generates can show you how to
produce a simple document of your own.

\section{Ordinary Text}      % Produces section heading.  Lower-level
                             % sections are begun with similar
                             % \subsection and \subsubsection commands.

The ends  of words and sentences are marked
  by   spaces. It  doesn't matter how many
spaces    you type; one is as good as 100.  The
end of   a line counts as a space.

One   or more   blank lines denote the  end
of  a paragraph.

Since any number of consecutive spaces are treated
like a single one, the formatting of the input
file makes no difference to
      \LaTeX,                % The \LaTeX command generates the LaTeX logo.
but it makes a difference to you.  When you use
\LaTeX, making your input file as easy to read
as possible will be a great help as you write
your document and when you change it.  This sample
file shows how you can add comments to your own input
file.

Because printing is different from typewriting,
there are a number of things that you have to do
differently when preparing an input file than if
you were just typing the document directly.
Quotation marks like
       ``this''
have to be handled specially, as do quotes within
quotes:
       ``\,`this'            % \, separates the double and single quote.
        is what I just
        wrote, not  `that'\,''.

Dashes come in three sizes: an
       intra-word
dash, a medium dash for number ranges like
       1--2,
and a punctuation
       dash---like
this.

A sentence-ending space should be larger than the
space between words within a sentence.  You
sometimes have to type special commands in
conjunction with punctuation characters to get
this right, as in the following sentence.
       Gnats, gnus, etc.\ all  % `\ ' makes an inter-word space.
       begin with G\@.         % \@ marks end-of-sentence punctuation.
You should check the spaces after periods when
reading your output to make sure you haven't
forgotten any special cases.  Generating an
ellipsis
       \ldots\               % `\ ' is needed after `\ldots' because TeX
                             % ignores spaces after command names like \ldots
                             % made from \ + letters.
                             %
                             % Note how a `%' character causes TeX to ignore
                             % the end of the input line, so these blank lines
                             % do not start a new paragraph.
                             %
with the right spacing around the periods requires
a special command.

\LaTeX\ interprets some common characters as
commands, so you must type special commands to
generate them.  These characters include the
following:
       \$ \& \% \# \{ and \}.

In printing, text is usually emphasized with an
       \emph{italic}
type style.

\begin{em}
   A long segment of text can also be emphasized
   in this way.  Text within such a segment can be
   given \emph{additional} emphasis.
\end{em}

It is sometimes necessary to prevent \LaTeX\ from
breaking a line where it might otherwise do so.
This may be at a space, as between the ``Mr.''\ and
``Jones'' in
       ``Mr.~Jones'',        % ~ produces an unbreakable interword space.
or within a word---especially when the word is a
symbol like
       \mbox{\emph{itemnum}}
that makes little sense when hyphenated across
lines.

Footnotes\footnote{This is an example of a footnote.}
pose no problem.

\LaTeX\ is good at typesetting mathematical formulas
like
       \( x-3y + z = 7 \)
or
       \( a_{1} > x^{2n} + y^{2n} > x' \)
or
       \( \ip{A}{B} = \sum_{i} a_{i} b_{i} \).
The spaces you type in a formula are
ignored.  Remember that a letter like
       $x$                   % $ ... $  and  \( ... \)  are equivalent
is a formula when it denotes a mathematical
symbol, and it should be typed as one.

\section{Displayed Text}

Text is displayed by indenting it from the left
margin.  Quotations are commonly displayed.  There
are short quotations
\begin{quote}
   This is a short quotation.  It consists of a
   single paragraph of text.  See how it is formatted.
\end{quote}
and longer ones.
\begin{quotation}
   This is a longer quotation.  It consists of two
   paragraphs of text, neither of which are
   particularly interesting.

   This is the second paragraph of the quotation.  It
   is just as dull as the first paragraph.
\end{quotation}
Another frequently-displayed structure is a list.
The following is an example of an \emph{itemized}
list.
\begin{itemize}
   \item This is the first item of an itemized list.
         Each item in the list is marked with a ``tick''.
         You don't have to worry about what kind of tick
         mark is used.

   \item This is the second item of the list.  It
         contains another list nested inside it.  The inner
         list is an \emph{enumerated} list.
         \begin{enumerate}
            \item This is the first item of an enumerated
                  list that is nested within the itemized list.

            \item This is the second item of the inner list.
                  \LaTeX\ allows you to nest lists deeper than
                  you really should.
         \end{enumerate}
         This is the rest of the second item of the outer
         list.  It is no more interesting than any other
         part of the item.
   \item This is the third item of the list.
\end{itemize}
You can even display poetry.
\begin{verse}
   There is an environment
    for verse \\             % The \\ command separates lines
   Whose features some poets % within a stanza.
   will curse.

                             % One or more blank lines separate stanzas.

   For instead of making\\
   Them do \emph{all} line breaking, \\
   It allows them to put too many words on a line when they'd rather be
   forced to be terse.
\end{verse}

Mathematical formulas may also be displayed.  A
displayed formula
is
one-line long; multiline
formulas require special formatting instructions.
   \[  \ip{\Gamma}{\psi'} = x'' + y^{2} + z_{i}^{n}\]
Don't start a paragraph with a displayed equation,
nor make one a paragraph by itself.

\end{document}               % End of document.

\section{Bench Section 17}
% synthetic bench section 17
% This is a sample LaTeX input file.  (Version of 12 August 2004.)
%
% A '%' character causes TeX to ignore all remaining text on the line,
% and is used for comments like this one.

\documentclass{article}      % Specifies the document class

                             % The preamble begins here.
\title{An Example Document}  % Declares the document's title.
\author{Leslie Lamport}      % Declares the author's name.
\date{January 21, 1994}      % Deleting this command produces today's date.

\newcommand{\ip}[2]{(#1, #2)}
                             % Defines \ip{arg1}{arg2} to mean
                             % (arg1, arg2).

%\newcommand{\ip}[2]{\langle #1 | #2\rangle}
                             % This is an alternative definition of
                             % \ip that is commented out.

\begin{document}             % End of preamble and beginning of text.

\maketitle                   % Produces the title.

This is an example input file.  Comparing it with
the output it generates can show you how to
produce a simple document of your own.

\section{Ordinary Text}      % Produces section heading.  Lower-level
                             % sections are begun with similar
                             % \subsection and \subsubsection commands.

The ends  of words and sentences are marked
  by   spaces. It  doesn't matter how many
spaces    you type; one is as good as 100.  The
end of   a line counts as a space.

One   or more   blank lines denote the  end
of  a paragraph.

Since any number of consecutive spaces are treated
like a single one, the formatting of the input
file makes no difference to
      \LaTeX,                % The \LaTeX command generates the LaTeX logo.
but it makes a difference to you.  When you use
\LaTeX, making your input file as easy to read
as possible will be a great help as you write
your document and when you change it.  This sample
file shows how you can add comments to your own input
file.

Because printing is different from typewriting,
there are a number of things that you have to do
differently when preparing an input file than if
you were just typing the document directly.
Quotation marks like
       ``this''
have to be handled specially, as do quotes within
quotes:
       ``\,`this'            % \, separates the double and single quote.
        is what I just
        wrote, not  `that'\,''.

Dashes come in three sizes: an
       intra-word
dash, a medium dash for number ranges like
       1--2,
and a punctuation
       dash---like
this.

A sentence-ending space should be larger than the
space between words within a sentence.  You
sometimes have to type special commands in
conjunction with punctuation characters to get
this right, as in the following sentence.
       Gnats, gnus, etc.\ all  % `\ ' makes an inter-word space.
       begin with G\@.         % \@ marks end-of-sentence punctuation.
You should check the spaces after periods when
reading your output to make sure you haven't
forgotten any special cases.  Generating an
ellipsis
       \ldots\               % `\ ' is needed after `\ldots' because TeX
                             % ignores spaces after command names like \ldots
                             % made from \ + letters.
                             %
                             % Note how a `%' character causes TeX to ignore
                             % the end of the input line, so these blank lines
                             % do not start a new paragraph.
                             %
with the right spacing around the periods requires
a special command.

\LaTeX\ interprets some common characters as
commands, so you must type special commands to
generate them.  These characters include the
following:
       \$ \& \% \# \{ and \}.

In printing, text is usually emphasized with an
       \emph{italic}
type style.

\begin{em}
   A long segment of text can also be emphasized
   in this way.  Text within such a segment can be
   given \emph{additional} emphasis.
\end{em}

It is sometimes necessary to prevent \LaTeX\ from
breaking a line where it might otherwise do so.
This may be at a space, as between the ``Mr.''\ and
``Jones'' in
       ``Mr.~Jones'',        % ~ produces an unbreakable interword space.
or within a word---especially when the word is a
symbol like
       \mbox{\emph{itemnum}}
that makes little sense when hyphenated across
lines.

Footnotes\footnote{This is an example of a footnote.}
pose no problem.

\LaTeX\ is good at typesetting mathematical formulas
like
       \( x-3y + z = 7 \)
or
       \( a_{1} > x^{2n} + y^{2n} > x' \)
or
       \( \ip{A}{B} = \sum_{i} a_{i} b_{i} \).
The spaces you type in a formula are
ignored.  Remember that a letter like
       $x$                   % $ ... $  and  \( ... \)  are equivalent
is a formula when it denotes a mathematical
symbol, and it should be typed as one.

\section{Displayed Text}

Text is displayed by indenting it from the left
margin.  Quotations are commonly displayed.  There
are short quotations
\begin{quote}
   This is a short quotation.  It consists of a
   single paragraph of text.  See how it is formatted.
\end{quote}
and longer ones.
\begin{quotation}
   This is a longer quotation.  It consists of two
   paragraphs of text, neither of which are
   particularly interesting.

   This is the second paragraph of the quotation.  It
   is just as dull as the first paragraph.
\end{quotation}
Another frequently-displayed structure is a list.
The following is an example of an \emph{itemized}
list.
\begin{itemize}
   \item This is the first item of an itemized list.
         Each item in the list is marked with a ``tick''.
         You don't have to worry about what kind of tick
         mark is used.

   \item This is the second item of the list.  It
         contains another list nested inside it.  The inner
         list is an \emph{enumerated} list.
         \begin{enumerate}
            \item This is the first item of an enumerated
                  list that is nested within the itemized list.

            \item This is the second item of the inner list.
                  \LaTeX\ allows you to nest lists deeper than
                  you really should.
         \end{enumerate}
         This is the rest of the second item of the outer
         list.  It is no more interesting than any other
         part of the item.
   \item This is the third item of the list.
\end{itemize}
You can even display poetry.
\begin{verse}
   There is an environment
    for verse \\             % The \\ command separates lines
   Whose features some poets % within a stanza.
   will curse.

                             % One or more blank lines separate stanzas.

   For instead of making\\
   Them do \emph{all} line breaking, \\
   It allows them to put too many words on a line when they'd rather be
   forced to be terse.
\end{verse}

Mathematical formulas may also be displayed.  A
displayed formula
is
one-line long; multiline
formulas require special formatting instructions.
   \[  \ip{\Gamma}{\psi'} = x'' + y^{2} + z_{i}^{n}\]
Don't start a paragraph with a displayed equation,
nor make one a paragraph by itself.

\end{document}               % End of document.

\section{Bench Section 18}
% synthetic bench section 18
% This is a sample LaTeX input file.  (Version of 12 August 2004.)
%
% A '%' character causes TeX to ignore all remaining text on the line,
% and is used for comments like this one.

\documentclass{article}      % Specifies the document class

                             % The preamble begins here.
\title{An Example Document}  % Declares the document's title.
\author{Leslie Lamport}      % Declares the author's name.
\date{January 21, 1994}      % Deleting this command produces today's date.

\newcommand{\ip}[2]{(#1, #2)}
                             % Defines \ip{arg1}{arg2} to mean
                             % (arg1, arg2).

%\newcommand{\ip}[2]{\langle #1 | #2\rangle}
                             % This is an alternative definition of
                             % \ip that is commented out.

\begin{document}             % End of preamble and beginning of text.

\maketitle                   % Produces the title.

This is an example input file.  Comparing it with
the output it generates can show you how to
produce a simple document of your own.

\section{Ordinary Text}      % Produces section heading.  Lower-level
                             % sections are begun with similar
                             % \subsection and \subsubsection commands.

The ends  of words and sentences are marked
  by   spaces. It  doesn't matter how many
spaces    you type; one is as good as 100.  The
end of   a line counts as a space.

One   or more   blank lines denote the  end
of  a paragraph.

Since any number of consecutive spaces are treated
like a single one, the formatting of the input
file makes no difference to
      \LaTeX,                % The \LaTeX command generates the LaTeX logo.
but it makes a difference to you.  When you use
\LaTeX, making your input file as easy to read
as possible will be a great help as you write
your document and when you change it.  This sample
file shows how you can add comments to your own input
file.

Because printing is different from typewriting,
there are a number of things that you have to do
differently when preparing an input file than if
you were just typing the document directly.
Quotation marks like
       ``this''
have to be handled specially, as do quotes within
quotes:
       ``\,`this'            % \, separates the double and single quote.
        is what I just
        wrote, not  `that'\,''.

Dashes come in three sizes: an
       intra-word
dash, a medium dash for number ranges like
       1--2,
and a punctuation
       dash---like
this.

A sentence-ending space should be larger than the
space between words within a sentence.  You
sometimes have to type special commands in
conjunction with punctuation characters to get
this right, as in the following sentence.
       Gnats, gnus, etc.\ all  % `\ ' makes an inter-word space.
       begin with G\@.         % \@ marks end-of-sentence punctuation.
You should check the spaces after periods when
reading your output to make sure you haven't
forgotten any special cases.  Generating an
ellipsis
       \ldots\               % `\ ' is needed after `\ldots' because TeX
                             % ignores spaces after command names like \ldots
                             % made from \ + letters.
                             %
                             % Note how a `%' character causes TeX to ignore
                             % the end of the input line, so these blank lines
                             % do not start a new paragraph.
                             %
with the right spacing around the periods requires
a special command.

\LaTeX\ interprets some common characters as
commands, so you must type special commands to
generate them.  These characters include the
following:
       \$ \& \% \# \{ and \}.

In printing, text is usually emphasized with an
       \emph{italic}
type style.

\begin{em}
   A long segment of text can also be emphasized
   in this way.  Text within such a segment can be
   given \emph{additional} emphasis.
\end{em}

It is sometimes necessary to prevent \LaTeX\ from
breaking a line where it might otherwise do so.
This may be at a space, as between the ``Mr.''\ and
``Jones'' in
       ``Mr.~Jones'',        % ~ produces an unbreakable interword space.
or within a word---especially when the word is a
symbol like
       \mbox{\emph{itemnum}}
that makes little sense when hyphenated across
lines.

Footnotes\footnote{This is an example of a footnote.}
pose no problem.

\LaTeX\ is good at typesetting mathematical formulas
like
       \( x-3y + z = 7 \)
or
       \( a_{1} > x^{2n} + y^{2n} > x' \)
or
       \( \ip{A}{B} = \sum_{i} a_{i} b_{i} \).
The spaces you type in a formula are
ignored.  Remember that a letter like
       $x$                   % $ ... $  and  \( ... \)  are equivalent
is a formula when it denotes a mathematical
symbol, and it should be typed as one.

\section{Displayed Text}

Text is displayed by indenting it from the left
margin.  Quotations are commonly displayed.  There
are short quotations
\begin{quote}
   This is a short quotation.  It consists of a
   single paragraph of text.  See how it is formatted.
\end{quote}
and longer ones.
\begin{quotation}
   This is a longer quotation.  It consists of two
   paragraphs of text, neither of which are
   particularly interesting.

   This is the second paragraph of the quotation.  It
   is just as dull as the first paragraph.
\end{quotation}
Another frequently-displayed structure is a list.
The following is an example of an \emph{itemized}
list.
\begin{itemize}
   \item This is the first item of an itemized list.
         Each item in the list is marked with a ``tick''.
         You don't have to worry about what kind of tick
         mark is used.

   \item This is the second item of the list.  It
         contains another list nested inside it.  The inner
         list is an \emph{enumerated} list.
         \begin{enumerate}
            \item This is the first item of an enumerated
                  list that is nested within the itemized list.

            \item This is the second item of the inner list.
                  \LaTeX\ allows you to nest lists deeper than
                  you really should.
         \end{enumerate}
         This is the rest of the second item of the outer
         list.  It is no more interesting than any other
         part of the item.
   \item This is the third item of the list.
\end{itemize}
You can even display poetry.
\begin{verse}
   There is an environment
    for verse \\             % The \\ command separates lines
   Whose features some poets % within a stanza.
   will curse.

                             % One or more blank lines separate stanzas.

   For instead of making\\
   Them do \emph{all} line breaking, \\
   It allows them to put too many words on a line when they'd rather be
   forced to be terse.
\end{verse}

Mathematical formulas may also be displayed.  A
displayed formula
is
one-line long; multiline
formulas require special formatting instructions.
   \[  \ip{\Gamma}{\psi'} = x'' + y^{2} + z_{i}^{n}\]
Don't start a paragraph with a displayed equation,
nor make one a paragraph by itself.

\end{document}               % End of document.

\section{Bench Section 19}
% synthetic bench section 19
% This is a sample LaTeX input file.  (Version of 12 August 2004.)
%
% A '%' character causes TeX to ignore all remaining text on the line,
% and is used for comments like this one.

\documentclass{article}      % Specifies the document class

                             % The preamble begins here.
\title{An Example Document}  % Declares the document's title.
\author{Leslie Lamport}      % Declares the author's name.
\date{January 21, 1994}      % Deleting this command produces today's date.

\newcommand{\ip}[2]{(#1, #2)}
                             % Defines \ip{arg1}{arg2} to mean
                             % (arg1, arg2).

%\newcommand{\ip}[2]{\langle #1 | #2\rangle}
                             % This is an alternative definition of
                             % \ip that is commented out.

\begin{document}             % End of preamble and beginning of text.

\maketitle                   % Produces the title.

This is an example input file.  Comparing it with
the output it generates can show you how to
produce a simple document of your own.

\section{Ordinary Text}      % Produces section heading.  Lower-level
                             % sections are begun with similar
                             % \subsection and \subsubsection commands.

The ends  of words and sentences are marked
  by   spaces. It  doesn't matter how many
spaces    you type; one is as good as 100.  The
end of   a line counts as a space.

One   or more   blank lines denote the  end
of  a paragraph.

Since any number of consecutive spaces are treated
like a single one, the formatting of the input
file makes no difference to
      \LaTeX,                % The \LaTeX command generates the LaTeX logo.
but it makes a difference to you.  When you use
\LaTeX, making your input file as easy to read
as possible will be a great help as you write
your document and when you change it.  This sample
file shows how you can add comments to your own input
file.

Because printing is different from typewriting,
there are a number of things that you have to do
differently when preparing an input file than if
you were just typing the document directly.
Quotation marks like
       ``this''
have to be handled specially, as do quotes within
quotes:
       ``\,`this'            % \, separates the double and single quote.
        is what I just
        wrote, not  `that'\,''.

Dashes come in three sizes: an
       intra-word
dash, a medium dash for number ranges like
       1--2,
and a punctuation
       dash---like
this.

A sentence-ending space should be larger than the
space between words within a sentence.  You
sometimes have to type special commands in
conjunction with punctuation characters to get
this right, as in the following sentence.
       Gnats, gnus, etc.\ all  % `\ ' makes an inter-word space.
       begin with G\@.         % \@ marks end-of-sentence punctuation.
You should check the spaces after periods when
reading your output to make sure you haven't
forgotten any special cases.  Generating an
ellipsis
       \ldots\               % `\ ' is needed after `\ldots' because TeX
                             % ignores spaces after command names like \ldots
                             % made from \ + letters.
                             %
                             % Note how a `%' character causes TeX to ignore
                             % the end of the input line, so these blank lines
                             % do not start a new paragraph.
                             %
with the right spacing around the periods requires
a special command.

\LaTeX\ interprets some common characters as
commands, so you must type special commands to
generate them.  These characters include the
following:
       \$ \& \% \# \{ and \}.

In printing, text is usually emphasized with an
       \emph{italic}
type style.

\begin{em}
   A long segment of text can also be emphasized
   in this way.  Text within such a segment can be
   given \emph{additional} emphasis.
\end{em}

It is sometimes necessary to prevent \LaTeX\ from
breaking a line where it might otherwise do so.
This may be at a space, as between the ``Mr.''\ and
``Jones'' in
       ``Mr.~Jones'',        % ~ produces an unbreakable interword space.
or within a word---especially when the word is a
symbol like
       \mbox{\emph{itemnum}}
that makes little sense when hyphenated across
lines.

Footnotes\footnote{This is an example of a footnote.}
pose no problem.

\LaTeX\ is good at typesetting mathematical formulas
like
       \( x-3y + z = 7 \)
or
       \( a_{1} > x^{2n} + y^{2n} > x' \)
or
       \( \ip{A}{B} = \sum_{i} a_{i} b_{i} \).
The spaces you type in a formula are
ignored.  Remember that a letter like
       $x$                   % $ ... $  and  \( ... \)  are equivalent
is a formula when it denotes a mathematical
symbol, and it should be typed as one.

\section{Displayed Text}

Text is displayed by indenting it from the left
margin.  Quotations are commonly displayed.  There
are short quotations
\begin{quote}
   This is a short quotation.  It consists of a
   single paragraph of text.  See how it is formatted.
\end{quote}
and longer ones.
\begin{quotation}
   This is a longer quotation.  It consists of two
   paragraphs of text, neither of which are
   particularly interesting.

   This is the second paragraph of the quotation.  It
   is just as dull as the first paragraph.
\end{quotation}
Another frequently-displayed structure is a list.
The following is an example of an \emph{itemized}
list.
\begin{itemize}
   \item This is the first item of an itemized list.
         Each item in the list is marked with a ``tick''.
         You don't have to worry about what kind of tick
         mark is used.

   \item This is the second item of the list.  It
         contains another list nested inside it.  The inner
         list is an \emph{enumerated} list.
         \begin{enumerate}
            \item This is the first item of an enumerated
                  list that is nested within the itemized list.

            \item This is the second item of the inner list.
                  \LaTeX\ allows you to nest lists deeper than
                  you really should.
         \end{enumerate}
         This is the rest of the second item of the outer
         list.  It is no more interesting than any other
         part of the item.
   \item This is the third item of the list.
\end{itemize}
You can even display poetry.
\begin{verse}
   There is an environment
    for verse \\             % The \\ command separates lines
   Whose features some poets % within a stanza.
   will curse.

                             % One or more blank lines separate stanzas.

   For instead of making\\
   Them do \emph{all} line breaking, \\
   It allows them to put too many words on a line when they'd rather be
   forced to be terse.
\end{verse}

Mathematical formulas may also be displayed.  A
displayed formula
is
one-line long; multiline
formulas require special formatting instructions.
   \[  \ip{\Gamma}{\psi'} = x'' + y^{2} + z_{i}^{n}\]
Don't start a paragraph with a displayed equation,
nor make one a paragraph by itself.

\end{document}               % End of document.

\section{Bench Section 20}
% synthetic bench section 20
% This is a sample LaTeX input file.  (Version of 12 August 2004.)
%
% A '%' character causes TeX to ignore all remaining text on the line,
% and is used for comments like this one.

\documentclass{article}      % Specifies the document class

                             % The preamble begins here.
\title{An Example Document}  % Declares the document's title.
\author{Leslie Lamport}      % Declares the author's name.
\date{January 21, 1994}      % Deleting this command produces today's date.

\newcommand{\ip}[2]{(#1, #2)}
                             % Defines \ip{arg1}{arg2} to mean
                             % (arg1, arg2).

%\newcommand{\ip}[2]{\langle #1 | #2\rangle}
                             % This is an alternative definition of
                             % \ip that is commented out.

\begin{document}             % End of preamble and beginning of text.

\maketitle                   % Produces the title.

This is an example input file.  Comparing it with
the output it generates can show you how to
produce a simple document of your own.

\section{Ordinary Text}      % Produces section heading.  Lower-level
                             % sections are begun with similar
                             % \subsection and \subsubsection commands.

The ends  of words and sentences are marked
  by   spaces. It  doesn't matter how many
spaces    you type; one is as good as 100.  The
end of   a line counts as a space.

One   or more   blank lines denote the  end
of  a paragraph.

Since any number of consecutive spaces are treated
like a single one, the formatting of the input
file makes no difference to
      \LaTeX,                % The \LaTeX command generates the LaTeX logo.
but it makes a difference to you.  When you use
\LaTeX, making your input file as easy to read
as possible will be a great help as you write
your document and when you change it.  This sample
file shows how you can add comments to your own input
file.

Because printing is different from typewriting,
there are a number of things that you have to do
differently when preparing an input file than if
you were just typing the document directly.
Quotation marks like
       ``this''
have to be handled specially, as do quotes within
quotes:
       ``\,`this'            % \, separates the double and single quote.
        is what I just
        wrote, not  `that'\,''.

Dashes come in three sizes: an
       intra-word
dash, a medium dash for number ranges like
       1--2,
and a punctuation
       dash---like
this.

A sentence-ending space should be larger than the
space between words within a sentence.  You
sometimes have to type special commands in
conjunction with punctuation characters to get
this right, as in the following sentence.
       Gnats, gnus, etc.\ all  % `\ ' makes an inter-word space.
       begin with G\@.         % \@ marks end-of-sentence punctuation.
You should check the spaces after periods when
reading your output to make sure you haven't
forgotten any special cases.  Generating an
ellipsis
       \ldots\               % `\ ' is needed after `\ldots' because TeX
                             % ignores spaces after command names like \ldots
                             % made from \ + letters.
                             %
                             % Note how a `%' character causes TeX to ignore
                             % the end of the input line, so these blank lines
                             % do not start a new paragraph.
                             %
with the right spacing around the periods requires
a special command.

\LaTeX\ interprets some common characters as
commands, so you must type special commands to
generate them.  These characters include the
following:
       \$ \& \% \# \{ and \}.

In printing, text is usually emphasized with an
       \emph{italic}
type style.

\begin{em}
   A long segment of text can also be emphasized
   in this way.  Text within such a segment can be
   given \emph{additional} emphasis.
\end{em}

It is sometimes necessary to prevent \LaTeX\ from
breaking a line where it might otherwise do so.
This may be at a space, as between the ``Mr.''\ and
``Jones'' in
       ``Mr.~Jones'',        % ~ produces an unbreakable interword space.
or within a word---especially when the word is a
symbol like
       \mbox{\emph{itemnum}}
that makes little sense when hyphenated across
lines.

Footnotes\footnote{This is an example of a footnote.}
pose no problem.

\LaTeX\ is good at typesetting mathematical formulas
like
       \( x-3y + z = 7 \)
or
       \( a_{1} > x^{2n} + y^{2n} > x' \)
or
       \( \ip{A}{B} = \sum_{i} a_{i} b_{i} \).
The spaces you type in a formula are
ignored.  Remember that a letter like
       $x$                   % $ ... $  and  \( ... \)  are equivalent
is a formula when it denotes a mathematical
symbol, and it should be typed as one.

\section{Displayed Text}

Text is displayed by indenting it from the left
margin.  Quotations are commonly displayed.  There
are short quotations
\begin{quote}
   This is a short quotation.  It consists of a
   single paragraph of text.  See how it is formatted.
\end{quote}
and longer ones.
\begin{quotation}
   This is a longer quotation.  It consists of two
   paragraphs of text, neither of which are
   particularly interesting.

   This is the second paragraph of the quotation.  It
   is just as dull as the first paragraph.
\end{quotation}
Another frequently-displayed structure is a list.
The following is an example of an \emph{itemized}
list.
\begin{itemize}
   \item This is the first item of an itemized list.
         Each item in the list is marked with a ``tick''.
         You don't have to worry about what kind of tick
         mark is used.

   \item This is the second item of the list.  It
         contains another list nested inside it.  The inner
         list is an \emph{enumerated} list.
         \begin{enumerate}
            \item This is the first item of an enumerated
                  list that is nested within the itemized list.

            \item This is the second item of the inner list.
                  \LaTeX\ allows you to nest lists deeper than
                  you really should.
         \end{enumerate}
         This is the rest of the second item of the outer
         list.  It is no more interesting than any other
         part of the item.
   \item This is the third item of the list.
\end{itemize}
You can even display poetry.
\begin{verse}
   There is an environment
    for verse \\             % The \\ command separates lines
   Whose features some poets % within a stanza.
   will curse.

                             % One or more blank lines separate stanzas.

   For instead of making\\
   Them do \emph{all} line breaking, \\
   It allows them to put too many words on a line when they'd rather be
   forced to be terse.
\end{verse}

Mathematical formulas may also be displayed.  A
displayed formula
is
one-line long; multiline
formulas require special formatting instructions.
   \[  \ip{\Gamma}{\psi'} = x'' + y^{2} + z_{i}^{n}\]
Don't start a paragraph with a displayed equation,
nor make one a paragraph by itself.

\end{document}               % End of document.

\section{Bench Section 21}
% synthetic bench section 21
% This is a sample LaTeX input file.  (Version of 12 August 2004.)
%
% A '%' character causes TeX to ignore all remaining text on the line,
% and is used for comments like this one.

\documentclass{article}      % Specifies the document class

                             % The preamble begins here.
\title{An Example Document}  % Declares the document's title.
\author{Leslie Lamport}      % Declares the author's name.
\date{January 21, 1994}      % Deleting this command produces today's date.

\newcommand{\ip}[2]{(#1, #2)}
                             % Defines \ip{arg1}{arg2} to mean
                             % (arg1, arg2).

%\newcommand{\ip}[2]{\langle #1 | #2\rangle}
                             % This is an alternative definition of
                             % \ip that is commented out.

\begin{document}             % End of preamble and beginning of text.

\maketitle                   % Produces the title.

This is an example input file.  Comparing it with
the output it generates can show you how to
produce a simple document of your own.

\section{Ordinary Text}      % Produces section heading.  Lower-level
                             % sections are begun with similar
                             % \subsection and \subsubsection commands.

The ends  of words and sentences are marked
  by   spaces. It  doesn't matter how many
spaces    you type; one is as good as 100.  The
end of   a line counts as a space.

One   or more   blank lines denote the  end
of  a paragraph.

Since any number of consecutive spaces are treated
like a single one, the formatting of the input
file makes no difference to
      \LaTeX,                % The \LaTeX command generates the LaTeX logo.
but it makes a difference to you.  When you use
\LaTeX, making your input file as easy to read
as possible will be a great help as you write
your document and when you change it.  This sample
file shows how you can add comments to your own input
file.

Because printing is different from typewriting,
there are a number of things that you have to do
differently when preparing an input file than if
you were just typing the document directly.
Quotation marks like
       ``this''
have to be handled specially, as do quotes within
quotes:
       ``\,`this'            % \, separates the double and single quote.
        is what I just
        wrote, not  `that'\,''.

Dashes come in three sizes: an
       intra-word
dash, a medium dash for number ranges like
       1--2,
and a punctuation
       dash---like
this.

A sentence-ending space should be larger than the
space between words within a sentence.  You
sometimes have to type special commands in
conjunction with punctuation characters to get
this right, as in the following sentence.
       Gnats, gnus, etc.\ all  % `\ ' makes an inter-word space.
       begin with G\@.         % \@ marks end-of-sentence punctuation.
You should check the spaces after periods when
reading your output to make sure you haven't
forgotten any special cases.  Generating an
ellipsis
       \ldots\               % `\ ' is needed after `\ldots' because TeX
                             % ignores spaces after command names like \ldots
                             % made from \ + letters.
                             %
                             % Note how a `%' character causes TeX to ignore
                             % the end of the input line, so these blank lines
                             % do not start a new paragraph.
                             %
with the right spacing around the periods requires
a special command.

\LaTeX\ interprets some common characters as
commands, so you must type special commands to
generate them.  These characters include the
following:
       \$ \& \% \# \{ and \}.

In printing, text is usually emphasized with an
       \emph{italic}
type style.

\begin{em}
   A long segment of text can also be emphasized
   in this way.  Text within such a segment can be
   given \emph{additional} emphasis.
\end{em}

It is sometimes necessary to prevent \LaTeX\ from
breaking a line where it might otherwise do so.
This may be at a space, as between the ``Mr.''\ and
``Jones'' in
       ``Mr.~Jones'',        % ~ produces an unbreakable interword space.
or within a word---especially when the word is a
symbol like
       \mbox{\emph{itemnum}}
that makes little sense when hyphenated across
lines.

Footnotes\footnote{This is an example of a footnote.}
pose no problem.

\LaTeX\ is good at typesetting mathematical formulas
like
       \( x-3y + z = 7 \)
or
       \( a_{1} > x^{2n} + y^{2n} > x' \)
or
       \( \ip{A}{B} = \sum_{i} a_{i} b_{i} \).
The spaces you type in a formula are
ignored.  Remember that a letter like
       $x$                   % $ ... $  and  \( ... \)  are equivalent
is a formula when it denotes a mathematical
symbol, and it should be typed as one.

\section{Displayed Text}

Text is displayed by indenting it from the left
margin.  Quotations are commonly displayed.  There
are short quotations
\begin{quote}
   This is a short quotation.  It consists of a
   single paragraph of text.  See how it is formatted.
\end{quote}
and longer ones.
\begin{quotation}
   This is a longer quotation.  It consists of two
   paragraphs of text, neither of which are
   particularly interesting.

   This is the second paragraph of the quotation.  It
   is just as dull as the first paragraph.
\end{quotation}
Another frequently-displayed structure is a list.
The following is an example of an \emph{itemized}
list.
\begin{itemize}
   \item This is the first item of an itemized list.
         Each item in the list is marked with a ``tick''.
         You don't have to worry about what kind of tick
         mark is used.

   \item This is the second item of the list.  It
         contains another list nested inside it.  The inner
         list is an \emph{enumerated} list.
         \begin{enumerate}
            \item This is the first item of an enumerated
                  list that is nested within the itemized list.

            \item This is the second item of the inner list.
                  \LaTeX\ allows you to nest lists deeper than
                  you really should.
         \end{enumerate}
         This is the rest of the second item of the outer
         list.  It is no more interesting than any other
         part of the item.
   \item This is the third item of the list.
\end{itemize}
You can even display poetry.
\begin{verse}
   There is an environment
    for verse \\             % The \\ command separates lines
   Whose features some poets % within a stanza.
   will curse.

                             % One or more blank lines separate stanzas.

   For instead of making\\
   Them do \emph{all} line breaking, \\
   It allows them to put too many words on a line when they'd rather be
   forced to be terse.
\end{verse}

Mathematical formulas may also be displayed.  A
displayed formula
is
one-line long; multiline
formulas require special formatting instructions.
   \[  \ip{\Gamma}{\psi'} = x'' + y^{2} + z_{i}^{n}\]
Don't start a paragraph with a displayed equation,
nor make one a paragraph by itself.

\end{document}               % End of document.

\section{Bench Section 22}
% synthetic bench section 22
% This is a sample LaTeX input file.  (Version of 12 August 2004.)
%
% A '%' character causes TeX to ignore all remaining text on the line,
% and is used for comments like this one.

\documentclass{article}      % Specifies the document class

                             % The preamble begins here.
\title{An Example Document}  % Declares the document's title.
\author{Leslie Lamport}      % Declares the author's name.
\date{January 21, 1994}      % Deleting this command produces today's date.

\newcommand{\ip}[2]{(#1, #2)}
                             % Defines \ip{arg1}{arg2} to mean
                             % (arg1, arg2).

%\newcommand{\ip}[2]{\langle #1 | #2\rangle}
                             % This is an alternative definition of
                             % \ip that is commented out.

\begin{document}             % End of preamble and beginning of text.

\maketitle                   % Produces the title.

This is an example input file.  Comparing it with
the output it generates can show you how to
produce a simple document of your own.

\section{Ordinary Text}      % Produces section heading.  Lower-level
                             % sections are begun with similar
                             % \subsection and \subsubsection commands.

The ends  of words and sentences are marked
  by   spaces. It  doesn't matter how many
spaces    you type; one is as good as 100.  The
end of   a line counts as a space.

One   or more   blank lines denote the  end
of  a paragraph.

Since any number of consecutive spaces are treated
like a single one, the formatting of the input
file makes no difference to
      \LaTeX,                % The \LaTeX command generates the LaTeX logo.
but it makes a difference to you.  When you use
\LaTeX, making your input file as easy to read
as possible will be a great help as you write
your document and when you change it.  This sample
file shows how you can add comments to your own input
file.

Because printing is different from typewriting,
there are a number of things that you have to do
differently when preparing an input file than if
you were just typing the document directly.
Quotation marks like
       ``this''
have to be handled specially, as do quotes within
quotes:
       ``\,`this'            % \, separates the double and single quote.
        is what I just
        wrote, not  `that'\,''.

Dashes come in three sizes: an
       intra-word
dash, a medium dash for number ranges like
       1--2,
and a punctuation
       dash---like
this.

A sentence-ending space should be larger than the
space between words within a sentence.  You
sometimes have to type special commands in
conjunction with punctuation characters to get
this right, as in the following sentence.
       Gnats, gnus, etc.\ all  % `\ ' makes an inter-word space.
       begin with G\@.         % \@ marks end-of-sentence punctuation.
You should check the spaces after periods when
reading your output to make sure you haven't
forgotten any special cases.  Generating an
ellipsis
       \ldots\               % `\ ' is needed after `\ldots' because TeX
                             % ignores spaces after command names like \ldots
                             % made from \ + letters.
                             %
                             % Note how a `%' character causes TeX to ignore
                             % the end of the input line, so these blank lines
                             % do not start a new paragraph.
                             %
with the right spacing around the periods requires
a special command.

\LaTeX\ interprets some common characters as
commands, so you must type special commands to
generate them.  These characters include the
following:
       \$ \& \% \# \{ and \}.

In printing, text is usually emphasized with an
       \emph{italic}
type style.

\begin{em}
   A long segment of text can also be emphasized
   in this way.  Text within such a segment can be
   given \emph{additional} emphasis.
\end{em}

It is sometimes necessary to prevent \LaTeX\ from
breaking a line where it might otherwise do so.
This may be at a space, as between the ``Mr.''\ and
``Jones'' in
       ``Mr.~Jones'',        % ~ produces an unbreakable interword space.
or within a word---especially when the word is a
symbol like
       \mbox{\emph{itemnum}}
that makes little sense when hyphenated across
lines.

Footnotes\footnote{This is an example of a footnote.}
pose no problem.

\LaTeX\ is good at typesetting mathematical formulas
like
       \( x-3y + z = 7 \)
or
       \( a_{1} > x^{2n} + y^{2n} > x' \)
or
       \( \ip{A}{B} = \sum_{i} a_{i} b_{i} \).
The spaces you type in a formula are
ignored.  Remember that a letter like
       $x$                   % $ ... $  and  \( ... \)  are equivalent
is a formula when it denotes a mathematical
symbol, and it should be typed as one.

\section{Displayed Text}

Text is displayed by indenting it from the left
margin.  Quotations are commonly displayed.  There
are short quotations
\begin{quote}
   This is a short quotation.  It consists of a
   single paragraph of text.  See how it is formatted.
\end{quote}
and longer ones.
\begin{quotation}
   This is a longer quotation.  It consists of two
   paragraphs of text, neither of which are
   particularly interesting.

   This is the second paragraph of the quotation.  It
   is just as dull as the first paragraph.
\end{quotation}
Another frequently-displayed structure is a list.
The following is an example of an \emph{itemized}
list.
\begin{itemize}
   \item This is the first item of an itemized list.
         Each item in the list is marked with a ``tick''.
         You don't have to worry about what kind of tick
         mark is used.

   \item This is the second item of the list.  It
         contains another list nested inside it.  The inner
         list is an \emph{enumerated} list.
         \begin{enumerate}
            \item This is the first item of an enumerated
                  list that is nested within the itemized list.

            \item This is the second item of the inner list.
                  \LaTeX\ allows you to nest lists deeper than
                  you really should.
         \end{enumerate}
         This is the rest of the second item of the outer
         list.  It is no more interesting than any other
         part of the item.
   \item This is the third item of the list.
\end{itemize}
You can even display poetry.
\begin{verse}
   There is an environment
    for verse \\             % The \\ command separates lines
   Whose features some poets % within a stanza.
   will curse.

                             % One or more blank lines separate stanzas.

   For instead of making\\
   Them do \emph{all} line breaking, \\
   It allows them to put too many words on a line when they'd rather be
   forced to be terse.
\end{verse}

Mathematical formulas may also be displayed.  A
displayed formula
is
one-line long; multiline
formulas require special formatting instructions.
   \[  \ip{\Gamma}{\psi'} = x'' + y^{2} + z_{i}^{n}\]
Don't start a paragraph with a displayed equation,
nor make one a paragraph by itself.

\end{document}               % End of document.

\section{Bench Section 23}
% synthetic bench section 23
% This is a sample LaTeX input file.  (Version of 12 August 2004.)
%
% A '%' character causes TeX to ignore all remaining text on the line,
% and is used for comments like this one.

\documentclass{article}      % Specifies the document class

                             % The preamble begins here.
\title{An Example Document}  % Declares the document's title.
\author{Leslie Lamport}      % Declares the author's name.
\date{January 21, 1994}      % Deleting this command produces today's date.

\newcommand{\ip}[2]{(#1, #2)}
                             % Defines \ip{arg1}{arg2} to mean
                             % (arg1, arg2).

%\newcommand{\ip}[2]{\langle #1 | #2\rangle}
                             % This is an alternative definition of
                             % \ip that is commented out.

\begin{document}             % End of preamble and beginning of text.

\maketitle                   % Produces the title.

This is an example input file.  Comparing it with
the output it generates can show you how to
produce a simple document of your own.

\section{Ordinary Text}      % Produces section heading.  Lower-level
                             % sections are begun with similar
                             % \subsection and \subsubsection commands.

The ends  of words and sentences are marked
  by   spaces. It  doesn't matter how many
spaces    you type; one is as good as 100.  The
end of   a line counts as a space.

One   or more   blank lines denote the  end
of  a paragraph.

Since any number of consecutive spaces are treated
like a single one, the formatting of the input
file makes no difference to
      \LaTeX,                % The \LaTeX command generates the LaTeX logo.
but it makes a difference to you.  When you use
\LaTeX, making your input file as easy to read
as possible will be a great help as you write
your document and when you change it.  This sample
file shows how you can add comments to your own input
file.

Because printing is different from typewriting,
there are a number of things that you have to do
differently when preparing an input file than if
you were just typing the document directly.
Quotation marks like
       ``this''
have to be handled specially, as do quotes within
quotes:
       ``\,`this'            % \, separates the double and single quote.
        is what I just
        wrote, not  `that'\,''.

Dashes come in three sizes: an
       intra-word
dash, a medium dash for number ranges like
       1--2,
and a punctuation
       dash---like
this.

A sentence-ending space should be larger than the
space between words within a sentence.  You
sometimes have to type special commands in
conjunction with punctuation characters to get
this right, as in the following sentence.
       Gnats, gnus, etc.\ all  % `\ ' makes an inter-word space.
       begin with G\@.         % \@ marks end-of-sentence punctuation.
You should check the spaces after periods when
reading your output to make sure you haven't
forgotten any special cases.  Generating an
ellipsis
       \ldots\               % `\ ' is needed after `\ldots' because TeX
                             % ignores spaces after command names like \ldots
                             % made from \ + letters.
                             %
                             % Note how a `%' character causes TeX to ignore
                             % the end of the input line, so these blank lines
                             % do not start a new paragraph.
                             %
with the right spacing around the periods requires
a special command.

\LaTeX\ interprets some common characters as
commands, so you must type special commands to
generate them.  These characters include the
following:
       \$ \& \% \# \{ and \}.

In printing, text is usually emphasized with an
       \emph{italic}
type style.

\begin{em}
   A long segment of text can also be emphasized
   in this way.  Text within such a segment can be
   given \emph{additional} emphasis.
\end{em}

It is sometimes necessary to prevent \LaTeX\ from
breaking a line where it might otherwise do so.
This may be at a space, as between the ``Mr.''\ and
``Jones'' in
       ``Mr.~Jones'',        % ~ produces an unbreakable interword space.
or within a word---especially when the word is a
symbol like
       \mbox{\emph{itemnum}}
that makes little sense when hyphenated across
lines.

Footnotes\footnote{This is an example of a footnote.}
pose no problem.

\LaTeX\ is good at typesetting mathematical formulas
like
       \( x-3y + z = 7 \)
or
       \( a_{1} > x^{2n} + y^{2n} > x' \)
or
       \( \ip{A}{B} = \sum_{i} a_{i} b_{i} \).
The spaces you type in a formula are
ignored.  Remember that a letter like
       $x$                   % $ ... $  and  \( ... \)  are equivalent
is a formula when it denotes a mathematical
symbol, and it should be typed as one.

\section{Displayed Text}

Text is displayed by indenting it from the left
margin.  Quotations are commonly displayed.  There
are short quotations
\begin{quote}
   This is a short quotation.  It consists of a
   single paragraph of text.  See how it is formatted.
\end{quote}
and longer ones.
\begin{quotation}
   This is a longer quotation.  It consists of two
   paragraphs of text, neither of which are
   particularly interesting.

   This is the second paragraph of the quotation.  It
   is just as dull as the first paragraph.
\end{quotation}
Another frequently-displayed structure is a list.
The following is an example of an \emph{itemized}
list.
\begin{itemize}
   \item This is the first item of an itemized list.
         Each item in the list is marked with a ``tick''.
         You don't have to worry about what kind of tick
         mark is used.

   \item This is the second item of the list.  It
         contains another list nested inside it.  The inner
         list is an \emph{enumerated} list.
         \begin{enumerate}
            \item This is the first item of an enumerated
                  list that is nested within the itemized list.

            \item This is the second item of the inner list.
                  \LaTeX\ allows you to nest lists deeper than
                  you really should.
         \end{enumerate}
         This is the rest of the second item of the outer
         list.  It is no more interesting than any other
         part of the item.
   \item This is the third item of the list.
\end{itemize}
You can even display poetry.
\begin{verse}
   There is an environment
    for verse \\             % The \\ command separates lines
   Whose features some poets % within a stanza.
   will curse.

                             % One or more blank lines separate stanzas.

   For instead of making\\
   Them do \emph{all} line breaking, \\
   It allows them to put too many words on a line when they'd rather be
   forced to be terse.
\end{verse}

Mathematical formulas may also be displayed.  A
displayed formula
is
one-line long; multiline
formulas require special formatting instructions.
   \[  \ip{\Gamma}{\psi'} = x'' + y^{2} + z_{i}^{n}\]
Don't start a paragraph with a displayed equation,
nor make one a paragraph by itself.

\end{document}               % End of document.

\section{Bench Section 24}
% synthetic bench section 24
% This is a sample LaTeX input file.  (Version of 12 August 2004.)
%
% A '%' character causes TeX to ignore all remaining text on the line,
% and is used for comments like this one.

\documentclass{article}      % Specifies the document class

                             % The preamble begins here.
\title{An Example Document}  % Declares the document's title.
\author{Leslie Lamport}      % Declares the author's name.
\date{January 21, 1994}      % Deleting this command produces today's date.

\newcommand{\ip}[2]{(#1, #2)}
                             % Defines \ip{arg1}{arg2} to mean
                             % (arg1, arg2).

%\newcommand{\ip}[2]{\langle #1 | #2\rangle}
                             % This is an alternative definition of
                             % \ip that is commented out.

\begin{document}             % End of preamble and beginning of text.

\maketitle                   % Produces the title.

This is an example input file.  Comparing it with
the output it generates can show you how to
produce a simple document of your own.

\section{Ordinary Text}      % Produces section heading.  Lower-level
                             % sections are begun with similar
                             % \subsection and \subsubsection commands.

The ends  of words and sentences are marked
  by   spaces. It  doesn't matter how many
spaces    you type; one is as good as 100.  The
end of   a line counts as a space.

One   or more   blank lines denote the  end
of  a paragraph.

Since any number of consecutive spaces are treated
like a single one, the formatting of the input
file makes no difference to
      \LaTeX,                % The \LaTeX command generates the LaTeX logo.
but it makes a difference to you.  When you use
\LaTeX, making your input file as easy to read
as possible will be a great help as you write
your document and when you change it.  This sample
file shows how you can add comments to your own input
file.

Because printing is different from typewriting,
there are a number of things that you have to do
differently when preparing an input file than if
you were just typing the document directly.
Quotation marks like
       ``this''
have to be handled specially, as do quotes within
quotes:
       ``\,`this'            % \, separates the double and single quote.
        is what I just
        wrote, not  `that'\,''.

Dashes come in three sizes: an
       intra-word
dash, a medium dash for number ranges like
       1--2,
and a punctuation
       dash---like
this.

A sentence-ending space should be larger than the
space between words within a sentence.  You
sometimes have to type special commands in
conjunction with punctuation characters to get
this right, as in the following sentence.
       Gnats, gnus, etc.\ all  % `\ ' makes an inter-word space.
       begin with G\@.         % \@ marks end-of-sentence punctuation.
You should check the spaces after periods when
reading your output to make sure you haven't
forgotten any special cases.  Generating an
ellipsis
       \ldots\               % `\ ' is needed after `\ldots' because TeX
                             % ignores spaces after command names like \ldots
                             % made from \ + letters.
                             %
                             % Note how a `%' character causes TeX to ignore
                             % the end of the input line, so these blank lines
                             % do not start a new paragraph.
                             %
with the right spacing around the periods requires
a special command.

\LaTeX\ interprets some common characters as
commands, so you must type special commands to
generate them.  These characters include the
following:
       \$ \& \% \# \{ and \}.

In printing, text is usually emphasized with an
       \emph{italic}
type style.

\begin{em}
   A long segment of text can also be emphasized
   in this way.  Text within such a segment can be
   given \emph{additional} emphasis.
\end{em}

It is sometimes necessary to prevent \LaTeX\ from
breaking a line where it might otherwise do so.
This may be at a space, as between the ``Mr.''\ and
``Jones'' in
       ``Mr.~Jones'',        % ~ produces an unbreakable interword space.
or within a word---especially when the word is a
symbol like
       \mbox{\emph{itemnum}}
that makes little sense when hyphenated across
lines.

Footnotes\footnote{This is an example of a footnote.}
pose no problem.

\LaTeX\ is good at typesetting mathematical formulas
like
       \( x-3y + z = 7 \)
or
       \( a_{1} > x^{2n} + y^{2n} > x' \)
or
       \( \ip{A}{B} = \sum_{i} a_{i} b_{i} \).
The spaces you type in a formula are
ignored.  Remember that a letter like
       $x$                   % $ ... $  and  \( ... \)  are equivalent
is a formula when it denotes a mathematical
symbol, and it should be typed as one.

\section{Displayed Text}

Text is displayed by indenting it from the left
margin.  Quotations are commonly displayed.  There
are short quotations
\begin{quote}
   This is a short quotation.  It consists of a
   single paragraph of text.  See how it is formatted.
\end{quote}
and longer ones.
\begin{quotation}
   This is a longer quotation.  It consists of two
   paragraphs of text, neither of which are
   particularly interesting.

   This is the second paragraph of the quotation.  It
   is just as dull as the first paragraph.
\end{quotation}
Another frequently-displayed structure is a list.
The following is an example of an \emph{itemized}
list.
\begin{itemize}
   \item This is the first item of an itemized list.
         Each item in the list is marked with a ``tick''.
         You don't have to worry about what kind of tick
         mark is used.

   \item This is the second item of the list.  It
         contains another list nested inside it.  The inner
         list is an \emph{enumerated} list.
         \begin{enumerate}
            \item This is the first item of an enumerated
                  list that is nested within the itemized list.

            \item This is the second item of the inner list.
                  \LaTeX\ allows you to nest lists deeper than
                  you really should.
         \end{enumerate}
         This is the rest of the second item of the outer
         list.  It is no more interesting than any other
         part of the item.
   \item This is the third item of the list.
\end{itemize}
You can even display poetry.
\begin{verse}
   There is an environment
    for verse \\             % The \\ command separates lines
   Whose features some poets % within a stanza.
   will curse.

                             % One or more blank lines separate stanzas.

   For instead of making\\
   Them do \emph{all} line breaking, \\
   It allows them to put too many words on a line when they'd rather be
   forced to be terse.
\end{verse}

Mathematical formulas may also be displayed.  A
displayed formula
is
one-line long; multiline
formulas require special formatting instructions.
   \[  \ip{\Gamma}{\psi'} = x'' + y^{2} + z_{i}^{n}\]
Don't start a paragraph with a displayed equation,
nor make one a paragraph by itself.

\end{document}               % End of document.

\section{Bench Section 25}
% synthetic bench section 25
% This is a sample LaTeX input file.  (Version of 12 August 2004.)
%
% A '%' character causes TeX to ignore all remaining text on the line,
% and is used for comments like this one.

\documentclass{article}      % Specifies the document class

                             % The preamble begins here.
\title{An Example Document}  % Declares the document's title.
\author{Leslie Lamport}      % Declares the author's name.
\date{January 21, 1994}      % Deleting this command produces today's date.

\newcommand{\ip}[2]{(#1, #2)}
                             % Defines \ip{arg1}{arg2} to mean
                             % (arg1, arg2).

%\newcommand{\ip}[2]{\langle #1 | #2\rangle}
                             % This is an alternative definition of
                             % \ip that is commented out.

\begin{document}             % End of preamble and beginning of text.

\maketitle                   % Produces the title.

This is an example input file.  Comparing it with
the output it generates can show you how to
produce a simple document of your own.

\section{Ordinary Text}      % Produces section heading.  Lower-level
                             % sections are begun with similar
                             % \subsection and \subsubsection commands.

The ends  of words and sentences are marked
  by   spaces. It  doesn't matter how many
spaces    you type; one is as good as 100.  The
end of   a line counts as a space.

One   or more   blank lines denote the  end
of  a paragraph.

Since any number of consecutive spaces are treated
like a single one, the formatting of the input
file makes no difference to
      \LaTeX,                % The \LaTeX command generates the LaTeX logo.
but it makes a difference to you.  When you use
\LaTeX, making your input file as easy to read
as possible will be a great help as you write
your document and when you change it.  This sample
file shows how you can add comments to your own input
file.

Because printing is different from typewriting,
there are a number of things that you have to do
differently when preparing an input file than if
you were just typing the document directly.
Quotation marks like
       ``this''
have to be handled specially, as do quotes within
quotes:
       ``\,`this'            % \, separates the double and single quote.
        is what I just
        wrote, not  `that'\,''.

Dashes come in three sizes: an
       intra-word
dash, a medium dash for number ranges like
       1--2,
and a punctuation
       dash---like
this.

A sentence-ending space should be larger than the
space between words within a sentence.  You
sometimes have to type special commands in
conjunction with punctuation characters to get
this right, as in the following sentence.
       Gnats, gnus, etc.\ all  % `\ ' makes an inter-word space.
       begin with G\@.         % \@ marks end-of-sentence punctuation.
You should check the spaces after periods when
reading your output to make sure you haven't
forgotten any special cases.  Generating an
ellipsis
       \ldots\               % `\ ' is needed after `\ldots' because TeX
                             % ignores spaces after command names like \ldots
                             % made from \ + letters.
                             %
                             % Note how a `%' character causes TeX to ignore
                             % the end of the input line, so these blank lines
                             % do not start a new paragraph.
                             %
with the right spacing around the periods requires
a special command.

\LaTeX\ interprets some common characters as
commands, so you must type special commands to
generate them.  These characters include the
following:
       \$ \& \% \# \{ and \}.

In printing, text is usually emphasized with an
       \emph{italic}
type style.

\begin{em}
   A long segment of text can also be emphasized
   in this way.  Text within such a segment can be
   given \emph{additional} emphasis.
\end{em}

It is sometimes necessary to prevent \LaTeX\ from
breaking a line where it might otherwise do so.
This may be at a space, as between the ``Mr.''\ and
``Jones'' in
       ``Mr.~Jones'',        % ~ produces an unbreakable interword space.
or within a word---especially when the word is a
symbol like
       \mbox{\emph{itemnum}}
that makes little sense when hyphenated across
lines.

Footnotes\footnote{This is an example of a footnote.}
pose no problem.

\LaTeX\ is good at typesetting mathematical formulas
like
       \( x-3y + z = 7 \)
or
       \( a_{1} > x^{2n} + y^{2n} > x' \)
or
       \( \ip{A}{B} = \sum_{i} a_{i} b_{i} \).
The spaces you type in a formula are
ignored.  Remember that a letter like
       $x$                   % $ ... $  and  \( ... \)  are equivalent
is a formula when it denotes a mathematical
symbol, and it should be typed as one.

\section{Displayed Text}

Text is displayed by indenting it from the left
margin.  Quotations are commonly displayed.  There
are short quotations
\begin{quote}
   This is a short quotation.  It consists of a
   single paragraph of text.  See how it is formatted.
\end{quote}
and longer ones.
\begin{quotation}
   This is a longer quotation.  It consists of two
   paragraphs of text, neither of which are
   particularly interesting.

   This is the second paragraph of the quotation.  It
   is just as dull as the first paragraph.
\end{quotation}
Another frequently-displayed structure is a list.
The following is an example of an \emph{itemized}
list.
\begin{itemize}
   \item This is the first item of an itemized list.
         Each item in the list is marked with a ``tick''.
         You don't have to worry about what kind of tick
         mark is used.

   \item This is the second item of the list.  It
         contains another list nested inside it.  The inner
         list is an \emph{enumerated} list.
         \begin{enumerate}
            \item This is the first item of an enumerated
                  list that is nested within the itemized list.

            \item This is the second item of the inner list.
                  \LaTeX\ allows you to nest lists deeper than
                  you really should.
         \end{enumerate}
         This is the rest of the second item of the outer
         list.  It is no more interesting than any other
         part of the item.
   \item This is the third item of the list.
\end{itemize}
You can even display poetry.
\begin{verse}
   There is an environment
    for verse \\             % The \\ command separates lines
   Whose features some poets % within a stanza.
   will curse.

                             % One or more blank lines separate stanzas.

   For instead of making\\
   Them do \emph{all} line breaking, \\
   It allows them to put too many words on a line when they'd rather be
   forced to be terse.
\end{verse}

Mathematical formulas may also be displayed.  A
displayed formula
is
one-line long; multiline
formulas require special formatting instructions.
   \[  \ip{\Gamma}{\psi'} = x'' + y^{2} + z_{i}^{n}\]
Don't start a paragraph with a displayed equation,
nor make one a paragraph by itself.

\end{document}               % End of document.

\section{Bench Section 26}
% synthetic bench section 26
% This is a sample LaTeX input file.  (Version of 12 August 2004.)
%
% A '%' character causes TeX to ignore all remaining text on the line,
% and is used for comments like this one.

\documentclass{article}      % Specifies the document class

                             % The preamble begins here.
\title{An Example Document}  % Declares the document's title.
\author{Leslie Lamport}      % Declares the author's name.
\date{January 21, 1994}      % Deleting this command produces today's date.

\newcommand{\ip}[2]{(#1, #2)}
                             % Defines \ip{arg1}{arg2} to mean
                             % (arg1, arg2).

%\newcommand{\ip}[2]{\langle #1 | #2\rangle}
                             % This is an alternative definition of
                             % \ip that is commented out.

\begin{document}             % End of preamble and beginning of text.

\maketitle                   % Produces the title.

This is an example input file.  Comparing it with
the output it generates can show you how to
produce a simple document of your own.

\section{Ordinary Text}      % Produces section heading.  Lower-level
                             % sections are begun with similar
                             % \subsection and \subsubsection commands.

The ends  of words and sentences are marked
  by   spaces. It  doesn't matter how many
spaces    you type; one is as good as 100.  The
end of   a line counts as a space.

One   or more   blank lines denote the  end
of  a paragraph.

Since any number of consecutive spaces are treated
like a single one, the formatting of the input
file makes no difference to
      \LaTeX,                % The \LaTeX command generates the LaTeX logo.
but it makes a difference to you.  When you use
\LaTeX, making your input file as easy to read
as possible will be a great help as you write
your document and when you change it.  This sample
file shows how you can add comments to your own input
file.

Because printing is different from typewriting,
there are a number of things that you have to do
differently when preparing an input file than if
you were just typing the document directly.
Quotation marks like
       ``this''
have to be handled specially, as do quotes within
quotes:
       ``\,`this'            % \, separates the double and single quote.
        is what I just
        wrote, not  `that'\,''.

Dashes come in three sizes: an
       intra-word
dash, a medium dash for number ranges like
       1--2,
and a punctuation
       dash---like
this.

A sentence-ending space should be larger than the
space between words within a sentence.  You
sometimes have to type special commands in
conjunction with punctuation characters to get
this right, as in the following sentence.
       Gnats, gnus, etc.\ all  % `\ ' makes an inter-word space.
       begin with G\@.         % \@ marks end-of-sentence punctuation.
You should check the spaces after periods when
reading your output to make sure you haven't
forgotten any special cases.  Generating an
ellipsis
       \ldots\               % `\ ' is needed after `\ldots' because TeX
                             % ignores spaces after command names like \ldots
                             % made from \ + letters.
                             %
                             % Note how a `%' character causes TeX to ignore
                             % the end of the input line, so these blank lines
                             % do not start a new paragraph.
                             %
with the right spacing around the periods requires
a special command.

\LaTeX\ interprets some common characters as
commands, so you must type special commands to
generate them.  These characters include the
following:
       \$ \& \% \# \{ and \}.

In printing, text is usually emphasized with an
       \emph{italic}
type style.

\begin{em}
   A long segment of text can also be emphasized
   in this way.  Text within such a segment can be
   given \emph{additional} emphasis.
\end{em}

It is sometimes necessary to prevent \LaTeX\ from
breaking a line where it might otherwise do so.
This may be at a space, as between the ``Mr.''\ and
``Jones'' in
       ``Mr.~Jones'',        % ~ produces an unbreakable interword space.
or within a word---especially when the word is a
symbol like
       \mbox{\emph{itemnum}}
that makes little sense when hyphenated across
lines.

Footnotes\footnote{This is an example of a footnote.}
pose no problem.

\LaTeX\ is good at typesetting mathematical formulas
like
       \( x-3y + z = 7 \)
or
       \( a_{1} > x^{2n} + y^{2n} > x' \)
or
       \( \ip{A}{B} = \sum_{i} a_{i} b_{i} \).
The spaces you type in a formula are
ignored.  Remember that a letter like
       $x$                   % $ ... $  and  \( ... \)  are equivalent
is a formula when it denotes a mathematical
symbol, and it should be typed as one.

\section{Displayed Text}

Text is displayed by indenting it from the left
margin.  Quotations are commonly displayed.  There
are short quotations
\begin{quote}
   This is a short quotation.  It consists of a
   single paragraph of text.  See how it is formatted.
\end{quote}
and longer ones.
\begin{quotation}
   This is a longer quotation.  It consists of two
   paragraphs of text, neither of which are
   particularly interesting.

   This is the second paragraph of the quotation.  It
   is just as dull as the first paragraph.
\end{quotation}
Another frequently-displayed structure is a list.
The following is an example of an \emph{itemized}
list.
\begin{itemize}
   \item This is the first item of an itemized list.
         Each item in the list is marked with a ``tick''.
         You don't have to worry about what kind of tick
         mark is used.

   \item This is the second item of the list.  It
         contains another list nested inside it.  The inner
         list is an \emph{enumerated} list.
         \begin{enumerate}
            \item This is the first item of an enumerated
                  list that is nested within the itemized list.

            \item This is the second item of the inner list.
                  \LaTeX\ allows you to nest lists deeper than
                  you really should.
         \end{enumerate}
         This is the rest of the second item of the outer
         list.  It is no more interesting than any other
         part of the item.
   \item This is the third item of the list.
\end{itemize}
You can even display poetry.
\begin{verse}
   There is an environment
    for verse \\             % The \\ command separates lines
   Whose features some poets % within a stanza.
   will curse.

                             % One or more blank lines separate stanzas.

   For instead of making\\
   Them do \emph{all} line breaking, \\
   It allows them to put too many words on a line when they'd rather be
   forced to be terse.
\end{verse}

Mathematical formulas may also be displayed.  A
displayed formula
is
one-line long; multiline
formulas require special formatting instructions.
   \[  \ip{\Gamma}{\psi'} = x'' + y^{2} + z_{i}^{n}\]
Don't start a paragraph with a displayed equation,
nor make one a paragraph by itself.

\end{document}               % End of document.

\section{Bench Section 27}
% synthetic bench section 27
% This is a sample LaTeX input file.  (Version of 12 August 2004.)
%
% A '%' character causes TeX to ignore all remaining text on the line,
% and is used for comments like this one.

\documentclass{article}      % Specifies the document class

                             % The preamble begins here.
\title{An Example Document}  % Declares the document's title.
\author{Leslie Lamport}      % Declares the author's name.
\date{January 21, 1994}      % Deleting this command produces today's date.

\newcommand{\ip}[2]{(#1, #2)}
                             % Defines \ip{arg1}{arg2} to mean
                             % (arg1, arg2).

%\newcommand{\ip}[2]{\langle #1 | #2\rangle}
                             % This is an alternative definition of
                             % \ip that is commented out.

\begin{document}             % End of preamble and beginning of text.

\maketitle                   % Produces the title.

This is an example input file.  Comparing it with
the output it generates can show you how to
produce a simple document of your own.

\section{Ordinary Text}      % Produces section heading.  Lower-level
                             % sections are begun with similar
                             % \subsection and \subsubsection commands.

The ends  of words and sentences are marked
  by   spaces. It  doesn't matter how many
spaces    you type; one is as good as 100.  The
end of   a line counts as a space.

One   or more   blank lines denote the  end
of  a paragraph.

Since any number of consecutive spaces are treated
like a single one, the formatting of the input
file makes no difference to
      \LaTeX,                % The \LaTeX command generates the LaTeX logo.
but it makes a difference to you.  When you use
\LaTeX, making your input file as easy to read
as possible will be a great help as you write
your document and when you change it.  This sample
file shows how you can add comments to your own input
file.

Because printing is different from typewriting,
there are a number of things that you have to do
differently when preparing an input file than if
you were just typing the document directly.
Quotation marks like
       ``this''
have to be handled specially, as do quotes within
quotes:
       ``\,`this'            % \, separates the double and single quote.
        is what I just
        wrote, not  `that'\,''.

Dashes come in three sizes: an
       intra-word
dash, a medium dash for number ranges like
       1--2,
and a punctuation
       dash---like
this.

A sentence-ending space should be larger than the
space between words within a sentence.  You
sometimes have to type special commands in
conjunction with punctuation characters to get
this right, as in the following sentence.
       Gnats, gnus, etc.\ all  % `\ ' makes an inter-word space.
       begin with G\@.         % \@ marks end-of-sentence punctuation.
You should check the spaces after periods when
reading your output to make sure you haven't
forgotten any special cases.  Generating an
ellipsis
       \ldots\               % `\ ' is needed after `\ldots' because TeX
                             % ignores spaces after command names like \ldots
                             % made from \ + letters.
                             %
                             % Note how a `%' character causes TeX to ignore
                             % the end of the input line, so these blank lines
                             % do not start a new paragraph.
                             %
with the right spacing around the periods requires
a special command.

\LaTeX\ interprets some common characters as
commands, so you must type special commands to
generate them.  These characters include the
following:
       \$ \& \% \# \{ and \}.

In printing, text is usually emphasized with an
       \emph{italic}
type style.

\begin{em}
   A long segment of text can also be emphasized
   in this way.  Text within such a segment can be
   given \emph{additional} emphasis.
\end{em}

It is sometimes necessary to prevent \LaTeX\ from
breaking a line where it might otherwise do so.
This may be at a space, as between the ``Mr.''\ and
``Jones'' in
       ``Mr.~Jones'',        % ~ produces an unbreakable interword space.
or within a word---especially when the word is a
symbol like
       \mbox{\emph{itemnum}}
that makes little sense when hyphenated across
lines.

Footnotes\footnote{This is an example of a footnote.}
pose no problem.

\LaTeX\ is good at typesetting mathematical formulas
like
       \( x-3y + z = 7 \)
or
       \( a_{1} > x^{2n} + y^{2n} > x' \)
or
       \( \ip{A}{B} = \sum_{i} a_{i} b_{i} \).
The spaces you type in a formula are
ignored.  Remember that a letter like
       $x$                   % $ ... $  and  \( ... \)  are equivalent
is a formula when it denotes a mathematical
symbol, and it should be typed as one.

\section{Displayed Text}

Text is displayed by indenting it from the left
margin.  Quotations are commonly displayed.  There
are short quotations
\begin{quote}
   This is a short quotation.  It consists of a
   single paragraph of text.  See how it is formatted.
\end{quote}
and longer ones.
\begin{quotation}
   This is a longer quotation.  It consists of two
   paragraphs of text, neither of which are
   particularly interesting.

   This is the second paragraph of the quotation.  It
   is just as dull as the first paragraph.
\end{quotation}
Another frequently-displayed structure is a list.
The following is an example of an \emph{itemized}
list.
\begin{itemize}
   \item This is the first item of an itemized list.
         Each item in the list is marked with a ``tick''.
         You don't have to worry about what kind of tick
         mark is used.

   \item This is the second item of the list.  It
         contains another list nested inside it.  The inner
         list is an \emph{enumerated} list.
         \begin{enumerate}
            \item This is the first item of an enumerated
                  list that is nested within the itemized list.

            \item This is the second item of the inner list.
                  \LaTeX\ allows you to nest lists deeper than
                  you really should.
         \end{enumerate}
         This is the rest of the second item of the outer
         list.  It is no more interesting than any other
         part of the item.
   \item This is the third item of the list.
\end{itemize}
You can even display poetry.
\begin{verse}
   There is an environment
    for verse \\             % The \\ command separates lines
   Whose features some poets % within a stanza.
   will curse.

                             % One or more blank lines separate stanzas.

   For instead of making\\
   Them do \emph{all} line breaking, \\
   It allows them to put too many words on a line when they'd rather be
   forced to be terse.
\end{verse}

Mathematical formulas may also be displayed.  A
displayed formula
is
one-line long; multiline
formulas require special formatting instructions.
   \[  \ip{\Gamma}{\psi'} = x'' + y^{2} + z_{i}^{n}\]
Don't start a paragraph with a displayed equation,
nor make one a paragraph by itself.

\end{document}               % End of document.

\section{Bench Section 28}
% synthetic bench section 28
% This is a sample LaTeX input file.  (Version of 12 August 2004.)
%
% A '%' character causes TeX to ignore all remaining text on the line,
% and is used for comments like this one.

\documentclass{article}      % Specifies the document class

                             % The preamble begins here.
\title{An Example Document}  % Declares the document's title.
\author{Leslie Lamport}      % Declares the author's name.
\date{January 21, 1994}      % Deleting this command produces today's date.

\newcommand{\ip}[2]{(#1, #2)}
                             % Defines \ip{arg1}{arg2} to mean
                             % (arg1, arg2).

%\newcommand{\ip}[2]{\langle #1 | #2\rangle}
                             % This is an alternative definition of
                             % \ip that is commented out.

\begin{document}             % End of preamble and beginning of text.

\maketitle                   % Produces the title.

This is an example input file.  Comparing it with
the output it generates can show you how to
produce a simple document of your own.

\section{Ordinary Text}      % Produces section heading.  Lower-level
                             % sections are begun with similar
                             % \subsection and \subsubsection commands.

The ends  of words and sentences are marked
  by   spaces. It  doesn't matter how many
spaces    you type; one is as good as 100.  The
end of   a line counts as a space.

One   or more   blank lines denote the  end
of  a paragraph.

Since any number of consecutive spaces are treated
like a single one, the formatting of the input
file makes no difference to
      \LaTeX,                % The \LaTeX command generates the LaTeX logo.
but it makes a difference to you.  When you use
\LaTeX, making your input file as easy to read
as possible will be a great help as you write
your document and when you change it.  This sample
file shows how you can add comments to your own input
file.

Because printing is different from typewriting,
there are a number of things that you have to do
differently when preparing an input file than if
you were just typing the document directly.
Quotation marks like
       ``this''
have to be handled specially, as do quotes within
quotes:
       ``\,`this'            % \, separates the double and single quote.
        is what I just
        wrote, not  `that'\,''.

Dashes come in three sizes: an
       intra-word
dash, a medium dash for number ranges like
       1--2,
and a punctuation
       dash---like
this.

A sentence-ending space should be larger than the
space between words within a sentence.  You
sometimes have to type special commands in
conjunction with punctuation characters to get
this right, as in the following sentence.
       Gnats, gnus, etc.\ all  % `\ ' makes an inter-word space.
       begin with G\@.         % \@ marks end-of-sentence punctuation.
You should check the spaces after periods when
reading your output to make sure you haven't
forgotten any special cases.  Generating an
ellipsis
       \ldots\               % `\ ' is needed after `\ldots' because TeX
                             % ignores spaces after command names like \ldots
                             % made from \ + letters.
                             %
                             % Note how a `%' character causes TeX to ignore
                             % the end of the input line, so these blank lines
                             % do not start a new paragraph.
                             %
with the right spacing around the periods requires
a special command.

\LaTeX\ interprets some common characters as
commands, so you must type special commands to
generate them.  These characters include the
following:
       \$ \& \% \# \{ and \}.

In printing, text is usually emphasized with an
       \emph{italic}
type style.

\begin{em}
   A long segment of text can also be emphasized
   in this way.  Text within such a segment can be
   given \emph{additional} emphasis.
\end{em}

It is sometimes necessary to prevent \LaTeX\ from
breaking a line where it might otherwise do so.
This may be at a space, as between the ``Mr.''\ and
``Jones'' in
       ``Mr.~Jones'',        % ~ produces an unbreakable interword space.
or within a word---especially when the word is a
symbol like
       \mbox{\emph{itemnum}}
that makes little sense when hyphenated across
lines.

Footnotes\footnote{This is an example of a footnote.}
pose no problem.

\LaTeX\ is good at typesetting mathematical formulas
like
       \( x-3y + z = 7 \)
or
       \( a_{1} > x^{2n} + y^{2n} > x' \)
or
       \( \ip{A}{B} = \sum_{i} a_{i} b_{i} \).
The spaces you type in a formula are
ignored.  Remember that a letter like
       $x$                   % $ ... $  and  \( ... \)  are equivalent
is a formula when it denotes a mathematical
symbol, and it should be typed as one.

\section{Displayed Text}

Text is displayed by indenting it from the left
margin.  Quotations are commonly displayed.  There
are short quotations
\begin{quote}
   This is a short quotation.  It consists of a
   single paragraph of text.  See how it is formatted.
\end{quote}
and longer ones.
\begin{quotation}
   This is a longer quotation.  It consists of two
   paragraphs of text, neither of which are
   particularly interesting.

   This is the second paragraph of the quotation.  It
   is just as dull as the first paragraph.
\end{quotation}
Another frequently-displayed structure is a list.
The following is an example of an \emph{itemized}
list.
\begin{itemize}
   \item This is the first item of an itemized list.
         Each item in the list is marked with a ``tick''.
         You don't have to worry about what kind of tick
         mark is used.

   \item This is the second item of the list.  It
         contains another list nested inside it.  The inner
         list is an \emph{enumerated} list.
         \begin{enumerate}
            \item This is the first item of an enumerated
                  list that is nested within the itemized list.

            \item This is the second item of the inner list.
                  \LaTeX\ allows you to nest lists deeper than
                  you really should.
         \end{enumerate}
         This is the rest of the second item of the outer
         list.  It is no more interesting than any other
         part of the item.
   \item This is the third item of the list.
\end{itemize}
You can even display poetry.
\begin{verse}
   There is an environment
    for verse \\             % The \\ command separates lines
   Whose features some poets % within a stanza.
   will curse.

                             % One or more blank lines separate stanzas.

   For instead of making\\
   Them do \emph{all} line breaking, \\
   It allows them to put too many words on a line when they'd rather be
   forced to be terse.
\end{verse}

Mathematical formulas may also be displayed.  A
displayed formula
is
one-line long; multiline
formulas require special formatting instructions.
   \[  \ip{\Gamma}{\psi'} = x'' + y^{2} + z_{i}^{n}\]
Don't start a paragraph with a displayed equation,
nor make one a paragraph by itself.

\end{document}               % End of document.

\section{Bench Section 29}
% synthetic bench section 29
% This is a sample LaTeX input file.  (Version of 12 August 2004.)
%
% A '%' character causes TeX to ignore all remaining text on the line,
% and is used for comments like this one.

\documentclass{article}      % Specifies the document class

                             % The preamble begins here.
\title{An Example Document}  % Declares the document's title.
\author{Leslie Lamport}      % Declares the author's name.
\date{January 21, 1994}      % Deleting this command produces today's date.

\newcommand{\ip}[2]{(#1, #2)}
                             % Defines \ip{arg1}{arg2} to mean
                             % (arg1, arg2).

%\newcommand{\ip}[2]{\langle #1 | #2\rangle}
                             % This is an alternative definition of
                             % \ip that is commented out.

\begin{document}             % End of preamble and beginning of text.

\maketitle                   % Produces the title.

This is an example input file.  Comparing it with
the output it generates can show you how to
produce a simple document of your own.

\section{Ordinary Text}      % Produces section heading.  Lower-level
                             % sections are begun with similar
                             % \subsection and \subsubsection commands.

The ends  of words and sentences are marked
  by   spaces. It  doesn't matter how many
spaces    you type; one is as good as 100.  The
end of   a line counts as a space.

One   or more   blank lines denote the  end
of  a paragraph.

Since any number of consecutive spaces are treated
like a single one, the formatting of the input
file makes no difference to
      \LaTeX,                % The \LaTeX command generates the LaTeX logo.
but it makes a difference to you.  When you use
\LaTeX, making your input file as easy to read
as possible will be a great help as you write
your document and when you change it.  This sample
file shows how you can add comments to your own input
file.

Because printing is different from typewriting,
there are a number of things that you have to do
differently when preparing an input file than if
you were just typing the document directly.
Quotation marks like
       ``this''
have to be handled specially, as do quotes within
quotes:
       ``\,`this'            % \, separates the double and single quote.
        is what I just
        wrote, not  `that'\,''.

Dashes come in three sizes: an
       intra-word
dash, a medium dash for number ranges like
       1--2,
and a punctuation
       dash---like
this.

A sentence-ending space should be larger than the
space between words within a sentence.  You
sometimes have to type special commands in
conjunction with punctuation characters to get
this right, as in the following sentence.
       Gnats, gnus, etc.\ all  % `\ ' makes an inter-word space.
       begin with G\@.         % \@ marks end-of-sentence punctuation.
You should check the spaces after periods when
reading your output to make sure you haven't
forgotten any special cases.  Generating an
ellipsis
       \ldots\               % `\ ' is needed after `\ldots' because TeX
                             % ignores spaces after command names like \ldots
                             % made from \ + letters.
                             %
                             % Note how a `%' character causes TeX to ignore
                             % the end of the input line, so these blank lines
                             % do not start a new paragraph.
                             %
with the right spacing around the periods requires
a special command.

\LaTeX\ interprets some common characters as
commands, so you must type special commands to
generate them.  These characters include the
following:
       \$ \& \% \# \{ and \}.

In printing, text is usually emphasized with an
       \emph{italic}
type style.

\begin{em}
   A long segment of text can also be emphasized
   in this way.  Text within such a segment can be
   given \emph{additional} emphasis.
\end{em}

It is sometimes necessary to prevent \LaTeX\ from
breaking a line where it might otherwise do so.
This may be at a space, as between the ``Mr.''\ and
``Jones'' in
       ``Mr.~Jones'',        % ~ produces an unbreakable interword space.
or within a word---especially when the word is a
symbol like
       \mbox{\emph{itemnum}}
that makes little sense when hyphenated across
lines.

Footnotes\footnote{This is an example of a footnote.}
pose no problem.

\LaTeX\ is good at typesetting mathematical formulas
like
       \( x-3y + z = 7 \)
or
       \( a_{1} > x^{2n} + y^{2n} > x' \)
or
       \( \ip{A}{B} = \sum_{i} a_{i} b_{i} \).
The spaces you type in a formula are
ignored.  Remember that a letter like
       $x$                   % $ ... $  and  \( ... \)  are equivalent
is a formula when it denotes a mathematical
symbol, and it should be typed as one.

\section{Displayed Text}

Text is displayed by indenting it from the left
margin.  Quotations are commonly displayed.  There
are short quotations
\begin{quote}
   This is a short quotation.  It consists of a
   single paragraph of text.  See how it is formatted.
\end{quote}
and longer ones.
\begin{quotation}
   This is a longer quotation.  It consists of two
   paragraphs of text, neither of which are
   particularly interesting.

   This is the second paragraph of the quotation.  It
   is just as dull as the first paragraph.
\end{quotation}
Another frequently-displayed structure is a list.
The following is an example of an \emph{itemized}
list.
\begin{itemize}
   \item This is the first item of an itemized list.
         Each item in the list is marked with a ``tick''.
         You don't have to worry about what kind of tick
         mark is used.

   \item This is the second item of the list.  It
         contains another list nested inside it.  The inner
         list is an \emph{enumerated} list.
         \begin{enumerate}
            \item This is the first item of an enumerated
                  list that is nested within the itemized list.

            \item This is the second item of the inner list.
                  \LaTeX\ allows you to nest lists deeper than
                  you really should.
         \end{enumerate}
         This is the rest of the second item of the outer
         list.  It is no more interesting than any other
         part of the item.
   \item This is the third item of the list.
\end{itemize}
You can even display poetry.
\begin{verse}
   There is an environment
    for verse \\             % The \\ command separates lines
   Whose features some poets % within a stanza.
   will curse.

                             % One or more blank lines separate stanzas.

   For instead of making\\
   Them do \emph{all} line breaking, \\
   It allows them to put too many words on a line when they'd rather be
   forced to be terse.
\end{verse}

Mathematical formulas may also be displayed.  A
displayed formula
is
one-line long; multiline
formulas require special formatting instructions.
   \[  \ip{\Gamma}{\psi'} = x'' + y^{2} + z_{i}^{n}\]
Don't start a paragraph with a displayed equation,
nor make one a paragraph by itself.

\end{document}               % End of document.

\section{Bench Section 30}
% synthetic bench section 30
% This is a sample LaTeX input file.  (Version of 12 August 2004.)
%
% A '%' character causes TeX to ignore all remaining text on the line,
% and is used for comments like this one.

\documentclass{article}      % Specifies the document class

                             % The preamble begins here.
\title{An Example Document}  % Declares the document's title.
\author{Leslie Lamport}      % Declares the author's name.
\date{January 21, 1994}      % Deleting this command produces today's date.

\newcommand{\ip}[2]{(#1, #2)}
                             % Defines \ip{arg1}{arg2} to mean
                             % (arg1, arg2).

%\newcommand{\ip}[2]{\langle #1 | #2\rangle}
                             % This is an alternative definition of
                             % \ip that is commented out.

\begin{document}             % End of preamble and beginning of text.

\maketitle                   % Produces the title.

This is an example input file.  Comparing it with
the output it generates can show you how to
produce a simple document of your own.

\section{Ordinary Text}      % Produces section heading.  Lower-level
                             % sections are begun with similar
                             % \subsection and \subsubsection commands.

The ends  of words and sentences are marked
  by   spaces. It  doesn't matter how many
spaces    you type; one is as good as 100.  The
end of   a line counts as a space.

One   or more   blank lines denote the  end
of  a paragraph.

Since any number of consecutive spaces are treated
like a single one, the formatting of the input
file makes no difference to
      \LaTeX,                % The \LaTeX command generates the LaTeX logo.
but it makes a difference to you.  When you use
\LaTeX, making your input file as easy to read
as possible will be a great help as you write
your document and when you change it.  This sample
file shows how you can add comments to your own input
file.

Because printing is different from typewriting,
there are a number of things that you have to do
differently when preparing an input file than if
you were just typing the document directly.
Quotation marks like
       ``this''
have to be handled specially, as do quotes within
quotes:
       ``\,`this'            % \, separates the double and single quote.
        is what I just
        wrote, not  `that'\,''.

Dashes come in three sizes: an
       intra-word
dash, a medium dash for number ranges like
       1--2,
and a punctuation
       dash---like
this.

A sentence-ending space should be larger than the
space between words within a sentence.  You
sometimes have to type special commands in
conjunction with punctuation characters to get
this right, as in the following sentence.
       Gnats, gnus, etc.\ all  % `\ ' makes an inter-word space.
       begin with G\@.         % \@ marks end-of-sentence punctuation.
You should check the spaces after periods when
reading your output to make sure you haven't
forgotten any special cases.  Generating an
ellipsis
       \ldots\               % `\ ' is needed after `\ldots' because TeX
                             % ignores spaces after command names like \ldots
                             % made from \ + letters.
                             %
                             % Note how a `%' character causes TeX to ignore
                             % the end of the input line, so these blank lines
                             % do not start a new paragraph.
                             %
with the right spacing around the periods requires
a special command.

\LaTeX\ interprets some common characters as
commands, so you must type special commands to
generate them.  These characters include the
following:
       \$ \& \% \# \{ and \}.

In printing, text is usually emphasized with an
       \emph{italic}
type style.

\begin{em}
   A long segment of text can also be emphasized
   in this way.  Text within such a segment can be
   given \emph{additional} emphasis.
\end{em}

It is sometimes necessary to prevent \LaTeX\ from
breaking a line where it might otherwise do so.
This may be at a space, as between the ``Mr.''\ and
``Jones'' in
       ``Mr.~Jones'',        % ~ produces an unbreakable interword space.
or within a word---especially when the word is a
symbol like
       \mbox{\emph{itemnum}}
that makes little sense when hyphenated across
lines.

Footnotes\footnote{This is an example of a footnote.}
pose no problem.

\LaTeX\ is good at typesetting mathematical formulas
like
       \( x-3y + z = 7 \)
or
       \( a_{1} > x^{2n} + y^{2n} > x' \)
or
       \( \ip{A}{B} = \sum_{i} a_{i} b_{i} \).
The spaces you type in a formula are
ignored.  Remember that a letter like
       $x$                   % $ ... $  and  \( ... \)  are equivalent
is a formula when it denotes a mathematical
symbol, and it should be typed as one.

\section{Displayed Text}

Text is displayed by indenting it from the left
margin.  Quotations are commonly displayed.  There
are short quotations
\begin{quote}
   This is a short quotation.  It consists of a
   single paragraph of text.  See how it is formatted.
\end{quote}
and longer ones.
\begin{quotation}
   This is a longer quotation.  It consists of two
   paragraphs of text, neither of which are
   particularly interesting.

   This is the second paragraph of the quotation.  It
   is just as dull as the first paragraph.
\end{quotation}
Another frequently-displayed structure is a list.
The following is an example of an \emph{itemized}
list.
\begin{itemize}
   \item This is the first item of an itemized list.
         Each item in the list is marked with a ``tick''.
         You don't have to worry about what kind of tick
         mark is used.

   \item This is the second item of the list.  It
         contains another list nested inside it.  The inner
         list is an \emph{enumerated} list.
         \begin{enumerate}
            \item This is the first item of an enumerated
                  list that is nested within the itemized list.

            \item This is the second item of the inner list.
                  \LaTeX\ allows you to nest lists deeper than
                  you really should.
         \end{enumerate}
         This is the rest of the second item of the outer
         list.  It is no more interesting than any other
         part of the item.
   \item This is the third item of the list.
\end{itemize}
You can even display poetry.
\begin{verse}
   There is an environment
    for verse \\             % The \\ command separates lines
   Whose features some poets % within a stanza.
   will curse.

                             % One or more blank lines separate stanzas.

   For instead of making\\
   Them do \emph{all} line breaking, \\
   It allows them to put too many words on a line when they'd rather be
   forced to be terse.
\end{verse}

Mathematical formulas may also be displayed.  A
displayed formula
is
one-line long; multiline
formulas require special formatting instructions.
   \[  \ip{\Gamma}{\psi'} = x'' + y^{2} + z_{i}^{n}\]
Don't start a paragraph with a displayed equation,
nor make one a paragraph by itself.

\end{document}               % End of document.

\section{Bench Section 31}
% synthetic bench section 31
% This is a sample LaTeX input file.  (Version of 12 August 2004.)
%
% A '%' character causes TeX to ignore all remaining text on the line,
% and is used for comments like this one.

\documentclass{article}      % Specifies the document class

                             % The preamble begins here.
\title{An Example Document}  % Declares the document's title.
\author{Leslie Lamport}      % Declares the author's name.
\date{January 21, 1994}      % Deleting this command produces today's date.

\newcommand{\ip}[2]{(#1, #2)}
                             % Defines \ip{arg1}{arg2} to mean
                             % (arg1, arg2).

%\newcommand{\ip}[2]{\langle #1 | #2\rangle}
                             % This is an alternative definition of
                             % \ip that is commented out.

\begin{document}             % End of preamble and beginning of text.

\maketitle                   % Produces the title.

This is an example input file.  Comparing it with
the output it generates can show you how to
produce a simple document of your own.

\section{Ordinary Text}      % Produces section heading.  Lower-level
                             % sections are begun with similar
                             % \subsection and \subsubsection commands.

The ends  of words and sentences are marked
  by   spaces. It  doesn't matter how many
spaces    you type; one is as good as 100.  The
end of   a line counts as a space.

One   or more   blank lines denote the  end
of  a paragraph.

Since any number of consecutive spaces are treated
like a single one, the formatting of the input
file makes no difference to
      \LaTeX,                % The \LaTeX command generates the LaTeX logo.
but it makes a difference to you.  When you use
\LaTeX, making your input file as easy to read
as possible will be a great help as you write
your document and when you change it.  This sample
file shows how you can add comments to your own input
file.

Because printing is different from typewriting,
there are a number of things that you have to do
differently when preparing an input file than if
you were just typing the document directly.
Quotation marks like
       ``this''
have to be handled specially, as do quotes within
quotes:
       ``\,`this'            % \, separates the double and single quote.
        is what I just
        wrote, not  `that'\,''.

Dashes come in three sizes: an
       intra-word
dash, a medium dash for number ranges like
       1--2,
and a punctuation
       dash---like
this.

A sentence-ending space should be larger than the
space between words within a sentence.  You
sometimes have to type special commands in
conjunction with punctuation characters to get
this right, as in the following sentence.
       Gnats, gnus, etc.\ all  % `\ ' makes an inter-word space.
       begin with G\@.         % \@ marks end-of-sentence punctuation.
You should check the spaces after periods when
reading your output to make sure you haven't
forgotten any special cases.  Generating an
ellipsis
       \ldots\               % `\ ' is needed after `\ldots' because TeX
                             % ignores spaces after command names like \ldots
                             % made from \ + letters.
                             %
                             % Note how a `%' character causes TeX to ignore
                             % the end of the input line, so these blank lines
                             % do not start a new paragraph.
                             %
with the right spacing around the periods requires
a special command.

\LaTeX\ interprets some common characters as
commands, so you must type special commands to
generate them.  These characters include the
following:
       \$ \& \% \# \{ and \}.

In printing, text is usually emphasized with an
       \emph{italic}
type style.

\begin{em}
   A long segment of text can also be emphasized
   in this way.  Text within such a segment can be
   given \emph{additional} emphasis.
\end{em}

It is sometimes necessary to prevent \LaTeX\ from
breaking a line where it might otherwise do so.
This may be at a space, as between the ``Mr.''\ and
``Jones'' in
       ``Mr.~Jones'',        % ~ produces an unbreakable interword space.
or within a word---especially when the word is a
symbol like
       \mbox{\emph{itemnum}}
that makes little sense when hyphenated across
lines.

Footnotes\footnote{This is an example of a footnote.}
pose no problem.

\LaTeX\ is good at typesetting mathematical formulas
like
       \( x-3y + z = 7 \)
or
       \( a_{1} > x^{2n} + y^{2n} > x' \)
or
       \( \ip{A}{B} = \sum_{i} a_{i} b_{i} \).
The spaces you type in a formula are
ignored.  Remember that a letter like
       $x$                   % $ ... $  and  \( ... \)  are equivalent
is a formula when it denotes a mathematical
symbol, and it should be typed as one.

\section{Displayed Text}

Text is displayed by indenting it from the left
margin.  Quotations are commonly displayed.  There
are short quotations
\begin{quote}
   This is a short quotation.  It consists of a
   single paragraph of text.  See how it is formatted.
\end{quote}
and longer ones.
\begin{quotation}
   This is a longer quotation.  It consists of two
   paragraphs of text, neither of which are
   particularly interesting.

   This is the second paragraph of the quotation.  It
   is just as dull as the first paragraph.
\end{quotation}
Another frequently-displayed structure is a list.
The following is an example of an \emph{itemized}
list.
\begin{itemize}
   \item This is the first item of an itemized list.
         Each item in the list is marked with a ``tick''.
         You don't have to worry about what kind of tick
         mark is used.

   \item This is the second item of the list.  It
         contains another list nested inside it.  The inner
         list is an \emph{enumerated} list.
         \begin{enumerate}
            \item This is the first item of an enumerated
                  list that is nested within the itemized list.

            \item This is the second item of the inner list.
                  \LaTeX\ allows you to nest lists deeper than
                  you really should.
         \end{enumerate}
         This is the rest of the second item of the outer
         list.  It is no more interesting than any other
         part of the item.
   \item This is the third item of the list.
\end{itemize}
You can even display poetry.
\begin{verse}
   There is an environment
    for verse \\             % The \\ command separates lines
   Whose features some poets % within a stanza.
   will curse.

                             % One or more blank lines separate stanzas.

   For instead of making\\
   Them do \emph{all} line breaking, \\
   It allows them to put too many words on a line when they'd rather be
   forced to be terse.
\end{verse}

Mathematical formulas may also be displayed.  A
displayed formula
is
one-line long; multiline
formulas require special formatting instructions.
   \[  \ip{\Gamma}{\psi'} = x'' + y^{2} + z_{i}^{n}\]
Don't start a paragraph with a displayed equation,
nor make one a paragraph by itself.

\end{document}               % End of document.

\section{Bench Section 32}
% synthetic bench section 32
% This is a sample LaTeX input file.  (Version of 12 August 2004.)
%
% A '%' character causes TeX to ignore all remaining text on the line,
% and is used for comments like this one.

\documentclass{article}      % Specifies the document class

                             % The preamble begins here.
\title{An Example Document}  % Declares the document's title.
\author{Leslie Lamport}      % Declares the author's name.
\date{January 21, 1994}      % Deleting this command produces today's date.

\newcommand{\ip}[2]{(#1, #2)}
                             % Defines \ip{arg1}{arg2} to mean
                             % (arg1, arg2).

%\newcommand{\ip}[2]{\langle #1 | #2\rangle}
                             % This is an alternative definition of
                             % \ip that is commented out.

\begin{document}             % End of preamble and beginning of text.

\maketitle                   % Produces the title.

This is an example input file.  Comparing it with
the output it generates can show you how to
produce a simple document of your own.

\section{Ordinary Text}      % Produces section heading.  Lower-level
                             % sections are begun with similar
                             % \subsection and \subsubsection commands.

The ends  of words and sentences are marked
  by   spaces. It  doesn't matter how many
spaces    you type; one is as good as 100.  The
end of   a line counts as a space.

One   or more   blank lines denote the  end
of  a paragraph.

Since any number of consecutive spaces are treated
like a single one, the formatting of the input
file makes no difference to
      \LaTeX,                % The \LaTeX command generates the LaTeX logo.
but it makes a difference to you.  When you use
\LaTeX, making your input file as easy to read
as possible will be a great help as you write
your document and when you change it.  This sample
file shows how you can add comments to your own input
file.

Because printing is different from typewriting,
there are a number of things that you have to do
differently when preparing an input file than if
you were just typing the document directly.
Quotation marks like
       ``this''
have to be handled specially, as do quotes within
quotes:
       ``\,`this'            % \, separates the double and single quote.
        is what I just
        wrote, not  `that'\,''.

Dashes come in three sizes: an
       intra-word
dash, a medium dash for number ranges like
       1--2,
and a punctuation
       dash---like
this.

A sentence-ending space should be larger than the
space between words within a sentence.  You
sometimes have to type special commands in
conjunction with punctuation characters to get
this right, as in the following sentence.
       Gnats, gnus, etc.\ all  % `\ ' makes an inter-word space.
       begin with G\@.         % \@ marks end-of-sentence punctuation.
You should check the spaces after periods when
reading your output to make sure you haven't
forgotten any special cases.  Generating an
ellipsis
       \ldots\               % `\ ' is needed after `\ldots' because TeX
                             % ignores spaces after command names like \ldots
                             % made from \ + letters.
                             %
                             % Note how a `%' character causes TeX to ignore
                             % the end of the input line, so these blank lines
                             % do not start a new paragraph.
                             %
with the right spacing around the periods requires
a special command.

\LaTeX\ interprets some common characters as
commands, so you must type special commands to
generate them.  These characters include the
following:
       \$ \& \% \# \{ and \}.

In printing, text is usually emphasized with an
       \emph{italic}
type style.

\begin{em}
   A long segment of text can also be emphasized
   in this way.  Text within such a segment can be
   given \emph{additional} emphasis.
\end{em}

It is sometimes necessary to prevent \LaTeX\ from
breaking a line where it might otherwise do so.
This may be at a space, as between the ``Mr.''\ and
``Jones'' in
       ``Mr.~Jones'',        % ~ produces an unbreakable interword space.
or within a word---especially when the word is a
symbol like
       \mbox{\emph{itemnum}}
that makes little sense when hyphenated across
lines.

Footnotes\footnote{This is an example of a footnote.}
pose no problem.

\LaTeX\ is good at typesetting mathematical formulas
like
       \( x-3y + z = 7 \)
or
       \( a_{1} > x^{2n} + y^{2n} > x' \)
or
       \( \ip{A}{B} = \sum_{i} a_{i} b_{i} \).
The spaces you type in a formula are
ignored.  Remember that a letter like
       $x$                   % $ ... $  and  \( ... \)  are equivalent
is a formula when it denotes a mathematical
symbol, and it should be typed as one.

\section{Displayed Text}

Text is displayed by indenting it from the left
margin.  Quotations are commonly displayed.  There
are short quotations
\begin{quote}
   This is a short quotation.  It consists of a
   single paragraph of text.  See how it is formatted.
\end{quote}
and longer ones.
\begin{quotation}
   This is a longer quotation.  It consists of two
   paragraphs of text, neither of which are
   particularly interesting.

   This is the second paragraph of the quotation.  It
   is just as dull as the first paragraph.
\end{quotation}
Another frequently-displayed structure is a list.
The following is an example of an \emph{itemized}
list.
\begin{itemize}
   \item This is the first item of an itemized list.
         Each item in the list is marked with a ``tick''.
         You don't have to worry about what kind of tick
         mark is used.

   \item This is the second item of the list.  It
         contains another list nested inside it.  The inner
         list is an \emph{enumerated} list.
         \begin{enumerate}
            \item This is the first item of an enumerated
                  list that is nested within the itemized list.

            \item This is the second item of the inner list.
                  \LaTeX\ allows you to nest lists deeper than
                  you really should.
         \end{enumerate}
         This is the rest of the second item of the outer
         list.  It is no more interesting than any other
         part of the item.
   \item This is the third item of the list.
\end{itemize}
You can even display poetry.
\begin{verse}
   There is an environment
    for verse \\             % The \\ command separates lines
   Whose features some poets % within a stanza.
   will curse.

                             % One or more blank lines separate stanzas.

   For instead of making\\
   Them do \emph{all} line breaking, \\
   It allows them to put too many words on a line when they'd rather be
   forced to be terse.
\end{verse}

Mathematical formulas may also be displayed.  A
displayed formula
is
one-line long; multiline
formulas require special formatting instructions.
   \[  \ip{\Gamma}{\psi'} = x'' + y^{2} + z_{i}^{n}\]
Don't start a paragraph with a displayed equation,
nor make one a paragraph by itself.

\end{document}               % End of document.

\section{Bench Section 33}
% synthetic bench section 33
% This is a sample LaTeX input file.  (Version of 12 August 2004.)
%
% A '%' character causes TeX to ignore all remaining text on the line,
% and is used for comments like this one.

\documentclass{article}      % Specifies the document class

                             % The preamble begins here.
\title{An Example Document}  % Declares the document's title.
\author{Leslie Lamport}      % Declares the author's name.
\date{January 21, 1994}      % Deleting this command produces today's date.

\newcommand{\ip}[2]{(#1, #2)}
                             % Defines \ip{arg1}{arg2} to mean
                             % (arg1, arg2).

%\newcommand{\ip}[2]{\langle #1 | #2\rangle}
                             % This is an alternative definition of
                             % \ip that is commented out.

\begin{document}             % End of preamble and beginning of text.

\maketitle                   % Produces the title.

This is an example input file.  Comparing it with
the output it generates can show you how to
produce a simple document of your own.

\section{Ordinary Text}      % Produces section heading.  Lower-level
                             % sections are begun with similar
                             % \subsection and \subsubsection commands.

The ends  of words and sentences are marked
  by   spaces. It  doesn't matter how many
spaces    you type; one is as good as 100.  The
end of   a line counts as a space.

One   or more   blank lines denote the  end
of  a paragraph.

Since any number of consecutive spaces are treated
like a single one, the formatting of the input
file makes no difference to
      \LaTeX,                % The \LaTeX command generates the LaTeX logo.
but it makes a difference to you.  When you use
\LaTeX, making your input file as easy to read
as possible will be a great help as you write
your document and when you change it.  This sample
file shows how you can add comments to your own input
file.

Because printing is different from typewriting,
there are a number of things that you have to do
differently when preparing an input file than if
you were just typing the document directly.
Quotation marks like
       ``this''
have to be handled specially, as do quotes within
quotes:
       ``\,`this'            % \, separates the double and single quote.
        is what I just
        wrote, not  `that'\,''.

Dashes come in three sizes: an
       intra-word
dash, a medium dash for number ranges like
       1--2,
and a punctuation
       dash---like
this.

A sentence-ending space should be larger than the
space between words within a sentence.  You
sometimes have to type special commands in
conjunction with punctuation characters to get
this right, as in the following sentence.
       Gnats, gnus, etc.\ all  % `\ ' makes an inter-word space.
       begin with G\@.         % \@ marks end-of-sentence punctuation.
You should check the spaces after periods when
reading your output to make sure you haven't
forgotten any special cases.  Generating an
ellipsis
       \ldots\               % `\ ' is needed after `\ldots' because TeX
                             % ignores spaces after command names like \ldots
                             % made from \ + letters.
                             %
                             % Note how a `%' character causes TeX to ignore
                             % the end of the input line, so these blank lines
                             % do not start a new paragraph.
                             %
with the right spacing around the periods requires
a special command.

\LaTeX\ interprets some common characters as
commands, so you must type special commands to
generate them.  These characters include the
following:
       \$ \& \% \# \{ and \}.

In printing, text is usually emphasized with an
       \emph{italic}
type style.

\begin{em}
   A long segment of text can also be emphasized
   in this way.  Text within such a segment can be
   given \emph{additional} emphasis.
\end{em}

It is sometimes necessary to prevent \LaTeX\ from
breaking a line where it might otherwise do so.
This may be at a space, as between the ``Mr.''\ and
``Jones'' in
       ``Mr.~Jones'',        % ~ produces an unbreakable interword space.
or within a word---especially when the word is a
symbol like
       \mbox{\emph{itemnum}}
that makes little sense when hyphenated across
lines.

Footnotes\footnote{This is an example of a footnote.}
pose no problem.

\LaTeX\ is good at typesetting mathematical formulas
like
       \( x-3y + z = 7 \)
or
       \( a_{1} > x^{2n} + y^{2n} > x' \)
or
       \( \ip{A}{B} = \sum_{i} a_{i} b_{i} \).
The spaces you type in a formula are
ignored.  Remember that a letter like
       $x$                   % $ ... $  and  \( ... \)  are equivalent
is a formula when it denotes a mathematical
symbol, and it should be typed as one.

\section{Displayed Text}

Text is displayed by indenting it from the left
margin.  Quotations are commonly displayed.  There
are short quotations
\begin{quote}
   This is a short quotation.  It consists of a
   single paragraph of text.  See how it is formatted.
\end{quote}
and longer ones.
\begin{quotation}
   This is a longer quotation.  It consists of two
   paragraphs of text, neither of which are
   particularly interesting.

   This is the second paragraph of the quotation.  It
   is just as dull as the first paragraph.
\end{quotation}
Another frequently-displayed structure is a list.
The following is an example of an \emph{itemized}
list.
\begin{itemize}
   \item This is the first item of an itemized list.
         Each item in the list is marked with a ``tick''.
         You don't have to worry about what kind of tick
         mark is used.

   \item This is the second item of the list.  It
         contains another list nested inside it.  The inner
         list is an \emph{enumerated} list.
         \begin{enumerate}
            \item This is the first item of an enumerated
                  list that is nested within the itemized list.

            \item This is the second item of the inner list.
                  \LaTeX\ allows you to nest lists deeper than
                  you really should.
         \end{enumerate}
         This is the rest of the second item of the outer
         list.  It is no more interesting than any other
         part of the item.
   \item This is the third item of the list.
\end{itemize}
You can even display poetry.
\begin{verse}
   There is an environment
    for verse \\             % The \\ command separates lines
   Whose features some poets % within a stanza.
   will curse.

                             % One or more blank lines separate stanzas.

   For instead of making\\
   Them do \emph{all} line breaking, \\
   It allows them to put too many words on a line when they'd rather be
   forced to be terse.
\end{verse}

Mathematical formulas may also be displayed.  A
displayed formula
is
one-line long; multiline
formulas require special formatting instructions.
   \[  \ip{\Gamma}{\psi'} = x'' + y^{2} + z_{i}^{n}\]
Don't start a paragraph with a displayed equation,
nor make one a paragraph by itself.

\end{document}               % End of document.

\section{Bench Section 34}
% synthetic bench section 34
% This is a sample LaTeX input file.  (Version of 12 August 2004.)
%
% A '%' character causes TeX to ignore all remaining text on the line,
% and is used for comments like this one.

\documentclass{article}      % Specifies the document class

                             % The preamble begins here.
\title{An Example Document}  % Declares the document's title.
\author{Leslie Lamport}      % Declares the author's name.
\date{January 21, 1994}      % Deleting this command produces today's date.

\newcommand{\ip}[2]{(#1, #2)}
                             % Defines \ip{arg1}{arg2} to mean
                             % (arg1, arg2).

%\newcommand{\ip}[2]{\langle #1 | #2\rangle}
                             % This is an alternative definition of
                             % \ip that is commented out.

\begin{document}             % End of preamble and beginning of text.

\maketitle                   % Produces the title.

This is an example input file.  Comparing it with
the output it generates can show you how to
produce a simple document of your own.

\section{Ordinary Text}      % Produces section heading.  Lower-level
                             % sections are begun with similar
                             % \subsection and \subsubsection commands.

The ends  of words and sentences are marked
  by   spaces. It  doesn't matter how many
spaces    you type; one is as good as 100.  The
end of   a line counts as a space.

One   or more   blank lines denote the  end
of  a paragraph.

Since any number of consecutive spaces are treated
like a single one, the formatting of the input
file makes no difference to
      \LaTeX,                % The \LaTeX command generates the LaTeX logo.
but it makes a difference to you.  When you use
\LaTeX, making your input file as easy to read
as possible will be a great help as you write
your document and when you change it.  This sample
file shows how you can add comments to your own input
file.

Because printing is different from typewriting,
there are a number of things that you have to do
differently when preparing an input file than if
you were just typing the document directly.
Quotation marks like
       ``this''
have to be handled specially, as do quotes within
quotes:
       ``\,`this'            % \, separates the double and single quote.
        is what I just
        wrote, not  `that'\,''.

Dashes come in three sizes: an
       intra-word
dash, a medium dash for number ranges like
       1--2,
and a punctuation
       dash---like
this.

A sentence-ending space should be larger than the
space between words within a sentence.  You
sometimes have to type special commands in
conjunction with punctuation characters to get
this right, as in the following sentence.
       Gnats, gnus, etc.\ all  % `\ ' makes an inter-word space.
       begin with G\@.         % \@ marks end-of-sentence punctuation.
You should check the spaces after periods when
reading your output to make sure you haven't
forgotten any special cases.  Generating an
ellipsis
       \ldots\               % `\ ' is needed after `\ldots' because TeX
                             % ignores spaces after command names like \ldots
                             % made from \ + letters.
                             %
                             % Note how a `%' character causes TeX to ignore
                             % the end of the input line, so these blank lines
                             % do not start a new paragraph.
                             %
with the right spacing around the periods requires
a special command.

\LaTeX\ interprets some common characters as
commands, so you must type special commands to
generate them.  These characters include the
following:
       \$ \& \% \# \{ and \}.

In printing, text is usually emphasized with an
       \emph{italic}
type style.

\begin{em}
   A long segment of text can also be emphasized
   in this way.  Text within such a segment can be
   given \emph{additional} emphasis.
\end{em}

It is sometimes necessary to prevent \LaTeX\ from
breaking a line where it might otherwise do so.
This may be at a space, as between the ``Mr.''\ and
``Jones'' in
       ``Mr.~Jones'',        % ~ produces an unbreakable interword space.
or within a word---especially when the word is a
symbol like
       \mbox{\emph{itemnum}}
that makes little sense when hyphenated across
lines.

Footnotes\footnote{This is an example of a footnote.}
pose no problem.

\LaTeX\ is good at typesetting mathematical formulas
like
       \( x-3y + z = 7 \)
or
       \( a_{1} > x^{2n} + y^{2n} > x' \)
or
       \( \ip{A}{B} = \sum_{i} a_{i} b_{i} \).
The spaces you type in a formula are
ignored.  Remember that a letter like
       $x$                   % $ ... $  and  \( ... \)  are equivalent
is a formula when it denotes a mathematical
symbol, and it should be typed as one.

\section{Displayed Text}

Text is displayed by indenting it from the left
margin.  Quotations are commonly displayed.  There
are short quotations
\begin{quote}
   This is a short quotation.  It consists of a
   single paragraph of text.  See how it is formatted.
\end{quote}
and longer ones.
\begin{quotation}
   This is a longer quotation.  It consists of two
   paragraphs of text, neither of which are
   particularly interesting.

   This is the second paragraph of the quotation.  It
   is just as dull as the first paragraph.
\end{quotation}
Another frequently-displayed structure is a list.
The following is an example of an \emph{itemized}
list.
\begin{itemize}
   \item This is the first item of an itemized list.
         Each item in the list is marked with a ``tick''.
         You don't have to worry about what kind of tick
         mark is used.

   \item This is the second item of the list.  It
         contains another list nested inside it.  The inner
         list is an \emph{enumerated} list.
         \begin{enumerate}
            \item This is the first item of an enumerated
                  list that is nested within the itemized list.

            \item This is the second item of the inner list.
                  \LaTeX\ allows you to nest lists deeper than
                  you really should.
         \end{enumerate}
         This is the rest of the second item of the outer
         list.  It is no more interesting than any other
         part of the item.
   \item This is the third item of the list.
\end{itemize}
You can even display poetry.
\begin{verse}
   There is an environment
    for verse \\             % The \\ command separates lines
   Whose features some poets % within a stanza.
   will curse.

                             % One or more blank lines separate stanzas.

   For instead of making\\
   Them do \emph{all} line breaking, \\
   It allows them to put too many words on a line when they'd rather be
   forced to be terse.
\end{verse}

Mathematical formulas may also be displayed.  A
displayed formula
is
one-line long; multiline
formulas require special formatting instructions.
   \[  \ip{\Gamma}{\psi'} = x'' + y^{2} + z_{i}^{n}\]
Don't start a paragraph with a displayed equation,
nor make one a paragraph by itself.

\end{document}               % End of document.

\section{Bench Section 35}
% synthetic bench section 35
% This is a sample LaTeX input file.  (Version of 12 August 2004.)
%
% A '%' character causes TeX to ignore all remaining text on the line,
% and is used for comments like this one.

\documentclass{article}      % Specifies the document class

                             % The preamble begins here.
\title{An Example Document}  % Declares the document's title.
\author{Leslie Lamport}      % Declares the author's name.
\date{January 21, 1994}      % Deleting this command produces today's date.

\newcommand{\ip}[2]{(#1, #2)}
                             % Defines \ip{arg1}{arg2} to mean
                             % (arg1, arg2).

%\newcommand{\ip}[2]{\langle #1 | #2\rangle}
                             % This is an alternative definition of
                             % \ip that is commented out.

\begin{document}             % End of preamble and beginning of text.

\maketitle                   % Produces the title.

This is an example input file.  Comparing it with
the output it generates can show you how to
produce a simple document of your own.

\section{Ordinary Text}      % Produces section heading.  Lower-level
                             % sections are begun with similar
                             % \subsection and \subsubsection commands.

The ends  of words and sentences are marked
  by   spaces. It  doesn't matter how many
spaces    you type; one is as good as 100.  The
end of   a line counts as a space.

One   or more   blank lines denote the  end
of  a paragraph.

Since any number of consecutive spaces are treated
like a single one, the formatting of the input
file makes no difference to
      \LaTeX,                % The \LaTeX command generates the LaTeX logo.
but it makes a difference to you.  When you use
\LaTeX, making your input file as easy to read
as possible will be a great help as you write
your document and when you change it.  This sample
file shows how you can add comments to your own input
file.

Because printing is different from typewriting,
there are a number of things that you have to do
differently when preparing an input file than if
you were just typing the document directly.
Quotation marks like
       ``this''
have to be handled specially, as do quotes within
quotes:
       ``\,`this'            % \, separates the double and single quote.
        is what I just
        wrote, not  `that'\,''.

Dashes come in three sizes: an
       intra-word
dash, a medium dash for number ranges like
       1--2,
and a punctuation
       dash---like
this.

A sentence-ending space should be larger than the
space between words within a sentence.  You
sometimes have to type special commands in
conjunction with punctuation characters to get
this right, as in the following sentence.
       Gnats, gnus, etc.\ all  % `\ ' makes an inter-word space.
       begin with G\@.         % \@ marks end-of-sentence punctuation.
You should check the spaces after periods when
reading your output to make sure you haven't
forgotten any special cases.  Generating an
ellipsis
       \ldots\               % `\ ' is needed after `\ldots' because TeX
                             % ignores spaces after command names like \ldots
                             % made from \ + letters.
                             %
                             % Note how a `%' character causes TeX to ignore
                             % the end of the input line, so these blank lines
                             % do not start a new paragraph.
                             %
with the right spacing around the periods requires
a special command.

\LaTeX\ interprets some common characters as
commands, so you must type special commands to
generate them.  These characters include the
following:
       \$ \& \% \# \{ and \}.

In printing, text is usually emphasized with an
       \emph{italic}
type style.

\begin{em}
   A long segment of text can also be emphasized
   in this way.  Text within such a segment can be
   given \emph{additional} emphasis.
\end{em}

It is sometimes necessary to prevent \LaTeX\ from
breaking a line where it might otherwise do so.
This may be at a space, as between the ``Mr.''\ and
``Jones'' in
       ``Mr.~Jones'',        % ~ produces an unbreakable interword space.
or within a word---especially when the word is a
symbol like
       \mbox{\emph{itemnum}}
that makes little sense when hyphenated across
lines.

Footnotes\footnote{This is an example of a footnote.}
pose no problem.

\LaTeX\ is good at typesetting mathematical formulas
like
       \( x-3y + z = 7 \)
or
       \( a_{1} > x^{2n} + y^{2n} > x' \)
or
       \( \ip{A}{B} = \sum_{i} a_{i} b_{i} \).
The spaces you type in a formula are
ignored.  Remember that a letter like
       $x$                   % $ ... $  and  \( ... \)  are equivalent
is a formula when it denotes a mathematical
symbol, and it should be typed as one.

\section{Displayed Text}

Text is displayed by indenting it from the left
margin.  Quotations are commonly displayed.  There
are short quotations
\begin{quote}
   This is a short quotation.  It consists of a
   single paragraph of text.  See how it is formatted.
\end{quote}
and longer ones.
\begin{quotation}
   This is a longer quotation.  It consists of two
   paragraphs of text, neither of which are
   particularly interesting.

   This is the second paragraph of the quotation.  It
   is just as dull as the first paragraph.
\end{quotation}
Another frequently-displayed structure is a list.
The following is an example of an \emph{itemized}
list.
\begin{itemize}
   \item This is the first item of an itemized list.
         Each item in the list is marked with a ``tick''.
         You don't have to worry about what kind of tick
         mark is used.

   \item This is the second item of the list.  It
         contains another list nested inside it.  The inner
         list is an \emph{enumerated} list.
         \begin{enumerate}
            \item This is the first item of an enumerated
                  list that is nested within the itemized list.

            \item This is the second item of the inner list.
                  \LaTeX\ allows you to nest lists deeper than
                  you really should.
         \end{enumerate}
         This is the rest of the second item of the outer
         list.  It is no more interesting than any other
         part of the item.
   \item This is the third item of the list.
\end{itemize}
You can even display poetry.
\begin{verse}
   There is an environment
    for verse \\             % The \\ command separates lines
   Whose features some poets % within a stanza.
   will curse.

                             % One or more blank lines separate stanzas.

   For instead of making\\
   Them do \emph{all} line breaking, \\
   It allows them to put too many words on a line when they'd rather be
   forced to be terse.
\end{verse}

Mathematical formulas may also be displayed.  A
displayed formula
is
one-line long; multiline
formulas require special formatting instructions.
   \[  \ip{\Gamma}{\psi'} = x'' + y^{2} + z_{i}^{n}\]
Don't start a paragraph with a displayed equation,
nor make one a paragraph by itself.

\end{document}               % End of document.

\section{Bench Section 36}
% synthetic bench section 36
% This is a sample LaTeX input file.  (Version of 12 August 2004.)
%
% A '%' character causes TeX to ignore all remaining text on the line,
% and is used for comments like this one.

\documentclass{article}      % Specifies the document class

                             % The preamble begins here.
\title{An Example Document}  % Declares the document's title.
\author{Leslie Lamport}      % Declares the author's name.
\date{January 21, 1994}      % Deleting this command produces today's date.

\newcommand{\ip}[2]{(#1, #2)}
                             % Defines \ip{arg1}{arg2} to mean
                             % (arg1, arg2).

%\newcommand{\ip}[2]{\langle #1 | #2\rangle}
                             % This is an alternative definition of
                             % \ip that is commented out.

\begin{document}             % End of preamble and beginning of text.

\maketitle                   % Produces the title.

This is an example input file.  Comparing it with
the output it generates can show you how to
produce a simple document of your own.

\section{Ordinary Text}      % Produces section heading.  Lower-level
                             % sections are begun with similar
                             % \subsection and \subsubsection commands.

The ends  of words and sentences are marked
  by   spaces. It  doesn't matter how many
spaces    you type; one is as good as 100.  The
end of   a line counts as a space.

One   or more   blank lines denote the  end
of  a paragraph.

Since any number of consecutive spaces are treated
like a single one, the formatting of the input
file makes no difference to
      \LaTeX,                % The \LaTeX command generates the LaTeX logo.
but it makes a difference to you.  When you use
\LaTeX, making your input file as easy to read
as possible will be a great help as you write
your document and when you change it.  This sample
file shows how you can add comments to your own input
file.

Because printing is different from typewriting,
there are a number of things that you have to do
differently when preparing an input file than if
you were just typing the document directly.
Quotation marks like
       ``this''
have to be handled specially, as do quotes within
quotes:
       ``\,`this'            % \, separates the double and single quote.
        is what I just
        wrote, not  `that'\,''.

Dashes come in three sizes: an
       intra-word
dash, a medium dash for number ranges like
       1--2,
and a punctuation
       dash---like
this.

A sentence-ending space should be larger than the
space between words within a sentence.  You
sometimes have to type special commands in
conjunction with punctuation characters to get
this right, as in the following sentence.
       Gnats, gnus, etc.\ all  % `\ ' makes an inter-word space.
       begin with G\@.         % \@ marks end-of-sentence punctuation.
You should check the spaces after periods when
reading your output to make sure you haven't
forgotten any special cases.  Generating an
ellipsis
       \ldots\               % `\ ' is needed after `\ldots' because TeX
                             % ignores spaces after command names like \ldots
                             % made from \ + letters.
                             %
                             % Note how a `%' character causes TeX to ignore
                             % the end of the input line, so these blank lines
                             % do not start a new paragraph.
                             %
with the right spacing around the periods requires
a special command.

\LaTeX\ interprets some common characters as
commands, so you must type special commands to
generate them.  These characters include the
following:
       \$ \& \% \# \{ and \}.

In printing, text is usually emphasized with an
       \emph{italic}
type style.

\begin{em}
   A long segment of text can also be emphasized
   in this way.  Text within such a segment can be
   given \emph{additional} emphasis.
\end{em}

It is sometimes necessary to prevent \LaTeX\ from
breaking a line where it might otherwise do so.
This may be at a space, as between the ``Mr.''\ and
``Jones'' in
       ``Mr.~Jones'',        % ~ produces an unbreakable interword space.
or within a word---especially when the word is a
symbol like
       \mbox{\emph{itemnum}}
that makes little sense when hyphenated across
lines.

Footnotes\footnote{This is an example of a footnote.}
pose no problem.

\LaTeX\ is good at typesetting mathematical formulas
like
       \( x-3y + z = 7 \)
or
       \( a_{1} > x^{2n} + y^{2n} > x' \)
or
       \( \ip{A}{B} = \sum_{i} a_{i} b_{i} \).
The spaces you type in a formula are
ignored.  Remember that a letter like
       $x$                   % $ ... $  and  \( ... \)  are equivalent
is a formula when it denotes a mathematical
symbol, and it should be typed as one.

\section{Displayed Text}

Text is displayed by indenting it from the left
margin.  Quotations are commonly displayed.  There
are short quotations
\begin{quote}
   This is a short quotation.  It consists of a
   single paragraph of text.  See how it is formatted.
\end{quote}
and longer ones.
\begin{quotation}
   This is a longer quotation.  It consists of two
   paragraphs of text, neither of which are
   particularly interesting.

   This is the second paragraph of the quotation.  It
   is just as dull as the first paragraph.
\end{quotation}
Another frequently-displayed structure is a list.
The following is an example of an \emph{itemized}
list.
\begin{itemize}
   \item This is the first item of an itemized list.
         Each item in the list is marked with a ``tick''.
         You don't have to worry about what kind of tick
         mark is used.

   \item This is the second item of the list.  It
         contains another list nested inside it.  The inner
         list is an \emph{enumerated} list.
         \begin{enumerate}
            \item This is the first item of an enumerated
                  list that is nested within the itemized list.

            \item This is the second item of the inner list.
                  \LaTeX\ allows you to nest lists deeper than
                  you really should.
         \end{enumerate}
         This is the rest of the second item of the outer
         list.  It is no more interesting than any other
         part of the item.
   \item This is the third item of the list.
\end{itemize}
You can even display poetry.
\begin{verse}
   There is an environment
    for verse \\             % The \\ command separates lines
   Whose features some poets % within a stanza.
   will curse.

                             % One or more blank lines separate stanzas.

   For instead of making\\
   Them do \emph{all} line breaking, \\
   It allows them to put too many words on a line when they'd rather be
   forced to be terse.
\end{verse}

Mathematical formulas may also be displayed.  A
displayed formula
is
one-line long; multiline
formulas require special formatting instructions.
   \[  \ip{\Gamma}{\psi'} = x'' + y^{2} + z_{i}^{n}\]
Don't start a paragraph with a displayed equation,
nor make one a paragraph by itself.

\end{document}               % End of document.

\section{Bench Section 37}
% synthetic bench section 37
% This is a sample LaTeX input file.  (Version of 12 August 2004.)
%
% A '%' character causes TeX to ignore all remaining text on the line,
% and is used for comments like this one.

\documentclass{article}      % Specifies the document class

                             % The preamble begins here.
\title{An Example Document}  % Declares the document's title.
\author{Leslie Lamport}      % Declares the author's name.
\date{January 21, 1994}      % Deleting this command produces today's date.

\newcommand{\ip}[2]{(#1, #2)}
                             % Defines \ip{arg1}{arg2} to mean
                             % (arg1, arg2).

%\newcommand{\ip}[2]{\langle #1 | #2\rangle}
                             % This is an alternative definition of
                             % \ip that is commented out.

\begin{document}             % End of preamble and beginning of text.

\maketitle                   % Produces the title.

This is an example input file.  Comparing it with
the output it generates can show you how to
produce a simple document of your own.

\section{Ordinary Text}      % Produces section heading.  Lower-level
                             % sections are begun with similar
                             % \subsection and \subsubsection commands.

The ends  of words and sentences are marked
  by   spaces. It  doesn't matter how many
spaces    you type; one is as good as 100.  The
end of   a line counts as a space.

One   or more   blank lines denote the  end
of  a paragraph.

Since any number of consecutive spaces are treated
like a single one, the formatting of the input
file makes no difference to
      \LaTeX,                % The \LaTeX command generates the LaTeX logo.
but it makes a difference to you.  When you use
\LaTeX, making your input file as easy to read
as possible will be a great help as you write
your document and when you change it.  This sample
file shows how you can add comments to your own input
file.

Because printing is different from typewriting,
there are a number of things that you have to do
differently when preparing an input file than if
you were just typing the document directly.
Quotation marks like
       ``this''
have to be handled specially, as do quotes within
quotes:
       ``\,`this'            % \, separates the double and single quote.
        is what I just
        wrote, not  `that'\,''.

Dashes come in three sizes: an
       intra-word
dash, a medium dash for number ranges like
       1--2,
and a punctuation
       dash---like
this.

A sentence-ending space should be larger than the
space between words within a sentence.  You
sometimes have to type special commands in
conjunction with punctuation characters to get
this right, as in the following sentence.
       Gnats, gnus, etc.\ all  % `\ ' makes an inter-word space.
       begin with G\@.         % \@ marks end-of-sentence punctuation.
You should check the spaces after periods when
reading your output to make sure you haven't
forgotten any special cases.  Generating an
ellipsis
       \ldots\               % `\ ' is needed after `\ldots' because TeX
                             % ignores spaces after command names like \ldots
                             % made from \ + letters.
                             %
                             % Note how a `%' character causes TeX to ignore
                             % the end of the input line, so these blank lines
                             % do not start a new paragraph.
                             %
with the right spacing around the periods requires
a special command.

\LaTeX\ interprets some common characters as
commands, so you must type special commands to
generate them.  These characters include the
following:
       \$ \& \% \# \{ and \}.

In printing, text is usually emphasized with an
       \emph{italic}
type style.

\begin{em}
   A long segment of text can also be emphasized
   in this way.  Text within such a segment can be
   given \emph{additional} emphasis.
\end{em}

It is sometimes necessary to prevent \LaTeX\ from
breaking a line where it might otherwise do so.
This may be at a space, as between the ``Mr.''\ and
``Jones'' in
       ``Mr.~Jones'',        % ~ produces an unbreakable interword space.
or within a word---especially when the word is a
symbol like
       \mbox{\emph{itemnum}}
that makes little sense when hyphenated across
lines.

Footnotes\footnote{This is an example of a footnote.}
pose no problem.

\LaTeX\ is good at typesetting mathematical formulas
like
       \( x-3y + z = 7 \)
or
       \( a_{1} > x^{2n} + y^{2n} > x' \)
or
       \( \ip{A}{B} = \sum_{i} a_{i} b_{i} \).
The spaces you type in a formula are
ignored.  Remember that a letter like
       $x$                   % $ ... $  and  \( ... \)  are equivalent
is a formula when it denotes a mathematical
symbol, and it should be typed as one.

\section{Displayed Text}

Text is displayed by indenting it from the left
margin.  Quotations are commonly displayed.  There
are short quotations
\begin{quote}
   This is a short quotation.  It consists of a
   single paragraph of text.  See how it is formatted.
\end{quote}
and longer ones.
\begin{quotation}
   This is a longer quotation.  It consists of two
   paragraphs of text, neither of which are
   particularly interesting.

   This is the second paragraph of the quotation.  It
   is just as dull as the first paragraph.
\end{quotation}
Another frequently-displayed structure is a list.
The following is an example of an \emph{itemized}
list.
\begin{itemize}
   \item This is the first item of an itemized list.
         Each item in the list is marked with a ``tick''.
         You don't have to worry about what kind of tick
         mark is used.

   \item This is the second item of the list.  It
         contains another list nested inside it.  The inner
         list is an \emph{enumerated} list.
         \begin{enumerate}
            \item This is the first item of an enumerated
                  list that is nested within the itemized list.

            \item This is the second item of the inner list.
                  \LaTeX\ allows you to nest lists deeper than
                  you really should.
         \end{enumerate}
         This is the rest of the second item of the outer
         list.  It is no more interesting than any other
         part of the item.
   \item This is the third item of the list.
\end{itemize}
You can even display poetry.
\begin{verse}
   There is an environment
    for verse \\             % The \\ command separates lines
   Whose features some poets % within a stanza.
   will curse.

                             % One or more blank lines separate stanzas.

   For instead of making\\
   Them do \emph{all} line breaking, \\
   It allows them to put too many words on a line when they'd rather be
   forced to be terse.
\end{verse}

Mathematical formulas may also be displayed.  A
displayed formula
is
one-line long; multiline
formulas require special formatting instructions.
   \[  \ip{\Gamma}{\psi'} = x'' + y^{2} + z_{i}^{n}\]
Don't start a paragraph with a displayed equation,
nor make one a paragraph by itself.

\end{document}               % End of document.

\section{Bench Section 38}
% synthetic bench section 38
% This is a sample LaTeX input file.  (Version of 12 August 2004.)
%
% A '%' character causes TeX to ignore all remaining text on the line,
% and is used for comments like this one.

\documentclass{article}      % Specifies the document class

                             % The preamble begins here.
\title{An Example Document}  % Declares the document's title.
\author{Leslie Lamport}      % Declares the author's name.
\date{January 21, 1994}      % Deleting this command produces today's date.

\newcommand{\ip}[2]{(#1, #2)}
                             % Defines \ip{arg1}{arg2} to mean
                             % (arg1, arg2).

%\newcommand{\ip}[2]{\langle #1 | #2\rangle}
                             % This is an alternative definition of
                             % \ip that is commented out.

\begin{document}             % End of preamble and beginning of text.

\maketitle                   % Produces the title.

This is an example input file.  Comparing it with
the output it generates can show you how to
produce a simple document of your own.

\section{Ordinary Text}      % Produces section heading.  Lower-level
                             % sections are begun with similar
                             % \subsection and \subsubsection commands.

The ends  of words and sentences are marked
  by   spaces. It  doesn't matter how many
spaces    you type; one is as good as 100.  The
end of   a line counts as a space.

One   or more   blank lines denote the  end
of  a paragraph.

Since any number of consecutive spaces are treated
like a single one, the formatting of the input
file makes no difference to
      \LaTeX,                % The \LaTeX command generates the LaTeX logo.
but it makes a difference to you.  When you use
\LaTeX, making your input file as easy to read
as possible will be a great help as you write
your document and when you change it.  This sample
file shows how you can add comments to your own input
file.

Because printing is different from typewriting,
there are a number of things that you have to do
differently when preparing an input file than if
you were just typing the document directly.
Quotation marks like
       ``this''
have to be handled specially, as do quotes within
quotes:
       ``\,`this'            % \, separates the double and single quote.
        is what I just
        wrote, not  `that'\,''.

Dashes come in three sizes: an
       intra-word
dash, a medium dash for number ranges like
       1--2,
and a punctuation
       dash---like
this.

A sentence-ending space should be larger than the
space between words within a sentence.  You
sometimes have to type special commands in
conjunction with punctuation characters to get
this right, as in the following sentence.
       Gnats, gnus, etc.\ all  % `\ ' makes an inter-word space.
       begin with G\@.         % \@ marks end-of-sentence punctuation.
You should check the spaces after periods when
reading your output to make sure you haven't
forgotten any special cases.  Generating an
ellipsis
       \ldots\               % `\ ' is needed after `\ldots' because TeX
                             % ignores spaces after command names like \ldots
                             % made from \ + letters.
                             %
                             % Note how a `%' character causes TeX to ignore
                             % the end of the input line, so these blank lines
                             % do not start a new paragraph.
                             %
with the right spacing around the periods requires
a special command.

\LaTeX\ interprets some common characters as
commands, so you must type special commands to
generate them.  These characters include the
following:
       \$ \& \% \# \{ and \}.

In printing, text is usually emphasized with an
       \emph{italic}
type style.

\begin{em}
   A long segment of text can also be emphasized
   in this way.  Text within such a segment can be
   given \emph{additional} emphasis.
\end{em}

It is sometimes necessary to prevent \LaTeX\ from
breaking a line where it might otherwise do so.
This may be at a space, as between the ``Mr.''\ and
``Jones'' in
       ``Mr.~Jones'',        % ~ produces an unbreakable interword space.
or within a word---especially when the word is a
symbol like
       \mbox{\emph{itemnum}}
that makes little sense when hyphenated across
lines.

Footnotes\footnote{This is an example of a footnote.}
pose no problem.

\LaTeX\ is good at typesetting mathematical formulas
like
       \( x-3y + z = 7 \)
or
       \( a_{1} > x^{2n} + y^{2n} > x' \)
or
       \( \ip{A}{B} = \sum_{i} a_{i} b_{i} \).
The spaces you type in a formula are
ignored.  Remember that a letter like
       $x$                   % $ ... $  and  \( ... \)  are equivalent
is a formula when it denotes a mathematical
symbol, and it should be typed as one.

\section{Displayed Text}

Text is displayed by indenting it from the left
margin.  Quotations are commonly displayed.  There
are short quotations
\begin{quote}
   This is a short quotation.  It consists of a
   single paragraph of text.  See how it is formatted.
\end{quote}
and longer ones.
\begin{quotation}
   This is a longer quotation.  It consists of two
   paragraphs of text, neither of which are
   particularly interesting.

   This is the second paragraph of the quotation.  It
   is just as dull as the first paragraph.
\end{quotation}
Another frequently-displayed structure is a list.
The following is an example of an \emph{itemized}
list.
\begin{itemize}
   \item This is the first item of an itemized list.
         Each item in the list is marked with a ``tick''.
         You don't have to worry about what kind of tick
         mark is used.

   \item This is the second item of the list.  It
         contains another list nested inside it.  The inner
         list is an \emph{enumerated} list.
         \begin{enumerate}
            \item This is the first item of an enumerated
                  list that is nested within the itemized list.

            \item This is the second item of the inner list.
                  \LaTeX\ allows you to nest lists deeper than
                  you really should.
         \end{enumerate}
         This is the rest of the second item of the outer
         list.  It is no more interesting than any other
         part of the item.
   \item This is the third item of the list.
\end{itemize}
You can even display poetry.
\begin{verse}
   There is an environment
    for verse \\             % The \\ command separates lines
   Whose features some poets % within a stanza.
   will curse.

                             % One or more blank lines separate stanzas.

   For instead of making\\
   Them do \emph{all} line breaking, \\
   It allows them to put too many words on a line when they'd rather be
   forced to be terse.
\end{verse}

Mathematical formulas may also be displayed.  A
displayed formula
is
one-line long; multiline
formulas require special formatting instructions.
   \[  \ip{\Gamma}{\psi'} = x'' + y^{2} + z_{i}^{n}\]
Don't start a paragraph with a displayed equation,
nor make one a paragraph by itself.

\end{document}               % End of document.

\section{Bench Section 39}
% synthetic bench section 39
% This is a sample LaTeX input file.  (Version of 12 August 2004.)
%
% A '%' character causes TeX to ignore all remaining text on the line,
% and is used for comments like this one.

\documentclass{article}      % Specifies the document class

                             % The preamble begins here.
\title{An Example Document}  % Declares the document's title.
\author{Leslie Lamport}      % Declares the author's name.
\date{January 21, 1994}      % Deleting this command produces today's date.

\newcommand{\ip}[2]{(#1, #2)}
                             % Defines \ip{arg1}{arg2} to mean
                             % (arg1, arg2).

%\newcommand{\ip}[2]{\langle #1 | #2\rangle}
                             % This is an alternative definition of
                             % \ip that is commented out.

\begin{document}             % End of preamble and beginning of text.

\maketitle                   % Produces the title.

This is an example input file.  Comparing it with
the output it generates can show you how to
produce a simple document of your own.

\section{Ordinary Text}      % Produces section heading.  Lower-level
                             % sections are begun with similar
                             % \subsection and \subsubsection commands.

The ends  of words and sentences are marked
  by   spaces. It  doesn't matter how many
spaces    you type; one is as good as 100.  The
end of   a line counts as a space.

One   or more   blank lines denote the  end
of  a paragraph.

Since any number of consecutive spaces are treated
like a single one, the formatting of the input
file makes no difference to
      \LaTeX,                % The \LaTeX command generates the LaTeX logo.
but it makes a difference to you.  When you use
\LaTeX, making your input file as easy to read
as possible will be a great help as you write
your document and when you change it.  This sample
file shows how you can add comments to your own input
file.

Because printing is different from typewriting,
there are a number of things that you have to do
differently when preparing an input file than if
you were just typing the document directly.
Quotation marks like
       ``this''
have to be handled specially, as do quotes within
quotes:
       ``\,`this'            % \, separates the double and single quote.
        is what I just
        wrote, not  `that'\,''.

Dashes come in three sizes: an
       intra-word
dash, a medium dash for number ranges like
       1--2,
and a punctuation
       dash---like
this.

A sentence-ending space should be larger than the
space between words within a sentence.  You
sometimes have to type special commands in
conjunction with punctuation characters to get
this right, as in the following sentence.
       Gnats, gnus, etc.\ all  % `\ ' makes an inter-word space.
       begin with G\@.         % \@ marks end-of-sentence punctuation.
You should check the spaces after periods when
reading your output to make sure you haven't
forgotten any special cases.  Generating an
ellipsis
       \ldots\               % `\ ' is needed after `\ldots' because TeX
                             % ignores spaces after command names like \ldots
                             % made from \ + letters.
                             %
                             % Note how a `%' character causes TeX to ignore
                             % the end of the input line, so these blank lines
                             % do not start a new paragraph.
                             %
with the right spacing around the periods requires
a special command.

\LaTeX\ interprets some common characters as
commands, so you must type special commands to
generate them.  These characters include the
following:
       \$ \& \% \# \{ and \}.

In printing, text is usually emphasized with an
       \emph{italic}
type style.

\begin{em}
   A long segment of text can also be emphasized
   in this way.  Text within such a segment can be
   given \emph{additional} emphasis.
\end{em}

It is sometimes necessary to prevent \LaTeX\ from
breaking a line where it might otherwise do so.
This may be at a space, as between the ``Mr.''\ and
``Jones'' in
       ``Mr.~Jones'',        % ~ produces an unbreakable interword space.
or within a word---especially when the word is a
symbol like
       \mbox{\emph{itemnum}}
that makes little sense when hyphenated across
lines.

Footnotes\footnote{This is an example of a footnote.}
pose no problem.

\LaTeX\ is good at typesetting mathematical formulas
like
       \( x-3y + z = 7 \)
or
       \( a_{1} > x^{2n} + y^{2n} > x' \)
or
       \( \ip{A}{B} = \sum_{i} a_{i} b_{i} \).
The spaces you type in a formula are
ignored.  Remember that a letter like
       $x$                   % $ ... $  and  \( ... \)  are equivalent
is a formula when it denotes a mathematical
symbol, and it should be typed as one.

\section{Displayed Text}

Text is displayed by indenting it from the left
margin.  Quotations are commonly displayed.  There
are short quotations
\begin{quote}
   This is a short quotation.  It consists of a
   single paragraph of text.  See how it is formatted.
\end{quote}
and longer ones.
\begin{quotation}
   This is a longer quotation.  It consists of two
   paragraphs of text, neither of which are
   particularly interesting.

   This is the second paragraph of the quotation.  It
   is just as dull as the first paragraph.
\end{quotation}
Another frequently-displayed structure is a list.
The following is an example of an \emph{itemized}
list.
\begin{itemize}
   \item This is the first item of an itemized list.
         Each item in the list is marked with a ``tick''.
         You don't have to worry about what kind of tick
         mark is used.

   \item This is the second item of the list.  It
         contains another list nested inside it.  The inner
         list is an \emph{enumerated} list.
         \begin{enumerate}
            \item This is the first item of an enumerated
                  list that is nested within the itemized list.

            \item This is the second item of the inner list.
                  \LaTeX\ allows you to nest lists deeper than
                  you really should.
         \end{enumerate}
         This is the rest of the second item of the outer
         list.  It is no more interesting than any other
         part of the item.
   \item This is the third item of the list.
\end{itemize}
You can even display poetry.
\begin{verse}
   There is an environment
    for verse \\             % The \\ command separates lines
   Whose features some poets % within a stanza.
   will curse.

                             % One or more blank lines separate stanzas.

   For instead of making\\
   Them do \emph{all} line breaking, \\
   It allows them to put too many words on a line when they'd rather be
   forced to be terse.
\end{verse}

Mathematical formulas may also be displayed.  A
displayed formula
is
one-line long; multiline
formulas require special formatting instructions.
   \[  \ip{\Gamma}{\psi'} = x'' + y^{2} + z_{i}^{n}\]
Don't start a paragraph with a displayed equation,
nor make one a paragraph by itself.

\end{document}               % End of document.

\section{Bench Section 40}
% synthetic bench section 40
% This is a sample LaTeX input file.  (Version of 12 August 2004.)
%
% A '%' character causes TeX to ignore all remaining text on the line,
% and is used for comments like this one.

\documentclass{article}      % Specifies the document class

                             % The preamble begins here.
\title{An Example Document}  % Declares the document's title.
\author{Leslie Lamport}      % Declares the author's name.
\date{January 21, 1994}      % Deleting this command produces today's date.

\newcommand{\ip}[2]{(#1, #2)}
                             % Defines \ip{arg1}{arg2} to mean
                             % (arg1, arg2).

%\newcommand{\ip}[2]{\langle #1 | #2\rangle}
                             % This is an alternative definition of
                             % \ip that is commented out.

\begin{document}             % End of preamble and beginning of text.

\maketitle                   % Produces the title.

This is an example input file.  Comparing it with
the output it generates can show you how to
produce a simple document of your own.

\section{Ordinary Text}      % Produces section heading.  Lower-level
                             % sections are begun with similar
                             % \subsection and \subsubsection commands.

The ends  of words and sentences are marked
  by   spaces. It  doesn't matter how many
spaces    you type; one is as good as 100.  The
end of   a line counts as a space.

One   or more   blank lines denote the  end
of  a paragraph.

Since any number of consecutive spaces are treated
like a single one, the formatting of the input
file makes no difference to
      \LaTeX,                % The \LaTeX command generates the LaTeX logo.
but it makes a difference to you.  When you use
\LaTeX, making your input file as easy to read
as possible will be a great help as you write
your document and when you change it.  This sample
file shows how you can add comments to your own input
file.

Because printing is different from typewriting,
there are a number of things that you have to do
differently when preparing an input file than if
you were just typing the document directly.
Quotation marks like
       ``this''
have to be handled specially, as do quotes within
quotes:
       ``\,`this'            % \, separates the double and single quote.
        is what I just
        wrote, not  `that'\,''.

Dashes come in three sizes: an
       intra-word
dash, a medium dash for number ranges like
       1--2,
and a punctuation
       dash---like
this.

A sentence-ending space should be larger than the
space between words within a sentence.  You
sometimes have to type special commands in
conjunction with punctuation characters to get
this right, as in the following sentence.
       Gnats, gnus, etc.\ all  % `\ ' makes an inter-word space.
       begin with G\@.         % \@ marks end-of-sentence punctuation.
You should check the spaces after periods when
reading your output to make sure you haven't
forgotten any special cases.  Generating an
ellipsis
       \ldots\               % `\ ' is needed after `\ldots' because TeX
                             % ignores spaces after command names like \ldots
                             % made from \ + letters.
                             %
                             % Note how a `%' character causes TeX to ignore
                             % the end of the input line, so these blank lines
                             % do not start a new paragraph.
                             %
with the right spacing around the periods requires
a special command.

\LaTeX\ interprets some common characters as
commands, so you must type special commands to
generate them.  These characters include the
following:
       \$ \& \% \# \{ and \}.

In printing, text is usually emphasized with an
       \emph{italic}
type style.

\begin{em}
   A long segment of text can also be emphasized
   in this way.  Text within such a segment can be
   given \emph{additional} emphasis.
\end{em}

It is sometimes necessary to prevent \LaTeX\ from
breaking a line where it might otherwise do so.
This may be at a space, as between the ``Mr.''\ and
``Jones'' in
       ``Mr.~Jones'',        % ~ produces an unbreakable interword space.
or within a word---especially when the word is a
symbol like
       \mbox{\emph{itemnum}}
that makes little sense when hyphenated across
lines.

Footnotes\footnote{This is an example of a footnote.}
pose no problem.

\LaTeX\ is good at typesetting mathematical formulas
like
       \( x-3y + z = 7 \)
or
       \( a_{1} > x^{2n} + y^{2n} > x' \)
or
       \( \ip{A}{B} = \sum_{i} a_{i} b_{i} \).
The spaces you type in a formula are
ignored.  Remember that a letter like
       $x$                   % $ ... $  and  \( ... \)  are equivalent
is a formula when it denotes a mathematical
symbol, and it should be typed as one.

\section{Displayed Text}

Text is displayed by indenting it from the left
margin.  Quotations are commonly displayed.  There
are short quotations
\begin{quote}
   This is a short quotation.  It consists of a
   single paragraph of text.  See how it is formatted.
\end{quote}
and longer ones.
\begin{quotation}
   This is a longer quotation.  It consists of two
   paragraphs of text, neither of which are
   particularly interesting.

   This is the second paragraph of the quotation.  It
   is just as dull as the first paragraph.
\end{quotation}
Another frequently-displayed structure is a list.
The following is an example of an \emph{itemized}
list.
\begin{itemize}
   \item This is the first item of an itemized list.
         Each item in the list is marked with a ``tick''.
         You don't have to worry about what kind of tick
         mark is used.

   \item This is the second item of the list.  It
         contains another list nested inside it.  The inner
         list is an \emph{enumerated} list.
         \begin{enumerate}
            \item This is the first item of an enumerated
                  list that is nested within the itemized list.

            \item This is the second item of the inner list.
                  \LaTeX\ allows you to nest lists deeper than
                  you really should.
         \end{enumerate}
         This is the rest of the second item of the outer
         list.  It is no more interesting than any other
         part of the item.
   \item This is the third item of the list.
\end{itemize}
You can even display poetry.
\begin{verse}
   There is an environment
    for verse \\             % The \\ command separates lines
   Whose features some poets % within a stanza.
   will curse.

                             % One or more blank lines separate stanzas.

   For instead of making\\
   Them do \emph{all} line breaking, \\
   It allows them to put too many words on a line when they'd rather be
   forced to be terse.
\end{verse}

Mathematical formulas may also be displayed.  A
displayed formula
is
one-line long; multiline
formulas require special formatting instructions.
   \[  \ip{\Gamma}{\psi'} = x'' + y^{2} + z_{i}^{n}\]
Don't start a paragraph with a displayed equation,
nor make one a paragraph by itself.

\end{document}               % End of document.

\section{Bench Section 41}
% synthetic bench section 41
% This is a sample LaTeX input file.  (Version of 12 August 2004.)
%
% A '%' character causes TeX to ignore all remaining text on the line,
% and is used for comments like this one.

\documentclass{article}      % Specifies the document class

                             % The preamble begins here.
\title{An Example Document}  % Declares the document's title.
\author{Leslie Lamport}      % Declares the author's name.
\date{January 21, 1994}      % Deleting this command produces today's date.

\newcommand{\ip}[2]{(#1, #2)}
                             % Defines \ip{arg1}{arg2} to mean
                             % (arg1, arg2).

%\newcommand{\ip}[2]{\langle #1 | #2\rangle}
                             % This is an alternative definition of
                             % \ip that is commented out.

\begin{document}             % End of preamble and beginning of text.

\maketitle                   % Produces the title.

This is an example input file.  Comparing it with
the output it generates can show you how to
produce a simple document of your own.

\section{Ordinary Text}      % Produces section heading.  Lower-level
                             % sections are begun with similar
                             % \subsection and \subsubsection commands.

The ends  of words and sentences are marked
  by   spaces. It  doesn't matter how many
spaces    you type; one is as good as 100.  The
end of   a line counts as a space.

One   or more   blank lines denote the  end
of  a paragraph.

Since any number of consecutive spaces are treated
like a single one, the formatting of the input
file makes no difference to
      \LaTeX,                % The \LaTeX command generates the LaTeX logo.
but it makes a difference to you.  When you use
\LaTeX, making your input file as easy to read
as possible will be a great help as you write
your document and when you change it.  This sample
file shows how you can add comments to your own input
file.

Because printing is different from typewriting,
there are a number of things that you have to do
differently when preparing an input file than if
you were just typing the document directly.
Quotation marks like
       ``this''
have to be handled specially, as do quotes within
quotes:
       ``\,`this'            % \, separates the double and single quote.
        is what I just
        wrote, not  `that'\,''.

Dashes come in three sizes: an
       intra-word
dash, a medium dash for number ranges like
       1--2,
and a punctuation
       dash---like
this.

A sentence-ending space should be larger than the
space between words within a sentence.  You
sometimes have to type special commands in
conjunction with punctuation characters to get
this right, as in the following sentence.
       Gnats, gnus, etc.\ all  % `\ ' makes an inter-word space.
       begin with G\@.         % \@ marks end-of-sentence punctuation.
You should check the spaces after periods when
reading your output to make sure you haven't
forgotten any special cases.  Generating an
ellipsis
       \ldots\               % `\ ' is needed after `\ldots' because TeX
                             % ignores spaces after command names like \ldots
                             % made from \ + letters.
                             %
                             % Note how a `%' character causes TeX to ignore
                             % the end of the input line, so these blank lines
                             % do not start a new paragraph.
                             %
with the right spacing around the periods requires
a special command.

\LaTeX\ interprets some common characters as
commands, so you must type special commands to
generate them.  These characters include the
following:
       \$ \& \% \# \{ and \}.

In printing, text is usually emphasized with an
       \emph{italic}
type style.

\begin{em}
   A long segment of text can also be emphasized
   in this way.  Text within such a segment can be
   given \emph{additional} emphasis.
\end{em}

It is sometimes necessary to prevent \LaTeX\ from
breaking a line where it might otherwise do so.
This may be at a space, as between the ``Mr.''\ and
``Jones'' in
       ``Mr.~Jones'',        % ~ produces an unbreakable interword space.
or within a word---especially when the word is a
symbol like
       \mbox{\emph{itemnum}}
that makes little sense when hyphenated across
lines.

Footnotes\footnote{This is an example of a footnote.}
pose no problem.

\LaTeX\ is good at typesetting mathematical formulas
like
       \( x-3y + z = 7 \)
or
       \( a_{1} > x^{2n} + y^{2n} > x' \)
or
       \( \ip{A}{B} = \sum_{i} a_{i} b_{i} \).
The spaces you type in a formula are
ignored.  Remember that a letter like
       $x$                   % $ ... $  and  \( ... \)  are equivalent
is a formula when it denotes a mathematical
symbol, and it should be typed as one.

\section{Displayed Text}

Text is displayed by indenting it from the left
margin.  Quotations are commonly displayed.  There
are short quotations
\begin{quote}
   This is a short quotation.  It consists of a
   single paragraph of text.  See how it is formatted.
\end{quote}
and longer ones.
\begin{quotation}
   This is a longer quotation.  It consists of two
   paragraphs of text, neither of which are
   particularly interesting.

   This is the second paragraph of the quotation.  It
   is just as dull as the first paragraph.
\end{quotation}
Another frequently-displayed structure is a list.
The following is an example of an \emph{itemized}
list.
\begin{itemize}
   \item This is the first item of an itemized list.
         Each item in the list is marked with a ``tick''.
         You don't have to worry about what kind of tick
         mark is used.

   \item This is the second item of the list.  It
         contains another list nested inside it.  The inner
         list is an \emph{enumerated} list.
         \begin{enumerate}
            \item This is the first item of an enumerated
                  list that is nested within the itemized list.

            \item This is the second item of the inner list.
                  \LaTeX\ allows you to nest lists deeper than
                  you really should.
         \end{enumerate}
         This is the rest of the second item of the outer
         list.  It is no more interesting than any other
         part of the item.
   \item This is the third item of the list.
\end{itemize}
You can even display poetry.
\begin{verse}
   There is an environment
    for verse \\             % The \\ command separates lines
   Whose features some poets % within a stanza.
   will curse.

                             % One or more blank lines separate stanzas.

   For instead of making\\
   Them do \emph{all} line breaking, \\
   It allows them to put too many words on a line when they'd rather be
   forced to be terse.
\end{verse}

Mathematical formulas may also be displayed.  A
displayed formula
is
one-line long; multiline
formulas require special formatting instructions.
   \[  \ip{\Gamma}{\psi'} = x'' + y^{2} + z_{i}^{n}\]
Don't start a paragraph with a displayed equation,
nor make one a paragraph by itself.

\end{document}               % End of document.

\section{Bench Section 42}
% synthetic bench section 42
% This is a sample LaTeX input file.  (Version of 12 August 2004.)
%
% A '%' character causes TeX to ignore all remaining text on the line,
% and is used for comments like this one.

\documentclass{article}      % Specifies the document class

                             % The preamble begins here.
\title{An Example Document}  % Declares the document's title.
\author{Leslie Lamport}      % Declares the author's name.
\date{January 21, 1994}      % Deleting this command produces today's date.

\newcommand{\ip}[2]{(#1, #2)}
                             % Defines \ip{arg1}{arg2} to mean
                             % (arg1, arg2).

%\newcommand{\ip}[2]{\langle #1 | #2\rangle}
                             % This is an alternative definition of
                             % \ip that is commented out.

\begin{document}             % End of preamble and beginning of text.

\maketitle                   % Produces the title.

This is an example input file.  Comparing it with
the output it generates can show you how to
produce a simple document of your own.

\section{Ordinary Text}      % Produces section heading.  Lower-level
                             % sections are begun with similar
                             % \subsection and \subsubsection commands.

The ends  of words and sentences are marked
  by   spaces. It  doesn't matter how many
spaces    you type; one is as good as 100.  The
end of   a line counts as a space.

One   or more   blank lines denote the  end
of  a paragraph.

Since any number of consecutive spaces are treated
like a single one, the formatting of the input
file makes no difference to
      \LaTeX,                % The \LaTeX command generates the LaTeX logo.
but it makes a difference to you.  When you use
\LaTeX, making your input file as easy to read
as possible will be a great help as you write
your document and when you change it.  This sample
file shows how you can add comments to your own input
file.

Because printing is different from typewriting,
there are a number of things that you have to do
differently when preparing an input file than if
you were just typing the document directly.
Quotation marks like
       ``this''
have to be handled specially, as do quotes within
quotes:
       ``\,`this'            % \, separates the double and single quote.
        is what I just
        wrote, not  `that'\,''.

Dashes come in three sizes: an
       intra-word
dash, a medium dash for number ranges like
       1--2,
and a punctuation
       dash---like
this.

A sentence-ending space should be larger than the
space between words within a sentence.  You
sometimes have to type special commands in
conjunction with punctuation characters to get
this right, as in the following sentence.
       Gnats, gnus, etc.\ all  % `\ ' makes an inter-word space.
       begin with G\@.         % \@ marks end-of-sentence punctuation.
You should check the spaces after periods when
reading your output to make sure you haven't
forgotten any special cases.  Generating an
ellipsis
       \ldots\               % `\ ' is needed after `\ldots' because TeX
                             % ignores spaces after command names like \ldots
                             % made from \ + letters.
                             %
                             % Note how a `%' character causes TeX to ignore
                             % the end of the input line, so these blank lines
                             % do not start a new paragraph.
                             %
with the right spacing around the periods requires
a special command.

\LaTeX\ interprets some common characters as
commands, so you must type special commands to
generate them.  These characters include the
following:
       \$ \& \% \# \{ and \}.

In printing, text is usually emphasized with an
       \emph{italic}
type style.

\begin{em}
   A long segment of text can also be emphasized
   in this way.  Text within such a segment can be
   given \emph{additional} emphasis.
\end{em}

It is sometimes necessary to prevent \LaTeX\ from
breaking a line where it might otherwise do so.
This may be at a space, as between the ``Mr.''\ and
``Jones'' in
       ``Mr.~Jones'',        % ~ produces an unbreakable interword space.
or within a word---especially when the word is a
symbol like
       \mbox{\emph{itemnum}}
that makes little sense when hyphenated across
lines.

Footnotes\footnote{This is an example of a footnote.}
pose no problem.

\LaTeX\ is good at typesetting mathematical formulas
like
       \( x-3y + z = 7 \)
or
       \( a_{1} > x^{2n} + y^{2n} > x' \)
or
       \( \ip{A}{B} = \sum_{i} a_{i} b_{i} \).
The spaces you type in a formula are
ignored.  Remember that a letter like
       $x$                   % $ ... $  and  \( ... \)  are equivalent
is a formula when it denotes a mathematical
symbol, and it should be typed as one.

\section{Displayed Text}

Text is displayed by indenting it from the left
margin.  Quotations are commonly displayed.  There
are short quotations
\begin{quote}
   This is a short quotation.  It consists of a
   single paragraph of text.  See how it is formatted.
\end{quote}
and longer ones.
\begin{quotation}
   This is a longer quotation.  It consists of two
   paragraphs of text, neither of which are
   particularly interesting.

   This is the second paragraph of the quotation.  It
   is just as dull as the first paragraph.
\end{quotation}
Another frequently-displayed structure is a list.
The following is an example of an \emph{itemized}
list.
\begin{itemize}
   \item This is the first item of an itemized list.
         Each item in the list is marked with a ``tick''.
         You don't have to worry about what kind of tick
         mark is used.

   \item This is the second item of the list.  It
         contains another list nested inside it.  The inner
         list is an \emph{enumerated} list.
         \begin{enumerate}
            \item This is the first item of an enumerated
                  list that is nested within the itemized list.

            \item This is the second item of the inner list.
                  \LaTeX\ allows you to nest lists deeper than
                  you really should.
         \end{enumerate}
         This is the rest of the second item of the outer
         list.  It is no more interesting than any other
         part of the item.
   \item This is the third item of the list.
\end{itemize}
You can even display poetry.
\begin{verse}
   There is an environment
    for verse \\             % The \\ command separates lines
   Whose features some poets % within a stanza.
   will curse.

                             % One or more blank lines separate stanzas.

   For instead of making\\
   Them do \emph{all} line breaking, \\
   It allows them to put too many words on a line when they'd rather be
   forced to be terse.
\end{verse}

Mathematical formulas may also be displayed.  A
displayed formula
is
one-line long; multiline
formulas require special formatting instructions.
   \[  \ip{\Gamma}{\psi'} = x'' + y^{2} + z_{i}^{n}\]
Don't start a paragraph with a displayed equation,
nor make one a paragraph by itself.

\end{document}               % End of document.

\section{Bench Section 43}
% synthetic bench section 43
% This is a sample LaTeX input file.  (Version of 12 August 2004.)
%
% A '%' character causes TeX to ignore all remaining text on the line,
% and is used for comments like this one.

\documentclass{article}      % Specifies the document class

                             % The preamble begins here.
\title{An Example Document}  % Declares the document's title.
\author{Leslie Lamport}      % Declares the author's name.
\date{January 21, 1994}      % Deleting this command produces today's date.

\newcommand{\ip}[2]{(#1, #2)}
                             % Defines \ip{arg1}{arg2} to mean
                             % (arg1, arg2).

%\newcommand{\ip}[2]{\langle #1 | #2\rangle}
                             % This is an alternative definition of
                             % \ip that is commented out.

\begin{document}             % End of preamble and beginning of text.

\maketitle                   % Produces the title.

This is an example input file.  Comparing it with
the output it generates can show you how to
produce a simple document of your own.

\section{Ordinary Text}      % Produces section heading.  Lower-level
                             % sections are begun with similar
                             % \subsection and \subsubsection commands.

The ends  of words and sentences are marked
  by   spaces. It  doesn't matter how many
spaces    you type; one is as good as 100.  The
end of   a line counts as a space.

One   or more   blank lines denote the  end
of  a paragraph.

Since any number of consecutive spaces are treated
like a single one, the formatting of the input
file makes no difference to
      \LaTeX,                % The \LaTeX command generates the LaTeX logo.
but it makes a difference to you.  When you use
\LaTeX, making your input file as easy to read
as possible will be a great help as you write
your document and when you change it.  This sample
file shows how you can add comments to your own input
file.

Because printing is different from typewriting,
there are a number of things that you have to do
differently when preparing an input file than if
you were just typing the document directly.
Quotation marks like
       ``this''
have to be handled specially, as do quotes within
quotes:
       ``\,`this'            % \, separates the double and single quote.
        is what I just
        wrote, not  `that'\,''.

Dashes come in three sizes: an
       intra-word
dash, a medium dash for number ranges like
       1--2,
and a punctuation
       dash---like
this.

A sentence-ending space should be larger than the
space between words within a sentence.  You
sometimes have to type special commands in
conjunction with punctuation characters to get
this right, as in the following sentence.
       Gnats, gnus, etc.\ all  % `\ ' makes an inter-word space.
       begin with G\@.         % \@ marks end-of-sentence punctuation.
You should check the spaces after periods when
reading your output to make sure you haven't
forgotten any special cases.  Generating an
ellipsis
       \ldots\               % `\ ' is needed after `\ldots' because TeX
                             % ignores spaces after command names like \ldots
                             % made from \ + letters.
                             %
                             % Note how a `%' character causes TeX to ignore
                             % the end of the input line, so these blank lines
                             % do not start a new paragraph.
                             %
with the right spacing around the periods requires
a special command.

\LaTeX\ interprets some common characters as
commands, so you must type special commands to
generate them.  These characters include the
following:
       \$ \& \% \# \{ and \}.

In printing, text is usually emphasized with an
       \emph{italic}
type style.

\begin{em}
   A long segment of text can also be emphasized
   in this way.  Text within such a segment can be
   given \emph{additional} emphasis.
\end{em}

It is sometimes necessary to prevent \LaTeX\ from
breaking a line where it might otherwise do so.
This may be at a space, as between the ``Mr.''\ and
``Jones'' in
       ``Mr.~Jones'',        % ~ produces an unbreakable interword space.
or within a word---especially when the word is a
symbol like
       \mbox{\emph{itemnum}}
that makes little sense when hyphenated across
lines.

Footnotes\footnote{This is an example of a footnote.}
pose no problem.

\LaTeX\ is good at typesetting mathematical formulas
like
       \( x-3y + z = 7 \)
or
       \( a_{1} > x^{2n} + y^{2n} > x' \)
or
       \( \ip{A}{B} = \sum_{i} a_{i} b_{i} \).
The spaces you type in a formula are
ignored.  Remember that a letter like
       $x$                   % $ ... $  and  \( ... \)  are equivalent
is a formula when it denotes a mathematical
symbol, and it should be typed as one.

\section{Displayed Text}

Text is displayed by indenting it from the left
margin.  Quotations are commonly displayed.  There
are short quotations
\begin{quote}
   This is a short quotation.  It consists of a
   single paragraph of text.  See how it is formatted.
\end{quote}
and longer ones.
\begin{quotation}
   This is a longer quotation.  It consists of two
   paragraphs of text, neither of which are
   particularly interesting.

   This is the second paragraph of the quotation.  It
   is just as dull as the first paragraph.
\end{quotation}
Another frequently-displayed structure is a list.
The following is an example of an \emph{itemized}
list.
\begin{itemize}
   \item This is the first item of an itemized list.
         Each item in the list is marked with a ``tick''.
         You don't have to worry about what kind of tick
         mark is used.

   \item This is the second item of the list.  It
         contains another list nested inside it.  The inner
         list is an \emph{enumerated} list.
         \begin{enumerate}
            \item This is the first item of an enumerated
                  list that is nested within the itemized list.

            \item This is the second item of the inner list.
                  \LaTeX\ allows you to nest lists deeper than
                  you really should.
         \end{enumerate}
         This is the rest of the second item of the outer
         list.  It is no more interesting than any other
         part of the item.
   \item This is the third item of the list.
\end{itemize}
You can even display poetry.
\begin{verse}
   There is an environment
    for verse \\             % The \\ command separates lines
   Whose features some poets % within a stanza.
   will curse.

                             % One or more blank lines separate stanzas.

   For instead of making\\
   Them do \emph{all} line breaking, \\
   It allows them to put too many words on a line when they'd rather be
   forced to be terse.
\end{verse}

Mathematical formulas may also be displayed.  A
displayed formula
is
one-line long; multiline
formulas require special formatting instructions.
   \[  \ip{\Gamma}{\psi'} = x'' + y^{2} + z_{i}^{n}\]
Don't start a paragraph with a displayed equation,
nor make one a paragraph by itself.

\end{document}               % End of document.

\section{Bench Section 44}
% synthetic bench section 44
% This is a sample LaTeX input file.  (Version of 12 August 2004.)
%
% A '%' character causes TeX to ignore all remaining text on the line,
% and is used for comments like this one.

\documentclass{article}      % Specifies the document class

                             % The preamble begins here.
\title{An Example Document}  % Declares the document's title.
\author{Leslie Lamport}      % Declares the author's name.
\date{January 21, 1994}      % Deleting this command produces today's date.

\newcommand{\ip}[2]{(#1, #2)}
                             % Defines \ip{arg1}{arg2} to mean
                             % (arg1, arg2).

%\newcommand{\ip}[2]{\langle #1 | #2\rangle}
                             % This is an alternative definition of
                             % \ip that is commented out.

\begin{document}             % End of preamble and beginning of text.

\maketitle                   % Produces the title.

This is an example input file.  Comparing it with
the output it generates can show you how to
produce a simple document of your own.

\section{Ordinary Text}      % Produces section heading.  Lower-level
                             % sections are begun with similar
                             % \subsection and \subsubsection commands.

The ends  of words and sentences are marked
  by   spaces. It  doesn't matter how many
spaces    you type; one is as good as 100.  The
end of   a line counts as a space.

One   or more   blank lines denote the  end
of  a paragraph.

Since any number of consecutive spaces are treated
like a single one, the formatting of the input
file makes no difference to
      \LaTeX,                % The \LaTeX command generates the LaTeX logo.
but it makes a difference to you.  When you use
\LaTeX, making your input file as easy to read
as possible will be a great help as you write
your document and when you change it.  This sample
file shows how you can add comments to your own input
file.

Because printing is different from typewriting,
there are a number of things that you have to do
differently when preparing an input file than if
you were just typing the document directly.
Quotation marks like
       ``this''
have to be handled specially, as do quotes within
quotes:
       ``\,`this'            % \, separates the double and single quote.
        is what I just
        wrote, not  `that'\,''.

Dashes come in three sizes: an
       intra-word
dash, a medium dash for number ranges like
       1--2,
and a punctuation
       dash---like
this.

A sentence-ending space should be larger than the
space between words within a sentence.  You
sometimes have to type special commands in
conjunction with punctuation characters to get
this right, as in the following sentence.
       Gnats, gnus, etc.\ all  % `\ ' makes an inter-word space.
       begin with G\@.         % \@ marks end-of-sentence punctuation.
You should check the spaces after periods when
reading your output to make sure you haven't
forgotten any special cases.  Generating an
ellipsis
       \ldots\               % `\ ' is needed after `\ldots' because TeX
                             % ignores spaces after command names like \ldots
                             % made from \ + letters.
                             %
                             % Note how a `%' character causes TeX to ignore
                             % the end of the input line, so these blank lines
                             % do not start a new paragraph.
                             %
with the right spacing around the periods requires
a special command.

\LaTeX\ interprets some common characters as
commands, so you must type special commands to
generate them.  These characters include the
following:
       \$ \& \% \# \{ and \}.

In printing, text is usually emphasized with an
       \emph{italic}
type style.

\begin{em}
   A long segment of text can also be emphasized
   in this way.  Text within such a segment can be
   given \emph{additional} emphasis.
\end{em}

It is sometimes necessary to prevent \LaTeX\ from
breaking a line where it might otherwise do so.
This may be at a space, as between the ``Mr.''\ and
``Jones'' in
       ``Mr.~Jones'',        % ~ produces an unbreakable interword space.
or within a word---especially when the word is a
symbol like
       \mbox{\emph{itemnum}}
that makes little sense when hyphenated across
lines.

Footnotes\footnote{This is an example of a footnote.}
pose no problem.

\LaTeX\ is good at typesetting mathematical formulas
like
       \( x-3y + z = 7 \)
or
       \( a_{1} > x^{2n} + y^{2n} > x' \)
or
       \( \ip{A}{B} = \sum_{i} a_{i} b_{i} \).
The spaces you type in a formula are
ignored.  Remember that a letter like
       $x$                   % $ ... $  and  \( ... \)  are equivalent
is a formula when it denotes a mathematical
symbol, and it should be typed as one.

\section{Displayed Text}

Text is displayed by indenting it from the left
margin.  Quotations are commonly displayed.  There
are short quotations
\begin{quote}
   This is a short quotation.  It consists of a
   single paragraph of text.  See how it is formatted.
\end{quote}
and longer ones.
\begin{quotation}
   This is a longer quotation.  It consists of two
   paragraphs of text, neither of which are
   particularly interesting.

   This is the second paragraph of the quotation.  It
   is just as dull as the first paragraph.
\end{quotation}
Another frequently-displayed structure is a list.
The following is an example of an \emph{itemized}
list.
\begin{itemize}
   \item This is the first item of an itemized list.
         Each item in the list is marked with a ``tick''.
         You don't have to worry about what kind of tick
         mark is used.

   \item This is the second item of the list.  It
         contains another list nested inside it.  The inner
         list is an \emph{enumerated} list.
         \begin{enumerate}
            \item This is the first item of an enumerated
                  list that is nested within the itemized list.

            \item This is the second item of the inner list.
                  \LaTeX\ allows you to nest lists deeper than
                  you really should.
         \end{enumerate}
         This is the rest of the second item of the outer
         list.  It is no more interesting than any other
         part of the item.
   \item This is the third item of the list.
\end{itemize}
You can even display poetry.
\begin{verse}
   There is an environment
    for verse \\             % The \\ command separates lines
   Whose features some poets % within a stanza.
   will curse.

                             % One or more blank lines separate stanzas.

   For instead of making\\
   Them do \emph{all} line breaking, \\
   It allows them to put too many words on a line when they'd rather be
   forced to be terse.
\end{verse}

Mathematical formulas may also be displayed.  A
displayed formula
is
one-line long; multiline
formulas require special formatting instructions.
   \[  \ip{\Gamma}{\psi'} = x'' + y^{2} + z_{i}^{n}\]
Don't start a paragraph with a displayed equation,
nor make one a paragraph by itself.

\end{document}               % End of document.

\section{Bench Section 45}
% synthetic bench section 45
% This is a sample LaTeX input file.  (Version of 12 August 2004.)
%
% A '%' character causes TeX to ignore all remaining text on the line,
% and is used for comments like this one.

\documentclass{article}      % Specifies the document class

                             % The preamble begins here.
\title{An Example Document}  % Declares the document's title.
\author{Leslie Lamport}      % Declares the author's name.
\date{January 21, 1994}      % Deleting this command produces today's date.

\newcommand{\ip}[2]{(#1, #2)}
                             % Defines \ip{arg1}{arg2} to mean
                             % (arg1, arg2).

%\newcommand{\ip}[2]{\langle #1 | #2\rangle}
                             % This is an alternative definition of
                             % \ip that is commented out.

\begin{document}             % End of preamble and beginning of text.

\maketitle                   % Produces the title.

This is an example input file.  Comparing it with
the output it generates can show you how to
produce a simple document of your own.

\section{Ordinary Text}      % Produces section heading.  Lower-level
                             % sections are begun with similar
                             % \subsection and \subsubsection commands.

The ends  of words and sentences are marked
  by   spaces. It  doesn't matter how many
spaces    you type; one is as good as 100.  The
end of   a line counts as a space.

One   or more   blank lines denote the  end
of  a paragraph.

Since any number of consecutive spaces are treated
like a single one, the formatting of the input
file makes no difference to
      \LaTeX,                % The \LaTeX command generates the LaTeX logo.
but it makes a difference to you.  When you use
\LaTeX, making your input file as easy to read
as possible will be a great help as you write
your document and when you change it.  This sample
file shows how you can add comments to your own input
file.

Because printing is different from typewriting,
there are a number of things that you have to do
differently when preparing an input file than if
you were just typing the document directly.
Quotation marks like
       ``this''
have to be handled specially, as do quotes within
quotes:
       ``\,`this'            % \, separates the double and single quote.
        is what I just
        wrote, not  `that'\,''.

Dashes come in three sizes: an
       intra-word
dash, a medium dash for number ranges like
       1--2,
and a punctuation
       dash---like
this.

A sentence-ending space should be larger than the
space between words within a sentence.  You
sometimes have to type special commands in
conjunction with punctuation characters to get
this right, as in the following sentence.
       Gnats, gnus, etc.\ all  % `\ ' makes an inter-word space.
       begin with G\@.         % \@ marks end-of-sentence punctuation.
You should check the spaces after periods when
reading your output to make sure you haven't
forgotten any special cases.  Generating an
ellipsis
       \ldots\               % `\ ' is needed after `\ldots' because TeX
                             % ignores spaces after command names like \ldots
                             % made from \ + letters.
                             %
                             % Note how a `%' character causes TeX to ignore
                             % the end of the input line, so these blank lines
                             % do not start a new paragraph.
                             %
with the right spacing around the periods requires
a special command.

\LaTeX\ interprets some common characters as
commands, so you must type special commands to
generate them.  These characters include the
following:
       \$ \& \% \# \{ and \}.

In printing, text is usually emphasized with an
       \emph{italic}
type style.

\begin{em}
   A long segment of text can also be emphasized
   in this way.  Text within such a segment can be
   given \emph{additional} emphasis.
\end{em}

It is sometimes necessary to prevent \LaTeX\ from
breaking a line where it might otherwise do so.
This may be at a space, as between the ``Mr.''\ and
``Jones'' in
       ``Mr.~Jones'',        % ~ produces an unbreakable interword space.
or within a word---especially when the word is a
symbol like
       \mbox{\emph{itemnum}}
that makes little sense when hyphenated across
lines.

Footnotes\footnote{This is an example of a footnote.}
pose no problem.

\LaTeX\ is good at typesetting mathematical formulas
like
       \( x-3y + z = 7 \)
or
       \( a_{1} > x^{2n} + y^{2n} > x' \)
or
       \( \ip{A}{B} = \sum_{i} a_{i} b_{i} \).
The spaces you type in a formula are
ignored.  Remember that a letter like
       $x$                   % $ ... $  and  \( ... \)  are equivalent
is a formula when it denotes a mathematical
symbol, and it should be typed as one.

\section{Displayed Text}

Text is displayed by indenting it from the left
margin.  Quotations are commonly displayed.  There
are short quotations
\begin{quote}
   This is a short quotation.  It consists of a
   single paragraph of text.  See how it is formatted.
\end{quote}
and longer ones.
\begin{quotation}
   This is a longer quotation.  It consists of two
   paragraphs of text, neither of which are
   particularly interesting.

   This is the second paragraph of the quotation.  It
   is just as dull as the first paragraph.
\end{quotation}
Another frequently-displayed structure is a list.
The following is an example of an \emph{itemized}
list.
\begin{itemize}
   \item This is the first item of an itemized list.
         Each item in the list is marked with a ``tick''.
         You don't have to worry about what kind of tick
         mark is used.

   \item This is the second item of the list.  It
         contains another list nested inside it.  The inner
         list is an \emph{enumerated} list.
         \begin{enumerate}
            \item This is the first item of an enumerated
                  list that is nested within the itemized list.

            \item This is the second item of the inner list.
                  \LaTeX\ allows you to nest lists deeper than
                  you really should.
         \end{enumerate}
         This is the rest of the second item of the outer
         list.  It is no more interesting than any other
         part of the item.
   \item This is the third item of the list.
\end{itemize}
You can even display poetry.
\begin{verse}
   There is an environment
    for verse \\             % The \\ command separates lines
   Whose features some poets % within a stanza.
   will curse.

                             % One or more blank lines separate stanzas.

   For instead of making\\
   Them do \emph{all} line breaking, \\
   It allows them to put too many words on a line when they'd rather be
   forced to be terse.
\end{verse}

Mathematical formulas may also be displayed.  A
displayed formula
is
one-line long; multiline
formulas require special formatting instructions.
   \[  \ip{\Gamma}{\psi'} = x'' + y^{2} + z_{i}^{n}\]
Don't start a paragraph with a displayed equation,
nor make one a paragraph by itself.

\end{document}               % End of document.

\section{Bench Section 46}
% synthetic bench section 46
% This is a sample LaTeX input file.  (Version of 12 August 2004.)
%
% A '%' character causes TeX to ignore all remaining text on the line,
% and is used for comments like this one.

\documentclass{article}      % Specifies the document class

                             % The preamble begins here.
\title{An Example Document}  % Declares the document's title.
\author{Leslie Lamport}      % Declares the author's name.
\date{January 21, 1994}      % Deleting this command produces today's date.

\newcommand{\ip}[2]{(#1, #2)}
                             % Defines \ip{arg1}{arg2} to mean
                             % (arg1, arg2).

%\newcommand{\ip}[2]{\langle #1 | #2\rangle}
                             % This is an alternative definition of
                             % \ip that is commented out.

\begin{document}             % End of preamble and beginning of text.

\maketitle                   % Produces the title.

This is an example input file.  Comparing it with
the output it generates can show you how to
produce a simple document of your own.

\section{Ordinary Text}      % Produces section heading.  Lower-level
                             % sections are begun with similar
                             % \subsection and \subsubsection commands.

The ends  of words and sentences are marked
  by   spaces. It  doesn't matter how many
spaces    you type; one is as good as 100.  The
end of   a line counts as a space.

One   or more   blank lines denote the  end
of  a paragraph.

Since any number of consecutive spaces are treated
like a single one, the formatting of the input
file makes no difference to
      \LaTeX,                % The \LaTeX command generates the LaTeX logo.
but it makes a difference to you.  When you use
\LaTeX, making your input file as easy to read
as possible will be a great help as you write
your document and when you change it.  This sample
file shows how you can add comments to your own input
file.

Because printing is different from typewriting,
there are a number of things that you have to do
differently when preparing an input file than if
you were just typing the document directly.
Quotation marks like
       ``this''
have to be handled specially, as do quotes within
quotes:
       ``\,`this'            % \, separates the double and single quote.
        is what I just
        wrote, not  `that'\,''.

Dashes come in three sizes: an
       intra-word
dash, a medium dash for number ranges like
       1--2,
and a punctuation
       dash---like
this.

A sentence-ending space should be larger than the
space between words within a sentence.  You
sometimes have to type special commands in
conjunction with punctuation characters to get
this right, as in the following sentence.
       Gnats, gnus, etc.\ all  % `\ ' makes an inter-word space.
       begin with G\@.         % \@ marks end-of-sentence punctuation.
You should check the spaces after periods when
reading your output to make sure you haven't
forgotten any special cases.  Generating an
ellipsis
       \ldots\               % `\ ' is needed after `\ldots' because TeX
                             % ignores spaces after command names like \ldots
                             % made from \ + letters.
                             %
                             % Note how a `%' character causes TeX to ignore
                             % the end of the input line, so these blank lines
                             % do not start a new paragraph.
                             %
with the right spacing around the periods requires
a special command.

\LaTeX\ interprets some common characters as
commands, so you must type special commands to
generate them.  These characters include the
following:
       \$ \& \% \# \{ and \}.

In printing, text is usually emphasized with an
       \emph{italic}
type style.

\begin{em}
   A long segment of text can also be emphasized
   in this way.  Text within such a segment can be
   given \emph{additional} emphasis.
\end{em}

It is sometimes necessary to prevent \LaTeX\ from
breaking a line where it might otherwise do so.
This may be at a space, as between the ``Mr.''\ and
``Jones'' in
       ``Mr.~Jones'',        % ~ produces an unbreakable interword space.
or within a word---especially when the word is a
symbol like
       \mbox{\emph{itemnum}}
that makes little sense when hyphenated across
lines.

Footnotes\footnote{This is an example of a footnote.}
pose no problem.

\LaTeX\ is good at typesetting mathematical formulas
like
       \( x-3y + z = 7 \)
or
       \( a_{1} > x^{2n} + y^{2n} > x' \)
or
       \( \ip{A}{B} = \sum_{i} a_{i} b_{i} \).
The spaces you type in a formula are
ignored.  Remember that a letter like
       $x$                   % $ ... $  and  \( ... \)  are equivalent
is a formula when it denotes a mathematical
symbol, and it should be typed as one.

\section{Displayed Text}

Text is displayed by indenting it from the left
margin.  Quotations are commonly displayed.  There
are short quotations
\begin{quote}
   This is a short quotation.  It consists of a
   single paragraph of text.  See how it is formatted.
\end{quote}
and longer ones.
\begin{quotation}
   This is a longer quotation.  It consists of two
   paragraphs of text, neither of which are
   particularly interesting.

   This is the second paragraph of the quotation.  It
   is just as dull as the first paragraph.
\end{quotation}
Another frequently-displayed structure is a list.
The following is an example of an \emph{itemized}
list.
\begin{itemize}
   \item This is the first item of an itemized list.
         Each item in the list is marked with a ``tick''.
         You don't have to worry about what kind of tick
         mark is used.

   \item This is the second item of the list.  It
         contains another list nested inside it.  The inner
         list is an \emph{enumerated} list.
         \begin{enumerate}
            \item This is the first item of an enumerated
                  list that is nested within the itemized list.

            \item This is the second item of the inner list.
                  \LaTeX\ allows you to nest lists deeper than
                  you really should.
         \end{enumerate}
         This is the rest of the second item of the outer
         list.  It is no more interesting than any other
         part of the item.
   \item This is the third item of the list.
\end{itemize}
You can even display poetry.
\begin{verse}
   There is an environment
    for verse \\             % The \\ command separates lines
   Whose features some poets % within a stanza.
   will curse.

                             % One or more blank lines separate stanzas.

   For instead of making\\
   Them do \emph{all} line breaking, \\
   It allows them to put too many words on a line when they'd rather be
   forced to be terse.
\end{verse}

Mathematical formulas may also be displayed.  A
displayed formula
is
one-line long; multiline
formulas require special formatting instructions.
   \[  \ip{\Gamma}{\psi'} = x'' + y^{2} + z_{i}^{n}\]
Don't start a paragraph with a displayed equation,
nor make one a paragraph by itself.

\end{document}               % End of document.

\section{Bench Section 47}
% synthetic bench section 47
% This is a sample LaTeX input file.  (Version of 12 August 2004.)
%
% A '%' character causes TeX to ignore all remaining text on the line,
% and is used for comments like this one.

\documentclass{article}      % Specifies the document class

                             % The preamble begins here.
\title{An Example Document}  % Declares the document's title.
\author{Leslie Lamport}      % Declares the author's name.
\date{January 21, 1994}      % Deleting this command produces today's date.

\newcommand{\ip}[2]{(#1, #2)}
                             % Defines \ip{arg1}{arg2} to mean
                             % (arg1, arg2).

%\newcommand{\ip}[2]{\langle #1 | #2\rangle}
                             % This is an alternative definition of
                             % \ip that is commented out.

\begin{document}             % End of preamble and beginning of text.

\maketitle                   % Produces the title.

This is an example input file.  Comparing it with
the output it generates can show you how to
produce a simple document of your own.

\section{Ordinary Text}      % Produces section heading.  Lower-level
                             % sections are begun with similar
                             % \subsection and \subsubsection commands.

The ends  of words and sentences are marked
  by   spaces. It  doesn't matter how many
spaces    you type; one is as good as 100.  The
end of   a line counts as a space.

One   or more   blank lines denote the  end
of  a paragraph.

Since any number of consecutive spaces are treated
like a single one, the formatting of the input
file makes no difference to
      \LaTeX,                % The \LaTeX command generates the LaTeX logo.
but it makes a difference to you.  When you use
\LaTeX, making your input file as easy to read
as possible will be a great help as you write
your document and when you change it.  This sample
file shows how you can add comments to your own input
file.

Because printing is different from typewriting,
there are a number of things that you have to do
differently when preparing an input file than if
you were just typing the document directly.
Quotation marks like
       ``this''
have to be handled specially, as do quotes within
quotes:
       ``\,`this'            % \, separates the double and single quote.
        is what I just
        wrote, not  `that'\,''.

Dashes come in three sizes: an
       intra-word
dash, a medium dash for number ranges like
       1--2,
and a punctuation
       dash---like
this.

A sentence-ending space should be larger than the
space between words within a sentence.  You
sometimes have to type special commands in
conjunction with punctuation characters to get
this right, as in the following sentence.
       Gnats, gnus, etc.\ all  % `\ ' makes an inter-word space.
       begin with G\@.         % \@ marks end-of-sentence punctuation.
You should check the spaces after periods when
reading your output to make sure you haven't
forgotten any special cases.  Generating an
ellipsis
       \ldots\               % `\ ' is needed after `\ldots' because TeX
                             % ignores spaces after command names like \ldots
                             % made from \ + letters.
                             %
                             % Note how a `%' character causes TeX to ignore
                             % the end of the input line, so these blank lines
                             % do not start a new paragraph.
                             %
with the right spacing around the periods requires
a special command.

\LaTeX\ interprets some common characters as
commands, so you must type special commands to
generate them.  These characters include the
following:
       \$ \& \% \# \{ and \}.

In printing, text is usually emphasized with an
       \emph{italic}
type style.

\begin{em}
   A long segment of text can also be emphasized
   in this way.  Text within such a segment can be
   given \emph{additional} emphasis.
\end{em}

It is sometimes necessary to prevent \LaTeX\ from
breaking a line where it might otherwise do so.
This may be at a space, as between the ``Mr.''\ and
``Jones'' in
       ``Mr.~Jones'',        % ~ produces an unbreakable interword space.
or within a word---especially when the word is a
symbol like
       \mbox{\emph{itemnum}}
that makes little sense when hyphenated across
lines.

Footnotes\footnote{This is an example of a footnote.}
pose no problem.

\LaTeX\ is good at typesetting mathematical formulas
like
       \( x-3y + z = 7 \)
or
       \( a_{1} > x^{2n} + y^{2n} > x' \)
or
       \( \ip{A}{B} = \sum_{i} a_{i} b_{i} \).
The spaces you type in a formula are
ignored.  Remember that a letter like
       $x$                   % $ ... $  and  \( ... \)  are equivalent
is a formula when it denotes a mathematical
symbol, and it should be typed as one.

\section{Displayed Text}

Text is displayed by indenting it from the left
margin.  Quotations are commonly displayed.  There
are short quotations
\begin{quote}
   This is a short quotation.  It consists of a
   single paragraph of text.  See how it is formatted.
\end{quote}
and longer ones.
\begin{quotation}
   This is a longer quotation.  It consists of two
   paragraphs of text, neither of which are
   particularly interesting.

   This is the second paragraph of the quotation.  It
   is just as dull as the first paragraph.
\end{quotation}
Another frequently-displayed structure is a list.
The following is an example of an \emph{itemized}
list.
\begin{itemize}
   \item This is the first item of an itemized list.
         Each item in the list is marked with a ``tick''.
         You don't have to worry about what kind of tick
         mark is used.

   \item This is the second item of the list.  It
         contains another list nested inside it.  The inner
         list is an \emph{enumerated} list.
         \begin{enumerate}
            \item This is the first item of an enumerated
                  list that is nested within the itemized list.

            \item This is the second item of the inner list.
                  \LaTeX\ allows you to nest lists deeper than
                  you really should.
         \end{enumerate}
         This is the rest of the second item of the outer
         list.  It is no more interesting than any other
         part of the item.
   \item This is the third item of the list.
\end{itemize}
You can even display poetry.
\begin{verse}
   There is an environment
    for verse \\             % The \\ command separates lines
   Whose features some poets % within a stanza.
   will curse.

                             % One or more blank lines separate stanzas.

   For instead of making\\
   Them do \emph{all} line breaking, \\
   It allows them to put too many words on a line when they'd rather be
   forced to be terse.
\end{verse}

Mathematical formulas may also be displayed.  A
displayed formula
is
one-line long; multiline
formulas require special formatting instructions.
   \[  \ip{\Gamma}{\psi'} = x'' + y^{2} + z_{i}^{n}\]
Don't start a paragraph with a displayed equation,
nor make one a paragraph by itself.

\end{document}               % End of document.

\section{Bench Section 48}
% synthetic bench section 48
% This is a sample LaTeX input file.  (Version of 12 August 2004.)
%
% A '%' character causes TeX to ignore all remaining text on the line,
% and is used for comments like this one.

\documentclass{article}      % Specifies the document class

                             % The preamble begins here.
\title{An Example Document}  % Declares the document's title.
\author{Leslie Lamport}      % Declares the author's name.
\date{January 21, 1994}      % Deleting this command produces today's date.

\newcommand{\ip}[2]{(#1, #2)}
                             % Defines \ip{arg1}{arg2} to mean
                             % (arg1, arg2).

%\newcommand{\ip}[2]{\langle #1 | #2\rangle}
                             % This is an alternative definition of
                             % \ip that is commented out.

\begin{document}             % End of preamble and beginning of text.

\maketitle                   % Produces the title.

This is an example input file.  Comparing it with
the output it generates can show you how to
produce a simple document of your own.

\section{Ordinary Text}      % Produces section heading.  Lower-level
                             % sections are begun with similar
                             % \subsection and \subsubsection commands.

The ends  of words and sentences are marked
  by   spaces. It  doesn't matter how many
spaces    you type; one is as good as 100.  The
end of   a line counts as a space.

One   or more   blank lines denote the  end
of  a paragraph.

Since any number of consecutive spaces are treated
like a single one, the formatting of the input
file makes no difference to
      \LaTeX,                % The \LaTeX command generates the LaTeX logo.
but it makes a difference to you.  When you use
\LaTeX, making your input file as easy to read
as possible will be a great help as you write
your document and when you change it.  This sample
file shows how you can add comments to your own input
file.

Because printing is different from typewriting,
there are a number of things that you have to do
differently when preparing an input file than if
you were just typing the document directly.
Quotation marks like
       ``this''
have to be handled specially, as do quotes within
quotes:
       ``\,`this'            % \, separates the double and single quote.
        is what I just
        wrote, not  `that'\,''.

Dashes come in three sizes: an
       intra-word
dash, a medium dash for number ranges like
       1--2,
and a punctuation
       dash---like
this.

A sentence-ending space should be larger than the
space between words within a sentence.  You
sometimes have to type special commands in
conjunction with punctuation characters to get
this right, as in the following sentence.
       Gnats, gnus, etc.\ all  % `\ ' makes an inter-word space.
       begin with G\@.         % \@ marks end-of-sentence punctuation.
You should check the spaces after periods when
reading your output to make sure you haven't
forgotten any special cases.  Generating an
ellipsis
       \ldots\               % `\ ' is needed after `\ldots' because TeX
                             % ignores spaces after command names like \ldots
                             % made from \ + letters.
                             %
                             % Note how a `%' character causes TeX to ignore
                             % the end of the input line, so these blank lines
                             % do not start a new paragraph.
                             %
with the right spacing around the periods requires
a special command.

\LaTeX\ interprets some common characters as
commands, so you must type special commands to
generate them.  These characters include the
following:
       \$ \& \% \# \{ and \}.

In printing, text is usually emphasized with an
       \emph{italic}
type style.

\begin{em}
   A long segment of text can also be emphasized
   in this way.  Text within such a segment can be
   given \emph{additional} emphasis.
\end{em}

It is sometimes necessary to prevent \LaTeX\ from
breaking a line where it might otherwise do so.
This may be at a space, as between the ``Mr.''\ and
``Jones'' in
       ``Mr.~Jones'',        % ~ produces an unbreakable interword space.
or within a word---especially when the word is a
symbol like
       \mbox{\emph{itemnum}}
that makes little sense when hyphenated across
lines.

Footnotes\footnote{This is an example of a footnote.}
pose no problem.

\LaTeX\ is good at typesetting mathematical formulas
like
       \( x-3y + z = 7 \)
or
       \( a_{1} > x^{2n} + y^{2n} > x' \)
or
       \( \ip{A}{B} = \sum_{i} a_{i} b_{i} \).
The spaces you type in a formula are
ignored.  Remember that a letter like
       $x$                   % $ ... $  and  \( ... \)  are equivalent
is a formula when it denotes a mathematical
symbol, and it should be typed as one.

\section{Displayed Text}

Text is displayed by indenting it from the left
margin.  Quotations are commonly displayed.  There
are short quotations
\begin{quote}
   This is a short quotation.  It consists of a
   single paragraph of text.  See how it is formatted.
\end{quote}
and longer ones.
\begin{quotation}
   This is a longer quotation.  It consists of two
   paragraphs of text, neither of which are
   particularly interesting.

   This is the second paragraph of the quotation.  It
   is just as dull as the first paragraph.
\end{quotation}
Another frequently-displayed structure is a list.
The following is an example of an \emph{itemized}
list.
\begin{itemize}
   \item This is the first item of an itemized list.
         Each item in the list is marked with a ``tick''.
         You don't have to worry about what kind of tick
         mark is used.

   \item This is the second item of the list.  It
         contains another list nested inside it.  The inner
         list is an \emph{enumerated} list.
         \begin{enumerate}
            \item This is the first item of an enumerated
                  list that is nested within the itemized list.

            \item This is the second item of the inner list.
                  \LaTeX\ allows you to nest lists deeper than
                  you really should.
         \end{enumerate}
         This is the rest of the second item of the outer
         list.  It is no more interesting than any other
         part of the item.
   \item This is the third item of the list.
\end{itemize}
You can even display poetry.
\begin{verse}
   There is an environment
    for verse \\             % The \\ command separates lines
   Whose features some poets % within a stanza.
   will curse.

                             % One or more blank lines separate stanzas.

   For instead of making\\
   Them do \emph{all} line breaking, \\
   It allows them to put too many words on a line when they'd rather be
   forced to be terse.
\end{verse}

Mathematical formulas may also be displayed.  A
displayed formula
is
one-line long; multiline
formulas require special formatting instructions.
   \[  \ip{\Gamma}{\psi'} = x'' + y^{2} + z_{i}^{n}\]
Don't start a paragraph with a displayed equation,
nor make one a paragraph by itself.

\end{document}               % End of document.

\section{Bench Section 49}
% synthetic bench section 49
% This is a sample LaTeX input file.  (Version of 12 August 2004.)
%
% A '%' character causes TeX to ignore all remaining text on the line,
% and is used for comments like this one.

\documentclass{article}      % Specifies the document class

                             % The preamble begins here.
\title{An Example Document}  % Declares the document's title.
\author{Leslie Lamport}      % Declares the author's name.
\date{January 21, 1994}      % Deleting this command produces today's date.

\newcommand{\ip}[2]{(#1, #2)}
                             % Defines \ip{arg1}{arg2} to mean
                             % (arg1, arg2).

%\newcommand{\ip}[2]{\langle #1 | #2\rangle}
                             % This is an alternative definition of
                             % \ip that is commented out.

\begin{document}             % End of preamble and beginning of text.

\maketitle                   % Produces the title.

This is an example input file.  Comparing it with
the output it generates can show you how to
produce a simple document of your own.

\section{Ordinary Text}      % Produces section heading.  Lower-level
                             % sections are begun with similar
                             % \subsection and \subsubsection commands.

The ends  of words and sentences are marked
  by   spaces. It  doesn't matter how many
spaces    you type; one is as good as 100.  The
end of   a line counts as a space.

One   or more   blank lines denote the  end
of  a paragraph.

Since any number of consecutive spaces are treated
like a single one, the formatting of the input
file makes no difference to
      \LaTeX,                % The \LaTeX command generates the LaTeX logo.
but it makes a difference to you.  When you use
\LaTeX, making your input file as easy to read
as possible will be a great help as you write
your document and when you change it.  This sample
file shows how you can add comments to your own input
file.

Because printing is different from typewriting,
there are a number of things that you have to do
differently when preparing an input file than if
you were just typing the document directly.
Quotation marks like
       ``this''
have to be handled specially, as do quotes within
quotes:
       ``\,`this'            % \, separates the double and single quote.
        is what I just
        wrote, not  `that'\,''.

Dashes come in three sizes: an
       intra-word
dash, a medium dash for number ranges like
       1--2,
and a punctuation
       dash---like
this.

A sentence-ending space should be larger than the
space between words within a sentence.  You
sometimes have to type special commands in
conjunction with punctuation characters to get
this right, as in the following sentence.
       Gnats, gnus, etc.\ all  % `\ ' makes an inter-word space.
       begin with G\@.         % \@ marks end-of-sentence punctuation.
You should check the spaces after periods when
reading your output to make sure you haven't
forgotten any special cases.  Generating an
ellipsis
       \ldots\               % `\ ' is needed after `\ldots' because TeX
                             % ignores spaces after command names like \ldots
                             % made from \ + letters.
                             %
                             % Note how a `%' character causes TeX to ignore
                             % the end of the input line, so these blank lines
                             % do not start a new paragraph.
                             %
with the right spacing around the periods requires
a special command.

\LaTeX\ interprets some common characters as
commands, so you must type special commands to
generate them.  These characters include the
following:
       \$ \& \% \# \{ and \}.

In printing, text is usually emphasized with an
       \emph{italic}
type style.

\begin{em}
   A long segment of text can also be emphasized
   in this way.  Text within such a segment can be
   given \emph{additional} emphasis.
\end{em}

It is sometimes necessary to prevent \LaTeX\ from
breaking a line where it might otherwise do so.
This may be at a space, as between the ``Mr.''\ and
``Jones'' in
       ``Mr.~Jones'',        % ~ produces an unbreakable interword space.
or within a word---especially when the word is a
symbol like
       \mbox{\emph{itemnum}}
that makes little sense when hyphenated across
lines.

Footnotes\footnote{This is an example of a footnote.}
pose no problem.

\LaTeX\ is good at typesetting mathematical formulas
like
       \( x-3y + z = 7 \)
or
       \( a_{1} > x^{2n} + y^{2n} > x' \)
or
       \( \ip{A}{B} = \sum_{i} a_{i} b_{i} \).
The spaces you type in a formula are
ignored.  Remember that a letter like
       $x$                   % $ ... $  and  \( ... \)  are equivalent
is a formula when it denotes a mathematical
symbol, and it should be typed as one.

\section{Displayed Text}

Text is displayed by indenting it from the left
margin.  Quotations are commonly displayed.  There
are short quotations
\begin{quote}
   This is a short quotation.  It consists of a
   single paragraph of text.  See how it is formatted.
\end{quote}
and longer ones.
\begin{quotation}
   This is a longer quotation.  It consists of two
   paragraphs of text, neither of which are
   particularly interesting.

   This is the second paragraph of the quotation.  It
   is just as dull as the first paragraph.
\end{quotation}
Another frequently-displayed structure is a list.
The following is an example of an \emph{itemized}
list.
\begin{itemize}
   \item This is the first item of an itemized list.
         Each item in the list is marked with a ``tick''.
         You don't have to worry about what kind of tick
         mark is used.

   \item This is the second item of the list.  It
         contains another list nested inside it.  The inner
         list is an \emph{enumerated} list.
         \begin{enumerate}
            \item This is the first item of an enumerated
                  list that is nested within the itemized list.

            \item This is the second item of the inner list.
                  \LaTeX\ allows you to nest lists deeper than
                  you really should.
         \end{enumerate}
         This is the rest of the second item of the outer
         list.  It is no more interesting than any other
         part of the item.
   \item This is the third item of the list.
\end{itemize}
You can even display poetry.
\begin{verse}
   There is an environment
    for verse \\             % The \\ command separates lines
   Whose features some poets % within a stanza.
   will curse.

                             % One or more blank lines separate stanzas.

   For instead of making\\
   Them do \emph{all} line breaking, \\
   It allows them to put too many words on a line when they'd rather be
   forced to be terse.
\end{verse}

Mathematical formulas may also be displayed.  A
displayed formula
is
one-line long; multiline
formulas require special formatting instructions.
   \[  \ip{\Gamma}{\psi'} = x'' + y^{2} + z_{i}^{n}\]
Don't start a paragraph with a displayed equation,
nor make one a paragraph by itself.

\end{document}               % End of document.

\section{Bench Section 50}
% synthetic bench section 50
% This is a sample LaTeX input file.  (Version of 12 August 2004.)
%
% A '%' character causes TeX to ignore all remaining text on the line,
% and is used for comments like this one.

\documentclass{article}      % Specifies the document class

                             % The preamble begins here.
\title{An Example Document}  % Declares the document's title.
\author{Leslie Lamport}      % Declares the author's name.
\date{January 21, 1994}      % Deleting this command produces today's date.

\newcommand{\ip}[2]{(#1, #2)}
                             % Defines \ip{arg1}{arg2} to mean
                             % (arg1, arg2).

%\newcommand{\ip}[2]{\langle #1 | #2\rangle}
                             % This is an alternative definition of
                             % \ip that is commented out.

\begin{document}             % End of preamble and beginning of text.

\maketitle                   % Produces the title.

This is an example input file.  Comparing it with
the output it generates can show you how to
produce a simple document of your own.

\section{Ordinary Text}      % Produces section heading.  Lower-level
                             % sections are begun with similar
                             % \subsection and \subsubsection commands.

The ends  of words and sentences are marked
  by   spaces. It  doesn't matter how many
spaces    you type; one is as good as 100.  The
end of   a line counts as a space.

One   or more   blank lines denote the  end
of  a paragraph.

Since any number of consecutive spaces are treated
like a single one, the formatting of the input
file makes no difference to
      \LaTeX,                % The \LaTeX command generates the LaTeX logo.
but it makes a difference to you.  When you use
\LaTeX, making your input file as easy to read
as possible will be a great help as you write
your document and when you change it.  This sample
file shows how you can add comments to your own input
file.

Because printing is different from typewriting,
there are a number of things that you have to do
differently when preparing an input file than if
you were just typing the document directly.
Quotation marks like
       ``this''
have to be handled specially, as do quotes within
quotes:
       ``\,`this'            % \, separates the double and single quote.
        is what I just
        wrote, not  `that'\,''.

Dashes come in three sizes: an
       intra-word
dash, a medium dash for number ranges like
       1--2,
and a punctuation
       dash---like
this.

A sentence-ending space should be larger than the
space between words within a sentence.  You
sometimes have to type special commands in
conjunction with punctuation characters to get
this right, as in the following sentence.
       Gnats, gnus, etc.\ all  % `\ ' makes an inter-word space.
       begin with G\@.         % \@ marks end-of-sentence punctuation.
You should check the spaces after periods when
reading your output to make sure you haven't
forgotten any special cases.  Generating an
ellipsis
       \ldots\               % `\ ' is needed after `\ldots' because TeX
                             % ignores spaces after command names like \ldots
                             % made from \ + letters.
                             %
                             % Note how a `%' character causes TeX to ignore
                             % the end of the input line, so these blank lines
                             % do not start a new paragraph.
                             %
with the right spacing around the periods requires
a special command.

\LaTeX\ interprets some common characters as
commands, so you must type special commands to
generate them.  These characters include the
following:
       \$ \& \% \# \{ and \}.

In printing, text is usually emphasized with an
       \emph{italic}
type style.

\begin{em}
   A long segment of text can also be emphasized
   in this way.  Text within such a segment can be
   given \emph{additional} emphasis.
\end{em}

It is sometimes necessary to prevent \LaTeX\ from
breaking a line where it might otherwise do so.
This may be at a space, as between the ``Mr.''\ and
``Jones'' in
       ``Mr.~Jones'',        % ~ produces an unbreakable interword space.
or within a word---especially when the word is a
symbol like
       \mbox{\emph{itemnum}}
that makes little sense when hyphenated across
lines.

Footnotes\footnote{This is an example of a footnote.}
pose no problem.

\LaTeX\ is good at typesetting mathematical formulas
like
       \( x-3y + z = 7 \)
or
       \( a_{1} > x^{2n} + y^{2n} > x' \)
or
       \( \ip{A}{B} = \sum_{i} a_{i} b_{i} \).
The spaces you type in a formula are
ignored.  Remember that a letter like
       $x$                   % $ ... $  and  \( ... \)  are equivalent
is a formula when it denotes a mathematical
symbol, and it should be typed as one.

\section{Displayed Text}

Text is displayed by indenting it from the left
margin.  Quotations are commonly displayed.  There
are short quotations
\begin{quote}
   This is a short quotation.  It consists of a
   single paragraph of text.  See how it is formatted.
\end{quote}
and longer ones.
\begin{quotation}
   This is a longer quotation.  It consists of two
   paragraphs of text, neither of which are
   particularly interesting.

   This is the second paragraph of the quotation.  It
   is just as dull as the first paragraph.
\end{quotation}
Another frequently-displayed structure is a list.
The following is an example of an \emph{itemized}
list.
\begin{itemize}
   \item This is the first item of an itemized list.
         Each item in the list is marked with a ``tick''.
         You don't have to worry about what kind of tick
         mark is used.

   \item This is the second item of the list.  It
         contains another list nested inside it.  The inner
         list is an \emph{enumerated} list.
         \begin{enumerate}
            \item This is the first item of an enumerated
                  list that is nested within the itemized list.

            \item This is the second item of the inner list.
                  \LaTeX\ allows you to nest lists deeper than
                  you really should.
         \end{enumerate}
         This is the rest of the second item of the outer
         list.  It is no more interesting than any other
         part of the item.
   \item This is the third item of the list.
\end{itemize}
You can even display poetry.
\begin{verse}
   There is an environment
    for verse \\             % The \\ command separates lines
   Whose features some poets % within a stanza.
   will curse.

                             % One or more blank lines separate stanzas.

   For instead of making\\
   Them do \emph{all} line breaking, \\
   It allows them to put too many words on a line when they'd rather be
   forced to be terse.
\end{verse}

Mathematical formulas may also be displayed.  A
displayed formula
is
one-line long; multiline
formulas require special formatting instructions.
   \[  \ip{\Gamma}{\psi'} = x'' + y^{2} + z_{i}^{n}\]
Don't start a paragraph with a displayed equation,
nor make one a paragraph by itself.

\end{document}               % End of document.

\section{Bench Section 51}
% synthetic bench section 51
% This is a sample LaTeX input file.  (Version of 12 August 2004.)
%
% A '%' character causes TeX to ignore all remaining text on the line,
% and is used for comments like this one.

\documentclass{article}      % Specifies the document class

                             % The preamble begins here.
\title{An Example Document}  % Declares the document's title.
\author{Leslie Lamport}      % Declares the author's name.
\date{January 21, 1994}      % Deleting this command produces today's date.

\newcommand{\ip}[2]{(#1, #2)}
                             % Defines \ip{arg1}{arg2} to mean
                             % (arg1, arg2).

%\newcommand{\ip}[2]{\langle #1 | #2\rangle}
                             % This is an alternative definition of
                             % \ip that is commented out.

\begin{document}             % End of preamble and beginning of text.

\maketitle                   % Produces the title.

This is an example input file.  Comparing it with
the output it generates can show you how to
produce a simple document of your own.

\section{Ordinary Text}      % Produces section heading.  Lower-level
                             % sections are begun with similar
                             % \subsection and \subsubsection commands.

The ends  of words and sentences are marked
  by   spaces. It  doesn't matter how many
spaces    you type; one is as good as 100.  The
end of   a line counts as a space.

One   or more   blank lines denote the  end
of  a paragraph.

Since any number of consecutive spaces are treated
like a single one, the formatting of the input
file makes no difference to
      \LaTeX,                % The \LaTeX command generates the LaTeX logo.
but it makes a difference to you.  When you use
\LaTeX, making your input file as easy to read
as possible will be a great help as you write
your document and when you change it.  This sample
file shows how you can add comments to your own input
file.

Because printing is different from typewriting,
there are a number of things that you have to do
differently when preparing an input file than if
you were just typing the document directly.
Quotation marks like
       ``this''
have to be handled specially, as do quotes within
quotes:
       ``\,`this'            % \, separates the double and single quote.
        is what I just
        wrote, not  `that'\,''.

Dashes come in three sizes: an
       intra-word
dash, a medium dash for number ranges like
       1--2,
and a punctuation
       dash---like
this.

A sentence-ending space should be larger than the
space between words within a sentence.  You
sometimes have to type special commands in
conjunction with punctuation characters to get
this right, as in the following sentence.
       Gnats, gnus, etc.\ all  % `\ ' makes an inter-word space.
       begin with G\@.         % \@ marks end-of-sentence punctuation.
You should check the spaces after periods when
reading your output to make sure you haven't
forgotten any special cases.  Generating an
ellipsis
       \ldots\               % `\ ' is needed after `\ldots' because TeX
                             % ignores spaces after command names like \ldots
                             % made from \ + letters.
                             %
                             % Note how a `%' character causes TeX to ignore
                             % the end of the input line, so these blank lines
                             % do not start a new paragraph.
                             %
with the right spacing around the periods requires
a special command.

\LaTeX\ interprets some common characters as
commands, so you must type special commands to
generate them.  These characters include the
following:
       \$ \& \% \# \{ and \}.

In printing, text is usually emphasized with an
       \emph{italic}
type style.

\begin{em}
   A long segment of text can also be emphasized
   in this way.  Text within such a segment can be
   given \emph{additional} emphasis.
\end{em}

It is sometimes necessary to prevent \LaTeX\ from
breaking a line where it might otherwise do so.
This may be at a space, as between the ``Mr.''\ and
``Jones'' in
       ``Mr.~Jones'',        % ~ produces an unbreakable interword space.
or within a word---especially when the word is a
symbol like
       \mbox{\emph{itemnum}}
that makes little sense when hyphenated across
lines.

Footnotes\footnote{This is an example of a footnote.}
pose no problem.

\LaTeX\ is good at typesetting mathematical formulas
like
       \( x-3y + z = 7 \)
or
       \( a_{1} > x^{2n} + y^{2n} > x' \)
or
       \( \ip{A}{B} = \sum_{i} a_{i} b_{i} \).
The spaces you type in a formula are
ignored.  Remember that a letter like
       $x$                   % $ ... $  and  \( ... \)  are equivalent
is a formula when it denotes a mathematical
symbol, and it should be typed as one.

\section{Displayed Text}

Text is displayed by indenting it from the left
margin.  Quotations are commonly displayed.  There
are short quotations
\begin{quote}
   This is a short quotation.  It consists of a
   single paragraph of text.  See how it is formatted.
\end{quote}
and longer ones.
\begin{quotation}
   This is a longer quotation.  It consists of two
   paragraphs of text, neither of which are
   particularly interesting.

   This is the second paragraph of the quotation.  It
   is just as dull as the first paragraph.
\end{quotation}
Another frequently-displayed structure is a list.
The following is an example of an \emph{itemized}
list.
\begin{itemize}
   \item This is the first item of an itemized list.
         Each item in the list is marked with a ``tick''.
         You don't have to worry about what kind of tick
         mark is used.

   \item This is the second item of the list.  It
         contains another list nested inside it.  The inner
         list is an \emph{enumerated} list.
         \begin{enumerate}
            \item This is the first item of an enumerated
                  list that is nested within the itemized list.

            \item This is the second item of the inner list.
                  \LaTeX\ allows you to nest lists deeper than
                  you really should.
         \end{enumerate}
         This is the rest of the second item of the outer
         list.  It is no more interesting than any other
         part of the item.
   \item This is the third item of the list.
\end{itemize}
You can even display poetry.
\begin{verse}
   There is an environment
    for verse \\             % The \\ command separates lines
   Whose features some poets % within a stanza.
   will curse.

                             % One or more blank lines separate stanzas.

   For instead of making\\
   Them do \emph{all} line breaking, \\
   It allows them to put too many words on a line when they'd rather be
   forced to be terse.
\end{verse}

Mathematical formulas may also be displayed.  A
displayed formula
is
one-line long; multiline
formulas require special formatting instructions.
   \[  \ip{\Gamma}{\psi'} = x'' + y^{2} + z_{i}^{n}\]
Don't start a paragraph with a displayed equation,
nor make one a paragraph by itself.

\end{document}               % End of document.

\section{Bench Section 52}
% synthetic bench section 52
% This is a sample LaTeX input file.  (Version of 12 August 2004.)
%
% A '%' character causes TeX to ignore all remaining text on the line,
% and is used for comments like this one.

\documentclass{article}      % Specifies the document class

                             % The preamble begins here.
\title{An Example Document}  % Declares the document's title.
\author{Leslie Lamport}      % Declares the author's name.
\date{January 21, 1994}      % Deleting this command produces today's date.

\newcommand{\ip}[2]{(#1, #2)}
                             % Defines \ip{arg1}{arg2} to mean
                             % (arg1, arg2).

%\newcommand{\ip}[2]{\langle #1 | #2\rangle}
                             % This is an alternative definition of
                             % \ip that is commented out.

\begin{document}             % End of preamble and beginning of text.

\maketitle                   % Produces the title.

This is an example input file.  Comparing it with
the output it generates can show you how to
produce a simple document of your own.

\section{Ordinary Text}      % Produces section heading.  Lower-level
                             % sections are begun with similar
                             % \subsection and \subsubsection commands.

The ends  of words and sentences are marked
  by   spaces. It  doesn't matter how many
spaces    you type; one is as good as 100.  The
end of   a line counts as a space.

One   or more   blank lines denote the  end
of  a paragraph.

Since any number of consecutive spaces are treated
like a single one, the formatting of the input
file makes no difference to
      \LaTeX,                % The \LaTeX command generates the LaTeX logo.
but it makes a difference to you.  When you use
\LaTeX, making your input file as easy to read
as possible will be a great help as you write
your document and when you change it.  This sample
file shows how you can add comments to your own input
file.

Because printing is different from typewriting,
there are a number of things that you have to do
differently when preparing an input file than if
you were just typing the document directly.
Quotation marks like
       ``this''
have to be handled specially, as do quotes within
quotes:
       ``\,`this'            % \, separates the double and single quote.
        is what I just
        wrote, not  `that'\,''.

Dashes come in three sizes: an
       intra-word
dash, a medium dash for number ranges like
       1--2,
and a punctuation
       dash---like
this.

A sentence-ending space should be larger than the
space between words within a sentence.  You
sometimes have to type special commands in
conjunction with punctuation characters to get
this right, as in the following sentence.
       Gnats, gnus, etc.\ all  % `\ ' makes an inter-word space.
       begin with G\@.         % \@ marks end-of-sentence punctuation.
You should check the spaces after periods when
reading your output to make sure you haven't
forgotten any special cases.  Generating an
ellipsis
       \ldots\               % `\ ' is needed after `\ldots' because TeX
                             % ignores spaces after command names like \ldots
                             % made from \ + letters.
                             %
                             % Note how a `%' character causes TeX to ignore
                             % the end of the input line, so these blank lines
                             % do not start a new paragraph.
                             %
with the right spacing around the periods requires
a special command.

\LaTeX\ interprets some common characters as
commands, so you must type special commands to
generate them.  These characters include the
following:
       \$ \& \% \# \{ and \}.

In printing, text is usually emphasized with an
       \emph{italic}
type style.

\begin{em}
   A long segment of text can also be emphasized
   in this way.  Text within such a segment can be
   given \emph{additional} emphasis.
\end{em}

It is sometimes necessary to prevent \LaTeX\ from
breaking a line where it might otherwise do so.
This may be at a space, as between the ``Mr.''\ and
``Jones'' in
       ``Mr.~Jones'',        % ~ produces an unbreakable interword space.
or within a word---especially when the word is a
symbol like
       \mbox{\emph{itemnum}}
that makes little sense when hyphenated across
lines.

Footnotes\footnote{This is an example of a footnote.}
pose no problem.

\LaTeX\ is good at typesetting mathematical formulas
like
       \( x-3y + z = 7 \)
or
       \( a_{1} > x^{2n} + y^{2n} > x' \)
or
       \( \ip{A}{B} = \sum_{i} a_{i} b_{i} \).
The spaces you type in a formula are
ignored.  Remember that a letter like
       $x$                   % $ ... $  and  \( ... \)  are equivalent
is a formula when it denotes a mathematical
symbol, and it should be typed as one.

\section{Displayed Text}

Text is displayed by indenting it from the left
margin.  Quotations are commonly displayed.  There
are short quotations
\begin{quote}
   This is a short quotation.  It consists of a
   single paragraph of text.  See how it is formatted.
\end{quote}
and longer ones.
\begin{quotation}
   This is a longer quotation.  It consists of two
   paragraphs of text, neither of which are
   particularly interesting.

   This is the second paragraph of the quotation.  It
   is just as dull as the first paragraph.
\end{quotation}
Another frequently-displayed structure is a list.
The following is an example of an \emph{itemized}
list.
\begin{itemize}
   \item This is the first item of an itemized list.
         Each item in the list is marked with a ``tick''.
         You don't have to worry about what kind of tick
         mark is used.

   \item This is the second item of the list.  It
         contains another list nested inside it.  The inner
         list is an \emph{enumerated} list.
         \begin{enumerate}
            \item This is the first item of an enumerated
                  list that is nested within the itemized list.

            \item This is the second item of the inner list.
                  \LaTeX\ allows you to nest lists deeper than
                  you really should.
         \end{enumerate}
         This is the rest of the second item of the outer
         list.  It is no more interesting than any other
         part of the item.
   \item This is the third item of the list.
\end{itemize}
You can even display poetry.
\begin{verse}
   There is an environment
    for verse \\             % The \\ command separates lines
   Whose features some poets % within a stanza.
   will curse.

                             % One or more blank lines separate stanzas.

   For instead of making\\
   Them do \emph{all} line breaking, \\
   It allows them to put too many words on a line when they'd rather be
   forced to be terse.
\end{verse}

Mathematical formulas may also be displayed.  A
displayed formula
is
one-line long; multiline
formulas require special formatting instructions.
   \[  \ip{\Gamma}{\psi'} = x'' + y^{2} + z_{i}^{n}\]
Don't start a paragraph with a displayed equation,
nor make one a paragraph by itself.

\end{document}               % End of document.

\section{Bench Section 53}
% synthetic bench section 53
% This is a sample LaTeX input file.  (Version of 12 August 2004.)
%
% A '%' character causes TeX to ignore all remaining text on the line,
% and is used for comments like this one.

\documentclass{article}      % Specifies the document class

                             % The preamble begins here.
\title{An Example Document}  % Declares the document's title.
\author{Leslie Lamport}      % Declares the author's name.
\date{January 21, 1994}      % Deleting this command produces today's date.

\newcommand{\ip}[2]{(#1, #2)}
                             % Defines \ip{arg1}{arg2} to mean
                             % (arg1, arg2).

%\newcommand{\ip}[2]{\langle #1 | #2\rangle}
                             % This is an alternative definition of
                             % \ip that is commented out.

\begin{document}             % End of preamble and beginning of text.

\maketitle                   % Produces the title.

This is an example input file.  Comparing it with
the output it generates can show you how to
produce a simple document of your own.

\section{Ordinary Text}      % Produces section heading.  Lower-level
                             % sections are begun with similar
                             % \subsection and \subsubsection commands.

The ends  of words and sentences are marked
  by   spaces. It  doesn't matter how many
spaces    you type; one is as good as 100.  The
end of   a line counts as a space.

One   or more   blank lines denote the  end
of  a paragraph.

Since any number of consecutive spaces are treated
like a single one, the formatting of the input
file makes no difference to
      \LaTeX,                % The \LaTeX command generates the LaTeX logo.
but it makes a difference to you.  When you use
\LaTeX, making your input file as easy to read
as possible will be a great help as you write
your document and when you change it.  This sample
file shows how you can add comments to your own input
file.

Because printing is different from typewriting,
there are a number of things that you have to do
differently when preparing an input file than if
you were just typing the document directly.
Quotation marks like
       ``this''
have to be handled specially, as do quotes within
quotes:
       ``\,`this'            % \, separates the double and single quote.
        is what I just
        wrote, not  `that'\,''.

Dashes come in three sizes: an
       intra-word
dash, a medium dash for number ranges like
       1--2,
and a punctuation
       dash---like
this.

A sentence-ending space should be larger than the
space between words within a sentence.  You
sometimes have to type special commands in
conjunction with punctuation characters to get
this right, as in the following sentence.
       Gnats, gnus, etc.\ all  % `\ ' makes an inter-word space.
       begin with G\@.         % \@ marks end-of-sentence punctuation.
You should check the spaces after periods when
reading your output to make sure you haven't
forgotten any special cases.  Generating an
ellipsis
       \ldots\               % `\ ' is needed after `\ldots' because TeX
                             % ignores spaces after command names like \ldots
                             % made from \ + letters.
                             %
                             % Note how a `%' character causes TeX to ignore
                             % the end of the input line, so these blank lines
                             % do not start a new paragraph.
                             %
with the right spacing around the periods requires
a special command.

\LaTeX\ interprets some common characters as
commands, so you must type special commands to
generate them.  These characters include the
following:
       \$ \& \% \# \{ and \}.

In printing, text is usually emphasized with an
       \emph{italic}
type style.

\begin{em}
   A long segment of text can also be emphasized
   in this way.  Text within such a segment can be
   given \emph{additional} emphasis.
\end{em}

It is sometimes necessary to prevent \LaTeX\ from
breaking a line where it might otherwise do so.
This may be at a space, as between the ``Mr.''\ and
``Jones'' in
       ``Mr.~Jones'',        % ~ produces an unbreakable interword space.
or within a word---especially when the word is a
symbol like
       \mbox{\emph{itemnum}}
that makes little sense when hyphenated across
lines.

Footnotes\footnote{This is an example of a footnote.}
pose no problem.

\LaTeX\ is good at typesetting mathematical formulas
like
       \( x-3y + z = 7 \)
or
       \( a_{1} > x^{2n} + y^{2n} > x' \)
or
       \( \ip{A}{B} = \sum_{i} a_{i} b_{i} \).
The spaces you type in a formula are
ignored.  Remember that a letter like
       $x$                   % $ ... $  and  \( ... \)  are equivalent
is a formula when it denotes a mathematical
symbol, and it should be typed as one.

\section{Displayed Text}

Text is displayed by indenting it from the left
margin.  Quotations are commonly displayed.  There
are short quotations
\begin{quote}
   This is a short quotation.  It consists of a
   single paragraph of text.  See how it is formatted.
\end{quote}
and longer ones.
\begin{quotation}
   This is a longer quotation.  It consists of two
   paragraphs of text, neither of which are
   particularly interesting.

   This is the second paragraph of the quotation.  It
   is just as dull as the first paragraph.
\end{quotation}
Another frequently-displayed structure is a list.
The following is an example of an \emph{itemized}
list.
\begin{itemize}
   \item This is the first item of an itemized list.
         Each item in the list is marked with a ``tick''.
         You don't have to worry about what kind of tick
         mark is used.

   \item This is the second item of the list.  It
         contains another list nested inside it.  The inner
         list is an \emph{enumerated} list.
         \begin{enumerate}
            \item This is the first item of an enumerated
                  list that is nested within the itemized list.

            \item This is the second item of the inner list.
                  \LaTeX\ allows you to nest lists deeper than
                  you really should.
         \end{enumerate}
         This is the rest of the second item of the outer
         list.  It is no more interesting than any other
         part of the item.
   \item This is the third item of the list.
\end{itemize}
You can even display poetry.
\begin{verse}
   There is an environment
    for verse \\             % The \\ command separates lines
   Whose features some poets % within a stanza.
   will curse.

                             % One or more blank lines separate stanzas.

   For instead of making\\
   Them do \emph{all} line breaking, \\
   It allows them to put too many words on a line when they'd rather be
   forced to be terse.
\end{verse}

Mathematical formulas may also be displayed.  A
displayed formula
is
one-line long; multiline
formulas require special formatting instructions.
   \[  \ip{\Gamma}{\psi'} = x'' + y^{2} + z_{i}^{n}\]
Don't start a paragraph with a displayed equation,
nor make one a paragraph by itself.

\end{document}               % End of document.

\section{Bench Section 54}
% synthetic bench section 54
% This is a sample LaTeX input file.  (Version of 12 August 2004.)
%
% A '%' character causes TeX to ignore all remaining text on the line,
% and is used for comments like this one.

\documentclass{article}      % Specifies the document class

                             % The preamble begins here.
\title{An Example Document}  % Declares the document's title.
\author{Leslie Lamport}      % Declares the author's name.
\date{January 21, 1994}      % Deleting this command produces today's date.

\newcommand{\ip}[2]{(#1, #2)}
                             % Defines \ip{arg1}{arg2} to mean
                             % (arg1, arg2).

%\newcommand{\ip}[2]{\langle #1 | #2\rangle}
                             % This is an alternative definition of
                             % \ip that is commented out.

\begin{document}             % End of preamble and beginning of text.

\maketitle                   % Produces the title.

This is an example input file.  Comparing it with
the output it generates can show you how to
produce a simple document of your own.

\section{Ordinary Text}      % Produces section heading.  Lower-level
                             % sections are begun with similar
                             % \subsection and \subsubsection commands.

The ends  of words and sentences are marked
  by   spaces. It  doesn't matter how many
spaces    you type; one is as good as 100.  The
end of   a line counts as a space.

One   or more   blank lines denote the  end
of  a paragraph.

Since any number of consecutive spaces are treated
like a single one, the formatting of the input
file makes no difference to
      \LaTeX,                % The \LaTeX command generates the LaTeX logo.
but it makes a difference to you.  When you use
\LaTeX, making your input file as easy to read
as possible will be a great help as you write
your document and when you change it.  This sample
file shows how you can add comments to your own input
file.

Because printing is different from typewriting,
there are a number of things that you have to do
differently when preparing an input file than if
you were just typing the document directly.
Quotation marks like
       ``this''
have to be handled specially, as do quotes within
quotes:
       ``\,`this'            % \, separates the double and single quote.
        is what I just
        wrote, not  `that'\,''.

Dashes come in three sizes: an
       intra-word
dash, a medium dash for number ranges like
       1--2,
and a punctuation
       dash---like
this.

A sentence-ending space should be larger than the
space between words within a sentence.  You
sometimes have to type special commands in
conjunction with punctuation characters to get
this right, as in the following sentence.
       Gnats, gnus, etc.\ all  % `\ ' makes an inter-word space.
       begin with G\@.         % \@ marks end-of-sentence punctuation.
You should check the spaces after periods when
reading your output to make sure you haven't
forgotten any special cases.  Generating an
ellipsis
       \ldots\               % `\ ' is needed after `\ldots' because TeX
                             % ignores spaces after command names like \ldots
                             % made from \ + letters.
                             %
                             % Note how a `%' character causes TeX to ignore
                             % the end of the input line, so these blank lines
                             % do not start a new paragraph.
                             %
with the right spacing around the periods requires
a special command.

\LaTeX\ interprets some common characters as
commands, so you must type special commands to
generate them.  These characters include the
following:
       \$ \& \% \# \{ and \}.

In printing, text is usually emphasized with an
       \emph{italic}
type style.

\begin{em}
   A long segment of text can also be emphasized
   in this way.  Text within such a segment can be
   given \emph{additional} emphasis.
\end{em}

It is sometimes necessary to prevent \LaTeX\ from
breaking a line where it might otherwise do so.
This may be at a space, as between the ``Mr.''\ and
``Jones'' in
       ``Mr.~Jones'',        % ~ produces an unbreakable interword space.
or within a word---especially when the word is a
symbol like
       \mbox{\emph{itemnum}}
that makes little sense when hyphenated across
lines.

Footnotes\footnote{This is an example of a footnote.}
pose no problem.

\LaTeX\ is good at typesetting mathematical formulas
like
       \( x-3y + z = 7 \)
or
       \( a_{1} > x^{2n} + y^{2n} > x' \)
or
       \( \ip{A}{B} = \sum_{i} a_{i} b_{i} \).
The spaces you type in a formula are
ignored.  Remember that a letter like
       $x$                   % $ ... $  and  \( ... \)  are equivalent
is a formula when it denotes a mathematical
symbol, and it should be typed as one.

\section{Displayed Text}

Text is displayed by indenting it from the left
margin.  Quotations are commonly displayed.  There
are short quotations
\begin{quote}
   This is a short quotation.  It consists of a
   single paragraph of text.  See how it is formatted.
\end{quote}
and longer ones.
\begin{quotation}
   This is a longer quotation.  It consists of two
   paragraphs of text, neither of which are
   particularly interesting.

   This is the second paragraph of the quotation.  It
   is just as dull as the first paragraph.
\end{quotation}
Another frequently-displayed structure is a list.
The following is an example of an \emph{itemized}
list.
\begin{itemize}
   \item This is the first item of an itemized list.
         Each item in the list is marked with a ``tick''.
         You don't have to worry about what kind of tick
         mark is used.

   \item This is the second item of the list.  It
         contains another list nested inside it.  The inner
         list is an \emph{enumerated} list.
         \begin{enumerate}
            \item This is the first item of an enumerated
                  list that is nested within the itemized list.

            \item This is the second item of the inner list.
                  \LaTeX\ allows you to nest lists deeper than
                  you really should.
         \end{enumerate}
         This is the rest of the second item of the outer
         list.  It is no more interesting than any other
         part of the item.
   \item This is the third item of the list.
\end{itemize}
You can even display poetry.
\begin{verse}
   There is an environment
    for verse \\             % The \\ command separates lines
   Whose features some poets % within a stanza.
   will curse.

                             % One or more blank lines separate stanzas.

   For instead of making\\
   Them do \emph{all} line breaking, \\
   It allows them to put too many words on a line when they'd rather be
   forced to be terse.
\end{verse}

Mathematical formulas may also be displayed.  A
displayed formula
is
one-line long; multiline
formulas require special formatting instructions.
   \[  \ip{\Gamma}{\psi'} = x'' + y^{2} + z_{i}^{n}\]
Don't start a paragraph with a displayed equation,
nor make one a paragraph by itself.

\end{document}               % End of document.

\section{Bench Section 55}
% synthetic bench section 55
% This is a sample LaTeX input file.  (Version of 12 August 2004.)
%
% A '%' character causes TeX to ignore all remaining text on the line,
% and is used for comments like this one.

\documentclass{article}      % Specifies the document class

                             % The preamble begins here.
\title{An Example Document}  % Declares the document's title.
\author{Leslie Lamport}      % Declares the author's name.
\date{January 21, 1994}      % Deleting this command produces today's date.

\newcommand{\ip}[2]{(#1, #2)}
                             % Defines \ip{arg1}{arg2} to mean
                             % (arg1, arg2).

%\newcommand{\ip}[2]{\langle #1 | #2\rangle}
                             % This is an alternative definition of
                             % \ip that is commented out.

\begin{document}             % End of preamble and beginning of text.

\maketitle                   % Produces the title.

This is an example input file.  Comparing it with
the output it generates can show you how to
produce a simple document of your own.

\section{Ordinary Text}      % Produces section heading.  Lower-level
                             % sections are begun with similar
                             % \subsection and \subsubsection commands.

The ends  of words and sentences are marked
  by   spaces. It  doesn't matter how many
spaces    you type; one is as good as 100.  The
end of   a line counts as a space.

One   or more   blank lines denote the  end
of  a paragraph.

Since any number of consecutive spaces are treated
like a single one, the formatting of the input
file makes no difference to
      \LaTeX,                % The \LaTeX command generates the LaTeX logo.
but it makes a difference to you.  When you use
\LaTeX, making your input file as easy to read
as possible will be a great help as you write
your document and when you change it.  This sample
file shows how you can add comments to your own input
file.

Because printing is different from typewriting,
there are a number of things that you have to do
differently when preparing an input file than if
you were just typing the document directly.
Quotation marks like
       ``this''
have to be handled specially, as do quotes within
quotes:
       ``\,`this'            % \, separates the double and single quote.
        is what I just
        wrote, not  `that'\,''.

Dashes come in three sizes: an
       intra-word
dash, a medium dash for number ranges like
       1--2,
and a punctuation
       dash---like
this.

A sentence-ending space should be larger than the
space between words within a sentence.  You
sometimes have to type special commands in
conjunction with punctuation characters to get
this right, as in the following sentence.
       Gnats, gnus, etc.\ all  % `\ ' makes an inter-word space.
       begin with G\@.         % \@ marks end-of-sentence punctuation.
You should check the spaces after periods when
reading your output to make sure you haven't
forgotten any special cases.  Generating an
ellipsis
       \ldots\               % `\ ' is needed after `\ldots' because TeX
                             % ignores spaces after command names like \ldots
                             % made from \ + letters.
                             %
                             % Note how a `%' character causes TeX to ignore
                             % the end of the input line, so these blank lines
                             % do not start a new paragraph.
                             %
with the right spacing around the periods requires
a special command.

\LaTeX\ interprets some common characters as
commands, so you must type special commands to
generate them.  These characters include the
following:
       \$ \& \% \# \{ and \}.

In printing, text is usually emphasized with an
       \emph{italic}
type style.

\begin{em}
   A long segment of text can also be emphasized
   in this way.  Text within such a segment can be
   given \emph{additional} emphasis.
\end{em}

It is sometimes necessary to prevent \LaTeX\ from
breaking a line where it might otherwise do so.
This may be at a space, as between the ``Mr.''\ and
``Jones'' in
       ``Mr.~Jones'',        % ~ produces an unbreakable interword space.
or within a word---especially when the word is a
symbol like
       \mbox{\emph{itemnum}}
that makes little sense when hyphenated across
lines.

Footnotes\footnote{This is an example of a footnote.}
pose no problem.

\LaTeX\ is good at typesetting mathematical formulas
like
       \( x-3y + z = 7 \)
or
       \( a_{1} > x^{2n} + y^{2n} > x' \)
or
       \( \ip{A}{B} = \sum_{i} a_{i} b_{i} \).
The spaces you type in a formula are
ignored.  Remember that a letter like
       $x$                   % $ ... $  and  \( ... \)  are equivalent
is a formula when it denotes a mathematical
symbol, and it should be typed as one.

\section{Displayed Text}

Text is displayed by indenting it from the left
margin.  Quotations are commonly displayed.  There
are short quotations
\begin{quote}
   This is a short quotation.  It consists of a
   single paragraph of text.  See how it is formatted.
\end{quote}
and longer ones.
\begin{quotation}
   This is a longer quotation.  It consists of two
   paragraphs of text, neither of which are
   particularly interesting.

   This is the second paragraph of the quotation.  It
   is just as dull as the first paragraph.
\end{quotation}
Another frequently-displayed structure is a list.
The following is an example of an \emph{itemized}
list.
\begin{itemize}
   \item This is the first item of an itemized list.
         Each item in the list is marked with a ``tick''.
         You don't have to worry about what kind of tick
         mark is used.

   \item This is the second item of the list.  It
         contains another list nested inside it.  The inner
         list is an \emph{enumerated} list.
         \begin{enumerate}
            \item This is the first item of an enumerated
                  list that is nested within the itemized list.

            \item This is the second item of the inner list.
                  \LaTeX\ allows you to nest lists deeper than
                  you really should.
         \end{enumerate}
         This is the rest of the second item of the outer
         list.  It is no more interesting than any other
         part of the item.
   \item This is the third item of the list.
\end{itemize}
You can even display poetry.
\begin{verse}
   There is an environment
    for verse \\             % The \\ command separates lines
   Whose features some poets % within a stanza.
   will curse.

                             % One or more blank lines separate stanzas.

   For instead of making\\
   Them do \emph{all} line breaking, \\
   It allows them to put too many words on a line when they'd rather be
   forced to be terse.
\end{verse}

Mathematical formulas may also be displayed.  A
displayed formula
is
one-line long; multiline
formulas require special formatting instructions.
   \[  \ip{\Gamma}{\psi'} = x'' + y^{2} + z_{i}^{n}\]
Don't start a paragraph with a displayed equation,
nor make one a paragraph by itself.

\end{document}               % End of document.

\section{Bench Section 56}
% synthetic bench section 56
% This is a sample LaTeX input file.  (Version of 12 August 2004.)
%
% A '%' character causes TeX to ignore all remaining text on the line,
% and is used for comments like this one.

\documentclass{article}      % Specifies the document class

                             % The preamble begins here.
\title{An Example Document}  % Declares the document's title.
\author{Leslie Lamport}      % Declares the author's name.
\date{January 21, 1994}      % Deleting this command produces today's date.

\newcommand{\ip}[2]{(#1, #2)}
                             % Defines \ip{arg1}{arg2} to mean
                             % (arg1, arg2).

%\newcommand{\ip}[2]{\langle #1 | #2\rangle}
                             % This is an alternative definition of
                             % \ip that is commented out.

\begin{document}             % End of preamble and beginning of text.

\maketitle                   % Produces the title.

This is an example input file.  Comparing it with
the output it generates can show you how to
produce a simple document of your own.

\section{Ordinary Text}      % Produces section heading.  Lower-level
                             % sections are begun with similar
                             % \subsection and \subsubsection commands.

The ends  of words and sentences are marked
  by   spaces. It  doesn't matter how many
spaces    you type; one is as good as 100.  The
end of   a line counts as a space.

One   or more   blank lines denote the  end
of  a paragraph.

Since any number of consecutive spaces are treated
like a single one, the formatting of the input
file makes no difference to
      \LaTeX,                % The \LaTeX command generates the LaTeX logo.
but it makes a difference to you.  When you use
\LaTeX, making your input file as easy to read
as possible will be a great help as you write
your document and when you change it.  This sample
file shows how you can add comments to your own input
file.

Because printing is different from typewriting,
there are a number of things that you have to do
differently when preparing an input file than if
you were just typing the document directly.
Quotation marks like
       ``this''
have to be handled specially, as do quotes within
quotes:
       ``\,`this'            % \, separates the double and single quote.
        is what I just
        wrote, not  `that'\,''.

Dashes come in three sizes: an
       intra-word
dash, a medium dash for number ranges like
       1--2,
and a punctuation
       dash---like
this.

A sentence-ending space should be larger than the
space between words within a sentence.  You
sometimes have to type special commands in
conjunction with punctuation characters to get
this right, as in the following sentence.
       Gnats, gnus, etc.\ all  % `\ ' makes an inter-word space.
       begin with G\@.         % \@ marks end-of-sentence punctuation.
You should check the spaces after periods when
reading your output to make sure you haven't
forgotten any special cases.  Generating an
ellipsis
       \ldots\               % `\ ' is needed after `\ldots' because TeX
                             % ignores spaces after command names like \ldots
                             % made from \ + letters.
                             %
                             % Note how a `%' character causes TeX to ignore
                             % the end of the input line, so these blank lines
                             % do not start a new paragraph.
                             %
with the right spacing around the periods requires
a special command.

\LaTeX\ interprets some common characters as
commands, so you must type special commands to
generate them.  These characters include the
following:
       \$ \& \% \# \{ and \}.

In printing, text is usually emphasized with an
       \emph{italic}
type style.

\begin{em}
   A long segment of text can also be emphasized
   in this way.  Text within such a segment can be
   given \emph{additional} emphasis.
\end{em}

It is sometimes necessary to prevent \LaTeX\ from
breaking a line where it might otherwise do so.
This may be at a space, as between the ``Mr.''\ and
``Jones'' in
       ``Mr.~Jones'',        % ~ produces an unbreakable interword space.
or within a word---especially when the word is a
symbol like
       \mbox{\emph{itemnum}}
that makes little sense when hyphenated across
lines.

Footnotes\footnote{This is an example of a footnote.}
pose no problem.

\LaTeX\ is good at typesetting mathematical formulas
like
       \( x-3y + z = 7 \)
or
       \( a_{1} > x^{2n} + y^{2n} > x' \)
or
       \( \ip{A}{B} = \sum_{i} a_{i} b_{i} \).
The spaces you type in a formula are
ignored.  Remember that a letter like
       $x$                   % $ ... $  and  \( ... \)  are equivalent
is a formula when it denotes a mathematical
symbol, and it should be typed as one.

\section{Displayed Text}

Text is displayed by indenting it from the left
margin.  Quotations are commonly displayed.  There
are short quotations
\begin{quote}
   This is a short quotation.  It consists of a
   single paragraph of text.  See how it is formatted.
\end{quote}
and longer ones.
\begin{quotation}
   This is a longer quotation.  It consists of two
   paragraphs of text, neither of which are
   particularly interesting.

   This is the second paragraph of the quotation.  It
   is just as dull as the first paragraph.
\end{quotation}
Another frequently-displayed structure is a list.
The following is an example of an \emph{itemized}
list.
\begin{itemize}
   \item This is the first item of an itemized list.
         Each item in the list is marked with a ``tick''.
         You don't have to worry about what kind of tick
         mark is used.

   \item This is the second item of the list.  It
         contains another list nested inside it.  The inner
         list is an \emph{enumerated} list.
         \begin{enumerate}
            \item This is the first item of an enumerated
                  list that is nested within the itemized list.

            \item This is the second item of the inner list.
                  \LaTeX\ allows you to nest lists deeper than
                  you really should.
         \end{enumerate}
         This is the rest of the second item of the outer
         list.  It is no more interesting than any other
         part of the item.
   \item This is the third item of the list.
\end{itemize}
You can even display poetry.
\begin{verse}
   There is an environment
    for verse \\             % The \\ command separates lines
   Whose features some poets % within a stanza.
   will curse.

                             % One or more blank lines separate stanzas.

   For instead of making\\
   Them do \emph{all} line breaking, \\
   It allows them to put too many words on a line when they'd rather be
   forced to be terse.
\end{verse}

Mathematical formulas may also be displayed.  A
displayed formula
is
one-line long; multiline
formulas require special formatting instructions.
   \[  \ip{\Gamma}{\psi'} = x'' + y^{2} + z_{i}^{n}\]
Don't start a paragraph with a displayed equation,
nor make one a paragraph by itself.

\end{document}               % End of document.

\section{Bench Section 57}
% synthetic bench section 57
% This is a sample LaTeX input file.  (Version of 12 August 2004.)
%
% A '%' character causes TeX to ignore all remaining text on the line,
% and is used for comments like this one.

\documentclass{article}      % Specifies the document class

                             % The preamble begins here.
\title{An Example Document}  % Declares the document's title.
\author{Leslie Lamport}      % Declares the author's name.
\date{January 21, 1994}      % Deleting this command produces today's date.

\newcommand{\ip}[2]{(#1, #2)}
                             % Defines \ip{arg1}{arg2} to mean
                             % (arg1, arg2).

%\newcommand{\ip}[2]{\langle #1 | #2\rangle}
                             % This is an alternative definition of
                             % \ip that is commented out.

\begin{document}             % End of preamble and beginning of text.

\maketitle                   % Produces the title.

This is an example input file.  Comparing it with
the output it generates can show you how to
produce a simple document of your own.

\section{Ordinary Text}      % Produces section heading.  Lower-level
                             % sections are begun with similar
                             % \subsection and \subsubsection commands.

The ends  of words and sentences are marked
  by   spaces. It  doesn't matter how many
spaces    you type; one is as good as 100.  The
end of   a line counts as a space.

One   or more   blank lines denote the  end
of  a paragraph.

Since any number of consecutive spaces are treated
like a single one, the formatting of the input
file makes no difference to
      \LaTeX,                % The \LaTeX command generates the LaTeX logo.
but it makes a difference to you.  When you use
\LaTeX, making your input file as easy to read
as possible will be a great help as you write
your document and when you change it.  This sample
file shows how you can add comments to your own input
file.

Because printing is different from typewriting,
there are a number of things that you have to do
differently when preparing an input file than if
you were just typing the document directly.
Quotation marks like
       ``this''
have to be handled specially, as do quotes within
quotes:
       ``\,`this'            % \, separates the double and single quote.
        is what I just
        wrote, not  `that'\,''.

Dashes come in three sizes: an
       intra-word
dash, a medium dash for number ranges like
       1--2,
and a punctuation
       dash---like
this.

A sentence-ending space should be larger than the
space between words within a sentence.  You
sometimes have to type special commands in
conjunction with punctuation characters to get
this right, as in the following sentence.
       Gnats, gnus, etc.\ all  % `\ ' makes an inter-word space.
       begin with G\@.         % \@ marks end-of-sentence punctuation.
You should check the spaces after periods when
reading your output to make sure you haven't
forgotten any special cases.  Generating an
ellipsis
       \ldots\               % `\ ' is needed after `\ldots' because TeX
                             % ignores spaces after command names like \ldots
                             % made from \ + letters.
                             %
                             % Note how a `%' character causes TeX to ignore
                             % the end of the input line, so these blank lines
                             % do not start a new paragraph.
                             %
with the right spacing around the periods requires
a special command.

\LaTeX\ interprets some common characters as
commands, so you must type special commands to
generate them.  These characters include the
following:
       \$ \& \% \# \{ and \}.

In printing, text is usually emphasized with an
       \emph{italic}
type style.

\begin{em}
   A long segment of text can also be emphasized
   in this way.  Text within such a segment can be
   given \emph{additional} emphasis.
\end{em}

It is sometimes necessary to prevent \LaTeX\ from
breaking a line where it might otherwise do so.
This may be at a space, as between the ``Mr.''\ and
``Jones'' in
       ``Mr.~Jones'',        % ~ produces an unbreakable interword space.
or within a word---especially when the word is a
symbol like
       \mbox{\emph{itemnum}}
that makes little sense when hyphenated across
lines.

Footnotes\footnote{This is an example of a footnote.}
pose no problem.

\LaTeX\ is good at typesetting mathematical formulas
like
       \( x-3y + z = 7 \)
or
       \( a_{1} > x^{2n} + y^{2n} > x' \)
or
       \( \ip{A}{B} = \sum_{i} a_{i} b_{i} \).
The spaces you type in a formula are
ignored.  Remember that a letter like
       $x$                   % $ ... $  and  \( ... \)  are equivalent
is a formula when it denotes a mathematical
symbol, and it should be typed as one.

\section{Displayed Text}

Text is displayed by indenting it from the left
margin.  Quotations are commonly displayed.  There
are short quotations
\begin{quote}
   This is a short quotation.  It consists of a
   single paragraph of text.  See how it is formatted.
\end{quote}
and longer ones.
\begin{quotation}
   This is a longer quotation.  It consists of two
   paragraphs of text, neither of which are
   particularly interesting.

   This is the second paragraph of the quotation.  It
   is just as dull as the first paragraph.
\end{quotation}
Another frequently-displayed structure is a list.
The following is an example of an \emph{itemized}
list.
\begin{itemize}
   \item This is the first item of an itemized list.
         Each item in the list is marked with a ``tick''.
         You don't have to worry about what kind of tick
         mark is used.

   \item This is the second item of the list.  It
         contains another list nested inside it.  The inner
         list is an \emph{enumerated} list.
         \begin{enumerate}
            \item This is the first item of an enumerated
                  list that is nested within the itemized list.

            \item This is the second item of the inner list.
                  \LaTeX\ allows you to nest lists deeper than
                  you really should.
         \end{enumerate}
         This is the rest of the second item of the outer
         list.  It is no more interesting than any other
         part of the item.
   \item This is the third item of the list.
\end{itemize}
You can even display poetry.
\begin{verse}
   There is an environment
    for verse \\             % The \\ command separates lines
   Whose features some poets % within a stanza.
   will curse.

                             % One or more blank lines separate stanzas.

   For instead of making\\
   Them do \emph{all} line breaking, \\
   It allows them to put too many words on a line when they'd rather be
   forced to be terse.
\end{verse}

Mathematical formulas may also be displayed.  A
displayed formula
is
one-line long; multiline
formulas require special formatting instructions.
   \[  \ip{\Gamma}{\psi'} = x'' + y^{2} + z_{i}^{n}\]
Don't start a paragraph with a displayed equation,
nor make one a paragraph by itself.

\end{document}               % End of document.

\section{Bench Section 58}
% synthetic bench section 58
% This is a sample LaTeX input file.  (Version of 12 August 2004.)
%
% A '%' character causes TeX to ignore all remaining text on the line,
% and is used for comments like this one.

\documentclass{article}      % Specifies the document class

                             % The preamble begins here.
\title{An Example Document}  % Declares the document's title.
\author{Leslie Lamport}      % Declares the author's name.
\date{January 21, 1994}      % Deleting this command produces today's date.

\newcommand{\ip}[2]{(#1, #2)}
                             % Defines \ip{arg1}{arg2} to mean
                             % (arg1, arg2).

%\newcommand{\ip}[2]{\langle #1 | #2\rangle}
                             % This is an alternative definition of
                             % \ip that is commented out.

\begin{document}             % End of preamble and beginning of text.

\maketitle                   % Produces the title.

This is an example input file.  Comparing it with
the output it generates can show you how to
produce a simple document of your own.

\section{Ordinary Text}      % Produces section heading.  Lower-level
                             % sections are begun with similar
                             % \subsection and \subsubsection commands.

The ends  of words and sentences are marked
  by   spaces. It  doesn't matter how many
spaces    you type; one is as good as 100.  The
end of   a line counts as a space.

One   or more   blank lines denote the  end
of  a paragraph.

Since any number of consecutive spaces are treated
like a single one, the formatting of the input
file makes no difference to
      \LaTeX,                % The \LaTeX command generates the LaTeX logo.
but it makes a difference to you.  When you use
\LaTeX, making your input file as easy to read
as possible will be a great help as you write
your document and when you change it.  This sample
file shows how you can add comments to your own input
file.

Because printing is different from typewriting,
there are a number of things that you have to do
differently when preparing an input file than if
you were just typing the document directly.
Quotation marks like
       ``this''
have to be handled specially, as do quotes within
quotes:
       ``\,`this'            % \, separates the double and single quote.
        is what I just
        wrote, not  `that'\,''.

Dashes come in three sizes: an
       intra-word
dash, a medium dash for number ranges like
       1--2,
and a punctuation
       dash---like
this.

A sentence-ending space should be larger than the
space between words within a sentence.  You
sometimes have to type special commands in
conjunction with punctuation characters to get
this right, as in the following sentence.
       Gnats, gnus, etc.\ all  % `\ ' makes an inter-word space.
       begin with G\@.         % \@ marks end-of-sentence punctuation.
You should check the spaces after periods when
reading your output to make sure you haven't
forgotten any special cases.  Generating an
ellipsis
       \ldots\               % `\ ' is needed after `\ldots' because TeX
                             % ignores spaces after command names like \ldots
                             % made from \ + letters.
                             %
                             % Note how a `%' character causes TeX to ignore
                             % the end of the input line, so these blank lines
                             % do not start a new paragraph.
                             %
with the right spacing around the periods requires
a special command.

\LaTeX\ interprets some common characters as
commands, so you must type special commands to
generate them.  These characters include the
following:
       \$ \& \% \# \{ and \}.

In printing, text is usually emphasized with an
       \emph{italic}
type style.

\begin{em}
   A long segment of text can also be emphasized
   in this way.  Text within such a segment can be
   given \emph{additional} emphasis.
\end{em}

It is sometimes necessary to prevent \LaTeX\ from
breaking a line where it might otherwise do so.
This may be at a space, as between the ``Mr.''\ and
``Jones'' in
       ``Mr.~Jones'',        % ~ produces an unbreakable interword space.
or within a word---especially when the word is a
symbol like
       \mbox{\emph{itemnum}}
that makes little sense when hyphenated across
lines.

Footnotes\footnote{This is an example of a footnote.}
pose no problem.

\LaTeX\ is good at typesetting mathematical formulas
like
       \( x-3y + z = 7 \)
or
       \( a_{1} > x^{2n} + y^{2n} > x' \)
or
       \( \ip{A}{B} = \sum_{i} a_{i} b_{i} \).
The spaces you type in a formula are
ignored.  Remember that a letter like
       $x$                   % $ ... $  and  \( ... \)  are equivalent
is a formula when it denotes a mathematical
symbol, and it should be typed as one.

\section{Displayed Text}

Text is displayed by indenting it from the left
margin.  Quotations are commonly displayed.  There
are short quotations
\begin{quote}
   This is a short quotation.  It consists of a
   single paragraph of text.  See how it is formatted.
\end{quote}
and longer ones.
\begin{quotation}
   This is a longer quotation.  It consists of two
   paragraphs of text, neither of which are
   particularly interesting.

   This is the second paragraph of the quotation.  It
   is just as dull as the first paragraph.
\end{quotation}
Another frequently-displayed structure is a list.
The following is an example of an \emph{itemized}
list.
\begin{itemize}
   \item This is the first item of an itemized list.
         Each item in the list is marked with a ``tick''.
         You don't have to worry about what kind of tick
         mark is used.

   \item This is the second item of the list.  It
         contains another list nested inside it.  The inner
         list is an \emph{enumerated} list.
         \begin{enumerate}
            \item This is the first item of an enumerated
                  list that is nested within the itemized list.

            \item This is the second item of the inner list.
                  \LaTeX\ allows you to nest lists deeper than
                  you really should.
         \end{enumerate}
         This is the rest of the second item of the outer
         list.  It is no more interesting than any other
         part of the item.
   \item This is the third item of the list.
\end{itemize}
You can even display poetry.
\begin{verse}
   There is an environment
    for verse \\             % The \\ command separates lines
   Whose features some poets % within a stanza.
   will curse.

                             % One or more blank lines separate stanzas.

   For instead of making\\
   Them do \emph{all} line breaking, \\
   It allows them to put too many words on a line when they'd rather be
   forced to be terse.
\end{verse}

Mathematical formulas may also be displayed.  A
displayed formula
is
one-line long; multiline
formulas require special formatting instructions.
   \[  \ip{\Gamma}{\psi'} = x'' + y^{2} + z_{i}^{n}\]
Don't start a paragraph with a displayed equation,
nor make one a paragraph by itself.

\end{document}               % End of document.

\section{Bench Section 59}
% synthetic bench section 59
% This is a sample LaTeX input file.  (Version of 12 August 2004.)
%
% A '%' character causes TeX to ignore all remaining text on the line,
% and is used for comments like this one.

\documentclass{article}      % Specifies the document class

                             % The preamble begins here.
\title{An Example Document}  % Declares the document's title.
\author{Leslie Lamport}      % Declares the author's name.
\date{January 21, 1994}      % Deleting this command produces today's date.

\newcommand{\ip}[2]{(#1, #2)}
                             % Defines \ip{arg1}{arg2} to mean
                             % (arg1, arg2).

%\newcommand{\ip}[2]{\langle #1 | #2\rangle}
                             % This is an alternative definition of
                             % \ip that is commented out.

\begin{document}             % End of preamble and beginning of text.

\maketitle                   % Produces the title.

This is an example input file.  Comparing it with
the output it generates can show you how to
produce a simple document of your own.

\section{Ordinary Text}      % Produces section heading.  Lower-level
                             % sections are begun with similar
                             % \subsection and \subsubsection commands.

The ends  of words and sentences are marked
  by   spaces. It  doesn't matter how many
spaces    you type; one is as good as 100.  The
end of   a line counts as a space.

One   or more   blank lines denote the  end
of  a paragraph.

Since any number of consecutive spaces are treated
like a single one, the formatting of the input
file makes no difference to
      \LaTeX,                % The \LaTeX command generates the LaTeX logo.
but it makes a difference to you.  When you use
\LaTeX, making your input file as easy to read
as possible will be a great help as you write
your document and when you change it.  This sample
file shows how you can add comments to your own input
file.

Because printing is different from typewriting,
there are a number of things that you have to do
differently when preparing an input file than if
you were just typing the document directly.
Quotation marks like
       ``this''
have to be handled specially, as do quotes within
quotes:
       ``\,`this'            % \, separates the double and single quote.
        is what I just
        wrote, not  `that'\,''.

Dashes come in three sizes: an
       intra-word
dash, a medium dash for number ranges like
       1--2,
and a punctuation
       dash---like
this.

A sentence-ending space should be larger than the
space between words within a sentence.  You
sometimes have to type special commands in
conjunction with punctuation characters to get
this right, as in the following sentence.
       Gnats, gnus, etc.\ all  % `\ ' makes an inter-word space.
       begin with G\@.         % \@ marks end-of-sentence punctuation.
You should check the spaces after periods when
reading your output to make sure you haven't
forgotten any special cases.  Generating an
ellipsis
       \ldots\               % `\ ' is needed after `\ldots' because TeX
                             % ignores spaces after command names like \ldots
                             % made from \ + letters.
                             %
                             % Note how a `%' character causes TeX to ignore
                             % the end of the input line, so these blank lines
                             % do not start a new paragraph.
                             %
with the right spacing around the periods requires
a special command.

\LaTeX\ interprets some common characters as
commands, so you must type special commands to
generate them.  These characters include the
following:
       \$ \& \% \# \{ and \}.

In printing, text is usually emphasized with an
       \emph{italic}
type style.

\begin{em}
   A long segment of text can also be emphasized
   in this way.  Text within such a segment can be
   given \emph{additional} emphasis.
\end{em}

It is sometimes necessary to prevent \LaTeX\ from
breaking a line where it might otherwise do so.
This may be at a space, as between the ``Mr.''\ and
``Jones'' in
       ``Mr.~Jones'',        % ~ produces an unbreakable interword space.
or within a word---especially when the word is a
symbol like
       \mbox{\emph{itemnum}}
that makes little sense when hyphenated across
lines.

Footnotes\footnote{This is an example of a footnote.}
pose no problem.

\LaTeX\ is good at typesetting mathematical formulas
like
       \( x-3y + z = 7 \)
or
       \( a_{1} > x^{2n} + y^{2n} > x' \)
or
       \( \ip{A}{B} = \sum_{i} a_{i} b_{i} \).
The spaces you type in a formula are
ignored.  Remember that a letter like
       $x$                   % $ ... $  and  \( ... \)  are equivalent
is a formula when it denotes a mathematical
symbol, and it should be typed as one.

\section{Displayed Text}

Text is displayed by indenting it from the left
margin.  Quotations are commonly displayed.  There
are short quotations
\begin{quote}
   This is a short quotation.  It consists of a
   single paragraph of text.  See how it is formatted.
\end{quote}
and longer ones.
\begin{quotation}
   This is a longer quotation.  It consists of two
   paragraphs of text, neither of which are
   particularly interesting.

   This is the second paragraph of the quotation.  It
   is just as dull as the first paragraph.
\end{quotation}
Another frequently-displayed structure is a list.
The following is an example of an \emph{itemized}
list.
\begin{itemize}
   \item This is the first item of an itemized list.
         Each item in the list is marked with a ``tick''.
         You don't have to worry about what kind of tick
         mark is used.

   \item This is the second item of the list.  It
         contains another list nested inside it.  The inner
         list is an \emph{enumerated} list.
         \begin{enumerate}
            \item This is the first item of an enumerated
                  list that is nested within the itemized list.

            \item This is the second item of the inner list.
                  \LaTeX\ allows you to nest lists deeper than
                  you really should.
         \end{enumerate}
         This is the rest of the second item of the outer
         list.  It is no more interesting than any other
         part of the item.
   \item This is the third item of the list.
\end{itemize}
You can even display poetry.
\begin{verse}
   There is an environment
    for verse \\             % The \\ command separates lines
   Whose features some poets % within a stanza.
   will curse.

                             % One or more blank lines separate stanzas.

   For instead of making\\
   Them do \emph{all} line breaking, \\
   It allows them to put too many words on a line when they'd rather be
   forced to be terse.
\end{verse}

Mathematical formulas may also be displayed.  A
displayed formula
is
one-line long; multiline
formulas require special formatting instructions.
   \[  \ip{\Gamma}{\psi'} = x'' + y^{2} + z_{i}^{n}\]
Don't start a paragraph with a displayed equation,
nor make one a paragraph by itself.

\end{document}               % End of document.

\section{Bench Section 60}
% synthetic bench section 60
% This is a sample LaTeX input file.  (Version of 12 August 2004.)
%
% A '%' character causes TeX to ignore all remaining text on the line,
% and is used for comments like this one.

\documentclass{article}      % Specifies the document class

                             % The preamble begins here.
\title{An Example Document}  % Declares the document's title.
\author{Leslie Lamport}      % Declares the author's name.
\date{January 21, 1994}      % Deleting this command produces today's date.

\newcommand{\ip}[2]{(#1, #2)}
                             % Defines \ip{arg1}{arg2} to mean
                             % (arg1, arg2).

%\newcommand{\ip}[2]{\langle #1 | #2\rangle}
                             % This is an alternative definition of
                             % \ip that is commented out.

\begin{document}             % End of preamble and beginning of text.

\maketitle                   % Produces the title.

This is an example input file.  Comparing it with
the output it generates can show you how to
produce a simple document of your own.

\section{Ordinary Text}      % Produces section heading.  Lower-level
                             % sections are begun with similar
                             % \subsection and \subsubsection commands.

The ends  of words and sentences are marked
  by   spaces. It  doesn't matter how many
spaces    you type; one is as good as 100.  The
end of   a line counts as a space.

One   or more   blank lines denote the  end
of  a paragraph.

Since any number of consecutive spaces are treated
like a single one, the formatting of the input
file makes no difference to
      \LaTeX,                % The \LaTeX command generates the LaTeX logo.
but it makes a difference to you.  When you use
\LaTeX, making your input file as easy to read
as possible will be a great help as you write
your document and when you change it.  This sample
file shows how you can add comments to your own input
file.

Because printing is different from typewriting,
there are a number of things that you have to do
differently when preparing an input file than if
you were just typing the document directly.
Quotation marks like
       ``this''
have to be handled specially, as do quotes within
quotes:
       ``\,`this'            % \, separates the double and single quote.
        is what I just
        wrote, not  `that'\,''.

Dashes come in three sizes: an
       intra-word
dash, a medium dash for number ranges like
       1--2,
and a punctuation
       dash---like
this.

A sentence-ending space should be larger than the
space between words within a sentence.  You
sometimes have to type special commands in
conjunction with punctuation characters to get
this right, as in the following sentence.
       Gnats, gnus, etc.\ all  % `\ ' makes an inter-word space.
       begin with G\@.         % \@ marks end-of-sentence punctuation.
You should check the spaces after periods when
reading your output to make sure you haven't
forgotten any special cases.  Generating an
ellipsis
       \ldots\               % `\ ' is needed after `\ldots' because TeX
                             % ignores spaces after command names like \ldots
                             % made from \ + letters.
                             %
                             % Note how a `%' character causes TeX to ignore
                             % the end of the input line, so these blank lines
                             % do not start a new paragraph.
                             %
with the right spacing around the periods requires
a special command.

\LaTeX\ interprets some common characters as
commands, so you must type special commands to
generate them.  These characters include the
following:
       \$ \& \% \# \{ and \}.

In printing, text is usually emphasized with an
       \emph{italic}
type style.

\begin{em}
   A long segment of text can also be emphasized
   in this way.  Text within such a segment can be
   given \emph{additional} emphasis.
\end{em}

It is sometimes necessary to prevent \LaTeX\ from
breaking a line where it might otherwise do so.
This may be at a space, as between the ``Mr.''\ and
``Jones'' in
       ``Mr.~Jones'',        % ~ produces an unbreakable interword space.
or within a word---especially when the word is a
symbol like
       \mbox{\emph{itemnum}}
that makes little sense when hyphenated across
lines.

Footnotes\footnote{This is an example of a footnote.}
pose no problem.

\LaTeX\ is good at typesetting mathematical formulas
like
       \( x-3y + z = 7 \)
or
       \( a_{1} > x^{2n} + y^{2n} > x' \)
or
       \( \ip{A}{B} = \sum_{i} a_{i} b_{i} \).
The spaces you type in a formula are
ignored.  Remember that a letter like
       $x$                   % $ ... $  and  \( ... \)  are equivalent
is a formula when it denotes a mathematical
symbol, and it should be typed as one.

\section{Displayed Text}

Text is displayed by indenting it from the left
margin.  Quotations are commonly displayed.  There
are short quotations
\begin{quote}
   This is a short quotation.  It consists of a
   single paragraph of text.  See how it is formatted.
\end{quote}
and longer ones.
\begin{quotation}
   This is a longer quotation.  It consists of two
   paragraphs of text, neither of which are
   particularly interesting.

   This is the second paragraph of the quotation.  It
   is just as dull as the first paragraph.
\end{quotation}
Another frequently-displayed structure is a list.
The following is an example of an \emph{itemized}
list.
\begin{itemize}
   \item This is the first item of an itemized list.
         Each item in the list is marked with a ``tick''.
         You don't have to worry about what kind of tick
         mark is used.

   \item This is the second item of the list.  It
         contains another list nested inside it.  The inner
         list is an \emph{enumerated} list.
         \begin{enumerate}
            \item This is the first item of an enumerated
                  list that is nested within the itemized list.

            \item This is the second item of the inner list.
                  \LaTeX\ allows you to nest lists deeper than
                  you really should.
         \end{enumerate}
         This is the rest of the second item of the outer
         list.  It is no more interesting than any other
         part of the item.
   \item This is the third item of the list.
\end{itemize}
You can even display poetry.
\begin{verse}
   There is an environment
    for verse \\             % The \\ command separates lines
   Whose features some poets % within a stanza.
   will curse.

                             % One or more blank lines separate stanzas.

   For instead of making\\
   Them do \emph{all} line breaking, \\
   It allows them to put too many words on a line when they'd rather be
   forced to be terse.
\end{verse}

Mathematical formulas may also be displayed.  A
displayed formula
is
one-line long; multiline
formulas require special formatting instructions.
   \[  \ip{\Gamma}{\psi'} = x'' + y^{2} + z_{i}^{n}\]
Don't start a paragraph with a displayed equation,
nor make one a paragraph by itself.

\end{document}               % End of document.

\section{Bench Section 61}
% synthetic bench section 61
% This is a sample LaTeX input file.  (Version of 12 August 2004.)
%
% A '%' character causes TeX to ignore all remaining text on the line,
% and is used for comments like this one.

\documentclass{article}      % Specifies the document class

                             % The preamble begins here.
\title{An Example Document}  % Declares the document's title.
\author{Leslie Lamport}      % Declares the author's name.
\date{January 21, 1994}      % Deleting this command produces today's date.

\newcommand{\ip}[2]{(#1, #2)}
                             % Defines \ip{arg1}{arg2} to mean
                             % (arg1, arg2).

%\newcommand{\ip}[2]{\langle #1 | #2\rangle}
                             % This is an alternative definition of
                             % \ip that is commented out.

\begin{document}             % End of preamble and beginning of text.

\maketitle                   % Produces the title.

This is an example input file.  Comparing it with
the output it generates can show you how to
produce a simple document of your own.

\section{Ordinary Text}      % Produces section heading.  Lower-level
                             % sections are begun with similar
                             % \subsection and \subsubsection commands.

The ends  of words and sentences are marked
  by   spaces. It  doesn't matter how many
spaces    you type; one is as good as 100.  The
end of   a line counts as a space.

One   or more   blank lines denote the  end
of  a paragraph.

Since any number of consecutive spaces are treated
like a single one, the formatting of the input
file makes no difference to
      \LaTeX,                % The \LaTeX command generates the LaTeX logo.
but it makes a difference to you.  When you use
\LaTeX, making your input file as easy to read
as possible will be a great help as you write
your document and when you change it.  This sample
file shows how you can add comments to your own input
file.

Because printing is different from typewriting,
there are a number of things that you have to do
differently when preparing an input file than if
you were just typing the document directly.
Quotation marks like
       ``this''
have to be handled specially, as do quotes within
quotes:
       ``\,`this'            % \, separates the double and single quote.
        is what I just
        wrote, not  `that'\,''.

Dashes come in three sizes: an
       intra-word
dash, a medium dash for number ranges like
       1--2,
and a punctuation
       dash---like
this.

A sentence-ending space should be larger than the
space between words within a sentence.  You
sometimes have to type special commands in
conjunction with punctuation characters to get
this right, as in the following sentence.
       Gnats, gnus, etc.\ all  % `\ ' makes an inter-word space.
       begin with G\@.         % \@ marks end-of-sentence punctuation.
You should check the spaces after periods when
reading your output to make sure you haven't
forgotten any special cases.  Generating an
ellipsis
       \ldots\               % `\ ' is needed after `\ldots' because TeX
                             % ignores spaces after command names like \ldots
                             % made from \ + letters.
                             %
                             % Note how a `%' character causes TeX to ignore
                             % the end of the input line, so these blank lines
                             % do not start a new paragraph.
                             %
with the right spacing around the periods requires
a special command.

\LaTeX\ interprets some common characters as
commands, so you must type special commands to
generate them.  These characters include the
following:
       \$ \& \% \# \{ and \}.

In printing, text is usually emphasized with an
       \emph{italic}
type style.

\begin{em}
   A long segment of text can also be emphasized
   in this way.  Text within such a segment can be
   given \emph{additional} emphasis.
\end{em}

It is sometimes necessary to prevent \LaTeX\ from
breaking a line where it might otherwise do so.
This may be at a space, as between the ``Mr.''\ and
``Jones'' in
       ``Mr.~Jones'',        % ~ produces an unbreakable interword space.
or within a word---especially when the word is a
symbol like
       \mbox{\emph{itemnum}}
that makes little sense when hyphenated across
lines.

Footnotes\footnote{This is an example of a footnote.}
pose no problem.

\LaTeX\ is good at typesetting mathematical formulas
like
       \( x-3y + z = 7 \)
or
       \( a_{1} > x^{2n} + y^{2n} > x' \)
or
       \( \ip{A}{B} = \sum_{i} a_{i} b_{i} \).
The spaces you type in a formula are
ignored.  Remember that a letter like
       $x$                   % $ ... $  and  \( ... \)  are equivalent
is a formula when it denotes a mathematical
symbol, and it should be typed as one.

\section{Displayed Text}

Text is displayed by indenting it from the left
margin.  Quotations are commonly displayed.  There
are short quotations
\begin{quote}
   This is a short quotation.  It consists of a
   single paragraph of text.  See how it is formatted.
\end{quote}
and longer ones.
\begin{quotation}
   This is a longer quotation.  It consists of two
   paragraphs of text, neither of which are
   particularly interesting.

   This is the second paragraph of the quotation.  It
   is just as dull as the first paragraph.
\end{quotation}
Another frequently-displayed structure is a list.
The following is an example of an \emph{itemized}
list.
\begin{itemize}
   \item This is the first item of an itemized list.
         Each item in the list is marked with a ``tick''.
         You don't have to worry about what kind of tick
         mark is used.

   \item This is the second item of the list.  It
         contains another list nested inside it.  The inner
         list is an \emph{enumerated} list.
         \begin{enumerate}
            \item This is the first item of an enumerated
                  list that is nested within the itemized list.

            \item This is the second item of the inner list.
                  \LaTeX\ allows you to nest lists deeper than
                  you really should.
         \end{enumerate}
         This is the rest of the second item of the outer
         list.  It is no more interesting than any other
         part of the item.
   \item This is the third item of the list.
\end{itemize}
You can even display poetry.
\begin{verse}
   There is an environment
    for verse \\             % The \\ command separates lines
   Whose features some poets % within a stanza.
   will curse.

                             % One or more blank lines separate stanzas.

   For instead of making\\
   Them do \emph{all} line breaking, \\
   It allows them to put too many words on a line when they'd rather be
   forced to be terse.
\end{verse}

Mathematical formulas may also be displayed.  A
displayed formula
is
one-line long; multiline
formulas require special formatting instructions.
   \[  \ip{\Gamma}{\psi'} = x'' + y^{2} + z_{i}^{n}\]
Don't start a paragraph with a displayed equation,
nor make one a paragraph by itself.

\end{document}               % End of document.

\section{Bench Section 62}
% synthetic bench section 62
% This is a sample LaTeX input file.  (Version of 12 August 2004.)
%
% A '%' character causes TeX to ignore all remaining text on the line,
% and is used for comments like this one.

\documentclass{article}      % Specifies the document class

                             % The preamble begins here.
\title{An Example Document}  % Declares the document's title.
\author{Leslie Lamport}      % Declares the author's name.
\date{January 21, 1994}      % Deleting this command produces today's date.

\newcommand{\ip}[2]{(#1, #2)}
                             % Defines \ip{arg1}{arg2} to mean
                             % (arg1, arg2).

%\newcommand{\ip}[2]{\langle #1 | #2\rangle}
                             % This is an alternative definition of
                             % \ip that is commented out.

\begin{document}             % End of preamble and beginning of text.

\maketitle                   % Produces the title.

This is an example input file.  Comparing it with
the output it generates can show you how to
produce a simple document of your own.

\section{Ordinary Text}      % Produces section heading.  Lower-level
                             % sections are begun with similar
                             % \subsection and \subsubsection commands.

The ends  of words and sentences are marked
  by   spaces. It  doesn't matter how many
spaces    you type; one is as good as 100.  The
end of   a line counts as a space.

One   or more   blank lines denote the  end
of  a paragraph.

Since any number of consecutive spaces are treated
like a single one, the formatting of the input
file makes no difference to
      \LaTeX,                % The \LaTeX command generates the LaTeX logo.
but it makes a difference to you.  When you use
\LaTeX, making your input file as easy to read
as possible will be a great help as you write
your document and when you change it.  This sample
file shows how you can add comments to your own input
file.

Because printing is different from typewriting,
there are a number of things that you have to do
differently when preparing an input file than if
you were just typing the document directly.
Quotation marks like
       ``this''
have to be handled specially, as do quotes within
quotes:
       ``\,`this'            % \, separates the double and single quote.
        is what I just
        wrote, not  `that'\,''.

Dashes come in three sizes: an
       intra-word
dash, a medium dash for number ranges like
       1--2,
and a punctuation
       dash---like
this.

A sentence-ending space should be larger than the
space between words within a sentence.  You
sometimes have to type special commands in
conjunction with punctuation characters to get
this right, as in the following sentence.
       Gnats, gnus, etc.\ all  % `\ ' makes an inter-word space.
       begin with G\@.         % \@ marks end-of-sentence punctuation.
You should check the spaces after periods when
reading your output to make sure you haven't
forgotten any special cases.  Generating an
ellipsis
       \ldots\               % `\ ' is needed after `\ldots' because TeX
                             % ignores spaces after command names like \ldots
                             % made from \ + letters.
                             %
                             % Note how a `%' character causes TeX to ignore
                             % the end of the input line, so these blank lines
                             % do not start a new paragraph.
                             %
with the right spacing around the periods requires
a special command.

\LaTeX\ interprets some common characters as
commands, so you must type special commands to
generate them.  These characters include the
following:
       \$ \& \% \# \{ and \}.

In printing, text is usually emphasized with an
       \emph{italic}
type style.

\begin{em}
   A long segment of text can also be emphasized
   in this way.  Text within such a segment can be
   given \emph{additional} emphasis.
\end{em}

It is sometimes necessary to prevent \LaTeX\ from
breaking a line where it might otherwise do so.
This may be at a space, as between the ``Mr.''\ and
``Jones'' in
       ``Mr.~Jones'',        % ~ produces an unbreakable interword space.
or within a word---especially when the word is a
symbol like
       \mbox{\emph{itemnum}}
that makes little sense when hyphenated across
lines.

Footnotes\footnote{This is an example of a footnote.}
pose no problem.

\LaTeX\ is good at typesetting mathematical formulas
like
       \( x-3y + z = 7 \)
or
       \( a_{1} > x^{2n} + y^{2n} > x' \)
or
       \( \ip{A}{B} = \sum_{i} a_{i} b_{i} \).
The spaces you type in a formula are
ignored.  Remember that a letter like
       $x$                   % $ ... $  and  \( ... \)  are equivalent
is a formula when it denotes a mathematical
symbol, and it should be typed as one.

\section{Displayed Text}

Text is displayed by indenting it from the left
margin.  Quotations are commonly displayed.  There
are short quotations
\begin{quote}
   This is a short quotation.  It consists of a
   single paragraph of text.  See how it is formatted.
\end{quote}
and longer ones.
\begin{quotation}
   This is a longer quotation.  It consists of two
   paragraphs of text, neither of which are
   particularly interesting.

   This is the second paragraph of the quotation.  It
   is just as dull as the first paragraph.
\end{quotation}
Another frequently-displayed structure is a list.
The following is an example of an \emph{itemized}
list.
\begin{itemize}
   \item This is the first item of an itemized list.
         Each item in the list is marked with a ``tick''.
         You don't have to worry about what kind of tick
         mark is used.

   \item This is the second item of the list.  It
         contains another list nested inside it.  The inner
         list is an \emph{enumerated} list.
         \begin{enumerate}
            \item This is the first item of an enumerated
                  list that is nested within the itemized list.

            \item This is the second item of the inner list.
                  \LaTeX\ allows you to nest lists deeper than
                  you really should.
         \end{enumerate}
         This is the rest of the second item of the outer
         list.  It is no more interesting than any other
         part of the item.
   \item This is the third item of the list.
\end{itemize}
You can even display poetry.
\begin{verse}
   There is an environment
    for verse \\             % The \\ command separates lines
   Whose features some poets % within a stanza.
   will curse.

                             % One or more blank lines separate stanzas.

   For instead of making\\
   Them do \emph{all} line breaking, \\
   It allows them to put too many words on a line when they'd rather be
   forced to be terse.
\end{verse}

Mathematical formulas may also be displayed.  A
displayed formula
is
one-line long; multiline
formulas require special formatting instructions.
   \[  \ip{\Gamma}{\psi'} = x'' + y^{2} + z_{i}^{n}\]
Don't start a paragraph with a displayed equation,
nor make one a paragraph by itself.

\end{document}               % End of document.

\section{Bench Section 63}
% synthetic bench section 63
% This is a sample LaTeX input file.  (Version of 12 August 2004.)
%
% A '%' character causes TeX to ignore all remaining text on the line,
% and is used for comments like this one.

\documentclass{article}      % Specifies the document class

                             % The preamble begins here.
\title{An Example Document}  % Declares the document's title.
\author{Leslie Lamport}      % Declares the author's name.
\date{January 21, 1994}      % Deleting this command produces today's date.

\newcommand{\ip}[2]{(#1, #2)}
                             % Defines \ip{arg1}{arg2} to mean
                             % (arg1, arg2).

%\newcommand{\ip}[2]{\langle #1 | #2\rangle}
                             % This is an alternative definition of
                             % \ip that is commented out.

\begin{document}             % End of preamble and beginning of text.

\maketitle                   % Produces the title.

This is an example input file.  Comparing it with
the output it generates can show you how to
produce a simple document of your own.

\section{Ordinary Text}      % Produces section heading.  Lower-level
                             % sections are begun with similar
                             % \subsection and \subsubsection commands.

The ends  of words and sentences are marked
  by   spaces. It  doesn't matter how many
spaces    you type; one is as good as 100.  The
end of   a line counts as a space.

One   or more   blank lines denote the  end
of  a paragraph.

Since any number of consecutive spaces are treated
like a single one, the formatting of the input
file makes no difference to
      \LaTeX,                % The \LaTeX command generates the LaTeX logo.
but it makes a difference to you.  When you use
\LaTeX, making your input file as easy to read
as possible will be a great help as you write
your document and when you change it.  This sample
file shows how you can add comments to your own input
file.

Because printing is different from typewriting,
there are a number of things that you have to do
differently when preparing an input file than if
you were just typing the document directly.
Quotation marks like
       ``this''
have to be handled specially, as do quotes within
quotes:
       ``\,`this'            % \, separates the double and single quote.
        is what I just
        wrote, not  `that'\,''.

Dashes come in three sizes: an
       intra-word
dash, a medium dash for number ranges like
       1--2,
and a punctuation
       dash---like
this.

A sentence-ending space should be larger than the
space between words within a sentence.  You
sometimes have to type special commands in
conjunction with punctuation characters to get
this right, as in the following sentence.
       Gnats, gnus, etc.\ all  % `\ ' makes an inter-word space.
       begin with G\@.         % \@ marks end-of-sentence punctuation.
You should check the spaces after periods when
reading your output to make sure you haven't
forgotten any special cases.  Generating an
ellipsis
       \ldots\               % `\ ' is needed after `\ldots' because TeX
                             % ignores spaces after command names like \ldots
                             % made from \ + letters.
                             %
                             % Note how a `%' character causes TeX to ignore
                             % the end of the input line, so these blank lines
                             % do not start a new paragraph.
                             %
with the right spacing around the periods requires
a special command.

\LaTeX\ interprets some common characters as
commands, so you must type special commands to
generate them.  These characters include the
following:
       \$ \& \% \# \{ and \}.

In printing, text is usually emphasized with an
       \emph{italic}
type style.

\begin{em}
   A long segment of text can also be emphasized
   in this way.  Text within such a segment can be
   given \emph{additional} emphasis.
\end{em}

It is sometimes necessary to prevent \LaTeX\ from
breaking a line where it might otherwise do so.
This may be at a space, as between the ``Mr.''\ and
``Jones'' in
       ``Mr.~Jones'',        % ~ produces an unbreakable interword space.
or within a word---especially when the word is a
symbol like
       \mbox{\emph{itemnum}}
that makes little sense when hyphenated across
lines.

Footnotes\footnote{This is an example of a footnote.}
pose no problem.

\LaTeX\ is good at typesetting mathematical formulas
like
       \( x-3y + z = 7 \)
or
       \( a_{1} > x^{2n} + y^{2n} > x' \)
or
       \( \ip{A}{B} = \sum_{i} a_{i} b_{i} \).
The spaces you type in a formula are
ignored.  Remember that a letter like
       $x$                   % $ ... $  and  \( ... \)  are equivalent
is a formula when it denotes a mathematical
symbol, and it should be typed as one.

\section{Displayed Text}

Text is displayed by indenting it from the left
margin.  Quotations are commonly displayed.  There
are short quotations
\begin{quote}
   This is a short quotation.  It consists of a
   single paragraph of text.  See how it is formatted.
\end{quote}
and longer ones.
\begin{quotation}
   This is a longer quotation.  It consists of two
   paragraphs of text, neither of which are
   particularly interesting.

   This is the second paragraph of the quotation.  It
   is just as dull as the first paragraph.
\end{quotation}
Another frequently-displayed structure is a list.
The following is an example of an \emph{itemized}
list.
\begin{itemize}
   \item This is the first item of an itemized list.
         Each item in the list is marked with a ``tick''.
         You don't have to worry about what kind of tick
         mark is used.

   \item This is the second item of the list.  It
         contains another list nested inside it.  The inner
         list is an \emph{enumerated} list.
         \begin{enumerate}
            \item This is the first item of an enumerated
                  list that is nested within the itemized list.

            \item This is the second item of the inner list.
                  \LaTeX\ allows you to nest lists deeper than
                  you really should.
         \end{enumerate}
         This is the rest of the second item of the outer
         list.  It is no more interesting than any other
         part of the item.
   \item This is the third item of the list.
\end{itemize}
You can even display poetry.
\begin{verse}
   There is an environment
    for verse \\             % The \\ command separates lines
   Whose features some poets % within a stanza.
   will curse.

                             % One or more blank lines separate stanzas.

   For instead of making\\
   Them do \emph{all} line breaking, \\
   It allows them to put too many words on a line when they'd rather be
   forced to be terse.
\end{verse}

Mathematical formulas may also be displayed.  A
displayed formula
is
one-line long; multiline
formulas require special formatting instructions.
   \[  \ip{\Gamma}{\psi'} = x'' + y^{2} + z_{i}^{n}\]
Don't start a paragraph with a displayed equation,
nor make one a paragraph by itself.

\end{document}               % End of document.

\section{Bench Section 64}
% synthetic bench section 64
% This is a sample LaTeX input file.  (Version of 12 August 2004.)
%
% A '%' character causes TeX to ignore all remaining text on the line,
% and is used for comments like this one.

\documentclass{article}      % Specifies the document class

                             % The preamble begins here.
\title{An Example Document}  % Declares the document's title.
\author{Leslie Lamport}      % Declares the author's name.
\date{January 21, 1994}      % Deleting this command produces today's date.

\newcommand{\ip}[2]{(#1, #2)}
                             % Defines \ip{arg1}{arg2} to mean
                             % (arg1, arg2).

%\newcommand{\ip}[2]{\langle #1 | #2\rangle}
                             % This is an alternative definition of
                             % \ip that is commented out.

\begin{document}             % End of preamble and beginning of text.

\maketitle                   % Produces the title.

This is an example input file.  Comparing it with
the output it generates can show you how to
produce a simple document of your own.

\section{Ordinary Text}      % Produces section heading.  Lower-level
                             % sections are begun with similar
                             % \subsection and \subsubsection commands.

The ends  of words and sentences are marked
  by   spaces. It  doesn't matter how many
spaces    you type; one is as good as 100.  The
end of   a line counts as a space.

One   or more   blank lines denote the  end
of  a paragraph.

Since any number of consecutive spaces are treated
like a single one, the formatting of the input
file makes no difference to
      \LaTeX,                % The \LaTeX command generates the LaTeX logo.
but it makes a difference to you.  When you use
\LaTeX, making your input file as easy to read
as possible will be a great help as you write
your document and when you change it.  This sample
file shows how you can add comments to your own input
file.

Because printing is different from typewriting,
there are a number of things that you have to do
differently when preparing an input file than if
you were just typing the document directly.
Quotation marks like
       ``this''
have to be handled specially, as do quotes within
quotes:
       ``\,`this'            % \, separates the double and single quote.
        is what I just
        wrote, not  `that'\,''.

Dashes come in three sizes: an
       intra-word
dash, a medium dash for number ranges like
       1--2,
and a punctuation
       dash---like
this.

A sentence-ending space should be larger than the
space between words within a sentence.  You
sometimes have to type special commands in
conjunction with punctuation characters to get
this right, as in the following sentence.
       Gnats, gnus, etc.\ all  % `\ ' makes an inter-word space.
       begin with G\@.         % \@ marks end-of-sentence punctuation.
You should check the spaces after periods when
reading your output to make sure you haven't
forgotten any special cases.  Generating an
ellipsis
       \ldots\               % `\ ' is needed after `\ldots' because TeX
                             % ignores spaces after command names like \ldots
                             % made from \ + letters.
                             %
                             % Note how a `%' character causes TeX to ignore
                             % the end of the input line, so these blank lines
                             % do not start a new paragraph.
                             %
with the right spacing around the periods requires
a special command.

\LaTeX\ interprets some common characters as
commands, so you must type special commands to
generate them.  These characters include the
following:
       \$ \& \% \# \{ and \}.

In printing, text is usually emphasized with an
       \emph{italic}
type style.

\begin{em}
   A long segment of text can also be emphasized
   in this way.  Text within such a segment can be
   given \emph{additional} emphasis.
\end{em}

It is sometimes necessary to prevent \LaTeX\ from
breaking a line where it might otherwise do so.
This may be at a space, as between the ``Mr.''\ and
``Jones'' in
       ``Mr.~Jones'',        % ~ produces an unbreakable interword space.
or within a word---especially when the word is a
symbol like
       \mbox{\emph{itemnum}}
that makes little sense when hyphenated across
lines.

Footnotes\footnote{This is an example of a footnote.}
pose no problem.

\LaTeX\ is good at typesetting mathematical formulas
like
       \( x-3y + z = 7 \)
or
       \( a_{1} > x^{2n} + y^{2n} > x' \)
or
       \( \ip{A}{B} = \sum_{i} a_{i} b_{i} \).
The spaces you type in a formula are
ignored.  Remember that a letter like
       $x$                   % $ ... $  and  \( ... \)  are equivalent
is a formula when it denotes a mathematical
symbol, and it should be typed as one.

\section{Displayed Text}

Text is displayed by indenting it from the left
margin.  Quotations are commonly displayed.  There
are short quotations
\begin{quote}
   This is a short quotation.  It consists of a
   single paragraph of text.  See how it is formatted.
\end{quote}
and longer ones.
\begin{quotation}
   This is a longer quotation.  It consists of two
   paragraphs of text, neither of which are
   particularly interesting.

   This is the second paragraph of the quotation.  It
   is just as dull as the first paragraph.
\end{quotation}
Another frequently-displayed structure is a list.
The following is an example of an \emph{itemized}
list.
\begin{itemize}
   \item This is the first item of an itemized list.
         Each item in the list is marked with a ``tick''.
         You don't have to worry about what kind of tick
         mark is used.

   \item This is the second item of the list.  It
         contains another list nested inside it.  The inner
         list is an \emph{enumerated} list.
         \begin{enumerate}
            \item This is the first item of an enumerated
                  list that is nested within the itemized list.

            \item This is the second item of the inner list.
                  \LaTeX\ allows you to nest lists deeper than
                  you really should.
         \end{enumerate}
         This is the rest of the second item of the outer
         list.  It is no more interesting than any other
         part of the item.
   \item This is the third item of the list.
\end{itemize}
You can even display poetry.
\begin{verse}
   There is an environment
    for verse \\             % The \\ command separates lines
   Whose features some poets % within a stanza.
   will curse.

                             % One or more blank lines separate stanzas.

   For instead of making\\
   Them do \emph{all} line breaking, \\
   It allows them to put too many words on a line when they'd rather be
   forced to be terse.
\end{verse}

Mathematical formulas may also be displayed.  A
displayed formula
is
one-line long; multiline
formulas require special formatting instructions.
   \[  \ip{\Gamma}{\psi'} = x'' + y^{2} + z_{i}^{n}\]
Don't start a paragraph with a displayed equation,
nor make one a paragraph by itself.

\end{document}               % End of document.

\section{Bench Section 65}
% synthetic bench section 65
% This is a sample LaTeX input file.  (Version of 12 August 2004.)
%
% A '%' character causes TeX to ignore all remaining text on the line,
% and is used for comments like this one.

\documentclass{article}      % Specifies the document class

                             % The preamble begins here.
\title{An Example Document}  % Declares the document's title.
\author{Leslie Lamport}      % Declares the author's name.
\date{January 21, 1994}      % Deleting this command produces today's date.

\newcommand{\ip}[2]{(#1, #2)}
                             % Defines \ip{arg1}{arg2} to mean
                             % (arg1, arg2).

%\newcommand{\ip}[2]{\langle #1 | #2\rangle}
                             % This is an alternative definition of
                             % \ip that is commented out.

\begin{document}             % End of preamble and beginning of text.

\maketitle                   % Produces the title.

This is an example input file.  Comparing it with
the output it generates can show you how to
produce a simple document of your own.

\section{Ordinary Text}      % Produces section heading.  Lower-level
                             % sections are begun with similar
                             % \subsection and \subsubsection commands.

The ends  of words and sentences are marked
  by   spaces. It  doesn't matter how many
spaces    you type; one is as good as 100.  The
end of   a line counts as a space.

One   or more   blank lines denote the  end
of  a paragraph.

Since any number of consecutive spaces are treated
like a single one, the formatting of the input
file makes no difference to
      \LaTeX,                % The \LaTeX command generates the LaTeX logo.
but it makes a difference to you.  When you use
\LaTeX, making your input file as easy to read
as possible will be a great help as you write
your document and when you change it.  This sample
file shows how you can add comments to your own input
file.

Because printing is different from typewriting,
there are a number of things that you have to do
differently when preparing an input file than if
you were just typing the document directly.
Quotation marks like
       ``this''
have to be handled specially, as do quotes within
quotes:
       ``\,`this'            % \, separates the double and single quote.
        is what I just
        wrote, not  `that'\,''.

Dashes come in three sizes: an
       intra-word
dash, a medium dash for number ranges like
       1--2,
and a punctuation
       dash---like
this.

A sentence-ending space should be larger than the
space between words within a sentence.  You
sometimes have to type special commands in
conjunction with punctuation characters to get
this right, as in the following sentence.
       Gnats, gnus, etc.\ all  % `\ ' makes an inter-word space.
       begin with G\@.         % \@ marks end-of-sentence punctuation.
You should check the spaces after periods when
reading your output to make sure you haven't
forgotten any special cases.  Generating an
ellipsis
       \ldots\               % `\ ' is needed after `\ldots' because TeX
                             % ignores spaces after command names like \ldots
                             % made from \ + letters.
                             %
                             % Note how a `%' character causes TeX to ignore
                             % the end of the input line, so these blank lines
                             % do not start a new paragraph.
                             %
with the right spacing around the periods requires
a special command.

\LaTeX\ interprets some common characters as
commands, so you must type special commands to
generate them.  These characters include the
following:
       \$ \& \% \# \{ and \}.

In printing, text is usually emphasized with an
       \emph{italic}
type style.

\begin{em}
   A long segment of text can also be emphasized
   in this way.  Text within such a segment can be
   given \emph{additional} emphasis.
\end{em}

It is sometimes necessary to prevent \LaTeX\ from
breaking a line where it might otherwise do so.
This may be at a space, as between the ``Mr.''\ and
``Jones'' in
       ``Mr.~Jones'',        % ~ produces an unbreakable interword space.
or within a word---especially when the word is a
symbol like
       \mbox{\emph{itemnum}}
that makes little sense when hyphenated across
lines.

Footnotes\footnote{This is an example of a footnote.}
pose no problem.

\LaTeX\ is good at typesetting mathematical formulas
like
       \( x-3y + z = 7 \)
or
       \( a_{1} > x^{2n} + y^{2n} > x' \)
or
       \( \ip{A}{B} = \sum_{i} a_{i} b_{i} \).
The spaces you type in a formula are
ignored.  Remember that a letter like
       $x$                   % $ ... $  and  \( ... \)  are equivalent
is a formula when it denotes a mathematical
symbol, and it should be typed as one.

\section{Displayed Text}

Text is displayed by indenting it from the left
margin.  Quotations are commonly displayed.  There
are short quotations
\begin{quote}
   This is a short quotation.  It consists of a
   single paragraph of text.  See how it is formatted.
\end{quote}
and longer ones.
\begin{quotation}
   This is a longer quotation.  It consists of two
   paragraphs of text, neither of which are
   particularly interesting.

   This is the second paragraph of the quotation.  It
   is just as dull as the first paragraph.
\end{quotation}
Another frequently-displayed structure is a list.
The following is an example of an \emph{itemized}
list.
\begin{itemize}
   \item This is the first item of an itemized list.
         Each item in the list is marked with a ``tick''.
         You don't have to worry about what kind of tick
         mark is used.

   \item This is the second item of the list.  It
         contains another list nested inside it.  The inner
         list is an \emph{enumerated} list.
         \begin{enumerate}
            \item This is the first item of an enumerated
                  list that is nested within the itemized list.

            \item This is the second item of the inner list.
                  \LaTeX\ allows you to nest lists deeper than
                  you really should.
         \end{enumerate}
         This is the rest of the second item of the outer
         list.  It is no more interesting than any other
         part of the item.
   \item This is the third item of the list.
\end{itemize}
You can even display poetry.
\begin{verse}
   There is an environment
    for verse \\             % The \\ command separates lines
   Whose features some poets % within a stanza.
   will curse.

                             % One or more blank lines separate stanzas.

   For instead of making\\
   Them do \emph{all} line breaking, \\
   It allows them to put too many words on a line when they'd rather be
   forced to be terse.
\end{verse}

Mathematical formulas may also be displayed.  A
displayed formula
is
one-line long; multiline
formulas require special formatting instructions.
   \[  \ip{\Gamma}{\psi'} = x'' + y^{2} + z_{i}^{n}\]
Don't start a paragraph with a displayed equation,
nor make one a paragraph by itself.

\end{document}               % End of document.

\section{Bench Section 66}
% synthetic bench section 66
% This is a sample LaTeX input file.  (Version of 12 August 2004.)
%
% A '%' character causes TeX to ignore all remaining text on the line,
% and is used for comments like this one.

\documentclass{article}      % Specifies the document class

                             % The preamble begins here.
\title{An Example Document}  % Declares the document's title.
\author{Leslie Lamport}      % Declares the author's name.
\date{January 21, 1994}      % Deleting this command produces today's date.

\newcommand{\ip}[2]{(#1, #2)}
                             % Defines \ip{arg1}{arg2} to mean
                             % (arg1, arg2).

%\newcommand{\ip}[2]{\langle #1 | #2\rangle}
                             % This is an alternative definition of
                             % \ip that is commented out.

\begin{document}             % End of preamble and beginning of text.

\maketitle                   % Produces the title.

This is an example input file.  Comparing it with
the output it generates can show you how to
produce a simple document of your own.

\section{Ordinary Text}      % Produces section heading.  Lower-level
                             % sections are begun with similar
                             % \subsection and \subsubsection commands.

The ends  of words and sentences are marked
  by   spaces. It  doesn't matter how many
spaces    you type; one is as good as 100.  The
end of   a line counts as a space.

One   or more   blank lines denote the  end
of  a paragraph.

Since any number of consecutive spaces are treated
like a single one, the formatting of the input
file makes no difference to
      \LaTeX,                % The \LaTeX command generates the LaTeX logo.
but it makes a difference to you.  When you use
\LaTeX, making your input file as easy to read
as possible will be a great help as you write
your document and when you change it.  This sample
file shows how you can add comments to your own input
file.

Because printing is different from typewriting,
there are a number of things that you have to do
differently when preparing an input file than if
you were just typing the document directly.
Quotation marks like
       ``this''
have to be handled specially, as do quotes within
quotes:
       ``\,`this'            % \, separates the double and single quote.
        is what I just
        wrote, not  `that'\,''.

Dashes come in three sizes: an
       intra-word
dash, a medium dash for number ranges like
       1--2,
and a punctuation
       dash---like
this.

A sentence-ending space should be larger than the
space between words within a sentence.  You
sometimes have to type special commands in
conjunction with punctuation characters to get
this right, as in the following sentence.
       Gnats, gnus, etc.\ all  % `\ ' makes an inter-word space.
       begin with G\@.         % \@ marks end-of-sentence punctuation.
You should check the spaces after periods when
reading your output to make sure you haven't
forgotten any special cases.  Generating an
ellipsis
       \ldots\               % `\ ' is needed after `\ldots' because TeX
                             % ignores spaces after command names like \ldots
                             % made from \ + letters.
                             %
                             % Note how a `%' character causes TeX to ignore
                             % the end of the input line, so these blank lines
                             % do not start a new paragraph.
                             %
with the right spacing around the periods requires
a special command.

\LaTeX\ interprets some common characters as
commands, so you must type special commands to
generate them.  These characters include the
following:
       \$ \& \% \# \{ and \}.

In printing, text is usually emphasized with an
       \emph{italic}
type style.

\begin{em}
   A long segment of text can also be emphasized
   in this way.  Text within such a segment can be
   given \emph{additional} emphasis.
\end{em}

It is sometimes necessary to prevent \LaTeX\ from
breaking a line where it might otherwise do so.
This may be at a space, as between the ``Mr.''\ and
``Jones'' in
       ``Mr.~Jones'',        % ~ produces an unbreakable interword space.
or within a word---especially when the word is a
symbol like
       \mbox{\emph{itemnum}}
that makes little sense when hyphenated across
lines.

Footnotes\footnote{This is an example of a footnote.}
pose no problem.

\LaTeX\ is good at typesetting mathematical formulas
like
       \( x-3y + z = 7 \)
or
       \( a_{1} > x^{2n} + y^{2n} > x' \)
or
       \( \ip{A}{B} = \sum_{i} a_{i} b_{i} \).
The spaces you type in a formula are
ignored.  Remember that a letter like
       $x$                   % $ ... $  and  \( ... \)  are equivalent
is a formula when it denotes a mathematical
symbol, and it should be typed as one.

\section{Displayed Text}

Text is displayed by indenting it from the left
margin.  Quotations are commonly displayed.  There
are short quotations
\begin{quote}
   This is a short quotation.  It consists of a
   single paragraph of text.  See how it is formatted.
\end{quote}
and longer ones.
\begin{quotation}
   This is a longer quotation.  It consists of two
   paragraphs of text, neither of which are
   particularly interesting.

   This is the second paragraph of the quotation.  It
   is just as dull as the first paragraph.
\end{quotation}
Another frequently-displayed structure is a list.
The following is an example of an \emph{itemized}
list.
\begin{itemize}
   \item This is the first item of an itemized list.
         Each item in the list is marked with a ``tick''.
         You don't have to worry about what kind of tick
         mark is used.

   \item This is the second item of the list.  It
         contains another list nested inside it.  The inner
         list is an \emph{enumerated} list.
         \begin{enumerate}
            \item This is the first item of an enumerated
                  list that is nested within the itemized list.

            \item This is the second item of the inner list.
                  \LaTeX\ allows you to nest lists deeper than
                  you really should.
         \end{enumerate}
         This is the rest of the second item of the outer
         list.  It is no more interesting than any other
         part of the item.
   \item This is the third item of the list.
\end{itemize}
You can even display poetry.
\begin{verse}
   There is an environment
    for verse \\             % The \\ command separates lines
   Whose features some poets % within a stanza.
   will curse.

                             % One or more blank lines separate stanzas.

   For instead of making\\
   Them do \emph{all} line breaking, \\
   It allows them to put too many words on a line when they'd rather be
   forced to be terse.
\end{verse}

Mathematical formulas may also be displayed.  A
displayed formula
is
one-line long; multiline
formulas require special formatting instructions.
   \[  \ip{\Gamma}{\psi'} = x'' + y^{2} + z_{i}^{n}\]
Don't start a paragraph with a displayed equation,
nor make one a paragraph by itself.

\end{document}               % End of document.

\section{Bench Section 67}
% synthetic bench section 67
% This is a sample LaTeX input file.  (Version of 12 August 2004.)
%
% A '%' character causes TeX to ignore all remaining text on the line,
% and is used for comments like this one.

\documentclass{article}      % Specifies the document class

                             % The preamble begins here.
\title{An Example Document}  % Declares the document's title.
\author{Leslie Lamport}      % Declares the author's name.
\date{January 21, 1994}      % Deleting this command produces today's date.

\newcommand{\ip}[2]{(#1, #2)}
                             % Defines \ip{arg1}{arg2} to mean
                             % (arg1, arg2).

%\newcommand{\ip}[2]{\langle #1 | #2\rangle}
                             % This is an alternative definition of
                             % \ip that is commented out.

\begin{document}             % End of preamble and beginning of text.

\maketitle                   % Produces the title.

This is an example input file.  Comparing it with
the output it generates can show you how to
produce a simple document of your own.

\section{Ordinary Text}      % Produces section heading.  Lower-level
                             % sections are begun with similar
                             % \subsection and \subsubsection commands.

The ends  of words and sentences are marked
  by   spaces. It  doesn't matter how many
spaces    you type; one is as good as 100.  The
end of   a line counts as a space.

One   or more   blank lines denote the  end
of  a paragraph.

Since any number of consecutive spaces are treated
like a single one, the formatting of the input
file makes no difference to
      \LaTeX,                % The \LaTeX command generates the LaTeX logo.
but it makes a difference to you.  When you use
\LaTeX, making your input file as easy to read
as possible will be a great help as you write
your document and when you change it.  This sample
file shows how you can add comments to your own input
file.

Because printing is different from typewriting,
there are a number of things that you have to do
differently when preparing an input file than if
you were just typing the document directly.
Quotation marks like
       ``this''
have to be handled specially, as do quotes within
quotes:
       ``\,`this'            % \, separates the double and single quote.
        is what I just
        wrote, not  `that'\,''.

Dashes come in three sizes: an
       intra-word
dash, a medium dash for number ranges like
       1--2,
and a punctuation
       dash---like
this.

A sentence-ending space should be larger than the
space between words within a sentence.  You
sometimes have to type special commands in
conjunction with punctuation characters to get
this right, as in the following sentence.
       Gnats, gnus, etc.\ all  % `\ ' makes an inter-word space.
       begin with G\@.         % \@ marks end-of-sentence punctuation.
You should check the spaces after periods when
reading your output to make sure you haven't
forgotten any special cases.  Generating an
ellipsis
       \ldots\               % `\ ' is needed after `\ldots' because TeX
                             % ignores spaces after command names like \ldots
                             % made from \ + letters.
                             %
                             % Note how a `%' character causes TeX to ignore
                             % the end of the input line, so these blank lines
                             % do not start a new paragraph.
                             %
with the right spacing around the periods requires
a special command.

\LaTeX\ interprets some common characters as
commands, so you must type special commands to
generate them.  These characters include the
following:
       \$ \& \% \# \{ and \}.

In printing, text is usually emphasized with an
       \emph{italic}
type style.

\begin{em}
   A long segment of text can also be emphasized
   in this way.  Text within such a segment can be
   given \emph{additional} emphasis.
\end{em}

It is sometimes necessary to prevent \LaTeX\ from
breaking a line where it might otherwise do so.
This may be at a space, as between the ``Mr.''\ and
``Jones'' in
       ``Mr.~Jones'',        % ~ produces an unbreakable interword space.
or within a word---especially when the word is a
symbol like
       \mbox{\emph{itemnum}}
that makes little sense when hyphenated across
lines.

Footnotes\footnote{This is an example of a footnote.}
pose no problem.

\LaTeX\ is good at typesetting mathematical formulas
like
       \( x-3y + z = 7 \)
or
       \( a_{1} > x^{2n} + y^{2n} > x' \)
or
       \( \ip{A}{B} = \sum_{i} a_{i} b_{i} \).
The spaces you type in a formula are
ignored.  Remember that a letter like
       $x$                   % $ ... $  and  \( ... \)  are equivalent
is a formula when it denotes a mathematical
symbol, and it should be typed as one.

\section{Displayed Text}

Text is displayed by indenting it from the left
margin.  Quotations are commonly displayed.  There
are short quotations
\begin{quote}
   This is a short quotation.  It consists of a
   single paragraph of text.  See how it is formatted.
\end{quote}
and longer ones.
\begin{quotation}
   This is a longer quotation.  It consists of two
   paragraphs of text, neither of which are
   particularly interesting.

   This is the second paragraph of the quotation.  It
   is just as dull as the first paragraph.
\end{quotation}
Another frequently-displayed structure is a list.
The following is an example of an \emph{itemized}
list.
\begin{itemize}
   \item This is the first item of an itemized list.
         Each item in the list is marked with a ``tick''.
         You don't have to worry about what kind of tick
         mark is used.

   \item This is the second item of the list.  It
         contains another list nested inside it.  The inner
         list is an \emph{enumerated} list.
         \begin{enumerate}
            \item This is the first item of an enumerated
                  list that is nested within the itemized list.

            \item This is the second item of the inner list.
                  \LaTeX\ allows you to nest lists deeper than
                  you really should.
         \end{enumerate}
         This is the rest of the second item of the outer
         list.  It is no more interesting than any other
         part of the item.
   \item This is the third item of the list.
\end{itemize}
You can even display poetry.
\begin{verse}
   There is an environment
    for verse \\             % The \\ command separates lines
   Whose features some poets % within a stanza.
   will curse.

                             % One or more blank lines separate stanzas.

   For instead of making\\
   Them do \emph{all} line breaking, \\
   It allows them to put too many words on a line when they'd rather be
   forced to be terse.
\end{verse}

Mathematical formulas may also be displayed.  A
displayed formula
is
one-line long; multiline
formulas require special formatting instructions.
   \[  \ip{\Gamma}{\psi'} = x'' + y^{2} + z_{i}^{n}\]
Don't start a paragraph with a displayed equation,
nor make one a paragraph by itself.

\end{document}               % End of document.

\section{Bench Section 68}
% synthetic bench section 68
% This is a sample LaTeX input file.  (Version of 12 August 2004.)
%
% A '%' character causes TeX to ignore all remaining text on the line,
% and is used for comments like this one.

\documentclass{article}      % Specifies the document class

                             % The preamble begins here.
\title{An Example Document}  % Declares the document's title.
\author{Leslie Lamport}      % Declares the author's name.
\date{January 21, 1994}      % Deleting this command produces today's date.

\newcommand{\ip}[2]{(#1, #2)}
                             % Defines \ip{arg1}{arg2} to mean
                             % (arg1, arg2).

%\newcommand{\ip}[2]{\langle #1 | #2\rangle}
                             % This is an alternative definition of
                             % \ip that is commented out.

\begin{document}             % End of preamble and beginning of text.

\maketitle                   % Produces the title.

This is an example input file.  Comparing it with
the output it generates can show you how to
produce a simple document of your own.

\section{Ordinary Text}      % Produces section heading.  Lower-level
                             % sections are begun with similar
                             % \subsection and \subsubsection commands.

The ends  of words and sentences are marked
  by   spaces. It  doesn't matter how many
spaces    you type; one is as good as 100.  The
end of   a line counts as a space.

One   or more   blank lines denote the  end
of  a paragraph.

Since any number of consecutive spaces are treated
like a single one, the formatting of the input
file makes no difference to
      \LaTeX,                % The \LaTeX command generates the LaTeX logo.
but it makes a difference to you.  When you use
\LaTeX, making your input file as easy to read
as possible will be a great help as you write
your document and when you change it.  This sample
file shows how you can add comments to your own input
file.

Because printing is different from typewriting,
there are a number of things that you have to do
differently when preparing an input file than if
you were just typing the document directly.
Quotation marks like
       ``this''
have to be handled specially, as do quotes within
quotes:
       ``\,`this'            % \, separates the double and single quote.
        is what I just
        wrote, not  `that'\,''.

Dashes come in three sizes: an
       intra-word
dash, a medium dash for number ranges like
       1--2,
and a punctuation
       dash---like
this.

A sentence-ending space should be larger than the
space between words within a sentence.  You
sometimes have to type special commands in
conjunction with punctuation characters to get
this right, as in the following sentence.
       Gnats, gnus, etc.\ all  % `\ ' makes an inter-word space.
       begin with G\@.         % \@ marks end-of-sentence punctuation.
You should check the spaces after periods when
reading your output to make sure you haven't
forgotten any special cases.  Generating an
ellipsis
       \ldots\               % `\ ' is needed after `\ldots' because TeX
                             % ignores spaces after command names like \ldots
                             % made from \ + letters.
                             %
                             % Note how a `%' character causes TeX to ignore
                             % the end of the input line, so these blank lines
                             % do not start a new paragraph.
                             %
with the right spacing around the periods requires
a special command.

\LaTeX\ interprets some common characters as
commands, so you must type special commands to
generate them.  These characters include the
following:
       \$ \& \% \# \{ and \}.

In printing, text is usually emphasized with an
       \emph{italic}
type style.

\begin{em}
   A long segment of text can also be emphasized
   in this way.  Text within such a segment can be
   given \emph{additional} emphasis.
\end{em}

It is sometimes necessary to prevent \LaTeX\ from
breaking a line where it might otherwise do so.
This may be at a space, as between the ``Mr.''\ and
``Jones'' in
       ``Mr.~Jones'',        % ~ produces an unbreakable interword space.
or within a word---especially when the word is a
symbol like
       \mbox{\emph{itemnum}}
that makes little sense when hyphenated across
lines.

Footnotes\footnote{This is an example of a footnote.}
pose no problem.

\LaTeX\ is good at typesetting mathematical formulas
like
       \( x-3y + z = 7 \)
or
       \( a_{1} > x^{2n} + y^{2n} > x' \)
or
       \( \ip{A}{B} = \sum_{i} a_{i} b_{i} \).
The spaces you type in a formula are
ignored.  Remember that a letter like
       $x$                   % $ ... $  and  \( ... \)  are equivalent
is a formula when it denotes a mathematical
symbol, and it should be typed as one.

\section{Displayed Text}

Text is displayed by indenting it from the left
margin.  Quotations are commonly displayed.  There
are short quotations
\begin{quote}
   This is a short quotation.  It consists of a
   single paragraph of text.  See how it is formatted.
\end{quote}
and longer ones.
\begin{quotation}
   This is a longer quotation.  It consists of two
   paragraphs of text, neither of which are
   particularly interesting.

   This is the second paragraph of the quotation.  It
   is just as dull as the first paragraph.
\end{quotation}
Another frequently-displayed structure is a list.
The following is an example of an \emph{itemized}
list.
\begin{itemize}
   \item This is the first item of an itemized list.
         Each item in the list is marked with a ``tick''.
         You don't have to worry about what kind of tick
         mark is used.

   \item This is the second item of the list.  It
         contains another list nested inside it.  The inner
         list is an \emph{enumerated} list.
         \begin{enumerate}
            \item This is the first item of an enumerated
                  list that is nested within the itemized list.

            \item This is the second item of the inner list.
                  \LaTeX\ allows you to nest lists deeper than
                  you really should.
         \end{enumerate}
         This is the rest of the second item of the outer
         list.  It is no more interesting than any other
         part of the item.
   \item This is the third item of the list.
\end{itemize}
You can even display poetry.
\begin{verse}
   There is an environment
    for verse \\             % The \\ command separates lines
   Whose features some poets % within a stanza.
   will curse.

                             % One or more blank lines separate stanzas.

   For instead of making\\
   Them do \emph{all} line breaking, \\
   It allows them to put too many words on a line when they'd rather be
   forced to be terse.
\end{verse}

Mathematical formulas may also be displayed.  A
displayed formula
is
one-line long; multiline
formulas require special formatting instructions.
   \[  \ip{\Gamma}{\psi'} = x'' + y^{2} + z_{i}^{n}\]
Don't start a paragraph with a displayed equation,
nor make one a paragraph by itself.

\end{document}               % End of document.

\section{Bench Section 69}
% synthetic bench section 69
% This is a sample LaTeX input file.  (Version of 12 August 2004.)
%
% A '%' character causes TeX to ignore all remaining text on the line,
% and is used for comments like this one.

\documentclass{article}      % Specifies the document class

                             % The preamble begins here.
\title{An Example Document}  % Declares the document's title.
\author{Leslie Lamport}      % Declares the author's name.
\date{January 21, 1994}      % Deleting this command produces today's date.

\newcommand{\ip}[2]{(#1, #2)}
                             % Defines \ip{arg1}{arg2} to mean
                             % (arg1, arg2).

%\newcommand{\ip}[2]{\langle #1 | #2\rangle}
                             % This is an alternative definition of
                             % \ip that is commented out.

\begin{document}             % End of preamble and beginning of text.

\maketitle                   % Produces the title.

This is an example input file.  Comparing it with
the output it generates can show you how to
produce a simple document of your own.

\section{Ordinary Text}      % Produces section heading.  Lower-level
                             % sections are begun with similar
                             % \subsection and \subsubsection commands.

The ends  of words and sentences are marked
  by   spaces. It  doesn't matter how many
spaces    you type; one is as good as 100.  The
end of   a line counts as a space.

One   or more   blank lines denote the  end
of  a paragraph.

Since any number of consecutive spaces are treated
like a single one, the formatting of the input
file makes no difference to
      \LaTeX,                % The \LaTeX command generates the LaTeX logo.
but it makes a difference to you.  When you use
\LaTeX, making your input file as easy to read
as possible will be a great help as you write
your document and when you change it.  This sample
file shows how you can add comments to your own input
file.

Because printing is different from typewriting,
there are a number of things that you have to do
differently when preparing an input file than if
you were just typing the document directly.
Quotation marks like
       ``this''
have to be handled specially, as do quotes within
quotes:
       ``\,`this'            % \, separates the double and single quote.
        is what I just
        wrote, not  `that'\,''.

Dashes come in three sizes: an
       intra-word
dash, a medium dash for number ranges like
       1--2,
and a punctuation
       dash---like
this.

A sentence-ending space should be larger than the
space between words within a sentence.  You
sometimes have to type special commands in
conjunction with punctuation characters to get
this right, as in the following sentence.
       Gnats, gnus, etc.\ all  % `\ ' makes an inter-word space.
       begin with G\@.         % \@ marks end-of-sentence punctuation.
You should check the spaces after periods when
reading your output to make sure you haven't
forgotten any special cases.  Generating an
ellipsis
       \ldots\               % `\ ' is needed after `\ldots' because TeX
                             % ignores spaces after command names like \ldots
                             % made from \ + letters.
                             %
                             % Note how a `%' character causes TeX to ignore
                             % the end of the input line, so these blank lines
                             % do not start a new paragraph.
                             %
with the right spacing around the periods requires
a special command.

\LaTeX\ interprets some common characters as
commands, so you must type special commands to
generate them.  These characters include the
following:
       \$ \& \% \# \{ and \}.

In printing, text is usually emphasized with an
       \emph{italic}
type style.

\begin{em}
   A long segment of text can also be emphasized
   in this way.  Text within such a segment can be
   given \emph{additional} emphasis.
\end{em}

It is sometimes necessary to prevent \LaTeX\ from
breaking a line where it might otherwise do so.
This may be at a space, as between the ``Mr.''\ and
``Jones'' in
       ``Mr.~Jones'',        % ~ produces an unbreakable interword space.
or within a word---especially when the word is a
symbol like
       \mbox{\emph{itemnum}}
that makes little sense when hyphenated across
lines.

Footnotes\footnote{This is an example of a footnote.}
pose no problem.

\LaTeX\ is good at typesetting mathematical formulas
like
       \( x-3y + z = 7 \)
or
       \( a_{1} > x^{2n} + y^{2n} > x' \)
or
       \( \ip{A}{B} = \sum_{i} a_{i} b_{i} \).
The spaces you type in a formula are
ignored.  Remember that a letter like
       $x$                   % $ ... $  and  \( ... \)  are equivalent
is a formula when it denotes a mathematical
symbol, and it should be typed as one.

\section{Displayed Text}

Text is displayed by indenting it from the left
margin.  Quotations are commonly displayed.  There
are short quotations
\begin{quote}
   This is a short quotation.  It consists of a
   single paragraph of text.  See how it is formatted.
\end{quote}
and longer ones.
\begin{quotation}
   This is a longer quotation.  It consists of two
   paragraphs of text, neither of which are
   particularly interesting.

   This is the second paragraph of the quotation.  It
   is just as dull as the first paragraph.
\end{quotation}
Another frequently-displayed structure is a list.
The following is an example of an \emph{itemized}
list.
\begin{itemize}
   \item This is the first item of an itemized list.
         Each item in the list is marked with a ``tick''.
         You don't have to worry about what kind of tick
         mark is used.

   \item This is the second item of the list.  It
         contains another list nested inside it.  The inner
         list is an \emph{enumerated} list.
         \begin{enumerate}
            \item This is the first item of an enumerated
                  list that is nested within the itemized list.

            \item This is the second item of the inner list.
                  \LaTeX\ allows you to nest lists deeper than
                  you really should.
         \end{enumerate}
         This is the rest of the second item of the outer
         list.  It is no more interesting than any other
         part of the item.
   \item This is the third item of the list.
\end{itemize}
You can even display poetry.
\begin{verse}
   There is an environment
    for verse \\             % The \\ command separates lines
   Whose features some poets % within a stanza.
   will curse.

                             % One or more blank lines separate stanzas.

   For instead of making\\
   Them do \emph{all} line breaking, \\
   It allows them to put too many words on a line when they'd rather be
   forced to be terse.
\end{verse}

Mathematical formulas may also be displayed.  A
displayed formula
is
one-line long; multiline
formulas require special formatting instructions.
   \[  \ip{\Gamma}{\psi'} = x'' + y^{2} + z_{i}^{n}\]
Don't start a paragraph with a displayed equation,
nor make one a paragraph by itself.

\end{document}               % End of document.

\section{Bench Section 70}
% synthetic bench section 70
% This is a sample LaTeX input file.  (Version of 12 August 2004.)
%
% A '%' character causes TeX to ignore all remaining text on the line,
% and is used for comments like this one.

\documentclass{article}      % Specifies the document class

                             % The preamble begins here.
\title{An Example Document}  % Declares the document's title.
\author{Leslie Lamport}      % Declares the author's name.
\date{January 21, 1994}      % Deleting this command produces today's date.

\newcommand{\ip}[2]{(#1, #2)}
                             % Defines \ip{arg1}{arg2} to mean
                             % (arg1, arg2).

%\newcommand{\ip}[2]{\langle #1 | #2\rangle}
                             % This is an alternative definition of
                             % \ip that is commented out.

\begin{document}             % End of preamble and beginning of text.

\maketitle                   % Produces the title.

This is an example input file.  Comparing it with
the output it generates can show you how to
produce a simple document of your own.

\section{Ordinary Text}      % Produces section heading.  Lower-level
                             % sections are begun with similar
                             % \subsection and \subsubsection commands.

The ends  of words and sentences are marked
  by   spaces. It  doesn't matter how many
spaces    you type; one is as good as 100.  The
end of   a line counts as a space.

One   or more   blank lines denote the  end
of  a paragraph.

Since any number of consecutive spaces are treated
like a single one, the formatting of the input
file makes no difference to
      \LaTeX,                % The \LaTeX command generates the LaTeX logo.
but it makes a difference to you.  When you use
\LaTeX, making your input file as easy to read
as possible will be a great help as you write
your document and when you change it.  This sample
file shows how you can add comments to your own input
file.

Because printing is different from typewriting,
there are a number of things that you have to do
differently when preparing an input file than if
you were just typing the document directly.
Quotation marks like
       ``this''
have to be handled specially, as do quotes within
quotes:
       ``\,`this'            % \, separates the double and single quote.
        is what I just
        wrote, not  `that'\,''.

Dashes come in three sizes: an
       intra-word
dash, a medium dash for number ranges like
       1--2,
and a punctuation
       dash---like
this.

A sentence-ending space should be larger than the
space between words within a sentence.  You
sometimes have to type special commands in
conjunction with punctuation characters to get
this right, as in the following sentence.
       Gnats, gnus, etc.\ all  % `\ ' makes an inter-word space.
       begin with G\@.         % \@ marks end-of-sentence punctuation.
You should check the spaces after periods when
reading your output to make sure you haven't
forgotten any special cases.  Generating an
ellipsis
       \ldots\               % `\ ' is needed after `\ldots' because TeX
                             % ignores spaces after command names like \ldots
                             % made from \ + letters.
                             %
                             % Note how a `%' character causes TeX to ignore
                             % the end of the input line, so these blank lines
                             % do not start a new paragraph.
                             %
with the right spacing around the periods requires
a special command.

\LaTeX\ interprets some common characters as
commands, so you must type special commands to
generate them.  These characters include the
following:
       \$ \& \% \# \{ and \}.

In printing, text is usually emphasized with an
       \emph{italic}
type style.

\begin{em}
   A long segment of text can also be emphasized
   in this way.  Text within such a segment can be
   given \emph{additional} emphasis.
\end{em}

It is sometimes necessary to prevent \LaTeX\ from
breaking a line where it might otherwise do so.
This may be at a space, as between the ``Mr.''\ and
``Jones'' in
       ``Mr.~Jones'',        % ~ produces an unbreakable interword space.
or within a word---especially when the word is a
symbol like
       \mbox{\emph{itemnum}}
that makes little sense when hyphenated across
lines.

Footnotes\footnote{This is an example of a footnote.}
pose no problem.

\LaTeX\ is good at typesetting mathematical formulas
like
       \( x-3y + z = 7 \)
or
       \( a_{1} > x^{2n} + y^{2n} > x' \)
or
       \( \ip{A}{B} = \sum_{i} a_{i} b_{i} \).
The spaces you type in a formula are
ignored.  Remember that a letter like
       $x$                   % $ ... $  and  \( ... \)  are equivalent
is a formula when it denotes a mathematical
symbol, and it should be typed as one.

\section{Displayed Text}

Text is displayed by indenting it from the left
margin.  Quotations are commonly displayed.  There
are short quotations
\begin{quote}
   This is a short quotation.  It consists of a
   single paragraph of text.  See how it is formatted.
\end{quote}
and longer ones.
\begin{quotation}
   This is a longer quotation.  It consists of two
   paragraphs of text, neither of which are
   particularly interesting.

   This is the second paragraph of the quotation.  It
   is just as dull as the first paragraph.
\end{quotation}
Another frequently-displayed structure is a list.
The following is an example of an \emph{itemized}
list.
\begin{itemize}
   \item This is the first item of an itemized list.
         Each item in the list is marked with a ``tick''.
         You don't have to worry about what kind of tick
         mark is used.

   \item This is the second item of the list.  It
         contains another list nested inside it.  The inner
         list is an \emph{enumerated} list.
         \begin{enumerate}
            \item This is the first item of an enumerated
                  list that is nested within the itemized list.

            \item This is the second item of the inner list.
                  \LaTeX\ allows you to nest lists deeper than
                  you really should.
         \end{enumerate}
         This is the rest of the second item of the outer
         list.  It is no more interesting than any other
         part of the item.
   \item This is the third item of the list.
\end{itemize}
You can even display poetry.
\begin{verse}
   There is an environment
    for verse \\             % The \\ command separates lines
   Whose features some poets % within a stanza.
   will curse.

                             % One or more blank lines separate stanzas.

   For instead of making\\
   Them do \emph{all} line breaking, \\
   It allows them to put too many words on a line when they'd rather be
   forced to be terse.
\end{verse}

Mathematical formulas may also be displayed.  A
displayed formula
is
one-line long; multiline
formulas require special formatting instructions.
   \[  \ip{\Gamma}{\psi'} = x'' + y^{2} + z_{i}^{n}\]
Don't start a paragraph with a displayed equation,
nor make one a paragraph by itself.

\end{document}               % End of document.

\section{Bench Section 71}
% synthetic bench section 71
% This is a sample LaTeX input file.  (Version of 12 August 2004.)
%
% A '%' character causes TeX to ignore all remaining text on the line,
% and is used for comments like this one.

\documentclass{article}      % Specifies the document class

                             % The preamble begins here.
\title{An Example Document}  % Declares the document's title.
\author{Leslie Lamport}      % Declares the author's name.
\date{January 21, 1994}      % Deleting this command produces today's date.

\newcommand{\ip}[2]{(#1, #2)}
                             % Defines \ip{arg1}{arg2} to mean
                             % (arg1, arg2).

%\newcommand{\ip}[2]{\langle #1 | #2\rangle}
                             % This is an alternative definition of
                             % \ip that is commented out.

\begin{document}             % End of preamble and beginning of text.

\maketitle                   % Produces the title.

This is an example input file.  Comparing it with
the output it generates can show you how to
produce a simple document of your own.

\section{Ordinary Text}      % Produces section heading.  Lower-level
                             % sections are begun with similar
                             % \subsection and \subsubsection commands.

The ends  of words and sentences are marked
  by   spaces. It  doesn't matter how many
spaces    you type; one is as good as 100.  The
end of   a line counts as a space.

One   or more   blank lines denote the  end
of  a paragraph.

Since any number of consecutive spaces are treated
like a single one, the formatting of the input
file makes no difference to
      \LaTeX,                % The \LaTeX command generates the LaTeX logo.
but it makes a difference to you.  When you use
\LaTeX, making your input file as easy to read
as possible will be a great help as you write
your document and when you change it.  This sample
file shows how you can add comments to your own input
file.

Because printing is different from typewriting,
there are a number of things that you have to do
differently when preparing an input file than if
you were just typing the document directly.
Quotation marks like
       ``this''
have to be handled specially, as do quotes within
quotes:
       ``\,`this'            % \, separates the double and single quote.
        is what I just
        wrote, not  `that'\,''.

Dashes come in three sizes: an
       intra-word
dash, a medium dash for number ranges like
       1--2,
and a punctuation
       dash---like
this.

A sentence-ending space should be larger than the
space between words within a sentence.  You
sometimes have to type special commands in
conjunction with punctuation characters to get
this right, as in the following sentence.
       Gnats, gnus, etc.\ all  % `\ ' makes an inter-word space.
       begin with G\@.         % \@ marks end-of-sentence punctuation.
You should check the spaces after periods when
reading your output to make sure you haven't
forgotten any special cases.  Generating an
ellipsis
       \ldots\               % `\ ' is needed after `\ldots' because TeX
                             % ignores spaces after command names like \ldots
                             % made from \ + letters.
                             %
                             % Note how a `%' character causes TeX to ignore
                             % the end of the input line, so these blank lines
                             % do not start a new paragraph.
                             %
with the right spacing around the periods requires
a special command.

\LaTeX\ interprets some common characters as
commands, so you must type special commands to
generate them.  These characters include the
following:
       \$ \& \% \# \{ and \}.

In printing, text is usually emphasized with an
       \emph{italic}
type style.

\begin{em}
   A long segment of text can also be emphasized
   in this way.  Text within such a segment can be
   given \emph{additional} emphasis.
\end{em}

It is sometimes necessary to prevent \LaTeX\ from
breaking a line where it might otherwise do so.
This may be at a space, as between the ``Mr.''\ and
``Jones'' in
       ``Mr.~Jones'',        % ~ produces an unbreakable interword space.
or within a word---especially when the word is a
symbol like
       \mbox{\emph{itemnum}}
that makes little sense when hyphenated across
lines.

Footnotes\footnote{This is an example of a footnote.}
pose no problem.

\LaTeX\ is good at typesetting mathematical formulas
like
       \( x-3y + z = 7 \)
or
       \( a_{1} > x^{2n} + y^{2n} > x' \)
or
       \( \ip{A}{B} = \sum_{i} a_{i} b_{i} \).
The spaces you type in a formula are
ignored.  Remember that a letter like
       $x$                   % $ ... $  and  \( ... \)  are equivalent
is a formula when it denotes a mathematical
symbol, and it should be typed as one.

\section{Displayed Text}

Text is displayed by indenting it from the left
margin.  Quotations are commonly displayed.  There
are short quotations
\begin{quote}
   This is a short quotation.  It consists of a
   single paragraph of text.  See how it is formatted.
\end{quote}
and longer ones.
\begin{quotation}
   This is a longer quotation.  It consists of two
   paragraphs of text, neither of which are
   particularly interesting.

   This is the second paragraph of the quotation.  It
   is just as dull as the first paragraph.
\end{quotation}
Another frequently-displayed structure is a list.
The following is an example of an \emph{itemized}
list.
\begin{itemize}
   \item This is the first item of an itemized list.
         Each item in the list is marked with a ``tick''.
         You don't have to worry about what kind of tick
         mark is used.

   \item This is the second item of the list.  It
         contains another list nested inside it.  The inner
         list is an \emph{enumerated} list.
         \begin{enumerate}
            \item This is the first item of an enumerated
                  list that is nested within the itemized list.

            \item This is the second item of the inner list.
                  \LaTeX\ allows you to nest lists deeper than
                  you really should.
         \end{enumerate}
         This is the rest of the second item of the outer
         list.  It is no more interesting than any other
         part of the item.
   \item This is the third item of the list.
\end{itemize}
You can even display poetry.
\begin{verse}
   There is an environment
    for verse \\             % The \\ command separates lines
   Whose features some poets % within a stanza.
   will curse.

                             % One or more blank lines separate stanzas.

   For instead of making\\
   Them do \emph{all} line breaking, \\
   It allows them to put too many words on a line when they'd rather be
   forced to be terse.
\end{verse}

Mathematical formulas may also be displayed.  A
displayed formula
is
one-line long; multiline
formulas require special formatting instructions.
   \[  \ip{\Gamma}{\psi'} = x'' + y^{2} + z_{i}^{n}\]
Don't start a paragraph with a displayed equation,
nor make one a paragraph by itself.

\end{document}               % End of document.

\section{Bench Section 72}
% synthetic bench section 72
% This is a sample LaTeX input file.  (Version of 12 August 2004.)
%
% A '%' character causes TeX to ignore all remaining text on the line,
% and is used for comments like this one.

\documentclass{article}      % Specifies the document class

                             % The preamble begins here.
\title{An Example Document}  % Declares the document's title.
\author{Leslie Lamport}      % Declares the author's name.
\date{January 21, 1994}      % Deleting this command produces today's date.

\newcommand{\ip}[2]{(#1, #2)}
                             % Defines \ip{arg1}{arg2} to mean
                             % (arg1, arg2).

%\newcommand{\ip}[2]{\langle #1 | #2\rangle}
                             % This is an alternative definition of
                             % \ip that is commented out.

\begin{document}             % End of preamble and beginning of text.

\maketitle                   % Produces the title.

This is an example input file.  Comparing it with
the output it generates can show you how to
produce a simple document of your own.

\section{Ordinary Text}      % Produces section heading.  Lower-level
                             % sections are begun with similar
                             % \subsection and \subsubsection commands.

The ends  of words and sentences are marked
  by   spaces. It  doesn't matter how many
spaces    you type; one is as good as 100.  The
end of   a line counts as a space.

One   or more   blank lines denote the  end
of  a paragraph.

Since any number of consecutive spaces are treated
like a single one, the formatting of the input
file makes no difference to
      \LaTeX,                % The \LaTeX command generates the LaTeX logo.
but it makes a difference to you.  When you use
\LaTeX, making your input file as easy to read
as possible will be a great help as you write
your document and when you change it.  This sample
file shows how you can add comments to your own input
file.

Because printing is different from typewriting,
there are a number of things that you have to do
differently when preparing an input file than if
you were just typing the document directly.
Quotation marks like
       ``this''
have to be handled specially, as do quotes within
quotes:
       ``\,`this'            % \, separates the double and single quote.
        is what I just
        wrote, not  `that'\,''.

Dashes come in three sizes: an
       intra-word
dash, a medium dash for number ranges like
       1--2,
and a punctuation
       dash---like
this.

A sentence-ending space should be larger than the
space between words within a sentence.  You
sometimes have to type special commands in
conjunction with punctuation characters to get
this right, as in the following sentence.
       Gnats, gnus, etc.\ all  % `\ ' makes an inter-word space.
       begin with G\@.         % \@ marks end-of-sentence punctuation.
You should check the spaces after periods when
reading your output to make sure you haven't
forgotten any special cases.  Generating an
ellipsis
       \ldots\               % `\ ' is needed after `\ldots' because TeX
                             % ignores spaces after command names like \ldots
                             % made from \ + letters.
                             %
                             % Note how a `%' character causes TeX to ignore
                             % the end of the input line, so these blank lines
                             % do not start a new paragraph.
                             %
with the right spacing around the periods requires
a special command.

\LaTeX\ interprets some common characters as
commands, so you must type special commands to
generate them.  These characters include the
following:
       \$ \& \% \# \{ and \}.

In printing, text is usually emphasized with an
       \emph{italic}
type style.

\begin{em}
   A long segment of text can also be emphasized
   in this way.  Text within such a segment can be
   given \emph{additional} emphasis.
\end{em}

It is sometimes necessary to prevent \LaTeX\ from
breaking a line where it might otherwise do so.
This may be at a space, as between the ``Mr.''\ and
``Jones'' in
       ``Mr.~Jones'',        % ~ produces an unbreakable interword space.
or within a word---especially when the word is a
symbol like
       \mbox{\emph{itemnum}}
that makes little sense when hyphenated across
lines.

Footnotes\footnote{This is an example of a footnote.}
pose no problem.

\LaTeX\ is good at typesetting mathematical formulas
like
       \( x-3y + z = 7 \)
or
       \( a_{1} > x^{2n} + y^{2n} > x' \)
or
       \( \ip{A}{B} = \sum_{i} a_{i} b_{i} \).
The spaces you type in a formula are
ignored.  Remember that a letter like
       $x$                   % $ ... $  and  \( ... \)  are equivalent
is a formula when it denotes a mathematical
symbol, and it should be typed as one.

\section{Displayed Text}

Text is displayed by indenting it from the left
margin.  Quotations are commonly displayed.  There
are short quotations
\begin{quote}
   This is a short quotation.  It consists of a
   single paragraph of text.  See how it is formatted.
\end{quote}
and longer ones.
\begin{quotation}
   This is a longer quotation.  It consists of two
   paragraphs of text, neither of which are
   particularly interesting.

   This is the second paragraph of the quotation.  It
   is just as dull as the first paragraph.
\end{quotation}
Another frequently-displayed structure is a list.
The following is an example of an \emph{itemized}
list.
\begin{itemize}
   \item This is the first item of an itemized list.
         Each item in the list is marked with a ``tick''.
         You don't have to worry about what kind of tick
         mark is used.

   \item This is the second item of the list.  It
         contains another list nested inside it.  The inner
         list is an \emph{enumerated} list.
         \begin{enumerate}
            \item This is the first item of an enumerated
                  list that is nested within the itemized list.

            \item This is the second item of the inner list.
                  \LaTeX\ allows you to nest lists deeper than
                  you really should.
         \end{enumerate}
         This is the rest of the second item of the outer
         list.  It is no more interesting than any other
         part of the item.
   \item This is the third item of the list.
\end{itemize}
You can even display poetry.
\begin{verse}
   There is an environment
    for verse \\             % The \\ command separates lines
   Whose features some poets % within a stanza.
   will curse.

                             % One or more blank lines separate stanzas.

   For instead of making\\
   Them do \emph{all} line breaking, \\
   It allows them to put too many words on a line when they'd rather be
   forced to be terse.
\end{verse}

Mathematical formulas may also be displayed.  A
displayed formula
is
one-line long; multiline
formulas require special formatting instructions.
   \[  \ip{\Gamma}{\psi'} = x'' + y^{2} + z_{i}^{n}\]
Don't start a paragraph with a displayed equation,
nor make one a paragraph by itself.

\end{document}               % End of document.

\section{Bench Section 73}
% synthetic bench section 73
% This is a sample LaTeX input file.  (Version of 12 August 2004.)
%
% A '%' character causes TeX to ignore all remaining text on the line,
% and is used for comments like this one.

\documentclass{article}      % Specifies the document class

                             % The preamble begins here.
\title{An Example Document}  % Declares the document's title.
\author{Leslie Lamport}      % Declares the author's name.
\date{January 21, 1994}      % Deleting this command produces today's date.

\newcommand{\ip}[2]{(#1, #2)}
                             % Defines \ip{arg1}{arg2} to mean
                             % (arg1, arg2).

%\newcommand{\ip}[2]{\langle #1 | #2\rangle}
                             % This is an alternative definition of
                             % \ip that is commented out.

\begin{document}             % End of preamble and beginning of text.

\maketitle                   % Produces the title.

This is an example input file.  Comparing it with
the output it generates can show you how to
produce a simple document of your own.

\section{Ordinary Text}      % Produces section heading.  Lower-level
                             % sections are begun with similar
                             % \subsection and \subsubsection commands.

The ends  of words and sentences are marked
  by   spaces. It  doesn't matter how many
spaces    you type; one is as good as 100.  The
end of   a line counts as a space.

One   or more   blank lines denote the  end
of  a paragraph.

Since any number of consecutive spaces are treated
like a single one, the formatting of the input
file makes no difference to
      \LaTeX,                % The \LaTeX command generates the LaTeX logo.
but it makes a difference to you.  When you use
\LaTeX, making your input file as easy to read
as possible will be a great help as you write
your document and when you change it.  This sample
file shows how you can add comments to your own input
file.

Because printing is different from typewriting,
there are a number of things that you have to do
differently when preparing an input file than if
you were just typing the document directly.
Quotation marks like
       ``this''
have to be handled specially, as do quotes within
quotes:
       ``\,`this'            % \, separates the double and single quote.
        is what I just
        wrote, not  `that'\,''.

Dashes come in three sizes: an
       intra-word
dash, a medium dash for number ranges like
       1--2,
and a punctuation
       dash---like
this.

A sentence-ending space should be larger than the
space between words within a sentence.  You
sometimes have to type special commands in
conjunction with punctuation characters to get
this right, as in the following sentence.
       Gnats, gnus, etc.\ all  % `\ ' makes an inter-word space.
       begin with G\@.         % \@ marks end-of-sentence punctuation.
You should check the spaces after periods when
reading your output to make sure you haven't
forgotten any special cases.  Generating an
ellipsis
       \ldots\               % `\ ' is needed after `\ldots' because TeX
                             % ignores spaces after command names like \ldots
                             % made from \ + letters.
                             %
                             % Note how a `%' character causes TeX to ignore
                             % the end of the input line, so these blank lines
                             % do not start a new paragraph.
                             %
with the right spacing around the periods requires
a special command.

\LaTeX\ interprets some common characters as
commands, so you must type special commands to
generate them.  These characters include the
following:
       \$ \& \% \# \{ and \}.

In printing, text is usually emphasized with an
       \emph{italic}
type style.

\begin{em}
   A long segment of text can also be emphasized
   in this way.  Text within such a segment can be
   given \emph{additional} emphasis.
\end{em}

It is sometimes necessary to prevent \LaTeX\ from
breaking a line where it might otherwise do so.
This may be at a space, as between the ``Mr.''\ and
``Jones'' in
       ``Mr.~Jones'',        % ~ produces an unbreakable interword space.
or within a word---especially when the word is a
symbol like
       \mbox{\emph{itemnum}}
that makes little sense when hyphenated across
lines.

Footnotes\footnote{This is an example of a footnote.}
pose no problem.

\LaTeX\ is good at typesetting mathematical formulas
like
       \( x-3y + z = 7 \)
or
       \( a_{1} > x^{2n} + y^{2n} > x' \)
or
       \( \ip{A}{B} = \sum_{i} a_{i} b_{i} \).
The spaces you type in a formula are
ignored.  Remember that a letter like
       $x$                   % $ ... $  and  \( ... \)  are equivalent
is a formula when it denotes a mathematical
symbol, and it should be typed as one.

\section{Displayed Text}

Text is displayed by indenting it from the left
margin.  Quotations are commonly displayed.  There
are short quotations
\begin{quote}
   This is a short quotation.  It consists of a
   single paragraph of text.  See how it is formatted.
\end{quote}
and longer ones.
\begin{quotation}
   This is a longer quotation.  It consists of two
   paragraphs of text, neither of which are
   particularly interesting.

   This is the second paragraph of the quotation.  It
   is just as dull as the first paragraph.
\end{quotation}
Another frequently-displayed structure is a list.
The following is an example of an \emph{itemized}
list.
\begin{itemize}
   \item This is the first item of an itemized list.
         Each item in the list is marked with a ``tick''.
         You don't have to worry about what kind of tick
         mark is used.

   \item This is the second item of the list.  It
         contains another list nested inside it.  The inner
         list is an \emph{enumerated} list.
         \begin{enumerate}
            \item This is the first item of an enumerated
                  list that is nested within the itemized list.

            \item This is the second item of the inner list.
                  \LaTeX\ allows you to nest lists deeper than
                  you really should.
         \end{enumerate}
         This is the rest of the second item of the outer
         list.  It is no more interesting than any other
         part of the item.
   \item This is the third item of the list.
\end{itemize}
You can even display poetry.
\begin{verse}
   There is an environment
    for verse \\             % The \\ command separates lines
   Whose features some poets % within a stanza.
   will curse.

                             % One or more blank lines separate stanzas.

   For instead of making\\
   Them do \emph{all} line breaking, \\
   It allows them to put too many words on a line when they'd rather be
   forced to be terse.
\end{verse}

Mathematical formulas may also be displayed.  A
displayed formula
is
one-line long; multiline
formulas require special formatting instructions.
   \[  \ip{\Gamma}{\psi'} = x'' + y^{2} + z_{i}^{n}\]
Don't start a paragraph with a displayed equation,
nor make one a paragraph by itself.

\end{document}               % End of document.

\section{Bench Section 74}
% synthetic bench section 74
% This is a sample LaTeX input file.  (Version of 12 August 2004.)
%
% A '%' character causes TeX to ignore all remaining text on the line,
% and is used for comments like this one.

\documentclass{article}      % Specifies the document class

                             % The preamble begins here.
\title{An Example Document}  % Declares the document's title.
\author{Leslie Lamport}      % Declares the author's name.
\date{January 21, 1994}      % Deleting this command produces today's date.

\newcommand{\ip}[2]{(#1, #2)}
                             % Defines \ip{arg1}{arg2} to mean
                             % (arg1, arg2).

%\newcommand{\ip}[2]{\langle #1 | #2\rangle}
                             % This is an alternative definition of
                             % \ip that is commented out.

\begin{document}             % End of preamble and beginning of text.

\maketitle                   % Produces the title.

This is an example input file.  Comparing it with
the output it generates can show you how to
produce a simple document of your own.

\section{Ordinary Text}      % Produces section heading.  Lower-level
                             % sections are begun with similar
                             % \subsection and \subsubsection commands.

The ends  of words and sentences are marked
  by   spaces. It  doesn't matter how many
spaces    you type; one is as good as 100.  The
end of   a line counts as a space.

One   or more   blank lines denote the  end
of  a paragraph.

Since any number of consecutive spaces are treated
like a single one, the formatting of the input
file makes no difference to
      \LaTeX,                % The \LaTeX command generates the LaTeX logo.
but it makes a difference to you.  When you use
\LaTeX, making your input file as easy to read
as possible will be a great help as you write
your document and when you change it.  This sample
file shows how you can add comments to your own input
file.

Because printing is different from typewriting,
there are a number of things that you have to do
differently when preparing an input file than if
you were just typing the document directly.
Quotation marks like
       ``this''
have to be handled specially, as do quotes within
quotes:
       ``\,`this'            % \, separates the double and single quote.
        is what I just
        wrote, not  `that'\,''.

Dashes come in three sizes: an
       intra-word
dash, a medium dash for number ranges like
       1--2,
and a punctuation
       dash---like
this.

A sentence-ending space should be larger than the
space between words within a sentence.  You
sometimes have to type special commands in
conjunction with punctuation characters to get
this right, as in the following sentence.
       Gnats, gnus, etc.\ all  % `\ ' makes an inter-word space.
       begin with G\@.         % \@ marks end-of-sentence punctuation.
You should check the spaces after periods when
reading your output to make sure you haven't
forgotten any special cases.  Generating an
ellipsis
       \ldots\               % `\ ' is needed after `\ldots' because TeX
                             % ignores spaces after command names like \ldots
                             % made from \ + letters.
                             %
                             % Note how a `%' character causes TeX to ignore
                             % the end of the input line, so these blank lines
                             % do not start a new paragraph.
                             %
with the right spacing around the periods requires
a special command.

\LaTeX\ interprets some common characters as
commands, so you must type special commands to
generate them.  These characters include the
following:
       \$ \& \% \# \{ and \}.

In printing, text is usually emphasized with an
       \emph{italic}
type style.

\begin{em}
   A long segment of text can also be emphasized
   in this way.  Text within such a segment can be
   given \emph{additional} emphasis.
\end{em}

It is sometimes necessary to prevent \LaTeX\ from
breaking a line where it might otherwise do so.
This may be at a space, as between the ``Mr.''\ and
``Jones'' in
       ``Mr.~Jones'',        % ~ produces an unbreakable interword space.
or within a word---especially when the word is a
symbol like
       \mbox{\emph{itemnum}}
that makes little sense when hyphenated across
lines.

Footnotes\footnote{This is an example of a footnote.}
pose no problem.

\LaTeX\ is good at typesetting mathematical formulas
like
       \( x-3y + z = 7 \)
or
       \( a_{1} > x^{2n} + y^{2n} > x' \)
or
       \( \ip{A}{B} = \sum_{i} a_{i} b_{i} \).
The spaces you type in a formula are
ignored.  Remember that a letter like
       $x$                   % $ ... $  and  \( ... \)  are equivalent
is a formula when it denotes a mathematical
symbol, and it should be typed as one.

\section{Displayed Text}

Text is displayed by indenting it from the left
margin.  Quotations are commonly displayed.  There
are short quotations
\begin{quote}
   This is a short quotation.  It consists of a
   single paragraph of text.  See how it is formatted.
\end{quote}
and longer ones.
\begin{quotation}
   This is a longer quotation.  It consists of two
   paragraphs of text, neither of which are
   particularly interesting.

   This is the second paragraph of the quotation.  It
   is just as dull as the first paragraph.
\end{quotation}
Another frequently-displayed structure is a list.
The following is an example of an \emph{itemized}
list.
\begin{itemize}
   \item This is the first item of an itemized list.
         Each item in the list is marked with a ``tick''.
         You don't have to worry about what kind of tick
         mark is used.

   \item This is the second item of the list.  It
         contains another list nested inside it.  The inner
         list is an \emph{enumerated} list.
         \begin{enumerate}
            \item This is the first item of an enumerated
                  list that is nested within the itemized list.

            \item This is the second item of the inner list.
                  \LaTeX\ allows you to nest lists deeper than
                  you really should.
         \end{enumerate}
         This is the rest of the second item of the outer
         list.  It is no more interesting than any other
         part of the item.
   \item This is the third item of the list.
\end{itemize}
You can even display poetry.
\begin{verse}
   There is an environment
    for verse \\             % The \\ command separates lines
   Whose features some poets % within a stanza.
   will curse.

                             % One or more blank lines separate stanzas.

   For instead of making\\
   Them do \emph{all} line breaking, \\
   It allows them to put too many words on a line when they'd rather be
   forced to be terse.
\end{verse}

Mathematical formulas may also be displayed.  A
displayed formula
is
one-line long; multiline
formulas require special formatting instructions.
   \[  \ip{\Gamma}{\psi'} = x'' + y^{2} + z_{i}^{n}\]
Don't start a paragraph with a displayed equation,
nor make one a paragraph by itself.

\end{document}               % End of document.

\section{Bench Section 75}
% synthetic bench section 75
% This is a sample LaTeX input file.  (Version of 12 August 2004.)
%
% A '%' character causes TeX to ignore all remaining text on the line,
% and is used for comments like this one.

\documentclass{article}      % Specifies the document class

                             % The preamble begins here.
\title{An Example Document}  % Declares the document's title.
\author{Leslie Lamport}      % Declares the author's name.
\date{January 21, 1994}      % Deleting this command produces today's date.

\newcommand{\ip}[2]{(#1, #2)}
                             % Defines \ip{arg1}{arg2} to mean
                             % (arg1, arg2).

%\newcommand{\ip}[2]{\langle #1 | #2\rangle}
                             % This is an alternative definition of
                             % \ip that is commented out.

\begin{document}             % End of preamble and beginning of text.

\maketitle                   % Produces the title.

This is an example input file.  Comparing it with
the output it generates can show you how to
produce a simple document of your own.

\section{Ordinary Text}      % Produces section heading.  Lower-level
                             % sections are begun with similar
                             % \subsection and \subsubsection commands.

The ends  of words and sentences are marked
  by   spaces. It  doesn't matter how many
spaces    you type; one is as good as 100.  The
end of   a line counts as a space.

One   or more   blank lines denote the  end
of  a paragraph.

Since any number of consecutive spaces are treated
like a single one, the formatting of the input
file makes no difference to
      \LaTeX,                % The \LaTeX command generates the LaTeX logo.
but it makes a difference to you.  When you use
\LaTeX, making your input file as easy to read
as possible will be a great help as you write
your document and when you change it.  This sample
file shows how you can add comments to your own input
file.

Because printing is different from typewriting,
there are a number of things that you have to do
differently when preparing an input file than if
you were just typing the document directly.
Quotation marks like
       ``this''
have to be handled specially, as do quotes within
quotes:
       ``\,`this'            % \, separates the double and single quote.
        is what I just
        wrote, not  `that'\,''.

Dashes come in three sizes: an
       intra-word
dash, a medium dash for number ranges like
       1--2,
and a punctuation
       dash---like
this.

A sentence-ending space should be larger than the
space between words within a sentence.  You
sometimes have to type special commands in
conjunction with punctuation characters to get
this right, as in the following sentence.
       Gnats, gnus, etc.\ all  % `\ ' makes an inter-word space.
       begin with G\@.         % \@ marks end-of-sentence punctuation.
You should check the spaces after periods when
reading your output to make sure you haven't
forgotten any special cases.  Generating an
ellipsis
       \ldots\               % `\ ' is needed after `\ldots' because TeX
                             % ignores spaces after command names like \ldots
                             % made from \ + letters.
                             %
                             % Note how a `%' character causes TeX to ignore
                             % the end of the input line, so these blank lines
                             % do not start a new paragraph.
                             %
with the right spacing around the periods requires
a special command.

\LaTeX\ interprets some common characters as
commands, so you must type special commands to
generate them.  These characters include the
following:
       \$ \& \% \# \{ and \}.

In printing, text is usually emphasized with an
       \emph{italic}
type style.

\begin{em}
   A long segment of text can also be emphasized
   in this way.  Text within such a segment can be
   given \emph{additional} emphasis.
\end{em}

It is sometimes necessary to prevent \LaTeX\ from
breaking a line where it might otherwise do so.
This may be at a space, as between the ``Mr.''\ and
``Jones'' in
       ``Mr.~Jones'',        % ~ produces an unbreakable interword space.
or within a word---especially when the word is a
symbol like
       \mbox{\emph{itemnum}}
that makes little sense when hyphenated across
lines.

Footnotes\footnote{This is an example of a footnote.}
pose no problem.

\LaTeX\ is good at typesetting mathematical formulas
like
       \( x-3y + z = 7 \)
or
       \( a_{1} > x^{2n} + y^{2n} > x' \)
or
       \( \ip{A}{B} = \sum_{i} a_{i} b_{i} \).
The spaces you type in a formula are
ignored.  Remember that a letter like
       $x$                   % $ ... $  and  \( ... \)  are equivalent
is a formula when it denotes a mathematical
symbol, and it should be typed as one.

\section{Displayed Text}

Text is displayed by indenting it from the left
margin.  Quotations are commonly displayed.  There
are short quotations
\begin{quote}
   This is a short quotation.  It consists of a
   single paragraph of text.  See how it is formatted.
\end{quote}
and longer ones.
\begin{quotation}
   This is a longer quotation.  It consists of two
   paragraphs of text, neither of which are
   particularly interesting.

   This is the second paragraph of the quotation.  It
   is just as dull as the first paragraph.
\end{quotation}
Another frequently-displayed structure is a list.
The following is an example of an \emph{itemized}
list.
\begin{itemize}
   \item This is the first item of an itemized list.
         Each item in the list is marked with a ``tick''.
         You don't have to worry about what kind of tick
         mark is used.

   \item This is the second item of the list.  It
         contains another list nested inside it.  The inner
         list is an \emph{enumerated} list.
         \begin{enumerate}
            \item This is the first item of an enumerated
                  list that is nested within the itemized list.

            \item This is the second item of the inner list.
                  \LaTeX\ allows you to nest lists deeper than
                  you really should.
         \end{enumerate}
         This is the rest of the second item of the outer
         list.  It is no more interesting than any other
         part of the item.
   \item This is the third item of the list.
\end{itemize}
You can even display poetry.
\begin{verse}
   There is an environment
    for verse \\             % The \\ command separates lines
   Whose features some poets % within a stanza.
   will curse.

                             % One or more blank lines separate stanzas.

   For instead of making\\
   Them do \emph{all} line breaking, \\
   It allows them to put too many words on a line when they'd rather be
   forced to be terse.
\end{verse}

Mathematical formulas may also be displayed.  A
displayed formula
is
one-line long; multiline
formulas require special formatting instructions.
   \[  \ip{\Gamma}{\psi'} = x'' + y^{2} + z_{i}^{n}\]
Don't start a paragraph with a displayed equation,
nor make one a paragraph by itself.

\end{document}               % End of document.

\section{Bench Section 76}
% synthetic bench section 76
% This is a sample LaTeX input file.  (Version of 12 August 2004.)
%
% A '%' character causes TeX to ignore all remaining text on the line,
% and is used for comments like this one.

\documentclass{article}      % Specifies the document class

                             % The preamble begins here.
\title{An Example Document}  % Declares the document's title.
\author{Leslie Lamport}      % Declares the author's name.
\date{January 21, 1994}      % Deleting this command produces today's date.

\newcommand{\ip}[2]{(#1, #2)}
                             % Defines \ip{arg1}{arg2} to mean
                             % (arg1, arg2).

%\newcommand{\ip}[2]{\langle #1 | #2\rangle}
                             % This is an alternative definition of
                             % \ip that is commented out.

\begin{document}             % End of preamble and beginning of text.

\maketitle                   % Produces the title.

This is an example input file.  Comparing it with
the output it generates can show you how to
produce a simple document of your own.

\section{Ordinary Text}      % Produces section heading.  Lower-level
                             % sections are begun with similar
                             % \subsection and \subsubsection commands.

The ends  of words and sentences are marked
  by   spaces. It  doesn't matter how many
spaces    you type; one is as good as 100.  The
end of   a line counts as a space.

One   or more   blank lines denote the  end
of  a paragraph.

Since any number of consecutive spaces are treated
like a single one, the formatting of the input
file makes no difference to
      \LaTeX,                % The \LaTeX command generates the LaTeX logo.
but it makes a difference to you.  When you use
\LaTeX, making your input file as easy to read
as possible will be a great help as you write
your document and when you change it.  This sample
file shows how you can add comments to your own input
file.

Because printing is different from typewriting,
there are a number of things that you have to do
differently when preparing an input file than if
you were just typing the document directly.
Quotation marks like
       ``this''
have to be handled specially, as do quotes within
quotes:
       ``\,`this'            % \, separates the double and single quote.
        is what I just
        wrote, not  `that'\,''.

Dashes come in three sizes: an
       intra-word
dash, a medium dash for number ranges like
       1--2,
and a punctuation
       dash---like
this.

A sentence-ending space should be larger than the
space between words within a sentence.  You
sometimes have to type special commands in
conjunction with punctuation characters to get
this right, as in the following sentence.
       Gnats, gnus, etc.\ all  % `\ ' makes an inter-word space.
       begin with G\@.         % \@ marks end-of-sentence punctuation.
You should check the spaces after periods when
reading your output to make sure you haven't
forgotten any special cases.  Generating an
ellipsis
       \ldots\               % `\ ' is needed after `\ldots' because TeX
                             % ignores spaces after command names like \ldots
                             % made from \ + letters.
                             %
                             % Note how a `%' character causes TeX to ignore
                             % the end of the input line, so these blank lines
                             % do not start a new paragraph.
                             %
with the right spacing around the periods requires
a special command.

\LaTeX\ interprets some common characters as
commands, so you must type special commands to
generate them.  These characters include the
following:
       \$ \& \% \# \{ and \}.

In printing, text is usually emphasized with an
       \emph{italic}
type style.

\begin{em}
   A long segment of text can also be emphasized
   in this way.  Text within such a segment can be
   given \emph{additional} emphasis.
\end{em}

It is sometimes necessary to prevent \LaTeX\ from
breaking a line where it might otherwise do so.
This may be at a space, as between the ``Mr.''\ and
``Jones'' in
       ``Mr.~Jones'',        % ~ produces an unbreakable interword space.
or within a word---especially when the word is a
symbol like
       \mbox{\emph{itemnum}}
that makes little sense when hyphenated across
lines.

Footnotes\footnote{This is an example of a footnote.}
pose no problem.

\LaTeX\ is good at typesetting mathematical formulas
like
       \( x-3y + z = 7 \)
or
       \( a_{1} > x^{2n} + y^{2n} > x' \)
or
       \( \ip{A}{B} = \sum_{i} a_{i} b_{i} \).
The spaces you type in a formula are
ignored.  Remember that a letter like
       $x$                   % $ ... $  and  \( ... \)  are equivalent
is a formula when it denotes a mathematical
symbol, and it should be typed as one.

\section{Displayed Text}

Text is displayed by indenting it from the left
margin.  Quotations are commonly displayed.  There
are short quotations
\begin{quote}
   This is a short quotation.  It consists of a
   single paragraph of text.  See how it is formatted.
\end{quote}
and longer ones.
\begin{quotation}
   This is a longer quotation.  It consists of two
   paragraphs of text, neither of which are
   particularly interesting.

   This is the second paragraph of the quotation.  It
   is just as dull as the first paragraph.
\end{quotation}
Another frequently-displayed structure is a list.
The following is an example of an \emph{itemized}
list.
\begin{itemize}
   \item This is the first item of an itemized list.
         Each item in the list is marked with a ``tick''.
         You don't have to worry about what kind of tick
         mark is used.

   \item This is the second item of the list.  It
         contains another list nested inside it.  The inner
         list is an \emph{enumerated} list.
         \begin{enumerate}
            \item This is the first item of an enumerated
                  list that is nested within the itemized list.

            \item This is the second item of the inner list.
                  \LaTeX\ allows you to nest lists deeper than
                  you really should.
         \end{enumerate}
         This is the rest of the second item of the outer
         list.  It is no more interesting than any other
         part of the item.
   \item This is the third item of the list.
\end{itemize}
You can even display poetry.
\begin{verse}
   There is an environment
    for verse \\             % The \\ command separates lines
   Whose features some poets % within a stanza.
   will curse.

                             % One or more blank lines separate stanzas.

   For instead of making\\
   Them do \emph{all} line breaking, \\
   It allows them to put too many words on a line when they'd rather be
   forced to be terse.
\end{verse}

Mathematical formulas may also be displayed.  A
displayed formula
is
one-line long; multiline
formulas require special formatting instructions.
   \[  \ip{\Gamma}{\psi'} = x'' + y^{2} + z_{i}^{n}\]
Don't start a paragraph with a displayed equation,
nor make one a paragraph by itself.

\end{document}               % End of document.

\section{Bench Section 77}
% synthetic bench section 77
% This is a sample LaTeX input file.  (Version of 12 August 2004.)
%
% A '%' character causes TeX to ignore all remaining text on the line,
% and is used for comments like this one.

\documentclass{article}      % Specifies the document class

                             % The preamble begins here.
\title{An Example Document}  % Declares the document's title.
\author{Leslie Lamport}      % Declares the author's name.
\date{January 21, 1994}      % Deleting this command produces today's date.

\newcommand{\ip}[2]{(#1, #2)}
                             % Defines \ip{arg1}{arg2} to mean
                             % (arg1, arg2).

%\newcommand{\ip}[2]{\langle #1 | #2\rangle}
                             % This is an alternative definition of
                             % \ip that is commented out.

\begin{document}             % End of preamble and beginning of text.

\maketitle                   % Produces the title.

This is an example input file.  Comparing it with
the output it generates can show you how to
produce a simple document of your own.

\section{Ordinary Text}      % Produces section heading.  Lower-level
                             % sections are begun with similar
                             % \subsection and \subsubsection commands.

The ends  of words and sentences are marked
  by   spaces. It  doesn't matter how many
spaces    you type; one is as good as 100.  The
end of   a line counts as a space.

One   or more   blank lines denote the  end
of  a paragraph.

Since any number of consecutive spaces are treated
like a single one, the formatting of the input
file makes no difference to
      \LaTeX,                % The \LaTeX command generates the LaTeX logo.
but it makes a difference to you.  When you use
\LaTeX, making your input file as easy to read
as possible will be a great help as you write
your document and when you change it.  This sample
file shows how you can add comments to your own input
file.

Because printing is different from typewriting,
there are a number of things that you have to do
differently when preparing an input file than if
you were just typing the document directly.
Quotation marks like
       ``this''
have to be handled specially, as do quotes within
quotes:
       ``\,`this'            % \, separates the double and single quote.
        is what I just
        wrote, not  `that'\,''.

Dashes come in three sizes: an
       intra-word
dash, a medium dash for number ranges like
       1--2,
and a punctuation
       dash---like
this.

A sentence-ending space should be larger than the
space between words within a sentence.  You
sometimes have to type special commands in
conjunction with punctuation characters to get
this right, as in the following sentence.
       Gnats, gnus, etc.\ all  % `\ ' makes an inter-word space.
       begin with G\@.         % \@ marks end-of-sentence punctuation.
You should check the spaces after periods when
reading your output to make sure you haven't
forgotten any special cases.  Generating an
ellipsis
       \ldots\               % `\ ' is needed after `\ldots' because TeX
                             % ignores spaces after command names like \ldots
                             % made from \ + letters.
                             %
                             % Note how a `%' character causes TeX to ignore
                             % the end of the input line, so these blank lines
                             % do not start a new paragraph.
                             %
with the right spacing around the periods requires
a special command.

\LaTeX\ interprets some common characters as
commands, so you must type special commands to
generate them.  These characters include the
following:
       \$ \& \% \# \{ and \}.

In printing, text is usually emphasized with an
       \emph{italic}
type style.

\begin{em}
   A long segment of text can also be emphasized
   in this way.  Text within such a segment can be
   given \emph{additional} emphasis.
\end{em}

It is sometimes necessary to prevent \LaTeX\ from
breaking a line where it might otherwise do so.
This may be at a space, as between the ``Mr.''\ and
``Jones'' in
       ``Mr.~Jones'',        % ~ produces an unbreakable interword space.
or within a word---especially when the word is a
symbol like
       \mbox{\emph{itemnum}}
that makes little sense when hyphenated across
lines.

Footnotes\footnote{This is an example of a footnote.}
pose no problem.

\LaTeX\ is good at typesetting mathematical formulas
like
       \( x-3y + z = 7 \)
or
       \( a_{1} > x^{2n} + y^{2n} > x' \)
or
       \( \ip{A}{B} = \sum_{i} a_{i} b_{i} \).
The spaces you type in a formula are
ignored.  Remember that a letter like
       $x$                   % $ ... $  and  \( ... \)  are equivalent
is a formula when it denotes a mathematical
symbol, and it should be typed as one.

\section{Displayed Text}

Text is displayed by indenting it from the left
margin.  Quotations are commonly displayed.  There
are short quotations
\begin{quote}
   This is a short quotation.  It consists of a
   single paragraph of text.  See how it is formatted.
\end{quote}
and longer ones.
\begin{quotation}
   This is a longer quotation.  It consists of two
   paragraphs of text, neither of which are
   particularly interesting.

   This is the second paragraph of the quotation.  It
   is just as dull as the first paragraph.
\end{quotation}
Another frequently-displayed structure is a list.
The following is an example of an \emph{itemized}
list.
\begin{itemize}
   \item This is the first item of an itemized list.
         Each item in the list is marked with a ``tick''.
         You don't have to worry about what kind of tick
         mark is used.

   \item This is the second item of the list.  It
         contains another list nested inside it.  The inner
         list is an \emph{enumerated} list.
         \begin{enumerate}
            \item This is the first item of an enumerated
                  list that is nested within the itemized list.

            \item This is the second item of the inner list.
                  \LaTeX\ allows you to nest lists deeper than
                  you really should.
         \end{enumerate}
         This is the rest of the second item of the outer
         list.  It is no more interesting than any other
         part of the item.
   \item This is the third item of the list.
\end{itemize}
You can even display poetry.
\begin{verse}
   There is an environment
    for verse \\             % The \\ command separates lines
   Whose features some poets % within a stanza.
   will curse.

                             % One or more blank lines separate stanzas.

   For instead of making\\
   Them do \emph{all} line breaking, \\
   It allows them to put too many words on a line when they'd rather be
   forced to be terse.
\end{verse}

Mathematical formulas may also be displayed.  A
displayed formula
is
one-line long; multiline
formulas require special formatting instructions.
   \[  \ip{\Gamma}{\psi'} = x'' + y^{2} + z_{i}^{n}\]
Don't start a paragraph with a displayed equation,
nor make one a paragraph by itself.

\end{document}               % End of document.

\section{Bench Section 78}
% synthetic bench section 78
% This is a sample LaTeX input file.  (Version of 12 August 2004.)
%
% A '%' character causes TeX to ignore all remaining text on the line,
% and is used for comments like this one.

\documentclass{article}      % Specifies the document class

                             % The preamble begins here.
\title{An Example Document}  % Declares the document's title.
\author{Leslie Lamport}      % Declares the author's name.
\date{January 21, 1994}      % Deleting this command produces today's date.

\newcommand{\ip}[2]{(#1, #2)}
                             % Defines \ip{arg1}{arg2} to mean
                             % (arg1, arg2).

%\newcommand{\ip}[2]{\langle #1 | #2\rangle}
                             % This is an alternative definition of
                             % \ip that is commented out.

\begin{document}             % End of preamble and beginning of text.

\maketitle                   % Produces the title.

This is an example input file.  Comparing it with
the output it generates can show you how to
produce a simple document of your own.

\section{Ordinary Text}      % Produces section heading.  Lower-level
                             % sections are begun with similar
                             % \subsection and \subsubsection commands.

The ends  of words and sentences are marked
  by   spaces. It  doesn't matter how many
spaces    you type; one is as good as 100.  The
end of   a line counts as a space.

One   or more   blank lines denote the  end
of  a paragraph.

Since any number of consecutive spaces are treated
like a single one, the formatting of the input
file makes no difference to
      \LaTeX,                % The \LaTeX command generates the LaTeX logo.
but it makes a difference to you.  When you use
\LaTeX, making your input file as easy to read
as possible will be a great help as you write
your document and when you change it.  This sample
file shows how you can add comments to your own input
file.

Because printing is different from typewriting,
there are a number of things that you have to do
differently when preparing an input file than if
you were just typing the document directly.
Quotation marks like
       ``this''
have to be handled specially, as do quotes within
quotes:
       ``\,`this'            % \, separates the double and single quote.
        is what I just
        wrote, not  `that'\,''.

Dashes come in three sizes: an
       intra-word
dash, a medium dash for number ranges like
       1--2,
and a punctuation
       dash---like
this.

A sentence-ending space should be larger than the
space between words within a sentence.  You
sometimes have to type special commands in
conjunction with punctuation characters to get
this right, as in the following sentence.
       Gnats, gnus, etc.\ all  % `\ ' makes an inter-word space.
       begin with G\@.         % \@ marks end-of-sentence punctuation.
You should check the spaces after periods when
reading your output to make sure you haven't
forgotten any special cases.  Generating an
ellipsis
       \ldots\               % `\ ' is needed after `\ldots' because TeX
                             % ignores spaces after command names like \ldots
                             % made from \ + letters.
                             %
                             % Note how a `%' character causes TeX to ignore
                             % the end of the input line, so these blank lines
                             % do not start a new paragraph.
                             %
with the right spacing around the periods requires
a special command.

\LaTeX\ interprets some common characters as
commands, so you must type special commands to
generate them.  These characters include the
following:
       \$ \& \% \# \{ and \}.

In printing, text is usually emphasized with an
       \emph{italic}
type style.

\begin{em}
   A long segment of text can also be emphasized
   in this way.  Text within such a segment can be
   given \emph{additional} emphasis.
\end{em}

It is sometimes necessary to prevent \LaTeX\ from
breaking a line where it might otherwise do so.
This may be at a space, as between the ``Mr.''\ and
``Jones'' in
       ``Mr.~Jones'',        % ~ produces an unbreakable interword space.
or within a word---especially when the word is a
symbol like
       \mbox{\emph{itemnum}}
that makes little sense when hyphenated across
lines.

Footnotes\footnote{This is an example of a footnote.}
pose no problem.

\LaTeX\ is good at typesetting mathematical formulas
like
       \( x-3y + z = 7 \)
or
       \( a_{1} > x^{2n} + y^{2n} > x' \)
or
       \( \ip{A}{B} = \sum_{i} a_{i} b_{i} \).
The spaces you type in a formula are
ignored.  Remember that a letter like
       $x$                   % $ ... $  and  \( ... \)  are equivalent
is a formula when it denotes a mathematical
symbol, and it should be typed as one.

\section{Displayed Text}

Text is displayed by indenting it from the left
margin.  Quotations are commonly displayed.  There
are short quotations
\begin{quote}
   This is a short quotation.  It consists of a
   single paragraph of text.  See how it is formatted.
\end{quote}
and longer ones.
\begin{quotation}
   This is a longer quotation.  It consists of two
   paragraphs of text, neither of which are
   particularly interesting.

   This is the second paragraph of the quotation.  It
   is just as dull as the first paragraph.
\end{quotation}
Another frequently-displayed structure is a list.
The following is an example of an \emph{itemized}
list.
\begin{itemize}
   \item This is the first item of an itemized list.
         Each item in the list is marked with a ``tick''.
         You don't have to worry about what kind of tick
         mark is used.

   \item This is the second item of the list.  It
         contains another list nested inside it.  The inner
         list is an \emph{enumerated} list.
         \begin{enumerate}
            \item This is the first item of an enumerated
                  list that is nested within the itemized list.

            \item This is the second item of the inner list.
                  \LaTeX\ allows you to nest lists deeper than
                  you really should.
         \end{enumerate}
         This is the rest of the second item of the outer
         list.  It is no more interesting than any other
         part of the item.
   \item This is the third item of the list.
\end{itemize}
You can even display poetry.
\begin{verse}
   There is an environment
    for verse \\             % The \\ command separates lines
   Whose features some poets % within a stanza.
   will curse.

                             % One or more blank lines separate stanzas.

   For instead of making\\
   Them do \emph{all} line breaking, \\
   It allows them to put too many words on a line when they'd rather be
   forced to be terse.
\end{verse}

Mathematical formulas may also be displayed.  A
displayed formula
is
one-line long; multiline
formulas require special formatting instructions.
   \[  \ip{\Gamma}{\psi'} = x'' + y^{2} + z_{i}^{n}\]
Don't start a paragraph with a displayed equation,
nor make one a paragraph by itself.

\end{document}               % End of document.

\section{Bench Section 79}
% synthetic bench section 79
% This is a sample LaTeX input file.  (Version of 12 August 2004.)
%
% A '%' character causes TeX to ignore all remaining text on the line,
% and is used for comments like this one.

\documentclass{article}      % Specifies the document class

                             % The preamble begins here.
\title{An Example Document}  % Declares the document's title.
\author{Leslie Lamport}      % Declares the author's name.
\date{January 21, 1994}      % Deleting this command produces today's date.

\newcommand{\ip}[2]{(#1, #2)}
                             % Defines \ip{arg1}{arg2} to mean
                             % (arg1, arg2).

%\newcommand{\ip}[2]{\langle #1 | #2\rangle}
                             % This is an alternative definition of
                             % \ip that is commented out.

\begin{document}             % End of preamble and beginning of text.

\maketitle                   % Produces the title.

This is an example input file.  Comparing it with
the output it generates can show you how to
produce a simple document of your own.

\section{Ordinary Text}      % Produces section heading.  Lower-level
                             % sections are begun with similar
                             % \subsection and \subsubsection commands.

The ends  of words and sentences are marked
  by   spaces. It  doesn't matter how many
spaces    you type; one is as good as 100.  The
end of   a line counts as a space.

One   or more   blank lines denote the  end
of  a paragraph.

Since any number of consecutive spaces are treated
like a single one, the formatting of the input
file makes no difference to
      \LaTeX,                % The \LaTeX command generates the LaTeX logo.
but it makes a difference to you.  When you use
\LaTeX, making your input file as easy to read
as possible will be a great help as you write
your document and when you change it.  This sample
file shows how you can add comments to your own input
file.

Because printing is different from typewriting,
there are a number of things that you have to do
differently when preparing an input file than if
you were just typing the document directly.
Quotation marks like
       ``this''
have to be handled specially, as do quotes within
quotes:
       ``\,`this'            % \, separates the double and single quote.
        is what I just
        wrote, not  `that'\,''.

Dashes come in three sizes: an
       intra-word
dash, a medium dash for number ranges like
       1--2,
and a punctuation
       dash---like
this.

A sentence-ending space should be larger than the
space between words within a sentence.  You
sometimes have to type special commands in
conjunction with punctuation characters to get
this right, as in the following sentence.
       Gnats, gnus, etc.\ all  % `\ ' makes an inter-word space.
       begin with G\@.         % \@ marks end-of-sentence punctuation.
You should check the spaces after periods when
reading your output to make sure you haven't
forgotten any special cases.  Generating an
ellipsis
       \ldots\               % `\ ' is needed after `\ldots' because TeX
                             % ignores spaces after command names like \ldots
                             % made from \ + letters.
                             %
                             % Note how a `%' character causes TeX to ignore
                             % the end of the input line, so these blank lines
                             % do not start a new paragraph.
                             %
with the right spacing around the periods requires
a special command.

\LaTeX\ interprets some common characters as
commands, so you must type special commands to
generate them.  These characters include the
following:
       \$ \& \% \# \{ and \}.

In printing, text is usually emphasized with an
       \emph{italic}
type style.

\begin{em}
   A long segment of text can also be emphasized
   in this way.  Text within such a segment can be
   given \emph{additional} emphasis.
\end{em}

It is sometimes necessary to prevent \LaTeX\ from
breaking a line where it might otherwise do so.
This may be at a space, as between the ``Mr.''\ and
``Jones'' in
       ``Mr.~Jones'',        % ~ produces an unbreakable interword space.
or within a word---especially when the word is a
symbol like
       \mbox{\emph{itemnum}}
that makes little sense when hyphenated across
lines.

Footnotes\footnote{This is an example of a footnote.}
pose no problem.

\LaTeX\ is good at typesetting mathematical formulas
like
       \( x-3y + z = 7 \)
or
       \( a_{1} > x^{2n} + y^{2n} > x' \)
or
       \( \ip{A}{B} = \sum_{i} a_{i} b_{i} \).
The spaces you type in a formula are
ignored.  Remember that a letter like
       $x$                   % $ ... $  and  \( ... \)  are equivalent
is a formula when it denotes a mathematical
symbol, and it should be typed as one.

\section{Displayed Text}

Text is displayed by indenting it from the left
margin.  Quotations are commonly displayed.  There
are short quotations
\begin{quote}
   This is a short quotation.  It consists of a
   single paragraph of text.  See how it is formatted.
\end{quote}
and longer ones.
\begin{quotation}
   This is a longer quotation.  It consists of two
   paragraphs of text, neither of which are
   particularly interesting.

   This is the second paragraph of the quotation.  It
   is just as dull as the first paragraph.
\end{quotation}
Another frequently-displayed structure is a list.
The following is an example of an \emph{itemized}
list.
\begin{itemize}
   \item This is the first item of an itemized list.
         Each item in the list is marked with a ``tick''.
         You don't have to worry about what kind of tick
         mark is used.

   \item This is the second item of the list.  It
         contains another list nested inside it.  The inner
         list is an \emph{enumerated} list.
         \begin{enumerate}
            \item This is the first item of an enumerated
                  list that is nested within the itemized list.

            \item This is the second item of the inner list.
                  \LaTeX\ allows you to nest lists deeper than
                  you really should.
         \end{enumerate}
         This is the rest of the second item of the outer
         list.  It is no more interesting than any other
         part of the item.
   \item This is the third item of the list.
\end{itemize}
You can even display poetry.
\begin{verse}
   There is an environment
    for verse \\             % The \\ command separates lines
   Whose features some poets % within a stanza.
   will curse.

                             % One or more blank lines separate stanzas.

   For instead of making\\
   Them do \emph{all} line breaking, \\
   It allows them to put too many words on a line when they'd rather be
   forced to be terse.
\end{verse}

Mathematical formulas may also be displayed.  A
displayed formula
is
one-line long; multiline
formulas require special formatting instructions.
   \[  \ip{\Gamma}{\psi'} = x'' + y^{2} + z_{i}^{n}\]
Don't start a paragraph with a displayed equation,
nor make one a paragraph by itself.

\end{document}               % End of document.

\section{Bench Section 80}
% synthetic bench section 80
% This is a sample LaTeX input file.  (Version of 12 August 2004.)
%
% A '%' character causes TeX to ignore all remaining text on the line,
% and is used for comments like this one.

\documentclass{article}      % Specifies the document class

                             % The preamble begins here.
\title{An Example Document}  % Declares the document's title.
\author{Leslie Lamport}      % Declares the author's name.
\date{January 21, 1994}      % Deleting this command produces today's date.

\newcommand{\ip}[2]{(#1, #2)}
                             % Defines \ip{arg1}{arg2} to mean
                             % (arg1, arg2).

%\newcommand{\ip}[2]{\langle #1 | #2\rangle}
                             % This is an alternative definition of
                             % \ip that is commented out.

\begin{document}             % End of preamble and beginning of text.

\maketitle                   % Produces the title.

This is an example input file.  Comparing it with
the output it generates can show you how to
produce a simple document of your own.

\section{Ordinary Text}      % Produces section heading.  Lower-level
                             % sections are begun with similar
                             % \subsection and \subsubsection commands.

The ends  of words and sentences are marked
  by   spaces. It  doesn't matter how many
spaces    you type; one is as good as 100.  The
end of   a line counts as a space.

One   or more   blank lines denote the  end
of  a paragraph.

Since any number of consecutive spaces are treated
like a single one, the formatting of the input
file makes no difference to
      \LaTeX,                % The \LaTeX command generates the LaTeX logo.
but it makes a difference to you.  When you use
\LaTeX, making your input file as easy to read
as possible will be a great help as you write
your document and when you change it.  This sample
file shows how you can add comments to your own input
file.

Because printing is different from typewriting,
there are a number of things that you have to do
differently when preparing an input file than if
you were just typing the document directly.
Quotation marks like
       ``this''
have to be handled specially, as do quotes within
quotes:
       ``\,`this'            % \, separates the double and single quote.
        is what I just
        wrote, not  `that'\,''.

Dashes come in three sizes: an
       intra-word
dash, a medium dash for number ranges like
       1--2,
and a punctuation
       dash---like
this.

A sentence-ending space should be larger than the
space between words within a sentence.  You
sometimes have to type special commands in
conjunction with punctuation characters to get
this right, as in the following sentence.
       Gnats, gnus, etc.\ all  % `\ ' makes an inter-word space.
       begin with G\@.         % \@ marks end-of-sentence punctuation.
You should check the spaces after periods when
reading your output to make sure you haven't
forgotten any special cases.  Generating an
ellipsis
       \ldots\               % `\ ' is needed after `\ldots' because TeX
                             % ignores spaces after command names like \ldots
                             % made from \ + letters.
                             %
                             % Note how a `%' character causes TeX to ignore
                             % the end of the input line, so these blank lines
                             % do not start a new paragraph.
                             %
with the right spacing around the periods requires
a special command.

\LaTeX\ interprets some common characters as
commands, so you must type special commands to
generate them.  These characters include the
following:
       \$ \& \% \# \{ and \}.

In printing, text is usually emphasized with an
       \emph{italic}
type style.

\begin{em}
   A long segment of text can also be emphasized
   in this way.  Text within such a segment can be
   given \emph{additional} emphasis.
\end{em}

It is sometimes necessary to prevent \LaTeX\ from
breaking a line where it might otherwise do so.
This may be at a space, as between the ``Mr.''\ and
``Jones'' in
       ``Mr.~Jones'',        % ~ produces an unbreakable interword space.
or within a word---especially when the word is a
symbol like
       \mbox{\emph{itemnum}}
that makes little sense when hyphenated across
lines.

Footnotes\footnote{This is an example of a footnote.}
pose no problem.

\LaTeX\ is good at typesetting mathematical formulas
like
       \( x-3y + z = 7 \)
or
       \( a_{1} > x^{2n} + y^{2n} > x' \)
or
       \( \ip{A}{B} = \sum_{i} a_{i} b_{i} \).
The spaces you type in a formula are
ignored.  Remember that a letter like
       $x$                   % $ ... $  and  \( ... \)  are equivalent
is a formula when it denotes a mathematical
symbol, and it should be typed as one.

\section{Displayed Text}

Text is displayed by indenting it from the left
margin.  Quotations are commonly displayed.  There
are short quotations
\begin{quote}
   This is a short quotation.  It consists of a
   single paragraph of text.  See how it is formatted.
\end{quote}
and longer ones.
\begin{quotation}
   This is a longer quotation.  It consists of two
   paragraphs of text, neither of which are
   particularly interesting.

   This is the second paragraph of the quotation.  It
   is just as dull as the first paragraph.
\end{quotation}
Another frequently-displayed structure is a list.
The following is an example of an \emph{itemized}
list.
\begin{itemize}
   \item This is the first item of an itemized list.
         Each item in the list is marked with a ``tick''.
         You don't have to worry about what kind of tick
         mark is used.

   \item This is the second item of the list.  It
         contains another list nested inside it.  The inner
         list is an \emph{enumerated} list.
         \begin{enumerate}
            \item This is the first item of an enumerated
                  list that is nested within the itemized list.

            \item This is the second item of the inner list.
                  \LaTeX\ allows you to nest lists deeper than
                  you really should.
         \end{enumerate}
         This is the rest of the second item of the outer
         list.  It is no more interesting than any other
         part of the item.
   \item This is the third item of the list.
\end{itemize}
You can even display poetry.
\begin{verse}
   There is an environment
    for verse \\             % The \\ command separates lines
   Whose features some poets % within a stanza.
   will curse.

                             % One or more blank lines separate stanzas.

   For instead of making\\
   Them do \emph{all} line breaking, \\
   It allows them to put too many words on a line when they'd rather be
   forced to be terse.
\end{verse}

Mathematical formulas may also be displayed.  A
displayed formula
is
one-line long; multiline
formulas require special formatting instructions.
   \[  \ip{\Gamma}{\psi'} = x'' + y^{2} + z_{i}^{n}\]
Don't start a paragraph with a displayed equation,
nor make one a paragraph by itself.

\end{document}               % End of document.

\section{Bench Section 81}
% synthetic bench section 81
% This is a sample LaTeX input file.  (Version of 12 August 2004.)
%
% A '%' character causes TeX to ignore all remaining text on the line,
% and is used for comments like this one.

\documentclass{article}      % Specifies the document class

                             % The preamble begins here.
\title{An Example Document}  % Declares the document's title.
\author{Leslie Lamport}      % Declares the author's name.
\date{January 21, 1994}      % Deleting this command produces today's date.

\newcommand{\ip}[2]{(#1, #2)}
                             % Defines \ip{arg1}{arg2} to mean
                             % (arg1, arg2).

%\newcommand{\ip}[2]{\langle #1 | #2\rangle}
                             % This is an alternative definition of
                             % \ip that is commented out.

\begin{document}             % End of preamble and beginning of text.

\maketitle                   % Produces the title.

This is an example input file.  Comparing it with
the output it generates can show you how to
produce a simple document of your own.

\section{Ordinary Text}      % Produces section heading.  Lower-level
                             % sections are begun with similar
                             % \subsection and \subsubsection commands.

The ends  of words and sentences are marked
  by   spaces. It  doesn't matter how many
spaces    you type; one is as good as 100.  The
end of   a line counts as a space.

One   or more   blank lines denote the  end
of  a paragraph.

Since any number of consecutive spaces are treated
like a single one, the formatting of the input
file makes no difference to
      \LaTeX,                % The \LaTeX command generates the LaTeX logo.
but it makes a difference to you.  When you use
\LaTeX, making your input file as easy to read
as possible will be a great help as you write
your document and when you change it.  This sample
file shows how you can add comments to your own input
file.

Because printing is different from typewriting,
there are a number of things that you have to do
differently when preparing an input file than if
you were just typing the document directly.
Quotation marks like
       ``this''
have to be handled specially, as do quotes within
quotes:
       ``\,`this'            % \, separates the double and single quote.
        is what I just
        wrote, not  `that'\,''.

Dashes come in three sizes: an
       intra-word
dash, a medium dash for number ranges like
       1--2,
and a punctuation
       dash---like
this.

A sentence-ending space should be larger than the
space between words within a sentence.  You
sometimes have to type special commands in
conjunction with punctuation characters to get
this right, as in the following sentence.
       Gnats, gnus, etc.\ all  % `\ ' makes an inter-word space.
       begin with G\@.         % \@ marks end-of-sentence punctuation.
You should check the spaces after periods when
reading your output to make sure you haven't
forgotten any special cases.  Generating an
ellipsis
       \ldots\               % `\ ' is needed after `\ldots' because TeX
                             % ignores spaces after command names like \ldots
                             % made from \ + letters.
                             %
                             % Note how a `%' character causes TeX to ignore
                             % the end of the input line, so these blank lines
                             % do not start a new paragraph.
                             %
with the right spacing around the periods requires
a special command.

\LaTeX\ interprets some common characters as
commands, so you must type special commands to
generate them.  These characters include the
following:
       \$ \& \% \# \{ and \}.

In printing, text is usually emphasized with an
       \emph{italic}
type style.

\begin{em}
   A long segment of text can also be emphasized
   in this way.  Text within such a segment can be
   given \emph{additional} emphasis.
\end{em}

It is sometimes necessary to prevent \LaTeX\ from
breaking a line where it might otherwise do so.
This may be at a space, as between the ``Mr.''\ and
``Jones'' in
       ``Mr.~Jones'',        % ~ produces an unbreakable interword space.
or within a word---especially when the word is a
symbol like
       \mbox{\emph{itemnum}}
that makes little sense when hyphenated across
lines.

Footnotes\footnote{This is an example of a footnote.}
pose no problem.

\LaTeX\ is good at typesetting mathematical formulas
like
       \( x-3y + z = 7 \)
or
       \( a_{1} > x^{2n} + y^{2n} > x' \)
or
       \( \ip{A}{B} = \sum_{i} a_{i} b_{i} \).
The spaces you type in a formula are
ignored.  Remember that a letter like
       $x$                   % $ ... $  and  \( ... \)  are equivalent
is a formula when it denotes a mathematical
symbol, and it should be typed as one.

\section{Displayed Text}

Text is displayed by indenting it from the left
margin.  Quotations are commonly displayed.  There
are short quotations
\begin{quote}
   This is a short quotation.  It consists of a
   single paragraph of text.  See how it is formatted.
\end{quote}
and longer ones.
\begin{quotation}
   This is a longer quotation.  It consists of two
   paragraphs of text, neither of which are
   particularly interesting.

   This is the second paragraph of the quotation.  It
   is just as dull as the first paragraph.
\end{quotation}
Another frequently-displayed structure is a list.
The following is an example of an \emph{itemized}
list.
\begin{itemize}
   \item This is the first item of an itemized list.
         Each item in the list is marked with a ``tick''.
         You don't have to worry about what kind of tick
         mark is used.

   \item This is the second item of the list.  It
         contains another list nested inside it.  The inner
         list is an \emph{enumerated} list.
         \begin{enumerate}
            \item This is the first item of an enumerated
                  list that is nested within the itemized list.

            \item This is the second item of the inner list.
                  \LaTeX\ allows you to nest lists deeper than
                  you really should.
         \end{enumerate}
         This is the rest of the second item of the outer
         list.  It is no more interesting than any other
         part of the item.
   \item This is the third item of the list.
\end{itemize}
You can even display poetry.
\begin{verse}
   There is an environment
    for verse \\             % The \\ command separates lines
   Whose features some poets % within a stanza.
   will curse.

                             % One or more blank lines separate stanzas.

   For instead of making\\
   Them do \emph{all} line breaking, \\
   It allows them to put too many words on a line when they'd rather be
   forced to be terse.
\end{verse}

Mathematical formulas may also be displayed.  A
displayed formula
is
one-line long; multiline
formulas require special formatting instructions.
   \[  \ip{\Gamma}{\psi'} = x'' + y^{2} + z_{i}^{n}\]
Don't start a paragraph with a displayed equation,
nor make one a paragraph by itself.

\end{document}               % End of document.

\section{Bench Section 82}
% synthetic bench section 82
% This is a sample LaTeX input file.  (Version of 12 August 2004.)
%
% A '%' character causes TeX to ignore all remaining text on the line,
% and is used for comments like this one.

\documentclass{article}      % Specifies the document class

                             % The preamble begins here.
\title{An Example Document}  % Declares the document's title.
\author{Leslie Lamport}      % Declares the author's name.
\date{January 21, 1994}      % Deleting this command produces today's date.

\newcommand{\ip}[2]{(#1, #2)}
                             % Defines \ip{arg1}{arg2} to mean
                             % (arg1, arg2).

%\newcommand{\ip}[2]{\langle #1 | #2\rangle}
                             % This is an alternative definition of
                             % \ip that is commented out.

\begin{document}             % End of preamble and beginning of text.

\maketitle                   % Produces the title.

This is an example input file.  Comparing it with
the output it generates can show you how to
produce a simple document of your own.

\section{Ordinary Text}      % Produces section heading.  Lower-level
                             % sections are begun with similar
                             % \subsection and \subsubsection commands.

The ends  of words and sentences are marked
  by   spaces. It  doesn't matter how many
spaces    you type; one is as good as 100.  The
end of   a line counts as a space.

One   or more   blank lines denote the  end
of  a paragraph.

Since any number of consecutive spaces are treated
like a single one, the formatting of the input
file makes no difference to
      \LaTeX,                % The \LaTeX command generates the LaTeX logo.
but it makes a difference to you.  When you use
\LaTeX, making your input file as easy to read
as possible will be a great help as you write
your document and when you change it.  This sample
file shows how you can add comments to your own input
file.

Because printing is different from typewriting,
there are a number of things that you have to do
differently when preparing an input file than if
you were just typing the document directly.
Quotation marks like
       ``this''
have to be handled specially, as do quotes within
quotes:
       ``\,`this'            % \, separates the double and single quote.
        is what I just
        wrote, not  `that'\,''.

Dashes come in three sizes: an
       intra-word
dash, a medium dash for number ranges like
       1--2,
and a punctuation
       dash---like
this.

A sentence-ending space should be larger than the
space between words within a sentence.  You
sometimes have to type special commands in
conjunction with punctuation characters to get
this right, as in the following sentence.
       Gnats, gnus, etc.\ all  % `\ ' makes an inter-word space.
       begin with G\@.         % \@ marks end-of-sentence punctuation.
You should check the spaces after periods when
reading your output to make sure you haven't
forgotten any special cases.  Generating an
ellipsis
       \ldots\               % `\ ' is needed after `\ldots' because TeX
                             % ignores spaces after command names like \ldots
                             % made from \ + letters.
                             %
                             % Note how a `%' character causes TeX to ignore
                             % the end of the input line, so these blank lines
                             % do not start a new paragraph.
                             %
with the right spacing around the periods requires
a special command.

\LaTeX\ interprets some common characters as
commands, so you must type special commands to
generate them.  These characters include the
following:
       \$ \& \% \# \{ and \}.

In printing, text is usually emphasized with an
       \emph{italic}
type style.

\begin{em}
   A long segment of text can also be emphasized
   in this way.  Text within such a segment can be
   given \emph{additional} emphasis.
\end{em}

It is sometimes necessary to prevent \LaTeX\ from
breaking a line where it might otherwise do so.
This may be at a space, as between the ``Mr.''\ and
``Jones'' in
       ``Mr.~Jones'',        % ~ produces an unbreakable interword space.
or within a word---especially when the word is a
symbol like
       \mbox{\emph{itemnum}}
that makes little sense when hyphenated across
lines.

Footnotes\footnote{This is an example of a footnote.}
pose no problem.

\LaTeX\ is good at typesetting mathematical formulas
like
       \( x-3y + z = 7 \)
or
       \( a_{1} > x^{2n} + y^{2n} > x' \)
or
       \( \ip{A}{B} = \sum_{i} a_{i} b_{i} \).
The spaces you type in a formula are
ignored.  Remember that a letter like
       $x$                   % $ ... $  and  \( ... \)  are equivalent
is a formula when it denotes a mathematical
symbol, and it should be typed as one.

\section{Displayed Text}

Text is displayed by indenting it from the left
margin.  Quotations are commonly displayed.  There
are short quotations
\begin{quote}
   This is a short quotation.  It consists of a
   single paragraph of text.  See how it is formatted.
\end{quote}
and longer ones.
\begin{quotation}
   This is a longer quotation.  It consists of two
   paragraphs of text, neither of which are
   particularly interesting.

   This is the second paragraph of the quotation.  It
   is just as dull as the first paragraph.
\end{quotation}
Another frequently-displayed structure is a list.
The following is an example of an \emph{itemized}
list.
\begin{itemize}
   \item This is the first item of an itemized list.
         Each item in the list is marked with a ``tick''.
         You don't have to worry about what kind of tick
         mark is used.

   \item This is the second item of the list.  It
         contains another list nested inside it.  The inner
         list is an \emph{enumerated} list.
         \begin{enumerate}
            \item This is the first item of an enumerated
                  list that is nested within the itemized list.

            \item This is the second item of the inner list.
                  \LaTeX\ allows you to nest lists deeper than
                  you really should.
         \end{enumerate}
         This is the rest of the second item of the outer
         list.  It is no more interesting than any other
         part of the item.
   \item This is the third item of the list.
\end{itemize}
You can even display poetry.
\begin{verse}
   There is an environment
    for verse \\             % The \\ command separates lines
   Whose features some poets % within a stanza.
   will curse.

                             % One or more blank lines separate stanzas.

   For instead of making\\
   Them do \emph{all} line breaking, \\
   It allows them to put too many words on a line when they'd rather be
   forced to be terse.
\end{verse}

Mathematical formulas may also be displayed.  A
displayed formula
is
one-line long; multiline
formulas require special formatting instructions.
   \[  \ip{\Gamma}{\psi'} = x'' + y^{2} + z_{i}^{n}\]
Don't start a paragraph with a displayed equation,
nor make one a paragraph by itself.

\end{document}               % End of document.

\section{Bench Section 83}
% synthetic bench section 83
% This is a sample LaTeX input file.  (Version of 12 August 2004.)
%
% A '%' character causes TeX to ignore all remaining text on the line,
% and is used for comments like this one.

\documentclass{article}      % Specifies the document class

                             % The preamble begins here.
\title{An Example Document}  % Declares the document's title.
\author{Leslie Lamport}      % Declares the author's name.
\date{January 21, 1994}      % Deleting this command produces today's date.

\newcommand{\ip}[2]{(#1, #2)}
                             % Defines \ip{arg1}{arg2} to mean
                             % (arg1, arg2).

%\newcommand{\ip}[2]{\langle #1 | #2\rangle}
                             % This is an alternative definition of
                             % \ip that is commented out.

\begin{document}             % End of preamble and beginning of text.

\maketitle                   % Produces the title.

This is an example input file.  Comparing it with
the output it generates can show you how to
produce a simple document of your own.

\section{Ordinary Text}      % Produces section heading.  Lower-level
                             % sections are begun with similar
                             % \subsection and \subsubsection commands.

The ends  of words and sentences are marked
  by   spaces. It  doesn't matter how many
spaces    you type; one is as good as 100.  The
end of   a line counts as a space.

One   or more   blank lines denote the  end
of  a paragraph.

Since any number of consecutive spaces are treated
like a single one, the formatting of the input
file makes no difference to
      \LaTeX,                % The \LaTeX command generates the LaTeX logo.
but it makes a difference to you.  When you use
\LaTeX, making your input file as easy to read
as possible will be a great help as you write
your document and when you change it.  This sample
file shows how you can add comments to your own input
file.

Because printing is different from typewriting,
there are a number of things that you have to do
differently when preparing an input file than if
you were just typing the document directly.
Quotation marks like
       ``this''
have to be handled specially, as do quotes within
quotes:
       ``\,`this'            % \, separates the double and single quote.
        is what I just
        wrote, not  `that'\,''.

Dashes come in three sizes: an
       intra-word
dash, a medium dash for number ranges like
       1--2,
and a punctuation
       dash---like
this.

A sentence-ending space should be larger than the
space between words within a sentence.  You
sometimes have to type special commands in
conjunction with punctuation characters to get
this right, as in the following sentence.
       Gnats, gnus, etc.\ all  % `\ ' makes an inter-word space.
       begin with G\@.         % \@ marks end-of-sentence punctuation.
You should check the spaces after periods when
reading your output to make sure you haven't
forgotten any special cases.  Generating an
ellipsis
       \ldots\               % `\ ' is needed after `\ldots' because TeX
                             % ignores spaces after command names like \ldots
                             % made from \ + letters.
                             %
                             % Note how a `%' character causes TeX to ignore
                             % the end of the input line, so these blank lines
                             % do not start a new paragraph.
                             %
with the right spacing around the periods requires
a special command.

\LaTeX\ interprets some common characters as
commands, so you must type special commands to
generate them.  These characters include the
following:
       \$ \& \% \# \{ and \}.

In printing, text is usually emphasized with an
       \emph{italic}
type style.

\begin{em}
   A long segment of text can also be emphasized
   in this way.  Text within such a segment can be
   given \emph{additional} emphasis.
\end{em}

It is sometimes necessary to prevent \LaTeX\ from
breaking a line where it might otherwise do so.
This may be at a space, as between the ``Mr.''\ and
``Jones'' in
       ``Mr.~Jones'',        % ~ produces an unbreakable interword space.
or within a word---especially when the word is a
symbol like
       \mbox{\emph{itemnum}}
that makes little sense when hyphenated across
lines.

Footnotes\footnote{This is an example of a footnote.}
pose no problem.

\LaTeX\ is good at typesetting mathematical formulas
like
       \( x-3y + z = 7 \)
or
       \( a_{1} > x^{2n} + y^{2n} > x' \)
or
       \( \ip{A}{B} = \sum_{i} a_{i} b_{i} \).
The spaces you type in a formula are
ignored.  Remember that a letter like
       $x$                   % $ ... $  and  \( ... \)  are equivalent
is a formula when it denotes a mathematical
symbol, and it should be typed as one.

\section{Displayed Text}

Text is displayed by indenting it from the left
margin.  Quotations are commonly displayed.  There
are short quotations
\begin{quote}
   This is a short quotation.  It consists of a
   single paragraph of text.  See how it is formatted.
\end{quote}
and longer ones.
\begin{quotation}
   This is a longer quotation.  It consists of two
   paragraphs of text, neither of which are
   particularly interesting.

   This is the second paragraph of the quotation.  It
   is just as dull as the first paragraph.
\end{quotation}
Another frequently-displayed structure is a list.
The following is an example of an \emph{itemized}
list.
\begin{itemize}
   \item This is the first item of an itemized list.
         Each item in the list is marked with a ``tick''.
         You don't have to worry about what kind of tick
         mark is used.

   \item This is the second item of the list.  It
         contains another list nested inside it.  The inner
         list is an \emph{enumerated} list.
         \begin{enumerate}
            \item This is the first item of an enumerated
                  list that is nested within the itemized list.

            \item This is the second item of the inner list.
                  \LaTeX\ allows you to nest lists deeper than
                  you really should.
         \end{enumerate}
         This is the rest of the second item of the outer
         list.  It is no more interesting than any other
         part of the item.
   \item This is the third item of the list.
\end{itemize}
You can even display poetry.
\begin{verse}
   There is an environment
    for verse \\             % The \\ command separates lines
   Whose features some poets % within a stanza.
   will curse.

                             % One or more blank lines separate stanzas.

   For instead of making\\
   Them do \emph{all} line breaking, \\
   It allows them to put too many words on a line when they'd rather be
   forced to be terse.
\end{verse}

Mathematical formulas may also be displayed.  A
displayed formula
is
one-line long; multiline
formulas require special formatting instructions.
   \[  \ip{\Gamma}{\psi'} = x'' + y^{2} + z_{i}^{n}\]
Don't start a paragraph with a displayed equation,
nor make one a paragraph by itself.

\end{document}               % End of document.

\section{Bench Section 84}
% synthetic bench section 84
% This is a sample LaTeX input file.  (Version of 12 August 2004.)
%
% A '%' character causes TeX to ignore all remaining text on the line,
% and is used for comments like this one.

\documentclass{article}      % Specifies the document class

                             % The preamble begins here.
\title{An Example Document}  % Declares the document's title.
\author{Leslie Lamport}      % Declares the author's name.
\date{January 21, 1994}      % Deleting this command produces today's date.

\newcommand{\ip}[2]{(#1, #2)}
                             % Defines \ip{arg1}{arg2} to mean
                             % (arg1, arg2).

%\newcommand{\ip}[2]{\langle #1 | #2\rangle}
                             % This is an alternative definition of
                             % \ip that is commented out.

\begin{document}             % End of preamble and beginning of text.

\maketitle                   % Produces the title.

This is an example input file.  Comparing it with
the output it generates can show you how to
produce a simple document of your own.

\section{Ordinary Text}      % Produces section heading.  Lower-level
                             % sections are begun with similar
                             % \subsection and \subsubsection commands.

The ends  of words and sentences are marked
  by   spaces. It  doesn't matter how many
spaces    you type; one is as good as 100.  The
end of   a line counts as a space.

One   or more   blank lines denote the  end
of  a paragraph.

Since any number of consecutive spaces are treated
like a single one, the formatting of the input
file makes no difference to
      \LaTeX,                % The \LaTeX command generates the LaTeX logo.
but it makes a difference to you.  When you use
\LaTeX, making your input file as easy to read
as possible will be a great help as you write
your document and when you change it.  This sample
file shows how you can add comments to your own input
file.

Because printing is different from typewriting,
there are a number of things that you have to do
differently when preparing an input file than if
you were just typing the document directly.
Quotation marks like
       ``this''
have to be handled specially, as do quotes within
quotes:
       ``\,`this'            % \, separates the double and single quote.
        is what I just
        wrote, not  `that'\,''.

Dashes come in three sizes: an
       intra-word
dash, a medium dash for number ranges like
       1--2,
and a punctuation
       dash---like
this.

A sentence-ending space should be larger than the
space between words within a sentence.  You
sometimes have to type special commands in
conjunction with punctuation characters to get
this right, as in the following sentence.
       Gnats, gnus, etc.\ all  % `\ ' makes an inter-word space.
       begin with G\@.         % \@ marks end-of-sentence punctuation.
You should check the spaces after periods when
reading your output to make sure you haven't
forgotten any special cases.  Generating an
ellipsis
       \ldots\               % `\ ' is needed after `\ldots' because TeX
                             % ignores spaces after command names like \ldots
                             % made from \ + letters.
                             %
                             % Note how a `%' character causes TeX to ignore
                             % the end of the input line, so these blank lines
                             % do not start a new paragraph.
                             %
with the right spacing around the periods requires
a special command.

\LaTeX\ interprets some common characters as
commands, so you must type special commands to
generate them.  These characters include the
following:
       \$ \& \% \# \{ and \}.

In printing, text is usually emphasized with an
       \emph{italic}
type style.

\begin{em}
   A long segment of text can also be emphasized
   in this way.  Text within such a segment can be
   given \emph{additional} emphasis.
\end{em}

It is sometimes necessary to prevent \LaTeX\ from
breaking a line where it might otherwise do so.
This may be at a space, as between the ``Mr.''\ and
``Jones'' in
       ``Mr.~Jones'',        % ~ produces an unbreakable interword space.
or within a word---especially when the word is a
symbol like
       \mbox{\emph{itemnum}}
that makes little sense when hyphenated across
lines.

Footnotes\footnote{This is an example of a footnote.}
pose no problem.

\LaTeX\ is good at typesetting mathematical formulas
like
       \( x-3y + z = 7 \)
or
       \( a_{1} > x^{2n} + y^{2n} > x' \)
or
       \( \ip{A}{B} = \sum_{i} a_{i} b_{i} \).
The spaces you type in a formula are
ignored.  Remember that a letter like
       $x$                   % $ ... $  and  \( ... \)  are equivalent
is a formula when it denotes a mathematical
symbol, and it should be typed as one.

\section{Displayed Text}

Text is displayed by indenting it from the left
margin.  Quotations are commonly displayed.  There
are short quotations
\begin{quote}
   This is a short quotation.  It consists of a
   single paragraph of text.  See how it is formatted.
\end{quote}
and longer ones.
\begin{quotation}
   This is a longer quotation.  It consists of two
   paragraphs of text, neither of which are
   particularly interesting.

   This is the second paragraph of the quotation.  It
   is just as dull as the first paragraph.
\end{quotation}
Another frequently-displayed structure is a list.
The following is an example of an \emph{itemized}
list.
\begin{itemize}
   \item This is the first item of an itemized list.
         Each item in the list is marked with a ``tick''.
         You don't have to worry about what kind of tick
         mark is used.

   \item This is the second item of the list.  It
         contains another list nested inside it.  The inner
         list is an \emph{enumerated} list.
         \begin{enumerate}
            \item This is the first item of an enumerated
                  list that is nested within the itemized list.

            \item This is the second item of the inner list.
                  \LaTeX\ allows you to nest lists deeper than
                  you really should.
         \end{enumerate}
         This is the rest of the second item of the outer
         list.  It is no more interesting than any other
         part of the item.
   \item This is the third item of the list.
\end{itemize}
You can even display poetry.
\begin{verse}
   There is an environment
    for verse \\             % The \\ command separates lines
   Whose features some poets % within a stanza.
   will curse.

                             % One or more blank lines separate stanzas.

   For instead of making\\
   Them do \emph{all} line breaking, \\
   It allows them to put too many words on a line when they'd rather be
   forced to be terse.
\end{verse}

Mathematical formulas may also be displayed.  A
displayed formula
is
one-line long; multiline
formulas require special formatting instructions.
   \[  \ip{\Gamma}{\psi'} = x'' + y^{2} + z_{i}^{n}\]
Don't start a paragraph with a displayed equation,
nor make one a paragraph by itself.

\end{document}               % End of document.

\section{Bench Section 85}
% synthetic bench section 85
% This is a sample LaTeX input file.  (Version of 12 August 2004.)
%
% A '%' character causes TeX to ignore all remaining text on the line,
% and is used for comments like this one.

\documentclass{article}      % Specifies the document class

                             % The preamble begins here.
\title{An Example Document}  % Declares the document's title.
\author{Leslie Lamport}      % Declares the author's name.
\date{January 21, 1994}      % Deleting this command produces today's date.

\newcommand{\ip}[2]{(#1, #2)}
                             % Defines \ip{arg1}{arg2} to mean
                             % (arg1, arg2).

%\newcommand{\ip}[2]{\langle #1 | #2\rangle}
                             % This is an alternative definition of
                             % \ip that is commented out.

\begin{document}             % End of preamble and beginning of text.

\maketitle                   % Produces the title.

This is an example input file.  Comparing it with
the output it generates can show you how to
produce a simple document of your own.

\section{Ordinary Text}      % Produces section heading.  Lower-level
                             % sections are begun with similar
                             % \subsection and \subsubsection commands.

The ends  of words and sentences are marked
  by   spaces. It  doesn't matter how many
spaces    you type; one is as good as 100.  The
end of   a line counts as a space.

One   or more   blank lines denote the  end
of  a paragraph.

Since any number of consecutive spaces are treated
like a single one, the formatting of the input
file makes no difference to
      \LaTeX,                % The \LaTeX command generates the LaTeX logo.
but it makes a difference to you.  When you use
\LaTeX, making your input file as easy to read
as possible will be a great help as you write
your document and when you change it.  This sample
file shows how you can add comments to your own input
file.

Because printing is different from typewriting,
there are a number of things that you have to do
differently when preparing an input file than if
you were just typing the document directly.
Quotation marks like
       ``this''
have to be handled specially, as do quotes within
quotes:
       ``\,`this'            % \, separates the double and single quote.
        is what I just
        wrote, not  `that'\,''.

Dashes come in three sizes: an
       intra-word
dash, a medium dash for number ranges like
       1--2,
and a punctuation
       dash---like
this.

A sentence-ending space should be larger than the
space between words within a sentence.  You
sometimes have to type special commands in
conjunction with punctuation characters to get
this right, as in the following sentence.
       Gnats, gnus, etc.\ all  % `\ ' makes an inter-word space.
       begin with G\@.         % \@ marks end-of-sentence punctuation.
You should check the spaces after periods when
reading your output to make sure you haven't
forgotten any special cases.  Generating an
ellipsis
       \ldots\               % `\ ' is needed after `\ldots' because TeX
                             % ignores spaces after command names like \ldots
                             % made from \ + letters.
                             %
                             % Note how a `%' character causes TeX to ignore
                             % the end of the input line, so these blank lines
                             % do not start a new paragraph.
                             %
with the right spacing around the periods requires
a special command.

\LaTeX\ interprets some common characters as
commands, so you must type special commands to
generate them.  These characters include the
following:
       \$ \& \% \# \{ and \}.

In printing, text is usually emphasized with an
       \emph{italic}
type style.

\begin{em}
   A long segment of text can also be emphasized
   in this way.  Text within such a segment can be
   given \emph{additional} emphasis.
\end{em}

It is sometimes necessary to prevent \LaTeX\ from
breaking a line where it might otherwise do so.
This may be at a space, as between the ``Mr.''\ and
``Jones'' in
       ``Mr.~Jones'',        % ~ produces an unbreakable interword space.
or within a word---especially when the word is a
symbol like
       \mbox{\emph{itemnum}}
that makes little sense when hyphenated across
lines.

Footnotes\footnote{This is an example of a footnote.}
pose no problem.

\LaTeX\ is good at typesetting mathematical formulas
like
       \( x-3y + z = 7 \)
or
       \( a_{1} > x^{2n} + y^{2n} > x' \)
or
       \( \ip{A}{B} = \sum_{i} a_{i} b_{i} \).
The spaces you type in a formula are
ignored.  Remember that a letter like
       $x$                   % $ ... $  and  \( ... \)  are equivalent
is a formula when it denotes a mathematical
symbol, and it should be typed as one.

\section{Displayed Text}

Text is displayed by indenting it from the left
margin.  Quotations are commonly displayed.  There
are short quotations
\begin{quote}
   This is a short quotation.  It consists of a
   single paragraph of text.  See how it is formatted.
\end{quote}
and longer ones.
\begin{quotation}
   This is a longer quotation.  It consists of two
   paragraphs of text, neither of which are
   particularly interesting.

   This is the second paragraph of the quotation.  It
   is just as dull as the first paragraph.
\end{quotation}
Another frequently-displayed structure is a list.
The following is an example of an \emph{itemized}
list.
\begin{itemize}
   \item This is the first item of an itemized list.
         Each item in the list is marked with a ``tick''.
         You don't have to worry about what kind of tick
         mark is used.

   \item This is the second item of the list.  It
         contains another list nested inside it.  The inner
         list is an \emph{enumerated} list.
         \begin{enumerate}
            \item This is the first item of an enumerated
                  list that is nested within the itemized list.

            \item This is the second item of the inner list.
                  \LaTeX\ allows you to nest lists deeper than
                  you really should.
         \end{enumerate}
         This is the rest of the second item of the outer
         list.  It is no more interesting than any other
         part of the item.
   \item This is the third item of the list.
\end{itemize}
You can even display poetry.
\begin{verse}
   There is an environment
    for verse \\             % The \\ command separates lines
   Whose features some poets % within a stanza.
   will curse.

                             % One or more blank lines separate stanzas.

   For instead of making\\
   Them do \emph{all} line breaking, \\
   It allows them to put too many words on a line when they'd rather be
   forced to be terse.
\end{verse}

Mathematical formulas may also be displayed.  A
displayed formula
is
one-line long; multiline
formulas require special formatting instructions.
   \[  \ip{\Gamma}{\psi'} = x'' + y^{2} + z_{i}^{n}\]
Don't start a paragraph with a displayed equation,
nor make one a paragraph by itself.

\end{document}               % End of document.

\section{Bench Section 86}
% synthetic bench section 86
% This is a sample LaTeX input file.  (Version of 12 August 2004.)
%
% A '%' character causes TeX to ignore all remaining text on the line,
% and is used for comments like this one.

\documentclass{article}      % Specifies the document class

                             % The preamble begins here.
\title{An Example Document}  % Declares the document's title.
\author{Leslie Lamport}      % Declares the author's name.
\date{January 21, 1994}      % Deleting this command produces today's date.

\newcommand{\ip}[2]{(#1, #2)}
                             % Defines \ip{arg1}{arg2} to mean
                             % (arg1, arg2).

%\newcommand{\ip}[2]{\langle #1 | #2\rangle}
                             % This is an alternative definition of
                             % \ip that is commented out.

\begin{document}             % End of preamble and beginning of text.

\maketitle                   % Produces the title.

This is an example input file.  Comparing it with
the output it generates can show you how to
produce a simple document of your own.

\section{Ordinary Text}      % Produces section heading.  Lower-level
                             % sections are begun with similar
                             % \subsection and \subsubsection commands.

The ends  of words and sentences are marked
  by   spaces. It  doesn't matter how many
spaces    you type; one is as good as 100.  The
end of   a line counts as a space.

One   or more   blank lines denote the  end
of  a paragraph.

Since any number of consecutive spaces are treated
like a single one, the formatting of the input
file makes no difference to
      \LaTeX,                % The \LaTeX command generates the LaTeX logo.
but it makes a difference to you.  When you use
\LaTeX, making your input file as easy to read
as possible will be a great help as you write
your document and when you change it.  This sample
file shows how you can add comments to your own input
file.

Because printing is different from typewriting,
there are a number of things that you have to do
differently when preparing an input file than if
you were just typing the document directly.
Quotation marks like
       ``this''
have to be handled specially, as do quotes within
quotes:
       ``\,`this'            % \, separates the double and single quote.
        is what I just
        wrote, not  `that'\,''.

Dashes come in three sizes: an
       intra-word
dash, a medium dash for number ranges like
       1--2,
and a punctuation
       dash---like
this.

A sentence-ending space should be larger than the
space between words within a sentence.  You
sometimes have to type special commands in
conjunction with punctuation characters to get
this right, as in the following sentence.
       Gnats, gnus, etc.\ all  % `\ ' makes an inter-word space.
       begin with G\@.         % \@ marks end-of-sentence punctuation.
You should check the spaces after periods when
reading your output to make sure you haven't
forgotten any special cases.  Generating an
ellipsis
       \ldots\               % `\ ' is needed after `\ldots' because TeX
                             % ignores spaces after command names like \ldots
                             % made from \ + letters.
                             %
                             % Note how a `%' character causes TeX to ignore
                             % the end of the input line, so these blank lines
                             % do not start a new paragraph.
                             %
with the right spacing around the periods requires
a special command.

\LaTeX\ interprets some common characters as
commands, so you must type special commands to
generate them.  These characters include the
following:
       \$ \& \% \# \{ and \}.

In printing, text is usually emphasized with an
       \emph{italic}
type style.

\begin{em}
   A long segment of text can also be emphasized
   in this way.  Text within such a segment can be
   given \emph{additional} emphasis.
\end{em}

It is sometimes necessary to prevent \LaTeX\ from
breaking a line where it might otherwise do so.
This may be at a space, as between the ``Mr.''\ and
``Jones'' in
       ``Mr.~Jones'',        % ~ produces an unbreakable interword space.
or within a word---especially when the word is a
symbol like
       \mbox{\emph{itemnum}}
that makes little sense when hyphenated across
lines.

Footnotes\footnote{This is an example of a footnote.}
pose no problem.

\LaTeX\ is good at typesetting mathematical formulas
like
       \( x-3y + z = 7 \)
or
       \( a_{1} > x^{2n} + y^{2n} > x' \)
or
       \( \ip{A}{B} = \sum_{i} a_{i} b_{i} \).
The spaces you type in a formula are
ignored.  Remember that a letter like
       $x$                   % $ ... $  and  \( ... \)  are equivalent
is a formula when it denotes a mathematical
symbol, and it should be typed as one.

\section{Displayed Text}

Text is displayed by indenting it from the left
margin.  Quotations are commonly displayed.  There
are short quotations
\begin{quote}
   This is a short quotation.  It consists of a
   single paragraph of text.  See how it is formatted.
\end{quote}
and longer ones.
\begin{quotation}
   This is a longer quotation.  It consists of two
   paragraphs of text, neither of which are
   particularly interesting.

   This is the second paragraph of the quotation.  It
   is just as dull as the first paragraph.
\end{quotation}
Another frequently-displayed structure is a list.
The following is an example of an \emph{itemized}
list.
\begin{itemize}
   \item This is the first item of an itemized list.
         Each item in the list is marked with a ``tick''.
         You don't have to worry about what kind of tick
         mark is used.

   \item This is the second item of the list.  It
         contains another list nested inside it.  The inner
         list is an \emph{enumerated} list.
         \begin{enumerate}
            \item This is the first item of an enumerated
                  list that is nested within the itemized list.

            \item This is the second item of the inner list.
                  \LaTeX\ allows you to nest lists deeper than
                  you really should.
         \end{enumerate}
         This is the rest of the second item of the outer
         list.  It is no more interesting than any other
         part of the item.
   \item This is the third item of the list.
\end{itemize}
You can even display poetry.
\begin{verse}
   There is an environment
    for verse \\             % The \\ command separates lines
   Whose features some poets % within a stanza.
   will curse.

                             % One or more blank lines separate stanzas.

   For instead of making\\
   Them do \emph{all} line breaking, \\
   It allows them to put too many words on a line when they'd rather be
   forced to be terse.
\end{verse}

Mathematical formulas may also be displayed.  A
displayed formula
is
one-line long; multiline
formulas require special formatting instructions.
   \[  \ip{\Gamma}{\psi'} = x'' + y^{2} + z_{i}^{n}\]
Don't start a paragraph with a displayed equation,
nor make one a paragraph by itself.

\end{document}               % End of document.

\section{Bench Section 87}
% synthetic bench section 87
% This is a sample LaTeX input file.  (Version of 12 August 2004.)
%
% A '%' character causes TeX to ignore all remaining text on the line,
% and is used for comments like this one.

\documentclass{article}      % Specifies the document class

                             % The preamble begins here.
\title{An Example Document}  % Declares the document's title.
\author{Leslie Lamport}      % Declares the author's name.
\date{January 21, 1994}      % Deleting this command produces today's date.

\newcommand{\ip}[2]{(#1, #2)}
                             % Defines \ip{arg1}{arg2} to mean
                             % (arg1, arg2).

%\newcommand{\ip}[2]{\langle #1 | #2\rangle}
                             % This is an alternative definition of
                             % \ip that is commented out.

\begin{document}             % End of preamble and beginning of text.

\maketitle                   % Produces the title.

This is an example input file.  Comparing it with
the output it generates can show you how to
produce a simple document of your own.

\section{Ordinary Text}      % Produces section heading.  Lower-level
                             % sections are begun with similar
                             % \subsection and \subsubsection commands.

The ends  of words and sentences are marked
  by   spaces. It  doesn't matter how many
spaces    you type; one is as good as 100.  The
end of   a line counts as a space.

One   or more   blank lines denote the  end
of  a paragraph.

Since any number of consecutive spaces are treated
like a single one, the formatting of the input
file makes no difference to
      \LaTeX,                % The \LaTeX command generates the LaTeX logo.
but it makes a difference to you.  When you use
\LaTeX, making your input file as easy to read
as possible will be a great help as you write
your document and when you change it.  This sample
file shows how you can add comments to your own input
file.

Because printing is different from typewriting,
there are a number of things that you have to do
differently when preparing an input file than if
you were just typing the document directly.
Quotation marks like
       ``this''
have to be handled specially, as do quotes within
quotes:
       ``\,`this'            % \, separates the double and single quote.
        is what I just
        wrote, not  `that'\,''.

Dashes come in three sizes: an
       intra-word
dash, a medium dash for number ranges like
       1--2,
and a punctuation
       dash---like
this.

A sentence-ending space should be larger than the
space between words within a sentence.  You
sometimes have to type special commands in
conjunction with punctuation characters to get
this right, as in the following sentence.
       Gnats, gnus, etc.\ all  % `\ ' makes an inter-word space.
       begin with G\@.         % \@ marks end-of-sentence punctuation.
You should check the spaces after periods when
reading your output to make sure you haven't
forgotten any special cases.  Generating an
ellipsis
       \ldots\               % `\ ' is needed after `\ldots' because TeX
                             % ignores spaces after command names like \ldots
                             % made from \ + letters.
                             %
                             % Note how a `%' character causes TeX to ignore
                             % the end of the input line, so these blank lines
                             % do not start a new paragraph.
                             %
with the right spacing around the periods requires
a special command.

\LaTeX\ interprets some common characters as
commands, so you must type special commands to
generate them.  These characters include the
following:
       \$ \& \% \# \{ and \}.

In printing, text is usually emphasized with an
       \emph{italic}
type style.

\begin{em}
   A long segment of text can also be emphasized
   in this way.  Text within such a segment can be
   given \emph{additional} emphasis.
\end{em}

It is sometimes necessary to prevent \LaTeX\ from
breaking a line where it might otherwise do so.
This may be at a space, as between the ``Mr.''\ and
``Jones'' in
       ``Mr.~Jones'',        % ~ produces an unbreakable interword space.
or within a word---especially when the word is a
symbol like
       \mbox{\emph{itemnum}}
that makes little sense when hyphenated across
lines.

Footnotes\footnote{This is an example of a footnote.}
pose no problem.

\LaTeX\ is good at typesetting mathematical formulas
like
       \( x-3y + z = 7 \)
or
       \( a_{1} > x^{2n} + y^{2n} > x' \)
or
       \( \ip{A}{B} = \sum_{i} a_{i} b_{i} \).
The spaces you type in a formula are
ignored.  Remember that a letter like
       $x$                   % $ ... $  and  \( ... \)  are equivalent
is a formula when it denotes a mathematical
symbol, and it should be typed as one.

\section{Displayed Text}

Text is displayed by indenting it from the left
margin.  Quotations are commonly displayed.  There
are short quotations
\begin{quote}
   This is a short quotation.  It consists of a
   single paragraph of text.  See how it is formatted.
\end{quote}
and longer ones.
\begin{quotation}
   This is a longer quotation.  It consists of two
   paragraphs of text, neither of which are
   particularly interesting.

   This is the second paragraph of the quotation.  It
   is just as dull as the first paragraph.
\end{quotation}
Another frequently-displayed structure is a list.
The following is an example of an \emph{itemized}
list.
\begin{itemize}
   \item This is the first item of an itemized list.
         Each item in the list is marked with a ``tick''.
         You don't have to worry about what kind of tick
         mark is used.

   \item This is the second item of the list.  It
         contains another list nested inside it.  The inner
         list is an \emph{enumerated} list.
         \begin{enumerate}
            \item This is the first item of an enumerated
                  list that is nested within the itemized list.

            \item This is the second item of the inner list.
                  \LaTeX\ allows you to nest lists deeper than
                  you really should.
         \end{enumerate}
         This is the rest of the second item of the outer
         list.  It is no more interesting than any other
         part of the item.
   \item This is the third item of the list.
\end{itemize}
You can even display poetry.
\begin{verse}
   There is an environment
    for verse \\             % The \\ command separates lines
   Whose features some poets % within a stanza.
   will curse.

                             % One or more blank lines separate stanzas.

   For instead of making\\
   Them do \emph{all} line breaking, \\
   It allows them to put too many words on a line when they'd rather be
   forced to be terse.
\end{verse}

Mathematical formulas may also be displayed.  A
displayed formula
is
one-line long; multiline
formulas require special formatting instructions.
   \[  \ip{\Gamma}{\psi'} = x'' + y^{2} + z_{i}^{n}\]
Don't start a paragraph with a displayed equation,
nor make one a paragraph by itself.

\end{document}               % End of document.

\section{Bench Section 88}
% synthetic bench section 88
% This is a sample LaTeX input file.  (Version of 12 August 2004.)
%
% A '%' character causes TeX to ignore all remaining text on the line,
% and is used for comments like this one.

\documentclass{article}      % Specifies the document class

                             % The preamble begins here.
\title{An Example Document}  % Declares the document's title.
\author{Leslie Lamport}      % Declares the author's name.
\date{January 21, 1994}      % Deleting this command produces today's date.

\newcommand{\ip}[2]{(#1, #2)}
                             % Defines \ip{arg1}{arg2} to mean
                             % (arg1, arg2).

%\newcommand{\ip}[2]{\langle #1 | #2\rangle}
                             % This is an alternative definition of
                             % \ip that is commented out.

\begin{document}             % End of preamble and beginning of text.

\maketitle                   % Produces the title.

This is an example input file.  Comparing it with
the output it generates can show you how to
produce a simple document of your own.

\section{Ordinary Text}      % Produces section heading.  Lower-level
                             % sections are begun with similar
                             % \subsection and \subsubsection commands.

The ends  of words and sentences are marked
  by   spaces. It  doesn't matter how many
spaces    you type; one is as good as 100.  The
end of   a line counts as a space.

One   or more   blank lines denote the  end
of  a paragraph.

Since any number of consecutive spaces are treated
like a single one, the formatting of the input
file makes no difference to
      \LaTeX,                % The \LaTeX command generates the LaTeX logo.
but it makes a difference to you.  When you use
\LaTeX, making your input file as easy to read
as possible will be a great help as you write
your document and when you change it.  This sample
file shows how you can add comments to your own input
file.

Because printing is different from typewriting,
there are a number of things that you have to do
differently when preparing an input file than if
you were just typing the document directly.
Quotation marks like
       ``this''
have to be handled specially, as do quotes within
quotes:
       ``\,`this'            % \, separates the double and single quote.
        is what I just
        wrote, not  `that'\,''.

Dashes come in three sizes: an
       intra-word
dash, a medium dash for number ranges like
       1--2,
and a punctuation
       dash---like
this.

A sentence-ending space should be larger than the
space between words within a sentence.  You
sometimes have to type special commands in
conjunction with punctuation characters to get
this right, as in the following sentence.
       Gnats, gnus, etc.\ all  % `\ ' makes an inter-word space.
       begin with G\@.         % \@ marks end-of-sentence punctuation.
You should check the spaces after periods when
reading your output to make sure you haven't
forgotten any special cases.  Generating an
ellipsis
       \ldots\               % `\ ' is needed after `\ldots' because TeX
                             % ignores spaces after command names like \ldots
                             % made from \ + letters.
                             %
                             % Note how a `%' character causes TeX to ignore
                             % the end of the input line, so these blank lines
                             % do not start a new paragraph.
                             %
with the right spacing around the periods requires
a special command.

\LaTeX\ interprets some common characters as
commands, so you must type special commands to
generate them.  These characters include the
following:
       \$ \& \% \# \{ and \}.

In printing, text is usually emphasized with an
       \emph{italic}
type style.

\begin{em}
   A long segment of text can also be emphasized
   in this way.  Text within such a segment can be
   given \emph{additional} emphasis.
\end{em}

It is sometimes necessary to prevent \LaTeX\ from
breaking a line where it might otherwise do so.
This may be at a space, as between the ``Mr.''\ and
``Jones'' in
       ``Mr.~Jones'',        % ~ produces an unbreakable interword space.
or within a word---especially when the word is a
symbol like
       \mbox{\emph{itemnum}}
that makes little sense when hyphenated across
lines.

Footnotes\footnote{This is an example of a footnote.}
pose no problem.

\LaTeX\ is good at typesetting mathematical formulas
like
       \( x-3y + z = 7 \)
or
       \( a_{1} > x^{2n} + y^{2n} > x' \)
or
       \( \ip{A}{B} = \sum_{i} a_{i} b_{i} \).
The spaces you type in a formula are
ignored.  Remember that a letter like
       $x$                   % $ ... $  and  \( ... \)  are equivalent
is a formula when it denotes a mathematical
symbol, and it should be typed as one.

\section{Displayed Text}

Text is displayed by indenting it from the left
margin.  Quotations are commonly displayed.  There
are short quotations
\begin{quote}
   This is a short quotation.  It consists of a
   single paragraph of text.  See how it is formatted.
\end{quote}
and longer ones.
\begin{quotation}
   This is a longer quotation.  It consists of two
   paragraphs of text, neither of which are
   particularly interesting.

   This is the second paragraph of the quotation.  It
   is just as dull as the first paragraph.
\end{quotation}
Another frequently-displayed structure is a list.
The following is an example of an \emph{itemized}
list.
\begin{itemize}
   \item This is the first item of an itemized list.
         Each item in the list is marked with a ``tick''.
         You don't have to worry about what kind of tick
         mark is used.

   \item This is the second item of the list.  It
         contains another list nested inside it.  The inner
         list is an \emph{enumerated} list.
         \begin{enumerate}
            \item This is the first item of an enumerated
                  list that is nested within the itemized list.

            \item This is the second item of the inner list.
                  \LaTeX\ allows you to nest lists deeper than
                  you really should.
         \end{enumerate}
         This is the rest of the second item of the outer
         list.  It is no more interesting than any other
         part of the item.
   \item This is the third item of the list.
\end{itemize}
You can even display poetry.
\begin{verse}
   There is an environment
    for verse \\             % The \\ command separates lines
   Whose features some poets % within a stanza.
   will curse.

                             % One or more blank lines separate stanzas.

   For instead of making\\
   Them do \emph{all} line breaking, \\
   It allows them to put too many words on a line when they'd rather be
   forced to be terse.
\end{verse}

Mathematical formulas may also be displayed.  A
displayed formula
is
one-line long; multiline
formulas require special formatting instructions.
   \[  \ip{\Gamma}{\psi'} = x'' + y^{2} + z_{i}^{n}\]
Don't start a paragraph with a displayed equation,
nor make one a paragraph by itself.

\end{document}               % End of document.

\section{Bench Section 89}
% synthetic bench section 89
% This is a sample LaTeX input file.  (Version of 12 August 2004.)
%
% A '%' character causes TeX to ignore all remaining text on the line,
% and is used for comments like this one.

\documentclass{article}      % Specifies the document class

                             % The preamble begins here.
\title{An Example Document}  % Declares the document's title.
\author{Leslie Lamport}      % Declares the author's name.
\date{January 21, 1994}      % Deleting this command produces today's date.

\newcommand{\ip}[2]{(#1, #2)}
                             % Defines \ip{arg1}{arg2} to mean
                             % (arg1, arg2).

%\newcommand{\ip}[2]{\langle #1 | #2\rangle}
                             % This is an alternative definition of
                             % \ip that is commented out.

\begin{document}             % End of preamble and beginning of text.

\maketitle                   % Produces the title.

This is an example input file.  Comparing it with
the output it generates can show you how to
produce a simple document of your own.

\section{Ordinary Text}      % Produces section heading.  Lower-level
                             % sections are begun with similar
                             % \subsection and \subsubsection commands.

The ends  of words and sentences are marked
  by   spaces. It  doesn't matter how many
spaces    you type; one is as good as 100.  The
end of   a line counts as a space.

One   or more   blank lines denote the  end
of  a paragraph.

Since any number of consecutive spaces are treated
like a single one, the formatting of the input
file makes no difference to
      \LaTeX,                % The \LaTeX command generates the LaTeX logo.
but it makes a difference to you.  When you use
\LaTeX, making your input file as easy to read
as possible will be a great help as you write
your document and when you change it.  This sample
file shows how you can add comments to your own input
file.

Because printing is different from typewriting,
there are a number of things that you have to do
differently when preparing an input file than if
you were just typing the document directly.
Quotation marks like
       ``this''
have to be handled specially, as do quotes within
quotes:
       ``\,`this'            % \, separates the double and single quote.
        is what I just
        wrote, not  `that'\,''.

Dashes come in three sizes: an
       intra-word
dash, a medium dash for number ranges like
       1--2,
and a punctuation
       dash---like
this.

A sentence-ending space should be larger than the
space between words within a sentence.  You
sometimes have to type special commands in
conjunction with punctuation characters to get
this right, as in the following sentence.
       Gnats, gnus, etc.\ all  % `\ ' makes an inter-word space.
       begin with G\@.         % \@ marks end-of-sentence punctuation.
You should check the spaces after periods when
reading your output to make sure you haven't
forgotten any special cases.  Generating an
ellipsis
       \ldots\               % `\ ' is needed after `\ldots' because TeX
                             % ignores spaces after command names like \ldots
                             % made from \ + letters.
                             %
                             % Note how a `%' character causes TeX to ignore
                             % the end of the input line, so these blank lines
                             % do not start a new paragraph.
                             %
with the right spacing around the periods requires
a special command.

\LaTeX\ interprets some common characters as
commands, so you must type special commands to
generate them.  These characters include the
following:
       \$ \& \% \# \{ and \}.

In printing, text is usually emphasized with an
       \emph{italic}
type style.

\begin{em}
   A long segment of text can also be emphasized
   in this way.  Text within such a segment can be
   given \emph{additional} emphasis.
\end{em}

It is sometimes necessary to prevent \LaTeX\ from
breaking a line where it might otherwise do so.
This may be at a space, as between the ``Mr.''\ and
``Jones'' in
       ``Mr.~Jones'',        % ~ produces an unbreakable interword space.
or within a word---especially when the word is a
symbol like
       \mbox{\emph{itemnum}}
that makes little sense when hyphenated across
lines.

Footnotes\footnote{This is an example of a footnote.}
pose no problem.

\LaTeX\ is good at typesetting mathematical formulas
like
       \( x-3y + z = 7 \)
or
       \( a_{1} > x^{2n} + y^{2n} > x' \)
or
       \( \ip{A}{B} = \sum_{i} a_{i} b_{i} \).
The spaces you type in a formula are
ignored.  Remember that a letter like
       $x$                   % $ ... $  and  \( ... \)  are equivalent
is a formula when it denotes a mathematical
symbol, and it should be typed as one.

\section{Displayed Text}

Text is displayed by indenting it from the left
margin.  Quotations are commonly displayed.  There
are short quotations
\begin{quote}
   This is a short quotation.  It consists of a
   single paragraph of text.  See how it is formatted.
\end{quote}
and longer ones.
\begin{quotation}
   This is a longer quotation.  It consists of two
   paragraphs of text, neither of which are
   particularly interesting.

   This is the second paragraph of the quotation.  It
   is just as dull as the first paragraph.
\end{quotation}
Another frequently-displayed structure is a list.
The following is an example of an \emph{itemized}
list.
\begin{itemize}
   \item This is the first item of an itemized list.
         Each item in the list is marked with a ``tick''.
         You don't have to worry about what kind of tick
         mark is used.

   \item This is the second item of the list.  It
         contains another list nested inside it.  The inner
         list is an \emph{enumerated} list.
         \begin{enumerate}
            \item This is the first item of an enumerated
                  list that is nested within the itemized list.

            \item This is the second item of the inner list.
                  \LaTeX\ allows you to nest lists deeper than
                  you really should.
         \end{enumerate}
         This is the rest of the second item of the outer
         list.  It is no more interesting than any other
         part of the item.
   \item This is the third item of the list.
\end{itemize}
You can even display poetry.
\begin{verse}
   There is an environment
    for verse \\             % The \\ command separates lines
   Whose features some poets % within a stanza.
   will curse.

                             % One or more blank lines separate stanzas.

   For instead of making\\
   Them do \emph{all} line breaking, \\
   It allows them to put too many words on a line when they'd rather be
   forced to be terse.
\end{verse}

Mathematical formulas may also be displayed.  A
displayed formula
is
one-line long; multiline
formulas require special formatting instructions.
   \[  \ip{\Gamma}{\psi'} = x'' + y^{2} + z_{i}^{n}\]
Don't start a paragraph with a displayed equation,
nor make one a paragraph by itself.

\end{document}               % End of document.

\section{Bench Section 90}
% synthetic bench section 90
% This is a sample LaTeX input file.  (Version of 12 August 2004.)
%
% A '%' character causes TeX to ignore all remaining text on the line,
% and is used for comments like this one.

\documentclass{article}      % Specifies the document class

                             % The preamble begins here.
\title{An Example Document}  % Declares the document's title.
\author{Leslie Lamport}      % Declares the author's name.
\date{January 21, 1994}      % Deleting this command produces today's date.

\newcommand{\ip}[2]{(#1, #2)}
                             % Defines \ip{arg1}{arg2} to mean
                             % (arg1, arg2).

%\newcommand{\ip}[2]{\langle #1 | #2\rangle}
                             % This is an alternative definition of
                             % \ip that is commented out.

\begin{document}             % End of preamble and beginning of text.

\maketitle                   % Produces the title.

This is an example input file.  Comparing it with
the output it generates can show you how to
produce a simple document of your own.

\section{Ordinary Text}      % Produces section heading.  Lower-level
                             % sections are begun with similar
                             % \subsection and \subsubsection commands.

The ends  of words and sentences are marked
  by   spaces. It  doesn't matter how many
spaces    you type; one is as good as 100.  The
end of   a line counts as a space.

One   or more   blank lines denote the  end
of  a paragraph.

Since any number of consecutive spaces are treated
like a single one, the formatting of the input
file makes no difference to
      \LaTeX,                % The \LaTeX command generates the LaTeX logo.
but it makes a difference to you.  When you use
\LaTeX, making your input file as easy to read
as possible will be a great help as you write
your document and when you change it.  This sample
file shows how you can add comments to your own input
file.

Because printing is different from typewriting,
there are a number of things that you have to do
differently when preparing an input file than if
you were just typing the document directly.
Quotation marks like
       ``this''
have to be handled specially, as do quotes within
quotes:
       ``\,`this'            % \, separates the double and single quote.
        is what I just
        wrote, not  `that'\,''.

Dashes come in three sizes: an
       intra-word
dash, a medium dash for number ranges like
       1--2,
and a punctuation
       dash---like
this.

A sentence-ending space should be larger than the
space between words within a sentence.  You
sometimes have to type special commands in
conjunction with punctuation characters to get
this right, as in the following sentence.
       Gnats, gnus, etc.\ all  % `\ ' makes an inter-word space.
       begin with G\@.         % \@ marks end-of-sentence punctuation.
You should check the spaces after periods when
reading your output to make sure you haven't
forgotten any special cases.  Generating an
ellipsis
       \ldots\               % `\ ' is needed after `\ldots' because TeX
                             % ignores spaces after command names like \ldots
                             % made from \ + letters.
                             %
                             % Note how a `%' character causes TeX to ignore
                             % the end of the input line, so these blank lines
                             % do not start a new paragraph.
                             %
with the right spacing around the periods requires
a special command.

\LaTeX\ interprets some common characters as
commands, so you must type special commands to
generate them.  These characters include the
following:
       \$ \& \% \# \{ and \}.

In printing, text is usually emphasized with an
       \emph{italic}
type style.

\begin{em}
   A long segment of text can also be emphasized
   in this way.  Text within such a segment can be
   given \emph{additional} emphasis.
\end{em}

It is sometimes necessary to prevent \LaTeX\ from
breaking a line where it might otherwise do so.
This may be at a space, as between the ``Mr.''\ and
``Jones'' in
       ``Mr.~Jones'',        % ~ produces an unbreakable interword space.
or within a word---especially when the word is a
symbol like
       \mbox{\emph{itemnum}}
that makes little sense when hyphenated across
lines.

Footnotes\footnote{This is an example of a footnote.}
pose no problem.

\LaTeX\ is good at typesetting mathematical formulas
like
       \( x-3y + z = 7 \)
or
       \( a_{1} > x^{2n} + y^{2n} > x' \)
or
       \( \ip{A}{B} = \sum_{i} a_{i} b_{i} \).
The spaces you type in a formula are
ignored.  Remember that a letter like
       $x$                   % $ ... $  and  \( ... \)  are equivalent
is a formula when it denotes a mathematical
symbol, and it should be typed as one.

\section{Displayed Text}

Text is displayed by indenting it from the left
margin.  Quotations are commonly displayed.  There
are short quotations
\begin{quote}
   This is a short quotation.  It consists of a
   single paragraph of text.  See how it is formatted.
\end{quote}
and longer ones.
\begin{quotation}
   This is a longer quotation.  It consists of two
   paragraphs of text, neither of which are
   particularly interesting.

   This is the second paragraph of the quotation.  It
   is just as dull as the first paragraph.
\end{quotation}
Another frequently-displayed structure is a list.
The following is an example of an \emph{itemized}
list.
\begin{itemize}
   \item This is the first item of an itemized list.
         Each item in the list is marked with a ``tick''.
         You don't have to worry about what kind of tick
         mark is used.

   \item This is the second item of the list.  It
         contains another list nested inside it.  The inner
         list is an \emph{enumerated} list.
         \begin{enumerate}
            \item This is the first item of an enumerated
                  list that is nested within the itemized list.

            \item This is the second item of the inner list.
                  \LaTeX\ allows you to nest lists deeper than
                  you really should.
         \end{enumerate}
         This is the rest of the second item of the outer
         list.  It is no more interesting than any other
         part of the item.
   \item This is the third item of the list.
\end{itemize}
You can even display poetry.
\begin{verse}
   There is an environment
    for verse \\             % The \\ command separates lines
   Whose features some poets % within a stanza.
   will curse.

                             % One or more blank lines separate stanzas.

   For instead of making\\
   Them do \emph{all} line breaking, \\
   It allows them to put too many words on a line when they'd rather be
   forced to be terse.
\end{verse}

Mathematical formulas may also be displayed.  A
displayed formula
is
one-line long; multiline
formulas require special formatting instructions.
   \[  \ip{\Gamma}{\psi'} = x'' + y^{2} + z_{i}^{n}\]
Don't start a paragraph with a displayed equation,
nor make one a paragraph by itself.

\end{document}               % End of document.

\section{Bench Section 91}
% synthetic bench section 91
% This is a sample LaTeX input file.  (Version of 12 August 2004.)
%
% A '%' character causes TeX to ignore all remaining text on the line,
% and is used for comments like this one.

\documentclass{article}      % Specifies the document class

                             % The preamble begins here.
\title{An Example Document}  % Declares the document's title.
\author{Leslie Lamport}      % Declares the author's name.
\date{January 21, 1994}      % Deleting this command produces today's date.

\newcommand{\ip}[2]{(#1, #2)}
                             % Defines \ip{arg1}{arg2} to mean
                             % (arg1, arg2).

%\newcommand{\ip}[2]{\langle #1 | #2\rangle}
                             % This is an alternative definition of
                             % \ip that is commented out.

\begin{document}             % End of preamble and beginning of text.

\maketitle                   % Produces the title.

This is an example input file.  Comparing it with
the output it generates can show you how to
produce a simple document of your own.

\section{Ordinary Text}      % Produces section heading.  Lower-level
                             % sections are begun with similar
                             % \subsection and \subsubsection commands.

The ends  of words and sentences are marked
  by   spaces. It  doesn't matter how many
spaces    you type; one is as good as 100.  The
end of   a line counts as a space.

One   or more   blank lines denote the  end
of  a paragraph.

Since any number of consecutive spaces are treated
like a single one, the formatting of the input
file makes no difference to
      \LaTeX,                % The \LaTeX command generates the LaTeX logo.
but it makes a difference to you.  When you use
\LaTeX, making your input file as easy to read
as possible will be a great help as you write
your document and when you change it.  This sample
file shows how you can add comments to your own input
file.

Because printing is different from typewriting,
there are a number of things that you have to do
differently when preparing an input file than if
you were just typing the document directly.
Quotation marks like
       ``this''
have to be handled specially, as do quotes within
quotes:
       ``\,`this'            % \, separates the double and single quote.
        is what I just
        wrote, not  `that'\,''.

Dashes come in three sizes: an
       intra-word
dash, a medium dash for number ranges like
       1--2,
and a punctuation
       dash---like
this.

A sentence-ending space should be larger than the
space between words within a sentence.  You
sometimes have to type special commands in
conjunction with punctuation characters to get
this right, as in the following sentence.
       Gnats, gnus, etc.\ all  % `\ ' makes an inter-word space.
       begin with G\@.         % \@ marks end-of-sentence punctuation.
You should check the spaces after periods when
reading your output to make sure you haven't
forgotten any special cases.  Generating an
ellipsis
       \ldots\               % `\ ' is needed after `\ldots' because TeX
                             % ignores spaces after command names like \ldots
                             % made from \ + letters.
                             %
                             % Note how a `%' character causes TeX to ignore
                             % the end of the input line, so these blank lines
                             % do not start a new paragraph.
                             %
with the right spacing around the periods requires
a special command.

\LaTeX\ interprets some common characters as
commands, so you must type special commands to
generate them.  These characters include the
following:
       \$ \& \% \# \{ and \}.

In printing, text is usually emphasized with an
       \emph{italic}
type style.

\begin{em}
   A long segment of text can also be emphasized
   in this way.  Text within such a segment can be
   given \emph{additional} emphasis.
\end{em}

It is sometimes necessary to prevent \LaTeX\ from
breaking a line where it might otherwise do so.
This may be at a space, as between the ``Mr.''\ and
``Jones'' in
       ``Mr.~Jones'',        % ~ produces an unbreakable interword space.
or within a word---especially when the word is a
symbol like
       \mbox{\emph{itemnum}}
that makes little sense when hyphenated across
lines.

Footnotes\footnote{This is an example of a footnote.}
pose no problem.

\LaTeX\ is good at typesetting mathematical formulas
like
       \( x-3y + z = 7 \)
or
       \( a_{1} > x^{2n} + y^{2n} > x' \)
or
       \( \ip{A}{B} = \sum_{i} a_{i} b_{i} \).
The spaces you type in a formula are
ignored.  Remember that a letter like
       $x$                   % $ ... $  and  \( ... \)  are equivalent
is a formula when it denotes a mathematical
symbol, and it should be typed as one.

\section{Displayed Text}

Text is displayed by indenting it from the left
margin.  Quotations are commonly displayed.  There
are short quotations
\begin{quote}
   This is a short quotation.  It consists of a
   single paragraph of text.  See how it is formatted.
\end{quote}
and longer ones.
\begin{quotation}
   This is a longer quotation.  It consists of two
   paragraphs of text, neither of which are
   particularly interesting.

   This is the second paragraph of the quotation.  It
   is just as dull as the first paragraph.
\end{quotation}
Another frequently-displayed structure is a list.
The following is an example of an \emph{itemized}
list.
\begin{itemize}
   \item This is the first item of an itemized list.
         Each item in the list is marked with a ``tick''.
         You don't have to worry about what kind of tick
         mark is used.

   \item This is the second item of the list.  It
         contains another list nested inside it.  The inner
         list is an \emph{enumerated} list.
         \begin{enumerate}
            \item This is the first item of an enumerated
                  list that is nested within the itemized list.

            \item This is the second item of the inner list.
                  \LaTeX\ allows you to nest lists deeper than
                  you really should.
         \end{enumerate}
         This is the rest of the second item of the outer
         list.  It is no more interesting than any other
         part of the item.
   \item This is the third item of the list.
\end{itemize}
You can even display poetry.
\begin{verse}
   There is an environment
    for verse \\             % The \\ command separates lines
   Whose features some poets % within a stanza.
   will curse.

                             % One or more blank lines separate stanzas.

   For instead of making\\
   Them do \emph{all} line breaking, \\
   It allows them to put too many words on a line when they'd rather be
   forced to be terse.
\end{verse}

Mathematical formulas may also be displayed.  A
displayed formula
is
one-line long; multiline
formulas require special formatting instructions.
   \[  \ip{\Gamma}{\psi'} = x'' + y^{2} + z_{i}^{n}\]
Don't start a paragraph with a displayed equation,
nor make one a paragraph by itself.

\end{document}               % End of document.

\section{Bench Section 92}
% synthetic bench section 92
% This is a sample LaTeX input file.  (Version of 12 August 2004.)
%
% A '%' character causes TeX to ignore all remaining text on the line,
% and is used for comments like this one.

\documentclass{article}      % Specifies the document class

                             % The preamble begins here.
\title{An Example Document}  % Declares the document's title.
\author{Leslie Lamport}      % Declares the author's name.
\date{January 21, 1994}      % Deleting this command produces today's date.

\newcommand{\ip}[2]{(#1, #2)}
                             % Defines \ip{arg1}{arg2} to mean
                             % (arg1, arg2).

%\newcommand{\ip}[2]{\langle #1 | #2\rangle}
                             % This is an alternative definition of
                             % \ip that is commented out.

\begin{document}             % End of preamble and beginning of text.

\maketitle                   % Produces the title.

This is an example input file.  Comparing it with
the output it generates can show you how to
produce a simple document of your own.

\section{Ordinary Text}      % Produces section heading.  Lower-level
                             % sections are begun with similar
                             % \subsection and \subsubsection commands.

The ends  of words and sentences are marked
  by   spaces. It  doesn't matter how many
spaces    you type; one is as good as 100.  The
end of   a line counts as a space.

One   or more   blank lines denote the  end
of  a paragraph.

Since any number of consecutive spaces are treated
like a single one, the formatting of the input
file makes no difference to
      \LaTeX,                % The \LaTeX command generates the LaTeX logo.
but it makes a difference to you.  When you use
\LaTeX, making your input file as easy to read
as possible will be a great help as you write
your document and when you change it.  This sample
file shows how you can add comments to your own input
file.

Because printing is different from typewriting,
there are a number of things that you have to do
differently when preparing an input file than if
you were just typing the document directly.
Quotation marks like
       ``this''
have to be handled specially, as do quotes within
quotes:
       ``\,`this'            % \, separates the double and single quote.
        is what I just
        wrote, not  `that'\,''.

Dashes come in three sizes: an
       intra-word
dash, a medium dash for number ranges like
       1--2,
and a punctuation
       dash---like
this.

A sentence-ending space should be larger than the
space between words within a sentence.  You
sometimes have to type special commands in
conjunction with punctuation characters to get
this right, as in the following sentence.
       Gnats, gnus, etc.\ all  % `\ ' makes an inter-word space.
       begin with G\@.         % \@ marks end-of-sentence punctuation.
You should check the spaces after periods when
reading your output to make sure you haven't
forgotten any special cases.  Generating an
ellipsis
       \ldots\               % `\ ' is needed after `\ldots' because TeX
                             % ignores spaces after command names like \ldots
                             % made from \ + letters.
                             %
                             % Note how a `%' character causes TeX to ignore
                             % the end of the input line, so these blank lines
                             % do not start a new paragraph.
                             %
with the right spacing around the periods requires
a special command.

\LaTeX\ interprets some common characters as
commands, so you must type special commands to
generate them.  These characters include the
following:
       \$ \& \% \# \{ and \}.

In printing, text is usually emphasized with an
       \emph{italic}
type style.

\begin{em}
   A long segment of text can also be emphasized
   in this way.  Text within such a segment can be
   given \emph{additional} emphasis.
\end{em}

It is sometimes necessary to prevent \LaTeX\ from
breaking a line where it might otherwise do so.
This may be at a space, as between the ``Mr.''\ and
``Jones'' in
       ``Mr.~Jones'',        % ~ produces an unbreakable interword space.
or within a word---especially when the word is a
symbol like
       \mbox{\emph{itemnum}}
that makes little sense when hyphenated across
lines.

Footnotes\footnote{This is an example of a footnote.}
pose no problem.

\LaTeX\ is good at typesetting mathematical formulas
like
       \( x-3y + z = 7 \)
or
       \( a_{1} > x^{2n} + y^{2n} > x' \)
or
       \( \ip{A}{B} = \sum_{i} a_{i} b_{i} \).
The spaces you type in a formula are
ignored.  Remember that a letter like
       $x$                   % $ ... $  and  \( ... \)  are equivalent
is a formula when it denotes a mathematical
symbol, and it should be typed as one.

\section{Displayed Text}

Text is displayed by indenting it from the left
margin.  Quotations are commonly displayed.  There
are short quotations
\begin{quote}
   This is a short quotation.  It consists of a
   single paragraph of text.  See how it is formatted.
\end{quote}
and longer ones.
\begin{quotation}
   This is a longer quotation.  It consists of two
   paragraphs of text, neither of which are
   particularly interesting.

   This is the second paragraph of the quotation.  It
   is just as dull as the first paragraph.
\end{quotation}
Another frequently-displayed structure is a list.
The following is an example of an \emph{itemized}
list.
\begin{itemize}
   \item This is the first item of an itemized list.
         Each item in the list is marked with a ``tick''.
         You don't have to worry about what kind of tick
         mark is used.

   \item This is the second item of the list.  It
         contains another list nested inside it.  The inner
         list is an \emph{enumerated} list.
         \begin{enumerate}
            \item This is the first item of an enumerated
                  list that is nested within the itemized list.

            \item This is the second item of the inner list.
                  \LaTeX\ allows you to nest lists deeper than
                  you really should.
         \end{enumerate}
         This is the rest of the second item of the outer
         list.  It is no more interesting than any other
         part of the item.
   \item This is the third item of the list.
\end{itemize}
You can even display poetry.
\begin{verse}
   There is an environment
    for verse \\             % The \\ command separates lines
   Whose features some poets % within a stanza.
   will curse.

                             % One or more blank lines separate stanzas.

   For instead of making\\
   Them do \emph{all} line breaking, \\
   It allows them to put too many words on a line when they'd rather be
   forced to be terse.
\end{verse}

Mathematical formulas may also be displayed.  A
displayed formula
is
one-line long; multiline
formulas require special formatting instructions.
   \[  \ip{\Gamma}{\psi'} = x'' + y^{2} + z_{i}^{n}\]
Don't start a paragraph with a displayed equation,
nor make one a paragraph by itself.

\end{document}               % End of document.

\section{Bench Section 93}
% synthetic bench section 93
% This is a sample LaTeX input file.  (Version of 12 August 2004.)
%
% A '%' character causes TeX to ignore all remaining text on the line,
% and is used for comments like this one.

\documentclass{article}      % Specifies the document class

                             % The preamble begins here.
\title{An Example Document}  % Declares the document's title.
\author{Leslie Lamport}      % Declares the author's name.
\date{January 21, 1994}      % Deleting this command produces today's date.

\newcommand{\ip}[2]{(#1, #2)}
                             % Defines \ip{arg1}{arg2} to mean
                             % (arg1, arg2).

%\newcommand{\ip}[2]{\langle #1 | #2\rangle}
                             % This is an alternative definition of
                             % \ip that is commented out.

\begin{document}             % End of preamble and beginning of text.

\maketitle                   % Produces the title.

This is an example input file.  Comparing it with
the output it generates can show you how to
produce a simple document of your own.

\section{Ordinary Text}      % Produces section heading.  Lower-level
                             % sections are begun with similar
                             % \subsection and \subsubsection commands.

The ends  of words and sentences are marked
  by   spaces. It  doesn't matter how many
spaces    you type; one is as good as 100.  The
end of   a line counts as a space.

One   or more   blank lines denote the  end
of  a paragraph.

Since any number of consecutive spaces are treated
like a single one, the formatting of the input
file makes no difference to
      \LaTeX,                % The \LaTeX command generates the LaTeX logo.
but it makes a difference to you.  When you use
\LaTeX, making your input file as easy to read
as possible will be a great help as you write
your document and when you change it.  This sample
file shows how you can add comments to your own input
file.

Because printing is different from typewriting,
there are a number of things that you have to do
differently when preparing an input file than if
you were just typing the document directly.
Quotation marks like
       ``this''
have to be handled specially, as do quotes within
quotes:
       ``\,`this'            % \, separates the double and single quote.
        is what I just
        wrote, not  `that'\,''.

Dashes come in three sizes: an
       intra-word
dash, a medium dash for number ranges like
       1--2,
and a punctuation
       dash---like
this.

A sentence-ending space should be larger than the
space between words within a sentence.  You
sometimes have to type special commands in
conjunction with punctuation characters to get
this right, as in the following sentence.
       Gnats, gnus, etc.\ all  % `\ ' makes an inter-word space.
       begin with G\@.         % \@ marks end-of-sentence punctuation.
You should check the spaces after periods when
reading your output to make sure you haven't
forgotten any special cases.  Generating an
ellipsis
       \ldots\               % `\ ' is needed after `\ldots' because TeX
                             % ignores spaces after command names like \ldots
                             % made from \ + letters.
                             %
                             % Note how a `%' character causes TeX to ignore
                             % the end of the input line, so these blank lines
                             % do not start a new paragraph.
                             %
with the right spacing around the periods requires
a special command.

\LaTeX\ interprets some common characters as
commands, so you must type special commands to
generate them.  These characters include the
following:
       \$ \& \% \# \{ and \}.

In printing, text is usually emphasized with an
       \emph{italic}
type style.

\begin{em}
   A long segment of text can also be emphasized
   in this way.  Text within such a segment can be
   given \emph{additional} emphasis.
\end{em}

It is sometimes necessary to prevent \LaTeX\ from
breaking a line where it might otherwise do so.
This may be at a space, as between the ``Mr.''\ and
``Jones'' in
       ``Mr.~Jones'',        % ~ produces an unbreakable interword space.
or within a word---especially when the word is a
symbol like
       \mbox{\emph{itemnum}}
that makes little sense when hyphenated across
lines.

Footnotes\footnote{This is an example of a footnote.}
pose no problem.

\LaTeX\ is good at typesetting mathematical formulas
like
       \( x-3y + z = 7 \)
or
       \( a_{1} > x^{2n} + y^{2n} > x' \)
or
       \( \ip{A}{B} = \sum_{i} a_{i} b_{i} \).
The spaces you type in a formula are
ignored.  Remember that a letter like
       $x$                   % $ ... $  and  \( ... \)  are equivalent
is a formula when it denotes a mathematical
symbol, and it should be typed as one.

\section{Displayed Text}

Text is displayed by indenting it from the left
margin.  Quotations are commonly displayed.  There
are short quotations
\begin{quote}
   This is a short quotation.  It consists of a
   single paragraph of text.  See how it is formatted.
\end{quote}
and longer ones.
\begin{quotation}
   This is a longer quotation.  It consists of two
   paragraphs of text, neither of which are
   particularly interesting.

   This is the second paragraph of the quotation.  It
   is just as dull as the first paragraph.
\end{quotation}
Another frequently-displayed structure is a list.
The following is an example of an \emph{itemized}
list.
\begin{itemize}
   \item This is the first item of an itemized list.
         Each item in the list is marked with a ``tick''.
         You don't have to worry about what kind of tick
         mark is used.

   \item This is the second item of the list.  It
         contains another list nested inside it.  The inner
         list is an \emph{enumerated} list.
         \begin{enumerate}
            \item This is the first item of an enumerated
                  list that is nested within the itemized list.

            \item This is the second item of the inner list.
                  \LaTeX\ allows you to nest lists deeper than
                  you really should.
         \end{enumerate}
         This is the rest of the second item of the outer
         list.  It is no more interesting than any other
         part of the item.
   \item This is the third item of the list.
\end{itemize}
You can even display poetry.
\begin{verse}
   There is an environment
    for verse \\             % The \\ command separates lines
   Whose features some poets % within a stanza.
   will curse.

                             % One or more blank lines separate stanzas.

   For instead of making\\
   Them do \emph{all} line breaking, \\
   It allows them to put too many words on a line when they'd rather be
   forced to be terse.
\end{verse}

Mathematical formulas may also be displayed.  A
displayed formula
is
one-line long; multiline
formulas require special formatting instructions.
   \[  \ip{\Gamma}{\psi'} = x'' + y^{2} + z_{i}^{n}\]
Don't start a paragraph with a displayed equation,
nor make one a paragraph by itself.

\end{document}               % End of document.

\section{Bench Section 94}
% synthetic bench section 94
% This is a sample LaTeX input file.  (Version of 12 August 2004.)
%
% A '%' character causes TeX to ignore all remaining text on the line,
% and is used for comments like this one.

\documentclass{article}      % Specifies the document class

                             % The preamble begins here.
\title{An Example Document}  % Declares the document's title.
\author{Leslie Lamport}      % Declares the author's name.
\date{January 21, 1994}      % Deleting this command produces today's date.

\newcommand{\ip}[2]{(#1, #2)}
                             % Defines \ip{arg1}{arg2} to mean
                             % (arg1, arg2).

%\newcommand{\ip}[2]{\langle #1 | #2\rangle}
                             % This is an alternative definition of
                             % \ip that is commented out.

\begin{document}             % End of preamble and beginning of text.

\maketitle                   % Produces the title.

This is an example input file.  Comparing it with
the output it generates can show you how to
produce a simple document of your own.

\section{Ordinary Text}      % Produces section heading.  Lower-level
                             % sections are begun with similar
                             % \subsection and \subsubsection commands.

The ends  of words and sentences are marked
  by   spaces. It  doesn't matter how many
spaces    you type; one is as good as 100.  The
end of   a line counts as a space.

One   or more   blank lines denote the  end
of  a paragraph.

Since any number of consecutive spaces are treated
like a single one, the formatting of the input
file makes no difference to
      \LaTeX,                % The \LaTeX command generates the LaTeX logo.
but it makes a difference to you.  When you use
\LaTeX, making your input file as easy to read
as possible will be a great help as you write
your document and when you change it.  This sample
file shows how you can add comments to your own input
file.

Because printing is different from typewriting,
there are a number of things that you have to do
differently when preparing an input file than if
you were just typing the document directly.
Quotation marks like
       ``this''
have to be handled specially, as do quotes within
quotes:
       ``\,`this'            % \, separates the double and single quote.
        is what I just
        wrote, not  `that'\,''.

Dashes come in three sizes: an
       intra-word
dash, a medium dash for number ranges like
       1--2,
and a punctuation
       dash---like
this.

A sentence-ending space should be larger than the
space between words within a sentence.  You
sometimes have to type special commands in
conjunction with punctuation characters to get
this right, as in the following sentence.
       Gnats, gnus, etc.\ all  % `\ ' makes an inter-word space.
       begin with G\@.         % \@ marks end-of-sentence punctuation.
You should check the spaces after periods when
reading your output to make sure you haven't
forgotten any special cases.  Generating an
ellipsis
       \ldots\               % `\ ' is needed after `\ldots' because TeX
                             % ignores spaces after command names like \ldots
                             % made from \ + letters.
                             %
                             % Note how a `%' character causes TeX to ignore
                             % the end of the input line, so these blank lines
                             % do not start a new paragraph.
                             %
with the right spacing around the periods requires
a special command.

\LaTeX\ interprets some common characters as
commands, so you must type special commands to
generate them.  These characters include the
following:
       \$ \& \% \# \{ and \}.

In printing, text is usually emphasized with an
       \emph{italic}
type style.

\begin{em}
   A long segment of text can also be emphasized
   in this way.  Text within such a segment can be
   given \emph{additional} emphasis.
\end{em}

It is sometimes necessary to prevent \LaTeX\ from
breaking a line where it might otherwise do so.
This may be at a space, as between the ``Mr.''\ and
``Jones'' in
       ``Mr.~Jones'',        % ~ produces an unbreakable interword space.
or within a word---especially when the word is a
symbol like
       \mbox{\emph{itemnum}}
that makes little sense when hyphenated across
lines.

Footnotes\footnote{This is an example of a footnote.}
pose no problem.

\LaTeX\ is good at typesetting mathematical formulas
like
       \( x-3y + z = 7 \)
or
       \( a_{1} > x^{2n} + y^{2n} > x' \)
or
       \( \ip{A}{B} = \sum_{i} a_{i} b_{i} \).
The spaces you type in a formula are
ignored.  Remember that a letter like
       $x$                   % $ ... $  and  \( ... \)  are equivalent
is a formula when it denotes a mathematical
symbol, and it should be typed as one.

\section{Displayed Text}

Text is displayed by indenting it from the left
margin.  Quotations are commonly displayed.  There
are short quotations
\begin{quote}
   This is a short quotation.  It consists of a
   single paragraph of text.  See how it is formatted.
\end{quote}
and longer ones.
\begin{quotation}
   This is a longer quotation.  It consists of two
   paragraphs of text, neither of which are
   particularly interesting.

   This is the second paragraph of the quotation.  It
   is just as dull as the first paragraph.
\end{quotation}
Another frequently-displayed structure is a list.
The following is an example of an \emph{itemized}
list.
\begin{itemize}
   \item This is the first item of an itemized list.
         Each item in the list is marked with a ``tick''.
         You don't have to worry about what kind of tick
         mark is used.

   \item This is the second item of the list.  It
         contains another list nested inside it.  The inner
         list is an \emph{enumerated} list.
         \begin{enumerate}
            \item This is the first item of an enumerated
                  list that is nested within the itemized list.

            \item This is the second item of the inner list.
                  \LaTeX\ allows you to nest lists deeper than
                  you really should.
         \end{enumerate}
         This is the rest of the second item of the outer
         list.  It is no more interesting than any other
         part of the item.
   \item This is the third item of the list.
\end{itemize}
You can even display poetry.
\begin{verse}
   There is an environment
    for verse \\             % The \\ command separates lines
   Whose features some poets % within a stanza.
   will curse.

                             % One or more blank lines separate stanzas.

   For instead of making\\
   Them do \emph{all} line breaking, \\
   It allows them to put too many words on a line when they'd rather be
   forced to be terse.
\end{verse}

Mathematical formulas may also be displayed.  A
displayed formula
is
one-line long; multiline
formulas require special formatting instructions.
   \[  \ip{\Gamma}{\psi'} = x'' + y^{2} + z_{i}^{n}\]
Don't start a paragraph with a displayed equation,
nor make one a paragraph by itself.

\end{document}               % End of document.

\section{Bench Section 95}
% synthetic bench section 95
% This is a sample LaTeX input file.  (Version of 12 August 2004.)
%
% A '%' character causes TeX to ignore all remaining text on the line,
% and is used for comments like this one.

\documentclass{article}      % Specifies the document class

                             % The preamble begins here.
\title{An Example Document}  % Declares the document's title.
\author{Leslie Lamport}      % Declares the author's name.
\date{January 21, 1994}      % Deleting this command produces today's date.

\newcommand{\ip}[2]{(#1, #2)}
                             % Defines \ip{arg1}{arg2} to mean
                             % (arg1, arg2).

%\newcommand{\ip}[2]{\langle #1 | #2\rangle}
                             % This is an alternative definition of
                             % \ip that is commented out.

\begin{document}             % End of preamble and beginning of text.

\maketitle                   % Produces the title.

This is an example input file.  Comparing it with
the output it generates can show you how to
produce a simple document of your own.

\section{Ordinary Text}      % Produces section heading.  Lower-level
                             % sections are begun with similar
                             % \subsection and \subsubsection commands.

The ends  of words and sentences are marked
  by   spaces. It  doesn't matter how many
spaces    you type; one is as good as 100.  The
end of   a line counts as a space.

One   or more   blank lines denote the  end
of  a paragraph.

Since any number of consecutive spaces are treated
like a single one, the formatting of the input
file makes no difference to
      \LaTeX,                % The \LaTeX command generates the LaTeX logo.
but it makes a difference to you.  When you use
\LaTeX, making your input file as easy to read
as possible will be a great help as you write
your document and when you change it.  This sample
file shows how you can add comments to your own input
file.

Because printing is different from typewriting,
there are a number of things that you have to do
differently when preparing an input file than if
you were just typing the document directly.
Quotation marks like
       ``this''
have to be handled specially, as do quotes within
quotes:
       ``\,`this'            % \, separates the double and single quote.
        is what I just
        wrote, not  `that'\,''.

Dashes come in three sizes: an
       intra-word
dash, a medium dash for number ranges like
       1--2,
and a punctuation
       dash---like
this.

A sentence-ending space should be larger than the
space between words within a sentence.  You
sometimes have to type special commands in
conjunction with punctuation characters to get
this right, as in the following sentence.
       Gnats, gnus, etc.\ all  % `\ ' makes an inter-word space.
       begin with G\@.         % \@ marks end-of-sentence punctuation.
You should check the spaces after periods when
reading your output to make sure you haven't
forgotten any special cases.  Generating an
ellipsis
       \ldots\               % `\ ' is needed after `\ldots' because TeX
                             % ignores spaces after command names like \ldots
                             % made from \ + letters.
                             %
                             % Note how a `%' character causes TeX to ignore
                             % the end of the input line, so these blank lines
                             % do not start a new paragraph.
                             %
with the right spacing around the periods requires
a special command.

\LaTeX\ interprets some common characters as
commands, so you must type special commands to
generate them.  These characters include the
following:
       \$ \& \% \# \{ and \}.

In printing, text is usually emphasized with an
       \emph{italic}
type style.

\begin{em}
   A long segment of text can also be emphasized
   in this way.  Text within such a segment can be
   given \emph{additional} emphasis.
\end{em}

It is sometimes necessary to prevent \LaTeX\ from
breaking a line where it might otherwise do so.
This may be at a space, as between the ``Mr.''\ and
``Jones'' in
       ``Mr.~Jones'',        % ~ produces an unbreakable interword space.
or within a word---especially when the word is a
symbol like
       \mbox{\emph{itemnum}}
that makes little sense when hyphenated across
lines.

Footnotes\footnote{This is an example of a footnote.}
pose no problem.

\LaTeX\ is good at typesetting mathematical formulas
like
       \( x-3y + z = 7 \)
or
       \( a_{1} > x^{2n} + y^{2n} > x' \)
or
       \( \ip{A}{B} = \sum_{i} a_{i} b_{i} \).
The spaces you type in a formula are
ignored.  Remember that a letter like
       $x$                   % $ ... $  and  \( ... \)  are equivalent
is a formula when it denotes a mathematical
symbol, and it should be typed as one.

\section{Displayed Text}

Text is displayed by indenting it from the left
margin.  Quotations are commonly displayed.  There
are short quotations
\begin{quote}
   This is a short quotation.  It consists of a
   single paragraph of text.  See how it is formatted.
\end{quote}
and longer ones.
\begin{quotation}
   This is a longer quotation.  It consists of two
   paragraphs of text, neither of which are
   particularly interesting.

   This is the second paragraph of the quotation.  It
   is just as dull as the first paragraph.
\end{quotation}
Another frequently-displayed structure is a list.
The following is an example of an \emph{itemized}
list.
\begin{itemize}
   \item This is the first item of an itemized list.
         Each item in the list is marked with a ``tick''.
         You don't have to worry about what kind of tick
         mark is used.

   \item This is the second item of the list.  It
         contains another list nested inside it.  The inner
         list is an \emph{enumerated} list.
         \begin{enumerate}
            \item This is the first item of an enumerated
                  list that is nested within the itemized list.

            \item This is the second item of the inner list.
                  \LaTeX\ allows you to nest lists deeper than
                  you really should.
         \end{enumerate}
         This is the rest of the second item of the outer
         list.  It is no more interesting than any other
         part of the item.
   \item This is the third item of the list.
\end{itemize}
You can even display poetry.
\begin{verse}
   There is an environment
    for verse \\             % The \\ command separates lines
   Whose features some poets % within a stanza.
   will curse.

                             % One or more blank lines separate stanzas.

   For instead of making\\
   Them do \emph{all} line breaking, \\
   It allows them to put too many words on a line when they'd rather be
   forced to be terse.
\end{verse}

Mathematical formulas may also be displayed.  A
displayed formula
is
one-line long; multiline
formulas require special formatting instructions.
   \[  \ip{\Gamma}{\psi'} = x'' + y^{2} + z_{i}^{n}\]
Don't start a paragraph with a displayed equation,
nor make one a paragraph by itself.

\end{document}               % End of document.

\section{Bench Section 96}
% synthetic bench section 96
% This is a sample LaTeX input file.  (Version of 12 August 2004.)
%
% A '%' character causes TeX to ignore all remaining text on the line,
% and is used for comments like this one.

\documentclass{article}      % Specifies the document class

                             % The preamble begins here.
\title{An Example Document}  % Declares the document's title.
\author{Leslie Lamport}      % Declares the author's name.
\date{January 21, 1994}      % Deleting this command produces today's date.

\newcommand{\ip}[2]{(#1, #2)}
                             % Defines \ip{arg1}{arg2} to mean
                             % (arg1, arg2).

%\newcommand{\ip}[2]{\langle #1 | #2\rangle}
                             % This is an alternative definition of
                             % \ip that is commented out.

\begin{document}             % End of preamble and beginning of text.

\maketitle                   % Produces the title.

This is an example input file.  Comparing it with
the output it generates can show you how to
produce a simple document of your own.

\section{Ordinary Text}      % Produces section heading.  Lower-level
                             % sections are begun with similar
                             % \subsection and \subsubsection commands.

The ends  of words and sentences are marked
  by   spaces. It  doesn't matter how many
spaces    you type; one is as good as 100.  The
end of   a line counts as a space.

One   or more   blank lines denote the  end
of  a paragraph.

Since any number of consecutive spaces are treated
like a single one, the formatting of the input
file makes no difference to
      \LaTeX,                % The \LaTeX command generates the LaTeX logo.
but it makes a difference to you.  When you use
\LaTeX, making your input file as easy to read
as possible will be a great help as you write
your document and when you change it.  This sample
file shows how you can add comments to your own input
file.

Because printing is different from typewriting,
there are a number of things that you have to do
differently when preparing an input file than if
you were just typing the document directly.
Quotation marks like
       ``this''
have to be handled specially, as do quotes within
quotes:
       ``\,`this'            % \, separates the double and single quote.
        is what I just
        wrote, not  `that'\,''.

Dashes come in three sizes: an
       intra-word
dash, a medium dash for number ranges like
       1--2,
and a punctuation
       dash---like
this.

A sentence-ending space should be larger than the
space between words within a sentence.  You
sometimes have to type special commands in
conjunction with punctuation characters to get
this right, as in the following sentence.
       Gnats, gnus, etc.\ all  % `\ ' makes an inter-word space.
       begin with G\@.         % \@ marks end-of-sentence punctuation.
You should check the spaces after periods when
reading your output to make sure you haven't
forgotten any special cases.  Generating an
ellipsis
       \ldots\               % `\ ' is needed after `\ldots' because TeX
                             % ignores spaces after command names like \ldots
                             % made from \ + letters.
                             %
                             % Note how a `%' character causes TeX to ignore
                             % the end of the input line, so these blank lines
                             % do not start a new paragraph.
                             %
with the right spacing around the periods requires
a special command.

\LaTeX\ interprets some common characters as
commands, so you must type special commands to
generate them.  These characters include the
following:
       \$ \& \% \# \{ and \}.

In printing, text is usually emphasized with an
       \emph{italic}
type style.

\begin{em}
   A long segment of text can also be emphasized
   in this way.  Text within such a segment can be
   given \emph{additional} emphasis.
\end{em}

It is sometimes necessary to prevent \LaTeX\ from
breaking a line where it might otherwise do so.
This may be at a space, as between the ``Mr.''\ and
``Jones'' in
       ``Mr.~Jones'',        % ~ produces an unbreakable interword space.
or within a word---especially when the word is a
symbol like
       \mbox{\emph{itemnum}}
that makes little sense when hyphenated across
lines.

Footnotes\footnote{This is an example of a footnote.}
pose no problem.

\LaTeX\ is good at typesetting mathematical formulas
like
       \( x-3y + z = 7 \)
or
       \( a_{1} > x^{2n} + y^{2n} > x' \)
or
       \( \ip{A}{B} = \sum_{i} a_{i} b_{i} \).
The spaces you type in a formula are
ignored.  Remember that a letter like
       $x$                   % $ ... $  and  \( ... \)  are equivalent
is a formula when it denotes a mathematical
symbol, and it should be typed as one.

\section{Displayed Text}

Text is displayed by indenting it from the left
margin.  Quotations are commonly displayed.  There
are short quotations
\begin{quote}
   This is a short quotation.  It consists of a
   single paragraph of text.  See how it is formatted.
\end{quote}
and longer ones.
\begin{quotation}
   This is a longer quotation.  It consists of two
   paragraphs of text, neither of which are
   particularly interesting.

   This is the second paragraph of the quotation.  It
   is just as dull as the first paragraph.
\end{quotation}
Another frequently-displayed structure is a list.
The following is an example of an \emph{itemized}
list.
\begin{itemize}
   \item This is the first item of an itemized list.
         Each item in the list is marked with a ``tick''.
         You don't have to worry about what kind of tick
         mark is used.

   \item This is the second item of the list.  It
         contains another list nested inside it.  The inner
         list is an \emph{enumerated} list.
         \begin{enumerate}
            \item This is the first item of an enumerated
                  list that is nested within the itemized list.

            \item This is the second item of the inner list.
                  \LaTeX\ allows you to nest lists deeper than
                  you really should.
         \end{enumerate}
         This is the rest of the second item of the outer
         list.  It is no more interesting than any other
         part of the item.
   \item This is the third item of the list.
\end{itemize}
You can even display poetry.
\begin{verse}
   There is an environment
    for verse \\             % The \\ command separates lines
   Whose features some poets % within a stanza.
   will curse.

                             % One or more blank lines separate stanzas.

   For instead of making\\
   Them do \emph{all} line breaking, \\
   It allows them to put too many words on a line when they'd rather be
   forced to be terse.
\end{verse}

Mathematical formulas may also be displayed.  A
displayed formula
is
one-line long; multiline
formulas require special formatting instructions.
   \[  \ip{\Gamma}{\psi'} = x'' + y^{2} + z_{i}^{n}\]
Don't start a paragraph with a displayed equation,
nor make one a paragraph by itself.

\end{document}               % End of document.

\section{Bench Section 97}
% synthetic bench section 97
% This is a sample LaTeX input file.  (Version of 12 August 2004.)
%
% A '%' character causes TeX to ignore all remaining text on the line,
% and is used for comments like this one.

\documentclass{article}      % Specifies the document class

                             % The preamble begins here.
\title{An Example Document}  % Declares the document's title.
\author{Leslie Lamport}      % Declares the author's name.
\date{January 21, 1994}      % Deleting this command produces today's date.

\newcommand{\ip}[2]{(#1, #2)}
                             % Defines \ip{arg1}{arg2} to mean
                             % (arg1, arg2).

%\newcommand{\ip}[2]{\langle #1 | #2\rangle}
                             % This is an alternative definition of
                             % \ip that is commented out.

\begin{document}             % End of preamble and beginning of text.

\maketitle                   % Produces the title.

This is an example input file.  Comparing it with
the output it generates can show you how to
produce a simple document of your own.

\section{Ordinary Text}      % Produces section heading.  Lower-level
                             % sections are begun with similar
                             % \subsection and \subsubsection commands.

The ends  of words and sentences are marked
  by   spaces. It  doesn't matter how many
spaces    you type; one is as good as 100.  The
end of   a line counts as a space.

One   or more   blank lines denote the  end
of  a paragraph.

Since any number of consecutive spaces are treated
like a single one, the formatting of the input
file makes no difference to
      \LaTeX,                % The \LaTeX command generates the LaTeX logo.
but it makes a difference to you.  When you use
\LaTeX, making your input file as easy to read
as possible will be a great help as you write
your document and when you change it.  This sample
file shows how you can add comments to your own input
file.

Because printing is different from typewriting,
there are a number of things that you have to do
differently when preparing an input file than if
you were just typing the document directly.
Quotation marks like
       ``this''
have to be handled specially, as do quotes within
quotes:
       ``\,`this'            % \, separates the double and single quote.
        is what I just
        wrote, not  `that'\,''.

Dashes come in three sizes: an
       intra-word
dash, a medium dash for number ranges like
       1--2,
and a punctuation
       dash---like
this.

A sentence-ending space should be larger than the
space between words within a sentence.  You
sometimes have to type special commands in
conjunction with punctuation characters to get
this right, as in the following sentence.
       Gnats, gnus, etc.\ all  % `\ ' makes an inter-word space.
       begin with G\@.         % \@ marks end-of-sentence punctuation.
You should check the spaces after periods when
reading your output to make sure you haven't
forgotten any special cases.  Generating an
ellipsis
       \ldots\               % `\ ' is needed after `\ldots' because TeX
                             % ignores spaces after command names like \ldots
                             % made from \ + letters.
                             %
                             % Note how a `%' character causes TeX to ignore
                             % the end of the input line, so these blank lines
                             % do not start a new paragraph.
                             %
with the right spacing around the periods requires
a special command.

\LaTeX\ interprets some common characters as
commands, so you must type special commands to
generate them.  These characters include the
following:
       \$ \& \% \# \{ and \}.

In printing, text is usually emphasized with an
       \emph{italic}
type style.

\begin{em}
   A long segment of text can also be emphasized
   in this way.  Text within such a segment can be
   given \emph{additional} emphasis.
\end{em}

It is sometimes necessary to prevent \LaTeX\ from
breaking a line where it might otherwise do so.
This may be at a space, as between the ``Mr.''\ and
``Jones'' in
       ``Mr.~Jones'',        % ~ produces an unbreakable interword space.
or within a word---especially when the word is a
symbol like
       \mbox{\emph{itemnum}}
that makes little sense when hyphenated across
lines.

Footnotes\footnote{This is an example of a footnote.}
pose no problem.

\LaTeX\ is good at typesetting mathematical formulas
like
       \( x-3y + z = 7 \)
or
       \( a_{1} > x^{2n} + y^{2n} > x' \)
or
       \( \ip{A}{B} = \sum_{i} a_{i} b_{i} \).
The spaces you type in a formula are
ignored.  Remember that a letter like
       $x$                   % $ ... $  and  \( ... \)  are equivalent
is a formula when it denotes a mathematical
symbol, and it should be typed as one.

\section{Displayed Text}

Text is displayed by indenting it from the left
margin.  Quotations are commonly displayed.  There
are short quotations
\begin{quote}
   This is a short quotation.  It consists of a
   single paragraph of text.  See how it is formatted.
\end{quote}
and longer ones.
\begin{quotation}
   This is a longer quotation.  It consists of two
   paragraphs of text, neither of which are
   particularly interesting.

   This is the second paragraph of the quotation.  It
   is just as dull as the first paragraph.
\end{quotation}
Another frequently-displayed structure is a list.
The following is an example of an \emph{itemized}
list.
\begin{itemize}
   \item This is the first item of an itemized list.
         Each item in the list is marked with a ``tick''.
         You don't have to worry about what kind of tick
         mark is used.

   \item This is the second item of the list.  It
         contains another list nested inside it.  The inner
         list is an \emph{enumerated} list.
         \begin{enumerate}
            \item This is the first item of an enumerated
                  list that is nested within the itemized list.

            \item This is the second item of the inner list.
                  \LaTeX\ allows you to nest lists deeper than
                  you really should.
         \end{enumerate}
         This is the rest of the second item of the outer
         list.  It is no more interesting than any other
         part of the item.
   \item This is the third item of the list.
\end{itemize}
You can even display poetry.
\begin{verse}
   There is an environment
    for verse \\             % The \\ command separates lines
   Whose features some poets % within a stanza.
   will curse.

                             % One or more blank lines separate stanzas.

   For instead of making\\
   Them do \emph{all} line breaking, \\
   It allows them to put too many words on a line when they'd rather be
   forced to be terse.
\end{verse}

Mathematical formulas may also be displayed.  A
displayed formula
is
one-line long; multiline
formulas require special formatting instructions.
   \[  \ip{\Gamma}{\psi'} = x'' + y^{2} + z_{i}^{n}\]
Don't start a paragraph with a displayed equation,
nor make one a paragraph by itself.

\end{document}               % End of document.

\section{Bench Section 98}
% synthetic bench section 98
% This is a sample LaTeX input file.  (Version of 12 August 2004.)
%
% A '%' character causes TeX to ignore all remaining text on the line,
% and is used for comments like this one.

\documentclass{article}      % Specifies the document class

                             % The preamble begins here.
\title{An Example Document}  % Declares the document's title.
\author{Leslie Lamport}      % Declares the author's name.
\date{January 21, 1994}      % Deleting this command produces today's date.

\newcommand{\ip}[2]{(#1, #2)}
                             % Defines \ip{arg1}{arg2} to mean
                             % (arg1, arg2).

%\newcommand{\ip}[2]{\langle #1 | #2\rangle}
                             % This is an alternative definition of
                             % \ip that is commented out.

\begin{document}             % End of preamble and beginning of text.

\maketitle                   % Produces the title.

This is an example input file.  Comparing it with
the output it generates can show you how to
produce a simple document of your own.

\section{Ordinary Text}      % Produces section heading.  Lower-level
                             % sections are begun with similar
                             % \subsection and \subsubsection commands.

The ends  of words and sentences are marked
  by   spaces. It  doesn't matter how many
spaces    you type; one is as good as 100.  The
end of   a line counts as a space.

One   or more   blank lines denote the  end
of  a paragraph.

Since any number of consecutive spaces are treated
like a single one, the formatting of the input
file makes no difference to
      \LaTeX,                % The \LaTeX command generates the LaTeX logo.
but it makes a difference to you.  When you use
\LaTeX, making your input file as easy to read
as possible will be a great help as you write
your document and when you change it.  This sample
file shows how you can add comments to your own input
file.

Because printing is different from typewriting,
there are a number of things that you have to do
differently when preparing an input file than if
you were just typing the document directly.
Quotation marks like
       ``this''
have to be handled specially, as do quotes within
quotes:
       ``\,`this'            % \, separates the double and single quote.
        is what I just
        wrote, not  `that'\,''.

Dashes come in three sizes: an
       intra-word
dash, a medium dash for number ranges like
       1--2,
and a punctuation
       dash---like
this.

A sentence-ending space should be larger than the
space between words within a sentence.  You
sometimes have to type special commands in
conjunction with punctuation characters to get
this right, as in the following sentence.
       Gnats, gnus, etc.\ all  % `\ ' makes an inter-word space.
       begin with G\@.         % \@ marks end-of-sentence punctuation.
You should check the spaces after periods when
reading your output to make sure you haven't
forgotten any special cases.  Generating an
ellipsis
       \ldots\               % `\ ' is needed after `\ldots' because TeX
                             % ignores spaces after command names like \ldots
                             % made from \ + letters.
                             %
                             % Note how a `%' character causes TeX to ignore
                             % the end of the input line, so these blank lines
                             % do not start a new paragraph.
                             %
with the right spacing around the periods requires
a special command.

\LaTeX\ interprets some common characters as
commands, so you must type special commands to
generate them.  These characters include the
following:
       \$ \& \% \# \{ and \}.

In printing, text is usually emphasized with an
       \emph{italic}
type style.

\begin{em}
   A long segment of text can also be emphasized
   in this way.  Text within such a segment can be
   given \emph{additional} emphasis.
\end{em}

It is sometimes necessary to prevent \LaTeX\ from
breaking a line where it might otherwise do so.
This may be at a space, as between the ``Mr.''\ and
``Jones'' in
       ``Mr.~Jones'',        % ~ produces an unbreakable interword space.
or within a word---especially when the word is a
symbol like
       \mbox{\emph{itemnum}}
that makes little sense when hyphenated across
lines.

Footnotes\footnote{This is an example of a footnote.}
pose no problem.

\LaTeX\ is good at typesetting mathematical formulas
like
       \( x-3y + z = 7 \)
or
       \( a_{1} > x^{2n} + y^{2n} > x' \)
or
       \( \ip{A}{B} = \sum_{i} a_{i} b_{i} \).
The spaces you type in a formula are
ignored.  Remember that a letter like
       $x$                   % $ ... $  and  \( ... \)  are equivalent
is a formula when it denotes a mathematical
symbol, and it should be typed as one.

\section{Displayed Text}

Text is displayed by indenting it from the left
margin.  Quotations are commonly displayed.  There
are short quotations
\begin{quote}
   This is a short quotation.  It consists of a
   single paragraph of text.  See how it is formatted.
\end{quote}
and longer ones.
\begin{quotation}
   This is a longer quotation.  It consists of two
   paragraphs of text, neither of which are
   particularly interesting.

   This is the second paragraph of the quotation.  It
   is just as dull as the first paragraph.
\end{quotation}
Another frequently-displayed structure is a list.
The following is an example of an \emph{itemized}
list.
\begin{itemize}
   \item This is the first item of an itemized list.
         Each item in the list is marked with a ``tick''.
         You don't have to worry about what kind of tick
         mark is used.

   \item This is the second item of the list.  It
         contains another list nested inside it.  The inner
         list is an \emph{enumerated} list.
         \begin{enumerate}
            \item This is the first item of an enumerated
                  list that is nested within the itemized list.

            \item This is the second item of the inner list.
                  \LaTeX\ allows you to nest lists deeper than
                  you really should.
         \end{enumerate}
         This is the rest of the second item of the outer
         list.  It is no more interesting than any other
         part of the item.
   \item This is the third item of the list.
\end{itemize}
You can even display poetry.
\begin{verse}
   There is an environment
    for verse \\             % The \\ command separates lines
   Whose features some poets % within a stanza.
   will curse.

                             % One or more blank lines separate stanzas.

   For instead of making\\
   Them do \emph{all} line breaking, \\
   It allows them to put too many words on a line when they'd rather be
   forced to be terse.
\end{verse}

Mathematical formulas may also be displayed.  A
displayed formula
is
one-line long; multiline
formulas require special formatting instructions.
   \[  \ip{\Gamma}{\psi'} = x'' + y^{2} + z_{i}^{n}\]
Don't start a paragraph with a displayed equation,
nor make one a paragraph by itself.

\end{document}               % End of document.

\section{Bench Section 99}
% synthetic bench section 99
% This is a sample LaTeX input file.  (Version of 12 August 2004.)
%
% A '%' character causes TeX to ignore all remaining text on the line,
% and is used for comments like this one.

\documentclass{article}      % Specifies the document class

                             % The preamble begins here.
\title{An Example Document}  % Declares the document's title.
\author{Leslie Lamport}      % Declares the author's name.
\date{January 21, 1994}      % Deleting this command produces today's date.

\newcommand{\ip}[2]{(#1, #2)}
                             % Defines \ip{arg1}{arg2} to mean
                             % (arg1, arg2).

%\newcommand{\ip}[2]{\langle #1 | #2\rangle}
                             % This is an alternative definition of
                             % \ip that is commented out.

\begin{document}             % End of preamble and beginning of text.

\maketitle                   % Produces the title.

This is an example input file.  Comparing it with
the output it generates can show you how to
produce a simple document of your own.

\section{Ordinary Text}      % Produces section heading.  Lower-level
                             % sections are begun with similar
                             % \subsection and \subsubsection commands.

The ends  of words and sentences are marked
  by   spaces. It  doesn't matter how many
spaces    you type; one is as good as 100.  The
end of   a line counts as a space.

One   or more   blank lines denote the  end
of  a paragraph.

Since any number of consecutive spaces are treated
like a single one, the formatting of the input
file makes no difference to
      \LaTeX,                % The \LaTeX command generates the LaTeX logo.
but it makes a difference to you.  When you use
\LaTeX, making your input file as easy to read
as possible will be a great help as you write
your document and when you change it.  This sample
file shows how you can add comments to your own input
file.

Because printing is different from typewriting,
there are a number of things that you have to do
differently when preparing an input file than if
you were just typing the document directly.
Quotation marks like
       ``this''
have to be handled specially, as do quotes within
quotes:
       ``\,`this'            % \, separates the double and single quote.
        is what I just
        wrote, not  `that'\,''.

Dashes come in three sizes: an
       intra-word
dash, a medium dash for number ranges like
       1--2,
and a punctuation
       dash---like
this.

A sentence-ending space should be larger than the
space between words within a sentence.  You
sometimes have to type special commands in
conjunction with punctuation characters to get
this right, as in the following sentence.
       Gnats, gnus, etc.\ all  % `\ ' makes an inter-word space.
       begin with G\@.         % \@ marks end-of-sentence punctuation.
You should check the spaces after periods when
reading your output to make sure you haven't
forgotten any special cases.  Generating an
ellipsis
       \ldots\               % `\ ' is needed after `\ldots' because TeX
                             % ignores spaces after command names like \ldots
                             % made from \ + letters.
                             %
                             % Note how a `%' character causes TeX to ignore
                             % the end of the input line, so these blank lines
                             % do not start a new paragraph.
                             %
with the right spacing around the periods requires
a special command.

\LaTeX\ interprets some common characters as
commands, so you must type special commands to
generate them.  These characters include the
following:
       \$ \& \% \# \{ and \}.

In printing, text is usually emphasized with an
       \emph{italic}
type style.

\begin{em}
   A long segment of text can also be emphasized
   in this way.  Text within such a segment can be
   given \emph{additional} emphasis.
\end{em}

It is sometimes necessary to prevent \LaTeX\ from
breaking a line where it might otherwise do so.
This may be at a space, as between the ``Mr.''\ and
``Jones'' in
       ``Mr.~Jones'',        % ~ produces an unbreakable interword space.
or within a word---especially when the word is a
symbol like
       \mbox{\emph{itemnum}}
that makes little sense when hyphenated across
lines.

Footnotes\footnote{This is an example of a footnote.}
pose no problem.

\LaTeX\ is good at typesetting mathematical formulas
like
       \( x-3y + z = 7 \)
or
       \( a_{1} > x^{2n} + y^{2n} > x' \)
or
       \( \ip{A}{B} = \sum_{i} a_{i} b_{i} \).
The spaces you type in a formula are
ignored.  Remember that a letter like
       $x$                   % $ ... $  and  \( ... \)  are equivalent
is a formula when it denotes a mathematical
symbol, and it should be typed as one.

\section{Displayed Text}

Text is displayed by indenting it from the left
margin.  Quotations are commonly displayed.  There
are short quotations
\begin{quote}
   This is a short quotation.  It consists of a
   single paragraph of text.  See how it is formatted.
\end{quote}
and longer ones.
\begin{quotation}
   This is a longer quotation.  It consists of two
   paragraphs of text, neither of which are
   particularly interesting.

   This is the second paragraph of the quotation.  It
   is just as dull as the first paragraph.
\end{quotation}
Another frequently-displayed structure is a list.
The following is an example of an \emph{itemized}
list.
\begin{itemize}
   \item This is the first item of an itemized list.
         Each item in the list is marked with a ``tick''.
         You don't have to worry about what kind of tick
         mark is used.

   \item This is the second item of the list.  It
         contains another list nested inside it.  The inner
         list is an \emph{enumerated} list.
         \begin{enumerate}
            \item This is the first item of an enumerated
                  list that is nested within the itemized list.

            \item This is the second item of the inner list.
                  \LaTeX\ allows you to nest lists deeper than
                  you really should.
         \end{enumerate}
         This is the rest of the second item of the outer
         list.  It is no more interesting than any other
         part of the item.
   \item This is the third item of the list.
\end{itemize}
You can even display poetry.
\begin{verse}
   There is an environment
    for verse \\             % The \\ command separates lines
   Whose features some poets % within a stanza.
   will curse.

                             % One or more blank lines separate stanzas.

   For instead of making\\
   Them do \emph{all} line breaking, \\
   It allows them to put too many words on a line when they'd rather be
   forced to be terse.
\end{verse}

Mathematical formulas may also be displayed.  A
displayed formula
is
one-line long; multiline
formulas require special formatting instructions.
   \[  \ip{\Gamma}{\psi'} = x'' + y^{2} + z_{i}^{n}\]
Don't start a paragraph with a displayed equation,
nor make one a paragraph by itself.

\end{document}               % End of document.

\section{Bench Section 100}
% synthetic bench section 100
% This is a sample LaTeX input file.  (Version of 12 August 2004.)
%
% A '%' character causes TeX to ignore all remaining text on the line,
% and is used for comments like this one.

\documentclass{article}      % Specifies the document class

                             % The preamble begins here.
\title{An Example Document}  % Declares the document's title.
\author{Leslie Lamport}      % Declares the author's name.
\date{January 21, 1994}      % Deleting this command produces today's date.

\newcommand{\ip}[2]{(#1, #2)}
                             % Defines \ip{arg1}{arg2} to mean
                             % (arg1, arg2).

%\newcommand{\ip}[2]{\langle #1 | #2\rangle}
                             % This is an alternative definition of
                             % \ip that is commented out.

\begin{document}             % End of preamble and beginning of text.

\maketitle                   % Produces the title.

This is an example input file.  Comparing it with
the output it generates can show you how to
produce a simple document of your own.

\section{Ordinary Text}      % Produces section heading.  Lower-level
                             % sections are begun with similar
                             % \subsection and \subsubsection commands.

The ends  of words and sentences are marked
  by   spaces. It  doesn't matter how many
spaces    you type; one is as good as 100.  The
end of   a line counts as a space.

One   or more   blank lines denote the  end
of  a paragraph.

Since any number of consecutive spaces are treated
like a single one, the formatting of the input
file makes no difference to
      \LaTeX,                % The \LaTeX command generates the LaTeX logo.
but it makes a difference to you.  When you use
\LaTeX, making your input file as easy to read
as possible will be a great help as you write
your document and when you change it.  This sample
file shows how you can add comments to your own input
file.

Because printing is different from typewriting,
there are a number of things that you have to do
differently when preparing an input file than if
you were just typing the document directly.
Quotation marks like
       ``this''
have to be handled specially, as do quotes within
quotes:
       ``\,`this'            % \, separates the double and single quote.
        is what I just
        wrote, not  `that'\,''.

Dashes come in three sizes: an
       intra-word
dash, a medium dash for number ranges like
       1--2,
and a punctuation
       dash---like
this.

A sentence-ending space should be larger than the
space between words within a sentence.  You
sometimes have to type special commands in
conjunction with punctuation characters to get
this right, as in the following sentence.
       Gnats, gnus, etc.\ all  % `\ ' makes an inter-word space.
       begin with G\@.         % \@ marks end-of-sentence punctuation.
You should check the spaces after periods when
reading your output to make sure you haven't
forgotten any special cases.  Generating an
ellipsis
       \ldots\               % `\ ' is needed after `\ldots' because TeX
                             % ignores spaces after command names like \ldots
                             % made from \ + letters.
                             %
                             % Note how a `%' character causes TeX to ignore
                             % the end of the input line, so these blank lines
                             % do not start a new paragraph.
                             %
with the right spacing around the periods requires
a special command.

\LaTeX\ interprets some common characters as
commands, so you must type special commands to
generate them.  These characters include the
following:
       \$ \& \% \# \{ and \}.

In printing, text is usually emphasized with an
       \emph{italic}
type style.

\begin{em}
   A long segment of text can also be emphasized
   in this way.  Text within such a segment can be
   given \emph{additional} emphasis.
\end{em}

It is sometimes necessary to prevent \LaTeX\ from
breaking a line where it might otherwise do so.
This may be at a space, as between the ``Mr.''\ and
``Jones'' in
       ``Mr.~Jones'',        % ~ produces an unbreakable interword space.
or within a word---especially when the word is a
symbol like
       \mbox{\emph{itemnum}}
that makes little sense when hyphenated across
lines.

Footnotes\footnote{This is an example of a footnote.}
pose no problem.

\LaTeX\ is good at typesetting mathematical formulas
like
       \( x-3y + z = 7 \)
or
       \( a_{1} > x^{2n} + y^{2n} > x' \)
or
       \( \ip{A}{B} = \sum_{i} a_{i} b_{i} \).
The spaces you type in a formula are
ignored.  Remember that a letter like
       $x$                   % $ ... $  and  \( ... \)  are equivalent
is a formula when it denotes a mathematical
symbol, and it should be typed as one.

\section{Displayed Text}

Text is displayed by indenting it from the left
margin.  Quotations are commonly displayed.  There
are short quotations
\begin{quote}
   This is a short quotation.  It consists of a
   single paragraph of text.  See how it is formatted.
\end{quote}
and longer ones.
\begin{quotation}
   This is a longer quotation.  It consists of two
   paragraphs of text, neither of which are
   particularly interesting.

   This is the second paragraph of the quotation.  It
   is just as dull as the first paragraph.
\end{quotation}
Another frequently-displayed structure is a list.
The following is an example of an \emph{itemized}
list.
\begin{itemize}
   \item This is the first item of an itemized list.
         Each item in the list is marked with a ``tick''.
         You don't have to worry about what kind of tick
         mark is used.

   \item This is the second item of the list.  It
         contains another list nested inside it.  The inner
         list is an \emph{enumerated} list.
         \begin{enumerate}
            \item This is the first item of an enumerated
                  list that is nested within the itemized list.

            \item This is the second item of the inner list.
                  \LaTeX\ allows you to nest lists deeper than
                  you really should.
         \end{enumerate}
         This is the rest of the second item of the outer
         list.  It is no more interesting than any other
         part of the item.
   \item This is the third item of the list.
\end{itemize}
You can even display poetry.
\begin{verse}
   There is an environment
    for verse \\             % The \\ command separates lines
   Whose features some poets % within a stanza.
   will curse.

                             % One or more blank lines separate stanzas.

   For instead of making\\
   Them do \emph{all} line breaking, \\
   It allows them to put too many words on a line when they'd rather be
   forced to be terse.
\end{verse}

Mathematical formulas may also be displayed.  A
displayed formula
is
one-line long; multiline
formulas require special formatting instructions.
   \[  \ip{\Gamma}{\psi'} = x'' + y^{2} + z_{i}^{n}\]
Don't start a paragraph with a displayed equation,
nor make one a paragraph by itself.

\end{document}               % End of document.

\section{Bench Section 101}
% synthetic bench section 101
% This is a sample LaTeX input file.  (Version of 12 August 2004.)
%
% A '%' character causes TeX to ignore all remaining text on the line,
% and is used for comments like this one.

\documentclass{article}      % Specifies the document class

                             % The preamble begins here.
\title{An Example Document}  % Declares the document's title.
\author{Leslie Lamport}      % Declares the author's name.
\date{January 21, 1994}      % Deleting this command produces today's date.

\newcommand{\ip}[2]{(#1, #2)}
                             % Defines \ip{arg1}{arg2} to mean
                             % (arg1, arg2).

%\newcommand{\ip}[2]{\langle #1 | #2\rangle}
                             % This is an alternative definition of
                             % \ip that is commented out.

\begin{document}             % End of preamble and beginning of text.

\maketitle                   % Produces the title.

This is an example input file.  Comparing it with
the output it generates can show you how to
produce a simple document of your own.

\section{Ordinary Text}      % Produces section heading.  Lower-level
                             % sections are begun with similar
                             % \subsection and \subsubsection commands.

The ends  of words and sentences are marked
  by   spaces. It  doesn't matter how many
spaces    you type; one is as good as 100.  The
end of   a line counts as a space.

One   or more   blank lines denote the  end
of  a paragraph.

Since any number of consecutive spaces are treated
like a single one, the formatting of the input
file makes no difference to
      \LaTeX,                % The \LaTeX command generates the LaTeX logo.
but it makes a difference to you.  When you use
\LaTeX, making your input file as easy to read
as possible will be a great help as you write
your document and when you change it.  This sample
file shows how you can add comments to your own input
file.

Because printing is different from typewriting,
there are a number of things that you have to do
differently when preparing an input file than if
you were just typing the document directly.
Quotation marks like
       ``this''
have to be handled specially, as do quotes within
quotes:
       ``\,`this'            % \, separates the double and single quote.
        is what I just
        wrote, not  `that'\,''.

Dashes come in three sizes: an
       intra-word
dash, a medium dash for number ranges like
       1--2,
and a punctuation
       dash---like
this.

A sentence-ending space should be larger than the
space between words within a sentence.  You
sometimes have to type special commands in
conjunction with punctuation characters to get
this right, as in the following sentence.
       Gnats, gnus, etc.\ all  % `\ ' makes an inter-word space.
       begin with G\@.         % \@ marks end-of-sentence punctuation.
You should check the spaces after periods when
reading your output to make sure you haven't
forgotten any special cases.  Generating an
ellipsis
       \ldots\               % `\ ' is needed after `\ldots' because TeX
                             % ignores spaces after command names like \ldots
                             % made from \ + letters.
                             %
                             % Note how a `%' character causes TeX to ignore
                             % the end of the input line, so these blank lines
                             % do not start a new paragraph.
                             %
with the right spacing around the periods requires
a special command.

\LaTeX\ interprets some common characters as
commands, so you must type special commands to
generate them.  These characters include the
following:
       \$ \& \% \# \{ and \}.

In printing, text is usually emphasized with an
       \emph{italic}
type style.

\begin{em}
   A long segment of text can also be emphasized
   in this way.  Text within such a segment can be
   given \emph{additional} emphasis.
\end{em}

It is sometimes necessary to prevent \LaTeX\ from
breaking a line where it might otherwise do so.
This may be at a space, as between the ``Mr.''\ and
``Jones'' in
       ``Mr.~Jones'',        % ~ produces an unbreakable interword space.
or within a word---especially when the word is a
symbol like
       \mbox{\emph{itemnum}}
that makes little sense when hyphenated across
lines.

Footnotes\footnote{This is an example of a footnote.}
pose no problem.

\LaTeX\ is good at typesetting mathematical formulas
like
       \( x-3y + z = 7 \)
or
       \( a_{1} > x^{2n} + y^{2n} > x' \)
or
       \( \ip{A}{B} = \sum_{i} a_{i} b_{i} \).
The spaces you type in a formula are
ignored.  Remember that a letter like
       $x$                   % $ ... $  and  \( ... \)  are equivalent
is a formula when it denotes a mathematical
symbol, and it should be typed as one.

\section{Displayed Text}

Text is displayed by indenting it from the left
margin.  Quotations are commonly displayed.  There
are short quotations
\begin{quote}
   This is a short quotation.  It consists of a
   single paragraph of text.  See how it is formatted.
\end{quote}
and longer ones.
\begin{quotation}
   This is a longer quotation.  It consists of two
   paragraphs of text, neither of which are
   particularly interesting.

   This is the second paragraph of the quotation.  It
   is just as dull as the first paragraph.
\end{quotation}
Another frequently-displayed structure is a list.
The following is an example of an \emph{itemized}
list.
\begin{itemize}
   \item This is the first item of an itemized list.
         Each item in the list is marked with a ``tick''.
         You don't have to worry about what kind of tick
         mark is used.

   \item This is the second item of the list.  It
         contains another list nested inside it.  The inner
         list is an \emph{enumerated} list.
         \begin{enumerate}
            \item This is the first item of an enumerated
                  list that is nested within the itemized list.

            \item This is the second item of the inner list.
                  \LaTeX\ allows you to nest lists deeper than
                  you really should.
         \end{enumerate}
         This is the rest of the second item of the outer
         list.  It is no more interesting than any other
         part of the item.
   \item This is the third item of the list.
\end{itemize}
You can even display poetry.
\begin{verse}
   There is an environment
    for verse \\             % The \\ command separates lines
   Whose features some poets % within a stanza.
   will curse.

                             % One or more blank lines separate stanzas.

   For instead of making\\
   Them do \emph{all} line breaking, \\
   It allows them to put too many words on a line when they'd rather be
   forced to be terse.
\end{verse}

Mathematical formulas may also be displayed.  A
displayed formula
is
one-line long; multiline
formulas require special formatting instructions.
   \[  \ip{\Gamma}{\psi'} = x'' + y^{2} + z_{i}^{n}\]
Don't start a paragraph with a displayed equation,
nor make one a paragraph by itself.

\end{document}               % End of document.

\section{Bench Section 102}
% synthetic bench section 102
% This is a sample LaTeX input file.  (Version of 12 August 2004.)
%
% A '%' character causes TeX to ignore all remaining text on the line,
% and is used for comments like this one.

\documentclass{article}      % Specifies the document class

                             % The preamble begins here.
\title{An Example Document}  % Declares the document's title.
\author{Leslie Lamport}      % Declares the author's name.
\date{January 21, 1994}      % Deleting this command produces today's date.

\newcommand{\ip}[2]{(#1, #2)}
                             % Defines \ip{arg1}{arg2} to mean
                             % (arg1, arg2).

%\newcommand{\ip}[2]{\langle #1 | #2\rangle}
                             % This is an alternative definition of
                             % \ip that is commented out.

\begin{document}             % End of preamble and beginning of text.

\maketitle                   % Produces the title.

This is an example input file.  Comparing it with
the output it generates can show you how to
produce a simple document of your own.

\section{Ordinary Text}      % Produces section heading.  Lower-level
                             % sections are begun with similar
                             % \subsection and \subsubsection commands.

The ends  of words and sentences are marked
  by   spaces. It  doesn't matter how many
spaces    you type; one is as good as 100.  The
end of   a line counts as a space.

One   or more   blank lines denote the  end
of  a paragraph.

Since any number of consecutive spaces are treated
like a single one, the formatting of the input
file makes no difference to
      \LaTeX,                % The \LaTeX command generates the LaTeX logo.
but it makes a difference to you.  When you use
\LaTeX, making your input file as easy to read
as possible will be a great help as you write
your document and when you change it.  This sample
file shows how you can add comments to your own input
file.

Because printing is different from typewriting,
there are a number of things that you have to do
differently when preparing an input file than if
you were just typing the document directly.
Quotation marks like
       ``this''
have to be handled specially, as do quotes within
quotes:
       ``\,`this'            % \, separates the double and single quote.
        is what I just
        wrote, not  `that'\,''.

Dashes come in three sizes: an
       intra-word
dash, a medium dash for number ranges like
       1--2,
and a punctuation
       dash---like
this.

A sentence-ending space should be larger than the
space between words within a sentence.  You
sometimes have to type special commands in
conjunction with punctuation characters to get
this right, as in the following sentence.
       Gnats, gnus, etc.\ all  % `\ ' makes an inter-word space.
       begin with G\@.         % \@ marks end-of-sentence punctuation.
You should check the spaces after periods when
reading your output to make sure you haven't
forgotten any special cases.  Generating an
ellipsis
       \ldots\               % `\ ' is needed after `\ldots' because TeX
                             % ignores spaces after command names like \ldots
                             % made from \ + letters.
                             %
                             % Note how a `%' character causes TeX to ignore
                             % the end of the input line, so these blank lines
                             % do not start a new paragraph.
                             %
with the right spacing around the periods requires
a special command.

\LaTeX\ interprets some common characters as
commands, so you must type special commands to
generate them.  These characters include the
following:
       \$ \& \% \# \{ and \}.

In printing, text is usually emphasized with an
       \emph{italic}
type style.

\begin{em}
   A long segment of text can also be emphasized
   in this way.  Text within such a segment can be
   given \emph{additional} emphasis.
\end{em}

It is sometimes necessary to prevent \LaTeX\ from
breaking a line where it might otherwise do so.
This may be at a space, as between the ``Mr.''\ and
``Jones'' in
       ``Mr.~Jones'',        % ~ produces an unbreakable interword space.
or within a word---especially when the word is a
symbol like
       \mbox{\emph{itemnum}}
that makes little sense when hyphenated across
lines.

Footnotes\footnote{This is an example of a footnote.}
pose no problem.

\LaTeX\ is good at typesetting mathematical formulas
like
       \( x-3y + z = 7 \)
or
       \( a_{1} > x^{2n} + y^{2n} > x' \)
or
       \( \ip{A}{B} = \sum_{i} a_{i} b_{i} \).
The spaces you type in a formula are
ignored.  Remember that a letter like
       $x$                   % $ ... $  and  \( ... \)  are equivalent
is a formula when it denotes a mathematical
symbol, and it should be typed as one.

\section{Displayed Text}

Text is displayed by indenting it from the left
margin.  Quotations are commonly displayed.  There
are short quotations
\begin{quote}
   This is a short quotation.  It consists of a
   single paragraph of text.  See how it is formatted.
\end{quote}
and longer ones.
\begin{quotation}
   This is a longer quotation.  It consists of two
   paragraphs of text, neither of which are
   particularly interesting.

   This is the second paragraph of the quotation.  It
   is just as dull as the first paragraph.
\end{quotation}
Another frequently-displayed structure is a list.
The following is an example of an \emph{itemized}
list.
\begin{itemize}
   \item This is the first item of an itemized list.
         Each item in the list is marked with a ``tick''.
         You don't have to worry about what kind of tick
         mark is used.

   \item This is the second item of the list.  It
         contains another list nested inside it.  The inner
         list is an \emph{enumerated} list.
         \begin{enumerate}
            \item This is the first item of an enumerated
                  list that is nested within the itemized list.

            \item This is the second item of the inner list.
                  \LaTeX\ allows you to nest lists deeper than
                  you really should.
         \end{enumerate}
         This is the rest of the second item of the outer
         list.  It is no more interesting than any other
         part of the item.
   \item This is the third item of the list.
\end{itemize}
You can even display poetry.
\begin{verse}
   There is an environment
    for verse \\             % The \\ command separates lines
   Whose features some poets % within a stanza.
   will curse.

                             % One or more blank lines separate stanzas.

   For instead of making\\
   Them do \emph{all} line breaking, \\
   It allows them to put too many words on a line when they'd rather be
   forced to be terse.
\end{verse}

Mathematical formulas may also be displayed.  A
displayed formula
is
one-line long; multiline
formulas require special formatting instructions.
   \[  \ip{\Gamma}{\psi'} = x'' + y^{2} + z_{i}^{n}\]
Don't start a paragraph with a displayed equation,
nor make one a paragraph by itself.

\end{document}               % End of document.

\section{Bench Section 103}
% synthetic bench section 103
% This is a sample LaTeX input file.  (Version of 12 August 2004.)
%
% A '%' character causes TeX to ignore all remaining text on the line,
% and is used for comments like this one.

\documentclass{article}      % Specifies the document class

                             % The preamble begins here.
\title{An Example Document}  % Declares the document's title.
\author{Leslie Lamport}      % Declares the author's name.
\date{January 21, 1994}      % Deleting this command produces today's date.

\newcommand{\ip}[2]{(#1, #2)}
                             % Defines \ip{arg1}{arg2} to mean
                             % (arg1, arg2).

%\newcommand{\ip}[2]{\langle #1 | #2\rangle}
                             % This is an alternative definition of
                             % \ip that is commented out.

\begin{document}             % End of preamble and beginning of text.

\maketitle                   % Produces the title.

This is an example input file.  Comparing it with
the output it generates can show you how to
produce a simple document of your own.

\section{Ordinary Text}      % Produces section heading.  Lower-level
                             % sections are begun with similar
                             % \subsection and \subsubsection commands.

The ends  of words and sentences are marked
  by   spaces. It  doesn't matter how many
spaces    you type; one is as good as 100.  The
end of   a line counts as a space.

One   or more   blank lines denote the  end
of  a paragraph.

Since any number of consecutive spaces are treated
like a single one, the formatting of the input
file makes no difference to
      \LaTeX,                % The \LaTeX command generates the LaTeX logo.
but it makes a difference to you.  When you use
\LaTeX, making your input file as easy to read
as possible will be a great help as you write
your document and when you change it.  This sample
file shows how you can add comments to your own input
file.

Because printing is different from typewriting,
there are a number of things that you have to do
differently when preparing an input file than if
you were just typing the document directly.
Quotation marks like
       ``this''
have to be handled specially, as do quotes within
quotes:
       ``\,`this'            % \, separates the double and single quote.
        is what I just
        wrote, not  `that'\,''.

Dashes come in three sizes: an
       intra-word
dash, a medium dash for number ranges like
       1--2,
and a punctuation
       dash---like
this.

A sentence-ending space should be larger than the
space between words within a sentence.  You
sometimes have to type special commands in
conjunction with punctuation characters to get
this right, as in the following sentence.
       Gnats, gnus, etc.\ all  % `\ ' makes an inter-word space.
       begin with G\@.         % \@ marks end-of-sentence punctuation.
You should check the spaces after periods when
reading your output to make sure you haven't
forgotten any special cases.  Generating an
ellipsis
       \ldots\               % `\ ' is needed after `\ldots' because TeX
                             % ignores spaces after command names like \ldots
                             % made from \ + letters.
                             %
                             % Note how a `%' character causes TeX to ignore
                             % the end of the input line, so these blank lines
                             % do not start a new paragraph.
                             %
with the right spacing around the periods requires
a special command.

\LaTeX\ interprets some common characters as
commands, so you must type special commands to
generate them.  These characters include the
following:
       \$ \& \% \# \{ and \}.

In printing, text is usually emphasized with an
       \emph{italic}
type style.

\begin{em}
   A long segment of text can also be emphasized
   in this way.  Text within such a segment can be
   given \emph{additional} emphasis.
\end{em}

It is sometimes necessary to prevent \LaTeX\ from
breaking a line where it might otherwise do so.
This may be at a space, as between the ``Mr.''\ and
``Jones'' in
       ``Mr.~Jones'',        % ~ produces an unbreakable interword space.
or within a word---especially when the word is a
symbol like
       \mbox{\emph{itemnum}}
that makes little sense when hyphenated across
lines.

Footnotes\footnote{This is an example of a footnote.}
pose no problem.

\LaTeX\ is good at typesetting mathematical formulas
like
       \( x-3y + z = 7 \)
or
       \( a_{1} > x^{2n} + y^{2n} > x' \)
or
       \( \ip{A}{B} = \sum_{i} a_{i} b_{i} \).
The spaces you type in a formula are
ignored.  Remember that a letter like
       $x$                   % $ ... $  and  \( ... \)  are equivalent
is a formula when it denotes a mathematical
symbol, and it should be typed as one.

\section{Displayed Text}

Text is displayed by indenting it from the left
margin.  Quotations are commonly displayed.  There
are short quotations
\begin{quote}
   This is a short quotation.  It consists of a
   single paragraph of text.  See how it is formatted.
\end{quote}
and longer ones.
\begin{quotation}
   This is a longer quotation.  It consists of two
   paragraphs of text, neither of which are
   particularly interesting.

   This is the second paragraph of the quotation.  It
   is just as dull as the first paragraph.
\end{quotation}
Another frequently-displayed structure is a list.
The following is an example of an \emph{itemized}
list.
\begin{itemize}
   \item This is the first item of an itemized list.
         Each item in the list is marked with a ``tick''.
         You don't have to worry about what kind of tick
         mark is used.

   \item This is the second item of the list.  It
         contains another list nested inside it.  The inner
         list is an \emph{enumerated} list.
         \begin{enumerate}
            \item This is the first item of an enumerated
                  list that is nested within the itemized list.

            \item This is the second item of the inner list.
                  \LaTeX\ allows you to nest lists deeper than
                  you really should.
         \end{enumerate}
         This is the rest of the second item of the outer
         list.  It is no more interesting than any other
         part of the item.
   \item This is the third item of the list.
\end{itemize}
You can even display poetry.
\begin{verse}
   There is an environment
    for verse \\             % The \\ command separates lines
   Whose features some poets % within a stanza.
   will curse.

                             % One or more blank lines separate stanzas.

   For instead of making\\
   Them do \emph{all} line breaking, \\
   It allows them to put too many words on a line when they'd rather be
   forced to be terse.
\end{verse}

Mathematical formulas may also be displayed.  A
displayed formula
is
one-line long; multiline
formulas require special formatting instructions.
   \[  \ip{\Gamma}{\psi'} = x'' + y^{2} + z_{i}^{n}\]
Don't start a paragraph with a displayed equation,
nor make one a paragraph by itself.

\end{document}               % End of document.

\section{Bench Section 104}
% synthetic bench section 104
% This is a sample LaTeX input file.  (Version of 12 August 2004.)
%
% A '%' character causes TeX to ignore all remaining text on the line,
% and is used for comments like this one.

\documentclass{article}      % Specifies the document class

                             % The preamble begins here.
\title{An Example Document}  % Declares the document's title.
\author{Leslie Lamport}      % Declares the author's name.
\date{January 21, 1994}      % Deleting this command produces today's date.

\newcommand{\ip}[2]{(#1, #2)}
                             % Defines \ip{arg1}{arg2} to mean
                             % (arg1, arg2).

%\newcommand{\ip}[2]{\langle #1 | #2\rangle}
                             % This is an alternative definition of
                             % \ip that is commented out.

\begin{document}             % End of preamble and beginning of text.

\maketitle                   % Produces the title.

This is an example input file.  Comparing it with
the output it generates can show you how to
produce a simple document of your own.

\section{Ordinary Text}      % Produces section heading.  Lower-level
                             % sections are begun with similar
                             % \subsection and \subsubsection commands.

The ends  of words and sentences are marked
  by   spaces. It  doesn't matter how many
spaces    you type; one is as good as 100.  The
end of   a line counts as a space.

One   or more   blank lines denote the  end
of  a paragraph.

Since any number of consecutive spaces are treated
like a single one, the formatting of the input
file makes no difference to
      \LaTeX,                % The \LaTeX command generates the LaTeX logo.
but it makes a difference to you.  When you use
\LaTeX, making your input file as easy to read
as possible will be a great help as you write
your document and when you change it.  This sample
file shows how you can add comments to your own input
file.

Because printing is different from typewriting,
there are a number of things that you have to do
differently when preparing an input file than if
you were just typing the document directly.
Quotation marks like
       ``this''
have to be handled specially, as do quotes within
quotes:
       ``\,`this'            % \, separates the double and single quote.
        is what I just
        wrote, not  `that'\,''.

Dashes come in three sizes: an
       intra-word
dash, a medium dash for number ranges like
       1--2,
and a punctuation
       dash---like
this.

A sentence-ending space should be larger than the
space between words within a sentence.  You
sometimes have to type special commands in
conjunction with punctuation characters to get
this right, as in the following sentence.
       Gnats, gnus, etc.\ all  % `\ ' makes an inter-word space.
       begin with G\@.         % \@ marks end-of-sentence punctuation.
You should check the spaces after periods when
reading your output to make sure you haven't
forgotten any special cases.  Generating an
ellipsis
       \ldots\               % `\ ' is needed after `\ldots' because TeX
                             % ignores spaces after command names like \ldots
                             % made from \ + letters.
                             %
                             % Note how a `%' character causes TeX to ignore
                             % the end of the input line, so these blank lines
                             % do not start a new paragraph.
                             %
with the right spacing around the periods requires
a special command.

\LaTeX\ interprets some common characters as
commands, so you must type special commands to
generate them.  These characters include the
following:
       \$ \& \% \# \{ and \}.

In printing, text is usually emphasized with an
       \emph{italic}
type style.

\begin{em}
   A long segment of text can also be emphasized
   in this way.  Text within such a segment can be
   given \emph{additional} emphasis.
\end{em}

It is sometimes necessary to prevent \LaTeX\ from
breaking a line where it might otherwise do so.
This may be at a space, as between the ``Mr.''\ and
``Jones'' in
       ``Mr.~Jones'',        % ~ produces an unbreakable interword space.
or within a word---especially when the word is a
symbol like
       \mbox{\emph{itemnum}}
that makes little sense when hyphenated across
lines.

Footnotes\footnote{This is an example of a footnote.}
pose no problem.

\LaTeX\ is good at typesetting mathematical formulas
like
       \( x-3y + z = 7 \)
or
       \( a_{1} > x^{2n} + y^{2n} > x' \)
or
       \( \ip{A}{B} = \sum_{i} a_{i} b_{i} \).
The spaces you type in a formula are
ignored.  Remember that a letter like
       $x$                   % $ ... $  and  \( ... \)  are equivalent
is a formula when it denotes a mathematical
symbol, and it should be typed as one.

\section{Displayed Text}

Text is displayed by indenting it from the left
margin.  Quotations are commonly displayed.  There
are short quotations
\begin{quote}
   This is a short quotation.  It consists of a
   single paragraph of text.  See how it is formatted.
\end{quote}
and longer ones.
\begin{quotation}
   This is a longer quotation.  It consists of two
   paragraphs of text, neither of which are
   particularly interesting.

   This is the second paragraph of the quotation.  It
   is just as dull as the first paragraph.
\end{quotation}
Another frequently-displayed structure is a list.
The following is an example of an \emph{itemized}
list.
\begin{itemize}
   \item This is the first item of an itemized list.
         Each item in the list is marked with a ``tick''.
         You don't have to worry about what kind of tick
         mark is used.

   \item This is the second item of the list.  It
         contains another list nested inside it.  The inner
         list is an \emph{enumerated} list.
         \begin{enumerate}
            \item This is the first item of an enumerated
                  list that is nested within the itemized list.

            \item This is the second item of the inner list.
                  \LaTeX\ allows you to nest lists deeper than
                  you really should.
         \end{enumerate}
         This is the rest of the second item of the outer
         list.  It is no more interesting than any other
         part of the item.
   \item This is the third item of the list.
\end{itemize}
You can even display poetry.
\begin{verse}
   There is an environment
    for verse \\             % The \\ command separates lines
   Whose features some poets % within a stanza.
   will curse.

                             % One or more blank lines separate stanzas.

   For instead of making\\
   Them do \emph{all} line breaking, \\
   It allows them to put too many words on a line when they'd rather be
   forced to be terse.
\end{verse}

Mathematical formulas may also be displayed.  A
displayed formula
is
one-line long; multiline
formulas require special formatting instructions.
   \[  \ip{\Gamma}{\psi'} = x'' + y^{2} + z_{i}^{n}\]
Don't start a paragraph with a displayed equation,
nor make one a paragraph by itself.

\end{document}               % End of document.

\section{Bench Section 105}
% synthetic bench section 105
% This is a sample LaTeX input file.  (Version of 12 August 2004.)
%
% A '%' character causes TeX to ignore all remaining text on the line,
% and is used for comments like this one.

\documentclass{article}      % Specifies the document class

                             % The preamble begins here.
\title{An Example Document}  % Declares the document's title.
\author{Leslie Lamport}      % Declares the author's name.
\date{January 21, 1994}      % Deleting this command produces today's date.

\newcommand{\ip}[2]{(#1, #2)}
                             % Defines \ip{arg1}{arg2} to mean
                             % (arg1, arg2).

%\newcommand{\ip}[2]{\langle #1 | #2\rangle}
                             % This is an alternative definition of
                             % \ip that is commented out.

\begin{document}             % End of preamble and beginning of text.

\maketitle                   % Produces the title.

This is an example input file.  Comparing it with
the output it generates can show you how to
produce a simple document of your own.

\section{Ordinary Text}      % Produces section heading.  Lower-level
                             % sections are begun with similar
                             % \subsection and \subsubsection commands.

The ends  of words and sentences are marked
  by   spaces. It  doesn't matter how many
spaces    you type; one is as good as 100.  The
end of   a line counts as a space.

One   or more   blank lines denote the  end
of  a paragraph.

Since any number of consecutive spaces are treated
like a single one, the formatting of the input
file makes no difference to
      \LaTeX,                % The \LaTeX command generates the LaTeX logo.
but it makes a difference to you.  When you use
\LaTeX, making your input file as easy to read
as possible will be a great help as you write
your document and when you change it.  This sample
file shows how you can add comments to your own input
file.

Because printing is different from typewriting,
there are a number of things that you have to do
differently when preparing an input file than if
you were just typing the document directly.
Quotation marks like
       ``this''
have to be handled specially, as do quotes within
quotes:
       ``\,`this'            % \, separates the double and single quote.
        is what I just
        wrote, not  `that'\,''.

Dashes come in three sizes: an
       intra-word
dash, a medium dash for number ranges like
       1--2,
and a punctuation
       dash---like
this.

A sentence-ending space should be larger than the
space between words within a sentence.  You
sometimes have to type special commands in
conjunction with punctuation characters to get
this right, as in the following sentence.
       Gnats, gnus, etc.\ all  % `\ ' makes an inter-word space.
       begin with G\@.         % \@ marks end-of-sentence punctuation.
You should check the spaces after periods when
reading your output to make sure you haven't
forgotten any special cases.  Generating an
ellipsis
       \ldots\               % `\ ' is needed after `\ldots' because TeX
                             % ignores spaces after command names like \ldots
                             % made from \ + letters.
                             %
                             % Note how a `%' character causes TeX to ignore
                             % the end of the input line, so these blank lines
                             % do not start a new paragraph.
                             %
with the right spacing around the periods requires
a special command.

\LaTeX\ interprets some common characters as
commands, so you must type special commands to
generate them.  These characters include the
following:
       \$ \& \% \# \{ and \}.

In printing, text is usually emphasized with an
       \emph{italic}
type style.

\begin{em}
   A long segment of text can also be emphasized
   in this way.  Text within such a segment can be
   given \emph{additional} emphasis.
\end{em}

It is sometimes necessary to prevent \LaTeX\ from
breaking a line where it might otherwise do so.
This may be at a space, as between the ``Mr.''\ and
``Jones'' in
       ``Mr.~Jones'',        % ~ produces an unbreakable interword space.
or within a word---especially when the word is a
symbol like
       \mbox{\emph{itemnum}}
that makes little sense when hyphenated across
lines.

Footnotes\footnote{This is an example of a footnote.}
pose no problem.

\LaTeX\ is good at typesetting mathematical formulas
like
       \( x-3y + z = 7 \)
or
       \( a_{1} > x^{2n} + y^{2n} > x' \)
or
       \( \ip{A}{B} = \sum_{i} a_{i} b_{i} \).
The spaces you type in a formula are
ignored.  Remember that a letter like
       $x$                   % $ ... $  and  \( ... \)  are equivalent
is a formula when it denotes a mathematical
symbol, and it should be typed as one.

\section{Displayed Text}

Text is displayed by indenting it from the left
margin.  Quotations are commonly displayed.  There
are short quotations
\begin{quote}
   This is a short quotation.  It consists of a
   single paragraph of text.  See how it is formatted.
\end{quote}
and longer ones.
\begin{quotation}
   This is a longer quotation.  It consists of two
   paragraphs of text, neither of which are
   particularly interesting.

   This is the second paragraph of the quotation.  It
   is just as dull as the first paragraph.
\end{quotation}
Another frequently-displayed structure is a list.
The following is an example of an \emph{itemized}
list.
\begin{itemize}
   \item This is the first item of an itemized list.
         Each item in the list is marked with a ``tick''.
         You don't have to worry about what kind of tick
         mark is used.

   \item This is the second item of the list.  It
         contains another list nested inside it.  The inner
         list is an \emph{enumerated} list.
         \begin{enumerate}
            \item This is the first item of an enumerated
                  list that is nested within the itemized list.

            \item This is the second item of the inner list.
                  \LaTeX\ allows you to nest lists deeper than
                  you really should.
         \end{enumerate}
         This is the rest of the second item of the outer
         list.  It is no more interesting than any other
         part of the item.
   \item This is the third item of the list.
\end{itemize}
You can even display poetry.
\begin{verse}
   There is an environment
    for verse \\             % The \\ command separates lines
   Whose features some poets % within a stanza.
   will curse.

                             % One or more blank lines separate stanzas.

   For instead of making\\
   Them do \emph{all} line breaking, \\
   It allows them to put too many words on a line when they'd rather be
   forced to be terse.
\end{verse}

Mathematical formulas may also be displayed.  A
displayed formula
is
one-line long; multiline
formulas require special formatting instructions.
   \[  \ip{\Gamma}{\psi'} = x'' + y^{2} + z_{i}^{n}\]
Don't start a paragraph with a displayed equation,
nor make one a paragraph by itself.

\end{document}               % End of document.

\section{Bench Section 106}
% synthetic bench section 106
% This is a sample LaTeX input file.  (Version of 12 August 2004.)
%
% A '%' character causes TeX to ignore all remaining text on the line,
% and is used for comments like this one.

\documentclass{article}      % Specifies the document class

                             % The preamble begins here.
\title{An Example Document}  % Declares the document's title.
\author{Leslie Lamport}      % Declares the author's name.
\date{January 21, 1994}      % Deleting this command produces today's date.

\newcommand{\ip}[2]{(#1, #2)}
                             % Defines \ip{arg1}{arg2} to mean
                             % (arg1, arg2).

%\newcommand{\ip}[2]{\langle #1 | #2\rangle}
                             % This is an alternative definition of
                             % \ip that is commented out.

\begin{document}             % End of preamble and beginning of text.

\maketitle                   % Produces the title.

This is an example input file.  Comparing it with
the output it generates can show you how to
produce a simple document of your own.

\section{Ordinary Text}      % Produces section heading.  Lower-level
                             % sections are begun with similar
                             % \subsection and \subsubsection commands.

The ends  of words and sentences are marked
  by   spaces. It  doesn't matter how many
spaces    you type; one is as good as 100.  The
end of   a line counts as a space.

One   or more   blank lines denote the  end
of  a paragraph.

Since any number of consecutive spaces are treated
like a single one, the formatting of the input
file makes no difference to
      \LaTeX,                % The \LaTeX command generates the LaTeX logo.
but it makes a difference to you.  When you use
\LaTeX, making your input file as easy to read
as possible will be a great help as you write
your document and when you change it.  This sample
file shows how you can add comments to your own input
file.

Because printing is different from typewriting,
there are a number of things that you have to do
differently when preparing an input file than if
you were just typing the document directly.
Quotation marks like
       ``this''
have to be handled specially, as do quotes within
quotes:
       ``\,`this'            % \, separates the double and single quote.
        is what I just
        wrote, not  `that'\,''.

Dashes come in three sizes: an
       intra-word
dash, a medium dash for number ranges like
       1--2,
and a punctuation
       dash---like
this.

A sentence-ending space should be larger than the
space between words within a sentence.  You
sometimes have to type special commands in
conjunction with punctuation characters to get
this right, as in the following sentence.
       Gnats, gnus, etc.\ all  % `\ ' makes an inter-word space.
       begin with G\@.         % \@ marks end-of-sentence punctuation.
You should check the spaces after periods when
reading your output to make sure you haven't
forgotten any special cases.  Generating an
ellipsis
       \ldots\               % `\ ' is needed after `\ldots' because TeX
                             % ignores spaces after command names like \ldots
                             % made from \ + letters.
                             %
                             % Note how a `%' character causes TeX to ignore
                             % the end of the input line, so these blank lines
                             % do not start a new paragraph.
                             %
with the right spacing around the periods requires
a special command.

\LaTeX\ interprets some common characters as
commands, so you must type special commands to
generate them.  These characters include the
following:
       \$ \& \% \# \{ and \}.

In printing, text is usually emphasized with an
       \emph{italic}
type style.

\begin{em}
   A long segment of text can also be emphasized
   in this way.  Text within such a segment can be
   given \emph{additional} emphasis.
\end{em}

It is sometimes necessary to prevent \LaTeX\ from
breaking a line where it might otherwise do so.
This may be at a space, as between the ``Mr.''\ and
``Jones'' in
       ``Mr.~Jones'',        % ~ produces an unbreakable interword space.
or within a word---especially when the word is a
symbol like
       \mbox{\emph{itemnum}}
that makes little sense when hyphenated across
lines.

Footnotes\footnote{This is an example of a footnote.}
pose no problem.

\LaTeX\ is good at typesetting mathematical formulas
like
       \( x-3y + z = 7 \)
or
       \( a_{1} > x^{2n} + y^{2n} > x' \)
or
       \( \ip{A}{B} = \sum_{i} a_{i} b_{i} \).
The spaces you type in a formula are
ignored.  Remember that a letter like
       $x$                   % $ ... $  and  \( ... \)  are equivalent
is a formula when it denotes a mathematical
symbol, and it should be typed as one.

\section{Displayed Text}

Text is displayed by indenting it from the left
margin.  Quotations are commonly displayed.  There
are short quotations
\begin{quote}
   This is a short quotation.  It consists of a
   single paragraph of text.  See how it is formatted.
\end{quote}
and longer ones.
\begin{quotation}
   This is a longer quotation.  It consists of two
   paragraphs of text, neither of which are
   particularly interesting.

   This is the second paragraph of the quotation.  It
   is just as dull as the first paragraph.
\end{quotation}
Another frequently-displayed structure is a list.
The following is an example of an \emph{itemized}
list.
\begin{itemize}
   \item This is the first item of an itemized list.
         Each item in the list is marked with a ``tick''.
         You don't have to worry about what kind of tick
         mark is used.

   \item This is the second item of the list.  It
         contains another list nested inside it.  The inner
         list is an \emph{enumerated} list.
         \begin{enumerate}
            \item This is the first item of an enumerated
                  list that is nested within the itemized list.

            \item This is the second item of the inner list.
                  \LaTeX\ allows you to nest lists deeper than
                  you really should.
         \end{enumerate}
         This is the rest of the second item of the outer
         list.  It is no more interesting than any other
         part of the item.
   \item This is the third item of the list.
\end{itemize}
You can even display poetry.
\begin{verse}
   There is an environment
    for verse \\             % The \\ command separates lines
   Whose features some poets % within a stanza.
   will curse.

                             % One or more blank lines separate stanzas.

   For instead of making\\
   Them do \emph{all} line breaking, \\
   It allows them to put too many words on a line when they'd rather be
   forced to be terse.
\end{verse}

Mathematical formulas may also be displayed.  A
displayed formula
is
one-line long; multiline
formulas require special formatting instructions.
   \[  \ip{\Gamma}{\psi'} = x'' + y^{2} + z_{i}^{n}\]
Don't start a paragraph with a displayed equation,
nor make one a paragraph by itself.

\end{document}               % End of document.

\section{Bench Section 107}
% synthetic bench section 107
% This is a sample LaTeX input file.  (Version of 12 August 2004.)
%
% A '%' character causes TeX to ignore all remaining text on the line,
% and is used for comments like this one.

\documentclass{article}      % Specifies the document class

                             % The preamble begins here.
\title{An Example Document}  % Declares the document's title.
\author{Leslie Lamport}      % Declares the author's name.
\date{January 21, 1994}      % Deleting this command produces today's date.

\newcommand{\ip}[2]{(#1, #2)}
                             % Defines \ip{arg1}{arg2} to mean
                             % (arg1, arg2).

%\newcommand{\ip}[2]{\langle #1 | #2\rangle}
                             % This is an alternative definition of
                             % \ip that is commented out.

\begin{document}             % End of preamble and beginning of text.

\maketitle                   % Produces the title.

This is an example input file.  Comparing it with
the output it generates can show you how to
produce a simple document of your own.

\section{Ordinary Text}      % Produces section heading.  Lower-level
                             % sections are begun with similar
                             % \subsection and \subsubsection commands.

The ends  of words and sentences are marked
  by   spaces. It  doesn't matter how many
spaces    you type; one is as good as 100.  The
end of   a line counts as a space.

One   or more   blank lines denote the  end
of  a paragraph.

Since any number of consecutive spaces are treated
like a single one, the formatting of the input
file makes no difference to
      \LaTeX,                % The \LaTeX command generates the LaTeX logo.
but it makes a difference to you.  When you use
\LaTeX, making your input file as easy to read
as possible will be a great help as you write
your document and when you change it.  This sample
file shows how you can add comments to your own input
file.

Because printing is different from typewriting,
there are a number of things that you have to do
differently when preparing an input file than if
you were just typing the document directly.
Quotation marks like
       ``this''
have to be handled specially, as do quotes within
quotes:
       ``\,`this'            % \, separates the double and single quote.
        is what I just
        wrote, not  `that'\,''.

Dashes come in three sizes: an
       intra-word
dash, a medium dash for number ranges like
       1--2,
and a punctuation
       dash---like
this.

A sentence-ending space should be larger than the
space between words within a sentence.  You
sometimes have to type special commands in
conjunction with punctuation characters to get
this right, as in the following sentence.
       Gnats, gnus, etc.\ all  % `\ ' makes an inter-word space.
       begin with G\@.         % \@ marks end-of-sentence punctuation.
You should check the spaces after periods when
reading your output to make sure you haven't
forgotten any special cases.  Generating an
ellipsis
       \ldots\               % `\ ' is needed after `\ldots' because TeX
                             % ignores spaces after command names like \ldots
                             % made from \ + letters.
                             %
                             % Note how a `%' character causes TeX to ignore
                             % the end of the input line, so these blank lines
                             % do not start a new paragraph.
                             %
with the right spacing around the periods requires
a special command.

\LaTeX\ interprets some common characters as
commands, so you must type special commands to
generate them.  These characters include the
following:
       \$ \& \% \# \{ and \}.

In printing, text is usually emphasized with an
       \emph{italic}
type style.

\begin{em}
   A long segment of text can also be emphasized
   in this way.  Text within such a segment can be
   given \emph{additional} emphasis.
\end{em}

It is sometimes necessary to prevent \LaTeX\ from
breaking a line where it might otherwise do so.
This may be at a space, as between the ``Mr.''\ and
``Jones'' in
       ``Mr.~Jones'',        % ~ produces an unbreakable interword space.
or within a word---especially when the word is a
symbol like
       \mbox{\emph{itemnum}}
that makes little sense when hyphenated across
lines.

Footnotes\footnote{This is an example of a footnote.}
pose no problem.

\LaTeX\ is good at typesetting mathematical formulas
like
       \( x-3y + z = 7 \)
or
       \( a_{1} > x^{2n} + y^{2n} > x' \)
or
       \( \ip{A}{B} = \sum_{i} a_{i} b_{i} \).
The spaces you type in a formula are
ignored.  Remember that a letter like
       $x$                   % $ ... $  and  \( ... \)  are equivalent
is a formula when it denotes a mathematical
symbol, and it should be typed as one.

\section{Displayed Text}

Text is displayed by indenting it from the left
margin.  Quotations are commonly displayed.  There
are short quotations
\begin{quote}
   This is a short quotation.  It consists of a
   single paragraph of text.  See how it is formatted.
\end{quote}
and longer ones.
\begin{quotation}
   This is a longer quotation.  It consists of two
   paragraphs of text, neither of which are
   particularly interesting.

   This is the second paragraph of the quotation.  It
   is just as dull as the first paragraph.
\end{quotation}
Another frequently-displayed structure is a list.
The following is an example of an \emph{itemized}
list.
\begin{itemize}
   \item This is the first item of an itemized list.
         Each item in the list is marked with a ``tick''.
         You don't have to worry about what kind of tick
         mark is used.

   \item This is the second item of the list.  It
         contains another list nested inside it.  The inner
         list is an \emph{enumerated} list.
         \begin{enumerate}
            \item This is the first item of an enumerated
                  list that is nested within the itemized list.

            \item This is the second item of the inner list.
                  \LaTeX\ allows you to nest lists deeper than
                  you really should.
         \end{enumerate}
         This is the rest of the second item of the outer
         list.  It is no more interesting than any other
         part of the item.
   \item This is the third item of the list.
\end{itemize}
You can even display poetry.
\begin{verse}
   There is an environment
    for verse \\             % The \\ command separates lines
   Whose features some poets % within a stanza.
   will curse.

                             % One or more blank lines separate stanzas.

   For instead of making\\
   Them do \emph{all} line breaking, \\
   It allows them to put too many words on a line when they'd rather be
   forced to be terse.
\end{verse}

Mathematical formulas may also be displayed.  A
displayed formula
is
one-line long; multiline
formulas require special formatting instructions.
   \[  \ip{\Gamma}{\psi'} = x'' + y^{2} + z_{i}^{n}\]
Don't start a paragraph with a displayed equation,
nor make one a paragraph by itself.

\end{document}               % End of document.

\section{Bench Section 108}
% synthetic bench section 108
% This is a sample LaTeX input file.  (Version of 12 August 2004.)
%
% A '%' character causes TeX to ignore all remaining text on the line,
% and is used for comments like this one.

\documentclass{article}      % Specifies the document class

                             % The preamble begins here.
\title{An Example Document}  % Declares the document's title.
\author{Leslie Lamport}      % Declares the author's name.
\date{January 21, 1994}      % Deleting this command produces today's date.

\newcommand{\ip}[2]{(#1, #2)}
                             % Defines \ip{arg1}{arg2} to mean
                             % (arg1, arg2).

%\newcommand{\ip}[2]{\langle #1 | #2\rangle}
                             % This is an alternative definition of
                             % \ip that is commented out.

\begin{document}             % End of preamble and beginning of text.

\maketitle                   % Produces the title.

This is an example input file.  Comparing it with
the output it generates can show you how to
produce a simple document of your own.

\section{Ordinary Text}      % Produces section heading.  Lower-level
                             % sections are begun with similar
                             % \subsection and \subsubsection commands.

The ends  of words and sentences are marked
  by   spaces. It  doesn't matter how many
spaces    you type; one is as good as 100.  The
end of   a line counts as a space.

One   or more   blank lines denote the  end
of  a paragraph.

Since any number of consecutive spaces are treated
like a single one, the formatting of the input
file makes no difference to
      \LaTeX,                % The \LaTeX command generates the LaTeX logo.
but it makes a difference to you.  When you use
\LaTeX, making your input file as easy to read
as possible will be a great help as you write
your document and when you change it.  This sample
file shows how you can add comments to your own input
file.

Because printing is different from typewriting,
there are a number of things that you have to do
differently when preparing an input file than if
you were just typing the document directly.
Quotation marks like
       ``this''
have to be handled specially, as do quotes within
quotes:
       ``\,`this'            % \, separates the double and single quote.
        is what I just
        wrote, not  `that'\,''.

Dashes come in three sizes: an
       intra-word
dash, a medium dash for number ranges like
       1--2,
and a punctuation
       dash---like
this.

A sentence-ending space should be larger than the
space between words within a sentence.  You
sometimes have to type special commands in
conjunction with punctuation characters to get
this right, as in the following sentence.
       Gnats, gnus, etc.\ all  % `\ ' makes an inter-word space.
       begin with G\@.         % \@ marks end-of-sentence punctuation.
You should check the spaces after periods when
reading your output to make sure you haven't
forgotten any special cases.  Generating an
ellipsis
       \ldots\               % `\ ' is needed after `\ldots' because TeX
                             % ignores spaces after command names like \ldots
                             % made from \ + letters.
                             %
                             % Note how a `%' character causes TeX to ignore
                             % the end of the input line, so these blank lines
                             % do not start a new paragraph.
                             %
with the right spacing around the periods requires
a special command.

\LaTeX\ interprets some common characters as
commands, so you must type special commands to
generate them.  These characters include the
following:
       \$ \& \% \# \{ and \}.

In printing, text is usually emphasized with an
       \emph{italic}
type style.

\begin{em}
   A long segment of text can also be emphasized
   in this way.  Text within such a segment can be
   given \emph{additional} emphasis.
\end{em}

It is sometimes necessary to prevent \LaTeX\ from
breaking a line where it might otherwise do so.
This may be at a space, as between the ``Mr.''\ and
``Jones'' in
       ``Mr.~Jones'',        % ~ produces an unbreakable interword space.
or within a word---especially when the word is a
symbol like
       \mbox{\emph{itemnum}}
that makes little sense when hyphenated across
lines.

Footnotes\footnote{This is an example of a footnote.}
pose no problem.

\LaTeX\ is good at typesetting mathematical formulas
like
       \( x-3y + z = 7 \)
or
       \( a_{1} > x^{2n} + y^{2n} > x' \)
or
       \( \ip{A}{B} = \sum_{i} a_{i} b_{i} \).
The spaces you type in a formula are
ignored.  Remember that a letter like
       $x$                   % $ ... $  and  \( ... \)  are equivalent
is a formula when it denotes a mathematical
symbol, and it should be typed as one.

\section{Displayed Text}

Text is displayed by indenting it from the left
margin.  Quotations are commonly displayed.  There
are short quotations
\begin{quote}
   This is a short quotation.  It consists of a
   single paragraph of text.  See how it is formatted.
\end{quote}
and longer ones.
\begin{quotation}
   This is a longer quotation.  It consists of two
   paragraphs of text, neither of which are
   particularly interesting.

   This is the second paragraph of the quotation.  It
   is just as dull as the first paragraph.
\end{quotation}
Another frequently-displayed structure is a list.
The following is an example of an \emph{itemized}
list.
\begin{itemize}
   \item This is the first item of an itemized list.
         Each item in the list is marked with a ``tick''.
         You don't have to worry about what kind of tick
         mark is used.

   \item This is the second item of the list.  It
         contains another list nested inside it.  The inner
         list is an \emph{enumerated} list.
         \begin{enumerate}
            \item This is the first item of an enumerated
                  list that is nested within the itemized list.

            \item This is the second item of the inner list.
                  \LaTeX\ allows you to nest lists deeper than
                  you really should.
         \end{enumerate}
         This is the rest of the second item of the outer
         list.  It is no more interesting than any other
         part of the item.
   \item This is the third item of the list.
\end{itemize}
You can even display poetry.
\begin{verse}
   There is an environment
    for verse \\             % The \\ command separates lines
   Whose features some poets % within a stanza.
   will curse.

                             % One or more blank lines separate stanzas.

   For instead of making\\
   Them do \emph{all} line breaking, \\
   It allows them to put too many words on a line when they'd rather be
   forced to be terse.
\end{verse}

Mathematical formulas may also be displayed.  A
displayed formula
is
one-line long; multiline
formulas require special formatting instructions.
   \[  \ip{\Gamma}{\psi'} = x'' + y^{2} + z_{i}^{n}\]
Don't start a paragraph with a displayed equation,
nor make one a paragraph by itself.

\end{document}               % End of document.

\section{Bench Section 109}
% synthetic bench section 109
% This is a sample LaTeX input file.  (Version of 12 August 2004.)
%
% A '%' character causes TeX to ignore all remaining text on the line,
% and is used for comments like this one.

\documentclass{article}      % Specifies the document class

                             % The preamble begins here.
\title{An Example Document}  % Declares the document's title.
\author{Leslie Lamport}      % Declares the author's name.
\date{January 21, 1994}      % Deleting this command produces today's date.

\newcommand{\ip}[2]{(#1, #2)}
                             % Defines \ip{arg1}{arg2} to mean
                             % (arg1, arg2).

%\newcommand{\ip}[2]{\langle #1 | #2\rangle}
                             % This is an alternative definition of
                             % \ip that is commented out.

\begin{document}             % End of preamble and beginning of text.

\maketitle                   % Produces the title.

This is an example input file.  Comparing it with
the output it generates can show you how to
produce a simple document of your own.

\section{Ordinary Text}      % Produces section heading.  Lower-level
                             % sections are begun with similar
                             % \subsection and \subsubsection commands.

The ends  of words and sentences are marked
  by   spaces. It  doesn't matter how many
spaces    you type; one is as good as 100.  The
end of   a line counts as a space.

One   or more   blank lines denote the  end
of  a paragraph.

Since any number of consecutive spaces are treated
like a single one, the formatting of the input
file makes no difference to
      \LaTeX,                % The \LaTeX command generates the LaTeX logo.
but it makes a difference to you.  When you use
\LaTeX, making your input file as easy to read
as possible will be a great help as you write
your document and when you change it.  This sample
file shows how you can add comments to your own input
file.

Because printing is different from typewriting,
there are a number of things that you have to do
differently when preparing an input file than if
you were just typing the document directly.
Quotation marks like
       ``this''
have to be handled specially, as do quotes within
quotes:
       ``\,`this'            % \, separates the double and single quote.
        is what I just
        wrote, not  `that'\,''.

Dashes come in three sizes: an
       intra-word
dash, a medium dash for number ranges like
       1--2,
and a punctuation
       dash---like
this.

A sentence-ending space should be larger than the
space between words within a sentence.  You
sometimes have to type special commands in
conjunction with punctuation characters to get
this right, as in the following sentence.
       Gnats, gnus, etc.\ all  % `\ ' makes an inter-word space.
       begin with G\@.         % \@ marks end-of-sentence punctuation.
You should check the spaces after periods when
reading your output to make sure you haven't
forgotten any special cases.  Generating an
ellipsis
       \ldots\               % `\ ' is needed after `\ldots' because TeX
                             % ignores spaces after command names like \ldots
                             % made from \ + letters.
                             %
                             % Note how a `%' character causes TeX to ignore
                             % the end of the input line, so these blank lines
                             % do not start a new paragraph.
                             %
with the right spacing around the periods requires
a special command.

\LaTeX\ interprets some common characters as
commands, so you must type special commands to
generate them.  These characters include the
following:
       \$ \& \% \# \{ and \}.

In printing, text is usually emphasized with an
       \emph{italic}
type style.

\begin{em}
   A long segment of text can also be emphasized
   in this way.  Text within such a segment can be
   given \emph{additional} emphasis.
\end{em}

It is sometimes necessary to prevent \LaTeX\ from
breaking a line where it might otherwise do so.
This may be at a space, as between the ``Mr.''\ and
``Jones'' in
       ``Mr.~Jones'',        % ~ produces an unbreakable interword space.
or within a word---especially when the word is a
symbol like
       \mbox{\emph{itemnum}}
that makes little sense when hyphenated across
lines.

Footnotes\footnote{This is an example of a footnote.}
pose no problem.

\LaTeX\ is good at typesetting mathematical formulas
like
       \( x-3y + z = 7 \)
or
       \( a_{1} > x^{2n} + y^{2n} > x' \)
or
       \( \ip{A}{B} = \sum_{i} a_{i} b_{i} \).
The spaces you type in a formula are
ignored.  Remember that a letter like
       $x$                   % $ ... $  and  \( ... \)  are equivalent
is a formula when it denotes a mathematical
symbol, and it should be typed as one.

\section{Displayed Text}

Text is displayed by indenting it from the left
margin.  Quotations are commonly displayed.  There
are short quotations
\begin{quote}
   This is a short quotation.  It consists of a
   single paragraph of text.  See how it is formatted.
\end{quote}
and longer ones.
\begin{quotation}
   This is a longer quotation.  It consists of two
   paragraphs of text, neither of which are
   particularly interesting.

   This is the second paragraph of the quotation.  It
   is just as dull as the first paragraph.
\end{quotation}
Another frequently-displayed structure is a list.
The following is an example of an \emph{itemized}
list.
\begin{itemize}
   \item This is the first item of an itemized list.
         Each item in the list is marked with a ``tick''.
         You don't have to worry about what kind of tick
         mark is used.

   \item This is the second item of the list.  It
         contains another list nested inside it.  The inner
         list is an \emph{enumerated} list.
         \begin{enumerate}
            \item This is the first item of an enumerated
                  list that is nested within the itemized list.

            \item This is the second item of the inner list.
                  \LaTeX\ allows you to nest lists deeper than
                  you really should.
         \end{enumerate}
         This is the rest of the second item of the outer
         list.  It is no more interesting than any other
         part of the item.
   \item This is the third item of the list.
\end{itemize}
You can even display poetry.
\begin{verse}
   There is an environment
    for verse \\             % The \\ command separates lines
   Whose features some poets % within a stanza.
   will curse.

                             % One or more blank lines separate stanzas.

   For instead of making\\
   Them do \emph{all} line breaking, \\
   It allows them to put too many words on a line when they'd rather be
   forced to be terse.
\end{verse}

Mathematical formulas may also be displayed.  A
displayed formula
is
one-line long; multiline
formulas require special formatting instructions.
   \[  \ip{\Gamma}{\psi'} = x'' + y^{2} + z_{i}^{n}\]
Don't start a paragraph with a displayed equation,
nor make one a paragraph by itself.

\end{document}               % End of document.

\section{Bench Section 110}
% synthetic bench section 110
% This is a sample LaTeX input file.  (Version of 12 August 2004.)
%
% A '%' character causes TeX to ignore all remaining text on the line,
% and is used for comments like this one.

\documentclass{article}      % Specifies the document class

                             % The preamble begins here.
\title{An Example Document}  % Declares the document's title.
\author{Leslie Lamport}      % Declares the author's name.
\date{January 21, 1994}      % Deleting this command produces today's date.

\newcommand{\ip}[2]{(#1, #2)}
                             % Defines \ip{arg1}{arg2} to mean
                             % (arg1, arg2).

%\newcommand{\ip}[2]{\langle #1 | #2\rangle}
                             % This is an alternative definition of
                             % \ip that is commented out.

\begin{document}             % End of preamble and beginning of text.

\maketitle                   % Produces the title.

This is an example input file.  Comparing it with
the output it generates can show you how to
produce a simple document of your own.

\section{Ordinary Text}      % Produces section heading.  Lower-level
                             % sections are begun with similar
                             % \subsection and \subsubsection commands.

The ends  of words and sentences are marked
  by   spaces. It  doesn't matter how many
spaces    you type; one is as good as 100.  The
end of   a line counts as a space.

One   or more   blank lines denote the  end
of  a paragraph.

Since any number of consecutive spaces are treated
like a single one, the formatting of the input
file makes no difference to
      \LaTeX,                % The \LaTeX command generates the LaTeX logo.
but it makes a difference to you.  When you use
\LaTeX, making your input file as easy to read
as possible will be a great help as you write
your document and when you change it.  This sample
file shows how you can add comments to your own input
file.

Because printing is different from typewriting,
there are a number of things that you have to do
differently when preparing an input file than if
you were just typing the document directly.
Quotation marks like
       ``this''
have to be handled specially, as do quotes within
quotes:
       ``\,`this'            % \, separates the double and single quote.
        is what I just
        wrote, not  `that'\,''.

Dashes come in three sizes: an
       intra-word
dash, a medium dash for number ranges like
       1--2,
and a punctuation
       dash---like
this.

A sentence-ending space should be larger than the
space between words within a sentence.  You
sometimes have to type special commands in
conjunction with punctuation characters to get
this right, as in the following sentence.
       Gnats, gnus, etc.\ all  % `\ ' makes an inter-word space.
       begin with G\@.         % \@ marks end-of-sentence punctuation.
You should check the spaces after periods when
reading your output to make sure you haven't
forgotten any special cases.  Generating an
ellipsis
       \ldots\               % `\ ' is needed after `\ldots' because TeX
                             % ignores spaces after command names like \ldots
                             % made from \ + letters.
                             %
                             % Note how a `%' character causes TeX to ignore
                             % the end of the input line, so these blank lines
                             % do not start a new paragraph.
                             %
with the right spacing around the periods requires
a special command.

\LaTeX\ interprets some common characters as
commands, so you must type special commands to
generate them.  These characters include the
following:
       \$ \& \% \# \{ and \}.

In printing, text is usually emphasized with an
       \emph{italic}
type style.

\begin{em}
   A long segment of text can also be emphasized
   in this way.  Text within such a segment can be
   given \emph{additional} emphasis.
\end{em}

It is sometimes necessary to prevent \LaTeX\ from
breaking a line where it might otherwise do so.
This may be at a space, as between the ``Mr.''\ and
``Jones'' in
       ``Mr.~Jones'',        % ~ produces an unbreakable interword space.
or within a word---especially when the word is a
symbol like
       \mbox{\emph{itemnum}}
that makes little sense when hyphenated across
lines.

Footnotes\footnote{This is an example of a footnote.}
pose no problem.

\LaTeX\ is good at typesetting mathematical formulas
like
       \( x-3y + z = 7 \)
or
       \( a_{1} > x^{2n} + y^{2n} > x' \)
or
       \( \ip{A}{B} = \sum_{i} a_{i} b_{i} \).
The spaces you type in a formula are
ignored.  Remember that a letter like
       $x$                   % $ ... $  and  \( ... \)  are equivalent
is a formula when it denotes a mathematical
symbol, and it should be typed as one.

\section{Displayed Text}

Text is displayed by indenting it from the left
margin.  Quotations are commonly displayed.  There
are short quotations
\begin{quote}
   This is a short quotation.  It consists of a
   single paragraph of text.  See how it is formatted.
\end{quote}
and longer ones.
\begin{quotation}
   This is a longer quotation.  It consists of two
   paragraphs of text, neither of which are
   particularly interesting.

   This is the second paragraph of the quotation.  It
   is just as dull as the first paragraph.
\end{quotation}
Another frequently-displayed structure is a list.
The following is an example of an \emph{itemized}
list.
\begin{itemize}
   \item This is the first item of an itemized list.
         Each item in the list is marked with a ``tick''.
         You don't have to worry about what kind of tick
         mark is used.

   \item This is the second item of the list.  It
         contains another list nested inside it.  The inner
         list is an \emph{enumerated} list.
         \begin{enumerate}
            \item This is the first item of an enumerated
                  list that is nested within the itemized list.

            \item This is the second item of the inner list.
                  \LaTeX\ allows you to nest lists deeper than
                  you really should.
         \end{enumerate}
         This is the rest of the second item of the outer
         list.  It is no more interesting than any other
         part of the item.
   \item This is the third item of the list.
\end{itemize}
You can even display poetry.
\begin{verse}
   There is an environment
    for verse \\             % The \\ command separates lines
   Whose features some poets % within a stanza.
   will curse.

                             % One or more blank lines separate stanzas.

   For instead of making\\
   Them do \emph{all} line breaking, \\
   It allows them to put too many words on a line when they'd rather be
   forced to be terse.
\end{verse}

Mathematical formulas may also be displayed.  A
displayed formula
is
one-line long; multiline
formulas require special formatting instructions.
   \[  \ip{\Gamma}{\psi'} = x'' + y^{2} + z_{i}^{n}\]
Don't start a paragraph with a displayed equation,
nor make one a paragraph by itself.

\end{document}               % End of document.

\section{Bench Section 111}
% synthetic bench section 111
% This is a sample LaTeX input file.  (Version of 12 August 2004.)
%
% A '%' character causes TeX to ignore all remaining text on the line,
% and is used for comments like this one.

\documentclass{article}      % Specifies the document class

                             % The preamble begins here.
\title{An Example Document}  % Declares the document's title.
\author{Leslie Lamport}      % Declares the author's name.
\date{January 21, 1994}      % Deleting this command produces today's date.

\newcommand{\ip}[2]{(#1, #2)}
                             % Defines \ip{arg1}{arg2} to mean
                             % (arg1, arg2).

%\newcommand{\ip}[2]{\langle #1 | #2\rangle}
                             % This is an alternative definition of
                             % \ip that is commented out.

\begin{document}             % End of preamble and beginning of text.

\maketitle                   % Produces the title.

This is an example input file.  Comparing it with
the output it generates can show you how to
produce a simple document of your own.

\section{Ordinary Text}      % Produces section heading.  Lower-level
                             % sections are begun with similar
                             % \subsection and \subsubsection commands.

The ends  of words and sentences are marked
  by   spaces. It  doesn't matter how many
spaces    you type; one is as good as 100.  The
end of   a line counts as a space.

One   or more   blank lines denote the  end
of  a paragraph.

Since any number of consecutive spaces are treated
like a single one, the formatting of the input
file makes no difference to
      \LaTeX,                % The \LaTeX command generates the LaTeX logo.
but it makes a difference to you.  When you use
\LaTeX, making your input file as easy to read
as possible will be a great help as you write
your document and when you change it.  This sample
file shows how you can add comments to your own input
file.

Because printing is different from typewriting,
there are a number of things that you have to do
differently when preparing an input file than if
you were just typing the document directly.
Quotation marks like
       ``this''
have to be handled specially, as do quotes within
quotes:
       ``\,`this'            % \, separates the double and single quote.
        is what I just
        wrote, not  `that'\,''.

Dashes come in three sizes: an
       intra-word
dash, a medium dash for number ranges like
       1--2,
and a punctuation
       dash---like
this.

A sentence-ending space should be larger than the
space between words within a sentence.  You
sometimes have to type special commands in
conjunction with punctuation characters to get
this right, as in the following sentence.
       Gnats, gnus, etc.\ all  % `\ ' makes an inter-word space.
       begin with G\@.         % \@ marks end-of-sentence punctuation.
You should check the spaces after periods when
reading your output to make sure you haven't
forgotten any special cases.  Generating an
ellipsis
       \ldots\               % `\ ' is needed after `\ldots' because TeX
                             % ignores spaces after command names like \ldots
                             % made from \ + letters.
                             %
                             % Note how a `%' character causes TeX to ignore
                             % the end of the input line, so these blank lines
                             % do not start a new paragraph.
                             %
with the right spacing around the periods requires
a special command.

\LaTeX\ interprets some common characters as
commands, so you must type special commands to
generate them.  These characters include the
following:
       \$ \& \% \# \{ and \}.

In printing, text is usually emphasized with an
       \emph{italic}
type style.

\begin{em}
   A long segment of text can also be emphasized
   in this way.  Text within such a segment can be
   given \emph{additional} emphasis.
\end{em}

It is sometimes necessary to prevent \LaTeX\ from
breaking a line where it might otherwise do so.
This may be at a space, as between the ``Mr.''\ and
``Jones'' in
       ``Mr.~Jones'',        % ~ produces an unbreakable interword space.
or within a word---especially when the word is a
symbol like
       \mbox{\emph{itemnum}}
that makes little sense when hyphenated across
lines.

Footnotes\footnote{This is an example of a footnote.}
pose no problem.

\LaTeX\ is good at typesetting mathematical formulas
like
       \( x-3y + z = 7 \)
or
       \( a_{1} > x^{2n} + y^{2n} > x' \)
or
       \( \ip{A}{B} = \sum_{i} a_{i} b_{i} \).
The spaces you type in a formula are
ignored.  Remember that a letter like
       $x$                   % $ ... $  and  \( ... \)  are equivalent
is a formula when it denotes a mathematical
symbol, and it should be typed as one.

\section{Displayed Text}

Text is displayed by indenting it from the left
margin.  Quotations are commonly displayed.  There
are short quotations
\begin{quote}
   This is a short quotation.  It consists of a
   single paragraph of text.  See how it is formatted.
\end{quote}
and longer ones.
\begin{quotation}
   This is a longer quotation.  It consists of two
   paragraphs of text, neither of which are
   particularly interesting.

   This is the second paragraph of the quotation.  It
   is just as dull as the first paragraph.
\end{quotation}
Another frequently-displayed structure is a list.
The following is an example of an \emph{itemized}
list.
\begin{itemize}
   \item This is the first item of an itemized list.
         Each item in the list is marked with a ``tick''.
         You don't have to worry about what kind of tick
         mark is used.

   \item This is the second item of the list.  It
         contains another list nested inside it.  The inner
         list is an \emph{enumerated} list.
         \begin{enumerate}
            \item This is the first item of an enumerated
                  list that is nested within the itemized list.

            \item This is the second item of the inner list.
                  \LaTeX\ allows you to nest lists deeper than
                  you really should.
         \end{enumerate}
         This is the rest of the second item of the outer
         list.  It is no more interesting than any other
         part of the item.
   \item This is the third item of the list.
\end{itemize}
You can even display poetry.
\begin{verse}
   There is an environment
    for verse \\             % The \\ command separates lines
   Whose features some poets % within a stanza.
   will curse.

                             % One or more blank lines separate stanzas.

   For instead of making\\
   Them do \emph{all} line breaking, \\
   It allows them to put too many words on a line when they'd rather be
   forced to be terse.
\end{verse}

Mathematical formulas may also be displayed.  A
displayed formula
is
one-line long; multiline
formulas require special formatting instructions.
   \[  \ip{\Gamma}{\psi'} = x'' + y^{2} + z_{i}^{n}\]
Don't start a paragraph with a displayed equation,
nor make one a paragraph by itself.

\end{document}               % End of document.

\section{Bench Section 112}
% synthetic bench section 112
% This is a sample LaTeX input file.  (Version of 12 August 2004.)
%
% A '%' character causes TeX to ignore all remaining text on the line,
% and is used for comments like this one.

\documentclass{article}      % Specifies the document class

                             % The preamble begins here.
\title{An Example Document}  % Declares the document's title.
\author{Leslie Lamport}      % Declares the author's name.
\date{January 21, 1994}      % Deleting this command produces today's date.

\newcommand{\ip}[2]{(#1, #2)}
                             % Defines \ip{arg1}{arg2} to mean
                             % (arg1, arg2).

%\newcommand{\ip}[2]{\langle #1 | #2\rangle}
                             % This is an alternative definition of
                             % \ip that is commented out.

\begin{document}             % End of preamble and beginning of text.

\maketitle                   % Produces the title.

This is an example input file.  Comparing it with
the output it generates can show you how to
produce a simple document of your own.

\section{Ordinary Text}      % Produces section heading.  Lower-level
                             % sections are begun with similar
                             % \subsection and \subsubsection commands.

The ends  of words and sentences are marked
  by   spaces. It  doesn't matter how many
spaces    you type; one is as good as 100.  The
end of   a line counts as a space.

One   or more   blank lines denote the  end
of  a paragraph.

Since any number of consecutive spaces are treated
like a single one, the formatting of the input
file makes no difference to
      \LaTeX,                % The \LaTeX command generates the LaTeX logo.
but it makes a difference to you.  When you use
\LaTeX, making your input file as easy to read
as possible will be a great help as you write
your document and when you change it.  This sample
file shows how you can add comments to your own input
file.

Because printing is different from typewriting,
there are a number of things that you have to do
differently when preparing an input file than if
you were just typing the document directly.
Quotation marks like
       ``this''
have to be handled specially, as do quotes within
quotes:
       ``\,`this'            % \, separates the double and single quote.
        is what I just
        wrote, not  `that'\,''.

Dashes come in three sizes: an
       intra-word
dash, a medium dash for number ranges like
       1--2,
and a punctuation
       dash---like
this.

A sentence-ending space should be larger than the
space between words within a sentence.  You
sometimes have to type special commands in
conjunction with punctuation characters to get
this right, as in the following sentence.
       Gnats, gnus, etc.\ all  % `\ ' makes an inter-word space.
       begin with G\@.         % \@ marks end-of-sentence punctuation.
You should check the spaces after periods when
reading your output to make sure you haven't
forgotten any special cases.  Generating an
ellipsis
       \ldots\               % `\ ' is needed after `\ldots' because TeX
                             % ignores spaces after command names like \ldots
                             % made from \ + letters.
                             %
                             % Note how a `%' character causes TeX to ignore
                             % the end of the input line, so these blank lines
                             % do not start a new paragraph.
                             %
with the right spacing around the periods requires
a special command.

\LaTeX\ interprets some common characters as
commands, so you must type special commands to
generate them.  These characters include the
following:
       \$ \& \% \# \{ and \}.

In printing, text is usually emphasized with an
       \emph{italic}
type style.

\begin{em}
   A long segment of text can also be emphasized
   in this way.  Text within such a segment can be
   given \emph{additional} emphasis.
\end{em}

It is sometimes necessary to prevent \LaTeX\ from
breaking a line where it might otherwise do so.
This may be at a space, as between the ``Mr.''\ and
``Jones'' in
       ``Mr.~Jones'',        % ~ produces an unbreakable interword space.
or within a word---especially when the word is a
symbol like
       \mbox{\emph{itemnum}}
that makes little sense when hyphenated across
lines.

Footnotes\footnote{This is an example of a footnote.}
pose no problem.

\LaTeX\ is good at typesetting mathematical formulas
like
       \( x-3y + z = 7 \)
or
       \( a_{1} > x^{2n} + y^{2n} > x' \)
or
       \( \ip{A}{B} = \sum_{i} a_{i} b_{i} \).
The spaces you type in a formula are
ignored.  Remember that a letter like
       $x$                   % $ ... $  and  \( ... \)  are equivalent
is a formula when it denotes a mathematical
symbol, and it should be typed as one.

\section{Displayed Text}

Text is displayed by indenting it from the left
margin.  Quotations are commonly displayed.  There
are short quotations
\begin{quote}
   This is a short quotation.  It consists of a
   single paragraph of text.  See how it is formatted.
\end{quote}
and longer ones.
\begin{quotation}
   This is a longer quotation.  It consists of two
   paragraphs of text, neither of which are
   particularly interesting.

   This is the second paragraph of the quotation.  It
   is just as dull as the first paragraph.
\end{quotation}
Another frequently-displayed structure is a list.
The following is an example of an \emph{itemized}
list.
\begin{itemize}
   \item This is the first item of an itemized list.
         Each item in the list is marked with a ``tick''.
         You don't have to worry about what kind of tick
         mark is used.

   \item This is the second item of the list.  It
         contains another list nested inside it.  The inner
         list is an \emph{enumerated} list.
         \begin{enumerate}
            \item This is the first item of an enumerated
                  list that is nested within the itemized list.

            \item This is the second item of the inner list.
                  \LaTeX\ allows you to nest lists deeper than
                  you really should.
         \end{enumerate}
         This is the rest of the second item of the outer
         list.  It is no more interesting than any other
         part of the item.
   \item This is the third item of the list.
\end{itemize}
You can even display poetry.
\begin{verse}
   There is an environment
    for verse \\             % The \\ command separates lines
   Whose features some poets % within a stanza.
   will curse.

                             % One or more blank lines separate stanzas.

   For instead of making\\
   Them do \emph{all} line breaking, \\
   It allows them to put too many words on a line when they'd rather be
   forced to be terse.
\end{verse}

Mathematical formulas may also be displayed.  A
displayed formula
is
one-line long; multiline
formulas require special formatting instructions.
   \[  \ip{\Gamma}{\psi'} = x'' + y^{2} + z_{i}^{n}\]
Don't start a paragraph with a displayed equation,
nor make one a paragraph by itself.

\end{document}               % End of document.

\section{Bench Section 113}
% synthetic bench section 113
% This is a sample LaTeX input file.  (Version of 12 August 2004.)
%
% A '%' character causes TeX to ignore all remaining text on the line,
% and is used for comments like this one.

\documentclass{article}      % Specifies the document class

                             % The preamble begins here.
\title{An Example Document}  % Declares the document's title.
\author{Leslie Lamport}      % Declares the author's name.
\date{January 21, 1994}      % Deleting this command produces today's date.

\newcommand{\ip}[2]{(#1, #2)}
                             % Defines \ip{arg1}{arg2} to mean
                             % (arg1, arg2).

%\newcommand{\ip}[2]{\langle #1 | #2\rangle}
                             % This is an alternative definition of
                             % \ip that is commented out.

\begin{document}             % End of preamble and beginning of text.

\maketitle                   % Produces the title.

This is an example input file.  Comparing it with
the output it generates can show you how to
produce a simple document of your own.

\section{Ordinary Text}      % Produces section heading.  Lower-level
                             % sections are begun with similar
                             % \subsection and \subsubsection commands.

The ends  of words and sentences are marked
  by   spaces. It  doesn't matter how many
spaces    you type; one is as good as 100.  The
end of   a line counts as a space.

One   or more   blank lines denote the  end
of  a paragraph.

Since any number of consecutive spaces are treated
like a single one, the formatting of the input
file makes no difference to
      \LaTeX,                % The \LaTeX command generates the LaTeX logo.
but it makes a difference to you.  When you use
\LaTeX, making your input file as easy to read
as possible will be a great help as you write
your document and when you change it.  This sample
file shows how you can add comments to your own input
file.

Because printing is different from typewriting,
there are a number of things that you have to do
differently when preparing an input file than if
you were just typing the document directly.
Quotation marks like
       ``this''
have to be handled specially, as do quotes within
quotes:
       ``\,`this'            % \, separates the double and single quote.
        is what I just
        wrote, not  `that'\,''.

Dashes come in three sizes: an
       intra-word
dash, a medium dash for number ranges like
       1--2,
and a punctuation
       dash---like
this.

A sentence-ending space should be larger than the
space between words within a sentence.  You
sometimes have to type special commands in
conjunction with punctuation characters to get
this right, as in the following sentence.
       Gnats, gnus, etc.\ all  % `\ ' makes an inter-word space.
       begin with G\@.         % \@ marks end-of-sentence punctuation.
You should check the spaces after periods when
reading your output to make sure you haven't
forgotten any special cases.  Generating an
ellipsis
       \ldots\               % `\ ' is needed after `\ldots' because TeX
                             % ignores spaces after command names like \ldots
                             % made from \ + letters.
                             %
                             % Note how a `%' character causes TeX to ignore
                             % the end of the input line, so these blank lines
                             % do not start a new paragraph.
                             %
with the right spacing around the periods requires
a special command.

\LaTeX\ interprets some common characters as
commands, so you must type special commands to
generate them.  These characters include the
following:
       \$ \& \% \# \{ and \}.

In printing, text is usually emphasized with an
       \emph{italic}
type style.

\begin{em}
   A long segment of text can also be emphasized
   in this way.  Text within such a segment can be
   given \emph{additional} emphasis.
\end{em}

It is sometimes necessary to prevent \LaTeX\ from
breaking a line where it might otherwise do so.
This may be at a space, as between the ``Mr.''\ and
``Jones'' in
       ``Mr.~Jones'',        % ~ produces an unbreakable interword space.
or within a word---especially when the word is a
symbol like
       \mbox{\emph{itemnum}}
that makes little sense when hyphenated across
lines.

Footnotes\footnote{This is an example of a footnote.}
pose no problem.

\LaTeX\ is good at typesetting mathematical formulas
like
       \( x-3y + z = 7 \)
or
       \( a_{1} > x^{2n} + y^{2n} > x' \)
or
       \( \ip{A}{B} = \sum_{i} a_{i} b_{i} \).
The spaces you type in a formula are
ignored.  Remember that a letter like
       $x$                   % $ ... $  and  \( ... \)  are equivalent
is a formula when it denotes a mathematical
symbol, and it should be typed as one.

\section{Displayed Text}

Text is displayed by indenting it from the left
margin.  Quotations are commonly displayed.  There
are short quotations
\begin{quote}
   This is a short quotation.  It consists of a
   single paragraph of text.  See how it is formatted.
\end{quote}
and longer ones.
\begin{quotation}
   This is a longer quotation.  It consists of two
   paragraphs of text, neither of which are
   particularly interesting.

   This is the second paragraph of the quotation.  It
   is just as dull as the first paragraph.
\end{quotation}
Another frequently-displayed structure is a list.
The following is an example of an \emph{itemized}
list.
\begin{itemize}
   \item This is the first item of an itemized list.
         Each item in the list is marked with a ``tick''.
         You don't have to worry about what kind of tick
         mark is used.

   \item This is the second item of the list.  It
         contains another list nested inside it.  The inner
         list is an \emph{enumerated} list.
         \begin{enumerate}
            \item This is the first item of an enumerated
                  list that is nested within the itemized list.

            \item This is the second item of the inner list.
                  \LaTeX\ allows you to nest lists deeper than
                  you really should.
         \end{enumerate}
         This is the rest of the second item of the outer
         list.  It is no more interesting than any other
         part of the item.
   \item This is the third item of the list.
\end{itemize}
You can even display poetry.
\begin{verse}
   There is an environment
    for verse \\             % The \\ command separates lines
   Whose features some poets % within a stanza.
   will curse.

                             % One or more blank lines separate stanzas.

   For instead of making\\
   Them do \emph{all} line breaking, \\
   It allows them to put too many words on a line when they'd rather be
   forced to be terse.
\end{verse}

Mathematical formulas may also be displayed.  A
displayed formula
is
one-line long; multiline
formulas require special formatting instructions.
   \[  \ip{\Gamma}{\psi'} = x'' + y^{2} + z_{i}^{n}\]
Don't start a paragraph with a displayed equation,
nor make one a paragraph by itself.

\end{document}               % End of document.

\section{Bench Section 114}
% synthetic bench section 114
% This is a sample LaTeX input file.  (Version of 12 August 2004.)
%
% A '%' character causes TeX to ignore all remaining text on the line,
% and is used for comments like this one.

\documentclass{article}      % Specifies the document class

                             % The preamble begins here.
\title{An Example Document}  % Declares the document's title.
\author{Leslie Lamport}      % Declares the author's name.
\date{January 21, 1994}      % Deleting this command produces today's date.

\newcommand{\ip}[2]{(#1, #2)}
                             % Defines \ip{arg1}{arg2} to mean
                             % (arg1, arg2).

%\newcommand{\ip}[2]{\langle #1 | #2\rangle}
                             % This is an alternative definition of
                             % \ip that is commented out.

\begin{document}             % End of preamble and beginning of text.

\maketitle                   % Produces the title.

This is an example input file.  Comparing it with
the output it generates can show you how to
produce a simple document of your own.

\section{Ordinary Text}      % Produces section heading.  Lower-level
                             % sections are begun with similar
                             % \subsection and \subsubsection commands.

The ends  of words and sentences are marked
  by   spaces. It  doesn't matter how many
spaces    you type; one is as good as 100.  The
end of   a line counts as a space.

One   or more   blank lines denote the  end
of  a paragraph.

Since any number of consecutive spaces are treated
like a single one, the formatting of the input
file makes no difference to
      \LaTeX,                % The \LaTeX command generates the LaTeX logo.
but it makes a difference to you.  When you use
\LaTeX, making your input file as easy to read
as possible will be a great help as you write
your document and when you change it.  This sample
file shows how you can add comments to your own input
file.

Because printing is different from typewriting,
there are a number of things that you have to do
differently when preparing an input file than if
you were just typing the document directly.
Quotation marks like
       ``this''
have to be handled specially, as do quotes within
quotes:
       ``\,`this'            % \, separates the double and single quote.
        is what I just
        wrote, not  `that'\,''.

Dashes come in three sizes: an
       intra-word
dash, a medium dash for number ranges like
       1--2,
and a punctuation
       dash---like
this.

A sentence-ending space should be larger than the
space between words within a sentence.  You
sometimes have to type special commands in
conjunction with punctuation characters to get
this right, as in the following sentence.
       Gnats, gnus, etc.\ all  % `\ ' makes an inter-word space.
       begin with G\@.         % \@ marks end-of-sentence punctuation.
You should check the spaces after periods when
reading your output to make sure you haven't
forgotten any special cases.  Generating an
ellipsis
       \ldots\               % `\ ' is needed after `\ldots' because TeX
                             % ignores spaces after command names like \ldots
                             % made from \ + letters.
                             %
                             % Note how a `%' character causes TeX to ignore
                             % the end of the input line, so these blank lines
                             % do not start a new paragraph.
                             %
with the right spacing around the periods requires
a special command.

\LaTeX\ interprets some common characters as
commands, so you must type special commands to
generate them.  These characters include the
following:
       \$ \& \% \# \{ and \}.

In printing, text is usually emphasized with an
       \emph{italic}
type style.

\begin{em}
   A long segment of text can also be emphasized
   in this way.  Text within such a segment can be
   given \emph{additional} emphasis.
\end{em}

It is sometimes necessary to prevent \LaTeX\ from
breaking a line where it might otherwise do so.
This may be at a space, as between the ``Mr.''\ and
``Jones'' in
       ``Mr.~Jones'',        % ~ produces an unbreakable interword space.
or within a word---especially when the word is a
symbol like
       \mbox{\emph{itemnum}}
that makes little sense when hyphenated across
lines.

Footnotes\footnote{This is an example of a footnote.}
pose no problem.

\LaTeX\ is good at typesetting mathematical formulas
like
       \( x-3y + z = 7 \)
or
       \( a_{1} > x^{2n} + y^{2n} > x' \)
or
       \( \ip{A}{B} = \sum_{i} a_{i} b_{i} \).
The spaces you type in a formula are
ignored.  Remember that a letter like
       $x$                   % $ ... $  and  \( ... \)  are equivalent
is a formula when it denotes a mathematical
symbol, and it should be typed as one.

\section{Displayed Text}

Text is displayed by indenting it from the left
margin.  Quotations are commonly displayed.  There
are short quotations
\begin{quote}
   This is a short quotation.  It consists of a
   single paragraph of text.  See how it is formatted.
\end{quote}
and longer ones.
\begin{quotation}
   This is a longer quotation.  It consists of two
   paragraphs of text, neither of which are
   particularly interesting.

   This is the second paragraph of the quotation.  It
   is just as dull as the first paragraph.
\end{quotation}
Another frequently-displayed structure is a list.
The following is an example of an \emph{itemized}
list.
\begin{itemize}
   \item This is the first item of an itemized list.
         Each item in the list is marked with a ``tick''.
         You don't have to worry about what kind of tick
         mark is used.

   \item This is the second item of the list.  It
         contains another list nested inside it.  The inner
         list is an \emph{enumerated} list.
         \begin{enumerate}
            \item This is the first item of an enumerated
                  list that is nested within the itemized list.

            \item This is the second item of the inner list.
                  \LaTeX\ allows you to nest lists deeper than
                  you really should.
         \end{enumerate}
         This is the rest of the second item of the outer
         list.  It is no more interesting than any other
         part of the item.
   \item This is the third item of the list.
\end{itemize}
You can even display poetry.
\begin{verse}
   There is an environment
    for verse \\             % The \\ command separates lines
   Whose features some poets % within a stanza.
   will curse.

                             % One or more blank lines separate stanzas.

   For instead of making\\
   Them do \emph{all} line breaking, \\
   It allows them to put too many words on a line when they'd rather be
   forced to be terse.
\end{verse}

Mathematical formulas may also be displayed.  A
displayed formula
is
one-line long; multiline
formulas require special formatting instructions.
   \[  \ip{\Gamma}{\psi'} = x'' + y^{2} + z_{i}^{n}\]
Don't start a paragraph with a displayed equation,
nor make one a paragraph by itself.

\end{document}               % End of document.

\section{Bench Section 115}
% synthetic bench section 115
% This is a sample LaTeX input file.  (Version of 12 August 2004.)
%
% A '%' character causes TeX to ignore all remaining text on the line,
% and is used for comments like this one.

\documentclass{article}      % Specifies the document class

                             % The preamble begins here.
\title{An Example Document}  % Declares the document's title.
\author{Leslie Lamport}      % Declares the author's name.
\date{January 21, 1994}      % Deleting this command produces today's date.

\newcommand{\ip}[2]{(#1, #2)}
                             % Defines \ip{arg1}{arg2} to mean
                             % (arg1, arg2).

%\newcommand{\ip}[2]{\langle #1 | #2\rangle}
                             % This is an alternative definition of
                             % \ip that is commented out.

\begin{document}             % End of preamble and beginning of text.

\maketitle                   % Produces the title.

This is an example input file.  Comparing it with
the output it generates can show you how to
produce a simple document of your own.

\section{Ordinary Text}      % Produces section heading.  Lower-level
                             % sections are begun with similar
                             % \subsection and \subsubsection commands.

The ends  of words and sentences are marked
  by   spaces. It  doesn't matter how many
spaces    you type; one is as good as 100.  The
end of   a line counts as a space.

One   or more   blank lines denote the  end
of  a paragraph.

Since any number of consecutive spaces are treated
like a single one, the formatting of the input
file makes no difference to
      \LaTeX,                % The \LaTeX command generates the LaTeX logo.
but it makes a difference to you.  When you use
\LaTeX, making your input file as easy to read
as possible will be a great help as you write
your document and when you change it.  This sample
file shows how you can add comments to your own input
file.

Because printing is different from typewriting,
there are a number of things that you have to do
differently when preparing an input file than if
you were just typing the document directly.
Quotation marks like
       ``this''
have to be handled specially, as do quotes within
quotes:
       ``\,`this'            % \, separates the double and single quote.
        is what I just
        wrote, not  `that'\,''.

Dashes come in three sizes: an
       intra-word
dash, a medium dash for number ranges like
       1--2,
and a punctuation
       dash---like
this.

A sentence-ending space should be larger than the
space between words within a sentence.  You
sometimes have to type special commands in
conjunction with punctuation characters to get
this right, as in the following sentence.
       Gnats, gnus, etc.\ all  % `\ ' makes an inter-word space.
       begin with G\@.         % \@ marks end-of-sentence punctuation.
You should check the spaces after periods when
reading your output to make sure you haven't
forgotten any special cases.  Generating an
ellipsis
       \ldots\               % `\ ' is needed after `\ldots' because TeX
                             % ignores spaces after command names like \ldots
                             % made from \ + letters.
                             %
                             % Note how a `%' character causes TeX to ignore
                             % the end of the input line, so these blank lines
                             % do not start a new paragraph.
                             %
with the right spacing around the periods requires
a special command.

\LaTeX\ interprets some common characters as
commands, so you must type special commands to
generate them.  These characters include the
following:
       \$ \& \% \# \{ and \}.

In printing, text is usually emphasized with an
       \emph{italic}
type style.

\begin{em}
   A long segment of text can also be emphasized
   in this way.  Text within such a segment can be
   given \emph{additional} emphasis.
\end{em}

It is sometimes necessary to prevent \LaTeX\ from
breaking a line where it might otherwise do so.
This may be at a space, as between the ``Mr.''\ and
``Jones'' in
       ``Mr.~Jones'',        % ~ produces an unbreakable interword space.
or within a word---especially when the word is a
symbol like
       \mbox{\emph{itemnum}}
that makes little sense when hyphenated across
lines.

Footnotes\footnote{This is an example of a footnote.}
pose no problem.

\LaTeX\ is good at typesetting mathematical formulas
like
       \( x-3y + z = 7 \)
or
       \( a_{1} > x^{2n} + y^{2n} > x' \)
or
       \( \ip{A}{B} = \sum_{i} a_{i} b_{i} \).
The spaces you type in a formula are
ignored.  Remember that a letter like
       $x$                   % $ ... $  and  \( ... \)  are equivalent
is a formula when it denotes a mathematical
symbol, and it should be typed as one.

\section{Displayed Text}

Text is displayed by indenting it from the left
margin.  Quotations are commonly displayed.  There
are short quotations
\begin{quote}
   This is a short quotation.  It consists of a
   single paragraph of text.  See how it is formatted.
\end{quote}
and longer ones.
\begin{quotation}
   This is a longer quotation.  It consists of two
   paragraphs of text, neither of which are
   particularly interesting.

   This is the second paragraph of the quotation.  It
   is just as dull as the first paragraph.
\end{quotation}
Another frequently-displayed structure is a list.
The following is an example of an \emph{itemized}
list.
\begin{itemize}
   \item This is the first item of an itemized list.
         Each item in the list is marked with a ``tick''.
         You don't have to worry about what kind of tick
         mark is used.

   \item This is the second item of the list.  It
         contains another list nested inside it.  The inner
         list is an \emph{enumerated} list.
         \begin{enumerate}
            \item This is the first item of an enumerated
                  list that is nested within the itemized list.

            \item This is the second item of the inner list.
                  \LaTeX\ allows you to nest lists deeper than
                  you really should.
         \end{enumerate}
         This is the rest of the second item of the outer
         list.  It is no more interesting than any other
         part of the item.
   \item This is the third item of the list.
\end{itemize}
You can even display poetry.
\begin{verse}
   There is an environment
    for verse \\             % The \\ command separates lines
   Whose features some poets % within a stanza.
   will curse.

                             % One or more blank lines separate stanzas.

   For instead of making\\
   Them do \emph{all} line breaking, \\
   It allows them to put too many words on a line when they'd rather be
   forced to be terse.
\end{verse}

Mathematical formulas may also be displayed.  A
displayed formula
is
one-line long; multiline
formulas require special formatting instructions.
   \[  \ip{\Gamma}{\psi'} = x'' + y^{2} + z_{i}^{n}\]
Don't start a paragraph with a displayed equation,
nor make one a paragraph by itself.

\end{document}               % End of document.

\section{Bench Section 116}
% synthetic bench section 116
% This is a sample LaTeX input file.  (Version of 12 August 2004.)
%
% A '%' character causes TeX to ignore all remaining text on the line,
% and is used for comments like this one.

\documentclass{article}      % Specifies the document class

                             % The preamble begins here.
\title{An Example Document}  % Declares the document's title.
\author{Leslie Lamport}      % Declares the author's name.
\date{January 21, 1994}      % Deleting this command produces today's date.

\newcommand{\ip}[2]{(#1, #2)}
                             % Defines \ip{arg1}{arg2} to mean
                             % (arg1, arg2).

%\newcommand{\ip}[2]{\langle #1 | #2\rangle}
                             % This is an alternative definition of
                             % \ip that is commented out.

\begin{document}             % End of preamble and beginning of text.

\maketitle                   % Produces the title.

This is an example input file.  Comparing it with
the output it generates can show you how to
produce a simple document of your own.

\section{Ordinary Text}      % Produces section heading.  Lower-level
                             % sections are begun with similar
                             % \subsection and \subsubsection commands.

The ends  of words and sentences are marked
  by   spaces. It  doesn't matter how many
spaces    you type; one is as good as 100.  The
end of   a line counts as a space.

One   or more   blank lines denote the  end
of  a paragraph.

Since any number of consecutive spaces are treated
like a single one, the formatting of the input
file makes no difference to
      \LaTeX,                % The \LaTeX command generates the LaTeX logo.
but it makes a difference to you.  When you use
\LaTeX, making your input file as easy to read
as possible will be a great help as you write
your document and when you change it.  This sample
file shows how you can add comments to your own input
file.

Because printing is different from typewriting,
there are a number of things that you have to do
differently when preparing an input file than if
you were just typing the document directly.
Quotation marks like
       ``this''
have to be handled specially, as do quotes within
quotes:
       ``\,`this'            % \, separates the double and single quote.
        is what I just
        wrote, not  `that'\,''.

Dashes come in three sizes: an
       intra-word
dash, a medium dash for number ranges like
       1--2,
and a punctuation
       dash---like
this.

A sentence-ending space should be larger than the
space between words within a sentence.  You
sometimes have to type special commands in
conjunction with punctuation characters to get
this right, as in the following sentence.
       Gnats, gnus, etc.\ all  % `\ ' makes an inter-word space.
       begin with G\@.         % \@ marks end-of-sentence punctuation.
You should check the spaces after periods when
reading your output to make sure you haven't
forgotten any special cases.  Generating an
ellipsis
       \ldots\               % `\ ' is needed after `\ldots' because TeX
                             % ignores spaces after command names like \ldots
                             % made from \ + letters.
                             %
                             % Note how a `%' character causes TeX to ignore
                             % the end of the input line, so these blank lines
                             % do not start a new paragraph.
                             %
with the right spacing around the periods requires
a special command.

\LaTeX\ interprets some common characters as
commands, so you must type special commands to
generate them.  These characters include the
following:
       \$ \& \% \# \{ and \}.

In printing, text is usually emphasized with an
       \emph{italic}
type style.

\begin{em}
   A long segment of text can also be emphasized
   in this way.  Text within such a segment can be
   given \emph{additional} emphasis.
\end{em}

It is sometimes necessary to prevent \LaTeX\ from
breaking a line where it might otherwise do so.
This may be at a space, as between the ``Mr.''\ and
``Jones'' in
       ``Mr.~Jones'',        % ~ produces an unbreakable interword space.
or within a word---especially when the word is a
symbol like
       \mbox{\emph{itemnum}}
that makes little sense when hyphenated across
lines.

Footnotes\footnote{This is an example of a footnote.}
pose no problem.

\LaTeX\ is good at typesetting mathematical formulas
like
       \( x-3y + z = 7 \)
or
       \( a_{1} > x^{2n} + y^{2n} > x' \)
or
       \( \ip{A}{B} = \sum_{i} a_{i} b_{i} \).
The spaces you type in a formula are
ignored.  Remember that a letter like
       $x$                   % $ ... $  and  \( ... \)  are equivalent
is a formula when it denotes a mathematical
symbol, and it should be typed as one.

\section{Displayed Text}

Text is displayed by indenting it from the left
margin.  Quotations are commonly displayed.  There
are short quotations
\begin{quote}
   This is a short quotation.  It consists of a
   single paragraph of text.  See how it is formatted.
\end{quote}
and longer ones.
\begin{quotation}
   This is a longer quotation.  It consists of two
   paragraphs of text, neither of which are
   particularly interesting.

   This is the second paragraph of the quotation.  It
   is just as dull as the first paragraph.
\end{quotation}
Another frequently-displayed structure is a list.
The following is an example of an \emph{itemized}
list.
\begin{itemize}
   \item This is the first item of an itemized list.
         Each item in the list is marked with a ``tick''.
         You don't have to worry about what kind of tick
         mark is used.

   \item This is the second item of the list.  It
         contains another list nested inside it.  The inner
         list is an \emph{enumerated} list.
         \begin{enumerate}
            \item This is the first item of an enumerated
                  list that is nested within the itemized list.

            \item This is the second item of the inner list.
                  \LaTeX\ allows you to nest lists deeper than
                  you really should.
         \end{enumerate}
         This is the rest of the second item of the outer
         list.  It is no more interesting than any other
         part of the item.
   \item This is the third item of the list.
\end{itemize}
You can even display poetry.
\begin{verse}
   There is an environment
    for verse \\             % The \\ command separates lines
   Whose features some poets % within a stanza.
   will curse.

                             % One or more blank lines separate stanzas.

   For instead of making\\
   Them do \emph{all} line breaking, \\
   It allows them to put too many words on a line when they'd rather be
   forced to be terse.
\end{verse}

Mathematical formulas may also be displayed.  A
displayed formula
is
one-line long; multiline
formulas require special formatting instructions.
   \[  \ip{\Gamma}{\psi'} = x'' + y^{2} + z_{i}^{n}\]
Don't start a paragraph with a displayed equation,
nor make one a paragraph by itself.

\end{document}               % End of document.

\section{Bench Section 117}
% synthetic bench section 117
% This is a sample LaTeX input file.  (Version of 12 August 2004.)
%
% A '%' character causes TeX to ignore all remaining text on the line,
% and is used for comments like this one.

\documentclass{article}      % Specifies the document class

                             % The preamble begins here.
\title{An Example Document}  % Declares the document's title.
\author{Leslie Lamport}      % Declares the author's name.
\date{January 21, 1994}      % Deleting this command produces today's date.

\newcommand{\ip}[2]{(#1, #2)}
                             % Defines \ip{arg1}{arg2} to mean
                             % (arg1, arg2).

%\newcommand{\ip}[2]{\langle #1 | #2\rangle}
                             % This is an alternative definition of
                             % \ip that is commented out.

\begin{document}             % End of preamble and beginning of text.

\maketitle                   % Produces the title.

This is an example input file.  Comparing it with
the output it generates can show you how to
produce a simple document of your own.

\section{Ordinary Text}      % Produces section heading.  Lower-level
                             % sections are begun with similar
                             % \subsection and \subsubsection commands.

The ends  of words and sentences are marked
  by   spaces. It  doesn't matter how many
spaces    you type; one is as good as 100.  The
end of   a line counts as a space.

One   or more   blank lines denote the  end
of  a paragraph.

Since any number of consecutive spaces are treated
like a single one, the formatting of the input
file makes no difference to
      \LaTeX,                % The \LaTeX command generates the LaTeX logo.
but it makes a difference to you.  When you use
\LaTeX, making your input file as easy to read
as possible will be a great help as you write
your document and when you change it.  This sample
file shows how you can add comments to your own input
file.

Because printing is different from typewriting,
there are a number of things that you have to do
differently when preparing an input file than if
you were just typing the document directly.
Quotation marks like
       ``this''
have to be handled specially, as do quotes within
quotes:
       ``\,`this'            % \, separates the double and single quote.
        is what I just
        wrote, not  `that'\,''.

Dashes come in three sizes: an
       intra-word
dash, a medium dash for number ranges like
       1--2,
and a punctuation
       dash---like
this.

A sentence-ending space should be larger than the
space between words within a sentence.  You
sometimes have to type special commands in
conjunction with punctuation characters to get
this right, as in the following sentence.
       Gnats, gnus, etc.\ all  % `\ ' makes an inter-word space.
       begin with G\@.         % \@ marks end-of-sentence punctuation.
You should check the spaces after periods when
reading your output to make sure you haven't
forgotten any special cases.  Generating an
ellipsis
       \ldots\               % `\ ' is needed after `\ldots' because TeX
                             % ignores spaces after command names like \ldots
                             % made from \ + letters.
                             %
                             % Note how a `%' character causes TeX to ignore
                             % the end of the input line, so these blank lines
                             % do not start a new paragraph.
                             %
with the right spacing around the periods requires
a special command.

\LaTeX\ interprets some common characters as
commands, so you must type special commands to
generate them.  These characters include the
following:
       \$ \& \% \# \{ and \}.

In printing, text is usually emphasized with an
       \emph{italic}
type style.

\begin{em}
   A long segment of text can also be emphasized
   in this way.  Text within such a segment can be
   given \emph{additional} emphasis.
\end{em}

It is sometimes necessary to prevent \LaTeX\ from
breaking a line where it might otherwise do so.
This may be at a space, as between the ``Mr.''\ and
``Jones'' in
       ``Mr.~Jones'',        % ~ produces an unbreakable interword space.
or within a word---especially when the word is a
symbol like
       \mbox{\emph{itemnum}}
that makes little sense when hyphenated across
lines.

Footnotes\footnote{This is an example of a footnote.}
pose no problem.

\LaTeX\ is good at typesetting mathematical formulas
like
       \( x-3y + z = 7 \)
or
       \( a_{1} > x^{2n} + y^{2n} > x' \)
or
       \( \ip{A}{B} = \sum_{i} a_{i} b_{i} \).
The spaces you type in a formula are
ignored.  Remember that a letter like
       $x$                   % $ ... $  and  \( ... \)  are equivalent
is a formula when it denotes a mathematical
symbol, and it should be typed as one.

\section{Displayed Text}

Text is displayed by indenting it from the left
margin.  Quotations are commonly displayed.  There
are short quotations
\begin{quote}
   This is a short quotation.  It consists of a
   single paragraph of text.  See how it is formatted.
\end{quote}
and longer ones.
\begin{quotation}
   This is a longer quotation.  It consists of two
   paragraphs of text, neither of which are
   particularly interesting.

   This is the second paragraph of the quotation.  It
   is just as dull as the first paragraph.
\end{quotation}
Another frequently-displayed structure is a list.
The following is an example of an \emph{itemized}
list.
\begin{itemize}
   \item This is the first item of an itemized list.
         Each item in the list is marked with a ``tick''.
         You don't have to worry about what kind of tick
         mark is used.

   \item This is the second item of the list.  It
         contains another list nested inside it.  The inner
         list is an \emph{enumerated} list.
         \begin{enumerate}
            \item This is the first item of an enumerated
                  list that is nested within the itemized list.

            \item This is the second item of the inner list.
                  \LaTeX\ allows you to nest lists deeper than
                  you really should.
         \end{enumerate}
         This is the rest of the second item of the outer
         list.  It is no more interesting than any other
         part of the item.
   \item This is the third item of the list.
\end{itemize}
You can even display poetry.
\begin{verse}
   There is an environment
    for verse \\             % The \\ command separates lines
   Whose features some poets % within a stanza.
   will curse.

                             % One or more blank lines separate stanzas.

   For instead of making\\
   Them do \emph{all} line breaking, \\
   It allows them to put too many words on a line when they'd rather be
   forced to be terse.
\end{verse}

Mathematical formulas may also be displayed.  A
displayed formula
is
one-line long; multiline
formulas require special formatting instructions.
   \[  \ip{\Gamma}{\psi'} = x'' + y^{2} + z_{i}^{n}\]
Don't start a paragraph with a displayed equation,
nor make one a paragraph by itself.

\end{document}               % End of document.

\section{Bench Section 118}
% synthetic bench section 118
% This is a sample LaTeX input file.  (Version of 12 August 2004.)
%
% A '%' character causes TeX to ignore all remaining text on the line,
% and is used for comments like this one.

\documentclass{article}      % Specifies the document class

                             % The preamble begins here.
\title{An Example Document}  % Declares the document's title.
\author{Leslie Lamport}      % Declares the author's name.
\date{January 21, 1994}      % Deleting this command produces today's date.

\newcommand{\ip}[2]{(#1, #2)}
                             % Defines \ip{arg1}{arg2} to mean
                             % (arg1, arg2).

%\newcommand{\ip}[2]{\langle #1 | #2\rangle}
                             % This is an alternative definition of
                             % \ip that is commented out.

\begin{document}             % End of preamble and beginning of text.

\maketitle                   % Produces the title.

This is an example input file.  Comparing it with
the output it generates can show you how to
produce a simple document of your own.

\section{Ordinary Text}      % Produces section heading.  Lower-level
                             % sections are begun with similar
                             % \subsection and \subsubsection commands.

The ends  of words and sentences are marked
  by   spaces. It  doesn't matter how many
spaces    you type; one is as good as 100.  The
end of   a line counts as a space.

One   or more   blank lines denote the  end
of  a paragraph.

Since any number of consecutive spaces are treated
like a single one, the formatting of the input
file makes no difference to
      \LaTeX,                % The \LaTeX command generates the LaTeX logo.
but it makes a difference to you.  When you use
\LaTeX, making your input file as easy to read
as possible will be a great help as you write
your document and when you change it.  This sample
file shows how you can add comments to your own input
file.

Because printing is different from typewriting,
there are a number of things that you have to do
differently when preparing an input file than if
you were just typing the document directly.
Quotation marks like
       ``this''
have to be handled specially, as do quotes within
quotes:
       ``\,`this'            % \, separates the double and single quote.
        is what I just
        wrote, not  `that'\,''.

Dashes come in three sizes: an
       intra-word
dash, a medium dash for number ranges like
       1--2,
and a punctuation
       dash---like
this.

A sentence-ending space should be larger than the
space between words within a sentence.  You
sometimes have to type special commands in
conjunction with punctuation characters to get
this right, as in the following sentence.
       Gnats, gnus, etc.\ all  % `\ ' makes an inter-word space.
       begin with G\@.         % \@ marks end-of-sentence punctuation.
You should check the spaces after periods when
reading your output to make sure you haven't
forgotten any special cases.  Generating an
ellipsis
       \ldots\               % `\ ' is needed after `\ldots' because TeX
                             % ignores spaces after command names like \ldots
                             % made from \ + letters.
                             %
                             % Note how a `%' character causes TeX to ignore
                             % the end of the input line, so these blank lines
                             % do not start a new paragraph.
                             %
with the right spacing around the periods requires
a special command.

\LaTeX\ interprets some common characters as
commands, so you must type special commands to
generate them.  These characters include the
following:
       \$ \& \% \# \{ and \}.

In printing, text is usually emphasized with an
       \emph{italic}
type style.

\begin{em}
   A long segment of text can also be emphasized
   in this way.  Text within such a segment can be
   given \emph{additional} emphasis.
\end{em}

It is sometimes necessary to prevent \LaTeX\ from
breaking a line where it might otherwise do so.
This may be at a space, as between the ``Mr.''\ and
``Jones'' in
       ``Mr.~Jones'',        % ~ produces an unbreakable interword space.
or within a word---especially when the word is a
symbol like
       \mbox{\emph{itemnum}}
that makes little sense when hyphenated across
lines.

Footnotes\footnote{This is an example of a footnote.}
pose no problem.

\LaTeX\ is good at typesetting mathematical formulas
like
       \( x-3y + z = 7 \)
or
       \( a_{1} > x^{2n} + y^{2n} > x' \)
or
       \( \ip{A}{B} = \sum_{i} a_{i} b_{i} \).
The spaces you type in a formula are
ignored.  Remember that a letter like
       $x$                   % $ ... $  and  \( ... \)  are equivalent
is a formula when it denotes a mathematical
symbol, and it should be typed as one.

\section{Displayed Text}

Text is displayed by indenting it from the left
margin.  Quotations are commonly displayed.  There
are short quotations
\begin{quote}
   This is a short quotation.  It consists of a
   single paragraph of text.  See how it is formatted.
\end{quote}
and longer ones.
\begin{quotation}
   This is a longer quotation.  It consists of two
   paragraphs of text, neither of which are
   particularly interesting.

   This is the second paragraph of the quotation.  It
   is just as dull as the first paragraph.
\end{quotation}
Another frequently-displayed structure is a list.
The following is an example of an \emph{itemized}
list.
\begin{itemize}
   \item This is the first item of an itemized list.
         Each item in the list is marked with a ``tick''.
         You don't have to worry about what kind of tick
         mark is used.

   \item This is the second item of the list.  It
         contains another list nested inside it.  The inner
         list is an \emph{enumerated} list.
         \begin{enumerate}
            \item This is the first item of an enumerated
                  list that is nested within the itemized list.

            \item This is the second item of the inner list.
                  \LaTeX\ allows you to nest lists deeper than
                  you really should.
         \end{enumerate}
         This is the rest of the second item of the outer
         list.  It is no more interesting than any other
         part of the item.
   \item This is the third item of the list.
\end{itemize}
You can even display poetry.
\begin{verse}
   There is an environment
    for verse \\             % The \\ command separates lines
   Whose features some poets % within a stanza.
   will curse.

                             % One or more blank lines separate stanzas.

   For instead of making\\
   Them do \emph{all} line breaking, \\
   It allows them to put too many words on a line when they'd rather be
   forced to be terse.
\end{verse}

Mathematical formulas may also be displayed.  A
displayed formula
is
one-line long; multiline
formulas require special formatting instructions.
   \[  \ip{\Gamma}{\psi'} = x'' + y^{2} + z_{i}^{n}\]
Don't start a paragraph with a displayed equation,
nor make one a paragraph by itself.

\end{document}               % End of document.

\section{Bench Section 119}
% synthetic bench section 119
% This is a sample LaTeX input file.  (Version of 12 August 2004.)
%
% A '%' character causes TeX to ignore all remaining text on the line,
% and is used for comments like this one.

\documentclass{article}      % Specifies the document class

                             % The preamble begins here.
\title{An Example Document}  % Declares the document's title.
\author{Leslie Lamport}      % Declares the author's name.
\date{January 21, 1994}      % Deleting this command produces today's date.

\newcommand{\ip}[2]{(#1, #2)}
                             % Defines \ip{arg1}{arg2} to mean
                             % (arg1, arg2).

%\newcommand{\ip}[2]{\langle #1 | #2\rangle}
                             % This is an alternative definition of
                             % \ip that is commented out.

\begin{document}             % End of preamble and beginning of text.

\maketitle                   % Produces the title.

This is an example input file.  Comparing it with
the output it generates can show you how to
produce a simple document of your own.

\section{Ordinary Text}      % Produces section heading.  Lower-level
                             % sections are begun with similar
                             % \subsection and \subsubsection commands.

The ends  of words and sentences are marked
  by   spaces. It  doesn't matter how many
spaces    you type; one is as good as 100.  The
end of   a line counts as a space.

One   or more   blank lines denote the  end
of  a paragraph.

Since any number of consecutive spaces are treated
like a single one, the formatting of the input
file makes no difference to
      \LaTeX,                % The \LaTeX command generates the LaTeX logo.
but it makes a difference to you.  When you use
\LaTeX, making your input file as easy to read
as possible will be a great help as you write
your document and when you change it.  This sample
file shows how you can add comments to your own input
file.

Because printing is different from typewriting,
there are a number of things that you have to do
differently when preparing an input file than if
you were just typing the document directly.
Quotation marks like
       ``this''
have to be handled specially, as do quotes within
quotes:
       ``\,`this'            % \, separates the double and single quote.
        is what I just
        wrote, not  `that'\,''.

Dashes come in three sizes: an
       intra-word
dash, a medium dash for number ranges like
       1--2,
and a punctuation
       dash---like
this.

A sentence-ending space should be larger than the
space between words within a sentence.  You
sometimes have to type special commands in
conjunction with punctuation characters to get
this right, as in the following sentence.
       Gnats, gnus, etc.\ all  % `\ ' makes an inter-word space.
       begin with G\@.         % \@ marks end-of-sentence punctuation.
You should check the spaces after periods when
reading your output to make sure you haven't
forgotten any special cases.  Generating an
ellipsis
       \ldots\               % `\ ' is needed after `\ldots' because TeX
                             % ignores spaces after command names like \ldots
                             % made from \ + letters.
                             %
                             % Note how a `%' character causes TeX to ignore
                             % the end of the input line, so these blank lines
                             % do not start a new paragraph.
                             %
with the right spacing around the periods requires
a special command.

\LaTeX\ interprets some common characters as
commands, so you must type special commands to
generate them.  These characters include the
following:
       \$ \& \% \# \{ and \}.

In printing, text is usually emphasized with an
       \emph{italic}
type style.

\begin{em}
   A long segment of text can also be emphasized
   in this way.  Text within such a segment can be
   given \emph{additional} emphasis.
\end{em}

It is sometimes necessary to prevent \LaTeX\ from
breaking a line where it might otherwise do so.
This may be at a space, as between the ``Mr.''\ and
``Jones'' in
       ``Mr.~Jones'',        % ~ produces an unbreakable interword space.
or within a word---especially when the word is a
symbol like
       \mbox{\emph{itemnum}}
that makes little sense when hyphenated across
lines.

Footnotes\footnote{This is an example of a footnote.}
pose no problem.

\LaTeX\ is good at typesetting mathematical formulas
like
       \( x-3y + z = 7 \)
or
       \( a_{1} > x^{2n} + y^{2n} > x' \)
or
       \( \ip{A}{B} = \sum_{i} a_{i} b_{i} \).
The spaces you type in a formula are
ignored.  Remember that a letter like
       $x$                   % $ ... $  and  \( ... \)  are equivalent
is a formula when it denotes a mathematical
symbol, and it should be typed as one.

\section{Displayed Text}

Text is displayed by indenting it from the left
margin.  Quotations are commonly displayed.  There
are short quotations
\begin{quote}
   This is a short quotation.  It consists of a
   single paragraph of text.  See how it is formatted.
\end{quote}
and longer ones.
\begin{quotation}
   This is a longer quotation.  It consists of two
   paragraphs of text, neither of which are
   particularly interesting.

   This is the second paragraph of the quotation.  It
   is just as dull as the first paragraph.
\end{quotation}
Another frequently-displayed structure is a list.
The following is an example of an \emph{itemized}
list.
\begin{itemize}
   \item This is the first item of an itemized list.
         Each item in the list is marked with a ``tick''.
         You don't have to worry about what kind of tick
         mark is used.

   \item This is the second item of the list.  It
         contains another list nested inside it.  The inner
         list is an \emph{enumerated} list.
         \begin{enumerate}
            \item This is the first item of an enumerated
                  list that is nested within the itemized list.

            \item This is the second item of the inner list.
                  \LaTeX\ allows you to nest lists deeper than
                  you really should.
         \end{enumerate}
         This is the rest of the second item of the outer
         list.  It is no more interesting than any other
         part of the item.
   \item This is the third item of the list.
\end{itemize}
You can even display poetry.
\begin{verse}
   There is an environment
    for verse \\             % The \\ command separates lines
   Whose features some poets % within a stanza.
   will curse.

                             % One or more blank lines separate stanzas.

   For instead of making\\
   Them do \emph{all} line breaking, \\
   It allows them to put too many words on a line when they'd rather be
   forced to be terse.
\end{verse}

Mathematical formulas may also be displayed.  A
displayed formula
is
one-line long; multiline
formulas require special formatting instructions.
   \[  \ip{\Gamma}{\psi'} = x'' + y^{2} + z_{i}^{n}\]
Don't start a paragraph with a displayed equation,
nor make one a paragraph by itself.

\end{document}               % End of document.

\section{Bench Section 120}
% synthetic bench section 120
% This is a sample LaTeX input file.  (Version of 12 August 2004.)
%
% A '%' character causes TeX to ignore all remaining text on the line,
% and is used for comments like this one.

\documentclass{article}      % Specifies the document class

                             % The preamble begins here.
\title{An Example Document}  % Declares the document's title.
\author{Leslie Lamport}      % Declares the author's name.
\date{January 21, 1994}      % Deleting this command produces today's date.

\newcommand{\ip}[2]{(#1, #2)}
                             % Defines \ip{arg1}{arg2} to mean
                             % (arg1, arg2).

%\newcommand{\ip}[2]{\langle #1 | #2\rangle}
                             % This is an alternative definition of
                             % \ip that is commented out.

\begin{document}             % End of preamble and beginning of text.

\maketitle                   % Produces the title.

This is an example input file.  Comparing it with
the output it generates can show you how to
produce a simple document of your own.

\section{Ordinary Text}      % Produces section heading.  Lower-level
                             % sections are begun with similar
                             % \subsection and \subsubsection commands.

The ends  of words and sentences are marked
  by   spaces. It  doesn't matter how many
spaces    you type; one is as good as 100.  The
end of   a line counts as a space.

One   or more   blank lines denote the  end
of  a paragraph.

Since any number of consecutive spaces are treated
like a single one, the formatting of the input
file makes no difference to
      \LaTeX,                % The \LaTeX command generates the LaTeX logo.
but it makes a difference to you.  When you use
\LaTeX, making your input file as easy to read
as possible will be a great help as you write
your document and when you change it.  This sample
file shows how you can add comments to your own input
file.

Because printing is different from typewriting,
there are a number of things that you have to do
differently when preparing an input file than if
you were just typing the document directly.
Quotation marks like
       ``this''
have to be handled specially, as do quotes within
quotes:
       ``\,`this'            % \, separates the double and single quote.
        is what I just
        wrote, not  `that'\,''.

Dashes come in three sizes: an
       intra-word
dash, a medium dash for number ranges like
       1--2,
and a punctuation
       dash---like
this.

A sentence-ending space should be larger than the
space between words within a sentence.  You
sometimes have to type special commands in
conjunction with punctuation characters to get
this right, as in the following sentence.
       Gnats, gnus, etc.\ all  % `\ ' makes an inter-word space.
       begin with G\@.         % \@ marks end-of-sentence punctuation.
You should check the spaces after periods when
reading your output to make sure you haven't
forgotten any special cases.  Generating an
ellipsis
       \ldots\               % `\ ' is needed after `\ldots' because TeX
                             % ignores spaces after command names like \ldots
                             % made from \ + letters.
                             %
                             % Note how a `%' character causes TeX to ignore
                             % the end of the input line, so these blank lines
                             % do not start a new paragraph.
                             %
with the right spacing around the periods requires
a special command.

\LaTeX\ interprets some common characters as
commands, so you must type special commands to
generate them.  These characters include the
following:
       \$ \& \% \# \{ and \}.

In printing, text is usually emphasized with an
       \emph{italic}
type style.

\begin{em}
   A long segment of text can also be emphasized
   in this way.  Text within such a segment can be
   given \emph{additional} emphasis.
\end{em}

It is sometimes necessary to prevent \LaTeX\ from
breaking a line where it might otherwise do so.
This may be at a space, as between the ``Mr.''\ and
``Jones'' in
       ``Mr.~Jones'',        % ~ produces an unbreakable interword space.
or within a word---especially when the word is a
symbol like
       \mbox{\emph{itemnum}}
that makes little sense when hyphenated across
lines.

Footnotes\footnote{This is an example of a footnote.}
pose no problem.

\LaTeX\ is good at typesetting mathematical formulas
like
       \( x-3y + z = 7 \)
or
       \( a_{1} > x^{2n} + y^{2n} > x' \)
or
       \( \ip{A}{B} = \sum_{i} a_{i} b_{i} \).
The spaces you type in a formula are
ignored.  Remember that a letter like
       $x$                   % $ ... $  and  \( ... \)  are equivalent
is a formula when it denotes a mathematical
symbol, and it should be typed as one.

\section{Displayed Text}

Text is displayed by indenting it from the left
margin.  Quotations are commonly displayed.  There
are short quotations
\begin{quote}
   This is a short quotation.  It consists of a
   single paragraph of text.  See how it is formatted.
\end{quote}
and longer ones.
\begin{quotation}
   This is a longer quotation.  It consists of two
   paragraphs of text, neither of which are
   particularly interesting.

   This is the second paragraph of the quotation.  It
   is just as dull as the first paragraph.
\end{quotation}
Another frequently-displayed structure is a list.
The following is an example of an \emph{itemized}
list.
\begin{itemize}
   \item This is the first item of an itemized list.
         Each item in the list is marked with a ``tick''.
         You don't have to worry about what kind of tick
         mark is used.

   \item This is the second item of the list.  It
         contains another list nested inside it.  The inner
         list is an \emph{enumerated} list.
         \begin{enumerate}
            \item This is the first item of an enumerated
                  list that is nested within the itemized list.

            \item This is the second item of the inner list.
                  \LaTeX\ allows you to nest lists deeper than
                  you really should.
         \end{enumerate}
         This is the rest of the second item of the outer
         list.  It is no more interesting than any other
         part of the item.
   \item This is the third item of the list.
\end{itemize}
You can even display poetry.
\begin{verse}
   There is an environment
    for verse \\             % The \\ command separates lines
   Whose features some poets % within a stanza.
   will curse.

                             % One or more blank lines separate stanzas.

   For instead of making\\
   Them do \emph{all} line breaking, \\
   It allows them to put too many words on a line when they'd rather be
   forced to be terse.
\end{verse}

Mathematical formulas may also be displayed.  A
displayed formula
is
one-line long; multiline
formulas require special formatting instructions.
   \[  \ip{\Gamma}{\psi'} = x'' + y^{2} + z_{i}^{n}\]
Don't start a paragraph with a displayed equation,
nor make one a paragraph by itself.

\end{document}               % End of document.

\section{Bench Section 121}
% synthetic bench section 121
% This is a sample LaTeX input file.  (Version of 12 August 2004.)
%
% A '%' character causes TeX to ignore all remaining text on the line,
% and is used for comments like this one.

\documentclass{article}      % Specifies the document class

                             % The preamble begins here.
\title{An Example Document}  % Declares the document's title.
\author{Leslie Lamport}      % Declares the author's name.
\date{January 21, 1994}      % Deleting this command produces today's date.

\newcommand{\ip}[2]{(#1, #2)}
                             % Defines \ip{arg1}{arg2} to mean
                             % (arg1, arg2).

%\newcommand{\ip}[2]{\langle #1 | #2\rangle}
                             % This is an alternative definition of
                             % \ip that is commented out.

\begin{document}             % End of preamble and beginning of text.

\maketitle                   % Produces the title.

This is an example input file.  Comparing it with
the output it generates can show you how to
produce a simple document of your own.

\section{Ordinary Text}      % Produces section heading.  Lower-level
                             % sections are begun with similar
                             % \subsection and \subsubsection commands.

The ends  of words and sentences are marked
  by   spaces. It  doesn't matter how many
spaces    you type; one is as good as 100.  The
end of   a line counts as a space.

One   or more   blank lines denote the  end
of  a paragraph.

Since any number of consecutive spaces are treated
like a single one, the formatting of the input
file makes no difference to
      \LaTeX,                % The \LaTeX command generates the LaTeX logo.
but it makes a difference to you.  When you use
\LaTeX, making your input file as easy to read
as possible will be a great help as you write
your document and when you change it.  This sample
file shows how you can add comments to your own input
file.

Because printing is different from typewriting,
there are a number of things that you have to do
differently when preparing an input file than if
you were just typing the document directly.
Quotation marks like
       ``this''
have to be handled specially, as do quotes within
quotes:
       ``\,`this'            % \, separates the double and single quote.
        is what I just
        wrote, not  `that'\,''.

Dashes come in three sizes: an
       intra-word
dash, a medium dash for number ranges like
       1--2,
and a punctuation
       dash---like
this.

A sentence-ending space should be larger than the
space between words within a sentence.  You
sometimes have to type special commands in
conjunction with punctuation characters to get
this right, as in the following sentence.
       Gnats, gnus, etc.\ all  % `\ ' makes an inter-word space.
       begin with G\@.         % \@ marks end-of-sentence punctuation.
You should check the spaces after periods when
reading your output to make sure you haven't
forgotten any special cases.  Generating an
ellipsis
       \ldots\               % `\ ' is needed after `\ldots' because TeX
                             % ignores spaces after command names like \ldots
                             % made from \ + letters.
                             %
                             % Note how a `%' character causes TeX to ignore
                             % the end of the input line, so these blank lines
                             % do not start a new paragraph.
                             %
with the right spacing around the periods requires
a special command.

\LaTeX\ interprets some common characters as
commands, so you must type special commands to
generate them.  These characters include the
following:
       \$ \& \% \# \{ and \}.

In printing, text is usually emphasized with an
       \emph{italic}
type style.

\begin{em}
   A long segment of text can also be emphasized
   in this way.  Text within such a segment can be
   given \emph{additional} emphasis.
\end{em}

It is sometimes necessary to prevent \LaTeX\ from
breaking a line where it might otherwise do so.
This may be at a space, as between the ``Mr.''\ and
``Jones'' in
       ``Mr.~Jones'',        % ~ produces an unbreakable interword space.
or within a word---especially when the word is a
symbol like
       \mbox{\emph{itemnum}}
that makes little sense when hyphenated across
lines.

Footnotes\footnote{This is an example of a footnote.}
pose no problem.

\LaTeX\ is good at typesetting mathematical formulas
like
       \( x-3y + z = 7 \)
or
       \( a_{1} > x^{2n} + y^{2n} > x' \)
or
       \( \ip{A}{B} = \sum_{i} a_{i} b_{i} \).
The spaces you type in a formula are
ignored.  Remember that a letter like
       $x$                   % $ ... $  and  \( ... \)  are equivalent
is a formula when it denotes a mathematical
symbol, and it should be typed as one.

\section{Displayed Text}

Text is displayed by indenting it from the left
margin.  Quotations are commonly displayed.  There
are short quotations
\begin{quote}
   This is a short quotation.  It consists of a
   single paragraph of text.  See how it is formatted.
\end{quote}
and longer ones.
\begin{quotation}
   This is a longer quotation.  It consists of two
   paragraphs of text, neither of which are
   particularly interesting.

   This is the second paragraph of the quotation.  It
   is just as dull as the first paragraph.
\end{quotation}
Another frequently-displayed structure is a list.
The following is an example of an \emph{itemized}
list.
\begin{itemize}
   \item This is the first item of an itemized list.
         Each item in the list is marked with a ``tick''.
         You don't have to worry about what kind of tick
         mark is used.

   \item This is the second item of the list.  It
         contains another list nested inside it.  The inner
         list is an \emph{enumerated} list.
         \begin{enumerate}
            \item This is the first item of an enumerated
                  list that is nested within the itemized list.

            \item This is the second item of the inner list.
                  \LaTeX\ allows you to nest lists deeper than
                  you really should.
         \end{enumerate}
         This is the rest of the second item of the outer
         list.  It is no more interesting than any other
         part of the item.
   \item This is the third item of the list.
\end{itemize}
You can even display poetry.
\begin{verse}
   There is an environment
    for verse \\             % The \\ command separates lines
   Whose features some poets % within a stanza.
   will curse.

                             % One or more blank lines separate stanzas.

   For instead of making\\
   Them do \emph{all} line breaking, \\
   It allows them to put too many words on a line when they'd rather be
   forced to be terse.
\end{verse}

Mathematical formulas may also be displayed.  A
displayed formula
is
one-line long; multiline
formulas require special formatting instructions.
   \[  \ip{\Gamma}{\psi'} = x'' + y^{2} + z_{i}^{n}\]
Don't start a paragraph with a displayed equation,
nor make one a paragraph by itself.

\end{document}               % End of document.

\section{Bench Section 122}
% synthetic bench section 122
% This is a sample LaTeX input file.  (Version of 12 August 2004.)
%
% A '%' character causes TeX to ignore all remaining text on the line,
% and is used for comments like this one.

\documentclass{article}      % Specifies the document class

                             % The preamble begins here.
\title{An Example Document}  % Declares the document's title.
\author{Leslie Lamport}      % Declares the author's name.
\date{January 21, 1994}      % Deleting this command produces today's date.

\newcommand{\ip}[2]{(#1, #2)}
                             % Defines \ip{arg1}{arg2} to mean
                             % (arg1, arg2).

%\newcommand{\ip}[2]{\langle #1 | #2\rangle}
                             % This is an alternative definition of
                             % \ip that is commented out.

\begin{document}             % End of preamble and beginning of text.

\maketitle                   % Produces the title.

This is an example input file.  Comparing it with
the output it generates can show you how to
produce a simple document of your own.

\section{Ordinary Text}      % Produces section heading.  Lower-level
                             % sections are begun with similar
                             % \subsection and \subsubsection commands.

The ends  of words and sentences are marked
  by   spaces. It  doesn't matter how many
spaces    you type; one is as good as 100.  The
end of   a line counts as a space.

One   or more   blank lines denote the  end
of  a paragraph.

Since any number of consecutive spaces are treated
like a single one, the formatting of the input
file makes no difference to
      \LaTeX,                % The \LaTeX command generates the LaTeX logo.
but it makes a difference to you.  When you use
\LaTeX, making your input file as easy to read
as possible will be a great help as you write
your document and when you change it.  This sample
file shows how you can add comments to your own input
file.

Because printing is different from typewriting,
there are a number of things that you have to do
differently when preparing an input file than if
you were just typing the document directly.
Quotation marks like
       ``this''
have to be handled specially, as do quotes within
quotes:
       ``\,`this'            % \, separates the double and single quote.
        is what I just
        wrote, not  `that'\,''.

Dashes come in three sizes: an
       intra-word
dash, a medium dash for number ranges like
       1--2,
and a punctuation
       dash---like
this.

A sentence-ending space should be larger than the
space between words within a sentence.  You
sometimes have to type special commands in
conjunction with punctuation characters to get
this right, as in the following sentence.
       Gnats, gnus, etc.\ all  % `\ ' makes an inter-word space.
       begin with G\@.         % \@ marks end-of-sentence punctuation.
You should check the spaces after periods when
reading your output to make sure you haven't
forgotten any special cases.  Generating an
ellipsis
       \ldots\               % `\ ' is needed after `\ldots' because TeX
                             % ignores spaces after command names like \ldots
                             % made from \ + letters.
                             %
                             % Note how a `%' character causes TeX to ignore
                             % the end of the input line, so these blank lines
                             % do not start a new paragraph.
                             %
with the right spacing around the periods requires
a special command.

\LaTeX\ interprets some common characters as
commands, so you must type special commands to
generate them.  These characters include the
following:
       \$ \& \% \# \{ and \}.

In printing, text is usually emphasized with an
       \emph{italic}
type style.

\begin{em}
   A long segment of text can also be emphasized
   in this way.  Text within such a segment can be
   given \emph{additional} emphasis.
\end{em}

It is sometimes necessary to prevent \LaTeX\ from
breaking a line where it might otherwise do so.
This may be at a space, as between the ``Mr.''\ and
``Jones'' in
       ``Mr.~Jones'',        % ~ produces an unbreakable interword space.
or within a word---especially when the word is a
symbol like
       \mbox{\emph{itemnum}}
that makes little sense when hyphenated across
lines.

Footnotes\footnote{This is an example of a footnote.}
pose no problem.

\LaTeX\ is good at typesetting mathematical formulas
like
       \( x-3y + z = 7 \)
or
       \( a_{1} > x^{2n} + y^{2n} > x' \)
or
       \( \ip{A}{B} = \sum_{i} a_{i} b_{i} \).
The spaces you type in a formula are
ignored.  Remember that a letter like
       $x$                   % $ ... $  and  \( ... \)  are equivalent
is a formula when it denotes a mathematical
symbol, and it should be typed as one.

\section{Displayed Text}

Text is displayed by indenting it from the left
margin.  Quotations are commonly displayed.  There
are short quotations
\begin{quote}
   This is a short quotation.  It consists of a
   single paragraph of text.  See how it is formatted.
\end{quote}
and longer ones.
\begin{quotation}
   This is a longer quotation.  It consists of two
   paragraphs of text, neither of which are
   particularly interesting.

   This is the second paragraph of the quotation.  It
   is just as dull as the first paragraph.
\end{quotation}
Another frequently-displayed structure is a list.
The following is an example of an \emph{itemized}
list.
\begin{itemize}
   \item This is the first item of an itemized list.
         Each item in the list is marked with a ``tick''.
         You don't have to worry about what kind of tick
         mark is used.

   \item This is the second item of the list.  It
         contains another list nested inside it.  The inner
         list is an \emph{enumerated} list.
         \begin{enumerate}
            \item This is the first item of an enumerated
                  list that is nested within the itemized list.

            \item This is the second item of the inner list.
                  \LaTeX\ allows you to nest lists deeper than
                  you really should.
         \end{enumerate}
         This is the rest of the second item of the outer
         list.  It is no more interesting than any other
         part of the item.
   \item This is the third item of the list.
\end{itemize}
You can even display poetry.
\begin{verse}
   There is an environment
    for verse \\             % The \\ command separates lines
   Whose features some poets % within a stanza.
   will curse.

                             % One or more blank lines separate stanzas.

   For instead of making\\
   Them do \emph{all} line breaking, \\
   It allows them to put too many words on a line when they'd rather be
   forced to be terse.
\end{verse}

Mathematical formulas may also be displayed.  A
displayed formula
is
one-line long; multiline
formulas require special formatting instructions.
   \[  \ip{\Gamma}{\psi'} = x'' + y^{2} + z_{i}^{n}\]
Don't start a paragraph with a displayed equation,
nor make one a paragraph by itself.

\end{document}               % End of document.

\section{Bench Section 123}
% synthetic bench section 123
% This is a sample LaTeX input file.  (Version of 12 August 2004.)
%
% A '%' character causes TeX to ignore all remaining text on the line,
% and is used for comments like this one.

\documentclass{article}      % Specifies the document class

                             % The preamble begins here.
\title{An Example Document}  % Declares the document's title.
\author{Leslie Lamport}      % Declares the author's name.
\date{January 21, 1994}      % Deleting this command produces today's date.

\newcommand{\ip}[2]{(#1, #2)}
                             % Defines \ip{arg1}{arg2} to mean
                             % (arg1, arg2).

%\newcommand{\ip}[2]{\langle #1 | #2\rangle}
                             % This is an alternative definition of
                             % \ip that is commented out.

\begin{document}             % End of preamble and beginning of text.

\maketitle                   % Produces the title.

This is an example input file.  Comparing it with
the output it generates can show you how to
produce a simple document of your own.

\section{Ordinary Text}      % Produces section heading.  Lower-level
                             % sections are begun with similar
                             % \subsection and \subsubsection commands.

The ends  of words and sentences are marked
  by   spaces. It  doesn't matter how many
spaces    you type; one is as good as 100.  The
end of   a line counts as a space.

One   or more   blank lines denote the  end
of  a paragraph.

Since any number of consecutive spaces are treated
like a single one, the formatting of the input
file makes no difference to
      \LaTeX,                % The \LaTeX command generates the LaTeX logo.
but it makes a difference to you.  When you use
\LaTeX, making your input file as easy to read
as possible will be a great help as you write
your document and when you change it.  This sample
file shows how you can add comments to your own input
file.

Because printing is different from typewriting,
there are a number of things that you have to do
differently when preparing an input file than if
you were just typing the document directly.
Quotation marks like
       ``this''
have to be handled specially, as do quotes within
quotes:
       ``\,`this'            % \, separates the double and single quote.
        is what I just
        wrote, not  `that'\,''.

Dashes come in three sizes: an
       intra-word
dash, a medium dash for number ranges like
       1--2,
and a punctuation
       dash---like
this.

A sentence-ending space should be larger than the
space between words within a sentence.  You
sometimes have to type special commands in
conjunction with punctuation characters to get
this right, as in the following sentence.
       Gnats, gnus, etc.\ all  % `\ ' makes an inter-word space.
       begin with G\@.         % \@ marks end-of-sentence punctuation.
You should check the spaces after periods when
reading your output to make sure you haven't
forgotten any special cases.  Generating an
ellipsis
       \ldots\               % `\ ' is needed after `\ldots' because TeX
                             % ignores spaces after command names like \ldots
                             % made from \ + letters.
                             %
                             % Note how a `%' character causes TeX to ignore
                             % the end of the input line, so these blank lines
                             % do not start a new paragraph.
                             %
with the right spacing around the periods requires
a special command.

\LaTeX\ interprets some common characters as
commands, so you must type special commands to
generate them.  These characters include the
following:
       \$ \& \% \# \{ and \}.

In printing, text is usually emphasized with an
       \emph{italic}
type style.

\begin{em}
   A long segment of text can also be emphasized
   in this way.  Text within such a segment can be
   given \emph{additional} emphasis.
\end{em}

It is sometimes necessary to prevent \LaTeX\ from
breaking a line where it might otherwise do so.
This may be at a space, as between the ``Mr.''\ and
``Jones'' in
       ``Mr.~Jones'',        % ~ produces an unbreakable interword space.
or within a word---especially when the word is a
symbol like
       \mbox{\emph{itemnum}}
that makes little sense when hyphenated across
lines.

Footnotes\footnote{This is an example of a footnote.}
pose no problem.

\LaTeX\ is good at typesetting mathematical formulas
like
       \( x-3y + z = 7 \)
or
       \( a_{1} > x^{2n} + y^{2n} > x' \)
or
       \( \ip{A}{B} = \sum_{i} a_{i} b_{i} \).
The spaces you type in a formula are
ignored.  Remember that a letter like
       $x$                   % $ ... $  and  \( ... \)  are equivalent
is a formula when it denotes a mathematical
symbol, and it should be typed as one.

\section{Displayed Text}

Text is displayed by indenting it from the left
margin.  Quotations are commonly displayed.  There
are short quotations
\begin{quote}
   This is a short quotation.  It consists of a
   single paragraph of text.  See how it is formatted.
\end{quote}
and longer ones.
\begin{quotation}
   This is a longer quotation.  It consists of two
   paragraphs of text, neither of which are
   particularly interesting.

   This is the second paragraph of the quotation.  It
   is just as dull as the first paragraph.
\end{quotation}
Another frequently-displayed structure is a list.
The following is an example of an \emph{itemized}
list.
\begin{itemize}
   \item This is the first item of an itemized list.
         Each item in the list is marked with a ``tick''.
         You don't have to worry about what kind of tick
         mark is used.

   \item This is the second item of the list.  It
         contains another list nested inside it.  The inner
         list is an \emph{enumerated} list.
         \begin{enumerate}
            \item This is the first item of an enumerated
                  list that is nested within the itemized list.

            \item This is the second item of the inner list.
                  \LaTeX\ allows you to nest lists deeper than
                  you really should.
         \end{enumerate}
         This is the rest of the second item of the outer
         list.  It is no more interesting than any other
         part of the item.
   \item This is the third item of the list.
\end{itemize}
You can even display poetry.
\begin{verse}
   There is an environment
    for verse \\             % The \\ command separates lines
   Whose features some poets % within a stanza.
   will curse.

                             % One or more blank lines separate stanzas.

   For instead of making\\
   Them do \emph{all} line breaking, \\
   It allows them to put too many words on a line when they'd rather be
   forced to be terse.
\end{verse}

Mathematical formulas may also be displayed.  A
displayed formula
is
one-line long; multiline
formulas require special formatting instructions.
   \[  \ip{\Gamma}{\psi'} = x'' + y^{2} + z_{i}^{n}\]
Don't start a paragraph with a displayed equation,
nor make one a paragraph by itself.

\end{document}               % End of document.

\section{Bench Section 124}
% synthetic bench section 124
% This is a sample LaTeX input file.  (Version of 12 August 2004.)
%
% A '%' character causes TeX to ignore all remaining text on the line,
% and is used for comments like this one.

\documentclass{article}      % Specifies the document class

                             % The preamble begins here.
\title{An Example Document}  % Declares the document's title.
\author{Leslie Lamport}      % Declares the author's name.
\date{January 21, 1994}      % Deleting this command produces today's date.

\newcommand{\ip}[2]{(#1, #2)}
                             % Defines \ip{arg1}{arg2} to mean
                             % (arg1, arg2).

%\newcommand{\ip}[2]{\langle #1 | #2\rangle}
                             % This is an alternative definition of
                             % \ip that is commented out.

\begin{document}             % End of preamble and beginning of text.

\maketitle                   % Produces the title.

This is an example input file.  Comparing it with
the output it generates can show you how to
produce a simple document of your own.

\section{Ordinary Text}      % Produces section heading.  Lower-level
                             % sections are begun with similar
                             % \subsection and \subsubsection commands.

The ends  of words and sentences are marked
  by   spaces. It  doesn't matter how many
spaces    you type; one is as good as 100.  The
end of   a line counts as a space.

One   or more   blank lines denote the  end
of  a paragraph.

Since any number of consecutive spaces are treated
like a single one, the formatting of the input
file makes no difference to
      \LaTeX,                % The \LaTeX command generates the LaTeX logo.
but it makes a difference to you.  When you use
\LaTeX, making your input file as easy to read
as possible will be a great help as you write
your document and when you change it.  This sample
file shows how you can add comments to your own input
file.

Because printing is different from typewriting,
there are a number of things that you have to do
differently when preparing an input file than if
you were just typing the document directly.
Quotation marks like
       ``this''
have to be handled specially, as do quotes within
quotes:
       ``\,`this'            % \, separates the double and single quote.
        is what I just
        wrote, not  `that'\,''.

Dashes come in three sizes: an
       intra-word
dash, a medium dash for number ranges like
       1--2,
and a punctuation
       dash---like
this.

A sentence-ending space should be larger than the
space between words within a sentence.  You
sometimes have to type special commands in
conjunction with punctuation characters to get
this right, as in the following sentence.
       Gnats, gnus, etc.\ all  % `\ ' makes an inter-word space.
       begin with G\@.         % \@ marks end-of-sentence punctuation.
You should check the spaces after periods when
reading your output to make sure you haven't
forgotten any special cases.  Generating an
ellipsis
       \ldots\               % `\ ' is needed after `\ldots' because TeX
                             % ignores spaces after command names like \ldots
                             % made from \ + letters.
                             %
                             % Note how a `%' character causes TeX to ignore
                             % the end of the input line, so these blank lines
                             % do not start a new paragraph.
                             %
with the right spacing around the periods requires
a special command.

\LaTeX\ interprets some common characters as
commands, so you must type special commands to
generate them.  These characters include the
following:
       \$ \& \% \# \{ and \}.

In printing, text is usually emphasized with an
       \emph{italic}
type style.

\begin{em}
   A long segment of text can also be emphasized
   in this way.  Text within such a segment can be
   given \emph{additional} emphasis.
\end{em}

It is sometimes necessary to prevent \LaTeX\ from
breaking a line where it might otherwise do so.
This may be at a space, as between the ``Mr.''\ and
``Jones'' in
       ``Mr.~Jones'',        % ~ produces an unbreakable interword space.
or within a word---especially when the word is a
symbol like
       \mbox{\emph{itemnum}}
that makes little sense when hyphenated across
lines.

Footnotes\footnote{This is an example of a footnote.}
pose no problem.

\LaTeX\ is good at typesetting mathematical formulas
like
       \( x-3y + z = 7 \)
or
       \( a_{1} > x^{2n} + y^{2n} > x' \)
or
       \( \ip{A}{B} = \sum_{i} a_{i} b_{i} \).
The spaces you type in a formula are
ignored.  Remember that a letter like
       $x$                   % $ ... $  and  \( ... \)  are equivalent
is a formula when it denotes a mathematical
symbol, and it should be typed as one.

\section{Displayed Text}

Text is displayed by indenting it from the left
margin.  Quotations are commonly displayed.  There
are short quotations
\begin{quote}
   This is a short quotation.  It consists of a
   single paragraph of text.  See how it is formatted.
\end{quote}
and longer ones.
\begin{quotation}
   This is a longer quotation.  It consists of two
   paragraphs of text, neither of which are
   particularly interesting.

   This is the second paragraph of the quotation.  It
   is just as dull as the first paragraph.
\end{quotation}
Another frequently-displayed structure is a list.
The following is an example of an \emph{itemized}
list.
\begin{itemize}
   \item This is the first item of an itemized list.
         Each item in the list is marked with a ``tick''.
         You don't have to worry about what kind of tick
         mark is used.

   \item This is the second item of the list.  It
         contains another list nested inside it.  The inner
         list is an \emph{enumerated} list.
         \begin{enumerate}
            \item This is the first item of an enumerated
                  list that is nested within the itemized list.

            \item This is the second item of the inner list.
                  \LaTeX\ allows you to nest lists deeper than
                  you really should.
         \end{enumerate}
         This is the rest of the second item of the outer
         list.  It is no more interesting than any other
         part of the item.
   \item This is the third item of the list.
\end{itemize}
You can even display poetry.
\begin{verse}
   There is an environment
    for verse \\             % The \\ command separates lines
   Whose features some poets % within a stanza.
   will curse.

                             % One or more blank lines separate stanzas.

   For instead of making\\
   Them do \emph{all} line breaking, \\
   It allows them to put too many words on a line when they'd rather be
   forced to be terse.
\end{verse}

Mathematical formulas may also be displayed.  A
displayed formula
is
one-line long; multiline
formulas require special formatting instructions.
   \[  \ip{\Gamma}{\psi'} = x'' + y^{2} + z_{i}^{n}\]
Don't start a paragraph with a displayed equation,
nor make one a paragraph by itself.

\end{document}               % End of document.

\section{Bench Section 125}
% synthetic bench section 125
% This is a sample LaTeX input file.  (Version of 12 August 2004.)
%
% A '%' character causes TeX to ignore all remaining text on the line,
% and is used for comments like this one.

\documentclass{article}      % Specifies the document class

                             % The preamble begins here.
\title{An Example Document}  % Declares the document's title.
\author{Leslie Lamport}      % Declares the author's name.
\date{January 21, 1994}      % Deleting this command produces today's date.

\newcommand{\ip}[2]{(#1, #2)}
                             % Defines \ip{arg1}{arg2} to mean
                             % (arg1, arg2).

%\newcommand{\ip}[2]{\langle #1 | #2\rangle}
                             % This is an alternative definition of
                             % \ip that is commented out.

\begin{document}             % End of preamble and beginning of text.

\maketitle                   % Produces the title.

This is an example input file.  Comparing it with
the output it generates can show you how to
produce a simple document of your own.

\section{Ordinary Text}      % Produces section heading.  Lower-level
                             % sections are begun with similar
                             % \subsection and \subsubsection commands.

The ends  of words and sentences are marked
  by   spaces. It  doesn't matter how many
spaces    you type; one is as good as 100.  The
end of   a line counts as a space.

One   or more   blank lines denote the  end
of  a paragraph.

Since any number of consecutive spaces are treated
like a single one, the formatting of the input
file makes no difference to
      \LaTeX,                % The \LaTeX command generates the LaTeX logo.
but it makes a difference to you.  When you use
\LaTeX, making your input file as easy to read
as possible will be a great help as you write
your document and when you change it.  This sample
file shows how you can add comments to your own input
file.

Because printing is different from typewriting,
there are a number of things that you have to do
differently when preparing an input file than if
you were just typing the document directly.
Quotation marks like
       ``this''
have to be handled specially, as do quotes within
quotes:
       ``\,`this'            % \, separates the double and single quote.
        is what I just
        wrote, not  `that'\,''.

Dashes come in three sizes: an
       intra-word
dash, a medium dash for number ranges like
       1--2,
and a punctuation
       dash---like
this.

A sentence-ending space should be larger than the
space between words within a sentence.  You
sometimes have to type special commands in
conjunction with punctuation characters to get
this right, as in the following sentence.
       Gnats, gnus, etc.\ all  % `\ ' makes an inter-word space.
       begin with G\@.         % \@ marks end-of-sentence punctuation.
You should check the spaces after periods when
reading your output to make sure you haven't
forgotten any special cases.  Generating an
ellipsis
       \ldots\               % `\ ' is needed after `\ldots' because TeX
                             % ignores spaces after command names like \ldots
                             % made from \ + letters.
                             %
                             % Note how a `%' character causes TeX to ignore
                             % the end of the input line, so these blank lines
                             % do not start a new paragraph.
                             %
with the right spacing around the periods requires
a special command.

\LaTeX\ interprets some common characters as
commands, so you must type special commands to
generate them.  These characters include the
following:
       \$ \& \% \# \{ and \}.

In printing, text is usually emphasized with an
       \emph{italic}
type style.

\begin{em}
   A long segment of text can also be emphasized
   in this way.  Text within such a segment can be
   given \emph{additional} emphasis.
\end{em}

It is sometimes necessary to prevent \LaTeX\ from
breaking a line where it might otherwise do so.
This may be at a space, as between the ``Mr.''\ and
``Jones'' in
       ``Mr.~Jones'',        % ~ produces an unbreakable interword space.
or within a word---especially when the word is a
symbol like
       \mbox{\emph{itemnum}}
that makes little sense when hyphenated across
lines.

Footnotes\footnote{This is an example of a footnote.}
pose no problem.

\LaTeX\ is good at typesetting mathematical formulas
like
       \( x-3y + z = 7 \)
or
       \( a_{1} > x^{2n} + y^{2n} > x' \)
or
       \( \ip{A}{B} = \sum_{i} a_{i} b_{i} \).
The spaces you type in a formula are
ignored.  Remember that a letter like
       $x$                   % $ ... $  and  \( ... \)  are equivalent
is a formula when it denotes a mathematical
symbol, and it should be typed as one.

\section{Displayed Text}

Text is displayed by indenting it from the left
margin.  Quotations are commonly displayed.  There
are short quotations
\begin{quote}
   This is a short quotation.  It consists of a
   single paragraph of text.  See how it is formatted.
\end{quote}
and longer ones.
\begin{quotation}
   This is a longer quotation.  It consists of two
   paragraphs of text, neither of which are
   particularly interesting.

   This is the second paragraph of the quotation.  It
   is just as dull as the first paragraph.
\end{quotation}
Another frequently-displayed structure is a list.
The following is an example of an \emph{itemized}
list.
\begin{itemize}
   \item This is the first item of an itemized list.
         Each item in the list is marked with a ``tick''.
         You don't have to worry about what kind of tick
         mark is used.

   \item This is the second item of the list.  It
         contains another list nested inside it.  The inner
         list is an \emph{enumerated} list.
         \begin{enumerate}
            \item This is the first item of an enumerated
                  list that is nested within the itemized list.

            \item This is the second item of the inner list.
                  \LaTeX\ allows you to nest lists deeper than
                  you really should.
         \end{enumerate}
         This is the rest of the second item of the outer
         list.  It is no more interesting than any other
         part of the item.
   \item This is the third item of the list.
\end{itemize}
You can even display poetry.
\begin{verse}
   There is an environment
    for verse \\             % The \\ command separates lines
   Whose features some poets % within a stanza.
   will curse.

                             % One or more blank lines separate stanzas.

   For instead of making\\
   Them do \emph{all} line breaking, \\
   It allows them to put too many words on a line when they'd rather be
   forced to be terse.
\end{verse}

Mathematical formulas may also be displayed.  A
displayed formula
is
one-line long; multiline
formulas require special formatting instructions.
   \[  \ip{\Gamma}{\psi'} = x'' + y^{2} + z_{i}^{n}\]
Don't start a paragraph with a displayed equation,
nor make one a paragraph by itself.

\end{document}               % End of document.

\section{Bench Section 126}
% synthetic bench section 126
% This is a sample LaTeX input file.  (Version of 12 August 2004.)
%
% A '%' character causes TeX to ignore all remaining text on the line,
% and is used for comments like this one.

\documentclass{article}      % Specifies the document class

                             % The preamble begins here.
\title{An Example Document}  % Declares the document's title.
\author{Leslie Lamport}      % Declares the author's name.
\date{January 21, 1994}      % Deleting this command produces today's date.

\newcommand{\ip}[2]{(#1, #2)}
                             % Defines \ip{arg1}{arg2} to mean
                             % (arg1, arg2).

%\newcommand{\ip}[2]{\langle #1 | #2\rangle}
                             % This is an alternative definition of
                             % \ip that is commented out.

\begin{document}             % End of preamble and beginning of text.

\maketitle                   % Produces the title.

This is an example input file.  Comparing it with
the output it generates can show you how to
produce a simple document of your own.

\section{Ordinary Text}      % Produces section heading.  Lower-level
                             % sections are begun with similar
                             % \subsection and \subsubsection commands.

The ends  of words and sentences are marked
  by   spaces. It  doesn't matter how many
spaces    you type; one is as good as 100.  The
end of   a line counts as a space.

One   or more   blank lines denote the  end
of  a paragraph.

Since any number of consecutive spaces are treated
like a single one, the formatting of the input
file makes no difference to
      \LaTeX,                % The \LaTeX command generates the LaTeX logo.
but it makes a difference to you.  When you use
\LaTeX, making your input file as easy to read
as possible will be a great help as you write
your document and when you change it.  This sample
file shows how you can add comments to your own input
file.

Because printing is different from typewriting,
there are a number of things that you have to do
differently when preparing an input file than if
you were just typing the document directly.
Quotation marks like
       ``this''
have to be handled specially, as do quotes within
quotes:
       ``\,`this'            % \, separates the double and single quote.
        is what I just
        wrote, not  `that'\,''.

Dashes come in three sizes: an
       intra-word
dash, a medium dash for number ranges like
       1--2,
and a punctuation
       dash---like
this.

A sentence-ending space should be larger than the
space between words within a sentence.  You
sometimes have to type special commands in
conjunction with punctuation characters to get
this right, as in the following sentence.
       Gnats, gnus, etc.\ all  % `\ ' makes an inter-word space.
       begin with G\@.         % \@ marks end-of-sentence punctuation.
You should check the spaces after periods when
reading your output to make sure you haven't
forgotten any special cases.  Generating an
ellipsis
       \ldots\               % `\ ' is needed after `\ldots' because TeX
                             % ignores spaces after command names like \ldots
                             % made from \ + letters.
                             %
                             % Note how a `%' character causes TeX to ignore
                             % the end of the input line, so these blank lines
                             % do not start a new paragraph.
                             %
with the right spacing around the periods requires
a special command.

\LaTeX\ interprets some common characters as
commands, so you must type special commands to
generate them.  These characters include the
following:
       \$ \& \% \# \{ and \}.

In printing, text is usually emphasized with an
       \emph{italic}
type style.

\begin{em}
   A long segment of text can also be emphasized
   in this way.  Text within such a segment can be
   given \emph{additional} emphasis.
\end{em}

It is sometimes necessary to prevent \LaTeX\ from
breaking a line where it might otherwise do so.
This may be at a space, as between the ``Mr.''\ and
``Jones'' in
       ``Mr.~Jones'',        % ~ produces an unbreakable interword space.
or within a word---especially when the word is a
symbol like
       \mbox{\emph{itemnum}}
that makes little sense when hyphenated across
lines.

Footnotes\footnote{This is an example of a footnote.}
pose no problem.

\LaTeX\ is good at typesetting mathematical formulas
like
       \( x-3y + z = 7 \)
or
       \( a_{1} > x^{2n} + y^{2n} > x' \)
or
       \( \ip{A}{B} = \sum_{i} a_{i} b_{i} \).
The spaces you type in a formula are
ignored.  Remember that a letter like
       $x$                   % $ ... $  and  \( ... \)  are equivalent
is a formula when it denotes a mathematical
symbol, and it should be typed as one.

\section{Displayed Text}

Text is displayed by indenting it from the left
margin.  Quotations are commonly displayed.  There
are short quotations
\begin{quote}
   This is a short quotation.  It consists of a
   single paragraph of text.  See how it is formatted.
\end{quote}
and longer ones.
\begin{quotation}
   This is a longer quotation.  It consists of two
   paragraphs of text, neither of which are
   particularly interesting.

   This is the second paragraph of the quotation.  It
   is just as dull as the first paragraph.
\end{quotation}
Another frequently-displayed structure is a list.
The following is an example of an \emph{itemized}
list.
\begin{itemize}
   \item This is the first item of an itemized list.
         Each item in the list is marked with a ``tick''.
         You don't have to worry about what kind of tick
         mark is used.

   \item This is the second item of the list.  It
         contains another list nested inside it.  The inner
         list is an \emph{enumerated} list.
         \begin{enumerate}
            \item This is the first item of an enumerated
                  list that is nested within the itemized list.

            \item This is the second item of the inner list.
                  \LaTeX\ allows you to nest lists deeper than
                  you really should.
         \end{enumerate}
         This is the rest of the second item of the outer
         list.  It is no more interesting than any other
         part of the item.
   \item This is the third item of the list.
\end{itemize}
You can even display poetry.
\begin{verse}
   There is an environment
    for verse \\             % The \\ command separates lines
   Whose features some poets % within a stanza.
   will curse.

                             % One or more blank lines separate stanzas.

   For instead of making\\
   Them do \emph{all} line breaking, \\
   It allows them to put too many words on a line when they'd rather be
   forced to be terse.
\end{verse}

Mathematical formulas may also be displayed.  A
displayed formula
is
one-line long; multiline
formulas require special formatting instructions.
   \[  \ip{\Gamma}{\psi'} = x'' + y^{2} + z_{i}^{n}\]
Don't start a paragraph with a displayed equation,
nor make one a paragraph by itself.

\end{document}               % End of document.

\section{Bench Section 127}
% synthetic bench section 127
% This is a sample LaTeX input file.  (Version of 12 August 2004.)
%
% A '%' character causes TeX to ignore all remaining text on the line,
% and is used for comments like this one.

\documentclass{article}      % Specifies the document class

                             % The preamble begins here.
\title{An Example Document}  % Declares the document's title.
\author{Leslie Lamport}      % Declares the author's name.
\date{January 21, 1994}      % Deleting this command produces today's date.

\newcommand{\ip}[2]{(#1, #2)}
                             % Defines \ip{arg1}{arg2} to mean
                             % (arg1, arg2).

%\newcommand{\ip}[2]{\langle #1 | #2\rangle}
                             % This is an alternative definition of
                             % \ip that is commented out.

\begin{document}             % End of preamble and beginning of text.

\maketitle                   % Produces the title.

This is an example input file.  Comparing it with
the output it generates can show you how to
produce a simple document of your own.

\section{Ordinary Text}      % Produces section heading.  Lower-level
                             % sections are begun with similar
                             % \subsection and \subsubsection commands.

The ends  of words and sentences are marked
  by   spaces. It  doesn't matter how many
spaces    you type; one is as good as 100.  The
end of   a line counts as a space.

One   or more   blank lines denote the  end
of  a paragraph.

Since any number of consecutive spaces are treated
like a single one, the formatting of the input
file makes no difference to
      \LaTeX,                % The \LaTeX command generates the LaTeX logo.
but it makes a difference to you.  When you use
\LaTeX, making your input file as easy to read
as possible will be a great help as you write
your document and when you change it.  This sample
file shows how you can add comments to your own input
file.

Because printing is different from typewriting,
there are a number of things that you have to do
differently when preparing an input file than if
you were just typing the document directly.
Quotation marks like
       ``this''
have to be handled specially, as do quotes within
quotes:
       ``\,`this'            % \, separates the double and single quote.
        is what I just
        wrote, not  `that'\,''.

Dashes come in three sizes: an
       intra-word
dash, a medium dash for number ranges like
       1--2,
and a punctuation
       dash---like
this.

A sentence-ending space should be larger than the
space between words within a sentence.  You
sometimes have to type special commands in
conjunction with punctuation characters to get
this right, as in the following sentence.
       Gnats, gnus, etc.\ all  % `\ ' makes an inter-word space.
       begin with G\@.         % \@ marks end-of-sentence punctuation.
You should check the spaces after periods when
reading your output to make sure you haven't
forgotten any special cases.  Generating an
ellipsis
       \ldots\               % `\ ' is needed after `\ldots' because TeX
                             % ignores spaces after command names like \ldots
                             % made from \ + letters.
                             %
                             % Note how a `%' character causes TeX to ignore
                             % the end of the input line, so these blank lines
                             % do not start a new paragraph.
                             %
with the right spacing around the periods requires
a special command.

\LaTeX\ interprets some common characters as
commands, so you must type special commands to
generate them.  These characters include the
following:
       \$ \& \% \# \{ and \}.

In printing, text is usually emphasized with an
       \emph{italic}
type style.

\begin{em}
   A long segment of text can also be emphasized
   in this way.  Text within such a segment can be
   given \emph{additional} emphasis.
\end{em}

It is sometimes necessary to prevent \LaTeX\ from
breaking a line where it might otherwise do so.
This may be at a space, as between the ``Mr.''\ and
``Jones'' in
       ``Mr.~Jones'',        % ~ produces an unbreakable interword space.
or within a word---especially when the word is a
symbol like
       \mbox{\emph{itemnum}}
that makes little sense when hyphenated across
lines.

Footnotes\footnote{This is an example of a footnote.}
pose no problem.

\LaTeX\ is good at typesetting mathematical formulas
like
       \( x-3y + z = 7 \)
or
       \( a_{1} > x^{2n} + y^{2n} > x' \)
or
       \( \ip{A}{B} = \sum_{i} a_{i} b_{i} \).
The spaces you type in a formula are
ignored.  Remember that a letter like
       $x$                   % $ ... $  and  \( ... \)  are equivalent
is a formula when it denotes a mathematical
symbol, and it should be typed as one.

\section{Displayed Text}

Text is displayed by indenting it from the left
margin.  Quotations are commonly displayed.  There
are short quotations
\begin{quote}
   This is a short quotation.  It consists of a
   single paragraph of text.  See how it is formatted.
\end{quote}
and longer ones.
\begin{quotation}
   This is a longer quotation.  It consists of two
   paragraphs of text, neither of which are
   particularly interesting.

   This is the second paragraph of the quotation.  It
   is just as dull as the first paragraph.
\end{quotation}
Another frequently-displayed structure is a list.
The following is an example of an \emph{itemized}
list.
\begin{itemize}
   \item This is the first item of an itemized list.
         Each item in the list is marked with a ``tick''.
         You don't have to worry about what kind of tick
         mark is used.

   \item This is the second item of the list.  It
         contains another list nested inside it.  The inner
         list is an \emph{enumerated} list.
         \begin{enumerate}
            \item This is the first item of an enumerated
                  list that is nested within the itemized list.

            \item This is the second item of the inner list.
                  \LaTeX\ allows you to nest lists deeper than
                  you really should.
         \end{enumerate}
         This is the rest of the second item of the outer
         list.  It is no more interesting than any other
         part of the item.
   \item This is the third item of the list.
\end{itemize}
You can even display poetry.
\begin{verse}
   There is an environment
    for verse \\             % The \\ command separates lines
   Whose features some poets % within a stanza.
   will curse.

                             % One or more blank lines separate stanzas.

   For instead of making\\
   Them do \emph{all} line breaking, \\
   It allows them to put too many words on a line when they'd rather be
   forced to be terse.
\end{verse}

Mathematical formulas may also be displayed.  A
displayed formula
is
one-line long; multiline
formulas require special formatting instructions.
   \[  \ip{\Gamma}{\psi'} = x'' + y^{2} + z_{i}^{n}\]
Don't start a paragraph with a displayed equation,
nor make one a paragraph by itself.

\end{document}               % End of document.

\section{Bench Section 128}
% synthetic bench section 128
% This is a sample LaTeX input file.  (Version of 12 August 2004.)
%
% A '%' character causes TeX to ignore all remaining text on the line,
% and is used for comments like this one.

\documentclass{article}      % Specifies the document class

                             % The preamble begins here.
\title{An Example Document}  % Declares the document's title.
\author{Leslie Lamport}      % Declares the author's name.
\date{January 21, 1994}      % Deleting this command produces today's date.

\newcommand{\ip}[2]{(#1, #2)}
                             % Defines \ip{arg1}{arg2} to mean
                             % (arg1, arg2).

%\newcommand{\ip}[2]{\langle #1 | #2\rangle}
                             % This is an alternative definition of
                             % \ip that is commented out.

\begin{document}             % End of preamble and beginning of text.

\maketitle                   % Produces the title.

This is an example input file.  Comparing it with
the output it generates can show you how to
produce a simple document of your own.

\section{Ordinary Text}      % Produces section heading.  Lower-level
                             % sections are begun with similar
                             % \subsection and \subsubsection commands.

The ends  of words and sentences are marked
  by   spaces. It  doesn't matter how many
spaces    you type; one is as good as 100.  The
end of   a line counts as a space.

One   or more   blank lines denote the  end
of  a paragraph.

Since any number of consecutive spaces are treated
like a single one, the formatting of the input
file makes no difference to
      \LaTeX,                % The \LaTeX command generates the LaTeX logo.
but it makes a difference to you.  When you use
\LaTeX, making your input file as easy to read
as possible will be a great help as you write
your document and when you change it.  This sample
file shows how you can add comments to your own input
file.

Because printing is different from typewriting,
there are a number of things that you have to do
differently when preparing an input file than if
you were just typing the document directly.
Quotation marks like
       ``this''
have to be handled specially, as do quotes within
quotes:
       ``\,`this'            % \, separates the double and single quote.
        is what I just
        wrote, not  `that'\,''.

Dashes come in three sizes: an
       intra-word
dash, a medium dash for number ranges like
       1--2,
and a punctuation
       dash---like
this.

A sentence-ending space should be larger than the
space between words within a sentence.  You
sometimes have to type special commands in
conjunction with punctuation characters to get
this right, as in the following sentence.
       Gnats, gnus, etc.\ all  % `\ ' makes an inter-word space.
       begin with G\@.         % \@ marks end-of-sentence punctuation.
You should check the spaces after periods when
reading your output to make sure you haven't
forgotten any special cases.  Generating an
ellipsis
       \ldots\               % `\ ' is needed after `\ldots' because TeX
                             % ignores spaces after command names like \ldots
                             % made from \ + letters.
                             %
                             % Note how a `%' character causes TeX to ignore
                             % the end of the input line, so these blank lines
                             % do not start a new paragraph.
                             %
with the right spacing around the periods requires
a special command.

\LaTeX\ interprets some common characters as
commands, so you must type special commands to
generate them.  These characters include the
following:
       \$ \& \% \# \{ and \}.

In printing, text is usually emphasized with an
       \emph{italic}
type style.

\begin{em}
   A long segment of text can also be emphasized
   in this way.  Text within such a segment can be
   given \emph{additional} emphasis.
\end{em}

It is sometimes necessary to prevent \LaTeX\ from
breaking a line where it might otherwise do so.
This may be at a space, as between the ``Mr.''\ and
``Jones'' in
       ``Mr.~Jones'',        % ~ produces an unbreakable interword space.
or within a word---especially when the word is a
symbol like
       \mbox{\emph{itemnum}}
that makes little sense when hyphenated across
lines.

Footnotes\footnote{This is an example of a footnote.}
pose no problem.

\LaTeX\ is good at typesetting mathematical formulas
like
       \( x-3y + z = 7 \)
or
       \( a_{1} > x^{2n} + y^{2n} > x' \)
or
       \( \ip{A}{B} = \sum_{i} a_{i} b_{i} \).
The spaces you type in a formula are
ignored.  Remember that a letter like
       $x$                   % $ ... $  and  \( ... \)  are equivalent
is a formula when it denotes a mathematical
symbol, and it should be typed as one.

\section{Displayed Text}

Text is displayed by indenting it from the left
margin.  Quotations are commonly displayed.  There
are short quotations
\begin{quote}
   This is a short quotation.  It consists of a
   single paragraph of text.  See how it is formatted.
\end{quote}
and longer ones.
\begin{quotation}
   This is a longer quotation.  It consists of two
   paragraphs of text, neither of which are
   particularly interesting.

   This is the second paragraph of the quotation.  It
   is just as dull as the first paragraph.
\end{quotation}
Another frequently-displayed structure is a list.
The following is an example of an \emph{itemized}
list.
\begin{itemize}
   \item This is the first item of an itemized list.
         Each item in the list is marked with a ``tick''.
         You don't have to worry about what kind of tick
         mark is used.

   \item This is the second item of the list.  It
         contains another list nested inside it.  The inner
         list is an \emph{enumerated} list.
         \begin{enumerate}
            \item This is the first item of an enumerated
                  list that is nested within the itemized list.

            \item This is the second item of the inner list.
                  \LaTeX\ allows you to nest lists deeper than
                  you really should.
         \end{enumerate}
         This is the rest of the second item of the outer
         list.  It is no more interesting than any other
         part of the item.
   \item This is the third item of the list.
\end{itemize}
You can even display poetry.
\begin{verse}
   There is an environment
    for verse \\             % The \\ command separates lines
   Whose features some poets % within a stanza.
   will curse.

                             % One or more blank lines separate stanzas.

   For instead of making\\
   Them do \emph{all} line breaking, \\
   It allows them to put too many words on a line when they'd rather be
   forced to be terse.
\end{verse}

Mathematical formulas may also be displayed.  A
displayed formula
is
one-line long; multiline
formulas require special formatting instructions.
   \[  \ip{\Gamma}{\psi'} = x'' + y^{2} + z_{i}^{n}\]
Don't start a paragraph with a displayed equation,
nor make one a paragraph by itself.

\end{document}               % End of document.

\section{Bench Section 129}
% synthetic bench section 129
% This is a sample LaTeX input file.  (Version of 12 August 2004.)
%
% A '%' character causes TeX to ignore all remaining text on the line,
% and is used for comments like this one.

\documentclass{article}      % Specifies the document class

                             % The preamble begins here.
\title{An Example Document}  % Declares the document's title.
\author{Leslie Lamport}      % Declares the author's name.
\date{January 21, 1994}      % Deleting this command produces today's date.

\newcommand{\ip}[2]{(#1, #2)}
                             % Defines \ip{arg1}{arg2} to mean
                             % (arg1, arg2).

%\newcommand{\ip}[2]{\langle #1 | #2\rangle}
                             % This is an alternative definition of
                             % \ip that is commented out.

\begin{document}             % End of preamble and beginning of text.

\maketitle                   % Produces the title.

This is an example input file.  Comparing it with
the output it generates can show you how to
produce a simple document of your own.

\section{Ordinary Text}      % Produces section heading.  Lower-level
                             % sections are begun with similar
                             % \subsection and \subsubsection commands.

The ends  of words and sentences are marked
  by   spaces. It  doesn't matter how many
spaces    you type; one is as good as 100.  The
end of   a line counts as a space.

One   or more   blank lines denote the  end
of  a paragraph.

Since any number of consecutive spaces are treated
like a single one, the formatting of the input
file makes no difference to
      \LaTeX,                % The \LaTeX command generates the LaTeX logo.
but it makes a difference to you.  When you use
\LaTeX, making your input file as easy to read
as possible will be a great help as you write
your document and when you change it.  This sample
file shows how you can add comments to your own input
file.

Because printing is different from typewriting,
there are a number of things that you have to do
differently when preparing an input file than if
you were just typing the document directly.
Quotation marks like
       ``this''
have to be handled specially, as do quotes within
quotes:
       ``\,`this'            % \, separates the double and single quote.
        is what I just
        wrote, not  `that'\,''.

Dashes come in three sizes: an
       intra-word
dash, a medium dash for number ranges like
       1--2,
and a punctuation
       dash---like
this.

A sentence-ending space should be larger than the
space between words within a sentence.  You
sometimes have to type special commands in
conjunction with punctuation characters to get
this right, as in the following sentence.
       Gnats, gnus, etc.\ all  % `\ ' makes an inter-word space.
       begin with G\@.         % \@ marks end-of-sentence punctuation.
You should check the spaces after periods when
reading your output to make sure you haven't
forgotten any special cases.  Generating an
ellipsis
       \ldots\               % `\ ' is needed after `\ldots' because TeX
                             % ignores spaces after command names like \ldots
                             % made from \ + letters.
                             %
                             % Note how a `%' character causes TeX to ignore
                             % the end of the input line, so these blank lines
                             % do not start a new paragraph.
                             %
with the right spacing around the periods requires
a special command.

\LaTeX\ interprets some common characters as
commands, so you must type special commands to
generate them.  These characters include the
following:
       \$ \& \% \# \{ and \}.

In printing, text is usually emphasized with an
       \emph{italic}
type style.

\begin{em}
   A long segment of text can also be emphasized
   in this way.  Text within such a segment can be
   given \emph{additional} emphasis.
\end{em}

It is sometimes necessary to prevent \LaTeX\ from
breaking a line where it might otherwise do so.
This may be at a space, as between the ``Mr.''\ and
``Jones'' in
       ``Mr.~Jones'',        % ~ produces an unbreakable interword space.
or within a word---especially when the word is a
symbol like
       \mbox{\emph{itemnum}}
that makes little sense when hyphenated across
lines.

Footnotes\footnote{This is an example of a footnote.}
pose no problem.

\LaTeX\ is good at typesetting mathematical formulas
like
       \( x-3y + z = 7 \)
or
       \( a_{1} > x^{2n} + y^{2n} > x' \)
or
       \( \ip{A}{B} = \sum_{i} a_{i} b_{i} \).
The spaces you type in a formula are
ignored.  Remember that a letter like
       $x$                   % $ ... $  and  \( ... \)  are equivalent
is a formula when it denotes a mathematical
symbol, and it should be typed as one.

\section{Displayed Text}

Text is displayed by indenting it from the left
margin.  Quotations are commonly displayed.  There
are short quotations
\begin{quote}
   This is a short quotation.  It consists of a
   single paragraph of text.  See how it is formatted.
\end{quote}
and longer ones.
\begin{quotation}
   This is a longer quotation.  It consists of two
   paragraphs of text, neither of which are
   particularly interesting.

   This is the second paragraph of the quotation.  It
   is just as dull as the first paragraph.
\end{quotation}
Another frequently-displayed structure is a list.
The following is an example of an \emph{itemized}
list.
\begin{itemize}
   \item This is the first item of an itemized list.
         Each item in the list is marked with a ``tick''.
         You don't have to worry about what kind of tick
         mark is used.

   \item This is the second item of the list.  It
         contains another list nested inside it.  The inner
         list is an \emph{enumerated} list.
         \begin{enumerate}
            \item This is the first item of an enumerated
                  list that is nested within the itemized list.

            \item This is the second item of the inner list.
                  \LaTeX\ allows you to nest lists deeper than
                  you really should.
         \end{enumerate}
         This is the rest of the second item of the outer
         list.  It is no more interesting than any other
         part of the item.
   \item This is the third item of the list.
\end{itemize}
You can even display poetry.
\begin{verse}
   There is an environment
    for verse \\             % The \\ command separates lines
   Whose features some poets % within a stanza.
   will curse.

                             % One or more blank lines separate stanzas.

   For instead of making\\
   Them do \emph{all} line breaking, \\
   It allows them to put too many words on a line when they'd rather be
   forced to be terse.
\end{verse}

Mathematical formulas may also be displayed.  A
displayed formula
is
one-line long; multiline
formulas require special formatting instructions.
   \[  \ip{\Gamma}{\psi'} = x'' + y^{2} + z_{i}^{n}\]
Don't start a paragraph with a displayed equation,
nor make one a paragraph by itself.

\end{document}               % End of document.

\section{Bench Section 130}
% synthetic bench section 130
% This is a sample LaTeX input file.  (Version of 12 August 2004.)
%
% A '%' character causes TeX to ignore all remaining text on the line,
% and is used for comments like this one.

\documentclass{article}      % Specifies the document class

                             % The preamble begins here.
\title{An Example Document}  % Declares the document's title.
\author{Leslie Lamport}      % Declares the author's name.
\date{January 21, 1994}      % Deleting this command produces today's date.

\newcommand{\ip}[2]{(#1, #2)}
                             % Defines \ip{arg1}{arg2} to mean
                             % (arg1, arg2).

%\newcommand{\ip}[2]{\langle #1 | #2\rangle}
                             % This is an alternative definition of
                             % \ip that is commented out.

\begin{document}             % End of preamble and beginning of text.

\maketitle                   % Produces the title.

This is an example input file.  Comparing it with
the output it generates can show you how to
produce a simple document of your own.

\section{Ordinary Text}      % Produces section heading.  Lower-level
                             % sections are begun with similar
                             % \subsection and \subsubsection commands.

The ends  of words and sentences are marked
  by   spaces. It  doesn't matter how many
spaces    you type; one is as good as 100.  The
end of   a line counts as a space.

One   or more   blank lines denote the  end
of  a paragraph.

Since any number of consecutive spaces are treated
like a single one, the formatting of the input
file makes no difference to
      \LaTeX,                % The \LaTeX command generates the LaTeX logo.
but it makes a difference to you.  When you use
\LaTeX, making your input file as easy to read
as possible will be a great help as you write
your document and when you change it.  This sample
file shows how you can add comments to your own input
file.

Because printing is different from typewriting,
there are a number of things that you have to do
differently when preparing an input file than if
you were just typing the document directly.
Quotation marks like
       ``this''
have to be handled specially, as do quotes within
quotes:
       ``\,`this'            % \, separates the double and single quote.
        is what I just
        wrote, not  `that'\,''.

Dashes come in three sizes: an
       intra-word
dash, a medium dash for number ranges like
       1--2,
and a punctuation
       dash---like
this.

A sentence-ending space should be larger than the
space between words within a sentence.  You
sometimes have to type special commands in
conjunction with punctuation characters to get
this right, as in the following sentence.
       Gnats, gnus, etc.\ all  % `\ ' makes an inter-word space.
       begin with G\@.         % \@ marks end-of-sentence punctuation.
You should check the spaces after periods when
reading your output to make sure you haven't
forgotten any special cases.  Generating an
ellipsis
       \ldots\               % `\ ' is needed after `\ldots' because TeX
                             % ignores spaces after command names like \ldots
                             % made from \ + letters.
                             %
                             % Note how a `%' character causes TeX to ignore
                             % the end of the input line, so these blank lines
                             % do not start a new paragraph.
                             %
with the right spacing around the periods requires
a special command.

\LaTeX\ interprets some common characters as
commands, so you must type special commands to
generate them.  These characters include the
following:
       \$ \& \% \# \{ and \}.

In printing, text is usually emphasized with an
       \emph{italic}
type style.

\begin{em}
   A long segment of text can also be emphasized
   in this way.  Text within such a segment can be
   given \emph{additional} emphasis.
\end{em}

It is sometimes necessary to prevent \LaTeX\ from
breaking a line where it might otherwise do so.
This may be at a space, as between the ``Mr.''\ and
``Jones'' in
       ``Mr.~Jones'',        % ~ produces an unbreakable interword space.
or within a word---especially when the word is a
symbol like
       \mbox{\emph{itemnum}}
that makes little sense when hyphenated across
lines.

Footnotes\footnote{This is an example of a footnote.}
pose no problem.

\LaTeX\ is good at typesetting mathematical formulas
like
       \( x-3y + z = 7 \)
or
       \( a_{1} > x^{2n} + y^{2n} > x' \)
or
       \( \ip{A}{B} = \sum_{i} a_{i} b_{i} \).
The spaces you type in a formula are
ignored.  Remember that a letter like
       $x$                   % $ ... $  and  \( ... \)  are equivalent
is a formula when it denotes a mathematical
symbol, and it should be typed as one.

\section{Displayed Text}

Text is displayed by indenting it from the left
margin.  Quotations are commonly displayed.  There
are short quotations
\begin{quote}
   This is a short quotation.  It consists of a
   single paragraph of text.  See how it is formatted.
\end{quote}
and longer ones.
\begin{quotation}
   This is a longer quotation.  It consists of two
   paragraphs of text, neither of which are
   particularly interesting.

   This is the second paragraph of the quotation.  It
   is just as dull as the first paragraph.
\end{quotation}
Another frequently-displayed structure is a list.
The following is an example of an \emph{itemized}
list.
\begin{itemize}
   \item This is the first item of an itemized list.
         Each item in the list is marked with a ``tick''.
         You don't have to worry about what kind of tick
         mark is used.

   \item This is the second item of the list.  It
         contains another list nested inside it.  The inner
         list is an \emph{enumerated} list.
         \begin{enumerate}
            \item This is the first item of an enumerated
                  list that is nested within the itemized list.

            \item This is the second item of the inner list.
                  \LaTeX\ allows you to nest lists deeper than
                  you really should.
         \end{enumerate}
         This is the rest of the second item of the outer
         list.  It is no more interesting than any other
         part of the item.
   \item This is the third item of the list.
\end{itemize}
You can even display poetry.
\begin{verse}
   There is an environment
    for verse \\             % The \\ command separates lines
   Whose features some poets % within a stanza.
   will curse.

                             % One or more blank lines separate stanzas.

   For instead of making\\
   Them do \emph{all} line breaking, \\
   It allows them to put too many words on a line when they'd rather be
   forced to be terse.
\end{verse}

Mathematical formulas may also be displayed.  A
displayed formula
is
one-line long; multiline
formulas require special formatting instructions.
   \[  \ip{\Gamma}{\psi'} = x'' + y^{2} + z_{i}^{n}\]
Don't start a paragraph with a displayed equation,
nor make one a paragraph by itself.

\end{document}               % End of document.

\section{Bench Section 131}
% synthetic bench section 131
% This is a sample LaTeX input file.  (Version of 12 August 2004.)
%
% A '%' character causes TeX to ignore all remaining text on the line,
% and is used for comments like this one.

\documentclass{article}      % Specifies the document class

                             % The preamble begins here.
\title{An Example Document}  % Declares the document's title.
\author{Leslie Lamport}      % Declares the author's name.
\date{January 21, 1994}      % Deleting this command produces today's date.

\newcommand{\ip}[2]{(#1, #2)}
                             % Defines \ip{arg1}{arg2} to mean
                             % (arg1, arg2).

%\newcommand{\ip}[2]{\langle #1 | #2\rangle}
                             % This is an alternative definition of
                             % \ip that is commented out.

\begin{document}             % End of preamble and beginning of text.

\maketitle                   % Produces the title.

This is an example input file.  Comparing it with
the output it generates can show you how to
produce a simple document of your own.

\section{Ordinary Text}      % Produces section heading.  Lower-level
                             % sections are begun with similar
                             % \subsection and \subsubsection commands.

The ends  of words and sentences are marked
  by   spaces. It  doesn't matter how many
spaces    you type; one is as good as 100.  The
end of   a line counts as a space.

One   or more   blank lines denote the  end
of  a paragraph.

Since any number of consecutive spaces are treated
like a single one, the formatting of the input
file makes no difference to
      \LaTeX,                % The \LaTeX command generates the LaTeX logo.
but it makes a difference to you.  When you use
\LaTeX, making your input file as easy to read
as possible will be a great help as you write
your document and when you change it.  This sample
file shows how you can add comments to your own input
file.

Because printing is different from typewriting,
there are a number of things that you have to do
differently when preparing an input file than if
you were just typing the document directly.
Quotation marks like
       ``this''
have to be handled specially, as do quotes within
quotes:
       ``\,`this'            % \, separates the double and single quote.
        is what I just
        wrote, not  `that'\,''.

Dashes come in three sizes: an
       intra-word
dash, a medium dash for number ranges like
       1--2,
and a punctuation
       dash---like
this.

A sentence-ending space should be larger than the
space between words within a sentence.  You
sometimes have to type special commands in
conjunction with punctuation characters to get
this right, as in the following sentence.
       Gnats, gnus, etc.\ all  % `\ ' makes an inter-word space.
       begin with G\@.         % \@ marks end-of-sentence punctuation.
You should check the spaces after periods when
reading your output to make sure you haven't
forgotten any special cases.  Generating an
ellipsis
       \ldots\               % `\ ' is needed after `\ldots' because TeX
                             % ignores spaces after command names like \ldots
                             % made from \ + letters.
                             %
                             % Note how a `%' character causes TeX to ignore
                             % the end of the input line, so these blank lines
                             % do not start a new paragraph.
                             %
with the right spacing around the periods requires
a special command.

\LaTeX\ interprets some common characters as
commands, so you must type special commands to
generate them.  These characters include the
following:
       \$ \& \% \# \{ and \}.

In printing, text is usually emphasized with an
       \emph{italic}
type style.

\begin{em}
   A long segment of text can also be emphasized
   in this way.  Text within such a segment can be
   given \emph{additional} emphasis.
\end{em}

It is sometimes necessary to prevent \LaTeX\ from
breaking a line where it might otherwise do so.
This may be at a space, as between the ``Mr.''\ and
``Jones'' in
       ``Mr.~Jones'',        % ~ produces an unbreakable interword space.
or within a word---especially when the word is a
symbol like
       \mbox{\emph{itemnum}}
that makes little sense when hyphenated across
lines.

Footnotes\footnote{This is an example of a footnote.}
pose no problem.

\LaTeX\ is good at typesetting mathematical formulas
like
       \( x-3y + z = 7 \)
or
       \( a_{1} > x^{2n} + y^{2n} > x' \)
or
       \( \ip{A}{B} = \sum_{i} a_{i} b_{i} \).
The spaces you type in a formula are
ignored.  Remember that a letter like
       $x$                   % $ ... $  and  \( ... \)  are equivalent
is a formula when it denotes a mathematical
symbol, and it should be typed as one.

\section{Displayed Text}

Text is displayed by indenting it from the left
margin.  Quotations are commonly displayed.  There
are short quotations
\begin{quote}
   This is a short quotation.  It consists of a
   single paragraph of text.  See how it is formatted.
\end{quote}
and longer ones.
\begin{quotation}
   This is a longer quotation.  It consists of two
   paragraphs of text, neither of which are
   particularly interesting.

   This is the second paragraph of the quotation.  It
   is just as dull as the first paragraph.
\end{quotation}
Another frequently-displayed structure is a list.
The following is an example of an \emph{itemized}
list.
\begin{itemize}
   \item This is the first item of an itemized list.
         Each item in the list is marked with a ``tick''.
         You don't have to worry about what kind of tick
         mark is used.

   \item This is the second item of the list.  It
         contains another list nested inside it.  The inner
         list is an \emph{enumerated} list.
         \begin{enumerate}
            \item This is the first item of an enumerated
                  list that is nested within the itemized list.

            \item This is the second item of the inner list.
                  \LaTeX\ allows you to nest lists deeper than
                  you really should.
         \end{enumerate}
         This is the rest of the second item of the outer
         list.  It is no more interesting than any other
         part of the item.
   \item This is the third item of the list.
\end{itemize}
You can even display poetry.
\begin{verse}
   There is an environment
    for verse \\             % The \\ command separates lines
   Whose features some poets % within a stanza.
   will curse.

                             % One or more blank lines separate stanzas.

   For instead of making\\
   Them do \emph{all} line breaking, \\
   It allows them to put too many words on a line when they'd rather be
   forced to be terse.
\end{verse}

Mathematical formulas may also be displayed.  A
displayed formula
is
one-line long; multiline
formulas require special formatting instructions.
   \[  \ip{\Gamma}{\psi'} = x'' + y^{2} + z_{i}^{n}\]
Don't start a paragraph with a displayed equation,
nor make one a paragraph by itself.

\end{document}               % End of document.

\section{Bench Section 132}
% synthetic bench section 132
% This is a sample LaTeX input file.  (Version of 12 August 2004.)
%
% A '%' character causes TeX to ignore all remaining text on the line,
% and is used for comments like this one.

\documentclass{article}      % Specifies the document class

                             % The preamble begins here.
\title{An Example Document}  % Declares the document's title.
\author{Leslie Lamport}      % Declares the author's name.
\date{January 21, 1994}      % Deleting this command produces today's date.

\newcommand{\ip}[2]{(#1, #2)}
                             % Defines \ip{arg1}{arg2} to mean
                             % (arg1, arg2).

%\newcommand{\ip}[2]{\langle #1 | #2\rangle}
                             % This is an alternative definition of
                             % \ip that is commented out.

\begin{document}             % End of preamble and beginning of text.

\maketitle                   % Produces the title.

This is an example input file.  Comparing it with
the output it generates can show you how to
produce a simple document of your own.

\section{Ordinary Text}      % Produces section heading.  Lower-level
                             % sections are begun with similar
                             % \subsection and \subsubsection commands.

The ends  of words and sentences are marked
  by   spaces. It  doesn't matter how many
spaces    you type; one is as good as 100.  The
end of   a line counts as a space.

One   or more   blank lines denote the  end
of  a paragraph.

Since any number of consecutive spaces are treated
like a single one, the formatting of the input
file makes no difference to
      \LaTeX,                % The \LaTeX command generates the LaTeX logo.
but it makes a difference to you.  When you use
\LaTeX, making your input file as easy to read
as possible will be a great help as you write
your document and when you change it.  This sample
file shows how you can add comments to your own input
file.

Because printing is different from typewriting,
there are a number of things that you have to do
differently when preparing an input file than if
you were just typing the document directly.
Quotation marks like
       ``this''
have to be handled specially, as do quotes within
quotes:
       ``\,`this'            % \, separates the double and single quote.
        is what I just
        wrote, not  `that'\,''.

Dashes come in three sizes: an
       intra-word
dash, a medium dash for number ranges like
       1--2,
and a punctuation
       dash---like
this.

A sentence-ending space should be larger than the
space between words within a sentence.  You
sometimes have to type special commands in
conjunction with punctuation characters to get
this right, as in the following sentence.
       Gnats, gnus, etc.\ all  % `\ ' makes an inter-word space.
       begin with G\@.         % \@ marks end-of-sentence punctuation.
You should check the spaces after periods when
reading your output to make sure you haven't
forgotten any special cases.  Generating an
ellipsis
       \ldots\               % `\ ' is needed after `\ldots' because TeX
                             % ignores spaces after command names like \ldots
                             % made from \ + letters.
                             %
                             % Note how a `%' character causes TeX to ignore
                             % the end of the input line, so these blank lines
                             % do not start a new paragraph.
                             %
with the right spacing around the periods requires
a special command.

\LaTeX\ interprets some common characters as
commands, so you must type special commands to
generate them.  These characters include the
following:
       \$ \& \% \# \{ and \}.

In printing, text is usually emphasized with an
       \emph{italic}
type style.

\begin{em}
   A long segment of text can also be emphasized
   in this way.  Text within such a segment can be
   given \emph{additional} emphasis.
\end{em}

It is sometimes necessary to prevent \LaTeX\ from
breaking a line where it might otherwise do so.
This may be at a space, as between the ``Mr.''\ and
``Jones'' in
       ``Mr.~Jones'',        % ~ produces an unbreakable interword space.
or within a word---especially when the word is a
symbol like
       \mbox{\emph{itemnum}}
that makes little sense when hyphenated across
lines.

Footnotes\footnote{This is an example of a footnote.}
pose no problem.

\LaTeX\ is good at typesetting mathematical formulas
like
       \( x-3y + z = 7 \)
or
       \( a_{1} > x^{2n} + y^{2n} > x' \)
or
       \( \ip{A}{B} = \sum_{i} a_{i} b_{i} \).
The spaces you type in a formula are
ignored.  Remember that a letter like
       $x$                   % $ ... $  and  \( ... \)  are equivalent
is a formula when it denotes a mathematical
symbol, and it should be typed as one.

\section{Displayed Text}

Text is displayed by indenting it from the left
margin.  Quotations are commonly displayed.  There
are short quotations
\begin{quote}
   This is a short quotation.  It consists of a
   single paragraph of text.  See how it is formatted.
\end{quote}
and longer ones.
\begin{quotation}
   This is a longer quotation.  It consists of two
   paragraphs of text, neither of which are
   particularly interesting.

   This is the second paragraph of the quotation.  It
   is just as dull as the first paragraph.
\end{quotation}
Another frequently-displayed structure is a list.
The following is an example of an \emph{itemized}
list.
\begin{itemize}
   \item This is the first item of an itemized list.
         Each item in the list is marked with a ``tick''.
         You don't have to worry about what kind of tick
         mark is used.

   \item This is the second item of the list.  It
         contains another list nested inside it.  The inner
         list is an \emph{enumerated} list.
         \begin{enumerate}
            \item This is the first item of an enumerated
                  list that is nested within the itemized list.

            \item This is the second item of the inner list.
                  \LaTeX\ allows you to nest lists deeper than
                  you really should.
         \end{enumerate}
         This is the rest of the second item of the outer
         list.  It is no more interesting than any other
         part of the item.
   \item This is the third item of the list.
\end{itemize}
You can even display poetry.
\begin{verse}
   There is an environment
    for verse \\             % The \\ command separates lines
   Whose features some poets % within a stanza.
   will curse.

                             % One or more blank lines separate stanzas.

   For instead of making\\
   Them do \emph{all} line breaking, \\
   It allows them to put too many words on a line when they'd rather be
   forced to be terse.
\end{verse}

Mathematical formulas may also be displayed.  A
displayed formula
is
one-line long; multiline
formulas require special formatting instructions.
   \[  \ip{\Gamma}{\psi'} = x'' + y^{2} + z_{i}^{n}\]
Don't start a paragraph with a displayed equation,
nor make one a paragraph by itself.

\end{document}               % End of document.

\section{Bench Section 133}
% synthetic bench section 133
% This is a sample LaTeX input file.  (Version of 12 August 2004.)
%
% A '%' character causes TeX to ignore all remaining text on the line,
% and is used for comments like this one.

\documentclass{article}      % Specifies the document class

                             % The preamble begins here.
\title{An Example Document}  % Declares the document's title.
\author{Leslie Lamport}      % Declares the author's name.
\date{January 21, 1994}      % Deleting this command produces today's date.

\newcommand{\ip}[2]{(#1, #2)}
                             % Defines \ip{arg1}{arg2} to mean
                             % (arg1, arg2).

%\newcommand{\ip}[2]{\langle #1 | #2\rangle}
                             % This is an alternative definition of
                             % \ip that is commented out.

\begin{document}             % End of preamble and beginning of text.

\maketitle                   % Produces the title.

This is an example input file.  Comparing it with
the output it generates can show you how to
produce a simple document of your own.

\section{Ordinary Text}      % Produces section heading.  Lower-level
                             % sections are begun with similar
                             % \subsection and \subsubsection commands.

The ends  of words and sentences are marked
  by   spaces. It  doesn't matter how many
spaces    you type; one is as good as 100.  The
end of   a line counts as a space.

One   or more   blank lines denote the  end
of  a paragraph.

Since any number of consecutive spaces are treated
like a single one, the formatting of the input
file makes no difference to
      \LaTeX,                % The \LaTeX command generates the LaTeX logo.
but it makes a difference to you.  When you use
\LaTeX, making your input file as easy to read
as possible will be a great help as you write
your document and when you change it.  This sample
file shows how you can add comments to your own input
file.

Because printing is different from typewriting,
there are a number of things that you have to do
differently when preparing an input file than if
you were just typing the document directly.
Quotation marks like
       ``this''
have to be handled specially, as do quotes within
quotes:
       ``\,`this'            % \, separates the double and single quote.
        is what I just
        wrote, not  `that'\,''.

Dashes come in three sizes: an
       intra-word
dash, a medium dash for number ranges like
       1--2,
and a punctuation
       dash---like
this.

A sentence-ending space should be larger than the
space between words within a sentence.  You
sometimes have to type special commands in
conjunction with punctuation characters to get
this right, as in the following sentence.
       Gnats, gnus, etc.\ all  % `\ ' makes an inter-word space.
       begin with G\@.         % \@ marks end-of-sentence punctuation.
You should check the spaces after periods when
reading your output to make sure you haven't
forgotten any special cases.  Generating an
ellipsis
       \ldots\               % `\ ' is needed after `\ldots' because TeX
                             % ignores spaces after command names like \ldots
                             % made from \ + letters.
                             %
                             % Note how a `%' character causes TeX to ignore
                             % the end of the input line, so these blank lines
                             % do not start a new paragraph.
                             %
with the right spacing around the periods requires
a special command.

\LaTeX\ interprets some common characters as
commands, so you must type special commands to
generate them.  These characters include the
following:
       \$ \& \% \# \{ and \}.

In printing, text is usually emphasized with an
       \emph{italic}
type style.

\begin{em}
   A long segment of text can also be emphasized
   in this way.  Text within such a segment can be
   given \emph{additional} emphasis.
\end{em}

It is sometimes necessary to prevent \LaTeX\ from
breaking a line where it might otherwise do so.
This may be at a space, as between the ``Mr.''\ and
``Jones'' in
       ``Mr.~Jones'',        % ~ produces an unbreakable interword space.
or within a word---especially when the word is a
symbol like
       \mbox{\emph{itemnum}}
that makes little sense when hyphenated across
lines.

Footnotes\footnote{This is an example of a footnote.}
pose no problem.

\LaTeX\ is good at typesetting mathematical formulas
like
       \( x-3y + z = 7 \)
or
       \( a_{1} > x^{2n} + y^{2n} > x' \)
or
       \( \ip{A}{B} = \sum_{i} a_{i} b_{i} \).
The spaces you type in a formula are
ignored.  Remember that a letter like
       $x$                   % $ ... $  and  \( ... \)  are equivalent
is a formula when it denotes a mathematical
symbol, and it should be typed as one.

\section{Displayed Text}

Text is displayed by indenting it from the left
margin.  Quotations are commonly displayed.  There
are short quotations
\begin{quote}
   This is a short quotation.  It consists of a
   single paragraph of text.  See how it is formatted.
\end{quote}
and longer ones.
\begin{quotation}
   This is a longer quotation.  It consists of two
   paragraphs of text, neither of which are
   particularly interesting.

   This is the second paragraph of the quotation.  It
   is just as dull as the first paragraph.
\end{quotation}
Another frequently-displayed structure is a list.
The following is an example of an \emph{itemized}
list.
\begin{itemize}
   \item This is the first item of an itemized list.
         Each item in the list is marked with a ``tick''.
         You don't have to worry about what kind of tick
         mark is used.

   \item This is the second item of the list.  It
         contains another list nested inside it.  The inner
         list is an \emph{enumerated} list.
         \begin{enumerate}
            \item This is the first item of an enumerated
                  list that is nested within the itemized list.

            \item This is the second item of the inner list.
                  \LaTeX\ allows you to nest lists deeper than
                  you really should.
         \end{enumerate}
         This is the rest of the second item of the outer
         list.  It is no more interesting than any other
         part of the item.
   \item This is the third item of the list.
\end{itemize}
You can even display poetry.
\begin{verse}
   There is an environment
    for verse \\             % The \\ command separates lines
   Whose features some poets % within a stanza.
   will curse.

                             % One or more blank lines separate stanzas.

   For instead of making\\
   Them do \emph{all} line breaking, \\
   It allows them to put too many words on a line when they'd rather be
   forced to be terse.
\end{verse}

Mathematical formulas may also be displayed.  A
displayed formula
is
one-line long; multiline
formulas require special formatting instructions.
   \[  \ip{\Gamma}{\psi'} = x'' + y^{2} + z_{i}^{n}\]
Don't start a paragraph with a displayed equation,
nor make one a paragraph by itself.

\end{document}               % End of document.

\section{Bench Section 134}
% synthetic bench section 134
% This is a sample LaTeX input file.  (Version of 12 August 2004.)
%
% A '%' character causes TeX to ignore all remaining text on the line,
% and is used for comments like this one.

\documentclass{article}      % Specifies the document class

                             % The preamble begins here.
\title{An Example Document}  % Declares the document's title.
\author{Leslie Lamport}      % Declares the author's name.
\date{January 21, 1994}      % Deleting this command produces today's date.

\newcommand{\ip}[2]{(#1, #2)}
                             % Defines \ip{arg1}{arg2} to mean
                             % (arg1, arg2).

%\newcommand{\ip}[2]{\langle #1 | #2\rangle}
                             % This is an alternative definition of
                             % \ip that is commented out.

\begin{document}             % End of preamble and beginning of text.

\maketitle                   % Produces the title.

This is an example input file.  Comparing it with
the output it generates can show you how to
produce a simple document of your own.

\section{Ordinary Text}      % Produces section heading.  Lower-level
                             % sections are begun with similar
                             % \subsection and \subsubsection commands.

The ends  of words and sentences are marked
  by   spaces. It  doesn't matter how many
spaces    you type; one is as good as 100.  The
end of   a line counts as a space.

One   or more   blank lines denote the  end
of  a paragraph.

Since any number of consecutive spaces are treated
like a single one, the formatting of the input
file makes no difference to
      \LaTeX,                % The \LaTeX command generates the LaTeX logo.
but it makes a difference to you.  When you use
\LaTeX, making your input file as easy to read
as possible will be a great help as you write
your document and when you change it.  This sample
file shows how you can add comments to your own input
file.

Because printing is different from typewriting,
there are a number of things that you have to do
differently when preparing an input file than if
you were just typing the document directly.
Quotation marks like
       ``this''
have to be handled specially, as do quotes within
quotes:
       ``\,`this'            % \, separates the double and single quote.
        is what I just
        wrote, not  `that'\,''.

Dashes come in three sizes: an
       intra-word
dash, a medium dash for number ranges like
       1--2,
and a punctuation
       dash---like
this.

A sentence-ending space should be larger than the
space between words within a sentence.  You
sometimes have to type special commands in
conjunction with punctuation characters to get
this right, as in the following sentence.
       Gnats, gnus, etc.\ all  % `\ ' makes an inter-word space.
       begin with G\@.         % \@ marks end-of-sentence punctuation.
You should check the spaces after periods when
reading your output to make sure you haven't
forgotten any special cases.  Generating an
ellipsis
       \ldots\               % `\ ' is needed after `\ldots' because TeX
                             % ignores spaces after command names like \ldots
                             % made from \ + letters.
                             %
                             % Note how a `%' character causes TeX to ignore
                             % the end of the input line, so these blank lines
                             % do not start a new paragraph.
                             %
with the right spacing around the periods requires
a special command.

\LaTeX\ interprets some common characters as
commands, so you must type special commands to
generate them.  These characters include the
following:
       \$ \& \% \# \{ and \}.

In printing, text is usually emphasized with an
       \emph{italic}
type style.

\begin{em}
   A long segment of text can also be emphasized
   in this way.  Text within such a segment can be
   given \emph{additional} emphasis.
\end{em}

It is sometimes necessary to prevent \LaTeX\ from
breaking a line where it might otherwise do so.
This may be at a space, as between the ``Mr.''\ and
``Jones'' in
       ``Mr.~Jones'',        % ~ produces an unbreakable interword space.
or within a word---especially when the word is a
symbol like
       \mbox{\emph{itemnum}}
that makes little sense when hyphenated across
lines.

Footnotes\footnote{This is an example of a footnote.}
pose no problem.

\LaTeX\ is good at typesetting mathematical formulas
like
       \( x-3y + z = 7 \)
or
       \( a_{1} > x^{2n} + y^{2n} > x' \)
or
       \( \ip{A}{B} = \sum_{i} a_{i} b_{i} \).
The spaces you type in a formula are
ignored.  Remember that a letter like
       $x$                   % $ ... $  and  \( ... \)  are equivalent
is a formula when it denotes a mathematical
symbol, and it should be typed as one.

\section{Displayed Text}

Text is displayed by indenting it from the left
margin.  Quotations are commonly displayed.  There
are short quotations
\begin{quote}
   This is a short quotation.  It consists of a
   single paragraph of text.  See how it is formatted.
\end{quote}
and longer ones.
\begin{quotation}
   This is a longer quotation.  It consists of two
   paragraphs of text, neither of which are
   particularly interesting.

   This is the second paragraph of the quotation.  It
   is just as dull as the first paragraph.
\end{quotation}
Another frequently-displayed structure is a list.
The following is an example of an \emph{itemized}
list.
\begin{itemize}
   \item This is the first item of an itemized list.
         Each item in the list is marked with a ``tick''.
         You don't have to worry about what kind of tick
         mark is used.

   \item This is the second item of the list.  It
         contains another list nested inside it.  The inner
         list is an \emph{enumerated} list.
         \begin{enumerate}
            \item This is the first item of an enumerated
                  list that is nested within the itemized list.

            \item This is the second item of the inner list.
                  \LaTeX\ allows you to nest lists deeper than
                  you really should.
         \end{enumerate}
         This is the rest of the second item of the outer
         list.  It is no more interesting than any other
         part of the item.
   \item This is the third item of the list.
\end{itemize}
You can even display poetry.
\begin{verse}
   There is an environment
    for verse \\             % The \\ command separates lines
   Whose features some poets % within a stanza.
   will curse.

                             % One or more blank lines separate stanzas.

   For instead of making\\
   Them do \emph{all} line breaking, \\
   It allows them to put too many words on a line when they'd rather be
   forced to be terse.
\end{verse}

Mathematical formulas may also be displayed.  A
displayed formula
is
one-line long; multiline
formulas require special formatting instructions.
   \[  \ip{\Gamma}{\psi'} = x'' + y^{2} + z_{i}^{n}\]
Don't start a paragraph with a displayed equation,
nor make one a paragraph by itself.

\end{document}               % End of document.

\section{Bench Section 135}
% synthetic bench section 135
% This is a sample LaTeX input file.  (Version of 12 August 2004.)
%
% A '%' character causes TeX to ignore all remaining text on the line,
% and is used for comments like this one.

\documentclass{article}      % Specifies the document class

                             % The preamble begins here.
\title{An Example Document}  % Declares the document's title.
\author{Leslie Lamport}      % Declares the author's name.
\date{January 21, 1994}      % Deleting this command produces today's date.

\newcommand{\ip}[2]{(#1, #2)}
                             % Defines \ip{arg1}{arg2} to mean
                             % (arg1, arg2).

%\newcommand{\ip}[2]{\langle #1 | #2\rangle}
                             % This is an alternative definition of
                             % \ip that is commented out.

\begin{document}             % End of preamble and beginning of text.

\maketitle                   % Produces the title.

This is an example input file.  Comparing it with
the output it generates can show you how to
produce a simple document of your own.

\section{Ordinary Text}      % Produces section heading.  Lower-level
                             % sections are begun with similar
                             % \subsection and \subsubsection commands.

The ends  of words and sentences are marked
  by   spaces. It  doesn't matter how many
spaces    you type; one is as good as 100.  The
end of   a line counts as a space.

One   or more   blank lines denote the  end
of  a paragraph.

Since any number of consecutive spaces are treated
like a single one, the formatting of the input
file makes no difference to
      \LaTeX,                % The \LaTeX command generates the LaTeX logo.
but it makes a difference to you.  When you use
\LaTeX, making your input file as easy to read
as possible will be a great help as you write
your document and when you change it.  This sample
file shows how you can add comments to your own input
file.

Because printing is different from typewriting,
there are a number of things that you have to do
differently when preparing an input file than if
you were just typing the document directly.
Quotation marks like
       ``this''
have to be handled specially, as do quotes within
quotes:
       ``\,`this'            % \, separates the double and single quote.
        is what I just
        wrote, not  `that'\,''.

Dashes come in three sizes: an
       intra-word
dash, a medium dash for number ranges like
       1--2,
and a punctuation
       dash---like
this.

A sentence-ending space should be larger than the
space between words within a sentence.  You
sometimes have to type special commands in
conjunction with punctuation characters to get
this right, as in the following sentence.
       Gnats, gnus, etc.\ all  % `\ ' makes an inter-word space.
       begin with G\@.         % \@ marks end-of-sentence punctuation.
You should check the spaces after periods when
reading your output to make sure you haven't
forgotten any special cases.  Generating an
ellipsis
       \ldots\               % `\ ' is needed after `\ldots' because TeX
                             % ignores spaces after command names like \ldots
                             % made from \ + letters.
                             %
                             % Note how a `%' character causes TeX to ignore
                             % the end of the input line, so these blank lines
                             % do not start a new paragraph.
                             %
with the right spacing around the periods requires
a special command.

\LaTeX\ interprets some common characters as
commands, so you must type special commands to
generate them.  These characters include the
following:
       \$ \& \% \# \{ and \}.

In printing, text is usually emphasized with an
       \emph{italic}
type style.

\begin{em}
   A long segment of text can also be emphasized
   in this way.  Text within such a segment can be
   given \emph{additional} emphasis.
\end{em}

It is sometimes necessary to prevent \LaTeX\ from
breaking a line where it might otherwise do so.
This may be at a space, as between the ``Mr.''\ and
``Jones'' in
       ``Mr.~Jones'',        % ~ produces an unbreakable interword space.
or within a word---especially when the word is a
symbol like
       \mbox{\emph{itemnum}}
that makes little sense when hyphenated across
lines.

Footnotes\footnote{This is an example of a footnote.}
pose no problem.

\LaTeX\ is good at typesetting mathematical formulas
like
       \( x-3y + z = 7 \)
or
       \( a_{1} > x^{2n} + y^{2n} > x' \)
or
       \( \ip{A}{B} = \sum_{i} a_{i} b_{i} \).
The spaces you type in a formula are
ignored.  Remember that a letter like
       $x$                   % $ ... $  and  \( ... \)  are equivalent
is a formula when it denotes a mathematical
symbol, and it should be typed as one.

\section{Displayed Text}

Text is displayed by indenting it from the left
margin.  Quotations are commonly displayed.  There
are short quotations
\begin{quote}
   This is a short quotation.  It consists of a
   single paragraph of text.  See how it is formatted.
\end{quote}
and longer ones.
\begin{quotation}
   This is a longer quotation.  It consists of two
   paragraphs of text, neither of which are
   particularly interesting.

   This is the second paragraph of the quotation.  It
   is just as dull as the first paragraph.
\end{quotation}
Another frequently-displayed structure is a list.
The following is an example of an \emph{itemized}
list.
\begin{itemize}
   \item This is the first item of an itemized list.
         Each item in the list is marked with a ``tick''.
         You don't have to worry about what kind of tick
         mark is used.

   \item This is the second item of the list.  It
         contains another list nested inside it.  The inner
         list is an \emph{enumerated} list.
         \begin{enumerate}
            \item This is the first item of an enumerated
                  list that is nested within the itemized list.

            \item This is the second item of the inner list.
                  \LaTeX\ allows you to nest lists deeper than
                  you really should.
         \end{enumerate}
         This is the rest of the second item of the outer
         list.  It is no more interesting than any other
         part of the item.
   \item This is the third item of the list.
\end{itemize}
You can even display poetry.
\begin{verse}
   There is an environment
    for verse \\             % The \\ command separates lines
   Whose features some poets % within a stanza.
   will curse.

                             % One or more blank lines separate stanzas.

   For instead of making\\
   Them do \emph{all} line breaking, \\
   It allows them to put too many words on a line when they'd rather be
   forced to be terse.
\end{verse}

Mathematical formulas may also be displayed.  A
displayed formula
is
one-line long; multiline
formulas require special formatting instructions.
   \[  \ip{\Gamma}{\psi'} = x'' + y^{2} + z_{i}^{n}\]
Don't start a paragraph with a displayed equation,
nor make one a paragraph by itself.

\end{document}               % End of document.

\section{Bench Section 136}
% synthetic bench section 136
% This is a sample LaTeX input file.  (Version of 12 August 2004.)
%
% A '%' character causes TeX to ignore all remaining text on the line,
% and is used for comments like this one.

\documentclass{article}      % Specifies the document class

                             % The preamble begins here.
\title{An Example Document}  % Declares the document's title.
\author{Leslie Lamport}      % Declares the author's name.
\date{January 21, 1994}      % Deleting this command produces today's date.

\newcommand{\ip}[2]{(#1, #2)}
                             % Defines \ip{arg1}{arg2} to mean
                             % (arg1, arg2).

%\newcommand{\ip}[2]{\langle #1 | #2\rangle}
                             % This is an alternative definition of
                             % \ip that is commented out.

\begin{document}             % End of preamble and beginning of text.

\maketitle                   % Produces the title.

This is an example input file.  Comparing it with
the output it generates can show you how to
produce a simple document of your own.

\section{Ordinary Text}      % Produces section heading.  Lower-level
                             % sections are begun with similar
                             % \subsection and \subsubsection commands.

The ends  of words and sentences are marked
  by   spaces. It  doesn't matter how many
spaces    you type; one is as good as 100.  The
end of   a line counts as a space.

One   or more   blank lines denote the  end
of  a paragraph.

Since any number of consecutive spaces are treated
like a single one, the formatting of the input
file makes no difference to
      \LaTeX,                % The \LaTeX command generates the LaTeX logo.
but it makes a difference to you.  When you use
\LaTeX, making your input file as easy to read
as possible will be a great help as you write
your document and when you change it.  This sample
file shows how you can add comments to your own input
file.

Because printing is different from typewriting,
there are a number of things that you have to do
differently when preparing an input file than if
you were just typing the document directly.
Quotation marks like
       ``this''
have to be handled specially, as do quotes within
quotes:
       ``\,`this'            % \, separates the double and single quote.
        is what I just
        wrote, not  `that'\,''.

Dashes come in three sizes: an
       intra-word
dash, a medium dash for number ranges like
       1--2,
and a punctuation
       dash---like
this.

A sentence-ending space should be larger than the
space between words within a sentence.  You
sometimes have to type special commands in
conjunction with punctuation characters to get
this right, as in the following sentence.
       Gnats, gnus, etc.\ all  % `\ ' makes an inter-word space.
       begin with G\@.         % \@ marks end-of-sentence punctuation.
You should check the spaces after periods when
reading your output to make sure you haven't
forgotten any special cases.  Generating an
ellipsis
       \ldots\               % `\ ' is needed after `\ldots' because TeX
                             % ignores spaces after command names like \ldots
                             % made from \ + letters.
                             %
                             % Note how a `%' character causes TeX to ignore
                             % the end of the input line, so these blank lines
                             % do not start a new paragraph.
                             %
with the right spacing around the periods requires
a special command.

\LaTeX\ interprets some common characters as
commands, so you must type special commands to
generate them.  These characters include the
following:
       \$ \& \% \# \{ and \}.

In printing, text is usually emphasized with an
       \emph{italic}
type style.

\begin{em}
   A long segment of text can also be emphasized
   in this way.  Text within such a segment can be
   given \emph{additional} emphasis.
\end{em}

It is sometimes necessary to prevent \LaTeX\ from
breaking a line where it might otherwise do so.
This may be at a space, as between the ``Mr.''\ and
``Jones'' in
       ``Mr.~Jones'',        % ~ produces an unbreakable interword space.
or within a word---especially when the word is a
symbol like
       \mbox{\emph{itemnum}}
that makes little sense when hyphenated across
lines.

Footnotes\footnote{This is an example of a footnote.}
pose no problem.

\LaTeX\ is good at typesetting mathematical formulas
like
       \( x-3y + z = 7 \)
or
       \( a_{1} > x^{2n} + y^{2n} > x' \)
or
       \( \ip{A}{B} = \sum_{i} a_{i} b_{i} \).
The spaces you type in a formula are
ignored.  Remember that a letter like
       $x$                   % $ ... $  and  \( ... \)  are equivalent
is a formula when it denotes a mathematical
symbol, and it should be typed as one.

\section{Displayed Text}

Text is displayed by indenting it from the left
margin.  Quotations are commonly displayed.  There
are short quotations
\begin{quote}
   This is a short quotation.  It consists of a
   single paragraph of text.  See how it is formatted.
\end{quote}
and longer ones.
\begin{quotation}
   This is a longer quotation.  It consists of two
   paragraphs of text, neither of which are
   particularly interesting.

   This is the second paragraph of the quotation.  It
   is just as dull as the first paragraph.
\end{quotation}
Another frequently-displayed structure is a list.
The following is an example of an \emph{itemized}
list.
\begin{itemize}
   \item This is the first item of an itemized list.
         Each item in the list is marked with a ``tick''.
         You don't have to worry about what kind of tick
         mark is used.

   \item This is the second item of the list.  It
         contains another list nested inside it.  The inner
         list is an \emph{enumerated} list.
         \begin{enumerate}
            \item This is the first item of an enumerated
                  list that is nested within the itemized list.

            \item This is the second item of the inner list.
                  \LaTeX\ allows you to nest lists deeper than
                  you really should.
         \end{enumerate}
         This is the rest of the second item of the outer
         list.  It is no more interesting than any other
         part of the item.
   \item This is the third item of the list.
\end{itemize}
You can even display poetry.
\begin{verse}
   There is an environment
    for verse \\             % The \\ command separates lines
   Whose features some poets % within a stanza.
   will curse.

                             % One or more blank lines separate stanzas.

   For instead of making\\
   Them do \emph{all} line breaking, \\
   It allows them to put too many words on a line when they'd rather be
   forced to be terse.
\end{verse}

Mathematical formulas may also be displayed.  A
displayed formula
is
one-line long; multiline
formulas require special formatting instructions.
   \[  \ip{\Gamma}{\psi'} = x'' + y^{2} + z_{i}^{n}\]
Don't start a paragraph with a displayed equation,
nor make one a paragraph by itself.

\end{document}               % End of document.

\section{Bench Section 137}
% synthetic bench section 137
% This is a sample LaTeX input file.  (Version of 12 August 2004.)
%
% A '%' character causes TeX to ignore all remaining text on the line,
% and is used for comments like this one.

\documentclass{article}      % Specifies the document class

                             % The preamble begins here.
\title{An Example Document}  % Declares the document's title.
\author{Leslie Lamport}      % Declares the author's name.
\date{January 21, 1994}      % Deleting this command produces today's date.

\newcommand{\ip}[2]{(#1, #2)}
                             % Defines \ip{arg1}{arg2} to mean
                             % (arg1, arg2).

%\newcommand{\ip}[2]{\langle #1 | #2\rangle}
                             % This is an alternative definition of
                             % \ip that is commented out.

\begin{document}             % End of preamble and beginning of text.

\maketitle                   % Produces the title.

This is an example input file.  Comparing it with
the output it generates can show you how to
produce a simple document of your own.

\section{Ordinary Text}      % Produces section heading.  Lower-level
                             % sections are begun with similar
                             % \subsection and \subsubsection commands.

The ends  of words and sentences are marked
  by   spaces. It  doesn't matter how many
spaces    you type; one is as good as 100.  The
end of   a line counts as a space.

One   or more   blank lines denote the  end
of  a paragraph.

Since any number of consecutive spaces are treated
like a single one, the formatting of the input
file makes no difference to
      \LaTeX,                % The \LaTeX command generates the LaTeX logo.
but it makes a difference to you.  When you use
\LaTeX, making your input file as easy to read
as possible will be a great help as you write
your document and when you change it.  This sample
file shows how you can add comments to your own input
file.

Because printing is different from typewriting,
there are a number of things that you have to do
differently when preparing an input file than if
you were just typing the document directly.
Quotation marks like
       ``this''
have to be handled specially, as do quotes within
quotes:
       ``\,`this'            % \, separates the double and single quote.
        is what I just
        wrote, not  `that'\,''.

Dashes come in three sizes: an
       intra-word
dash, a medium dash for number ranges like
       1--2,
and a punctuation
       dash---like
this.

A sentence-ending space should be larger than the
space between words within a sentence.  You
sometimes have to type special commands in
conjunction with punctuation characters to get
this right, as in the following sentence.
       Gnats, gnus, etc.\ all  % `\ ' makes an inter-word space.
       begin with G\@.         % \@ marks end-of-sentence punctuation.
You should check the spaces after periods when
reading your output to make sure you haven't
forgotten any special cases.  Generating an
ellipsis
       \ldots\               % `\ ' is needed after `\ldots' because TeX
                             % ignores spaces after command names like \ldots
                             % made from \ + letters.
                             %
                             % Note how a `%' character causes TeX to ignore
                             % the end of the input line, so these blank lines
                             % do not start a new paragraph.
                             %
with the right spacing around the periods requires
a special command.

\LaTeX\ interprets some common characters as
commands, so you must type special commands to
generate them.  These characters include the
following:
       \$ \& \% \# \{ and \}.

In printing, text is usually emphasized with an
       \emph{italic}
type style.

\begin{em}
   A long segment of text can also be emphasized
   in this way.  Text within such a segment can be
   given \emph{additional} emphasis.
\end{em}

It is sometimes necessary to prevent \LaTeX\ from
breaking a line where it might otherwise do so.
This may be at a space, as between the ``Mr.''\ and
``Jones'' in
       ``Mr.~Jones'',        % ~ produces an unbreakable interword space.
or within a word---especially when the word is a
symbol like
       \mbox{\emph{itemnum}}
that makes little sense when hyphenated across
lines.

Footnotes\footnote{This is an example of a footnote.}
pose no problem.

\LaTeX\ is good at typesetting mathematical formulas
like
       \( x-3y + z = 7 \)
or
       \( a_{1} > x^{2n} + y^{2n} > x' \)
or
       \( \ip{A}{B} = \sum_{i} a_{i} b_{i} \).
The spaces you type in a formula are
ignored.  Remember that a letter like
       $x$                   % $ ... $  and  \( ... \)  are equivalent
is a formula when it denotes a mathematical
symbol, and it should be typed as one.

\section{Displayed Text}

Text is displayed by indenting it from the left
margin.  Quotations are commonly displayed.  There
are short quotations
\begin{quote}
   This is a short quotation.  It consists of a
   single paragraph of text.  See how it is formatted.
\end{quote}
and longer ones.
\begin{quotation}
   This is a longer quotation.  It consists of two
   paragraphs of text, neither of which are
   particularly interesting.

   This is the second paragraph of the quotation.  It
   is just as dull as the first paragraph.
\end{quotation}
Another frequently-displayed structure is a list.
The following is an example of an \emph{itemized}
list.
\begin{itemize}
   \item This is the first item of an itemized list.
         Each item in the list is marked with a ``tick''.
         You don't have to worry about what kind of tick
         mark is used.

   \item This is the second item of the list.  It
         contains another list nested inside it.  The inner
         list is an \emph{enumerated} list.
         \begin{enumerate}
            \item This is the first item of an enumerated
                  list that is nested within the itemized list.

            \item This is the second item of the inner list.
                  \LaTeX\ allows you to nest lists deeper than
                  you really should.
         \end{enumerate}
         This is the rest of the second item of the outer
         list.  It is no more interesting than any other
         part of the item.
   \item This is the third item of the list.
\end{itemize}
You can even display poetry.
\begin{verse}
   There is an environment
    for verse \\             % The \\ command separates lines
   Whose features some poets % within a stanza.
   will curse.

                             % One or more blank lines separate stanzas.

   For instead of making\\
   Them do \emph{all} line breaking, \\
   It allows them to put too many words on a line when they'd rather be
   forced to be terse.
\end{verse}

Mathematical formulas may also be displayed.  A
displayed formula
is
one-line long; multiline
formulas require special formatting instructions.
   \[  \ip{\Gamma}{\psi'} = x'' + y^{2} + z_{i}^{n}\]
Don't start a paragraph with a displayed equation,
nor make one a paragraph by itself.

\end{document}               % End of document.

\section{Bench Section 138}
% synthetic bench section 138
% This is a sample LaTeX input file.  (Version of 12 August 2004.)
%
% A '%' character causes TeX to ignore all remaining text on the line,
% and is used for comments like this one.

\documentclass{article}      % Specifies the document class

                             % The preamble begins here.
\title{An Example Document}  % Declares the document's title.
\author{Leslie Lamport}      % Declares the author's name.
\date{January 21, 1994}      % Deleting this command produces today's date.

\newcommand{\ip}[2]{(#1, #2)}
                             % Defines \ip{arg1}{arg2} to mean
                             % (arg1, arg2).

%\newcommand{\ip}[2]{\langle #1 | #2\rangle}
                             % This is an alternative definition of
                             % \ip that is commented out.

\begin{document}             % End of preamble and beginning of text.

\maketitle                   % Produces the title.

This is an example input file.  Comparing it with
the output it generates can show you how to
produce a simple document of your own.

\section{Ordinary Text}      % Produces section heading.  Lower-level
                             % sections are begun with similar
                             % \subsection and \subsubsection commands.

The ends  of words and sentences are marked
  by   spaces. It  doesn't matter how many
spaces    you type; one is as good as 100.  The
end of   a line counts as a space.

One   or more   blank lines denote the  end
of  a paragraph.

Since any number of consecutive spaces are treated
like a single one, the formatting of the input
file makes no difference to
      \LaTeX,                % The \LaTeX command generates the LaTeX logo.
but it makes a difference to you.  When you use
\LaTeX, making your input file as easy to read
as possible will be a great help as you write
your document and when you change it.  This sample
file shows how you can add comments to your own input
file.

Because printing is different from typewriting,
there are a number of things that you have to do
differently when preparing an input file than if
you were just typing the document directly.
Quotation marks like
       ``this''
have to be handled specially, as do quotes within
quotes:
       ``\,`this'            % \, separates the double and single quote.
        is what I just
        wrote, not  `that'\,''.

Dashes come in three sizes: an
       intra-word
dash, a medium dash for number ranges like
       1--2,
and a punctuation
       dash---like
this.

A sentence-ending space should be larger than the
space between words within a sentence.  You
sometimes have to type special commands in
conjunction with punctuation characters to get
this right, as in the following sentence.
       Gnats, gnus, etc.\ all  % `\ ' makes an inter-word space.
       begin with G\@.         % \@ marks end-of-sentence punctuation.
You should check the spaces after periods when
reading your output to make sure you haven't
forgotten any special cases.  Generating an
ellipsis
       \ldots\               % `\ ' is needed after `\ldots' because TeX
                             % ignores spaces after command names like \ldots
                             % made from \ + letters.
                             %
                             % Note how a `%' character causes TeX to ignore
                             % the end of the input line, so these blank lines
                             % do not start a new paragraph.
                             %
with the right spacing around the periods requires
a special command.

\LaTeX\ interprets some common characters as
commands, so you must type special commands to
generate them.  These characters include the
following:
       \$ \& \% \# \{ and \}.

In printing, text is usually emphasized with an
       \emph{italic}
type style.

\begin{em}
   A long segment of text can also be emphasized
   in this way.  Text within such a segment can be
   given \emph{additional} emphasis.
\end{em}

It is sometimes necessary to prevent \LaTeX\ from
breaking a line where it might otherwise do so.
This may be at a space, as between the ``Mr.''\ and
``Jones'' in
       ``Mr.~Jones'',        % ~ produces an unbreakable interword space.
or within a word---especially when the word is a
symbol like
       \mbox{\emph{itemnum}}
that makes little sense when hyphenated across
lines.

Footnotes\footnote{This is an example of a footnote.}
pose no problem.

\LaTeX\ is good at typesetting mathematical formulas
like
       \( x-3y + z = 7 \)
or
       \( a_{1} > x^{2n} + y^{2n} > x' \)
or
       \( \ip{A}{B} = \sum_{i} a_{i} b_{i} \).
The spaces you type in a formula are
ignored.  Remember that a letter like
       $x$                   % $ ... $  and  \( ... \)  are equivalent
is a formula when it denotes a mathematical
symbol, and it should be typed as one.

\section{Displayed Text}

Text is displayed by indenting it from the left
margin.  Quotations are commonly displayed.  There
are short quotations
\begin{quote}
   This is a short quotation.  It consists of a
   single paragraph of text.  See how it is formatted.
\end{quote}
and longer ones.
\begin{quotation}
   This is a longer quotation.  It consists of two
   paragraphs of text, neither of which are
   particularly interesting.

   This is the second paragraph of the quotation.  It
   is just as dull as the first paragraph.
\end{quotation}
Another frequently-displayed structure is a list.
The following is an example of an \emph{itemized}
list.
\begin{itemize}
   \item This is the first item of an itemized list.
         Each item in the list is marked with a ``tick''.
         You don't have to worry about what kind of tick
         mark is used.

   \item This is the second item of the list.  It
         contains another list nested inside it.  The inner
         list is an \emph{enumerated} list.
         \begin{enumerate}
            \item This is the first item of an enumerated
                  list that is nested within the itemized list.

            \item This is the second item of the inner list.
                  \LaTeX\ allows you to nest lists deeper than
                  you really should.
         \end{enumerate}
         This is the rest of the second item of the outer
         list.  It is no more interesting than any other
         part of the item.
   \item This is the third item of the list.
\end{itemize}
You can even display poetry.
\begin{verse}
   There is an environment
    for verse \\             % The \\ command separates lines
   Whose features some poets % within a stanza.
   will curse.

                             % One or more blank lines separate stanzas.

   For instead of making\\
   Them do \emph{all} line breaking, \\
   It allows them to put too many words on a line when they'd rather be
   forced to be terse.
\end{verse}

Mathematical formulas may also be displayed.  A
displayed formula
is
one-line long; multiline
formulas require special formatting instructions.
   \[  \ip{\Gamma}{\psi'} = x'' + y^{2} + z_{i}^{n}\]
Don't start a paragraph with a displayed equation,
nor make one a paragraph by itself.

\end{document}               % End of document.

\section{Bench Section 139}
% synthetic bench section 139
% This is a sample LaTeX input file.  (Version of 12 August 2004.)
%
% A '%' character causes TeX to ignore all remaining text on the line,
% and is used for comments like this one.

\documentclass{article}      % Specifies the document class

                             % The preamble begins here.
\title{An Example Document}  % Declares the document's title.
\author{Leslie Lamport}      % Declares the author's name.
\date{January 21, 1994}      % Deleting this command produces today's date.

\newcommand{\ip}[2]{(#1, #2)}
                             % Defines \ip{arg1}{arg2} to mean
                             % (arg1, arg2).

%\newcommand{\ip}[2]{\langle #1 | #2\rangle}
                             % This is an alternative definition of
                             % \ip that is commented out.

\begin{document}             % End of preamble and beginning of text.

\maketitle                   % Produces the title.

This is an example input file.  Comparing it with
the output it generates can show you how to
produce a simple document of your own.

\section{Ordinary Text}      % Produces section heading.  Lower-level
                             % sections are begun with similar
                             % \subsection and \subsubsection commands.

The ends  of words and sentences are marked
  by   spaces. It  doesn't matter how many
spaces    you type; one is as good as 100.  The
end of   a line counts as a space.

One   or more   blank lines denote the  end
of  a paragraph.

Since any number of consecutive spaces are treated
like a single one, the formatting of the input
file makes no difference to
      \LaTeX,                % The \LaTeX command generates the LaTeX logo.
but it makes a difference to you.  When you use
\LaTeX, making your input file as easy to read
as possible will be a great help as you write
your document and when you change it.  This sample
file shows how you can add comments to your own input
file.

Because printing is different from typewriting,
there are a number of things that you have to do
differently when preparing an input file than if
you were just typing the document directly.
Quotation marks like
       ``this''
have to be handled specially, as do quotes within
quotes:
       ``\,`this'            % \, separates the double and single quote.
        is what I just
        wrote, not  `that'\,''.

Dashes come in three sizes: an
       intra-word
dash, a medium dash for number ranges like
       1--2,
and a punctuation
       dash---like
this.

A sentence-ending space should be larger than the
space between words within a sentence.  You
sometimes have to type special commands in
conjunction with punctuation characters to get
this right, as in the following sentence.
       Gnats, gnus, etc.\ all  % `\ ' makes an inter-word space.
       begin with G\@.         % \@ marks end-of-sentence punctuation.
You should check the spaces after periods when
reading your output to make sure you haven't
forgotten any special cases.  Generating an
ellipsis
       \ldots\               % `\ ' is needed after `\ldots' because TeX
                             % ignores spaces after command names like \ldots
                             % made from \ + letters.
                             %
                             % Note how a `%' character causes TeX to ignore
                             % the end of the input line, so these blank lines
                             % do not start a new paragraph.
                             %
with the right spacing around the periods requires
a special command.

\LaTeX\ interprets some common characters as
commands, so you must type special commands to
generate them.  These characters include the
following:
       \$ \& \% \# \{ and \}.

In printing, text is usually emphasized with an
       \emph{italic}
type style.

\begin{em}
   A long segment of text can also be emphasized
   in this way.  Text within such a segment can be
   given \emph{additional} emphasis.
\end{em}

It is sometimes necessary to prevent \LaTeX\ from
breaking a line where it might otherwise do so.
This may be at a space, as between the ``Mr.''\ and
``Jones'' in
       ``Mr.~Jones'',        % ~ produces an unbreakable interword space.
or within a word---especially when the word is a
symbol like
       \mbox{\emph{itemnum}}
that makes little sense when hyphenated across
lines.

Footnotes\footnote{This is an example of a footnote.}
pose no problem.

\LaTeX\ is good at typesetting mathematical formulas
like
       \( x-3y + z = 7 \)
or
       \( a_{1} > x^{2n} + y^{2n} > x' \)
or
       \( \ip{A}{B} = \sum_{i} a_{i} b_{i} \).
The spaces you type in a formula are
ignored.  Remember that a letter like
       $x$                   % $ ... $  and  \( ... \)  are equivalent
is a formula when it denotes a mathematical
symbol, and it should be typed as one.

\section{Displayed Text}

Text is displayed by indenting it from the left
margin.  Quotations are commonly displayed.  There
are short quotations
\begin{quote}
   This is a short quotation.  It consists of a
   single paragraph of text.  See how it is formatted.
\end{quote}
and longer ones.
\begin{quotation}
   This is a longer quotation.  It consists of two
   paragraphs of text, neither of which are
   particularly interesting.

   This is the second paragraph of the quotation.  It
   is just as dull as the first paragraph.
\end{quotation}
Another frequently-displayed structure is a list.
The following is an example of an \emph{itemized}
list.
\begin{itemize}
   \item This is the first item of an itemized list.
         Each item in the list is marked with a ``tick''.
         You don't have to worry about what kind of tick
         mark is used.

   \item This is the second item of the list.  It
         contains another list nested inside it.  The inner
         list is an \emph{enumerated} list.
         \begin{enumerate}
            \item This is the first item of an enumerated
                  list that is nested within the itemized list.

            \item This is the second item of the inner list.
                  \LaTeX\ allows you to nest lists deeper than
                  you really should.
         \end{enumerate}
         This is the rest of the second item of the outer
         list.  It is no more interesting than any other
         part of the item.
   \item This is the third item of the list.
\end{itemize}
You can even display poetry.
\begin{verse}
   There is an environment
    for verse \\             % The \\ command separates lines
   Whose features some poets % within a stanza.
   will curse.

                             % One or more blank lines separate stanzas.

   For instead of making\\
   Them do \emph{all} line breaking, \\
   It allows them to put too many words on a line when they'd rather be
   forced to be terse.
\end{verse}

Mathematical formulas may also be displayed.  A
displayed formula
is
one-line long; multiline
formulas require special formatting instructions.
   \[  \ip{\Gamma}{\psi'} = x'' + y^{2} + z_{i}^{n}\]
Don't start a paragraph with a displayed equation,
nor make one a paragraph by itself.

\end{document}               % End of document.

\section{Bench Section 140}
% synthetic bench section 140
% This is a sample LaTeX input file.  (Version of 12 August 2004.)
%
% A '%' character causes TeX to ignore all remaining text on the line,
% and is used for comments like this one.

\documentclass{article}      % Specifies the document class

                             % The preamble begins here.
\title{An Example Document}  % Declares the document's title.
\author{Leslie Lamport}      % Declares the author's name.
\date{January 21, 1994}      % Deleting this command produces today's date.

\newcommand{\ip}[2]{(#1, #2)}
                             % Defines \ip{arg1}{arg2} to mean
                             % (arg1, arg2).

%\newcommand{\ip}[2]{\langle #1 | #2\rangle}
                             % This is an alternative definition of
                             % \ip that is commented out.

\begin{document}             % End of preamble and beginning of text.

\maketitle                   % Produces the title.

This is an example input file.  Comparing it with
the output it generates can show you how to
produce a simple document of your own.

\section{Ordinary Text}      % Produces section heading.  Lower-level
                             % sections are begun with similar
                             % \subsection and \subsubsection commands.

The ends  of words and sentences are marked
  by   spaces. It  doesn't matter how many
spaces    you type; one is as good as 100.  The
end of   a line counts as a space.

One   or more   blank lines denote the  end
of  a paragraph.

Since any number of consecutive spaces are treated
like a single one, the formatting of the input
file makes no difference to
      \LaTeX,                % The \LaTeX command generates the LaTeX logo.
but it makes a difference to you.  When you use
\LaTeX, making your input file as easy to read
as possible will be a great help as you write
your document and when you change it.  This sample
file shows how you can add comments to your own input
file.

Because printing is different from typewriting,
there are a number of things that you have to do
differently when preparing an input file than if
you were just typing the document directly.
Quotation marks like
       ``this''
have to be handled specially, as do quotes within
quotes:
       ``\,`this'            % \, separates the double and single quote.
        is what I just
        wrote, not  `that'\,''.

Dashes come in three sizes: an
       intra-word
dash, a medium dash for number ranges like
       1--2,
and a punctuation
       dash---like
this.

A sentence-ending space should be larger than the
space between words within a sentence.  You
sometimes have to type special commands in
conjunction with punctuation characters to get
this right, as in the following sentence.
       Gnats, gnus, etc.\ all  % `\ ' makes an inter-word space.
       begin with G\@.         % \@ marks end-of-sentence punctuation.
You should check the spaces after periods when
reading your output to make sure you haven't
forgotten any special cases.  Generating an
ellipsis
       \ldots\               % `\ ' is needed after `\ldots' because TeX
                             % ignores spaces after command names like \ldots
                             % made from \ + letters.
                             %
                             % Note how a `%' character causes TeX to ignore
                             % the end of the input line, so these blank lines
                             % do not start a new paragraph.
                             %
with the right spacing around the periods requires
a special command.

\LaTeX\ interprets some common characters as
commands, so you must type special commands to
generate them.  These characters include the
following:
       \$ \& \% \# \{ and \}.

In printing, text is usually emphasized with an
       \emph{italic}
type style.

\begin{em}
   A long segment of text can also be emphasized
   in this way.  Text within such a segment can be
   given \emph{additional} emphasis.
\end{em}

It is sometimes necessary to prevent \LaTeX\ from
breaking a line where it might otherwise do so.
This may be at a space, as between the ``Mr.''\ and
``Jones'' in
       ``Mr.~Jones'',        % ~ produces an unbreakable interword space.
or within a word---especially when the word is a
symbol like
       \mbox{\emph{itemnum}}
that makes little sense when hyphenated across
lines.

Footnotes\footnote{This is an example of a footnote.}
pose no problem.

\LaTeX\ is good at typesetting mathematical formulas
like
       \( x-3y + z = 7 \)
or
       \( a_{1} > x^{2n} + y^{2n} > x' \)
or
       \( \ip{A}{B} = \sum_{i} a_{i} b_{i} \).
The spaces you type in a formula are
ignored.  Remember that a letter like
       $x$                   % $ ... $  and  \( ... \)  are equivalent
is a formula when it denotes a mathematical
symbol, and it should be typed as one.

\section{Displayed Text}

Text is displayed by indenting it from the left
margin.  Quotations are commonly displayed.  There
are short quotations
\begin{quote}
   This is a short quotation.  It consists of a
   single paragraph of text.  See how it is formatted.
\end{quote}
and longer ones.
\begin{quotation}
   This is a longer quotation.  It consists of two
   paragraphs of text, neither of which are
   particularly interesting.

   This is the second paragraph of the quotation.  It
   is just as dull as the first paragraph.
\end{quotation}
Another frequently-displayed structure is a list.
The following is an example of an \emph{itemized}
list.
\begin{itemize}
   \item This is the first item of an itemized list.
         Each item in the list is marked with a ``tick''.
         You don't have to worry about what kind of tick
         mark is used.

   \item This is the second item of the list.  It
         contains another list nested inside it.  The inner
         list is an \emph{enumerated} list.
         \begin{enumerate}
            \item This is the first item of an enumerated
                  list that is nested within the itemized list.

            \item This is the second item of the inner list.
                  \LaTeX\ allows you to nest lists deeper than
                  you really should.
         \end{enumerate}
         This is the rest of the second item of the outer
         list.  It is no more interesting than any other
         part of the item.
   \item This is the third item of the list.
\end{itemize}
You can even display poetry.
\begin{verse}
   There is an environment
    for verse \\             % The \\ command separates lines
   Whose features some poets % within a stanza.
   will curse.

                             % One or more blank lines separate stanzas.

   For instead of making\\
   Them do \emph{all} line breaking, \\
   It allows them to put too many words on a line when they'd rather be
   forced to be terse.
\end{verse}

Mathematical formulas may also be displayed.  A
displayed formula
is
one-line long; multiline
formulas require special formatting instructions.
   \[  \ip{\Gamma}{\psi'} = x'' + y^{2} + z_{i}^{n}\]
Don't start a paragraph with a displayed equation,
nor make one a paragraph by itself.

\end{document}               % End of document.

\section{Bench Section 141}
% synthetic bench section 141
% This is a sample LaTeX input file.  (Version of 12 August 2004.)
%
% A '%' character causes TeX to ignore all remaining text on the line,
% and is used for comments like this one.

\documentclass{article}      % Specifies the document class

                             % The preamble begins here.
\title{An Example Document}  % Declares the document's title.
\author{Leslie Lamport}      % Declares the author's name.
\date{January 21, 1994}      % Deleting this command produces today's date.

\newcommand{\ip}[2]{(#1, #2)}
                             % Defines \ip{arg1}{arg2} to mean
                             % (arg1, arg2).

%\newcommand{\ip}[2]{\langle #1 | #2\rangle}
                             % This is an alternative definition of
                             % \ip that is commented out.

\begin{document}             % End of preamble and beginning of text.

\maketitle                   % Produces the title.

This is an example input file.  Comparing it with
the output it generates can show you how to
produce a simple document of your own.

\section{Ordinary Text}      % Produces section heading.  Lower-level
                             % sections are begun with similar
                             % \subsection and \subsubsection commands.

The ends  of words and sentences are marked
  by   spaces. It  doesn't matter how many
spaces    you type; one is as good as 100.  The
end of   a line counts as a space.

One   or more   blank lines denote the  end
of  a paragraph.

Since any number of consecutive spaces are treated
like a single one, the formatting of the input
file makes no difference to
      \LaTeX,                % The \LaTeX command generates the LaTeX logo.
but it makes a difference to you.  When you use
\LaTeX, making your input file as easy to read
as possible will be a great help as you write
your document and when you change it.  This sample
file shows how you can add comments to your own input
file.

Because printing is different from typewriting,
there are a number of things that you have to do
differently when preparing an input file than if
you were just typing the document directly.
Quotation marks like
       ``this''
have to be handled specially, as do quotes within
quotes:
       ``\,`this'            % \, separates the double and single quote.
        is what I just
        wrote, not  `that'\,''.

Dashes come in three sizes: an
       intra-word
dash, a medium dash for number ranges like
       1--2,
and a punctuation
       dash---like
this.

A sentence-ending space should be larger than the
space between words within a sentence.  You
sometimes have to type special commands in
conjunction with punctuation characters to get
this right, as in the following sentence.
       Gnats, gnus, etc.\ all  % `\ ' makes an inter-word space.
       begin with G\@.         % \@ marks end-of-sentence punctuation.
You should check the spaces after periods when
reading your output to make sure you haven't
forgotten any special cases.  Generating an
ellipsis
       \ldots\               % `\ ' is needed after `\ldots' because TeX
                             % ignores spaces after command names like \ldots
                             % made from \ + letters.
                             %
                             % Note how a `%' character causes TeX to ignore
                             % the end of the input line, so these blank lines
                             % do not start a new paragraph.
                             %
with the right spacing around the periods requires
a special command.

\LaTeX\ interprets some common characters as
commands, so you must type special commands to
generate them.  These characters include the
following:
       \$ \& \% \# \{ and \}.

In printing, text is usually emphasized with an
       \emph{italic}
type style.

\begin{em}
   A long segment of text can also be emphasized
   in this way.  Text within such a segment can be
   given \emph{additional} emphasis.
\end{em}

It is sometimes necessary to prevent \LaTeX\ from
breaking a line where it might otherwise do so.
This may be at a space, as between the ``Mr.''\ and
``Jones'' in
       ``Mr.~Jones'',        % ~ produces an unbreakable interword space.
or within a word---especially when the word is a
symbol like
       \mbox{\emph{itemnum}}
that makes little sense when hyphenated across
lines.

Footnotes\footnote{This is an example of a footnote.}
pose no problem.

\LaTeX\ is good at typesetting mathematical formulas
like
       \( x-3y + z = 7 \)
or
       \( a_{1} > x^{2n} + y^{2n} > x' \)
or
       \( \ip{A}{B} = \sum_{i} a_{i} b_{i} \).
The spaces you type in a formula are
ignored.  Remember that a letter like
       $x$                   % $ ... $  and  \( ... \)  are equivalent
is a formula when it denotes a mathematical
symbol, and it should be typed as one.

\section{Displayed Text}

Text is displayed by indenting it from the left
margin.  Quotations are commonly displayed.  There
are short quotations
\begin{quote}
   This is a short quotation.  It consists of a
   single paragraph of text.  See how it is formatted.
\end{quote}
and longer ones.
\begin{quotation}
   This is a longer quotation.  It consists of two
   paragraphs of text, neither of which are
   particularly interesting.

   This is the second paragraph of the quotation.  It
   is just as dull as the first paragraph.
\end{quotation}
Another frequently-displayed structure is a list.
The following is an example of an \emph{itemized}
list.
\begin{itemize}
   \item This is the first item of an itemized list.
         Each item in the list is marked with a ``tick''.
         You don't have to worry about what kind of tick
         mark is used.

   \item This is the second item of the list.  It
         contains another list nested inside it.  The inner
         list is an \emph{enumerated} list.
         \begin{enumerate}
            \item This is the first item of an enumerated
                  list that is nested within the itemized list.

            \item This is the second item of the inner list.
                  \LaTeX\ allows you to nest lists deeper than
                  you really should.
         \end{enumerate}
         This is the rest of the second item of the outer
         list.  It is no more interesting than any other
         part of the item.
   \item This is the third item of the list.
\end{itemize}
You can even display poetry.
\begin{verse}
   There is an environment
    for verse \\             % The \\ command separates lines
   Whose features some poets % within a stanza.
   will curse.

                             % One or more blank lines separate stanzas.

   For instead of making\\
   Them do \emph{all} line breaking, \\
   It allows them to put too many words on a line when they'd rather be
   forced to be terse.
\end{verse}

Mathematical formulas may also be displayed.  A
displayed formula
is
one-line long; multiline
formulas require special formatting instructions.
   \[  \ip{\Gamma}{\psi'} = x'' + y^{2} + z_{i}^{n}\]
Don't start a paragraph with a displayed equation,
nor make one a paragraph by itself.

\end{document}               % End of document.

\section{Bench Section 142}
% synthetic bench section 142
% This is a sample LaTeX input file.  (Version of 12 August 2004.)
%
% A '%' character causes TeX to ignore all remaining text on the line,
% and is used for comments like this one.

\documentclass{article}      % Specifies the document class

                             % The preamble begins here.
\title{An Example Document}  % Declares the document's title.
\author{Leslie Lamport}      % Declares the author's name.
\date{January 21, 1994}      % Deleting this command produces today's date.

\newcommand{\ip}[2]{(#1, #2)}
                             % Defines \ip{arg1}{arg2} to mean
                             % (arg1, arg2).

%\newcommand{\ip}[2]{\langle #1 | #2\rangle}
                             % This is an alternative definition of
                             % \ip that is commented out.

\begin{document}             % End of preamble and beginning of text.

\maketitle                   % Produces the title.

This is an example input file.  Comparing it with
the output it generates can show you how to
produce a simple document of your own.

\section{Ordinary Text}      % Produces section heading.  Lower-level
                             % sections are begun with similar
                             % \subsection and \subsubsection commands.

The ends  of words and sentences are marked
  by   spaces. It  doesn't matter how many
spaces    you type; one is as good as 100.  The
end of   a line counts as a space.

One   or more   blank lines denote the  end
of  a paragraph.

Since any number of consecutive spaces are treated
like a single one, the formatting of the input
file makes no difference to
      \LaTeX,                % The \LaTeX command generates the LaTeX logo.
but it makes a difference to you.  When you use
\LaTeX, making your input file as easy to read
as possible will be a great help as you write
your document and when you change it.  This sample
file shows how you can add comments to your own input
file.

Because printing is different from typewriting,
there are a number of things that you have to do
differently when preparing an input file than if
you were just typing the document directly.
Quotation marks like
       ``this''
have to be handled specially, as do quotes within
quotes:
       ``\,`this'            % \, separates the double and single quote.
        is what I just
        wrote, not  `that'\,''.

Dashes come in three sizes: an
       intra-word
dash, a medium dash for number ranges like
       1--2,
and a punctuation
       dash---like
this.

A sentence-ending space should be larger than the
space between words within a sentence.  You
sometimes have to type special commands in
conjunction with punctuation characters to get
this right, as in the following sentence.
       Gnats, gnus, etc.\ all  % `\ ' makes an inter-word space.
       begin with G\@.         % \@ marks end-of-sentence punctuation.
You should check the spaces after periods when
reading your output to make sure you haven't
forgotten any special cases.  Generating an
ellipsis
       \ldots\               % `\ ' is needed after `\ldots' because TeX
                             % ignores spaces after command names like \ldots
                             % made from \ + letters.
                             %
                             % Note how a `%' character causes TeX to ignore
                             % the end of the input line, so these blank lines
                             % do not start a new paragraph.
                             %
with the right spacing around the periods requires
a special command.

\LaTeX\ interprets some common characters as
commands, so you must type special commands to
generate them.  These characters include the
following:
       \$ \& \% \# \{ and \}.

In printing, text is usually emphasized with an
       \emph{italic}
type style.

\begin{em}
   A long segment of text can also be emphasized
   in this way.  Text within such a segment can be
   given \emph{additional} emphasis.
\end{em}

It is sometimes necessary to prevent \LaTeX\ from
breaking a line where it might otherwise do so.
This may be at a space, as between the ``Mr.''\ and
``Jones'' in
       ``Mr.~Jones'',        % ~ produces an unbreakable interword space.
or within a word---especially when the word is a
symbol like
       \mbox{\emph{itemnum}}
that makes little sense when hyphenated across
lines.

Footnotes\footnote{This is an example of a footnote.}
pose no problem.

\LaTeX\ is good at typesetting mathematical formulas
like
       \( x-3y + z = 7 \)
or
       \( a_{1} > x^{2n} + y^{2n} > x' \)
or
       \( \ip{A}{B} = \sum_{i} a_{i} b_{i} \).
The spaces you type in a formula are
ignored.  Remember that a letter like
       $x$                   % $ ... $  and  \( ... \)  are equivalent
is a formula when it denotes a mathematical
symbol, and it should be typed as one.

\section{Displayed Text}

Text is displayed by indenting it from the left
margin.  Quotations are commonly displayed.  There
are short quotations
\begin{quote}
   This is a short quotation.  It consists of a
   single paragraph of text.  See how it is formatted.
\end{quote}
and longer ones.
\begin{quotation}
   This is a longer quotation.  It consists of two
   paragraphs of text, neither of which are
   particularly interesting.

   This is the second paragraph of the quotation.  It
   is just as dull as the first paragraph.
\end{quotation}
Another frequently-displayed structure is a list.
The following is an example of an \emph{itemized}
list.
\begin{itemize}
   \item This is the first item of an itemized list.
         Each item in the list is marked with a ``tick''.
         You don't have to worry about what kind of tick
         mark is used.

   \item This is the second item of the list.  It
         contains another list nested inside it.  The inner
         list is an \emph{enumerated} list.
         \begin{enumerate}
            \item This is the first item of an enumerated
                  list that is nested within the itemized list.

            \item This is the second item of the inner list.
                  \LaTeX\ allows you to nest lists deeper than
                  you really should.
         \end{enumerate}
         This is the rest of the second item of the outer
         list.  It is no more interesting than any other
         part of the item.
   \item This is the third item of the list.
\end{itemize}
You can even display poetry.
\begin{verse}
   There is an environment
    for verse \\             % The \\ command separates lines
   Whose features some poets % within a stanza.
   will curse.

                             % One or more blank lines separate stanzas.

   For instead of making\\
   Them do \emph{all} line breaking, \\
   It allows them to put too many words on a line when they'd rather be
   forced to be terse.
\end{verse}

Mathematical formulas may also be displayed.  A
displayed formula
is
one-line long; multiline
formulas require special formatting instructions.
   \[  \ip{\Gamma}{\psi'} = x'' + y^{2} + z_{i}^{n}\]
Don't start a paragraph with a displayed equation,
nor make one a paragraph by itself.

\end{document}               % End of document.

\section{Bench Section 143}
% synthetic bench section 143
% This is a sample LaTeX input file.  (Version of 12 August 2004.)
%
% A '%' character causes TeX to ignore all remaining text on the line,
% and is used for comments like this one.

\documentclass{article}      % Specifies the document class

                             % The preamble begins here.
\title{An Example Document}  % Declares the document's title.
\author{Leslie Lamport}      % Declares the author's name.
\date{January 21, 1994}      % Deleting this command produces today's date.

\newcommand{\ip}[2]{(#1, #2)}
                             % Defines \ip{arg1}{arg2} to mean
                             % (arg1, arg2).

%\newcommand{\ip}[2]{\langle #1 | #2\rangle}
                             % This is an alternative definition of
                             % \ip that is commented out.

\begin{document}             % End of preamble and beginning of text.

\maketitle                   % Produces the title.

This is an example input file.  Comparing it with
the output it generates can show you how to
produce a simple document of your own.

\section{Ordinary Text}      % Produces section heading.  Lower-level
                             % sections are begun with similar
                             % \subsection and \subsubsection commands.

The ends  of words and sentences are marked
  by   spaces. It  doesn't matter how many
spaces    you type; one is as good as 100.  The
end of   a line counts as a space.

One   or more   blank lines denote the  end
of  a paragraph.

Since any number of consecutive spaces are treated
like a single one, the formatting of the input
file makes no difference to
      \LaTeX,                % The \LaTeX command generates the LaTeX logo.
but it makes a difference to you.  When you use
\LaTeX, making your input file as easy to read
as possible will be a great help as you write
your document and when you change it.  This sample
file shows how you can add comments to your own input
file.

Because printing is different from typewriting,
there are a number of things that you have to do
differently when preparing an input file than if
you were just typing the document directly.
Quotation marks like
       ``this''
have to be handled specially, as do quotes within
quotes:
       ``\,`this'            % \, separates the double and single quote.
        is what I just
        wrote, not  `that'\,''.

Dashes come in three sizes: an
       intra-word
dash, a medium dash for number ranges like
       1--2,
and a punctuation
       dash---like
this.

A sentence-ending space should be larger than the
space between words within a sentence.  You
sometimes have to type special commands in
conjunction with punctuation characters to get
this right, as in the following sentence.
       Gnats, gnus, etc.\ all  % `\ ' makes an inter-word space.
       begin with G\@.         % \@ marks end-of-sentence punctuation.
You should check the spaces after periods when
reading your output to make sure you haven't
forgotten any special cases.  Generating an
ellipsis
       \ldots\               % `\ ' is needed after `\ldots' because TeX
                             % ignores spaces after command names like \ldots
                             % made from \ + letters.
                             %
                             % Note how a `%' character causes TeX to ignore
                             % the end of the input line, so these blank lines
                             % do not start a new paragraph.
                             %
with the right spacing around the periods requires
a special command.

\LaTeX\ interprets some common characters as
commands, so you must type special commands to
generate them.  These characters include the
following:
       \$ \& \% \# \{ and \}.

In printing, text is usually emphasized with an
       \emph{italic}
type style.

\begin{em}
   A long segment of text can also be emphasized
   in this way.  Text within such a segment can be
   given \emph{additional} emphasis.
\end{em}

It is sometimes necessary to prevent \LaTeX\ from
breaking a line where it might otherwise do so.
This may be at a space, as between the ``Mr.''\ and
``Jones'' in
       ``Mr.~Jones'',        % ~ produces an unbreakable interword space.
or within a word---especially when the word is a
symbol like
       \mbox{\emph{itemnum}}
that makes little sense when hyphenated across
lines.

Footnotes\footnote{This is an example of a footnote.}
pose no problem.

\LaTeX\ is good at typesetting mathematical formulas
like
       \( x-3y + z = 7 \)
or
       \( a_{1} > x^{2n} + y^{2n} > x' \)
or
       \( \ip{A}{B} = \sum_{i} a_{i} b_{i} \).
The spaces you type in a formula are
ignored.  Remember that a letter like
       $x$                   % $ ... $  and  \( ... \)  are equivalent
is a formula when it denotes a mathematical
symbol, and it should be typed as one.

\section{Displayed Text}

Text is displayed by indenting it from the left
margin.  Quotations are commonly displayed.  There
are short quotations
\begin{quote}
   This is a short quotation.  It consists of a
   single paragraph of text.  See how it is formatted.
\end{quote}
and longer ones.
\begin{quotation}
   This is a longer quotation.  It consists of two
   paragraphs of text, neither of which are
   particularly interesting.

   This is the second paragraph of the quotation.  It
   is just as dull as the first paragraph.
\end{quotation}
Another frequently-displayed structure is a list.
The following is an example of an \emph{itemized}
list.
\begin{itemize}
   \item This is the first item of an itemized list.
         Each item in the list is marked with a ``tick''.
         You don't have to worry about what kind of tick
         mark is used.

   \item This is the second item of the list.  It
         contains another list nested inside it.  The inner
         list is an \emph{enumerated} list.
         \begin{enumerate}
            \item This is the first item of an enumerated
                  list that is nested within the itemized list.

            \item This is the second item of the inner list.
                  \LaTeX\ allows you to nest lists deeper than
                  you really should.
         \end{enumerate}
         This is the rest of the second item of the outer
         list.  It is no more interesting than any other
         part of the item.
   \item This is the third item of the list.
\end{itemize}
You can even display poetry.
\begin{verse}
   There is an environment
    for verse \\             % The \\ command separates lines
   Whose features some poets % within a stanza.
   will curse.

                             % One or more blank lines separate stanzas.

   For instead of making\\
   Them do \emph{all} line breaking, \\
   It allows them to put too many words on a line when they'd rather be
   forced to be terse.
\end{verse}

Mathematical formulas may also be displayed.  A
displayed formula
is
one-line long; multiline
formulas require special formatting instructions.
   \[  \ip{\Gamma}{\psi'} = x'' + y^{2} + z_{i}^{n}\]
Don't start a paragraph with a displayed equation,
nor make one a paragraph by itself.

\end{document}               % End of document.

\section{Bench Section 144}
% synthetic bench section 144
% This is a sample LaTeX input file.  (Version of 12 August 2004.)
%
% A '%' character causes TeX to ignore all remaining text on the line,
% and is used for comments like this one.

\documentclass{article}      % Specifies the document class

                             % The preamble begins here.
\title{An Example Document}  % Declares the document's title.
\author{Leslie Lamport}      % Declares the author's name.
\date{January 21, 1994}      % Deleting this command produces today's date.

\newcommand{\ip}[2]{(#1, #2)}
                             % Defines \ip{arg1}{arg2} to mean
                             % (arg1, arg2).

%\newcommand{\ip}[2]{\langle #1 | #2\rangle}
                             % This is an alternative definition of
                             % \ip that is commented out.

\begin{document}             % End of preamble and beginning of text.

\maketitle                   % Produces the title.

This is an example input file.  Comparing it with
the output it generates can show you how to
produce a simple document of your own.

\section{Ordinary Text}      % Produces section heading.  Lower-level
                             % sections are begun with similar
                             % \subsection and \subsubsection commands.

The ends  of words and sentences are marked
  by   spaces. It  doesn't matter how many
spaces    you type; one is as good as 100.  The
end of   a line counts as a space.

One   or more   blank lines denote the  end
of  a paragraph.

Since any number of consecutive spaces are treated
like a single one, the formatting of the input
file makes no difference to
      \LaTeX,                % The \LaTeX command generates the LaTeX logo.
but it makes a difference to you.  When you use
\LaTeX, making your input file as easy to read
as possible will be a great help as you write
your document and when you change it.  This sample
file shows how you can add comments to your own input
file.

Because printing is different from typewriting,
there are a number of things that you have to do
differently when preparing an input file than if
you were just typing the document directly.
Quotation marks like
       ``this''
have to be handled specially, as do quotes within
quotes:
       ``\,`this'            % \, separates the double and single quote.
        is what I just
        wrote, not  `that'\,''.

Dashes come in three sizes: an
       intra-word
dash, a medium dash for number ranges like
       1--2,
and a punctuation
       dash---like
this.

A sentence-ending space should be larger than the
space between words within a sentence.  You
sometimes have to type special commands in
conjunction with punctuation characters to get
this right, as in the following sentence.
       Gnats, gnus, etc.\ all  % `\ ' makes an inter-word space.
       begin with G\@.         % \@ marks end-of-sentence punctuation.
You should check the spaces after periods when
reading your output to make sure you haven't
forgotten any special cases.  Generating an
ellipsis
       \ldots\               % `\ ' is needed after `\ldots' because TeX
                             % ignores spaces after command names like \ldots
                             % made from \ + letters.
                             %
                             % Note how a `%' character causes TeX to ignore
                             % the end of the input line, so these blank lines
                             % do not start a new paragraph.
                             %
with the right spacing around the periods requires
a special command.

\LaTeX\ interprets some common characters as
commands, so you must type special commands to
generate them.  These characters include the
following:
       \$ \& \% \# \{ and \}.

In printing, text is usually emphasized with an
       \emph{italic}
type style.

\begin{em}
   A long segment of text can also be emphasized
   in this way.  Text within such a segment can be
   given \emph{additional} emphasis.
\end{em}

It is sometimes necessary to prevent \LaTeX\ from
breaking a line where it might otherwise do so.
This may be at a space, as between the ``Mr.''\ and
``Jones'' in
       ``Mr.~Jones'',        % ~ produces an unbreakable interword space.
or within a word---especially when the word is a
symbol like
       \mbox{\emph{itemnum}}
that makes little sense when hyphenated across
lines.

Footnotes\footnote{This is an example of a footnote.}
pose no problem.

\LaTeX\ is good at typesetting mathematical formulas
like
       \( x-3y + z = 7 \)
or
       \( a_{1} > x^{2n} + y^{2n} > x' \)
or
       \( \ip{A}{B} = \sum_{i} a_{i} b_{i} \).
The spaces you type in a formula are
ignored.  Remember that a letter like
       $x$                   % $ ... $  and  \( ... \)  are equivalent
is a formula when it denotes a mathematical
symbol, and it should be typed as one.

\section{Displayed Text}

Text is displayed by indenting it from the left
margin.  Quotations are commonly displayed.  There
are short quotations
\begin{quote}
   This is a short quotation.  It consists of a
   single paragraph of text.  See how it is formatted.
\end{quote}
and longer ones.
\begin{quotation}
   This is a longer quotation.  It consists of two
   paragraphs of text, neither of which are
   particularly interesting.

   This is the second paragraph of the quotation.  It
   is just as dull as the first paragraph.
\end{quotation}
Another frequently-displayed structure is a list.
The following is an example of an \emph{itemized}
list.
\begin{itemize}
   \item This is the first item of an itemized list.
         Each item in the list is marked with a ``tick''.
         You don't have to worry about what kind of tick
         mark is used.

   \item This is the second item of the list.  It
         contains another list nested inside it.  The inner
         list is an \emph{enumerated} list.
         \begin{enumerate}
            \item This is the first item of an enumerated
                  list that is nested within the itemized list.

            \item This is the second item of the inner list.
                  \LaTeX\ allows you to nest lists deeper than
                  you really should.
         \end{enumerate}
         This is the rest of the second item of the outer
         list.  It is no more interesting than any other
         part of the item.
   \item This is the third item of the list.
\end{itemize}
You can even display poetry.
\begin{verse}
   There is an environment
    for verse \\             % The \\ command separates lines
   Whose features some poets % within a stanza.
   will curse.

                             % One or more blank lines separate stanzas.

   For instead of making\\
   Them do \emph{all} line breaking, \\
   It allows them to put too many words on a line when they'd rather be
   forced to be terse.
\end{verse}

Mathematical formulas may also be displayed.  A
displayed formula
is
one-line long; multiline
formulas require special formatting instructions.
   \[  \ip{\Gamma}{\psi'} = x'' + y^{2} + z_{i}^{n}\]
Don't start a paragraph with a displayed equation,
nor make one a paragraph by itself.

\end{document}               % End of document.

\section{Bench Section 145}
% synthetic bench section 145
% This is a sample LaTeX input file.  (Version of 12 August 2004.)
%
% A '%' character causes TeX to ignore all remaining text on the line,
% and is used for comments like this one.

\documentclass{article}      % Specifies the document class

                             % The preamble begins here.
\title{An Example Document}  % Declares the document's title.
\author{Leslie Lamport}      % Declares the author's name.
\date{January 21, 1994}      % Deleting this command produces today's date.

\newcommand{\ip}[2]{(#1, #2)}
                             % Defines \ip{arg1}{arg2} to mean
                             % (arg1, arg2).

%\newcommand{\ip}[2]{\langle #1 | #2\rangle}
                             % This is an alternative definition of
                             % \ip that is commented out.

\begin{document}             % End of preamble and beginning of text.

\maketitle                   % Produces the title.

This is an example input file.  Comparing it with
the output it generates can show you how to
produce a simple document of your own.

\section{Ordinary Text}      % Produces section heading.  Lower-level
                             % sections are begun with similar
                             % \subsection and \subsubsection commands.

The ends  of words and sentences are marked
  by   spaces. It  doesn't matter how many
spaces    you type; one is as good as 100.  The
end of   a line counts as a space.

One   or more   blank lines denote the  end
of  a paragraph.

Since any number of consecutive spaces are treated
like a single one, the formatting of the input
file makes no difference to
      \LaTeX,                % The \LaTeX command generates the LaTeX logo.
but it makes a difference to you.  When you use
\LaTeX, making your input file as easy to read
as possible will be a great help as you write
your document and when you change it.  This sample
file shows how you can add comments to your own input
file.

Because printing is different from typewriting,
there are a number of things that you have to do
differently when preparing an input file than if
you were just typing the document directly.
Quotation marks like
       ``this''
have to be handled specially, as do quotes within
quotes:
       ``\,`this'            % \, separates the double and single quote.
        is what I just
        wrote, not  `that'\,''.

Dashes come in three sizes: an
       intra-word
dash, a medium dash for number ranges like
       1--2,
and a punctuation
       dash---like
this.

A sentence-ending space should be larger than the
space between words within a sentence.  You
sometimes have to type special commands in
conjunction with punctuation characters to get
this right, as in the following sentence.
       Gnats, gnus, etc.\ all  % `\ ' makes an inter-word space.
       begin with G\@.         % \@ marks end-of-sentence punctuation.
You should check the spaces after periods when
reading your output to make sure you haven't
forgotten any special cases.  Generating an
ellipsis
       \ldots\               % `\ ' is needed after `\ldots' because TeX
                             % ignores spaces after command names like \ldots
                             % made from \ + letters.
                             %
                             % Note how a `%' character causes TeX to ignore
                             % the end of the input line, so these blank lines
                             % do not start a new paragraph.
                             %
with the right spacing around the periods requires
a special command.

\LaTeX\ interprets some common characters as
commands, so you must type special commands to
generate them.  These characters include the
following:
       \$ \& \% \# \{ and \}.

In printing, text is usually emphasized with an
       \emph{italic}
type style.

\begin{em}
   A long segment of text can also be emphasized
   in this way.  Text within such a segment can be
   given \emph{additional} emphasis.
\end{em}

It is sometimes necessary to prevent \LaTeX\ from
breaking a line where it might otherwise do so.
This may be at a space, as between the ``Mr.''\ and
``Jones'' in
       ``Mr.~Jones'',        % ~ produces an unbreakable interword space.
or within a word---especially when the word is a
symbol like
       \mbox{\emph{itemnum}}
that makes little sense when hyphenated across
lines.

Footnotes\footnote{This is an example of a footnote.}
pose no problem.

\LaTeX\ is good at typesetting mathematical formulas
like
       \( x-3y + z = 7 \)
or
       \( a_{1} > x^{2n} + y^{2n} > x' \)
or
       \( \ip{A}{B} = \sum_{i} a_{i} b_{i} \).
The spaces you type in a formula are
ignored.  Remember that a letter like
       $x$                   % $ ... $  and  \( ... \)  are equivalent
is a formula when it denotes a mathematical
symbol, and it should be typed as one.

\section{Displayed Text}

Text is displayed by indenting it from the left
margin.  Quotations are commonly displayed.  There
are short quotations
\begin{quote}
   This is a short quotation.  It consists of a
   single paragraph of text.  See how it is formatted.
\end{quote}
and longer ones.
\begin{quotation}
   This is a longer quotation.  It consists of two
   paragraphs of text, neither of which are
   particularly interesting.

   This is the second paragraph of the quotation.  It
   is just as dull as the first paragraph.
\end{quotation}
Another frequently-displayed structure is a list.
The following is an example of an \emph{itemized}
list.
\begin{itemize}
   \item This is the first item of an itemized list.
         Each item in the list is marked with a ``tick''.
         You don't have to worry about what kind of tick
         mark is used.

   \item This is the second item of the list.  It
         contains another list nested inside it.  The inner
         list is an \emph{enumerated} list.
         \begin{enumerate}
            \item This is the first item of an enumerated
                  list that is nested within the itemized list.

            \item This is the second item of the inner list.
                  \LaTeX\ allows you to nest lists deeper than
                  you really should.
         \end{enumerate}
         This is the rest of the second item of the outer
         list.  It is no more interesting than any other
         part of the item.
   \item This is the third item of the list.
\end{itemize}
You can even display poetry.
\begin{verse}
   There is an environment
    for verse \\             % The \\ command separates lines
   Whose features some poets % within a stanza.
   will curse.

                             % One or more blank lines separate stanzas.

   For instead of making\\
   Them do \emph{all} line breaking, \\
   It allows them to put too many words on a line when they'd rather be
   forced to be terse.
\end{verse}

Mathematical formulas may also be displayed.  A
displayed formula
is
one-line long; multiline
formulas require special formatting instructions.
   \[  \ip{\Gamma}{\psi'} = x'' + y^{2} + z_{i}^{n}\]
Don't start a paragraph with a displayed equation,
nor make one a paragraph by itself.

\end{document}               % End of document.

\section{Bench Section 146}
% synthetic bench section 146
% This is a sample LaTeX input file.  (Version of 12 August 2004.)
%
% A '%' character causes TeX to ignore all remaining text on the line,
% and is used for comments like this one.

\documentclass{article}      % Specifies the document class

                             % The preamble begins here.
\title{An Example Document}  % Declares the document's title.
\author{Leslie Lamport}      % Declares the author's name.
\date{January 21, 1994}      % Deleting this command produces today's date.

\newcommand{\ip}[2]{(#1, #2)}
                             % Defines \ip{arg1}{arg2} to mean
                             % (arg1, arg2).

%\newcommand{\ip}[2]{\langle #1 | #2\rangle}
                             % This is an alternative definition of
                             % \ip that is commented out.

\begin{document}             % End of preamble and beginning of text.

\maketitle                   % Produces the title.

This is an example input file.  Comparing it with
the output it generates can show you how to
produce a simple document of your own.

\section{Ordinary Text}      % Produces section heading.  Lower-level
                             % sections are begun with similar
                             % \subsection and \subsubsection commands.

The ends  of words and sentences are marked
  by   spaces. It  doesn't matter how many
spaces    you type; one is as good as 100.  The
end of   a line counts as a space.

One   or more   blank lines denote the  end
of  a paragraph.

Since any number of consecutive spaces are treated
like a single one, the formatting of the input
file makes no difference to
      \LaTeX,                % The \LaTeX command generates the LaTeX logo.
but it makes a difference to you.  When you use
\LaTeX, making your input file as easy to read
as possible will be a great help as you write
your document and when you change it.  This sample
file shows how you can add comments to your own input
file.

Because printing is different from typewriting,
there are a number of things that you have to do
differently when preparing an input file than if
you were just typing the document directly.
Quotation marks like
       ``this''
have to be handled specially, as do quotes within
quotes:
       ``\,`this'            % \, separates the double and single quote.
        is what I just
        wrote, not  `that'\,''.

Dashes come in three sizes: an
       intra-word
dash, a medium dash for number ranges like
       1--2,
and a punctuation
       dash---like
this.

A sentence-ending space should be larger than the
space between words within a sentence.  You
sometimes have to type special commands in
conjunction with punctuation characters to get
this right, as in the following sentence.
       Gnats, gnus, etc.\ all  % `\ ' makes an inter-word space.
       begin with G\@.         % \@ marks end-of-sentence punctuation.
You should check the spaces after periods when
reading your output to make sure you haven't
forgotten any special cases.  Generating an
ellipsis
       \ldots\               % `\ ' is needed after `\ldots' because TeX
                             % ignores spaces after command names like \ldots
                             % made from \ + letters.
                             %
                             % Note how a `%' character causes TeX to ignore
                             % the end of the input line, so these blank lines
                             % do not start a new paragraph.
                             %
with the right spacing around the periods requires
a special command.

\LaTeX\ interprets some common characters as
commands, so you must type special commands to
generate them.  These characters include the
following:
       \$ \& \% \# \{ and \}.

In printing, text is usually emphasized with an
       \emph{italic}
type style.

\begin{em}
   A long segment of text can also be emphasized
   in this way.  Text within such a segment can be
   given \emph{additional} emphasis.
\end{em}

It is sometimes necessary to prevent \LaTeX\ from
breaking a line where it might otherwise do so.
This may be at a space, as between the ``Mr.''\ and
``Jones'' in
       ``Mr.~Jones'',        % ~ produces an unbreakable interword space.
or within a word---especially when the word is a
symbol like
       \mbox{\emph{itemnum}}
that makes little sense when hyphenated across
lines.

Footnotes\footnote{This is an example of a footnote.}
pose no problem.

\LaTeX\ is good at typesetting mathematical formulas
like
       \( x-3y + z = 7 \)
or
       \( a_{1} > x^{2n} + y^{2n} > x' \)
or
       \( \ip{A}{B} = \sum_{i} a_{i} b_{i} \).
The spaces you type in a formula are
ignored.  Remember that a letter like
       $x$                   % $ ... $  and  \( ... \)  are equivalent
is a formula when it denotes a mathematical
symbol, and it should be typed as one.

\section{Displayed Text}

Text is displayed by indenting it from the left
margin.  Quotations are commonly displayed.  There
are short quotations
\begin{quote}
   This is a short quotation.  It consists of a
   single paragraph of text.  See how it is formatted.
\end{quote}
and longer ones.
\begin{quotation}
   This is a longer quotation.  It consists of two
   paragraphs of text, neither of which are
   particularly interesting.

   This is the second paragraph of the quotation.  It
   is just as dull as the first paragraph.
\end{quotation}
Another frequently-displayed structure is a list.
The following is an example of an \emph{itemized}
list.
\begin{itemize}
   \item This is the first item of an itemized list.
         Each item in the list is marked with a ``tick''.
         You don't have to worry about what kind of tick
         mark is used.

   \item This is the second item of the list.  It
         contains another list nested inside it.  The inner
         list is an \emph{enumerated} list.
         \begin{enumerate}
            \item This is the first item of an enumerated
                  list that is nested within the itemized list.

            \item This is the second item of the inner list.
                  \LaTeX\ allows you to nest lists deeper than
                  you really should.
         \end{enumerate}
         This is the rest of the second item of the outer
         list.  It is no more interesting than any other
         part of the item.
   \item This is the third item of the list.
\end{itemize}
You can even display poetry.
\begin{verse}
   There is an environment
    for verse \\             % The \\ command separates lines
   Whose features some poets % within a stanza.
   will curse.

                             % One or more blank lines separate stanzas.

   For instead of making\\
   Them do \emph{all} line breaking, \\
   It allows them to put too many words on a line when they'd rather be
   forced to be terse.
\end{verse}

Mathematical formulas may also be displayed.  A
displayed formula
is
one-line long; multiline
formulas require special formatting instructions.
   \[  \ip{\Gamma}{\psi'} = x'' + y^{2} + z_{i}^{n}\]
Don't start a paragraph with a displayed equation,
nor make one a paragraph by itself.

\end{document}               % End of document.

\section{Bench Section 147}
% synthetic bench section 147
% This is a sample LaTeX input file.  (Version of 12 August 2004.)
%
% A '%' character causes TeX to ignore all remaining text on the line,
% and is used for comments like this one.

\documentclass{article}      % Specifies the document class

                             % The preamble begins here.
\title{An Example Document}  % Declares the document's title.
\author{Leslie Lamport}      % Declares the author's name.
\date{January 21, 1994}      % Deleting this command produces today's date.

\newcommand{\ip}[2]{(#1, #2)}
                             % Defines \ip{arg1}{arg2} to mean
                             % (arg1, arg2).

%\newcommand{\ip}[2]{\langle #1 | #2\rangle}
                             % This is an alternative definition of
                             % \ip that is commented out.

\begin{document}             % End of preamble and beginning of text.

\maketitle                   % Produces the title.

This is an example input file.  Comparing it with
the output it generates can show you how to
produce a simple document of your own.

\section{Ordinary Text}      % Produces section heading.  Lower-level
                             % sections are begun with similar
                             % \subsection and \subsubsection commands.

The ends  of words and sentences are marked
  by   spaces. It  doesn't matter how many
spaces    you type; one is as good as 100.  The
end of   a line counts as a space.

One   or more   blank lines denote the  end
of  a paragraph.

Since any number of consecutive spaces are treated
like a single one, the formatting of the input
file makes no difference to
      \LaTeX,                % The \LaTeX command generates the LaTeX logo.
but it makes a difference to you.  When you use
\LaTeX, making your input file as easy to read
as possible will be a great help as you write
your document and when you change it.  This sample
file shows how you can add comments to your own input
file.

Because printing is different from typewriting,
there are a number of things that you have to do
differently when preparing an input file than if
you were just typing the document directly.
Quotation marks like
       ``this''
have to be handled specially, as do quotes within
quotes:
       ``\,`this'            % \, separates the double and single quote.
        is what I just
        wrote, not  `that'\,''.

Dashes come in three sizes: an
       intra-word
dash, a medium dash for number ranges like
       1--2,
and a punctuation
       dash---like
this.

A sentence-ending space should be larger than the
space between words within a sentence.  You
sometimes have to type special commands in
conjunction with punctuation characters to get
this right, as in the following sentence.
       Gnats, gnus, etc.\ all  % `\ ' makes an inter-word space.
       begin with G\@.         % \@ marks end-of-sentence punctuation.
You should check the spaces after periods when
reading your output to make sure you haven't
forgotten any special cases.  Generating an
ellipsis
       \ldots\               % `\ ' is needed after `\ldots' because TeX
                             % ignores spaces after command names like \ldots
                             % made from \ + letters.
                             %
                             % Note how a `%' character causes TeX to ignore
                             % the end of the input line, so these blank lines
                             % do not start a new paragraph.
                             %
with the right spacing around the periods requires
a special command.

\LaTeX\ interprets some common characters as
commands, so you must type special commands to
generate them.  These characters include the
following:
       \$ \& \% \# \{ and \}.

In printing, text is usually emphasized with an
       \emph{italic}
type style.

\begin{em}
   A long segment of text can also be emphasized
   in this way.  Text within such a segment can be
   given \emph{additional} emphasis.
\end{em}

It is sometimes necessary to prevent \LaTeX\ from
breaking a line where it might otherwise do so.
This may be at a space, as between the ``Mr.''\ and
``Jones'' in
       ``Mr.~Jones'',        % ~ produces an unbreakable interword space.
or within a word---especially when the word is a
symbol like
       \mbox{\emph{itemnum}}
that makes little sense when hyphenated across
lines.

Footnotes\footnote{This is an example of a footnote.}
pose no problem.

\LaTeX\ is good at typesetting mathematical formulas
like
       \( x-3y + z = 7 \)
or
       \( a_{1} > x^{2n} + y^{2n} > x' \)
or
       \( \ip{A}{B} = \sum_{i} a_{i} b_{i} \).
The spaces you type in a formula are
ignored.  Remember that a letter like
       $x$                   % $ ... $  and  \( ... \)  are equivalent
is a formula when it denotes a mathematical
symbol, and it should be typed as one.

\section{Displayed Text}

Text is displayed by indenting it from the left
margin.  Quotations are commonly displayed.  There
are short quotations
\begin{quote}
   This is a short quotation.  It consists of a
   single paragraph of text.  See how it is formatted.
\end{quote}
and longer ones.
\begin{quotation}
   This is a longer quotation.  It consists of two
   paragraphs of text, neither of which are
   particularly interesting.

   This is the second paragraph of the quotation.  It
   is just as dull as the first paragraph.
\end{quotation}
Another frequently-displayed structure is a list.
The following is an example of an \emph{itemized}
list.
\begin{itemize}
   \item This is the first item of an itemized list.
         Each item in the list is marked with a ``tick''.
         You don't have to worry about what kind of tick
         mark is used.

   \item This is the second item of the list.  It
         contains another list nested inside it.  The inner
         list is an \emph{enumerated} list.
         \begin{enumerate}
            \item This is the first item of an enumerated
                  list that is nested within the itemized list.

            \item This is the second item of the inner list.
                  \LaTeX\ allows you to nest lists deeper than
                  you really should.
         \end{enumerate}
         This is the rest of the second item of the outer
         list.  It is no more interesting than any other
         part of the item.
   \item This is the third item of the list.
\end{itemize}
You can even display poetry.
\begin{verse}
   There is an environment
    for verse \\             % The \\ command separates lines
   Whose features some poets % within a stanza.
   will curse.

                             % One or more blank lines separate stanzas.

   For instead of making\\
   Them do \emph{all} line breaking, \\
   It allows them to put too many words on a line when they'd rather be
   forced to be terse.
\end{verse}

Mathematical formulas may also be displayed.  A
displayed formula
is
one-line long; multiline
formulas require special formatting instructions.
   \[  \ip{\Gamma}{\psi'} = x'' + y^{2} + z_{i}^{n}\]
Don't start a paragraph with a displayed equation,
nor make one a paragraph by itself.

\end{document}               % End of document.

\section{Bench Section 148}
% synthetic bench section 148
% This is a sample LaTeX input file.  (Version of 12 August 2004.)
%
% A '%' character causes TeX to ignore all remaining text on the line,
% and is used for comments like this one.

\documentclass{article}      % Specifies the document class

                             % The preamble begins here.
\title{An Example Document}  % Declares the document's title.
\author{Leslie Lamport}      % Declares the author's name.
\date{January 21, 1994}      % Deleting this command produces today's date.

\newcommand{\ip}[2]{(#1, #2)}
                             % Defines \ip{arg1}{arg2} to mean
                             % (arg1, arg2).

%\newcommand{\ip}[2]{\langle #1 | #2\rangle}
                             % This is an alternative definition of
                             % \ip that is commented out.

\begin{document}             % End of preamble and beginning of text.

\maketitle                   % Produces the title.

This is an example input file.  Comparing it with
the output it generates can show you how to
produce a simple document of your own.

\section{Ordinary Text}      % Produces section heading.  Lower-level
                             % sections are begun with similar
                             % \subsection and \subsubsection commands.

The ends  of words and sentences are marked
  by   spaces. It  doesn't matter how many
spaces    you type; one is as good as 100.  The
end of   a line counts as a space.

One   or more   blank lines denote the  end
of  a paragraph.

Since any number of consecutive spaces are treated
like a single one, the formatting of the input
file makes no difference to
      \LaTeX,                % The \LaTeX command generates the LaTeX logo.
but it makes a difference to you.  When you use
\LaTeX, making your input file as easy to read
as possible will be a great help as you write
your document and when you change it.  This sample
file shows how you can add comments to your own input
file.

Because printing is different from typewriting,
there are a number of things that you have to do
differently when preparing an input file than if
you were just typing the document directly.
Quotation marks like
       ``this''
have to be handled specially, as do quotes within
quotes:
       ``\,`this'            % \, separates the double and single quote.
        is what I just
        wrote, not  `that'\,''.

Dashes come in three sizes: an
       intra-word
dash, a medium dash for number ranges like
       1--2,
and a punctuation
       dash---like
this.

A sentence-ending space should be larger than the
space between words within a sentence.  You
sometimes have to type special commands in
conjunction with punctuation characters to get
this right, as in the following sentence.
       Gnats, gnus, etc.\ all  % `\ ' makes an inter-word space.
       begin with G\@.         % \@ marks end-of-sentence punctuation.
You should check the spaces after periods when
reading your output to make sure you haven't
forgotten any special cases.  Generating an
ellipsis
       \ldots\               % `\ ' is needed after `\ldots' because TeX
                             % ignores spaces after command names like \ldots
                             % made from \ + letters.
                             %
                             % Note how a `%' character causes TeX to ignore
                             % the end of the input line, so these blank lines
                             % do not start a new paragraph.
                             %
with the right spacing around the periods requires
a special command.

\LaTeX\ interprets some common characters as
commands, so you must type special commands to
generate them.  These characters include the
following:
       \$ \& \% \# \{ and \}.

In printing, text is usually emphasized with an
       \emph{italic}
type style.

\begin{em}
   A long segment of text can also be emphasized
   in this way.  Text within such a segment can be
   given \emph{additional} emphasis.
\end{em}

It is sometimes necessary to prevent \LaTeX\ from
breaking a line where it might otherwise do so.
This may be at a space, as between the ``Mr.''\ and
``Jones'' in
       ``Mr.~Jones'',        % ~ produces an unbreakable interword space.
or within a word---especially when the word is a
symbol like
       \mbox{\emph{itemnum}}
that makes little sense when hyphenated across
lines.

Footnotes\footnote{This is an example of a footnote.}
pose no problem.

\LaTeX\ is good at typesetting mathematical formulas
like
       \( x-3y + z = 7 \)
or
       \( a_{1} > x^{2n} + y^{2n} > x' \)
or
       \( \ip{A}{B} = \sum_{i} a_{i} b_{i} \).
The spaces you type in a formula are
ignored.  Remember that a letter like
       $x$                   % $ ... $  and  \( ... \)  are equivalent
is a formula when it denotes a mathematical
symbol, and it should be typed as one.

\section{Displayed Text}

Text is displayed by indenting it from the left
margin.  Quotations are commonly displayed.  There
are short quotations
\begin{quote}
   This is a short quotation.  It consists of a
   single paragraph of text.  See how it is formatted.
\end{quote}
and longer ones.
\begin{quotation}
   This is a longer quotation.  It consists of two
   paragraphs of text, neither of which are
   particularly interesting.

   This is the second paragraph of the quotation.  It
   is just as dull as the first paragraph.
\end{quotation}
Another frequently-displayed structure is a list.
The following is an example of an \emph{itemized}
list.
\begin{itemize}
   \item This is the first item of an itemized list.
         Each item in the list is marked with a ``tick''.
         You don't have to worry about what kind of tick
         mark is used.

   \item This is the second item of the list.  It
         contains another list nested inside it.  The inner
         list is an \emph{enumerated} list.
         \begin{enumerate}
            \item This is the first item of an enumerated
                  list that is nested within the itemized list.

            \item This is the second item of the inner list.
                  \LaTeX\ allows you to nest lists deeper than
                  you really should.
         \end{enumerate}
         This is the rest of the second item of the outer
         list.  It is no more interesting than any other
         part of the item.
   \item This is the third item of the list.
\end{itemize}
You can even display poetry.
\begin{verse}
   There is an environment
    for verse \\             % The \\ command separates lines
   Whose features some poets % within a stanza.
   will curse.

                             % One or more blank lines separate stanzas.

   For instead of making\\
   Them do \emph{all} line breaking, \\
   It allows them to put too many words on a line when they'd rather be
   forced to be terse.
\end{verse}

Mathematical formulas may also be displayed.  A
displayed formula
is
one-line long; multiline
formulas require special formatting instructions.
   \[  \ip{\Gamma}{\psi'} = x'' + y^{2} + z_{i}^{n}\]
Don't start a paragraph with a displayed equation,
nor make one a paragraph by itself.

\end{document}               % End of document.

\section{Bench Section 149}
% synthetic bench section 149
% This is a sample LaTeX input file.  (Version of 12 August 2004.)
%
% A '%' character causes TeX to ignore all remaining text on the line,
% and is used for comments like this one.

\documentclass{article}      % Specifies the document class

                             % The preamble begins here.
\title{An Example Document}  % Declares the document's title.
\author{Leslie Lamport}      % Declares the author's name.
\date{January 21, 1994}      % Deleting this command produces today's date.

\newcommand{\ip}[2]{(#1, #2)}
                             % Defines \ip{arg1}{arg2} to mean
                             % (arg1, arg2).

%\newcommand{\ip}[2]{\langle #1 | #2\rangle}
                             % This is an alternative definition of
                             % \ip that is commented out.

\begin{document}             % End of preamble and beginning of text.

\maketitle                   % Produces the title.

This is an example input file.  Comparing it with
the output it generates can show you how to
produce a simple document of your own.

\section{Ordinary Text}      % Produces section heading.  Lower-level
                             % sections are begun with similar
                             % \subsection and \subsubsection commands.

The ends  of words and sentences are marked
  by   spaces. It  doesn't matter how many
spaces    you type; one is as good as 100.  The
end of   a line counts as a space.

One   or more   blank lines denote the  end
of  a paragraph.

Since any number of consecutive spaces are treated
like a single one, the formatting of the input
file makes no difference to
      \LaTeX,                % The \LaTeX command generates the LaTeX logo.
but it makes a difference to you.  When you use
\LaTeX, making your input file as easy to read
as possible will be a great help as you write
your document and when you change it.  This sample
file shows how you can add comments to your own input
file.

Because printing is different from typewriting,
there are a number of things that you have to do
differently when preparing an input file than if
you were just typing the document directly.
Quotation marks like
       ``this''
have to be handled specially, as do quotes within
quotes:
       ``\,`this'            % \, separates the double and single quote.
        is what I just
        wrote, not  `that'\,''.

Dashes come in three sizes: an
       intra-word
dash, a medium dash for number ranges like
       1--2,
and a punctuation
       dash---like
this.

A sentence-ending space should be larger than the
space between words within a sentence.  You
sometimes have to type special commands in
conjunction with punctuation characters to get
this right, as in the following sentence.
       Gnats, gnus, etc.\ all  % `\ ' makes an inter-word space.
       begin with G\@.         % \@ marks end-of-sentence punctuation.
You should check the spaces after periods when
reading your output to make sure you haven't
forgotten any special cases.  Generating an
ellipsis
       \ldots\               % `\ ' is needed after `\ldots' because TeX
                             % ignores spaces after command names like \ldots
                             % made from \ + letters.
                             %
                             % Note how a `%' character causes TeX to ignore
                             % the end of the input line, so these blank lines
                             % do not start a new paragraph.
                             %
with the right spacing around the periods requires
a special command.

\LaTeX\ interprets some common characters as
commands, so you must type special commands to
generate them.  These characters include the
following:
       \$ \& \% \# \{ and \}.

In printing, text is usually emphasized with an
       \emph{italic}
type style.

\begin{em}
   A long segment of text can also be emphasized
   in this way.  Text within such a segment can be
   given \emph{additional} emphasis.
\end{em}

It is sometimes necessary to prevent \LaTeX\ from
breaking a line where it might otherwise do so.
This may be at a space, as between the ``Mr.''\ and
``Jones'' in
       ``Mr.~Jones'',        % ~ produces an unbreakable interword space.
or within a word---especially when the word is a
symbol like
       \mbox{\emph{itemnum}}
that makes little sense when hyphenated across
lines.

Footnotes\footnote{This is an example of a footnote.}
pose no problem.

\LaTeX\ is good at typesetting mathematical formulas
like
       \( x-3y + z = 7 \)
or
       \( a_{1} > x^{2n} + y^{2n} > x' \)
or
       \( \ip{A}{B} = \sum_{i} a_{i} b_{i} \).
The spaces you type in a formula are
ignored.  Remember that a letter like
       $x$                   % $ ... $  and  \( ... \)  are equivalent
is a formula when it denotes a mathematical
symbol, and it should be typed as one.

\section{Displayed Text}

Text is displayed by indenting it from the left
margin.  Quotations are commonly displayed.  There
are short quotations
\begin{quote}
   This is a short quotation.  It consists of a
   single paragraph of text.  See how it is formatted.
\end{quote}
and longer ones.
\begin{quotation}
   This is a longer quotation.  It consists of two
   paragraphs of text, neither of which are
   particularly interesting.

   This is the second paragraph of the quotation.  It
   is just as dull as the first paragraph.
\end{quotation}
Another frequently-displayed structure is a list.
The following is an example of an \emph{itemized}
list.
\begin{itemize}
   \item This is the first item of an itemized list.
         Each item in the list is marked with a ``tick''.
         You don't have to worry about what kind of tick
         mark is used.

   \item This is the second item of the list.  It
         contains another list nested inside it.  The inner
         list is an \emph{enumerated} list.
         \begin{enumerate}
            \item This is the first item of an enumerated
                  list that is nested within the itemized list.

            \item This is the second item of the inner list.
                  \LaTeX\ allows you to nest lists deeper than
                  you really should.
         \end{enumerate}
         This is the rest of the second item of the outer
         list.  It is no more interesting than any other
         part of the item.
   \item This is the third item of the list.
\end{itemize}
You can even display poetry.
\begin{verse}
   There is an environment
    for verse \\             % The \\ command separates lines
   Whose features some poets % within a stanza.
   will curse.

                             % One or more blank lines separate stanzas.

   For instead of making\\
   Them do \emph{all} line breaking, \\
   It allows them to put too many words on a line when they'd rather be
   forced to be terse.
\end{verse}

Mathematical formulas may also be displayed.  A
displayed formula
is
one-line long; multiline
formulas require special formatting instructions.
   \[  \ip{\Gamma}{\psi'} = x'' + y^{2} + z_{i}^{n}\]
Don't start a paragraph with a displayed equation,
nor make one a paragraph by itself.

\end{document}               % End of document.

\section{Bench Section 150}
% synthetic bench section 150
% This is a sample LaTeX input file.  (Version of 12 August 2004.)
%
% A '%' character causes TeX to ignore all remaining text on the line,
% and is used for comments like this one.

\documentclass{article}      % Specifies the document class

                             % The preamble begins here.
\title{An Example Document}  % Declares the document's title.
\author{Leslie Lamport}      % Declares the author's name.
\date{January 21, 1994}      % Deleting this command produces today's date.

\newcommand{\ip}[2]{(#1, #2)}
                             % Defines \ip{arg1}{arg2} to mean
                             % (arg1, arg2).

%\newcommand{\ip}[2]{\langle #1 | #2\rangle}
                             % This is an alternative definition of
                             % \ip that is commented out.

\begin{document}             % End of preamble and beginning of text.

\maketitle                   % Produces the title.

This is an example input file.  Comparing it with
the output it generates can show you how to
produce a simple document of your own.

\section{Ordinary Text}      % Produces section heading.  Lower-level
                             % sections are begun with similar
                             % \subsection and \subsubsection commands.

The ends  of words and sentences are marked
  by   spaces. It  doesn't matter how many
spaces    you type; one is as good as 100.  The
end of   a line counts as a space.

One   or more   blank lines denote the  end
of  a paragraph.

Since any number of consecutive spaces are treated
like a single one, the formatting of the input
file makes no difference to
      \LaTeX,                % The \LaTeX command generates the LaTeX logo.
but it makes a difference to you.  When you use
\LaTeX, making your input file as easy to read
as possible will be a great help as you write
your document and when you change it.  This sample
file shows how you can add comments to your own input
file.

Because printing is different from typewriting,
there are a number of things that you have to do
differently when preparing an input file than if
you were just typing the document directly.
Quotation marks like
       ``this''
have to be handled specially, as do quotes within
quotes:
       ``\,`this'            % \, separates the double and single quote.
        is what I just
        wrote, not  `that'\,''.

Dashes come in three sizes: an
       intra-word
dash, a medium dash for number ranges like
       1--2,
and a punctuation
       dash---like
this.

A sentence-ending space should be larger than the
space between words within a sentence.  You
sometimes have to type special commands in
conjunction with punctuation characters to get
this right, as in the following sentence.
       Gnats, gnus, etc.\ all  % `\ ' makes an inter-word space.
       begin with G\@.         % \@ marks end-of-sentence punctuation.
You should check the spaces after periods when
reading your output to make sure you haven't
forgotten any special cases.  Generating an
ellipsis
       \ldots\               % `\ ' is needed after `\ldots' because TeX
                             % ignores spaces after command names like \ldots
                             % made from \ + letters.
                             %
                             % Note how a `%' character causes TeX to ignore
                             % the end of the input line, so these blank lines
                             % do not start a new paragraph.
                             %
with the right spacing around the periods requires
a special command.

\LaTeX\ interprets some common characters as
commands, so you must type special commands to
generate them.  These characters include the
following:
       \$ \& \% \# \{ and \}.

In printing, text is usually emphasized with an
       \emph{italic}
type style.

\begin{em}
   A long segment of text can also be emphasized
   in this way.  Text within such a segment can be
   given \emph{additional} emphasis.
\end{em}

It is sometimes necessary to prevent \LaTeX\ from
breaking a line where it might otherwise do so.
This may be at a space, as between the ``Mr.''\ and
``Jones'' in
       ``Mr.~Jones'',        % ~ produces an unbreakable interword space.
or within a word---especially when the word is a
symbol like
       \mbox{\emph{itemnum}}
that makes little sense when hyphenated across
lines.

Footnotes\footnote{This is an example of a footnote.}
pose no problem.

\LaTeX\ is good at typesetting mathematical formulas
like
       \( x-3y + z = 7 \)
or
       \( a_{1} > x^{2n} + y^{2n} > x' \)
or
       \( \ip{A}{B} = \sum_{i} a_{i} b_{i} \).
The spaces you type in a formula are
ignored.  Remember that a letter like
       $x$                   % $ ... $  and  \( ... \)  are equivalent
is a formula when it denotes a mathematical
symbol, and it should be typed as one.

\section{Displayed Text}

Text is displayed by indenting it from the left
margin.  Quotations are commonly displayed.  There
are short quotations
\begin{quote}
   This is a short quotation.  It consists of a
   single paragraph of text.  See how it is formatted.
\end{quote}
and longer ones.
\begin{quotation}
   This is a longer quotation.  It consists of two
   paragraphs of text, neither of which are
   particularly interesting.

   This is the second paragraph of the quotation.  It
   is just as dull as the first paragraph.
\end{quotation}
Another frequently-displayed structure is a list.
The following is an example of an \emph{itemized}
list.
\begin{itemize}
   \item This is the first item of an itemized list.
         Each item in the list is marked with a ``tick''.
         You don't have to worry about what kind of tick
         mark is used.

   \item This is the second item of the list.  It
         contains another list nested inside it.  The inner
         list is an \emph{enumerated} list.
         \begin{enumerate}
            \item This is the first item of an enumerated
                  list that is nested within the itemized list.

            \item This is the second item of the inner list.
                  \LaTeX\ allows you to nest lists deeper than
                  you really should.
         \end{enumerate}
         This is the rest of the second item of the outer
         list.  It is no more interesting than any other
         part of the item.
   \item This is the third item of the list.
\end{itemize}
You can even display poetry.
\begin{verse}
   There is an environment
    for verse \\             % The \\ command separates lines
   Whose features some poets % within a stanza.
   will curse.

                             % One or more blank lines separate stanzas.

   For instead of making\\
   Them do \emph{all} line breaking, \\
   It allows them to put too many words on a line when they'd rather be
   forced to be terse.
\end{verse}

Mathematical formulas may also be displayed.  A
displayed formula
is
one-line long; multiline
formulas require special formatting instructions.
   \[  \ip{\Gamma}{\psi'} = x'' + y^{2} + z_{i}^{n}\]
Don't start a paragraph with a displayed equation,
nor make one a paragraph by itself.

\end{document}               % End of document.

\section{Bench Section 151}
% synthetic bench section 151
% This is a sample LaTeX input file.  (Version of 12 August 2004.)
%
% A '%' character causes TeX to ignore all remaining text on the line,
% and is used for comments like this one.

\documentclass{article}      % Specifies the document class

                             % The preamble begins here.
\title{An Example Document}  % Declares the document's title.
\author{Leslie Lamport}      % Declares the author's name.
\date{January 21, 1994}      % Deleting this command produces today's date.

\newcommand{\ip}[2]{(#1, #2)}
                             % Defines \ip{arg1}{arg2} to mean
                             % (arg1, arg2).

%\newcommand{\ip}[2]{\langle #1 | #2\rangle}
                             % This is an alternative definition of
                             % \ip that is commented out.

\begin{document}             % End of preamble and beginning of text.

\maketitle                   % Produces the title.

This is an example input file.  Comparing it with
the output it generates can show you how to
produce a simple document of your own.

\section{Ordinary Text}      % Produces section heading.  Lower-level
                             % sections are begun with similar
                             % \subsection and \subsubsection commands.

The ends  of words and sentences are marked
  by   spaces. It  doesn't matter how many
spaces    you type; one is as good as 100.  The
end of   a line counts as a space.

One   or more   blank lines denote the  end
of  a paragraph.

Since any number of consecutive spaces are treated
like a single one, the formatting of the input
file makes no difference to
      \LaTeX,                % The \LaTeX command generates the LaTeX logo.
but it makes a difference to you.  When you use
\LaTeX, making your input file as easy to read
as possible will be a great help as you write
your document and when you change it.  This sample
file shows how you can add comments to your own input
file.

Because printing is different from typewriting,
there are a number of things that you have to do
differently when preparing an input file than if
you were just typing the document directly.
Quotation marks like
       ``this''
have to be handled specially, as do quotes within
quotes:
       ``\,`this'            % \, separates the double and single quote.
        is what I just
        wrote, not  `that'\,''.

Dashes come in three sizes: an
       intra-word
dash, a medium dash for number ranges like
       1--2,
and a punctuation
       dash---like
this.

A sentence-ending space should be larger than the
space between words within a sentence.  You
sometimes have to type special commands in
conjunction with punctuation characters to get
this right, as in the following sentence.
       Gnats, gnus, etc.\ all  % `\ ' makes an inter-word space.
       begin with G\@.         % \@ marks end-of-sentence punctuation.
You should check the spaces after periods when
reading your output to make sure you haven't
forgotten any special cases.  Generating an
ellipsis
       \ldots\               % `\ ' is needed after `\ldots' because TeX
                             % ignores spaces after command names like \ldots
                             % made from \ + letters.
                             %
                             % Note how a `%' character causes TeX to ignore
                             % the end of the input line, so these blank lines
                             % do not start a new paragraph.
                             %
with the right spacing around the periods requires
a special command.

\LaTeX\ interprets some common characters as
commands, so you must type special commands to
generate them.  These characters include the
following:
       \$ \& \% \# \{ and \}.

In printing, text is usually emphasized with an
       \emph{italic}
type style.

\begin{em}
   A long segment of text can also be emphasized
   in this way.  Text within such a segment can be
   given \emph{additional} emphasis.
\end{em}

It is sometimes necessary to prevent \LaTeX\ from
breaking a line where it might otherwise do so.
This may be at a space, as between the ``Mr.''\ and
``Jones'' in
       ``Mr.~Jones'',        % ~ produces an unbreakable interword space.
or within a word---especially when the word is a
symbol like
       \mbox{\emph{itemnum}}
that makes little sense when hyphenated across
lines.

Footnotes\footnote{This is an example of a footnote.}
pose no problem.

\LaTeX\ is good at typesetting mathematical formulas
like
       \( x-3y + z = 7 \)
or
       \( a_{1} > x^{2n} + y^{2n} > x' \)
or
       \( \ip{A}{B} = \sum_{i} a_{i} b_{i} \).
The spaces you type in a formula are
ignored.  Remember that a letter like
       $x$                   % $ ... $  and  \( ... \)  are equivalent
is a formula when it denotes a mathematical
symbol, and it should be typed as one.

\section{Displayed Text}

Text is displayed by indenting it from the left
margin.  Quotations are commonly displayed.  There
are short quotations
\begin{quote}
   This is a short quotation.  It consists of a
   single paragraph of text.  See how it is formatted.
\end{quote}
and longer ones.
\begin{quotation}
   This is a longer quotation.  It consists of two
   paragraphs of text, neither of which are
   particularly interesting.

   This is the second paragraph of the quotation.  It
   is just as dull as the first paragraph.
\end{quotation}
Another frequently-displayed structure is a list.
The following is an example of an \emph{itemized}
list.
\begin{itemize}
   \item This is the first item of an itemized list.
         Each item in the list is marked with a ``tick''.
         You don't have to worry about what kind of tick
         mark is used.

   \item This is the second item of the list.  It
         contains another list nested inside it.  The inner
         list is an \emph{enumerated} list.
         \begin{enumerate}
            \item This is the first item of an enumerated
                  list that is nested within the itemized list.

            \item This is the second item of the inner list.
                  \LaTeX\ allows you to nest lists deeper than
                  you really should.
         \end{enumerate}
         This is the rest of the second item of the outer
         list.  It is no more interesting than any other
         part of the item.
   \item This is the third item of the list.
\end{itemize}
You can even display poetry.
\begin{verse}
   There is an environment
    for verse \\             % The \\ command separates lines
   Whose features some poets % within a stanza.
   will curse.

                             % One or more blank lines separate stanzas.

   For instead of making\\
   Them do \emph{all} line breaking, \\
   It allows them to put too many words on a line when they'd rather be
   forced to be terse.
\end{verse}

Mathematical formulas may also be displayed.  A
displayed formula
is
one-line long; multiline
formulas require special formatting instructions.
   \[  \ip{\Gamma}{\psi'} = x'' + y^{2} + z_{i}^{n}\]
Don't start a paragraph with a displayed equation,
nor make one a paragraph by itself.

\end{document}               % End of document.

\section{Bench Section 152}
% synthetic bench section 152
% This is a sample LaTeX input file.  (Version of 12 August 2004.)
%
% A '%' character causes TeX to ignore all remaining text on the line,
% and is used for comments like this one.

\documentclass{article}      % Specifies the document class

                             % The preamble begins here.
\title{An Example Document}  % Declares the document's title.
\author{Leslie Lamport}      % Declares the author's name.
\date{January 21, 1994}      % Deleting this command produces today's date.

\newcommand{\ip}[2]{(#1, #2)}
                             % Defines \ip{arg1}{arg2} to mean
                             % (arg1, arg2).

%\newcommand{\ip}[2]{\langle #1 | #2\rangle}
                             % This is an alternative definition of
                             % \ip that is commented out.

\begin{document}             % End of preamble and beginning of text.

\maketitle                   % Produces the title.

This is an example input file.  Comparing it with
the output it generates can show you how to
produce a simple document of your own.

\section{Ordinary Text}      % Produces section heading.  Lower-level
                             % sections are begun with similar
                             % \subsection and \subsubsection commands.

The ends  of words and sentences are marked
  by   spaces. It  doesn't matter how many
spaces    you type; one is as good as 100.  The
end of   a line counts as a space.

One   or more   blank lines denote the  end
of  a paragraph.

Since any number of consecutive spaces are treated
like a single one, the formatting of the input
file makes no difference to
      \LaTeX,                % The \LaTeX command generates the LaTeX logo.
but it makes a difference to you.  When you use
\LaTeX, making your input file as easy to read
as possible will be a great help as you write
your document and when you change it.  This sample
file shows how you can add comments to your own input
file.

Because printing is different from typewriting,
there are a number of things that you have to do
differently when preparing an input file than if
you were just typing the document directly.
Quotation marks like
       ``this''
have to be handled specially, as do quotes within
quotes:
       ``\,`this'            % \, separates the double and single quote.
        is what I just
        wrote, not  `that'\,''.

Dashes come in three sizes: an
       intra-word
dash, a medium dash for number ranges like
       1--2,
and a punctuation
       dash---like
this.

A sentence-ending space should be larger than the
space between words within a sentence.  You
sometimes have to type special commands in
conjunction with punctuation characters to get
this right, as in the following sentence.
       Gnats, gnus, etc.\ all  % `\ ' makes an inter-word space.
       begin with G\@.         % \@ marks end-of-sentence punctuation.
You should check the spaces after periods when
reading your output to make sure you haven't
forgotten any special cases.  Generating an
ellipsis
       \ldots\               % `\ ' is needed after `\ldots' because TeX
                             % ignores spaces after command names like \ldots
                             % made from \ + letters.
                             %
                             % Note how a `%' character causes TeX to ignore
                             % the end of the input line, so these blank lines
                             % do not start a new paragraph.
                             %
with the right spacing around the periods requires
a special command.

\LaTeX\ interprets some common characters as
commands, so you must type special commands to
generate them.  These characters include the
following:
       \$ \& \% \# \{ and \}.

In printing, text is usually emphasized with an
       \emph{italic}
type style.

\begin{em}
   A long segment of text can also be emphasized
   in this way.  Text within such a segment can be
   given \emph{additional} emphasis.
\end{em}

It is sometimes necessary to prevent \LaTeX\ from
breaking a line where it might otherwise do so.
This may be at a space, as between the ``Mr.''\ and
``Jones'' in
       ``Mr.~Jones'',        % ~ produces an unbreakable interword space.
or within a word---especially when the word is a
symbol like
       \mbox{\emph{itemnum}}
that makes little sense when hyphenated across
lines.

Footnotes\footnote{This is an example of a footnote.}
pose no problem.

\LaTeX\ is good at typesetting mathematical formulas
like
       \( x-3y + z = 7 \)
or
       \( a_{1} > x^{2n} + y^{2n} > x' \)
or
       \( \ip{A}{B} = \sum_{i} a_{i} b_{i} \).
The spaces you type in a formula are
ignored.  Remember that a letter like
       $x$                   % $ ... $  and  \( ... \)  are equivalent
is a formula when it denotes a mathematical
symbol, and it should be typed as one.

\section{Displayed Text}

Text is displayed by indenting it from the left
margin.  Quotations are commonly displayed.  There
are short quotations
\begin{quote}
   This is a short quotation.  It consists of a
   single paragraph of text.  See how it is formatted.
\end{quote}
and longer ones.
\begin{quotation}
   This is a longer quotation.  It consists of two
   paragraphs of text, neither of which are
   particularly interesting.

   This is the second paragraph of the quotation.  It
   is just as dull as the first paragraph.
\end{quotation}
Another frequently-displayed structure is a list.
The following is an example of an \emph{itemized}
list.
\begin{itemize}
   \item This is the first item of an itemized list.
         Each item in the list is marked with a ``tick''.
         You don't have to worry about what kind of tick
         mark is used.

   \item This is the second item of the list.  It
         contains another list nested inside it.  The inner
         list is an \emph{enumerated} list.
         \begin{enumerate}
            \item This is the first item of an enumerated
                  list that is nested within the itemized list.

            \item This is the second item of the inner list.
                  \LaTeX\ allows you to nest lists deeper than
                  you really should.
         \end{enumerate}
         This is the rest of the second item of the outer
         list.  It is no more interesting than any other
         part of the item.
   \item This is the third item of the list.
\end{itemize}
You can even display poetry.
\begin{verse}
   There is an environment
    for verse \\             % The \\ command separates lines
   Whose features some poets % within a stanza.
   will curse.

                             % One or more blank lines separate stanzas.

   For instead of making\\
   Them do \emph{all} line breaking, \\
   It allows them to put too many words on a line when they'd rather be
   forced to be terse.
\end{verse}

Mathematical formulas may also be displayed.  A
displayed formula
is
one-line long; multiline
formulas require special formatting instructions.
   \[  \ip{\Gamma}{\psi'} = x'' + y^{2} + z_{i}^{n}\]
Don't start a paragraph with a displayed equation,
nor make one a paragraph by itself.

\end{document}               % End of document.

\section{Bench Section 153}
% synthetic bench section 153
% This is a sample LaTeX input file.  (Version of 12 August 2004.)
%
% A '%' character causes TeX to ignore all remaining text on the line,
% and is used for comments like this one.

\documentclass{article}      % Specifies the document class

                             % The preamble begins here.
\title{An Example Document}  % Declares the document's title.
\author{Leslie Lamport}      % Declares the author's name.
\date{January 21, 1994}      % Deleting this command produces today's date.

\newcommand{\ip}[2]{(#1, #2)}
                             % Defines \ip{arg1}{arg2} to mean
                             % (arg1, arg2).

%\newcommand{\ip}[2]{\langle #1 | #2\rangle}
                             % This is an alternative definition of
                             % \ip that is commented out.

\begin{document}             % End of preamble and beginning of text.

\maketitle                   % Produces the title.

This is an example input file.  Comparing it with
the output it generates can show you how to
produce a simple document of your own.

\section{Ordinary Text}      % Produces section heading.  Lower-level
                             % sections are begun with similar
                             % \subsection and \subsubsection commands.

The ends  of words and sentences are marked
  by   spaces. It  doesn't matter how many
spaces    you type; one is as good as 100.  The
end of   a line counts as a space.

One   or more   blank lines denote the  end
of  a paragraph.

Since any number of consecutive spaces are treated
like a single one, the formatting of the input
file makes no difference to
      \LaTeX,                % The \LaTeX command generates the LaTeX logo.
but it makes a difference to you.  When you use
\LaTeX, making your input file as easy to read
as possible will be a great help as you write
your document and when you change it.  This sample
file shows how you can add comments to your own input
file.

Because printing is different from typewriting,
there are a number of things that you have to do
differently when preparing an input file than if
you were just typing the document directly.
Quotation marks like
       ``this''
have to be handled specially, as do quotes within
quotes:
       ``\,`this'            % \, separates the double and single quote.
        is what I just
        wrote, not  `that'\,''.

Dashes come in three sizes: an
       intra-word
dash, a medium dash for number ranges like
       1--2,
and a punctuation
       dash---like
this.

A sentence-ending space should be larger than the
space between words within a sentence.  You
sometimes have to type special commands in
conjunction with punctuation characters to get
this right, as in the following sentence.
       Gnats, gnus, etc.\ all  % `\ ' makes an inter-word space.
       begin with G\@.         % \@ marks end-of-sentence punctuation.
You should check the spaces after periods when
reading your output to make sure you haven't
forgotten any special cases.  Generating an
ellipsis
       \ldots\               % `\ ' is needed after `\ldots' because TeX
                             % ignores spaces after command names like \ldots
                             % made from \ + letters.
                             %
                             % Note how a `%' character causes TeX to ignore
                             % the end of the input line, so these blank lines
                             % do not start a new paragraph.
                             %
with the right spacing around the periods requires
a special command.

\LaTeX\ interprets some common characters as
commands, so you must type special commands to
generate them.  These characters include the
following:
       \$ \& \% \# \{ and \}.

In printing, text is usually emphasized with an
       \emph{italic}
type style.

\begin{em}
   A long segment of text can also be emphasized
   in this way.  Text within such a segment can be
   given \emph{additional} emphasis.
\end{em}

It is sometimes necessary to prevent \LaTeX\ from
breaking a line where it might otherwise do so.
This may be at a space, as between the ``Mr.''\ and
``Jones'' in
       ``Mr.~Jones'',        % ~ produces an unbreakable interword space.
or within a word---especially when the word is a
symbol like
       \mbox{\emph{itemnum}}
that makes little sense when hyphenated across
lines.

Footnotes\footnote{This is an example of a footnote.}
pose no problem.

\LaTeX\ is good at typesetting mathematical formulas
like
       \( x-3y + z = 7 \)
or
       \( a_{1} > x^{2n} + y^{2n} > x' \)
or
       \( \ip{A}{B} = \sum_{i} a_{i} b_{i} \).
The spaces you type in a formula are
ignored.  Remember that a letter like
       $x$                   % $ ... $  and  \( ... \)  are equivalent
is a formula when it denotes a mathematical
symbol, and it should be typed as one.

\section{Displayed Text}

Text is displayed by indenting it from the left
margin.  Quotations are commonly displayed.  There
are short quotations
\begin{quote}
   This is a short quotation.  It consists of a
   single paragraph of text.  See how it is formatted.
\end{quote}
and longer ones.
\begin{quotation}
   This is a longer quotation.  It consists of two
   paragraphs of text, neither of which are
   particularly interesting.

   This is the second paragraph of the quotation.  It
   is just as dull as the first paragraph.
\end{quotation}
Another frequently-displayed structure is a list.
The following is an example of an \emph{itemized}
list.
\begin{itemize}
   \item This is the first item of an itemized list.
         Each item in the list is marked with a ``tick''.
         You don't have to worry about what kind of tick
         mark is used.

   \item This is the second item of the list.  It
         contains another list nested inside it.  The inner
         list is an \emph{enumerated} list.
         \begin{enumerate}
            \item This is the first item of an enumerated
                  list that is nested within the itemized list.

            \item This is the second item of the inner list.
                  \LaTeX\ allows you to nest lists deeper than
                  you really should.
         \end{enumerate}
         This is the rest of the second item of the outer
         list.  It is no more interesting than any other
         part of the item.
   \item This is the third item of the list.
\end{itemize}
You can even display poetry.
\begin{verse}
   There is an environment
    for verse \\             % The \\ command separates lines
   Whose features some poets % within a stanza.
   will curse.

                             % One or more blank lines separate stanzas.

   For instead of making\\
   Them do \emph{all} line breaking, \\
   It allows them to put too many words on a line when they'd rather be
   forced to be terse.
\end{verse}

Mathematical formulas may also be displayed.  A
displayed formula
is
one-line long; multiline
formulas require special formatting instructions.
   \[  \ip{\Gamma}{\psi'} = x'' + y^{2} + z_{i}^{n}\]
Don't start a paragraph with a displayed equation,
nor make one a paragraph by itself.

\end{document}               % End of document.

\section{Bench Section 154}
% synthetic bench section 154
% This is a sample LaTeX input file.  (Version of 12 August 2004.)
%
% A '%' character causes TeX to ignore all remaining text on the line,
% and is used for comments like this one.

\documentclass{article}      % Specifies the document class

                             % The preamble begins here.
\title{An Example Document}  % Declares the document's title.
\author{Leslie Lamport}      % Declares the author's name.
\date{January 21, 1994}      % Deleting this command produces today's date.

\newcommand{\ip}[2]{(#1, #2)}
                             % Defines \ip{arg1}{arg2} to mean
                             % (arg1, arg2).

%\newcommand{\ip}[2]{\langle #1 | #2\rangle}
                             % This is an alternative definition of
                             % \ip that is commented out.

\begin{document}             % End of preamble and beginning of text.

\maketitle                   % Produces the title.

This is an example input file.  Comparing it with
the output it generates can show you how to
produce a simple document of your own.

\section{Ordinary Text}      % Produces section heading.  Lower-level
                             % sections are begun with similar
                             % \subsection and \subsubsection commands.

The ends  of words and sentences are marked
  by   spaces. It  doesn't matter how many
spaces    you type; one is as good as 100.  The
end of   a line counts as a space.

One   or more   blank lines denote the  end
of  a paragraph.

Since any number of consecutive spaces are treated
like a single one, the formatting of the input
file makes no difference to
      \LaTeX,                % The \LaTeX command generates the LaTeX logo.
but it makes a difference to you.  When you use
\LaTeX, making your input file as easy to read
as possible will be a great help as you write
your document and when you change it.  This sample
file shows how you can add comments to your own input
file.

Because printing is different from typewriting,
there are a number of things that you have to do
differently when preparing an input file than if
you were just typing the document directly.
Quotation marks like
       ``this''
have to be handled specially, as do quotes within
quotes:
       ``\,`this'            % \, separates the double and single quote.
        is what I just
        wrote, not  `that'\,''.

Dashes come in three sizes: an
       intra-word
dash, a medium dash for number ranges like
       1--2,
and a punctuation
       dash---like
this.

A sentence-ending space should be larger than the
space between words within a sentence.  You
sometimes have to type special commands in
conjunction with punctuation characters to get
this right, as in the following sentence.
       Gnats, gnus, etc.\ all  % `\ ' makes an inter-word space.
       begin with G\@.         % \@ marks end-of-sentence punctuation.
You should check the spaces after periods when
reading your output to make sure you haven't
forgotten any special cases.  Generating an
ellipsis
       \ldots\               % `\ ' is needed after `\ldots' because TeX
                             % ignores spaces after command names like \ldots
                             % made from \ + letters.
                             %
                             % Note how a `%' character causes TeX to ignore
                             % the end of the input line, so these blank lines
                             % do not start a new paragraph.
                             %
with the right spacing around the periods requires
a special command.

\LaTeX\ interprets some common characters as
commands, so you must type special commands to
generate them.  These characters include the
following:
       \$ \& \% \# \{ and \}.

In printing, text is usually emphasized with an
       \emph{italic}
type style.

\begin{em}
   A long segment of text can also be emphasized
   in this way.  Text within such a segment can be
   given \emph{additional} emphasis.
\end{em}

It is sometimes necessary to prevent \LaTeX\ from
breaking a line where it might otherwise do so.
This may be at a space, as between the ``Mr.''\ and
``Jones'' in
       ``Mr.~Jones'',        % ~ produces an unbreakable interword space.
or within a word---especially when the word is a
symbol like
       \mbox{\emph{itemnum}}
that makes little sense when hyphenated across
lines.

Footnotes\footnote{This is an example of a footnote.}
pose no problem.

\LaTeX\ is good at typesetting mathematical formulas
like
       \( x-3y + z = 7 \)
or
       \( a_{1} > x^{2n} + y^{2n} > x' \)
or
       \( \ip{A}{B} = \sum_{i} a_{i} b_{i} \).
The spaces you type in a formula are
ignored.  Remember that a letter like
       $x$                   % $ ... $  and  \( ... \)  are equivalent
is a formula when it denotes a mathematical
symbol, and it should be typed as one.

\section{Displayed Text}

Text is displayed by indenting it from the left
margin.  Quotations are commonly displayed.  There
are short quotations
\begin{quote}
   This is a short quotation.  It consists of a
   single paragraph of text.  See how it is formatted.
\end{quote}
and longer ones.
\begin{quotation}
   This is a longer quotation.  It consists of two
   paragraphs of text, neither of which are
   particularly interesting.

   This is the second paragraph of the quotation.  It
   is just as dull as the first paragraph.
\end{quotation}
Another frequently-displayed structure is a list.
The following is an example of an \emph{itemized}
list.
\begin{itemize}
   \item This is the first item of an itemized list.
         Each item in the list is marked with a ``tick''.
         You don't have to worry about what kind of tick
         mark is used.

   \item This is the second item of the list.  It
         contains another list nested inside it.  The inner
         list is an \emph{enumerated} list.
         \begin{enumerate}
            \item This is the first item of an enumerated
                  list that is nested within the itemized list.

            \item This is the second item of the inner list.
                  \LaTeX\ allows you to nest lists deeper than
                  you really should.
         \end{enumerate}
         This is the rest of the second item of the outer
         list.  It is no more interesting than any other
         part of the item.
   \item This is the third item of the list.
\end{itemize}
You can even display poetry.
\begin{verse}
   There is an environment
    for verse \\             % The \\ command separates lines
   Whose features some poets % within a stanza.
   will curse.

                             % One or more blank lines separate stanzas.

   For instead of making\\
   Them do \emph{all} line breaking, \\
   It allows them to put too many words on a line when they'd rather be
   forced to be terse.
\end{verse}

Mathematical formulas may also be displayed.  A
displayed formula
is
one-line long; multiline
formulas require special formatting instructions.
   \[  \ip{\Gamma}{\psi'} = x'' + y^{2} + z_{i}^{n}\]
Don't start a paragraph with a displayed equation,
nor make one a paragraph by itself.

\end{document}               % End of document.

\section{Bench Section 155}
% synthetic bench section 155
% This is a sample LaTeX input file.  (Version of 12 August 2004.)
%
% A '%' character causes TeX to ignore all remaining text on the line,
% and is used for comments like this one.

\documentclass{article}      % Specifies the document class

                             % The preamble begins here.
\title{An Example Document}  % Declares the document's title.
\author{Leslie Lamport}      % Declares the author's name.
\date{January 21, 1994}      % Deleting this command produces today's date.

\newcommand{\ip}[2]{(#1, #2)}
                             % Defines \ip{arg1}{arg2} to mean
                             % (arg1, arg2).

%\newcommand{\ip}[2]{\langle #1 | #2\rangle}
                             % This is an alternative definition of
                             % \ip that is commented out.

\begin{document}             % End of preamble and beginning of text.

\maketitle                   % Produces the title.

This is an example input file.  Comparing it with
the output it generates can show you how to
produce a simple document of your own.

\section{Ordinary Text}      % Produces section heading.  Lower-level
                             % sections are begun with similar
                             % \subsection and \subsubsection commands.

The ends  of words and sentences are marked
  by   spaces. It  doesn't matter how many
spaces    you type; one is as good as 100.  The
end of   a line counts as a space.

One   or more   blank lines denote the  end
of  a paragraph.

Since any number of consecutive spaces are treated
like a single one, the formatting of the input
file makes no difference to
      \LaTeX,                % The \LaTeX command generates the LaTeX logo.
but it makes a difference to you.  When you use
\LaTeX, making your input file as easy to read
as possible will be a great help as you write
your document and when you change it.  This sample
file shows how you can add comments to your own input
file.

Because printing is different from typewriting,
there are a number of things that you have to do
differently when preparing an input file than if
you were just typing the document directly.
Quotation marks like
       ``this''
have to be handled specially, as do quotes within
quotes:
       ``\,`this'            % \, separates the double and single quote.
        is what I just
        wrote, not  `that'\,''.

Dashes come in three sizes: an
       intra-word
dash, a medium dash for number ranges like
       1--2,
and a punctuation
       dash---like
this.

A sentence-ending space should be larger than the
space between words within a sentence.  You
sometimes have to type special commands in
conjunction with punctuation characters to get
this right, as in the following sentence.
       Gnats, gnus, etc.\ all  % `\ ' makes an inter-word space.
       begin with G\@.         % \@ marks end-of-sentence punctuation.
You should check the spaces after periods when
reading your output to make sure you haven't
forgotten any special cases.  Generating an
ellipsis
       \ldots\               % `\ ' is needed after `\ldots' because TeX
                             % ignores spaces after command names like \ldots
                             % made from \ + letters.
                             %
                             % Note how a `%' character causes TeX to ignore
                             % the end of the input line, so these blank lines
                             % do not start a new paragraph.
                             %
with the right spacing around the periods requires
a special command.

\LaTeX\ interprets some common characters as
commands, so you must type special commands to
generate them.  These characters include the
following:
       \$ \& \% \# \{ and \}.

In printing, text is usually emphasized with an
       \emph{italic}
type style.

\begin{em}
   A long segment of text can also be emphasized
   in this way.  Text within such a segment can be
   given \emph{additional} emphasis.
\end{em}

It is sometimes necessary to prevent \LaTeX\ from
breaking a line where it might otherwise do so.
This may be at a space, as between the ``Mr.''\ and
``Jones'' in
       ``Mr.~Jones'',        % ~ produces an unbreakable interword space.
or within a word---especially when the word is a
symbol like
       \mbox{\emph{itemnum}}
that makes little sense when hyphenated across
lines.

Footnotes\footnote{This is an example of a footnote.}
pose no problem.

\LaTeX\ is good at typesetting mathematical formulas
like
       \( x-3y + z = 7 \)
or
       \( a_{1} > x^{2n} + y^{2n} > x' \)
or
       \( \ip{A}{B} = \sum_{i} a_{i} b_{i} \).
The spaces you type in a formula are
ignored.  Remember that a letter like
       $x$                   % $ ... $  and  \( ... \)  are equivalent
is a formula when it denotes a mathematical
symbol, and it should be typed as one.

\section{Displayed Text}

Text is displayed by indenting it from the left
margin.  Quotations are commonly displayed.  There
are short quotations
\begin{quote}
   This is a short quotation.  It consists of a
   single paragraph of text.  See how it is formatted.
\end{quote}
and longer ones.
\begin{quotation}
   This is a longer quotation.  It consists of two
   paragraphs of text, neither of which are
   particularly interesting.

   This is the second paragraph of the quotation.  It
   is just as dull as the first paragraph.
\end{quotation}
Another frequently-displayed structure is a list.
The following is an example of an \emph{itemized}
list.
\begin{itemize}
   \item This is the first item of an itemized list.
         Each item in the list is marked with a ``tick''.
         You don't have to worry about what kind of tick
         mark is used.

   \item This is the second item of the list.  It
         contains another list nested inside it.  The inner
         list is an \emph{enumerated} list.
         \begin{enumerate}
            \item This is the first item of an enumerated
                  list that is nested within the itemized list.

            \item This is the second item of the inner list.
                  \LaTeX\ allows you to nest lists deeper than
                  you really should.
         \end{enumerate}
         This is the rest of the second item of the outer
         list.  It is no more interesting than any other
         part of the item.
   \item This is the third item of the list.
\end{itemize}
You can even display poetry.
\begin{verse}
   There is an environment
    for verse \\             % The \\ command separates lines
   Whose features some poets % within a stanza.
   will curse.

                             % One or more blank lines separate stanzas.

   For instead of making\\
   Them do \emph{all} line breaking, \\
   It allows them to put too many words on a line when they'd rather be
   forced to be terse.
\end{verse}

Mathematical formulas may also be displayed.  A
displayed formula
is
one-line long; multiline
formulas require special formatting instructions.
   \[  \ip{\Gamma}{\psi'} = x'' + y^{2} + z_{i}^{n}\]
Don't start a paragraph with a displayed equation,
nor make one a paragraph by itself.

\end{document}               % End of document.

\section{Bench Section 156}
% synthetic bench section 156
% This is a sample LaTeX input file.  (Version of 12 August 2004.)
%
% A '%' character causes TeX to ignore all remaining text on the line,
% and is used for comments like this one.

\documentclass{article}      % Specifies the document class

                             % The preamble begins here.
\title{An Example Document}  % Declares the document's title.
\author{Leslie Lamport}      % Declares the author's name.
\date{January 21, 1994}      % Deleting this command produces today's date.

\newcommand{\ip}[2]{(#1, #2)}
                             % Defines \ip{arg1}{arg2} to mean
                             % (arg1, arg2).

%\newcommand{\ip}[2]{\langle #1 | #2\rangle}
                             % This is an alternative definition of
                             % \ip that is commented out.

\begin{document}             % End of preamble and beginning of text.

\maketitle                   % Produces the title.

This is an example input file.  Comparing it with
the output it generates can show you how to
produce a simple document of your own.

\section{Ordinary Text}      % Produces section heading.  Lower-level
                             % sections are begun with similar
                             % \subsection and \subsubsection commands.

The ends  of words and sentences are marked
  by   spaces. It  doesn't matter how many
spaces    you type; one is as good as 100.  The
end of   a line counts as a space.

One   or more   blank lines denote the  end
of  a paragraph.

Since any number of consecutive spaces are treated
like a single one, the formatting of the input
file makes no difference to
      \LaTeX,                % The \LaTeX command generates the LaTeX logo.
but it makes a difference to you.  When you use
\LaTeX, making your input file as easy to read
as possible will be a great help as you write
your document and when you change it.  This sample
file shows how you can add comments to your own input
file.

Because printing is different from typewriting,
there are a number of things that you have to do
differently when preparing an input file than if
you were just typing the document directly.
Quotation marks like
       ``this''
have to be handled specially, as do quotes within
quotes:
       ``\,`this'            % \, separates the double and single quote.
        is what I just
        wrote, not  `that'\,''.

Dashes come in three sizes: an
       intra-word
dash, a medium dash for number ranges like
       1--2,
and a punctuation
       dash---like
this.

A sentence-ending space should be larger than the
space between words within a sentence.  You
sometimes have to type special commands in
conjunction with punctuation characters to get
this right, as in the following sentence.
       Gnats, gnus, etc.\ all  % `\ ' makes an inter-word space.
       begin with G\@.         % \@ marks end-of-sentence punctuation.
You should check the spaces after periods when
reading your output to make sure you haven't
forgotten any special cases.  Generating an
ellipsis
       \ldots\               % `\ ' is needed after `\ldots' because TeX
                             % ignores spaces after command names like \ldots
                             % made from \ + letters.
                             %
                             % Note how a `%' character causes TeX to ignore
                             % the end of the input line, so these blank lines
                             % do not start a new paragraph.
                             %
with the right spacing around the periods requires
a special command.

\LaTeX\ interprets some common characters as
commands, so you must type special commands to
generate them.  These characters include the
following:
       \$ \& \% \# \{ and \}.

In printing, text is usually emphasized with an
       \emph{italic}
type style.

\begin{em}
   A long segment of text can also be emphasized
   in this way.  Text within such a segment can be
   given \emph{additional} emphasis.
\end{em}

It is sometimes necessary to prevent \LaTeX\ from
breaking a line where it might otherwise do so.
This may be at a space, as between the ``Mr.''\ and
``Jones'' in
       ``Mr.~Jones'',        % ~ produces an unbreakable interword space.
or within a word---especially when the word is a
symbol like
       \mbox{\emph{itemnum}}
that makes little sense when hyphenated across
lines.

Footnotes\footnote{This is an example of a footnote.}
pose no problem.

\LaTeX\ is good at typesetting mathematical formulas
like
       \( x-3y + z = 7 \)
or
       \( a_{1} > x^{2n} + y^{2n} > x' \)
or
       \( \ip{A}{B} = \sum_{i} a_{i} b_{i} \).
The spaces you type in a formula are
ignored.  Remember that a letter like
       $x$                   % $ ... $  and  \( ... \)  are equivalent
is a formula when it denotes a mathematical
symbol, and it should be typed as one.

\section{Displayed Text}

Text is displayed by indenting it from the left
margin.  Quotations are commonly displayed.  There
are short quotations
\begin{quote}
   This is a short quotation.  It consists of a
   single paragraph of text.  See how it is formatted.
\end{quote}
and longer ones.
\begin{quotation}
   This is a longer quotation.  It consists of two
   paragraphs of text, neither of which are
   particularly interesting.

   This is the second paragraph of the quotation.  It
   is just as dull as the first paragraph.
\end{quotation}
Another frequently-displayed structure is a list.
The following is an example of an \emph{itemized}
list.
\begin{itemize}
   \item This is the first item of an itemized list.
         Each item in the list is marked with a ``tick''.
         You don't have to worry about what kind of tick
         mark is used.

   \item This is the second item of the list.  It
         contains another list nested inside it.  The inner
         list is an \emph{enumerated} list.
         \begin{enumerate}
            \item This is the first item of an enumerated
                  list that is nested within the itemized list.

            \item This is the second item of the inner list.
                  \LaTeX\ allows you to nest lists deeper than
                  you really should.
         \end{enumerate}
         This is the rest of the second item of the outer
         list.  It is no more interesting than any other
         part of the item.
   \item This is the third item of the list.
\end{itemize}
You can even display poetry.
\begin{verse}
   There is an environment
    for verse \\             % The \\ command separates lines
   Whose features some poets % within a stanza.
   will curse.

                             % One or more blank lines separate stanzas.

   For instead of making\\
   Them do \emph{all} line breaking, \\
   It allows them to put too many words on a line when they'd rather be
   forced to be terse.
\end{verse}

Mathematical formulas may also be displayed.  A
displayed formula
is
one-line long; multiline
formulas require special formatting instructions.
   \[  \ip{\Gamma}{\psi'} = x'' + y^{2} + z_{i}^{n}\]
Don't start a paragraph with a displayed equation,
nor make one a paragraph by itself.

\end{document}               % End of document.

\section{Bench Section 157}
% synthetic bench section 157
% This is a sample LaTeX input file.  (Version of 12 August 2004.)
%
% A '%' character causes TeX to ignore all remaining text on the line,
% and is used for comments like this one.

\documentclass{article}      % Specifies the document class

                             % The preamble begins here.
\title{An Example Document}  % Declares the document's title.
\author{Leslie Lamport}      % Declares the author's name.
\date{January 21, 1994}      % Deleting this command produces today's date.

\newcommand{\ip}[2]{(#1, #2)}
                             % Defines \ip{arg1}{arg2} to mean
                             % (arg1, arg2).

%\newcommand{\ip}[2]{\langle #1 | #2\rangle}
                             % This is an alternative definition of
                             % \ip that is commented out.

\begin{document}             % End of preamble and beginning of text.

\maketitle                   % Produces the title.

This is an example input file.  Comparing it with
the output it generates can show you how to
produce a simple document of your own.

\section{Ordinary Text}      % Produces section heading.  Lower-level
                             % sections are begun with similar
                             % \subsection and \subsubsection commands.

The ends  of words and sentences are marked
  by   spaces. It  doesn't matter how many
spaces    you type; one is as good as 100.  The
end of   a line counts as a space.

One   or more   blank lines denote the  end
of  a paragraph.

Since any number of consecutive spaces are treated
like a single one, the formatting of the input
file makes no difference to
      \LaTeX,                % The \LaTeX command generates the LaTeX logo.
but it makes a difference to you.  When you use
\LaTeX, making your input file as easy to read
as possible will be a great help as you write
your document and when you change it.  This sample
file shows how you can add comments to your own input
file.

Because printing is different from typewriting,
there are a number of things that you have to do
differently when preparing an input file than if
you were just typing the document directly.
Quotation marks like
       ``this''
have to be handled specially, as do quotes within
quotes:
       ``\,`this'            % \, separates the double and single quote.
        is what I just
        wrote, not  `that'\,''.

Dashes come in three sizes: an
       intra-word
dash, a medium dash for number ranges like
       1--2,
and a punctuation
       dash---like
this.

A sentence-ending space should be larger than the
space between words within a sentence.  You
sometimes have to type special commands in
conjunction with punctuation characters to get
this right, as in the following sentence.
       Gnats, gnus, etc.\ all  % `\ ' makes an inter-word space.
       begin with G\@.         % \@ marks end-of-sentence punctuation.
You should check the spaces after periods when
reading your output to make sure you haven't
forgotten any special cases.  Generating an
ellipsis
       \ldots\               % `\ ' is needed after `\ldots' because TeX
                             % ignores spaces after command names like \ldots
                             % made from \ + letters.
                             %
                             % Note how a `%' character causes TeX to ignore
                             % the end of the input line, so these blank lines
                             % do not start a new paragraph.
                             %
with the right spacing around the periods requires
a special command.

\LaTeX\ interprets some common characters as
commands, so you must type special commands to
generate them.  These characters include the
following:
       \$ \& \% \# \{ and \}.

In printing, text is usually emphasized with an
       \emph{italic}
type style.

\begin{em}
   A long segment of text can also be emphasized
   in this way.  Text within such a segment can be
   given \emph{additional} emphasis.
\end{em}

It is sometimes necessary to prevent \LaTeX\ from
breaking a line where it might otherwise do so.
This may be at a space, as between the ``Mr.''\ and
``Jones'' in
       ``Mr.~Jones'',        % ~ produces an unbreakable interword space.
or within a word---especially when the word is a
symbol like
       \mbox{\emph{itemnum}}
that makes little sense when hyphenated across
lines.

Footnotes\footnote{This is an example of a footnote.}
pose no problem.

\LaTeX\ is good at typesetting mathematical formulas
like
       \( x-3y + z = 7 \)
or
       \( a_{1} > x^{2n} + y^{2n} > x' \)
or
       \( \ip{A}{B} = \sum_{i} a_{i} b_{i} \).
The spaces you type in a formula are
ignored.  Remember that a letter like
       $x$                   % $ ... $  and  \( ... \)  are equivalent
is a formula when it denotes a mathematical
symbol, and it should be typed as one.

\section{Displayed Text}

Text is displayed by indenting it from the left
margin.  Quotations are commonly displayed.  There
are short quotations
\begin{quote}
   This is a short quotation.  It consists of a
   single paragraph of text.  See how it is formatted.
\end{quote}
and longer ones.
\begin{quotation}
   This is a longer quotation.  It consists of two
   paragraphs of text, neither of which are
   particularly interesting.

   This is the second paragraph of the quotation.  It
   is just as dull as the first paragraph.
\end{quotation}
Another frequently-displayed structure is a list.
The following is an example of an \emph{itemized}
list.
\begin{itemize}
   \item This is the first item of an itemized list.
         Each item in the list is marked with a ``tick''.
         You don't have to worry about what kind of tick
         mark is used.

   \item This is the second item of the list.  It
         contains another list nested inside it.  The inner
         list is an \emph{enumerated} list.
         \begin{enumerate}
            \item This is the first item of an enumerated
                  list that is nested within the itemized list.

            \item This is the second item of the inner list.
                  \LaTeX\ allows you to nest lists deeper than
                  you really should.
         \end{enumerate}
         This is the rest of the second item of the outer
         list.  It is no more interesting than any other
         part of the item.
   \item This is the third item of the list.
\end{itemize}
You can even display poetry.
\begin{verse}
   There is an environment
    for verse \\             % The \\ command separates lines
   Whose features some poets % within a stanza.
   will curse.

                             % One or more blank lines separate stanzas.

   For instead of making\\
   Them do \emph{all} line breaking, \\
   It allows them to put too many words on a line when they'd rather be
   forced to be terse.
\end{verse}

Mathematical formulas may also be displayed.  A
displayed formula
is
one-line long; multiline
formulas require special formatting instructions.
   \[  \ip{\Gamma}{\psi'} = x'' + y^{2} + z_{i}^{n}\]
Don't start a paragraph with a displayed equation,
nor make one a paragraph by itself.

\end{document}               % End of document.

\section{Bench Section 158}
% synthetic bench section 158
% This is a sample LaTeX input file.  (Version of 12 August 2004.)
%
% A '%' character causes TeX to ignore all remaining text on the line,
% and is used for comments like this one.

\documentclass{article}      % Specifies the document class

                             % The preamble begins here.
\title{An Example Document}  % Declares the document's title.
\author{Leslie Lamport}      % Declares the author's name.
\date{January 21, 1994}      % Deleting this command produces today's date.

\newcommand{\ip}[2]{(#1, #2)}
                             % Defines \ip{arg1}{arg2} to mean
                             % (arg1, arg2).

%\newcommand{\ip}[2]{\langle #1 | #2\rangle}
                             % This is an alternative definition of
                             % \ip that is commented out.

\begin{document}             % End of preamble and beginning of text.

\maketitle                   % Produces the title.

This is an example input file.  Comparing it with
the output it generates can show you how to
produce a simple document of your own.

\section{Ordinary Text}      % Produces section heading.  Lower-level
                             % sections are begun with similar
                             % \subsection and \subsubsection commands.

The ends  of words and sentences are marked
  by   spaces. It  doesn't matter how many
spaces    you type; one is as good as 100.  The
end of   a line counts as a space.

One   or more   blank lines denote the  end
of  a paragraph.

Since any number of consecutive spaces are treated
like a single one, the formatting of the input
file makes no difference to
      \LaTeX,                % The \LaTeX command generates the LaTeX logo.
but it makes a difference to you.  When you use
\LaTeX, making your input file as easy to read
as possible will be a great help as you write
your document and when you change it.  This sample
file shows how you can add comments to your own input
file.

Because printing is different from typewriting,
there are a number of things that you have to do
differently when preparing an input file than if
you were just typing the document directly.
Quotation marks like
       ``this''
have to be handled specially, as do quotes within
quotes:
       ``\,`this'            % \, separates the double and single quote.
        is what I just
        wrote, not  `that'\,''.

Dashes come in three sizes: an
       intra-word
dash, a medium dash for number ranges like
       1--2,
and a punctuation
       dash---like
this.

A sentence-ending space should be larger than the
space between words within a sentence.  You
sometimes have to type special commands in
conjunction with punctuation characters to get
this right, as in the following sentence.
       Gnats, gnus, etc.\ all  % `\ ' makes an inter-word space.
       begin with G\@.         % \@ marks end-of-sentence punctuation.
You should check the spaces after periods when
reading your output to make sure you haven't
forgotten any special cases.  Generating an
ellipsis
       \ldots\               % `\ ' is needed after `\ldots' because TeX
                             % ignores spaces after command names like \ldots
                             % made from \ + letters.
                             %
                             % Note how a `%' character causes TeX to ignore
                             % the end of the input line, so these blank lines
                             % do not start a new paragraph.
                             %
with the right spacing around the periods requires
a special command.

\LaTeX\ interprets some common characters as
commands, so you must type special commands to
generate them.  These characters include the
following:
       \$ \& \% \# \{ and \}.

In printing, text is usually emphasized with an
       \emph{italic}
type style.

\begin{em}
   A long segment of text can also be emphasized
   in this way.  Text within such a segment can be
   given \emph{additional} emphasis.
\end{em}

It is sometimes necessary to prevent \LaTeX\ from
breaking a line where it might otherwise do so.
This may be at a space, as between the ``Mr.''\ and
``Jones'' in
       ``Mr.~Jones'',        % ~ produces an unbreakable interword space.
or within a word---especially when the word is a
symbol like
       \mbox{\emph{itemnum}}
that makes little sense when hyphenated across
lines.

Footnotes\footnote{This is an example of a footnote.}
pose no problem.

\LaTeX\ is good at typesetting mathematical formulas
like
       \( x-3y + z = 7 \)
or
       \( a_{1} > x^{2n} + y^{2n} > x' \)
or
       \( \ip{A}{B} = \sum_{i} a_{i} b_{i} \).
The spaces you type in a formula are
ignored.  Remember that a letter like
       $x$                   % $ ... $  and  \( ... \)  are equivalent
is a formula when it denotes a mathematical
symbol, and it should be typed as one.

\section{Displayed Text}

Text is displayed by indenting it from the left
margin.  Quotations are commonly displayed.  There
are short quotations
\begin{quote}
   This is a short quotation.  It consists of a
   single paragraph of text.  See how it is formatted.
\end{quote}
and longer ones.
\begin{quotation}
   This is a longer quotation.  It consists of two
   paragraphs of text, neither of which are
   particularly interesting.

   This is the second paragraph of the quotation.  It
   is just as dull as the first paragraph.
\end{quotation}
Another frequently-displayed structure is a list.
The following is an example of an \emph{itemized}
list.
\begin{itemize}
   \item This is the first item of an itemized list.
         Each item in the list is marked with a ``tick''.
         You don't have to worry about what kind of tick
         mark is used.

   \item This is the second item of the list.  It
         contains another list nested inside it.  The inner
         list is an \emph{enumerated} list.
         \begin{enumerate}
            \item This is the first item of an enumerated
                  list that is nested within the itemized list.

            \item This is the second item of the inner list.
                  \LaTeX\ allows you to nest lists deeper than
                  you really should.
         \end{enumerate}
         This is the rest of the second item of the outer
         list.  It is no more interesting than any other
         part of the item.
   \item This is the third item of the list.
\end{itemize}
You can even display poetry.
\begin{verse}
   There is an environment
    for verse \\             % The \\ command separates lines
   Whose features some poets % within a stanza.
   will curse.

                             % One or more blank lines separate stanzas.

   For instead of making\\
   Them do \emph{all} line breaking, \\
   It allows them to put too many words on a line when they'd rather be
   forced to be terse.
\end{verse}

Mathematical formulas may also be displayed.  A
displayed formula
is
one-line long; multiline
formulas require special formatting instructions.
   \[  \ip{\Gamma}{\psi'} = x'' + y^{2} + z_{i}^{n}\]
Don't start a paragraph with a displayed equation,
nor make one a paragraph by itself.

\end{document}               % End of document.

\section{Bench Section 159}
% synthetic bench section 159
% This is a sample LaTeX input file.  (Version of 12 August 2004.)
%
% A '%' character causes TeX to ignore all remaining text on the line,
% and is used for comments like this one.

\documentclass{article}      % Specifies the document class

                             % The preamble begins here.
\title{An Example Document}  % Declares the document's title.
\author{Leslie Lamport}      % Declares the author's name.
\date{January 21, 1994}      % Deleting this command produces today's date.

\newcommand{\ip}[2]{(#1, #2)}
                             % Defines \ip{arg1}{arg2} to mean
                             % (arg1, arg2).

%\newcommand{\ip}[2]{\langle #1 | #2\rangle}
                             % This is an alternative definition of
                             % \ip that is commented out.

\begin{document}             % End of preamble and beginning of text.

\maketitle                   % Produces the title.

This is an example input file.  Comparing it with
the output it generates can show you how to
produce a simple document of your own.

\section{Ordinary Text}      % Produces section heading.  Lower-level
                             % sections are begun with similar
                             % \subsection and \subsubsection commands.

The ends  of words and sentences are marked
  by   spaces. It  doesn't matter how many
spaces    you type; one is as good as 100.  The
end of   a line counts as a space.

One   or more   blank lines denote the  end
of  a paragraph.

Since any number of consecutive spaces are treated
like a single one, the formatting of the input
file makes no difference to
      \LaTeX,                % The \LaTeX command generates the LaTeX logo.
but it makes a difference to you.  When you use
\LaTeX, making your input file as easy to read
as possible will be a great help as you write
your document and when you change it.  This sample
file shows how you can add comments to your own input
file.

Because printing is different from typewriting,
there are a number of things that you have to do
differently when preparing an input file than if
you were just typing the document directly.
Quotation marks like
       ``this''
have to be handled specially, as do quotes within
quotes:
       ``\,`this'            % \, separates the double and single quote.
        is what I just
        wrote, not  `that'\,''.

Dashes come in three sizes: an
       intra-word
dash, a medium dash for number ranges like
       1--2,
and a punctuation
       dash---like
this.

A sentence-ending space should be larger than the
space between words within a sentence.  You
sometimes have to type special commands in
conjunction with punctuation characters to get
this right, as in the following sentence.
       Gnats, gnus, etc.\ all  % `\ ' makes an inter-word space.
       begin with G\@.         % \@ marks end-of-sentence punctuation.
You should check the spaces after periods when
reading your output to make sure you haven't
forgotten any special cases.  Generating an
ellipsis
       \ldots\               % `\ ' is needed after `\ldots' because TeX
                             % ignores spaces after command names like \ldots
                             % made from \ + letters.
                             %
                             % Note how a `%' character causes TeX to ignore
                             % the end of the input line, so these blank lines
                             % do not start a new paragraph.
                             %
with the right spacing around the periods requires
a special command.

\LaTeX\ interprets some common characters as
commands, so you must type special commands to
generate them.  These characters include the
following:
       \$ \& \% \# \{ and \}.

In printing, text is usually emphasized with an
       \emph{italic}
type style.

\begin{em}
   A long segment of text can also be emphasized
   in this way.  Text within such a segment can be
   given \emph{additional} emphasis.
\end{em}

It is sometimes necessary to prevent \LaTeX\ from
breaking a line where it might otherwise do so.
This may be at a space, as between the ``Mr.''\ and
``Jones'' in
       ``Mr.~Jones'',        % ~ produces an unbreakable interword space.
or within a word---especially when the word is a
symbol like
       \mbox{\emph{itemnum}}
that makes little sense when hyphenated across
lines.

Footnotes\footnote{This is an example of a footnote.}
pose no problem.

\LaTeX\ is good at typesetting mathematical formulas
like
       \( x-3y + z = 7 \)
or
       \( a_{1} > x^{2n} + y^{2n} > x' \)
or
       \( \ip{A}{B} = \sum_{i} a_{i} b_{i} \).
The spaces you type in a formula are
ignored.  Remember that a letter like
       $x$                   % $ ... $  and  \( ... \)  are equivalent
is a formula when it denotes a mathematical
symbol, and it should be typed as one.

\section{Displayed Text}

Text is displayed by indenting it from the left
margin.  Quotations are commonly displayed.  There
are short quotations
\begin{quote}
   This is a short quotation.  It consists of a
   single paragraph of text.  See how it is formatted.
\end{quote}
and longer ones.
\begin{quotation}
   This is a longer quotation.  It consists of two
   paragraphs of text, neither of which are
   particularly interesting.

   This is the second paragraph of the quotation.  It
   is just as dull as the first paragraph.
\end{quotation}
Another frequently-displayed structure is a list.
The following is an example of an \emph{itemized}
list.
\begin{itemize}
   \item This is the first item of an itemized list.
         Each item in the list is marked with a ``tick''.
         You don't have to worry about what kind of tick
         mark is used.

   \item This is the second item of the list.  It
         contains another list nested inside it.  The inner
         list is an \emph{enumerated} list.
         \begin{enumerate}
            \item This is the first item of an enumerated
                  list that is nested within the itemized list.

            \item This is the second item of the inner list.
                  \LaTeX\ allows you to nest lists deeper than
                  you really should.
         \end{enumerate}
         This is the rest of the second item of the outer
         list.  It is no more interesting than any other
         part of the item.
   \item This is the third item of the list.
\end{itemize}
You can even display poetry.
\begin{verse}
   There is an environment
    for verse \\             % The \\ command separates lines
   Whose features some poets % within a stanza.
   will curse.

                             % One or more blank lines separate stanzas.

   For instead of making\\
   Them do \emph{all} line breaking, \\
   It allows them to put too many words on a line when they'd rather be
   forced to be terse.
\end{verse}

Mathematical formulas may also be displayed.  A
displayed formula
is
one-line long; multiline
formulas require special formatting instructions.
   \[  \ip{\Gamma}{\psi'} = x'' + y^{2} + z_{i}^{n}\]
Don't start a paragraph with a displayed equation,
nor make one a paragraph by itself.

\end{document}               % End of document.

\section{Bench Section 160}
% synthetic bench section 160
% This is a sample LaTeX input file.  (Version of 12 August 2004.)
%
% A '%' character causes TeX to ignore all remaining text on the line,
% and is used for comments like this one.

\documentclass{article}      % Specifies the document class

                             % The preamble begins here.
\title{An Example Document}  % Declares the document's title.
\author{Leslie Lamport}      % Declares the author's name.
\date{January 21, 1994}      % Deleting this command produces today's date.

\newcommand{\ip}[2]{(#1, #2)}
                             % Defines \ip{arg1}{arg2} to mean
                             % (arg1, arg2).

%\newcommand{\ip}[2]{\langle #1 | #2\rangle}
                             % This is an alternative definition of
                             % \ip that is commented out.

\begin{document}             % End of preamble and beginning of text.

\maketitle                   % Produces the title.

This is an example input file.  Comparing it with
the output it generates can show you how to
produce a simple document of your own.

\section{Ordinary Text}      % Produces section heading.  Lower-level
                             % sections are begun with similar
                             % \subsection and \subsubsection commands.

The ends  of words and sentences are marked
  by   spaces. It  doesn't matter how many
spaces    you type; one is as good as 100.  The
end of   a line counts as a space.

One   or more   blank lines denote the  end
of  a paragraph.

Since any number of consecutive spaces are treated
like a single one, the formatting of the input
file makes no difference to
      \LaTeX,                % The \LaTeX command generates the LaTeX logo.
but it makes a difference to you.  When you use
\LaTeX, making your input file as easy to read
as possible will be a great help as you write
your document and when you change it.  This sample
file shows how you can add comments to your own input
file.

Because printing is different from typewriting,
there are a number of things that you have to do
differently when preparing an input file than if
you were just typing the document directly.
Quotation marks like
       ``this''
have to be handled specially, as do quotes within
quotes:
       ``\,`this'            % \, separates the double and single quote.
        is what I just
        wrote, not  `that'\,''.

Dashes come in three sizes: an
       intra-word
dash, a medium dash for number ranges like
       1--2,
and a punctuation
       dash---like
this.

A sentence-ending space should be larger than the
space between words within a sentence.  You
sometimes have to type special commands in
conjunction with punctuation characters to get
this right, as in the following sentence.
       Gnats, gnus, etc.\ all  % `\ ' makes an inter-word space.
       begin with G\@.         % \@ marks end-of-sentence punctuation.
You should check the spaces after periods when
reading your output to make sure you haven't
forgotten any special cases.  Generating an
ellipsis
       \ldots\               % `\ ' is needed after `\ldots' because TeX
                             % ignores spaces after command names like \ldots
                             % made from \ + letters.
                             %
                             % Note how a `%' character causes TeX to ignore
                             % the end of the input line, so these blank lines
                             % do not start a new paragraph.
                             %
with the right spacing around the periods requires
a special command.

\LaTeX\ interprets some common characters as
commands, so you must type special commands to
generate them.  These characters include the
following:
       \$ \& \% \# \{ and \}.

In printing, text is usually emphasized with an
       \emph{italic}
type style.

\begin{em}
   A long segment of text can also be emphasized
   in this way.  Text within such a segment can be
   given \emph{additional} emphasis.
\end{em}

It is sometimes necessary to prevent \LaTeX\ from
breaking a line where it might otherwise do so.
This may be at a space, as between the ``Mr.''\ and
``Jones'' in
       ``Mr.~Jones'',        % ~ produces an unbreakable interword space.
or within a word---especially when the word is a
symbol like
       \mbox{\emph{itemnum}}
that makes little sense when hyphenated across
lines.

Footnotes\footnote{This is an example of a footnote.}
pose no problem.

\LaTeX\ is good at typesetting mathematical formulas
like
       \( x-3y + z = 7 \)
or
       \( a_{1} > x^{2n} + y^{2n} > x' \)
or
       \( \ip{A}{B} = \sum_{i} a_{i} b_{i} \).
The spaces you type in a formula are
ignored.  Remember that a letter like
       $x$                   % $ ... $  and  \( ... \)  are equivalent
is a formula when it denotes a mathematical
symbol, and it should be typed as one.

\section{Displayed Text}

Text is displayed by indenting it from the left
margin.  Quotations are commonly displayed.  There
are short quotations
\begin{quote}
   This is a short quotation.  It consists of a
   single paragraph of text.  See how it is formatted.
\end{quote}
and longer ones.
\begin{quotation}
   This is a longer quotation.  It consists of two
   paragraphs of text, neither of which are
   particularly interesting.

   This is the second paragraph of the quotation.  It
   is just as dull as the first paragraph.
\end{quotation}
Another frequently-displayed structure is a list.
The following is an example of an \emph{itemized}
list.
\begin{itemize}
   \item This is the first item of an itemized list.
         Each item in the list is marked with a ``tick''.
         You don't have to worry about what kind of tick
         mark is used.

   \item This is the second item of the list.  It
         contains another list nested inside it.  The inner
         list is an \emph{enumerated} list.
         \begin{enumerate}
            \item This is the first item of an enumerated
                  list that is nested within the itemized list.

            \item This is the second item of the inner list.
                  \LaTeX\ allows you to nest lists deeper than
                  you really should.
         \end{enumerate}
         This is the rest of the second item of the outer
         list.  It is no more interesting than any other
         part of the item.
   \item This is the third item of the list.
\end{itemize}
You can even display poetry.
\begin{verse}
   There is an environment
    for verse \\             % The \\ command separates lines
   Whose features some poets % within a stanza.
   will curse.

                             % One or more blank lines separate stanzas.

   For instead of making\\
   Them do \emph{all} line breaking, \\
   It allows them to put too many words on a line when they'd rather be
   forced to be terse.
\end{verse}

Mathematical formulas may also be displayed.  A
displayed formula
is
one-line long; multiline
formulas require special formatting instructions.
   \[  \ip{\Gamma}{\psi'} = x'' + y^{2} + z_{i}^{n}\]
Don't start a paragraph with a displayed equation,
nor make one a paragraph by itself.

\end{document}               % End of document.

\section{Bench Section 161}
% synthetic bench section 161
% This is a sample LaTeX input file.  (Version of 12 August 2004.)
%
% A '%' character causes TeX to ignore all remaining text on the line,
% and is used for comments like this one.

\documentclass{article}      % Specifies the document class

                             % The preamble begins here.
\title{An Example Document}  % Declares the document's title.
\author{Leslie Lamport}      % Declares the author's name.
\date{January 21, 1994}      % Deleting this command produces today's date.

\newcommand{\ip}[2]{(#1, #2)}
                             % Defines \ip{arg1}{arg2} to mean
                             % (arg1, arg2).

%\newcommand{\ip}[2]{\langle #1 | #2\rangle}
                             % This is an alternative definition of
                             % \ip that is commented out.

\begin{document}             % End of preamble and beginning of text.

\maketitle                   % Produces the title.

This is an example input file.  Comparing it with
the output it generates can show you how to
produce a simple document of your own.

\section{Ordinary Text}      % Produces section heading.  Lower-level
                             % sections are begun with similar
                             % \subsection and \subsubsection commands.

The ends  of words and sentences are marked
  by   spaces. It  doesn't matter how many
spaces    you type; one is as good as 100.  The
end of   a line counts as a space.

One   or more   blank lines denote the  end
of  a paragraph.

Since any number of consecutive spaces are treated
like a single one, the formatting of the input
file makes no difference to
      \LaTeX,                % The \LaTeX command generates the LaTeX logo.
but it makes a difference to you.  When you use
\LaTeX, making your input file as easy to read
as possible will be a great help as you write
your document and when you change it.  This sample
file shows how you can add comments to your own input
file.

Because printing is different from typewriting,
there are a number of things that you have to do
differently when preparing an input file than if
you were just typing the document directly.
Quotation marks like
       ``this''
have to be handled specially, as do quotes within
quotes:
       ``\,`this'            % \, separates the double and single quote.
        is what I just
        wrote, not  `that'\,''.

Dashes come in three sizes: an
       intra-word
dash, a medium dash for number ranges like
       1--2,
and a punctuation
       dash---like
this.

A sentence-ending space should be larger than the
space between words within a sentence.  You
sometimes have to type special commands in
conjunction with punctuation characters to get
this right, as in the following sentence.
       Gnats, gnus, etc.\ all  % `\ ' makes an inter-word space.
       begin with G\@.         % \@ marks end-of-sentence punctuation.
You should check the spaces after periods when
reading your output to make sure you haven't
forgotten any special cases.  Generating an
ellipsis
       \ldots\               % `\ ' is needed after `\ldots' because TeX
                             % ignores spaces after command names like \ldots
                             % made from \ + letters.
                             %
                             % Note how a `%' character causes TeX to ignore
                             % the end of the input line, so these blank lines
                             % do not start a new paragraph.
                             %
with the right spacing around the periods requires
a special command.

\LaTeX\ interprets some common characters as
commands, so you must type special commands to
generate them.  These characters include the
following:
       \$ \& \% \# \{ and \}.

In printing, text is usually emphasized with an
       \emph{italic}
type style.

\begin{em}
   A long segment of text can also be emphasized
   in this way.  Text within such a segment can be
   given \emph{additional} emphasis.
\end{em}

It is sometimes necessary to prevent \LaTeX\ from
breaking a line where it might otherwise do so.
This may be at a space, as between the ``Mr.''\ and
``Jones'' in
       ``Mr.~Jones'',        % ~ produces an unbreakable interword space.
or within a word---especially when the word is a
symbol like
       \mbox{\emph{itemnum}}
that makes little sense when hyphenated across
lines.

Footnotes\footnote{This is an example of a footnote.}
pose no problem.

\LaTeX\ is good at typesetting mathematical formulas
like
       \( x-3y + z = 7 \)
or
       \( a_{1} > x^{2n} + y^{2n} > x' \)
or
       \( \ip{A}{B} = \sum_{i} a_{i} b_{i} \).
The spaces you type in a formula are
ignored.  Remember that a letter like
       $x$                   % $ ... $  and  \( ... \)  are equivalent
is a formula when it denotes a mathematical
symbol, and it should be typed as one.

\section{Displayed Text}

Text is displayed by indenting it from the left
margin.  Quotations are commonly displayed.  There
are short quotations
\begin{quote}
   This is a short quotation.  It consists of a
   single paragraph of text.  See how it is formatted.
\end{quote}
and longer ones.
\begin{quotation}
   This is a longer quotation.  It consists of two
   paragraphs of text, neither of which are
   particularly interesting.

   This is the second paragraph of the quotation.  It
   is just as dull as the first paragraph.
\end{quotation}
Another frequently-displayed structure is a list.
The following is an example of an \emph{itemized}
list.
\begin{itemize}
   \item This is the first item of an itemized list.
         Each item in the list is marked with a ``tick''.
         You don't have to worry about what kind of tick
         mark is used.

   \item This is the second item of the list.  It
         contains another list nested inside it.  The inner
         list is an \emph{enumerated} list.
         \begin{enumerate}
            \item This is the first item of an enumerated
                  list that is nested within the itemized list.

            \item This is the second item of the inner list.
                  \LaTeX\ allows you to nest lists deeper than
                  you really should.
         \end{enumerate}
         This is the rest of the second item of the outer
         list.  It is no more interesting than any other
         part of the item.
   \item This is the third item of the list.
\end{itemize}
You can even display poetry.
\begin{verse}
   There is an environment
    for verse \\             % The \\ command separates lines
   Whose features some poets % within a stanza.
   will curse.

                             % One or more blank lines separate stanzas.

   For instead of making\\
   Them do \emph{all} line breaking, \\
   It allows them to put too many words on a line when they'd rather be
   forced to be terse.
\end{verse}

Mathematical formulas may also be displayed.  A
displayed formula
is
one-line long; multiline
formulas require special formatting instructions.
   \[  \ip{\Gamma}{\psi'} = x'' + y^{2} + z_{i}^{n}\]
Don't start a paragraph with a displayed equation,
nor make one a paragraph by itself.

\end{document}               % End of document.

\section{Bench Section 162}
% synthetic bench section 162
% This is a sample LaTeX input file.  (Version of 12 August 2004.)
%
% A '%' character causes TeX to ignore all remaining text on the line,
% and is used for comments like this one.

\documentclass{article}      % Specifies the document class

                             % The preamble begins here.
\title{An Example Document}  % Declares the document's title.
\author{Leslie Lamport}      % Declares the author's name.
\date{January 21, 1994}      % Deleting this command produces today's date.

\newcommand{\ip}[2]{(#1, #2)}
                             % Defines \ip{arg1}{arg2} to mean
                             % (arg1, arg2).

%\newcommand{\ip}[2]{\langle #1 | #2\rangle}
                             % This is an alternative definition of
                             % \ip that is commented out.

\begin{document}             % End of preamble and beginning of text.

\maketitle                   % Produces the title.

This is an example input file.  Comparing it with
the output it generates can show you how to
produce a simple document of your own.

\section{Ordinary Text}      % Produces section heading.  Lower-level
                             % sections are begun with similar
                             % \subsection and \subsubsection commands.

The ends  of words and sentences are marked
  by   spaces. It  doesn't matter how many
spaces    you type; one is as good as 100.  The
end of   a line counts as a space.

One   or more   blank lines denote the  end
of  a paragraph.

Since any number of consecutive spaces are treated
like a single one, the formatting of the input
file makes no difference to
      \LaTeX,                % The \LaTeX command generates the LaTeX logo.
but it makes a difference to you.  When you use
\LaTeX, making your input file as easy to read
as possible will be a great help as you write
your document and when you change it.  This sample
file shows how you can add comments to your own input
file.

Because printing is different from typewriting,
there are a number of things that you have to do
differently when preparing an input file than if
you were just typing the document directly.
Quotation marks like
       ``this''
have to be handled specially, as do quotes within
quotes:
       ``\,`this'            % \, separates the double and single quote.
        is what I just
        wrote, not  `that'\,''.

Dashes come in three sizes: an
       intra-word
dash, a medium dash for number ranges like
       1--2,
and a punctuation
       dash---like
this.

A sentence-ending space should be larger than the
space between words within a sentence.  You
sometimes have to type special commands in
conjunction with punctuation characters to get
this right, as in the following sentence.
       Gnats, gnus, etc.\ all  % `\ ' makes an inter-word space.
       begin with G\@.         % \@ marks end-of-sentence punctuation.
You should check the spaces after periods when
reading your output to make sure you haven't
forgotten any special cases.  Generating an
ellipsis
       \ldots\               % `\ ' is needed after `\ldots' because TeX
                             % ignores spaces after command names like \ldots
                             % made from \ + letters.
                             %
                             % Note how a `%' character causes TeX to ignore
                             % the end of the input line, so these blank lines
                             % do not start a new paragraph.
                             %
with the right spacing around the periods requires
a special command.

\LaTeX\ interprets some common characters as
commands, so you must type special commands to
generate them.  These characters include the
following:
       \$ \& \% \# \{ and \}.

In printing, text is usually emphasized with an
       \emph{italic}
type style.

\begin{em}
   A long segment of text can also be emphasized
   in this way.  Text within such a segment can be
   given \emph{additional} emphasis.
\end{em}

It is sometimes necessary to prevent \LaTeX\ from
breaking a line where it might otherwise do so.
This may be at a space, as between the ``Mr.''\ and
``Jones'' in
       ``Mr.~Jones'',        % ~ produces an unbreakable interword space.
or within a word---especially when the word is a
symbol like
       \mbox{\emph{itemnum}}
that makes little sense when hyphenated across
lines.

Footnotes\footnote{This is an example of a footnote.}
pose no problem.

\LaTeX\ is good at typesetting mathematical formulas
like
       \( x-3y + z = 7 \)
or
       \( a_{1} > x^{2n} + y^{2n} > x' \)
or
       \( \ip{A}{B} = \sum_{i} a_{i} b_{i} \).
The spaces you type in a formula are
ignored.  Remember that a letter like
       $x$                   % $ ... $  and  \( ... \)  are equivalent
is a formula when it denotes a mathematical
symbol, and it should be typed as one.

\section{Displayed Text}

Text is displayed by indenting it from the left
margin.  Quotations are commonly displayed.  There
are short quotations
\begin{quote}
   This is a short quotation.  It consists of a
   single paragraph of text.  See how it is formatted.
\end{quote}
and longer ones.
\begin{quotation}
   This is a longer quotation.  It consists of two
   paragraphs of text, neither of which are
   particularly interesting.

   This is the second paragraph of the quotation.  It
   is just as dull as the first paragraph.
\end{quotation}
Another frequently-displayed structure is a list.
The following is an example of an \emph{itemized}
list.
\begin{itemize}
   \item This is the first item of an itemized list.
         Each item in the list is marked with a ``tick''.
         You don't have to worry about what kind of tick
         mark is used.

   \item This is the second item of the list.  It
         contains another list nested inside it.  The inner
         list is an \emph{enumerated} list.
         \begin{enumerate}
            \item This is the first item of an enumerated
                  list that is nested within the itemized list.

            \item This is the second item of the inner list.
                  \LaTeX\ allows you to nest lists deeper than
                  you really should.
         \end{enumerate}
         This is the rest of the second item of the outer
         list.  It is no more interesting than any other
         part of the item.
   \item This is the third item of the list.
\end{itemize}
You can even display poetry.
\begin{verse}
   There is an environment
    for verse \\             % The \\ command separates lines
   Whose features some poets % within a stanza.
   will curse.

                             % One or more blank lines separate stanzas.

   For instead of making\\
   Them do \emph{all} line breaking, \\
   It allows them to put too many words on a line when they'd rather be
   forced to be terse.
\end{verse}

Mathematical formulas may also be displayed.  A
displayed formula
is
one-line long; multiline
formulas require special formatting instructions.
   \[  \ip{\Gamma}{\psi'} = x'' + y^{2} + z_{i}^{n}\]
Don't start a paragraph with a displayed equation,
nor make one a paragraph by itself.

\end{document}               % End of document.

\section{Bench Section 163}
% synthetic bench section 163
% This is a sample LaTeX input file.  (Version of 12 August 2004.)
%
% A '%' character causes TeX to ignore all remaining text on the line,
% and is used for comments like this one.

\documentclass{article}      % Specifies the document class

                             % The preamble begins here.
\title{An Example Document}  % Declares the document's title.
\author{Leslie Lamport}      % Declares the author's name.
\date{January 21, 1994}      % Deleting this command produces today's date.

\newcommand{\ip}[2]{(#1, #2)}
                             % Defines \ip{arg1}{arg2} to mean
                             % (arg1, arg2).

%\newcommand{\ip}[2]{\langle #1 | #2\rangle}
                             % This is an alternative definition of
                             % \ip that is commented out.

\begin{document}             % End of preamble and beginning of text.

\maketitle                   % Produces the title.

This is an example input file.  Comparing it with
the output it generates can show you how to
produce a simple document of your own.

\section{Ordinary Text}      % Produces section heading.  Lower-level
                             % sections are begun with similar
                             % \subsection and \subsubsection commands.

The ends  of words and sentences are marked
  by   spaces. It  doesn't matter how many
spaces    you type; one is as good as 100.  The
end of   a line counts as a space.

One   or more   blank lines denote the  end
of  a paragraph.

Since any number of consecutive spaces are treated
like a single one, the formatting of the input
file makes no difference to
      \LaTeX,                % The \LaTeX command generates the LaTeX logo.
but it makes a difference to you.  When you use
\LaTeX, making your input file as easy to read
as possible will be a great help as you write
your document and when you change it.  This sample
file shows how you can add comments to your own input
file.

Because printing is different from typewriting,
there are a number of things that you have to do
differently when preparing an input file than if
you were just typing the document directly.
Quotation marks like
       ``this''
have to be handled specially, as do quotes within
quotes:
       ``\,`this'            % \, separates the double and single quote.
        is what I just
        wrote, not  `that'\,''.

Dashes come in three sizes: an
       intra-word
dash, a medium dash for number ranges like
       1--2,
and a punctuation
       dash---like
this.

A sentence-ending space should be larger than the
space between words within a sentence.  You
sometimes have to type special commands in
conjunction with punctuation characters to get
this right, as in the following sentence.
       Gnats, gnus, etc.\ all  % `\ ' makes an inter-word space.
       begin with G\@.         % \@ marks end-of-sentence punctuation.
You should check the spaces after periods when
reading your output to make sure you haven't
forgotten any special cases.  Generating an
ellipsis
       \ldots\               % `\ ' is needed after `\ldots' because TeX
                             % ignores spaces after command names like \ldots
                             % made from \ + letters.
                             %
                             % Note how a `%' character causes TeX to ignore
                             % the end of the input line, so these blank lines
                             % do not start a new paragraph.
                             %
with the right spacing around the periods requires
a special command.

\LaTeX\ interprets some common characters as
commands, so you must type special commands to
generate them.  These characters include the
following:
       \$ \& \% \# \{ and \}.

In printing, text is usually emphasized with an
       \emph{italic}
type style.

\begin{em}
   A long segment of text can also be emphasized
   in this way.  Text within such a segment can be
   given \emph{additional} emphasis.
\end{em}

It is sometimes necessary to prevent \LaTeX\ from
breaking a line where it might otherwise do so.
This may be at a space, as between the ``Mr.''\ and
``Jones'' in
       ``Mr.~Jones'',        % ~ produces an unbreakable interword space.
or within a word---especially when the word is a
symbol like
       \mbox{\emph{itemnum}}
that makes little sense when hyphenated across
lines.

Footnotes\footnote{This is an example of a footnote.}
pose no problem.

\LaTeX\ is good at typesetting mathematical formulas
like
       \( x-3y + z = 7 \)
or
       \( a_{1} > x^{2n} + y^{2n} > x' \)
or
       \( \ip{A}{B} = \sum_{i} a_{i} b_{i} \).
The spaces you type in a formula are
ignored.  Remember that a letter like
       $x$                   % $ ... $  and  \( ... \)  are equivalent
is a formula when it denotes a mathematical
symbol, and it should be typed as one.

\section{Displayed Text}

Text is displayed by indenting it from the left
margin.  Quotations are commonly displayed.  There
are short quotations
\begin{quote}
   This is a short quotation.  It consists of a
   single paragraph of text.  See how it is formatted.
\end{quote}
and longer ones.
\begin{quotation}
   This is a longer quotation.  It consists of two
   paragraphs of text, neither of which are
   particularly interesting.

   This is the second paragraph of the quotation.  It
   is just as dull as the first paragraph.
\end{quotation}
Another frequently-displayed structure is a list.
The following is an example of an \emph{itemized}
list.
\begin{itemize}
   \item This is the first item of an itemized list.
         Each item in the list is marked with a ``tick''.
         You don't have to worry about what kind of tick
         mark is used.

   \item This is the second item of the list.  It
         contains another list nested inside it.  The inner
         list is an \emph{enumerated} list.
         \begin{enumerate}
            \item This is the first item of an enumerated
                  list that is nested within the itemized list.

            \item This is the second item of the inner list.
                  \LaTeX\ allows you to nest lists deeper than
                  you really should.
         \end{enumerate}
         This is the rest of the second item of the outer
         list.  It is no more interesting than any other
         part of the item.
   \item This is the third item of the list.
\end{itemize}
You can even display poetry.
\begin{verse}
   There is an environment
    for verse \\             % The \\ command separates lines
   Whose features some poets % within a stanza.
   will curse.

                             % One or more blank lines separate stanzas.

   For instead of making\\
   Them do \emph{all} line breaking, \\
   It allows them to put too many words on a line when they'd rather be
   forced to be terse.
\end{verse}

Mathematical formulas may also be displayed.  A
displayed formula
is
one-line long; multiline
formulas require special formatting instructions.
   \[  \ip{\Gamma}{\psi'} = x'' + y^{2} + z_{i}^{n}\]
Don't start a paragraph with a displayed equation,
nor make one a paragraph by itself.

\end{document}               % End of document.

\section{Bench Section 164}
% synthetic bench section 164
% This is a sample LaTeX input file.  (Version of 12 August 2004.)
%
% A '%' character causes TeX to ignore all remaining text on the line,
% and is used for comments like this one.

\documentclass{article}      % Specifies the document class

                             % The preamble begins here.
\title{An Example Document}  % Declares the document's title.
\author{Leslie Lamport}      % Declares the author's name.
\date{January 21, 1994}      % Deleting this command produces today's date.

\newcommand{\ip}[2]{(#1, #2)}
                             % Defines \ip{arg1}{arg2} to mean
                             % (arg1, arg2).

%\newcommand{\ip}[2]{\langle #1 | #2\rangle}
                             % This is an alternative definition of
                             % \ip that is commented out.

\begin{document}             % End of preamble and beginning of text.

\maketitle                   % Produces the title.

This is an example input file.  Comparing it with
the output it generates can show you how to
produce a simple document of your own.

\section{Ordinary Text}      % Produces section heading.  Lower-level
                             % sections are begun with similar
                             % \subsection and \subsubsection commands.

The ends  of words and sentences are marked
  by   spaces. It  doesn't matter how many
spaces    you type; one is as good as 100.  The
end of   a line counts as a space.

One   or more   blank lines denote the  end
of  a paragraph.

Since any number of consecutive spaces are treated
like a single one, the formatting of the input
file makes no difference to
      \LaTeX,                % The \LaTeX command generates the LaTeX logo.
but it makes a difference to you.  When you use
\LaTeX, making your input file as easy to read
as possible will be a great help as you write
your document and when you change it.  This sample
file shows how you can add comments to your own input
file.

Because printing is different from typewriting,
there are a number of things that you have to do
differently when preparing an input file than if
you were just typing the document directly.
Quotation marks like
       ``this''
have to be handled specially, as do quotes within
quotes:
       ``\,`this'            % \, separates the double and single quote.
        is what I just
        wrote, not  `that'\,''.

Dashes come in three sizes: an
       intra-word
dash, a medium dash for number ranges like
       1--2,
and a punctuation
       dash---like
this.

A sentence-ending space should be larger than the
space between words within a sentence.  You
sometimes have to type special commands in
conjunction with punctuation characters to get
this right, as in the following sentence.
       Gnats, gnus, etc.\ all  % `\ ' makes an inter-word space.
       begin with G\@.         % \@ marks end-of-sentence punctuation.
You should check the spaces after periods when
reading your output to make sure you haven't
forgotten any special cases.  Generating an
ellipsis
       \ldots\               % `\ ' is needed after `\ldots' because TeX
                             % ignores spaces after command names like \ldots
                             % made from \ + letters.
                             %
                             % Note how a `%' character causes TeX to ignore
                             % the end of the input line, so these blank lines
                             % do not start a new paragraph.
                             %
with the right spacing around the periods requires
a special command.

\LaTeX\ interprets some common characters as
commands, so you must type special commands to
generate them.  These characters include the
following:
       \$ \& \% \# \{ and \}.

In printing, text is usually emphasized with an
       \emph{italic}
type style.

\begin{em}
   A long segment of text can also be emphasized
   in this way.  Text within such a segment can be
   given \emph{additional} emphasis.
\end{em}

It is sometimes necessary to prevent \LaTeX\ from
breaking a line where it might otherwise do so.
This may be at a space, as between the ``Mr.''\ and
``Jones'' in
       ``Mr.~Jones'',        % ~ produces an unbreakable interword space.
or within a word---especially when the word is a
symbol like
       \mbox{\emph{itemnum}}
that makes little sense when hyphenated across
lines.

Footnotes\footnote{This is an example of a footnote.}
pose no problem.

\LaTeX\ is good at typesetting mathematical formulas
like
       \( x-3y + z = 7 \)
or
       \( a_{1} > x^{2n} + y^{2n} > x' \)
or
       \( \ip{A}{B} = \sum_{i} a_{i} b_{i} \).
The spaces you type in a formula are
ignored.  Remember that a letter like
       $x$                   % $ ... $  and  \( ... \)  are equivalent
is a formula when it denotes a mathematical
symbol, and it should be typed as one.

\section{Displayed Text}

Text is displayed by indenting it from the left
margin.  Quotations are commonly displayed.  There
are short quotations
\begin{quote}
   This is a short quotation.  It consists of a
   single paragraph of text.  See how it is formatted.
\end{quote}
and longer ones.
\begin{quotation}
   This is a longer quotation.  It consists of two
   paragraphs of text, neither of which are
   particularly interesting.

   This is the second paragraph of the quotation.  It
   is just as dull as the first paragraph.
\end{quotation}
Another frequently-displayed structure is a list.
The following is an example of an \emph{itemized}
list.
\begin{itemize}
   \item This is the first item of an itemized list.
         Each item in the list is marked with a ``tick''.
         You don't have to worry about what kind of tick
         mark is used.

   \item This is the second item of the list.  It
         contains another list nested inside it.  The inner
         list is an \emph{enumerated} list.
         \begin{enumerate}
            \item This is the first item of an enumerated
                  list that is nested within the itemized list.

            \item This is the second item of the inner list.
                  \LaTeX\ allows you to nest lists deeper than
                  you really should.
         \end{enumerate}
         This is the rest of the second item of the outer
         list.  It is no more interesting than any other
         part of the item.
   \item This is the third item of the list.
\end{itemize}
You can even display poetry.
\begin{verse}
   There is an environment
    for verse \\             % The \\ command separates lines
   Whose features some poets % within a stanza.
   will curse.

                             % One or more blank lines separate stanzas.

   For instead of making\\
   Them do \emph{all} line breaking, \\
   It allows them to put too many words on a line when they'd rather be
   forced to be terse.
\end{verse}

Mathematical formulas may also be displayed.  A
displayed formula
is
one-line long; multiline
formulas require special formatting instructions.
   \[  \ip{\Gamma}{\psi'} = x'' + y^{2} + z_{i}^{n}\]
Don't start a paragraph with a displayed equation,
nor make one a paragraph by itself.

\end{document}               % End of document.

\section{Bench Section 165}
% synthetic bench section 165
% This is a sample LaTeX input file.  (Version of 12 August 2004.)
%
% A '%' character causes TeX to ignore all remaining text on the line,
% and is used for comments like this one.

\documentclass{article}      % Specifies the document class

                             % The preamble begins here.
\title{An Example Document}  % Declares the document's title.
\author{Leslie Lamport}      % Declares the author's name.
\date{January 21, 1994}      % Deleting this command produces today's date.

\newcommand{\ip}[2]{(#1, #2)}
                             % Defines \ip{arg1}{arg2} to mean
                             % (arg1, arg2).

%\newcommand{\ip}[2]{\langle #1 | #2\rangle}
                             % This is an alternative definition of
                             % \ip that is commented out.

\begin{document}             % End of preamble and beginning of text.

\maketitle                   % Produces the title.

This is an example input file.  Comparing it with
the output it generates can show you how to
produce a simple document of your own.

\section{Ordinary Text}      % Produces section heading.  Lower-level
                             % sections are begun with similar
                             % \subsection and \subsubsection commands.

The ends  of words and sentences are marked
  by   spaces. It  doesn't matter how many
spaces    you type; one is as good as 100.  The
end of   a line counts as a space.

One   or more   blank lines denote the  end
of  a paragraph.

Since any number of consecutive spaces are treated
like a single one, the formatting of the input
file makes no difference to
      \LaTeX,                % The \LaTeX command generates the LaTeX logo.
but it makes a difference to you.  When you use
\LaTeX, making your input file as easy to read
as possible will be a great help as you write
your document and when you change it.  This sample
file shows how you can add comments to your own input
file.

Because printing is different from typewriting,
there are a number of things that you have to do
differently when preparing an input file than if
you were just typing the document directly.
Quotation marks like
       ``this''
have to be handled specially, as do quotes within
quotes:
       ``\,`this'            % \, separates the double and single quote.
        is what I just
        wrote, not  `that'\,''.

Dashes come in three sizes: an
       intra-word
dash, a medium dash for number ranges like
       1--2,
and a punctuation
       dash---like
this.

A sentence-ending space should be larger than the
space between words within a sentence.  You
sometimes have to type special commands in
conjunction with punctuation characters to get
this right, as in the following sentence.
       Gnats, gnus, etc.\ all  % `\ ' makes an inter-word space.
       begin with G\@.         % \@ marks end-of-sentence punctuation.
You should check the spaces after periods when
reading your output to make sure you haven't
forgotten any special cases.  Generating an
ellipsis
       \ldots\               % `\ ' is needed after `\ldots' because TeX
                             % ignores spaces after command names like \ldots
                             % made from \ + letters.
                             %
                             % Note how a `%' character causes TeX to ignore
                             % the end of the input line, so these blank lines
                             % do not start a new paragraph.
                             %
with the right spacing around the periods requires
a special command.

\LaTeX\ interprets some common characters as
commands, so you must type special commands to
generate them.  These characters include the
following:
       \$ \& \% \# \{ and \}.

In printing, text is usually emphasized with an
       \emph{italic}
type style.

\begin{em}
   A long segment of text can also be emphasized
   in this way.  Text within such a segment can be
   given \emph{additional} emphasis.
\end{em}

It is sometimes necessary to prevent \LaTeX\ from
breaking a line where it might otherwise do so.
This may be at a space, as between the ``Mr.''\ and
``Jones'' in
       ``Mr.~Jones'',        % ~ produces an unbreakable interword space.
or within a word---especially when the word is a
symbol like
       \mbox{\emph{itemnum}}
that makes little sense when hyphenated across
lines.

Footnotes\footnote{This is an example of a footnote.}
pose no problem.

\LaTeX\ is good at typesetting mathematical formulas
like
       \( x-3y + z = 7 \)
or
       \( a_{1} > x^{2n} + y^{2n} > x' \)
or
       \( \ip{A}{B} = \sum_{i} a_{i} b_{i} \).
The spaces you type in a formula are
ignored.  Remember that a letter like
       $x$                   % $ ... $  and  \( ... \)  are equivalent
is a formula when it denotes a mathematical
symbol, and it should be typed as one.

\section{Displayed Text}

Text is displayed by indenting it from the left
margin.  Quotations are commonly displayed.  There
are short quotations
\begin{quote}
   This is a short quotation.  It consists of a
   single paragraph of text.  See how it is formatted.
\end{quote}
and longer ones.
\begin{quotation}
   This is a longer quotation.  It consists of two
   paragraphs of text, neither of which are
   particularly interesting.

   This is the second paragraph of the quotation.  It
   is just as dull as the first paragraph.
\end{quotation}
Another frequently-displayed structure is a list.
The following is an example of an \emph{itemized}
list.
\begin{itemize}
   \item This is the first item of an itemized list.
         Each item in the list is marked with a ``tick''.
         You don't have to worry about what kind of tick
         mark is used.

   \item This is the second item of the list.  It
         contains another list nested inside it.  The inner
         list is an \emph{enumerated} list.
         \begin{enumerate}
            \item This is the first item of an enumerated
                  list that is nested within the itemized list.

            \item This is the second item of the inner list.
                  \LaTeX\ allows you to nest lists deeper than
                  you really should.
         \end{enumerate}
         This is the rest of the second item of the outer
         list.  It is no more interesting than any other
         part of the item.
   \item This is the third item of the list.
\end{itemize}
You can even display poetry.
\begin{verse}
   There is an environment
    for verse \\             % The \\ command separates lines
   Whose features some poets % within a stanza.
   will curse.

                             % One or more blank lines separate stanzas.

   For instead of making\\
   Them do \emph{all} line breaking, \\
   It allows them to put too many words on a line when they'd rather be
   forced to be terse.
\end{verse}

Mathematical formulas may also be displayed.  A
displayed formula
is
one-line long; multiline
formulas require special formatting instructions.
   \[  \ip{\Gamma}{\psi'} = x'' + y^{2} + z_{i}^{n}\]
Don't start a paragraph with a displayed equation,
nor make one a paragraph by itself.

\end{document}               % End of document.

\section{Bench Section 166}
% synthetic bench section 166
% This is a sample LaTeX input file.  (Version of 12 August 2004.)
%
% A '%' character causes TeX to ignore all remaining text on the line,
% and is used for comments like this one.

\documentclass{article}      % Specifies the document class

                             % The preamble begins here.
\title{An Example Document}  % Declares the document's title.
\author{Leslie Lamport}      % Declares the author's name.
\date{January 21, 1994}      % Deleting this command produces today's date.

\newcommand{\ip}[2]{(#1, #2)}
                             % Defines \ip{arg1}{arg2} to mean
                             % (arg1, arg2).

%\newcommand{\ip}[2]{\langle #1 | #2\rangle}
                             % This is an alternative definition of
                             % \ip that is commented out.

\begin{document}             % End of preamble and beginning of text.

\maketitle                   % Produces the title.

This is an example input file.  Comparing it with
the output it generates can show you how to
produce a simple document of your own.

\section{Ordinary Text}      % Produces section heading.  Lower-level
                             % sections are begun with similar
                             % \subsection and \subsubsection commands.

The ends  of words and sentences are marked
  by   spaces. It  doesn't matter how many
spaces    you type; one is as good as 100.  The
end of   a line counts as a space.

One   or more   blank lines denote the  end
of  a paragraph.

Since any number of consecutive spaces are treated
like a single one, the formatting of the input
file makes no difference to
      \LaTeX,                % The \LaTeX command generates the LaTeX logo.
but it makes a difference to you.  When you use
\LaTeX, making your input file as easy to read
as possible will be a great help as you write
your document and when you change it.  This sample
file shows how you can add comments to your own input
file.

Because printing is different from typewriting,
there are a number of things that you have to do
differently when preparing an input file than if
you were just typing the document directly.
Quotation marks like
       ``this''
have to be handled specially, as do quotes within
quotes:
       ``\,`this'            % \, separates the double and single quote.
        is what I just
        wrote, not  `that'\,''.

Dashes come in three sizes: an
       intra-word
dash, a medium dash for number ranges like
       1--2,
and a punctuation
       dash---like
this.

A sentence-ending space should be larger than the
space between words within a sentence.  You
sometimes have to type special commands in
conjunction with punctuation characters to get
this right, as in the following sentence.
       Gnats, gnus, etc.\ all  % `\ ' makes an inter-word space.
       begin with G\@.         % \@ marks end-of-sentence punctuation.
You should check the spaces after periods when
reading your output to make sure you haven't
forgotten any special cases.  Generating an
ellipsis
       \ldots\               % `\ ' is needed after `\ldots' because TeX
                             % ignores spaces after command names like \ldots
                             % made from \ + letters.
                             %
                             % Note how a `%' character causes TeX to ignore
                             % the end of the input line, so these blank lines
                             % do not start a new paragraph.
                             %
with the right spacing around the periods requires
a special command.

\LaTeX\ interprets some common characters as
commands, so you must type special commands to
generate them.  These characters include the
following:
       \$ \& \% \# \{ and \}.

In printing, text is usually emphasized with an
       \emph{italic}
type style.

\begin{em}
   A long segment of text can also be emphasized
   in this way.  Text within such a segment can be
   given \emph{additional} emphasis.
\end{em}

It is sometimes necessary to prevent \LaTeX\ from
breaking a line where it might otherwise do so.
This may be at a space, as between the ``Mr.''\ and
``Jones'' in
       ``Mr.~Jones'',        % ~ produces an unbreakable interword space.
or within a word---especially when the word is a
symbol like
       \mbox{\emph{itemnum}}
that makes little sense when hyphenated across
lines.

Footnotes\footnote{This is an example of a footnote.}
pose no problem.

\LaTeX\ is good at typesetting mathematical formulas
like
       \( x-3y + z = 7 \)
or
       \( a_{1} > x^{2n} + y^{2n} > x' \)
or
       \( \ip{A}{B} = \sum_{i} a_{i} b_{i} \).
The spaces you type in a formula are
ignored.  Remember that a letter like
       $x$                   % $ ... $  and  \( ... \)  are equivalent
is a formula when it denotes a mathematical
symbol, and it should be typed as one.

\section{Displayed Text}

Text is displayed by indenting it from the left
margin.  Quotations are commonly displayed.  There
are short quotations
\begin{quote}
   This is a short quotation.  It consists of a
   single paragraph of text.  See how it is formatted.
\end{quote}
and longer ones.
\begin{quotation}
   This is a longer quotation.  It consists of two
   paragraphs of text, neither of which are
   particularly interesting.

   This is the second paragraph of the quotation.  It
   is just as dull as the first paragraph.
\end{quotation}
Another frequently-displayed structure is a list.
The following is an example of an \emph{itemized}
list.
\begin{itemize}
   \item This is the first item of an itemized list.
         Each item in the list is marked with a ``tick''.
         You don't have to worry about what kind of tick
         mark is used.

   \item This is the second item of the list.  It
         contains another list nested inside it.  The inner
         list is an \emph{enumerated} list.
         \begin{enumerate}
            \item This is the first item of an enumerated
                  list that is nested within the itemized list.

            \item This is the second item of the inner list.
                  \LaTeX\ allows you to nest lists deeper than
                  you really should.
         \end{enumerate}
         This is the rest of the second item of the outer
         list.  It is no more interesting than any other
         part of the item.
   \item This is the third item of the list.
\end{itemize}
You can even display poetry.
\begin{verse}
   There is an environment
    for verse \\             % The \\ command separates lines
   Whose features some poets % within a stanza.
   will curse.

                             % One or more blank lines separate stanzas.

   For instead of making\\
   Them do \emph{all} line breaking, \\
   It allows them to put too many words on a line when they'd rather be
   forced to be terse.
\end{verse}

Mathematical formulas may also be displayed.  A
displayed formula
is
one-line long; multiline
formulas require special formatting instructions.
   \[  \ip{\Gamma}{\psi'} = x'' + y^{2} + z_{i}^{n}\]
Don't start a paragraph with a displayed equation,
nor make one a paragraph by itself.

\end{document}               % End of document.

\section{Bench Section 167}
% synthetic bench section 167
% This is a sample LaTeX input file.  (Version of 12 August 2004.)
%
% A '%' character causes TeX to ignore all remaining text on the line,
% and is used for comments like this one.

\documentclass{article}      % Specifies the document class

                             % The preamble begins here.
\title{An Example Document}  % Declares the document's title.
\author{Leslie Lamport}      % Declares the author's name.
\date{January 21, 1994}      % Deleting this command produces today's date.

\newcommand{\ip}[2]{(#1, #2)}
                             % Defines \ip{arg1}{arg2} to mean
                             % (arg1, arg2).

%\newcommand{\ip}[2]{\langle #1 | #2\rangle}
                             % This is an alternative definition of
                             % \ip that is commented out.

\begin{document}             % End of preamble and beginning of text.

\maketitle                   % Produces the title.

This is an example input file.  Comparing it with
the output it generates can show you how to
produce a simple document of your own.

\section{Ordinary Text}      % Produces section heading.  Lower-level
                             % sections are begun with similar
                             % \subsection and \subsubsection commands.

The ends  of words and sentences are marked
  by   spaces. It  doesn't matter how many
spaces    you type; one is as good as 100.  The
end of   a line counts as a space.

One   or more   blank lines denote the  end
of  a paragraph.

Since any number of consecutive spaces are treated
like a single one, the formatting of the input
file makes no difference to
      \LaTeX,                % The \LaTeX command generates the LaTeX logo.
but it makes a difference to you.  When you use
\LaTeX, making your input file as easy to read
as possible will be a great help as you write
your document and when you change it.  This sample
file shows how you can add comments to your own input
file.

Because printing is different from typewriting,
there are a number of things that you have to do
differently when preparing an input file than if
you were just typing the document directly.
Quotation marks like
       ``this''
have to be handled specially, as do quotes within
quotes:
       ``\,`this'            % \, separates the double and single quote.
        is what I just
        wrote, not  `that'\,''.

Dashes come in three sizes: an
       intra-word
dash, a medium dash for number ranges like
       1--2,
and a punctuation
       dash---like
this.

A sentence-ending space should be larger than the
space between words within a sentence.  You
sometimes have to type special commands in
conjunction with punctuation characters to get
this right, as in the following sentence.
       Gnats, gnus, etc.\ all  % `\ ' makes an inter-word space.
       begin with G\@.         % \@ marks end-of-sentence punctuation.
You should check the spaces after periods when
reading your output to make sure you haven't
forgotten any special cases.  Generating an
ellipsis
       \ldots\               % `\ ' is needed after `\ldots' because TeX
                             % ignores spaces after command names like \ldots
                             % made from \ + letters.
                             %
                             % Note how a `%' character causes TeX to ignore
                             % the end of the input line, so these blank lines
                             % do not start a new paragraph.
                             %
with the right spacing around the periods requires
a special command.

\LaTeX\ interprets some common characters as
commands, so you must type special commands to
generate them.  These characters include the
following:
       \$ \& \% \# \{ and \}.

In printing, text is usually emphasized with an
       \emph{italic}
type style.

\begin{em}
   A long segment of text can also be emphasized
   in this way.  Text within such a segment can be
   given \emph{additional} emphasis.
\end{em}

It is sometimes necessary to prevent \LaTeX\ from
breaking a line where it might otherwise do so.
This may be at a space, as between the ``Mr.''\ and
``Jones'' in
       ``Mr.~Jones'',        % ~ produces an unbreakable interword space.
or within a word---especially when the word is a
symbol like
       \mbox{\emph{itemnum}}
that makes little sense when hyphenated across
lines.

Footnotes\footnote{This is an example of a footnote.}
pose no problem.

\LaTeX\ is good at typesetting mathematical formulas
like
       \( x-3y + z = 7 \)
or
       \( a_{1} > x^{2n} + y^{2n} > x' \)
or
       \( \ip{A}{B} = \sum_{i} a_{i} b_{i} \).
The spaces you type in a formula are
ignored.  Remember that a letter like
       $x$                   % $ ... $  and  \( ... \)  are equivalent
is a formula when it denotes a mathematical
symbol, and it should be typed as one.

\section{Displayed Text}

Text is displayed by indenting it from the left
margin.  Quotations are commonly displayed.  There
are short quotations
\begin{quote}
   This is a short quotation.  It consists of a
   single paragraph of text.  See how it is formatted.
\end{quote}
and longer ones.
\begin{quotation}
   This is a longer quotation.  It consists of two
   paragraphs of text, neither of which are
   particularly interesting.

   This is the second paragraph of the quotation.  It
   is just as dull as the first paragraph.
\end{quotation}
Another frequently-displayed structure is a list.
The following is an example of an \emph{itemized}
list.
\begin{itemize}
   \item This is the first item of an itemized list.
         Each item in the list is marked with a ``tick''.
         You don't have to worry about what kind of tick
         mark is used.

   \item This is the second item of the list.  It
         contains another list nested inside it.  The inner
         list is an \emph{enumerated} list.
         \begin{enumerate}
            \item This is the first item of an enumerated
                  list that is nested within the itemized list.

            \item This is the second item of the inner list.
                  \LaTeX\ allows you to nest lists deeper than
                  you really should.
         \end{enumerate}
         This is the rest of the second item of the outer
         list.  It is no more interesting than any other
         part of the item.
   \item This is the third item of the list.
\end{itemize}
You can even display poetry.
\begin{verse}
   There is an environment
    for verse \\             % The \\ command separates lines
   Whose features some poets % within a stanza.
   will curse.

                             % One or more blank lines separate stanzas.

   For instead of making\\
   Them do \emph{all} line breaking, \\
   It allows them to put too many words on a line when they'd rather be
   forced to be terse.
\end{verse}

Mathematical formulas may also be displayed.  A
displayed formula
is
one-line long; multiline
formulas require special formatting instructions.
   \[  \ip{\Gamma}{\psi'} = x'' + y^{2} + z_{i}^{n}\]
Don't start a paragraph with a displayed equation,
nor make one a paragraph by itself.

\end{document}               % End of document.

\section{Bench Section 168}
% synthetic bench section 168
% This is a sample LaTeX input file.  (Version of 12 August 2004.)
%
% A '%' character causes TeX to ignore all remaining text on the line,
% and is used for comments like this one.

\documentclass{article}      % Specifies the document class

                             % The preamble begins here.
\title{An Example Document}  % Declares the document's title.
\author{Leslie Lamport}      % Declares the author's name.
\date{January 21, 1994}      % Deleting this command produces today's date.

\newcommand{\ip}[2]{(#1, #2)}
                             % Defines \ip{arg1}{arg2} to mean
                             % (arg1, arg2).

%\newcommand{\ip}[2]{\langle #1 | #2\rangle}
                             % This is an alternative definition of
                             % \ip that is commented out.

\begin{document}             % End of preamble and beginning of text.

\maketitle                   % Produces the title.

This is an example input file.  Comparing it with
the output it generates can show you how to
produce a simple document of your own.

\section{Ordinary Text}      % Produces section heading.  Lower-level
                             % sections are begun with similar
                             % \subsection and \subsubsection commands.

The ends  of words and sentences are marked
  by   spaces. It  doesn't matter how many
spaces    you type; one is as good as 100.  The
end of   a line counts as a space.

One   or more   blank lines denote the  end
of  a paragraph.

Since any number of consecutive spaces are treated
like a single one, the formatting of the input
file makes no difference to
      \LaTeX,                % The \LaTeX command generates the LaTeX logo.
but it makes a difference to you.  When you use
\LaTeX, making your input file as easy to read
as possible will be a great help as you write
your document and when you change it.  This sample
file shows how you can add comments to your own input
file.

Because printing is different from typewriting,
there are a number of things that you have to do
differently when preparing an input file than if
you were just typing the document directly.
Quotation marks like
       ``this''
have to be handled specially, as do quotes within
quotes:
       ``\,`this'            % \, separates the double and single quote.
        is what I just
        wrote, not  `that'\,''.

Dashes come in three sizes: an
       intra-word
dash, a medium dash for number ranges like
       1--2,
and a punctuation
       dash---like
this.

A sentence-ending space should be larger than the
space between words within a sentence.  You
sometimes have to type special commands in
conjunction with punctuation characters to get
this right, as in the following sentence.
       Gnats, gnus, etc.\ all  % `\ ' makes an inter-word space.
       begin with G\@.         % \@ marks end-of-sentence punctuation.
You should check the spaces after periods when
reading your output to make sure you haven't
forgotten any special cases.  Generating an
ellipsis
       \ldots\               % `\ ' is needed after `\ldots' because TeX
                             % ignores spaces after command names like \ldots
                             % made from \ + letters.
                             %
                             % Note how a `%' character causes TeX to ignore
                             % the end of the input line, so these blank lines
                             % do not start a new paragraph.
                             %
with the right spacing around the periods requires
a special command.

\LaTeX\ interprets some common characters as
commands, so you must type special commands to
generate them.  These characters include the
following:
       \$ \& \% \# \{ and \}.

In printing, text is usually emphasized with an
       \emph{italic}
type style.

\begin{em}
   A long segment of text can also be emphasized
   in this way.  Text within such a segment can be
   given \emph{additional} emphasis.
\end{em}

It is sometimes necessary to prevent \LaTeX\ from
breaking a line where it might otherwise do so.
This may be at a space, as between the ``Mr.''\ and
``Jones'' in
       ``Mr.~Jones'',        % ~ produces an unbreakable interword space.
or within a word---especially when the word is a
symbol like
       \mbox{\emph{itemnum}}
that makes little sense when hyphenated across
lines.

Footnotes\footnote{This is an example of a footnote.}
pose no problem.

\LaTeX\ is good at typesetting mathematical formulas
like
       \( x-3y + z = 7 \)
or
       \( a_{1} > x^{2n} + y^{2n} > x' \)
or
       \( \ip{A}{B} = \sum_{i} a_{i} b_{i} \).
The spaces you type in a formula are
ignored.  Remember that a letter like
       $x$                   % $ ... $  and  \( ... \)  are equivalent
is a formula when it denotes a mathematical
symbol, and it should be typed as one.

\section{Displayed Text}

Text is displayed by indenting it from the left
margin.  Quotations are commonly displayed.  There
are short quotations
\begin{quote}
   This is a short quotation.  It consists of a
   single paragraph of text.  See how it is formatted.
\end{quote}
and longer ones.
\begin{quotation}
   This is a longer quotation.  It consists of two
   paragraphs of text, neither of which are
   particularly interesting.

   This is the second paragraph of the quotation.  It
   is just as dull as the first paragraph.
\end{quotation}
Another frequently-displayed structure is a list.
The following is an example of an \emph{itemized}
list.
\begin{itemize}
   \item This is the first item of an itemized list.
         Each item in the list is marked with a ``tick''.
         You don't have to worry about what kind of tick
         mark is used.

   \item This is the second item of the list.  It
         contains another list nested inside it.  The inner
         list is an \emph{enumerated} list.
         \begin{enumerate}
            \item This is the first item of an enumerated
                  list that is nested within the itemized list.

            \item This is the second item of the inner list.
                  \LaTeX\ allows you to nest lists deeper than
                  you really should.
         \end{enumerate}
         This is the rest of the second item of the outer
         list.  It is no more interesting than any other
         part of the item.
   \item This is the third item of the list.
\end{itemize}
You can even display poetry.
\begin{verse}
   There is an environment
    for verse \\             % The \\ command separates lines
   Whose features some poets % within a stanza.
   will curse.

                             % One or more blank lines separate stanzas.

   For instead of making\\
   Them do \emph{all} line breaking, \\
   It allows them to put too many words on a line when they'd rather be
   forced to be terse.
\end{verse}

Mathematical formulas may also be displayed.  A
displayed formula
is
one-line long; multiline
formulas require special formatting instructions.
   \[  \ip{\Gamma}{\psi'} = x'' + y^{2} + z_{i}^{n}\]
Don't start a paragraph with a displayed equation,
nor make one a paragraph by itself.

\end{document}               % End of document.

\section{Bench Section 169}
% synthetic bench section 169
% This is a sample LaTeX input file.  (Version of 12 August 2004.)
%
% A '%' character causes TeX to ignore all remaining text on the line,
% and is used for comments like this one.

\documentclass{article}      % Specifies the document class

                             % The preamble begins here.
\title{An Example Document}  % Declares the document's title.
\author{Leslie Lamport}      % Declares the author's name.
\date{January 21, 1994}      % Deleting this command produces today's date.

\newcommand{\ip}[2]{(#1, #2)}
                             % Defines \ip{arg1}{arg2} to mean
                             % (arg1, arg2).

%\newcommand{\ip}[2]{\langle #1 | #2\rangle}
                             % This is an alternative definition of
                             % \ip that is commented out.

\begin{document}             % End of preamble and beginning of text.

\maketitle                   % Produces the title.

This is an example input file.  Comparing it with
the output it generates can show you how to
produce a simple document of your own.

\section{Ordinary Text}      % Produces section heading.  Lower-level
                             % sections are begun with similar
                             % \subsection and \subsubsection commands.

The ends  of words and sentences are marked
  by   spaces. It  doesn't matter how many
spaces    you type; one is as good as 100.  The
end of   a line counts as a space.

One   or more   blank lines denote the  end
of  a paragraph.

Since any number of consecutive spaces are treated
like a single one, the formatting of the input
file makes no difference to
      \LaTeX,                % The \LaTeX command generates the LaTeX logo.
but it makes a difference to you.  When you use
\LaTeX, making your input file as easy to read
as possible will be a great help as you write
your document and when you change it.  This sample
file shows how you can add comments to your own input
file.

Because printing is different from typewriting,
there are a number of things that you have to do
differently when preparing an input file than if
you were just typing the document directly.
Quotation marks like
       ``this''
have to be handled specially, as do quotes within
quotes:
       ``\,`this'            % \, separates the double and single quote.
        is what I just
        wrote, not  `that'\,''.

Dashes come in three sizes: an
       intra-word
dash, a medium dash for number ranges like
       1--2,
and a punctuation
       dash---like
this.

A sentence-ending space should be larger than the
space between words within a sentence.  You
sometimes have to type special commands in
conjunction with punctuation characters to get
this right, as in the following sentence.
       Gnats, gnus, etc.\ all  % `\ ' makes an inter-word space.
       begin with G\@.         % \@ marks end-of-sentence punctuation.
You should check the spaces after periods when
reading your output to make sure you haven't
forgotten any special cases.  Generating an
ellipsis
       \ldots\               % `\ ' is needed after `\ldots' because TeX
                             % ignores spaces after command names like \ldots
                             % made from \ + letters.
                             %
                             % Note how a `%' character causes TeX to ignore
                             % the end of the input line, so these blank lines
                             % do not start a new paragraph.
                             %
with the right spacing around the periods requires
a special command.

\LaTeX\ interprets some common characters as
commands, so you must type special commands to
generate them.  These characters include the
following:
       \$ \& \% \# \{ and \}.

In printing, text is usually emphasized with an
       \emph{italic}
type style.

\begin{em}
   A long segment of text can also be emphasized
   in this way.  Text within such a segment can be
   given \emph{additional} emphasis.
\end{em}

It is sometimes necessary to prevent \LaTeX\ from
breaking a line where it might otherwise do so.
This may be at a space, as between the ``Mr.''\ and
``Jones'' in
       ``Mr.~Jones'',        % ~ produces an unbreakable interword space.
or within a word---especially when the word is a
symbol like
       \mbox{\emph{itemnum}}
that makes little sense when hyphenated across
lines.

Footnotes\footnote{This is an example of a footnote.}
pose no problem.

\LaTeX\ is good at typesetting mathematical formulas
like
       \( x-3y + z = 7 \)
or
       \( a_{1} > x^{2n} + y^{2n} > x' \)
or
       \( \ip{A}{B} = \sum_{i} a_{i} b_{i} \).
The spaces you type in a formula are
ignored.  Remember that a letter like
       $x$                   % $ ... $  and  \( ... \)  are equivalent
is a formula when it denotes a mathematical
symbol, and it should be typed as one.

\section{Displayed Text}

Text is displayed by indenting it from the left
margin.  Quotations are commonly displayed.  There
are short quotations
\begin{quote}
   This is a short quotation.  It consists of a
   single paragraph of text.  See how it is formatted.
\end{quote}
and longer ones.
\begin{quotation}
   This is a longer quotation.  It consists of two
   paragraphs of text, neither of which are
   particularly interesting.

   This is the second paragraph of the quotation.  It
   is just as dull as the first paragraph.
\end{quotation}
Another frequently-displayed structure is a list.
The following is an example of an \emph{itemized}
list.
\begin{itemize}
   \item This is the first item of an itemized list.
         Each item in the list is marked with a ``tick''.
         You don't have to worry about what kind of tick
         mark is used.

   \item This is the second item of the list.  It
         contains another list nested inside it.  The inner
         list is an \emph{enumerated} list.
         \begin{enumerate}
            \item This is the first item of an enumerated
                  list that is nested within the itemized list.

            \item This is the second item of the inner list.
                  \LaTeX\ allows you to nest lists deeper than
                  you really should.
         \end{enumerate}
         This is the rest of the second item of the outer
         list.  It is no more interesting than any other
         part of the item.
   \item This is the third item of the list.
\end{itemize}
You can even display poetry.
\begin{verse}
   There is an environment
    for verse \\             % The \\ command separates lines
   Whose features some poets % within a stanza.
   will curse.

                             % One or more blank lines separate stanzas.

   For instead of making\\
   Them do \emph{all} line breaking, \\
   It allows them to put too many words on a line when they'd rather be
   forced to be terse.
\end{verse}

Mathematical formulas may also be displayed.  A
displayed formula
is
one-line long; multiline
formulas require special formatting instructions.
   \[  \ip{\Gamma}{\psi'} = x'' + y^{2} + z_{i}^{n}\]
Don't start a paragraph with a displayed equation,
nor make one a paragraph by itself.

\end{document}               % End of document.

\section{Bench Section 170}
% synthetic bench section 170
% This is a sample LaTeX input file.  (Version of 12 August 2004.)
%
% A '%' character causes TeX to ignore all remaining text on the line,
% and is used for comments like this one.

\documentclass{article}      % Specifies the document class

                             % The preamble begins here.
\title{An Example Document}  % Declares the document's title.
\author{Leslie Lamport}      % Declares the author's name.
\date{January 21, 1994}      % Deleting this command produces today's date.

\newcommand{\ip}[2]{(#1, #2)}
                             % Defines \ip{arg1}{arg2} to mean
                             % (arg1, arg2).

%\newcommand{\ip}[2]{\langle #1 | #2\rangle}
                             % This is an alternative definition of
                             % \ip that is commented out.

\begin{document}             % End of preamble and beginning of text.

\maketitle                   % Produces the title.

This is an example input file.  Comparing it with
the output it generates can show you how to
produce a simple document of your own.

\section{Ordinary Text}      % Produces section heading.  Lower-level
                             % sections are begun with similar
                             % \subsection and \subsubsection commands.

The ends  of words and sentences are marked
  by   spaces. It  doesn't matter how many
spaces    you type; one is as good as 100.  The
end of   a line counts as a space.

One   or more   blank lines denote the  end
of  a paragraph.

Since any number of consecutive spaces are treated
like a single one, the formatting of the input
file makes no difference to
      \LaTeX,                % The \LaTeX command generates the LaTeX logo.
but it makes a difference to you.  When you use
\LaTeX, making your input file as easy to read
as possible will be a great help as you write
your document and when you change it.  This sample
file shows how you can add comments to your own input
file.

Because printing is different from typewriting,
there are a number of things that you have to do
differently when preparing an input file than if
you were just typing the document directly.
Quotation marks like
       ``this''
have to be handled specially, as do quotes within
quotes:
       ``\,`this'            % \, separates the double and single quote.
        is what I just
        wrote, not  `that'\,''.

Dashes come in three sizes: an
       intra-word
dash, a medium dash for number ranges like
       1--2,
and a punctuation
       dash---like
this.

A sentence-ending space should be larger than the
space between words within a sentence.  You
sometimes have to type special commands in
conjunction with punctuation characters to get
this right, as in the following sentence.
       Gnats, gnus, etc.\ all  % `\ ' makes an inter-word space.
       begin with G\@.         % \@ marks end-of-sentence punctuation.
You should check the spaces after periods when
reading your output to make sure you haven't
forgotten any special cases.  Generating an
ellipsis
       \ldots\               % `\ ' is needed after `\ldots' because TeX
                             % ignores spaces after command names like \ldots
                             % made from \ + letters.
                             %
                             % Note how a `%' character causes TeX to ignore
                             % the end of the input line, so these blank lines
                             % do not start a new paragraph.
                             %
with the right spacing around the periods requires
a special command.

\LaTeX\ interprets some common characters as
commands, so you must type special commands to
generate them.  These characters include the
following:
       \$ \& \% \# \{ and \}.

In printing, text is usually emphasized with an
       \emph{italic}
type style.

\begin{em}
   A long segment of text can also be emphasized
   in this way.  Text within such a segment can be
   given \emph{additional} emphasis.
\end{em}

It is sometimes necessary to prevent \LaTeX\ from
breaking a line where it might otherwise do so.
This may be at a space, as between the ``Mr.''\ and
``Jones'' in
       ``Mr.~Jones'',        % ~ produces an unbreakable interword space.
or within a word---especially when the word is a
symbol like
       \mbox{\emph{itemnum}}
that makes little sense when hyphenated across
lines.

Footnotes\footnote{This is an example of a footnote.}
pose no problem.

\LaTeX\ is good at typesetting mathematical formulas
like
       \( x-3y + z = 7 \)
or
       \( a_{1} > x^{2n} + y^{2n} > x' \)
or
       \( \ip{A}{B} = \sum_{i} a_{i} b_{i} \).
The spaces you type in a formula are
ignored.  Remember that a letter like
       $x$                   % $ ... $  and  \( ... \)  are equivalent
is a formula when it denotes a mathematical
symbol, and it should be typed as one.

\section{Displayed Text}

Text is displayed by indenting it from the left
margin.  Quotations are commonly displayed.  There
are short quotations
\begin{quote}
   This is a short quotation.  It consists of a
   single paragraph of text.  See how it is formatted.
\end{quote}
and longer ones.
\begin{quotation}
   This is a longer quotation.  It consists of two
   paragraphs of text, neither of which are
   particularly interesting.

   This is the second paragraph of the quotation.  It
   is just as dull as the first paragraph.
\end{quotation}
Another frequently-displayed structure is a list.
The following is an example of an \emph{itemized}
list.
\begin{itemize}
   \item This is the first item of an itemized list.
         Each item in the list is marked with a ``tick''.
         You don't have to worry about what kind of tick
         mark is used.

   \item This is the second item of the list.  It
         contains another list nested inside it.  The inner
         list is an \emph{enumerated} list.
         \begin{enumerate}
            \item This is the first item of an enumerated
                  list that is nested within the itemized list.

            \item This is the second item of the inner list.
                  \LaTeX\ allows you to nest lists deeper than
                  you really should.
         \end{enumerate}
         This is the rest of the second item of the outer
         list.  It is no more interesting than any other
         part of the item.
   \item This is the third item of the list.
\end{itemize}
You can even display poetry.
\begin{verse}
   There is an environment
    for verse \\             % The \\ command separates lines
   Whose features some poets % within a stanza.
   will curse.

                             % One or more blank lines separate stanzas.

   For instead of making\\
   Them do \emph{all} line breaking, \\
   It allows them to put too many words on a line when they'd rather be
   forced to be terse.
\end{verse}

Mathematical formulas may also be displayed.  A
displayed formula
is
one-line long; multiline
formulas require special formatting instructions.
   \[  \ip{\Gamma}{\psi'} = x'' + y^{2} + z_{i}^{n}\]
Don't start a paragraph with a displayed equation,
nor make one a paragraph by itself.

\end{document}               % End of document.

\section{Bench Section 171}
% synthetic bench section 171
% This is a sample LaTeX input file.  (Version of 12 August 2004.)
%
% A '%' character causes TeX to ignore all remaining text on the line,
% and is used for comments like this one.

\documentclass{article}      % Specifies the document class

                             % The preamble begins here.
\title{An Example Document}  % Declares the document's title.
\author{Leslie Lamport}      % Declares the author's name.
\date{January 21, 1994}      % Deleting this command produces today's date.

\newcommand{\ip}[2]{(#1, #2)}
                             % Defines \ip{arg1}{arg2} to mean
                             % (arg1, arg2).

%\newcommand{\ip}[2]{\langle #1 | #2\rangle}
                             % This is an alternative definition of
                             % \ip that is commented out.

\begin{document}             % End of preamble and beginning of text.

\maketitle                   % Produces the title.

This is an example input file.  Comparing it with
the output it generates can show you how to
produce a simple document of your own.

\section{Ordinary Text}      % Produces section heading.  Lower-level
                             % sections are begun with similar
                             % \subsection and \subsubsection commands.

The ends  of words and sentences are marked
  by   spaces. It  doesn't matter how many
spaces    you type; one is as good as 100.  The
end of   a line counts as a space.

One   or more   blank lines denote the  end
of  a paragraph.

Since any number of consecutive spaces are treated
like a single one, the formatting of the input
file makes no difference to
      \LaTeX,                % The \LaTeX command generates the LaTeX logo.
but it makes a difference to you.  When you use
\LaTeX, making your input file as easy to read
as possible will be a great help as you write
your document and when you change it.  This sample
file shows how you can add comments to your own input
file.

Because printing is different from typewriting,
there are a number of things that you have to do
differently when preparing an input file than if
you were just typing the document directly.
Quotation marks like
       ``this''
have to be handled specially, as do quotes within
quotes:
       ``\,`this'            % \, separates the double and single quote.
        is what I just
        wrote, not  `that'\,''.

Dashes come in three sizes: an
       intra-word
dash, a medium dash for number ranges like
       1--2,
and a punctuation
       dash---like
this.

A sentence-ending space should be larger than the
space between words within a sentence.  You
sometimes have to type special commands in
conjunction with punctuation characters to get
this right, as in the following sentence.
       Gnats, gnus, etc.\ all  % `\ ' makes an inter-word space.
       begin with G\@.         % \@ marks end-of-sentence punctuation.
You should check the spaces after periods when
reading your output to make sure you haven't
forgotten any special cases.  Generating an
ellipsis
       \ldots\               % `\ ' is needed after `\ldots' because TeX
                             % ignores spaces after command names like \ldots
                             % made from \ + letters.
                             %
                             % Note how a `%' character causes TeX to ignore
                             % the end of the input line, so these blank lines
                             % do not start a new paragraph.
                             %
with the right spacing around the periods requires
a special command.

\LaTeX\ interprets some common characters as
commands, so you must type special commands to
generate them.  These characters include the
following:
       \$ \& \% \# \{ and \}.

In printing, text is usually emphasized with an
       \emph{italic}
type style.

\begin{em}
   A long segment of text can also be emphasized
   in this way.  Text within such a segment can be
   given \emph{additional} emphasis.
\end{em}

It is sometimes necessary to prevent \LaTeX\ from
breaking a line where it might otherwise do so.
This may be at a space, as between the ``Mr.''\ and
``Jones'' in
       ``Mr.~Jones'',        % ~ produces an unbreakable interword space.
or within a word---especially when the word is a
symbol like
       \mbox{\emph{itemnum}}
that makes little sense when hyphenated across
lines.

Footnotes\footnote{This is an example of a footnote.}
pose no problem.

\LaTeX\ is good at typesetting mathematical formulas
like
       \( x-3y + z = 7 \)
or
       \( a_{1} > x^{2n} + y^{2n} > x' \)
or
       \( \ip{A}{B} = \sum_{i} a_{i} b_{i} \).
The spaces you type in a formula are
ignored.  Remember that a letter like
       $x$                   % $ ... $  and  \( ... \)  are equivalent
is a formula when it denotes a mathematical
symbol, and it should be typed as one.

\section{Displayed Text}

Text is displayed by indenting it from the left
margin.  Quotations are commonly displayed.  There
are short quotations
\begin{quote}
   This is a short quotation.  It consists of a
   single paragraph of text.  See how it is formatted.
\end{quote}
and longer ones.
\begin{quotation}
   This is a longer quotation.  It consists of two
   paragraphs of text, neither of which are
   particularly interesting.

   This is the second paragraph of the quotation.  It
   is just as dull as the first paragraph.
\end{quotation}
Another frequently-displayed structure is a list.
The following is an example of an \emph{itemized}
list.
\begin{itemize}
   \item This is the first item of an itemized list.
         Each item in the list is marked with a ``tick''.
         You don't have to worry about what kind of tick
         mark is used.

   \item This is the second item of the list.  It
         contains another list nested inside it.  The inner
         list is an \emph{enumerated} list.
         \begin{enumerate}
            \item This is the first item of an enumerated
                  list that is nested within the itemized list.

            \item This is the second item of the inner list.
                  \LaTeX\ allows you to nest lists deeper than
                  you really should.
         \end{enumerate}
         This is the rest of the second item of the outer
         list.  It is no more interesting than any other
         part of the item.
   \item This is the third item of the list.
\end{itemize}
You can even display poetry.
\begin{verse}
   There is an environment
    for verse \\             % The \\ command separates lines
   Whose features some poets % within a stanza.
   will curse.

                             % One or more blank lines separate stanzas.

   For instead of making\\
   Them do \emph{all} line breaking, \\
   It allows them to put too many words on a line when they'd rather be
   forced to be terse.
\end{verse}

Mathematical formulas may also be displayed.  A
displayed formula
is
one-line long; multiline
formulas require special formatting instructions.
   \[  \ip{\Gamma}{\psi'} = x'' + y^{2} + z_{i}^{n}\]
Don't start a paragraph with a displayed equation,
nor make one a paragraph by itself.

\end{document}               % End of document.

\section{Bench Section 172}
% synthetic bench section 172
% This is a sample LaTeX input file.  (Version of 12 August 2004.)
%
% A '%' character causes TeX to ignore all remaining text on the line,
% and is used for comments like this one.

\documentclass{article}      % Specifies the document class

                             % The preamble begins here.
\title{An Example Document}  % Declares the document's title.
\author{Leslie Lamport}      % Declares the author's name.
\date{January 21, 1994}      % Deleting this command produces today's date.

\newcommand{\ip}[2]{(#1, #2)}
                             % Defines \ip{arg1}{arg2} to mean
                             % (arg1, arg2).

%\newcommand{\ip}[2]{\langle #1 | #2\rangle}
                             % This is an alternative definition of
                             % \ip that is commented out.

\begin{document}             % End of preamble and beginning of text.

\maketitle                   % Produces the title.

This is an example input file.  Comparing it with
the output it generates can show you how to
produce a simple document of your own.

\section{Ordinary Text}      % Produces section heading.  Lower-level
                             % sections are begun with similar
                             % \subsection and \subsubsection commands.

The ends  of words and sentences are marked
  by   spaces. It  doesn't matter how many
spaces    you type; one is as good as 100.  The
end of   a line counts as a space.

One   or more   blank lines denote the  end
of  a paragraph.

Since any number of consecutive spaces are treated
like a single one, the formatting of the input
file makes no difference to
      \LaTeX,                % The \LaTeX command generates the LaTeX logo.
but it makes a difference to you.  When you use
\LaTeX, making your input file as easy to read
as possible will be a great help as you write
your document and when you change it.  This sample
file shows how you can add comments to your own input
file.

Because printing is different from typewriting,
there are a number of things that you have to do
differently when preparing an input file than if
you were just typing the document directly.
Quotation marks like
       ``this''
have to be handled specially, as do quotes within
quotes:
       ``\,`this'            % \, separates the double and single quote.
        is what I just
        wrote, not  `that'\,''.

Dashes come in three sizes: an
       intra-word
dash, a medium dash for number ranges like
       1--2,
and a punctuation
       dash---like
this.

A sentence-ending space should be larger than the
space between words within a sentence.  You
sometimes have to type special commands in
conjunction with punctuation characters to get
this right, as in the following sentence.
       Gnats, gnus, etc.\ all  % `\ ' makes an inter-word space.
       begin with G\@.         % \@ marks end-of-sentence punctuation.
You should check the spaces after periods when
reading your output to make sure you haven't
forgotten any special cases.  Generating an
ellipsis
       \ldots\               % `\ ' is needed after `\ldots' because TeX
                             % ignores spaces after command names like \ldots
                             % made from \ + letters.
                             %
                             % Note how a `%' character causes TeX to ignore
                             % the end of the input line, so these blank lines
                             % do not start a new paragraph.
                             %
with the right spacing around the periods requires
a special command.

\LaTeX\ interprets some common characters as
commands, so you must type special commands to
generate them.  These characters include the
following:
       \$ \& \% \# \{ and \}.

In printing, text is usually emphasized with an
       \emph{italic}
type style.

\begin{em}
   A long segment of text can also be emphasized
   in this way.  Text within such a segment can be
   given \emph{additional} emphasis.
\end{em}

It is sometimes necessary to prevent \LaTeX\ from
breaking a line where it might otherwise do so.
This may be at a space, as between the ``Mr.''\ and
``Jones'' in
       ``Mr.~Jones'',        % ~ produces an unbreakable interword space.
or within a word---especially when the word is a
symbol like
       \mbox{\emph{itemnum}}
that makes little sense when hyphenated across
lines.

Footnotes\footnote{This is an example of a footnote.}
pose no problem.

\LaTeX\ is good at typesetting mathematical formulas
like
       \( x-3y + z = 7 \)
or
       \( a_{1} > x^{2n} + y^{2n} > x' \)
or
       \( \ip{A}{B} = \sum_{i} a_{i} b_{i} \).
The spaces you type in a formula are
ignored.  Remember that a letter like
       $x$                   % $ ... $  and  \( ... \)  are equivalent
is a formula when it denotes a mathematical
symbol, and it should be typed as one.

\section{Displayed Text}

Text is displayed by indenting it from the left
margin.  Quotations are commonly displayed.  There
are short quotations
\begin{quote}
   This is a short quotation.  It consists of a
   single paragraph of text.  See how it is formatted.
\end{quote}
and longer ones.
\begin{quotation}
   This is a longer quotation.  It consists of two
   paragraphs of text, neither of which are
   particularly interesting.

   This is the second paragraph of the quotation.  It
   is just as dull as the first paragraph.
\end{quotation}
Another frequently-displayed structure is a list.
The following is an example of an \emph{itemized}
list.
\begin{itemize}
   \item This is the first item of an itemized list.
         Each item in the list is marked with a ``tick''.
         You don't have to worry about what kind of tick
         mark is used.

   \item This is the second item of the list.  It
         contains another list nested inside it.  The inner
         list is an \emph{enumerated} list.
         \begin{enumerate}
            \item This is the first item of an enumerated
                  list that is nested within the itemized list.

            \item This is the second item of the inner list.
                  \LaTeX\ allows you to nest lists deeper than
                  you really should.
         \end{enumerate}
         This is the rest of the second item of the outer
         list.  It is no more interesting than any other
         part of the item.
   \item This is the third item of the list.
\end{itemize}
You can even display poetry.
\begin{verse}
   There is an environment
    for verse \\             % The \\ command separates lines
   Whose features some poets % within a stanza.
   will curse.

                             % One or more blank lines separate stanzas.

   For instead of making\\
   Them do \emph{all} line breaking, \\
   It allows them to put too many words on a line when they'd rather be
   forced to be terse.
\end{verse}

Mathematical formulas may also be displayed.  A
displayed formula
is
one-line long; multiline
formulas require special formatting instructions.
   \[  \ip{\Gamma}{\psi'} = x'' + y^{2} + z_{i}^{n}\]
Don't start a paragraph with a displayed equation,
nor make one a paragraph by itself.

\end{document}               % End of document.

\section{Bench Section 173}
% synthetic bench section 173
% This is a sample LaTeX input file.  (Version of 12 August 2004.)
%
% A '%' character causes TeX to ignore all remaining text on the line,
% and is used for comments like this one.

\documentclass{article}      % Specifies the document class

                             % The preamble begins here.
\title{An Example Document}  % Declares the document's title.
\author{Leslie Lamport}      % Declares the author's name.
\date{January 21, 1994}      % Deleting this command produces today's date.

\newcommand{\ip}[2]{(#1, #2)}
                             % Defines \ip{arg1}{arg2} to mean
                             % (arg1, arg2).

%\newcommand{\ip}[2]{\langle #1 | #2\rangle}
                             % This is an alternative definition of
                             % \ip that is commented out.

\begin{document}             % End of preamble and beginning of text.

\maketitle                   % Produces the title.

This is an example input file.  Comparing it with
the output it generates can show you how to
produce a simple document of your own.

\section{Ordinary Text}      % Produces section heading.  Lower-level
                             % sections are begun with similar
                             % \subsection and \subsubsection commands.

The ends  of words and sentences are marked
  by   spaces. It  doesn't matter how many
spaces    you type; one is as good as 100.  The
end of   a line counts as a space.

One   or more   blank lines denote the  end
of  a paragraph.

Since any number of consecutive spaces are treated
like a single one, the formatting of the input
file makes no difference to
      \LaTeX,                % The \LaTeX command generates the LaTeX logo.
but it makes a difference to you.  When you use
\LaTeX, making your input file as easy to read
as possible will be a great help as you write
your document and when you change it.  This sample
file shows how you can add comments to your own input
file.

Because printing is different from typewriting,
there are a number of things that you have to do
differently when preparing an input file than if
you were just typing the document directly.
Quotation marks like
       ``this''
have to be handled specially, as do quotes within
quotes:
       ``\,`this'            % \, separates the double and single quote.
        is what I just
        wrote, not  `that'\,''.

Dashes come in three sizes: an
       intra-word
dash, a medium dash for number ranges like
       1--2,
and a punctuation
       dash---like
this.

A sentence-ending space should be larger than the
space between words within a sentence.  You
sometimes have to type special commands in
conjunction with punctuation characters to get
this right, as in the following sentence.
       Gnats, gnus, etc.\ all  % `\ ' makes an inter-word space.
       begin with G\@.         % \@ marks end-of-sentence punctuation.
You should check the spaces after periods when
reading your output to make sure you haven't
forgotten any special cases.  Generating an
ellipsis
       \ldots\               % `\ ' is needed after `\ldots' because TeX
                             % ignores spaces after command names like \ldots
                             % made from \ + letters.
                             %
                             % Note how a `%' character causes TeX to ignore
                             % the end of the input line, so these blank lines
                             % do not start a new paragraph.
                             %
with the right spacing around the periods requires
a special command.

\LaTeX\ interprets some common characters as
commands, so you must type special commands to
generate them.  These characters include the
following:
       \$ \& \% \# \{ and \}.

In printing, text is usually emphasized with an
       \emph{italic}
type style.

\begin{em}
   A long segment of text can also be emphasized
   in this way.  Text within such a segment can be
   given \emph{additional} emphasis.
\end{em}

It is sometimes necessary to prevent \LaTeX\ from
breaking a line where it might otherwise do so.
This may be at a space, as between the ``Mr.''\ and
``Jones'' in
       ``Mr.~Jones'',        % ~ produces an unbreakable interword space.
or within a word---especially when the word is a
symbol like
       \mbox{\emph{itemnum}}
that makes little sense when hyphenated across
lines.

Footnotes\footnote{This is an example of a footnote.}
pose no problem.

\LaTeX\ is good at typesetting mathematical formulas
like
       \( x-3y + z = 7 \)
or
       \( a_{1} > x^{2n} + y^{2n} > x' \)
or
       \( \ip{A}{B} = \sum_{i} a_{i} b_{i} \).
The spaces you type in a formula are
ignored.  Remember that a letter like
       $x$                   % $ ... $  and  \( ... \)  are equivalent
is a formula when it denotes a mathematical
symbol, and it should be typed as one.

\section{Displayed Text}

Text is displayed by indenting it from the left
margin.  Quotations are commonly displayed.  There
are short quotations
\begin{quote}
   This is a short quotation.  It consists of a
   single paragraph of text.  See how it is formatted.
\end{quote}
and longer ones.
\begin{quotation}
   This is a longer quotation.  It consists of two
   paragraphs of text, neither of which are
   particularly interesting.

   This is the second paragraph of the quotation.  It
   is just as dull as the first paragraph.
\end{quotation}
Another frequently-displayed structure is a list.
The following is an example of an \emph{itemized}
list.
\begin{itemize}
   \item This is the first item of an itemized list.
         Each item in the list is marked with a ``tick''.
         You don't have to worry about what kind of tick
         mark is used.

   \item This is the second item of the list.  It
         contains another list nested inside it.  The inner
         list is an \emph{enumerated} list.
         \begin{enumerate}
            \item This is the first item of an enumerated
                  list that is nested within the itemized list.

            \item This is the second item of the inner list.
                  \LaTeX\ allows you to nest lists deeper than
                  you really should.
         \end{enumerate}
         This is the rest of the second item of the outer
         list.  It is no more interesting than any other
         part of the item.
   \item This is the third item of the list.
\end{itemize}
You can even display poetry.
\begin{verse}
   There is an environment
    for verse \\             % The \\ command separates lines
   Whose features some poets % within a stanza.
   will curse.

                             % One or more blank lines separate stanzas.

   For instead of making\\
   Them do \emph{all} line breaking, \\
   It allows them to put too many words on a line when they'd rather be
   forced to be terse.
\end{verse}

Mathematical formulas may also be displayed.  A
displayed formula
is
one-line long; multiline
formulas require special formatting instructions.
   \[  \ip{\Gamma}{\psi'} = x'' + y^{2} + z_{i}^{n}\]
Don't start a paragraph with a displayed equation,
nor make one a paragraph by itself.

\end{document}               % End of document.

\section{Bench Section 174}
% synthetic bench section 174
% This is a sample LaTeX input file.  (Version of 12 August 2004.)
%
% A '%' character causes TeX to ignore all remaining text on the line,
% and is used for comments like this one.

\documentclass{article}      % Specifies the document class

                             % The preamble begins here.
\title{An Example Document}  % Declares the document's title.
\author{Leslie Lamport}      % Declares the author's name.
\date{January 21, 1994}      % Deleting this command produces today's date.

\newcommand{\ip}[2]{(#1, #2)}
                             % Defines \ip{arg1}{arg2} to mean
                             % (arg1, arg2).

%\newcommand{\ip}[2]{\langle #1 | #2\rangle}
                             % This is an alternative definition of
                             % \ip that is commented out.

\begin{document}             % End of preamble and beginning of text.

\maketitle                   % Produces the title.

This is an example input file.  Comparing it with
the output it generates can show you how to
produce a simple document of your own.

\section{Ordinary Text}      % Produces section heading.  Lower-level
                             % sections are begun with similar
                             % \subsection and \subsubsection commands.

The ends  of words and sentences are marked
  by   spaces. It  doesn't matter how many
spaces    you type; one is as good as 100.  The
end of   a line counts as a space.

One   or more   blank lines denote the  end
of  a paragraph.

Since any number of consecutive spaces are treated
like a single one, the formatting of the input
file makes no difference to
      \LaTeX,                % The \LaTeX command generates the LaTeX logo.
but it makes a difference to you.  When you use
\LaTeX, making your input file as easy to read
as possible will be a great help as you write
your document and when you change it.  This sample
file shows how you can add comments to your own input
file.

Because printing is different from typewriting,
there are a number of things that you have to do
differently when preparing an input file than if
you were just typing the document directly.
Quotation marks like
       ``this''
have to be handled specially, as do quotes within
quotes:
       ``\,`this'            % \, separates the double and single quote.
        is what I just
        wrote, not  `that'\,''.

Dashes come in three sizes: an
       intra-word
dash, a medium dash for number ranges like
       1--2,
and a punctuation
       dash---like
this.

A sentence-ending space should be larger than the
space between words within a sentence.  You
sometimes have to type special commands in
conjunction with punctuation characters to get
this right, as in the following sentence.
       Gnats, gnus, etc.\ all  % `\ ' makes an inter-word space.
       begin with G\@.         % \@ marks end-of-sentence punctuation.
You should check the spaces after periods when
reading your output to make sure you haven't
forgotten any special cases.  Generating an
ellipsis
       \ldots\               % `\ ' is needed after `\ldots' because TeX
                             % ignores spaces after command names like \ldots
                             % made from \ + letters.
                             %
                             % Note how a `%' character causes TeX to ignore
                             % the end of the input line, so these blank lines
                             % do not start a new paragraph.
                             %
with the right spacing around the periods requires
a special command.

\LaTeX\ interprets some common characters as
commands, so you must type special commands to
generate them.  These characters include the
following:
       \$ \& \% \# \{ and \}.

In printing, text is usually emphasized with an
       \emph{italic}
type style.

\begin{em}
   A long segment of text can also be emphasized
   in this way.  Text within such a segment can be
   given \emph{additional} emphasis.
\end{em}

It is sometimes necessary to prevent \LaTeX\ from
breaking a line where it might otherwise do so.
This may be at a space, as between the ``Mr.''\ and
``Jones'' in
       ``Mr.~Jones'',        % ~ produces an unbreakable interword space.
or within a word---especially when the word is a
symbol like
       \mbox{\emph{itemnum}}
that makes little sense when hyphenated across
lines.

Footnotes\footnote{This is an example of a footnote.}
pose no problem.

\LaTeX\ is good at typesetting mathematical formulas
like
       \( x-3y + z = 7 \)
or
       \( a_{1} > x^{2n} + y^{2n} > x' \)
or
       \( \ip{A}{B} = \sum_{i} a_{i} b_{i} \).
The spaces you type in a formula are
ignored.  Remember that a letter like
       $x$                   % $ ... $  and  \( ... \)  are equivalent
is a formula when it denotes a mathematical
symbol, and it should be typed as one.

\section{Displayed Text}

Text is displayed by indenting it from the left
margin.  Quotations are commonly displayed.  There
are short quotations
\begin{quote}
   This is a short quotation.  It consists of a
   single paragraph of text.  See how it is formatted.
\end{quote}
and longer ones.
\begin{quotation}
   This is a longer quotation.  It consists of two
   paragraphs of text, neither of which are
   particularly interesting.

   This is the second paragraph of the quotation.  It
   is just as dull as the first paragraph.
\end{quotation}
Another frequently-displayed structure is a list.
The following is an example of an \emph{itemized}
list.
\begin{itemize}
   \item This is the first item of an itemized list.
         Each item in the list is marked with a ``tick''.
         You don't have to worry about what kind of tick
         mark is used.

   \item This is the second item of the list.  It
         contains another list nested inside it.  The inner
         list is an \emph{enumerated} list.
         \begin{enumerate}
            \item This is the first item of an enumerated
                  list that is nested within the itemized list.

            \item This is the second item of the inner list.
                  \LaTeX\ allows you to nest lists deeper than
                  you really should.
         \end{enumerate}
         This is the rest of the second item of the outer
         list.  It is no more interesting than any other
         part of the item.
   \item This is the third item of the list.
\end{itemize}
You can even display poetry.
\begin{verse}
   There is an environment
    for verse \\             % The \\ command separates lines
   Whose features some poets % within a stanza.
   will curse.

                             % One or more blank lines separate stanzas.

   For instead of making\\
   Them do \emph{all} line breaking, \\
   It allows them to put too many words on a line when they'd rather be
   forced to be terse.
\end{verse}

Mathematical formulas may also be displayed.  A
displayed formula
is
one-line long; multiline
formulas require special formatting instructions.
   \[  \ip{\Gamma}{\psi'} = x'' + y^{2} + z_{i}^{n}\]
Don't start a paragraph with a displayed equation,
nor make one a paragraph by itself.

\end{document}               % End of document.

\section{Bench Section 175}
% synthetic bench section 175
% This is a sample LaTeX input file.  (Version of 12 August 2004.)
%
% A '%' character causes TeX to ignore all remaining text on the line,
% and is used for comments like this one.

\documentclass{article}      % Specifies the document class

                             % The preamble begins here.
\title{An Example Document}  % Declares the document's title.
\author{Leslie Lamport}      % Declares the author's name.
\date{January 21, 1994}      % Deleting this command produces today's date.

\newcommand{\ip}[2]{(#1, #2)}
                             % Defines \ip{arg1}{arg2} to mean
                             % (arg1, arg2).

%\newcommand{\ip}[2]{\langle #1 | #2\rangle}
                             % This is an alternative definition of
                             % \ip that is commented out.

\begin{document}             % End of preamble and beginning of text.

\maketitle                   % Produces the title.

This is an example input file.  Comparing it with
the output it generates can show you how to
produce a simple document of your own.

\section{Ordinary Text}      % Produces section heading.  Lower-level
                             % sections are begun with similar
                             % \subsection and \subsubsection commands.

The ends  of words and sentences are marked
  by   spaces. It  doesn't matter how many
spaces    you type; one is as good as 100.  The
end of   a line counts as a space.

One   or more   blank lines denote the  end
of  a paragraph.

Since any number of consecutive spaces are treated
like a single one, the formatting of the input
file makes no difference to
      \LaTeX,                % The \LaTeX command generates the LaTeX logo.
but it makes a difference to you.  When you use
\LaTeX, making your input file as easy to read
as possible will be a great help as you write
your document and when you change it.  This sample
file shows how you can add comments to your own input
file.

Because printing is different from typewriting,
there are a number of things that you have to do
differently when preparing an input file than if
you were just typing the document directly.
Quotation marks like
       ``this''
have to be handled specially, as do quotes within
quotes:
       ``\,`this'            % \, separates the double and single quote.
        is what I just
        wrote, not  `that'\,''.

Dashes come in three sizes: an
       intra-word
dash, a medium dash for number ranges like
       1--2,
and a punctuation
       dash---like
this.

A sentence-ending space should be larger than the
space between words within a sentence.  You
sometimes have to type special commands in
conjunction with punctuation characters to get
this right, as in the following sentence.
       Gnats, gnus, etc.\ all  % `\ ' makes an inter-word space.
       begin with G\@.         % \@ marks end-of-sentence punctuation.
You should check the spaces after periods when
reading your output to make sure you haven't
forgotten any special cases.  Generating an
ellipsis
       \ldots\               % `\ ' is needed after `\ldots' because TeX
                             % ignores spaces after command names like \ldots
                             % made from \ + letters.
                             %
                             % Note how a `%' character causes TeX to ignore
                             % the end of the input line, so these blank lines
                             % do not start a new paragraph.
                             %
with the right spacing around the periods requires
a special command.

\LaTeX\ interprets some common characters as
commands, so you must type special commands to
generate them.  These characters include the
following:
       \$ \& \% \# \{ and \}.

In printing, text is usually emphasized with an
       \emph{italic}
type style.

\begin{em}
   A long segment of text can also be emphasized
   in this way.  Text within such a segment can be
   given \emph{additional} emphasis.
\end{em}

It is sometimes necessary to prevent \LaTeX\ from
breaking a line where it might otherwise do so.
This may be at a space, as between the ``Mr.''\ and
``Jones'' in
       ``Mr.~Jones'',        % ~ produces an unbreakable interword space.
or within a word---especially when the word is a
symbol like
       \mbox{\emph{itemnum}}
that makes little sense when hyphenated across
lines.

Footnotes\footnote{This is an example of a footnote.}
pose no problem.

\LaTeX\ is good at typesetting mathematical formulas
like
       \( x-3y + z = 7 \)
or
       \( a_{1} > x^{2n} + y^{2n} > x' \)
or
       \( \ip{A}{B} = \sum_{i} a_{i} b_{i} \).
The spaces you type in a formula are
ignored.  Remember that a letter like
       $x$                   % $ ... $  and  \( ... \)  are equivalent
is a formula when it denotes a mathematical
symbol, and it should be typed as one.

\section{Displayed Text}

Text is displayed by indenting it from the left
margin.  Quotations are commonly displayed.  There
are short quotations
\begin{quote}
   This is a short quotation.  It consists of a
   single paragraph of text.  See how it is formatted.
\end{quote}
and longer ones.
\begin{quotation}
   This is a longer quotation.  It consists of two
   paragraphs of text, neither of which are
   particularly interesting.

   This is the second paragraph of the quotation.  It
   is just as dull as the first paragraph.
\end{quotation}
Another frequently-displayed structure is a list.
The following is an example of an \emph{itemized}
list.
\begin{itemize}
   \item This is the first item of an itemized list.
         Each item in the list is marked with a ``tick''.
         You don't have to worry about what kind of tick
         mark is used.

   \item This is the second item of the list.  It
         contains another list nested inside it.  The inner
         list is an \emph{enumerated} list.
         \begin{enumerate}
            \item This is the first item of an enumerated
                  list that is nested within the itemized list.

            \item This is the second item of the inner list.
                  \LaTeX\ allows you to nest lists deeper than
                  you really should.
         \end{enumerate}
         This is the rest of the second item of the outer
         list.  It is no more interesting than any other
         part of the item.
   \item This is the third item of the list.
\end{itemize}
You can even display poetry.
\begin{verse}
   There is an environment
    for verse \\             % The \\ command separates lines
   Whose features some poets % within a stanza.
   will curse.

                             % One or more blank lines separate stanzas.

   For instead of making\\
   Them do \emph{all} line breaking, \\
   It allows them to put too many words on a line when they'd rather be
   forced to be terse.
\end{verse}

Mathematical formulas may also be displayed.  A
displayed formula
is
one-line long; multiline
formulas require special formatting instructions.
   \[  \ip{\Gamma}{\psi'} = x'' + y^{2} + z_{i}^{n}\]
Don't start a paragraph with a displayed equation,
nor make one a paragraph by itself.

\end{document}               % End of document.

\section{Bench Section 176}
% synthetic bench section 176
% This is a sample LaTeX input file.  (Version of 12 August 2004.)
%
% A '%' character causes TeX to ignore all remaining text on the line,
% and is used for comments like this one.

\documentclass{article}      % Specifies the document class

                             % The preamble begins here.
\title{An Example Document}  % Declares the document's title.
\author{Leslie Lamport}      % Declares the author's name.
\date{January 21, 1994}      % Deleting this command produces today's date.

\newcommand{\ip}[2]{(#1, #2)}
                             % Defines \ip{arg1}{arg2} to mean
                             % (arg1, arg2).

%\newcommand{\ip}[2]{\langle #1 | #2\rangle}
                             % This is an alternative definition of
                             % \ip that is commented out.

\begin{document}             % End of preamble and beginning of text.

\maketitle                   % Produces the title.

This is an example input file.  Comparing it with
the output it generates can show you how to
produce a simple document of your own.

\section{Ordinary Text}      % Produces section heading.  Lower-level
                             % sections are begun with similar
                             % \subsection and \subsubsection commands.

The ends  of words and sentences are marked
  by   spaces. It  doesn't matter how many
spaces    you type; one is as good as 100.  The
end of   a line counts as a space.

One   or more   blank lines denote the  end
of  a paragraph.

Since any number of consecutive spaces are treated
like a single one, the formatting of the input
file makes no difference to
      \LaTeX,                % The \LaTeX command generates the LaTeX logo.
but it makes a difference to you.  When you use
\LaTeX, making your input file as easy to read
as possible will be a great help as you write
your document and when you change it.  This sample
file shows how you can add comments to your own input
file.

Because printing is different from typewriting,
there are a number of things that you have to do
differently when preparing an input file than if
you were just typing the document directly.
Quotation marks like
       ``this''
have to be handled specially, as do quotes within
quotes:
       ``\,`this'            % \, separates the double and single quote.
        is what I just
        wrote, not  `that'\,''.

Dashes come in three sizes: an
       intra-word
dash, a medium dash for number ranges like
       1--2,
and a punctuation
       dash---like
this.

A sentence-ending space should be larger than the
space between words within a sentence.  You
sometimes have to type special commands in
conjunction with punctuation characters to get
this right, as in the following sentence.
       Gnats, gnus, etc.\ all  % `\ ' makes an inter-word space.
       begin with G\@.         % \@ marks end-of-sentence punctuation.
You should check the spaces after periods when
reading your output to make sure you haven't
forgotten any special cases.  Generating an
ellipsis
       \ldots\               % `\ ' is needed after `\ldots' because TeX
                             % ignores spaces after command names like \ldots
                             % made from \ + letters.
                             %
                             % Note how a `%' character causes TeX to ignore
                             % the end of the input line, so these blank lines
                             % do not start a new paragraph.
                             %
with the right spacing around the periods requires
a special command.

\LaTeX\ interprets some common characters as
commands, so you must type special commands to
generate them.  These characters include the
following:
       \$ \& \% \# \{ and \}.

In printing, text is usually emphasized with an
       \emph{italic}
type style.

\begin{em}
   A long segment of text can also be emphasized
   in this way.  Text within such a segment can be
   given \emph{additional} emphasis.
\end{em}

It is sometimes necessary to prevent \LaTeX\ from
breaking a line where it might otherwise do so.
This may be at a space, as between the ``Mr.''\ and
``Jones'' in
       ``Mr.~Jones'',        % ~ produces an unbreakable interword space.
or within a word---especially when the word is a
symbol like
       \mbox{\emph{itemnum}}
that makes little sense when hyphenated across
lines.

Footnotes\footnote{This is an example of a footnote.}
pose no problem.

\LaTeX\ is good at typesetting mathematical formulas
like
       \( x-3y + z = 7 \)
or
       \( a_{1} > x^{2n} + y^{2n} > x' \)
or
       \( \ip{A}{B} = \sum_{i} a_{i} b_{i} \).
The spaces you type in a formula are
ignored.  Remember that a letter like
       $x$                   % $ ... $  and  \( ... \)  are equivalent
is a formula when it denotes a mathematical
symbol, and it should be typed as one.

\section{Displayed Text}

Text is displayed by indenting it from the left
margin.  Quotations are commonly displayed.  There
are short quotations
\begin{quote}
   This is a short quotation.  It consists of a
   single paragraph of text.  See how it is formatted.
\end{quote}
and longer ones.
\begin{quotation}
   This is a longer quotation.  It consists of two
   paragraphs of text, neither of which are
   particularly interesting.

   This is the second paragraph of the quotation.  It
   is just as dull as the first paragraph.
\end{quotation}
Another frequently-displayed structure is a list.
The following is an example of an \emph{itemized}
list.
\begin{itemize}
   \item This is the first item of an itemized list.
         Each item in the list is marked with a ``tick''.
         You don't have to worry about what kind of tick
         mark is used.

   \item This is the second item of the list.  It
         contains another list nested inside it.  The inner
         list is an \emph{enumerated} list.
         \begin{enumerate}
            \item This is the first item of an enumerated
                  list that is nested within the itemized list.

            \item This is the second item of the inner list.
                  \LaTeX\ allows you to nest lists deeper than
                  you really should.
         \end{enumerate}
         This is the rest of the second item of the outer
         list.  It is no more interesting than any other
         part of the item.
   \item This is the third item of the list.
\end{itemize}
You can even display poetry.
\begin{verse}
   There is an environment
    for verse \\             % The \\ command separates lines
   Whose features some poets % within a stanza.
   will curse.

                             % One or more blank lines separate stanzas.

   For instead of making\\
   Them do \emph{all} line breaking, \\
   It allows them to put too many words on a line when they'd rather be
   forced to be terse.
\end{verse}

Mathematical formulas may also be displayed.  A
displayed formula
is
one-line long; multiline
formulas require special formatting instructions.
   \[  \ip{\Gamma}{\psi'} = x'' + y^{2} + z_{i}^{n}\]
Don't start a paragraph with a displayed equation,
nor make one a paragraph by itself.

\end{document}               % End of document.

\section{Bench Section 177}
% synthetic bench section 177
% This is a sample LaTeX input file.  (Version of 12 August 2004.)
%
% A '%' character causes TeX to ignore all remaining text on the line,
% and is used for comments like this one.

\documentclass{article}      % Specifies the document class

                             % The preamble begins here.
\title{An Example Document}  % Declares the document's title.
\author{Leslie Lamport}      % Declares the author's name.
\date{January 21, 1994}      % Deleting this command produces today's date.

\newcommand{\ip}[2]{(#1, #2)}
                             % Defines \ip{arg1}{arg2} to mean
                             % (arg1, arg2).

%\newcommand{\ip}[2]{\langle #1 | #2\rangle}
                             % This is an alternative definition of
                             % \ip that is commented out.

\begin{document}             % End of preamble and beginning of text.

\maketitle                   % Produces the title.

This is an example input file.  Comparing it with
the output it generates can show you how to
produce a simple document of your own.

\section{Ordinary Text}      % Produces section heading.  Lower-level
                             % sections are begun with similar
                             % \subsection and \subsubsection commands.

The ends  of words and sentences are marked
  by   spaces. It  doesn't matter how many
spaces    you type; one is as good as 100.  The
end of   a line counts as a space.

One   or more   blank lines denote the  end
of  a paragraph.

Since any number of consecutive spaces are treated
like a single one, the formatting of the input
file makes no difference to
      \LaTeX,                % The \LaTeX command generates the LaTeX logo.
but it makes a difference to you.  When you use
\LaTeX, making your input file as easy to read
as possible will be a great help as you write
your document and when you change it.  This sample
file shows how you can add comments to your own input
file.

Because printing is different from typewriting,
there are a number of things that you have to do
differently when preparing an input file than if
you were just typing the document directly.
Quotation marks like
       ``this''
have to be handled specially, as do quotes within
quotes:
       ``\,`this'            % \, separates the double and single quote.
        is what I just
        wrote, not  `that'\,''.

Dashes come in three sizes: an
       intra-word
dash, a medium dash for number ranges like
       1--2,
and a punctuation
       dash---like
this.

A sentence-ending space should be larger than the
space between words within a sentence.  You
sometimes have to type special commands in
conjunction with punctuation characters to get
this right, as in the following sentence.
       Gnats, gnus, etc.\ all  % `\ ' makes an inter-word space.
       begin with G\@.         % \@ marks end-of-sentence punctuation.
You should check the spaces after periods when
reading your output to make sure you haven't
forgotten any special cases.  Generating an
ellipsis
       \ldots\               % `\ ' is needed after `\ldots' because TeX
                             % ignores spaces after command names like \ldots
                             % made from \ + letters.
                             %
                             % Note how a `%' character causes TeX to ignore
                             % the end of the input line, so these blank lines
                             % do not start a new paragraph.
                             %
with the right spacing around the periods requires
a special command.

\LaTeX\ interprets some common characters as
commands, so you must type special commands to
generate them.  These characters include the
following:
       \$ \& \% \# \{ and \}.

In printing, text is usually emphasized with an
       \emph{italic}
type style.

\begin{em}
   A long segment of text can also be emphasized
   in this way.  Text within such a segment can be
   given \emph{additional} emphasis.
\end{em}

It is sometimes necessary to prevent \LaTeX\ from
breaking a line where it might otherwise do so.
This may be at a space, as between the ``Mr.''\ and
``Jones'' in
       ``Mr.~Jones'',        % ~ produces an unbreakable interword space.
or within a word---especially when the word is a
symbol like
       \mbox{\emph{itemnum}}
that makes little sense when hyphenated across
lines.

Footnotes\footnote{This is an example of a footnote.}
pose no problem.

\LaTeX\ is good at typesetting mathematical formulas
like
       \( x-3y + z = 7 \)
or
       \( a_{1} > x^{2n} + y^{2n} > x' \)
or
       \( \ip{A}{B} = \sum_{i} a_{i} b_{i} \).
The spaces you type in a formula are
ignored.  Remember that a letter like
       $x$                   % $ ... $  and  \( ... \)  are equivalent
is a formula when it denotes a mathematical
symbol, and it should be typed as one.

\section{Displayed Text}

Text is displayed by indenting it from the left
margin.  Quotations are commonly displayed.  There
are short quotations
\begin{quote}
   This is a short quotation.  It consists of a
   single paragraph of text.  See how it is formatted.
\end{quote}
and longer ones.
\begin{quotation}
   This is a longer quotation.  It consists of two
   paragraphs of text, neither of which are
   particularly interesting.

   This is the second paragraph of the quotation.  It
   is just as dull as the first paragraph.
\end{quotation}
Another frequently-displayed structure is a list.
The following is an example of an \emph{itemized}
list.
\begin{itemize}
   \item This is the first item of an itemized list.
         Each item in the list is marked with a ``tick''.
         You don't have to worry about what kind of tick
         mark is used.

   \item This is the second item of the list.  It
         contains another list nested inside it.  The inner
         list is an \emph{enumerated} list.
         \begin{enumerate}
            \item This is the first item of an enumerated
                  list that is nested within the itemized list.

            \item This is the second item of the inner list.
                  \LaTeX\ allows you to nest lists deeper than
                  you really should.
         \end{enumerate}
         This is the rest of the second item of the outer
         list.  It is no more interesting than any other
         part of the item.
   \item This is the third item of the list.
\end{itemize}
You can even display poetry.
\begin{verse}
   There is an environment
    for verse \\             % The \\ command separates lines
   Whose features some poets % within a stanza.
   will curse.

                             % One or more blank lines separate stanzas.

   For instead of making\\
   Them do \emph{all} line breaking, \\
   It allows them to put too many words on a line when they'd rather be
   forced to be terse.
\end{verse}

Mathematical formulas may also be displayed.  A
displayed formula
is
one-line long; multiline
formulas require special formatting instructions.
   \[  \ip{\Gamma}{\psi'} = x'' + y^{2} + z_{i}^{n}\]
Don't start a paragraph with a displayed equation,
nor make one a paragraph by itself.

\end{document}               % End of document.

\section{Bench Section 178}
% synthetic bench section 178
% This is a sample LaTeX input file.  (Version of 12 August 2004.)
%
% A '%' character causes TeX to ignore all remaining text on the line,
% and is used for comments like this one.

\documentclass{article}      % Specifies the document class

                             % The preamble begins here.
\title{An Example Document}  % Declares the document's title.
\author{Leslie Lamport}      % Declares the author's name.
\date{January 21, 1994}      % Deleting this command produces today's date.

\newcommand{\ip}[2]{(#1, #2)}
                             % Defines \ip{arg1}{arg2} to mean
                             % (arg1, arg2).

%\newcommand{\ip}[2]{\langle #1 | #2\rangle}
                             % This is an alternative definition of
                             % \ip that is commented out.

\begin{document}             % End of preamble and beginning of text.

\maketitle                   % Produces the title.

This is an example input file.  Comparing it with
the output it generates can show you how to
produce a simple document of your own.

\section{Ordinary Text}      % Produces section heading.  Lower-level
                             % sections are begun with similar
                             % \subsection and \subsubsection commands.

The ends  of words and sentences are marked
  by   spaces. It  doesn't matter how many
spaces    you type; one is as good as 100.  The
end of   a line counts as a space.

One   or more   blank lines denote the  end
of  a paragraph.

Since any number of consecutive spaces are treated
like a single one, the formatting of the input
file makes no difference to
      \LaTeX,                % The \LaTeX command generates the LaTeX logo.
but it makes a difference to you.  When you use
\LaTeX, making your input file as easy to read
as possible will be a great help as you write
your document and when you change it.  This sample
file shows how you can add comments to your own input
file.

Because printing is different from typewriting,
there are a number of things that you have to do
differently when preparing an input file than if
you were just typing the document directly.
Quotation marks like
       ``this''
have to be handled specially, as do quotes within
quotes:
       ``\,`this'            % \, separates the double and single quote.
        is what I just
        wrote, not  `that'\,''.

Dashes come in three sizes: an
       intra-word
dash, a medium dash for number ranges like
       1--2,
and a punctuation
       dash---like
this.

A sentence-ending space should be larger than the
space between words within a sentence.  You
sometimes have to type special commands in
conjunction with punctuation characters to get
this right, as in the following sentence.
       Gnats, gnus, etc.\ all  % `\ ' makes an inter-word space.
       begin with G\@.         % \@ marks end-of-sentence punctuation.
You should check the spaces after periods when
reading your output to make sure you haven't
forgotten any special cases.  Generating an
ellipsis
       \ldots\               % `\ ' is needed after `\ldots' because TeX
                             % ignores spaces after command names like \ldots
                             % made from \ + letters.
                             %
                             % Note how a `%' character causes TeX to ignore
                             % the end of the input line, so these blank lines
                             % do not start a new paragraph.
                             %
with the right spacing around the periods requires
a special command.

\LaTeX\ interprets some common characters as
commands, so you must type special commands to
generate them.  These characters include the
following:
       \$ \& \% \# \{ and \}.

In printing, text is usually emphasized with an
       \emph{italic}
type style.

\begin{em}
   A long segment of text can also be emphasized
   in this way.  Text within such a segment can be
   given \emph{additional} emphasis.
\end{em}

It is sometimes necessary to prevent \LaTeX\ from
breaking a line where it might otherwise do so.
This may be at a space, as between the ``Mr.''\ and
``Jones'' in
       ``Mr.~Jones'',        % ~ produces an unbreakable interword space.
or within a word---especially when the word is a
symbol like
       \mbox{\emph{itemnum}}
that makes little sense when hyphenated across
lines.

Footnotes\footnote{This is an example of a footnote.}
pose no problem.

\LaTeX\ is good at typesetting mathematical formulas
like
       \( x-3y + z = 7 \)
or
       \( a_{1} > x^{2n} + y^{2n} > x' \)
or
       \( \ip{A}{B} = \sum_{i} a_{i} b_{i} \).
The spaces you type in a formula are
ignored.  Remember that a letter like
       $x$                   % $ ... $  and  \( ... \)  are equivalent
is a formula when it denotes a mathematical
symbol, and it should be typed as one.

\section{Displayed Text}

Text is displayed by indenting it from the left
margin.  Quotations are commonly displayed.  There
are short quotations
\begin{quote}
   This is a short quotation.  It consists of a
   single paragraph of text.  See how it is formatted.
\end{quote}
and longer ones.
\begin{quotation}
   This is a longer quotation.  It consists of two
   paragraphs of text, neither of which are
   particularly interesting.

   This is the second paragraph of the quotation.  It
   is just as dull as the first paragraph.
\end{quotation}
Another frequently-displayed structure is a list.
The following is an example of an \emph{itemized}
list.
\begin{itemize}
   \item This is the first item of an itemized list.
         Each item in the list is marked with a ``tick''.
         You don't have to worry about what kind of tick
         mark is used.

   \item This is the second item of the list.  It
         contains another list nested inside it.  The inner
         list is an \emph{enumerated} list.
         \begin{enumerate}
            \item This is the first item of an enumerated
                  list that is nested within the itemized list.

            \item This is the second item of the inner list.
                  \LaTeX\ allows you to nest lists deeper than
                  you really should.
         \end{enumerate}
         This is the rest of the second item of the outer
         list.  It is no more interesting than any other
         part of the item.
   \item This is the third item of the list.
\end{itemize}
You can even display poetry.
\begin{verse}
   There is an environment
    for verse \\             % The \\ command separates lines
   Whose features some poets % within a stanza.
   will curse.

                             % One or more blank lines separate stanzas.

   For instead of making\\
   Them do \emph{all} line breaking, \\
   It allows them to put too many words on a line when they'd rather be
   forced to be terse.
\end{verse}

Mathematical formulas may also be displayed.  A
displayed formula
is
one-line long; multiline
formulas require special formatting instructions.
   \[  \ip{\Gamma}{\psi'} = x'' + y^{2} + z_{i}^{n}\]
Don't start a paragraph with a displayed equation,
nor make one a paragraph by itself.

\end{document}               % End of document.

\section{Bench Section 179}
% synthetic bench section 179
% This is a sample LaTeX input file.  (Version of 12 August 2004.)
%
% A '%' character causes TeX to ignore all remaining text on the line,
% and is used for comments like this one.

\documentclass{article}      % Specifies the document class

                             % The preamble begins here.
\title{An Example Document}  % Declares the document's title.
\author{Leslie Lamport}      % Declares the author's name.
\date{January 21, 1994}      % Deleting this command produces today's date.

\newcommand{\ip}[2]{(#1, #2)}
                             % Defines \ip{arg1}{arg2} to mean
                             % (arg1, arg2).

%\newcommand{\ip}[2]{\langle #1 | #2\rangle}
                             % This is an alternative definition of
                             % \ip that is commented out.

\begin{document}             % End of preamble and beginning of text.

\maketitle                   % Produces the title.

This is an example input file.  Comparing it with
the output it generates can show you how to
produce a simple document of your own.

\section{Ordinary Text}      % Produces section heading.  Lower-level
                             % sections are begun with similar
                             % \subsection and \subsubsection commands.

The ends  of words and sentences are marked
  by   spaces. It  doesn't matter how many
spaces    you type; one is as good as 100.  The
end of   a line counts as a space.

One   or more   blank lines denote the  end
of  a paragraph.

Since any number of consecutive spaces are treated
like a single one, the formatting of the input
file makes no difference to
      \LaTeX,                % The \LaTeX command generates the LaTeX logo.
but it makes a difference to you.  When you use
\LaTeX, making your input file as easy to read
as possible will be a great help as you write
your document and when you change it.  This sample
file shows how you can add comments to your own input
file.

Because printing is different from typewriting,
there are a number of things that you have to do
differently when preparing an input file than if
you were just typing the document directly.
Quotation marks like
       ``this''
have to be handled specially, as do quotes within
quotes:
       ``\,`this'            % \, separates the double and single quote.
        is what I just
        wrote, not  `that'\,''.

Dashes come in three sizes: an
       intra-word
dash, a medium dash for number ranges like
       1--2,
and a punctuation
       dash---like
this.

A sentence-ending space should be larger than the
space between words within a sentence.  You
sometimes have to type special commands in
conjunction with punctuation characters to get
this right, as in the following sentence.
       Gnats, gnus, etc.\ all  % `\ ' makes an inter-word space.
       begin with G\@.         % \@ marks end-of-sentence punctuation.
You should check the spaces after periods when
reading your output to make sure you haven't
forgotten any special cases.  Generating an
ellipsis
       \ldots\               % `\ ' is needed after `\ldots' because TeX
                             % ignores spaces after command names like \ldots
                             % made from \ + letters.
                             %
                             % Note how a `%' character causes TeX to ignore
                             % the end of the input line, so these blank lines
                             % do not start a new paragraph.
                             %
with the right spacing around the periods requires
a special command.

\LaTeX\ interprets some common characters as
commands, so you must type special commands to
generate them.  These characters include the
following:
       \$ \& \% \# \{ and \}.

In printing, text is usually emphasized with an
       \emph{italic}
type style.

\begin{em}
   A long segment of text can also be emphasized
   in this way.  Text within such a segment can be
   given \emph{additional} emphasis.
\end{em}

It is sometimes necessary to prevent \LaTeX\ from
breaking a line where it might otherwise do so.
This may be at a space, as between the ``Mr.''\ and
``Jones'' in
       ``Mr.~Jones'',        % ~ produces an unbreakable interword space.
or within a word---especially when the word is a
symbol like
       \mbox{\emph{itemnum}}
that makes little sense when hyphenated across
lines.

Footnotes\footnote{This is an example of a footnote.}
pose no problem.

\LaTeX\ is good at typesetting mathematical formulas
like
       \( x-3y + z = 7 \)
or
       \( a_{1} > x^{2n} + y^{2n} > x' \)
or
       \( \ip{A}{B} = \sum_{i} a_{i} b_{i} \).
The spaces you type in a formula are
ignored.  Remember that a letter like
       $x$                   % $ ... $  and  \( ... \)  are equivalent
is a formula when it denotes a mathematical
symbol, and it should be typed as one.

\section{Displayed Text}

Text is displayed by indenting it from the left
margin.  Quotations are commonly displayed.  There
are short quotations
\begin{quote}
   This is a short quotation.  It consists of a
   single paragraph of text.  See how it is formatted.
\end{quote}
and longer ones.
\begin{quotation}
   This is a longer quotation.  It consists of two
   paragraphs of text, neither of which are
   particularly interesting.

   This is the second paragraph of the quotation.  It
   is just as dull as the first paragraph.
\end{quotation}
Another frequently-displayed structure is a list.
The following is an example of an \emph{itemized}
list.
\begin{itemize}
   \item This is the first item of an itemized list.
         Each item in the list is marked with a ``tick''.
         You don't have to worry about what kind of tick
         mark is used.

   \item This is the second item of the list.  It
         contains another list nested inside it.  The inner
         list is an \emph{enumerated} list.
         \begin{enumerate}
            \item This is the first item of an enumerated
                  list that is nested within the itemized list.

            \item This is the second item of the inner list.
                  \LaTeX\ allows you to nest lists deeper than
                  you really should.
         \end{enumerate}
         This is the rest of the second item of the outer
         list.  It is no more interesting than any other
         part of the item.
   \item This is the third item of the list.
\end{itemize}
You can even display poetry.
\begin{verse}
   There is an environment
    for verse \\             % The \\ command separates lines
   Whose features some poets % within a stanza.
   will curse.

                             % One or more blank lines separate stanzas.

   For instead of making\\
   Them do \emph{all} line breaking, \\
   It allows them to put too many words on a line when they'd rather be
   forced to be terse.
\end{verse}

Mathematical formulas may also be displayed.  A
displayed formula
is
one-line long; multiline
formulas require special formatting instructions.
   \[  \ip{\Gamma}{\psi'} = x'' + y^{2} + z_{i}^{n}\]
Don't start a paragraph with a displayed equation,
nor make one a paragraph by itself.

\end{document}               % End of document.

\section{Bench Section 180}
% synthetic bench section 180
% This is a sample LaTeX input file.  (Version of 12 August 2004.)
%
% A '%' character causes TeX to ignore all remaining text on the line,
% and is used for comments like this one.

\documentclass{article}      % Specifies the document class

                             % The preamble begins here.
\title{An Example Document}  % Declares the document's title.
\author{Leslie Lamport}      % Declares the author's name.
\date{January 21, 1994}      % Deleting this command produces today's date.

\newcommand{\ip}[2]{(#1, #2)}
                             % Defines \ip{arg1}{arg2} to mean
                             % (arg1, arg2).

%\newcommand{\ip}[2]{\langle #1 | #2\rangle}
                             % This is an alternative definition of
                             % \ip that is commented out.

\begin{document}             % End of preamble and beginning of text.

\maketitle                   % Produces the title.

This is an example input file.  Comparing it with
the output it generates can show you how to
produce a simple document of your own.

\section{Ordinary Text}      % Produces section heading.  Lower-level
                             % sections are begun with similar
                             % \subsection and \subsubsection commands.

The ends  of words and sentences are marked
  by   spaces. It  doesn't matter how many
spaces    you type; one is as good as 100.  The
end of   a line counts as a space.

One   or more   blank lines denote the  end
of  a paragraph.

Since any number of consecutive spaces are treated
like a single one, the formatting of the input
file makes no difference to
      \LaTeX,                % The \LaTeX command generates the LaTeX logo.
but it makes a difference to you.  When you use
\LaTeX, making your input file as easy to read
as possible will be a great help as you write
your document and when you change it.  This sample
file shows how you can add comments to your own input
file.

Because printing is different from typewriting,
there are a number of things that you have to do
differently when preparing an input file than if
you were just typing the document directly.
Quotation marks like
       ``this''
have to be handled specially, as do quotes within
quotes:
       ``\,`this'            % \, separates the double and single quote.
        is what I just
        wrote, not  `that'\,''.

Dashes come in three sizes: an
       intra-word
dash, a medium dash for number ranges like
       1--2,
and a punctuation
       dash---like
this.

A sentence-ending space should be larger than the
space between words within a sentence.  You
sometimes have to type special commands in
conjunction with punctuation characters to get
this right, as in the following sentence.
       Gnats, gnus, etc.\ all  % `\ ' makes an inter-word space.
       begin with G\@.         % \@ marks end-of-sentence punctuation.
You should check the spaces after periods when
reading your output to make sure you haven't
forgotten any special cases.  Generating an
ellipsis
       \ldots\               % `\ ' is needed after `\ldots' because TeX
                             % ignores spaces after command names like \ldots
                             % made from \ + letters.
                             %
                             % Note how a `%' character causes TeX to ignore
                             % the end of the input line, so these blank lines
                             % do not start a new paragraph.
                             %
with the right spacing around the periods requires
a special command.

\LaTeX\ interprets some common characters as
commands, so you must type special commands to
generate them.  These characters include the
following:
       \$ \& \% \# \{ and \}.

In printing, text is usually emphasized with an
       \emph{italic}
type style.

\begin{em}
   A long segment of text can also be emphasized
   in this way.  Text within such a segment can be
   given \emph{additional} emphasis.
\end{em}

It is sometimes necessary to prevent \LaTeX\ from
breaking a line where it might otherwise do so.
This may be at a space, as between the ``Mr.''\ and
``Jones'' in
       ``Mr.~Jones'',        % ~ produces an unbreakable interword space.
or within a word---especially when the word is a
symbol like
       \mbox{\emph{itemnum}}
that makes little sense when hyphenated across
lines.

Footnotes\footnote{This is an example of a footnote.}
pose no problem.

\LaTeX\ is good at typesetting mathematical formulas
like
       \( x-3y + z = 7 \)
or
       \( a_{1} > x^{2n} + y^{2n} > x' \)
or
       \( \ip{A}{B} = \sum_{i} a_{i} b_{i} \).
The spaces you type in a formula are
ignored.  Remember that a letter like
       $x$                   % $ ... $  and  \( ... \)  are equivalent
is a formula when it denotes a mathematical
symbol, and it should be typed as one.

\section{Displayed Text}

Text is displayed by indenting it from the left
margin.  Quotations are commonly displayed.  There
are short quotations
\begin{quote}
   This is a short quotation.  It consists of a
   single paragraph of text.  See how it is formatted.
\end{quote}
and longer ones.
\begin{quotation}
   This is a longer quotation.  It consists of two
   paragraphs of text, neither of which are
   particularly interesting.

   This is the second paragraph of the quotation.  It
   is just as dull as the first paragraph.
\end{quotation}
Another frequently-displayed structure is a list.
The following is an example of an \emph{itemized}
list.
\begin{itemize}
   \item This is the first item of an itemized list.
         Each item in the list is marked with a ``tick''.
         You don't have to worry about what kind of tick
         mark is used.

   \item This is the second item of the list.  It
         contains another list nested inside it.  The inner
         list is an \emph{enumerated} list.
         \begin{enumerate}
            \item This is the first item of an enumerated
                  list that is nested within the itemized list.

            \item This is the second item of the inner list.
                  \LaTeX\ allows you to nest lists deeper than
                  you really should.
         \end{enumerate}
         This is the rest of the second item of the outer
         list.  It is no more interesting than any other
         part of the item.
   \item This is the third item of the list.
\end{itemize}
You can even display poetry.
\begin{verse}
   There is an environment
    for verse \\             % The \\ command separates lines
   Whose features some poets % within a stanza.
   will curse.

                             % One or more blank lines separate stanzas.

   For instead of making\\
   Them do \emph{all} line breaking, \\
   It allows them to put too many words on a line when they'd rather be
   forced to be terse.
\end{verse}

Mathematical formulas may also be displayed.  A
displayed formula
is
one-line long; multiline
formulas require special formatting instructions.
   \[  \ip{\Gamma}{\psi'} = x'' + y^{2} + z_{i}^{n}\]
Don't start a paragraph with a displayed equation,
nor make one a paragraph by itself.

\end{document}               % End of document.

\section{Bench Section 181}
% synthetic bench section 181
% This is a sample LaTeX input file.  (Version of 12 August 2004.)
%
% A '%' character causes TeX to ignore all remaining text on the line,
% and is used for comments like this one.

\documentclass{article}      % Specifies the document class

                             % The preamble begins here.
\title{An Example Document}  % Declares the document's title.
\author{Leslie Lamport}      % Declares the author's name.
\date{January 21, 1994}      % Deleting this command produces today's date.

\newcommand{\ip}[2]{(#1, #2)}
                             % Defines \ip{arg1}{arg2} to mean
                             % (arg1, arg2).

%\newcommand{\ip}[2]{\langle #1 | #2\rangle}
                             % This is an alternative definition of
                             % \ip that is commented out.

\begin{document}             % End of preamble and beginning of text.

\maketitle                   % Produces the title.

This is an example input file.  Comparing it with
the output it generates can show you how to
produce a simple document of your own.

\section{Ordinary Text}      % Produces section heading.  Lower-level
                             % sections are begun with similar
                             % \subsection and \subsubsection commands.

The ends  of words and sentences are marked
  by   spaces. It  doesn't matter how many
spaces    you type; one is as good as 100.  The
end of   a line counts as a space.

One   or more   blank lines denote the  end
of  a paragraph.

Since any number of consecutive spaces are treated
like a single one, the formatting of the input
file makes no difference to
      \LaTeX,                % The \LaTeX command generates the LaTeX logo.
but it makes a difference to you.  When you use
\LaTeX, making your input file as easy to read
as possible will be a great help as you write
your document and when you change it.  This sample
file shows how you can add comments to your own input
file.

Because printing is different from typewriting,
there are a number of things that you have to do
differently when preparing an input file than if
you were just typing the document directly.
Quotation marks like
       ``this''
have to be handled specially, as do quotes within
quotes:
       ``\,`this'            % \, separates the double and single quote.
        is what I just
        wrote, not  `that'\,''.

Dashes come in three sizes: an
       intra-word
dash, a medium dash for number ranges like
       1--2,
and a punctuation
       dash---like
this.

A sentence-ending space should be larger than the
space between words within a sentence.  You
sometimes have to type special commands in
conjunction with punctuation characters to get
this right, as in the following sentence.
       Gnats, gnus, etc.\ all  % `\ ' makes an inter-word space.
       begin with G\@.         % \@ marks end-of-sentence punctuation.
You should check the spaces after periods when
reading your output to make sure you haven't
forgotten any special cases.  Generating an
ellipsis
       \ldots\               % `\ ' is needed after `\ldots' because TeX
                             % ignores spaces after command names like \ldots
                             % made from \ + letters.
                             %
                             % Note how a `%' character causes TeX to ignore
                             % the end of the input line, so these blank lines
                             % do not start a new paragraph.
                             %
with the right spacing around the periods requires
a special command.

\LaTeX\ interprets some common characters as
commands, so you must type special commands to
generate them.  These characters include the
following:
       \$ \& \% \# \{ and \}.

In printing, text is usually emphasized with an
       \emph{italic}
type style.

\begin{em}
   A long segment of text can also be emphasized
   in this way.  Text within such a segment can be
   given \emph{additional} emphasis.
\end{em}

It is sometimes necessary to prevent \LaTeX\ from
breaking a line where it might otherwise do so.
This may be at a space, as between the ``Mr.''\ and
``Jones'' in
       ``Mr.~Jones'',        % ~ produces an unbreakable interword space.
or within a word---especially when the word is a
symbol like
       \mbox{\emph{itemnum}}
that makes little sense when hyphenated across
lines.

Footnotes\footnote{This is an example of a footnote.}
pose no problem.

\LaTeX\ is good at typesetting mathematical formulas
like
       \( x-3y + z = 7 \)
or
       \( a_{1} > x^{2n} + y^{2n} > x' \)
or
       \( \ip{A}{B} = \sum_{i} a_{i} b_{i} \).
The spaces you type in a formula are
ignored.  Remember that a letter like
       $x$                   % $ ... $  and  \( ... \)  are equivalent
is a formula when it denotes a mathematical
symbol, and it should be typed as one.

\section{Displayed Text}

Text is displayed by indenting it from the left
margin.  Quotations are commonly displayed.  There
are short quotations
\begin{quote}
   This is a short quotation.  It consists of a
   single paragraph of text.  See how it is formatted.
\end{quote}
and longer ones.
\begin{quotation}
   This is a longer quotation.  It consists of two
   paragraphs of text, neither of which are
   particularly interesting.

   This is the second paragraph of the quotation.  It
   is just as dull as the first paragraph.
\end{quotation}
Another frequently-displayed structure is a list.
The following is an example of an \emph{itemized}
list.
\begin{itemize}
   \item This is the first item of an itemized list.
         Each item in the list is marked with a ``tick''.
         You don't have to worry about what kind of tick
         mark is used.

   \item This is the second item of the list.  It
         contains another list nested inside it.  The inner
         list is an \emph{enumerated} list.
         \begin{enumerate}
            \item This is the first item of an enumerated
                  list that is nested within the itemized list.

            \item This is the second item of the inner list.
                  \LaTeX\ allows you to nest lists deeper than
                  you really should.
         \end{enumerate}
         This is the rest of the second item of the outer
         list.  It is no more interesting than any other
         part of the item.
   \item This is the third item of the list.
\end{itemize}
You can even display poetry.
\begin{verse}
   There is an environment
    for verse \\             % The \\ command separates lines
   Whose features some poets % within a stanza.
   will curse.

                             % One or more blank lines separate stanzas.

   For instead of making\\
   Them do \emph{all} line breaking, \\
   It allows them to put too many words on a line when they'd rather be
   forced to be terse.
\end{verse}

Mathematical formulas may also be displayed.  A
displayed formula
is
one-line long; multiline
formulas require special formatting instructions.
   \[  \ip{\Gamma}{\psi'} = x'' + y^{2} + z_{i}^{n}\]
Don't start a paragraph with a displayed equation,
nor make one a paragraph by itself.

\end{document}               % End of document.

\section{Bench Section 182}
% synthetic bench section 182
% This is a sample LaTeX input file.  (Version of 12 August 2004.)
%
% A '%' character causes TeX to ignore all remaining text on the line,
% and is used for comments like this one.

\documentclass{article}      % Specifies the document class

                             % The preamble begins here.
\title{An Example Document}  % Declares the document's title.
\author{Leslie Lamport}      % Declares the author's name.
\date{January 21, 1994}      % Deleting this command produces today's date.

\newcommand{\ip}[2]{(#1, #2)}
                             % Defines \ip{arg1}{arg2} to mean
                             % (arg1, arg2).

%\newcommand{\ip}[2]{\langle #1 | #2\rangle}
                             % This is an alternative definition of
                             % \ip that is commented out.

\begin{document}             % End of preamble and beginning of text.

\maketitle                   % Produces the title.

This is an example input file.  Comparing it with
the output it generates can show you how to
produce a simple document of your own.

\section{Ordinary Text}      % Produces section heading.  Lower-level
                             % sections are begun with similar
                             % \subsection and \subsubsection commands.

The ends  of words and sentences are marked
  by   spaces. It  doesn't matter how many
spaces    you type; one is as good as 100.  The
end of   a line counts as a space.

One   or more   blank lines denote the  end
of  a paragraph.

Since any number of consecutive spaces are treated
like a single one, the formatting of the input
file makes no difference to
      \LaTeX,                % The \LaTeX command generates the LaTeX logo.
but it makes a difference to you.  When you use
\LaTeX, making your input file as easy to read
as possible will be a great help as you write
your document and when you change it.  This sample
file shows how you can add comments to your own input
file.

Because printing is different from typewriting,
there are a number of things that you have to do
differently when preparing an input file than if
you were just typing the document directly.
Quotation marks like
       ``this''
have to be handled specially, as do quotes within
quotes:
       ``\,`this'            % \, separates the double and single quote.
        is what I just
        wrote, not  `that'\,''.

Dashes come in three sizes: an
       intra-word
dash, a medium dash for number ranges like
       1--2,
and a punctuation
       dash---like
this.

A sentence-ending space should be larger than the
space between words within a sentence.  You
sometimes have to type special commands in
conjunction with punctuation characters to get
this right, as in the following sentence.
       Gnats, gnus, etc.\ all  % `\ ' makes an inter-word space.
       begin with G\@.         % \@ marks end-of-sentence punctuation.
You should check the spaces after periods when
reading your output to make sure you haven't
forgotten any special cases.  Generating an
ellipsis
       \ldots\               % `\ ' is needed after `\ldots' because TeX
                             % ignores spaces after command names like \ldots
                             % made from \ + letters.
                             %
                             % Note how a `%' character causes TeX to ignore
                             % the end of the input line, so these blank lines
                             % do not start a new paragraph.
                             %
with the right spacing around the periods requires
a special command.

\LaTeX\ interprets some common characters as
commands, so you must type special commands to
generate them.  These characters include the
following:
       \$ \& \% \# \{ and \}.

In printing, text is usually emphasized with an
       \emph{italic}
type style.

\begin{em}
   A long segment of text can also be emphasized
   in this way.  Text within such a segment can be
   given \emph{additional} emphasis.
\end{em}

It is sometimes necessary to prevent \LaTeX\ from
breaking a line where it might otherwise do so.
This may be at a space, as between the ``Mr.''\ and
``Jones'' in
       ``Mr.~Jones'',        % ~ produces an unbreakable interword space.
or within a word---especially when the word is a
symbol like
       \mbox{\emph{itemnum}}
that makes little sense when hyphenated across
lines.

Footnotes\footnote{This is an example of a footnote.}
pose no problem.

\LaTeX\ is good at typesetting mathematical formulas
like
       \( x-3y + z = 7 \)
or
       \( a_{1} > x^{2n} + y^{2n} > x' \)
or
       \( \ip{A}{B} = \sum_{i} a_{i} b_{i} \).
The spaces you type in a formula are
ignored.  Remember that a letter like
       $x$                   % $ ... $  and  \( ... \)  are equivalent
is a formula when it denotes a mathematical
symbol, and it should be typed as one.

\section{Displayed Text}

Text is displayed by indenting it from the left
margin.  Quotations are commonly displayed.  There
are short quotations
\begin{quote}
   This is a short quotation.  It consists of a
   single paragraph of text.  See how it is formatted.
\end{quote}
and longer ones.
\begin{quotation}
   This is a longer quotation.  It consists of two
   paragraphs of text, neither of which are
   particularly interesting.

   This is the second paragraph of the quotation.  It
   is just as dull as the first paragraph.
\end{quotation}
Another frequently-displayed structure is a list.
The following is an example of an \emph{itemized}
list.
\begin{itemize}
   \item This is the first item of an itemized list.
         Each item in the list is marked with a ``tick''.
         You don't have to worry about what kind of tick
         mark is used.

   \item This is the second item of the list.  It
         contains another list nested inside it.  The inner
         list is an \emph{enumerated} list.
         \begin{enumerate}
            \item This is the first item of an enumerated
                  list that is nested within the itemized list.

            \item This is the second item of the inner list.
                  \LaTeX\ allows you to nest lists deeper than
                  you really should.
         \end{enumerate}
         This is the rest of the second item of the outer
         list.  It is no more interesting than any other
         part of the item.
   \item This is the third item of the list.
\end{itemize}
You can even display poetry.
\begin{verse}
   There is an environment
    for verse \\             % The \\ command separates lines
   Whose features some poets % within a stanza.
   will curse.

                             % One or more blank lines separate stanzas.

   For instead of making\\
   Them do \emph{all} line breaking, \\
   It allows them to put too many words on a line when they'd rather be
   forced to be terse.
\end{verse}

Mathematical formulas may also be displayed.  A
displayed formula
is
one-line long; multiline
formulas require special formatting instructions.
   \[  \ip{\Gamma}{\psi'} = x'' + y^{2} + z_{i}^{n}\]
Don't start a paragraph with a displayed equation,
nor make one a paragraph by itself.

\end{document}               % End of document.

\section{Bench Section 183}
% synthetic bench section 183
% This is a sample LaTeX input file.  (Version of 12 August 2004.)
%
% A '%' character causes TeX to ignore all remaining text on the line,
% and is used for comments like this one.

\documentclass{article}      % Specifies the document class

                             % The preamble begins here.
\title{An Example Document}  % Declares the document's title.
\author{Leslie Lamport}      % Declares the author's name.
\date{January 21, 1994}      % Deleting this command produces today's date.

\newcommand{\ip}[2]{(#1, #2)}
                             % Defines \ip{arg1}{arg2} to mean
                             % (arg1, arg2).

%\newcommand{\ip}[2]{\langle #1 | #2\rangle}
                             % This is an alternative definition of
                             % \ip that is commented out.

\begin{document}             % End of preamble and beginning of text.

\maketitle                   % Produces the title.

This is an example input file.  Comparing it with
the output it generates can show you how to
produce a simple document of your own.

\section{Ordinary Text}      % Produces section heading.  Lower-level
                             % sections are begun with similar
                             % \subsection and \subsubsection commands.

The ends  of words and sentences are marked
  by   spaces. It  doesn't matter how many
spaces    you type; one is as good as 100.  The
end of   a line counts as a space.

One   or more   blank lines denote the  end
of  a paragraph.

Since any number of consecutive spaces are treated
like a single one, the formatting of the input
file makes no difference to
      \LaTeX,                % The \LaTeX command generates the LaTeX logo.
but it makes a difference to you.  When you use
\LaTeX, making your input file as easy to read
as possible will be a great help as you write
your document and when you change it.  This sample
file shows how you can add comments to your own input
file.

Because printing is different from typewriting,
there are a number of things that you have to do
differently when preparing an input file than if
you were just typing the document directly.
Quotation marks like
       ``this''
have to be handled specially, as do quotes within
quotes:
       ``\,`this'            % \, separates the double and single quote.
        is what I just
        wrote, not  `that'\,''.

Dashes come in three sizes: an
       intra-word
dash, a medium dash for number ranges like
       1--2,
and a punctuation
       dash---like
this.

A sentence-ending space should be larger than the
space between words within a sentence.  You
sometimes have to type special commands in
conjunction with punctuation characters to get
this right, as in the following sentence.
       Gnats, gnus, etc.\ all  % `\ ' makes an inter-word space.
       begin with G\@.         % \@ marks end-of-sentence punctuation.
You should check the spaces after periods when
reading your output to make sure you haven't
forgotten any special cases.  Generating an
ellipsis
       \ldots\               % `\ ' is needed after `\ldots' because TeX
                             % ignores spaces after command names like \ldots
                             % made from \ + letters.
                             %
                             % Note how a `%' character causes TeX to ignore
                             % the end of the input line, so these blank lines
                             % do not start a new paragraph.
                             %
with the right spacing around the periods requires
a special command.

\LaTeX\ interprets some common characters as
commands, so you must type special commands to
generate them.  These characters include the
following:
       \$ \& \% \# \{ and \}.

In printing, text is usually emphasized with an
       \emph{italic}
type style.

\begin{em}
   A long segment of text can also be emphasized
   in this way.  Text within such a segment can be
   given \emph{additional} emphasis.
\end{em}

It is sometimes necessary to prevent \LaTeX\ from
breaking a line where it might otherwise do so.
This may be at a space, as between the ``Mr.''\ and
``Jones'' in
       ``Mr.~Jones'',        % ~ produces an unbreakable interword space.
or within a word---especially when the word is a
symbol like
       \mbox{\emph{itemnum}}
that makes little sense when hyphenated across
lines.

Footnotes\footnote{This is an example of a footnote.}
pose no problem.

\LaTeX\ is good at typesetting mathematical formulas
like
       \( x-3y + z = 7 \)
or
       \( a_{1} > x^{2n} + y^{2n} > x' \)
or
       \( \ip{A}{B} = \sum_{i} a_{i} b_{i} \).
The spaces you type in a formula are
ignored.  Remember that a letter like
       $x$                   % $ ... $  and  \( ... \)  are equivalent
is a formula when it denotes a mathematical
symbol, and it should be typed as one.

\section{Displayed Text}

Text is displayed by indenting it from the left
margin.  Quotations are commonly displayed.  There
are short quotations
\begin{quote}
   This is a short quotation.  It consists of a
   single paragraph of text.  See how it is formatted.
\end{quote}
and longer ones.
\begin{quotation}
   This is a longer quotation.  It consists of two
   paragraphs of text, neither of which are
   particularly interesting.

   This is the second paragraph of the quotation.  It
   is just as dull as the first paragraph.
\end{quotation}
Another frequently-displayed structure is a list.
The following is an example of an \emph{itemized}
list.
\begin{itemize}
   \item This is the first item of an itemized list.
         Each item in the list is marked with a ``tick''.
         You don't have to worry about what kind of tick
         mark is used.

   \item This is the second item of the list.  It
         contains another list nested inside it.  The inner
         list is an \emph{enumerated} list.
         \begin{enumerate}
            \item This is the first item of an enumerated
                  list that is nested within the itemized list.

            \item This is the second item of the inner list.
                  \LaTeX\ allows you to nest lists deeper than
                  you really should.
         \end{enumerate}
         This is the rest of the second item of the outer
         list.  It is no more interesting than any other
         part of the item.
   \item This is the third item of the list.
\end{itemize}
You can even display poetry.
\begin{verse}
   There is an environment
    for verse \\             % The \\ command separates lines
   Whose features some poets % within a stanza.
   will curse.

                             % One or more blank lines separate stanzas.

   For instead of making\\
   Them do \emph{all} line breaking, \\
   It allows them to put too many words on a line when they'd rather be
   forced to be terse.
\end{verse}

Mathematical formulas may also be displayed.  A
displayed formula
is
one-line long; multiline
formulas require special formatting instructions.
   \[  \ip{\Gamma}{\psi'} = x'' + y^{2} + z_{i}^{n}\]
Don't start a paragraph with a displayed equation,
nor make one a paragraph by itself.

\end{document}               % End of document.

\section{Bench Section 184}
% synthetic bench section 184
% This is a sample LaTeX input file.  (Version of 12 August 2004.)
%
% A '%' character causes TeX to ignore all remaining text on the line,
% and is used for comments like this one.

\documentclass{article}      % Specifies the document class

                             % The preamble begins here.
\title{An Example Document}  % Declares the document's title.
\author{Leslie Lamport}      % Declares the author's name.
\date{January 21, 1994}      % Deleting this command produces today's date.

\newcommand{\ip}[2]{(#1, #2)}
                             % Defines \ip{arg1}{arg2} to mean
                             % (arg1, arg2).

%\newcommand{\ip}[2]{\langle #1 | #2\rangle}
                             % This is an alternative definition of
                             % \ip that is commented out.

\begin{document}             % End of preamble and beginning of text.

\maketitle                   % Produces the title.

This is an example input file.  Comparing it with
the output it generates can show you how to
produce a simple document of your own.

\section{Ordinary Text}      % Produces section heading.  Lower-level
                             % sections are begun with similar
                             % \subsection and \subsubsection commands.

The ends  of words and sentences are marked
  by   spaces. It  doesn't matter how many
spaces    you type; one is as good as 100.  The
end of   a line counts as a space.

One   or more   blank lines denote the  end
of  a paragraph.

Since any number of consecutive spaces are treated
like a single one, the formatting of the input
file makes no difference to
      \LaTeX,                % The \LaTeX command generates the LaTeX logo.
but it makes a difference to you.  When you use
\LaTeX, making your input file as easy to read
as possible will be a great help as you write
your document and when you change it.  This sample
file shows how you can add comments to your own input
file.

Because printing is different from typewriting,
there are a number of things that you have to do
differently when preparing an input file than if
you were just typing the document directly.
Quotation marks like
       ``this''
have to be handled specially, as do quotes within
quotes:
       ``\,`this'            % \, separates the double and single quote.
        is what I just
        wrote, not  `that'\,''.

Dashes come in three sizes: an
       intra-word
dash, a medium dash for number ranges like
       1--2,
and a punctuation
       dash---like
this.

A sentence-ending space should be larger than the
space between words within a sentence.  You
sometimes have to type special commands in
conjunction with punctuation characters to get
this right, as in the following sentence.
       Gnats, gnus, etc.\ all  % `\ ' makes an inter-word space.
       begin with G\@.         % \@ marks end-of-sentence punctuation.
You should check the spaces after periods when
reading your output to make sure you haven't
forgotten any special cases.  Generating an
ellipsis
       \ldots\               % `\ ' is needed after `\ldots' because TeX
                             % ignores spaces after command names like \ldots
                             % made from \ + letters.
                             %
                             % Note how a `%' character causes TeX to ignore
                             % the end of the input line, so these blank lines
                             % do not start a new paragraph.
                             %
with the right spacing around the periods requires
a special command.

\LaTeX\ interprets some common characters as
commands, so you must type special commands to
generate them.  These characters include the
following:
       \$ \& \% \# \{ and \}.

In printing, text is usually emphasized with an
       \emph{italic}
type style.

\begin{em}
   A long segment of text can also be emphasized
   in this way.  Text within such a segment can be
   given \emph{additional} emphasis.
\end{em}

It is sometimes necessary to prevent \LaTeX\ from
breaking a line where it might otherwise do so.
This may be at a space, as between the ``Mr.''\ and
``Jones'' in
       ``Mr.~Jones'',        % ~ produces an unbreakable interword space.
or within a word---especially when the word is a
symbol like
       \mbox{\emph{itemnum}}
that makes little sense when hyphenated across
lines.

Footnotes\footnote{This is an example of a footnote.}
pose no problem.

\LaTeX\ is good at typesetting mathematical formulas
like
       \( x-3y + z = 7 \)
or
       \( a_{1} > x^{2n} + y^{2n} > x' \)
or
       \( \ip{A}{B} = \sum_{i} a_{i} b_{i} \).
The spaces you type in a formula are
ignored.  Remember that a letter like
       $x$                   % $ ... $  and  \( ... \)  are equivalent
is a formula when it denotes a mathematical
symbol, and it should be typed as one.

\section{Displayed Text}

Text is displayed by indenting it from the left
margin.  Quotations are commonly displayed.  There
are short quotations
\begin{quote}
   This is a short quotation.  It consists of a
   single paragraph of text.  See how it is formatted.
\end{quote}
and longer ones.
\begin{quotation}
   This is a longer quotation.  It consists of two
   paragraphs of text, neither of which are
   particularly interesting.

   This is the second paragraph of the quotation.  It
   is just as dull as the first paragraph.
\end{quotation}
Another frequently-displayed structure is a list.
The following is an example of an \emph{itemized}
list.
\begin{itemize}
   \item This is the first item of an itemized list.
         Each item in the list is marked with a ``tick''.
         You don't have to worry about what kind of tick
         mark is used.

   \item This is the second item of the list.  It
         contains another list nested inside it.  The inner
         list is an \emph{enumerated} list.
         \begin{enumerate}
            \item This is the first item of an enumerated
                  list that is nested within the itemized list.

            \item This is the second item of the inner list.
                  \LaTeX\ allows you to nest lists deeper than
                  you really should.
         \end{enumerate}
         This is the rest of the second item of the outer
         list.  It is no more interesting than any other
         part of the item.
   \item This is the third item of the list.
\end{itemize}
You can even display poetry.
\begin{verse}
   There is an environment
    for verse \\             % The \\ command separates lines
   Whose features some poets % within a stanza.
   will curse.

                             % One or more blank lines separate stanzas.

   For instead of making\\
   Them do \emph{all} line breaking, \\
   It allows them to put too many words on a line when they'd rather be
   forced to be terse.
\end{verse}

Mathematical formulas may also be displayed.  A
displayed formula
is
one-line long; multiline
formulas require special formatting instructions.
   \[  \ip{\Gamma}{\psi'} = x'' + y^{2} + z_{i}^{n}\]
Don't start a paragraph with a displayed equation,
nor make one a paragraph by itself.

\end{document}               % End of document.

\section{Bench Section 185}
% synthetic bench section 185
% This is a sample LaTeX input file.  (Version of 12 August 2004.)
%
% A '%' character causes TeX to ignore all remaining text on the line,
% and is used for comments like this one.

\documentclass{article}      % Specifies the document class

                             % The preamble begins here.
\title{An Example Document}  % Declares the document's title.
\author{Leslie Lamport}      % Declares the author's name.
\date{January 21, 1994}      % Deleting this command produces today's date.

\newcommand{\ip}[2]{(#1, #2)}
                             % Defines \ip{arg1}{arg2} to mean
                             % (arg1, arg2).

%\newcommand{\ip}[2]{\langle #1 | #2\rangle}
                             % This is an alternative definition of
                             % \ip that is commented out.

\begin{document}             % End of preamble and beginning of text.

\maketitle                   % Produces the title.

This is an example input file.  Comparing it with
the output it generates can show you how to
produce a simple document of your own.

\section{Ordinary Text}      % Produces section heading.  Lower-level
                             % sections are begun with similar
                             % \subsection and \subsubsection commands.

The ends  of words and sentences are marked
  by   spaces. It  doesn't matter how many
spaces    you type; one is as good as 100.  The
end of   a line counts as a space.

One   or more   blank lines denote the  end
of  a paragraph.

Since any number of consecutive spaces are treated
like a single one, the formatting of the input
file makes no difference to
      \LaTeX,                % The \LaTeX command generates the LaTeX logo.
but it makes a difference to you.  When you use
\LaTeX, making your input file as easy to read
as possible will be a great help as you write
your document and when you change it.  This sample
file shows how you can add comments to your own input
file.

Because printing is different from typewriting,
there are a number of things that you have to do
differently when preparing an input file than if
you were just typing the document directly.
Quotation marks like
       ``this''
have to be handled specially, as do quotes within
quotes:
       ``\,`this'            % \, separates the double and single quote.
        is what I just
        wrote, not  `that'\,''.

Dashes come in three sizes: an
       intra-word
dash, a medium dash for number ranges like
       1--2,
and a punctuation
       dash---like
this.

A sentence-ending space should be larger than the
space between words within a sentence.  You
sometimes have to type special commands in
conjunction with punctuation characters to get
this right, as in the following sentence.
       Gnats, gnus, etc.\ all  % `\ ' makes an inter-word space.
       begin with G\@.         % \@ marks end-of-sentence punctuation.
You should check the spaces after periods when
reading your output to make sure you haven't
forgotten any special cases.  Generating an
ellipsis
       \ldots\               % `\ ' is needed after `\ldots' because TeX
                             % ignores spaces after command names like \ldots
                             % made from \ + letters.
                             %
                             % Note how a `%' character causes TeX to ignore
                             % the end of the input line, so these blank lines
                             % do not start a new paragraph.
                             %
with the right spacing around the periods requires
a special command.

\LaTeX\ interprets some common characters as
commands, so you must type special commands to
generate them.  These characters include the
following:
       \$ \& \% \# \{ and \}.

In printing, text is usually emphasized with an
       \emph{italic}
type style.

\begin{em}
   A long segment of text can also be emphasized
   in this way.  Text within such a segment can be
   given \emph{additional} emphasis.
\end{em}

It is sometimes necessary to prevent \LaTeX\ from
breaking a line where it might otherwise do so.
This may be at a space, as between the ``Mr.''\ and
``Jones'' in
       ``Mr.~Jones'',        % ~ produces an unbreakable interword space.
or within a word---especially when the word is a
symbol like
       \mbox{\emph{itemnum}}
that makes little sense when hyphenated across
lines.

Footnotes\footnote{This is an example of a footnote.}
pose no problem.

\LaTeX\ is good at typesetting mathematical formulas
like
       \( x-3y + z = 7 \)
or
       \( a_{1} > x^{2n} + y^{2n} > x' \)
or
       \( \ip{A}{B} = \sum_{i} a_{i} b_{i} \).
The spaces you type in a formula are
ignored.  Remember that a letter like
       $x$                   % $ ... $  and  \( ... \)  are equivalent
is a formula when it denotes a mathematical
symbol, and it should be typed as one.

\section{Displayed Text}

Text is displayed by indenting it from the left
margin.  Quotations are commonly displayed.  There
are short quotations
\begin{quote}
   This is a short quotation.  It consists of a
   single paragraph of text.  See how it is formatted.
\end{quote}
and longer ones.
\begin{quotation}
   This is a longer quotation.  It consists of two
   paragraphs of text, neither of which are
   particularly interesting.

   This is the second paragraph of the quotation.  It
   is just as dull as the first paragraph.
\end{quotation}
Another frequently-displayed structure is a list.
The following is an example of an \emph{itemized}
list.
\begin{itemize}
   \item This is the first item of an itemized list.
         Each item in the list is marked with a ``tick''.
         You don't have to worry about what kind of tick
         mark is used.

   \item This is the second item of the list.  It
         contains another list nested inside it.  The inner
         list is an \emph{enumerated} list.
         \begin{enumerate}
            \item This is the first item of an enumerated
                  list that is nested within the itemized list.

            \item This is the second item of the inner list.
                  \LaTeX\ allows you to nest lists deeper than
                  you really should.
         \end{enumerate}
         This is the rest of the second item of the outer
         list.  It is no more interesting than any other
         part of the item.
   \item This is the third item of the list.
\end{itemize}
You can even display poetry.
\begin{verse}
   There is an environment
    for verse \\             % The \\ command separates lines
   Whose features some poets % within a stanza.
   will curse.

                             % One or more blank lines separate stanzas.

   For instead of making\\
   Them do \emph{all} line breaking, \\
   It allows them to put too many words on a line when they'd rather be
   forced to be terse.
\end{verse}

Mathematical formulas may also be displayed.  A
displayed formula
is
one-line long; multiline
formulas require special formatting instructions.
   \[  \ip{\Gamma}{\psi'} = x'' + y^{2} + z_{i}^{n}\]
Don't start a paragraph with a displayed equation,
nor make one a paragraph by itself.

\end{document}               % End of document.

\section{Bench Section 186}
% synthetic bench section 186
% This is a sample LaTeX input file.  (Version of 12 August 2004.)
%
% A '%' character causes TeX to ignore all remaining text on the line,
% and is used for comments like this one.

\documentclass{article}      % Specifies the document class

                             % The preamble begins here.
\title{An Example Document}  % Declares the document's title.
\author{Leslie Lamport}      % Declares the author's name.
\date{January 21, 1994}      % Deleting this command produces today's date.

\newcommand{\ip}[2]{(#1, #2)}
                             % Defines \ip{arg1}{arg2} to mean
                             % (arg1, arg2).

%\newcommand{\ip}[2]{\langle #1 | #2\rangle}
                             % This is an alternative definition of
                             % \ip that is commented out.

\begin{document}             % End of preamble and beginning of text.

\maketitle                   % Produces the title.

This is an example input file.  Comparing it with
the output it generates can show you how to
produce a simple document of your own.

\section{Ordinary Text}      % Produces section heading.  Lower-level
                             % sections are begun with similar
                             % \subsection and \subsubsection commands.

The ends  of words and sentences are marked
  by   spaces. It  doesn't matter how many
spaces    you type; one is as good as 100.  The
end of   a line counts as a space.

One   or more   blank lines denote the  end
of  a paragraph.

Since any number of consecutive spaces are treated
like a single one, the formatting of the input
file makes no difference to
      \LaTeX,                % The \LaTeX command generates the LaTeX logo.
but it makes a difference to you.  When you use
\LaTeX, making your input file as easy to read
as possible will be a great help as you write
your document and when you change it.  This sample
file shows how you can add comments to your own input
file.

Because printing is different from typewriting,
there are a number of things that you have to do
differently when preparing an input file than if
you were just typing the document directly.
Quotation marks like
       ``this''
have to be handled specially, as do quotes within
quotes:
       ``\,`this'            % \, separates the double and single quote.
        is what I just
        wrote, not  `that'\,''.

Dashes come in three sizes: an
       intra-word
dash, a medium dash for number ranges like
       1--2,
and a punctuation
       dash---like
this.

A sentence-ending space should be larger than the
space between words within a sentence.  You
sometimes have to type special commands in
conjunction with punctuation characters to get
this right, as in the following sentence.
       Gnats, gnus, etc.\ all  % `\ ' makes an inter-word space.
       begin with G\@.         % \@ marks end-of-sentence punctuation.
You should check the spaces after periods when
reading your output to make sure you haven't
forgotten any special cases.  Generating an
ellipsis
       \ldots\               % `\ ' is needed after `\ldots' because TeX
                             % ignores spaces after command names like \ldots
                             % made from \ + letters.
                             %
                             % Note how a `%' character causes TeX to ignore
                             % the end of the input line, so these blank lines
                             % do not start a new paragraph.
                             %
with the right spacing around the periods requires
a special command.

\LaTeX\ interprets some common characters as
commands, so you must type special commands to
generate them.  These characters include the
following:
       \$ \& \% \# \{ and \}.

In printing, text is usually emphasized with an
       \emph{italic}
type style.

\begin{em}
   A long segment of text can also be emphasized
   in this way.  Text within such a segment can be
   given \emph{additional} emphasis.
\end{em}

It is sometimes necessary to prevent \LaTeX\ from
breaking a line where it might otherwise do so.
This may be at a space, as between the ``Mr.''\ and
``Jones'' in
       ``Mr.~Jones'',        % ~ produces an unbreakable interword space.
or within a word---especially when the word is a
symbol like
       \mbox{\emph{itemnum}}
that makes little sense when hyphenated across
lines.

Footnotes\footnote{This is an example of a footnote.}
pose no problem.

\LaTeX\ is good at typesetting mathematical formulas
like
       \( x-3y + z = 7 \)
or
       \( a_{1} > x^{2n} + y^{2n} > x' \)
or
       \( \ip{A}{B} = \sum_{i} a_{i} b_{i} \).
The spaces you type in a formula are
ignored.  Remember that a letter like
       $x$                   % $ ... $  and  \( ... \)  are equivalent
is a formula when it denotes a mathematical
symbol, and it should be typed as one.

\section{Displayed Text}

Text is displayed by indenting it from the left
margin.  Quotations are commonly displayed.  There
are short quotations
\begin{quote}
   This is a short quotation.  It consists of a
   single paragraph of text.  See how it is formatted.
\end{quote}
and longer ones.
\begin{quotation}
   This is a longer quotation.  It consists of two
   paragraphs of text, neither of which are
   particularly interesting.

   This is the second paragraph of the quotation.  It
   is just as dull as the first paragraph.
\end{quotation}
Another frequently-displayed structure is a list.
The following is an example of an \emph{itemized}
list.
\begin{itemize}
   \item This is the first item of an itemized list.
         Each item in the list is marked with a ``tick''.
         You don't have to worry about what kind of tick
         mark is used.

   \item This is the second item of the list.  It
         contains another list nested inside it.  The inner
         list is an \emph{enumerated} list.
         \begin{enumerate}
            \item This is the first item of an enumerated
                  list that is nested within the itemized list.

            \item This is the second item of the inner list.
                  \LaTeX\ allows you to nest lists deeper than
                  you really should.
         \end{enumerate}
         This is the rest of the second item of the outer
         list.  It is no more interesting than any other
         part of the item.
   \item This is the third item of the list.
\end{itemize}
You can even display poetry.
\begin{verse}
   There is an environment
    for verse \\             % The \\ command separates lines
   Whose features some poets % within a stanza.
   will curse.

                             % One or more blank lines separate stanzas.

   For instead of making\\
   Them do \emph{all} line breaking, \\
   It allows them to put too many words on a line when they'd rather be
   forced to be terse.
\end{verse}

Mathematical formulas may also be displayed.  A
displayed formula
is
one-line long; multiline
formulas require special formatting instructions.
   \[  \ip{\Gamma}{\psi'} = x'' + y^{2} + z_{i}^{n}\]
Don't start a paragraph with a displayed equation,
nor make one a paragraph by itself.

\end{document}               % End of document.

\section{Bench Section 187}
% synthetic bench section 187
% This is a sample LaTeX input file.  (Version of 12 August 2004.)
%
% A '%' character causes TeX to ignore all remaining text on the line,
% and is used for comments like this one.

\documentclass{article}      % Specifies the document class

                             % The preamble begins here.
\title{An Example Document}  % Declares the document's title.
\author{Leslie Lamport}      % Declares the author's name.
\date{January 21, 1994}      % Deleting this command produces today's date.

\newcommand{\ip}[2]{(#1, #2)}
                             % Defines \ip{arg1}{arg2} to mean
                             % (arg1, arg2).

%\newcommand{\ip}[2]{\langle #1 | #2\rangle}
                             % This is an alternative definition of
                             % \ip that is commented out.

\begin{document}             % End of preamble and beginning of text.

\maketitle                   % Produces the title.

This is an example input file.  Comparing it with
the output it generates can show you how to
produce a simple document of your own.

\section{Ordinary Text}      % Produces section heading.  Lower-level
                             % sections are begun with similar
                             % \subsection and \subsubsection commands.

The ends  of words and sentences are marked
  by   spaces. It  doesn't matter how many
spaces    you type; one is as good as 100.  The
end of   a line counts as a space.

One   or more   blank lines denote the  end
of  a paragraph.

Since any number of consecutive spaces are treated
like a single one, the formatting of the input
file makes no difference to
      \LaTeX,                % The \LaTeX command generates the LaTeX logo.
but it makes a difference to you.  When you use
\LaTeX, making your input file as easy to read
as possible will be a great help as you write
your document and when you change it.  This sample
file shows how you can add comments to your own input
file.

Because printing is different from typewriting,
there are a number of things that you have to do
differently when preparing an input file than if
you were just typing the document directly.
Quotation marks like
       ``this''
have to be handled specially, as do quotes within
quotes:
       ``\,`this'            % \, separates the double and single quote.
        is what I just
        wrote, not  `that'\,''.

Dashes come in three sizes: an
       intra-word
dash, a medium dash for number ranges like
       1--2,
and a punctuation
       dash---like
this.

A sentence-ending space should be larger than the
space between words within a sentence.  You
sometimes have to type special commands in
conjunction with punctuation characters to get
this right, as in the following sentence.
       Gnats, gnus, etc.\ all  % `\ ' makes an inter-word space.
       begin with G\@.         % \@ marks end-of-sentence punctuation.
You should check the spaces after periods when
reading your output to make sure you haven't
forgotten any special cases.  Generating an
ellipsis
       \ldots\               % `\ ' is needed after `\ldots' because TeX
                             % ignores spaces after command names like \ldots
                             % made from \ + letters.
                             %
                             % Note how a `%' character causes TeX to ignore
                             % the end of the input line, so these blank lines
                             % do not start a new paragraph.
                             %
with the right spacing around the periods requires
a special command.

\LaTeX\ interprets some common characters as
commands, so you must type special commands to
generate them.  These characters include the
following:
       \$ \& \% \# \{ and \}.

In printing, text is usually emphasized with an
       \emph{italic}
type style.

\begin{em}
   A long segment of text can also be emphasized
   in this way.  Text within such a segment can be
   given \emph{additional} emphasis.
\end{em}

It is sometimes necessary to prevent \LaTeX\ from
breaking a line where it might otherwise do so.
This may be at a space, as between the ``Mr.''\ and
``Jones'' in
       ``Mr.~Jones'',        % ~ produces an unbreakable interword space.
or within a word---especially when the word is a
symbol like
       \mbox{\emph{itemnum}}
that makes little sense when hyphenated across
lines.

Footnotes\footnote{This is an example of a footnote.}
pose no problem.

\LaTeX\ is good at typesetting mathematical formulas
like
       \( x-3y + z = 7 \)
or
       \( a_{1} > x^{2n} + y^{2n} > x' \)
or
       \( \ip{A}{B} = \sum_{i} a_{i} b_{i} \).
The spaces you type in a formula are
ignored.  Remember that a letter like
       $x$                   % $ ... $  and  \( ... \)  are equivalent
is a formula when it denotes a mathematical
symbol, and it should be typed as one.

\section{Displayed Text}

Text is displayed by indenting it from the left
margin.  Quotations are commonly displayed.  There
are short quotations
\begin{quote}
   This is a short quotation.  It consists of a
   single paragraph of text.  See how it is formatted.
\end{quote}
and longer ones.
\begin{quotation}
   This is a longer quotation.  It consists of two
   paragraphs of text, neither of which are
   particularly interesting.

   This is the second paragraph of the quotation.  It
   is just as dull as the first paragraph.
\end{quotation}
Another frequently-displayed structure is a list.
The following is an example of an \emph{itemized}
list.
\begin{itemize}
   \item This is the first item of an itemized list.
         Each item in the list is marked with a ``tick''.
         You don't have to worry about what kind of tick
         mark is used.

   \item This is the second item of the list.  It
         contains another list nested inside it.  The inner
         list is an \emph{enumerated} list.
         \begin{enumerate}
            \item This is the first item of an enumerated
                  list that is nested within the itemized list.

            \item This is the second item of the inner list.
                  \LaTeX\ allows you to nest lists deeper than
                  you really should.
         \end{enumerate}
         This is the rest of the second item of the outer
         list.  It is no more interesting than any other
         part of the item.
   \item This is the third item of the list.
\end{itemize}
You can even display poetry.
\begin{verse}
   There is an environment
    for verse \\             % The \\ command separates lines
   Whose features some poets % within a stanza.
   will curse.

                             % One or more blank lines separate stanzas.

   For instead of making\\
   Them do \emph{all} line breaking, \\
   It allows them to put too many words on a line when they'd rather be
   forced to be terse.
\end{verse}

Mathematical formulas may also be displayed.  A
displayed formula
is
one-line long; multiline
formulas require special formatting instructions.
   \[  \ip{\Gamma}{\psi'} = x'' + y^{2} + z_{i}^{n}\]
Don't start a paragraph with a displayed equation,
nor make one a paragraph by itself.

\end{document}               % End of document.

\section{Bench Section 188}
% synthetic bench section 188
% This is a sample LaTeX input file.  (Version of 12 August 2004.)
%
% A '%' character causes TeX to ignore all remaining text on the line,
% and is used for comments like this one.

\documentclass{article}      % Specifies the document class

                             % The preamble begins here.
\title{An Example Document}  % Declares the document's title.
\author{Leslie Lamport}      % Declares the author's name.
\date{January 21, 1994}      % Deleting this command produces today's date.

\newcommand{\ip}[2]{(#1, #2)}
                             % Defines \ip{arg1}{arg2} to mean
                             % (arg1, arg2).

%\newcommand{\ip}[2]{\langle #1 | #2\rangle}
                             % This is an alternative definition of
                             % \ip that is commented out.

\begin{document}             % End of preamble and beginning of text.

\maketitle                   % Produces the title.

This is an example input file.  Comparing it with
the output it generates can show you how to
produce a simple document of your own.

\section{Ordinary Text}      % Produces section heading.  Lower-level
                             % sections are begun with similar
                             % \subsection and \subsubsection commands.

The ends  of words and sentences are marked
  by   spaces. It  doesn't matter how many
spaces    you type; one is as good as 100.  The
end of   a line counts as a space.

One   or more   blank lines denote the  end
of  a paragraph.

Since any number of consecutive spaces are treated
like a single one, the formatting of the input
file makes no difference to
      \LaTeX,                % The \LaTeX command generates the LaTeX logo.
but it makes a difference to you.  When you use
\LaTeX, making your input file as easy to read
as possible will be a great help as you write
your document and when you change it.  This sample
file shows how you can add comments to your own input
file.

Because printing is different from typewriting,
there are a number of things that you have to do
differently when preparing an input file than if
you were just typing the document directly.
Quotation marks like
       ``this''
have to be handled specially, as do quotes within
quotes:
       ``\,`this'            % \, separates the double and single quote.
        is what I just
        wrote, not  `that'\,''.

Dashes come in three sizes: an
       intra-word
dash, a medium dash for number ranges like
       1--2,
and a punctuation
       dash---like
this.

A sentence-ending space should be larger than the
space between words within a sentence.  You
sometimes have to type special commands in
conjunction with punctuation characters to get
this right, as in the following sentence.
       Gnats, gnus, etc.\ all  % `\ ' makes an inter-word space.
       begin with G\@.         % \@ marks end-of-sentence punctuation.
You should check the spaces after periods when
reading your output to make sure you haven't
forgotten any special cases.  Generating an
ellipsis
       \ldots\               % `\ ' is needed after `\ldots' because TeX
                             % ignores spaces after command names like \ldots
                             % made from \ + letters.
                             %
                             % Note how a `%' character causes TeX to ignore
                             % the end of the input line, so these blank lines
                             % do not start a new paragraph.
                             %
with the right spacing around the periods requires
a special command.

\LaTeX\ interprets some common characters as
commands, so you must type special commands to
generate them.  These characters include the
following:
       \$ \& \% \# \{ and \}.

In printing, text is usually emphasized with an
       \emph{italic}
type style.

\begin{em}
   A long segment of text can also be emphasized
   in this way.  Text within such a segment can be
   given \emph{additional} emphasis.
\end{em}

It is sometimes necessary to prevent \LaTeX\ from
breaking a line where it might otherwise do so.
This may be at a space, as between the ``Mr.''\ and
``Jones'' in
       ``Mr.~Jones'',        % ~ produces an unbreakable interword space.
or within a word---especially when the word is a
symbol like
       \mbox{\emph{itemnum}}
that makes little sense when hyphenated across
lines.

Footnotes\footnote{This is an example of a footnote.}
pose no problem.

\LaTeX\ is good at typesetting mathematical formulas
like
       \( x-3y + z = 7 \)
or
       \( a_{1} > x^{2n} + y^{2n} > x' \)
or
       \( \ip{A}{B} = \sum_{i} a_{i} b_{i} \).
The spaces you type in a formula are
ignored.  Remember that a letter like
       $x$                   % $ ... $  and  \( ... \)  are equivalent
is a formula when it denotes a mathematical
symbol, and it should be typed as one.

\section{Displayed Text}

Text is displayed by indenting it from the left
margin.  Quotations are commonly displayed.  There
are short quotations
\begin{quote}
   This is a short quotation.  It consists of a
   single paragraph of text.  See how it is formatted.
\end{quote}
and longer ones.
\begin{quotation}
   This is a longer quotation.  It consists of two
   paragraphs of text, neither of which are
   particularly interesting.

   This is the second paragraph of the quotation.  It
   is just as dull as the first paragraph.
\end{quotation}
Another frequently-displayed structure is a list.
The following is an example of an \emph{itemized}
list.
\begin{itemize}
   \item This is the first item of an itemized list.
         Each item in the list is marked with a ``tick''.
         You don't have to worry about what kind of tick
         mark is used.

   \item This is the second item of the list.  It
         contains another list nested inside it.  The inner
         list is an \emph{enumerated} list.
         \begin{enumerate}
            \item This is the first item of an enumerated
                  list that is nested within the itemized list.

            \item This is the second item of the inner list.
                  \LaTeX\ allows you to nest lists deeper than
                  you really should.
         \end{enumerate}
         This is the rest of the second item of the outer
         list.  It is no more interesting than any other
         part of the item.
   \item This is the third item of the list.
\end{itemize}
You can even display poetry.
\begin{verse}
   There is an environment
    for verse \\             % The \\ command separates lines
   Whose features some poets % within a stanza.
   will curse.

                             % One or more blank lines separate stanzas.

   For instead of making\\
   Them do \emph{all} line breaking, \\
   It allows them to put too many words on a line when they'd rather be
   forced to be terse.
\end{verse}

Mathematical formulas may also be displayed.  A
displayed formula
is
one-line long; multiline
formulas require special formatting instructions.
   \[  \ip{\Gamma}{\psi'} = x'' + y^{2} + z_{i}^{n}\]
Don't start a paragraph with a displayed equation,
nor make one a paragraph by itself.

\end{document}               % End of document.

\section{Bench Section 189}
% synthetic bench section 189
% This is a sample LaTeX input file.  (Version of 12 August 2004.)
%
% A '%' character causes TeX to ignore all remaining text on the line,
% and is used for comments like this one.

\documentclass{article}      % Specifies the document class

                             % The preamble begins here.
\title{An Example Document}  % Declares the document's title.
\author{Leslie Lamport}      % Declares the author's name.
\date{January 21, 1994}      % Deleting this command produces today's date.

\newcommand{\ip}[2]{(#1, #2)}
                             % Defines \ip{arg1}{arg2} to mean
                             % (arg1, arg2).

%\newcommand{\ip}[2]{\langle #1 | #2\rangle}
                             % This is an alternative definition of
                             % \ip that is commented out.

\begin{document}             % End of preamble and beginning of text.

\maketitle                   % Produces the title.

This is an example input file.  Comparing it with
the output it generates can show you how to
produce a simple document of your own.

\section{Ordinary Text}      % Produces section heading.  Lower-level
                             % sections are begun with similar
                             % \subsection and \subsubsection commands.

The ends  of words and sentences are marked
  by   spaces. It  doesn't matter how many
spaces    you type; one is as good as 100.  The
end of   a line counts as a space.

One   or more   blank lines denote the  end
of  a paragraph.

Since any number of consecutive spaces are treated
like a single one, the formatting of the input
file makes no difference to
      \LaTeX,                % The \LaTeX command generates the LaTeX logo.
but it makes a difference to you.  When you use
\LaTeX, making your input file as easy to read
as possible will be a great help as you write
your document and when you change it.  This sample
file shows how you can add comments to your own input
file.

Because printing is different from typewriting,
there are a number of things that you have to do
differently when preparing an input file than if
you were just typing the document directly.
Quotation marks like
       ``this''
have to be handled specially, as do quotes within
quotes:
       ``\,`this'            % \, separates the double and single quote.
        is what I just
        wrote, not  `that'\,''.

Dashes come in three sizes: an
       intra-word
dash, a medium dash for number ranges like
       1--2,
and a punctuation
       dash---like
this.

A sentence-ending space should be larger than the
space between words within a sentence.  You
sometimes have to type special commands in
conjunction with punctuation characters to get
this right, as in the following sentence.
       Gnats, gnus, etc.\ all  % `\ ' makes an inter-word space.
       begin with G\@.         % \@ marks end-of-sentence punctuation.
You should check the spaces after periods when
reading your output to make sure you haven't
forgotten any special cases.  Generating an
ellipsis
       \ldots\               % `\ ' is needed after `\ldots' because TeX
                             % ignores spaces after command names like \ldots
                             % made from \ + letters.
                             %
                             % Note how a `%' character causes TeX to ignore
                             % the end of the input line, so these blank lines
                             % do not start a new paragraph.
                             %
with the right spacing around the periods requires
a special command.

\LaTeX\ interprets some common characters as
commands, so you must type special commands to
generate them.  These characters include the
following:
       \$ \& \% \# \{ and \}.

In printing, text is usually emphasized with an
       \emph{italic}
type style.

\begin{em}
   A long segment of text can also be emphasized
   in this way.  Text within such a segment can be
   given \emph{additional} emphasis.
\end{em}

It is sometimes necessary to prevent \LaTeX\ from
breaking a line where it might otherwise do so.
This may be at a space, as between the ``Mr.''\ and
``Jones'' in
       ``Mr.~Jones'',        % ~ produces an unbreakable interword space.
or within a word---especially when the word is a
symbol like
       \mbox{\emph{itemnum}}
that makes little sense when hyphenated across
lines.

Footnotes\footnote{This is an example of a footnote.}
pose no problem.

\LaTeX\ is good at typesetting mathematical formulas
like
       \( x-3y + z = 7 \)
or
       \( a_{1} > x^{2n} + y^{2n} > x' \)
or
       \( \ip{A}{B} = \sum_{i} a_{i} b_{i} \).
The spaces you type in a formula are
ignored.  Remember that a letter like
       $x$                   % $ ... $  and  \( ... \)  are equivalent
is a formula when it denotes a mathematical
symbol, and it should be typed as one.

\section{Displayed Text}

Text is displayed by indenting it from the left
margin.  Quotations are commonly displayed.  There
are short quotations
\begin{quote}
   This is a short quotation.  It consists of a
   single paragraph of text.  See how it is formatted.
\end{quote}
and longer ones.
\begin{quotation}
   This is a longer quotation.  It consists of two
   paragraphs of text, neither of which are
   particularly interesting.

   This is the second paragraph of the quotation.  It
   is just as dull as the first paragraph.
\end{quotation}
Another frequently-displayed structure is a list.
The following is an example of an \emph{itemized}
list.
\begin{itemize}
   \item This is the first item of an itemized list.
         Each item in the list is marked with a ``tick''.
         You don't have to worry about what kind of tick
         mark is used.

   \item This is the second item of the list.  It
         contains another list nested inside it.  The inner
         list is an \emph{enumerated} list.
         \begin{enumerate}
            \item This is the first item of an enumerated
                  list that is nested within the itemized list.

            \item This is the second item of the inner list.
                  \LaTeX\ allows you to nest lists deeper than
                  you really should.
         \end{enumerate}
         This is the rest of the second item of the outer
         list.  It is no more interesting than any other
         part of the item.
   \item This is the third item of the list.
\end{itemize}
You can even display poetry.
\begin{verse}
   There is an environment
    for verse \\             % The \\ command separates lines
   Whose features some poets % within a stanza.
   will curse.

                             % One or more blank lines separate stanzas.

   For instead of making\\
   Them do \emph{all} line breaking, \\
   It allows them to put too many words on a line when they'd rather be
   forced to be terse.
\end{verse}

Mathematical formulas may also be displayed.  A
displayed formula
is
one-line long; multiline
formulas require special formatting instructions.
   \[  \ip{\Gamma}{\psi'} = x'' + y^{2} + z_{i}^{n}\]
Don't start a paragraph with a displayed equation,
nor make one a paragraph by itself.

\end{document}               % End of document.

\section{Bench Section 190}
% synthetic bench section 190
% This is a sample LaTeX input file.  (Version of 12 August 2004.)
%
% A '%' character causes TeX to ignore all remaining text on the line,
% and is used for comments like this one.

\documentclass{article}      % Specifies the document class

                             % The preamble begins here.
\title{An Example Document}  % Declares the document's title.
\author{Leslie Lamport}      % Declares the author's name.
\date{January 21, 1994}      % Deleting this command produces today's date.

\newcommand{\ip}[2]{(#1, #2)}
                             % Defines \ip{arg1}{arg2} to mean
                             % (arg1, arg2).

%\newcommand{\ip}[2]{\langle #1 | #2\rangle}
                             % This is an alternative definition of
                             % \ip that is commented out.

\begin{document}             % End of preamble and beginning of text.

\maketitle                   % Produces the title.

This is an example input file.  Comparing it with
the output it generates can show you how to
produce a simple document of your own.

\section{Ordinary Text}      % Produces section heading.  Lower-level
                             % sections are begun with similar
                             % \subsection and \subsubsection commands.

The ends  of words and sentences are marked
  by   spaces. It  doesn't matter how many
spaces    you type; one is as good as 100.  The
end of   a line counts as a space.

One   or more   blank lines denote the  end
of  a paragraph.

Since any number of consecutive spaces are treated
like a single one, the formatting of the input
file makes no difference to
      \LaTeX,                % The \LaTeX command generates the LaTeX logo.
but it makes a difference to you.  When you use
\LaTeX, making your input file as easy to read
as possible will be a great help as you write
your document and when you change it.  This sample
file shows how you can add comments to your own input
file.

Because printing is different from typewriting,
there are a number of things that you have to do
differently when preparing an input file than if
you were just typing the document directly.
Quotation marks like
       ``this''
have to be handled specially, as do quotes within
quotes:
       ``\,`this'            % \, separates the double and single quote.
        is what I just
        wrote, not  `that'\,''.

Dashes come in three sizes: an
       intra-word
dash, a medium dash for number ranges like
       1--2,
and a punctuation
       dash---like
this.

A sentence-ending space should be larger than the
space between words within a sentence.  You
sometimes have to type special commands in
conjunction with punctuation characters to get
this right, as in the following sentence.
       Gnats, gnus, etc.\ all  % `\ ' makes an inter-word space.
       begin with G\@.         % \@ marks end-of-sentence punctuation.
You should check the spaces after periods when
reading your output to make sure you haven't
forgotten any special cases.  Generating an
ellipsis
       \ldots\               % `\ ' is needed after `\ldots' because TeX
                             % ignores spaces after command names like \ldots
                             % made from \ + letters.
                             %
                             % Note how a `%' character causes TeX to ignore
                             % the end of the input line, so these blank lines
                             % do not start a new paragraph.
                             %
with the right spacing around the periods requires
a special command.

\LaTeX\ interprets some common characters as
commands, so you must type special commands to
generate them.  These characters include the
following:
       \$ \& \% \# \{ and \}.

In printing, text is usually emphasized with an
       \emph{italic}
type style.

\begin{em}
   A long segment of text can also be emphasized
   in this way.  Text within such a segment can be
   given \emph{additional} emphasis.
\end{em}

It is sometimes necessary to prevent \LaTeX\ from
breaking a line where it might otherwise do so.
This may be at a space, as between the ``Mr.''\ and
``Jones'' in
       ``Mr.~Jones'',        % ~ produces an unbreakable interword space.
or within a word---especially when the word is a
symbol like
       \mbox{\emph{itemnum}}
that makes little sense when hyphenated across
lines.

Footnotes\footnote{This is an example of a footnote.}
pose no problem.

\LaTeX\ is good at typesetting mathematical formulas
like
       \( x-3y + z = 7 \)
or
       \( a_{1} > x^{2n} + y^{2n} > x' \)
or
       \( \ip{A}{B} = \sum_{i} a_{i} b_{i} \).
The spaces you type in a formula are
ignored.  Remember that a letter like
       $x$                   % $ ... $  and  \( ... \)  are equivalent
is a formula when it denotes a mathematical
symbol, and it should be typed as one.

\section{Displayed Text}

Text is displayed by indenting it from the left
margin.  Quotations are commonly displayed.  There
are short quotations
\begin{quote}
   This is a short quotation.  It consists of a
   single paragraph of text.  See how it is formatted.
\end{quote}
and longer ones.
\begin{quotation}
   This is a longer quotation.  It consists of two
   paragraphs of text, neither of which are
   particularly interesting.

   This is the second paragraph of the quotation.  It
   is just as dull as the first paragraph.
\end{quotation}
Another frequently-displayed structure is a list.
The following is an example of an \emph{itemized}
list.
\begin{itemize}
   \item This is the first item of an itemized list.
         Each item in the list is marked with a ``tick''.
         You don't have to worry about what kind of tick
         mark is used.

   \item This is the second item of the list.  It
         contains another list nested inside it.  The inner
         list is an \emph{enumerated} list.
         \begin{enumerate}
            \item This is the first item of an enumerated
                  list that is nested within the itemized list.

            \item This is the second item of the inner list.
                  \LaTeX\ allows you to nest lists deeper than
                  you really should.
         \end{enumerate}
         This is the rest of the second item of the outer
         list.  It is no more interesting than any other
         part of the item.
   \item This is the third item of the list.
\end{itemize}
You can even display poetry.
\begin{verse}
   There is an environment
    for verse \\             % The \\ command separates lines
   Whose features some poets % within a stanza.
   will curse.

                             % One or more blank lines separate stanzas.

   For instead of making\\
   Them do \emph{all} line breaking, \\
   It allows them to put too many words on a line when they'd rather be
   forced to be terse.
\end{verse}

Mathematical formulas may also be displayed.  A
displayed formula
is
one-line long; multiline
formulas require special formatting instructions.
   \[  \ip{\Gamma}{\psi'} = x'' + y^{2} + z_{i}^{n}\]
Don't start a paragraph with a displayed equation,
nor make one a paragraph by itself.

\end{document}               % End of document.

\section{Bench Section 191}
% synthetic bench section 191
% This is a sample LaTeX input file.  (Version of 12 August 2004.)
%
% A '%' character causes TeX to ignore all remaining text on the line,
% and is used for comments like this one.

\documentclass{article}      % Specifies the document class

                             % The preamble begins here.
\title{An Example Document}  % Declares the document's title.
\author{Leslie Lamport}      % Declares the author's name.
\date{January 21, 1994}      % Deleting this command produces today's date.

\newcommand{\ip}[2]{(#1, #2)}
                             % Defines \ip{arg1}{arg2} to mean
                             % (arg1, arg2).

%\newcommand{\ip}[2]{\langle #1 | #2\rangle}
                             % This is an alternative definition of
                             % \ip that is commented out.

\begin{document}             % End of preamble and beginning of text.

\maketitle                   % Produces the title.

This is an example input file.  Comparing it with
the output it generates can show you how to
produce a simple document of your own.

\section{Ordinary Text}      % Produces section heading.  Lower-level
                             % sections are begun with similar
                             % \subsection and \subsubsection commands.

The ends  of words and sentences are marked
  by   spaces. It  doesn't matter how many
spaces    you type; one is as good as 100.  The
end of   a line counts as a space.

One   or more   blank lines denote the  end
of  a paragraph.

Since any number of consecutive spaces are treated
like a single one, the formatting of the input
file makes no difference to
      \LaTeX,                % The \LaTeX command generates the LaTeX logo.
but it makes a difference to you.  When you use
\LaTeX, making your input file as easy to read
as possible will be a great help as you write
your document and when you change it.  This sample
file shows how you can add comments to your own input
file.

Because printing is different from typewriting,
there are a number of things that you have to do
differently when preparing an input file than if
you were just typing the document directly.
Quotation marks like
       ``this''
have to be handled specially, as do quotes within
quotes:
       ``\,`this'            % \, separates the double and single quote.
        is what I just
        wrote, not  `that'\,''.

Dashes come in three sizes: an
       intra-word
dash, a medium dash for number ranges like
       1--2,
and a punctuation
       dash---like
this.

A sentence-ending space should be larger than the
space between words within a sentence.  You
sometimes have to type special commands in
conjunction with punctuation characters to get
this right, as in the following sentence.
       Gnats, gnus, etc.\ all  % `\ ' makes an inter-word space.
       begin with G\@.         % \@ marks end-of-sentence punctuation.
You should check the spaces after periods when
reading your output to make sure you haven't
forgotten any special cases.  Generating an
ellipsis
       \ldots\               % `\ ' is needed after `\ldots' because TeX
                             % ignores spaces after command names like \ldots
                             % made from \ + letters.
                             %
                             % Note how a `%' character causes TeX to ignore
                             % the end of the input line, so these blank lines
                             % do not start a new paragraph.
                             %
with the right spacing around the periods requires
a special command.

\LaTeX\ interprets some common characters as
commands, so you must type special commands to
generate them.  These characters include the
following:
       \$ \& \% \# \{ and \}.

In printing, text is usually emphasized with an
       \emph{italic}
type style.

\begin{em}
   A long segment of text can also be emphasized
   in this way.  Text within such a segment can be
   given \emph{additional} emphasis.
\end{em}

It is sometimes necessary to prevent \LaTeX\ from
breaking a line where it might otherwise do so.
This may be at a space, as between the ``Mr.''\ and
``Jones'' in
       ``Mr.~Jones'',        % ~ produces an unbreakable interword space.
or within a word---especially when the word is a
symbol like
       \mbox{\emph{itemnum}}
that makes little sense when hyphenated across
lines.

Footnotes\footnote{This is an example of a footnote.}
pose no problem.

\LaTeX\ is good at typesetting mathematical formulas
like
       \( x-3y + z = 7 \)
or
       \( a_{1} > x^{2n} + y^{2n} > x' \)
or
       \( \ip{A}{B} = \sum_{i} a_{i} b_{i} \).
The spaces you type in a formula are
ignored.  Remember that a letter like
       $x$                   % $ ... $  and  \( ... \)  are equivalent
is a formula when it denotes a mathematical
symbol, and it should be typed as one.

\section{Displayed Text}

Text is displayed by indenting it from the left
margin.  Quotations are commonly displayed.  There
are short quotations
\begin{quote}
   This is a short quotation.  It consists of a
   single paragraph of text.  See how it is formatted.
\end{quote}
and longer ones.
\begin{quotation}
   This is a longer quotation.  It consists of two
   paragraphs of text, neither of which are
   particularly interesting.

   This is the second paragraph of the quotation.  It
   is just as dull as the first paragraph.
\end{quotation}
Another frequently-displayed structure is a list.
The following is an example of an \emph{itemized}
list.
\begin{itemize}
   \item This is the first item of an itemized list.
         Each item in the list is marked with a ``tick''.
         You don't have to worry about what kind of tick
         mark is used.

   \item This is the second item of the list.  It
         contains another list nested inside it.  The inner
         list is an \emph{enumerated} list.
         \begin{enumerate}
            \item This is the first item of an enumerated
                  list that is nested within the itemized list.

            \item This is the second item of the inner list.
                  \LaTeX\ allows you to nest lists deeper than
                  you really should.
         \end{enumerate}
         This is the rest of the second item of the outer
         list.  It is no more interesting than any other
         part of the item.
   \item This is the third item of the list.
\end{itemize}
You can even display poetry.
\begin{verse}
   There is an environment
    for verse \\             % The \\ command separates lines
   Whose features some poets % within a stanza.
   will curse.

                             % One or more blank lines separate stanzas.

   For instead of making\\
   Them do \emph{all} line breaking, \\
   It allows them to put too many words on a line when they'd rather be
   forced to be terse.
\end{verse}

Mathematical formulas may also be displayed.  A
displayed formula
is
one-line long; multiline
formulas require special formatting instructions.
   \[  \ip{\Gamma}{\psi'} = x'' + y^{2} + z_{i}^{n}\]
Don't start a paragraph with a displayed equation,
nor make one a paragraph by itself.

\end{document}               % End of document.

\section{Bench Section 192}
% synthetic bench section 192
% This is a sample LaTeX input file.  (Version of 12 August 2004.)
%
% A '%' character causes TeX to ignore all remaining text on the line,
% and is used for comments like this one.

\documentclass{article}      % Specifies the document class

                             % The preamble begins here.
\title{An Example Document}  % Declares the document's title.
\author{Leslie Lamport}      % Declares the author's name.
\date{January 21, 1994}      % Deleting this command produces today's date.

\newcommand{\ip}[2]{(#1, #2)}
                             % Defines \ip{arg1}{arg2} to mean
                             % (arg1, arg2).

%\newcommand{\ip}[2]{\langle #1 | #2\rangle}
                             % This is an alternative definition of
                             % \ip that is commented out.

\begin{document}             % End of preamble and beginning of text.

\maketitle                   % Produces the title.

This is an example input file.  Comparing it with
the output it generates can show you how to
produce a simple document of your own.

\section{Ordinary Text}      % Produces section heading.  Lower-level
                             % sections are begun with similar
                             % \subsection and \subsubsection commands.

The ends  of words and sentences are marked
  by   spaces. It  doesn't matter how many
spaces    you type; one is as good as 100.  The
end of   a line counts as a space.

One   or more   blank lines denote the  end
of  a paragraph.

Since any number of consecutive spaces are treated
like a single one, the formatting of the input
file makes no difference to
      \LaTeX,                % The \LaTeX command generates the LaTeX logo.
but it makes a difference to you.  When you use
\LaTeX, making your input file as easy to read
as possible will be a great help as you write
your document and when you change it.  This sample
file shows how you can add comments to your own input
file.

Because printing is different from typewriting,
there are a number of things that you have to do
differently when preparing an input file than if
you were just typing the document directly.
Quotation marks like
       ``this''
have to be handled specially, as do quotes within
quotes:
       ``\,`this'            % \, separates the double and single quote.
        is what I just
        wrote, not  `that'\,''.

Dashes come in three sizes: an
       intra-word
dash, a medium dash for number ranges like
       1--2,
and a punctuation
       dash---like
this.

A sentence-ending space should be larger than the
space between words within a sentence.  You
sometimes have to type special commands in
conjunction with punctuation characters to get
this right, as in the following sentence.
       Gnats, gnus, etc.\ all  % `\ ' makes an inter-word space.
       begin with G\@.         % \@ marks end-of-sentence punctuation.
You should check the spaces after periods when
reading your output to make sure you haven't
forgotten any special cases.  Generating an
ellipsis
       \ldots\               % `\ ' is needed after `\ldots' because TeX
                             % ignores spaces after command names like \ldots
                             % made from \ + letters.
                             %
                             % Note how a `%' character causes TeX to ignore
                             % the end of the input line, so these blank lines
                             % do not start a new paragraph.
                             %
with the right spacing around the periods requires
a special command.

\LaTeX\ interprets some common characters as
commands, so you must type special commands to
generate them.  These characters include the
following:
       \$ \& \% \# \{ and \}.

In printing, text is usually emphasized with an
       \emph{italic}
type style.

\begin{em}
   A long segment of text can also be emphasized
   in this way.  Text within such a segment can be
   given \emph{additional} emphasis.
\end{em}

It is sometimes necessary to prevent \LaTeX\ from
breaking a line where it might otherwise do so.
This may be at a space, as between the ``Mr.''\ and
``Jones'' in
       ``Mr.~Jones'',        % ~ produces an unbreakable interword space.
or within a word---especially when the word is a
symbol like
       \mbox{\emph{itemnum}}
that makes little sense when hyphenated across
lines.

Footnotes\footnote{This is an example of a footnote.}
pose no problem.

\LaTeX\ is good at typesetting mathematical formulas
like
       \( x-3y + z = 7 \)
or
       \( a_{1} > x^{2n} + y^{2n} > x' \)
or
       \( \ip{A}{B} = \sum_{i} a_{i} b_{i} \).
The spaces you type in a formula are
ignored.  Remember that a letter like
       $x$                   % $ ... $  and  \( ... \)  are equivalent
is a formula when it denotes a mathematical
symbol, and it should be typed as one.

\section{Displayed Text}

Text is displayed by indenting it from the left
margin.  Quotations are commonly displayed.  There
are short quotations
\begin{quote}
   This is a short quotation.  It consists of a
   single paragraph of text.  See how it is formatted.
\end{quote}
and longer ones.
\begin{quotation}
   This is a longer quotation.  It consists of two
   paragraphs of text, neither of which are
   particularly interesting.

   This is the second paragraph of the quotation.  It
   is just as dull as the first paragraph.
\end{quotation}
Another frequently-displayed structure is a list.
The following is an example of an \emph{itemized}
list.
\begin{itemize}
   \item This is the first item of an itemized list.
         Each item in the list is marked with a ``tick''.
         You don't have to worry about what kind of tick
         mark is used.

   \item This is the second item of the list.  It
         contains another list nested inside it.  The inner
         list is an \emph{enumerated} list.
         \begin{enumerate}
            \item This is the first item of an enumerated
                  list that is nested within the itemized list.

            \item This is the second item of the inner list.
                  \LaTeX\ allows you to nest lists deeper than
                  you really should.
         \end{enumerate}
         This is the rest of the second item of the outer
         list.  It is no more interesting than any other
         part of the item.
   \item This is the third item of the list.
\end{itemize}
You can even display poetry.
\begin{verse}
   There is an environment
    for verse \\             % The \\ command separates lines
   Whose features some poets % within a stanza.
   will curse.

                             % One or more blank lines separate stanzas.

   For instead of making\\
   Them do \emph{all} line breaking, \\
   It allows them to put too many words on a line when they'd rather be
   forced to be terse.
\end{verse}

Mathematical formulas may also be displayed.  A
displayed formula
is
one-line long; multiline
formulas require special formatting instructions.
   \[  \ip{\Gamma}{\psi'} = x'' + y^{2} + z_{i}^{n}\]
Don't start a paragraph with a displayed equation,
nor make one a paragraph by itself.

\end{document}               % End of document.

\section{Bench Section 193}
% synthetic bench section 193
% This is a sample LaTeX input file.  (Version of 12 August 2004.)
%
% A '%' character causes TeX to ignore all remaining text on the line,
% and is used for comments like this one.

\documentclass{article}      % Specifies the document class

                             % The preamble begins here.
\title{An Example Document}  % Declares the document's title.
\author{Leslie Lamport}      % Declares the author's name.
\date{January 21, 1994}      % Deleting this command produces today's date.

\newcommand{\ip}[2]{(#1, #2)}
                             % Defines \ip{arg1}{arg2} to mean
                             % (arg1, arg2).

%\newcommand{\ip}[2]{\langle #1 | #2\rangle}
                             % This is an alternative definition of
                             % \ip that is commented out.

\begin{document}             % End of preamble and beginning of text.

\maketitle                   % Produces the title.

This is an example input file.  Comparing it with
the output it generates can show you how to
produce a simple document of your own.

\section{Ordinary Text}      % Produces section heading.  Lower-level
                             % sections are begun with similar
                             % \subsection and \subsubsection commands.

The ends  of words and sentences are marked
  by   spaces. It  doesn't matter how many
spaces    you type; one is as good as 100.  The
end of   a line counts as a space.

One   or more   blank lines denote the  end
of  a paragraph.

Since any number of consecutive spaces are treated
like a single one, the formatting of the input
file makes no difference to
      \LaTeX,                % The \LaTeX command generates the LaTeX logo.
but it makes a difference to you.  When you use
\LaTeX, making your input file as easy to read
as possible will be a great help as you write
your document and when you change it.  This sample
file shows how you can add comments to your own input
file.

Because printing is different from typewriting,
there are a number of things that you have to do
differently when preparing an input file than if
you were just typing the document directly.
Quotation marks like
       ``this''
have to be handled specially, as do quotes within
quotes:
       ``\,`this'            % \, separates the double and single quote.
        is what I just
        wrote, not  `that'\,''.

Dashes come in three sizes: an
       intra-word
dash, a medium dash for number ranges like
       1--2,
and a punctuation
       dash---like
this.

A sentence-ending space should be larger than the
space between words within a sentence.  You
sometimes have to type special commands in
conjunction with punctuation characters to get
this right, as in the following sentence.
       Gnats, gnus, etc.\ all  % `\ ' makes an inter-word space.
       begin with G\@.         % \@ marks end-of-sentence punctuation.
You should check the spaces after periods when
reading your output to make sure you haven't
forgotten any special cases.  Generating an
ellipsis
       \ldots\               % `\ ' is needed after `\ldots' because TeX
                             % ignores spaces after command names like \ldots
                             % made from \ + letters.
                             %
                             % Note how a `%' character causes TeX to ignore
                             % the end of the input line, so these blank lines
                             % do not start a new paragraph.
                             %
with the right spacing around the periods requires
a special command.

\LaTeX\ interprets some common characters as
commands, so you must type special commands to
generate them.  These characters include the
following:
       \$ \& \% \# \{ and \}.

In printing, text is usually emphasized with an
       \emph{italic}
type style.

\begin{em}
   A long segment of text can also be emphasized
   in this way.  Text within such a segment can be
   given \emph{additional} emphasis.
\end{em}

It is sometimes necessary to prevent \LaTeX\ from
breaking a line where it might otherwise do so.
This may be at a space, as between the ``Mr.''\ and
``Jones'' in
       ``Mr.~Jones'',        % ~ produces an unbreakable interword space.
or within a word---especially when the word is a
symbol like
       \mbox{\emph{itemnum}}
that makes little sense when hyphenated across
lines.

Footnotes\footnote{This is an example of a footnote.}
pose no problem.

\LaTeX\ is good at typesetting mathematical formulas
like
       \( x-3y + z = 7 \)
or
       \( a_{1} > x^{2n} + y^{2n} > x' \)
or
       \( \ip{A}{B} = \sum_{i} a_{i} b_{i} \).
The spaces you type in a formula are
ignored.  Remember that a letter like
       $x$                   % $ ... $  and  \( ... \)  are equivalent
is a formula when it denotes a mathematical
symbol, and it should be typed as one.

\section{Displayed Text}

Text is displayed by indenting it from the left
margin.  Quotations are commonly displayed.  There
are short quotations
\begin{quote}
   This is a short quotation.  It consists of a
   single paragraph of text.  See how it is formatted.
\end{quote}
and longer ones.
\begin{quotation}
   This is a longer quotation.  It consists of two
   paragraphs of text, neither of which are
   particularly interesting.

   This is the second paragraph of the quotation.  It
   is just as dull as the first paragraph.
\end{quotation}
Another frequently-displayed structure is a list.
The following is an example of an \emph{itemized}
list.
\begin{itemize}
   \item This is the first item of an itemized list.
         Each item in the list is marked with a ``tick''.
         You don't have to worry about what kind of tick
         mark is used.

   \item This is the second item of the list.  It
         contains another list nested inside it.  The inner
         list is an \emph{enumerated} list.
         \begin{enumerate}
            \item This is the first item of an enumerated
                  list that is nested within the itemized list.

            \item This is the second item of the inner list.
                  \LaTeX\ allows you to nest lists deeper than
                  you really should.
         \end{enumerate}
         This is the rest of the second item of the outer
         list.  It is no more interesting than any other
         part of the item.
   \item This is the third item of the list.
\end{itemize}
You can even display poetry.
\begin{verse}
   There is an environment
    for verse \\             % The \\ command separates lines
   Whose features some poets % within a stanza.
   will curse.

                             % One or more blank lines separate stanzas.

   For instead of making\\
   Them do \emph{all} line breaking, \\
   It allows them to put too many words on a line when they'd rather be
   forced to be terse.
\end{verse}

Mathematical formulas may also be displayed.  A
displayed formula
is
one-line long; multiline
formulas require special formatting instructions.
   \[  \ip{\Gamma}{\psi'} = x'' + y^{2} + z_{i}^{n}\]
Don't start a paragraph with a displayed equation,
nor make one a paragraph by itself.

\end{document}               % End of document.

\section{Bench Section 194}
% synthetic bench section 194
% This is a sample LaTeX input file.  (Version of 12 August 2004.)
%
% A '%' character causes TeX to ignore all remaining text on the line,
% and is used for comments like this one.

\documentclass{article}      % Specifies the document class

                             % The preamble begins here.
\title{An Example Document}  % Declares the document's title.
\author{Leslie Lamport}      % Declares the author's name.
\date{January 21, 1994}      % Deleting this command produces today's date.

\newcommand{\ip}[2]{(#1, #2)}
                             % Defines \ip{arg1}{arg2} to mean
                             % (arg1, arg2).

%\newcommand{\ip}[2]{\langle #1 | #2\rangle}
                             % This is an alternative definition of
                             % \ip that is commented out.

\begin{document}             % End of preamble and beginning of text.

\maketitle                   % Produces the title.

This is an example input file.  Comparing it with
the output it generates can show you how to
produce a simple document of your own.

\section{Ordinary Text}      % Produces section heading.  Lower-level
                             % sections are begun with similar
                             % \subsection and \subsubsection commands.

The ends  of words and sentences are marked
  by   spaces. It  doesn't matter how many
spaces    you type; one is as good as 100.  The
end of   a line counts as a space.

One   or more   blank lines denote the  end
of  a paragraph.

Since any number of consecutive spaces are treated
like a single one, the formatting of the input
file makes no difference to
      \LaTeX,                % The \LaTeX command generates the LaTeX logo.
but it makes a difference to you.  When you use
\LaTeX, making your input file as easy to read
as possible will be a great help as you write
your document and when you change it.  This sample
file shows how you can add comments to your own input
file.

Because printing is different from typewriting,
there are a number of things that you have to do
differently when preparing an input file than if
you were just typing the document directly.
Quotation marks like
       ``this''
have to be handled specially, as do quotes within
quotes:
       ``\,`this'            % \, separates the double and single quote.
        is what I just
        wrote, not  `that'\,''.

Dashes come in three sizes: an
       intra-word
dash, a medium dash for number ranges like
       1--2,
and a punctuation
       dash---like
this.

A sentence-ending space should be larger than the
space between words within a sentence.  You
sometimes have to type special commands in
conjunction with punctuation characters to get
this right, as in the following sentence.
       Gnats, gnus, etc.\ all  % `\ ' makes an inter-word space.
       begin with G\@.         % \@ marks end-of-sentence punctuation.
You should check the spaces after periods when
reading your output to make sure you haven't
forgotten any special cases.  Generating an
ellipsis
       \ldots\               % `\ ' is needed after `\ldots' because TeX
                             % ignores spaces after command names like \ldots
                             % made from \ + letters.
                             %
                             % Note how a `%' character causes TeX to ignore
                             % the end of the input line, so these blank lines
                             % do not start a new paragraph.
                             %
with the right spacing around the periods requires
a special command.

\LaTeX\ interprets some common characters as
commands, so you must type special commands to
generate them.  These characters include the
following:
       \$ \& \% \# \{ and \}.

In printing, text is usually emphasized with an
       \emph{italic}
type style.

\begin{em}
   A long segment of text can also be emphasized
   in this way.  Text within such a segment can be
   given \emph{additional} emphasis.
\end{em}

It is sometimes necessary to prevent \LaTeX\ from
breaking a line where it might otherwise do so.
This may be at a space, as between the ``Mr.''\ and
``Jones'' in
       ``Mr.~Jones'',        % ~ produces an unbreakable interword space.
or within a word---especially when the word is a
symbol like
       \mbox{\emph{itemnum}}
that makes little sense when hyphenated across
lines.

Footnotes\footnote{This is an example of a footnote.}
pose no problem.

\LaTeX\ is good at typesetting mathematical formulas
like
       \( x-3y + z = 7 \)
or
       \( a_{1} > x^{2n} + y^{2n} > x' \)
or
       \( \ip{A}{B} = \sum_{i} a_{i} b_{i} \).
The spaces you type in a formula are
ignored.  Remember that a letter like
       $x$                   % $ ... $  and  \( ... \)  are equivalent
is a formula when it denotes a mathematical
symbol, and it should be typed as one.

\section{Displayed Text}

Text is displayed by indenting it from the left
margin.  Quotations are commonly displayed.  There
are short quotations
\begin{quote}
   This is a short quotation.  It consists of a
   single paragraph of text.  See how it is formatted.
\end{quote}
and longer ones.
\begin{quotation}
   This is a longer quotation.  It consists of two
   paragraphs of text, neither of which are
   particularly interesting.

   This is the second paragraph of the quotation.  It
   is just as dull as the first paragraph.
\end{quotation}
Another frequently-displayed structure is a list.
The following is an example of an \emph{itemized}
list.
\begin{itemize}
   \item This is the first item of an itemized list.
         Each item in the list is marked with a ``tick''.
         You don't have to worry about what kind of tick
         mark is used.

   \item This is the second item of the list.  It
         contains another list nested inside it.  The inner
         list is an \emph{enumerated} list.
         \begin{enumerate}
            \item This is the first item of an enumerated
                  list that is nested within the itemized list.

            \item This is the second item of the inner list.
                  \LaTeX\ allows you to nest lists deeper than
                  you really should.
         \end{enumerate}
         This is the rest of the second item of the outer
         list.  It is no more interesting than any other
         part of the item.
   \item This is the third item of the list.
\end{itemize}
You can even display poetry.
\begin{verse}
   There is an environment
    for verse \\             % The \\ command separates lines
   Whose features some poets % within a stanza.
   will curse.

                             % One or more blank lines separate stanzas.

   For instead of making\\
   Them do \emph{all} line breaking, \\
   It allows them to put too many words on a line when they'd rather be
   forced to be terse.
\end{verse}

Mathematical formulas may also be displayed.  A
displayed formula
is
one-line long; multiline
formulas require special formatting instructions.
   \[  \ip{\Gamma}{\psi'} = x'' + y^{2} + z_{i}^{n}\]
Don't start a paragraph with a displayed equation,
nor make one a paragraph by itself.

\end{document}               % End of document.

\section{Bench Section 195}
% synthetic bench section 195
% This is a sample LaTeX input file.  (Version of 12 August 2004.)
%
% A '%' character causes TeX to ignore all remaining text on the line,
% and is used for comments like this one.

\documentclass{article}      % Specifies the document class

                             % The preamble begins here.
\title{An Example Document}  % Declares the document's title.
\author{Leslie Lamport}      % Declares the author's name.
\date{January 21, 1994}      % Deleting this command produces today's date.

\newcommand{\ip}[2]{(#1, #2)}
                             % Defines \ip{arg1}{arg2} to mean
                             % (arg1, arg2).

%\newcommand{\ip}[2]{\langle #1 | #2\rangle}
                             % This is an alternative definition of
                             % \ip that is commented out.

\begin{document}             % End of preamble and beginning of text.

\maketitle                   % Produces the title.

This is an example input file.  Comparing it with
the output it generates can show you how to
produce a simple document of your own.

\section{Ordinary Text}      % Produces section heading.  Lower-level
                             % sections are begun with similar
                             % \subsection and \subsubsection commands.

The ends  of words and sentences are marked
  by   spaces. It  doesn't matter how many
spaces    you type; one is as good as 100.  The
end of   a line counts as a space.

One   or more   blank lines denote the  end
of  a paragraph.

Since any number of consecutive spaces are treated
like a single one, the formatting of the input
file makes no difference to
      \LaTeX,                % The \LaTeX command generates the LaTeX logo.
but it makes a difference to you.  When you use
\LaTeX, making your input file as easy to read
as possible will be a great help as you write
your document and when you change it.  This sample
file shows how you can add comments to your own input
file.

Because printing is different from typewriting,
there are a number of things that you have to do
differently when preparing an input file than if
you were just typing the document directly.
Quotation marks like
       ``this''
have to be handled specially, as do quotes within
quotes:
       ``\,`this'            % \, separates the double and single quote.
        is what I just
        wrote, not  `that'\,''.

Dashes come in three sizes: an
       intra-word
dash, a medium dash for number ranges like
       1--2,
and a punctuation
       dash---like
this.

A sentence-ending space should be larger than the
space between words within a sentence.  You
sometimes have to type special commands in
conjunction with punctuation characters to get
this right, as in the following sentence.
       Gnats, gnus, etc.\ all  % `\ ' makes an inter-word space.
       begin with G\@.         % \@ marks end-of-sentence punctuation.
You should check the spaces after periods when
reading your output to make sure you haven't
forgotten any special cases.  Generating an
ellipsis
       \ldots\               % `\ ' is needed after `\ldots' because TeX
                             % ignores spaces after command names like \ldots
                             % made from \ + letters.
                             %
                             % Note how a `%' character causes TeX to ignore
                             % the end of the input line, so these blank lines
                             % do not start a new paragraph.
                             %
with the right spacing around the periods requires
a special command.

\LaTeX\ interprets some common characters as
commands, so you must type special commands to
generate them.  These characters include the
following:
       \$ \& \% \# \{ and \}.

In printing, text is usually emphasized with an
       \emph{italic}
type style.

\begin{em}
   A long segment of text can also be emphasized
   in this way.  Text within such a segment can be
   given \emph{additional} emphasis.
\end{em}

It is sometimes necessary to prevent \LaTeX\ from
breaking a line where it might otherwise do so.
This may be at a space, as between the ``Mr.''\ and
``Jones'' in
       ``Mr.~Jones'',        % ~ produces an unbreakable interword space.
or within a word---especially when the word is a
symbol like
       \mbox{\emph{itemnum}}
that makes little sense when hyphenated across
lines.

Footnotes\footnote{This is an example of a footnote.}
pose no problem.

\LaTeX\ is good at typesetting mathematical formulas
like
       \( x-3y + z = 7 \)
or
       \( a_{1} > x^{2n} + y^{2n} > x' \)
or
       \( \ip{A}{B} = \sum_{i} a_{i} b_{i} \).
The spaces you type in a formula are
ignored.  Remember that a letter like
       $x$                   % $ ... $  and  \( ... \)  are equivalent
is a formula when it denotes a mathematical
symbol, and it should be typed as one.

\section{Displayed Text}

Text is displayed by indenting it from the left
margin.  Quotations are commonly displayed.  There
are short quotations
\begin{quote}
   This is a short quotation.  It consists of a
   single paragraph of text.  See how it is formatted.
\end{quote}
and longer ones.
\begin{quotation}
   This is a longer quotation.  It consists of two
   paragraphs of text, neither of which are
   particularly interesting.

   This is the second paragraph of the quotation.  It
   is just as dull as the first paragraph.
\end{quotation}
Another frequently-displayed structure is a list.
The following is an example of an \emph{itemized}
list.
\begin{itemize}
   \item This is the first item of an itemized list.
         Each item in the list is marked with a ``tick''.
         You don't have to worry about what kind of tick
         mark is used.

   \item This is the second item of the list.  It
         contains another list nested inside it.  The inner
         list is an \emph{enumerated} list.
         \begin{enumerate}
            \item This is the first item of an enumerated
                  list that is nested within the itemized list.

            \item This is the second item of the inner list.
                  \LaTeX\ allows you to nest lists deeper than
                  you really should.
         \end{enumerate}
         This is the rest of the second item of the outer
         list.  It is no more interesting than any other
         part of the item.
   \item This is the third item of the list.
\end{itemize}
You can even display poetry.
\begin{verse}
   There is an environment
    for verse \\             % The \\ command separates lines
   Whose features some poets % within a stanza.
   will curse.

                             % One or more blank lines separate stanzas.

   For instead of making\\
   Them do \emph{all} line breaking, \\
   It allows them to put too many words on a line when they'd rather be
   forced to be terse.
\end{verse}

Mathematical formulas may also be displayed.  A
displayed formula
is
one-line long; multiline
formulas require special formatting instructions.
   \[  \ip{\Gamma}{\psi'} = x'' + y^{2} + z_{i}^{n}\]
Don't start a paragraph with a displayed equation,
nor make one a paragraph by itself.

\end{document}               % End of document.

\section{Bench Section 196}
% synthetic bench section 196
% This is a sample LaTeX input file.  (Version of 12 August 2004.)
%
% A '%' character causes TeX to ignore all remaining text on the line,
% and is used for comments like this one.

\documentclass{article}      % Specifies the document class

                             % The preamble begins here.
\title{An Example Document}  % Declares the document's title.
\author{Leslie Lamport}      % Declares the author's name.
\date{January 21, 1994}      % Deleting this command produces today's date.

\newcommand{\ip}[2]{(#1, #2)}
                             % Defines \ip{arg1}{arg2} to mean
                             % (arg1, arg2).

%\newcommand{\ip}[2]{\langle #1 | #2\rangle}
                             % This is an alternative definition of
                             % \ip that is commented out.

\begin{document}             % End of preamble and beginning of text.

\maketitle                   % Produces the title.

This is an example input file.  Comparing it with
the output it generates can show you how to
produce a simple document of your own.

\section{Ordinary Text}      % Produces section heading.  Lower-level
                             % sections are begun with similar
                             % \subsection and \subsubsection commands.

The ends  of words and sentences are marked
  by   spaces. It  doesn't matter how many
spaces    you type; one is as good as 100.  The
end of   a line counts as a space.

One   or more   blank lines denote the  end
of  a paragraph.

Since any number of consecutive spaces are treated
like a single one, the formatting of the input
file makes no difference to
      \LaTeX,                % The \LaTeX command generates the LaTeX logo.
but it makes a difference to you.  When you use
\LaTeX, making your input file as easy to read
as possible will be a great help as you write
your document and when you change it.  This sample
file shows how you can add comments to your own input
file.

Because printing is different from typewriting,
there are a number of things that you have to do
differently when preparing an input file than if
you were just typing the document directly.
Quotation marks like
       ``this''
have to be handled specially, as do quotes within
quotes:
       ``\,`this'            % \, separates the double and single quote.
        is what I just
        wrote, not  `that'\,''.

Dashes come in three sizes: an
       intra-word
dash, a medium dash for number ranges like
       1--2,
and a punctuation
       dash---like
this.

A sentence-ending space should be larger than the
space between words within a sentence.  You
sometimes have to type special commands in
conjunction with punctuation characters to get
this right, as in the following sentence.
       Gnats, gnus, etc.\ all  % `\ ' makes an inter-word space.
       begin with G\@.         % \@ marks end-of-sentence punctuation.
You should check the spaces after periods when
reading your output to make sure you haven't
forgotten any special cases.  Generating an
ellipsis
       \ldots\               % `\ ' is needed after `\ldots' because TeX
                             % ignores spaces after command names like \ldots
                             % made from \ + letters.
                             %
                             % Note how a `%' character causes TeX to ignore
                             % the end of the input line, so these blank lines
                             % do not start a new paragraph.
                             %
with the right spacing around the periods requires
a special command.

\LaTeX\ interprets some common characters as
commands, so you must type special commands to
generate them.  These characters include the
following:
       \$ \& \% \# \{ and \}.

In printing, text is usually emphasized with an
       \emph{italic}
type style.

\begin{em}
   A long segment of text can also be emphasized
   in this way.  Text within such a segment can be
   given \emph{additional} emphasis.
\end{em}

It is sometimes necessary to prevent \LaTeX\ from
breaking a line where it might otherwise do so.
This may be at a space, as between the ``Mr.''\ and
``Jones'' in
       ``Mr.~Jones'',        % ~ produces an unbreakable interword space.
or within a word---especially when the word is a
symbol like
       \mbox{\emph{itemnum}}
that makes little sense when hyphenated across
lines.

Footnotes\footnote{This is an example of a footnote.}
pose no problem.

\LaTeX\ is good at typesetting mathematical formulas
like
       \( x-3y + z = 7 \)
or
       \( a_{1} > x^{2n} + y^{2n} > x' \)
or
       \( \ip{A}{B} = \sum_{i} a_{i} b_{i} \).
The spaces you type in a formula are
ignored.  Remember that a letter like
       $x$                   % $ ... $  and  \( ... \)  are equivalent
is a formula when it denotes a mathematical
symbol, and it should be typed as one.

\section{Displayed Text}

Text is displayed by indenting it from the left
margin.  Quotations are commonly displayed.  There
are short quotations
\begin{quote}
   This is a short quotation.  It consists of a
   single paragraph of text.  See how it is formatted.
\end{quote}
and longer ones.
\begin{quotation}
   This is a longer quotation.  It consists of two
   paragraphs of text, neither of which are
   particularly interesting.

   This is the second paragraph of the quotation.  It
   is just as dull as the first paragraph.
\end{quotation}
Another frequently-displayed structure is a list.
The following is an example of an \emph{itemized}
list.
\begin{itemize}
   \item This is the first item of an itemized list.
         Each item in the list is marked with a ``tick''.
         You don't have to worry about what kind of tick
         mark is used.

   \item This is the second item of the list.  It
         contains another list nested inside it.  The inner
         list is an \emph{enumerated} list.
         \begin{enumerate}
            \item This is the first item of an enumerated
                  list that is nested within the itemized list.

            \item This is the second item of the inner list.
                  \LaTeX\ allows you to nest lists deeper than
                  you really should.
         \end{enumerate}
         This is the rest of the second item of the outer
         list.  It is no more interesting than any other
         part of the item.
   \item This is the third item of the list.
\end{itemize}
You can even display poetry.
\begin{verse}
   There is an environment
    for verse \\             % The \\ command separates lines
   Whose features some poets % within a stanza.
   will curse.

                             % One or more blank lines separate stanzas.

   For instead of making\\
   Them do \emph{all} line breaking, \\
   It allows them to put too many words on a line when they'd rather be
   forced to be terse.
\end{verse}

Mathematical formulas may also be displayed.  A
displayed formula
is
one-line long; multiline
formulas require special formatting instructions.
   \[  \ip{\Gamma}{\psi'} = x'' + y^{2} + z_{i}^{n}\]
Don't start a paragraph with a displayed equation,
nor make one a paragraph by itself.

\end{document}               % End of document.

\section{Bench Section 197}
% synthetic bench section 197
% This is a sample LaTeX input file.  (Version of 12 August 2004.)
%
% A '%' character causes TeX to ignore all remaining text on the line,
% and is used for comments like this one.

\documentclass{article}      % Specifies the document class

                             % The preamble begins here.
\title{An Example Document}  % Declares the document's title.
\author{Leslie Lamport}      % Declares the author's name.
\date{January 21, 1994}      % Deleting this command produces today's date.

\newcommand{\ip}[2]{(#1, #2)}
                             % Defines \ip{arg1}{arg2} to mean
                             % (arg1, arg2).

%\newcommand{\ip}[2]{\langle #1 | #2\rangle}
                             % This is an alternative definition of
                             % \ip that is commented out.

\begin{document}             % End of preamble and beginning of text.

\maketitle                   % Produces the title.

This is an example input file.  Comparing it with
the output it generates can show you how to
produce a simple document of your own.

\section{Ordinary Text}      % Produces section heading.  Lower-level
                             % sections are begun with similar
                             % \subsection and \subsubsection commands.

The ends  of words and sentences are marked
  by   spaces. It  doesn't matter how many
spaces    you type; one is as good as 100.  The
end of   a line counts as a space.

One   or more   blank lines denote the  end
of  a paragraph.

Since any number of consecutive spaces are treated
like a single one, the formatting of the input
file makes no difference to
      \LaTeX,                % The \LaTeX command generates the LaTeX logo.
but it makes a difference to you.  When you use
\LaTeX, making your input file as easy to read
as possible will be a great help as you write
your document and when you change it.  This sample
file shows how you can add comments to your own input
file.

Because printing is different from typewriting,
there are a number of things that you have to do
differently when preparing an input file than if
you were just typing the document directly.
Quotation marks like
       ``this''
have to be handled specially, as do quotes within
quotes:
       ``\,`this'            % \, separates the double and single quote.
        is what I just
        wrote, not  `that'\,''.

Dashes come in three sizes: an
       intra-word
dash, a medium dash for number ranges like
       1--2,
and a punctuation
       dash---like
this.

A sentence-ending space should be larger than the
space between words within a sentence.  You
sometimes have to type special commands in
conjunction with punctuation characters to get
this right, as in the following sentence.
       Gnats, gnus, etc.\ all  % `\ ' makes an inter-word space.
       begin with G\@.         % \@ marks end-of-sentence punctuation.
You should check the spaces after periods when
reading your output to make sure you haven't
forgotten any special cases.  Generating an
ellipsis
       \ldots\               % `\ ' is needed after `\ldots' because TeX
                             % ignores spaces after command names like \ldots
                             % made from \ + letters.
                             %
                             % Note how a `%' character causes TeX to ignore
                             % the end of the input line, so these blank lines
                             % do not start a new paragraph.
                             %
with the right spacing around the periods requires
a special command.

\LaTeX\ interprets some common characters as
commands, so you must type special commands to
generate them.  These characters include the
following:
       \$ \& \% \# \{ and \}.

In printing, text is usually emphasized with an
       \emph{italic}
type style.

\begin{em}
   A long segment of text can also be emphasized
   in this way.  Text within such a segment can be
   given \emph{additional} emphasis.
\end{em}

It is sometimes necessary to prevent \LaTeX\ from
breaking a line where it might otherwise do so.
This may be at a space, as between the ``Mr.''\ and
``Jones'' in
       ``Mr.~Jones'',        % ~ produces an unbreakable interword space.
or within a word---especially when the word is a
symbol like
       \mbox{\emph{itemnum}}
that makes little sense when hyphenated across
lines.

Footnotes\footnote{This is an example of a footnote.}
pose no problem.

\LaTeX\ is good at typesetting mathematical formulas
like
       \( x-3y + z = 7 \)
or
       \( a_{1} > x^{2n} + y^{2n} > x' \)
or
       \( \ip{A}{B} = \sum_{i} a_{i} b_{i} \).
The spaces you type in a formula are
ignored.  Remember that a letter like
       $x$                   % $ ... $  and  \( ... \)  are equivalent
is a formula when it denotes a mathematical
symbol, and it should be typed as one.

\section{Displayed Text}

Text is displayed by indenting it from the left
margin.  Quotations are commonly displayed.  There
are short quotations
\begin{quote}
   This is a short quotation.  It consists of a
   single paragraph of text.  See how it is formatted.
\end{quote}
and longer ones.
\begin{quotation}
   This is a longer quotation.  It consists of two
   paragraphs of text, neither of which are
   particularly interesting.

   This is the second paragraph of the quotation.  It
   is just as dull as the first paragraph.
\end{quotation}
Another frequently-displayed structure is a list.
The following is an example of an \emph{itemized}
list.
\begin{itemize}
   \item This is the first item of an itemized list.
         Each item in the list is marked with a ``tick''.
         You don't have to worry about what kind of tick
         mark is used.

   \item This is the second item of the list.  It
         contains another list nested inside it.  The inner
         list is an \emph{enumerated} list.
         \begin{enumerate}
            \item This is the first item of an enumerated
                  list that is nested within the itemized list.

            \item This is the second item of the inner list.
                  \LaTeX\ allows you to nest lists deeper than
                  you really should.
         \end{enumerate}
         This is the rest of the second item of the outer
         list.  It is no more interesting than any other
         part of the item.
   \item This is the third item of the list.
\end{itemize}
You can even display poetry.
\begin{verse}
   There is an environment
    for verse \\             % The \\ command separates lines
   Whose features some poets % within a stanza.
   will curse.

                             % One or more blank lines separate stanzas.

   For instead of making\\
   Them do \emph{all} line breaking, \\
   It allows them to put too many words on a line when they'd rather be
   forced to be terse.
\end{verse}

Mathematical formulas may also be displayed.  A
displayed formula
is
one-line long; multiline
formulas require special formatting instructions.
   \[  \ip{\Gamma}{\psi'} = x'' + y^{2} + z_{i}^{n}\]
Don't start a paragraph with a displayed equation,
nor make one a paragraph by itself.

\end{document}               % End of document.

\section{Bench Section 198}
% synthetic bench section 198
% This is a sample LaTeX input file.  (Version of 12 August 2004.)
%
% A '%' character causes TeX to ignore all remaining text on the line,
% and is used for comments like this one.

\documentclass{article}      % Specifies the document class

                             % The preamble begins here.
\title{An Example Document}  % Declares the document's title.
\author{Leslie Lamport}      % Declares the author's name.
\date{January 21, 1994}      % Deleting this command produces today's date.

\newcommand{\ip}[2]{(#1, #2)}
                             % Defines \ip{arg1}{arg2} to mean
                             % (arg1, arg2).

%\newcommand{\ip}[2]{\langle #1 | #2\rangle}
                             % This is an alternative definition of
                             % \ip that is commented out.

\begin{document}             % End of preamble and beginning of text.

\maketitle                   % Produces the title.

This is an example input file.  Comparing it with
the output it generates can show you how to
produce a simple document of your own.

\section{Ordinary Text}      % Produces section heading.  Lower-level
                             % sections are begun with similar
                             % \subsection and \subsubsection commands.

The ends  of words and sentences are marked
  by   spaces. It  doesn't matter how many
spaces    you type; one is as good as 100.  The
end of   a line counts as a space.

One   or more   blank lines denote the  end
of  a paragraph.

Since any number of consecutive spaces are treated
like a single one, the formatting of the input
file makes no difference to
      \LaTeX,                % The \LaTeX command generates the LaTeX logo.
but it makes a difference to you.  When you use
\LaTeX, making your input file as easy to read
as possible will be a great help as you write
your document and when you change it.  This sample
file shows how you can add comments to your own input
file.

Because printing is different from typewriting,
there are a number of things that you have to do
differently when preparing an input file than if
you were just typing the document directly.
Quotation marks like
       ``this''
have to be handled specially, as do quotes within
quotes:
       ``\,`this'            % \, separates the double and single quote.
        is what I just
        wrote, not  `that'\,''.

Dashes come in three sizes: an
       intra-word
dash, a medium dash for number ranges like
       1--2,
and a punctuation
       dash---like
this.

A sentence-ending space should be larger than the
space between words within a sentence.  You
sometimes have to type special commands in
conjunction with punctuation characters to get
this right, as in the following sentence.
       Gnats, gnus, etc.\ all  % `\ ' makes an inter-word space.
       begin with G\@.         % \@ marks end-of-sentence punctuation.
You should check the spaces after periods when
reading your output to make sure you haven't
forgotten any special cases.  Generating an
ellipsis
       \ldots\               % `\ ' is needed after `\ldots' because TeX
                             % ignores spaces after command names like \ldots
                             % made from \ + letters.
                             %
                             % Note how a `%' character causes TeX to ignore
                             % the end of the input line, so these blank lines
                             % do not start a new paragraph.
                             %
with the right spacing around the periods requires
a special command.

\LaTeX\ interprets some common characters as
commands, so you must type special commands to
generate them.  These characters include the
following:
       \$ \& \% \# \{ and \}.

In printing, text is usually emphasized with an
       \emph{italic}
type style.

\begin{em}
   A long segment of text can also be emphasized
   in this way.  Text within such a segment can be
   given \emph{additional} emphasis.
\end{em}

It is sometimes necessary to prevent \LaTeX\ from
breaking a line where it might otherwise do so.
This may be at a space, as between the ``Mr.''\ and
``Jones'' in
       ``Mr.~Jones'',        % ~ produces an unbreakable interword space.
or within a word---especially when the word is a
symbol like
       \mbox{\emph{itemnum}}
that makes little sense when hyphenated across
lines.

Footnotes\footnote{This is an example of a footnote.}
pose no problem.

\LaTeX\ is good at typesetting mathematical formulas
like
       \( x-3y + z = 7 \)
or
       \( a_{1} > x^{2n} + y^{2n} > x' \)
or
       \( \ip{A}{B} = \sum_{i} a_{i} b_{i} \).
The spaces you type in a formula are
ignored.  Remember that a letter like
       $x$                   % $ ... $  and  \( ... \)  are equivalent
is a formula when it denotes a mathematical
symbol, and it should be typed as one.

\section{Displayed Text}

Text is displayed by indenting it from the left
margin.  Quotations are commonly displayed.  There
are short quotations
\begin{quote}
   This is a short quotation.  It consists of a
   single paragraph of text.  See how it is formatted.
\end{quote}
and longer ones.
\begin{quotation}
   This is a longer quotation.  It consists of two
   paragraphs of text, neither of which are
   particularly interesting.

   This is the second paragraph of the quotation.  It
   is just as dull as the first paragraph.
\end{quotation}
Another frequently-displayed structure is a list.
The following is an example of an \emph{itemized}
list.
\begin{itemize}
   \item This is the first item of an itemized list.
         Each item in the list is marked with a ``tick''.
         You don't have to worry about what kind of tick
         mark is used.

   \item This is the second item of the list.  It
         contains another list nested inside it.  The inner
         list is an \emph{enumerated} list.
         \begin{enumerate}
            \item This is the first item of an enumerated
                  list that is nested within the itemized list.

            \item This is the second item of the inner list.
                  \LaTeX\ allows you to nest lists deeper than
                  you really should.
         \end{enumerate}
         This is the rest of the second item of the outer
         list.  It is no more interesting than any other
         part of the item.
   \item This is the third item of the list.
\end{itemize}
You can even display poetry.
\begin{verse}
   There is an environment
    for verse \\             % The \\ command separates lines
   Whose features some poets % within a stanza.
   will curse.

                             % One or more blank lines separate stanzas.

   For instead of making\\
   Them do \emph{all} line breaking, \\
   It allows them to put too many words on a line when they'd rather be
   forced to be terse.
\end{verse}

Mathematical formulas may also be displayed.  A
displayed formula
is
one-line long; multiline
formulas require special formatting instructions.
   \[  \ip{\Gamma}{\psi'} = x'' + y^{2} + z_{i}^{n}\]
Don't start a paragraph with a displayed equation,
nor make one a paragraph by itself.

\end{document}               % End of document.

\section{Bench Section 199}
% synthetic bench section 199
% This is a sample LaTeX input file.  (Version of 12 August 2004.)
%
% A '%' character causes TeX to ignore all remaining text on the line,
% and is used for comments like this one.

\documentclass{article}      % Specifies the document class

                             % The preamble begins here.
\title{An Example Document}  % Declares the document's title.
\author{Leslie Lamport}      % Declares the author's name.
\date{January 21, 1994}      % Deleting this command produces today's date.

\newcommand{\ip}[2]{(#1, #2)}
                             % Defines \ip{arg1}{arg2} to mean
                             % (arg1, arg2).

%\newcommand{\ip}[2]{\langle #1 | #2\rangle}
                             % This is an alternative definition of
                             % \ip that is commented out.

\begin{document}             % End of preamble and beginning of text.

\maketitle                   % Produces the title.

This is an example input file.  Comparing it with
the output it generates can show you how to
produce a simple document of your own.

\section{Ordinary Text}      % Produces section heading.  Lower-level
                             % sections are begun with similar
                             % \subsection and \subsubsection commands.

The ends  of words and sentences are marked
  by   spaces. It  doesn't matter how many
spaces    you type; one is as good as 100.  The
end of   a line counts as a space.

One   or more   blank lines denote the  end
of  a paragraph.

Since any number of consecutive spaces are treated
like a single one, the formatting of the input
file makes no difference to
      \LaTeX,                % The \LaTeX command generates the LaTeX logo.
but it makes a difference to you.  When you use
\LaTeX, making your input file as easy to read
as possible will be a great help as you write
your document and when you change it.  This sample
file shows how you can add comments to your own input
file.

Because printing is different from typewriting,
there are a number of things that you have to do
differently when preparing an input file than if
you were just typing the document directly.
Quotation marks like
       ``this''
have to be handled specially, as do quotes within
quotes:
       ``\,`this'            % \, separates the double and single quote.
        is what I just
        wrote, not  `that'\,''.

Dashes come in three sizes: an
       intra-word
dash, a medium dash for number ranges like
       1--2,
and a punctuation
       dash---like
this.

A sentence-ending space should be larger than the
space between words within a sentence.  You
sometimes have to type special commands in
conjunction with punctuation characters to get
this right, as in the following sentence.
       Gnats, gnus, etc.\ all  % `\ ' makes an inter-word space.
       begin with G\@.         % \@ marks end-of-sentence punctuation.
You should check the spaces after periods when
reading your output to make sure you haven't
forgotten any special cases.  Generating an
ellipsis
       \ldots\               % `\ ' is needed after `\ldots' because TeX
                             % ignores spaces after command names like \ldots
                             % made from \ + letters.
                             %
                             % Note how a `%' character causes TeX to ignore
                             % the end of the input line, so these blank lines
                             % do not start a new paragraph.
                             %
with the right spacing around the periods requires
a special command.

\LaTeX\ interprets some common characters as
commands, so you must type special commands to
generate them.  These characters include the
following:
       \$ \& \% \# \{ and \}.

In printing, text is usually emphasized with an
       \emph{italic}
type style.

\begin{em}
   A long segment of text can also be emphasized
   in this way.  Text within such a segment can be
   given \emph{additional} emphasis.
\end{em}

It is sometimes necessary to prevent \LaTeX\ from
breaking a line where it might otherwise do so.
This may be at a space, as between the ``Mr.''\ and
``Jones'' in
       ``Mr.~Jones'',        % ~ produces an unbreakable interword space.
or within a word---especially when the word is a
symbol like
       \mbox{\emph{itemnum}}
that makes little sense when hyphenated across
lines.

Footnotes\footnote{This is an example of a footnote.}
pose no problem.

\LaTeX\ is good at typesetting mathematical formulas
like
       \( x-3y + z = 7 \)
or
       \( a_{1} > x^{2n} + y^{2n} > x' \)
or
       \( \ip{A}{B} = \sum_{i} a_{i} b_{i} \).
The spaces you type in a formula are
ignored.  Remember that a letter like
       $x$                   % $ ... $  and  \( ... \)  are equivalent
is a formula when it denotes a mathematical
symbol, and it should be typed as one.

\section{Displayed Text}

Text is displayed by indenting it from the left
margin.  Quotations are commonly displayed.  There
are short quotations
\begin{quote}
   This is a short quotation.  It consists of a
   single paragraph of text.  See how it is formatted.
\end{quote}
and longer ones.
\begin{quotation}
   This is a longer quotation.  It consists of two
   paragraphs of text, neither of which are
   particularly interesting.

   This is the second paragraph of the quotation.  It
   is just as dull as the first paragraph.
\end{quotation}
Another frequently-displayed structure is a list.
The following is an example of an \emph{itemized}
list.
\begin{itemize}
   \item This is the first item of an itemized list.
         Each item in the list is marked with a ``tick''.
         You don't have to worry about what kind of tick
         mark is used.

   \item This is the second item of the list.  It
         contains another list nested inside it.  The inner
         list is an \emph{enumerated} list.
         \begin{enumerate}
            \item This is the first item of an enumerated
                  list that is nested within the itemized list.

            \item This is the second item of the inner list.
                  \LaTeX\ allows you to nest lists deeper than
                  you really should.
         \end{enumerate}
         This is the rest of the second item of the outer
         list.  It is no more interesting than any other
         part of the item.
   \item This is the third item of the list.
\end{itemize}
You can even display poetry.
\begin{verse}
   There is an environment
    for verse \\             % The \\ command separates lines
   Whose features some poets % within a stanza.
   will curse.

                             % One or more blank lines separate stanzas.

   For instead of making\\
   Them do \emph{all} line breaking, \\
   It allows them to put too many words on a line when they'd rather be
   forced to be terse.
\end{verse}

Mathematical formulas may also be displayed.  A
displayed formula
is
one-line long; multiline
formulas require special formatting instructions.
   \[  \ip{\Gamma}{\psi'} = x'' + y^{2} + z_{i}^{n}\]
Don't start a paragraph with a displayed equation,
nor make one a paragraph by itself.

\end{document}               % End of document.

\section{Bench Section 200}
% synthetic bench section 200
% This is a sample LaTeX input file.  (Version of 12 August 2004.)
%
% A '%' character causes TeX to ignore all remaining text on the line,
% and is used for comments like this one.

\documentclass{article}      % Specifies the document class

                             % The preamble begins here.
\title{An Example Document}  % Declares the document's title.
\author{Leslie Lamport}      % Declares the author's name.
\date{January 21, 1994}      % Deleting this command produces today's date.

\newcommand{\ip}[2]{(#1, #2)}
                             % Defines \ip{arg1}{arg2} to mean
                             % (arg1, arg2).

%\newcommand{\ip}[2]{\langle #1 | #2\rangle}
                             % This is an alternative definition of
                             % \ip that is commented out.

\begin{document}             % End of preamble and beginning of text.

\maketitle                   % Produces the title.

This is an example input file.  Comparing it with
the output it generates can show you how to
produce a simple document of your own.

\section{Ordinary Text}      % Produces section heading.  Lower-level
                             % sections are begun with similar
                             % \subsection and \subsubsection commands.

The ends  of words and sentences are marked
  by   spaces. It  doesn't matter how many
spaces    you type; one is as good as 100.  The
end of   a line counts as a space.

One   or more   blank lines denote the  end
of  a paragraph.

Since any number of consecutive spaces are treated
like a single one, the formatting of the input
file makes no difference to
      \LaTeX,                % The \LaTeX command generates the LaTeX logo.
but it makes a difference to you.  When you use
\LaTeX, making your input file as easy to read
as possible will be a great help as you write
your document and when you change it.  This sample
file shows how you can add comments to your own input
file.

Because printing is different from typewriting,
there are a number of things that you have to do
differently when preparing an input file than if
you were just typing the document directly.
Quotation marks like
       ``this''
have to be handled specially, as do quotes within
quotes:
       ``\,`this'            % \, separates the double and single quote.
        is what I just
        wrote, not  `that'\,''.

Dashes come in three sizes: an
       intra-word
dash, a medium dash for number ranges like
       1--2,
and a punctuation
       dash---like
this.

A sentence-ending space should be larger than the
space between words within a sentence.  You
sometimes have to type special commands in
conjunction with punctuation characters to get
this right, as in the following sentence.
       Gnats, gnus, etc.\ all  % `\ ' makes an inter-word space.
       begin with G\@.         % \@ marks end-of-sentence punctuation.
You should check the spaces after periods when
reading your output to make sure you haven't
forgotten any special cases.  Generating an
ellipsis
       \ldots\               % `\ ' is needed after `\ldots' because TeX
                             % ignores spaces after command names like \ldots
                             % made from \ + letters.
                             %
                             % Note how a `%' character causes TeX to ignore
                             % the end of the input line, so these blank lines
                             % do not start a new paragraph.
                             %
with the right spacing around the periods requires
a special command.

\LaTeX\ interprets some common characters as
commands, so you must type special commands to
generate them.  These characters include the
following:
       \$ \& \% \# \{ and \}.

In printing, text is usually emphasized with an
       \emph{italic}
type style.

\begin{em}
   A long segment of text can also be emphasized
   in this way.  Text within such a segment can be
   given \emph{additional} emphasis.
\end{em}

It is sometimes necessary to prevent \LaTeX\ from
breaking a line where it might otherwise do so.
This may be at a space, as between the ``Mr.''\ and
``Jones'' in
       ``Mr.~Jones'',        % ~ produces an unbreakable interword space.
or within a word---especially when the word is a
symbol like
       \mbox{\emph{itemnum}}
that makes little sense when hyphenated across
lines.

Footnotes\footnote{This is an example of a footnote.}
pose no problem.

\LaTeX\ is good at typesetting mathematical formulas
like
       \( x-3y + z = 7 \)
or
       \( a_{1} > x^{2n} + y^{2n} > x' \)
or
       \( \ip{A}{B} = \sum_{i} a_{i} b_{i} \).
The spaces you type in a formula are
ignored.  Remember that a letter like
       $x$                   % $ ... $  and  \( ... \)  are equivalent
is a formula when it denotes a mathematical
symbol, and it should be typed as one.

\section{Displayed Text}

Text is displayed by indenting it from the left
margin.  Quotations are commonly displayed.  There
are short quotations
\begin{quote}
   This is a short quotation.  It consists of a
   single paragraph of text.  See how it is formatted.
\end{quote}
and longer ones.
\begin{quotation}
   This is a longer quotation.  It consists of two
   paragraphs of text, neither of which are
   particularly interesting.

   This is the second paragraph of the quotation.  It
   is just as dull as the first paragraph.
\end{quotation}
Another frequently-displayed structure is a list.
The following is an example of an \emph{itemized}
list.
\begin{itemize}
   \item This is the first item of an itemized list.
         Each item in the list is marked with a ``tick''.
         You don't have to worry about what kind of tick
         mark is used.

   \item This is the second item of the list.  It
         contains another list nested inside it.  The inner
         list is an \emph{enumerated} list.
         \begin{enumerate}
            \item This is the first item of an enumerated
                  list that is nested within the itemized list.

            \item This is the second item of the inner list.
                  \LaTeX\ allows you to nest lists deeper than
                  you really should.
         \end{enumerate}
         This is the rest of the second item of the outer
         list.  It is no more interesting than any other
         part of the item.
   \item This is the third item of the list.
\end{itemize}
You can even display poetry.
\begin{verse}
   There is an environment
    for verse \\             % The \\ command separates lines
   Whose features some poets % within a stanza.
   will curse.

                             % One or more blank lines separate stanzas.

   For instead of making\\
   Them do \emph{all} line breaking, \\
   It allows them to put too many words on a line when they'd rather be
   forced to be terse.
\end{verse}

Mathematical formulas may also be displayed.  A
displayed formula
is
one-line long; multiline
formulas require special formatting instructions.
   \[  \ip{\Gamma}{\psi'} = x'' + y^{2} + z_{i}^{n}\]
Don't start a paragraph with a displayed equation,
nor make one a paragraph by itself.

\end{document}               % End of document.

\section{Bench Section 201}
% synthetic bench section 201
% This is a sample LaTeX input file.  (Version of 12 August 2004.)
%
% A '%' character causes TeX to ignore all remaining text on the line,
% and is used for comments like this one.

\documentclass{article}      % Specifies the document class

                             % The preamble begins here.
\title{An Example Document}  % Declares the document's title.
\author{Leslie Lamport}      % Declares the author's name.
\date{January 21, 1994}      % Deleting this command produces today's date.

\newcommand{\ip}[2]{(#1, #2)}
                             % Defines \ip{arg1}{arg2} to mean
                             % (arg1, arg2).

%\newcommand{\ip}[2]{\langle #1 | #2\rangle}
                             % This is an alternative definition of
                             % \ip that is commented out.

\begin{document}             % End of preamble and beginning of text.

\maketitle                   % Produces the title.

This is an example input file.  Comparing it with
the output it generates can show you how to
produce a simple document of your own.

\section{Ordinary Text}      % Produces section heading.  Lower-level
                             % sections are begun with similar
                             % \subsection and \subsubsection commands.

The ends  of words and sentences are marked
  by   spaces. It  doesn't matter how many
spaces    you type; one is as good as 100.  The
end of   a line counts as a space.

One   or more   blank lines denote the  end
of  a paragraph.

Since any number of consecutive spaces are treated
like a single one, the formatting of the input
file makes no difference to
      \LaTeX,                % The \LaTeX command generates the LaTeX logo.
but it makes a difference to you.  When you use
\LaTeX, making your input file as easy to read
as possible will be a great help as you write
your document and when you change it.  This sample
file shows how you can add comments to your own input
file.

Because printing is different from typewriting,
there are a number of things that you have to do
differently when preparing an input file than if
you were just typing the document directly.
Quotation marks like
       ``this''
have to be handled specially, as do quotes within
quotes:
       ``\,`this'            % \, separates the double and single quote.
        is what I just
        wrote, not  `that'\,''.

Dashes come in three sizes: an
       intra-word
dash, a medium dash for number ranges like
       1--2,
and a punctuation
       dash---like
this.

A sentence-ending space should be larger than the
space between words within a sentence.  You
sometimes have to type special commands in
conjunction with punctuation characters to get
this right, as in the following sentence.
       Gnats, gnus, etc.\ all  % `\ ' makes an inter-word space.
       begin with G\@.         % \@ marks end-of-sentence punctuation.
You should check the spaces after periods when
reading your output to make sure you haven't
forgotten any special cases.  Generating an
ellipsis
       \ldots\               % `\ ' is needed after `\ldots' because TeX
                             % ignores spaces after command names like \ldots
                             % made from \ + letters.
                             %
                             % Note how a `%' character causes TeX to ignore
                             % the end of the input line, so these blank lines
                             % do not start a new paragraph.
                             %
with the right spacing around the periods requires
a special command.

\LaTeX\ interprets some common characters as
commands, so you must type special commands to
generate them.  These characters include the
following:
       \$ \& \% \# \{ and \}.

In printing, text is usually emphasized with an
       \emph{italic}
type style.

\begin{em}
   A long segment of text can also be emphasized
   in this way.  Text within such a segment can be
   given \emph{additional} emphasis.
\end{em}

It is sometimes necessary to prevent \LaTeX\ from
breaking a line where it might otherwise do so.
This may be at a space, as between the ``Mr.''\ and
``Jones'' in
       ``Mr.~Jones'',        % ~ produces an unbreakable interword space.
or within a word---especially when the word is a
symbol like
       \mbox{\emph{itemnum}}
that makes little sense when hyphenated across
lines.

Footnotes\footnote{This is an example of a footnote.}
pose no problem.

\LaTeX\ is good at typesetting mathematical formulas
like
       \( x-3y + z = 7 \)
or
       \( a_{1} > x^{2n} + y^{2n} > x' \)
or
       \( \ip{A}{B} = \sum_{i} a_{i} b_{i} \).
The spaces you type in a formula are
ignored.  Remember that a letter like
       $x$                   % $ ... $  and  \( ... \)  are equivalent
is a formula when it denotes a mathematical
symbol, and it should be typed as one.

\section{Displayed Text}

Text is displayed by indenting it from the left
margin.  Quotations are commonly displayed.  There
are short quotations
\begin{quote}
   This is a short quotation.  It consists of a
   single paragraph of text.  See how it is formatted.
\end{quote}
and longer ones.
\begin{quotation}
   This is a longer quotation.  It consists of two
   paragraphs of text, neither of which are
   particularly interesting.

   This is the second paragraph of the quotation.  It
   is just as dull as the first paragraph.
\end{quotation}
Another frequently-displayed structure is a list.
The following is an example of an \emph{itemized}
list.
\begin{itemize}
   \item This is the first item of an itemized list.
         Each item in the list is marked with a ``tick''.
         You don't have to worry about what kind of tick
         mark is used.

   \item This is the second item of the list.  It
         contains another list nested inside it.  The inner
         list is an \emph{enumerated} list.
         \begin{enumerate}
            \item This is the first item of an enumerated
                  list that is nested within the itemized list.

            \item This is the second item of the inner list.
                  \LaTeX\ allows you to nest lists deeper than
                  you really should.
         \end{enumerate}
         This is the rest of the second item of the outer
         list.  It is no more interesting than any other
         part of the item.
   \item This is the third item of the list.
\end{itemize}
You can even display poetry.
\begin{verse}
   There is an environment
    for verse \\             % The \\ command separates lines
   Whose features some poets % within a stanza.
   will curse.

                             % One or more blank lines separate stanzas.

   For instead of making\\
   Them do \emph{all} line breaking, \\
   It allows them to put too many words on a line when they'd rather be
   forced to be terse.
\end{verse}

Mathematical formulas may also be displayed.  A
displayed formula
is
one-line long; multiline
formulas require special formatting instructions.
   \[  \ip{\Gamma}{\psi'} = x'' + y^{2} + z_{i}^{n}\]
Don't start a paragraph with a displayed equation,
nor make one a paragraph by itself.

\end{document}               % End of document.

\section{Bench Section 202}
% synthetic bench section 202
% This is a sample LaTeX input file.  (Version of 12 August 2004.)
%
% A '%' character causes TeX to ignore all remaining text on the line,
% and is used for comments like this one.

\documentclass{article}      % Specifies the document class

                             % The preamble begins here.
\title{An Example Document}  % Declares the document's title.
\author{Leslie Lamport}      % Declares the author's name.
\date{January 21, 1994}      % Deleting this command produces today's date.

\newcommand{\ip}[2]{(#1, #2)}
                             % Defines \ip{arg1}{arg2} to mean
                             % (arg1, arg2).

%\newcommand{\ip}[2]{\langle #1 | #2\rangle}
                             % This is an alternative definition of
                             % \ip that is commented out.

\begin{document}             % End of preamble and beginning of text.

\maketitle                   % Produces the title.

This is an example input file.  Comparing it with
the output it generates can show you how to
produce a simple document of your own.

\section{Ordinary Text}      % Produces section heading.  Lower-level
                             % sections are begun with similar
                             % \subsection and \subsubsection commands.

The ends  of words and sentences are marked
  by   spaces. It  doesn't matter how many
spaces    you type; one is as good as 100.  The
end of   a line counts as a space.

One   or more   blank lines denote the  end
of  a paragraph.

Since any number of consecutive spaces are treated
like a single one, the formatting of the input
file makes no difference to
      \LaTeX,                % The \LaTeX command generates the LaTeX logo.
but it makes a difference to you.  When you use
\LaTeX, making your input file as easy to read
as possible will be a great help as you write
your document and when you change it.  This sample
file shows how you can add comments to your own input
file.

Because printing is different from typewriting,
there are a number of things that you have to do
differently when preparing an input file than if
you were just typing the document directly.
Quotation marks like
       ``this''
have to be handled specially, as do quotes within
quotes:
       ``\,`this'            % \, separates the double and single quote.
        is what I just
        wrote, not  `that'\,''.

Dashes come in three sizes: an
       intra-word
dash, a medium dash for number ranges like
       1--2,
and a punctuation
       dash---like
this.

A sentence-ending space should be larger than the
space between words within a sentence.  You
sometimes have to type special commands in
conjunction with punctuation characters to get
this right, as in the following sentence.
       Gnats, gnus, etc.\ all  % `\ ' makes an inter-word space.
       begin with G\@.         % \@ marks end-of-sentence punctuation.
You should check the spaces after periods when
reading your output to make sure you haven't
forgotten any special cases.  Generating an
ellipsis
       \ldots\               % `\ ' is needed after `\ldots' because TeX
                             % ignores spaces after command names like \ldots
                             % made from \ + letters.
                             %
                             % Note how a `%' character causes TeX to ignore
                             % the end of the input line, so these blank lines
                             % do not start a new paragraph.
                             %
with the right spacing around the periods requires
a special command.

\LaTeX\ interprets some common characters as
commands, so you must type special commands to
generate them.  These characters include the
following:
       \$ \& \% \# \{ and \}.

In printing, text is usually emphasized with an
       \emph{italic}
type style.

\begin{em}
   A long segment of text can also be emphasized
   in this way.  Text within such a segment can be
   given \emph{additional} emphasis.
\end{em}

It is sometimes necessary to prevent \LaTeX\ from
breaking a line where it might otherwise do so.
This may be at a space, as between the ``Mr.''\ and
``Jones'' in
       ``Mr.~Jones'',        % ~ produces an unbreakable interword space.
or within a word---especially when the word is a
symbol like
       \mbox{\emph{itemnum}}
that makes little sense when hyphenated across
lines.

Footnotes\footnote{This is an example of a footnote.}
pose no problem.

\LaTeX\ is good at typesetting mathematical formulas
like
       \( x-3y + z = 7 \)
or
       \( a_{1} > x^{2n} + y^{2n} > x' \)
or
       \( \ip{A}{B} = \sum_{i} a_{i} b_{i} \).
The spaces you type in a formula are
ignored.  Remember that a letter like
       $x$                   % $ ... $  and  \( ... \)  are equivalent
is a formula when it denotes a mathematical
symbol, and it should be typed as one.

\section{Displayed Text}

Text is displayed by indenting it from the left
margin.  Quotations are commonly displayed.  There
are short quotations
\begin{quote}
   This is a short quotation.  It consists of a
   single paragraph of text.  See how it is formatted.
\end{quote}
and longer ones.
\begin{quotation}
   This is a longer quotation.  It consists of two
   paragraphs of text, neither of which are
   particularly interesting.

   This is the second paragraph of the quotation.  It
   is just as dull as the first paragraph.
\end{quotation}
Another frequently-displayed structure is a list.
The following is an example of an \emph{itemized}
list.
\begin{itemize}
   \item This is the first item of an itemized list.
         Each item in the list is marked with a ``tick''.
         You don't have to worry about what kind of tick
         mark is used.

   \item This is the second item of the list.  It
         contains another list nested inside it.  The inner
         list is an \emph{enumerated} list.
         \begin{enumerate}
            \item This is the first item of an enumerated
                  list that is nested within the itemized list.

            \item This is the second item of the inner list.
                  \LaTeX\ allows you to nest lists deeper than
                  you really should.
         \end{enumerate}
         This is the rest of the second item of the outer
         list.  It is no more interesting than any other
         part of the item.
   \item This is the third item of the list.
\end{itemize}
You can even display poetry.
\begin{verse}
   There is an environment
    for verse \\             % The \\ command separates lines
   Whose features some poets % within a stanza.
   will curse.

                             % One or more blank lines separate stanzas.

   For instead of making\\
   Them do \emph{all} line breaking, \\
   It allows them to put too many words on a line when they'd rather be
   forced to be terse.
\end{verse}

Mathematical formulas may also be displayed.  A
displayed formula
is
one-line long; multiline
formulas require special formatting instructions.
   \[  \ip{\Gamma}{\psi'} = x'' + y^{2} + z_{i}^{n}\]
Don't start a paragraph with a displayed equation,
nor make one a paragraph by itself.

\end{document}               % End of document.

\section{Bench Section 203}
% synthetic bench section 203
% This is a sample LaTeX input file.  (Version of 12 August 2004.)
%
% A '%' character causes TeX to ignore all remaining text on the line,
% and is used for comments like this one.

\documentclass{article}      % Specifies the document class

                             % The preamble begins here.
\title{An Example Document}  % Declares the document's title.
\author{Leslie Lamport}      % Declares the author's name.
\date{January 21, 1994}      % Deleting this command produces today's date.

\newcommand{\ip}[2]{(#1, #2)}
                             % Defines \ip{arg1}{arg2} to mean
                             % (arg1, arg2).

%\newcommand{\ip}[2]{\langle #1 | #2\rangle}
                             % This is an alternative definition of
                             % \ip that is commented out.

\begin{document}             % End of preamble and beginning of text.

\maketitle                   % Produces the title.

This is an example input file.  Comparing it with
the output it generates can show you how to
produce a simple document of your own.

\section{Ordinary Text}      % Produces section heading.  Lower-level
                             % sections are begun with similar
                             % \subsection and \subsubsection commands.

The ends  of words and sentences are marked
  by   spaces. It  doesn't matter how many
spaces    you type; one is as good as 100.  The
end of   a line counts as a space.

One   or more   blank lines denote the  end
of  a paragraph.

Since any number of consecutive spaces are treated
like a single one, the formatting of the input
file makes no difference to
      \LaTeX,                % The \LaTeX command generates the LaTeX logo.
but it makes a difference to you.  When you use
\LaTeX, making your input file as easy to read
as possible will be a great help as you write
your document and when you change it.  This sample
file shows how you can add comments to your own input
file.

Because printing is different from typewriting,
there are a number of things that you have to do
differently when preparing an input file than if
you were just typing the document directly.
Quotation marks like
       ``this''
have to be handled specially, as do quotes within
quotes:
       ``\,`this'            % \, separates the double and single quote.
        is what I just
        wrote, not  `that'\,''.

Dashes come in three sizes: an
       intra-word
dash, a medium dash for number ranges like
       1--2,
and a punctuation
       dash---like
this.

A sentence-ending space should be larger than the
space between words within a sentence.  You
sometimes have to type special commands in
conjunction with punctuation characters to get
this right, as in the following sentence.
       Gnats, gnus, etc.\ all  % `\ ' makes an inter-word space.
       begin with G\@.         % \@ marks end-of-sentence punctuation.
You should check the spaces after periods when
reading your output to make sure you haven't
forgotten any special cases.  Generating an
ellipsis
       \ldots\               % `\ ' is needed after `\ldots' because TeX
                             % ignores spaces after command names like \ldots
                             % made from \ + letters.
                             %
                             % Note how a `%' character causes TeX to ignore
                             % the end of the input line, so these blank lines
                             % do not start a new paragraph.
                             %
with the right spacing around the periods requires
a special command.

\LaTeX\ interprets some common characters as
commands, so you must type special commands to
generate them.  These characters include the
following:
       \$ \& \% \# \{ and \}.

In printing, text is usually emphasized with an
       \emph{italic}
type style.

\begin{em}
   A long segment of text can also be emphasized
   in this way.  Text within such a segment can be
   given \emph{additional} emphasis.
\end{em}

It is sometimes necessary to prevent \LaTeX\ from
breaking a line where it might otherwise do so.
This may be at a space, as between the ``Mr.''\ and
``Jones'' in
       ``Mr.~Jones'',        % ~ produces an unbreakable interword space.
or within a word---especially when the word is a
symbol like
       \mbox{\emph{itemnum}}
that makes little sense when hyphenated across
lines.

Footnotes\footnote{This is an example of a footnote.}
pose no problem.

\LaTeX\ is good at typesetting mathematical formulas
like
       \( x-3y + z = 7 \)
or
       \( a_{1} > x^{2n} + y^{2n} > x' \)
or
       \( \ip{A}{B} = \sum_{i} a_{i} b_{i} \).
The spaces you type in a formula are
ignored.  Remember that a letter like
       $x$                   % $ ... $  and  \( ... \)  are equivalent
is a formula when it denotes a mathematical
symbol, and it should be typed as one.

\section{Displayed Text}

Text is displayed by indenting it from the left
margin.  Quotations are commonly displayed.  There
are short quotations
\begin{quote}
   This is a short quotation.  It consists of a
   single paragraph of text.  See how it is formatted.
\end{quote}
and longer ones.
\begin{quotation}
   This is a longer quotation.  It consists of two
   paragraphs of text, neither of which are
   particularly interesting.

   This is the second paragraph of the quotation.  It
   is just as dull as the first paragraph.
\end{quotation}
Another frequently-displayed structure is a list.
The following is an example of an \emph{itemized}
list.
\begin{itemize}
   \item This is the first item of an itemized list.
         Each item in the list is marked with a ``tick''.
         You don't have to worry about what kind of tick
         mark is used.

   \item This is the second item of the list.  It
         contains another list nested inside it.  The inner
         list is an \emph{enumerated} list.
         \begin{enumerate}
            \item This is the first item of an enumerated
                  list that is nested within the itemized list.

            \item This is the second item of the inner list.
                  \LaTeX\ allows you to nest lists deeper than
                  you really should.
         \end{enumerate}
         This is the rest of the second item of the outer
         list.  It is no more interesting than any other
         part of the item.
   \item This is the third item of the list.
\end{itemize}
You can even display poetry.
\begin{verse}
   There is an environment
    for verse \\             % The \\ command separates lines
   Whose features some poets % within a stanza.
   will curse.

                             % One or more blank lines separate stanzas.

   For instead of making\\
   Them do \emph{all} line breaking, \\
   It allows them to put too many words on a line when they'd rather be
   forced to be terse.
\end{verse}

Mathematical formulas may also be displayed.  A
displayed formula
is
one-line long; multiline
formulas require special formatting instructions.
   \[  \ip{\Gamma}{\psi'} = x'' + y^{2} + z_{i}^{n}\]
Don't start a paragraph with a displayed equation,
nor make one a paragraph by itself.

\end{document}               % End of document.

\section{Bench Section 204}
% synthetic bench section 204
% This is a sample LaTeX input file.  (Version of 12 August 2004.)
%
% A '%' character causes TeX to ignore all remaining text on the line,
% and is used for comments like this one.

\documentclass{article}      % Specifies the document class

                             % The preamble begins here.
\title{An Example Document}  % Declares the document's title.
\author{Leslie Lamport}      % Declares the author's name.
\date{January 21, 1994}      % Deleting this command produces today's date.

\newcommand{\ip}[2]{(#1, #2)}
                             % Defines \ip{arg1}{arg2} to mean
                             % (arg1, arg2).

%\newcommand{\ip}[2]{\langle #1 | #2\rangle}
                             % This is an alternative definition of
                             % \ip that is commented out.

\begin{document}             % End of preamble and beginning of text.

\maketitle                   % Produces the title.

This is an example input file.  Comparing it with
the output it generates can show you how to
produce a simple document of your own.

\section{Ordinary Text}      % Produces section heading.  Lower-level
                             % sections are begun with similar
                             % \subsection and \subsubsection commands.

The ends  of words and sentences are marked
  by   spaces. It  doesn't matter how many
spaces    you type; one is as good as 100.  The
end of   a line counts as a space.

One   or more   blank lines denote the  end
of  a paragraph.

Since any number of consecutive spaces are treated
like a single one, the formatting of the input
file makes no difference to
      \LaTeX,                % The \LaTeX command generates the LaTeX logo.
but it makes a difference to you.  When you use
\LaTeX, making your input file as easy to read
as possible will be a great help as you write
your document and when you change it.  This sample
file shows how you can add comments to your own input
file.

Because printing is different from typewriting,
there are a number of things that you have to do
differently when preparing an input file than if
you were just typing the document directly.
Quotation marks like
       ``this''
have to be handled specially, as do quotes within
quotes:
       ``\,`this'            % \, separates the double and single quote.
        is what I just
        wrote, not  `that'\,''.

Dashes come in three sizes: an
       intra-word
dash, a medium dash for number ranges like
       1--2,
and a punctuation
       dash---like
this.

A sentence-ending space should be larger than the
space between words within a sentence.  You
sometimes have to type special commands in
conjunction with punctuation characters to get
this right, as in the following sentence.
       Gnats, gnus, etc.\ all  % `\ ' makes an inter-word space.
       begin with G\@.         % \@ marks end-of-sentence punctuation.
You should check the spaces after periods when
reading your output to make sure you haven't
forgotten any special cases.  Generating an
ellipsis
       \ldots\               % `\ ' is needed after `\ldots' because TeX
                             % ignores spaces after command names like \ldots
                             % made from \ + letters.
                             %
                             % Note how a `%' character causes TeX to ignore
                             % the end of the input line, so these blank lines
                             % do not start a new paragraph.
                             %
with the right spacing around the periods requires
a special command.

\LaTeX\ interprets some common characters as
commands, so you must type special commands to
generate them.  These characters include the
following:
       \$ \& \% \# \{ and \}.

In printing, text is usually emphasized with an
       \emph{italic}
type style.

\begin{em}
   A long segment of text can also be emphasized
   in this way.  Text within such a segment can be
   given \emph{additional} emphasis.
\end{em}

It is sometimes necessary to prevent \LaTeX\ from
breaking a line where it might otherwise do so.
This may be at a space, as between the ``Mr.''\ and
``Jones'' in
       ``Mr.~Jones'',        % ~ produces an unbreakable interword space.
or within a word---especially when the word is a
symbol like
       \mbox{\emph{itemnum}}
that makes little sense when hyphenated across
lines.

Footnotes\footnote{This is an example of a footnote.}
pose no problem.

\LaTeX\ is good at typesetting mathematical formulas
like
       \( x-3y + z = 7 \)
or
       \( a_{1} > x^{2n} + y^{2n} > x' \)
or
       \( \ip{A}{B} = \sum_{i} a_{i} b_{i} \).
The spaces you type in a formula are
ignored.  Remember that a letter like
       $x$                   % $ ... $  and  \( ... \)  are equivalent
is a formula when it denotes a mathematical
symbol, and it should be typed as one.

\section{Displayed Text}

Text is displayed by indenting it from the left
margin.  Quotations are commonly displayed.  There
are short quotations
\begin{quote}
   This is a short quotation.  It consists of a
   single paragraph of text.  See how it is formatted.
\end{quote}
and longer ones.
\begin{quotation}
   This is a longer quotation.  It consists of two
   paragraphs of text, neither of which are
   particularly interesting.

   This is the second paragraph of the quotation.  It
   is just as dull as the first paragraph.
\end{quotation}
Another frequently-displayed structure is a list.
The following is an example of an \emph{itemized}
list.
\begin{itemize}
   \item This is the first item of an itemized list.
         Each item in the list is marked with a ``tick''.
         You don't have to worry about what kind of tick
         mark is used.

   \item This is the second item of the list.  It
         contains another list nested inside it.  The inner
         list is an \emph{enumerated} list.
         \begin{enumerate}
            \item This is the first item of an enumerated
                  list that is nested within the itemized list.

            \item This is the second item of the inner list.
                  \LaTeX\ allows you to nest lists deeper than
                  you really should.
         \end{enumerate}
         This is the rest of the second item of the outer
         list.  It is no more interesting than any other
         part of the item.
   \item This is the third item of the list.
\end{itemize}
You can even display poetry.
\begin{verse}
   There is an environment
    for verse \\             % The \\ command separates lines
   Whose features some poets % within a stanza.
   will curse.

                             % One or more blank lines separate stanzas.

   For instead of making\\
   Them do \emph{all} line breaking, \\
   It allows them to put too many words on a line when they'd rather be
   forced to be terse.
\end{verse}

Mathematical formulas may also be displayed.  A
displayed formula
is
one-line long; multiline
formulas require special formatting instructions.
   \[  \ip{\Gamma}{\psi'} = x'' + y^{2} + z_{i}^{n}\]
Don't start a paragraph with a displayed equation,
nor make one a paragraph by itself.

\end{document}               % End of document.

\section{Bench Section 205}
% synthetic bench section 205
% This is a sample LaTeX input file.  (Version of 12 August 2004.)
%
% A '%' character causes TeX to ignore all remaining text on the line,
% and is used for comments like this one.

\documentclass{article}      % Specifies the document class

                             % The preamble begins here.
\title{An Example Document}  % Declares the document's title.
\author{Leslie Lamport}      % Declares the author's name.
\date{January 21, 1994}      % Deleting this command produces today's date.

\newcommand{\ip}[2]{(#1, #2)}
                             % Defines \ip{arg1}{arg2} to mean
                             % (arg1, arg2).

%\newcommand{\ip}[2]{\langle #1 | #2\rangle}
                             % This is an alternative definition of
                             % \ip that is commented out.

\begin{document}             % End of preamble and beginning of text.

\maketitle                   % Produces the title.

This is an example input file.  Comparing it with
the output it generates can show you how to
produce a simple document of your own.

\section{Ordinary Text}      % Produces section heading.  Lower-level
                             % sections are begun with similar
                             % \subsection and \subsubsection commands.

The ends  of words and sentences are marked
  by   spaces. It  doesn't matter how many
spaces    you type; one is as good as 100.  The
end of   a line counts as a space.

One   or more   blank lines denote the  end
of  a paragraph.

Since any number of consecutive spaces are treated
like a single one, the formatting of the input
file makes no difference to
      \LaTeX,                % The \LaTeX command generates the LaTeX logo.
but it makes a difference to you.  When you use
\LaTeX, making your input file as easy to read
as possible will be a great help as you write
your document and when you change it.  This sample
file shows how you can add comments to your own input
file.

Because printing is different from typewriting,
there are a number of things that you have to do
differently when preparing an input file than if
you were just typing the document directly.
Quotation marks like
       ``this''
have to be handled specially, as do quotes within
quotes:
       ``\,`this'            % \, separates the double and single quote.
        is what I just
        wrote, not  `that'\,''.

Dashes come in three sizes: an
       intra-word
dash, a medium dash for number ranges like
       1--2,
and a punctuation
       dash---like
this.

A sentence-ending space should be larger than the
space between words within a sentence.  You
sometimes have to type special commands in
conjunction with punctuation characters to get
this right, as in the following sentence.
       Gnats, gnus, etc.\ all  % `\ ' makes an inter-word space.
       begin with G\@.         % \@ marks end-of-sentence punctuation.
You should check the spaces after periods when
reading your output to make sure you haven't
forgotten any special cases.  Generating an
ellipsis
       \ldots\               % `\ ' is needed after `\ldots' because TeX
                             % ignores spaces after command names like \ldots
                             % made from \ + letters.
                             %
                             % Note how a `%' character causes TeX to ignore
                             % the end of the input line, so these blank lines
                             % do not start a new paragraph.
                             %
with the right spacing around the periods requires
a special command.

\LaTeX\ interprets some common characters as
commands, so you must type special commands to
generate them.  These characters include the
following:
       \$ \& \% \# \{ and \}.

In printing, text is usually emphasized with an
       \emph{italic}
type style.

\begin{em}
   A long segment of text can also be emphasized
   in this way.  Text within such a segment can be
   given \emph{additional} emphasis.
\end{em}

It is sometimes necessary to prevent \LaTeX\ from
breaking a line where it might otherwise do so.
This may be at a space, as between the ``Mr.''\ and
``Jones'' in
       ``Mr.~Jones'',        % ~ produces an unbreakable interword space.
or within a word---especially when the word is a
symbol like
       \mbox{\emph{itemnum}}
that makes little sense when hyphenated across
lines.

Footnotes\footnote{This is an example of a footnote.}
pose no problem.

\LaTeX\ is good at typesetting mathematical formulas
like
       \( x-3y + z = 7 \)
or
       \( a_{1} > x^{2n} + y^{2n} > x' \)
or
       \( \ip{A}{B} = \sum_{i} a_{i} b_{i} \).
The spaces you type in a formula are
ignored.  Remember that a letter like
       $x$                   % $ ... $  and  \( ... \)  are equivalent
is a formula when it denotes a mathematical
symbol, and it should be typed as one.

\section{Displayed Text}

Text is displayed by indenting it from the left
margin.  Quotations are commonly displayed.  There
are short quotations
\begin{quote}
   This is a short quotation.  It consists of a
   single paragraph of text.  See how it is formatted.
\end{quote}
and longer ones.
\begin{quotation}
   This is a longer quotation.  It consists of two
   paragraphs of text, neither of which are
   particularly interesting.

   This is the second paragraph of the quotation.  It
   is just as dull as the first paragraph.
\end{quotation}
Another frequently-displayed structure is a list.
The following is an example of an \emph{itemized}
list.
\begin{itemize}
   \item This is the first item of an itemized list.
         Each item in the list is marked with a ``tick''.
         You don't have to worry about what kind of tick
         mark is used.

   \item This is the second item of the list.  It
         contains another list nested inside it.  The inner
         list is an \emph{enumerated} list.
         \begin{enumerate}
            \item This is the first item of an enumerated
                  list that is nested within the itemized list.

            \item This is the second item of the inner list.
                  \LaTeX\ allows you to nest lists deeper than
                  you really should.
         \end{enumerate}
         This is the rest of the second item of the outer
         list.  It is no more interesting than any other
         part of the item.
   \item This is the third item of the list.
\end{itemize}
You can even display poetry.
\begin{verse}
   There is an environment
    for verse \\             % The \\ command separates lines
   Whose features some poets % within a stanza.
   will curse.

                             % One or more blank lines separate stanzas.

   For instead of making\\
   Them do \emph{all} line breaking, \\
   It allows them to put too many words on a line when they'd rather be
   forced to be terse.
\end{verse}

Mathematical formulas may also be displayed.  A
displayed formula
is
one-line long; multiline
formulas require special formatting instructions.
   \[  \ip{\Gamma}{\psi'} = x'' + y^{2} + z_{i}^{n}\]
Don't start a paragraph with a displayed equation,
nor make one a paragraph by itself.

\end{document}               % End of document.

\section{Bench Section 206}
% synthetic bench section 206
% This is a sample LaTeX input file.  (Version of 12 August 2004.)
%
% A '%' character causes TeX to ignore all remaining text on the line,
% and is used for comments like this one.

\documentclass{article}      % Specifies the document class

                             % The preamble begins here.
\title{An Example Document}  % Declares the document's title.
\author{Leslie Lamport}      % Declares the author's name.
\date{January 21, 1994}      % Deleting this command produces today's date.

\newcommand{\ip}[2]{(#1, #2)}
                             % Defines \ip{arg1}{arg2} to mean
                             % (arg1, arg2).

%\newcommand{\ip}[2]{\langle #1 | #2\rangle}
                             % This is an alternative definition of
                             % \ip that is commented out.

\begin{document}             % End of preamble and beginning of text.

\maketitle                   % Produces the title.

This is an example input file.  Comparing it with
the output it generates can show you how to
produce a simple document of your own.

\section{Ordinary Text}      % Produces section heading.  Lower-level
                             % sections are begun with similar
                             % \subsection and \subsubsection commands.

The ends  of words and sentences are marked
  by   spaces. It  doesn't matter how many
spaces    you type; one is as good as 100.  The
end of   a line counts as a space.

One   or more   blank lines denote the  end
of  a paragraph.

Since any number of consecutive spaces are treated
like a single one, the formatting of the input
file makes no difference to
      \LaTeX,                % The \LaTeX command generates the LaTeX logo.
but it makes a difference to you.  When you use
\LaTeX, making your input file as easy to read
as possible will be a great help as you write
your document and when you change it.  This sample
file shows how you can add comments to your own input
file.

Because printing is different from typewriting,
there are a number of things that you have to do
differently when preparing an input file than if
you were just typing the document directly.
Quotation marks like
       ``this''
have to be handled specially, as do quotes within
quotes:
       ``\,`this'            % \, separates the double and single quote.
        is what I just
        wrote, not  `that'\,''.

Dashes come in three sizes: an
       intra-word
dash, a medium dash for number ranges like
       1--2,
and a punctuation
       dash---like
this.

A sentence-ending space should be larger than the
space between words within a sentence.  You
sometimes have to type special commands in
conjunction with punctuation characters to get
this right, as in the following sentence.
       Gnats, gnus, etc.\ all  % `\ ' makes an inter-word space.
       begin with G\@.         % \@ marks end-of-sentence punctuation.
You should check the spaces after periods when
reading your output to make sure you haven't
forgotten any special cases.  Generating an
ellipsis
       \ldots\               % `\ ' is needed after `\ldots' because TeX
                             % ignores spaces after command names like \ldots
                             % made from \ + letters.
                             %
                             % Note how a `%' character causes TeX to ignore
                             % the end of the input line, so these blank lines
                             % do not start a new paragraph.
                             %
with the right spacing around the periods requires
a special command.

\LaTeX\ interprets some common characters as
commands, so you must type special commands to
generate them.  These characters include the
following:
       \$ \& \% \# \{ and \}.

In printing, text is usually emphasized with an
       \emph{italic}
type style.

\begin{em}
   A long segment of text can also be emphasized
   in this way.  Text within such a segment can be
   given \emph{additional} emphasis.
\end{em}

It is sometimes necessary to prevent \LaTeX\ from
breaking a line where it might otherwise do so.
This may be at a space, as between the ``Mr.''\ and
``Jones'' in
       ``Mr.~Jones'',        % ~ produces an unbreakable interword space.
or within a word---especially when the word is a
symbol like
       \mbox{\emph{itemnum}}
that makes little sense when hyphenated across
lines.

Footnotes\footnote{This is an example of a footnote.}
pose no problem.

\LaTeX\ is good at typesetting mathematical formulas
like
       \( x-3y + z = 7 \)
or
       \( a_{1} > x^{2n} + y^{2n} > x' \)
or
       \( \ip{A}{B} = \sum_{i} a_{i} b_{i} \).
The spaces you type in a formula are
ignored.  Remember that a letter like
       $x$                   % $ ... $  and  \( ... \)  are equivalent
is a formula when it denotes a mathematical
symbol, and it should be typed as one.

\section{Displayed Text}

Text is displayed by indenting it from the left
margin.  Quotations are commonly displayed.  There
are short quotations
\begin{quote}
   This is a short quotation.  It consists of a
   single paragraph of text.  See how it is formatted.
\end{quote}
and longer ones.
\begin{quotation}
   This is a longer quotation.  It consists of two
   paragraphs of text, neither of which are
   particularly interesting.

   This is the second paragraph of the quotation.  It
   is just as dull as the first paragraph.
\end{quotation}
Another frequently-displayed structure is a list.
The following is an example of an \emph{itemized}
list.
\begin{itemize}
   \item This is the first item of an itemized list.
         Each item in the list is marked with a ``tick''.
         You don't have to worry about what kind of tick
         mark is used.

   \item This is the second item of the list.  It
         contains another list nested inside it.  The inner
         list is an \emph{enumerated} list.
         \begin{enumerate}
            \item This is the first item of an enumerated
                  list that is nested within the itemized list.

            \item This is the second item of the inner list.
                  \LaTeX\ allows you to nest lists deeper than
                  you really should.
         \end{enumerate}
         This is the rest of the second item of the outer
         list.  It is no more interesting than any other
         part of the item.
   \item This is the third item of the list.
\end{itemize}
You can even display poetry.
\begin{verse}
   There is an environment
    for verse \\             % The \\ command separates lines
   Whose features some poets % within a stanza.
   will curse.

                             % One or more blank lines separate stanzas.

   For instead of making\\
   Them do \emph{all} line breaking, \\
   It allows them to put too many words on a line when they'd rather be
   forced to be terse.
\end{verse}

Mathematical formulas may also be displayed.  A
displayed formula
is
one-line long; multiline
formulas require special formatting instructions.
   \[  \ip{\Gamma}{\psi'} = x'' + y^{2} + z_{i}^{n}\]
Don't start a paragraph with a displayed equation,
nor make one a paragraph by itself.

\end{document}               % End of document.

\section{Bench Section 207}
% synthetic bench section 207
% This is a sample LaTeX input file.  (Version of 12 August 2004.)
%
% A '%' character causes TeX to ignore all remaining text on the line,
% and is used for comments like this one.

\documentclass{article}      % Specifies the document class

                             % The preamble begins here.
\title{An Example Document}  % Declares the document's title.
\author{Leslie Lamport}      % Declares the author's name.
\date{January 21, 1994}      % Deleting this command produces today's date.

\newcommand{\ip}[2]{(#1, #2)}
                             % Defines \ip{arg1}{arg2} to mean
                             % (arg1, arg2).

%\newcommand{\ip}[2]{\langle #1 | #2\rangle}
                             % This is an alternative definition of
                             % \ip that is commented out.

\begin{document}             % End of preamble and beginning of text.

\maketitle                   % Produces the title.

This is an example input file.  Comparing it with
the output it generates can show you how to
produce a simple document of your own.

\section{Ordinary Text}      % Produces section heading.  Lower-level
                             % sections are begun with similar
                             % \subsection and \subsubsection commands.

The ends  of words and sentences are marked
  by   spaces. It  doesn't matter how many
spaces    you type; one is as good as 100.  The
end of   a line counts as a space.

One   or more   blank lines denote the  end
of  a paragraph.

Since any number of consecutive spaces are treated
like a single one, the formatting of the input
file makes no difference to
      \LaTeX,                % The \LaTeX command generates the LaTeX logo.
but it makes a difference to you.  When you use
\LaTeX, making your input file as easy to read
as possible will be a great help as you write
your document and when you change it.  This sample
file shows how you can add comments to your own input
file.

Because printing is different from typewriting,
there are a number of things that you have to do
differently when preparing an input file than if
you were just typing the document directly.
Quotation marks like
       ``this''
have to be handled specially, as do quotes within
quotes:
       ``\,`this'            % \, separates the double and single quote.
        is what I just
        wrote, not  `that'\,''.

Dashes come in three sizes: an
       intra-word
dash, a medium dash for number ranges like
       1--2,
and a punctuation
       dash---like
this.

A sentence-ending space should be larger than the
space between words within a sentence.  You
sometimes have to type special commands in
conjunction with punctuation characters to get
this right, as in the following sentence.
       Gnats, gnus, etc.\ all  % `\ ' makes an inter-word space.
       begin with G\@.         % \@ marks end-of-sentence punctuation.
You should check the spaces after periods when
reading your output to make sure you haven't
forgotten any special cases.  Generating an
ellipsis
       \ldots\               % `\ ' is needed after `\ldots' because TeX
                             % ignores spaces after command names like \ldots
                             % made from \ + letters.
                             %
                             % Note how a `%' character causes TeX to ignore
                             % the end of the input line, so these blank lines
                             % do not start a new paragraph.
                             %
with the right spacing around the periods requires
a special command.

\LaTeX\ interprets some common characters as
commands, so you must type special commands to
generate them.  These characters include the
following:
       \$ \& \% \# \{ and \}.

In printing, text is usually emphasized with an
       \emph{italic}
type style.

\begin{em}
   A long segment of text can also be emphasized
   in this way.  Text within such a segment can be
   given \emph{additional} emphasis.
\end{em}

It is sometimes necessary to prevent \LaTeX\ from
breaking a line where it might otherwise do so.
This may be at a space, as between the ``Mr.''\ and
``Jones'' in
       ``Mr.~Jones'',        % ~ produces an unbreakable interword space.
or within a word---especially when the word is a
symbol like
       \mbox{\emph{itemnum}}
that makes little sense when hyphenated across
lines.

Footnotes\footnote{This is an example of a footnote.}
pose no problem.

\LaTeX\ is good at typesetting mathematical formulas
like
       \( x-3y + z = 7 \)
or
       \( a_{1} > x^{2n} + y^{2n} > x' \)
or
       \( \ip{A}{B} = \sum_{i} a_{i} b_{i} \).
The spaces you type in a formula are
ignored.  Remember that a letter like
       $x$                   % $ ... $  and  \( ... \)  are equivalent
is a formula when it denotes a mathematical
symbol, and it should be typed as one.

\section{Displayed Text}

Text is displayed by indenting it from the left
margin.  Quotations are commonly displayed.  There
are short quotations
\begin{quote}
   This is a short quotation.  It consists of a
   single paragraph of text.  See how it is formatted.
\end{quote}
and longer ones.
\begin{quotation}
   This is a longer quotation.  It consists of two
   paragraphs of text, neither of which are
   particularly interesting.

   This is the second paragraph of the quotation.  It
   is just as dull as the first paragraph.
\end{quotation}
Another frequently-displayed structure is a list.
The following is an example of an \emph{itemized}
list.
\begin{itemize}
   \item This is the first item of an itemized list.
         Each item in the list is marked with a ``tick''.
         You don't have to worry about what kind of tick
         mark is used.

   \item This is the second item of the list.  It
         contains another list nested inside it.  The inner
         list is an \emph{enumerated} list.
         \begin{enumerate}
            \item This is the first item of an enumerated
                  list that is nested within the itemized list.

            \item This is the second item of the inner list.
                  \LaTeX\ allows you to nest lists deeper than
                  you really should.
         \end{enumerate}
         This is the rest of the second item of the outer
         list.  It is no more interesting than any other
         part of the item.
   \item This is the third item of the list.
\end{itemize}
You can even display poetry.
\begin{verse}
   There is an environment
    for verse \\             % The \\ command separates lines
   Whose features some poets % within a stanza.
   will curse.

                             % One or more blank lines separate stanzas.

   For instead of making\\
   Them do \emph{all} line breaking, \\
   It allows them to put too many words on a line when they'd rather be
   forced to be terse.
\end{verse}

Mathematical formulas may also be displayed.  A
displayed formula
is
one-line long; multiline
formulas require special formatting instructions.
   \[  \ip{\Gamma}{\psi'} = x'' + y^{2} + z_{i}^{n}\]
Don't start a paragraph with a displayed equation,
nor make one a paragraph by itself.

\end{document}               % End of document.

\section{Bench Section 208}
% synthetic bench section 208
% This is a sample LaTeX input file.  (Version of 12 August 2004.)
%
% A '%' character causes TeX to ignore all remaining text on the line,
% and is used for comments like this one.

\documentclass{article}      % Specifies the document class

                             % The preamble begins here.
\title{An Example Document}  % Declares the document's title.
\author{Leslie Lamport}      % Declares the author's name.
\date{January 21, 1994}      % Deleting this command produces today's date.

\newcommand{\ip}[2]{(#1, #2)}
                             % Defines \ip{arg1}{arg2} to mean
                             % (arg1, arg2).

%\newcommand{\ip}[2]{\langle #1 | #2\rangle}
                             % This is an alternative definition of
                             % \ip that is commented out.

\begin{document}             % End of preamble and beginning of text.

\maketitle                   % Produces the title.

This is an example input file.  Comparing it with
the output it generates can show you how to
produce a simple document of your own.

\section{Ordinary Text}      % Produces section heading.  Lower-level
                             % sections are begun with similar
                             % \subsection and \subsubsection commands.

The ends  of words and sentences are marked
  by   spaces. It  doesn't matter how many
spaces    you type; one is as good as 100.  The
end of   a line counts as a space.

One   or more   blank lines denote the  end
of  a paragraph.

Since any number of consecutive spaces are treated
like a single one, the formatting of the input
file makes no difference to
      \LaTeX,                % The \LaTeX command generates the LaTeX logo.
but it makes a difference to you.  When you use
\LaTeX, making your input file as easy to read
as possible will be a great help as you write
your document and when you change it.  This sample
file shows how you can add comments to your own input
file.

Because printing is different from typewriting,
there are a number of things that you have to do
differently when preparing an input file than if
you were just typing the document directly.
Quotation marks like
       ``this''
have to be handled specially, as do quotes within
quotes:
       ``\,`this'            % \, separates the double and single quote.
        is what I just
        wrote, not  `that'\,''.

Dashes come in three sizes: an
       intra-word
dash, a medium dash for number ranges like
       1--2,
and a punctuation
       dash---like
this.

A sentence-ending space should be larger than the
space between words within a sentence.  You
sometimes have to type special commands in
conjunction with punctuation characters to get
this right, as in the following sentence.
       Gnats, gnus, etc.\ all  % `\ ' makes an inter-word space.
       begin with G\@.         % \@ marks end-of-sentence punctuation.
You should check the spaces after periods when
reading your output to make sure you haven't
forgotten any special cases.  Generating an
ellipsis
       \ldots\               % `\ ' is needed after `\ldots' because TeX
                             % ignores spaces after command names like \ldots
                             % made from \ + letters.
                             %
                             % Note how a `%' character causes TeX to ignore
                             % the end of the input line, so these blank lines
                             % do not start a new paragraph.
                             %
with the right spacing around the periods requires
a special command.

\LaTeX\ interprets some common characters as
commands, so you must type special commands to
generate them.  These characters include the
following:
       \$ \& \% \# \{ and \}.

In printing, text is usually emphasized with an
       \emph{italic}
type style.

\begin{em}
   A long segment of text can also be emphasized
   in this way.  Text within such a segment can be
   given \emph{additional} emphasis.
\end{em}

It is sometimes necessary to prevent \LaTeX\ from
breaking a line where it might otherwise do so.
This may be at a space, as between the ``Mr.''\ and
``Jones'' in
       ``Mr.~Jones'',        % ~ produces an unbreakable interword space.
or within a word---especially when the word is a
symbol like
       \mbox{\emph{itemnum}}
that makes little sense when hyphenated across
lines.

Footnotes\footnote{This is an example of a footnote.}
pose no problem.

\LaTeX\ is good at typesetting mathematical formulas
like
       \( x-3y + z = 7 \)
or
       \( a_{1} > x^{2n} + y^{2n} > x' \)
or
       \( \ip{A}{B} = \sum_{i} a_{i} b_{i} \).
The spaces you type in a formula are
ignored.  Remember that a letter like
       $x$                   % $ ... $  and  \( ... \)  are equivalent
is a formula when it denotes a mathematical
symbol, and it should be typed as one.

\section{Displayed Text}

Text is displayed by indenting it from the left
margin.  Quotations are commonly displayed.  There
are short quotations
\begin{quote}
   This is a short quotation.  It consists of a
   single paragraph of text.  See how it is formatted.
\end{quote}
and longer ones.
\begin{quotation}
   This is a longer quotation.  It consists of two
   paragraphs of text, neither of which are
   particularly interesting.

   This is the second paragraph of the quotation.  It
   is just as dull as the first paragraph.
\end{quotation}
Another frequently-displayed structure is a list.
The following is an example of an \emph{itemized}
list.
\begin{itemize}
   \item This is the first item of an itemized list.
         Each item in the list is marked with a ``tick''.
         You don't have to worry about what kind of tick
         mark is used.

   \item This is the second item of the list.  It
         contains another list nested inside it.  The inner
         list is an \emph{enumerated} list.
         \begin{enumerate}
            \item This is the first item of an enumerated
                  list that is nested within the itemized list.

            \item This is the second item of the inner list.
                  \LaTeX\ allows you to nest lists deeper than
                  you really should.
         \end{enumerate}
         This is the rest of the second item of the outer
         list.  It is no more interesting than any other
         part of the item.
   \item This is the third item of the list.
\end{itemize}
You can even display poetry.
\begin{verse}
   There is an environment
    for verse \\             % The \\ command separates lines
   Whose features some poets % within a stanza.
   will curse.

                             % One or more blank lines separate stanzas.

   For instead of making\\
   Them do \emph{all} line breaking, \\
   It allows them to put too many words on a line when they'd rather be
   forced to be terse.
\end{verse}

Mathematical formulas may also be displayed.  A
displayed formula
is
one-line long; multiline
formulas require special formatting instructions.
   \[  \ip{\Gamma}{\psi'} = x'' + y^{2} + z_{i}^{n}\]
Don't start a paragraph with a displayed equation,
nor make one a paragraph by itself.

\end{document}               % End of document.

\section{Bench Section 209}
% synthetic bench section 209
% This is a sample LaTeX input file.  (Version of 12 August 2004.)
%
% A '%' character causes TeX to ignore all remaining text on the line,
% and is used for comments like this one.

\documentclass{article}      % Specifies the document class

                             % The preamble begins here.
\title{An Example Document}  % Declares the document's title.
\author{Leslie Lamport}      % Declares the author's name.
\date{January 21, 1994}      % Deleting this command produces today's date.

\newcommand{\ip}[2]{(#1, #2)}
                             % Defines \ip{arg1}{arg2} to mean
                             % (arg1, arg2).

%\newcommand{\ip}[2]{\langle #1 | #2\rangle}
                             % This is an alternative definition of
                             % \ip that is commented out.

\begin{document}             % End of preamble and beginning of text.

\maketitle                   % Produces the title.

This is an example input file.  Comparing it with
the output it generates can show you how to
produce a simple document of your own.

\section{Ordinary Text}      % Produces section heading.  Lower-level
                             % sections are begun with similar
                             % \subsection and \subsubsection commands.

The ends  of words and sentences are marked
  by   spaces. It  doesn't matter how many
spaces    you type; one is as good as 100.  The
end of   a line counts as a space.

One   or more   blank lines denote the  end
of  a paragraph.

Since any number of consecutive spaces are treated
like a single one, the formatting of the input
file makes no difference to
      \LaTeX,                % The \LaTeX command generates the LaTeX logo.
but it makes a difference to you.  When you use
\LaTeX, making your input file as easy to read
as possible will be a great help as you write
your document and when you change it.  This sample
file shows how you can add comments to your own input
file.

Because printing is different from typewriting,
there are a number of things that you have to do
differently when preparing an input file than if
you were just typing the document directly.
Quotation marks like
       ``this''
have to be handled specially, as do quotes within
quotes:
       ``\,`this'            % \, separates the double and single quote.
        is what I just
        wrote, not  `that'\,''.

Dashes come in three sizes: an
       intra-word
dash, a medium dash for number ranges like
       1--2,
and a punctuation
       dash---like
this.

A sentence-ending space should be larger than the
space between words within a sentence.  You
sometimes have to type special commands in
conjunction with punctuation characters to get
this right, as in the following sentence.
       Gnats, gnus, etc.\ all  % `\ ' makes an inter-word space.
       begin with G\@.         % \@ marks end-of-sentence punctuation.
You should check the spaces after periods when
reading your output to make sure you haven't
forgotten any special cases.  Generating an
ellipsis
       \ldots\               % `\ ' is needed after `\ldots' because TeX
                             % ignores spaces after command names like \ldots
                             % made from \ + letters.
                             %
                             % Note how a `%' character causes TeX to ignore
                             % the end of the input line, so these blank lines
                             % do not start a new paragraph.
                             %
with the right spacing around the periods requires
a special command.

\LaTeX\ interprets some common characters as
commands, so you must type special commands to
generate them.  These characters include the
following:
       \$ \& \% \# \{ and \}.

In printing, text is usually emphasized with an
       \emph{italic}
type style.

\begin{em}
   A long segment of text can also be emphasized
   in this way.  Text within such a segment can be
   given \emph{additional} emphasis.
\end{em}

It is sometimes necessary to prevent \LaTeX\ from
breaking a line where it might otherwise do so.
This may be at a space, as between the ``Mr.''\ and
``Jones'' in
       ``Mr.~Jones'',        % ~ produces an unbreakable interword space.
or within a word---especially when the word is a
symbol like
       \mbox{\emph{itemnum}}
that makes little sense when hyphenated across
lines.

Footnotes\footnote{This is an example of a footnote.}
pose no problem.

\LaTeX\ is good at typesetting mathematical formulas
like
       \( x-3y + z = 7 \)
or
       \( a_{1} > x^{2n} + y^{2n} > x' \)
or
       \( \ip{A}{B} = \sum_{i} a_{i} b_{i} \).
The spaces you type in a formula are
ignored.  Remember that a letter like
       $x$                   % $ ... $  and  \( ... \)  are equivalent
is a formula when it denotes a mathematical
symbol, and it should be typed as one.

\section{Displayed Text}

Text is displayed by indenting it from the left
margin.  Quotations are commonly displayed.  There
are short quotations
\begin{quote}
   This is a short quotation.  It consists of a
   single paragraph of text.  See how it is formatted.
\end{quote}
and longer ones.
\begin{quotation}
   This is a longer quotation.  It consists of two
   paragraphs of text, neither of which are
   particularly interesting.

   This is the second paragraph of the quotation.  It
   is just as dull as the first paragraph.
\end{quotation}
Another frequently-displayed structure is a list.
The following is an example of an \emph{itemized}
list.
\begin{itemize}
   \item This is the first item of an itemized list.
         Each item in the list is marked with a ``tick''.
         You don't have to worry about what kind of tick
         mark is used.

   \item This is the second item of the list.  It
         contains another list nested inside it.  The inner
         list is an \emph{enumerated} list.
         \begin{enumerate}
            \item This is the first item of an enumerated
                  list that is nested within the itemized list.

            \item This is the second item of the inner list.
                  \LaTeX\ allows you to nest lists deeper than
                  you really should.
         \end{enumerate}
         This is the rest of the second item of the outer
         list.  It is no more interesting than any other
         part of the item.
   \item This is the third item of the list.
\end{itemize}
You can even display poetry.
\begin{verse}
   There is an environment
    for verse \\             % The \\ command separates lines
   Whose features some poets % within a stanza.
   will curse.

                             % One or more blank lines separate stanzas.

   For instead of making\\
   Them do \emph{all} line breaking, \\
   It allows them to put too many words on a line when they'd rather be
   forced to be terse.
\end{verse}

Mathematical formulas may also be displayed.  A
displayed formula
is
one-line long; multiline
formulas require special formatting instructions.
   \[  \ip{\Gamma}{\psi'} = x'' + y^{2} + z_{i}^{n}\]
Don't start a paragraph with a displayed equation,
nor make one a paragraph by itself.

\end{document}               % End of document.

\section{Bench Section 210}
% synthetic bench section 210
% This is a sample LaTeX input file.  (Version of 12 August 2004.)
%
% A '%' character causes TeX to ignore all remaining text on the line,
% and is used for comments like this one.

\documentclass{article}      % Specifies the document class

                             % The preamble begins here.
\title{An Example Document}  % Declares the document's title.
\author{Leslie Lamport}      % Declares the author's name.
\date{January 21, 1994}      % Deleting this command produces today's date.

\newcommand{\ip}[2]{(#1, #2)}
                             % Defines \ip{arg1}{arg2} to mean
                             % (arg1, arg2).

%\newcommand{\ip}[2]{\langle #1 | #2\rangle}
                             % This is an alternative definition of
                             % \ip that is commented out.

\begin{document}             % End of preamble and beginning of text.

\maketitle                   % Produces the title.

This is an example input file.  Comparing it with
the output it generates can show you how to
produce a simple document of your own.

\section{Ordinary Text}      % Produces section heading.  Lower-level
                             % sections are begun with similar
                             % \subsection and \subsubsection commands.

The ends  of words and sentences are marked
  by   spaces. It  doesn't matter how many
spaces    you type; one is as good as 100.  The
end of   a line counts as a space.

One   or more   blank lines denote the  end
of  a paragraph.

Since any number of consecutive spaces are treated
like a single one, the formatting of the input
file makes no difference to
      \LaTeX,                % The \LaTeX command generates the LaTeX logo.
but it makes a difference to you.  When you use
\LaTeX, making your input file as easy to read
as possible will be a great help as you write
your document and when you change it.  This sample
file shows how you can add comments to your own input
file.

Because printing is different from typewriting,
there are a number of things that you have to do
differently when preparing an input file than if
you were just typing the document directly.
Quotation marks like
       ``this''
have to be handled specially, as do quotes within
quotes:
       ``\,`this'            % \, separates the double and single quote.
        is what I just
        wrote, not  `that'\,''.

Dashes come in three sizes: an
       intra-word
dash, a medium dash for number ranges like
       1--2,
and a punctuation
       dash---like
this.

A sentence-ending space should be larger than the
space between words within a sentence.  You
sometimes have to type special commands in
conjunction with punctuation characters to get
this right, as in the following sentence.
       Gnats, gnus, etc.\ all  % `\ ' makes an inter-word space.
       begin with G\@.         % \@ marks end-of-sentence punctuation.
You should check the spaces after periods when
reading your output to make sure you haven't
forgotten any special cases.  Generating an
ellipsis
       \ldots\               % `\ ' is needed after `\ldots' because TeX
                             % ignores spaces after command names like \ldots
                             % made from \ + letters.
                             %
                             % Note how a `%' character causes TeX to ignore
                             % the end of the input line, so these blank lines
                             % do not start a new paragraph.
                             %
with the right spacing around the periods requires
a special command.

\LaTeX\ interprets some common characters as
commands, so you must type special commands to
generate them.  These characters include the
following:
       \$ \& \% \# \{ and \}.

In printing, text is usually emphasized with an
       \emph{italic}
type style.

\begin{em}
   A long segment of text can also be emphasized
   in this way.  Text within such a segment can be
   given \emph{additional} emphasis.
\end{em}

It is sometimes necessary to prevent \LaTeX\ from
breaking a line where it might otherwise do so.
This may be at a space, as between the ``Mr.''\ and
``Jones'' in
       ``Mr.~Jones'',        % ~ produces an unbreakable interword space.
or within a word---especially when the word is a
symbol like
       \mbox{\emph{itemnum}}
that makes little sense when hyphenated across
lines.

Footnotes\footnote{This is an example of a footnote.}
pose no problem.

\LaTeX\ is good at typesetting mathematical formulas
like
       \( x-3y + z = 7 \)
or
       \( a_{1} > x^{2n} + y^{2n} > x' \)
or
       \( \ip{A}{B} = \sum_{i} a_{i} b_{i} \).
The spaces you type in a formula are
ignored.  Remember that a letter like
       $x$                   % $ ... $  and  \( ... \)  are equivalent
is a formula when it denotes a mathematical
symbol, and it should be typed as one.

\section{Displayed Text}

Text is displayed by indenting it from the left
margin.  Quotations are commonly displayed.  There
are short quotations
\begin{quote}
   This is a short quotation.  It consists of a
   single paragraph of text.  See how it is formatted.
\end{quote}
and longer ones.
\begin{quotation}
   This is a longer quotation.  It consists of two
   paragraphs of text, neither of which are
   particularly interesting.

   This is the second paragraph of the quotation.  It
   is just as dull as the first paragraph.
\end{quotation}
Another frequently-displayed structure is a list.
The following is an example of an \emph{itemized}
list.
\begin{itemize}
   \item This is the first item of an itemized list.
         Each item in the list is marked with a ``tick''.
         You don't have to worry about what kind of tick
         mark is used.

   \item This is the second item of the list.  It
         contains another list nested inside it.  The inner
         list is an \emph{enumerated} list.
         \begin{enumerate}
            \item This is the first item of an enumerated
                  list that is nested within the itemized list.

            \item This is the second item of the inner list.
                  \LaTeX\ allows you to nest lists deeper than
                  you really should.
         \end{enumerate}
         This is the rest of the second item of the outer
         list.  It is no more interesting than any other
         part of the item.
   \item This is the third item of the list.
\end{itemize}
You can even display poetry.
\begin{verse}
   There is an environment
    for verse \\             % The \\ command separates lines
   Whose features some poets % within a stanza.
   will curse.

                             % One or more blank lines separate stanzas.

   For instead of making\\
   Them do \emph{all} line breaking, \\
   It allows them to put too many words on a line when they'd rather be
   forced to be terse.
\end{verse}

Mathematical formulas may also be displayed.  A
displayed formula
is
one-line long; multiline
formulas require special formatting instructions.
   \[  \ip{\Gamma}{\psi'} = x'' + y^{2} + z_{i}^{n}\]
Don't start a paragraph with a displayed equation,
nor make one a paragraph by itself.

\end{document}               % End of document.

\section{Bench Section 211}
% synthetic bench section 211
% This is a sample LaTeX input file.  (Version of 12 August 2004.)
%
% A '%' character causes TeX to ignore all remaining text on the line,
% and is used for comments like this one.

\documentclass{article}      % Specifies the document class

                             % The preamble begins here.
\title{An Example Document}  % Declares the document's title.
\author{Leslie Lamport}      % Declares the author's name.
\date{January 21, 1994}      % Deleting this command produces today's date.

\newcommand{\ip}[2]{(#1, #2)}
                             % Defines \ip{arg1}{arg2} to mean
                             % (arg1, arg2).

%\newcommand{\ip}[2]{\langle #1 | #2\rangle}
                             % This is an alternative definition of
                             % \ip that is commented out.

\begin{document}             % End of preamble and beginning of text.

\maketitle                   % Produces the title.

This is an example input file.  Comparing it with
the output it generates can show you how to
produce a simple document of your own.

\section{Ordinary Text}      % Produces section heading.  Lower-level
                             % sections are begun with similar
                             % \subsection and \subsubsection commands.

The ends  of words and sentences are marked
  by   spaces. It  doesn't matter how many
spaces    you type; one is as good as 100.  The
end of   a line counts as a space.

One   or more   blank lines denote the  end
of  a paragraph.

Since any number of consecutive spaces are treated
like a single one, the formatting of the input
file makes no difference to
      \LaTeX,                % The \LaTeX command generates the LaTeX logo.
but it makes a difference to you.  When you use
\LaTeX, making your input file as easy to read
as possible will be a great help as you write
your document and when you change it.  This sample
file shows how you can add comments to your own input
file.

Because printing is different from typewriting,
there are a number of things that you have to do
differently when preparing an input file than if
you were just typing the document directly.
Quotation marks like
       ``this''
have to be handled specially, as do quotes within
quotes:
       ``\,`this'            % \, separates the double and single quote.
        is what I just
        wrote, not  `that'\,''.

Dashes come in three sizes: an
       intra-word
dash, a medium dash for number ranges like
       1--2,
and a punctuation
       dash---like
this.

A sentence-ending space should be larger than the
space between words within a sentence.  You
sometimes have to type special commands in
conjunction with punctuation characters to get
this right, as in the following sentence.
       Gnats, gnus, etc.\ all  % `\ ' makes an inter-word space.
       begin with G\@.         % \@ marks end-of-sentence punctuation.
You should check the spaces after periods when
reading your output to make sure you haven't
forgotten any special cases.  Generating an
ellipsis
       \ldots\               % `\ ' is needed after `\ldots' because TeX
                             % ignores spaces after command names like \ldots
                             % made from \ + letters.
                             %
                             % Note how a `%' character causes TeX to ignore
                             % the end of the input line, so these blank lines
                             % do not start a new paragraph.
                             %
with the right spacing around the periods requires
a special command.

\LaTeX\ interprets some common characters as
commands, so you must type special commands to
generate them.  These characters include the
following:
       \$ \& \% \# \{ and \}.

In printing, text is usually emphasized with an
       \emph{italic}
type style.

\begin{em}
   A long segment of text can also be emphasized
   in this way.  Text within such a segment can be
   given \emph{additional} emphasis.
\end{em}

It is sometimes necessary to prevent \LaTeX\ from
breaking a line where it might otherwise do so.
This may be at a space, as between the ``Mr.''\ and
``Jones'' in
       ``Mr.~Jones'',        % ~ produces an unbreakable interword space.
or within a word---especially when the word is a
symbol like
       \mbox{\emph{itemnum}}
that makes little sense when hyphenated across
lines.

Footnotes\footnote{This is an example of a footnote.}
pose no problem.

\LaTeX\ is good at typesetting mathematical formulas
like
       \( x-3y + z = 7 \)
or
       \( a_{1} > x^{2n} + y^{2n} > x' \)
or
       \( \ip{A}{B} = \sum_{i} a_{i} b_{i} \).
The spaces you type in a formula are
ignored.  Remember that a letter like
       $x$                   % $ ... $  and  \( ... \)  are equivalent
is a formula when it denotes a mathematical
symbol, and it should be typed as one.

\section{Displayed Text}

Text is displayed by indenting it from the left
margin.  Quotations are commonly displayed.  There
are short quotations
\begin{quote}
   This is a short quotation.  It consists of a
   single paragraph of text.  See how it is formatted.
\end{quote}
and longer ones.
\begin{quotation}
   This is a longer quotation.  It consists of two
   paragraphs of text, neither of which are
   particularly interesting.

   This is the second paragraph of the quotation.  It
   is just as dull as the first paragraph.
\end{quotation}
Another frequently-displayed structure is a list.
The following is an example of an \emph{itemized}
list.
\begin{itemize}
   \item This is the first item of an itemized list.
         Each item in the list is marked with a ``tick''.
         You don't have to worry about what kind of tick
         mark is used.

   \item This is the second item of the list.  It
         contains another list nested inside it.  The inner
         list is an \emph{enumerated} list.
         \begin{enumerate}
            \item This is the first item of an enumerated
                  list that is nested within the itemized list.

            \item This is the second item of the inner list.
                  \LaTeX\ allows you to nest lists deeper than
                  you really should.
         \end{enumerate}
         This is the rest of the second item of the outer
         list.  It is no more interesting than any other
         part of the item.
   \item This is the third item of the list.
\end{itemize}
You can even display poetry.
\begin{verse}
   There is an environment
    for verse \\             % The \\ command separates lines
   Whose features some poets % within a stanza.
   will curse.

                             % One or more blank lines separate stanzas.

   For instead of making\\
   Them do \emph{all} line breaking, \\
   It allows them to put too many words on a line when they'd rather be
   forced to be terse.
\end{verse}

Mathematical formulas may also be displayed.  A
displayed formula
is
one-line long; multiline
formulas require special formatting instructions.
   \[  \ip{\Gamma}{\psi'} = x'' + y^{2} + z_{i}^{n}\]
Don't start a paragraph with a displayed equation,
nor make one a paragraph by itself.

\end{document}               % End of document.

\section{Bench Section 212}
% synthetic bench section 212
% This is a sample LaTeX input file.  (Version of 12 August 2004.)
%
% A '%' character causes TeX to ignore all remaining text on the line,
% and is used for comments like this one.

\documentclass{article}      % Specifies the document class

                             % The preamble begins here.
\title{An Example Document}  % Declares the document's title.
\author{Leslie Lamport}      % Declares the author's name.
\date{January 21, 1994}      % Deleting this command produces today's date.

\newcommand{\ip}[2]{(#1, #2)}
                             % Defines \ip{arg1}{arg2} to mean
                             % (arg1, arg2).

%\newcommand{\ip}[2]{\langle #1 | #2\rangle}
                             % This is an alternative definition of
                             % \ip that is commented out.

\begin{document}             % End of preamble and beginning of text.

\maketitle                   % Produces the title.

This is an example input file.  Comparing it with
the output it generates can show you how to
produce a simple document of your own.

\section{Ordinary Text}      % Produces section heading.  Lower-level
                             % sections are begun with similar
                             % \subsection and \subsubsection commands.

The ends  of words and sentences are marked
  by   spaces. It  doesn't matter how many
spaces    you type; one is as good as 100.  The
end of   a line counts as a space.

One   or more   blank lines denote the  end
of  a paragraph.

Since any number of consecutive spaces are treated
like a single one, the formatting of the input
file makes no difference to
      \LaTeX,                % The \LaTeX command generates the LaTeX logo.
but it makes a difference to you.  When you use
\LaTeX, making your input file as easy to read
as possible will be a great help as you write
your document and when you change it.  This sample
file shows how you can add comments to your own input
file.

Because printing is different from typewriting,
there are a number of things that you have to do
differently when preparing an input file than if
you were just typing the document directly.
Quotation marks like
       ``this''
have to be handled specially, as do quotes within
quotes:
       ``\,`this'            % \, separates the double and single quote.
        is what I just
        wrote, not  `that'\,''.

Dashes come in three sizes: an
       intra-word
dash, a medium dash for number ranges like
       1--2,
and a punctuation
       dash---like
this.

A sentence-ending space should be larger than the
space between words within a sentence.  You
sometimes have to type special commands in
conjunction with punctuation characters to get
this right, as in the following sentence.
       Gnats, gnus, etc.\ all  % `\ ' makes an inter-word space.
       begin with G\@.         % \@ marks end-of-sentence punctuation.
You should check the spaces after periods when
reading your output to make sure you haven't
forgotten any special cases.  Generating an
ellipsis
       \ldots\               % `\ ' is needed after `\ldots' because TeX
                             % ignores spaces after command names like \ldots
                             % made from \ + letters.
                             %
                             % Note how a `%' character causes TeX to ignore
                             % the end of the input line, so these blank lines
                             % do not start a new paragraph.
                             %
with the right spacing around the periods requires
a special command.

\LaTeX\ interprets some common characters as
commands, so you must type special commands to
generate them.  These characters include the
following:
       \$ \& \% \# \{ and \}.

In printing, text is usually emphasized with an
       \emph{italic}
type style.

\begin{em}
   A long segment of text can also be emphasized
   in this way.  Text within such a segment can be
   given \emph{additional} emphasis.
\end{em}

It is sometimes necessary to prevent \LaTeX\ from
breaking a line where it might otherwise do so.
This may be at a space, as between the ``Mr.''\ and
``Jones'' in
       ``Mr.~Jones'',        % ~ produces an unbreakable interword space.
or within a word---especially when the word is a
symbol like
       \mbox{\emph{itemnum}}
that makes little sense when hyphenated across
lines.

Footnotes\footnote{This is an example of a footnote.}
pose no problem.

\LaTeX\ is good at typesetting mathematical formulas
like
       \( x-3y + z = 7 \)
or
       \( a_{1} > x^{2n} + y^{2n} > x' \)
or
       \( \ip{A}{B} = \sum_{i} a_{i} b_{i} \).
The spaces you type in a formula are
ignored.  Remember that a letter like
       $x$                   % $ ... $  and  \( ... \)  are equivalent
is a formula when it denotes a mathematical
symbol, and it should be typed as one.

\section{Displayed Text}

Text is displayed by indenting it from the left
margin.  Quotations are commonly displayed.  There
are short quotations
\begin{quote}
   This is a short quotation.  It consists of a
   single paragraph of text.  See how it is formatted.
\end{quote}
and longer ones.
\begin{quotation}
   This is a longer quotation.  It consists of two
   paragraphs of text, neither of which are
   particularly interesting.

   This is the second paragraph of the quotation.  It
   is just as dull as the first paragraph.
\end{quotation}
Another frequently-displayed structure is a list.
The following is an example of an \emph{itemized}
list.
\begin{itemize}
   \item This is the first item of an itemized list.
         Each item in the list is marked with a ``tick''.
         You don't have to worry about what kind of tick
         mark is used.

   \item This is the second item of the list.  It
         contains another list nested inside it.  The inner
         list is an \emph{enumerated} list.
         \begin{enumerate}
            \item This is the first item of an enumerated
                  list that is nested within the itemized list.

            \item This is the second item of the inner list.
                  \LaTeX\ allows you to nest lists deeper than
                  you really should.
         \end{enumerate}
         This is the rest of the second item of the outer
         list.  It is no more interesting than any other
         part of the item.
   \item This is the third item of the list.
\end{itemize}
You can even display poetry.
\begin{verse}
   There is an environment
    for verse \\             % The \\ command separates lines
   Whose features some poets % within a stanza.
   will curse.

                             % One or more blank lines separate stanzas.

   For instead of making\\
   Them do \emph{all} line breaking, \\
   It allows them to put too many words on a line when they'd rather be
   forced to be terse.
\end{verse}

Mathematical formulas may also be displayed.  A
displayed formula
is
one-line long; multiline
formulas require special formatting instructions.
   \[  \ip{\Gamma}{\psi'} = x'' + y^{2} + z_{i}^{n}\]
Don't start a paragraph with a displayed equation,
nor make one a paragraph by itself.

\end{document}               % End of document.

\section{Bench Section 213}
% synthetic bench section 213
% This is a sample LaTeX input file.  (Version of 12 August 2004.)
%
% A '%' character causes TeX to ignore all remaining text on the line,
% and is used for comments like this one.

\documentclass{article}      % Specifies the document class

                             % The preamble begins here.
\title{An Example Document}  % Declares the document's title.
\author{Leslie Lamport}      % Declares the author's name.
\date{January 21, 1994}      % Deleting this command produces today's date.

\newcommand{\ip}[2]{(#1, #2)}
                             % Defines \ip{arg1}{arg2} to mean
                             % (arg1, arg2).

%\newcommand{\ip}[2]{\langle #1 | #2\rangle}
                             % This is an alternative definition of
                             % \ip that is commented out.

\begin{document}             % End of preamble and beginning of text.

\maketitle                   % Produces the title.

This is an example input file.  Comparing it with
the output it generates can show you how to
produce a simple document of your own.

\section{Ordinary Text}      % Produces section heading.  Lower-level
                             % sections are begun with similar
                             % \subsection and \subsubsection commands.

The ends  of words and sentences are marked
  by   spaces. It  doesn't matter how many
spaces    you type; one is as good as 100.  The
end of   a line counts as a space.

One   or more   blank lines denote the  end
of  a paragraph.

Since any number of consecutive spaces are treated
like a single one, the formatting of the input
file makes no difference to
      \LaTeX,                % The \LaTeX command generates the LaTeX logo.
but it makes a difference to you.  When you use
\LaTeX, making your input file as easy to read
as possible will be a great help as you write
your document and when you change it.  This sample
file shows how you can add comments to your own input
file.

Because printing is different from typewriting,
there are a number of things that you have to do
differently when preparing an input file than if
you were just typing the document directly.
Quotation marks like
       ``this''
have to be handled specially, as do quotes within
quotes:
       ``\,`this'            % \, separates the double and single quote.
        is what I just
        wrote, not  `that'\,''.

Dashes come in three sizes: an
       intra-word
dash, a medium dash for number ranges like
       1--2,
and a punctuation
       dash---like
this.

A sentence-ending space should be larger than the
space between words within a sentence.  You
sometimes have to type special commands in
conjunction with punctuation characters to get
this right, as in the following sentence.
       Gnats, gnus, etc.\ all  % `\ ' makes an inter-word space.
       begin with G\@.         % \@ marks end-of-sentence punctuation.
You should check the spaces after periods when
reading your output to make sure you haven't
forgotten any special cases.  Generating an
ellipsis
       \ldots\               % `\ ' is needed after `\ldots' because TeX
                             % ignores spaces after command names like \ldots
                             % made from \ + letters.
                             %
                             % Note how a `%' character causes TeX to ignore
                             % the end of the input line, so these blank lines
                             % do not start a new paragraph.
                             %
with the right spacing around the periods requires
a special command.

\LaTeX\ interprets some common characters as
commands, so you must type special commands to
generate them.  These characters include the
following:
       \$ \& \% \# \{ and \}.

In printing, text is usually emphasized with an
       \emph{italic}
type style.

\begin{em}
   A long segment of text can also be emphasized
   in this way.  Text within such a segment can be
   given \emph{additional} emphasis.
\end{em}

It is sometimes necessary to prevent \LaTeX\ from
breaking a line where it might otherwise do so.
This may be at a space, as between the ``Mr.''\ and
``Jones'' in
       ``Mr.~Jones'',        % ~ produces an unbreakable interword space.
or within a word---especially when the word is a
symbol like
       \mbox{\emph{itemnum}}
that makes little sense when hyphenated across
lines.

Footnotes\footnote{This is an example of a footnote.}
pose no problem.

\LaTeX\ is good at typesetting mathematical formulas
like
       \( x-3y + z = 7 \)
or
       \( a_{1} > x^{2n} + y^{2n} > x' \)
or
       \( \ip{A}{B} = \sum_{i} a_{i} b_{i} \).
The spaces you type in a formula are
ignored.  Remember that a letter like
       $x$                   % $ ... $  and  \( ... \)  are equivalent
is a formula when it denotes a mathematical
symbol, and it should be typed as one.

\section{Displayed Text}

Text is displayed by indenting it from the left
margin.  Quotations are commonly displayed.  There
are short quotations
\begin{quote}
   This is a short quotation.  It consists of a
   single paragraph of text.  See how it is formatted.
\end{quote}
and longer ones.
\begin{quotation}
   This is a longer quotation.  It consists of two
   paragraphs of text, neither of which are
   particularly interesting.

   This is the second paragraph of the quotation.  It
   is just as dull as the first paragraph.
\end{quotation}
Another frequently-displayed structure is a list.
The following is an example of an \emph{itemized}
list.
\begin{itemize}
   \item This is the first item of an itemized list.
         Each item in the list is marked with a ``tick''.
         You don't have to worry about what kind of tick
         mark is used.

   \item This is the second item of the list.  It
         contains another list nested inside it.  The inner
         list is an \emph{enumerated} list.
         \begin{enumerate}
            \item This is the first item of an enumerated
                  list that is nested within the itemized list.

            \item This is the second item of the inner list.
                  \LaTeX\ allows you to nest lists deeper than
                  you really should.
         \end{enumerate}
         This is the rest of the second item of the outer
         list.  It is no more interesting than any other
         part of the item.
   \item This is the third item of the list.
\end{itemize}
You can even display poetry.
\begin{verse}
   There is an environment
    for verse \\             % The \\ command separates lines
   Whose features some poets % within a stanza.
   will curse.

                             % One or more blank lines separate stanzas.

   For instead of making\\
   Them do \emph{all} line breaking, \\
   It allows them to put too many words on a line when they'd rather be
   forced to be terse.
\end{verse}

Mathematical formulas may also be displayed.  A
displayed formula
is
one-line long; multiline
formulas require special formatting instructions.
   \[  \ip{\Gamma}{\psi'} = x'' + y^{2} + z_{i}^{n}\]
Don't start a paragraph with a displayed equation,
nor make one a paragraph by itself.

\end{document}               % End of document.

\section{Bench Section 214}
% synthetic bench section 214
% This is a sample LaTeX input file.  (Version of 12 August 2004.)
%
% A '%' character causes TeX to ignore all remaining text on the line,
% and is used for comments like this one.

\documentclass{article}      % Specifies the document class

                             % The preamble begins here.
\title{An Example Document}  % Declares the document's title.
\author{Leslie Lamport}      % Declares the author's name.
\date{January 21, 1994}      % Deleting this command produces today's date.

\newcommand{\ip}[2]{(#1, #2)}
                             % Defines \ip{arg1}{arg2} to mean
                             % (arg1, arg2).

%\newcommand{\ip}[2]{\langle #1 | #2\rangle}
                             % This is an alternative definition of
                             % \ip that is commented out.

\begin{document}             % End of preamble and beginning of text.

\maketitle                   % Produces the title.

This is an example input file.  Comparing it with
the output it generates can show you how to
produce a simple document of your own.

\section{Ordinary Text}      % Produces section heading.  Lower-level
                             % sections are begun with similar
                             % \subsection and \subsubsection commands.

The ends  of words and sentences are marked
  by   spaces. It  doesn't matter how many
spaces    you type; one is as good as 100.  The
end of   a line counts as a space.

One   or more   blank lines denote the  end
of  a paragraph.

Since any number of consecutive spaces are treated
like a single one, the formatting of the input
file makes no difference to
      \LaTeX,                % The \LaTeX command generates the LaTeX logo.
but it makes a difference to you.  When you use
\LaTeX, making your input file as easy to read
as possible will be a great help as you write
your document and when you change it.  This sample
file shows how you can add comments to your own input
file.

Because printing is different from typewriting,
there are a number of things that you have to do
differently when preparing an input file than if
you were just typing the document directly.
Quotation marks like
       ``this''
have to be handled specially, as do quotes within
quotes:
       ``\,`this'            % \, separates the double and single quote.
        is what I just
        wrote, not  `that'\,''.

Dashes come in three sizes: an
       intra-word
dash, a medium dash for number ranges like
       1--2,
and a punctuation
       dash---like
this.

A sentence-ending space should be larger than the
space between words within a sentence.  You
sometimes have to type special commands in
conjunction with punctuation characters to get
this right, as in the following sentence.
       Gnats, gnus, etc.\ all  % `\ ' makes an inter-word space.
       begin with G\@.         % \@ marks end-of-sentence punctuation.
You should check the spaces after periods when
reading your output to make sure you haven't
forgotten any special cases.  Generating an
ellipsis
       \ldots\               % `\ ' is needed after `\ldots' because TeX
                             % ignores spaces after command names like \ldots
                             % made from \ + letters.
                             %
                             % Note how a `%' character causes TeX to ignore
                             % the end of the input line, so these blank lines
                             % do not start a new paragraph.
                             %
with the right spacing around the periods requires
a special command.

\LaTeX\ interprets some common characters as
commands, so you must type special commands to
generate them.  These characters include the
following:
       \$ \& \% \# \{ and \}.

In printing, text is usually emphasized with an
       \emph{italic}
type style.

\begin{em}
   A long segment of text can also be emphasized
   in this way.  Text within such a segment can be
   given \emph{additional} emphasis.
\end{em}

It is sometimes necessary to prevent \LaTeX\ from
breaking a line where it might otherwise do so.
This may be at a space, as between the ``Mr.''\ and
``Jones'' in
       ``Mr.~Jones'',        % ~ produces an unbreakable interword space.
or within a word---especially when the word is a
symbol like
       \mbox{\emph{itemnum}}
that makes little sense when hyphenated across
lines.

Footnotes\footnote{This is an example of a footnote.}
pose no problem.

\LaTeX\ is good at typesetting mathematical formulas
like
       \( x-3y + z = 7 \)
or
       \( a_{1} > x^{2n} + y^{2n} > x' \)
or
       \( \ip{A}{B} = \sum_{i} a_{i} b_{i} \).
The spaces you type in a formula are
ignored.  Remember that a letter like
       $x$                   % $ ... $  and  \( ... \)  are equivalent
is a formula when it denotes a mathematical
symbol, and it should be typed as one.

\section{Displayed Text}

Text is displayed by indenting it from the left
margin.  Quotations are commonly displayed.  There
are short quotations
\begin{quote}
   This is a short quotation.  It consists of a
   single paragraph of text.  See how it is formatted.
\end{quote}
and longer ones.
\begin{quotation}
   This is a longer quotation.  It consists of two
   paragraphs of text, neither of which are
   particularly interesting.

   This is the second paragraph of the quotation.  It
   is just as dull as the first paragraph.
\end{quotation}
Another frequently-displayed structure is a list.
The following is an example of an \emph{itemized}
list.
\begin{itemize}
   \item This is the first item of an itemized list.
         Each item in the list is marked with a ``tick''.
         You don't have to worry about what kind of tick
         mark is used.

   \item This is the second item of the list.  It
         contains another list nested inside it.  The inner
         list is an \emph{enumerated} list.
         \begin{enumerate}
            \item This is the first item of an enumerated
                  list that is nested within the itemized list.

            \item This is the second item of the inner list.
                  \LaTeX\ allows you to nest lists deeper than
                  you really should.
         \end{enumerate}
         This is the rest of the second item of the outer
         list.  It is no more interesting than any other
         part of the item.
   \item This is the third item of the list.
\end{itemize}
You can even display poetry.
\begin{verse}
   There is an environment
    for verse \\             % The \\ command separates lines
   Whose features some poets % within a stanza.
   will curse.

                             % One or more blank lines separate stanzas.

   For instead of making\\
   Them do \emph{all} line breaking, \\
   It allows them to put too many words on a line when they'd rather be
   forced to be terse.
\end{verse}

Mathematical formulas may also be displayed.  A
displayed formula
is
one-line long; multiline
formulas require special formatting instructions.
   \[  \ip{\Gamma}{\psi'} = x'' + y^{2} + z_{i}^{n}\]
Don't start a paragraph with a displayed equation,
nor make one a paragraph by itself.

\end{document}               % End of document.

\section{Bench Section 215}
% synthetic bench section 215
% This is a sample LaTeX input file.  (Version of 12 August 2004.)
%
% A '%' character causes TeX to ignore all remaining text on the line,
% and is used for comments like this one.

\documentclass{article}      % Specifies the document class

                             % The preamble begins here.
\title{An Example Document}  % Declares the document's title.
\author{Leslie Lamport}      % Declares the author's name.
\date{January 21, 1994}      % Deleting this command produces today's date.

\newcommand{\ip}[2]{(#1, #2)}
                             % Defines \ip{arg1}{arg2} to mean
                             % (arg1, arg2).

%\newcommand{\ip}[2]{\langle #1 | #2\rangle}
                             % This is an alternative definition of
                             % \ip that is commented out.

\begin{document}             % End of preamble and beginning of text.

\maketitle                   % Produces the title.

This is an example input file.  Comparing it with
the output it generates can show you how to
produce a simple document of your own.

\section{Ordinary Text}      % Produces section heading.  Lower-level
                             % sections are begun with similar
                             % \subsection and \subsubsection commands.

The ends  of words and sentences are marked
  by   spaces. It  doesn't matter how many
spaces    you type; one is as good as 100.  The
end of   a line counts as a space.

One   or more   blank lines denote the  end
of  a paragraph.

Since any number of consecutive spaces are treated
like a single one, the formatting of the input
file makes no difference to
      \LaTeX,                % The \LaTeX command generates the LaTeX logo.
but it makes a difference to you.  When you use
\LaTeX, making your input file as easy to read
as possible will be a great help as you write
your document and when you change it.  This sample
file shows how you can add comments to your own input
file.

Because printing is different from typewriting,
there are a number of things that you have to do
differently when preparing an input file than if
you were just typing the document directly.
Quotation marks like
       ``this''
have to be handled specially, as do quotes within
quotes:
       ``\,`this'            % \, separates the double and single quote.
        is what I just
        wrote, not  `that'\,''.

Dashes come in three sizes: an
       intra-word
dash, a medium dash for number ranges like
       1--2,
and a punctuation
       dash---like
this.

A sentence-ending space should be larger than the
space between words within a sentence.  You
sometimes have to type special commands in
conjunction with punctuation characters to get
this right, as in the following sentence.
       Gnats, gnus, etc.\ all  % `\ ' makes an inter-word space.
       begin with G\@.         % \@ marks end-of-sentence punctuation.
You should check the spaces after periods when
reading your output to make sure you haven't
forgotten any special cases.  Generating an
ellipsis
       \ldots\               % `\ ' is needed after `\ldots' because TeX
                             % ignores spaces after command names like \ldots
                             % made from \ + letters.
                             %
                             % Note how a `%' character causes TeX to ignore
                             % the end of the input line, so these blank lines
                             % do not start a new paragraph.
                             %
with the right spacing around the periods requires
a special command.

\LaTeX\ interprets some common characters as
commands, so you must type special commands to
generate them.  These characters include the
following:
       \$ \& \% \# \{ and \}.

In printing, text is usually emphasized with an
       \emph{italic}
type style.

\begin{em}
   A long segment of text can also be emphasized
   in this way.  Text within such a segment can be
   given \emph{additional} emphasis.
\end{em}

It is sometimes necessary to prevent \LaTeX\ from
breaking a line where it might otherwise do so.
This may be at a space, as between the ``Mr.''\ and
``Jones'' in
       ``Mr.~Jones'',        % ~ produces an unbreakable interword space.
or within a word---especially when the word is a
symbol like
       \mbox{\emph{itemnum}}
that makes little sense when hyphenated across
lines.

Footnotes\footnote{This is an example of a footnote.}
pose no problem.

\LaTeX\ is good at typesetting mathematical formulas
like
       \( x-3y + z = 7 \)
or
       \( a_{1} > x^{2n} + y^{2n} > x' \)
or
       \( \ip{A}{B} = \sum_{i} a_{i} b_{i} \).
The spaces you type in a formula are
ignored.  Remember that a letter like
       $x$                   % $ ... $  and  \( ... \)  are equivalent
is a formula when it denotes a mathematical
symbol, and it should be typed as one.

\section{Displayed Text}

Text is displayed by indenting it from the left
margin.  Quotations are commonly displayed.  There
are short quotations
\begin{quote}
   This is a short quotation.  It consists of a
   single paragraph of text.  See how it is formatted.
\end{quote}
and longer ones.
\begin{quotation}
   This is a longer quotation.  It consists of two
   paragraphs of text, neither of which are
   particularly interesting.

   This is the second paragraph of the quotation.  It
   is just as dull as the first paragraph.
\end{quotation}
Another frequently-displayed structure is a list.
The following is an example of an \emph{itemized}
list.
\begin{itemize}
   \item This is the first item of an itemized list.
         Each item in the list is marked with a ``tick''.
         You don't have to worry about what kind of tick
         mark is used.

   \item This is the second item of the list.  It
         contains another list nested inside it.  The inner
         list is an \emph{enumerated} list.
         \begin{enumerate}
            \item This is the first item of an enumerated
                  list that is nested within the itemized list.

            \item This is the second item of the inner list.
                  \LaTeX\ allows you to nest lists deeper than
                  you really should.
         \end{enumerate}
         This is the rest of the second item of the outer
         list.  It is no more interesting than any other
         part of the item.
   \item This is the third item of the list.
\end{itemize}
You can even display poetry.
\begin{verse}
   There is an environment
    for verse \\             % The \\ command separates lines
   Whose features some poets % within a stanza.
   will curse.

                             % One or more blank lines separate stanzas.

   For instead of making\\
   Them do \emph{all} line breaking, \\
   It allows them to put too many words on a line when they'd rather be
   forced to be terse.
\end{verse}

Mathematical formulas may also be displayed.  A
displayed formula
is
one-line long; multiline
formulas require special formatting instructions.
   \[  \ip{\Gamma}{\psi'} = x'' + y^{2} + z_{i}^{n}\]
Don't start a paragraph with a displayed equation,
nor make one a paragraph by itself.

\end{document}               % End of document.

\section{Bench Section 216}
% synthetic bench section 216
% This is a sample LaTeX input file.  (Version of 12 August 2004.)
%
% A '%' character causes TeX to ignore all remaining text on the line,
% and is used for comments like this one.

\documentclass{article}      % Specifies the document class

                             % The preamble begins here.
\title{An Example Document}  % Declares the document's title.
\author{Leslie Lamport}      % Declares the author's name.
\date{January 21, 1994}      % Deleting this command produces today's date.

\newcommand{\ip}[2]{(#1, #2)}
                             % Defines \ip{arg1}{arg2} to mean
                             % (arg1, arg2).

%\newcommand{\ip}[2]{\langle #1 | #2\rangle}
                             % This is an alternative definition of
                             % \ip that is commented out.

\begin{document}             % End of preamble and beginning of text.

\maketitle                   % Produces the title.

This is an example input file.  Comparing it with
the output it generates can show you how to
produce a simple document of your own.

\section{Ordinary Text}      % Produces section heading.  Lower-level
                             % sections are begun with similar
                             % \subsection and \subsubsection commands.

The ends  of words and sentences are marked
  by   spaces. It  doesn't matter how many
spaces    you type; one is as good as 100.  The
end of   a line counts as a space.

One   or more   blank lines denote the  end
of  a paragraph.

Since any number of consecutive spaces are treated
like a single one, the formatting of the input
file makes no difference to
      \LaTeX,                % The \LaTeX command generates the LaTeX logo.
but it makes a difference to you.  When you use
\LaTeX, making your input file as easy to read
as possible will be a great help as you write
your document and when you change it.  This sample
file shows how you can add comments to your own input
file.

Because printing is different from typewriting,
there are a number of things that you have to do
differently when preparing an input file than if
you were just typing the document directly.
Quotation marks like
       ``this''
have to be handled specially, as do quotes within
quotes:
       ``\,`this'            % \, separates the double and single quote.
        is what I just
        wrote, not  `that'\,''.

Dashes come in three sizes: an
       intra-word
dash, a medium dash for number ranges like
       1--2,
and a punctuation
       dash---like
this.

A sentence-ending space should be larger than the
space between words within a sentence.  You
sometimes have to type special commands in
conjunction with punctuation characters to get
this right, as in the following sentence.
       Gnats, gnus, etc.\ all  % `\ ' makes an inter-word space.
       begin with G\@.         % \@ marks end-of-sentence punctuation.
You should check the spaces after periods when
reading your output to make sure you haven't
forgotten any special cases.  Generating an
ellipsis
       \ldots\               % `\ ' is needed after `\ldots' because TeX
                             % ignores spaces after command names like \ldots
                             % made from \ + letters.
                             %
                             % Note how a `%' character causes TeX to ignore
                             % the end of the input line, so these blank lines
                             % do not start a new paragraph.
                             %
with the right spacing around the periods requires
a special command.

\LaTeX\ interprets some common characters as
commands, so you must type special commands to
generate them.  These characters include the
following:
       \$ \& \% \# \{ and \}.

In printing, text is usually emphasized with an
       \emph{italic}
type style.

\begin{em}
   A long segment of text can also be emphasized
   in this way.  Text within such a segment can be
   given \emph{additional} emphasis.
\end{em}

It is sometimes necessary to prevent \LaTeX\ from
breaking a line where it might otherwise do so.
This may be at a space, as between the ``Mr.''\ and
``Jones'' in
       ``Mr.~Jones'',        % ~ produces an unbreakable interword space.
or within a word---especially when the word is a
symbol like
       \mbox{\emph{itemnum}}
that makes little sense when hyphenated across
lines.

Footnotes\footnote{This is an example of a footnote.}
pose no problem.

\LaTeX\ is good at typesetting mathematical formulas
like
       \( x-3y + z = 7 \)
or
       \( a_{1} > x^{2n} + y^{2n} > x' \)
or
       \( \ip{A}{B} = \sum_{i} a_{i} b_{i} \).
The spaces you type in a formula are
ignored.  Remember that a letter like
       $x$                   % $ ... $  and  \( ... \)  are equivalent
is a formula when it denotes a mathematical
symbol, and it should be typed as one.

\section{Displayed Text}

Text is displayed by indenting it from the left
margin.  Quotations are commonly displayed.  There
are short quotations
\begin{quote}
   This is a short quotation.  It consists of a
   single paragraph of text.  See how it is formatted.
\end{quote}
and longer ones.
\begin{quotation}
   This is a longer quotation.  It consists of two
   paragraphs of text, neither of which are
   particularly interesting.

   This is the second paragraph of the quotation.  It
   is just as dull as the first paragraph.
\end{quotation}
Another frequently-displayed structure is a list.
The following is an example of an \emph{itemized}
list.
\begin{itemize}
   \item This is the first item of an itemized list.
         Each item in the list is marked with a ``tick''.
         You don't have to worry about what kind of tick
         mark is used.

   \item This is the second item of the list.  It
         contains another list nested inside it.  The inner
         list is an \emph{enumerated} list.
         \begin{enumerate}
            \item This is the first item of an enumerated
                  list that is nested within the itemized list.

            \item This is the second item of the inner list.
                  \LaTeX\ allows you to nest lists deeper than
                  you really should.
         \end{enumerate}
         This is the rest of the second item of the outer
         list.  It is no more interesting than any other
         part of the item.
   \item This is the third item of the list.
\end{itemize}
You can even display poetry.
\begin{verse}
   There is an environment
    for verse \\             % The \\ command separates lines
   Whose features some poets % within a stanza.
   will curse.

                             % One or more blank lines separate stanzas.

   For instead of making\\
   Them do \emph{all} line breaking, \\
   It allows them to put too many words on a line when they'd rather be
   forced to be terse.
\end{verse}

Mathematical formulas may also be displayed.  A
displayed formula
is
one-line long; multiline
formulas require special formatting instructions.
   \[  \ip{\Gamma}{\psi'} = x'' + y^{2} + z_{i}^{n}\]
Don't start a paragraph with a displayed equation,
nor make one a paragraph by itself.

\end{document}               % End of document.

\section{Bench Section 217}
% synthetic bench section 217
% This is a sample LaTeX input file.  (Version of 12 August 2004.)
%
% A '%' character causes TeX to ignore all remaining text on the line,
% and is used for comments like this one.

\documentclass{article}      % Specifies the document class

                             % The preamble begins here.
\title{An Example Document}  % Declares the document's title.
\author{Leslie Lamport}      % Declares the author's name.
\date{January 21, 1994}      % Deleting this command produces today's date.

\newcommand{\ip}[2]{(#1, #2)}
                             % Defines \ip{arg1}{arg2} to mean
                             % (arg1, arg2).

%\newcommand{\ip}[2]{\langle #1 | #2\rangle}
                             % This is an alternative definition of
                             % \ip that is commented out.

\begin{document}             % End of preamble and beginning of text.

\maketitle                   % Produces the title.

This is an example input file.  Comparing it with
the output it generates can show you how to
produce a simple document of your own.

\section{Ordinary Text}      % Produces section heading.  Lower-level
                             % sections are begun with similar
                             % \subsection and \subsubsection commands.

The ends  of words and sentences are marked
  by   spaces. It  doesn't matter how many
spaces    you type; one is as good as 100.  The
end of   a line counts as a space.

One   or more   blank lines denote the  end
of  a paragraph.

Since any number of consecutive spaces are treated
like a single one, the formatting of the input
file makes no difference to
      \LaTeX,                % The \LaTeX command generates the LaTeX logo.
but it makes a difference to you.  When you use
\LaTeX, making your input file as easy to read
as possible will be a great help as you write
your document and when you change it.  This sample
file shows how you can add comments to your own input
file.

Because printing is different from typewriting,
there are a number of things that you have to do
differently when preparing an input file than if
you were just typing the document directly.
Quotation marks like
       ``this''
have to be handled specially, as do quotes within
quotes:
       ``\,`this'            % \, separates the double and single quote.
        is what I just
        wrote, not  `that'\,''.

Dashes come in three sizes: an
       intra-word
dash, a medium dash for number ranges like
       1--2,
and a punctuation
       dash---like
this.

A sentence-ending space should be larger than the
space between words within a sentence.  You
sometimes have to type special commands in
conjunction with punctuation characters to get
this right, as in the following sentence.
       Gnats, gnus, etc.\ all  % `\ ' makes an inter-word space.
       begin with G\@.         % \@ marks end-of-sentence punctuation.
You should check the spaces after periods when
reading your output to make sure you haven't
forgotten any special cases.  Generating an
ellipsis
       \ldots\               % `\ ' is needed after `\ldots' because TeX
                             % ignores spaces after command names like \ldots
                             % made from \ + letters.
                             %
                             % Note how a `%' character causes TeX to ignore
                             % the end of the input line, so these blank lines
                             % do not start a new paragraph.
                             %
with the right spacing around the periods requires
a special command.

\LaTeX\ interprets some common characters as
commands, so you must type special commands to
generate them.  These characters include the
following:
       \$ \& \% \# \{ and \}.

In printing, text is usually emphasized with an
       \emph{italic}
type style.

\begin{em}
   A long segment of text can also be emphasized
   in this way.  Text within such a segment can be
   given \emph{additional} emphasis.
\end{em}

It is sometimes necessary to prevent \LaTeX\ from
breaking a line where it might otherwise do so.
This may be at a space, as between the ``Mr.''\ and
``Jones'' in
       ``Mr.~Jones'',        % ~ produces an unbreakable interword space.
or within a word---especially when the word is a
symbol like
       \mbox{\emph{itemnum}}
that makes little sense when hyphenated across
lines.

Footnotes\footnote{This is an example of a footnote.}
pose no problem.

\LaTeX\ is good at typesetting mathematical formulas
like
       \( x-3y + z = 7 \)
or
       \( a_{1} > x^{2n} + y^{2n} > x' \)
or
       \( \ip{A}{B} = \sum_{i} a_{i} b_{i} \).
The spaces you type in a formula are
ignored.  Remember that a letter like
       $x$                   % $ ... $  and  \( ... \)  are equivalent
is a formula when it denotes a mathematical
symbol, and it should be typed as one.

\section{Displayed Text}

Text is displayed by indenting it from the left
margin.  Quotations are commonly displayed.  There
are short quotations
\begin{quote}
   This is a short quotation.  It consists of a
   single paragraph of text.  See how it is formatted.
\end{quote}
and longer ones.
\begin{quotation}
   This is a longer quotation.  It consists of two
   paragraphs of text, neither of which are
   particularly interesting.

   This is the second paragraph of the quotation.  It
   is just as dull as the first paragraph.
\end{quotation}
Another frequently-displayed structure is a list.
The following is an example of an \emph{itemized}
list.
\begin{itemize}
   \item This is the first item of an itemized list.
         Each item in the list is marked with a ``tick''.
         You don't have to worry about what kind of tick
         mark is used.

   \item This is the second item of the list.  It
         contains another list nested inside it.  The inner
         list is an \emph{enumerated} list.
         \begin{enumerate}
            \item This is the first item of an enumerated
                  list that is nested within the itemized list.

            \item This is the second item of the inner list.
                  \LaTeX\ allows you to nest lists deeper than
                  you really should.
         \end{enumerate}
         This is the rest of the second item of the outer
         list.  It is no more interesting than any other
         part of the item.
   \item This is the third item of the list.
\end{itemize}
You can even display poetry.
\begin{verse}
   There is an environment
    for verse \\             % The \\ command separates lines
   Whose features some poets % within a stanza.
   will curse.

                             % One or more blank lines separate stanzas.

   For instead of making\\
   Them do \emph{all} line breaking, \\
   It allows them to put too many words on a line when they'd rather be
   forced to be terse.
\end{verse}

Mathematical formulas may also be displayed.  A
displayed formula
is
one-line long; multiline
formulas require special formatting instructions.
   \[  \ip{\Gamma}{\psi'} = x'' + y^{2} + z_{i}^{n}\]
Don't start a paragraph with a displayed equation,
nor make one a paragraph by itself.

\end{document}               % End of document.

\section{Bench Section 218}
% synthetic bench section 218
% This is a sample LaTeX input file.  (Version of 12 August 2004.)
%
% A '%' character causes TeX to ignore all remaining text on the line,
% and is used for comments like this one.

\documentclass{article}      % Specifies the document class

                             % The preamble begins here.
\title{An Example Document}  % Declares the document's title.
\author{Leslie Lamport}      % Declares the author's name.
\date{January 21, 1994}      % Deleting this command produces today's date.

\newcommand{\ip}[2]{(#1, #2)}
                             % Defines \ip{arg1}{arg2} to mean
                             % (arg1, arg2).

%\newcommand{\ip}[2]{\langle #1 | #2\rangle}
                             % This is an alternative definition of
                             % \ip that is commented out.

\begin{document}             % End of preamble and beginning of text.

\maketitle                   % Produces the title.

This is an example input file.  Comparing it with
the output it generates can show you how to
produce a simple document of your own.

\section{Ordinary Text}      % Produces section heading.  Lower-level
                             % sections are begun with similar
                             % \subsection and \subsubsection commands.

The ends  of words and sentences are marked
  by   spaces. It  doesn't matter how many
spaces    you type; one is as good as 100.  The
end of   a line counts as a space.

One   or more   blank lines denote the  end
of  a paragraph.

Since any number of consecutive spaces are treated
like a single one, the formatting of the input
file makes no difference to
      \LaTeX,                % The \LaTeX command generates the LaTeX logo.
but it makes a difference to you.  When you use
\LaTeX, making your input file as easy to read
as possible will be a great help as you write
your document and when you change it.  This sample
file shows how you can add comments to your own input
file.

Because printing is different from typewriting,
there are a number of things that you have to do
differently when preparing an input file than if
you were just typing the document directly.
Quotation marks like
       ``this''
have to be handled specially, as do quotes within
quotes:
       ``\,`this'            % \, separates the double and single quote.
        is what I just
        wrote, not  `that'\,''.

Dashes come in three sizes: an
       intra-word
dash, a medium dash for number ranges like
       1--2,
and a punctuation
       dash---like
this.

A sentence-ending space should be larger than the
space between words within a sentence.  You
sometimes have to type special commands in
conjunction with punctuation characters to get
this right, as in the following sentence.
       Gnats, gnus, etc.\ all  % `\ ' makes an inter-word space.
       begin with G\@.         % \@ marks end-of-sentence punctuation.
You should check the spaces after periods when
reading your output to make sure you haven't
forgotten any special cases.  Generating an
ellipsis
       \ldots\               % `\ ' is needed after `\ldots' because TeX
                             % ignores spaces after command names like \ldots
                             % made from \ + letters.
                             %
                             % Note how a `%' character causes TeX to ignore
                             % the end of the input line, so these blank lines
                             % do not start a new paragraph.
                             %
with the right spacing around the periods requires
a special command.

\LaTeX\ interprets some common characters as
commands, so you must type special commands to
generate them.  These characters include the
following:
       \$ \& \% \# \{ and \}.

In printing, text is usually emphasized with an
       \emph{italic}
type style.

\begin{em}
   A long segment of text can also be emphasized
   in this way.  Text within such a segment can be
   given \emph{additional} emphasis.
\end{em}

It is sometimes necessary to prevent \LaTeX\ from
breaking a line where it might otherwise do so.
This may be at a space, as between the ``Mr.''\ and
``Jones'' in
       ``Mr.~Jones'',        % ~ produces an unbreakable interword space.
or within a word---especially when the word is a
symbol like
       \mbox{\emph{itemnum}}
that makes little sense when hyphenated across
lines.

Footnotes\footnote{This is an example of a footnote.}
pose no problem.

\LaTeX\ is good at typesetting mathematical formulas
like
       \( x-3y + z = 7 \)
or
       \( a_{1} > x^{2n} + y^{2n} > x' \)
or
       \( \ip{A}{B} = \sum_{i} a_{i} b_{i} \).
The spaces you type in a formula are
ignored.  Remember that a letter like
       $x$                   % $ ... $  and  \( ... \)  are equivalent
is a formula when it denotes a mathematical
symbol, and it should be typed as one.

\section{Displayed Text}

Text is displayed by indenting it from the left
margin.  Quotations are commonly displayed.  There
are short quotations
\begin{quote}
   This is a short quotation.  It consists of a
   single paragraph of text.  See how it is formatted.
\end{quote}
and longer ones.
\begin{quotation}
   This is a longer quotation.  It consists of two
   paragraphs of text, neither of which are
   particularly interesting.

   This is the second paragraph of the quotation.  It
   is just as dull as the first paragraph.
\end{quotation}
Another frequently-displayed structure is a list.
The following is an example of an \emph{itemized}
list.
\begin{itemize}
   \item This is the first item of an itemized list.
         Each item in the list is marked with a ``tick''.
         You don't have to worry about what kind of tick
         mark is used.

   \item This is the second item of the list.  It
         contains another list nested inside it.  The inner
         list is an \emph{enumerated} list.
         \begin{enumerate}
            \item This is the first item of an enumerated
                  list that is nested within the itemized list.

            \item This is the second item of the inner list.
                  \LaTeX\ allows you to nest lists deeper than
                  you really should.
         \end{enumerate}
         This is the rest of the second item of the outer
         list.  It is no more interesting than any other
         part of the item.
   \item This is the third item of the list.
\end{itemize}
You can even display poetry.
\begin{verse}
   There is an environment
    for verse \\             % The \\ command separates lines
   Whose features some poets % within a stanza.
   will curse.

                             % One or more blank lines separate stanzas.

   For instead of making\\
   Them do \emph{all} line breaking, \\
   It allows them to put too many words on a line when they'd rather be
   forced to be terse.
\end{verse}

Mathematical formulas may also be displayed.  A
displayed formula
is
one-line long; multiline
formulas require special formatting instructions.
   \[  \ip{\Gamma}{\psi'} = x'' + y^{2} + z_{i}^{n}\]
Don't start a paragraph with a displayed equation,
nor make one a paragraph by itself.

\end{document}               % End of document.

\section{Bench Section 219}
% synthetic bench section 219
% This is a sample LaTeX input file.  (Version of 12 August 2004.)
%
% A '%' character causes TeX to ignore all remaining text on the line,
% and is used for comments like this one.

\documentclass{article}      % Specifies the document class

                             % The preamble begins here.
\title{An Example Document}  % Declares the document's title.
\author{Leslie Lamport}      % Declares the author's name.
\date{January 21, 1994}      % Deleting this command produces today's date.

\newcommand{\ip}[2]{(#1, #2)}
                             % Defines \ip{arg1}{arg2} to mean
                             % (arg1, arg2).

%\newcommand{\ip}[2]{\langle #1 | #2\rangle}
                             % This is an alternative definition of
                             % \ip that is commented out.

\begin{document}             % End of preamble and beginning of text.

\maketitle                   % Produces the title.

This is an example input file.  Comparing it with
the output it generates can show you how to
produce a simple document of your own.

\section{Ordinary Text}      % Produces section heading.  Lower-level
                             % sections are begun with similar
                             % \subsection and \subsubsection commands.

The ends  of words and sentences are marked
  by   spaces. It  doesn't matter how many
spaces    you type; one is as good as 100.  The
end of   a line counts as a space.

One   or more   blank lines denote the  end
of  a paragraph.

Since any number of consecutive spaces are treated
like a single one, the formatting of the input
file makes no difference to
      \LaTeX,                % The \LaTeX command generates the LaTeX logo.
but it makes a difference to you.  When you use
\LaTeX, making your input file as easy to read
as possible will be a great help as you write
your document and when you change it.  This sample
file shows how you can add comments to your own input
file.

Because printing is different from typewriting,
there are a number of things that you have to do
differently when preparing an input file than if
you were just typing the document directly.
Quotation marks like
       ``this''
have to be handled specially, as do quotes within
quotes:
       ``\,`this'            % \, separates the double and single quote.
        is what I just
        wrote, not  `that'\,''.

Dashes come in three sizes: an
       intra-word
dash, a medium dash for number ranges like
       1--2,
and a punctuation
       dash---like
this.

A sentence-ending space should be larger than the
space between words within a sentence.  You
sometimes have to type special commands in
conjunction with punctuation characters to get
this right, as in the following sentence.
       Gnats, gnus, etc.\ all  % `\ ' makes an inter-word space.
       begin with G\@.         % \@ marks end-of-sentence punctuation.
You should check the spaces after periods when
reading your output to make sure you haven't
forgotten any special cases.  Generating an
ellipsis
       \ldots\               % `\ ' is needed after `\ldots' because TeX
                             % ignores spaces after command names like \ldots
                             % made from \ + letters.
                             %
                             % Note how a `%' character causes TeX to ignore
                             % the end of the input line, so these blank lines
                             % do not start a new paragraph.
                             %
with the right spacing around the periods requires
a special command.

\LaTeX\ interprets some common characters as
commands, so you must type special commands to
generate them.  These characters include the
following:
       \$ \& \% \# \{ and \}.

In printing, text is usually emphasized with an
       \emph{italic}
type style.

\begin{em}
   A long segment of text can also be emphasized
   in this way.  Text within such a segment can be
   given \emph{additional} emphasis.
\end{em}

It is sometimes necessary to prevent \LaTeX\ from
breaking a line where it might otherwise do so.
This may be at a space, as between the ``Mr.''\ and
``Jones'' in
       ``Mr.~Jones'',        % ~ produces an unbreakable interword space.
or within a word---especially when the word is a
symbol like
       \mbox{\emph{itemnum}}
that makes little sense when hyphenated across
lines.

Footnotes\footnote{This is an example of a footnote.}
pose no problem.

\LaTeX\ is good at typesetting mathematical formulas
like
       \( x-3y + z = 7 \)
or
       \( a_{1} > x^{2n} + y^{2n} > x' \)
or
       \( \ip{A}{B} = \sum_{i} a_{i} b_{i} \).
The spaces you type in a formula are
ignored.  Remember that a letter like
       $x$                   % $ ... $  and  \( ... \)  are equivalent
is a formula when it denotes a mathematical
symbol, and it should be typed as one.

\section{Displayed Text}

Text is displayed by indenting it from the left
margin.  Quotations are commonly displayed.  There
are short quotations
\begin{quote}
   This is a short quotation.  It consists of a
   single paragraph of text.  See how it is formatted.
\end{quote}
and longer ones.
\begin{quotation}
   This is a longer quotation.  It consists of two
   paragraphs of text, neither of which are
   particularly interesting.

   This is the second paragraph of the quotation.  It
   is just as dull as the first paragraph.
\end{quotation}
Another frequently-displayed structure is a list.
The following is an example of an \emph{itemized}
list.
\begin{itemize}
   \item This is the first item of an itemized list.
         Each item in the list is marked with a ``tick''.
         You don't have to worry about what kind of tick
         mark is used.

   \item This is the second item of the list.  It
         contains another list nested inside it.  The inner
         list is an \emph{enumerated} list.
         \begin{enumerate}
            \item This is the first item of an enumerated
                  list that is nested within the itemized list.

            \item This is the second item of the inner list.
                  \LaTeX\ allows you to nest lists deeper than
                  you really should.
         \end{enumerate}
         This is the rest of the second item of the outer
         list.  It is no more interesting than any other
         part of the item.
   \item This is the third item of the list.
\end{itemize}
You can even display poetry.
\begin{verse}
   There is an environment
    for verse \\             % The \\ command separates lines
   Whose features some poets % within a stanza.
   will curse.

                             % One or more blank lines separate stanzas.

   For instead of making\\
   Them do \emph{all} line breaking, \\
   It allows them to put too many words on a line when they'd rather be
   forced to be terse.
\end{verse}

Mathematical formulas may also be displayed.  A
displayed formula
is
one-line long; multiline
formulas require special formatting instructions.
   \[  \ip{\Gamma}{\psi'} = x'' + y^{2} + z_{i}^{n}\]
Don't start a paragraph with a displayed equation,
nor make one a paragraph by itself.

\end{document}               % End of document.

\section{Bench Section 220}
% synthetic bench section 220
% This is a sample LaTeX input file.  (Version of 12 August 2004.)
%
% A '%' character causes TeX to ignore all remaining text on the line,
% and is used for comments like this one.

\documentclass{article}      % Specifies the document class

                             % The preamble begins here.
\title{An Example Document}  % Declares the document's title.
\author{Leslie Lamport}      % Declares the author's name.
\date{January 21, 1994}      % Deleting this command produces today's date.

\newcommand{\ip}[2]{(#1, #2)}
                             % Defines \ip{arg1}{arg2} to mean
                             % (arg1, arg2).

%\newcommand{\ip}[2]{\langle #1 | #2\rangle}
                             % This is an alternative definition of
                             % \ip that is commented out.

\begin{document}             % End of preamble and beginning of text.

\maketitle                   % Produces the title.

This is an example input file.  Comparing it with
the output it generates can show you how to
produce a simple document of your own.

\section{Ordinary Text}      % Produces section heading.  Lower-level
                             % sections are begun with similar
                             % \subsection and \subsubsection commands.

The ends  of words and sentences are marked
  by   spaces. It  doesn't matter how many
spaces    you type; one is as good as 100.  The
end of   a line counts as a space.

One   or more   blank lines denote the  end
of  a paragraph.

Since any number of consecutive spaces are treated
like a single one, the formatting of the input
file makes no difference to
      \LaTeX,                % The \LaTeX command generates the LaTeX logo.
but it makes a difference to you.  When you use
\LaTeX, making your input file as easy to read
as possible will be a great help as you write
your document and when you change it.  This sample
file shows how you can add comments to your own input
file.

Because printing is different from typewriting,
there are a number of things that you have to do
differently when preparing an input file than if
you were just typing the document directly.
Quotation marks like
       ``this''
have to be handled specially, as do quotes within
quotes:
       ``\,`this'            % \, separates the double and single quote.
        is what I just
        wrote, not  `that'\,''.

Dashes come in three sizes: an
       intra-word
dash, a medium dash for number ranges like
       1--2,
and a punctuation
       dash---like
this.

A sentence-ending space should be larger than the
space between words within a sentence.  You
sometimes have to type special commands in
conjunction with punctuation characters to get
this right, as in the following sentence.
       Gnats, gnus, etc.\ all  % `\ ' makes an inter-word space.
       begin with G\@.         % \@ marks end-of-sentence punctuation.
You should check the spaces after periods when
reading your output to make sure you haven't
forgotten any special cases.  Generating an
ellipsis
       \ldots\               % `\ ' is needed after `\ldots' because TeX
                             % ignores spaces after command names like \ldots
                             % made from \ + letters.
                             %
                             % Note how a `%' character causes TeX to ignore
                             % the end of the input line, so these blank lines
                             % do not start a new paragraph.
                             %
with the right spacing around the periods requires
a special command.

\LaTeX\ interprets some common characters as
commands, so you must type special commands to
generate them.  These characters include the
following:
       \$ \& \% \# \{ and \}.

In printing, text is usually emphasized with an
       \emph{italic}
type style.

\begin{em}
   A long segment of text can also be emphasized
   in this way.  Text within such a segment can be
   given \emph{additional} emphasis.
\end{em}

It is sometimes necessary to prevent \LaTeX\ from
breaking a line where it might otherwise do so.
This may be at a space, as between the ``Mr.''\ and
``Jones'' in
       ``Mr.~Jones'',        % ~ produces an unbreakable interword space.
or within a word---especially when the word is a
symbol like
       \mbox{\emph{itemnum}}
that makes little sense when hyphenated across
lines.

Footnotes\footnote{This is an example of a footnote.}
pose no problem.

\LaTeX\ is good at typesetting mathematical formulas
like
       \( x-3y + z = 7 \)
or
       \( a_{1} > x^{2n} + y^{2n} > x' \)
or
       \( \ip{A}{B} = \sum_{i} a_{i} b_{i} \).
The spaces you type in a formula are
ignored.  Remember that a letter like
       $x$                   % $ ... $  and  \( ... \)  are equivalent
is a formula when it denotes a mathematical
symbol, and it should be typed as one.

\section{Displayed Text}

Text is displayed by indenting it from the left
margin.  Quotations are commonly displayed.  There
are short quotations
\begin{quote}
   This is a short quotation.  It consists of a
   single paragraph of text.  See how it is formatted.
\end{quote}
and longer ones.
\begin{quotation}
   This is a longer quotation.  It consists of two
   paragraphs of text, neither of which are
   particularly interesting.

   This is the second paragraph of the quotation.  It
   is just as dull as the first paragraph.
\end{quotation}
Another frequently-displayed structure is a list.
The following is an example of an \emph{itemized}
list.
\begin{itemize}
   \item This is the first item of an itemized list.
         Each item in the list is marked with a ``tick''.
         You don't have to worry about what kind of tick
         mark is used.

   \item This is the second item of the list.  It
         contains another list nested inside it.  The inner
         list is an \emph{enumerated} list.
         \begin{enumerate}
            \item This is the first item of an enumerated
                  list that is nested within the itemized list.

            \item This is the second item of the inner list.
                  \LaTeX\ allows you to nest lists deeper than
                  you really should.
         \end{enumerate}
         This is the rest of the second item of the outer
         list.  It is no more interesting than any other
         part of the item.
   \item This is the third item of the list.
\end{itemize}
You can even display poetry.
\begin{verse}
   There is an environment
    for verse \\             % The \\ command separates lines
   Whose features some poets % within a stanza.
   will curse.

                             % One or more blank lines separate stanzas.

   For instead of making\\
   Them do \emph{all} line breaking, \\
   It allows them to put too many words on a line when they'd rather be
   forced to be terse.
\end{verse}

Mathematical formulas may also be displayed.  A
displayed formula
is
one-line long; multiline
formulas require special formatting instructions.
   \[  \ip{\Gamma}{\psi'} = x'' + y^{2} + z_{i}^{n}\]
Don't start a paragraph with a displayed equation,
nor make one a paragraph by itself.

\end{document}               % End of document.

\section{Bench Section 221}
% synthetic bench section 221
% This is a sample LaTeX input file.  (Version of 12 August 2004.)
%
% A '%' character causes TeX to ignore all remaining text on the line,
% and is used for comments like this one.

\documentclass{article}      % Specifies the document class

                             % The preamble begins here.
\title{An Example Document}  % Declares the document's title.
\author{Leslie Lamport}      % Declares the author's name.
\date{January 21, 1994}      % Deleting this command produces today's date.

\newcommand{\ip}[2]{(#1, #2)}
                             % Defines \ip{arg1}{arg2} to mean
                             % (arg1, arg2).

%\newcommand{\ip}[2]{\langle #1 | #2\rangle}
                             % This is an alternative definition of
                             % \ip that is commented out.

\begin{document}             % End of preamble and beginning of text.

\maketitle                   % Produces the title.

This is an example input file.  Comparing it with
the output it generates can show you how to
produce a simple document of your own.

\section{Ordinary Text}      % Produces section heading.  Lower-level
                             % sections are begun with similar
                             % \subsection and \subsubsection commands.

The ends  of words and sentences are marked
  by   spaces. It  doesn't matter how many
spaces    you type; one is as good as 100.  The
end of   a line counts as a space.

One   or more   blank lines denote the  end
of  a paragraph.

Since any number of consecutive spaces are treated
like a single one, the formatting of the input
file makes no difference to
      \LaTeX,                % The \LaTeX command generates the LaTeX logo.
but it makes a difference to you.  When you use
\LaTeX, making your input file as easy to read
as possible will be a great help as you write
your document and when you change it.  This sample
file shows how you can add comments to your own input
file.

Because printing is different from typewriting,
there are a number of things that you have to do
differently when preparing an input file than if
you were just typing the document directly.
Quotation marks like
       ``this''
have to be handled specially, as do quotes within
quotes:
       ``\,`this'            % \, separates the double and single quote.
        is what I just
        wrote, not  `that'\,''.

Dashes come in three sizes: an
       intra-word
dash, a medium dash for number ranges like
       1--2,
and a punctuation
       dash---like
this.

A sentence-ending space should be larger than the
space between words within a sentence.  You
sometimes have to type special commands in
conjunction with punctuation characters to get
this right, as in the following sentence.
       Gnats, gnus, etc.\ all  % `\ ' makes an inter-word space.
       begin with G\@.         % \@ marks end-of-sentence punctuation.
You should check the spaces after periods when
reading your output to make sure you haven't
forgotten any special cases.  Generating an
ellipsis
       \ldots\               % `\ ' is needed after `\ldots' because TeX
                             % ignores spaces after command names like \ldots
                             % made from \ + letters.
                             %
                             % Note how a `%' character causes TeX to ignore
                             % the end of the input line, so these blank lines
                             % do not start a new paragraph.
                             %
with the right spacing around the periods requires
a special command.

\LaTeX\ interprets some common characters as
commands, so you must type special commands to
generate them.  These characters include the
following:
       \$ \& \% \# \{ and \}.

In printing, text is usually emphasized with an
       \emph{italic}
type style.

\begin{em}
   A long segment of text can also be emphasized
   in this way.  Text within such a segment can be
   given \emph{additional} emphasis.
\end{em}

It is sometimes necessary to prevent \LaTeX\ from
breaking a line where it might otherwise do so.
This may be at a space, as between the ``Mr.''\ and
``Jones'' in
       ``Mr.~Jones'',        % ~ produces an unbreakable interword space.
or within a word---especially when the word is a
symbol like
       \mbox{\emph{itemnum}}
that makes little sense when hyphenated across
lines.

Footnotes\footnote{This is an example of a footnote.}
pose no problem.

\LaTeX\ is good at typesetting mathematical formulas
like
       \( x-3y + z = 7 \)
or
       \( a_{1} > x^{2n} + y^{2n} > x' \)
or
       \( \ip{A}{B} = \sum_{i} a_{i} b_{i} \).
The spaces you type in a formula are
ignored.  Remember that a letter like
       $x$                   % $ ... $  and  \( ... \)  are equivalent
is a formula when it denotes a mathematical
symbol, and it should be typed as one.

\section{Displayed Text}

Text is displayed by indenting it from the left
margin.  Quotations are commonly displayed.  There
are short quotations
\begin{quote}
   This is a short quotation.  It consists of a
   single paragraph of text.  See how it is formatted.
\end{quote}
and longer ones.
\begin{quotation}
   This is a longer quotation.  It consists of two
   paragraphs of text, neither of which are
   particularly interesting.

   This is the second paragraph of the quotation.  It
   is just as dull as the first paragraph.
\end{quotation}
Another frequently-displayed structure is a list.
The following is an example of an \emph{itemized}
list.
\begin{itemize}
   \item This is the first item of an itemized list.
         Each item in the list is marked with a ``tick''.
         You don't have to worry about what kind of tick
         mark is used.

   \item This is the second item of the list.  It
         contains another list nested inside it.  The inner
         list is an \emph{enumerated} list.
         \begin{enumerate}
            \item This is the first item of an enumerated
                  list that is nested within the itemized list.

            \item This is the second item of the inner list.
                  \LaTeX\ allows you to nest lists deeper than
                  you really should.
         \end{enumerate}
         This is the rest of the second item of the outer
         list.  It is no more interesting than any other
         part of the item.
   \item This is the third item of the list.
\end{itemize}
You can even display poetry.
\begin{verse}
   There is an environment
    for verse \\             % The \\ command separates lines
   Whose features some poets % within a stanza.
   will curse.

                             % One or more blank lines separate stanzas.

   For instead of making\\
   Them do \emph{all} line breaking, \\
   It allows them to put too many words on a line when they'd rather be
   forced to be terse.
\end{verse}

Mathematical formulas may also be displayed.  A
displayed formula
is
one-line long; multiline
formulas require special formatting instructions.
   \[  \ip{\Gamma}{\psi'} = x'' + y^{2} + z_{i}^{n}\]
Don't start a paragraph with a displayed equation,
nor make one a paragraph by itself.

\end{document}               % End of document.

\section{Bench Section 222}
% synthetic bench section 222
% This is a sample LaTeX input file.  (Version of 12 August 2004.)
%
% A '%' character causes TeX to ignore all remaining text on the line,
% and is used for comments like this one.

\documentclass{article}      % Specifies the document class

                             % The preamble begins here.
\title{An Example Document}  % Declares the document's title.
\author{Leslie Lamport}      % Declares the author's name.
\date{January 21, 1994}      % Deleting this command produces today's date.

\newcommand{\ip}[2]{(#1, #2)}
                             % Defines \ip{arg1}{arg2} to mean
                             % (arg1, arg2).

%\newcommand{\ip}[2]{\langle #1 | #2\rangle}
                             % This is an alternative definition of
                             % \ip that is commented out.

\begin{document}             % End of preamble and beginning of text.

\maketitle                   % Produces the title.

This is an example input file.  Comparing it with
the output it generates can show you how to
produce a simple document of your own.

\section{Ordinary Text}      % Produces section heading.  Lower-level
                             % sections are begun with similar
                             % \subsection and \subsubsection commands.

The ends  of words and sentences are marked
  by   spaces. It  doesn't matter how many
spaces    you type; one is as good as 100.  The
end of   a line counts as a space.

One   or more   blank lines denote the  end
of  a paragraph.

Since any number of consecutive spaces are treated
like a single one, the formatting of the input
file makes no difference to
      \LaTeX,                % The \LaTeX command generates the LaTeX logo.
but it makes a difference to you.  When you use
\LaTeX, making your input file as easy to read
as possible will be a great help as you write
your document and when you change it.  This sample
file shows how you can add comments to your own input
file.

Because printing is different from typewriting,
there are a number of things that you have to do
differently when preparing an input file than if
you were just typing the document directly.
Quotation marks like
       ``this''
have to be handled specially, as do quotes within
quotes:
       ``\,`this'            % \, separates the double and single quote.
        is what I just
        wrote, not  `that'\,''.

Dashes come in three sizes: an
       intra-word
dash, a medium dash for number ranges like
       1--2,
and a punctuation
       dash---like
this.

A sentence-ending space should be larger than the
space between words within a sentence.  You
sometimes have to type special commands in
conjunction with punctuation characters to get
this right, as in the following sentence.
       Gnats, gnus, etc.\ all  % `\ ' makes an inter-word space.
       begin with G\@.         % \@ marks end-of-sentence punctuation.
You should check the spaces after periods when
reading your output to make sure you haven't
forgotten any special cases.  Generating an
ellipsis
       \ldots\               % `\ ' is needed after `\ldots' because TeX
                             % ignores spaces after command names like \ldots
                             % made from \ + letters.
                             %
                             % Note how a `%' character causes TeX to ignore
                             % the end of the input line, so these blank lines
                             % do not start a new paragraph.
                             %
with the right spacing around the periods requires
a special command.

\LaTeX\ interprets some common characters as
commands, so you must type special commands to
generate them.  These characters include the
following:
       \$ \& \% \# \{ and \}.

In printing, text is usually emphasized with an
       \emph{italic}
type style.

\begin{em}
   A long segment of text can also be emphasized
   in this way.  Text within such a segment can be
   given \emph{additional} emphasis.
\end{em}

It is sometimes necessary to prevent \LaTeX\ from
breaking a line where it might otherwise do so.
This may be at a space, as between the ``Mr.''\ and
``Jones'' in
       ``Mr.~Jones'',        % ~ produces an unbreakable interword space.
or within a word---especially when the word is a
symbol like
       \mbox{\emph{itemnum}}
that makes little sense when hyphenated across
lines.

Footnotes\footnote{This is an example of a footnote.}
pose no problem.

\LaTeX\ is good at typesetting mathematical formulas
like
       \( x-3y + z = 7 \)
or
       \( a_{1} > x^{2n} + y^{2n} > x' \)
or
       \( \ip{A}{B} = \sum_{i} a_{i} b_{i} \).
The spaces you type in a formula are
ignored.  Remember that a letter like
       $x$                   % $ ... $  and  \( ... \)  are equivalent
is a formula when it denotes a mathematical
symbol, and it should be typed as one.

\section{Displayed Text}

Text is displayed by indenting it from the left
margin.  Quotations are commonly displayed.  There
are short quotations
\begin{quote}
   This is a short quotation.  It consists of a
   single paragraph of text.  See how it is formatted.
\end{quote}
and longer ones.
\begin{quotation}
   This is a longer quotation.  It consists of two
   paragraphs of text, neither of which are
   particularly interesting.

   This is the second paragraph of the quotation.  It
   is just as dull as the first paragraph.
\end{quotation}
Another frequently-displayed structure is a list.
The following is an example of an \emph{itemized}
list.
\begin{itemize}
   \item This is the first item of an itemized list.
         Each item in the list is marked with a ``tick''.
         You don't have to worry about what kind of tick
         mark is used.

   \item This is the second item of the list.  It
         contains another list nested inside it.  The inner
         list is an \emph{enumerated} list.
         \begin{enumerate}
            \item This is the first item of an enumerated
                  list that is nested within the itemized list.

            \item This is the second item of the inner list.
                  \LaTeX\ allows you to nest lists deeper than
                  you really should.
         \end{enumerate}
         This is the rest of the second item of the outer
         list.  It is no more interesting than any other
         part of the item.
   \item This is the third item of the list.
\end{itemize}
You can even display poetry.
\begin{verse}
   There is an environment
    for verse \\             % The \\ command separates lines
   Whose features some poets % within a stanza.
   will curse.

                             % One or more blank lines separate stanzas.

   For instead of making\\
   Them do \emph{all} line breaking, \\
   It allows them to put too many words on a line when they'd rather be
   forced to be terse.
\end{verse}

Mathematical formulas may also be displayed.  A
displayed formula
is
one-line long; multiline
formulas require special formatting instructions.
   \[  \ip{\Gamma}{\psi'} = x'' + y^{2} + z_{i}^{n}\]
Don't start a paragraph with a displayed equation,
nor make one a paragraph by itself.

\end{document}               % End of document.

\section{Bench Section 223}
% synthetic bench section 223
% This is a sample LaTeX input file.  (Version of 12 August 2004.)
%
% A '%' character causes TeX to ignore all remaining text on the line,
% and is used for comments like this one.

\documentclass{article}      % Specifies the document class

                             % The preamble begins here.
\title{An Example Document}  % Declares the document's title.
\author{Leslie Lamport}      % Declares the author's name.
\date{January 21, 1994}      % Deleting this command produces today's date.

\newcommand{\ip}[2]{(#1, #2)}
                             % Defines \ip{arg1}{arg2} to mean
                             % (arg1, arg2).

%\newcommand{\ip}[2]{\langle #1 | #2\rangle}
                             % This is an alternative definition of
                             % \ip that is commented out.

\begin{document}             % End of preamble and beginning of text.

\maketitle                   % Produces the title.

This is an example input file.  Comparing it with
the output it generates can show you how to
produce a simple document of your own.

\section{Ordinary Text}      % Produces section heading.  Lower-level
                             % sections are begun with similar
                             % \subsection and \subsubsection commands.

The ends  of words and sentences are marked
  by   spaces. It  doesn't matter how many
spaces    you type; one is as good as 100.  The
end of   a line counts as a space.

One   or more   blank lines denote the  end
of  a paragraph.

Since any number of consecutive spaces are treated
like a single one, the formatting of the input
file makes no difference to
      \LaTeX,                % The \LaTeX command generates the LaTeX logo.
but it makes a difference to you.  When you use
\LaTeX, making your input file as easy to read
as possible will be a great help as you write
your document and when you change it.  This sample
file shows how you can add comments to your own input
file.

Because printing is different from typewriting,
there are a number of things that you have to do
differently when preparing an input file than if
you were just typing the document directly.
Quotation marks like
       ``this''
have to be handled specially, as do quotes within
quotes:
       ``\,`this'            % \, separates the double and single quote.
        is what I just
        wrote, not  `that'\,''.

Dashes come in three sizes: an
       intra-word
dash, a medium dash for number ranges like
       1--2,
and a punctuation
       dash---like
this.

A sentence-ending space should be larger than the
space between words within a sentence.  You
sometimes have to type special commands in
conjunction with punctuation characters to get
this right, as in the following sentence.
       Gnats, gnus, etc.\ all  % `\ ' makes an inter-word space.
       begin with G\@.         % \@ marks end-of-sentence punctuation.
You should check the spaces after periods when
reading your output to make sure you haven't
forgotten any special cases.  Generating an
ellipsis
       \ldots\               % `\ ' is needed after `\ldots' because TeX
                             % ignores spaces after command names like \ldots
                             % made from \ + letters.
                             %
                             % Note how a `%' character causes TeX to ignore
                             % the end of the input line, so these blank lines
                             % do not start a new paragraph.
                             %
with the right spacing around the periods requires
a special command.

\LaTeX\ interprets some common characters as
commands, so you must type special commands to
generate them.  These characters include the
following:
       \$ \& \% \# \{ and \}.

In printing, text is usually emphasized with an
       \emph{italic}
type style.

\begin{em}
   A long segment of text can also be emphasized
   in this way.  Text within such a segment can be
   given \emph{additional} emphasis.
\end{em}

It is sometimes necessary to prevent \LaTeX\ from
breaking a line where it might otherwise do so.
This may be at a space, as between the ``Mr.''\ and
``Jones'' in
       ``Mr.~Jones'',        % ~ produces an unbreakable interword space.
or within a word---especially when the word is a
symbol like
       \mbox{\emph{itemnum}}
that makes little sense when hyphenated across
lines.

Footnotes\footnote{This is an example of a footnote.}
pose no problem.

\LaTeX\ is good at typesetting mathematical formulas
like
       \( x-3y + z = 7 \)
or
       \( a_{1} > x^{2n} + y^{2n} > x' \)
or
       \( \ip{A}{B} = \sum_{i} a_{i} b_{i} \).
The spaces you type in a formula are
ignored.  Remember that a letter like
       $x$                   % $ ... $  and  \( ... \)  are equivalent
is a formula when it denotes a mathematical
symbol, and it should be typed as one.

\section{Displayed Text}

Text is displayed by indenting it from the left
margin.  Quotations are commonly displayed.  There
are short quotations
\begin{quote}
   This is a short quotation.  It consists of a
   single paragraph of text.  See how it is formatted.
\end{quote}
and longer ones.
\begin{quotation}
   This is a longer quotation.  It consists of two
   paragraphs of text, neither of which are
   particularly interesting.

   This is the second paragraph of the quotation.  It
   is just as dull as the first paragraph.
\end{quotation}
Another frequently-displayed structure is a list.
The following is an example of an \emph{itemized}
list.
\begin{itemize}
   \item This is the first item of an itemized list.
         Each item in the list is marked with a ``tick''.
         You don't have to worry about what kind of tick
         mark is used.

   \item This is the second item of the list.  It
         contains another list nested inside it.  The inner
         list is an \emph{enumerated} list.
         \begin{enumerate}
            \item This is the first item of an enumerated
                  list that is nested within the itemized list.

            \item This is the second item of the inner list.
                  \LaTeX\ allows you to nest lists deeper than
                  you really should.
         \end{enumerate}
         This is the rest of the second item of the outer
         list.  It is no more interesting than any other
         part of the item.
   \item This is the third item of the list.
\end{itemize}
You can even display poetry.
\begin{verse}
   There is an environment
    for verse \\             % The \\ command separates lines
   Whose features some poets % within a stanza.
   will curse.

                             % One or more blank lines separate stanzas.

   For instead of making\\
   Them do \emph{all} line breaking, \\
   It allows them to put too many words on a line when they'd rather be
   forced to be terse.
\end{verse}

Mathematical formulas may also be displayed.  A
displayed formula
is
one-line long; multiline
formulas require special formatting instructions.
   \[  \ip{\Gamma}{\psi'} = x'' + y^{2} + z_{i}^{n}\]
Don't start a paragraph with a displayed equation,
nor make one a paragraph by itself.

\end{document}               % End of document.

\section{Bench Section 224}
% synthetic bench section 224
% This is a sample LaTeX input file.  (Version of 12 August 2004.)
%
% A '%' character causes TeX to ignore all remaining text on the line,
% and is used for comments like this one.

\documentclass{article}      % Specifies the document class

                             % The preamble begins here.
\title{An Example Document}  % Declares the document's title.
\author{Leslie Lamport}      % Declares the author's name.
\date{January 21, 1994}      % Deleting this command produces today's date.

\newcommand{\ip}[2]{(#1, #2)}
                             % Defines \ip{arg1}{arg2} to mean
                             % (arg1, arg2).

%\newcommand{\ip}[2]{\langle #1 | #2\rangle}
                             % This is an alternative definition of
                             % \ip that is commented out.

\begin{document}             % End of preamble and beginning of text.

\maketitle                   % Produces the title.

This is an example input file.  Comparing it with
the output it generates can show you how to
produce a simple document of your own.

\section{Ordinary Text}      % Produces section heading.  Lower-level
                             % sections are begun with similar
                             % \subsection and \subsubsection commands.

The ends  of words and sentences are marked
  by   spaces. It  doesn't matter how many
spaces    you type; one is as good as 100.  The
end of   a line counts as a space.

One   or more   blank lines denote the  end
of  a paragraph.

Since any number of consecutive spaces are treated
like a single one, the formatting of the input
file makes no difference to
      \LaTeX,                % The \LaTeX command generates the LaTeX logo.
but it makes a difference to you.  When you use
\LaTeX, making your input file as easy to read
as possible will be a great help as you write
your document and when you change it.  This sample
file shows how you can add comments to your own input
file.

Because printing is different from typewriting,
there are a number of things that you have to do
differently when preparing an input file than if
you were just typing the document directly.
Quotation marks like
       ``this''
have to be handled specially, as do quotes within
quotes:
       ``\,`this'            % \, separates the double and single quote.
        is what I just
        wrote, not  `that'\,''.

Dashes come in three sizes: an
       intra-word
dash, a medium dash for number ranges like
       1--2,
and a punctuation
       dash---like
this.

A sentence-ending space should be larger than the
space between words within a sentence.  You
sometimes have to type special commands in
conjunction with punctuation characters to get
this right, as in the following sentence.
       Gnats, gnus, etc.\ all  % `\ ' makes an inter-word space.
       begin with G\@.         % \@ marks end-of-sentence punctuation.
You should check the spaces after periods when
reading your output to make sure you haven't
forgotten any special cases.  Generating an
ellipsis
       \ldots\               % `\ ' is needed after `\ldots' because TeX
                             % ignores spaces after command names like \ldots
                             % made from \ + letters.
                             %
                             % Note how a `%' character causes TeX to ignore
                             % the end of the input line, so these blank lines
                             % do not start a new paragraph.
                             %
with the right spacing around the periods requires
a special command.

\LaTeX\ interprets some common characters as
commands, so you must type special commands to
generate them.  These characters include the
following:
       \$ \& \% \# \{ and \}.

In printing, text is usually emphasized with an
       \emph{italic}
type style.

\begin{em}
   A long segment of text can also be emphasized
   in this way.  Text within such a segment can be
   given \emph{additional} emphasis.
\end{em}

It is sometimes necessary to prevent \LaTeX\ from
breaking a line where it might otherwise do so.
This may be at a space, as between the ``Mr.''\ and
``Jones'' in
       ``Mr.~Jones'',        % ~ produces an unbreakable interword space.
or within a word---especially when the word is a
symbol like
       \mbox{\emph{itemnum}}
that makes little sense when hyphenated across
lines.

Footnotes\footnote{This is an example of a footnote.}
pose no problem.

\LaTeX\ is good at typesetting mathematical formulas
like
       \( x-3y + z = 7 \)
or
       \( a_{1} > x^{2n} + y^{2n} > x' \)
or
       \( \ip{A}{B} = \sum_{i} a_{i} b_{i} \).
The spaces you type in a formula are
ignored.  Remember that a letter like
       $x$                   % $ ... $  and  \( ... \)  are equivalent
is a formula when it denotes a mathematical
symbol, and it should be typed as one.

\section{Displayed Text}

Text is displayed by indenting it from the left
margin.  Quotations are commonly displayed.  There
are short quotations
\begin{quote}
   This is a short quotation.  It consists of a
   single paragraph of text.  See how it is formatted.
\end{quote}
and longer ones.
\begin{quotation}
   This is a longer quotation.  It consists of two
   paragraphs of text, neither of which are
   particularly interesting.

   This is the second paragraph of the quotation.  It
   is just as dull as the first paragraph.
\end{quotation}
Another frequently-displayed structure is a list.
The following is an example of an \emph{itemized}
list.
\begin{itemize}
   \item This is the first item of an itemized list.
         Each item in the list is marked with a ``tick''.
         You don't have to worry about what kind of tick
         mark is used.

   \item This is the second item of the list.  It
         contains another list nested inside it.  The inner
         list is an \emph{enumerated} list.
         \begin{enumerate}
            \item This is the first item of an enumerated
                  list that is nested within the itemized list.

            \item This is the second item of the inner list.
                  \LaTeX\ allows you to nest lists deeper than
                  you really should.
         \end{enumerate}
         This is the rest of the second item of the outer
         list.  It is no more interesting than any other
         part of the item.
   \item This is the third item of the list.
\end{itemize}
You can even display poetry.
\begin{verse}
   There is an environment
    for verse \\             % The \\ command separates lines
   Whose features some poets % within a stanza.
   will curse.

                             % One or more blank lines separate stanzas.

   For instead of making\\
   Them do \emph{all} line breaking, \\
   It allows them to put too many words on a line when they'd rather be
   forced to be terse.
\end{verse}

Mathematical formulas may also be displayed.  A
displayed formula
is
one-line long; multiline
formulas require special formatting instructions.
   \[  \ip{\Gamma}{\psi'} = x'' + y^{2} + z_{i}^{n}\]
Don't start a paragraph with a displayed equation,
nor make one a paragraph by itself.

\end{document}               % End of document.

\section{Bench Section 225}
% synthetic bench section 225
% This is a sample LaTeX input file.  (Version of 12 August 2004.)
%
% A '%' character causes TeX to ignore all remaining text on the line,
% and is used for comments like this one.

\documentclass{article}      % Specifies the document class

                             % The preamble begins here.
\title{An Example Document}  % Declares the document's title.
\author{Leslie Lamport}      % Declares the author's name.
\date{January 21, 1994}      % Deleting this command produces today's date.

\newcommand{\ip}[2]{(#1, #2)}
                             % Defines \ip{arg1}{arg2} to mean
                             % (arg1, arg2).

%\newcommand{\ip}[2]{\langle #1 | #2\rangle}
                             % This is an alternative definition of
                             % \ip that is commented out.

\begin{document}             % End of preamble and beginning of text.

\maketitle                   % Produces the title.

This is an example input file.  Comparing it with
the output it generates can show you how to
produce a simple document of your own.

\section{Ordinary Text}      % Produces section heading.  Lower-level
                             % sections are begun with similar
                             % \subsection and \subsubsection commands.

The ends  of words and sentences are marked
  by   spaces. It  doesn't matter how many
spaces    you type; one is as good as 100.  The
end of   a line counts as a space.

One   or more   blank lines denote the  end
of  a paragraph.

Since any number of consecutive spaces are treated
like a single one, the formatting of the input
file makes no difference to
      \LaTeX,                % The \LaTeX command generates the LaTeX logo.
but it makes a difference to you.  When you use
\LaTeX, making your input file as easy to read
as possible will be a great help as you write
your document and when you change it.  This sample
file shows how you can add comments to your own input
file.

Because printing is different from typewriting,
there are a number of things that you have to do
differently when preparing an input file than if
you were just typing the document directly.
Quotation marks like
       ``this''
have to be handled specially, as do quotes within
quotes:
       ``\,`this'            % \, separates the double and single quote.
        is what I just
        wrote, not  `that'\,''.

Dashes come in three sizes: an
       intra-word
dash, a medium dash for number ranges like
       1--2,
and a punctuation
       dash---like
this.

A sentence-ending space should be larger than the
space between words within a sentence.  You
sometimes have to type special commands in
conjunction with punctuation characters to get
this right, as in the following sentence.
       Gnats, gnus, etc.\ all  % `\ ' makes an inter-word space.
       begin with G\@.         % \@ marks end-of-sentence punctuation.
You should check the spaces after periods when
reading your output to make sure you haven't
forgotten any special cases.  Generating an
ellipsis
       \ldots\               % `\ ' is needed after `\ldots' because TeX
                             % ignores spaces after command names like \ldots
                             % made from \ + letters.
                             %
                             % Note how a `%' character causes TeX to ignore
                             % the end of the input line, so these blank lines
                             % do not start a new paragraph.
                             %
with the right spacing around the periods requires
a special command.

\LaTeX\ interprets some common characters as
commands, so you must type special commands to
generate them.  These characters include the
following:
       \$ \& \% \# \{ and \}.

In printing, text is usually emphasized with an
       \emph{italic}
type style.

\begin{em}
   A long segment of text can also be emphasized
   in this way.  Text within such a segment can be
   given \emph{additional} emphasis.
\end{em}

It is sometimes necessary to prevent \LaTeX\ from
breaking a line where it might otherwise do so.
This may be at a space, as between the ``Mr.''\ and
``Jones'' in
       ``Mr.~Jones'',        % ~ produces an unbreakable interword space.
or within a word---especially when the word is a
symbol like
       \mbox{\emph{itemnum}}
that makes little sense when hyphenated across
lines.

Footnotes\footnote{This is an example of a footnote.}
pose no problem.

\LaTeX\ is good at typesetting mathematical formulas
like
       \( x-3y + z = 7 \)
or
       \( a_{1} > x^{2n} + y^{2n} > x' \)
or
       \( \ip{A}{B} = \sum_{i} a_{i} b_{i} \).
The spaces you type in a formula are
ignored.  Remember that a letter like
       $x$                   % $ ... $  and  \( ... \)  are equivalent
is a formula when it denotes a mathematical
symbol, and it should be typed as one.

\section{Displayed Text}

Text is displayed by indenting it from the left
margin.  Quotations are commonly displayed.  There
are short quotations
\begin{quote}
   This is a short quotation.  It consists of a
   single paragraph of text.  See how it is formatted.
\end{quote}
and longer ones.
\begin{quotation}
   This is a longer quotation.  It consists of two
   paragraphs of text, neither of which are
   particularly interesting.

   This is the second paragraph of the quotation.  It
   is just as dull as the first paragraph.
\end{quotation}
Another frequently-displayed structure is a list.
The following is an example of an \emph{itemized}
list.
\begin{itemize}
   \item This is the first item of an itemized list.
         Each item in the list is marked with a ``tick''.
         You don't have to worry about what kind of tick
         mark is used.

   \item This is the second item of the list.  It
         contains another list nested inside it.  The inner
         list is an \emph{enumerated} list.
         \begin{enumerate}
            \item This is the first item of an enumerated
                  list that is nested within the itemized list.

            \item This is the second item of the inner list.
                  \LaTeX\ allows you to nest lists deeper than
                  you really should.
         \end{enumerate}
         This is the rest of the second item of the outer
         list.  It is no more interesting than any other
         part of the item.
   \item This is the third item of the list.
\end{itemize}
You can even display poetry.
\begin{verse}
   There is an environment
    for verse \\             % The \\ command separates lines
   Whose features some poets % within a stanza.
   will curse.

                             % One or more blank lines separate stanzas.

   For instead of making\\
   Them do \emph{all} line breaking, \\
   It allows them to put too many words on a line when they'd rather be
   forced to be terse.
\end{verse}

Mathematical formulas may also be displayed.  A
displayed formula
is
one-line long; multiline
formulas require special formatting instructions.
   \[  \ip{\Gamma}{\psi'} = x'' + y^{2} + z_{i}^{n}\]
Don't start a paragraph with a displayed equation,
nor make one a paragraph by itself.

\end{document}               % End of document.

\section{Bench Section 226}
% synthetic bench section 226
% This is a sample LaTeX input file.  (Version of 12 August 2004.)
%
% A '%' character causes TeX to ignore all remaining text on the line,
% and is used for comments like this one.

\documentclass{article}      % Specifies the document class

                             % The preamble begins here.
\title{An Example Document}  % Declares the document's title.
\author{Leslie Lamport}      % Declares the author's name.
\date{January 21, 1994}      % Deleting this command produces today's date.

\newcommand{\ip}[2]{(#1, #2)}
                             % Defines \ip{arg1}{arg2} to mean
                             % (arg1, arg2).

%\newcommand{\ip}[2]{\langle #1 | #2\rangle}
                             % This is an alternative definition of
                             % \ip that is commented out.

\begin{document}             % End of preamble and beginning of text.

\maketitle                   % Produces the title.

This is an example input file.  Comparing it with
the output it generates can show you how to
produce a simple document of your own.

\section{Ordinary Text}      % Produces section heading.  Lower-level
                             % sections are begun with similar
                             % \subsection and \subsubsection commands.

The ends  of words and sentences are marked
  by   spaces. It  doesn't matter how many
spaces    you type; one is as good as 100.  The
end of   a line counts as a space.

One   or more   blank lines denote the  end
of  a paragraph.

Since any number of consecutive spaces are treated
like a single one, the formatting of the input
file makes no difference to
      \LaTeX,                % The \LaTeX command generates the LaTeX logo.
but it makes a difference to you.  When you use
\LaTeX, making your input file as easy to read
as possible will be a great help as you write
your document and when you change it.  This sample
file shows how you can add comments to your own input
file.

Because printing is different from typewriting,
there are a number of things that you have to do
differently when preparing an input file than if
you were just typing the document directly.
Quotation marks like
       ``this''
have to be handled specially, as do quotes within
quotes:
       ``\,`this'            % \, separates the double and single quote.
        is what I just
        wrote, not  `that'\,''.

Dashes come in three sizes: an
       intra-word
dash, a medium dash for number ranges like
       1--2,
and a punctuation
       dash---like
this.

A sentence-ending space should be larger than the
space between words within a sentence.  You
sometimes have to type special commands in
conjunction with punctuation characters to get
this right, as in the following sentence.
       Gnats, gnus, etc.\ all  % `\ ' makes an inter-word space.
       begin with G\@.         % \@ marks end-of-sentence punctuation.
You should check the spaces after periods when
reading your output to make sure you haven't
forgotten any special cases.  Generating an
ellipsis
       \ldots\               % `\ ' is needed after `\ldots' because TeX
                             % ignores spaces after command names like \ldots
                             % made from \ + letters.
                             %
                             % Note how a `%' character causes TeX to ignore
                             % the end of the input line, so these blank lines
                             % do not start a new paragraph.
                             %
with the right spacing around the periods requires
a special command.

\LaTeX\ interprets some common characters as
commands, so you must type special commands to
generate them.  These characters include the
following:
       \$ \& \% \# \{ and \}.

In printing, text is usually emphasized with an
       \emph{italic}
type style.

\begin{em}
   A long segment of text can also be emphasized
   in this way.  Text within such a segment can be
   given \emph{additional} emphasis.
\end{em}

It is sometimes necessary to prevent \LaTeX\ from
breaking a line where it might otherwise do so.
This may be at a space, as between the ``Mr.''\ and
``Jones'' in
       ``Mr.~Jones'',        % ~ produces an unbreakable interword space.
or within a word---especially when the word is a
symbol like
       \mbox{\emph{itemnum}}
that makes little sense when hyphenated across
lines.

Footnotes\footnote{This is an example of a footnote.}
pose no problem.

\LaTeX\ is good at typesetting mathematical formulas
like
       \( x-3y + z = 7 \)
or
       \( a_{1} > x^{2n} + y^{2n} > x' \)
or
       \( \ip{A}{B} = \sum_{i} a_{i} b_{i} \).
The spaces you type in a formula are
ignored.  Remember that a letter like
       $x$                   % $ ... $  and  \( ... \)  are equivalent
is a formula when it denotes a mathematical
symbol, and it should be typed as one.

\section{Displayed Text}

Text is displayed by indenting it from the left
margin.  Quotations are commonly displayed.  There
are short quotations
\begin{quote}
   This is a short quotation.  It consists of a
   single paragraph of text.  See how it is formatted.
\end{quote}
and longer ones.
\begin{quotation}
   This is a longer quotation.  It consists of two
   paragraphs of text, neither of which are
   particularly interesting.

   This is the second paragraph of the quotation.  It
   is just as dull as the first paragraph.
\end{quotation}
Another frequently-displayed structure is a list.
The following is an example of an \emph{itemized}
list.
\begin{itemize}
   \item This is the first item of an itemized list.
         Each item in the list is marked with a ``tick''.
         You don't have to worry about what kind of tick
         mark is used.

   \item This is the second item of the list.  It
         contains another list nested inside it.  The inner
         list is an \emph{enumerated} list.
         \begin{enumerate}
            \item This is the first item of an enumerated
                  list that is nested within the itemized list.

            \item This is the second item of the inner list.
                  \LaTeX\ allows you to nest lists deeper than
                  you really should.
         \end{enumerate}
         This is the rest of the second item of the outer
         list.  It is no more interesting than any other
         part of the item.
   \item This is the third item of the list.
\end{itemize}
You can even display poetry.
\begin{verse}
   There is an environment
    for verse \\             % The \\ command separates lines
   Whose features some poets % within a stanza.
   will curse.

                             % One or more blank lines separate stanzas.

   For instead of making\\
   Them do \emph{all} line breaking, \\
   It allows them to put too many words on a line when they'd rather be
   forced to be terse.
\end{verse}

Mathematical formulas may also be displayed.  A
displayed formula
is
one-line long; multiline
formulas require special formatting instructions.
   \[  \ip{\Gamma}{\psi'} = x'' + y^{2} + z_{i}^{n}\]
Don't start a paragraph with a displayed equation,
nor make one a paragraph by itself.

\end{document}               % End of document.

\section{Bench Section 227}
% synthetic bench section 227
% This is a sample LaTeX input file.  (Version of 12 August 2004.)
%
% A '%' character causes TeX to ignore all remaining text on the line,
% and is used for comments like this one.

\documentclass{article}      % Specifies the document class

                             % The preamble begins here.
\title{An Example Document}  % Declares the document's title.
\author{Leslie Lamport}      % Declares the author's name.
\date{January 21, 1994}      % Deleting this command produces today's date.

\newcommand{\ip}[2]{(#1, #2)}
                             % Defines \ip{arg1}{arg2} to mean
                             % (arg1, arg2).

%\newcommand{\ip}[2]{\langle #1 | #2\rangle}
                             % This is an alternative definition of
                             % \ip that is commented out.

\begin{document}             % End of preamble and beginning of text.

\maketitle                   % Produces the title.

This is an example input file.  Comparing it with
the output it generates can show you how to
produce a simple document of your own.

\section{Ordinary Text}      % Produces section heading.  Lower-level
                             % sections are begun with similar
                             % \subsection and \subsubsection commands.

The ends  of words and sentences are marked
  by   spaces. It  doesn't matter how many
spaces    you type; one is as good as 100.  The
end of   a line counts as a space.

One   or more   blank lines denote the  end
of  a paragraph.

Since any number of consecutive spaces are treated
like a single one, the formatting of the input
file makes no difference to
      \LaTeX,                % The \LaTeX command generates the LaTeX logo.
but it makes a difference to you.  When you use
\LaTeX, making your input file as easy to read
as possible will be a great help as you write
your document and when you change it.  This sample
file shows how you can add comments to your own input
file.

Because printing is different from typewriting,
there are a number of things that you have to do
differently when preparing an input file than if
you were just typing the document directly.
Quotation marks like
       ``this''
have to be handled specially, as do quotes within
quotes:
       ``\,`this'            % \, separates the double and single quote.
        is what I just
        wrote, not  `that'\,''.

Dashes come in three sizes: an
       intra-word
dash, a medium dash for number ranges like
       1--2,
and a punctuation
       dash---like
this.

A sentence-ending space should be larger than the
space between words within a sentence.  You
sometimes have to type special commands in
conjunction with punctuation characters to get
this right, as in the following sentence.
       Gnats, gnus, etc.\ all  % `\ ' makes an inter-word space.
       begin with G\@.         % \@ marks end-of-sentence punctuation.
You should check the spaces after periods when
reading your output to make sure you haven't
forgotten any special cases.  Generating an
ellipsis
       \ldots\               % `\ ' is needed after `\ldots' because TeX
                             % ignores spaces after command names like \ldots
                             % made from \ + letters.
                             %
                             % Note how a `%' character causes TeX to ignore
                             % the end of the input line, so these blank lines
                             % do not start a new paragraph.
                             %
with the right spacing around the periods requires
a special command.

\LaTeX\ interprets some common characters as
commands, so you must type special commands to
generate them.  These characters include the
following:
       \$ \& \% \# \{ and \}.

In printing, text is usually emphasized with an
       \emph{italic}
type style.

\begin{em}
   A long segment of text can also be emphasized
   in this way.  Text within such a segment can be
   given \emph{additional} emphasis.
\end{em}

It is sometimes necessary to prevent \LaTeX\ from
breaking a line where it might otherwise do so.
This may be at a space, as between the ``Mr.''\ and
``Jones'' in
       ``Mr.~Jones'',        % ~ produces an unbreakable interword space.
or within a word---especially when the word is a
symbol like
       \mbox{\emph{itemnum}}
that makes little sense when hyphenated across
lines.

Footnotes\footnote{This is an example of a footnote.}
pose no problem.

\LaTeX\ is good at typesetting mathematical formulas
like
       \( x-3y + z = 7 \)
or
       \( a_{1} > x^{2n} + y^{2n} > x' \)
or
       \( \ip{A}{B} = \sum_{i} a_{i} b_{i} \).
The spaces you type in a formula are
ignored.  Remember that a letter like
       $x$                   % $ ... $  and  \( ... \)  are equivalent
is a formula when it denotes a mathematical
symbol, and it should be typed as one.

\section{Displayed Text}

Text is displayed by indenting it from the left
margin.  Quotations are commonly displayed.  There
are short quotations
\begin{quote}
   This is a short quotation.  It consists of a
   single paragraph of text.  See how it is formatted.
\end{quote}
and longer ones.
\begin{quotation}
   This is a longer quotation.  It consists of two
   paragraphs of text, neither of which are
   particularly interesting.

   This is the second paragraph of the quotation.  It
   is just as dull as the first paragraph.
\end{quotation}
Another frequently-displayed structure is a list.
The following is an example of an \emph{itemized}
list.
\begin{itemize}
   \item This is the first item of an itemized list.
         Each item in the list is marked with a ``tick''.
         You don't have to worry about what kind of tick
         mark is used.

   \item This is the second item of the list.  It
         contains another list nested inside it.  The inner
         list is an \emph{enumerated} list.
         \begin{enumerate}
            \item This is the first item of an enumerated
                  list that is nested within the itemized list.

            \item This is the second item of the inner list.
                  \LaTeX\ allows you to nest lists deeper than
                  you really should.
         \end{enumerate}
         This is the rest of the second item of the outer
         list.  It is no more interesting than any other
         part of the item.
   \item This is the third item of the list.
\end{itemize}
You can even display poetry.
\begin{verse}
   There is an environment
    for verse \\             % The \\ command separates lines
   Whose features some poets % within a stanza.
   will curse.

                             % One or more blank lines separate stanzas.

   For instead of making\\
   Them do \emph{all} line breaking, \\
   It allows them to put too many words on a line when they'd rather be
   forced to be terse.
\end{verse}

Mathematical formulas may also be displayed.  A
displayed formula
is
one-line long; multiline
formulas require special formatting instructions.
   \[  \ip{\Gamma}{\psi'} = x'' + y^{2} + z_{i}^{n}\]
Don't start a paragraph with a displayed equation,
nor make one a paragraph by itself.

\end{document}               % End of document.

\section{Bench Section 228}
% synthetic bench section 228
% This is a sample LaTeX input file.  (Version of 12 August 2004.)
%
% A '%' character causes TeX to ignore all remaining text on the line,
% and is used for comments like this one.

\documentclass{article}      % Specifies the document class

                             % The preamble begins here.
\title{An Example Document}  % Declares the document's title.
\author{Leslie Lamport}      % Declares the author's name.
\date{January 21, 1994}      % Deleting this command produces today's date.

\newcommand{\ip}[2]{(#1, #2)}
                             % Defines \ip{arg1}{arg2} to mean
                             % (arg1, arg2).

%\newcommand{\ip}[2]{\langle #1 | #2\rangle}
                             % This is an alternative definition of
                             % \ip that is commented out.

\begin{document}             % End of preamble and beginning of text.

\maketitle                   % Produces the title.

This is an example input file.  Comparing it with
the output it generates can show you how to
produce a simple document of your own.

\section{Ordinary Text}      % Produces section heading.  Lower-level
                             % sections are begun with similar
                             % \subsection and \subsubsection commands.

The ends  of words and sentences are marked
  by   spaces. It  doesn't matter how many
spaces    you type; one is as good as 100.  The
end of   a line counts as a space.

One   or more   blank lines denote the  end
of  a paragraph.

Since any number of consecutive spaces are treated
like a single one, the formatting of the input
file makes no difference to
      \LaTeX,                % The \LaTeX command generates the LaTeX logo.
but it makes a difference to you.  When you use
\LaTeX, making your input file as easy to read
as possible will be a great help as you write
your document and when you change it.  This sample
file shows how you can add comments to your own input
file.

Because printing is different from typewriting,
there are a number of things that you have to do
differently when preparing an input file than if
you were just typing the document directly.
Quotation marks like
       ``this''
have to be handled specially, as do quotes within
quotes:
       ``\,`this'            % \, separates the double and single quote.
        is what I just
        wrote, not  `that'\,''.

Dashes come in three sizes: an
       intra-word
dash, a medium dash for number ranges like
       1--2,
and a punctuation
       dash---like
this.

A sentence-ending space should be larger than the
space between words within a sentence.  You
sometimes have to type special commands in
conjunction with punctuation characters to get
this right, as in the following sentence.
       Gnats, gnus, etc.\ all  % `\ ' makes an inter-word space.
       begin with G\@.         % \@ marks end-of-sentence punctuation.
You should check the spaces after periods when
reading your output to make sure you haven't
forgotten any special cases.  Generating an
ellipsis
       \ldots\               % `\ ' is needed after `\ldots' because TeX
                             % ignores spaces after command names like \ldots
                             % made from \ + letters.
                             %
                             % Note how a `%' character causes TeX to ignore
                             % the end of the input line, so these blank lines
                             % do not start a new paragraph.
                             %
with the right spacing around the periods requires
a special command.

\LaTeX\ interprets some common characters as
commands, so you must type special commands to
generate them.  These characters include the
following:
       \$ \& \% \# \{ and \}.

In printing, text is usually emphasized with an
       \emph{italic}
type style.

\begin{em}
   A long segment of text can also be emphasized
   in this way.  Text within such a segment can be
   given \emph{additional} emphasis.
\end{em}

It is sometimes necessary to prevent \LaTeX\ from
breaking a line where it might otherwise do so.
This may be at a space, as between the ``Mr.''\ and
``Jones'' in
       ``Mr.~Jones'',        % ~ produces an unbreakable interword space.
or within a word---especially when the word is a
symbol like
       \mbox{\emph{itemnum}}
that makes little sense when hyphenated across
lines.

Footnotes\footnote{This is an example of a footnote.}
pose no problem.

\LaTeX\ is good at typesetting mathematical formulas
like
       \( x-3y + z = 7 \)
or
       \( a_{1} > x^{2n} + y^{2n} > x' \)
or
       \( \ip{A}{B} = \sum_{i} a_{i} b_{i} \).
The spaces you type in a formula are
ignored.  Remember that a letter like
       $x$                   % $ ... $  and  \( ... \)  are equivalent
is a formula when it denotes a mathematical
symbol, and it should be typed as one.

\section{Displayed Text}

Text is displayed by indenting it from the left
margin.  Quotations are commonly displayed.  There
are short quotations
\begin{quote}
   This is a short quotation.  It consists of a
   single paragraph of text.  See how it is formatted.
\end{quote}
and longer ones.
\begin{quotation}
   This is a longer quotation.  It consists of two
   paragraphs of text, neither of which are
   particularly interesting.

   This is the second paragraph of the quotation.  It
   is just as dull as the first paragraph.
\end{quotation}
Another frequently-displayed structure is a list.
The following is an example of an \emph{itemized}
list.
\begin{itemize}
   \item This is the first item of an itemized list.
         Each item in the list is marked with a ``tick''.
         You don't have to worry about what kind of tick
         mark is used.

   \item This is the second item of the list.  It
         contains another list nested inside it.  The inner
         list is an \emph{enumerated} list.
         \begin{enumerate}
            \item This is the first item of an enumerated
                  list that is nested within the itemized list.

            \item This is the second item of the inner list.
                  \LaTeX\ allows you to nest lists deeper than
                  you really should.
         \end{enumerate}
         This is the rest of the second item of the outer
         list.  It is no more interesting than any other
         part of the item.
   \item This is the third item of the list.
\end{itemize}
You can even display poetry.
\begin{verse}
   There is an environment
    for verse \\             % The \\ command separates lines
   Whose features some poets % within a stanza.
   will curse.

                             % One or more blank lines separate stanzas.

   For instead of making\\
   Them do \emph{all} line breaking, \\
   It allows them to put too many words on a line when they'd rather be
   forced to be terse.
\end{verse}

Mathematical formulas may also be displayed.  A
displayed formula
is
one-line long; multiline
formulas require special formatting instructions.
   \[  \ip{\Gamma}{\psi'} = x'' + y^{2} + z_{i}^{n}\]
Don't start a paragraph with a displayed equation,
nor make one a paragraph by itself.

\end{document}               % End of document.

\section{Bench Section 229}
% synthetic bench section 229
% This is a sample LaTeX input file.  (Version of 12 August 2004.)
%
% A '%' character causes TeX to ignore all remaining text on the line,
% and is used for comments like this one.

\documentclass{article}      % Specifies the document class

                             % The preamble begins here.
\title{An Example Document}  % Declares the document's title.
\author{Leslie Lamport}      % Declares the author's name.
\date{January 21, 1994}      % Deleting this command produces today's date.

\newcommand{\ip}[2]{(#1, #2)}
                             % Defines \ip{arg1}{arg2} to mean
                             % (arg1, arg2).

%\newcommand{\ip}[2]{\langle #1 | #2\rangle}
                             % This is an alternative definition of
                             % \ip that is commented out.

\begin{document}             % End of preamble and beginning of text.

\maketitle                   % Produces the title.

This is an example input file.  Comparing it with
the output it generates can show you how to
produce a simple document of your own.

\section{Ordinary Text}      % Produces section heading.  Lower-level
                             % sections are begun with similar
                             % \subsection and \subsubsection commands.

The ends  of words and sentences are marked
  by   spaces. It  doesn't matter how many
spaces    you type; one is as good as 100.  The
end of   a line counts as a space.

One   or more   blank lines denote the  end
of  a paragraph.

Since any number of consecutive spaces are treated
like a single one, the formatting of the input
file makes no difference to
      \LaTeX,                % The \LaTeX command generates the LaTeX logo.
but it makes a difference to you.  When you use
\LaTeX, making your input file as easy to read
as possible will be a great help as you write
your document and when you change it.  This sample
file shows how you can add comments to your own input
file.

Because printing is different from typewriting,
there are a number of things that you have to do
differently when preparing an input file than if
you were just typing the document directly.
Quotation marks like
       ``this''
have to be handled specially, as do quotes within
quotes:
       ``\,`this'            % \, separates the double and single quote.
        is what I just
        wrote, not  `that'\,''.

Dashes come in three sizes: an
       intra-word
dash, a medium dash for number ranges like
       1--2,
and a punctuation
       dash---like
this.

A sentence-ending space should be larger than the
space between words within a sentence.  You
sometimes have to type special commands in
conjunction with punctuation characters to get
this right, as in the following sentence.
       Gnats, gnus, etc.\ all  % `\ ' makes an inter-word space.
       begin with G\@.         % \@ marks end-of-sentence punctuation.
You should check the spaces after periods when
reading your output to make sure you haven't
forgotten any special cases.  Generating an
ellipsis
       \ldots\               % `\ ' is needed after `\ldots' because TeX
                             % ignores spaces after command names like \ldots
                             % made from \ + letters.
                             %
                             % Note how a `%' character causes TeX to ignore
                             % the end of the input line, so these blank lines
                             % do not start a new paragraph.
                             %
with the right spacing around the periods requires
a special command.

\LaTeX\ interprets some common characters as
commands, so you must type special commands to
generate them.  These characters include the
following:
       \$ \& \% \# \{ and \}.

In printing, text is usually emphasized with an
       \emph{italic}
type style.

\begin{em}
   A long segment of text can also be emphasized
   in this way.  Text within such a segment can be
   given \emph{additional} emphasis.
\end{em}

It is sometimes necessary to prevent \LaTeX\ from
breaking a line where it might otherwise do so.
This may be at a space, as between the ``Mr.''\ and
``Jones'' in
       ``Mr.~Jones'',        % ~ produces an unbreakable interword space.
or within a word---especially when the word is a
symbol like
       \mbox{\emph{itemnum}}
that makes little sense when hyphenated across
lines.

Footnotes\footnote{This is an example of a footnote.}
pose no problem.

\LaTeX\ is good at typesetting mathematical formulas
like
       \( x-3y + z = 7 \)
or
       \( a_{1} > x^{2n} + y^{2n} > x' \)
or
       \( \ip{A}{B} = \sum_{i} a_{i} b_{i} \).
The spaces you type in a formula are
ignored.  Remember that a letter like
       $x$                   % $ ... $  and  \( ... \)  are equivalent
is a formula when it denotes a mathematical
symbol, and it should be typed as one.

\section{Displayed Text}

Text is displayed by indenting it from the left
margin.  Quotations are commonly displayed.  There
are short quotations
\begin{quote}
   This is a short quotation.  It consists of a
   single paragraph of text.  See how it is formatted.
\end{quote}
and longer ones.
\begin{quotation}
   This is a longer quotation.  It consists of two
   paragraphs of text, neither of which are
   particularly interesting.

   This is the second paragraph of the quotation.  It
   is just as dull as the first paragraph.
\end{quotation}
Another frequently-displayed structure is a list.
The following is an example of an \emph{itemized}
list.
\begin{itemize}
   \item This is the first item of an itemized list.
         Each item in the list is marked with a ``tick''.
         You don't have to worry about what kind of tick
         mark is used.

   \item This is the second item of the list.  It
         contains another list nested inside it.  The inner
         list is an \emph{enumerated} list.
         \begin{enumerate}
            \item This is the first item of an enumerated
                  list that is nested within the itemized list.

            \item This is the second item of the inner list.
                  \LaTeX\ allows you to nest lists deeper than
                  you really should.
         \end{enumerate}
         This is the rest of the second item of the outer
         list.  It is no more interesting than any other
         part of the item.
   \item This is the third item of the list.
\end{itemize}
You can even display poetry.
\begin{verse}
   There is an environment
    for verse \\             % The \\ command separates lines
   Whose features some poets % within a stanza.
   will curse.

                             % One or more blank lines separate stanzas.

   For instead of making\\
   Them do \emph{all} line breaking, \\
   It allows them to put too many words on a line when they'd rather be
   forced to be terse.
\end{verse}

Mathematical formulas may also be displayed.  A
displayed formula
is
one-line long; multiline
formulas require special formatting instructions.
   \[  \ip{\Gamma}{\psi'} = x'' + y^{2} + z_{i}^{n}\]
Don't start a paragraph with a displayed equation,
nor make one a paragraph by itself.

\end{document}               % End of document.

\section{Bench Section 230}
% synthetic bench section 230
% This is a sample LaTeX input file.  (Version of 12 August 2004.)
%
% A '%' character causes TeX to ignore all remaining text on the line,
% and is used for comments like this one.

\documentclass{article}      % Specifies the document class

                             % The preamble begins here.
\title{An Example Document}  % Declares the document's title.
\author{Leslie Lamport}      % Declares the author's name.
\date{January 21, 1994}      % Deleting this command produces today's date.

\newcommand{\ip}[2]{(#1, #2)}
                             % Defines \ip{arg1}{arg2} to mean
                             % (arg1, arg2).

%\newcommand{\ip}[2]{\langle #1 | #2\rangle}
                             % This is an alternative definition of
                             % \ip that is commented out.

\begin{document}             % End of preamble and beginning of text.

\maketitle                   % Produces the title.

This is an example input file.  Comparing it with
the output it generates can show you how to
produce a simple document of your own.

\section{Ordinary Text}      % Produces section heading.  Lower-level
                             % sections are begun with similar
                             % \subsection and \subsubsection commands.

The ends  of words and sentences are marked
  by   spaces. It  doesn't matter how many
spaces    you type; one is as good as 100.  The
end of   a line counts as a space.

One   or more   blank lines denote the  end
of  a paragraph.

Since any number of consecutive spaces are treated
like a single one, the formatting of the input
file makes no difference to
      \LaTeX,                % The \LaTeX command generates the LaTeX logo.
but it makes a difference to you.  When you use
\LaTeX, making your input file as easy to read
as possible will be a great help as you write
your document and when you change it.  This sample
file shows how you can add comments to your own input
file.

Because printing is different from typewriting,
there are a number of things that you have to do
differently when preparing an input file than if
you were just typing the document directly.
Quotation marks like
       ``this''
have to be handled specially, as do quotes within
quotes:
       ``\,`this'            % \, separates the double and single quote.
        is what I just
        wrote, not  `that'\,''.

Dashes come in three sizes: an
       intra-word
dash, a medium dash for number ranges like
       1--2,
and a punctuation
       dash---like
this.

A sentence-ending space should be larger than the
space between words within a sentence.  You
sometimes have to type special commands in
conjunction with punctuation characters to get
this right, as in the following sentence.
       Gnats, gnus, etc.\ all  % `\ ' makes an inter-word space.
       begin with G\@.         % \@ marks end-of-sentence punctuation.
You should check the spaces after periods when
reading your output to make sure you haven't
forgotten any special cases.  Generating an
ellipsis
       \ldots\               % `\ ' is needed after `\ldots' because TeX
                             % ignores spaces after command names like \ldots
                             % made from \ + letters.
                             %
                             % Note how a `%' character causes TeX to ignore
                             % the end of the input line, so these blank lines
                             % do not start a new paragraph.
                             %
with the right spacing around the periods requires
a special command.

\LaTeX\ interprets some common characters as
commands, so you must type special commands to
generate them.  These characters include the
following:
       \$ \& \% \# \{ and \}.

In printing, text is usually emphasized with an
       \emph{italic}
type style.

\begin{em}
   A long segment of text can also be emphasized
   in this way.  Text within such a segment can be
   given \emph{additional} emphasis.
\end{em}

It is sometimes necessary to prevent \LaTeX\ from
breaking a line where it might otherwise do so.
This may be at a space, as between the ``Mr.''\ and
``Jones'' in
       ``Mr.~Jones'',        % ~ produces an unbreakable interword space.
or within a word---especially when the word is a
symbol like
       \mbox{\emph{itemnum}}
that makes little sense when hyphenated across
lines.

Footnotes\footnote{This is an example of a footnote.}
pose no problem.

\LaTeX\ is good at typesetting mathematical formulas
like
       \( x-3y + z = 7 \)
or
       \( a_{1} > x^{2n} + y^{2n} > x' \)
or
       \( \ip{A}{B} = \sum_{i} a_{i} b_{i} \).
The spaces you type in a formula are
ignored.  Remember that a letter like
       $x$                   % $ ... $  and  \( ... \)  are equivalent
is a formula when it denotes a mathematical
symbol, and it should be typed as one.

\section{Displayed Text}

Text is displayed by indenting it from the left
margin.  Quotations are commonly displayed.  There
are short quotations
\begin{quote}
   This is a short quotation.  It consists of a
   single paragraph of text.  See how it is formatted.
\end{quote}
and longer ones.
\begin{quotation}
   This is a longer quotation.  It consists of two
   paragraphs of text, neither of which are
   particularly interesting.

   This is the second paragraph of the quotation.  It
   is just as dull as the first paragraph.
\end{quotation}
Another frequently-displayed structure is a list.
The following is an example of an \emph{itemized}
list.
\begin{itemize}
   \item This is the first item of an itemized list.
         Each item in the list is marked with a ``tick''.
         You don't have to worry about what kind of tick
         mark is used.

   \item This is the second item of the list.  It
         contains another list nested inside it.  The inner
         list is an \emph{enumerated} list.
         \begin{enumerate}
            \item This is the first item of an enumerated
                  list that is nested within the itemized list.

            \item This is the second item of the inner list.
                  \LaTeX\ allows you to nest lists deeper than
                  you really should.
         \end{enumerate}
         This is the rest of the second item of the outer
         list.  It is no more interesting than any other
         part of the item.
   \item This is the third item of the list.
\end{itemize}
You can even display poetry.
\begin{verse}
   There is an environment
    for verse \\             % The \\ command separates lines
   Whose features some poets % within a stanza.
   will curse.

                             % One or more blank lines separate stanzas.

   For instead of making\\
   Them do \emph{all} line breaking, \\
   It allows them to put too many words on a line when they'd rather be
   forced to be terse.
\end{verse}

Mathematical formulas may also be displayed.  A
displayed formula
is
one-line long; multiline
formulas require special formatting instructions.
   \[  \ip{\Gamma}{\psi'} = x'' + y^{2} + z_{i}^{n}\]
Don't start a paragraph with a displayed equation,
nor make one a paragraph by itself.

\end{document}               % End of document.

\section{Bench Section 231}
% synthetic bench section 231
% This is a sample LaTeX input file.  (Version of 12 August 2004.)
%
% A '%' character causes TeX to ignore all remaining text on the line,
% and is used for comments like this one.

\documentclass{article}      % Specifies the document class

                             % The preamble begins here.
\title{An Example Document}  % Declares the document's title.
\author{Leslie Lamport}      % Declares the author's name.
\date{January 21, 1994}      % Deleting this command produces today's date.

\newcommand{\ip}[2]{(#1, #2)}
                             % Defines \ip{arg1}{arg2} to mean
                             % (arg1, arg2).

%\newcommand{\ip}[2]{\langle #1 | #2\rangle}
                             % This is an alternative definition of
                             % \ip that is commented out.

\begin{document}             % End of preamble and beginning of text.

\maketitle                   % Produces the title.

This is an example input file.  Comparing it with
the output it generates can show you how to
produce a simple document of your own.

\section{Ordinary Text}      % Produces section heading.  Lower-level
                             % sections are begun with similar
                             % \subsection and \subsubsection commands.

The ends  of words and sentences are marked
  by   spaces. It  doesn't matter how many
spaces    you type; one is as good as 100.  The
end of   a line counts as a space.

One   or more   blank lines denote the  end
of  a paragraph.

Since any number of consecutive spaces are treated
like a single one, the formatting of the input
file makes no difference to
      \LaTeX,                % The \LaTeX command generates the LaTeX logo.
but it makes a difference to you.  When you use
\LaTeX, making your input file as easy to read
as possible will be a great help as you write
your document and when you change it.  This sample
file shows how you can add comments to your own input
file.

Because printing is different from typewriting,
there are a number of things that you have to do
differently when preparing an input file than if
you were just typing the document directly.
Quotation marks like
       ``this''
have to be handled specially, as do quotes within
quotes:
       ``\,`this'            % \, separates the double and single quote.
        is what I just
        wrote, not  `that'\,''.

Dashes come in three sizes: an
       intra-word
dash, a medium dash for number ranges like
       1--2,
and a punctuation
       dash---like
this.

A sentence-ending space should be larger than the
space between words within a sentence.  You
sometimes have to type special commands in
conjunction with punctuation characters to get
this right, as in the following sentence.
       Gnats, gnus, etc.\ all  % `\ ' makes an inter-word space.
       begin with G\@.         % \@ marks end-of-sentence punctuation.
You should check the spaces after periods when
reading your output to make sure you haven't
forgotten any special cases.  Generating an
ellipsis
       \ldots\               % `\ ' is needed after `\ldots' because TeX
                             % ignores spaces after command names like \ldots
                             % made from \ + letters.
                             %
                             % Note how a `%' character causes TeX to ignore
                             % the end of the input line, so these blank lines
                             % do not start a new paragraph.
                             %
with the right spacing around the periods requires
a special command.

\LaTeX\ interprets some common characters as
commands, so you must type special commands to
generate them.  These characters include the
following:
       \$ \& \% \# \{ and \}.

In printing, text is usually emphasized with an
       \emph{italic}
type style.

\begin{em}
   A long segment of text can also be emphasized
   in this way.  Text within such a segment can be
   given \emph{additional} emphasis.
\end{em}

It is sometimes necessary to prevent \LaTeX\ from
breaking a line where it might otherwise do so.
This may be at a space, as between the ``Mr.''\ and
``Jones'' in
       ``Mr.~Jones'',        % ~ produces an unbreakable interword space.
or within a word---especially when the word is a
symbol like
       \mbox{\emph{itemnum}}
that makes little sense when hyphenated across
lines.

Footnotes\footnote{This is an example of a footnote.}
pose no problem.

\LaTeX\ is good at typesetting mathematical formulas
like
       \( x-3y + z = 7 \)
or
       \( a_{1} > x^{2n} + y^{2n} > x' \)
or
       \( \ip{A}{B} = \sum_{i} a_{i} b_{i} \).
The spaces you type in a formula are
ignored.  Remember that a letter like
       $x$                   % $ ... $  and  \( ... \)  are equivalent
is a formula when it denotes a mathematical
symbol, and it should be typed as one.

\section{Displayed Text}

Text is displayed by indenting it from the left
margin.  Quotations are commonly displayed.  There
are short quotations
\begin{quote}
   This is a short quotation.  It consists of a
   single paragraph of text.  See how it is formatted.
\end{quote}
and longer ones.
\begin{quotation}
   This is a longer quotation.  It consists of two
   paragraphs of text, neither of which are
   particularly interesting.

   This is the second paragraph of the quotation.  It
   is just as dull as the first paragraph.
\end{quotation}
Another frequently-displayed structure is a list.
The following is an example of an \emph{itemized}
list.
\begin{itemize}
   \item This is the first item of an itemized list.
         Each item in the list is marked with a ``tick''.
         You don't have to worry about what kind of tick
         mark is used.

   \item This is the second item of the list.  It
         contains another list nested inside it.  The inner
         list is an \emph{enumerated} list.
         \begin{enumerate}
            \item This is the first item of an enumerated
                  list that is nested within the itemized list.

            \item This is the second item of the inner list.
                  \LaTeX\ allows you to nest lists deeper than
                  you really should.
         \end{enumerate}
         This is the rest of the second item of the outer
         list.  It is no more interesting than any other
         part of the item.
   \item This is the third item of the list.
\end{itemize}
You can even display poetry.
\begin{verse}
   There is an environment
    for verse \\             % The \\ command separates lines
   Whose features some poets % within a stanza.
   will curse.

                             % One or more blank lines separate stanzas.

   For instead of making\\
   Them do \emph{all} line breaking, \\
   It allows them to put too many words on a line when they'd rather be
   forced to be terse.
\end{verse}

Mathematical formulas may also be displayed.  A
displayed formula
is
one-line long; multiline
formulas require special formatting instructions.
   \[  \ip{\Gamma}{\psi'} = x'' + y^{2} + z_{i}^{n}\]
Don't start a paragraph with a displayed equation,
nor make one a paragraph by itself.

\end{document}               % End of document.

\section{Bench Section 232}
% synthetic bench section 232
% This is a sample LaTeX input file.  (Version of 12 August 2004.)
%
% A '%' character causes TeX to ignore all remaining text on the line,
% and is used for comments like this one.

\documentclass{article}      % Specifies the document class

                             % The preamble begins here.
\title{An Example Document}  % Declares the document's title.
\author{Leslie Lamport}      % Declares the author's name.
\date{January 21, 1994}      % Deleting this command produces today's date.

\newcommand{\ip}[2]{(#1, #2)}
                             % Defines \ip{arg1}{arg2} to mean
                             % (arg1, arg2).

%\newcommand{\ip}[2]{\langle #1 | #2\rangle}
                             % This is an alternative definition of
                             % \ip that is commented out.

\begin{document}             % End of preamble and beginning of text.

\maketitle                   % Produces the title.

This is an example input file.  Comparing it with
the output it generates can show you how to
produce a simple document of your own.

\section{Ordinary Text}      % Produces section heading.  Lower-level
                             % sections are begun with similar
                             % \subsection and \subsubsection commands.

The ends  of words and sentences are marked
  by   spaces. It  doesn't matter how many
spaces    you type; one is as good as 100.  The
end of   a line counts as a space.

One   or more   blank lines denote the  end
of  a paragraph.

Since any number of consecutive spaces are treated
like a single one, the formatting of the input
file makes no difference to
      \LaTeX,                % The \LaTeX command generates the LaTeX logo.
but it makes a difference to you.  When you use
\LaTeX, making your input file as easy to read
as possible will be a great help as you write
your document and when you change it.  This sample
file shows how you can add comments to your own input
file.

Because printing is different from typewriting,
there are a number of things that you have to do
differently when preparing an input file than if
you were just typing the document directly.
Quotation marks like
       ``this''
have to be handled specially, as do quotes within
quotes:
       ``\,`this'            % \, separates the double and single quote.
        is what I just
        wrote, not  `that'\,''.

Dashes come in three sizes: an
       intra-word
dash, a medium dash for number ranges like
       1--2,
and a punctuation
       dash---like
this.

A sentence-ending space should be larger than the
space between words within a sentence.  You
sometimes have to type special commands in
conjunction with punctuation characters to get
this right, as in the following sentence.
       Gnats, gnus, etc.\ all  % `\ ' makes an inter-word space.
       begin with G\@.         % \@ marks end-of-sentence punctuation.
You should check the spaces after periods when
reading your output to make sure you haven't
forgotten any special cases.  Generating an
ellipsis
       \ldots\               % `\ ' is needed after `\ldots' because TeX
                             % ignores spaces after command names like \ldots
                             % made from \ + letters.
                             %
                             % Note how a `%' character causes TeX to ignore
                             % the end of the input line, so these blank lines
                             % do not start a new paragraph.
                             %
with the right spacing around the periods requires
a special command.

\LaTeX\ interprets some common characters as
commands, so you must type special commands to
generate them.  These characters include the
following:
       \$ \& \% \# \{ and \}.

In printing, text is usually emphasized with an
       \emph{italic}
type style.

\begin{em}
   A long segment of text can also be emphasized
   in this way.  Text within such a segment can be
   given \emph{additional} emphasis.
\end{em}

It is sometimes necessary to prevent \LaTeX\ from
breaking a line where it might otherwise do so.
This may be at a space, as between the ``Mr.''\ and
``Jones'' in
       ``Mr.~Jones'',        % ~ produces an unbreakable interword space.
or within a word---especially when the word is a
symbol like
       \mbox{\emph{itemnum}}
that makes little sense when hyphenated across
lines.

Footnotes\footnote{This is an example of a footnote.}
pose no problem.

\LaTeX\ is good at typesetting mathematical formulas
like
       \( x-3y + z = 7 \)
or
       \( a_{1} > x^{2n} + y^{2n} > x' \)
or
       \( \ip{A}{B} = \sum_{i} a_{i} b_{i} \).
The spaces you type in a formula are
ignored.  Remember that a letter like
       $x$                   % $ ... $  and  \( ... \)  are equivalent
is a formula when it denotes a mathematical
symbol, and it should be typed as one.

\section{Displayed Text}

Text is displayed by indenting it from the left
margin.  Quotations are commonly displayed.  There
are short quotations
\begin{quote}
   This is a short quotation.  It consists of a
   single paragraph of text.  See how it is formatted.
\end{quote}
and longer ones.
\begin{quotation}
   This is a longer quotation.  It consists of two
   paragraphs of text, neither of which are
   particularly interesting.

   This is the second paragraph of the quotation.  It
   is just as dull as the first paragraph.
\end{quotation}
Another frequently-displayed structure is a list.
The following is an example of an \emph{itemized}
list.
\begin{itemize}
   \item This is the first item of an itemized list.
         Each item in the list is marked with a ``tick''.
         You don't have to worry about what kind of tick
         mark is used.

   \item This is the second item of the list.  It
         contains another list nested inside it.  The inner
         list is an \emph{enumerated} list.
         \begin{enumerate}
            \item This is the first item of an enumerated
                  list that is nested within the itemized list.

            \item This is the second item of the inner list.
                  \LaTeX\ allows you to nest lists deeper than
                  you really should.
         \end{enumerate}
         This is the rest of the second item of the outer
         list.  It is no more interesting than any other
         part of the item.
   \item This is the third item of the list.
\end{itemize}
You can even display poetry.
\begin{verse}
   There is an environment
    for verse \\             % The \\ command separates lines
   Whose features some poets % within a stanza.
   will curse.

                             % One or more blank lines separate stanzas.

   For instead of making\\
   Them do \emph{all} line breaking, \\
   It allows them to put too many words on a line when they'd rather be
   forced to be terse.
\end{verse}

Mathematical formulas may also be displayed.  A
displayed formula
is
one-line long; multiline
formulas require special formatting instructions.
   \[  \ip{\Gamma}{\psi'} = x'' + y^{2} + z_{i}^{n}\]
Don't start a paragraph with a displayed equation,
nor make one a paragraph by itself.

\end{document}               % End of document.

\section{Bench Section 233}
% synthetic bench section 233
% This is a sample LaTeX input file.  (Version of 12 August 2004.)
%
% A '%' character causes TeX to ignore all remaining text on the line,
% and is used for comments like this one.

\documentclass{article}      % Specifies the document class

                             % The preamble begins here.
\title{An Example Document}  % Declares the document's title.
\author{Leslie Lamport}      % Declares the author's name.
\date{January 21, 1994}      % Deleting this command produces today's date.

\newcommand{\ip}[2]{(#1, #2)}
                             % Defines \ip{arg1}{arg2} to mean
                             % (arg1, arg2).

%\newcommand{\ip}[2]{\langle #1 | #2\rangle}
                             % This is an alternative definition of
                             % \ip that is commented out.

\begin{document}             % End of preamble and beginning of text.

\maketitle                   % Produces the title.

This is an example input file.  Comparing it with
the output it generates can show you how to
produce a simple document of your own.

\section{Ordinary Text}      % Produces section heading.  Lower-level
                             % sections are begun with similar
                             % \subsection and \subsubsection commands.

The ends  of words and sentences are marked
  by   spaces. It  doesn't matter how many
spaces    you type; one is as good as 100.  The
end of   a line counts as a space.

One   or more   blank lines denote the  end
of  a paragraph.

Since any number of consecutive spaces are treated
like a single one, the formatting of the input
file makes no difference to
      \LaTeX,                % The \LaTeX command generates the LaTeX logo.
but it makes a difference to you.  When you use
\LaTeX, making your input file as easy to read
as possible will be a great help as you write
your document and when you change it.  This sample
file shows how you can add comments to your own input
file.

Because printing is different from typewriting,
there are a number of things that you have to do
differently when preparing an input file than if
you were just typing the document directly.
Quotation marks like
       ``this''
have to be handled specially, as do quotes within
quotes:
       ``\,`this'            % \, separates the double and single quote.
        is what I just
        wrote, not  `that'\,''.

Dashes come in three sizes: an
       intra-word
dash, a medium dash for number ranges like
       1--2,
and a punctuation
       dash---like
this.

A sentence-ending space should be larger than the
space between words within a sentence.  You
sometimes have to type special commands in
conjunction with punctuation characters to get
this right, as in the following sentence.
       Gnats, gnus, etc.\ all  % `\ ' makes an inter-word space.
       begin with G\@.         % \@ marks end-of-sentence punctuation.
You should check the spaces after periods when
reading your output to make sure you haven't
forgotten any special cases.  Generating an
ellipsis
       \ldots\               % `\ ' is needed after `\ldots' because TeX
                             % ignores spaces after command names like \ldots
                             % made from \ + letters.
                             %
                             % Note how a `%' character causes TeX to ignore
                             % the end of the input line, so these blank lines
                             % do not start a new paragraph.
                             %
with the right spacing around the periods requires
a special command.

\LaTeX\ interprets some common characters as
commands, so you must type special commands to
generate them.  These characters include the
following:
       \$ \& \% \# \{ and \}.

In printing, text is usually emphasized with an
       \emph{italic}
type style.

\begin{em}
   A long segment of text can also be emphasized
   in this way.  Text within such a segment can be
   given \emph{additional} emphasis.
\end{em}

It is sometimes necessary to prevent \LaTeX\ from
breaking a line where it might otherwise do so.
This may be at a space, as between the ``Mr.''\ and
``Jones'' in
       ``Mr.~Jones'',        % ~ produces an unbreakable interword space.
or within a word---especially when the word is a
symbol like
       \mbox{\emph{itemnum}}
that makes little sense when hyphenated across
lines.

Footnotes\footnote{This is an example of a footnote.}
pose no problem.

\LaTeX\ is good at typesetting mathematical formulas
like
       \( x-3y + z = 7 \)
or
       \( a_{1} > x^{2n} + y^{2n} > x' \)
or
       \( \ip{A}{B} = \sum_{i} a_{i} b_{i} \).
The spaces you type in a formula are
ignored.  Remember that a letter like
       $x$                   % $ ... $  and  \( ... \)  are equivalent
is a formula when it denotes a mathematical
symbol, and it should be typed as one.

\section{Displayed Text}

Text is displayed by indenting it from the left
margin.  Quotations are commonly displayed.  There
are short quotations
\begin{quote}
   This is a short quotation.  It consists of a
   single paragraph of text.  See how it is formatted.
\end{quote}
and longer ones.
\begin{quotation}
   This is a longer quotation.  It consists of two
   paragraphs of text, neither of which are
   particularly interesting.

   This is the second paragraph of the quotation.  It
   is just as dull as the first paragraph.
\end{quotation}
Another frequently-displayed structure is a list.
The following is an example of an \emph{itemized}
list.
\begin{itemize}
   \item This is the first item of an itemized list.
         Each item in the list is marked with a ``tick''.
         You don't have to worry about what kind of tick
         mark is used.

   \item This is the second item of the list.  It
         contains another list nested inside it.  The inner
         list is an \emph{enumerated} list.
         \begin{enumerate}
            \item This is the first item of an enumerated
                  list that is nested within the itemized list.

            \item This is the second item of the inner list.
                  \LaTeX\ allows you to nest lists deeper than
                  you really should.
         \end{enumerate}
         This is the rest of the second item of the outer
         list.  It is no more interesting than any other
         part of the item.
   \item This is the third item of the list.
\end{itemize}
You can even display poetry.
\begin{verse}
   There is an environment
    for verse \\             % The \\ command separates lines
   Whose features some poets % within a stanza.
   will curse.

                             % One or more blank lines separate stanzas.

   For instead of making\\
   Them do \emph{all} line breaking, \\
   It allows them to put too many words on a line when they'd rather be
   forced to be terse.
\end{verse}

Mathematical formulas may also be displayed.  A
displayed formula
is
one-line long; multiline
formulas require special formatting instructions.
   \[  \ip{\Gamma}{\psi'} = x'' + y^{2} + z_{i}^{n}\]
Don't start a paragraph with a displayed equation,
nor make one a paragraph by itself.

\end{document}               % End of document.

\section{Bench Section 234}
% synthetic bench section 234
% This is a sample LaTeX input file.  (Version of 12 August 2004.)
%
% A '%' character causes TeX to ignore all remaining text on the line,
% and is used for comments like this one.

\documentclass{article}      % Specifies the document class

                             % The preamble begins here.
\title{An Example Document}  % Declares the document's title.
\author{Leslie Lamport}      % Declares the author's name.
\date{January 21, 1994}      % Deleting this command produces today's date.

\newcommand{\ip}[2]{(#1, #2)}
                             % Defines \ip{arg1}{arg2} to mean
                             % (arg1, arg2).

%\newcommand{\ip}[2]{\langle #1 | #2\rangle}
                             % This is an alternative definition of
                             % \ip that is commented out.

\begin{document}             % End of preamble and beginning of text.

\maketitle                   % Produces the title.

This is an example input file.  Comparing it with
the output it generates can show you how to
produce a simple document of your own.

\section{Ordinary Text}      % Produces section heading.  Lower-level
                             % sections are begun with similar
                             % \subsection and \subsubsection commands.

The ends  of words and sentences are marked
  by   spaces. It  doesn't matter how many
spaces    you type; one is as good as 100.  The
end of   a line counts as a space.

One   or more   blank lines denote the  end
of  a paragraph.

Since any number of consecutive spaces are treated
like a single one, the formatting of the input
file makes no difference to
      \LaTeX,                % The \LaTeX command generates the LaTeX logo.
but it makes a difference to you.  When you use
\LaTeX, making your input file as easy to read
as possible will be a great help as you write
your document and when you change it.  This sample
file shows how you can add comments to your own input
file.

Because printing is different from typewriting,
there are a number of things that you have to do
differently when preparing an input file than if
you were just typing the document directly.
Quotation marks like
       ``this''
have to be handled specially, as do quotes within
quotes:
       ``\,`this'            % \, separates the double and single quote.
        is what I just
        wrote, not  `that'\,''.

Dashes come in three sizes: an
       intra-word
dash, a medium dash for number ranges like
       1--2,
and a punctuation
       dash---like
this.

A sentence-ending space should be larger than the
space between words within a sentence.  You
sometimes have to type special commands in
conjunction with punctuation characters to get
this right, as in the following sentence.
       Gnats, gnus, etc.\ all  % `\ ' makes an inter-word space.
       begin with G\@.         % \@ marks end-of-sentence punctuation.
You should check the spaces after periods when
reading your output to make sure you haven't
forgotten any special cases.  Generating an
ellipsis
       \ldots\               % `\ ' is needed after `\ldots' because TeX
                             % ignores spaces after command names like \ldots
                             % made from \ + letters.
                             %
                             % Note how a `%' character causes TeX to ignore
                             % the end of the input line, so these blank lines
                             % do not start a new paragraph.
                             %
with the right spacing around the periods requires
a special command.

\LaTeX\ interprets some common characters as
commands, so you must type special commands to
generate them.  These characters include the
following:
       \$ \& \% \# \{ and \}.

In printing, text is usually emphasized with an
       \emph{italic}
type style.

\begin{em}
   A long segment of text can also be emphasized
   in this way.  Text within such a segment can be
   given \emph{additional} emphasis.
\end{em}

It is sometimes necessary to prevent \LaTeX\ from
breaking a line where it might otherwise do so.
This may be at a space, as between the ``Mr.''\ and
``Jones'' in
       ``Mr.~Jones'',        % ~ produces an unbreakable interword space.
or within a word---especially when the word is a
symbol like
       \mbox{\emph{itemnum}}
that makes little sense when hyphenated across
lines.

Footnotes\footnote{This is an example of a footnote.}
pose no problem.

\LaTeX\ is good at typesetting mathematical formulas
like
       \( x-3y + z = 7 \)
or
       \( a_{1} > x^{2n} + y^{2n} > x' \)
or
       \( \ip{A}{B} = \sum_{i} a_{i} b_{i} \).
The spaces you type in a formula are
ignored.  Remember that a letter like
       $x$                   % $ ... $  and  \( ... \)  are equivalent
is a formula when it denotes a mathematical
symbol, and it should be typed as one.

\section{Displayed Text}

Text is displayed by indenting it from the left
margin.  Quotations are commonly displayed.  There
are short quotations
\begin{quote}
   This is a short quotation.  It consists of a
   single paragraph of text.  See how it is formatted.
\end{quote}
and longer ones.
\begin{quotation}
   This is a longer quotation.  It consists of two
   paragraphs of text, neither of which are
   particularly interesting.

   This is the second paragraph of the quotation.  It
   is just as dull as the first paragraph.
\end{quotation}
Another frequently-displayed structure is a list.
The following is an example of an \emph{itemized}
list.
\begin{itemize}
   \item This is the first item of an itemized list.
         Each item in the list is marked with a ``tick''.
         You don't have to worry about what kind of tick
         mark is used.

   \item This is the second item of the list.  It
         contains another list nested inside it.  The inner
         list is an \emph{enumerated} list.
         \begin{enumerate}
            \item This is the first item of an enumerated
                  list that is nested within the itemized list.

            \item This is the second item of the inner list.
                  \LaTeX\ allows you to nest lists deeper than
                  you really should.
         \end{enumerate}
         This is the rest of the second item of the outer
         list.  It is no more interesting than any other
         part of the item.
   \item This is the third item of the list.
\end{itemize}
You can even display poetry.
\begin{verse}
   There is an environment
    for verse \\             % The \\ command separates lines
   Whose features some poets % within a stanza.
   will curse.

                             % One or more blank lines separate stanzas.

   For instead of making\\
   Them do \emph{all} line breaking, \\
   It allows them to put too many words on a line when they'd rather be
   forced to be terse.
\end{verse}

Mathematical formulas may also be displayed.  A
displayed formula
is
one-line long; multiline
formulas require special formatting instructions.
   \[  \ip{\Gamma}{\psi'} = x'' + y^{2} + z_{i}^{n}\]
Don't start a paragraph with a displayed equation,
nor make one a paragraph by itself.

\end{document}               % End of document.

\section{Bench Section 235}
% synthetic bench section 235
% This is a sample LaTeX input file.  (Version of 12 August 2004.)
%
% A '%' character causes TeX to ignore all remaining text on the line,
% and is used for comments like this one.

\documentclass{article}      % Specifies the document class

                             % The preamble begins here.
\title{An Example Document}  % Declares the document's title.
\author{Leslie Lamport}      % Declares the author's name.
\date{January 21, 1994}      % Deleting this command produces today's date.

\newcommand{\ip}[2]{(#1, #2)}
                             % Defines \ip{arg1}{arg2} to mean
                             % (arg1, arg2).

%\newcommand{\ip}[2]{\langle #1 | #2\rangle}
                             % This is an alternative definition of
                             % \ip that is commented out.

\begin{document}             % End of preamble and beginning of text.

\maketitle                   % Produces the title.

This is an example input file.  Comparing it with
the output it generates can show you how to
produce a simple document of your own.

\section{Ordinary Text}      % Produces section heading.  Lower-level
                             % sections are begun with similar
                             % \subsection and \subsubsection commands.

The ends  of words and sentences are marked
  by   spaces. It  doesn't matter how many
spaces    you type; one is as good as 100.  The
end of   a line counts as a space.

One   or more   blank lines denote the  end
of  a paragraph.

Since any number of consecutive spaces are treated
like a single one, the formatting of the input
file makes no difference to
      \LaTeX,                % The \LaTeX command generates the LaTeX logo.
but it makes a difference to you.  When you use
\LaTeX, making your input file as easy to read
as possible will be a great help as you write
your document and when you change it.  This sample
file shows how you can add comments to your own input
file.

Because printing is different from typewriting,
there are a number of things that you have to do
differently when preparing an input file than if
you were just typing the document directly.
Quotation marks like
       ``this''
have to be handled specially, as do quotes within
quotes:
       ``\,`this'            % \, separates the double and single quote.
        is what I just
        wrote, not  `that'\,''.

Dashes come in three sizes: an
       intra-word
dash, a medium dash for number ranges like
       1--2,
and a punctuation
       dash---like
this.

A sentence-ending space should be larger than the
space between words within a sentence.  You
sometimes have to type special commands in
conjunction with punctuation characters to get
this right, as in the following sentence.
       Gnats, gnus, etc.\ all  % `\ ' makes an inter-word space.
       begin with G\@.         % \@ marks end-of-sentence punctuation.
You should check the spaces after periods when
reading your output to make sure you haven't
forgotten any special cases.  Generating an
ellipsis
       \ldots\               % `\ ' is needed after `\ldots' because TeX
                             % ignores spaces after command names like \ldots
                             % made from \ + letters.
                             %
                             % Note how a `%' character causes TeX to ignore
                             % the end of the input line, so these blank lines
                             % do not start a new paragraph.
                             %
with the right spacing around the periods requires
a special command.

\LaTeX\ interprets some common characters as
commands, so you must type special commands to
generate them.  These characters include the
following:
       \$ \& \% \# \{ and \}.

In printing, text is usually emphasized with an
       \emph{italic}
type style.

\begin{em}
   A long segment of text can also be emphasized
   in this way.  Text within such a segment can be
   given \emph{additional} emphasis.
\end{em}

It is sometimes necessary to prevent \LaTeX\ from
breaking a line where it might otherwise do so.
This may be at a space, as between the ``Mr.''\ and
``Jones'' in
       ``Mr.~Jones'',        % ~ produces an unbreakable interword space.
or within a word---especially when the word is a
symbol like
       \mbox{\emph{itemnum}}
that makes little sense when hyphenated across
lines.

Footnotes\footnote{This is an example of a footnote.}
pose no problem.

\LaTeX\ is good at typesetting mathematical formulas
like
       \( x-3y + z = 7 \)
or
       \( a_{1} > x^{2n} + y^{2n} > x' \)
or
       \( \ip{A}{B} = \sum_{i} a_{i} b_{i} \).
The spaces you type in a formula are
ignored.  Remember that a letter like
       $x$                   % $ ... $  and  \( ... \)  are equivalent
is a formula when it denotes a mathematical
symbol, and it should be typed as one.

\section{Displayed Text}

Text is displayed by indenting it from the left
margin.  Quotations are commonly displayed.  There
are short quotations
\begin{quote}
   This is a short quotation.  It consists of a
   single paragraph of text.  See how it is formatted.
\end{quote}
and longer ones.
\begin{quotation}
   This is a longer quotation.  It consists of two
   paragraphs of text, neither of which are
   particularly interesting.

   This is the second paragraph of the quotation.  It
   is just as dull as the first paragraph.
\end{quotation}
Another frequently-displayed structure is a list.
The following is an example of an \emph{itemized}
list.
\begin{itemize}
   \item This is the first item of an itemized list.
         Each item in the list is marked with a ``tick''.
         You don't have to worry about what kind of tick
         mark is used.

   \item This is the second item of the list.  It
         contains another list nested inside it.  The inner
         list is an \emph{enumerated} list.
         \begin{enumerate}
            \item This is the first item of an enumerated
                  list that is nested within the itemized list.

            \item This is the second item of the inner list.
                  \LaTeX\ allows you to nest lists deeper than
                  you really should.
         \end{enumerate}
         This is the rest of the second item of the outer
         list.  It is no more interesting than any other
         part of the item.
   \item This is the third item of the list.
\end{itemize}
You can even display poetry.
\begin{verse}
   There is an environment
    for verse \\             % The \\ command separates lines
   Whose features some poets % within a stanza.
   will curse.

                             % One or more blank lines separate stanzas.

   For instead of making\\
   Them do \emph{all} line breaking, \\
   It allows them to put too many words on a line when they'd rather be
   forced to be terse.
\end{verse}

Mathematical formulas may also be displayed.  A
displayed formula
is
one-line long; multiline
formulas require special formatting instructions.
   \[  \ip{\Gamma}{\psi'} = x'' + y^{2} + z_{i}^{n}\]
Don't start a paragraph with a displayed equation,
nor make one a paragraph by itself.

\end{document}               % End of document.

\section{Bench Section 236}
% synthetic bench section 236
% This is a sample LaTeX input file.  (Version of 12 August 2004.)
%
% A '%' character causes TeX to ignore all remaining text on the line,
% and is used for comments like this one.

\documentclass{article}      % Specifies the document class

                             % The preamble begins here.
\title{An Example Document}  % Declares the document's title.
\author{Leslie Lamport}      % Declares the author's name.
\date{January 21, 1994}      % Deleting this command produces today's date.

\newcommand{\ip}[2]{(#1, #2)}
                             % Defines \ip{arg1}{arg2} to mean
                             % (arg1, arg2).

%\newcommand{\ip}[2]{\langle #1 | #2\rangle}
                             % This is an alternative definition of
                             % \ip that is commented out.

\begin{document}             % End of preamble and beginning of text.

\maketitle                   % Produces the title.

This is an example input file.  Comparing it with
the output it generates can show you how to
produce a simple document of your own.

\section{Ordinary Text}      % Produces section heading.  Lower-level
                             % sections are begun with similar
                             % \subsection and \subsubsection commands.

The ends  of words and sentences are marked
  by   spaces. It  doesn't matter how many
spaces    you type; one is as good as 100.  The
end of   a line counts as a space.

One   or more   blank lines denote the  end
of  a paragraph.

Since any number of consecutive spaces are treated
like a single one, the formatting of the input
file makes no difference to
      \LaTeX,                % The \LaTeX command generates the LaTeX logo.
but it makes a difference to you.  When you use
\LaTeX, making your input file as easy to read
as possible will be a great help as you write
your document and when you change it.  This sample
file shows how you can add comments to your own input
file.

Because printing is different from typewriting,
there are a number of things that you have to do
differently when preparing an input file than if
you were just typing the document directly.
Quotation marks like
       ``this''
have to be handled specially, as do quotes within
quotes:
       ``\,`this'            % \, separates the double and single quote.
        is what I just
        wrote, not  `that'\,''.

Dashes come in three sizes: an
       intra-word
dash, a medium dash for number ranges like
       1--2,
and a punctuation
       dash---like
this.

A sentence-ending space should be larger than the
space between words within a sentence.  You
sometimes have to type special commands in
conjunction with punctuation characters to get
this right, as in the following sentence.
       Gnats, gnus, etc.\ all  % `\ ' makes an inter-word space.
       begin with G\@.         % \@ marks end-of-sentence punctuation.
You should check the spaces after periods when
reading your output to make sure you haven't
forgotten any special cases.  Generating an
ellipsis
       \ldots\               % `\ ' is needed after `\ldots' because TeX
                             % ignores spaces after command names like \ldots
                             % made from \ + letters.
                             %
                             % Note how a `%' character causes TeX to ignore
                             % the end of the input line, so these blank lines
                             % do not start a new paragraph.
                             %
with the right spacing around the periods requires
a special command.

\LaTeX\ interprets some common characters as
commands, so you must type special commands to
generate them.  These characters include the
following:
       \$ \& \% \# \{ and \}.

In printing, text is usually emphasized with an
       \emph{italic}
type style.

\begin{em}
   A long segment of text can also be emphasized
   in this way.  Text within such a segment can be
   given \emph{additional} emphasis.
\end{em}

It is sometimes necessary to prevent \LaTeX\ from
breaking a line where it might otherwise do so.
This may be at a space, as between the ``Mr.''\ and
``Jones'' in
       ``Mr.~Jones'',        % ~ produces an unbreakable interword space.
or within a word---especially when the word is a
symbol like
       \mbox{\emph{itemnum}}
that makes little sense when hyphenated across
lines.

Footnotes\footnote{This is an example of a footnote.}
pose no problem.

\LaTeX\ is good at typesetting mathematical formulas
like
       \( x-3y + z = 7 \)
or
       \( a_{1} > x^{2n} + y^{2n} > x' \)
or
       \( \ip{A}{B} = \sum_{i} a_{i} b_{i} \).
The spaces you type in a formula are
ignored.  Remember that a letter like
       $x$                   % $ ... $  and  \( ... \)  are equivalent
is a formula when it denotes a mathematical
symbol, and it should be typed as one.

\section{Displayed Text}

Text is displayed by indenting it from the left
margin.  Quotations are commonly displayed.  There
are short quotations
\begin{quote}
   This is a short quotation.  It consists of a
   single paragraph of text.  See how it is formatted.
\end{quote}
and longer ones.
\begin{quotation}
   This is a longer quotation.  It consists of two
   paragraphs of text, neither of which are
   particularly interesting.

   This is the second paragraph of the quotation.  It
   is just as dull as the first paragraph.
\end{quotation}
Another frequently-displayed structure is a list.
The following is an example of an \emph{itemized}
list.
\begin{itemize}
   \item This is the first item of an itemized list.
         Each item in the list is marked with a ``tick''.
         You don't have to worry about what kind of tick
         mark is used.

   \item This is the second item of the list.  It
         contains another list nested inside it.  The inner
         list is an \emph{enumerated} list.
         \begin{enumerate}
            \item This is the first item of an enumerated
                  list that is nested within the itemized list.

            \item This is the second item of the inner list.
                  \LaTeX\ allows you to nest lists deeper than
                  you really should.
         \end{enumerate}
         This is the rest of the second item of the outer
         list.  It is no more interesting than any other
         part of the item.
   \item This is the third item of the list.
\end{itemize}
You can even display poetry.
\begin{verse}
   There is an environment
    for verse \\             % The \\ command separates lines
   Whose features some poets % within a stanza.
   will curse.

                             % One or more blank lines separate stanzas.

   For instead of making\\
   Them do \emph{all} line breaking, \\
   It allows them to put too many words on a line when they'd rather be
   forced to be terse.
\end{verse}

Mathematical formulas may also be displayed.  A
displayed formula
is
one-line long; multiline
formulas require special formatting instructions.
   \[  \ip{\Gamma}{\psi'} = x'' + y^{2} + z_{i}^{n}\]
Don't start a paragraph with a displayed equation,
nor make one a paragraph by itself.

\end{document}               % End of document.

\section{Bench Section 237}
% synthetic bench section 237
% This is a sample LaTeX input file.  (Version of 12 August 2004.)
%
% A '%' character causes TeX to ignore all remaining text on the line,
% and is used for comments like this one.

\documentclass{article}      % Specifies the document class

                             % The preamble begins here.
\title{An Example Document}  % Declares the document's title.
\author{Leslie Lamport}      % Declares the author's name.
\date{January 21, 1994}      % Deleting this command produces today's date.

\newcommand{\ip}[2]{(#1, #2)}
                             % Defines \ip{arg1}{arg2} to mean
                             % (arg1, arg2).

%\newcommand{\ip}[2]{\langle #1 | #2\rangle}
                             % This is an alternative definition of
                             % \ip that is commented out.

\begin{document}             % End of preamble and beginning of text.

\maketitle                   % Produces the title.

This is an example input file.  Comparing it with
the output it generates can show you how to
produce a simple document of your own.

\section{Ordinary Text}      % Produces section heading.  Lower-level
                             % sections are begun with similar
                             % \subsection and \subsubsection commands.

The ends  of words and sentences are marked
  by   spaces. It  doesn't matter how many
spaces    you type; one is as good as 100.  The
end of   a line counts as a space.

One   or more   blank lines denote the  end
of  a paragraph.

Since any number of consecutive spaces are treated
like a single one, the formatting of the input
file makes no difference to
      \LaTeX,                % The \LaTeX command generates the LaTeX logo.
but it makes a difference to you.  When you use
\LaTeX, making your input file as easy to read
as possible will be a great help as you write
your document and when you change it.  This sample
file shows how you can add comments to your own input
file.

Because printing is different from typewriting,
there are a number of things that you have to do
differently when preparing an input file than if
you were just typing the document directly.
Quotation marks like
       ``this''
have to be handled specially, as do quotes within
quotes:
       ``\,`this'            % \, separates the double and single quote.
        is what I just
        wrote, not  `that'\,''.

Dashes come in three sizes: an
       intra-word
dash, a medium dash for number ranges like
       1--2,
and a punctuation
       dash---like
this.

A sentence-ending space should be larger than the
space between words within a sentence.  You
sometimes have to type special commands in
conjunction with punctuation characters to get
this right, as in the following sentence.
       Gnats, gnus, etc.\ all  % `\ ' makes an inter-word space.
       begin with G\@.         % \@ marks end-of-sentence punctuation.
You should check the spaces after periods when
reading your output to make sure you haven't
forgotten any special cases.  Generating an
ellipsis
       \ldots\               % `\ ' is needed after `\ldots' because TeX
                             % ignores spaces after command names like \ldots
                             % made from \ + letters.
                             %
                             % Note how a `%' character causes TeX to ignore
                             % the end of the input line, so these blank lines
                             % do not start a new paragraph.
                             %
with the right spacing around the periods requires
a special command.

\LaTeX\ interprets some common characters as
commands, so you must type special commands to
generate them.  These characters include the
following:
       \$ \& \% \# \{ and \}.

In printing, text is usually emphasized with an
       \emph{italic}
type style.

\begin{em}
   A long segment of text can also be emphasized
   in this way.  Text within such a segment can be
   given \emph{additional} emphasis.
\end{em}

It is sometimes necessary to prevent \LaTeX\ from
breaking a line where it might otherwise do so.
This may be at a space, as between the ``Mr.''\ and
``Jones'' in
       ``Mr.~Jones'',        % ~ produces an unbreakable interword space.
or within a word---especially when the word is a
symbol like
       \mbox{\emph{itemnum}}
that makes little sense when hyphenated across
lines.

Footnotes\footnote{This is an example of a footnote.}
pose no problem.

\LaTeX\ is good at typesetting mathematical formulas
like
       \( x-3y + z = 7 \)
or
       \( a_{1} > x^{2n} + y^{2n} > x' \)
or
       \( \ip{A}{B} = \sum_{i} a_{i} b_{i} \).
The spaces you type in a formula are
ignored.  Remember that a letter like
       $x$                   % $ ... $  and  \( ... \)  are equivalent
is a formula when it denotes a mathematical
symbol, and it should be typed as one.

\section{Displayed Text}

Text is displayed by indenting it from the left
margin.  Quotations are commonly displayed.  There
are short quotations
\begin{quote}
   This is a short quotation.  It consists of a
   single paragraph of text.  See how it is formatted.
\end{quote}
and longer ones.
\begin{quotation}
   This is a longer quotation.  It consists of two
   paragraphs of text, neither of which are
   particularly interesting.

   This is the second paragraph of the quotation.  It
   is just as dull as the first paragraph.
\end{quotation}
Another frequently-displayed structure is a list.
The following is an example of an \emph{itemized}
list.
\begin{itemize}
   \item This is the first item of an itemized list.
         Each item in the list is marked with a ``tick''.
         You don't have to worry about what kind of tick
         mark is used.

   \item This is the second item of the list.  It
         contains another list nested inside it.  The inner
         list is an \emph{enumerated} list.
         \begin{enumerate}
            \item This is the first item of an enumerated
                  list that is nested within the itemized list.

            \item This is the second item of the inner list.
                  \LaTeX\ allows you to nest lists deeper than
                  you really should.
         \end{enumerate}
         This is the rest of the second item of the outer
         list.  It is no more interesting than any other
         part of the item.
   \item This is the third item of the list.
\end{itemize}
You can even display poetry.
\begin{verse}
   There is an environment
    for verse \\             % The \\ command separates lines
   Whose features some poets % within a stanza.
   will curse.

                             % One or more blank lines separate stanzas.

   For instead of making\\
   Them do \emph{all} line breaking, \\
   It allows them to put too many words on a line when they'd rather be
   forced to be terse.
\end{verse}

Mathematical formulas may also be displayed.  A
displayed formula
is
one-line long; multiline
formulas require special formatting instructions.
   \[  \ip{\Gamma}{\psi'} = x'' + y^{2} + z_{i}^{n}\]
Don't start a paragraph with a displayed equation,
nor make one a paragraph by itself.

\end{document}               % End of document.

\section{Bench Section 238}
% synthetic bench section 238
% This is a sample LaTeX input file.  (Version of 12 August 2004.)
%
% A '%' character causes TeX to ignore all remaining text on the line,
% and is used for comments like this one.

\documentclass{article}      % Specifies the document class

                             % The preamble begins here.
\title{An Example Document}  % Declares the document's title.
\author{Leslie Lamport}      % Declares the author's name.
\date{January 21, 1994}      % Deleting this command produces today's date.

\newcommand{\ip}[2]{(#1, #2)}
                             % Defines \ip{arg1}{arg2} to mean
                             % (arg1, arg2).

%\newcommand{\ip}[2]{\langle #1 | #2\rangle}
                             % This is an alternative definition of
                             % \ip that is commented out.

\begin{document}             % End of preamble and beginning of text.

\maketitle                   % Produces the title.

This is an example input file.  Comparing it with
the output it generates can show you how to
produce a simple document of your own.

\section{Ordinary Text}      % Produces section heading.  Lower-level
                             % sections are begun with similar
                             % \subsection and \subsubsection commands.

The ends  of words and sentences are marked
  by   spaces. It  doesn't matter how many
spaces    you type; one is as good as 100.  The
end of   a line counts as a space.

One   or more   blank lines denote the  end
of  a paragraph.

Since any number of consecutive spaces are treated
like a single one, the formatting of the input
file makes no difference to
      \LaTeX,                % The \LaTeX command generates the LaTeX logo.
but it makes a difference to you.  When you use
\LaTeX, making your input file as easy to read
as possible will be a great help as you write
your document and when you change it.  This sample
file shows how you can add comments to your own input
file.

Because printing is different from typewriting,
there are a number of things that you have to do
differently when preparing an input file than if
you were just typing the document directly.
Quotation marks like
       ``this''
have to be handled specially, as do quotes within
quotes:
       ``\,`this'            % \, separates the double and single quote.
        is what I just
        wrote, not  `that'\,''.

Dashes come in three sizes: an
       intra-word
dash, a medium dash for number ranges like
       1--2,
and a punctuation
       dash---like
this.

A sentence-ending space should be larger than the
space between words within a sentence.  You
sometimes have to type special commands in
conjunction with punctuation characters to get
this right, as in the following sentence.
       Gnats, gnus, etc.\ all  % `\ ' makes an inter-word space.
       begin with G\@.         % \@ marks end-of-sentence punctuation.
You should check the spaces after periods when
reading your output to make sure you haven't
forgotten any special cases.  Generating an
ellipsis
       \ldots\               % `\ ' is needed after `\ldots' because TeX
                             % ignores spaces after command names like \ldots
                             % made from \ + letters.
                             %
                             % Note how a `%' character causes TeX to ignore
                             % the end of the input line, so these blank lines
                             % do not start a new paragraph.
                             %
with the right spacing around the periods requires
a special command.

\LaTeX\ interprets some common characters as
commands, so you must type special commands to
generate them.  These characters include the
following:
       \$ \& \% \# \{ and \}.

In printing, text is usually emphasized with an
       \emph{italic}
type style.

\begin{em}
   A long segment of text can also be emphasized
   in this way.  Text within such a segment can be
   given \emph{additional} emphasis.
\end{em}

It is sometimes necessary to prevent \LaTeX\ from
breaking a line where it might otherwise do so.
This may be at a space, as between the ``Mr.''\ and
``Jones'' in
       ``Mr.~Jones'',        % ~ produces an unbreakable interword space.
or within a word---especially when the word is a
symbol like
       \mbox{\emph{itemnum}}
that makes little sense when hyphenated across
lines.

Footnotes\footnote{This is an example of a footnote.}
pose no problem.

\LaTeX\ is good at typesetting mathematical formulas
like
       \( x-3y + z = 7 \)
or
       \( a_{1} > x^{2n} + y^{2n} > x' \)
or
       \( \ip{A}{B} = \sum_{i} a_{i} b_{i} \).
The spaces you type in a formula are
ignored.  Remember that a letter like
       $x$                   % $ ... $  and  \( ... \)  are equivalent
is a formula when it denotes a mathematical
symbol, and it should be typed as one.

\section{Displayed Text}

Text is displayed by indenting it from the left
margin.  Quotations are commonly displayed.  There
are short quotations
\begin{quote}
   This is a short quotation.  It consists of a
   single paragraph of text.  See how it is formatted.
\end{quote}
and longer ones.
\begin{quotation}
   This is a longer quotation.  It consists of two
   paragraphs of text, neither of which are
   particularly interesting.

   This is the second paragraph of the quotation.  It
   is just as dull as the first paragraph.
\end{quotation}
Another frequently-displayed structure is a list.
The following is an example of an \emph{itemized}
list.
\begin{itemize}
   \item This is the first item of an itemized list.
         Each item in the list is marked with a ``tick''.
         You don't have to worry about what kind of tick
         mark is used.

   \item This is the second item of the list.  It
         contains another list nested inside it.  The inner
         list is an \emph{enumerated} list.
         \begin{enumerate}
            \item This is the first item of an enumerated
                  list that is nested within the itemized list.

            \item This is the second item of the inner list.
                  \LaTeX\ allows you to nest lists deeper than
                  you really should.
         \end{enumerate}
         This is the rest of the second item of the outer
         list.  It is no more interesting than any other
         part of the item.
   \item This is the third item of the list.
\end{itemize}
You can even display poetry.
\begin{verse}
   There is an environment
    for verse \\             % The \\ command separates lines
   Whose features some poets % within a stanza.
   will curse.

                             % One or more blank lines separate stanzas.

   For instead of making\\
   Them do \emph{all} line breaking, \\
   It allows them to put too many words on a line when they'd rather be
   forced to be terse.
\end{verse}

Mathematical formulas may also be displayed.  A
displayed formula
is
one-line long; multiline
formulas require special formatting instructions.
   \[  \ip{\Gamma}{\psi'} = x'' + y^{2} + z_{i}^{n}\]
Don't start a paragraph with a displayed equation,
nor make one a paragraph by itself.

\end{document}               % End of document.

\section{Bench Section 239}
% synthetic bench section 239
% This is a sample LaTeX input file.  (Version of 12 August 2004.)
%
% A '%' character causes TeX to ignore all remaining text on the line,
% and is used for comments like this one.

\documentclass{article}      % Specifies the document class

                             % The preamble begins here.
\title{An Example Document}  % Declares the document's title.
\author{Leslie Lamport}      % Declares the author's name.
\date{January 21, 1994}      % Deleting this command produces today's date.

\newcommand{\ip}[2]{(#1, #2)}
                             % Defines \ip{arg1}{arg2} to mean
                             % (arg1, arg2).

%\newcommand{\ip}[2]{\langle #1 | #2\rangle}
                             % This is an alternative definition of
                             % \ip that is commented out.

\begin{document}             % End of preamble and beginning of text.

\maketitle                   % Produces the title.

This is an example input file.  Comparing it with
the output it generates can show you how to
produce a simple document of your own.

\section{Ordinary Text}      % Produces section heading.  Lower-level
                             % sections are begun with similar
                             % \subsection and \subsubsection commands.

The ends  of words and sentences are marked
  by   spaces. It  doesn't matter how many
spaces    you type; one is as good as 100.  The
end of   a line counts as a space.

One   or more   blank lines denote the  end
of  a paragraph.

Since any number of consecutive spaces are treated
like a single one, the formatting of the input
file makes no difference to
      \LaTeX,                % The \LaTeX command generates the LaTeX logo.
but it makes a difference to you.  When you use
\LaTeX, making your input file as easy to read
as possible will be a great help as you write
your document and when you change it.  This sample
file shows how you can add comments to your own input
file.

Because printing is different from typewriting,
there are a number of things that you have to do
differently when preparing an input file than if
you were just typing the document directly.
Quotation marks like
       ``this''
have to be handled specially, as do quotes within
quotes:
       ``\,`this'            % \, separates the double and single quote.
        is what I just
        wrote, not  `that'\,''.

Dashes come in three sizes: an
       intra-word
dash, a medium dash for number ranges like
       1--2,
and a punctuation
       dash---like
this.

A sentence-ending space should be larger than the
space between words within a sentence.  You
sometimes have to type special commands in
conjunction with punctuation characters to get
this right, as in the following sentence.
       Gnats, gnus, etc.\ all  % `\ ' makes an inter-word space.
       begin with G\@.         % \@ marks end-of-sentence punctuation.
You should check the spaces after periods when
reading your output to make sure you haven't
forgotten any special cases.  Generating an
ellipsis
       \ldots\               % `\ ' is needed after `\ldots' because TeX
                             % ignores spaces after command names like \ldots
                             % made from \ + letters.
                             %
                             % Note how a `%' character causes TeX to ignore
                             % the end of the input line, so these blank lines
                             % do not start a new paragraph.
                             %
with the right spacing around the periods requires
a special command.

\LaTeX\ interprets some common characters as
commands, so you must type special commands to
generate them.  These characters include the
following:
       \$ \& \% \# \{ and \}.

In printing, text is usually emphasized with an
       \emph{italic}
type style.

\begin{em}
   A long segment of text can also be emphasized
   in this way.  Text within such a segment can be
   given \emph{additional} emphasis.
\end{em}

It is sometimes necessary to prevent \LaTeX\ from
breaking a line where it might otherwise do so.
This may be at a space, as between the ``Mr.''\ and
``Jones'' in
       ``Mr.~Jones'',        % ~ produces an unbreakable interword space.
or within a word---especially when the word is a
symbol like
       \mbox{\emph{itemnum}}
that makes little sense when hyphenated across
lines.

Footnotes\footnote{This is an example of a footnote.}
pose no problem.

\LaTeX\ is good at typesetting mathematical formulas
like
       \( x-3y + z = 7 \)
or
       \( a_{1} > x^{2n} + y^{2n} > x' \)
or
       \( \ip{A}{B} = \sum_{i} a_{i} b_{i} \).
The spaces you type in a formula are
ignored.  Remember that a letter like
       $x$                   % $ ... $  and  \( ... \)  are equivalent
is a formula when it denotes a mathematical
symbol, and it should be typed as one.

\section{Displayed Text}

Text is displayed by indenting it from the left
margin.  Quotations are commonly displayed.  There
are short quotations
\begin{quote}
   This is a short quotation.  It consists of a
   single paragraph of text.  See how it is formatted.
\end{quote}
and longer ones.
\begin{quotation}
   This is a longer quotation.  It consists of two
   paragraphs of text, neither of which are
   particularly interesting.

   This is the second paragraph of the quotation.  It
   is just as dull as the first paragraph.
\end{quotation}
Another frequently-displayed structure is a list.
The following is an example of an \emph{itemized}
list.
\begin{itemize}
   \item This is the first item of an itemized list.
         Each item in the list is marked with a ``tick''.
         You don't have to worry about what kind of tick
         mark is used.

   \item This is the second item of the list.  It
         contains another list nested inside it.  The inner
         list is an \emph{enumerated} list.
         \begin{enumerate}
            \item This is the first item of an enumerated
                  list that is nested within the itemized list.

            \item This is the second item of the inner list.
                  \LaTeX\ allows you to nest lists deeper than
                  you really should.
         \end{enumerate}
         This is the rest of the second item of the outer
         list.  It is no more interesting than any other
         part of the item.
   \item This is the third item of the list.
\end{itemize}
You can even display poetry.
\begin{verse}
   There is an environment
    for verse \\             % The \\ command separates lines
   Whose features some poets % within a stanza.
   will curse.

                             % One or more blank lines separate stanzas.

   For instead of making\\
   Them do \emph{all} line breaking, \\
   It allows them to put too many words on a line when they'd rather be
   forced to be terse.
\end{verse}

Mathematical formulas may also be displayed.  A
displayed formula
is
one-line long; multiline
formulas require special formatting instructions.
   \[  \ip{\Gamma}{\psi'} = x'' + y^{2} + z_{i}^{n}\]
Don't start a paragraph with a displayed equation,
nor make one a paragraph by itself.

\end{document}               % End of document.

\section{Bench Section 240}
% synthetic bench section 240
% This is a sample LaTeX input file.  (Version of 12 August 2004.)
%
% A '%' character causes TeX to ignore all remaining text on the line,
% and is used for comments like this one.

\documentclass{article}      % Specifies the document class

                             % The preamble begins here.
\title{An Example Document}  % Declares the document's title.
\author{Leslie Lamport}      % Declares the author's name.
\date{January 21, 1994}      % Deleting this command produces today's date.

\newcommand{\ip}[2]{(#1, #2)}
                             % Defines \ip{arg1}{arg2} to mean
                             % (arg1, arg2).

%\newcommand{\ip}[2]{\langle #1 | #2\rangle}
                             % This is an alternative definition of
                             % \ip that is commented out.

\begin{document}             % End of preamble and beginning of text.

\maketitle                   % Produces the title.

This is an example input file.  Comparing it with
the output it generates can show you how to
produce a simple document of your own.

\section{Ordinary Text}      % Produces section heading.  Lower-level
                             % sections are begun with similar
                             % \subsection and \subsubsection commands.

The ends  of words and sentences are marked
  by   spaces. It  doesn't matter how many
spaces    you type; one is as good as 100.  The
end of   a line counts as a space.

One   or more   blank lines denote the  end
of  a paragraph.

Since any number of consecutive spaces are treated
like a single one, the formatting of the input
file makes no difference to
      \LaTeX,                % The \LaTeX command generates the LaTeX logo.
but it makes a difference to you.  When you use
\LaTeX, making your input file as easy to read
as possible will be a great help as you write
your document and when you change it.  This sample
file shows how you can add comments to your own input
file.

Because printing is different from typewriting,
there are a number of things that you have to do
differently when preparing an input file than if
you were just typing the document directly.
Quotation marks like
       ``this''
have to be handled specially, as do quotes within
quotes:
       ``\,`this'            % \, separates the double and single quote.
        is what I just
        wrote, not  `that'\,''.

Dashes come in three sizes: an
       intra-word
dash, a medium dash for number ranges like
       1--2,
and a punctuation
       dash---like
this.

A sentence-ending space should be larger than the
space between words within a sentence.  You
sometimes have to type special commands in
conjunction with punctuation characters to get
this right, as in the following sentence.
       Gnats, gnus, etc.\ all  % `\ ' makes an inter-word space.
       begin with G\@.         % \@ marks end-of-sentence punctuation.
You should check the spaces after periods when
reading your output to make sure you haven't
forgotten any special cases.  Generating an
ellipsis
       \ldots\               % `\ ' is needed after `\ldots' because TeX
                             % ignores spaces after command names like \ldots
                             % made from \ + letters.
                             %
                             % Note how a `%' character causes TeX to ignore
                             % the end of the input line, so these blank lines
                             % do not start a new paragraph.
                             %
with the right spacing around the periods requires
a special command.

\LaTeX\ interprets some common characters as
commands, so you must type special commands to
generate them.  These characters include the
following:
       \$ \& \% \# \{ and \}.

In printing, text is usually emphasized with an
       \emph{italic}
type style.

\begin{em}
   A long segment of text can also be emphasized
   in this way.  Text within such a segment can be
   given \emph{additional} emphasis.
\end{em}

It is sometimes necessary to prevent \LaTeX\ from
breaking a line where it might otherwise do so.
This may be at a space, as between the ``Mr.''\ and
``Jones'' in
       ``Mr.~Jones'',        % ~ produces an unbreakable interword space.
or within a word---especially when the word is a
symbol like
       \mbox{\emph{itemnum}}
that makes little sense when hyphenated across
lines.

Footnotes\footnote{This is an example of a footnote.}
pose no problem.

\LaTeX\ is good at typesetting mathematical formulas
like
       \( x-3y + z = 7 \)
or
       \( a_{1} > x^{2n} + y^{2n} > x' \)
or
       \( \ip{A}{B} = \sum_{i} a_{i} b_{i} \).
The spaces you type in a formula are
ignored.  Remember that a letter like
       $x$                   % $ ... $  and  \( ... \)  are equivalent
is a formula when it denotes a mathematical
symbol, and it should be typed as one.

\section{Displayed Text}

Text is displayed by indenting it from the left
margin.  Quotations are commonly displayed.  There
are short quotations
\begin{quote}
   This is a short quotation.  It consists of a
   single paragraph of text.  See how it is formatted.
\end{quote}
and longer ones.
\begin{quotation}
   This is a longer quotation.  It consists of two
   paragraphs of text, neither of which are
   particularly interesting.

   This is the second paragraph of the quotation.  It
   is just as dull as the first paragraph.
\end{quotation}
Another frequently-displayed structure is a list.
The following is an example of an \emph{itemized}
list.
\begin{itemize}
   \item This is the first item of an itemized list.
         Each item in the list is marked with a ``tick''.
         You don't have to worry about what kind of tick
         mark is used.

   \item This is the second item of the list.  It
         contains another list nested inside it.  The inner
         list is an \emph{enumerated} list.
         \begin{enumerate}
            \item This is the first item of an enumerated
                  list that is nested within the itemized list.

            \item This is the second item of the inner list.
                  \LaTeX\ allows you to nest lists deeper than
                  you really should.
         \end{enumerate}
         This is the rest of the second item of the outer
         list.  It is no more interesting than any other
         part of the item.
   \item This is the third item of the list.
\end{itemize}
You can even display poetry.
\begin{verse}
   There is an environment
    for verse \\             % The \\ command separates lines
   Whose features some poets % within a stanza.
   will curse.

                             % One or more blank lines separate stanzas.

   For instead of making\\
   Them do \emph{all} line breaking, \\
   It allows them to put too many words on a line when they'd rather be
   forced to be terse.
\end{verse}

Mathematical formulas may also be displayed.  A
displayed formula
is
one-line long; multiline
formulas require special formatting instructions.
   \[  \ip{\Gamma}{\psi'} = x'' + y^{2} + z_{i}^{n}\]
Don't start a paragraph with a displayed equation,
nor make one a paragraph by itself.

\end{document}               % End of document.

\section{Bench Section 241}
% synthetic bench section 241
% This is a sample LaTeX input file.  (Version of 12 August 2004.)
%
% A '%' character causes TeX to ignore all remaining text on the line,
% and is used for comments like this one.

\documentclass{article}      % Specifies the document class

                             % The preamble begins here.
\title{An Example Document}  % Declares the document's title.
\author{Leslie Lamport}      % Declares the author's name.
\date{January 21, 1994}      % Deleting this command produces today's date.

\newcommand{\ip}[2]{(#1, #2)}
                             % Defines \ip{arg1}{arg2} to mean
                             % (arg1, arg2).

%\newcommand{\ip}[2]{\langle #1 | #2\rangle}
                             % This is an alternative definition of
                             % \ip that is commented out.

\begin{document}             % End of preamble and beginning of text.

\maketitle                   % Produces the title.

This is an example input file.  Comparing it with
the output it generates can show you how to
produce a simple document of your own.

\section{Ordinary Text}      % Produces section heading.  Lower-level
                             % sections are begun with similar
                             % \subsection and \subsubsection commands.

The ends  of words and sentences are marked
  by   spaces. It  doesn't matter how many
spaces    you type; one is as good as 100.  The
end of   a line counts as a space.

One   or more   blank lines denote the  end
of  a paragraph.

Since any number of consecutive spaces are treated
like a single one, the formatting of the input
file makes no difference to
      \LaTeX,                % The \LaTeX command generates the LaTeX logo.
but it makes a difference to you.  When you use
\LaTeX, making your input file as easy to read
as possible will be a great help as you write
your document and when you change it.  This sample
file shows how you can add comments to your own input
file.

Because printing is different from typewriting,
there are a number of things that you have to do
differently when preparing an input file than if
you were just typing the document directly.
Quotation marks like
       ``this''
have to be handled specially, as do quotes within
quotes:
       ``\,`this'            % \, separates the double and single quote.
        is what I just
        wrote, not  `that'\,''.

Dashes come in three sizes: an
       intra-word
dash, a medium dash for number ranges like
       1--2,
and a punctuation
       dash---like
this.

A sentence-ending space should be larger than the
space between words within a sentence.  You
sometimes have to type special commands in
conjunction with punctuation characters to get
this right, as in the following sentence.
       Gnats, gnus, etc.\ all  % `\ ' makes an inter-word space.
       begin with G\@.         % \@ marks end-of-sentence punctuation.
You should check the spaces after periods when
reading your output to make sure you haven't
forgotten any special cases.  Generating an
ellipsis
       \ldots\               % `\ ' is needed after `\ldots' because TeX
                             % ignores spaces after command names like \ldots
                             % made from \ + letters.
                             %
                             % Note how a `%' character causes TeX to ignore
                             % the end of the input line, so these blank lines
                             % do not start a new paragraph.
                             %
with the right spacing around the periods requires
a special command.

\LaTeX\ interprets some common characters as
commands, so you must type special commands to
generate them.  These characters include the
following:
       \$ \& \% \# \{ and \}.

In printing, text is usually emphasized with an
       \emph{italic}
type style.

\begin{em}
   A long segment of text can also be emphasized
   in this way.  Text within such a segment can be
   given \emph{additional} emphasis.
\end{em}

It is sometimes necessary to prevent \LaTeX\ from
breaking a line where it might otherwise do so.
This may be at a space, as between the ``Mr.''\ and
``Jones'' in
       ``Mr.~Jones'',        % ~ produces an unbreakable interword space.
or within a word---especially when the word is a
symbol like
       \mbox{\emph{itemnum}}
that makes little sense when hyphenated across
lines.

Footnotes\footnote{This is an example of a footnote.}
pose no problem.

\LaTeX\ is good at typesetting mathematical formulas
like
       \( x-3y + z = 7 \)
or
       \( a_{1} > x^{2n} + y^{2n} > x' \)
or
       \( \ip{A}{B} = \sum_{i} a_{i} b_{i} \).
The spaces you type in a formula are
ignored.  Remember that a letter like
       $x$                   % $ ... $  and  \( ... \)  are equivalent
is a formula when it denotes a mathematical
symbol, and it should be typed as one.

\section{Displayed Text}

Text is displayed by indenting it from the left
margin.  Quotations are commonly displayed.  There
are short quotations
\begin{quote}
   This is a short quotation.  It consists of a
   single paragraph of text.  See how it is formatted.
\end{quote}
and longer ones.
\begin{quotation}
   This is a longer quotation.  It consists of two
   paragraphs of text, neither of which are
   particularly interesting.

   This is the second paragraph of the quotation.  It
   is just as dull as the first paragraph.
\end{quotation}
Another frequently-displayed structure is a list.
The following is an example of an \emph{itemized}
list.
\begin{itemize}
   \item This is the first item of an itemized list.
         Each item in the list is marked with a ``tick''.
         You don't have to worry about what kind of tick
         mark is used.

   \item This is the second item of the list.  It
         contains another list nested inside it.  The inner
         list is an \emph{enumerated} list.
         \begin{enumerate}
            \item This is the first item of an enumerated
                  list that is nested within the itemized list.

            \item This is the second item of the inner list.
                  \LaTeX\ allows you to nest lists deeper than
                  you really should.
         \end{enumerate}
         This is the rest of the second item of the outer
         list.  It is no more interesting than any other
         part of the item.
   \item This is the third item of the list.
\end{itemize}
You can even display poetry.
\begin{verse}
   There is an environment
    for verse \\             % The \\ command separates lines
   Whose features some poets % within a stanza.
   will curse.

                             % One or more blank lines separate stanzas.

   For instead of making\\
   Them do \emph{all} line breaking, \\
   It allows them to put too many words on a line when they'd rather be
   forced to be terse.
\end{verse}

Mathematical formulas may also be displayed.  A
displayed formula
is
one-line long; multiline
formulas require special formatting instructions.
   \[  \ip{\Gamma}{\psi'} = x'' + y^{2} + z_{i}^{n}\]
Don't start a paragraph with a displayed equation,
nor make one a paragraph by itself.

\end{document}               % End of document.

\section{Bench Section 242}
% synthetic bench section 242
% This is a sample LaTeX input file.  (Version of 12 August 2004.)
%
% A '%' character causes TeX to ignore all remaining text on the line,
% and is used for comments like this one.

\documentclass{article}      % Specifies the document class

                             % The preamble begins here.
\title{An Example Document}  % Declares the document's title.
\author{Leslie Lamport}      % Declares the author's name.
\date{January 21, 1994}      % Deleting this command produces today's date.

\newcommand{\ip}[2]{(#1, #2)}
                             % Defines \ip{arg1}{arg2} to mean
                             % (arg1, arg2).

%\newcommand{\ip}[2]{\langle #1 | #2\rangle}
                             % This is an alternative definition of
                             % \ip that is commented out.

\begin{document}             % End of preamble and beginning of text.

\maketitle                   % Produces the title.

This is an example input file.  Comparing it with
the output it generates can show you how to
produce a simple document of your own.

\section{Ordinary Text}      % Produces section heading.  Lower-level
                             % sections are begun with similar
                             % \subsection and \subsubsection commands.

The ends  of words and sentences are marked
  by   spaces. It  doesn't matter how many
spaces    you type; one is as good as 100.  The
end of   a line counts as a space.

One   or more   blank lines denote the  end
of  a paragraph.

Since any number of consecutive spaces are treated
like a single one, the formatting of the input
file makes no difference to
      \LaTeX,                % The \LaTeX command generates the LaTeX logo.
but it makes a difference to you.  When you use
\LaTeX, making your input file as easy to read
as possible will be a great help as you write
your document and when you change it.  This sample
file shows how you can add comments to your own input
file.

Because printing is different from typewriting,
there are a number of things that you have to do
differently when preparing an input file than if
you were just typing the document directly.
Quotation marks like
       ``this''
have to be handled specially, as do quotes within
quotes:
       ``\,`this'            % \, separates the double and single quote.
        is what I just
        wrote, not  `that'\,''.

Dashes come in three sizes: an
       intra-word
dash, a medium dash for number ranges like
       1--2,
and a punctuation
       dash---like
this.

A sentence-ending space should be larger than the
space between words within a sentence.  You
sometimes have to type special commands in
conjunction with punctuation characters to get
this right, as in the following sentence.
       Gnats, gnus, etc.\ all  % `\ ' makes an inter-word space.
       begin with G\@.         % \@ marks end-of-sentence punctuation.
You should check the spaces after periods when
reading your output to make sure you haven't
forgotten any special cases.  Generating an
ellipsis
       \ldots\               % `\ ' is needed after `\ldots' because TeX
                             % ignores spaces after command names like \ldots
                             % made from \ + letters.
                             %
                             % Note how a `%' character causes TeX to ignore
                             % the end of the input line, so these blank lines
                             % do not start a new paragraph.
                             %
with the right spacing around the periods requires
a special command.

\LaTeX\ interprets some common characters as
commands, so you must type special commands to
generate them.  These characters include the
following:
       \$ \& \% \# \{ and \}.

In printing, text is usually emphasized with an
       \emph{italic}
type style.

\begin{em}
   A long segment of text can also be emphasized
   in this way.  Text within such a segment can be
   given \emph{additional} emphasis.
\end{em}

It is sometimes necessary to prevent \LaTeX\ from
breaking a line where it might otherwise do so.
This may be at a space, as between the ``Mr.''\ and
``Jones'' in
       ``Mr.~Jones'',        % ~ produces an unbreakable interword space.
or within a word---especially when the word is a
symbol like
       \mbox{\emph{itemnum}}
that makes little sense when hyphenated across
lines.

Footnotes\footnote{This is an example of a footnote.}
pose no problem.

\LaTeX\ is good at typesetting mathematical formulas
like
       \( x-3y + z = 7 \)
or
       \( a_{1} > x^{2n} + y^{2n} > x' \)
or
       \( \ip{A}{B} = \sum_{i} a_{i} b_{i} \).
The spaces you type in a formula are
ignored.  Remember that a letter like
       $x$                   % $ ... $  and  \( ... \)  are equivalent
is a formula when it denotes a mathematical
symbol, and it should be typed as one.

\section{Displayed Text}

Text is displayed by indenting it from the left
margin.  Quotations are commonly displayed.  There
are short quotations
\begin{quote}
   This is a short quotation.  It consists of a
   single paragraph of text.  See how it is formatted.
\end{quote}
and longer ones.
\begin{quotation}
   This is a longer quotation.  It consists of two
   paragraphs of text, neither of which are
   particularly interesting.

   This is the second paragraph of the quotation.  It
   is just as dull as the first paragraph.
\end{quotation}
Another frequently-displayed structure is a list.
The following is an example of an \emph{itemized}
list.
\begin{itemize}
   \item This is the first item of an itemized list.
         Each item in the list is marked with a ``tick''.
         You don't have to worry about what kind of tick
         mark is used.

   \item This is the second item of the list.  It
         contains another list nested inside it.  The inner
         list is an \emph{enumerated} list.
         \begin{enumerate}
            \item This is the first item of an enumerated
                  list that is nested within the itemized list.

            \item This is the second item of the inner list.
                  \LaTeX\ allows you to nest lists deeper than
                  you really should.
         \end{enumerate}
         This is the rest of the second item of the outer
         list.  It is no more interesting than any other
         part of the item.
   \item This is the third item of the list.
\end{itemize}
You can even display poetry.
\begin{verse}
   There is an environment
    for verse \\             % The \\ command separates lines
   Whose features some poets % within a stanza.
   will curse.

                             % One or more blank lines separate stanzas.

   For instead of making\\
   Them do \emph{all} line breaking, \\
   It allows them to put too many words on a line when they'd rather be
   forced to be terse.
\end{verse}

Mathematical formulas may also be displayed.  A
displayed formula
is
one-line long; multiline
formulas require special formatting instructions.
   \[  \ip{\Gamma}{\psi'} = x'' + y^{2} + z_{i}^{n}\]
Don't start a paragraph with a displayed equation,
nor make one a paragraph by itself.

\end{document}               % End of document.

\section{Bench Section 243}
% synthetic bench section 243
% This is a sample LaTeX input file.  (Version of 12 August 2004.)
%
% A '%' character causes TeX to ignore all remaining text on the line,
% and is used for comments like this one.

\documentclass{article}      % Specifies the document class

                             % The preamble begins here.
\title{An Example Document}  % Declares the document's title.
\author{Leslie Lamport}      % Declares the author's name.
\date{January 21, 1994}      % Deleting this command produces today's date.

\newcommand{\ip}[2]{(#1, #2)}
                             % Defines \ip{arg1}{arg2} to mean
                             % (arg1, arg2).

%\newcommand{\ip}[2]{\langle #1 | #2\rangle}
                             % This is an alternative definition of
                             % \ip that is commented out.

\begin{document}             % End of preamble and beginning of text.

\maketitle                   % Produces the title.

This is an example input file.  Comparing it with
the output it generates can show you how to
produce a simple document of your own.

\section{Ordinary Text}      % Produces section heading.  Lower-level
                             % sections are begun with similar
                             % \subsection and \subsubsection commands.

The ends  of words and sentences are marked
  by   spaces. It  doesn't matter how many
spaces    you type; one is as good as 100.  The
end of   a line counts as a space.

One   or more   blank lines denote the  end
of  a paragraph.

Since any number of consecutive spaces are treated
like a single one, the formatting of the input
file makes no difference to
      \LaTeX,                % The \LaTeX command generates the LaTeX logo.
but it makes a difference to you.  When you use
\LaTeX, making your input file as easy to read
as possible will be a great help as you write
your document and when you change it.  This sample
file shows how you can add comments to your own input
file.

Because printing is different from typewriting,
there are a number of things that you have to do
differently when preparing an input file than if
you were just typing the document directly.
Quotation marks like
       ``this''
have to be handled specially, as do quotes within
quotes:
       ``\,`this'            % \, separates the double and single quote.
        is what I just
        wrote, not  `that'\,''.

Dashes come in three sizes: an
       intra-word
dash, a medium dash for number ranges like
       1--2,
and a punctuation
       dash---like
this.

A sentence-ending space should be larger than the
space between words within a sentence.  You
sometimes have to type special commands in
conjunction with punctuation characters to get
this right, as in the following sentence.
       Gnats, gnus, etc.\ all  % `\ ' makes an inter-word space.
       begin with G\@.         % \@ marks end-of-sentence punctuation.
You should check the spaces after periods when
reading your output to make sure you haven't
forgotten any special cases.  Generating an
ellipsis
       \ldots\               % `\ ' is needed after `\ldots' because TeX
                             % ignores spaces after command names like \ldots
                             % made from \ + letters.
                             %
                             % Note how a `%' character causes TeX to ignore
                             % the end of the input line, so these blank lines
                             % do not start a new paragraph.
                             %
with the right spacing around the periods requires
a special command.

\LaTeX\ interprets some common characters as
commands, so you must type special commands to
generate them.  These characters include the
following:
       \$ \& \% \# \{ and \}.

In printing, text is usually emphasized with an
       \emph{italic}
type style.

\begin{em}
   A long segment of text can also be emphasized
   in this way.  Text within such a segment can be
   given \emph{additional} emphasis.
\end{em}

It is sometimes necessary to prevent \LaTeX\ from
breaking a line where it might otherwise do so.
This may be at a space, as between the ``Mr.''\ and
``Jones'' in
       ``Mr.~Jones'',        % ~ produces an unbreakable interword space.
or within a word---especially when the word is a
symbol like
       \mbox{\emph{itemnum}}
that makes little sense when hyphenated across
lines.

Footnotes\footnote{This is an example of a footnote.}
pose no problem.

\LaTeX\ is good at typesetting mathematical formulas
like
       \( x-3y + z = 7 \)
or
       \( a_{1} > x^{2n} + y^{2n} > x' \)
or
       \( \ip{A}{B} = \sum_{i} a_{i} b_{i} \).
The spaces you type in a formula are
ignored.  Remember that a letter like
       $x$                   % $ ... $  and  \( ... \)  are equivalent
is a formula when it denotes a mathematical
symbol, and it should be typed as one.

\section{Displayed Text}

Text is displayed by indenting it from the left
margin.  Quotations are commonly displayed.  There
are short quotations
\begin{quote}
   This is a short quotation.  It consists of a
   single paragraph of text.  See how it is formatted.
\end{quote}
and longer ones.
\begin{quotation}
   This is a longer quotation.  It consists of two
   paragraphs of text, neither of which are
   particularly interesting.

   This is the second paragraph of the quotation.  It
   is just as dull as the first paragraph.
\end{quotation}
Another frequently-displayed structure is a list.
The following is an example of an \emph{itemized}
list.
\begin{itemize}
   \item This is the first item of an itemized list.
         Each item in the list is marked with a ``tick''.
         You don't have to worry about what kind of tick
         mark is used.

   \item This is the second item of the list.  It
         contains another list nested inside it.  The inner
         list is an \emph{enumerated} list.
         \begin{enumerate}
            \item This is the first item of an enumerated
                  list that is nested within the itemized list.

            \item This is the second item of the inner list.
                  \LaTeX\ allows you to nest lists deeper than
                  you really should.
         \end{enumerate}
         This is the rest of the second item of the outer
         list.  It is no more interesting than any other
         part of the item.
   \item This is the third item of the list.
\end{itemize}
You can even display poetry.
\begin{verse}
   There is an environment
    for verse \\             % The \\ command separates lines
   Whose features some poets % within a stanza.
   will curse.

                             % One or more blank lines separate stanzas.

   For instead of making\\
   Them do \emph{all} line breaking, \\
   It allows them to put too many words on a line when they'd rather be
   forced to be terse.
\end{verse}

Mathematical formulas may also be displayed.  A
displayed formula
is
one-line long; multiline
formulas require special formatting instructions.
   \[  \ip{\Gamma}{\psi'} = x'' + y^{2} + z_{i}^{n}\]
Don't start a paragraph with a displayed equation,
nor make one a paragraph by itself.

\end{document}               % End of document.

\section{Bench Section 244}
% synthetic bench section 244
% This is a sample LaTeX input file.  (Version of 12 August 2004.)
%
% A '%' character causes TeX to ignore all remaining text on the line,
% and is used for comments like this one.

\documentclass{article}      % Specifies the document class

                             % The preamble begins here.
\title{An Example Document}  % Declares the document's title.
\author{Leslie Lamport}      % Declares the author's name.
\date{January 21, 1994}      % Deleting this command produces today's date.

\newcommand{\ip}[2]{(#1, #2)}
                             % Defines \ip{arg1}{arg2} to mean
                             % (arg1, arg2).

%\newcommand{\ip}[2]{\langle #1 | #2\rangle}
                             % This is an alternative definition of
                             % \ip that is commented out.

\begin{document}             % End of preamble and beginning of text.

\maketitle                   % Produces the title.

This is an example input file.  Comparing it with
the output it generates can show you how to
produce a simple document of your own.

\section{Ordinary Text}      % Produces section heading.  Lower-level
                             % sections are begun with similar
                             % \subsection and \subsubsection commands.

The ends  of words and sentences are marked
  by   spaces. It  doesn't matter how many
spaces    you type; one is as good as 100.  The
end of   a line counts as a space.

One   or more   blank lines denote the  end
of  a paragraph.

Since any number of consecutive spaces are treated
like a single one, the formatting of the input
file makes no difference to
      \LaTeX,                % The \LaTeX command generates the LaTeX logo.
but it makes a difference to you.  When you use
\LaTeX, making your input file as easy to read
as possible will be a great help as you write
your document and when you change it.  This sample
file shows how you can add comments to your own input
file.

Because printing is different from typewriting,
there are a number of things that you have to do
differently when preparing an input file than if
you were just typing the document directly.
Quotation marks like
       ``this''
have to be handled specially, as do quotes within
quotes:
       ``\,`this'            % \, separates the double and single quote.
        is what I just
        wrote, not  `that'\,''.

Dashes come in three sizes: an
       intra-word
dash, a medium dash for number ranges like
       1--2,
and a punctuation
       dash---like
this.

A sentence-ending space should be larger than the
space between words within a sentence.  You
sometimes have to type special commands in
conjunction with punctuation characters to get
this right, as in the following sentence.
       Gnats, gnus, etc.\ all  % `\ ' makes an inter-word space.
       begin with G\@.         % \@ marks end-of-sentence punctuation.
You should check the spaces after periods when
reading your output to make sure you haven't
forgotten any special cases.  Generating an
ellipsis
       \ldots\               % `\ ' is needed after `\ldots' because TeX
                             % ignores spaces after command names like \ldots
                             % made from \ + letters.
                             %
                             % Note how a `%' character causes TeX to ignore
                             % the end of the input line, so these blank lines
                             % do not start a new paragraph.
                             %
with the right spacing around the periods requires
a special command.

\LaTeX\ interprets some common characters as
commands, so you must type special commands to
generate them.  These characters include the
following:
       \$ \& \% \# \{ and \}.

In printing, text is usually emphasized with an
       \emph{italic}
type style.

\begin{em}
   A long segment of text can also be emphasized
   in this way.  Text within such a segment can be
   given \emph{additional} emphasis.
\end{em}

It is sometimes necessary to prevent \LaTeX\ from
breaking a line where it might otherwise do so.
This may be at a space, as between the ``Mr.''\ and
``Jones'' in
       ``Mr.~Jones'',        % ~ produces an unbreakable interword space.
or within a word---especially when the word is a
symbol like
       \mbox{\emph{itemnum}}
that makes little sense when hyphenated across
lines.

Footnotes\footnote{This is an example of a footnote.}
pose no problem.

\LaTeX\ is good at typesetting mathematical formulas
like
       \( x-3y + z = 7 \)
or
       \( a_{1} > x^{2n} + y^{2n} > x' \)
or
       \( \ip{A}{B} = \sum_{i} a_{i} b_{i} \).
The spaces you type in a formula are
ignored.  Remember that a letter like
       $x$                   % $ ... $  and  \( ... \)  are equivalent
is a formula when it denotes a mathematical
symbol, and it should be typed as one.

\section{Displayed Text}

Text is displayed by indenting it from the left
margin.  Quotations are commonly displayed.  There
are short quotations
\begin{quote}
   This is a short quotation.  It consists of a
   single paragraph of text.  See how it is formatted.
\end{quote}
and longer ones.
\begin{quotation}
   This is a longer quotation.  It consists of two
   paragraphs of text, neither of which are
   particularly interesting.

   This is the second paragraph of the quotation.  It
   is just as dull as the first paragraph.
\end{quotation}
Another frequently-displayed structure is a list.
The following is an example of an \emph{itemized}
list.
\begin{itemize}
   \item This is the first item of an itemized list.
         Each item in the list is marked with a ``tick''.
         You don't have to worry about what kind of tick
         mark is used.

   \item This is the second item of the list.  It
         contains another list nested inside it.  The inner
         list is an \emph{enumerated} list.
         \begin{enumerate}
            \item This is the first item of an enumerated
                  list that is nested within the itemized list.

            \item This is the second item of the inner list.
                  \LaTeX\ allows you to nest lists deeper than
                  you really should.
         \end{enumerate}
         This is the rest of the second item of the outer
         list.  It is no more interesting than any other
         part of the item.
   \item This is the third item of the list.
\end{itemize}
You can even display poetry.
\begin{verse}
   There is an environment
    for verse \\             % The \\ command separates lines
   Whose features some poets % within a stanza.
   will curse.

                             % One or more blank lines separate stanzas.

   For instead of making\\
   Them do \emph{all} line breaking, \\
   It allows them to put too many words on a line when they'd rather be
   forced to be terse.
\end{verse}

Mathematical formulas may also be displayed.  A
displayed formula
is
one-line long; multiline
formulas require special formatting instructions.
   \[  \ip{\Gamma}{\psi'} = x'' + y^{2} + z_{i}^{n}\]
Don't start a paragraph with a displayed equation,
nor make one a paragraph by itself.

\end{document}               % End of document.

\section{Bench Section 245}
% synthetic bench section 245
% This is a sample LaTeX input file.  (Version of 12 August 2004.)
%
% A '%' character causes TeX to ignore all remaining text on the line,
% and is used for comments like this one.

\documentclass{article}      % Specifies the document class

                             % The preamble begins here.
\title{An Example Document}  % Declares the document's title.
\author{Leslie Lamport}      % Declares the author's name.
\date{January 21, 1994}      % Deleting this command produces today's date.

\newcommand{\ip}[2]{(#1, #2)}
                             % Defines \ip{arg1}{arg2} to mean
                             % (arg1, arg2).

%\newcommand{\ip}[2]{\langle #1 | #2\rangle}
                             % This is an alternative definition of
                             % \ip that is commented out.

\begin{document}             % End of preamble and beginning of text.

\maketitle                   % Produces the title.

This is an example input file.  Comparing it with
the output it generates can show you how to
produce a simple document of your own.

\section{Ordinary Text}      % Produces section heading.  Lower-level
                             % sections are begun with similar
                             % \subsection and \subsubsection commands.

The ends  of words and sentences are marked
  by   spaces. It  doesn't matter how many
spaces    you type; one is as good as 100.  The
end of   a line counts as a space.

One   or more   blank lines denote the  end
of  a paragraph.

Since any number of consecutive spaces are treated
like a single one, the formatting of the input
file makes no difference to
      \LaTeX,                % The \LaTeX command generates the LaTeX logo.
but it makes a difference to you.  When you use
\LaTeX, making your input file as easy to read
as possible will be a great help as you write
your document and when you change it.  This sample
file shows how you can add comments to your own input
file.

Because printing is different from typewriting,
there are a number of things that you have to do
differently when preparing an input file than if
you were just typing the document directly.
Quotation marks like
       ``this''
have to be handled specially, as do quotes within
quotes:
       ``\,`this'            % \, separates the double and single quote.
        is what I just
        wrote, not  `that'\,''.

Dashes come in three sizes: an
       intra-word
dash, a medium dash for number ranges like
       1--2,
and a punctuation
       dash---like
this.

A sentence-ending space should be larger than the
space between words within a sentence.  You
sometimes have to type special commands in
conjunction with punctuation characters to get
this right, as in the following sentence.
       Gnats, gnus, etc.\ all  % `\ ' makes an inter-word space.
       begin with G\@.         % \@ marks end-of-sentence punctuation.
You should check the spaces after periods when
reading your output to make sure you haven't
forgotten any special cases.  Generating an
ellipsis
       \ldots\               % `\ ' is needed after `\ldots' because TeX
                             % ignores spaces after command names like \ldots
                             % made from \ + letters.
                             %
                             % Note how a `%' character causes TeX to ignore
                             % the end of the input line, so these blank lines
                             % do not start a new paragraph.
                             %
with the right spacing around the periods requires
a special command.

\LaTeX\ interprets some common characters as
commands, so you must type special commands to
generate them.  These characters include the
following:
       \$ \& \% \# \{ and \}.

In printing, text is usually emphasized with an
       \emph{italic}
type style.

\begin{em}
   A long segment of text can also be emphasized
   in this way.  Text within such a segment can be
   given \emph{additional} emphasis.
\end{em}

It is sometimes necessary to prevent \LaTeX\ from
breaking a line where it might otherwise do so.
This may be at a space, as between the ``Mr.''\ and
``Jones'' in
       ``Mr.~Jones'',        % ~ produces an unbreakable interword space.
or within a word---especially when the word is a
symbol like
       \mbox{\emph{itemnum}}
that makes little sense when hyphenated across
lines.

Footnotes\footnote{This is an example of a footnote.}
pose no problem.

\LaTeX\ is good at typesetting mathematical formulas
like
       \( x-3y + z = 7 \)
or
       \( a_{1} > x^{2n} + y^{2n} > x' \)
or
       \( \ip{A}{B} = \sum_{i} a_{i} b_{i} \).
The spaces you type in a formula are
ignored.  Remember that a letter like
       $x$                   % $ ... $  and  \( ... \)  are equivalent
is a formula when it denotes a mathematical
symbol, and it should be typed as one.

\section{Displayed Text}

Text is displayed by indenting it from the left
margin.  Quotations are commonly displayed.  There
are short quotations
\begin{quote}
   This is a short quotation.  It consists of a
   single paragraph of text.  See how it is formatted.
\end{quote}
and longer ones.
\begin{quotation}
   This is a longer quotation.  It consists of two
   paragraphs of text, neither of which are
   particularly interesting.

   This is the second paragraph of the quotation.  It
   is just as dull as the first paragraph.
\end{quotation}
Another frequently-displayed structure is a list.
The following is an example of an \emph{itemized}
list.
\begin{itemize}
   \item This is the first item of an itemized list.
         Each item in the list is marked with a ``tick''.
         You don't have to worry about what kind of tick
         mark is used.

   \item This is the second item of the list.  It
         contains another list nested inside it.  The inner
         list is an \emph{enumerated} list.
         \begin{enumerate}
            \item This is the first item of an enumerated
                  list that is nested within the itemized list.

            \item This is the second item of the inner list.
                  \LaTeX\ allows you to nest lists deeper than
                  you really should.
         \end{enumerate}
         This is the rest of the second item of the outer
         list.  It is no more interesting than any other
         part of the item.
   \item This is the third item of the list.
\end{itemize}
You can even display poetry.
\begin{verse}
   There is an environment
    for verse \\             % The \\ command separates lines
   Whose features some poets % within a stanza.
   will curse.

                             % One or more blank lines separate stanzas.

   For instead of making\\
   Them do \emph{all} line breaking, \\
   It allows them to put too many words on a line when they'd rather be
   forced to be terse.
\end{verse}

Mathematical formulas may also be displayed.  A
displayed formula
is
one-line long; multiline
formulas require special formatting instructions.
   \[  \ip{\Gamma}{\psi'} = x'' + y^{2} + z_{i}^{n}\]
Don't start a paragraph with a displayed equation,
nor make one a paragraph by itself.

\end{document}               % End of document.

\section{Bench Section 246}
% synthetic bench section 246
% This is a sample LaTeX input file.  (Version of 12 August 2004.)
%
% A '%' character causes TeX to ignore all remaining text on the line,
% and is used for comments like this one.

\documentclass{article}      % Specifies the document class

                             % The preamble begins here.
\title{An Example Document}  % Declares the document's title.
\author{Leslie Lamport}      % Declares the author's name.
\date{January 21, 1994}      % Deleting this command produces today's date.

\newcommand{\ip}[2]{(#1, #2)}
                             % Defines \ip{arg1}{arg2} to mean
                             % (arg1, arg2).

%\newcommand{\ip}[2]{\langle #1 | #2\rangle}
                             % This is an alternative definition of
                             % \ip that is commented out.

\begin{document}             % End of preamble and beginning of text.

\maketitle                   % Produces the title.

This is an example input file.  Comparing it with
the output it generates can show you how to
produce a simple document of your own.

\section{Ordinary Text}      % Produces section heading.  Lower-level
                             % sections are begun with similar
                             % \subsection and \subsubsection commands.

The ends  of words and sentences are marked
  by   spaces. It  doesn't matter how many
spaces    you type; one is as good as 100.  The
end of   a line counts as a space.

One   or more   blank lines denote the  end
of  a paragraph.

Since any number of consecutive spaces are treated
like a single one, the formatting of the input
file makes no difference to
      \LaTeX,                % The \LaTeX command generates the LaTeX logo.
but it makes a difference to you.  When you use
\LaTeX, making your input file as easy to read
as possible will be a great help as you write
your document and when you change it.  This sample
file shows how you can add comments to your own input
file.

Because printing is different from typewriting,
there are a number of things that you have to do
differently when preparing an input file than if
you were just typing the document directly.
Quotation marks like
       ``this''
have to be handled specially, as do quotes within
quotes:
       ``\,`this'            % \, separates the double and single quote.
        is what I just
        wrote, not  `that'\,''.

Dashes come in three sizes: an
       intra-word
dash, a medium dash for number ranges like
       1--2,
and a punctuation
       dash---like
this.

A sentence-ending space should be larger than the
space between words within a sentence.  You
sometimes have to type special commands in
conjunction with punctuation characters to get
this right, as in the following sentence.
       Gnats, gnus, etc.\ all  % `\ ' makes an inter-word space.
       begin with G\@.         % \@ marks end-of-sentence punctuation.
You should check the spaces after periods when
reading your output to make sure you haven't
forgotten any special cases.  Generating an
ellipsis
       \ldots\               % `\ ' is needed after `\ldots' because TeX
                             % ignores spaces after command names like \ldots
                             % made from \ + letters.
                             %
                             % Note how a `%' character causes TeX to ignore
                             % the end of the input line, so these blank lines
                             % do not start a new paragraph.
                             %
with the right spacing around the periods requires
a special command.

\LaTeX\ interprets some common characters as
commands, so you must type special commands to
generate them.  These characters include the
following:
       \$ \& \% \# \{ and \}.

In printing, text is usually emphasized with an
       \emph{italic}
type style.

\begin{em}
   A long segment of text can also be emphasized
   in this way.  Text within such a segment can be
   given \emph{additional} emphasis.
\end{em}

It is sometimes necessary to prevent \LaTeX\ from
breaking a line where it might otherwise do so.
This may be at a space, as between the ``Mr.''\ and
``Jones'' in
       ``Mr.~Jones'',        % ~ produces an unbreakable interword space.
or within a word---especially when the word is a
symbol like
       \mbox{\emph{itemnum}}
that makes little sense when hyphenated across
lines.

Footnotes\footnote{This is an example of a footnote.}
pose no problem.

\LaTeX\ is good at typesetting mathematical formulas
like
       \( x-3y + z = 7 \)
or
       \( a_{1} > x^{2n} + y^{2n} > x' \)
or
       \( \ip{A}{B} = \sum_{i} a_{i} b_{i} \).
The spaces you type in a formula are
ignored.  Remember that a letter like
       $x$                   % $ ... $  and  \( ... \)  are equivalent
is a formula when it denotes a mathematical
symbol, and it should be typed as one.

\section{Displayed Text}

Text is displayed by indenting it from the left
margin.  Quotations are commonly displayed.  There
are short quotations
\begin{quote}
   This is a short quotation.  It consists of a
   single paragraph of text.  See how it is formatted.
\end{quote}
and longer ones.
\begin{quotation}
   This is a longer quotation.  It consists of two
   paragraphs of text, neither of which are
   particularly interesting.

   This is the second paragraph of the quotation.  It
   is just as dull as the first paragraph.
\end{quotation}
Another frequently-displayed structure is a list.
The following is an example of an \emph{itemized}
list.
\begin{itemize}
   \item This is the first item of an itemized list.
         Each item in the list is marked with a ``tick''.
         You don't have to worry about what kind of tick
         mark is used.

   \item This is the second item of the list.  It
         contains another list nested inside it.  The inner
         list is an \emph{enumerated} list.
         \begin{enumerate}
            \item This is the first item of an enumerated
                  list that is nested within the itemized list.

            \item This is the second item of the inner list.
                  \LaTeX\ allows you to nest lists deeper than
                  you really should.
         \end{enumerate}
         This is the rest of the second item of the outer
         list.  It is no more interesting than any other
         part of the item.
   \item This is the third item of the list.
\end{itemize}
You can even display poetry.
\begin{verse}
   There is an environment
    for verse \\             % The \\ command separates lines
   Whose features some poets % within a stanza.
   will curse.

                             % One or more blank lines separate stanzas.

   For instead of making\\
   Them do \emph{all} line breaking, \\
   It allows them to put too many words on a line when they'd rather be
   forced to be terse.
\end{verse}

Mathematical formulas may also be displayed.  A
displayed formula
is
one-line long; multiline
formulas require special formatting instructions.
   \[  \ip{\Gamma}{\psi'} = x'' + y^{2} + z_{i}^{n}\]
Don't start a paragraph with a displayed equation,
nor make one a paragraph by itself.

\end{document}               % End of document.

\section{Bench Section 247}
% synthetic bench section 247
% This is a sample LaTeX input file.  (Version of 12 August 2004.)
%
% A '%' character causes TeX to ignore all remaining text on the line,
% and is used for comments like this one.

\documentclass{article}      % Specifies the document class

                             % The preamble begins here.
\title{An Example Document}  % Declares the document's title.
\author{Leslie Lamport}      % Declares the author's name.
\date{January 21, 1994}      % Deleting this command produces today's date.

\newcommand{\ip}[2]{(#1, #2)}
                             % Defines \ip{arg1}{arg2} to mean
                             % (arg1, arg2).

%\newcommand{\ip}[2]{\langle #1 | #2\rangle}
                             % This is an alternative definition of
                             % \ip that is commented out.

\begin{document}             % End of preamble and beginning of text.

\maketitle                   % Produces the title.

This is an example input file.  Comparing it with
the output it generates can show you how to
produce a simple document of your own.

\section{Ordinary Text}      % Produces section heading.  Lower-level
                             % sections are begun with similar
                             % \subsection and \subsubsection commands.

The ends  of words and sentences are marked
  by   spaces. It  doesn't matter how many
spaces    you type; one is as good as 100.  The
end of   a line counts as a space.

One   or more   blank lines denote the  end
of  a paragraph.

Since any number of consecutive spaces are treated
like a single one, the formatting of the input
file makes no difference to
      \LaTeX,                % The \LaTeX command generates the LaTeX logo.
but it makes a difference to you.  When you use
\LaTeX, making your input file as easy to read
as possible will be a great help as you write
your document and when you change it.  This sample
file shows how you can add comments to your own input
file.

Because printing is different from typewriting,
there are a number of things that you have to do
differently when preparing an input file than if
you were just typing the document directly.
Quotation marks like
       ``this''
have to be handled specially, as do quotes within
quotes:
       ``\,`this'            % \, separates the double and single quote.
        is what I just
        wrote, not  `that'\,''.

Dashes come in three sizes: an
       intra-word
dash, a medium dash for number ranges like
       1--2,
and a punctuation
       dash---like
this.

A sentence-ending space should be larger than the
space between words within a sentence.  You
sometimes have to type special commands in
conjunction with punctuation characters to get
this right, as in the following sentence.
       Gnats, gnus, etc.\ all  % `\ ' makes an inter-word space.
       begin with G\@.         % \@ marks end-of-sentence punctuation.
You should check the spaces after periods when
reading your output to make sure you haven't
forgotten any special cases.  Generating an
ellipsis
       \ldots\               % `\ ' is needed after `\ldots' because TeX
                             % ignores spaces after command names like \ldots
                             % made from \ + letters.
                             %
                             % Note how a `%' character causes TeX to ignore
                             % the end of the input line, so these blank lines
                             % do not start a new paragraph.
                             %
with the right spacing around the periods requires
a special command.

\LaTeX\ interprets some common characters as
commands, so you must type special commands to
generate them.  These characters include the
following:
       \$ \& \% \# \{ and \}.

In printing, text is usually emphasized with an
       \emph{italic}
type style.

\begin{em}
   A long segment of text can also be emphasized
   in this way.  Text within such a segment can be
   given \emph{additional} emphasis.
\end{em}

It is sometimes necessary to prevent \LaTeX\ from
breaking a line where it might otherwise do so.
This may be at a space, as between the ``Mr.''\ and
``Jones'' in
       ``Mr.~Jones'',        % ~ produces an unbreakable interword space.
or within a word---especially when the word is a
symbol like
       \mbox{\emph{itemnum}}
that makes little sense when hyphenated across
lines.

Footnotes\footnote{This is an example of a footnote.}
pose no problem.

\LaTeX\ is good at typesetting mathematical formulas
like
       \( x-3y + z = 7 \)
or
       \( a_{1} > x^{2n} + y^{2n} > x' \)
or
       \( \ip{A}{B} = \sum_{i} a_{i} b_{i} \).
The spaces you type in a formula are
ignored.  Remember that a letter like
       $x$                   % $ ... $  and  \( ... \)  are equivalent
is a formula when it denotes a mathematical
symbol, and it should be typed as one.

\section{Displayed Text}

Text is displayed by indenting it from the left
margin.  Quotations are commonly displayed.  There
are short quotations
\begin{quote}
   This is a short quotation.  It consists of a
   single paragraph of text.  See how it is formatted.
\end{quote}
and longer ones.
\begin{quotation}
   This is a longer quotation.  It consists of two
   paragraphs of text, neither of which are
   particularly interesting.

   This is the second paragraph of the quotation.  It
   is just as dull as the first paragraph.
\end{quotation}
Another frequently-displayed structure is a list.
The following is an example of an \emph{itemized}
list.
\begin{itemize}
   \item This is the first item of an itemized list.
         Each item in the list is marked with a ``tick''.
         You don't have to worry about what kind of tick
         mark is used.

   \item This is the second item of the list.  It
         contains another list nested inside it.  The inner
         list is an \emph{enumerated} list.
         \begin{enumerate}
            \item This is the first item of an enumerated
                  list that is nested within the itemized list.

            \item This is the second item of the inner list.
                  \LaTeX\ allows you to nest lists deeper than
                  you really should.
         \end{enumerate}
         This is the rest of the second item of the outer
         list.  It is no more interesting than any other
         part of the item.
   \item This is the third item of the list.
\end{itemize}
You can even display poetry.
\begin{verse}
   There is an environment
    for verse \\             % The \\ command separates lines
   Whose features some poets % within a stanza.
   will curse.

                             % One or more blank lines separate stanzas.

   For instead of making\\
   Them do \emph{all} line breaking, \\
   It allows them to put too many words on a line when they'd rather be
   forced to be terse.
\end{verse}

Mathematical formulas may also be displayed.  A
displayed formula
is
one-line long; multiline
formulas require special formatting instructions.
   \[  \ip{\Gamma}{\psi'} = x'' + y^{2} + z_{i}^{n}\]
Don't start a paragraph with a displayed equation,
nor make one a paragraph by itself.

\end{document}               % End of document.

\section{Bench Section 248}
% synthetic bench section 248
% This is a sample LaTeX input file.  (Version of 12 August 2004.)
%
% A '%' character causes TeX to ignore all remaining text on the line,
% and is used for comments like this one.

\documentclass{article}      % Specifies the document class

                             % The preamble begins here.
\title{An Example Document}  % Declares the document's title.
\author{Leslie Lamport}      % Declares the author's name.
\date{January 21, 1994}      % Deleting this command produces today's date.

\newcommand{\ip}[2]{(#1, #2)}
                             % Defines \ip{arg1}{arg2} to mean
                             % (arg1, arg2).

%\newcommand{\ip}[2]{\langle #1 | #2\rangle}
                             % This is an alternative definition of
                             % \ip that is commented out.

\begin{document}             % End of preamble and beginning of text.

\maketitle                   % Produces the title.

This is an example input file.  Comparing it with
the output it generates can show you how to
produce a simple document of your own.

\section{Ordinary Text}      % Produces section heading.  Lower-level
                             % sections are begun with similar
                             % \subsection and \subsubsection commands.

The ends  of words and sentences are marked
  by   spaces. It  doesn't matter how many
spaces    you type; one is as good as 100.  The
end of   a line counts as a space.

One   or more   blank lines denote the  end
of  a paragraph.

Since any number of consecutive spaces are treated
like a single one, the formatting of the input
file makes no difference to
      \LaTeX,                % The \LaTeX command generates the LaTeX logo.
but it makes a difference to you.  When you use
\LaTeX, making your input file as easy to read
as possible will be a great help as you write
your document and when you change it.  This sample
file shows how you can add comments to your own input
file.

Because printing is different from typewriting,
there are a number of things that you have to do
differently when preparing an input file than if
you were just typing the document directly.
Quotation marks like
       ``this''
have to be handled specially, as do quotes within
quotes:
       ``\,`this'            % \, separates the double and single quote.
        is what I just
        wrote, not  `that'\,''.

Dashes come in three sizes: an
       intra-word
dash, a medium dash for number ranges like
       1--2,
and a punctuation
       dash---like
this.

A sentence-ending space should be larger than the
space between words within a sentence.  You
sometimes have to type special commands in
conjunction with punctuation characters to get
this right, as in the following sentence.
       Gnats, gnus, etc.\ all  % `\ ' makes an inter-word space.
       begin with G\@.         % \@ marks end-of-sentence punctuation.
You should check the spaces after periods when
reading your output to make sure you haven't
forgotten any special cases.  Generating an
ellipsis
       \ldots\               % `\ ' is needed after `\ldots' because TeX
                             % ignores spaces after command names like \ldots
                             % made from \ + letters.
                             %
                             % Note how a `%' character causes TeX to ignore
                             % the end of the input line, so these blank lines
                             % do not start a new paragraph.
                             %
with the right spacing around the periods requires
a special command.

\LaTeX\ interprets some common characters as
commands, so you must type special commands to
generate them.  These characters include the
following:
       \$ \& \% \# \{ and \}.

In printing, text is usually emphasized with an
       \emph{italic}
type style.

\begin{em}
   A long segment of text can also be emphasized
   in this way.  Text within such a segment can be
   given \emph{additional} emphasis.
\end{em}

It is sometimes necessary to prevent \LaTeX\ from
breaking a line where it might otherwise do so.
This may be at a space, as between the ``Mr.''\ and
``Jones'' in
       ``Mr.~Jones'',        % ~ produces an unbreakable interword space.
or within a word---especially when the word is a
symbol like
       \mbox{\emph{itemnum}}
that makes little sense when hyphenated across
lines.

Footnotes\footnote{This is an example of a footnote.}
pose no problem.

\LaTeX\ is good at typesetting mathematical formulas
like
       \( x-3y + z = 7 \)
or
       \( a_{1} > x^{2n} + y^{2n} > x' \)
or
       \( \ip{A}{B} = \sum_{i} a_{i} b_{i} \).
The spaces you type in a formula are
ignored.  Remember that a letter like
       $x$                   % $ ... $  and  \( ... \)  are equivalent
is a formula when it denotes a mathematical
symbol, and it should be typed as one.

\section{Displayed Text}

Text is displayed by indenting it from the left
margin.  Quotations are commonly displayed.  There
are short quotations
\begin{quote}
   This is a short quotation.  It consists of a
   single paragraph of text.  See how it is formatted.
\end{quote}
and longer ones.
\begin{quotation}
   This is a longer quotation.  It consists of two
   paragraphs of text, neither of which are
   particularly interesting.

   This is the second paragraph of the quotation.  It
   is just as dull as the first paragraph.
\end{quotation}
Another frequently-displayed structure is a list.
The following is an example of an \emph{itemized}
list.
\begin{itemize}
   \item This is the first item of an itemized list.
         Each item in the list is marked with a ``tick''.
         You don't have to worry about what kind of tick
         mark is used.

   \item This is the second item of the list.  It
         contains another list nested inside it.  The inner
         list is an \emph{enumerated} list.
         \begin{enumerate}
            \item This is the first item of an enumerated
                  list that is nested within the itemized list.

            \item This is the second item of the inner list.
                  \LaTeX\ allows you to nest lists deeper than
                  you really should.
         \end{enumerate}
         This is the rest of the second item of the outer
         list.  It is no more interesting than any other
         part of the item.
   \item This is the third item of the list.
\end{itemize}
You can even display poetry.
\begin{verse}
   There is an environment
    for verse \\             % The \\ command separates lines
   Whose features some poets % within a stanza.
   will curse.

                             % One or more blank lines separate stanzas.

   For instead of making\\
   Them do \emph{all} line breaking, \\
   It allows them to put too many words on a line when they'd rather be
   forced to be terse.
\end{verse}

Mathematical formulas may also be displayed.  A
displayed formula
is
one-line long; multiline
formulas require special formatting instructions.
   \[  \ip{\Gamma}{\psi'} = x'' + y^{2} + z_{i}^{n}\]
Don't start a paragraph with a displayed equation,
nor make one a paragraph by itself.

\end{document}               % End of document.

\section{Bench Section 249}
% synthetic bench section 249
% This is a sample LaTeX input file.  (Version of 12 August 2004.)
%
% A '%' character causes TeX to ignore all remaining text on the line,
% and is used for comments like this one.

\documentclass{article}      % Specifies the document class

                             % The preamble begins here.
\title{An Example Document}  % Declares the document's title.
\author{Leslie Lamport}      % Declares the author's name.
\date{January 21, 1994}      % Deleting this command produces today's date.

\newcommand{\ip}[2]{(#1, #2)}
                             % Defines \ip{arg1}{arg2} to mean
                             % (arg1, arg2).

%\newcommand{\ip}[2]{\langle #1 | #2\rangle}
                             % This is an alternative definition of
                             % \ip that is commented out.

\begin{document}             % End of preamble and beginning of text.

\maketitle                   % Produces the title.

This is an example input file.  Comparing it with
the output it generates can show you how to
produce a simple document of your own.

\section{Ordinary Text}      % Produces section heading.  Lower-level
                             % sections are begun with similar
                             % \subsection and \subsubsection commands.

The ends  of words and sentences are marked
  by   spaces. It  doesn't matter how many
spaces    you type; one is as good as 100.  The
end of   a line counts as a space.

One   or more   blank lines denote the  end
of  a paragraph.

Since any number of consecutive spaces are treated
like a single one, the formatting of the input
file makes no difference to
      \LaTeX,                % The \LaTeX command generates the LaTeX logo.
but it makes a difference to you.  When you use
\LaTeX, making your input file as easy to read
as possible will be a great help as you write
your document and when you change it.  This sample
file shows how you can add comments to your own input
file.

Because printing is different from typewriting,
there are a number of things that you have to do
differently when preparing an input file than if
you were just typing the document directly.
Quotation marks like
       ``this''
have to be handled specially, as do quotes within
quotes:
       ``\,`this'            % \, separates the double and single quote.
        is what I just
        wrote, not  `that'\,''.

Dashes come in three sizes: an
       intra-word
dash, a medium dash for number ranges like
       1--2,
and a punctuation
       dash---like
this.

A sentence-ending space should be larger than the
space between words within a sentence.  You
sometimes have to type special commands in
conjunction with punctuation characters to get
this right, as in the following sentence.
       Gnats, gnus, etc.\ all  % `\ ' makes an inter-word space.
       begin with G\@.         % \@ marks end-of-sentence punctuation.
You should check the spaces after periods when
reading your output to make sure you haven't
forgotten any special cases.  Generating an
ellipsis
       \ldots\               % `\ ' is needed after `\ldots' because TeX
                             % ignores spaces after command names like \ldots
                             % made from \ + letters.
                             %
                             % Note how a `%' character causes TeX to ignore
                             % the end of the input line, so these blank lines
                             % do not start a new paragraph.
                             %
with the right spacing around the periods requires
a special command.

\LaTeX\ interprets some common characters as
commands, so you must type special commands to
generate them.  These characters include the
following:
       \$ \& \% \# \{ and \}.

In printing, text is usually emphasized with an
       \emph{italic}
type style.

\begin{em}
   A long segment of text can also be emphasized
   in this way.  Text within such a segment can be
   given \emph{additional} emphasis.
\end{em}

It is sometimes necessary to prevent \LaTeX\ from
breaking a line where it might otherwise do so.
This may be at a space, as between the ``Mr.''\ and
``Jones'' in
       ``Mr.~Jones'',        % ~ produces an unbreakable interword space.
or within a word---especially when the word is a
symbol like
       \mbox{\emph{itemnum}}
that makes little sense when hyphenated across
lines.

Footnotes\footnote{This is an example of a footnote.}
pose no problem.

\LaTeX\ is good at typesetting mathematical formulas
like
       \( x-3y + z = 7 \)
or
       \( a_{1} > x^{2n} + y^{2n} > x' \)
or
       \( \ip{A}{B} = \sum_{i} a_{i} b_{i} \).
The spaces you type in a formula are
ignored.  Remember that a letter like
       $x$                   % $ ... $  and  \( ... \)  are equivalent
is a formula when it denotes a mathematical
symbol, and it should be typed as one.

\section{Displayed Text}

Text is displayed by indenting it from the left
margin.  Quotations are commonly displayed.  There
are short quotations
\begin{quote}
   This is a short quotation.  It consists of a
   single paragraph of text.  See how it is formatted.
\end{quote}
and longer ones.
\begin{quotation}
   This is a longer quotation.  It consists of two
   paragraphs of text, neither of which are
   particularly interesting.

   This is the second paragraph of the quotation.  It
   is just as dull as the first paragraph.
\end{quotation}
Another frequently-displayed structure is a list.
The following is an example of an \emph{itemized}
list.
\begin{itemize}
   \item This is the first item of an itemized list.
         Each item in the list is marked with a ``tick''.
         You don't have to worry about what kind of tick
         mark is used.

   \item This is the second item of the list.  It
         contains another list nested inside it.  The inner
         list is an \emph{enumerated} list.
         \begin{enumerate}
            \item This is the first item of an enumerated
                  list that is nested within the itemized list.

            \item This is the second item of the inner list.
                  \LaTeX\ allows you to nest lists deeper than
                  you really should.
         \end{enumerate}
         This is the rest of the second item of the outer
         list.  It is no more interesting than any other
         part of the item.
   \item This is the third item of the list.
\end{itemize}
You can even display poetry.
\begin{verse}
   There is an environment
    for verse \\             % The \\ command separates lines
   Whose features some poets % within a stanza.
   will curse.

                             % One or more blank lines separate stanzas.

   For instead of making\\
   Them do \emph{all} line breaking, \\
   It allows them to put too many words on a line when they'd rather be
   forced to be terse.
\end{verse}

Mathematical formulas may also be displayed.  A
displayed formula
is
one-line long; multiline
formulas require special formatting instructions.
   \[  \ip{\Gamma}{\psi'} = x'' + y^{2} + z_{i}^{n}\]
Don't start a paragraph with a displayed equation,
nor make one a paragraph by itself.

\end{document}               % End of document.

\section{Bench Section 250}
% synthetic bench section 250
% This is a sample LaTeX input file.  (Version of 12 August 2004.)
%
% A '%' character causes TeX to ignore all remaining text on the line,
% and is used for comments like this one.

\documentclass{article}      % Specifies the document class

                             % The preamble begins here.
\title{An Example Document}  % Declares the document's title.
\author{Leslie Lamport}      % Declares the author's name.
\date{January 21, 1994}      % Deleting this command produces today's date.

\newcommand{\ip}[2]{(#1, #2)}
                             % Defines \ip{arg1}{arg2} to mean
                             % (arg1, arg2).

%\newcommand{\ip}[2]{\langle #1 | #2\rangle}
                             % This is an alternative definition of
                             % \ip that is commented out.

\begin{document}             % End of preamble and beginning of text.

\maketitle                   % Produces the title.

This is an example input file.  Comparing it with
the output it generates can show you how to
produce a simple document of your own.

\section{Ordinary Text}      % Produces section heading.  Lower-level
                             % sections are begun with similar
                             % \subsection and \subsubsection commands.

The ends  of words and sentences are marked
  by   spaces. It  doesn't matter how many
spaces    you type; one is as good as 100.  The
end of   a line counts as a space.

One   or more   blank lines denote the  end
of  a paragraph.

Since any number of consecutive spaces are treated
like a single one, the formatting of the input
file makes no difference to
      \LaTeX,                % The \LaTeX command generates the LaTeX logo.
but it makes a difference to you.  When you use
\LaTeX, making your input file as easy to read
as possible will be a great help as you write
your document and when you change it.  This sample
file shows how you can add comments to your own input
file.

Because printing is different from typewriting,
there are a number of things that you have to do
differently when preparing an input file than if
you were just typing the document directly.
Quotation marks like
       ``this''
have to be handled specially, as do quotes within
quotes:
       ``\,`this'            % \, separates the double and single quote.
        is what I just
        wrote, not  `that'\,''.

Dashes come in three sizes: an
       intra-word
dash, a medium dash for number ranges like
       1--2,
and a punctuation
       dash---like
this.

A sentence-ending space should be larger than the
space between words within a sentence.  You
sometimes have to type special commands in
conjunction with punctuation characters to get
this right, as in the following sentence.
       Gnats, gnus, etc.\ all  % `\ ' makes an inter-word space.
       begin with G\@.         % \@ marks end-of-sentence punctuation.
You should check the spaces after periods when
reading your output to make sure you haven't
forgotten any special cases.  Generating an
ellipsis
       \ldots\               % `\ ' is needed after `\ldots' because TeX
                             % ignores spaces after command names like \ldots
                             % made from \ + letters.
                             %
                             % Note how a `%' character causes TeX to ignore
                             % the end of the input line, so these blank lines
                             % do not start a new paragraph.
                             %
with the right spacing around the periods requires
a special command.

\LaTeX\ interprets some common characters as
commands, so you must type special commands to
generate them.  These characters include the
following:
       \$ \& \% \# \{ and \}.

In printing, text is usually emphasized with an
       \emph{italic}
type style.

\begin{em}
   A long segment of text can also be emphasized
   in this way.  Text within such a segment can be
   given \emph{additional} emphasis.
\end{em}

It is sometimes necessary to prevent \LaTeX\ from
breaking a line where it might otherwise do so.
This may be at a space, as between the ``Mr.''\ and
``Jones'' in
       ``Mr.~Jones'',        % ~ produces an unbreakable interword space.
or within a word---especially when the word is a
symbol like
       \mbox{\emph{itemnum}}
that makes little sense when hyphenated across
lines.

Footnotes\footnote{This is an example of a footnote.}
pose no problem.

\LaTeX\ is good at typesetting mathematical formulas
like
       \( x-3y + z = 7 \)
or
       \( a_{1} > x^{2n} + y^{2n} > x' \)
or
       \( \ip{A}{B} = \sum_{i} a_{i} b_{i} \).
The spaces you type in a formula are
ignored.  Remember that a letter like
       $x$                   % $ ... $  and  \( ... \)  are equivalent
is a formula when it denotes a mathematical
symbol, and it should be typed as one.

\section{Displayed Text}

Text is displayed by indenting it from the left
margin.  Quotations are commonly displayed.  There
are short quotations
\begin{quote}
   This is a short quotation.  It consists of a
   single paragraph of text.  See how it is formatted.
\end{quote}
and longer ones.
\begin{quotation}
   This is a longer quotation.  It consists of two
   paragraphs of text, neither of which are
   particularly interesting.

   This is the second paragraph of the quotation.  It
   is just as dull as the first paragraph.
\end{quotation}
Another frequently-displayed structure is a list.
The following is an example of an \emph{itemized}
list.
\begin{itemize}
   \item This is the first item of an itemized list.
         Each item in the list is marked with a ``tick''.
         You don't have to worry about what kind of tick
         mark is used.

   \item This is the second item of the list.  It
         contains another list nested inside it.  The inner
         list is an \emph{enumerated} list.
         \begin{enumerate}
            \item This is the first item of an enumerated
                  list that is nested within the itemized list.

            \item This is the second item of the inner list.
                  \LaTeX\ allows you to nest lists deeper than
                  you really should.
         \end{enumerate}
         This is the rest of the second item of the outer
         list.  It is no more interesting than any other
         part of the item.
   \item This is the third item of the list.
\end{itemize}
You can even display poetry.
\begin{verse}
   There is an environment
    for verse \\             % The \\ command separates lines
   Whose features some poets % within a stanza.
   will curse.

                             % One or more blank lines separate stanzas.

   For instead of making\\
   Them do \emph{all} line breaking, \\
   It allows them to put too many words on a line when they'd rather be
   forced to be terse.
\end{verse}

Mathematical formulas may also be displayed.  A
displayed formula
is
one-line long; multiline
formulas require special formatting instructions.
   \[  \ip{\Gamma}{\psi'} = x'' + y^{2} + z_{i}^{n}\]
Don't start a paragraph with a displayed equation,
nor make one a paragraph by itself.

\end{document}               % End of document.

\section{Bench Section 251}
% synthetic bench section 251
% This is a sample LaTeX input file.  (Version of 12 August 2004.)
%
% A '%' character causes TeX to ignore all remaining text on the line,
% and is used for comments like this one.

\documentclass{article}      % Specifies the document class

                             % The preamble begins here.
\title{An Example Document}  % Declares the document's title.
\author{Leslie Lamport}      % Declares the author's name.
\date{January 21, 1994}      % Deleting this command produces today's date.

\newcommand{\ip}[2]{(#1, #2)}
                             % Defines \ip{arg1}{arg2} to mean
                             % (arg1, arg2).

%\newcommand{\ip}[2]{\langle #1 | #2\rangle}
                             % This is an alternative definition of
                             % \ip that is commented out.

\begin{document}             % End of preamble and beginning of text.

\maketitle                   % Produces the title.

This is an example input file.  Comparing it with
the output it generates can show you how to
produce a simple document of your own.

\section{Ordinary Text}      % Produces section heading.  Lower-level
                             % sections are begun with similar
                             % \subsection and \subsubsection commands.

The ends  of words and sentences are marked
  by   spaces. It  doesn't matter how many
spaces    you type; one is as good as 100.  The
end of   a line counts as a space.

One   or more   blank lines denote the  end
of  a paragraph.

Since any number of consecutive spaces are treated
like a single one, the formatting of the input
file makes no difference to
      \LaTeX,                % The \LaTeX command generates the LaTeX logo.
but it makes a difference to you.  When you use
\LaTeX, making your input file as easy to read
as possible will be a great help as you write
your document and when you change it.  This sample
file shows how you can add comments to your own input
file.

Because printing is different from typewriting,
there are a number of things that you have to do
differently when preparing an input file than if
you were just typing the document directly.
Quotation marks like
       ``this''
have to be handled specially, as do quotes within
quotes:
       ``\,`this'            % \, separates the double and single quote.
        is what I just
        wrote, not  `that'\,''.

Dashes come in three sizes: an
       intra-word
dash, a medium dash for number ranges like
       1--2,
and a punctuation
       dash---like
this.

A sentence-ending space should be larger than the
space between words within a sentence.  You
sometimes have to type special commands in
conjunction with punctuation characters to get
this right, as in the following sentence.
       Gnats, gnus, etc.\ all  % `\ ' makes an inter-word space.
       begin with G\@.         % \@ marks end-of-sentence punctuation.
You should check the spaces after periods when
reading your output to make sure you haven't
forgotten any special cases.  Generating an
ellipsis
       \ldots\               % `\ ' is needed after `\ldots' because TeX
                             % ignores spaces after command names like \ldots
                             % made from \ + letters.
                             %
                             % Note how a `%' character causes TeX to ignore
                             % the end of the input line, so these blank lines
                             % do not start a new paragraph.
                             %
with the right spacing around the periods requires
a special command.

\LaTeX\ interprets some common characters as
commands, so you must type special commands to
generate them.  These characters include the
following:
       \$ \& \% \# \{ and \}.

In printing, text is usually emphasized with an
       \emph{italic}
type style.

\begin{em}
   A long segment of text can also be emphasized
   in this way.  Text within such a segment can be
   given \emph{additional} emphasis.
\end{em}

It is sometimes necessary to prevent \LaTeX\ from
breaking a line where it might otherwise do so.
This may be at a space, as between the ``Mr.''\ and
``Jones'' in
       ``Mr.~Jones'',        % ~ produces an unbreakable interword space.
or within a word---especially when the word is a
symbol like
       \mbox{\emph{itemnum}}
that makes little sense when hyphenated across
lines.

Footnotes\footnote{This is an example of a footnote.}
pose no problem.

\LaTeX\ is good at typesetting mathematical formulas
like
       \( x-3y + z = 7 \)
or
       \( a_{1} > x^{2n} + y^{2n} > x' \)
or
       \( \ip{A}{B} = \sum_{i} a_{i} b_{i} \).
The spaces you type in a formula are
ignored.  Remember that a letter like
       $x$                   % $ ... $  and  \( ... \)  are equivalent
is a formula when it denotes a mathematical
symbol, and it should be typed as one.

\section{Displayed Text}

Text is displayed by indenting it from the left
margin.  Quotations are commonly displayed.  There
are short quotations
\begin{quote}
   This is a short quotation.  It consists of a
   single paragraph of text.  See how it is formatted.
\end{quote}
and longer ones.
\begin{quotation}
   This is a longer quotation.  It consists of two
   paragraphs of text, neither of which are
   particularly interesting.

   This is the second paragraph of the quotation.  It
   is just as dull as the first paragraph.
\end{quotation}
Another frequently-displayed structure is a list.
The following is an example of an \emph{itemized}
list.
\begin{itemize}
   \item This is the first item of an itemized list.
         Each item in the list is marked with a ``tick''.
         You don't have to worry about what kind of tick
         mark is used.

   \item This is the second item of the list.  It
         contains another list nested inside it.  The inner
         list is an \emph{enumerated} list.
         \begin{enumerate}
            \item This is the first item of an enumerated
                  list that is nested within the itemized list.

            \item This is the second item of the inner list.
                  \LaTeX\ allows you to nest lists deeper than
                  you really should.
         \end{enumerate}
         This is the rest of the second item of the outer
         list.  It is no more interesting than any other
         part of the item.
   \item This is the third item of the list.
\end{itemize}
You can even display poetry.
\begin{verse}
   There is an environment
    for verse \\             % The \\ command separates lines
   Whose features some poets % within a stanza.
   will curse.

                             % One or more blank lines separate stanzas.

   For instead of making\\
   Them do \emph{all} line breaking, \\
   It allows them to put too many words on a line when they'd rather be
   forced to be terse.
\end{verse}

Mathematical formulas may also be displayed.  A
displayed formula
is
one-line long; multiline
formulas require special formatting instructions.
   \[  \ip{\Gamma}{\psi'} = x'' + y^{2} + z_{i}^{n}\]
Don't start a paragraph with a displayed equation,
nor make one a paragraph by itself.

\end{document}               % End of document.

\section{Bench Section 252}
% synthetic bench section 252
% This is a sample LaTeX input file.  (Version of 12 August 2004.)
%
% A '%' character causes TeX to ignore all remaining text on the line,
% and is used for comments like this one.

\documentclass{article}      % Specifies the document class

                             % The preamble begins here.
\title{An Example Document}  % Declares the document's title.
\author{Leslie Lamport}      % Declares the author's name.
\date{January 21, 1994}      % Deleting this command produces today's date.

\newcommand{\ip}[2]{(#1, #2)}
                             % Defines \ip{arg1}{arg2} to mean
                             % (arg1, arg2).

%\newcommand{\ip}[2]{\langle #1 | #2\rangle}
                             % This is an alternative definition of
                             % \ip that is commented out.

\begin{document}             % End of preamble and beginning of text.

\maketitle                   % Produces the title.

This is an example input file.  Comparing it with
the output it generates can show you how to
produce a simple document of your own.

\section{Ordinary Text}      % Produces section heading.  Lower-level
                             % sections are begun with similar
                             % \subsection and \subsubsection commands.

The ends  of words and sentences are marked
  by   spaces. It  doesn't matter how many
spaces    you type; one is as good as 100.  The
end of   a line counts as a space.

One   or more   blank lines denote the  end
of  a paragraph.

Since any number of consecutive spaces are treated
like a single one, the formatting of the input
file makes no difference to
      \LaTeX,                % The \LaTeX command generates the LaTeX logo.
but it makes a difference to you.  When you use
\LaTeX, making your input file as easy to read
as possible will be a great help as you write
your document and when you change it.  This sample
file shows how you can add comments to your own input
file.

Because printing is different from typewriting,
there are a number of things that you have to do
differently when preparing an input file than if
you were just typing the document directly.
Quotation marks like
       ``this''
have to be handled specially, as do quotes within
quotes:
       ``\,`this'            % \, separates the double and single quote.
        is what I just
        wrote, not  `that'\,''.

Dashes come in three sizes: an
       intra-word
dash, a medium dash for number ranges like
       1--2,
and a punctuation
       dash---like
this.

A sentence-ending space should be larger than the
space between words within a sentence.  You
sometimes have to type special commands in
conjunction with punctuation characters to get
this right, as in the following sentence.
       Gnats, gnus, etc.\ all  % `\ ' makes an inter-word space.
       begin with G\@.         % \@ marks end-of-sentence punctuation.
You should check the spaces after periods when
reading your output to make sure you haven't
forgotten any special cases.  Generating an
ellipsis
       \ldots\               % `\ ' is needed after `\ldots' because TeX
                             % ignores spaces after command names like \ldots
                             % made from \ + letters.
                             %
                             % Note how a `%' character causes TeX to ignore
                             % the end of the input line, so these blank lines
                             % do not start a new paragraph.
                             %
with the right spacing around the periods requires
a special command.

\LaTeX\ interprets some common characters as
commands, so you must type special commands to
generate them.  These characters include the
following:
       \$ \& \% \# \{ and \}.

In printing, text is usually emphasized with an
       \emph{italic}
type style.

\begin{em}
   A long segment of text can also be emphasized
   in this way.  Text within such a segment can be
   given \emph{additional} emphasis.
\end{em}

It is sometimes necessary to prevent \LaTeX\ from
breaking a line where it might otherwise do so.
This may be at a space, as between the ``Mr.''\ and
``Jones'' in
       ``Mr.~Jones'',        % ~ produces an unbreakable interword space.
or within a word---especially when the word is a
symbol like
       \mbox{\emph{itemnum}}
that makes little sense when hyphenated across
lines.

Footnotes\footnote{This is an example of a footnote.}
pose no problem.

\LaTeX\ is good at typesetting mathematical formulas
like
       \( x-3y + z = 7 \)
or
       \( a_{1} > x^{2n} + y^{2n} > x' \)
or
       \( \ip{A}{B} = \sum_{i} a_{i} b_{i} \).
The spaces you type in a formula are
ignored.  Remember that a letter like
       $x$                   % $ ... $  and  \( ... \)  are equivalent
is a formula when it denotes a mathematical
symbol, and it should be typed as one.

\section{Displayed Text}

Text is displayed by indenting it from the left
margin.  Quotations are commonly displayed.  There
are short quotations
\begin{quote}
   This is a short quotation.  It consists of a
   single paragraph of text.  See how it is formatted.
\end{quote}
and longer ones.
\begin{quotation}
   This is a longer quotation.  It consists of two
   paragraphs of text, neither of which are
   particularly interesting.

   This is the second paragraph of the quotation.  It
   is just as dull as the first paragraph.
\end{quotation}
Another frequently-displayed structure is a list.
The following is an example of an \emph{itemized}
list.
\begin{itemize}
   \item This is the first item of an itemized list.
         Each item in the list is marked with a ``tick''.
         You don't have to worry about what kind of tick
         mark is used.

   \item This is the second item of the list.  It
         contains another list nested inside it.  The inner
         list is an \emph{enumerated} list.
         \begin{enumerate}
            \item This is the first item of an enumerated
                  list that is nested within the itemized list.

            \item This is the second item of the inner list.
                  \LaTeX\ allows you to nest lists deeper than
                  you really should.
         \end{enumerate}
         This is the rest of the second item of the outer
         list.  It is no more interesting than any other
         part of the item.
   \item This is the third item of the list.
\end{itemize}
You can even display poetry.
\begin{verse}
   There is an environment
    for verse \\             % The \\ command separates lines
   Whose features some poets % within a stanza.
   will curse.

                             % One or more blank lines separate stanzas.

   For instead of making\\
   Them do \emph{all} line breaking, \\
   It allows them to put too many words on a line when they'd rather be
   forced to be terse.
\end{verse}

Mathematical formulas may also be displayed.  A
displayed formula
is
one-line long; multiline
formulas require special formatting instructions.
   \[  \ip{\Gamma}{\psi'} = x'' + y^{2} + z_{i}^{n}\]
Don't start a paragraph with a displayed equation,
nor make one a paragraph by itself.

\end{document}               % End of document.

\section{Bench Section 253}
% synthetic bench section 253
% This is a sample LaTeX input file.  (Version of 12 August 2004.)
%
% A '%' character causes TeX to ignore all remaining text on the line,
% and is used for comments like this one.

\documentclass{article}      % Specifies the document class

                             % The preamble begins here.
\title{An Example Document}  % Declares the document's title.
\author{Leslie Lamport}      % Declares the author's name.
\date{January 21, 1994}      % Deleting this command produces today's date.

\newcommand{\ip}[2]{(#1, #2)}
                             % Defines \ip{arg1}{arg2} to mean
                             % (arg1, arg2).

%\newcommand{\ip}[2]{\langle #1 | #2\rangle}
                             % This is an alternative definition of
                             % \ip that is commented out.

\begin{document}             % End of preamble and beginning of text.

\maketitle                   % Produces the title.

This is an example input file.  Comparing it with
the output it generates can show you how to
produce a simple document of your own.

\section{Ordinary Text}      % Produces section heading.  Lower-level
                             % sections are begun with similar
                             % \subsection and \subsubsection commands.

The ends  of words and sentences are marked
  by   spaces. It  doesn't matter how many
spaces    you type; one is as good as 100.  The
end of   a line counts as a space.

One   or more   blank lines denote the  end
of  a paragraph.

Since any number of consecutive spaces are treated
like a single one, the formatting of the input
file makes no difference to
      \LaTeX,                % The \LaTeX command generates the LaTeX logo.
but it makes a difference to you.  When you use
\LaTeX, making your input file as easy to read
as possible will be a great help as you write
your document and when you change it.  This sample
file shows how you can add comments to your own input
file.

Because printing is different from typewriting,
there are a number of things that you have to do
differently when preparing an input file than if
you were just typing the document directly.
Quotation marks like
       ``this''
have to be handled specially, as do quotes within
quotes:
       ``\,`this'            % \, separates the double and single quote.
        is what I just
        wrote, not  `that'\,''.

Dashes come in three sizes: an
       intra-word
dash, a medium dash for number ranges like
       1--2,
and a punctuation
       dash---like
this.

A sentence-ending space should be larger than the
space between words within a sentence.  You
sometimes have to type special commands in
conjunction with punctuation characters to get
this right, as in the following sentence.
       Gnats, gnus, etc.\ all  % `\ ' makes an inter-word space.
       begin with G\@.         % \@ marks end-of-sentence punctuation.
You should check the spaces after periods when
reading your output to make sure you haven't
forgotten any special cases.  Generating an
ellipsis
       \ldots\               % `\ ' is needed after `\ldots' because TeX
                             % ignores spaces after command names like \ldots
                             % made from \ + letters.
                             %
                             % Note how a `%' character causes TeX to ignore
                             % the end of the input line, so these blank lines
                             % do not start a new paragraph.
                             %
with the right spacing around the periods requires
a special command.

\LaTeX\ interprets some common characters as
commands, so you must type special commands to
generate them.  These characters include the
following:
       \$ \& \% \# \{ and \}.

In printing, text is usually emphasized with an
       \emph{italic}
type style.

\begin{em}
   A long segment of text can also be emphasized
   in this way.  Text within such a segment can be
   given \emph{additional} emphasis.
\end{em}

It is sometimes necessary to prevent \LaTeX\ from
breaking a line where it might otherwise do so.
This may be at a space, as between the ``Mr.''\ and
``Jones'' in
       ``Mr.~Jones'',        % ~ produces an unbreakable interword space.
or within a word---especially when the word is a
symbol like
       \mbox{\emph{itemnum}}
that makes little sense when hyphenated across
lines.

Footnotes\footnote{This is an example of a footnote.}
pose no problem.

\LaTeX\ is good at typesetting mathematical formulas
like
       \( x-3y + z = 7 \)
or
       \( a_{1} > x^{2n} + y^{2n} > x' \)
or
       \( \ip{A}{B} = \sum_{i} a_{i} b_{i} \).
The spaces you type in a formula are
ignored.  Remember that a letter like
       $x$                   % $ ... $  and  \( ... \)  are equivalent
is a formula when it denotes a mathematical
symbol, and it should be typed as one.

\section{Displayed Text}

Text is displayed by indenting it from the left
margin.  Quotations are commonly displayed.  There
are short quotations
\begin{quote}
   This is a short quotation.  It consists of a
   single paragraph of text.  See how it is formatted.
\end{quote}
and longer ones.
\begin{quotation}
   This is a longer quotation.  It consists of two
   paragraphs of text, neither of which are
   particularly interesting.

   This is the second paragraph of the quotation.  It
   is just as dull as the first paragraph.
\end{quotation}
Another frequently-displayed structure is a list.
The following is an example of an \emph{itemized}
list.
\begin{itemize}
   \item This is the first item of an itemized list.
         Each item in the list is marked with a ``tick''.
         You don't have to worry about what kind of tick
         mark is used.

   \item This is the second item of the list.  It
         contains another list nested inside it.  The inner
         list is an \emph{enumerated} list.
         \begin{enumerate}
            \item This is the first item of an enumerated
                  list that is nested within the itemized list.

            \item This is the second item of the inner list.
                  \LaTeX\ allows you to nest lists deeper than
                  you really should.
         \end{enumerate}
         This is the rest of the second item of the outer
         list.  It is no more interesting than any other
         part of the item.
   \item This is the third item of the list.
\end{itemize}
You can even display poetry.
\begin{verse}
   There is an environment
    for verse \\             % The \\ command separates lines
   Whose features some poets % within a stanza.
   will curse.

                             % One or more blank lines separate stanzas.

   For instead of making\\
   Them do \emph{all} line breaking, \\
   It allows them to put too many words on a line when they'd rather be
   forced to be terse.
\end{verse}

Mathematical formulas may also be displayed.  A
displayed formula
is
one-line long; multiline
formulas require special formatting instructions.
   \[  \ip{\Gamma}{\psi'} = x'' + y^{2} + z_{i}^{n}\]
Don't start a paragraph with a displayed equation,
nor make one a paragraph by itself.

\end{document}               % End of document.

\section{Bench Section 254}
% synthetic bench section 254
% This is a sample LaTeX input file.  (Version of 12 August 2004.)
%
% A '%' character causes TeX to ignore all remaining text on the line,
% and is used for comments like this one.

\documentclass{article}      % Specifies the document class

                             % The preamble begins here.
\title{An Example Document}  % Declares the document's title.
\author{Leslie Lamport}      % Declares the author's name.
\date{January 21, 1994}      % Deleting this command produces today's date.

\newcommand{\ip}[2]{(#1, #2)}
                             % Defines \ip{arg1}{arg2} to mean
                             % (arg1, arg2).

%\newcommand{\ip}[2]{\langle #1 | #2\rangle}
                             % This is an alternative definition of
                             % \ip that is commented out.

\begin{document}             % End of preamble and beginning of text.

\maketitle                   % Produces the title.

This is an example input file.  Comparing it with
the output it generates can show you how to
produce a simple document of your own.

\section{Ordinary Text}      % Produces section heading.  Lower-level
                             % sections are begun with similar
                             % \subsection and \subsubsection commands.

The ends  of words and sentences are marked
  by   spaces. It  doesn't matter how many
spaces    you type; one is as good as 100.  The
end of   a line counts as a space.

One   or more   blank lines denote the  end
of  a paragraph.

Since any number of consecutive spaces are treated
like a single one, the formatting of the input
file makes no difference to
      \LaTeX,                % The \LaTeX command generates the LaTeX logo.
but it makes a difference to you.  When you use
\LaTeX, making your input file as easy to read
as possible will be a great help as you write
your document and when you change it.  This sample
file shows how you can add comments to your own input
file.

Because printing is different from typewriting,
there are a number of things that you have to do
differently when preparing an input file than if
you were just typing the document directly.
Quotation marks like
       ``this''
have to be handled specially, as do quotes within
quotes:
       ``\,`this'            % \, separates the double and single quote.
        is what I just
        wrote, not  `that'\,''.

Dashes come in three sizes: an
       intra-word
dash, a medium dash for number ranges like
       1--2,
and a punctuation
       dash---like
this.

A sentence-ending space should be larger than the
space between words within a sentence.  You
sometimes have to type special commands in
conjunction with punctuation characters to get
this right, as in the following sentence.
       Gnats, gnus, etc.\ all  % `\ ' makes an inter-word space.
       begin with G\@.         % \@ marks end-of-sentence punctuation.
You should check the spaces after periods when
reading your output to make sure you haven't
forgotten any special cases.  Generating an
ellipsis
       \ldots\               % `\ ' is needed after `\ldots' because TeX
                             % ignores spaces after command names like \ldots
                             % made from \ + letters.
                             %
                             % Note how a `%' character causes TeX to ignore
                             % the end of the input line, so these blank lines
                             % do not start a new paragraph.
                             %
with the right spacing around the periods requires
a special command.

\LaTeX\ interprets some common characters as
commands, so you must type special commands to
generate them.  These characters include the
following:
       \$ \& \% \# \{ and \}.

In printing, text is usually emphasized with an
       \emph{italic}
type style.

\begin{em}
   A long segment of text can also be emphasized
   in this way.  Text within such a segment can be
   given \emph{additional} emphasis.
\end{em}

It is sometimes necessary to prevent \LaTeX\ from
breaking a line where it might otherwise do so.
This may be at a space, as between the ``Mr.''\ and
``Jones'' in
       ``Mr.~Jones'',        % ~ produces an unbreakable interword space.
or within a word---especially when the word is a
symbol like
       \mbox{\emph{itemnum}}
that makes little sense when hyphenated across
lines.

Footnotes\footnote{This is an example of a footnote.}
pose no problem.

\LaTeX\ is good at typesetting mathematical formulas
like
       \( x-3y + z = 7 \)
or
       \( a_{1} > x^{2n} + y^{2n} > x' \)
or
       \( \ip{A}{B} = \sum_{i} a_{i} b_{i} \).
The spaces you type in a formula are
ignored.  Remember that a letter like
       $x$                   % $ ... $  and  \( ... \)  are equivalent
is a formula when it denotes a mathematical
symbol, and it should be typed as one.

\section{Displayed Text}

Text is displayed by indenting it from the left
margin.  Quotations are commonly displayed.  There
are short quotations
\begin{quote}
   This is a short quotation.  It consists of a
   single paragraph of text.  See how it is formatted.
\end{quote}
and longer ones.
\begin{quotation}
   This is a longer quotation.  It consists of two
   paragraphs of text, neither of which are
   particularly interesting.

   This is the second paragraph of the quotation.  It
   is just as dull as the first paragraph.
\end{quotation}
Another frequently-displayed structure is a list.
The following is an example of an \emph{itemized}
list.
\begin{itemize}
   \item This is the first item of an itemized list.
         Each item in the list is marked with a ``tick''.
         You don't have to worry about what kind of tick
         mark is used.

   \item This is the second item of the list.  It
         contains another list nested inside it.  The inner
         list is an \emph{enumerated} list.
         \begin{enumerate}
            \item This is the first item of an enumerated
                  list that is nested within the itemized list.

            \item This is the second item of the inner list.
                  \LaTeX\ allows you to nest lists deeper than
                  you really should.
         \end{enumerate}
         This is the rest of the second item of the outer
         list.  It is no more interesting than any other
         part of the item.
   \item This is the third item of the list.
\end{itemize}
You can even display poetry.
\begin{verse}
   There is an environment
    for verse \\             % The \\ command separates lines
   Whose features some poets % within a stanza.
   will curse.

                             % One or more blank lines separate stanzas.

   For instead of making\\
   Them do \emph{all} line breaking, \\
   It allows them to put too many words on a line when they'd rather be
   forced to be terse.
\end{verse}

Mathematical formulas may also be displayed.  A
displayed formula
is
one-line long; multiline
formulas require special formatting instructions.
   \[  \ip{\Gamma}{\psi'} = x'' + y^{2} + z_{i}^{n}\]
Don't start a paragraph with a displayed equation,
nor make one a paragraph by itself.

\end{document}               % End of document.

\section{Bench Section 255}
% synthetic bench section 255
% This is a sample LaTeX input file.  (Version of 12 August 2004.)
%
% A '%' character causes TeX to ignore all remaining text on the line,
% and is used for comments like this one.

\documentclass{article}      % Specifies the document class

                             % The preamble begins here.
\title{An Example Document}  % Declares the document's title.
\author{Leslie Lamport}      % Declares the author's name.
\date{January 21, 1994}      % Deleting this command produces today's date.

\newcommand{\ip}[2]{(#1, #2)}
                             % Defines \ip{arg1}{arg2} to mean
                             % (arg1, arg2).

%\newcommand{\ip}[2]{\langle #1 | #2\rangle}
                             % This is an alternative definition of
                             % \ip that is commented out.

\begin{document}             % End of preamble and beginning of text.

\maketitle                   % Produces the title.

This is an example input file.  Comparing it with
the output it generates can show you how to
produce a simple document of your own.

\section{Ordinary Text}      % Produces section heading.  Lower-level
                             % sections are begun with similar
                             % \subsection and \subsubsection commands.

The ends  of words and sentences are marked
  by   spaces. It  doesn't matter how many
spaces    you type; one is as good as 100.  The
end of   a line counts as a space.

One   or more   blank lines denote the  end
of  a paragraph.

Since any number of consecutive spaces are treated
like a single one, the formatting of the input
file makes no difference to
      \LaTeX,                % The \LaTeX command generates the LaTeX logo.
but it makes a difference to you.  When you use
\LaTeX, making your input file as easy to read
as possible will be a great help as you write
your document and when you change it.  This sample
file shows how you can add comments to your own input
file.

Because printing is different from typewriting,
there are a number of things that you have to do
differently when preparing an input file than if
you were just typing the document directly.
Quotation marks like
       ``this''
have to be handled specially, as do quotes within
quotes:
       ``\,`this'            % \, separates the double and single quote.
        is what I just
        wrote, not  `that'\,''.

Dashes come in three sizes: an
       intra-word
dash, a medium dash for number ranges like
       1--2,
and a punctuation
       dash---like
this.

A sentence-ending space should be larger than the
space between words within a sentence.  You
sometimes have to type special commands in
conjunction with punctuation characters to get
this right, as in the following sentence.
       Gnats, gnus, etc.\ all  % `\ ' makes an inter-word space.
       begin with G\@.         % \@ marks end-of-sentence punctuation.
You should check the spaces after periods when
reading your output to make sure you haven't
forgotten any special cases.  Generating an
ellipsis
       \ldots\               % `\ ' is needed after `\ldots' because TeX
                             % ignores spaces after command names like \ldots
                             % made from \ + letters.
                             %
                             % Note how a `%' character causes TeX to ignore
                             % the end of the input line, so these blank lines
                             % do not start a new paragraph.
                             %
with the right spacing around the periods requires
a special command.

\LaTeX\ interprets some common characters as
commands, so you must type special commands to
generate them.  These characters include the
following:
       \$ \& \% \# \{ and \}.

In printing, text is usually emphasized with an
       \emph{italic}
type style.

\begin{em}
   A long segment of text can also be emphasized
   in this way.  Text within such a segment can be
   given \emph{additional} emphasis.
\end{em}

It is sometimes necessary to prevent \LaTeX\ from
breaking a line where it might otherwise do so.
This may be at a space, as between the ``Mr.''\ and
``Jones'' in
       ``Mr.~Jones'',        % ~ produces an unbreakable interword space.
or within a word---especially when the word is a
symbol like
       \mbox{\emph{itemnum}}
that makes little sense when hyphenated across
lines.

Footnotes\footnote{This is an example of a footnote.}
pose no problem.

\LaTeX\ is good at typesetting mathematical formulas
like
       \( x-3y + z = 7 \)
or
       \( a_{1} > x^{2n} + y^{2n} > x' \)
or
       \( \ip{A}{B} = \sum_{i} a_{i} b_{i} \).
The spaces you type in a formula are
ignored.  Remember that a letter like
       $x$                   % $ ... $  and  \( ... \)  are equivalent
is a formula when it denotes a mathematical
symbol, and it should be typed as one.

\section{Displayed Text}

Text is displayed by indenting it from the left
margin.  Quotations are commonly displayed.  There
are short quotations
\begin{quote}
   This is a short quotation.  It consists of a
   single paragraph of text.  See how it is formatted.
\end{quote}
and longer ones.
\begin{quotation}
   This is a longer quotation.  It consists of two
   paragraphs of text, neither of which are
   particularly interesting.

   This is the second paragraph of the quotation.  It
   is just as dull as the first paragraph.
\end{quotation}
Another frequently-displayed structure is a list.
The following is an example of an \emph{itemized}
list.
\begin{itemize}
   \item This is the first item of an itemized list.
         Each item in the list is marked with a ``tick''.
         You don't have to worry about what kind of tick
         mark is used.

   \item This is the second item of the list.  It
         contains another list nested inside it.  The inner
         list is an \emph{enumerated} list.
         \begin{enumerate}
            \item This is the first item of an enumerated
                  list that is nested within the itemized list.

            \item This is the second item of the inner list.
                  \LaTeX\ allows you to nest lists deeper than
                  you really should.
         \end{enumerate}
         This is the rest of the second item of the outer
         list.  It is no more interesting than any other
         part of the item.
   \item This is the third item of the list.
\end{itemize}
You can even display poetry.
\begin{verse}
   There is an environment
    for verse \\             % The \\ command separates lines
   Whose features some poets % within a stanza.
   will curse.

                             % One or more blank lines separate stanzas.

   For instead of making\\
   Them do \emph{all} line breaking, \\
   It allows them to put too many words on a line when they'd rather be
   forced to be terse.
\end{verse}

Mathematical formulas may also be displayed.  A
displayed formula
is
one-line long; multiline
formulas require special formatting instructions.
   \[  \ip{\Gamma}{\psi'} = x'' + y^{2} + z_{i}^{n}\]
Don't start a paragraph with a displayed equation,
nor make one a paragraph by itself.

\end{document}               % End of document.

\section{Bench Section 256}
% synthetic bench section 256
% This is a sample LaTeX input file.  (Version of 12 August 2004.)
%
% A '%' character causes TeX to ignore all remaining text on the line,
% and is used for comments like this one.

\documentclass{article}      % Specifies the document class

                             % The preamble begins here.
\title{An Example Document}  % Declares the document's title.
\author{Leslie Lamport}      % Declares the author's name.
\date{January 21, 1994}      % Deleting this command produces today's date.

\newcommand{\ip}[2]{(#1, #2)}
                             % Defines \ip{arg1}{arg2} to mean
                             % (arg1, arg2).

%\newcommand{\ip}[2]{\langle #1 | #2\rangle}
                             % This is an alternative definition of
                             % \ip that is commented out.

\begin{document}             % End of preamble and beginning of text.

\maketitle                   % Produces the title.

This is an example input file.  Comparing it with
the output it generates can show you how to
produce a simple document of your own.

\section{Ordinary Text}      % Produces section heading.  Lower-level
                             % sections are begun with similar
                             % \subsection and \subsubsection commands.

The ends  of words and sentences are marked
  by   spaces. It  doesn't matter how many
spaces    you type; one is as good as 100.  The
end of   a line counts as a space.

One   or more   blank lines denote the  end
of  a paragraph.

Since any number of consecutive spaces are treated
like a single one, the formatting of the input
file makes no difference to
      \LaTeX,                % The \LaTeX command generates the LaTeX logo.
but it makes a difference to you.  When you use
\LaTeX, making your input file as easy to read
as possible will be a great help as you write
your document and when you change it.  This sample
file shows how you can add comments to your own input
file.

Because printing is different from typewriting,
there are a number of things that you have to do
differently when preparing an input file than if
you were just typing the document directly.
Quotation marks like
       ``this''
have to be handled specially, as do quotes within
quotes:
       ``\,`this'            % \, separates the double and single quote.
        is what I just
        wrote, not  `that'\,''.

Dashes come in three sizes: an
       intra-word
dash, a medium dash for number ranges like
       1--2,
and a punctuation
       dash---like
this.

A sentence-ending space should be larger than the
space between words within a sentence.  You
sometimes have to type special commands in
conjunction with punctuation characters to get
this right, as in the following sentence.
       Gnats, gnus, etc.\ all  % `\ ' makes an inter-word space.
       begin with G\@.         % \@ marks end-of-sentence punctuation.
You should check the spaces after periods when
reading your output to make sure you haven't
forgotten any special cases.  Generating an
ellipsis
       \ldots\               % `\ ' is needed after `\ldots' because TeX
                             % ignores spaces after command names like \ldots
                             % made from \ + letters.
                             %
                             % Note how a `%' character causes TeX to ignore
                             % the end of the input line, so these blank lines
                             % do not start a new paragraph.
                             %
with the right spacing around the periods requires
a special command.

\LaTeX\ interprets some common characters as
commands, so you must type special commands to
generate them.  These characters include the
following:
       \$ \& \% \# \{ and \}.

In printing, text is usually emphasized with an
       \emph{italic}
type style.

\begin{em}
   A long segment of text can also be emphasized
   in this way.  Text within such a segment can be
   given \emph{additional} emphasis.
\end{em}

It is sometimes necessary to prevent \LaTeX\ from
breaking a line where it might otherwise do so.
This may be at a space, as between the ``Mr.''\ and
``Jones'' in
       ``Mr.~Jones'',        % ~ produces an unbreakable interword space.
or within a word---especially when the word is a
symbol like
       \mbox{\emph{itemnum}}
that makes little sense when hyphenated across
lines.

Footnotes\footnote{This is an example of a footnote.}
pose no problem.

\LaTeX\ is good at typesetting mathematical formulas
like
       \( x-3y + z = 7 \)
or
       \( a_{1} > x^{2n} + y^{2n} > x' \)
or
       \( \ip{A}{B} = \sum_{i} a_{i} b_{i} \).
The spaces you type in a formula are
ignored.  Remember that a letter like
       $x$                   % $ ... $  and  \( ... \)  are equivalent
is a formula when it denotes a mathematical
symbol, and it should be typed as one.

\section{Displayed Text}

Text is displayed by indenting it from the left
margin.  Quotations are commonly displayed.  There
are short quotations
\begin{quote}
   This is a short quotation.  It consists of a
   single paragraph of text.  See how it is formatted.
\end{quote}
and longer ones.
\begin{quotation}
   This is a longer quotation.  It consists of two
   paragraphs of text, neither of which are
   particularly interesting.

   This is the second paragraph of the quotation.  It
   is just as dull as the first paragraph.
\end{quotation}
Another frequently-displayed structure is a list.
The following is an example of an \emph{itemized}
list.
\begin{itemize}
   \item This is the first item of an itemized list.
         Each item in the list is marked with a ``tick''.
         You don't have to worry about what kind of tick
         mark is used.

   \item This is the second item of the list.  It
         contains another list nested inside it.  The inner
         list is an \emph{enumerated} list.
         \begin{enumerate}
            \item This is the first item of an enumerated
                  list that is nested within the itemized list.

            \item This is the second item of the inner list.
                  \LaTeX\ allows you to nest lists deeper than
                  you really should.
         \end{enumerate}
         This is the rest of the second item of the outer
         list.  It is no more interesting than any other
         part of the item.
   \item This is the third item of the list.
\end{itemize}
You can even display poetry.
\begin{verse}
   There is an environment
    for verse \\             % The \\ command separates lines
   Whose features some poets % within a stanza.
   will curse.

                             % One or more blank lines separate stanzas.

   For instead of making\\
   Them do \emph{all} line breaking, \\
   It allows them to put too many words on a line when they'd rather be
   forced to be terse.
\end{verse}

Mathematical formulas may also be displayed.  A
displayed formula
is
one-line long; multiline
formulas require special formatting instructions.
   \[  \ip{\Gamma}{\psi'} = x'' + y^{2} + z_{i}^{n}\]
Don't start a paragraph with a displayed equation,
nor make one a paragraph by itself.

\end{document}               % End of document.

\section{Bench Section 257}
% synthetic bench section 257
% This is a sample LaTeX input file.  (Version of 12 August 2004.)
%
% A '%' character causes TeX to ignore all remaining text on the line,
% and is used for comments like this one.

\documentclass{article}      % Specifies the document class

                             % The preamble begins here.
\title{An Example Document}  % Declares the document's title.
\author{Leslie Lamport}      % Declares the author's name.
\date{January 21, 1994}      % Deleting this command produces today's date.

\newcommand{\ip}[2]{(#1, #2)}
                             % Defines \ip{arg1}{arg2} to mean
                             % (arg1, arg2).

%\newcommand{\ip}[2]{\langle #1 | #2\rangle}
                             % This is an alternative definition of
                             % \ip that is commented out.

\begin{document}             % End of preamble and beginning of text.

\maketitle                   % Produces the title.

This is an example input file.  Comparing it with
the output it generates can show you how to
produce a simple document of your own.

\section{Ordinary Text}      % Produces section heading.  Lower-level
                             % sections are begun with similar
                             % \subsection and \subsubsection commands.

The ends  of words and sentences are marked
  by   spaces. It  doesn't matter how many
spaces    you type; one is as good as 100.  The
end of   a line counts as a space.

One   or more   blank lines denote the  end
of  a paragraph.

Since any number of consecutive spaces are treated
like a single one, the formatting of the input
file makes no difference to
      \LaTeX,                % The \LaTeX command generates the LaTeX logo.
but it makes a difference to you.  When you use
\LaTeX, making your input file as easy to read
as possible will be a great help as you write
your document and when you change it.  This sample
file shows how you can add comments to your own input
file.

Because printing is different from typewriting,
there are a number of things that you have to do
differently when preparing an input file than if
you were just typing the document directly.
Quotation marks like
       ``this''
have to be handled specially, as do quotes within
quotes:
       ``\,`this'            % \, separates the double and single quote.
        is what I just
        wrote, not  `that'\,''.

Dashes come in three sizes: an
       intra-word
dash, a medium dash for number ranges like
       1--2,
and a punctuation
       dash---like
this.

A sentence-ending space should be larger than the
space between words within a sentence.  You
sometimes have to type special commands in
conjunction with punctuation characters to get
this right, as in the following sentence.
       Gnats, gnus, etc.\ all  % `\ ' makes an inter-word space.
       begin with G\@.         % \@ marks end-of-sentence punctuation.
You should check the spaces after periods when
reading your output to make sure you haven't
forgotten any special cases.  Generating an
ellipsis
       \ldots\               % `\ ' is needed after `\ldots' because TeX
                             % ignores spaces after command names like \ldots
                             % made from \ + letters.
                             %
                             % Note how a `%' character causes TeX to ignore
                             % the end of the input line, so these blank lines
                             % do not start a new paragraph.
                             %
with the right spacing around the periods requires
a special command.

\LaTeX\ interprets some common characters as
commands, so you must type special commands to
generate them.  These characters include the
following:
       \$ \& \% \# \{ and \}.

In printing, text is usually emphasized with an
       \emph{italic}
type style.

\begin{em}
   A long segment of text can also be emphasized
   in this way.  Text within such a segment can be
   given \emph{additional} emphasis.
\end{em}

It is sometimes necessary to prevent \LaTeX\ from
breaking a line where it might otherwise do so.
This may be at a space, as between the ``Mr.''\ and
``Jones'' in
       ``Mr.~Jones'',        % ~ produces an unbreakable interword space.
or within a word---especially when the word is a
symbol like
       \mbox{\emph{itemnum}}
that makes little sense when hyphenated across
lines.

Footnotes\footnote{This is an example of a footnote.}
pose no problem.

\LaTeX\ is good at typesetting mathematical formulas
like
       \( x-3y + z = 7 \)
or
       \( a_{1} > x^{2n} + y^{2n} > x' \)
or
       \( \ip{A}{B} = \sum_{i} a_{i} b_{i} \).
The spaces you type in a formula are
ignored.  Remember that a letter like
       $x$                   % $ ... $  and  \( ... \)  are equivalent
is a formula when it denotes a mathematical
symbol, and it should be typed as one.

\section{Displayed Text}

Text is displayed by indenting it from the left
margin.  Quotations are commonly displayed.  There
are short quotations
\begin{quote}
   This is a short quotation.  It consists of a
   single paragraph of text.  See how it is formatted.
\end{quote}
and longer ones.
\begin{quotation}
   This is a longer quotation.  It consists of two
   paragraphs of text, neither of which are
   particularly interesting.

   This is the second paragraph of the quotation.  It
   is just as dull as the first paragraph.
\end{quotation}
Another frequently-displayed structure is a list.
The following is an example of an \emph{itemized}
list.
\begin{itemize}
   \item This is the first item of an itemized list.
         Each item in the list is marked with a ``tick''.
         You don't have to worry about what kind of tick
         mark is used.

   \item This is the second item of the list.  It
         contains another list nested inside it.  The inner
         list is an \emph{enumerated} list.
         \begin{enumerate}
            \item This is the first item of an enumerated
                  list that is nested within the itemized list.

            \item This is the second item of the inner list.
                  \LaTeX\ allows you to nest lists deeper than
                  you really should.
         \end{enumerate}
         This is the rest of the second item of the outer
         list.  It is no more interesting than any other
         part of the item.
   \item This is the third item of the list.
\end{itemize}
You can even display poetry.
\begin{verse}
   There is an environment
    for verse \\             % The \\ command separates lines
   Whose features some poets % within a stanza.
   will curse.

                             % One or more blank lines separate stanzas.

   For instead of making\\
   Them do \emph{all} line breaking, \\
   It allows them to put too many words on a line when they'd rather be
   forced to be terse.
\end{verse}

Mathematical formulas may also be displayed.  A
displayed formula
is
one-line long; multiline
formulas require special formatting instructions.
   \[  \ip{\Gamma}{\psi'} = x'' + y^{2} + z_{i}^{n}\]
Don't start a paragraph with a displayed equation,
nor make one a paragraph by itself.

\end{document}               % End of document.

\section{Bench Section 258}
% synthetic bench section 258
% This is a sample LaTeX input file.  (Version of 12 August 2004.)
%
% A '%' character causes TeX to ignore all remaining text on the line,
% and is used for comments like this one.

\documentclass{article}      % Specifies the document class

                             % The preamble begins here.
\title{An Example Document}  % Declares the document's title.
\author{Leslie Lamport}      % Declares the author's name.
\date{January 21, 1994}      % Deleting this command produces today's date.

\newcommand{\ip}[2]{(#1, #2)}
                             % Defines \ip{arg1}{arg2} to mean
                             % (arg1, arg2).

%\newcommand{\ip}[2]{\langle #1 | #2\rangle}
                             % This is an alternative definition of
                             % \ip that is commented out.

\begin{document}             % End of preamble and beginning of text.

\maketitle                   % Produces the title.

This is an example input file.  Comparing it with
the output it generates can show you how to
produce a simple document of your own.

\section{Ordinary Text}      % Produces section heading.  Lower-level
                             % sections are begun with similar
                             % \subsection and \subsubsection commands.

The ends  of words and sentences are marked
  by   spaces. It  doesn't matter how many
spaces    you type; one is as good as 100.  The
end of   a line counts as a space.

One   or more   blank lines denote the  end
of  a paragraph.

Since any number of consecutive spaces are treated
like a single one, the formatting of the input
file makes no difference to
      \LaTeX,                % The \LaTeX command generates the LaTeX logo.
but it makes a difference to you.  When you use
\LaTeX, making your input file as easy to read
as possible will be a great help as you write
your document and when you change it.  This sample
file shows how you can add comments to your own input
file.

Because printing is different from typewriting,
there are a number of things that you have to do
differently when preparing an input file than if
you were just typing the document directly.
Quotation marks like
       ``this''
have to be handled specially, as do quotes within
quotes:
       ``\,`this'            % \, separates the double and single quote.
        is what I just
        wrote, not  `that'\,''.

Dashes come in three sizes: an
       intra-word
dash, a medium dash for number ranges like
       1--2,
and a punctuation
       dash---like
this.

A sentence-ending space should be larger than the
space between words within a sentence.  You
sometimes have to type special commands in
conjunction with punctuation characters to get
this right, as in the following sentence.
       Gnats, gnus, etc.\ all  % `\ ' makes an inter-word space.
       begin with G\@.         % \@ marks end-of-sentence punctuation.
You should check the spaces after periods when
reading your output to make sure you haven't
forgotten any special cases.  Generating an
ellipsis
       \ldots\               % `\ ' is needed after `\ldots' because TeX
                             % ignores spaces after command names like \ldots
                             % made from \ + letters.
                             %
                             % Note how a `%' character causes TeX to ignore
                             % the end of the input line, so these blank lines
                             % do not start a new paragraph.
                             %
with the right spacing around the periods requires
a special command.

\LaTeX\ interprets some common characters as
commands, so you must type special commands to
generate them.  These characters include the
following:
       \$ \& \% \# \{ and \}.

In printing, text is usually emphasized with an
       \emph{italic}
type style.

\begin{em}
   A long segment of text can also be emphasized
   in this way.  Text within such a segment can be
   given \emph{additional} emphasis.
\end{em}

It is sometimes necessary to prevent \LaTeX\ from
breaking a line where it might otherwise do so.
This may be at a space, as between the ``Mr.''\ and
``Jones'' in
       ``Mr.~Jones'',        % ~ produces an unbreakable interword space.
or within a word---especially when the word is a
symbol like
       \mbox{\emph{itemnum}}
that makes little sense when hyphenated across
lines.

Footnotes\footnote{This is an example of a footnote.}
pose no problem.

\LaTeX\ is good at typesetting mathematical formulas
like
       \( x-3y + z = 7 \)
or
       \( a_{1} > x^{2n} + y^{2n} > x' \)
or
       \( \ip{A}{B} = \sum_{i} a_{i} b_{i} \).
The spaces you type in a formula are
ignored.  Remember that a letter like
       $x$                   % $ ... $  and  \( ... \)  are equivalent
is a formula when it denotes a mathematical
symbol, and it should be typed as one.

\section{Displayed Text}

Text is displayed by indenting it from the left
margin.  Quotations are commonly displayed.  There
are short quotations
\begin{quote}
   This is a short quotation.  It consists of a
   single paragraph of text.  See how it is formatted.
\end{quote}
and longer ones.
\begin{quotation}
   This is a longer quotation.  It consists of two
   paragraphs of text, neither of which are
   particularly interesting.

   This is the second paragraph of the quotation.  It
   is just as dull as the first paragraph.
\end{quotation}
Another frequently-displayed structure is a list.
The following is an example of an \emph{itemized}
list.
\begin{itemize}
   \item This is the first item of an itemized list.
         Each item in the list is marked with a ``tick''.
         You don't have to worry about what kind of tick
         mark is used.

   \item This is the second item of the list.  It
         contains another list nested inside it.  The inner
         list is an \emph{enumerated} list.
         \begin{enumerate}
            \item This is the first item of an enumerated
                  list that is nested within the itemized list.

            \item This is the second item of the inner list.
                  \LaTeX\ allows you to nest lists deeper than
                  you really should.
         \end{enumerate}
         This is the rest of the second item of the outer
         list.  It is no more interesting than any other
         part of the item.
   \item This is the third item of the list.
\end{itemize}
You can even display poetry.
\begin{verse}
   There is an environment
    for verse \\             % The \\ command separates lines
   Whose features some poets % within a stanza.
   will curse.

                             % One or more blank lines separate stanzas.

   For instead of making\\
   Them do \emph{all} line breaking, \\
   It allows them to put too many words on a line when they'd rather be
   forced to be terse.
\end{verse}

Mathematical formulas may also be displayed.  A
displayed formula
is
one-line long; multiline
formulas require special formatting instructions.
   \[  \ip{\Gamma}{\psi'} = x'' + y^{2} + z_{i}^{n}\]
Don't start a paragraph with a displayed equation,
nor make one a paragraph by itself.

\end{document}               % End of document.

\section{Bench Section 259}
% synthetic bench section 259
% This is a sample LaTeX input file.  (Version of 12 August 2004.)
%
% A '%' character causes TeX to ignore all remaining text on the line,
% and is used for comments like this one.

\documentclass{article}      % Specifies the document class

                             % The preamble begins here.
\title{An Example Document}  % Declares the document's title.
\author{Leslie Lamport}      % Declares the author's name.
\date{January 21, 1994}      % Deleting this command produces today's date.

\newcommand{\ip}[2]{(#1, #2)}
                             % Defines \ip{arg1}{arg2} to mean
                             % (arg1, arg2).

%\newcommand{\ip}[2]{\langle #1 | #2\rangle}
                             % This is an alternative definition of
                             % \ip that is commented out.

\begin{document}             % End of preamble and beginning of text.

\maketitle                   % Produces the title.

This is an example input file.  Comparing it with
the output it generates can show you how to
produce a simple document of your own.

\section{Ordinary Text}      % Produces section heading.  Lower-level
                             % sections are begun with similar
                             % \subsection and \subsubsection commands.

The ends  of words and sentences are marked
  by   spaces. It  doesn't matter how many
spaces    you type; one is as good as 100.  The
end of   a line counts as a space.

One   or more   blank lines denote the  end
of  a paragraph.

Since any number of consecutive spaces are treated
like a single one, the formatting of the input
file makes no difference to
      \LaTeX,                % The \LaTeX command generates the LaTeX logo.
but it makes a difference to you.  When you use
\LaTeX, making your input file as easy to read
as possible will be a great help as you write
your document and when you change it.  This sample
file shows how you can add comments to your own input
file.

Because printing is different from typewriting,
there are a number of things that you have to do
differently when preparing an input file than if
you were just typing the document directly.
Quotation marks like
       ``this''
have to be handled specially, as do quotes within
quotes:
       ``\,`this'            % \, separates the double and single quote.
        is what I just
        wrote, not  `that'\,''.

Dashes come in three sizes: an
       intra-word
dash, a medium dash for number ranges like
       1--2,
and a punctuation
       dash---like
this.

A sentence-ending space should be larger than the
space between words within a sentence.  You
sometimes have to type special commands in
conjunction with punctuation characters to get
this right, as in the following sentence.
       Gnats, gnus, etc.\ all  % `\ ' makes an inter-word space.
       begin with G\@.         % \@ marks end-of-sentence punctuation.
You should check the spaces after periods when
reading your output to make sure you haven't
forgotten any special cases.  Generating an
ellipsis
       \ldots\               % `\ ' is needed after `\ldots' because TeX
                             % ignores spaces after command names like \ldots
                             % made from \ + letters.
                             %
                             % Note how a `%' character causes TeX to ignore
                             % the end of the input line, so these blank lines
                             % do not start a new paragraph.
                             %
with the right spacing around the periods requires
a special command.

\LaTeX\ interprets some common characters as
commands, so you must type special commands to
generate them.  These characters include the
following:
       \$ \& \% \# \{ and \}.

In printing, text is usually emphasized with an
       \emph{italic}
type style.

\begin{em}
   A long segment of text can also be emphasized
   in this way.  Text within such a segment can be
   given \emph{additional} emphasis.
\end{em}

It is sometimes necessary to prevent \LaTeX\ from
breaking a line where it might otherwise do so.
This may be at a space, as between the ``Mr.''\ and
``Jones'' in
       ``Mr.~Jones'',        % ~ produces an unbreakable interword space.
or within a word---especially when the word is a
symbol like
       \mbox{\emph{itemnum}}
that makes little sense when hyphenated across
lines.

Footnotes\footnote{This is an example of a footnote.}
pose no problem.

\LaTeX\ is good at typesetting mathematical formulas
like
       \( x-3y + z = 7 \)
or
       \( a_{1} > x^{2n} + y^{2n} > x' \)
or
       \( \ip{A}{B} = \sum_{i} a_{i} b_{i} \).
The spaces you type in a formula are
ignored.  Remember that a letter like
       $x$                   % $ ... $  and  \( ... \)  are equivalent
is a formula when it denotes a mathematical
symbol, and it should be typed as one.

\section{Displayed Text}

Text is displayed by indenting it from the left
margin.  Quotations are commonly displayed.  There
are short quotations
\begin{quote}
   This is a short quotation.  It consists of a
   single paragraph of text.  See how it is formatted.
\end{quote}
and longer ones.
\begin{quotation}
   This is a longer quotation.  It consists of two
   paragraphs of text, neither of which are
   particularly interesting.

   This is the second paragraph of the quotation.  It
   is just as dull as the first paragraph.
\end{quotation}
Another frequently-displayed structure is a list.
The following is an example of an \emph{itemized}
list.
\begin{itemize}
   \item This is the first item of an itemized list.
         Each item in the list is marked with a ``tick''.
         You don't have to worry about what kind of tick
         mark is used.

   \item This is the second item of the list.  It
         contains another list nested inside it.  The inner
         list is an \emph{enumerated} list.
         \begin{enumerate}
            \item This is the first item of an enumerated
                  list that is nested within the itemized list.

            \item This is the second item of the inner list.
                  \LaTeX\ allows you to nest lists deeper than
                  you really should.
         \end{enumerate}
         This is the rest of the second item of the outer
         list.  It is no more interesting than any other
         part of the item.
   \item This is the third item of the list.
\end{itemize}
You can even display poetry.
\begin{verse}
   There is an environment
    for verse \\             % The \\ command separates lines
   Whose features some poets % within a stanza.
   will curse.

                             % One or more blank lines separate stanzas.

   For instead of making\\
   Them do \emph{all} line breaking, \\
   It allows them to put too many words on a line when they'd rather be
   forced to be terse.
\end{verse}

Mathematical formulas may also be displayed.  A
displayed formula
is
one-line long; multiline
formulas require special formatting instructions.
   \[  \ip{\Gamma}{\psi'} = x'' + y^{2} + z_{i}^{n}\]
Don't start a paragraph with a displayed equation,
nor make one a paragraph by itself.

\end{document}               % End of document.

\section{Bench Section 260}
% synthetic bench section 260
% This is a sample LaTeX input file.  (Version of 12 August 2004.)
%
% A '%' character causes TeX to ignore all remaining text on the line,
% and is used for comments like this one.

\documentclass{article}      % Specifies the document class

                             % The preamble begins here.
\title{An Example Document}  % Declares the document's title.
\author{Leslie Lamport}      % Declares the author's name.
\date{January 21, 1994}      % Deleting this command produces today's date.

\newcommand{\ip}[2]{(#1, #2)}
                             % Defines \ip{arg1}{arg2} to mean
                             % (arg1, arg2).

%\newcommand{\ip}[2]{\langle #1 | #2\rangle}
                             % This is an alternative definition of
                             % \ip that is commented out.

\begin{document}             % End of preamble and beginning of text.

\maketitle                   % Produces the title.

This is an example input file.  Comparing it with
the output it generates can show you how to
produce a simple document of your own.

\section{Ordinary Text}      % Produces section heading.  Lower-level
                             % sections are begun with similar
                             % \subsection and \subsubsection commands.

The ends  of words and sentences are marked
  by   spaces. It  doesn't matter how many
spaces    you type; one is as good as 100.  The
end of   a line counts as a space.

One   or more   blank lines denote the  end
of  a paragraph.

Since any number of consecutive spaces are treated
like a single one, the formatting of the input
file makes no difference to
      \LaTeX,                % The \LaTeX command generates the LaTeX logo.
but it makes a difference to you.  When you use
\LaTeX, making your input file as easy to read
as possible will be a great help as you write
your document and when you change it.  This sample
file shows how you can add comments to your own input
file.

Because printing is different from typewriting,
there are a number of things that you have to do
differently when preparing an input file than if
you were just typing the document directly.
Quotation marks like
       ``this''
have to be handled specially, as do quotes within
quotes:
       ``\,`this'            % \, separates the double and single quote.
        is what I just
        wrote, not  `that'\,''.

Dashes come in three sizes: an
       intra-word
dash, a medium dash for number ranges like
       1--2,
and a punctuation
       dash---like
this.

A sentence-ending space should be larger than the
space between words within a sentence.  You
sometimes have to type special commands in
conjunction with punctuation characters to get
this right, as in the following sentence.
       Gnats, gnus, etc.\ all  % `\ ' makes an inter-word space.
       begin with G\@.         % \@ marks end-of-sentence punctuation.
You should check the spaces after periods when
reading your output to make sure you haven't
forgotten any special cases.  Generating an
ellipsis
       \ldots\               % `\ ' is needed after `\ldots' because TeX
                             % ignores spaces after command names like \ldots
                             % made from \ + letters.
                             %
                             % Note how a `%' character causes TeX to ignore
                             % the end of the input line, so these blank lines
                             % do not start a new paragraph.
                             %
with the right spacing around the periods requires
a special command.

\LaTeX\ interprets some common characters as
commands, so you must type special commands to
generate them.  These characters include the
following:
       \$ \& \% \# \{ and \}.

In printing, text is usually emphasized with an
       \emph{italic}
type style.

\begin{em}
   A long segment of text can also be emphasized
   in this way.  Text within such a segment can be
   given \emph{additional} emphasis.
\end{em}

It is sometimes necessary to prevent \LaTeX\ from
breaking a line where it might otherwise do so.
This may be at a space, as between the ``Mr.''\ and
``Jones'' in
       ``Mr.~Jones'',        % ~ produces an unbreakable interword space.
or within a word---especially when the word is a
symbol like
       \mbox{\emph{itemnum}}
that makes little sense when hyphenated across
lines.

Footnotes\footnote{This is an example of a footnote.}
pose no problem.

\LaTeX\ is good at typesetting mathematical formulas
like
       \( x-3y + z = 7 \)
or
       \( a_{1} > x^{2n} + y^{2n} > x' \)
or
       \( \ip{A}{B} = \sum_{i} a_{i} b_{i} \).
The spaces you type in a formula are
ignored.  Remember that a letter like
       $x$                   % $ ... $  and  \( ... \)  are equivalent
is a formula when it denotes a mathematical
symbol, and it should be typed as one.

\section{Displayed Text}

Text is displayed by indenting it from the left
margin.  Quotations are commonly displayed.  There
are short quotations
\begin{quote}
   This is a short quotation.  It consists of a
   single paragraph of text.  See how it is formatted.
\end{quote}
and longer ones.
\begin{quotation}
   This is a longer quotation.  It consists of two
   paragraphs of text, neither of which are
   particularly interesting.

   This is the second paragraph of the quotation.  It
   is just as dull as the first paragraph.
\end{quotation}
Another frequently-displayed structure is a list.
The following is an example of an \emph{itemized}
list.
\begin{itemize}
   \item This is the first item of an itemized list.
         Each item in the list is marked with a ``tick''.
         You don't have to worry about what kind of tick
         mark is used.

   \item This is the second item of the list.  It
         contains another list nested inside it.  The inner
         list is an \emph{enumerated} list.
         \begin{enumerate}
            \item This is the first item of an enumerated
                  list that is nested within the itemized list.

            \item This is the second item of the inner list.
                  \LaTeX\ allows you to nest lists deeper than
                  you really should.
         \end{enumerate}
         This is the rest of the second item of the outer
         list.  It is no more interesting than any other
         part of the item.
   \item This is the third item of the list.
\end{itemize}
You can even display poetry.
\begin{verse}
   There is an environment
    for verse \\             % The \\ command separates lines
   Whose features some poets % within a stanza.
   will curse.

                             % One or more blank lines separate stanzas.

   For instead of making\\
   Them do \emph{all} line breaking, \\
   It allows them to put too many words on a line when they'd rather be
   forced to be terse.
\end{verse}

Mathematical formulas may also be displayed.  A
displayed formula
is
one-line long; multiline
formulas require special formatting instructions.
   \[  \ip{\Gamma}{\psi'} = x'' + y^{2} + z_{i}^{n}\]
Don't start a paragraph with a displayed equation,
nor make one a paragraph by itself.

\end{document}               % End of document.

\section{Bench Section 261}
% synthetic bench section 261
% This is a sample LaTeX input file.  (Version of 12 August 2004.)
%
% A '%' character causes TeX to ignore all remaining text on the line,
% and is used for comments like this one.

\documentclass{article}      % Specifies the document class

                             % The preamble begins here.
\title{An Example Document}  % Declares the document's title.
\author{Leslie Lamport}      % Declares the author's name.
\date{January 21, 1994}      % Deleting this command produces today's date.

\newcommand{\ip}[2]{(#1, #2)}
                             % Defines \ip{arg1}{arg2} to mean
                             % (arg1, arg2).

%\newcommand{\ip}[2]{\langle #1 | #2\rangle}
                             % This is an alternative definition of
                             % \ip that is commented out.

\begin{document}             % End of preamble and beginning of text.

\maketitle                   % Produces the title.

This is an example input file.  Comparing it with
the output it generates can show you how to
produce a simple document of your own.

\section{Ordinary Text}      % Produces section heading.  Lower-level
                             % sections are begun with similar
                             % \subsection and \subsubsection commands.

The ends  of words and sentences are marked
  by   spaces. It  doesn't matter how many
spaces    you type; one is as good as 100.  The
end of   a line counts as a space.

One   or more   blank lines denote the  end
of  a paragraph.

Since any number of consecutive spaces are treated
like a single one, the formatting of the input
file makes no difference to
      \LaTeX,                % The \LaTeX command generates the LaTeX logo.
but it makes a difference to you.  When you use
\LaTeX, making your input file as easy to read
as possible will be a great help as you write
your document and when you change it.  This sample
file shows how you can add comments to your own input
file.

Because printing is different from typewriting,
there are a number of things that you have to do
differently when preparing an input file than if
you were just typing the document directly.
Quotation marks like
       ``this''
have to be handled specially, as do quotes within
quotes:
       ``\,`this'            % \, separates the double and single quote.
        is what I just
        wrote, not  `that'\,''.

Dashes come in three sizes: an
       intra-word
dash, a medium dash for number ranges like
       1--2,
and a punctuation
       dash---like
this.

A sentence-ending space should be larger than the
space between words within a sentence.  You
sometimes have to type special commands in
conjunction with punctuation characters to get
this right, as in the following sentence.
       Gnats, gnus, etc.\ all  % `\ ' makes an inter-word space.
       begin with G\@.         % \@ marks end-of-sentence punctuation.
You should check the spaces after periods when
reading your output to make sure you haven't
forgotten any special cases.  Generating an
ellipsis
       \ldots\               % `\ ' is needed after `\ldots' because TeX
                             % ignores spaces after command names like \ldots
                             % made from \ + letters.
                             %
                             % Note how a `%' character causes TeX to ignore
                             % the end of the input line, so these blank lines
                             % do not start a new paragraph.
                             %
with the right spacing around the periods requires
a special command.

\LaTeX\ interprets some common characters as
commands, so you must type special commands to
generate them.  These characters include the
following:
       \$ \& \% \# \{ and \}.

In printing, text is usually emphasized with an
       \emph{italic}
type style.

\begin{em}
   A long segment of text can also be emphasized
   in this way.  Text within such a segment can be
   given \emph{additional} emphasis.
\end{em}

It is sometimes necessary to prevent \LaTeX\ from
breaking a line where it might otherwise do so.
This may be at a space, as between the ``Mr.''\ and
``Jones'' in
       ``Mr.~Jones'',        % ~ produces an unbreakable interword space.
or within a word---especially when the word is a
symbol like
       \mbox{\emph{itemnum}}
that makes little sense when hyphenated across
lines.

Footnotes\footnote{This is an example of a footnote.}
pose no problem.

\LaTeX\ is good at typesetting mathematical formulas
like
       \( x-3y + z = 7 \)
or
       \( a_{1} > x^{2n} + y^{2n} > x' \)
or
       \( \ip{A}{B} = \sum_{i} a_{i} b_{i} \).
The spaces you type in a formula are
ignored.  Remember that a letter like
       $x$                   % $ ... $  and  \( ... \)  are equivalent
is a formula when it denotes a mathematical
symbol, and it should be typed as one.

\section{Displayed Text}

Text is displayed by indenting it from the left
margin.  Quotations are commonly displayed.  There
are short quotations
\begin{quote}
   This is a short quotation.  It consists of a
   single paragraph of text.  See how it is formatted.
\end{quote}
and longer ones.
\begin{quotation}
   This is a longer quotation.  It consists of two
   paragraphs of text, neither of which are
   particularly interesting.

   This is the second paragraph of the quotation.  It
   is just as dull as the first paragraph.
\end{quotation}
Another frequently-displayed structure is a list.
The following is an example of an \emph{itemized}
list.
\begin{itemize}
   \item This is the first item of an itemized list.
         Each item in the list is marked with a ``tick''.
         You don't have to worry about what kind of tick
         mark is used.

   \item This is the second item of the list.  It
         contains another list nested inside it.  The inner
         list is an \emph{enumerated} list.
         \begin{enumerate}
            \item This is the first item of an enumerated
                  list that is nested within the itemized list.

            \item This is the second item of the inner list.
                  \LaTeX\ allows you to nest lists deeper than
                  you really should.
         \end{enumerate}
         This is the rest of the second item of the outer
         list.  It is no more interesting than any other
         part of the item.
   \item This is the third item of the list.
\end{itemize}
You can even display poetry.
\begin{verse}
   There is an environment
    for verse \\             % The \\ command separates lines
   Whose features some poets % within a stanza.
   will curse.

                             % One or more blank lines separate stanzas.

   For instead of making\\
   Them do \emph{all} line breaking, \\
   It allows them to put too many words on a line when they'd rather be
   forced to be terse.
\end{verse}

Mathematical formulas may also be displayed.  A
displayed formula
is
one-line long; multiline
formulas require special formatting instructions.
   \[  \ip{\Gamma}{\psi'} = x'' + y^{2} + z_{i}^{n}\]
Don't start a paragraph with a displayed equation,
nor make one a paragraph by itself.

\end{document}               % End of document.

\section{Bench Section 262}
% synthetic bench section 262
% This is a sample LaTeX input file.  (Version of 12 August 2004.)
%
% A '%' character causes TeX to ignore all remaining text on the line,
% and is used for comments like this one.

\documentclass{article}      % Specifies the document class

                             % The preamble begins here.
\title{An Example Document}  % Declares the document's title.
\author{Leslie Lamport}      % Declares the author's name.
\date{January 21, 1994}      % Deleting this command produces today's date.

\newcommand{\ip}[2]{(#1, #2)}
                             % Defines \ip{arg1}{arg2} to mean
                             % (arg1, arg2).

%\newcommand{\ip}[2]{\langle #1 | #2\rangle}
                             % This is an alternative definition of
                             % \ip that is commented out.

\begin{document}             % End of preamble and beginning of text.

\maketitle                   % Produces the title.

This is an example input file.  Comparing it with
the output it generates can show you how to
produce a simple document of your own.

\section{Ordinary Text}      % Produces section heading.  Lower-level
                             % sections are begun with similar
                             % \subsection and \subsubsection commands.

The ends  of words and sentences are marked
  by   spaces. It  doesn't matter how many
spaces    you type; one is as good as 100.  The
end of   a line counts as a space.

One   or more   blank lines denote the  end
of  a paragraph.

Since any number of consecutive spaces are treated
like a single one, the formatting of the input
file makes no difference to
      \LaTeX,                % The \LaTeX command generates the LaTeX logo.
but it makes a difference to you.  When you use
\LaTeX, making your input file as easy to read
as possible will be a great help as you write
your document and when you change it.  This sample
file shows how you can add comments to your own input
file.

Because printing is different from typewriting,
there are a number of things that you have to do
differently when preparing an input file than if
you were just typing the document directly.
Quotation marks like
       ``this''
have to be handled specially, as do quotes within
quotes:
       ``\,`this'            % \, separates the double and single quote.
        is what I just
        wrote, not  `that'\,''.

Dashes come in three sizes: an
       intra-word
dash, a medium dash for number ranges like
       1--2,
and a punctuation
       dash---like
this.

A sentence-ending space should be larger than the
space between words within a sentence.  You
sometimes have to type special commands in
conjunction with punctuation characters to get
this right, as in the following sentence.
       Gnats, gnus, etc.\ all  % `\ ' makes an inter-word space.
       begin with G\@.         % \@ marks end-of-sentence punctuation.
You should check the spaces after periods when
reading your output to make sure you haven't
forgotten any special cases.  Generating an
ellipsis
       \ldots\               % `\ ' is needed after `\ldots' because TeX
                             % ignores spaces after command names like \ldots
                             % made from \ + letters.
                             %
                             % Note how a `%' character causes TeX to ignore
                             % the end of the input line, so these blank lines
                             % do not start a new paragraph.
                             %
with the right spacing around the periods requires
a special command.

\LaTeX\ interprets some common characters as
commands, so you must type special commands to
generate them.  These characters include the
following:
       \$ \& \% \# \{ and \}.

In printing, text is usually emphasized with an
       \emph{italic}
type style.

\begin{em}
   A long segment of text can also be emphasized
   in this way.  Text within such a segment can be
   given \emph{additional} emphasis.
\end{em}

It is sometimes necessary to prevent \LaTeX\ from
breaking a line where it might otherwise do so.
This may be at a space, as between the ``Mr.''\ and
``Jones'' in
       ``Mr.~Jones'',        % ~ produces an unbreakable interword space.
or within a word---especially when the word is a
symbol like
       \mbox{\emph{itemnum}}
that makes little sense when hyphenated across
lines.

Footnotes\footnote{This is an example of a footnote.}
pose no problem.

\LaTeX\ is good at typesetting mathematical formulas
like
       \( x-3y + z = 7 \)
or
       \( a_{1} > x^{2n} + y^{2n} > x' \)
or
       \( \ip{A}{B} = \sum_{i} a_{i} b_{i} \).
The spaces you type in a formula are
ignored.  Remember that a letter like
       $x$                   % $ ... $  and  \( ... \)  are equivalent
is a formula when it denotes a mathematical
symbol, and it should be typed as one.

\section{Displayed Text}

Text is displayed by indenting it from the left
margin.  Quotations are commonly displayed.  There
are short quotations
\begin{quote}
   This is a short quotation.  It consists of a
   single paragraph of text.  See how it is formatted.
\end{quote}
and longer ones.
\begin{quotation}
   This is a longer quotation.  It consists of two
   paragraphs of text, neither of which are
   particularly interesting.

   This is the second paragraph of the quotation.  It
   is just as dull as the first paragraph.
\end{quotation}
Another frequently-displayed structure is a list.
The following is an example of an \emph{itemized}
list.
\begin{itemize}
   \item This is the first item of an itemized list.
         Each item in the list is marked with a ``tick''.
         You don't have to worry about what kind of tick
         mark is used.

   \item This is the second item of the list.  It
         contains another list nested inside it.  The inner
         list is an \emph{enumerated} list.
         \begin{enumerate}
            \item This is the first item of an enumerated
                  list that is nested within the itemized list.

            \item This is the second item of the inner list.
                  \LaTeX\ allows you to nest lists deeper than
                  you really should.
         \end{enumerate}
         This is the rest of the second item of the outer
         list.  It is no more interesting than any other
         part of the item.
   \item This is the third item of the list.
\end{itemize}
You can even display poetry.
\begin{verse}
   There is an environment
    for verse \\             % The \\ command separates lines
   Whose features some poets % within a stanza.
   will curse.

                             % One or more blank lines separate stanzas.

   For instead of making\\
   Them do \emph{all} line breaking, \\
   It allows them to put too many words on a line when they'd rather be
   forced to be terse.
\end{verse}

Mathematical formulas may also be displayed.  A
displayed formula
is
one-line long; multiline
formulas require special formatting instructions.
   \[  \ip{\Gamma}{\psi'} = x'' + y^{2} + z_{i}^{n}\]
Don't start a paragraph with a displayed equation,
nor make one a paragraph by itself.

\end{document}               % End of document.

\section{Bench Section 263}
% synthetic bench section 263
% This is a sample LaTeX input file.  (Version of 12 August 2004.)
%
% A '%' character causes TeX to ignore all remaining text on the line,
% and is used for comments like this one.

\documentclass{article}      % Specifies the document class

                             % The preamble begins here.
\title{An Example Document}  % Declares the document's title.
\author{Leslie Lamport}      % Declares the author's name.
\date{January 21, 1994}      % Deleting this command produces today's date.

\newcommand{\ip}[2]{(#1, #2)}
                             % Defines \ip{arg1}{arg2} to mean
                             % (arg1, arg2).

%\newcommand{\ip}[2]{\langle #1 | #2\rangle}
                             % This is an alternative definition of
                             % \ip that is commented out.

\begin{document}             % End of preamble and beginning of text.

\maketitle                   % Produces the title.

This is an example input file.  Comparing it with
the output it generates can show you how to
produce a simple document of your own.

\section{Ordinary Text}      % Produces section heading.  Lower-level
                             % sections are begun with similar
                             % \subsection and \subsubsection commands.

The ends  of words and sentences are marked
  by   spaces. It  doesn't matter how many
spaces    you type; one is as good as 100.  The
end of   a line counts as a space.

One   or more   blank lines denote the  end
of  a paragraph.

Since any number of consecutive spaces are treated
like a single one, the formatting of the input
file makes no difference to
      \LaTeX,                % The \LaTeX command generates the LaTeX logo.
but it makes a difference to you.  When you use
\LaTeX, making your input file as easy to read
as possible will be a great help as you write
your document and when you change it.  This sample
file shows how you can add comments to your own input
file.

Because printing is different from typewriting,
there are a number of things that you have to do
differently when preparing an input file than if
you were just typing the document directly.
Quotation marks like
       ``this''
have to be handled specially, as do quotes within
quotes:
       ``\,`this'            % \, separates the double and single quote.
        is what I just
        wrote, not  `that'\,''.

Dashes come in three sizes: an
       intra-word
dash, a medium dash for number ranges like
       1--2,
and a punctuation
       dash---like
this.

A sentence-ending space should be larger than the
space between words within a sentence.  You
sometimes have to type special commands in
conjunction with punctuation characters to get
this right, as in the following sentence.
       Gnats, gnus, etc.\ all  % `\ ' makes an inter-word space.
       begin with G\@.         % \@ marks end-of-sentence punctuation.
You should check the spaces after periods when
reading your output to make sure you haven't
forgotten any special cases.  Generating an
ellipsis
       \ldots\               % `\ ' is needed after `\ldots' because TeX
                             % ignores spaces after command names like \ldots
                             % made from \ + letters.
                             %
                             % Note how a `%' character causes TeX to ignore
                             % the end of the input line, so these blank lines
                             % do not start a new paragraph.
                             %
with the right spacing around the periods requires
a special command.

\LaTeX\ interprets some common characters as
commands, so you must type special commands to
generate them.  These characters include the
following:
       \$ \& \% \# \{ and \}.

In printing, text is usually emphasized with an
       \emph{italic}
type style.

\begin{em}
   A long segment of text can also be emphasized
   in this way.  Text within such a segment can be
   given \emph{additional} emphasis.
\end{em}

It is sometimes necessary to prevent \LaTeX\ from
breaking a line where it might otherwise do so.
This may be at a space, as between the ``Mr.''\ and
``Jones'' in
       ``Mr.~Jones'',        % ~ produces an unbreakable interword space.
or within a word---especially when the word is a
symbol like
       \mbox{\emph{itemnum}}
that makes little sense when hyphenated across
lines.

Footnotes\footnote{This is an example of a footnote.}
pose no problem.

\LaTeX\ is good at typesetting mathematical formulas
like
       \( x-3y + z = 7 \)
or
       \( a_{1} > x^{2n} + y^{2n} > x' \)
or
       \( \ip{A}{B} = \sum_{i} a_{i} b_{i} \).
The spaces you type in a formula are
ignored.  Remember that a letter like
       $x$                   % $ ... $  and  \( ... \)  are equivalent
is a formula when it denotes a mathematical
symbol, and it should be typed as one.

\section{Displayed Text}

Text is displayed by indenting it from the left
margin.  Quotations are commonly displayed.  There
are short quotations
\begin{quote}
   This is a short quotation.  It consists of a
   single paragraph of text.  See how it is formatted.
\end{quote}
and longer ones.
\begin{quotation}
   This is a longer quotation.  It consists of two
   paragraphs of text, neither of which are
   particularly interesting.

   This is the second paragraph of the quotation.  It
   is just as dull as the first paragraph.
\end{quotation}
Another frequently-displayed structure is a list.
The following is an example of an \emph{itemized}
list.
\begin{itemize}
   \item This is the first item of an itemized list.
         Each item in the list is marked with a ``tick''.
         You don't have to worry about what kind of tick
         mark is used.

   \item This is the second item of the list.  It
         contains another list nested inside it.  The inner
         list is an \emph{enumerated} list.
         \begin{enumerate}
            \item This is the first item of an enumerated
                  list that is nested within the itemized list.

            \item This is the second item of the inner list.
                  \LaTeX\ allows you to nest lists deeper than
                  you really should.
         \end{enumerate}
         This is the rest of the second item of the outer
         list.  It is no more interesting than any other
         part of the item.
   \item This is the third item of the list.
\end{itemize}
You can even display poetry.
\begin{verse}
   There is an environment
    for verse \\             % The \\ command separates lines
   Whose features some poets % within a stanza.
   will curse.

                             % One or more blank lines separate stanzas.

   For instead of making\\
   Them do \emph{all} line breaking, \\
   It allows them to put too many words on a line when they'd rather be
   forced to be terse.
\end{verse}

Mathematical formulas may also be displayed.  A
displayed formula
is
one-line long; multiline
formulas require special formatting instructions.
   \[  \ip{\Gamma}{\psi'} = x'' + y^{2} + z_{i}^{n}\]
Don't start a paragraph with a displayed equation,
nor make one a paragraph by itself.

\end{document}               % End of document.

\section{Bench Section 264}
% synthetic bench section 264
% This is a sample LaTeX input file.  (Version of 12 August 2004.)
%
% A '%' character causes TeX to ignore all remaining text on the line,
% and is used for comments like this one.

\documentclass{article}      % Specifies the document class

                             % The preamble begins here.
\title{An Example Document}  % Declares the document's title.
\author{Leslie Lamport}      % Declares the author's name.
\date{January 21, 1994}      % Deleting this command produces today's date.

\newcommand{\ip}[2]{(#1, #2)}
                             % Defines \ip{arg1}{arg2} to mean
                             % (arg1, arg2).

%\newcommand{\ip}[2]{\langle #1 | #2\rangle}
                             % This is an alternative definition of
                             % \ip that is commented out.

\begin{document}             % End of preamble and beginning of text.

\maketitle                   % Produces the title.

This is an example input file.  Comparing it with
the output it generates can show you how to
produce a simple document of your own.

\section{Ordinary Text}      % Produces section heading.  Lower-level
                             % sections are begun with similar
                             % \subsection and \subsubsection commands.

The ends  of words and sentences are marked
  by   spaces. It  doesn't matter how many
spaces    you type; one is as good as 100.  The
end of   a line counts as a space.

One   or more   blank lines denote the  end
of  a paragraph.

Since any number of consecutive spaces are treated
like a single one, the formatting of the input
file makes no difference to
      \LaTeX,                % The \LaTeX command generates the LaTeX logo.
but it makes a difference to you.  When you use
\LaTeX, making your input file as easy to read
as possible will be a great help as you write
your document and when you change it.  This sample
file shows how you can add comments to your own input
file.

Because printing is different from typewriting,
there are a number of things that you have to do
differently when preparing an input file than if
you were just typing the document directly.
Quotation marks like
       ``this''
have to be handled specially, as do quotes within
quotes:
       ``\,`this'            % \, separates the double and single quote.
        is what I just
        wrote, not  `that'\,''.

Dashes come in three sizes: an
       intra-word
dash, a medium dash for number ranges like
       1--2,
and a punctuation
       dash---like
this.

A sentence-ending space should be larger than the
space between words within a sentence.  You
sometimes have to type special commands in
conjunction with punctuation characters to get
this right, as in the following sentence.
       Gnats, gnus, etc.\ all  % `\ ' makes an inter-word space.
       begin with G\@.         % \@ marks end-of-sentence punctuation.
You should check the spaces after periods when
reading your output to make sure you haven't
forgotten any special cases.  Generating an
ellipsis
       \ldots\               % `\ ' is needed after `\ldots' because TeX
                             % ignores spaces after command names like \ldots
                             % made from \ + letters.
                             %
                             % Note how a `%' character causes TeX to ignore
                             % the end of the input line, so these blank lines
                             % do not start a new paragraph.
                             %
with the right spacing around the periods requires
a special command.

\LaTeX\ interprets some common characters as
commands, so you must type special commands to
generate them.  These characters include the
following:
       \$ \& \% \# \{ and \}.

In printing, text is usually emphasized with an
       \emph{italic}
type style.

\begin{em}
   A long segment of text can also be emphasized
   in this way.  Text within such a segment can be
   given \emph{additional} emphasis.
\end{em}

It is sometimes necessary to prevent \LaTeX\ from
breaking a line where it might otherwise do so.
This may be at a space, as between the ``Mr.''\ and
``Jones'' in
       ``Mr.~Jones'',        % ~ produces an unbreakable interword space.
or within a word---especially when the word is a
symbol like
       \mbox{\emph{itemnum}}
that makes little sense when hyphenated across
lines.

Footnotes\footnote{This is an example of a footnote.}
pose no problem.

\LaTeX\ is good at typesetting mathematical formulas
like
       \( x-3y + z = 7 \)
or
       \( a_{1} > x^{2n} + y^{2n} > x' \)
or
       \( \ip{A}{B} = \sum_{i} a_{i} b_{i} \).
The spaces you type in a formula are
ignored.  Remember that a letter like
       $x$                   % $ ... $  and  \( ... \)  are equivalent
is a formula when it denotes a mathematical
symbol, and it should be typed as one.

\section{Displayed Text}

Text is displayed by indenting it from the left
margin.  Quotations are commonly displayed.  There
are short quotations
\begin{quote}
   This is a short quotation.  It consists of a
   single paragraph of text.  See how it is formatted.
\end{quote}
and longer ones.
\begin{quotation}
   This is a longer quotation.  It consists of two
   paragraphs of text, neither of which are
   particularly interesting.

   This is the second paragraph of the quotation.  It
   is just as dull as the first paragraph.
\end{quotation}
Another frequently-displayed structure is a list.
The following is an example of an \emph{itemized}
list.
\begin{itemize}
   \item This is the first item of an itemized list.
         Each item in the list is marked with a ``tick''.
         You don't have to worry about what kind of tick
         mark is used.

   \item This is the second item of the list.  It
         contains another list nested inside it.  The inner
         list is an \emph{enumerated} list.
         \begin{enumerate}
            \item This is the first item of an enumerated
                  list that is nested within the itemized list.

            \item This is the second item of the inner list.
                  \LaTeX\ allows you to nest lists deeper than
                  you really should.
         \end{enumerate}
         This is the rest of the second item of the outer
         list.  It is no more interesting than any other
         part of the item.
   \item This is the third item of the list.
\end{itemize}
You can even display poetry.
\begin{verse}
   There is an environment
    for verse \\             % The \\ command separates lines
   Whose features some poets % within a stanza.
   will curse.

                             % One or more blank lines separate stanzas.

   For instead of making\\
   Them do \emph{all} line breaking, \\
   It allows them to put too many words on a line when they'd rather be
   forced to be terse.
\end{verse}

Mathematical formulas may also be displayed.  A
displayed formula
is
one-line long; multiline
formulas require special formatting instructions.
   \[  \ip{\Gamma}{\psi'} = x'' + y^{2} + z_{i}^{n}\]
Don't start a paragraph with a displayed equation,
nor make one a paragraph by itself.

\end{document}               % End of document.

\section{Bench Section 265}
% synthetic bench section 265
% This is a sample LaTeX input file.  (Version of 12 August 2004.)
%
% A '%' character causes TeX to ignore all remaining text on the line,
% and is used for comments like this one.

\documentclass{article}      % Specifies the document class

                             % The preamble begins here.
\title{An Example Document}  % Declares the document's title.
\author{Leslie Lamport}      % Declares the author's name.
\date{January 21, 1994}      % Deleting this command produces today's date.

\newcommand{\ip}[2]{(#1, #2)}
                             % Defines \ip{arg1}{arg2} to mean
                             % (arg1, arg2).

%\newcommand{\ip}[2]{\langle #1 | #2\rangle}
                             % This is an alternative definition of
                             % \ip that is commented out.

\begin{document}             % End of preamble and beginning of text.

\maketitle                   % Produces the title.

This is an example input file.  Comparing it with
the output it generates can show you how to
produce a simple document of your own.

\section{Ordinary Text}      % Produces section heading.  Lower-level
                             % sections are begun with similar
                             % \subsection and \subsubsection commands.

The ends  of words and sentences are marked
  by   spaces. It  doesn't matter how many
spaces    you type; one is as good as 100.  The
end of   a line counts as a space.

One   or more   blank lines denote the  end
of  a paragraph.

Since any number of consecutive spaces are treated
like a single one, the formatting of the input
file makes no difference to
      \LaTeX,                % The \LaTeX command generates the LaTeX logo.
but it makes a difference to you.  When you use
\LaTeX, making your input file as easy to read
as possible will be a great help as you write
your document and when you change it.  This sample
file shows how you can add comments to your own input
file.

Because printing is different from typewriting,
there are a number of things that you have to do
differently when preparing an input file than if
you were just typing the document directly.
Quotation marks like
       ``this''
have to be handled specially, as do quotes within
quotes:
       ``\,`this'            % \, separates the double and single quote.
        is what I just
        wrote, not  `that'\,''.

Dashes come in three sizes: an
       intra-word
dash, a medium dash for number ranges like
       1--2,
and a punctuation
       dash---like
this.

A sentence-ending space should be larger than the
space between words within a sentence.  You
sometimes have to type special commands in
conjunction with punctuation characters to get
this right, as in the following sentence.
       Gnats, gnus, etc.\ all  % `\ ' makes an inter-word space.
       begin with G\@.         % \@ marks end-of-sentence punctuation.
You should check the spaces after periods when
reading your output to make sure you haven't
forgotten any special cases.  Generating an
ellipsis
       \ldots\               % `\ ' is needed after `\ldots' because TeX
                             % ignores spaces after command names like \ldots
                             % made from \ + letters.
                             %
                             % Note how a `%' character causes TeX to ignore
                             % the end of the input line, so these blank lines
                             % do not start a new paragraph.
                             %
with the right spacing around the periods requires
a special command.

\LaTeX\ interprets some common characters as
commands, so you must type special commands to
generate them.  These characters include the
following:
       \$ \& \% \# \{ and \}.

In printing, text is usually emphasized with an
       \emph{italic}
type style.

\begin{em}
   A long segment of text can also be emphasized
   in this way.  Text within such a segment can be
   given \emph{additional} emphasis.
\end{em}

It is sometimes necessary to prevent \LaTeX\ from
breaking a line where it might otherwise do so.
This may be at a space, as between the ``Mr.''\ and
``Jones'' in
       ``Mr.~Jones'',        % ~ produces an unbreakable interword space.
or within a word---especially when the word is a
symbol like
       \mbox{\emph{itemnum}}
that makes little sense when hyphenated across
lines.

Footnotes\footnote{This is an example of a footnote.}
pose no problem.

\LaTeX\ is good at typesetting mathematical formulas
like
       \( x-3y + z = 7 \)
or
       \( a_{1} > x^{2n} + y^{2n} > x' \)
or
       \( \ip{A}{B} = \sum_{i} a_{i} b_{i} \).
The spaces you type in a formula are
ignored.  Remember that a letter like
       $x$                   % $ ... $  and  \( ... \)  are equivalent
is a formula when it denotes a mathematical
symbol, and it should be typed as one.

\section{Displayed Text}

Text is displayed by indenting it from the left
margin.  Quotations are commonly displayed.  There
are short quotations
\begin{quote}
   This is a short quotation.  It consists of a
   single paragraph of text.  See how it is formatted.
\end{quote}
and longer ones.
\begin{quotation}
   This is a longer quotation.  It consists of two
   paragraphs of text, neither of which are
   particularly interesting.

   This is the second paragraph of the quotation.  It
   is just as dull as the first paragraph.
\end{quotation}
Another frequently-displayed structure is a list.
The following is an example of an \emph{itemized}
list.
\begin{itemize}
   \item This is the first item of an itemized list.
         Each item in the list is marked with a ``tick''.
         You don't have to worry about what kind of tick
         mark is used.

   \item This is the second item of the list.  It
         contains another list nested inside it.  The inner
         list is an \emph{enumerated} list.
         \begin{enumerate}
            \item This is the first item of an enumerated
                  list that is nested within the itemized list.

            \item This is the second item of the inner list.
                  \LaTeX\ allows you to nest lists deeper than
                  you really should.
         \end{enumerate}
         This is the rest of the second item of the outer
         list.  It is no more interesting than any other
         part of the item.
   \item This is the third item of the list.
\end{itemize}
You can even display poetry.
\begin{verse}
   There is an environment
    for verse \\             % The \\ command separates lines
   Whose features some poets % within a stanza.
   will curse.

                             % One or more blank lines separate stanzas.

   For instead of making\\
   Them do \emph{all} line breaking, \\
   It allows them to put too many words on a line when they'd rather be
   forced to be terse.
\end{verse}

Mathematical formulas may also be displayed.  A
displayed formula
is
one-line long; multiline
formulas require special formatting instructions.
   \[  \ip{\Gamma}{\psi'} = x'' + y^{2} + z_{i}^{n}\]
Don't start a paragraph with a displayed equation,
nor make one a paragraph by itself.

\end{document}               % End of document.

\section{Bench Section 266}
% synthetic bench section 266
% This is a sample LaTeX input file.  (Version of 12 August 2004.)
%
% A '%' character causes TeX to ignore all remaining text on the line,
% and is used for comments like this one.

\documentclass{article}      % Specifies the document class

                             % The preamble begins here.
\title{An Example Document}  % Declares the document's title.
\author{Leslie Lamport}      % Declares the author's name.
\date{January 21, 1994}      % Deleting this command produces today's date.

\newcommand{\ip}[2]{(#1, #2)}
                             % Defines \ip{arg1}{arg2} to mean
                             % (arg1, arg2).

%\newcommand{\ip}[2]{\langle #1 | #2\rangle}
                             % This is an alternative definition of
                             % \ip that is commented out.

\begin{document}             % End of preamble and beginning of text.

\maketitle                   % Produces the title.

This is an example input file.  Comparing it with
the output it generates can show you how to
produce a simple document of your own.

\section{Ordinary Text}      % Produces section heading.  Lower-level
                             % sections are begun with similar
                             % \subsection and \subsubsection commands.

The ends  of words and sentences are marked
  by   spaces. It  doesn't matter how many
spaces    you type; one is as good as 100.  The
end of   a line counts as a space.

One   or more   blank lines denote the  end
of  a paragraph.

Since any number of consecutive spaces are treated
like a single one, the formatting of the input
file makes no difference to
      \LaTeX,                % The \LaTeX command generates the LaTeX logo.
but it makes a difference to you.  When you use
\LaTeX, making your input file as easy to read
as possible will be a great help as you write
your document and when you change it.  This sample
file shows how you can add comments to your own input
file.

Because printing is different from typewriting,
there are a number of things that you have to do
differently when preparing an input file than if
you were just typing the document directly.
Quotation marks like
       ``this''
have to be handled specially, as do quotes within
quotes:
       ``\,`this'            % \, separates the double and single quote.
        is what I just
        wrote, not  `that'\,''.

Dashes come in three sizes: an
       intra-word
dash, a medium dash for number ranges like
       1--2,
and a punctuation
       dash---like
this.

A sentence-ending space should be larger than the
space between words within a sentence.  You
sometimes have to type special commands in
conjunction with punctuation characters to get
this right, as in the following sentence.
       Gnats, gnus, etc.\ all  % `\ ' makes an inter-word space.
       begin with G\@.         % \@ marks end-of-sentence punctuation.
You should check the spaces after periods when
reading your output to make sure you haven't
forgotten any special cases.  Generating an
ellipsis
       \ldots\               % `\ ' is needed after `\ldots' because TeX
                             % ignores spaces after command names like \ldots
                             % made from \ + letters.
                             %
                             % Note how a `%' character causes TeX to ignore
                             % the end of the input line, so these blank lines
                             % do not start a new paragraph.
                             %
with the right spacing around the periods requires
a special command.

\LaTeX\ interprets some common characters as
commands, so you must type special commands to
generate them.  These characters include the
following:
       \$ \& \% \# \{ and \}.

In printing, text is usually emphasized with an
       \emph{italic}
type style.

\begin{em}
   A long segment of text can also be emphasized
   in this way.  Text within such a segment can be
   given \emph{additional} emphasis.
\end{em}

It is sometimes necessary to prevent \LaTeX\ from
breaking a line where it might otherwise do so.
This may be at a space, as between the ``Mr.''\ and
``Jones'' in
       ``Mr.~Jones'',        % ~ produces an unbreakable interword space.
or within a word---especially when the word is a
symbol like
       \mbox{\emph{itemnum}}
that makes little sense when hyphenated across
lines.

Footnotes\footnote{This is an example of a footnote.}
pose no problem.

\LaTeX\ is good at typesetting mathematical formulas
like
       \( x-3y + z = 7 \)
or
       \( a_{1} > x^{2n} + y^{2n} > x' \)
or
       \( \ip{A}{B} = \sum_{i} a_{i} b_{i} \).
The spaces you type in a formula are
ignored.  Remember that a letter like
       $x$                   % $ ... $  and  \( ... \)  are equivalent
is a formula when it denotes a mathematical
symbol, and it should be typed as one.

\section{Displayed Text}

Text is displayed by indenting it from the left
margin.  Quotations are commonly displayed.  There
are short quotations
\begin{quote}
   This is a short quotation.  It consists of a
   single paragraph of text.  See how it is formatted.
\end{quote}
and longer ones.
\begin{quotation}
   This is a longer quotation.  It consists of two
   paragraphs of text, neither of which are
   particularly interesting.

   This is the second paragraph of the quotation.  It
   is just as dull as the first paragraph.
\end{quotation}
Another frequently-displayed structure is a list.
The following is an example of an \emph{itemized}
list.
\begin{itemize}
   \item This is the first item of an itemized list.
         Each item in the list is marked with a ``tick''.
         You don't have to worry about what kind of tick
         mark is used.

   \item This is the second item of the list.  It
         contains another list nested inside it.  The inner
         list is an \emph{enumerated} list.
         \begin{enumerate}
            \item This is the first item of an enumerated
                  list that is nested within the itemized list.

            \item This is the second item of the inner list.
                  \LaTeX\ allows you to nest lists deeper than
                  you really should.
         \end{enumerate}
         This is the rest of the second item of the outer
         list.  It is no more interesting than any other
         part of the item.
   \item This is the third item of the list.
\end{itemize}
You can even display poetry.
\begin{verse}
   There is an environment
    for verse \\             % The \\ command separates lines
   Whose features some poets % within a stanza.
   will curse.

                             % One or more blank lines separate stanzas.

   For instead of making\\
   Them do \emph{all} line breaking, \\
   It allows them to put too many words on a line when they'd rather be
   forced to be terse.
\end{verse}

Mathematical formulas may also be displayed.  A
displayed formula
is
one-line long; multiline
formulas require special formatting instructions.
   \[  \ip{\Gamma}{\psi'} = x'' + y^{2} + z_{i}^{n}\]
Don't start a paragraph with a displayed equation,
nor make one a paragraph by itself.

\end{document}               % End of document.

\section{Bench Section 267}
% synthetic bench section 267
% This is a sample LaTeX input file.  (Version of 12 August 2004.)
%
% A '%' character causes TeX to ignore all remaining text on the line,
% and is used for comments like this one.

\documentclass{article}      % Specifies the document class

                             % The preamble begins here.
\title{An Example Document}  % Declares the document's title.
\author{Leslie Lamport}      % Declares the author's name.
\date{January 21, 1994}      % Deleting this command produces today's date.

\newcommand{\ip}[2]{(#1, #2)}
                             % Defines \ip{arg1}{arg2} to mean
                             % (arg1, arg2).

%\newcommand{\ip}[2]{\langle #1 | #2\rangle}
                             % This is an alternative definition of
                             % \ip that is commented out.

\begin{document}             % End of preamble and beginning of text.

\maketitle                   % Produces the title.

This is an example input file.  Comparing it with
the output it generates can show you how to
produce a simple document of your own.

\section{Ordinary Text}      % Produces section heading.  Lower-level
                             % sections are begun with similar
                             % \subsection and \subsubsection commands.

The ends  of words and sentences are marked
  by   spaces. It  doesn't matter how many
spaces    you type; one is as good as 100.  The
end of   a line counts as a space.

One   or more   blank lines denote the  end
of  a paragraph.

Since any number of consecutive spaces are treated
like a single one, the formatting of the input
file makes no difference to
      \LaTeX,                % The \LaTeX command generates the LaTeX logo.
but it makes a difference to you.  When you use
\LaTeX, making your input file as easy to read
as possible will be a great help as you write
your document and when you change it.  This sample
file shows how you can add comments to your own input
file.

Because printing is different from typewriting,
there are a number of things that you have to do
differently when preparing an input file than if
you were just typing the document directly.
Quotation marks like
       ``this''
have to be handled specially, as do quotes within
quotes:
       ``\,`this'            % \, separates the double and single quote.
        is what I just
        wrote, not  `that'\,''.

Dashes come in three sizes: an
       intra-word
dash, a medium dash for number ranges like
       1--2,
and a punctuation
       dash---like
this.

A sentence-ending space should be larger than the
space between words within a sentence.  You
sometimes have to type special commands in
conjunction with punctuation characters to get
this right, as in the following sentence.
       Gnats, gnus, etc.\ all  % `\ ' makes an inter-word space.
       begin with G\@.         % \@ marks end-of-sentence punctuation.
You should check the spaces after periods when
reading your output to make sure you haven't
forgotten any special cases.  Generating an
ellipsis
       \ldots\               % `\ ' is needed after `\ldots' because TeX
                             % ignores spaces after command names like \ldots
                             % made from \ + letters.
                             %
                             % Note how a `%' character causes TeX to ignore
                             % the end of the input line, so these blank lines
                             % do not start a new paragraph.
                             %
with the right spacing around the periods requires
a special command.

\LaTeX\ interprets some common characters as
commands, so you must type special commands to
generate them.  These characters include the
following:
       \$ \& \% \# \{ and \}.

In printing, text is usually emphasized with an
       \emph{italic}
type style.

\begin{em}
   A long segment of text can also be emphasized
   in this way.  Text within such a segment can be
   given \emph{additional} emphasis.
\end{em}

It is sometimes necessary to prevent \LaTeX\ from
breaking a line where it might otherwise do so.
This may be at a space, as between the ``Mr.''\ and
``Jones'' in
       ``Mr.~Jones'',        % ~ produces an unbreakable interword space.
or within a word---especially when the word is a
symbol like
       \mbox{\emph{itemnum}}
that makes little sense when hyphenated across
lines.

Footnotes\footnote{This is an example of a footnote.}
pose no problem.

\LaTeX\ is good at typesetting mathematical formulas
like
       \( x-3y + z = 7 \)
or
       \( a_{1} > x^{2n} + y^{2n} > x' \)
or
       \( \ip{A}{B} = \sum_{i} a_{i} b_{i} \).
The spaces you type in a formula are
ignored.  Remember that a letter like
       $x$                   % $ ... $  and  \( ... \)  are equivalent
is a formula when it denotes a mathematical
symbol, and it should be typed as one.

\section{Displayed Text}

Text is displayed by indenting it from the left
margin.  Quotations are commonly displayed.  There
are short quotations
\begin{quote}
   This is a short quotation.  It consists of a
   single paragraph of text.  See how it is formatted.
\end{quote}
and longer ones.
\begin{quotation}
   This is a longer quotation.  It consists of two
   paragraphs of text, neither of which are
   particularly interesting.

   This is the second paragraph of the quotation.  It
   is just as dull as the first paragraph.
\end{quotation}
Another frequently-displayed structure is a list.
The following is an example of an \emph{itemized}
list.
\begin{itemize}
   \item This is the first item of an itemized list.
         Each item in the list is marked with a ``tick''.
         You don't have to worry about what kind of tick
         mark is used.

   \item This is the second item of the list.  It
         contains another list nested inside it.  The inner
         list is an \emph{enumerated} list.
         \begin{enumerate}
            \item This is the first item of an enumerated
                  list that is nested within the itemized list.

            \item This is the second item of the inner list.
                  \LaTeX\ allows you to nest lists deeper than
                  you really should.
         \end{enumerate}
         This is the rest of the second item of the outer
         list.  It is no more interesting than any other
         part of the item.
   \item This is the third item of the list.
\end{itemize}
You can even display poetry.
\begin{verse}
   There is an environment
    for verse \\             % The \\ command separates lines
   Whose features some poets % within a stanza.
   will curse.

                             % One or more blank lines separate stanzas.

   For instead of making\\
   Them do \emph{all} line breaking, \\
   It allows them to put too many words on a line when they'd rather be
   forced to be terse.
\end{verse}

Mathematical formulas may also be displayed.  A
displayed formula
is
one-line long; multiline
formulas require special formatting instructions.
   \[  \ip{\Gamma}{\psi'} = x'' + y^{2} + z_{i}^{n}\]
Don't start a paragraph with a displayed equation,
nor make one a paragraph by itself.

\end{document}               % End of document.

\section{Bench Section 268}
% synthetic bench section 268
% This is a sample LaTeX input file.  (Version of 12 August 2004.)
%
% A '%' character causes TeX to ignore all remaining text on the line,
% and is used for comments like this one.

\documentclass{article}      % Specifies the document class

                             % The preamble begins here.
\title{An Example Document}  % Declares the document's title.
\author{Leslie Lamport}      % Declares the author's name.
\date{January 21, 1994}      % Deleting this command produces today's date.

\newcommand{\ip}[2]{(#1, #2)}
                             % Defines \ip{arg1}{arg2} to mean
                             % (arg1, arg2).

%\newcommand{\ip}[2]{\langle #1 | #2\rangle}
                             % This is an alternative definition of
                             % \ip that is commented out.

\begin{document}             % End of preamble and beginning of text.

\maketitle                   % Produces the title.

This is an example input file.  Comparing it with
the output it generates can show you how to
produce a simple document of your own.

\section{Ordinary Text}      % Produces section heading.  Lower-level
                             % sections are begun with similar
                             % \subsection and \subsubsection commands.

The ends  of words and sentences are marked
  by   spaces. It  doesn't matter how many
spaces    you type; one is as good as 100.  The
end of   a line counts as a space.

One   or more   blank lines denote the  end
of  a paragraph.

Since any number of consecutive spaces are treated
like a single one, the formatting of the input
file makes no difference to
      \LaTeX,                % The \LaTeX command generates the LaTeX logo.
but it makes a difference to you.  When you use
\LaTeX, making your input file as easy to read
as possible will be a great help as you write
your document and when you change it.  This sample
file shows how you can add comments to your own input
file.

Because printing is different from typewriting,
there are a number of things that you have to do
differently when preparing an input file than if
you were just typing the document directly.
Quotation marks like
       ``this''
have to be handled specially, as do quotes within
quotes:
       ``\,`this'            % \, separates the double and single quote.
        is what I just
        wrote, not  `that'\,''.

Dashes come in three sizes: an
       intra-word
dash, a medium dash for number ranges like
       1--2,
and a punctuation
       dash---like
this.

A sentence-ending space should be larger than the
space between words within a sentence.  You
sometimes have to type special commands in
conjunction with punctuation characters to get
this right, as in the following sentence.
       Gnats, gnus, etc.\ all  % `\ ' makes an inter-word space.
       begin with G\@.         % \@ marks end-of-sentence punctuation.
You should check the spaces after periods when
reading your output to make sure you haven't
forgotten any special cases.  Generating an
ellipsis
       \ldots\               % `\ ' is needed after `\ldots' because TeX
                             % ignores spaces after command names like \ldots
                             % made from \ + letters.
                             %
                             % Note how a `%' character causes TeX to ignore
                             % the end of the input line, so these blank lines
                             % do not start a new paragraph.
                             %
with the right spacing around the periods requires
a special command.

\LaTeX\ interprets some common characters as
commands, so you must type special commands to
generate them.  These characters include the
following:
       \$ \& \% \# \{ and \}.

In printing, text is usually emphasized with an
       \emph{italic}
type style.

\begin{em}
   A long segment of text can also be emphasized
   in this way.  Text within such a segment can be
   given \emph{additional} emphasis.
\end{em}

It is sometimes necessary to prevent \LaTeX\ from
breaking a line where it might otherwise do so.
This may be at a space, as between the ``Mr.''\ and
``Jones'' in
       ``Mr.~Jones'',        % ~ produces an unbreakable interword space.
or within a word---especially when the word is a
symbol like
       \mbox{\emph{itemnum}}
that makes little sense when hyphenated across
lines.

Footnotes\footnote{This is an example of a footnote.}
pose no problem.

\LaTeX\ is good at typesetting mathematical formulas
like
       \( x-3y + z = 7 \)
or
       \( a_{1} > x^{2n} + y^{2n} > x' \)
or
       \( \ip{A}{B} = \sum_{i} a_{i} b_{i} \).
The spaces you type in a formula are
ignored.  Remember that a letter like
       $x$                   % $ ... $  and  \( ... \)  are equivalent
is a formula when it denotes a mathematical
symbol, and it should be typed as one.

\section{Displayed Text}

Text is displayed by indenting it from the left
margin.  Quotations are commonly displayed.  There
are short quotations
\begin{quote}
   This is a short quotation.  It consists of a
   single paragraph of text.  See how it is formatted.
\end{quote}
and longer ones.
\begin{quotation}
   This is a longer quotation.  It consists of two
   paragraphs of text, neither of which are
   particularly interesting.

   This is the second paragraph of the quotation.  It
   is just as dull as the first paragraph.
\end{quotation}
Another frequently-displayed structure is a list.
The following is an example of an \emph{itemized}
list.
\begin{itemize}
   \item This is the first item of an itemized list.
         Each item in the list is marked with a ``tick''.
         You don't have to worry about what kind of tick
         mark is used.

   \item This is the second item of the list.  It
         contains another list nested inside it.  The inner
         list is an \emph{enumerated} list.
         \begin{enumerate}
            \item This is the first item of an enumerated
                  list that is nested within the itemized list.

            \item This is the second item of the inner list.
                  \LaTeX\ allows you to nest lists deeper than
                  you really should.
         \end{enumerate}
         This is the rest of the second item of the outer
         list.  It is no more interesting than any other
         part of the item.
   \item This is the third item of the list.
\end{itemize}
You can even display poetry.
\begin{verse}
   There is an environment
    for verse \\             % The \\ command separates lines
   Whose features some poets % within a stanza.
   will curse.

                             % One or more blank lines separate stanzas.

   For instead of making\\
   Them do \emph{all} line breaking, \\
   It allows them to put too many words on a line when they'd rather be
   forced to be terse.
\end{verse}

Mathematical formulas may also be displayed.  A
displayed formula
is
one-line long; multiline
formulas require special formatting instructions.
   \[  \ip{\Gamma}{\psi'} = x'' + y^{2} + z_{i}^{n}\]
Don't start a paragraph with a displayed equation,
nor make one a paragraph by itself.

\end{document}               % End of document.

\section{Bench Section 269}
% synthetic bench section 269
% This is a sample LaTeX input file.  (Version of 12 August 2004.)
%
% A '%' character causes TeX to ignore all remaining text on the line,
% and is used for comments like this one.

\documentclass{article}      % Specifies the document class

                             % The preamble begins here.
\title{An Example Document}  % Declares the document's title.
\author{Leslie Lamport}      % Declares the author's name.
\date{January 21, 1994}      % Deleting this command produces today's date.

\newcommand{\ip}[2]{(#1, #2)}
                             % Defines \ip{arg1}{arg2} to mean
                             % (arg1, arg2).

%\newcommand{\ip}[2]{\langle #1 | #2\rangle}
                             % This is an alternative definition of
                             % \ip that is commented out.

\begin{document}             % End of preamble and beginning of text.

\maketitle                   % Produces the title.

This is an example input file.  Comparing it with
the output it generates can show you how to
produce a simple document of your own.

\section{Ordinary Text}      % Produces section heading.  Lower-level
                             % sections are begun with similar
                             % \subsection and \subsubsection commands.

The ends  of words and sentences are marked
  by   spaces. It  doesn't matter how many
spaces    you type; one is as good as 100.  The
end of   a line counts as a space.

One   or more   blank lines denote the  end
of  a paragraph.

Since any number of consecutive spaces are treated
like a single one, the formatting of the input
file makes no difference to
      \LaTeX,                % The \LaTeX command generates the LaTeX logo.
but it makes a difference to you.  When you use
\LaTeX, making your input file as easy to read
as possible will be a great help as you write
your document and when you change it.  This sample
file shows how you can add comments to your own input
file.

Because printing is different from typewriting,
there are a number of things that you have to do
differently when preparing an input file than if
you were just typing the document directly.
Quotation marks like
       ``this''
have to be handled specially, as do quotes within
quotes:
       ``\,`this'            % \, separates the double and single quote.
        is what I just
        wrote, not  `that'\,''.

Dashes come in three sizes: an
       intra-word
dash, a medium dash for number ranges like
       1--2,
and a punctuation
       dash---like
this.

A sentence-ending space should be larger than the
space between words within a sentence.  You
sometimes have to type special commands in
conjunction with punctuation characters to get
this right, as in the following sentence.
       Gnats, gnus, etc.\ all  % `\ ' makes an inter-word space.
       begin with G\@.         % \@ marks end-of-sentence punctuation.
You should check the spaces after periods when
reading your output to make sure you haven't
forgotten any special cases.  Generating an
ellipsis
       \ldots\               % `\ ' is needed after `\ldots' because TeX
                             % ignores spaces after command names like \ldots
                             % made from \ + letters.
                             %
                             % Note how a `%' character causes TeX to ignore
                             % the end of the input line, so these blank lines
                             % do not start a new paragraph.
                             %
with the right spacing around the periods requires
a special command.

\LaTeX\ interprets some common characters as
commands, so you must type special commands to
generate them.  These characters include the
following:
       \$ \& \% \# \{ and \}.

In printing, text is usually emphasized with an
       \emph{italic}
type style.

\begin{em}
   A long segment of text can also be emphasized
   in this way.  Text within such a segment can be
   given \emph{additional} emphasis.
\end{em}

It is sometimes necessary to prevent \LaTeX\ from
breaking a line where it might otherwise do so.
This may be at a space, as between the ``Mr.''\ and
``Jones'' in
       ``Mr.~Jones'',        % ~ produces an unbreakable interword space.
or within a word---especially when the word is a
symbol like
       \mbox{\emph{itemnum}}
that makes little sense when hyphenated across
lines.

Footnotes\footnote{This is an example of a footnote.}
pose no problem.

\LaTeX\ is good at typesetting mathematical formulas
like
       \( x-3y + z = 7 \)
or
       \( a_{1} > x^{2n} + y^{2n} > x' \)
or
       \( \ip{A}{B} = \sum_{i} a_{i} b_{i} \).
The spaces you type in a formula are
ignored.  Remember that a letter like
       $x$                   % $ ... $  and  \( ... \)  are equivalent
is a formula when it denotes a mathematical
symbol, and it should be typed as one.

\section{Displayed Text}

Text is displayed by indenting it from the left
margin.  Quotations are commonly displayed.  There
are short quotations
\begin{quote}
   This is a short quotation.  It consists of a
   single paragraph of text.  See how it is formatted.
\end{quote}
and longer ones.
\begin{quotation}
   This is a longer quotation.  It consists of two
   paragraphs of text, neither of which are
   particularly interesting.

   This is the second paragraph of the quotation.  It
   is just as dull as the first paragraph.
\end{quotation}
Another frequently-displayed structure is a list.
The following is an example of an \emph{itemized}
list.
\begin{itemize}
   \item This is the first item of an itemized list.
         Each item in the list is marked with a ``tick''.
         You don't have to worry about what kind of tick
         mark is used.

   \item This is the second item of the list.  It
         contains another list nested inside it.  The inner
         list is an \emph{enumerated} list.
         \begin{enumerate}
            \item This is the first item of an enumerated
                  list that is nested within the itemized list.

            \item This is the second item of the inner list.
                  \LaTeX\ allows you to nest lists deeper than
                  you really should.
         \end{enumerate}
         This is the rest of the second item of the outer
         list.  It is no more interesting than any other
         part of the item.
   \item This is the third item of the list.
\end{itemize}
You can even display poetry.
\begin{verse}
   There is an environment
    for verse \\             % The \\ command separates lines
   Whose features some poets % within a stanza.
   will curse.

                             % One or more blank lines separate stanzas.

   For instead of making\\
   Them do \emph{all} line breaking, \\
   It allows them to put too many words on a line when they'd rather be
   forced to be terse.
\end{verse}

Mathematical formulas may also be displayed.  A
displayed formula
is
one-line long; multiline
formulas require special formatting instructions.
   \[  \ip{\Gamma}{\psi'} = x'' + y^{2} + z_{i}^{n}\]
Don't start a paragraph with a displayed equation,
nor make one a paragraph by itself.

\end{document}               % End of document.

\section{Bench Section 270}
% synthetic bench section 270
% This is a sample LaTeX input file.  (Version of 12 August 2004.)
%
% A '%' character causes TeX to ignore all remaining text on the line,
% and is used for comments like this one.

\documentclass{article}      % Specifies the document class

                             % The preamble begins here.
\title{An Example Document}  % Declares the document's title.
\author{Leslie Lamport}      % Declares the author's name.
\date{January 21, 1994}      % Deleting this command produces today's date.

\newcommand{\ip}[2]{(#1, #2)}
                             % Defines \ip{arg1}{arg2} to mean
                             % (arg1, arg2).

%\newcommand{\ip}[2]{\langle #1 | #2\rangle}
                             % This is an alternative definition of
                             % \ip that is commented out.

\begin{document}             % End of preamble and beginning of text.

\maketitle                   % Produces the title.

This is an example input file.  Comparing it with
the output it generates can show you how to
produce a simple document of your own.

\section{Ordinary Text}      % Produces section heading.  Lower-level
                             % sections are begun with similar
                             % \subsection and \subsubsection commands.

The ends  of words and sentences are marked
  by   spaces. It  doesn't matter how many
spaces    you type; one is as good as 100.  The
end of   a line counts as a space.

One   or more   blank lines denote the  end
of  a paragraph.

Since any number of consecutive spaces are treated
like a single one, the formatting of the input
file makes no difference to
      \LaTeX,                % The \LaTeX command generates the LaTeX logo.
but it makes a difference to you.  When you use
\LaTeX, making your input file as easy to read
as possible will be a great help as you write
your document and when you change it.  This sample
file shows how you can add comments to your own input
file.

Because printing is different from typewriting,
there are a number of things that you have to do
differently when preparing an input file than if
you were just typing the document directly.
Quotation marks like
       ``this''
have to be handled specially, as do quotes within
quotes:
       ``\,`this'            % \, separates the double and single quote.
        is what I just
        wrote, not  `that'\,''.

Dashes come in three sizes: an
       intra-word
dash, a medium dash for number ranges like
       1--2,
and a punctuation
       dash---like
this.

A sentence-ending space should be larger than the
space between words within a sentence.  You
sometimes have to type special commands in
conjunction with punctuation characters to get
this right, as in the following sentence.
       Gnats, gnus, etc.\ all  % `\ ' makes an inter-word space.
       begin with G\@.         % \@ marks end-of-sentence punctuation.
You should check the spaces after periods when
reading your output to make sure you haven't
forgotten any special cases.  Generating an
ellipsis
       \ldots\               % `\ ' is needed after `\ldots' because TeX
                             % ignores spaces after command names like \ldots
                             % made from \ + letters.
                             %
                             % Note how a `%' character causes TeX to ignore
                             % the end of the input line, so these blank lines
                             % do not start a new paragraph.
                             %
with the right spacing around the periods requires
a special command.

\LaTeX\ interprets some common characters as
commands, so you must type special commands to
generate them.  These characters include the
following:
       \$ \& \% \# \{ and \}.

In printing, text is usually emphasized with an
       \emph{italic}
type style.

\begin{em}
   A long segment of text can also be emphasized
   in this way.  Text within such a segment can be
   given \emph{additional} emphasis.
\end{em}

It is sometimes necessary to prevent \LaTeX\ from
breaking a line where it might otherwise do so.
This may be at a space, as between the ``Mr.''\ and
``Jones'' in
       ``Mr.~Jones'',        % ~ produces an unbreakable interword space.
or within a word---especially when the word is a
symbol like
       \mbox{\emph{itemnum}}
that makes little sense when hyphenated across
lines.

Footnotes\footnote{This is an example of a footnote.}
pose no problem.

\LaTeX\ is good at typesetting mathematical formulas
like
       \( x-3y + z = 7 \)
or
       \( a_{1} > x^{2n} + y^{2n} > x' \)
or
       \( \ip{A}{B} = \sum_{i} a_{i} b_{i} \).
The spaces you type in a formula are
ignored.  Remember that a letter like
       $x$                   % $ ... $  and  \( ... \)  are equivalent
is a formula when it denotes a mathematical
symbol, and it should be typed as one.

\section{Displayed Text}

Text is displayed by indenting it from the left
margin.  Quotations are commonly displayed.  There
are short quotations
\begin{quote}
   This is a short quotation.  It consists of a
   single paragraph of text.  See how it is formatted.
\end{quote}
and longer ones.
\begin{quotation}
   This is a longer quotation.  It consists of two
   paragraphs of text, neither of which are
   particularly interesting.

   This is the second paragraph of the quotation.  It
   is just as dull as the first paragraph.
\end{quotation}
Another frequently-displayed structure is a list.
The following is an example of an \emph{itemized}
list.
\begin{itemize}
   \item This is the first item of an itemized list.
         Each item in the list is marked with a ``tick''.
         You don't have to worry about what kind of tick
         mark is used.

   \item This is the second item of the list.  It
         contains another list nested inside it.  The inner
         list is an \emph{enumerated} list.
         \begin{enumerate}
            \item This is the first item of an enumerated
                  list that is nested within the itemized list.

            \item This is the second item of the inner list.
                  \LaTeX\ allows you to nest lists deeper than
                  you really should.
         \end{enumerate}
         This is the rest of the second item of the outer
         list.  It is no more interesting than any other
         part of the item.
   \item This is the third item of the list.
\end{itemize}
You can even display poetry.
\begin{verse}
   There is an environment
    for verse \\             % The \\ command separates lines
   Whose features some poets % within a stanza.
   will curse.

                             % One or more blank lines separate stanzas.

   For instead of making\\
   Them do \emph{all} line breaking, \\
   It allows them to put too many words on a line when they'd rather be
   forced to be terse.
\end{verse}

Mathematical formulas may also be displayed.  A
displayed formula
is
one-line long; multiline
formulas require special formatting instructions.
   \[  \ip{\Gamma}{\psi'} = x'' + y^{2} + z_{i}^{n}\]
Don't start a paragraph with a displayed equation,
nor make one a paragraph by itself.

\end{document}               % End of document.

\section{Bench Section 271}
% synthetic bench section 271
% This is a sample LaTeX input file.  (Version of 12 August 2004.)
%
% A '%' character causes TeX to ignore all remaining text on the line,
% and is used for comments like this one.

\documentclass{article}      % Specifies the document class

                             % The preamble begins here.
\title{An Example Document}  % Declares the document's title.
\author{Leslie Lamport}      % Declares the author's name.
\date{January 21, 1994}      % Deleting this command produces today's date.

\newcommand{\ip}[2]{(#1, #2)}
                             % Defines \ip{arg1}{arg2} to mean
                             % (arg1, arg2).

%\newcommand{\ip}[2]{\langle #1 | #2\rangle}
                             % This is an alternative definition of
                             % \ip that is commented out.

\begin{document}             % End of preamble and beginning of text.

\maketitle                   % Produces the title.

This is an example input file.  Comparing it with
the output it generates can show you how to
produce a simple document of your own.

\section{Ordinary Text}      % Produces section heading.  Lower-level
                             % sections are begun with similar
                             % \subsection and \subsubsection commands.

The ends  of words and sentences are marked
  by   spaces. It  doesn't matter how many
spaces    you type; one is as good as 100.  The
end of   a line counts as a space.

One   or more   blank lines denote the  end
of  a paragraph.

Since any number of consecutive spaces are treated
like a single one, the formatting of the input
file makes no difference to
      \LaTeX,                % The \LaTeX command generates the LaTeX logo.
but it makes a difference to you.  When you use
\LaTeX, making your input file as easy to read
as possible will be a great help as you write
your document and when you change it.  This sample
file shows how you can add comments to your own input
file.

Because printing is different from typewriting,
there are a number of things that you have to do
differently when preparing an input file than if
you were just typing the document directly.
Quotation marks like
       ``this''
have to be handled specially, as do quotes within
quotes:
       ``\,`this'            % \, separates the double and single quote.
        is what I just
        wrote, not  `that'\,''.

Dashes come in three sizes: an
       intra-word
dash, a medium dash for number ranges like
       1--2,
and a punctuation
       dash---like
this.

A sentence-ending space should be larger than the
space between words within a sentence.  You
sometimes have to type special commands in
conjunction with punctuation characters to get
this right, as in the following sentence.
       Gnats, gnus, etc.\ all  % `\ ' makes an inter-word space.
       begin with G\@.         % \@ marks end-of-sentence punctuation.
You should check the spaces after periods when
reading your output to make sure you haven't
forgotten any special cases.  Generating an
ellipsis
       \ldots\               % `\ ' is needed after `\ldots' because TeX
                             % ignores spaces after command names like \ldots
                             % made from \ + letters.
                             %
                             % Note how a `%' character causes TeX to ignore
                             % the end of the input line, so these blank lines
                             % do not start a new paragraph.
                             %
with the right spacing around the periods requires
a special command.

\LaTeX\ interprets some common characters as
commands, so you must type special commands to
generate them.  These characters include the
following:
       \$ \& \% \# \{ and \}.

In printing, text is usually emphasized with an
       \emph{italic}
type style.

\begin{em}
   A long segment of text can also be emphasized
   in this way.  Text within such a segment can be
   given \emph{additional} emphasis.
\end{em}

It is sometimes necessary to prevent \LaTeX\ from
breaking a line where it might otherwise do so.
This may be at a space, as between the ``Mr.''\ and
``Jones'' in
       ``Mr.~Jones'',        % ~ produces an unbreakable interword space.
or within a word---especially when the word is a
symbol like
       \mbox{\emph{itemnum}}
that makes little sense when hyphenated across
lines.

Footnotes\footnote{This is an example of a footnote.}
pose no problem.

\LaTeX\ is good at typesetting mathematical formulas
like
       \( x-3y + z = 7 \)
or
       \( a_{1} > x^{2n} + y^{2n} > x' \)
or
       \( \ip{A}{B} = \sum_{i} a_{i} b_{i} \).
The spaces you type in a formula are
ignored.  Remember that a letter like
       $x$                   % $ ... $  and  \( ... \)  are equivalent
is a formula when it denotes a mathematical
symbol, and it should be typed as one.

\section{Displayed Text}

Text is displayed by indenting it from the left
margin.  Quotations are commonly displayed.  There
are short quotations
\begin{quote}
   This is a short quotation.  It consists of a
   single paragraph of text.  See how it is formatted.
\end{quote}
and longer ones.
\begin{quotation}
   This is a longer quotation.  It consists of two
   paragraphs of text, neither of which are
   particularly interesting.

   This is the second paragraph of the quotation.  It
   is just as dull as the first paragraph.
\end{quotation}
Another frequently-displayed structure is a list.
The following is an example of an \emph{itemized}
list.
\begin{itemize}
   \item This is the first item of an itemized list.
         Each item in the list is marked with a ``tick''.
         You don't have to worry about what kind of tick
         mark is used.

   \item This is the second item of the list.  It
         contains another list nested inside it.  The inner
         list is an \emph{enumerated} list.
         \begin{enumerate}
            \item This is the first item of an enumerated
                  list that is nested within the itemized list.

            \item This is the second item of the inner list.
                  \LaTeX\ allows you to nest lists deeper than
                  you really should.
         \end{enumerate}
         This is the rest of the second item of the outer
         list.  It is no more interesting than any other
         part of the item.
   \item This is the third item of the list.
\end{itemize}
You can even display poetry.
\begin{verse}
   There is an environment
    for verse \\             % The \\ command separates lines
   Whose features some poets % within a stanza.
   will curse.

                             % One or more blank lines separate stanzas.

   For instead of making\\
   Them do \emph{all} line breaking, \\
   It allows them to put too many words on a line when they'd rather be
   forced to be terse.
\end{verse}

Mathematical formulas may also be displayed.  A
displayed formula
is
one-line long; multiline
formulas require special formatting instructions.
   \[  \ip{\Gamma}{\psi'} = x'' + y^{2} + z_{i}^{n}\]
Don't start a paragraph with a displayed equation,
nor make one a paragraph by itself.

\end{document}               % End of document.

\section{Bench Section 272}
% synthetic bench section 272
% This is a sample LaTeX input file.  (Version of 12 August 2004.)
%
% A '%' character causes TeX to ignore all remaining text on the line,
% and is used for comments like this one.

\documentclass{article}      % Specifies the document class

                             % The preamble begins here.
\title{An Example Document}  % Declares the document's title.
\author{Leslie Lamport}      % Declares the author's name.
\date{January 21, 1994}      % Deleting this command produces today's date.

\newcommand{\ip}[2]{(#1, #2)}
                             % Defines \ip{arg1}{arg2} to mean
                             % (arg1, arg2).

%\newcommand{\ip}[2]{\langle #1 | #2\rangle}
                             % This is an alternative definition of
                             % \ip that is commented out.

\begin{document}             % End of preamble and beginning of text.

\maketitle                   % Produces the title.

This is an example input file.  Comparing it with
the output it generates can show you how to
produce a simple document of your own.

\section{Ordinary Text}      % Produces section heading.  Lower-level
                             % sections are begun with similar
                             % \subsection and \subsubsection commands.

The ends  of words and sentences are marked
  by   spaces. It  doesn't matter how many
spaces    you type; one is as good as 100.  The
end of   a line counts as a space.

One   or more   blank lines denote the  end
of  a paragraph.

Since any number of consecutive spaces are treated
like a single one, the formatting of the input
file makes no difference to
      \LaTeX,                % The \LaTeX command generates the LaTeX logo.
but it makes a difference to you.  When you use
\LaTeX, making your input file as easy to read
as possible will be a great help as you write
your document and when you change it.  This sample
file shows how you can add comments to your own input
file.

Because printing is different from typewriting,
there are a number of things that you have to do
differently when preparing an input file than if
you were just typing the document directly.
Quotation marks like
       ``this''
have to be handled specially, as do quotes within
quotes:
       ``\,`this'            % \, separates the double and single quote.
        is what I just
        wrote, not  `that'\,''.

Dashes come in three sizes: an
       intra-word
dash, a medium dash for number ranges like
       1--2,
and a punctuation
       dash---like
this.

A sentence-ending space should be larger than the
space between words within a sentence.  You
sometimes have to type special commands in
conjunction with punctuation characters to get
this right, as in the following sentence.
       Gnats, gnus, etc.\ all  % `\ ' makes an inter-word space.
       begin with G\@.         % \@ marks end-of-sentence punctuation.
You should check the spaces after periods when
reading your output to make sure you haven't
forgotten any special cases.  Generating an
ellipsis
       \ldots\               % `\ ' is needed after `\ldots' because TeX
                             % ignores spaces after command names like \ldots
                             % made from \ + letters.
                             %
                             % Note how a `%' character causes TeX to ignore
                             % the end of the input line, so these blank lines
                             % do not start a new paragraph.
                             %
with the right spacing around the periods requires
a special command.

\LaTeX\ interprets some common characters as
commands, so you must type special commands to
generate them.  These characters include the
following:
       \$ \& \% \# \{ and \}.

In printing, text is usually emphasized with an
       \emph{italic}
type style.

\begin{em}
   A long segment of text can also be emphasized
   in this way.  Text within such a segment can be
   given \emph{additional} emphasis.
\end{em}

It is sometimes necessary to prevent \LaTeX\ from
breaking a line where it might otherwise do so.
This may be at a space, as between the ``Mr.''\ and
``Jones'' in
       ``Mr.~Jones'',        % ~ produces an unbreakable interword space.
or within a word---especially when the word is a
symbol like
       \mbox{\emph{itemnum}}
that makes little sense when hyphenated across
lines.

Footnotes\footnote{This is an example of a footnote.}
pose no problem.

\LaTeX\ is good at typesetting mathematical formulas
like
       \( x-3y + z = 7 \)
or
       \( a_{1} > x^{2n} + y^{2n} > x' \)
or
       \( \ip{A}{B} = \sum_{i} a_{i} b_{i} \).
The spaces you type in a formula are
ignored.  Remember that a letter like
       $x$                   % $ ... $  and  \( ... \)  are equivalent
is a formula when it denotes a mathematical
symbol, and it should be typed as one.

\section{Displayed Text}

Text is displayed by indenting it from the left
margin.  Quotations are commonly displayed.  There
are short quotations
\begin{quote}
   This is a short quotation.  It consists of a
   single paragraph of text.  See how it is formatted.
\end{quote}
and longer ones.
\begin{quotation}
   This is a longer quotation.  It consists of two
   paragraphs of text, neither of which are
   particularly interesting.

   This is the second paragraph of the quotation.  It
   is just as dull as the first paragraph.
\end{quotation}
Another frequently-displayed structure is a list.
The following is an example of an \emph{itemized}
list.
\begin{itemize}
   \item This is the first item of an itemized list.
         Each item in the list is marked with a ``tick''.
         You don't have to worry about what kind of tick
         mark is used.

   \item This is the second item of the list.  It
         contains another list nested inside it.  The inner
         list is an \emph{enumerated} list.
         \begin{enumerate}
            \item This is the first item of an enumerated
                  list that is nested within the itemized list.

            \item This is the second item of the inner list.
                  \LaTeX\ allows you to nest lists deeper than
                  you really should.
         \end{enumerate}
         This is the rest of the second item of the outer
         list.  It is no more interesting than any other
         part of the item.
   \item This is the third item of the list.
\end{itemize}
You can even display poetry.
\begin{verse}
   There is an environment
    for verse \\             % The \\ command separates lines
   Whose features some poets % within a stanza.
   will curse.

                             % One or more blank lines separate stanzas.

   For instead of making\\
   Them do \emph{all} line breaking, \\
   It allows them to put too many words on a line when they'd rather be
   forced to be terse.
\end{verse}

Mathematical formulas may also be displayed.  A
displayed formula
is
one-line long; multiline
formulas require special formatting instructions.
   \[  \ip{\Gamma}{\psi'} = x'' + y^{2} + z_{i}^{n}\]
Don't start a paragraph with a displayed equation,
nor make one a paragraph by itself.

\end{document}               % End of document.

\section{Bench Section 273}
% synthetic bench section 273
% This is a sample LaTeX input file.  (Version of 12 August 2004.)
%
% A '%' character causes TeX to ignore all remaining text on the line,
% and is used for comments like this one.

\documentclass{article}      % Specifies the document class

                             % The preamble begins here.
\title{An Example Document}  % Declares the document's title.
\author{Leslie Lamport}      % Declares the author's name.
\date{January 21, 1994}      % Deleting this command produces today's date.

\newcommand{\ip}[2]{(#1, #2)}
                             % Defines \ip{arg1}{arg2} to mean
                             % (arg1, arg2).

%\newcommand{\ip}[2]{\langle #1 | #2\rangle}
                             % This is an alternative definition of
                             % \ip that is commented out.

\begin{document}             % End of preamble and beginning of text.

\maketitle                   % Produces the title.

This is an example input file.  Comparing it with
the output it generates can show you how to
produce a simple document of your own.

\section{Ordinary Text}      % Produces section heading.  Lower-level
                             % sections are begun with similar
                             % \subsection and \subsubsection commands.

The ends  of words and sentences are marked
  by   spaces. It  doesn't matter how many
spaces    you type; one is as good as 100.  The
end of   a line counts as a space.

One   or more   blank lines denote the  end
of  a paragraph.

Since any number of consecutive spaces are treated
like a single one, the formatting of the input
file makes no difference to
      \LaTeX,                % The \LaTeX command generates the LaTeX logo.
but it makes a difference to you.  When you use
\LaTeX, making your input file as easy to read
as possible will be a great help as you write
your document and when you change it.  This sample
file shows how you can add comments to your own input
file.

Because printing is different from typewriting,
there are a number of things that you have to do
differently when preparing an input file than if
you were just typing the document directly.
Quotation marks like
       ``this''
have to be handled specially, as do quotes within
quotes:
       ``\,`this'            % \, separates the double and single quote.
        is what I just
        wrote, not  `that'\,''.

Dashes come in three sizes: an
       intra-word
dash, a medium dash for number ranges like
       1--2,
and a punctuation
       dash---like
this.

A sentence-ending space should be larger than the
space between words within a sentence.  You
sometimes have to type special commands in
conjunction with punctuation characters to get
this right, as in the following sentence.
       Gnats, gnus, etc.\ all  % `\ ' makes an inter-word space.
       begin with G\@.         % \@ marks end-of-sentence punctuation.
You should check the spaces after periods when
reading your output to make sure you haven't
forgotten any special cases.  Generating an
ellipsis
       \ldots\               % `\ ' is needed after `\ldots' because TeX
                             % ignores spaces after command names like \ldots
                             % made from \ + letters.
                             %
                             % Note how a `%' character causes TeX to ignore
                             % the end of the input line, so these blank lines
                             % do not start a new paragraph.
                             %
with the right spacing around the periods requires
a special command.

\LaTeX\ interprets some common characters as
commands, so you must type special commands to
generate them.  These characters include the
following:
       \$ \& \% \# \{ and \}.

In printing, text is usually emphasized with an
       \emph{italic}
type style.

\begin{em}
   A long segment of text can also be emphasized
   in this way.  Text within such a segment can be
   given \emph{additional} emphasis.
\end{em}

It is sometimes necessary to prevent \LaTeX\ from
breaking a line where it might otherwise do so.
This may be at a space, as between the ``Mr.''\ and
``Jones'' in
       ``Mr.~Jones'',        % ~ produces an unbreakable interword space.
or within a word---especially when the word is a
symbol like
       \mbox{\emph{itemnum}}
that makes little sense when hyphenated across
lines.

Footnotes\footnote{This is an example of a footnote.}
pose no problem.

\LaTeX\ is good at typesetting mathematical formulas
like
       \( x-3y + z = 7 \)
or
       \( a_{1} > x^{2n} + y^{2n} > x' \)
or
       \( \ip{A}{B} = \sum_{i} a_{i} b_{i} \).
The spaces you type in a formula are
ignored.  Remember that a letter like
       $x$                   % $ ... $  and  \( ... \)  are equivalent
is a formula when it denotes a mathematical
symbol, and it should be typed as one.

\section{Displayed Text}

Text is displayed by indenting it from the left
margin.  Quotations are commonly displayed.  There
are short quotations
\begin{quote}
   This is a short quotation.  It consists of a
   single paragraph of text.  See how it is formatted.
\end{quote}
and longer ones.
\begin{quotation}
   This is a longer quotation.  It consists of two
   paragraphs of text, neither of which are
   particularly interesting.

   This is the second paragraph of the quotation.  It
   is just as dull as the first paragraph.
\end{quotation}
Another frequently-displayed structure is a list.
The following is an example of an \emph{itemized}
list.
\begin{itemize}
   \item This is the first item of an itemized list.
         Each item in the list is marked with a ``tick''.
         You don't have to worry about what kind of tick
         mark is used.

   \item This is the second item of the list.  It
         contains another list nested inside it.  The inner
         list is an \emph{enumerated} list.
         \begin{enumerate}
            \item This is the first item of an enumerated
                  list that is nested within the itemized list.

            \item This is the second item of the inner list.
                  \LaTeX\ allows you to nest lists deeper than
                  you really should.
         \end{enumerate}
         This is the rest of the second item of the outer
         list.  It is no more interesting than any other
         part of the item.
   \item This is the third item of the list.
\end{itemize}
You can even display poetry.
\begin{verse}
   There is an environment
    for verse \\             % The \\ command separates lines
   Whose features some poets % within a stanza.
   will curse.

                             % One or more blank lines separate stanzas.

   For instead of making\\
   Them do \emph{all} line breaking, \\
   It allows them to put too many words on a line when they'd rather be
   forced to be terse.
\end{verse}

Mathematical formulas may also be displayed.  A
displayed formula
is
one-line long; multiline
formulas require special formatting instructions.
   \[  \ip{\Gamma}{\psi'} = x'' + y^{2} + z_{i}^{n}\]
Don't start a paragraph with a displayed equation,
nor make one a paragraph by itself.

\end{document}               % End of document.

\section{Bench Section 274}
% synthetic bench section 274
% This is a sample LaTeX input file.  (Version of 12 August 2004.)
%
% A '%' character causes TeX to ignore all remaining text on the line,
% and is used for comments like this one.

\documentclass{article}      % Specifies the document class

                             % The preamble begins here.
\title{An Example Document}  % Declares the document's title.
\author{Leslie Lamport}      % Declares the author's name.
\date{January 21, 1994}      % Deleting this command produces today's date.

\newcommand{\ip}[2]{(#1, #2)}
                             % Defines \ip{arg1}{arg2} to mean
                             % (arg1, arg2).

%\newcommand{\ip}[2]{\langle #1 | #2\rangle}
                             % This is an alternative definition of
                             % \ip that is commented out.

\begin{document}             % End of preamble and beginning of text.

\maketitle                   % Produces the title.

This is an example input file.  Comparing it with
the output it generates can show you how to
produce a simple document of your own.

\section{Ordinary Text}      % Produces section heading.  Lower-level
                             % sections are begun with similar
                             % \subsection and \subsubsection commands.

The ends  of words and sentences are marked
  by   spaces. It  doesn't matter how many
spaces    you type; one is as good as 100.  The
end of   a line counts as a space.

One   or more   blank lines denote the  end
of  a paragraph.

Since any number of consecutive spaces are treated
like a single one, the formatting of the input
file makes no difference to
      \LaTeX,                % The \LaTeX command generates the LaTeX logo.
but it makes a difference to you.  When you use
\LaTeX, making your input file as easy to read
as possible will be a great help as you write
your document and when you change it.  This sample
file shows how you can add comments to your own input
file.

Because printing is different from typewriting,
there are a number of things that you have to do
differently when preparing an input file than if
you were just typing the document directly.
Quotation marks like
       ``this''
have to be handled specially, as do quotes within
quotes:
       ``\,`this'            % \, separates the double and single quote.
        is what I just
        wrote, not  `that'\,''.

Dashes come in three sizes: an
       intra-word
dash, a medium dash for number ranges like
       1--2,
and a punctuation
       dash---like
this.

A sentence-ending space should be larger than the
space between words within a sentence.  You
sometimes have to type special commands in
conjunction with punctuation characters to get
this right, as in the following sentence.
       Gnats, gnus, etc.\ all  % `\ ' makes an inter-word space.
       begin with G\@.         % \@ marks end-of-sentence punctuation.
You should check the spaces after periods when
reading your output to make sure you haven't
forgotten any special cases.  Generating an
ellipsis
       \ldots\               % `\ ' is needed after `\ldots' because TeX
                             % ignores spaces after command names like \ldots
                             % made from \ + letters.
                             %
                             % Note how a `%' character causes TeX to ignore
                             % the end of the input line, so these blank lines
                             % do not start a new paragraph.
                             %
with the right spacing around the periods requires
a special command.

\LaTeX\ interprets some common characters as
commands, so you must type special commands to
generate them.  These characters include the
following:
       \$ \& \% \# \{ and \}.

In printing, text is usually emphasized with an
       \emph{italic}
type style.

\begin{em}
   A long segment of text can also be emphasized
   in this way.  Text within such a segment can be
   given \emph{additional} emphasis.
\end{em}

It is sometimes necessary to prevent \LaTeX\ from
breaking a line where it might otherwise do so.
This may be at a space, as between the ``Mr.''\ and
``Jones'' in
       ``Mr.~Jones'',        % ~ produces an unbreakable interword space.
or within a word---especially when the word is a
symbol like
       \mbox{\emph{itemnum}}
that makes little sense when hyphenated across
lines.

Footnotes\footnote{This is an example of a footnote.}
pose no problem.

\LaTeX\ is good at typesetting mathematical formulas
like
       \( x-3y + z = 7 \)
or
       \( a_{1} > x^{2n} + y^{2n} > x' \)
or
       \( \ip{A}{B} = \sum_{i} a_{i} b_{i} \).
The spaces you type in a formula are
ignored.  Remember that a letter like
       $x$                   % $ ... $  and  \( ... \)  are equivalent
is a formula when it denotes a mathematical
symbol, and it should be typed as one.

\section{Displayed Text}

Text is displayed by indenting it from the left
margin.  Quotations are commonly displayed.  There
are short quotations
\begin{quote}
   This is a short quotation.  It consists of a
   single paragraph of text.  See how it is formatted.
\end{quote}
and longer ones.
\begin{quotation}
   This is a longer quotation.  It consists of two
   paragraphs of text, neither of which are
   particularly interesting.

   This is the second paragraph of the quotation.  It
   is just as dull as the first paragraph.
\end{quotation}
Another frequently-displayed structure is a list.
The following is an example of an \emph{itemized}
list.
\begin{itemize}
   \item This is the first item of an itemized list.
         Each item in the list is marked with a ``tick''.
         You don't have to worry about what kind of tick
         mark is used.

   \item This is the second item of the list.  It
         contains another list nested inside it.  The inner
         list is an \emph{enumerated} list.
         \begin{enumerate}
            \item This is the first item of an enumerated
                  list that is nested within the itemized list.

            \item This is the second item of the inner list.
                  \LaTeX\ allows you to nest lists deeper than
                  you really should.
         \end{enumerate}
         This is the rest of the second item of the outer
         list.  It is no more interesting than any other
         part of the item.
   \item This is the third item of the list.
\end{itemize}
You can even display poetry.
\begin{verse}
   There is an environment
    for verse \\             % The \\ command separates lines
   Whose features some poets % within a stanza.
   will curse.

                             % One or more blank lines separate stanzas.

   For instead of making\\
   Them do \emph{all} line breaking, \\
   It allows them to put too many words on a line when they'd rather be
   forced to be terse.
\end{verse}

Mathematical formulas may also be displayed.  A
displayed formula
is
one-line long; multiline
formulas require special formatting instructions.
   \[  \ip{\Gamma}{\psi'} = x'' + y^{2} + z_{i}^{n}\]
Don't start a paragraph with a displayed equation,
nor make one a paragraph by itself.

\end{document}               % End of document.

\section{Bench Section 275}
% synthetic bench section 275
% This is a sample LaTeX input file.  (Version of 12 August 2004.)
%
% A '%' character causes TeX to ignore all remaining text on the line,
% and is used for comments like this one.

\documentclass{article}      % Specifies the document class

                             % The preamble begins here.
\title{An Example Document}  % Declares the document's title.
\author{Leslie Lamport}      % Declares the author's name.
\date{January 21, 1994}      % Deleting this command produces today's date.

\newcommand{\ip}[2]{(#1, #2)}
                             % Defines \ip{arg1}{arg2} to mean
                             % (arg1, arg2).

%\newcommand{\ip}[2]{\langle #1 | #2\rangle}
                             % This is an alternative definition of
                             % \ip that is commented out.

\begin{document}             % End of preamble and beginning of text.

\maketitle                   % Produces the title.

This is an example input file.  Comparing it with
the output it generates can show you how to
produce a simple document of your own.

\section{Ordinary Text}      % Produces section heading.  Lower-level
                             % sections are begun with similar
                             % \subsection and \subsubsection commands.

The ends  of words and sentences are marked
  by   spaces. It  doesn't matter how many
spaces    you type; one is as good as 100.  The
end of   a line counts as a space.

One   or more   blank lines denote the  end
of  a paragraph.

Since any number of consecutive spaces are treated
like a single one, the formatting of the input
file makes no difference to
      \LaTeX,                % The \LaTeX command generates the LaTeX logo.
but it makes a difference to you.  When you use
\LaTeX, making your input file as easy to read
as possible will be a great help as you write
your document and when you change it.  This sample
file shows how you can add comments to your own input
file.

Because printing is different from typewriting,
there are a number of things that you have to do
differently when preparing an input file than if
you were just typing the document directly.
Quotation marks like
       ``this''
have to be handled specially, as do quotes within
quotes:
       ``\,`this'            % \, separates the double and single quote.
        is what I just
        wrote, not  `that'\,''.

Dashes come in three sizes: an
       intra-word
dash, a medium dash for number ranges like
       1--2,
and a punctuation
       dash---like
this.

A sentence-ending space should be larger than the
space between words within a sentence.  You
sometimes have to type special commands in
conjunction with punctuation characters to get
this right, as in the following sentence.
       Gnats, gnus, etc.\ all  % `\ ' makes an inter-word space.
       begin with G\@.         % \@ marks end-of-sentence punctuation.
You should check the spaces after periods when
reading your output to make sure you haven't
forgotten any special cases.  Generating an
ellipsis
       \ldots\               % `\ ' is needed after `\ldots' because TeX
                             % ignores spaces after command names like \ldots
                             % made from \ + letters.
                             %
                             % Note how a `%' character causes TeX to ignore
                             % the end of the input line, so these blank lines
                             % do not start a new paragraph.
                             %
with the right spacing around the periods requires
a special command.

\LaTeX\ interprets some common characters as
commands, so you must type special commands to
generate them.  These characters include the
following:
       \$ \& \% \# \{ and \}.

In printing, text is usually emphasized with an
       \emph{italic}
type style.

\begin{em}
   A long segment of text can also be emphasized
   in this way.  Text within such a segment can be
   given \emph{additional} emphasis.
\end{em}

It is sometimes necessary to prevent \LaTeX\ from
breaking a line where it might otherwise do so.
This may be at a space, as between the ``Mr.''\ and
``Jones'' in
       ``Mr.~Jones'',        % ~ produces an unbreakable interword space.
or within a word---especially when the word is a
symbol like
       \mbox{\emph{itemnum}}
that makes little sense when hyphenated across
lines.

Footnotes\footnote{This is an example of a footnote.}
pose no problem.

\LaTeX\ is good at typesetting mathematical formulas
like
       \( x-3y + z = 7 \)
or
       \( a_{1} > x^{2n} + y^{2n} > x' \)
or
       \( \ip{A}{B} = \sum_{i} a_{i} b_{i} \).
The spaces you type in a formula are
ignored.  Remember that a letter like
       $x$                   % $ ... $  and  \( ... \)  are equivalent
is a formula when it denotes a mathematical
symbol, and it should be typed as one.

\section{Displayed Text}

Text is displayed by indenting it from the left
margin.  Quotations are commonly displayed.  There
are short quotations
\begin{quote}
   This is a short quotation.  It consists of a
   single paragraph of text.  See how it is formatted.
\end{quote}
and longer ones.
\begin{quotation}
   This is a longer quotation.  It consists of two
   paragraphs of text, neither of which are
   particularly interesting.

   This is the second paragraph of the quotation.  It
   is just as dull as the first paragraph.
\end{quotation}
Another frequently-displayed structure is a list.
The following is an example of an \emph{itemized}
list.
\begin{itemize}
   \item This is the first item of an itemized list.
         Each item in the list is marked with a ``tick''.
         You don't have to worry about what kind of tick
         mark is used.

   \item This is the second item of the list.  It
         contains another list nested inside it.  The inner
         list is an \emph{enumerated} list.
         \begin{enumerate}
            \item This is the first item of an enumerated
                  list that is nested within the itemized list.

            \item This is the second item of the inner list.
                  \LaTeX\ allows you to nest lists deeper than
                  you really should.
         \end{enumerate}
         This is the rest of the second item of the outer
         list.  It is no more interesting than any other
         part of the item.
   \item This is the third item of the list.
\end{itemize}
You can even display poetry.
\begin{verse}
   There is an environment
    for verse \\             % The \\ command separates lines
   Whose features some poets % within a stanza.
   will curse.

                             % One or more blank lines separate stanzas.

   For instead of making\\
   Them do \emph{all} line breaking, \\
   It allows them to put too many words on a line when they'd rather be
   forced to be terse.
\end{verse}

Mathematical formulas may also be displayed.  A
displayed formula
is
one-line long; multiline
formulas require special formatting instructions.
   \[  \ip{\Gamma}{\psi'} = x'' + y^{2} + z_{i}^{n}\]
Don't start a paragraph with a displayed equation,
nor make one a paragraph by itself.

\end{document}               % End of document.

